% ############################################### PREAMBLE AND CONTROL PANEL #

% ########################################################### DOCUMENT CLASS #
\documentclass[  %   draft
             % , demo  % black boxes instead of figures
              , 8pt, twoside, openright]{extbook}

% ########################################################## AND CONTROL PANEL #

% true - hmarginration 1:1, urlcolor - darkgreen, version - for searching
\newif\ifscreen 	   \screentrue
% true - no apendix, no author of grammer, no backmatter ISBN, 
\newif\ifPDF               \PDFtrue
% true - rotate pages of morph verb tables
\newif\ifPDFrotate         \PDFrotatetrue
% true - adds white searchable value !
\newif\ifPDFsearchvalue    \PDFsearchvaluetrue
\newif\ifmakecovers        \makecoverstrue
\newif\ifmakefrontmatter   \makefrontmattertrue
\newif\ifinputintroduction \inputintroductiontrue
\newif\ifinputletters      \inputletterstrue
\newif\ifinputphon         \inputphontrue
\newif\ifinputmorpho       \inputmorphotrue
\newif\ifmakebackmatter    \makebackmattertrue
\newif\ifshowcrop           \showcroptrue
% corrections
\newif\ifC                  \Ctrue

\newif\ifPRINTHOUSE 	   \PRINTHOUSEfalse
\newif\ifPDFSEARCH  	   \PDFSEARCHfalse
\newif\ifPDFPRINT 	   \PDFPRINTfalse
\newif\ifPDFtest	   \PDFtesttrue

%versions
%PRINTHOUSE
\ifPRINTHOUSE
\makecoversfalse
\screenfalse
\PDFfalse
\PDFrotatefalse
\PDFsearchvaluefalse
\showcroptrue
\Ctrue
\fi

%PDFSEARCH
\ifPDFSEARCH
\screentrue
\PDFtrue
\PDFrotatetrue
\PDFsearchvaluetrue
\showcropfalse
\Ctrue
\fi

%PDFPRINT
\ifPDFPRINT
\screenfalse
\PDFtrue
\PDFrotatefalse
\PDFsearchvaluefalse
\showcropfalse
\Ctrue
\fi

%PDFtest
\ifPDFtest
\makecoversfalse
\makefrontmatterfalse
\inputintroductionfalse
\inputletterstrue
\inputphonfalse
\inputmorphofalse
\makebackmatterfalse
\screentrue
\PDFtrue
\PDFrotatefalse
\PDFsearchvaluefalse
\showcropfalse
\Cfalse
\fi

% 1. PDF (for printing)  
% a) searchable value, b) no apendix, c) version= PDF d) no ISBN and no ISBN bar
% e) automatic rotation of landscape pages, f) no author of Icelandic Grammer
% 2. PDF (for viewing offline and searching) + screen (horizontal ratio 1:1) 
% 3. PRINTHOUSE
% a) no searchable value, b) apendix, c) version ISBN d) ISBN bar
% e) no rotation of landscape pages, f)urllinkcolor black
% f) A4 paper dimentions, C5 layout, shows cropping marks

% ########################################################## PAGE GEOMETRIES #

\usepackage%[showframe]
           {geometry}
           

\ifshowcrop
% Basic geometry of document
\geometry
  { headsep    =   \baselineskip
  , textwidth  = 42\baselineskip
  , textheight = 60\baselineskip
  , hmarginratio = \ifscreen 1:1\else 2:3\fi
  , vmarginratio = 2:3
  , bindingoffset = 0cm
  , onecolumn
  , a4paper 
  , layout=c5paper,
  , layouthoffset=\dimexpr(\paperwidth-\csname Gm@layoutwidth\endcsname)/2\relax
   , layoutvoffset=\dimexpr(\paperheight-\csname Gm@layoutheight\endcsname)/2\relax
  , showcrop
  }
\else
\geometry
  { headsep    =   \baselineskip
  , textwidth  = 42\baselineskip
  , textheight = 60\baselineskip
  , hmarginratio = \ifscreen 1:1\else 2:3\fi
  , vmarginratio = 2:3
  , bindingoffset = 0cm
  , onecolumn,
  , c5paper
  }
\fi

\ifshowcrop
% Geometry of the dictionary section
\newcommand\dictionarygeometry{\newgeometry
  { textwidth  = 42\baselineskip
  , textheight = 60\baselineskip
  , hmarginratio = \ifscreen 1:1\else 2:3\fi
  , vmarginratio = 3:2
  , bindingoffset = 0cm
  , twocolumn
  , a4paper 
  , layout=c5paper
  , layouthoffset=\dimexpr(\paperwidth-\csname Gm@layoutwidth\endcsname)/2\relax
   , layoutvoffset=\dimexpr(\paperheight-\csname Gm@layoutheight\endcsname)/2\relax
  , showcrop
  }}
   % Geometry of the cover pages
\newcommand\covergeometry{\newgeometry
  { textwidth  = 42\baselineskip
  , textheight = 60\baselineskip
  , hmarginratio = 1:1
  , vmarginratio = 2:3
  , bindingoffset = 0cm
  , onecolumn
  , layout=c5paper
 , layouthoffset=\dimexpr(\paperwidth-\csname Gm@layoutwidth\endcsname)/2\relax
   , layoutvoffset=\dimexpr(\paperheight-\csname Gm@layoutheight\endcsname)/2\relax
  , showcrop
  }} 
\else
% Geometry of the dictionary section
\newcommand\dictionarygeometry{\newgeometry
  { textwidth  = 42\baselineskip
  , textheight = 60\baselineskip
  , hmarginratio = \ifscreen 1:1\else 2:3\fi
  , vmarginratio = 3:2
  , bindingoffset = 0cm
  , twocolumn,
  , c5paper
  }}
   % Geometry of the cover pages
\newcommand\covergeometry{\newgeometry
  { textwidth  = 42\baselineskip
  , textheight = 60\baselineskip
  , hmarginratio = 1:1
  , vmarginratio = 2:3
  , bindingoffset = 0cm
  , onecolumn
  , c5paper
  }} 
\fi

% ################################################################# LANGUAGES #

\usepackage[icelandic, latin, czech]{babel}
\usepackage{csquotes}


% ################################################################## ENCODING #

\usepackage[utf8]{inputenc}

% Smashing uppercase letters with diacritics
% ------------------ THIS IS EVIL AND SHOULD NEVER BE DONE -------------------
% Don't smash anything if we are on the first run to produce the index
\IfFileExists{figures.idx}{
  % Czech uppercase letters with diacritics:
  \DeclareUnicodeCharacter{00C1}{\strut\smash{\' A}} % Á
  \DeclareUnicodeCharacter{010C}{\strut\smash{\v C}} % Č
  \DeclareUnicodeCharacter{010E}{\strut\smash{\v D}} % Ď
  \DeclareUnicodeCharacter{00C9}{\strut\smash{\' E}} % É
  \DeclareUnicodeCharacter{011A}{\strut\smash{\v E}} % Ě
  \DeclareUnicodeCharacter{00CD}{\strut\smash{\' I}} % Í
  \DeclareUnicodeCharacter{0147}{\strut\smash{\v N}} % Ň
  \DeclareUnicodeCharacter{00D3}{\strut\smash{\' O}} % Ó
  \DeclareUnicodeCharacter{0158}{\strut\smash{\v R}} % Ř
  \DeclareUnicodeCharacter{0160}{\strut\smash{\v S}} % Š
  \DeclareUnicodeCharacter{0164}{\strut\smash{\v T}} % Ť
  \DeclareUnicodeCharacter{00DA}{\strut\smash{\' U}} % Ú
  \DeclareUnicodeCharacter{016E}{\strut\smash{\r U}} % Ů
  \DeclareUnicodeCharacter{00DD}{\strut\smash{\' Y}} % Ý
  \DeclareUnicodeCharacter{017D}{\strut\smash{\v Z}} % Ž
  % Icelandic uppercase letters with diacritics: ÁÉÍÓÚÝ (ÐÞÆ) and
  \DeclareUnicodeCharacter{00D6}{\strut\smash{\" O}} % Ö
}{}
% ------------------ THAT WAS EVIL AND SHOULD NEVER BE DONE ------------------

% Non-breakable space
\DeclareUnicodeCharacter{00A0}{~}

% Non-breakable hyphen
\newcommand*{\nobreakhyphen}{\mbox{-}}
\DeclareUnicodeCharacter{2011}{\nobreakhyphen}


% ############################################################ FONTS, SYMBOLS #

\usepackage[T1]{fontenc}

\usepackage{tgpagella}
\usepackage[scaled=0.88]{helvet}      % relative scale of the two fonts

\def\phvfamily{\fontfamily{phv}\selectfont}

%\usepackage{stmaryrd}
% Let's keep this minimal instead of using the package:
\DeclareSymbolFont{stmry}{U}{stmry}{m}{n}
%\SetSymbolFont{stmry}{bold}{U}{stmry}{b}{n}
\DeclareMathSymbol\shortrightarrow\mathrel{stmry}{"01}
\DeclareMathSymbol\shortuparrow\mathrel{stmry}{"02}

% Methinks these are not necessary
%\usepackage{dingbat}
%\usepackage{manfnt}
%\usepackage{latexsym}

\usepackage{pifont}




% #################################################################### LAYOUT #

% \normalsize should be {8pt}{9.6pt}
\def\HUGE {\fontsize{23.887872pt}{3\baselineskip}\selectfont}
\def\Huge {\fontsize{19.906560pt}{3\baselineskip}\selectfont}
\def\huge {\fontsize{16.588800pt}{3\baselineskip}\selectfont}
\def\LARGE{\fontsize{13.824000pt}{2\baselineskip}\selectfont}
\def\Large{\fontsize{11.520000pt}{2\baselineskip}\selectfont}
\def\large{\fontsize{ 9.600000pt}{2\baselineskip}\selectfont}

% Two columns layout ruler
\setlength\columnsep    {2\baselineskip}
\setlength\columnseprule{0.4pt}

% Necessary for baseline alignment
\topskip=\baselineskip
\raggedbottom
\setlength\parskip{0pt} % it's better to avoid glue

% Temporarily suppress warnings
\hbadness=5000
\vbadness=5000

\setlength\emergencystretch{17pt}

% Allow smaller emergencystretch in several cases
\newenvironment{xtolerant}[2]{%
  \par
  \ifx\empty#1\empty\else\tolerance=#1\relax\fi
  \ifx\empty#2\empty\else\emergencystretch=#2\relax\fi
}{%
  \par
}

% ######################################################### ASSORTED PACKAGES #

\usepackage{tikz}
\usetikzlibrary{calc}

\usepackage{tipa}                      % IPA phonetics

\usepackage{fourier-orns}              % ornaments
\usepackage{amsmath}                   % non-breakable dash
\usepackage{hyphenat}                  % no hyphen in abbreviations

\usepackage{hanging}
\usepackage{fix2col}
\usepackage{fixltx2e}                  % subscript
\usepackage{pagecolor}
\usepackage{afterpage}                 % Used in coloring trick for covers
\usepackage{pdfpages}

\usepackage{multirow}
\usepackage{multicol}
\usepackage{rotating}

\usepackage{etoolbox} % http://ctan.org/pkg/etoolbox
%\usepackage{newunicodechar} 		% new definition of middle dot
\usepackage[ISBN=978-80-260-2385-2,SC0]{ean13isbn}

\usepackage{pgffor}

\usepackage{newunicodechar}

\newunicodechar{·}{\chejnikcdot}
\makeatletter
\newcommand{\chejnikcdot}{%
  \ifnum\lastpenalty=10042
    \kern-0.06667em\relax % dots should always be kerned
  \fi
  \textperiodcentered
  \@ifnextchar|{\kern\chejnikkern\relax}{\penalty10042 }%
}

\newcommand{\closeupdotbar}{%
  \renewcommand\chejnikkern{-0.06667em}%
}
\newcommand\chejnikkern{0pt} % default

% Middle dot + | little bit closer, nicer
%\newunicodechar{·}{\advancedcdot}
%\makeatletter
%\newcommand*\advancedcdot{\textperiodcentered
%  \@ifnextchar|{\hspace{-0.03667em}}{\checkifcdotnext}}
%\newcommand*\checkifcdotnext[2]
%  {\def\tmpa{·}\def\tmpb{#1#2}\ifx\tmpa\tmpb\hspace{-0.06667em}\fi#1#2}
%\makeatother

% ############################################################ BASELINE GRID #

\usepackage{eso-pic}
\usepackage{tikzpagenodes}

% Command to draw a baseline grid
\newcommand\drawbaselinegrid{%
  \begin{tikzpicture}[overlay,remember picture]
    \draw [red!30!white, ultra thin, dashed]
      (0,0) grid [ ystep  = \baselineskip, xstep = \textwidth
                 , shift  = (current page text area.north west)
                 , yshift = -\dp\strutbox
                 ] ++(\textwidth,-\textheight);
    \draw [red!30!white, thin]
      (0,0) grid [ step = \baselineskip, xstep = \textwidth
                 , shift=(current page text area.north)
                 ] ++(0.5\textwidth,-\textheight)
      (0,0) grid [ step = \baselineskip, xstep = \textwidth
                 , shift=(current page text area.north)
                 ] ++(-0.5\textwidth,-\textheight);
    \draw[black!10!white]
      (current page text area.north west)
        rectangle (current page text area.south east);
  \end{tikzpicture}}

% Draw a baseline grid on every page
% \AddToShipoutPicture{\drawbaselinegrid}


% ################################################################## METADATA #

\title{%
  \textbf{Islandsko-český studijní slovník}%
  \thanks{Tato kniha byla vytvořena v \LaTeX{}u pod Ubuntu 14.04.%
    Poděkování patří všem autorům, kteří publikují pod svobodnými licencemi.}}
\author{Aleš Chejn, Ján Zaťko, Jón Gíslason}
\date{říjen 2011}


% ############################################################## BIBLIOGRAPHY #

\usepackage[style=authoryear]{biblatex}
\addbibresource{slovnik.bib}


% ##################################################################### INDEX #

\usepackage{imakeidx}

% Index of authors of photographs
\makeindex
  [ name = figures
  , title = {Seznam autorů fotografií a ilustrací}
  , intoc
  , columnseprule
  , columnsep = \columnsep
  , noautomatic                       % We handle this in make.sh
  ]

% icelandic rules start
\usepackage{filecontents}
\def\mygroup#1{{\bfseries#1}}
\begin{filecontents*}{style.xdy}
;; style.xdy
(markup-letter-group :open-head "~n\mygroup{" :close-head "}")
\end{filecontents*}
% icelandic rules end


% ############################################################### PAGE STYLES #

\usepackage{fancyhdr}

% -------------------------------------- BASIC PAGE STYLE
\fancypagestyle{plain}{
  \fancyhf\relax                       % Clear header/footer
  \renewcommand \headrulewidth {0.0pt} % No header rule
  \renewcommand \footrulewidth {0.0pt} % No footer rule
  \fancyfoot[C]{\thepage}              % Page number in footer, centred
}

% -------------------------------------- INDEX PAGE STYLE  
%\fancypagestyle{indexstyle}{
%  \fancyhf\relax                       % Clear header/footer
%  \renewcommand \headrulewidth {0.4pt} % Thin header rule
%  \renewcommand \footrulewidth {0.0pt} % No footer rule
%  \fancyhead[C]{\thepage}              % Page number in footer, centred
%}

% -------------------------------------- DICTIONARY PAGE STYLE
\fancypagestyle{myheadings}{
  \fancyhf\relax                       % Clear header/footer
  \renewcommand \headrulewidth {0.4pt} % Thin header rule
  \fancyhead[CO,CE]{\thepage}          % Page number in header, centred
  % NOTE: the page numbers will be printed when the dictionary is ready
  \fancyhead[LE,LO]{\phvfamily\bfseries\rightmark}
  \fancyhead[RE,RO]{\phvfamily\bfseries\leftmark}
}

\pagestyle{plain}  

% ##################################################################### LISTS #

\usepackage{expdlist}
\usepackage{enumitem}

\setlist{nosep, topsep=\baselineskip}
\setlist[itemize]{leftmargin=*,label=\scalebox{.6}{\textbullet}}


% #################################################################### FLOATS #

\usepackage{caption}
\captionsetup{labelformat=empty}

% Alter some LaTeX defaults for better treatment of figures:
% See p.105 of "TeX Unbound" for suggested values. 
% See pp. 199-200 of Lamport's "LaTeX" book for details.
% ------------------------------------- Parameters for ALL pages:
\renewcommand \topfraction    {0.9}   %   max fraction of floats at top
\renewcommand \bottomfraction {0.8}   %   max fraction of floats at bottom
% ------------------------------------- Parameters for TEXT (non float) pages:
\setcounter{topnumber}    {1}
\setcounter{bottomnumber} {1}
\setcounter{totalnumber}  {1}         %   2 may work better
\setcounter{dbltopnumber} {2}         %   for 2-column pages
\renewcommand \dbltopfraction   {0.9} %   fit big float above 2-col. text
\renewcommand \textfraction     {0.07}%   allow minimal text w. figs
% ------------------------------------- Parameters for FLOAT (non text) pages:
\renewcommand \floatpagefraction{0.7} %   require fuller float pages
% N.B.: floatpagefraction MUST be less than topfraction !!
\renewcommand{\dblfloatpagefraction}{0.7} % require fuller float pages
% remember to use [htp] or [htpb] for placement

\makeatletter
\setlength{\@fptop}{0pt}
\setlength{\@fpbot}{0pt plus 1fil}
\makeatother

% This kills spacing around floats
\setlength \intextsep        {0pt}
\setlength \floatsep         {0pt}
\setlength \textfloatsep     {0pt}
\setlength \dbltextfloatsep  {0pt}
\setlength \dblfloatsep      {0pt}

% This kills spacing around captions
\setlength \abovecaptionskip {0pt}
\setlength \belowcaptionskip {0pt}


% #################################################################### TABLES #

\usepackage{booktabs}
\usepackage{tabularx}
\usepackage{longtable}
\usepackage{ltxtable}
\usepackage{lscape}

\ifPDFrotate
\else
\makeatletter			       
\renewcommand\PLS@Rotate[1]{}          % this command prevents pdflscape 
\makeatother                           % to rotate pages
\fi

%\LetLtxMacro\oldtabularx\tabularx
%\LetLtxMacro\oldendtabularx\endtabularx
%\renewenvironment{tabularx}[1]
%  {\oldtabularx{\columnwidth}{#2}}
%  {\oldendtabularx}

\newcommand\LTXfw[1]{\LTXtable{\columnwidth}{#1}}

%%% \LetLtxMacro\oldtabularx\tabularx
%%% \LetLtxMacro\oldendtabularx\endtabularx
%%% \newenvironment{longtabularx}[2]
%%%   {\oldtabularx{#1}{#2}}
%%%   {\oldendtabularx}

% ######################################################### czech babel vs cmidrule
% ###################################### http://tex.stackexchange.com/a/112001/8380
\begingroup
  \makeatletter
  \catcode`\-=\active
  \AtBeginDocument{
  \def\@@@cmidrule[#1-#2]#3#4{\global\@cmidla#1\relax
    \global\advance\@cmidla\m@ne
    \ifnum\@cmidla>0\global\let\@gtempa\@cmidrulea\else
    \global\let\@gtempa\@cmidruleb\fi
    \global\@cmidlb#2\relax
    \global\advance\@cmidlb-\@cmidla
    \global\@thisrulewidth=#3
    \@setrulekerning{#4}
    \ifnum\@lastruleclass=\z@\vskip \aboverulesep\fi
    \ifnum0=`{\fi}\@gtempa
    \noalign{\ifnum0=`}\fi\futurenonspacelet\@tempa\@xcmidrule}
  }
\endgroup

% macro that supress cmidrule at page brake
% http://tex.stackexchange.com/questions/110548/using-longtable-with-cmidrule-page-breaking-issues
\makeatletter
\def\LT@output{%
  \ifnum\outputpenalty <-\@Mi
    \ifnum\outputpenalty > -\LT@end@pen
      \LT@err{floats and marginpars not allowed in a longtable}\@ehc
    \else
      \setbox\z@\vbox{\unvbox\@cclv}%
      \ifdim \ht\LT@lastfoot>\ht\LT@foot
        \dimen@\pagegoal
        \advance\dimen@-\ht\LT@lastfoot
        \ifdim\dimen@<\ht\z@
          \setbox\@cclv\vbox{\unvbox\z@
\copy\LT@foot\vss}%
          \@makecol
          \@outputpage
          \setbox\z@\vbox{\box\LT@head}%
        \fi
      \fi
      \global\@colroom\@colht
      \global\vsize\@colht
      \vbox
        {\unvbox\z@
\box\ifvoid\LT@lastfoot\LT@foot\else\LT@lastfoot\fi}%
    \fi
  \else
    \setbox\@cclv\vbox{\unvbox\@cclv

\setbox0\lastbox
\ifdim\ht0=0.29999pt \unskip % hope that was a cline we threw away
\else
\nointerlineskip
\box0     % put it back, whatever it was
\fi

\copy\LT@foot\vss}%
    \@makecol
    \@outputpage
      \global\vsize\@colroom
    \copy\LT@head\nobreak
  \fi}

\makeatother

% These values ensure that for top, mid and bottom rules
% space above + rule width + space below = baselineskip
\setlength \abovetopsep    {0.43\baselineskip}
\setlength \heavyrulewidth {0.10\baselineskip}
\setlength \belowbottomsep {0.43\baselineskip}
\setlength \aboverulesep   {0.47\baselineskip}
\setlength \lightrulewidth {0.06\baselineskip}
\setlength \belowrulesep   {0.47\baselineskip}

%\usepackage{letltxmacro}

%%% \LetLtxMacro\oldtabular\tabular
%%% \LetLtxMacro\oldendtabular\endtabular
%%% \renewenvironment{tabular}[2][t]
%%%   {\oldtabular[#1]{#2}}
%%%   {\oldendtabular}

%%% \LetLtxMacro\oldtable\table
%%% \LetLtxMacro\oldendtable\endtable
%%% \renewenvironment{table}[1][h]
%%%   {\leavevmode\oldtable[#1]\vspace{\dp\strutbox}}
%%%   {\vspace{-\dp\strutbox}\oldendtable}
%%%   % This aligns a [t] aligned tab-envo to baseline grid


%%% % declention and conjugation tables
%%% \usepackage{floatrow}
%%% \DeclareFloatFont{footnotesize}{\footnotesize}
%%% % "scriptsize" is defined by floatrow, "tiny" not
%%% \floatsetup[table]{font=footnotesize}

% ################################################################### FIGURES #

\usepackage{adjustbox} % loads also graphicx

\graphicspath{%
  {/home/chejnik/Dokumenty/hvalur.org/images/biolib/full/}%
  {/home/chejnik/Dokumenty/hvalur.org/images/uploaded_files/}}

\newlength\dicfigheight
\newlength\dicfigskip

\newcommand\dicFigure[3]{ % {filename}{caption}{index entry}
  \setlength\fboxsep{0pt}\setlength\fboxrule{0.5pt}
  \def\dicfigbox{\fbox{%
    \adjincludegraphics [ max height = 0.8\columnwidth
                        , max width  = 0.8\columnwidth ] {#1}}}
  \setlength\abovecaptionskip {\dp\strutbox}
  \setlength\belowcaptionskip {\ht\strutbox} 
  \setlength\dicfigheight {\heightof\dicfigbox+\fboxrule}
  \setlength\dicfigskip
    {\baselineskip*\numexpr1+\dicfigheight/\baselineskip\relax-\dicfigheight}
  \begin{figure}[ht]
    \centering\vspace*{\dicfigskip}\dicfigbox\caption{#2}\index[figures]{#3}
  \end{figure}
}


% #################################################################### COLORS #

\usepackage{color}

% http://www.colorschemer.com/online.html
% main color - grammar color
% #66023C
\definecolor {darkgreen}          {rgb} {0.40, 0.01, 0.24}

\definecolor {royalazure}         {rgb} {0.00, 0.22, 0.66}
\definecolor {brown}              {rgb} {0.40, 0.01, 0.24}

% COLORS FOR THUMB INDEXES
\definecolor {babyblueeyes}       {rgb} {0.63, 0.79, 0.95}
\definecolor {unitednationsblue}  {rgb} {0.36, 0.57, 0.90}
\definecolor {blue(ryb)}          {rgb} {0.01, 0.28, 1.00}
\definecolor {darkblue}           {rgb} {0.00, 0.00, 0.55}
\definecolor {screamingreen}      {rgb} {0.46, 1.00, 0.44}
\definecolor {limegreen}          {rgb} {0.20, 0.80, 0.20}
\definecolor {islamicgreen}       {rgb} {0.00, 0.56, 0.00}
\definecolor {upforestgreen}      {rgb} {0.00, 0.27, 0.13}
\definecolor {icterine}           {rgb} {0.99, 0.97, 0.37}
\definecolor {orange(colorwheel)} {rgb} {1.00, 0.50, 0.00}
\definecolor {orange-red}         {rgb} {1.00, 0.27, 0.00}
\definecolor {oucrimsonred}       {rgb} {0.60, 0.00, 0.00}
\definecolor {cottoncandy}        {rgb} {1.00, 0.74, 0.85}
\definecolor {orchid}             {rgb} {0.85, 0.44, 0.84}
\definecolor {vividcerise}        {rgb} {0.85, 0.11, 0.51}
\definecolor {patriarch}          {rgb} {0.50, 0.00, 0.50}

\definecolor {title}              {RGB} {  16,   13,   32}
\definecolor {golden}             {RGB} { 241,  184,   45}
\definecolor {freqstar}           {RGB} {  91,    3,   99}
\definecolor {newdev}             {RGB} {   3,   11,   99}

\definecolor {color1}             {RGB} { 182,   86,    0}
\definecolor {color2}             {RGB} { 143,    9,    6}
\definecolor {color3}             {RGB} {   3,   23,  118}
\definecolor {color4}             {RGB} {   0,   82,  168}
\definecolor {color5}             {RGB} {   0,   85,  142}
\definecolor {color6}             {RGB} {   0,  115,  162}
\definecolor {color7}             {RGB} {  34,  146,  186}
\definecolor {color8}             {RGB} {  40,  159,  153}
\definecolor {color9}             {RGB} {   0,  125,  111}
\definecolor {color10}            {RGB} {   4,  107,   60}
\definecolor {color11}            {RGB} {  71,  134,   81}
\definecolor {color12}            {RGB} { 109,  134,   42}
\definecolor {color13}            {RGB} { 205,  194,   18}
\definecolor {color14}            {RGB} { 204,  162,   24}

%\definecolor {color1}             {RGB} { 239,  222,  205}
%\definecolor {color2}             {RGB} { 251,  206,  177}
%\definecolor {color3}             {RGB} { 251,  185,  142}
%\definecolor {color4}             {RGB} { 248,  166,  112}
%\definecolor {color5}             {RGB} { 216,  231,  246}
%\definecolor {color6}             {RGB} { 186,  215,  244}
%\definecolor {color7}             {RGB} { 152,  197,  243}
%\definecolor {color8}             {RGB} { 118,  181,  245}
%\definecolor {color9}             {RGB} { 181,  249,  185}
%\definecolor {color10}            {RGB} { 160,  248,  167}
%\definecolor {color11}            {RGB} { 128,  247,  136}
%\definecolor {color12}            {RGB} {  87,  248,   98}
%\definecolor {color13}            {RGB} { 250,  241,  138}
%\definecolor {color14}            {RGB} { 250,  240,  111}
%\definecolor {color15}            {RGB} { 244,  239,   85}
%\definecolor {color16}            {RGB} { 245,  239,   57}

\definecolor {lightgrey}          {RGB} { 105,  105,  105}
\definecolor {darkgreen_real}     {rgb} {0.00, 0.50, 0.00}


% ############################################################# THUMB INDEXES #

% Thumb indexes' colors
%\newcommand\BoxColor{%
%  \ifcase\theletternum darkgreen!30\or babyblueeyes\or unitednationsblue\or blue(ryb)\or screamingreen\or limegreen\or islamicgreen\or upforestgreen\or icterine\or orange(colorwheel)\or orange-red%
%  \or oucrimsonred\or cottoncandy\or orchid\or vividcerise\or patriarch\or babyblueeyes\or unitednationsblue\or blue(ryb)\or screamingreen\or limegreen\or islamicgreen\or upforestgreen\or icterine\or orange(colorwheel)\or orange-red%
%  \or oucrimsonred\or cottoncandy\or orchid\or vividcerise\or patriarch\else darkgreen!30\fi}
  
% Thumb indexes' colors
%\newcommand\BoxColor{%
%  \ifcase\theletternum darkgreen!30\or color1\or color2\or color3\or color4\or color5\or color6\or color7\or color8\or color9\or color10%
%  \or color11\or color12\or color13\or color14\or color1\or color2\or color3\or color4\or color5\or color6\or color7\or color8\or color9%
%  \or color10\or color11\or color12\or color13\or color14\or color1\or color2\or color3\or color4\or color5\or color6\else color7\or color8\or color9\or color10\or color11\or color12\fi}

%\def\BoxColor#1{darkgreen!\the\numexpr103-#1-#1-#1\relax!blue}

\newcommand\BoxColor[1]{%
\ifcase#1 darkgreen!30\or color13\or color14\or color1\or color2\or color3\or color4\or color5\or color6\or color7\or color8\or color9\or color10%
\or color11\or color12\or color13\or color14\or color1\or color2\or color3\or color4\or color5\or color6\or color7\or color8\or color9%
\or color10\or color11\or color12\or color13\or color14\or color1\or color2\or color3\or color4\or color5\or color6\else color7\or color8\or color9\or color10\or color11\or color12\fi}

\ifshowcrop
%%%%%%%%%%%%%%%%%%%%%%%%%%%%%%%%%%%%%%%%%%%%%%%%%%%%%%%%%%%%%%%%%%%%%%%%%%%%%%%%

\newcounter{letternum}
\newcounter{lettersum}                 \setcounter{lettersum}{34}
% these margins separate top thumb from top of page and last thumb from
% bottom of page. BEWARE, it's from top of the (cut out) page, not the folio!
\newlength{\thumbtopmargin}            \setlength{\thumbtopmargin}{1cm}
\newlength{\thumbbottommargin}         \setlength{\thumbbottommargin}{1cm}
\newlength{\thumbwidth}                \setlength{\thumbwidth}{2.20cm}
\newlength{\thumbheight}
\pgfmathsetlength{\thumbheight}%
  {(\csname Gm@layoutheight\endcsname-\thumbtopmargin-\thumbbottommargin)/\value{lettersum}}
% inner xsep of letter from box margin
\newlength{\thumbsep}                  \setlength{\thumbsep}{.30cm}
% overlap over cutting line
\newlength{\thumboverlap}              \setlength{\thumboverlap}{3cm}

\tikzset{
  thumb/.style={
    text=white,
    outer sep=0pt,
    inner xsep=\thumbsep,
    minimum height=\thumbheight,
    text width=\thumbwidth-2\thumbsep,
    font=\sffamily\bfseries,},
  thumb east/.style={thumb,
    xshift=\dimexpr(+\thumboverlap-\paperwidth+\csname Gm@layoutwidth\endcsname)/2\relax,
    align=left,anchor=north east,},
  thumb west/.style={thumb,
    xshift=\dimexpr(-\thumboverlap+\paperwidth-\csname Gm@layoutwidth\endcsname)/2\relax,
    align=right,anchor=north west,},
}

\def\thumbnew{}
\def\thumbold{}
\usepackage{everypage}
\AddEverypageHook{\if\relax\thumbnew\relax\xdef\thumbnew{\thumbold}\fi}

\def\ethumbs#1,#2\relax%
  {\if\relax#1\relax\else\eventhumb{#1}\fi%
   \if\relax#2\relax\else\ethumbs#2\relax\fi%
   \gdef\thumbnew{}%
   \gdef\thumbold{#1,}}

\def\othumbs#1,#2\relax%
  {\if\relax#1\relax\else\oddthumb{#1}\fi%
   \if\relax#2\relax\else\othumbs#2\relax\fi%
   \gdef\thumbold{#1,}%
   \gdef\thumbnew{}}

\newcommand{\drawthumb}[2]{%
  \begin{tikzpicture}[remember picture, overlay]
    \node [thumb #2, fill=\BoxColor{#1}]
       at ($(current page.north #2)-%
            (0,\csname Gm@layoutvoffset\endcsname+\thumbtopmargin-\thumbheight+%
               #1*\thumbheight)$) {\csname Let#1\endcsname};
   \end{tikzpicture}}



%%%%%%%%%%%%%%%%%%%%%%%%%%%%%%%%%%%%%%%%%%%%%%%%%%%%%%%%%%%%%%%%%%%%%%%%%%%%%%%%
\else
% THUMB INDEXES
% new counter to hold the current number of the letter to determine the vertical position
\newcounter{letternum}
% newcounter for the sum of all letters to get the right height of a box
\newcounter{lettersum}
\setcounter{lettersum}{34}
% some margin settings
\newlength{\thumbtopmargin}
\setlength{\thumbtopmargin}{\ifshowcrop 3cm\else 1.5cm\fi} %1cm
\newlength{\thumbbottommargin}
\setlength{\thumbbottommargin}{\ifshowcrop 6cm\else 2cm\fi} %2.5cm
% calculate the box height by dividing the page height
\newlength{\thumbheight}
\pgfmathsetlength{\thumbheight}{%
(\paperheight-\thumbtopmargin-\thumbbottommargin)%
/%
\value{lettersum}
}
% box width
\newlength{\thumbwidth}
\setlength{\thumbwidth}{\ifshowcrop 2cm\else 0.5cm\fi} %0.5cm
% style the boxes
\tikzset{
thumb/.style={
   text=white,
   minimum height=\thumbheight,
   text width=\thumbwidth,
   outer sep=0pt,
   font=\sffamily\bfseries,
 }
 }

\def\thumbnew{}
\def\thumbold{}
\usepackage{everypage}
\AddEverypageHook{\if\relax\thumbnew\relax\xdef\thumbnew{\thumbold}\fi}

\def\ethumbs#1,#2\relax{\if\relax#1\relax\else\eventhumb{#1}\fi%
                        \if\relax#2\relax\else\ethumbs#2\relax\fi%
                        \gdef\thumbnew{}%
                        \gdef\thumbold{#1,}%
}

\def\othumbs#1,#2\relax{\if\relax#1\relax\else\oddthumb{#1}\fi%
                        \if\relax#2\relax\else\othumbs#2\relax\fi%
                        \gdef\thumbold{#1,}%
                        \gdef\thumbnew{}%
}

\newcommand{\drawthumb}[2]{%
  % see pgfmanual.pdf for more information about this part
  \begin{tikzpicture}[remember picture, overlay]
    \node [thumb, fill=\BoxColor{#1}, text centered, anchor=north #2]
       at ($(current page.north #2)-%
%            (0,\thumbtopmargin+\value{letternum}*\thumbheight)%
            (0,\thumbtopmargin+#1*\thumbheight)$) {\csname Let#1\endcsname};
   \end{tikzpicture}}
\fi

\newcommand{\oddthumb} [1]{\drawthumb{#1}{west}}
\newcommand{\eventhumb}[1]{\drawthumb{#1}{east}}

\newcommand{\lettergroup}[1]%
  {\refstepcounter{letternum}%
   \expandafter\gdef\csname Let\theletternum\endcsname{#1}%
   \xdef\thumbnew{\theletternum,\thumbnew}%
   \fancyhead[LO]{\phvfamily\bfseries\rightmark%
     \expandafter\ethumbs\thumbnew\relax\relax\relax}%
   \fancyhead[RE]{\phvfamily\bfseries\leftmark%
     \expandafter\othumbs\thumbnew\relax\relax\relax}}
     
% ############################################################## LOCALIZATION #

% renames the index name
%\addto\captionsczech{%
%  \renewcommand\indexname{{Seznam autorů fotografií a ilustrací}}}

% renames the contents name
\addto\captionsczech{%
  \renewcommand\contentsname{Obsah}}

\renewcommand \contentsname {Obsah}
\renewcommand \chaptername  {Kapitola}


% ################################################################ SECTIONING #
\usepackage[explicit]{titlesec}

\setlength{\beforetitleunit}{\baselineskip}
\setlength{\aftertitleunit} {\baselineskip}

\titlespacing*{\chapter}   {0em}{*0}{*3} % Remembere there's a topskip too
\titlespacing*{\section}   {0em}{*1}{*2}
\titlespacing*{\subsection}{0em}{*1}{*2}

\newcommand\ghostdrop[2]{%
  \raisebox{-#1\baselineskip}[0pt][0pt]{#2}}

\titleformat {\chapter} {} {} {0pt} % {#1}
  {\ghostdrop{1.0}{\parbox[t]{\columnwidth}
     {\LARGE\bfseries\centering\MakeUppercase{#1}}}}

\def\longchapterskip{\vspace{2\baselineskip}}
\def\longsectionskip{\vspace{2\baselineskip}}

\titleformat {\section} {\bfseries} {} {0pt}
  {\ghostdrop{1.0}{\parbox[t]{\columnwidth}
     {\Large\bfseries\centering#1}}}

\titleformat {\subsection} {\bfseries} {} {0pt}
  {\ghostdrop{1.0}{\hspace{1em}\parbox[t]{\columnwidth-1em}
     {\large\bfseries#1}}}


\setlength\multicolsep{0pt}
%\BeforeBeginEnvironment{multicols}{\vspace{-\dp\strutbox}}


% ######################################################## TABLE OF CONTENTS #
\usepackage{titletoc}

\titlecontents {chapter} [0pt] {\addvspace\baselineskip\bfseries}
  {} {} {\titlerule*[2\baselineskip]{}\contentspage}

\titlecontents {section} [\baselineskip] {} {} {}
  {\titlerule*[1\baselineskip]{.}\contentspage}


% #################################################### ASSORTED CUSTOM MACROS #

\def\textIS#1{\foreignlanguage{icelandic}{#1}}
\def\textCS#1{\foreignlanguage{czech}{#1}}
\def\textLA#1{\foreignlanguage{latin}{#1}}

\usepackage{accsupp}

\makeatletter                                                                                                          
  \def\ACCSUPP@bdc{\special{pdf:code \ACCSUPP@span BDC}}                                                               
  \def\ACCSUPP@emc{\special{pdf:code EMC}}                                                                             
\makeatother

\usepackage{pgffor}

\newcommand\blspace[1][1]{\vspace{#1\baselineskip}}

\usepackage{xparse}


\newcommand \dicsymFrequent  {\raisebox{-.2ex}{\color{darkgreen}\ding{167}}}
\newcommand \dicsymSee       {$\shortrightarrow$}
\newcommand \dicsymCompare   {$\shortuparrow$}
\newcommand \dicsymIdiom     {$\diamondsuit$}
\newcommand \dicsymExampleIS {$\triangleright$}
\newcommand \dicsymExampleCS {\guilsinglright}
\newcommand \dicsymProverb   {{\color{newdev}\ding{96}}}
\newcommand \dicsymPhoto
  {\includegraphics[keepaspectratio,height=0.65em]{photo_icon}}

%\NewDocumentCommand\dicTerm{mo}
 % {{\phvfamily\bfseries\textIS{#1}}\IfValueT{#2}{, \phvfamily\bfseries\textIS{#2}}}

\NewDocumentCommand\dicEntry{o}
  {\par\hangpara{0.25\baselineskip}{1}%
   \IfValueT{#1}{\markboth{#1}{#1}}%
    \ifPDFsearchvalue
    \settowidth{\eyja}{#1}%
  \makebox[\eyja]{\color{white}#1}\hspace{-\eyja}%
  \fi
   \ignorespaces}
   
   \NewDocumentCommand\dicEntryGuide{o}
  {\par\hangpara{0.29\baselineskip}{1}%
   \IfValueT{#1}{\markboth{#1}{#1}}%
    \ifPDFsearchvalue
    \settowidth{\eyja}{#1}%
  \makebox[\eyja]{\color{white}#1}\hspace{-\eyja}%
  \fi
   \ignorespaces}
   
\newcommand\dicSubEntry
  {\\\hspace*{-\hangindent}}

\NewDocumentCommand\textAlt{mm}
  {\BeginAccSupp{method=pdfstringdef,ActualText={#2}}#1\EndAccSupp{}}

\newcounter{terms}

\NewDocumentCommand\dicStyleIntTerm{m}{\textIS{{\phvfamily\bfseries#1}}}
\NewDocumentCommand\dicStyleExtTerm{m}{\textIS{\textit{#1}}}

   \NewDocumentCommand\dicTerm{mo}
  {{\closeupdotbar \phvfamily\bfseries\textIS{#1}}\IfValueT{#2}{, \phvfamily\bfseries\textIS{#2}}}

  \newcommand\dicTermList[1]
  {\foreach \t [count=\n] in {#1}{\ifnum\n=1\else, \fi\dicTerm{\t}}}
  
\def\dicIPA#1{\textipa{[#1]}}

\NewDocumentCommand\dicPos{smo}
  {{\color{darkgreen}\textbf{\small#2}%
   \IfValueT{#3}{\nobreak\textsubscript{#3}}}}

%\NewDocumentCommand\dicFlx{smo}
 % {{\color{darkgreen}\footnotesize#1\IfValueT{#2}{\textsubscript{#2}}}}

  \NewDocumentCommand\dicFlx{smo}
  {{\color{darkgreen}\footnotesize#2%
   \IfValueT{#3}{\nobreak\textsubscript{#3}}}}
  
\newcommand\dicExampleCS[1]{\dicsymExampleCS~\textit{#1}}
\newcommand\dicExampleIS[1]{\dicsymExampleIS~\textIS{\textit{#1}}}

\NewDocumentCommand\dicLink{sm}
  {\IfBooleanF#1{\dicsymSee\nobreak\hspace{.08667em plus .03333em}}\nobreak\hspace{.03667em}\dicTerm{#2}}
\NewDocumentCommand\dicSynonym{sm}
  {(\IfBooleanF#1{\dicsymSee\nobreak\hspace{.08667em plus .03333em}}\nobreak\hspace{.03667em}\textIS{\textit{#2}})}
\NewDocumentCommand\dicAntonym{sm}
  {(\IfBooleanF#1{\dicsymCompare\nobreak\hspace{.08667em plus .03333em}}\nobreak\hspace{.03667em}\textIS{\textit{#2}})}
\NewDocumentCommand\dicCompare{sm}
  {(\IfBooleanF#1{\dicsymCompare\nobreak\hspace{.08667em plus .03333em}}\nobreak\hspace{.03667em}\textIS{\textit{#2}})}


\newcommand\dicPhraseIS[1]{\textIS{\textbf{#1}}}

\NewDocumentCommand\dicIdiom{mo}
  {\dicSubEntry{\color{newdev}\dicStyleIntTerm{#1} %
     \IfValueTF{#2}{+ \dicStyleIntTerm{#2} \markboth{#1 + #2}{#1 + #2}}{\markboth{#1}{#1}}}\dicsymIdiom}
     
\newcommand\dicProverb
  {\dicSubEntry\dicsymProverb}



\newcommand\dicLangCat[1]{{\footnotesize#1}}
\newcommand\dicFieldCat[1]{{\footnotesize#1}}


\newcommand\dicDirectTranslationCS[1]{#1}
\newcommand\dicIndirectTranslationCS[1]{{\footnotesize#1}}

\usepackage{placeins}

\newcommand*{\dicLetter}[2]{%
  \FloatBarrier\newpage\phantomsection\pdfbookmark[1]{#1}{#2}%
  \noindent\parbox[b][9\baselineskip][c]{\columnwidth}
    {\centering\HUGE\strut\smash{\MakeUppercase{#1}~\MakeLowercase{#1}}}%
      %\par\xdef\tpd{\the\prevdepth}
      %\par\prevdepth\tpd%
  \expandafter\lettergroup{\MakeLowercase{#1}}}

  
%tabulky
\newcommand{\specialcell}[2][l]{%
\begin{tabular}[#1]{@{}l@{}}#2\end{tabular}}


%\newcommand{\devisionguide}[1]{\hspace*{0.3cm}{{{{\foreignlanguage{icelandic}{\color{newdev}{\phvfamily{\textbf{#1}}}}}}}}}
%\newcommand{\devisionguideindent}[1]{\hspace*{-0.1cm}{{{{\foreignlanguage{icelandic}{\color{newdev}{\phvfamily{\textbf{#1}}}}}}}}}

%\newcommand\n{$n$\nobreakdash-\hspace{0pt}} %nonbreakable dash in grammar endings
%\def \nobreakseq {\nobreak \hskip 0pt \hbox} %nolinebreak in case like (-s, -)

% rule line
\newcommand{\HRule}{\rule{\linewidth}{0.1mm}}

% width of pdf found box
\newlength\eyja

%\newcommand\strutsmash[1]{\strut\smash{#1}}


% Dictionary symbols


% Command synonyms
%\let\freqstar\dicsymFrequent
%\let\photoicon\dicsymPhoto

\newcommand{\vstretch}[1]{\vspace*{\stretch{#1}}}

\newcommand*{\addthin}{\hskip0.16667em\relax}
\newcommand*{\addthinS}{\hskip0.06667em\relax}
\newcommand*{\addthinM}{\hskip0.10007em\relax}
\newcommand*{\addthinSS}{\hskip0.00007em\relax}
% #################################################################### COVERS #

\def\dictnameCZ{islandsko-český studijní slovník}
\def\dictnameIS{íslensk-tékknesk stúdentaorðabók}

\newcommand*{\frontcoverimages}{%
  ds_image_blaberjalyng_0_2.jpg,
  ds_image_lundi_0_1.jpg,
  ds_image_hestur_0_1.jpg,
  ds_image_hrafnafifa_0_2.jpg,
  ds_image_baldursbra_0_1.jpg,
  20948.jpg,
  ds_image_hvalur_0_1.jpg,
  ds_image_blodberg_0_1.jpg}

\newcommand*{\backcoverimages}{%
  ds_image_spoi_0_1.jpg,
  ds_image_mosalyng_0_1.jpg,
  ds_image_blaklukkulyng_0_1.jpg,
  20787.jpg,
  ds_image_ufsi_0_2.jpg,
  %ds_image_gullbra_0_1.jpg,
  ds_image_vatnsberi_0_1.jpg,
  ds_image_islandsfifill_0_2.jpg,
  21137.jpg}

\newcommand\makecoverwith[1]{
 \pagecolor{title}\afterpage\nopagecolor
 \begin{center}
 \vspace*{0.3cm}%{\fill}
 {\Huge\scshape\color{white}%
   \dictnameCZ\\\dictnameIS}
 \vspace*{1em}                        % I think this is necessary
 \begin{figure}[ht]\centering%
   \setlength\fboxsep{0pt}\setlength\fboxrule{0.5pt}%
   \foreach \file in #1{
     \fbox{\includegraphics[width=5.5cm]{\file}}}
 \end{figure}
 \vspace*{\fill}
 \end{center}
}


% ###################################################### HYPERREF & PDF INFO #

% unicode neccessary so that national characters in hypersetup appear ok
\usepackage[pdftex, unicode, hyperfootnotes=false]{hyperref}

% hyperlinks in black
\makeatletter
\let\Hy@linktoc\Hy@linktoc@none
\makeatother

\hypersetup
  { pdftitle={Islandsko-český studijní slovník}
  , pdfauthor={Aleš Chejn, Jón Gíslason, Ján Zaťko}
  , pdfsubject={Islandsko-český studijní slovník}
  , pdfkeywords={Islandsko-český studijní slovník%
               , islandština, čeština, slovník}
  , bookmarks=true
  , colorlinks=true
  , citecolor=black
  , urlcolor=\ifscreen darkgreen\else black\fi
  , linkcolor=black}


% ################################################################# MICROTYPE #

\usepackage[final]{microtype}

% ############################################################### HYPHENATION #

% this allows to include brackets inside hyphenation rules f.e (pro)-dis-ku-to-vat
\lccode`\(`\(
\lccode`\)`\)
\lccode`\[`\[
\lccode`\]`\]

\begin{hyphenrules}{icelandic}
\hyphenation{
  frum-semja
  burm-neskur
  ó-breytan-le-gur
  pla-tó-nskt
  skatta-fram-tal
  skrúf-gang-ur
  strand-rauð-viður
  prent-sver-ta
  grun-nur
  uppi-skrop-pa
  ger-vi-mennska
  sö-lu-stað-ur
  sý-na-gó-ga
  mark-mið
  trau-st-ur
  ís-kyggi-leg-ur
  brjó-sta-höld
  kné-krjú-pa
  nel-li-ka
  kaffi-kvörn
  margs-konar
  skrúf-gang-ur
  uppi-skrop-pa
  veð-rað-ur
  vest-ur-átt
  mis-mun-andi
  Bar-ba-do-si
  tung-u-mál
  au-ga-brún
  há-sker-pu-sjón-varp
  rúnn-stykki
  ger-vi-tungl
  yfir-borðs-menn-ska
  gæ-fu-mun-ur
  hlí-fi-skil-di
  á-vís-un
  ó-um-flýj-an-leg-ur
  san-mar-ín-sk-ur
  gæ-ða-stað-all
}
\end{hyphenrules}
\begin{hyphenrules}{czech}
\hyphenation{
  teo-rie
  mi-liardy
  (z)-od-po-věd-nost
  tan-tié-ma
  alian-ce
  pe-dia-tr(ička)
  pa-cienty
  zá-chra-nář(ka)
  dia-man-to-vá
  poe-zie
  ma-te-riál-ní
  egoi-s-ta
  egoist-ka
  nai-vi-ta
  ne-upřím-nost
  pě-ti-úhel-ník
  še-s-ti-úhel-ník
  za-šmo-dr-chat
  ar-cheo-lo-gic-ký
  ne-omlu-vi-tel-ný
  (vy)-pro-du-ko-vat
  dia-gno-s-ti-ko-vat
  teo-lož-ka
  ateist-ka
  oceá-nu
  im-pe-ria-li-s-ta
  (vy)-zdvih-nout
  ne-úměr-ný
  kon-ti-nuál-ní
  neu-t-rum
  ne-opráv-ně-ně
  vig-vam
  pro-zaic-ký
  geo-lož-ka
  seiz-mic-ká
  teo-rém
  olean-dr
  poe-ta
  (vy)-fo-to-gra-fo-vat
  (vy)-ven-ti-lo-vat
  pneu-mo-nie
  Lu-cem-bur-čan-(ka)
  ma-te-riály
  am-bi-cióz-ní
  Myan-ma(r)
  Myan-ma-řan-(ka)
  mol-dav-šti-na
  si-tua-cí
  ne-ukáz-ně-ný
  he-te-ro-trof-ní
  ne-omlu-vi-tel-ná
  ne-upřím-ný
  dvou-bliz-ný
  ne-úmy-sl-ný
  ne-ustá-lý-mi
  ne-ohle-du-pl-ný
  ne-osob-ní
  ne-oso-le-ný
  ne-uvě-ři-tel-ný
  ne-ohro-že-ný
  ne-oby-čej-ný
  ne-útul-ný
  ne-únav-ný
  ne-ože-nit
  po-ukáz-ka
  za-úto-čit
  aero-dy-na-mic-ký
  hru-bián-ství
  (pro)-di-s-ku-to-vat
  (vy)-ča-lou-nit
  (na)-oč-ko-vá-ní
  Bhú-tán-ka
  (za)-šně-ro-vat
  soud-ní
  re-vo-lu-cio-nář(ka)
  á-zer-báj-d-žán
  čtyř-hran-ný
  ze-m-dle-lý
  klou-zač-ka
  Vest-mann-ské
  tří-stran-ný
  (za)-bu-rá-ce-ní
  os-tud-ný
  pro-du-cent-(ka)
  chvá-s-tal-(ka)
  apar-theid
  asis-ten-ce
  asis-tent-(ka)
  (u)-či-nit
  (u)-za-mknout
  past-vi-na
  (za)-zá-lo-ho-vat
  Af-ri-čan-(ka)
  (za)-an-ga-žo-vá-ní
  se-rióz-ní
  opo-nent-(ka)
  pro-tiv-ník
  zvlášt-ní
  (po)-div-ný
  ho-mo-se-xuál-ní
  ko-alič-ní
  si-tua-ce
  si-tua-cí
  egoi-s-mus
  spe-ciál-ní
  spe-cia-li-zo-va-ný
  spe-cia-li-zo-va-né
  sobě-stač-ný
  sa-mo-zvo-le-ný
  sa-mo-čin-ně
 chvá-s-ta-vost
 klient-(ka)
 ku-rióz-ní
 (na)-šrou-bo-vat
 (za)-šrou-bo-vat
 krea-ti-vi-ta
 pneu-ma-ti-ka
 toa-le-ta
 so-cia-list-ka
 (vy)-dez-in-fi-ko-va-ný
 vše-uměl
 vše-uměl-(kyně)
 (pře)-škrt-nout
 (vy)-svlék-nout
 vy-učo-va-cí
 vy-učo-va-cích
 (vy)-chla-s-tat
 (vy)-te-to-vat
 pa-cient-(ka)
 ne-opod-s-tat-ně-ný
 de-se-ti-úhel-ník
 pe-rio-di-kum
 (za)-mum-lá-ní
 ne-smě-lá
 ne-dů-vě-ři-vá
 tee-na-ger-(ka)
 (z)-de-ge-ne-ro-vá-ní
 mo-ci-chti-vý
 bez-úhon-nost
 zne-uctít
 vo-dě-odol-ný
 we-bo-vý
 web-de-sign
 in-že-nýr-(ka)
 zá-pa-do-is-land-ský
 klient-(ka)
 za-uja-tý
 vy-uží-va-né
 vy-uží-vat
 (za)-re-gi-s-t-ro-vat
 (z)-kom-pri-mo-va-ný
 na-cio-na-li-s-mus
 ma-k-ro-eko-no-mic-ká
 bez-bo-lest-ný
 viol-ka
 troj-úhel-ník
 (za)-bu-rá-ce-ní
 ne-ústup-nost
 (vy)-cvi-če-ný
 zprá-va
 ge-ria-tr-(ička)
 osud-ný
 ná-pra-vy
 vstup-ní
 obec-ně
 re-pub-li-ka
 (po)-div-ný
 pa-slun-ce
 vliv-ný
 ba-ham-ský
 pe-dia-tr-(ička)
 ha-sič-(ka)
 vnuč-ka
 bro-s-kvoň
 lo-di-vod-(ka)
 mi-s-tr-(yně)
 in-že-nýr-(ka)
 dr-ži-tel-(ka)
 to-pič-(ka)
 pro-mi-nent-ní
 takt-ní
 frag-men-tár-ní
 po-zo-ru-hod-ný
 úspě-š-ný
 trans-fu-ze
 bur-žoa-zie
 pře-stro-je-ní
 oh-ni-vzdor-ný
 dis-tan-ční
 de-frau-da-ce
 fi-s-kál-ní
 zběh-lý
 kvap-ný
 střed-ní
 sa-mo-hlá-s-ka
 pře-vo-di-tel-ný
 zban-k-ro-to-va-ný
 zá-clo-na
 zná-mý
 zá-klad-ní
 na-uka
 ro-š-ťa-č-ka
 lí-bán-ky
 sou-du
 zkou-š-ka
 mlý-nek
 prá-šiv-ka
 kul-ma
 la-hod-ný
 zá-meč-ni-ce
 ne-vý-znam-ný
 ovzdu-ší
 mno-ho-úhel-ní-ko-vý
 žra-vý
 zcest-ný
 esen-ce
 (při)-nu-ce-ný
 skr-b-lic-ký
 skle-s-lý
 bez-bož-ný
 ne-opráv-ně-ný
 ošu-mě-lý
 pru-der-ní
 chvá-s-ta-vý
 sou-drž-ný
 van-dal-(ka)
 sa-bo-tér-(ka)
 (po)-div-ně
 mo-k-řad
 plyt-ký
 na-spo-du
 izo-la-ce
 le-toun
 za-han-be-ný
 ří-ze-ní
 kom-plex-ní
 pr-kén-ko
 ba-vič-(ka)
 plyt-ký
 ne-mo-tor-ně
 šéf-(ka)
 přes-čas
 zji-s-ti-tel-ný
 ta-ko-vý
 kňou-ral-(ka)
 se-me-s-tr
 po-pra-s-ka-ný
 han-ba
 vzpří-me-ný
 chys-ta-ný
 lik-vi-da-ce
 or-ga-ni-za-ce
 ri-s-kant-ní
 (za)-pís-kat
 (za)-hvíz-dat
 vrou-cí
  }
\end{hyphenrules}


% ############################################################# CZECH TYPOGRAPHY #
% https://math.feld.cvut.cz/tkadlec/ftp/text/textypo.pdf



% ################################################################# DOCUMENT #

\begin{document}
\pagenumbering{gobble}

% -------------------------------------------------------------- Front cover -

\ifmakecovers
\covergeometry
  \makecoverwith\frontcoverimages
\restoregeometry  
\fi


% ============================================================= FRONT MATTER =

\ifmakefrontmatter

% ------------------------------------------------------------------ Patitul -

\cleardoublepage

\begin{center}
  %\vspace*{2cm}{\fontsize{20}{20}\selectfont\scshape\dictnameCZ\\\dictnameIS}
  \vspace*{\stretch 1}
  {\Huge\scshape\dictnameCZ\\\dictnameIS}
  \vspace{\stretch 4}
\end{center}

% -------------------------------------------------------------------- Titul -

\cleardoublepage

\begin{center}
  %\vspace*{2cm}{\fontsize{24}{24}\selectfont\scshape\dictnameCZ\\\dictnameIS}
  \vspace*{\stretch 1}
  {\HUGE\scshape\dictnameCZ\\\dictnameIS}
  \vspace{\stretch 2}
\end{center}

% ----------------------------------------------------------------- Impresum -

\clearpage

\noindent\begin{minipage}[t][\textheight][t]{\textwidth}
Autoři \textit{Islandsko-českého studijního slovníku}:\\
  \textbf{Aleš Chejn, Ján Zaťko, Jón Gíslason}\\[\baselineskip]
\ifPDF
\else
Autor \textit{Islandské gramatiky}:\\
  \textbf{Vojtěch Kupča}\\[\baselineskip]
\fi
Poděkování:\\
  \textbf{Dorotě Nierychlewské-Chejn, Jiřině Hanzlové, Jiřině Chejnové,
          Simoně Petrásové,\\ Petru Mikešovi, Vojtěchu Kupčovi, 
          Renatě Peškové Emilsson, Amiru Mulahumici, \\ Kristýně Antonové, 
          Martině Kašparové, Lucii Brabcové, Samoriele, Adamu Kožouškovi}\\[\baselineskip]
Finanční příspěvek na výtisk exemplářů určených knihovnám\\ v České republice,
Slovenské republice a na Islandu poskytli:\\
  \textbf{Petr Mikeš, Dorota Nierychlewska-Chejn, Jiřina Chejnová, Otta Chejn}
  \ifPDF
  \\[39\baselineskip]
  \else
  \\[36\baselineskip]
  \fi
\textbf{Hvalur.org - 2016}\\
Licence: \textbf{GNU Free Documentation License}\\
Sazba: \textbf{Aleš Chejn} v prostředí  \textbf\LaTeX\\
  Zvláštní poděkování patří \textbf{Paolu al-Imk\=an\=\i\ Brasolin} a uživatelům  \textbf{\TeX\ - \LaTeX\ Stack Exchange}\\[\baselineskip]
\ifPDF
Verze: \textbf{PDF verze} \textit{Islandsko-českého studijního slovníku} vhodná pro \textbf{\ifscreen offline prohlížení a vyhledávání\else vlastní tisk\fi}\\ volně dostupná na \url{www.hvalur.org}\\[\baselineskip]
\else
\textbf\ISBN\\[\baselineskip]
\fi
\ifPDF
\else
Tisk: \textbf{Tiskárna Melmen, Pardubice}\\
\fi
Heslová část vysázena patkovým fontem \TeX\ Gyre Pagella,
hesla bezpatkovým fontem {\phvfamily Helvetica}.
\end{minipage}

% -------------------------------------------------------- Table of contents -

\cleardoublepage 

\pdfbookmark\contentsname{toc}    
\tableofcontents
\thispagestyle{empty}
\cleardoublepage 

\fi

\pagenumbering{arabic}
\setcounter{page}{3}
% ============================================================= INTRODUCTION =

\ifinputintroduction

% ------------------------------------- The study Icelandic-Czech dictionary -

\chapter{O Islandsko-českém studijním slovníku}
Cílem vytvoření\textit{ Islandsko-českého studijního slovníku} je zpřístupnění islandského jazyka všem zájemcům o~tento jazyk. 

\textit{Islandsko-český studijní slovník} je určen především českým a~slovenským studentům islandštiny, začátečníkům i~pokročilým, kteří zde naleznou ucelené fonetické, morfologické a~syntaktické informace, jakož
i~slovní zásobu potřebnou k~aktivnímu i~pasivnímu užití s~rozsahem vhodným ke studiu islandštiny na vysoké škole.
Slovník je rovněž určen všem zájemcům o~islandštinu a~Island, jimž slovník nabízí encyklopedické informace, úvod do fonetiky a~gramatiky a~velké množství fotografií a~ilustrací. 
Třebaže se jedná o~slovník primárně zaměřený na potřeby českých a~slovenských uživatelů, obsahuje také řadu informací, které usnadní práci islandskému uživateli (např. použití více jak 10 000 synonym a~antonym, překlady příkladů ap.).

Slovní zásoba původně vychází z~\textit{Concise Icelandic-English Dictionary} (\cite {ic_en}), značně však byla rozšířena díky přístupu k~webovému projektu \textit{ISLEX} (\cite {int1}). Velké množství hesel bylo
přidáno na podnět uživatelů online verze \textit{Islandsko-českého studijního slovníku}.

Slovník obsahuje 30~575 hesel s~46~725 významy nebo slovními spojeními. Z~celkového počtu 30~575 je 28~097 svébytných hesel obsahujících alespoň jeden význam nebo slovní spojení. 
Zbylá hesla (v~počtu 2~478) jsou varianty, dublety nebo nepravidelné tvary.

Slovník je primárně online elektronický slovník, volně přístupný na internetových stránkách \url{www.hvalur.org}  (\cite {int14}). 
Tištěná verze slovníku je publikována na vlastní náklady v~počtu 30 kusů, přičemž většina exemplářů je určena knihovnám a~institucím v~České republice, Slovenské republice a~na Islandu
a~několik exemplářů je určeno autorům nebo zájemcům z~řad veřejnosti. Dvě PDF verze (vyhledávání\,/\addthin prohlížení offline a~vlastní tisk) jsou dostupné zdarma na internetových stránkách slovníku (viz výše).

\textit{Islandsko-český studijní slovník} je otevřený svobodný projekt, jehož se v~průběhu tvorby zúčastnilo velké množství spoluautorů (viz Historie tvorby IČSS), kterým tímto jménem autorů velmi děkuji.
Děkuji také všem uživatelům online verze slovníku za trpělivost, dobrou vůli a~veškeré připomínky a~návrhy, jmenovitě především Petru Mikešovi, Adamu Kožouškovi a~Samoriele za neutuchající víru v~náš projekt.
Poděkování patří také našim rodinám a~blízkým, díky kterým jsme se mohli soustavně věnovat tvorbě slovníku.
Srdečné poděkování patří mé ženě Dorotce za podporu a~nadšení potřebné k~tvorbě slovníku.

\blspace[5]

{\centering Aleš Chejn, Brwinów, 24. prosince 2015\par}


% ------------ The history of the formation Icelandic-Czech dictionary study -

\chapter{Historie tvorby Islandsko-českého studijního slovníku}
\longchapterskip
Projekt \textit{Islandsko-českého studijního slovníku} (IČSS) vznikl v~roce 2001 na nepřímý podnět doc. PhDr. Heleny Kadečkové, CSc. z~Univerzity Karlovy,
kdy byl vypracován seznam slov. 
Tento seznam obsahoval islandská slova a~jejich překlady, později byly přidány základní morfologické informace o~slovním druhu a~deklinační\,/\addthin konjugační koncovky podle vzoru z~\textit{Concise Icelandic-English Dictionary} (\cite {ic_en}). 
Seznam v~této době obsahoval okolo 1 200 slov. Byl zveřejněn pomocí statických webových stránek a~později byl umístěn na webu Univerzity Karlovy. 

V~roce 2006 byl slovník rozšířen o~nových 1 500 heslových slov a~stávající hesla byla rozšířena o~nové významy, zvláště významy frázových sloves. 
Ke konci roku 2006  byl slovník převeden z programu Excel do databáze MySQL a~byly zahájeny práce na vytvoření aplikace \textit{Dictionary System} (DS) (\cite {int19}), která by podporovala tvorbu IČSS. 
Obsah slovníku byl zveřejněn pod svobodnou licencí na webových stránkách \url{www.hvalur.org} (\cite {int14}), kde se nachází dodnes. 

V~roce 2007 byla vytvořena první verze aplikace DS. V~této době se k~projektu přidala Renata Pešková Emilsson, která se nyní věnuje podpoře výuky češtiny na Islandu. 
Ta slovník obohatila rozsáhlým seznamem slov z~praktického života na Islandu. Dále se slovník rozšiřoval o~heslová slova podle \textit{Concise Icelandic-English Dictionary} (\cite {ic_en}). 
V průběhu následujících třech let bylo doplněno přibližně 15~000 slov. V této době byla hesla rozšířena o~synonyma, příklady, frekvenci, oborovou a~stylovou charakteristiku. K projektu se na Islandu připojila na krátkou dobu Kristýna Antonová, studentka islandštiny na Islandské univerzitě, která 
nás upozornila na nesrovnalosti v~uživatelském rozhraní aplikace DS.

Během léta 2009 byl vytvořen skript na generování skloňování a~časování heslových slov. Výsledky generování byly porovnány s~\textit{Beygingarlýsing íslensks nútímamáls} (\cite {DalvikVM}) a~během následujících let konzultovány
s~Jónem Gíslasonem. Kristín Bjarnadóttir, vedoucí projektu \textit{Beygingarlýsing íslensks nútímamáls}, zodpověděla řadu našich dotazů týkajících se deklinací a~konjugací.

V~roce 2009 začala spolupráce s~\textit{Biolib.cz} (\cite {int5}), jmenovitě s~Ondřejem Zichou. Obdrželi jsme  databázi s~náhledy obrázků, autory a~licencemi a~latinskými názvy botanických a~zoologických druhů.

V~zimním semestru 2009 jsme požádali s~Amirem Mulahumicem, bosenským studentem islandštiny na Islandské univerzitě, Studentský fond Islandské univerzity o~grant na tvorbu zvukové databáze islandských slov. 
Studentský fond návrh přijal. Peníze z~grantu jsme určili pro islandského rodilého mluvčího, Jóna Gíslasona, lektora islandského jazyka na Islandské univerzitě v~Reykjavíku, který nahrál více jak 22~000 islandských slov do elektronické formy. 
S~Dorotou Nierychlewskou-Chejn jsme připravili dva seznamy heslových slov (jeden pro nahrávajícího a~druhý pro zpracování nahrávek). Na podporu nahrávek jsme vytvořili webové stránky nazvané \textit{Databáze s~výslovností islandských slov} (\cite {int13}) a~tyto stránky jsme přeložili do 6~jazyků (Jón Gíslason do islandštiny, Ján Zaťko do slovenštiny, francouzštiny a~angličtiny, Amir Mulamuhic do srbochorvatštiny a~Dorota Nierychlewska-Chejn do polštiny).

V~roce 2010 jsme přidali do databáze pravidla islandské výslovnosti, která jsou uvedena v~knize \textit{Handbók um íslenskan framburð} (\cite {is2}). Z~těchto pravidel jsme automaticky vygenerovali fonetický zápis výslovnosti všech hesel ve slovníku. 
Jón Gíslason ověřil správnost fonetického zápisu u~všech heslových slov.

Na konci roku 2011 jsme spojili webové stránky slovníku se stránkami aplikace a~sjednotili jsme a~vylepšili vzhled webových stránek. Cennými připomínkami se na vytvoření aplikace podílela
RNDr. Ingrid Nagyová, PhD., vysokoškolská pedagožka Ostravské Univerzity. Ján Zaťko přeložil nové webové stránky, Průvodce po slovníku do slovenštiny, francouzštiny a~angličtiny a~celý manuál aplikace do angličtiny. 
Dorota Nierychlewska-Chejn přeložila stránky do polštiny, Jón Gíslason do islandštiny a~Marcos Helena do portugalštiny.

Od prosince 2011 do března 2014 jsme realizovali první čtení, ve kterém jsme doplnili chybějící informace, přidali nová hesla a~opravili pravopisné chyby. 
Obsah slovníku jsme porovnávali především s~\textit{Íslensk orðabók} (\mbox{\cite {is}}) a~\textit{Islandsko-českým slovníkem} (\cite {is7}). S~Vojtěchem Kupčou jsme konzultovali mnohé nesnáze při překladech.

V~letech 2013 až 2015 byl vytvořen morfologický klíč, který rozšiřuje klasický systém zobrazení morfologické informace v~islandských slovnících o~přesné zařazení daného hesla do určité deklinační\,/\addthin konjugační skupiny, jejíž
skloňování\,/\addthin časování je ukázáno na příkladu. V~roce 2014 jsme obdrželi roční grant od Ostravské univerzity v~Ostravě na tvorbu morfologického klíče. S~realizací grantu nám velmi pomohla doc. PhDr. Diana Svobodová, Ph.D. a~Ing. Zuzana Havrlantová.

Od března 2014 do srpna 2015 jsme realizovali druhé finální čtení, ve kterém jsme porovnávali informace s~webovým slovníkem \textit{ISLEX} (\cite {int1}). 
Na požádání jsme obdrželi od Þórdís Úlfarsdóttir (vedoucí projektu \textit{ISLEX}) seznam víceslovných příslovečných frází. Takto poskytnutý seznam nám usnadnil opravy slovníku. 
S~Þórdís Úlfarsdóttir jsme konzultovali různé gramatické jevy a~záležitosti týkající se deklinací\,/\addthin konjugací.
Martina Kašparová, studentka islandštiny na Islandské univerzitě, prošla celý slovník, přičemž se soustředila mimo jiné na stylistické a~pravopisné korektury překladů. 
Lucie Brabcová udělala jazykové a~typografické korektury celého slovníku.

Vojtěch Kupča vydal v~roce 2014 \textit{Islandskou gramatiku} (\cite {is77}) a~souhlasil s~připojením této publikace ke~knižní verzi IČSS.
V~roce 2015 jsme s~přispěním Paola al-Imk\=an\=\i\ Brasolin vysázeli slovník v~typografickém prostředí {\LaTeX}. V~únoru 2016 byl slovník vytištěn v~Tiskárně MELMEN, s.r.o v~Pardubicích a svázán 
v~Knihařství Vojtíšková Hrubá v~Pardubicích. 

\blspace[5]

{\centering Všem spoluautorům děkujeme za spolupráci.\par}


% -------------------------------------------------  Guide to the dictionary -

\chapter{Průvodce po slovníku}
\begin{multicols}{2}
\def\tableA{\noindent\begin{tabularx}{\columnwidth}{>{\bfseries}lll}
  \toprule
  \textit{přípona} & \textit{přípona s~koncovkou} & \textit{příklad} \\
  \midrule
  -ing-   & -ingur   & rugl··ingur \\
          & -ing     & hreyf··ing \\
          & -ingi    & erf··ingi \\
          & -ingar   & lækn··ingar \\
  -ling-  & -lingur  & vett··lingur \\
  -ning-  & -ningur  & flut··ningur \\
          & -ning    & set··ning \\
          & -ningar  & kos··ningar \\
  -að-    & -aður    & mark··aður \\
  -nað-   & -naður   & skil··naður \\
  -ar-    & -ari     & dóm··ari \\
  -skap-  & -skapur  & kveð··skapur \\
  -dóm-   & -dómur   & lær··dómur \\
  -leik-  & -leikur  & sann··leikur \\
          & -leiki   & mögu··leiki \\
  -lát-   & -látur   & laus··látur \\
  -átt-   & -átta    & víð··átta \\
  -sam-   & -samur   & starf··samur \\
  -sem-   & -semi    & hamingju··semi \\
  -nesk-  & -neskur  & tékk··neskur \\
          & -neska   & tékk··neska \\
  -leg-   & -legur   & blíð··legur \\
          & -lega    & blíð··lega \\
  -ung-   & -ungur   & gler··ungur \\
          & -unga    & snopp··unga \\
          & -ungi    & sveit··ungi \\
          & -ung     & nýj··ung \\
  -ótt-   & -óttur   & blett··óttur \\
  -ræn-   & -rænn    & suð··rænn \\
          & -ræna    & nor··ræna \\
  -ern-   & -erni    & fað··erni \\
  -ind-   & -indi    & vís··indi \\
  -and-   & -andi    & gef··andi \\
  -ug-    & -ugur    & göf··ugur \\
          & -ugt     & stöð··ugt \\
  -ust-   & -usta    & þjón··usta \\
  -ald-   & -aldur   & far··aldur \\
  -uð-    & -uður    & söfn··uður \\
  -úð-    & -úð      & dul··úð \\
  -versk- & -verskur & ung··verskur \\
          & -verska  & ung··verska \\
  -ist-   & -isti    & hass··isti \\
  -ism-   & -ismi    & húman··ismi \\
  -ang-   & -angur   & far··angur \\
  \bottomrule
  \end{tabularx}
  \captionof{table}{Seznam přípon označených ve slovníku}
  \label{table1}}

\def\tableB{\noindent\begin{tabularx}{\columnwidth}{>{\bfseries}lX}
  \toprule
  \textbf{symbol} & \textbf{vysvětlivka} \\
  \midrule
  m                   & označuje slovní druh \\
  \textsubscript{1}   & označuje třídu \\
  (-a, -ar)           & označuje deklinační / konjugační koncovky hesla \\
  \textsubscript{8}   & označuje podtřídu\\
  \bottomrule
  \end{tabularx}
  \captionof{table}{Morfologická informace}
  \label{table2}}

\def\tableC{\noindent\begin{tabularx}{\columnwidth}
    {>{\bfseries}l>{\footnotesize\itshape}lllll}
  \toprule
  \multirow{3}{*}{{{\textbf{f{\textsubscript{1}}} \Large{\textbf{7}}}}}
    & nom & kisa & kisan    & kisur & kisurnar \\
    & acc & kisu & kisuna   & kisur & kisurnar \\
    & dat & kisu & kisunni  & kisum & kisunum \\
    & gen & kisu & kisunnar & kisa  & kisanna \\
  \bottomrule
  \end{tabularx}
  \captionof{table}{Skloňování hesla \textit{kisa}}
  \label{table3}}

\def\tableD{\noindent\begin{tabularx}{\columnwidth}{>{\bfseries}lllll}
  \toprule
  \textit{zkratka} & \textit{slovo} & \textit{pád} & \textit{číslo}
    & \specialcell{\textit{předmět}\\\textit{podmět}} \\
  \midrule
  e-að  & eitthvað   & 1. & sg & neživotný \\
  e-ð   & eitthvað   & 4. & sg & neživotný \\
  e-ir  & einhverjir & 1. & pl & životný \\
  e-ja  & einhverja  & 4. & pl & životný \\
  e-jum & einhverjum & 3. & pl & životný \\
  e-m   & einhverjum & 3. & sg & životný \\
  e-n   & einhvern   & 4. & sg & životný \\
  e-r   & einhver    & 1. & sg & životný \\
  e-rra & einhverra  & 2. & pl & životný \\
  e-rs  & einhvers   & 2. & sg & životný \\
  e-s   & einhvers  & 2. & sg & neživotný \\
  e-u   & einhverju  & 3. & sg & neživotný \\
  \bottomrule
  \end{tabularx}
  \captionof{table}{Seznam použitých zkratek neurčitých zájmen}
  \label{table4}}

%#############################################################################

\section{Heslo}

\subsection*{Řazení slov}

Heslová slova jsou zobrazena tučně a~jsou seřazena podle islandské abecedy:

\blspace
{\bfseries\centering
  a~á b c d ð e é f g h i~í j k~l m n\\
  o~ó p q r s~t u~ú v~w x y ý z~þ æ ö\par}
\blspace

K~snadnější orientaci v~heslové části je použit u~každého písmene tzv. palcový index, který se nachází na hřbetě slovníku. Každé písmeno je v~palcovém indexu barevně odlišeno.

\subsection*{Homonyma}

Homonyma (heslová slova stejného reprezentativního tvaru) jsou označena povýšenými arabskými číslicemi uváděnými za heslovým slovem.

\blspace
  \dicEntry \dicTerm{vor\textsuperscript{1}} \dicPos{n} \dicFlx{(-s, -)}
  \dicEntry \dicTerm{vor\textsuperscript{2}} \dicPos{pron} \dicFlx{pers gen pl}
\blspace

Homonyma jsou seřazena podle českého ustáleného řazení slovních druhů. Poté následují předpony a~zkratky.
V~případě dvou podstatných jmen je pořadí mužský rod, ženský rod, střední rod. V~případě dvou podstatných jmen stejného rodu má přednost slovo se slabým skloňováním před slovem se silným skloňováním. V~případě dvou sloves stejného reprezentativního tvaru má přednost sloveso se slabým časováním před slovesem se silným časováním.

\subsection*{Varianty a~dublety}

U~některých heslových slov jsou v~záhlaví hesla uvedeny a) varianty (obměny)

\blspace
  \dicEntry \dicTerm{ærsla··fenginn}, \dicTerm{ærslafullur} \dicPos{adj}
\blspace

nebo b) dublety (dvojtvary).

\blspace
  \dicEntry \dicTerm{{ævin··týri}, \dicTerm{ævintýr}} \dicPos{n}
\blspace

Varianty a~dublety jsou méně frekventované než samotné heslo. Lze je vyhledat ve slovníku jako hesla obsahující fonetické a~morfologické informace s~odkazem na původní heslové slovo.

\blspace
  \dicEntry \dicTerm{ærsla··|fullur} \dicIPA{{a}{i}{\textsubring{r}}{s}{\textsubring{d}}{l}{a}{f}{\textscy}{\textsubring{d}}{l}{\textscy}{\textsubring{r}}} \dicPos{adj} \dicFlx{(comp -fyllri, sup -fyllstur)} \dicLink{ærslafenginn}
  \dicEntry \dicTerm{ævin··týr} \dicIPA{{a}{i}{\textlengthmark}{v}{\textsci}{n}{t\textsuperscript{h}}{i}{\textsubring{r}}} \dicPos{n} \dicFlx{(-s, -)} \dicLink{ævintýri}
\blspace

\subsection*{Vybrané tvary jako hesla}

Ve slovníku jsou uvedeny vybrané tvary heslových slov, u~kterých by mohl mít uživatel potíže s~identifikací. Jedná se o~a) genitiv jednotného čísla ({\textbf{sg}}) a~nominativ množného čísla ({\textbf{pl}}) podstatných jmen, b) ženský a~střední rod přídavných jmen a~2. a~3. stupeň rodu mužského přídavných jmen, c) oznamovací způsob ({\textbf{ind}}) přítomného času ({\textbf{praes}}) a~minulého času ({\textbf{pf}}) 1. osoby ({\textbf{pers}}) sg a~pl a~spojovací způsob ({\textbf{con}}) minulého času 1. osoby sg a~jiné nepravidelné tvary.

\blspace
  \dicEntry \dicTerm{öttum} \dicPos{v} \dicFlx{ind pf pl 1 pers} \dicLink{etja}
\blspace

\subsection*{Dělení heslových slov}

K~dělení heslových slov je použita střední tečka \textbf{'·'}, která odděluje části složeného slova, přípony a~předpony. Platí následující pravidla:

\renewcommand{\labelenumi}{\alph{enumi})}
\begin{enumerate}
\item jediné dělení heslového slova je označeno dvěma středními tečkami \textbf{'··'}. Může se jednat o~dělení částí složeného slova
\blspace
\dicEntry \dicTerm{ráð··hús}
\blspace

nebo o~dělení mezi kořenem a~předponou / příponou.

\blspace
\dicEntry \dicTerm{þjón··usta}
\blspace
\item V~případě, že se v~hesle vyskytuje více dělení, jsou použity dvě střední tečky \textbf{'··'} pro významově poslední dělení a~jedna střední tečka \textbf{'·'} pro zbylé / zbylá dělení.
\blspace
\dicEntry \dicTerm{við·bótar··líf·eyris·sparn·aður}
\blspace
\item V~hesle s~předponou je označena jedna nebo i~více předpon (každá zvlášť).
\blspace
\dicEntry \dicTerm{ó··að·finnan·legur}
\blspace
\item V~hesle s~příponou je označena jedna nebo i~více přípon (každá zvlášť).
\blspace
\dicEntry \dicTerm{létt·úð··ugur}
\blspace
\item V~hesle s~příponou je označena pouze přípona / přípony poslední části složeného slova. Důvodem pro neoznačení přípon jiných částí složených slov je lepší čitelnost hesla. Příklad označení pouze přípony poslední části složeniny:
\blspace
\dicEntry \dicTerm{flutninga··miðl·ari}
\end{enumerate}

Následuje seznam přípon, které jsou ve slovníku označeny:

\tableA

\section{Frekvence slov}

Symbolem \dicsymFrequent\ je označeno 2~825 hesel patřících mezi velmi frekventovaná islandská slova. Tento seznam vychází z~korpusu \textit{Íslenskur Orðasjóður} (\cite{int7}) a~knihy \textit{Íslensk orðtíðnibók} (\cite{is2}).

\blspace
  \dicEntry \dicTerm{koma}\textsuperscript{2} \dicsymFrequent
\blspace

\section{Fonetický zápis}

V~hranatých závorkách v~záhlaví hesla je uveden fonetický zápis výslovnosti hesla v~mezinárodní fonetické abecedě (IPA, angl. International Phonetic Alphabet).

\blspace
  \dicEntry \dicTerm{öldu··rót} \dicIPA{{\oe}{l}{\textsubring{d}}{\textscy}{r}{ou}{\textsubring{d}}}
\blspace

Znak '{\textlengthmark}' označuje dlouhou samohlásku nebo souhlásku.

\blspace
  \dicEntry \dicTerm{ör}\textsuperscript{2} \dicIPA{{\oe}{\textlengthmark}{\textsubring{r}}}
\blspace

\section{Slovní druhy, předpony a~zkratky}

V~záhlaví hesla je uvedena morfologická informace popisující a) slovní druh hesla, b) předponu nebo c) zkratku. Pokud se heslo řadí do více slovních druhů, jsou jednotlivé zkratky pro slovní druhy odděleny lomítkem.

\blspace
  \dicEntry \dicTerm{af} \dicPos{prep / adv}
\blspace

Informace o~slovním druhu, předponě či zkratce je vysázena barevně a~tučně.

\subsection*{Podstatná jména}

Podstatná jména jsou uvedena v~1. pádě (\textbf{nom}) jednotného čísla (\textbf{sg}) nebo v~1. pádě množného čísla (\textbf{pl}), pokud se heslo v~jednotném čísle nevyskytuje. Podstatná jména jsou označena následovně: \textbf{m} (l. masculinum) - podstatné jméno rodu mužského,
\textbf{f} (l. femininum) - podstatné jméno rodu ženského, \textbf{n} (l. neutrum) - podstatné jméno rodu středního, \textbf{subs} (l. substantivum) - podstatné jméno bez jednoznačného mluvnického rodu.

\blspace
  \dicEntry \dicTerm{hestur} \dicPos{m}
  \dicEntry \dicTerm{Hanoí} \dicPos{subs}
\blspace

\subsection*{Přídavná jména}

Přídavná jména jsou uvedena v~1. pádě (\textbf{nom}) rodu mužského (\textbf{m}) jednotného čísla (\textbf{sg}) 1. stupně (\textbf{pos}) (l. positivus) a~jsou označena zkratkou \textbf{adj} (l. adjectivum).

\blspace
  \dicEntry \dicTerm{fallegur} \dicPos{adj}
\blspace

V~případě, že u~přídavného jména dochází v~rodě ženském (\textbf{f}) nebo středním (\textbf{n}) k~přehlásce anebo je tvar druhého stupně (\textbf{comp} (l. comparativus)) či třetího stupně (\textbf{sup} (l. superlativus)) nepravidelný, pak jsou tyto tvary uvedeny v~nominativu jako heslová slova s~odkazem na hlavní heslový tvar

\blspace
  \dicEntry \dicTerm{yngri} \dicPos{adj} \dicFlx{comp m} \dicLink{ungur}
\blspace

nebo jako svébytná hesla s~významy.

\blspace
  \dicEntry \dicTerm{fremri} \dicPos{adj} \dicFlx{comp (sup fremstur)}
\blspace

\subsection*{Zájmena}

Zájmena jsou označena zkratkou \textbf{pron} (l. pronomen) a~jsou tříděna na zájmena ukazovací \textbf{dem} (angl. demonstrative), přivlastňovací \textbf{poss} (angl. possessive), osobní \textbf{pers} (angl. personal), neurčitá \textbf{indef} (angl. indefinite), zvratná \textbf{refl} (angl. reflexive) a~vztažná \textbf{rel} (angl. relative).

\blspace
  \dicEntry \dicTerm{þessi} \dicPos{pron} \dicFlx{dem}
\blspace

U~zájmen přivlastňovacích je v~3. osobě označen rod.

\blspace
  \dicEntry \dicTerm{hans} \dicPos{pron} \dicFlx{poss m}
\blspace

U~zájmen osobních je označen pád a~číslo.

\blspace
  \dicEntry \dicTerm{mig} \dicPos{pron} \dicFlx{pers acc sg}
\blspace

\subsection*{Číslovky}

Číslovky jsou označeny zkratkou \textbf{num} (l. numerale). V~případě vybraných číslovek (1 až 4) jsou ve slovníku uvedeny jako heslová slova také tvary ženského i~středního rodu. V~takovém případě jsou označeny zkratkou \textbf{m} pro mužský rod, zkratkou \textbf{f} pro ženský rod a~zkratkou \textbf{n} pro střední rod. Zkratkou \textbf{ord} (angl. ordinal number) jsou označeny číslovky řadové.

\blspace
  \dicEntry \dicTerm{tut··tugu} \dicPos{num}
  \dicEntry \dicTerm{tveir} \dicPos{num} \dicFlx{pl m}
  \dicEntry \dicTerm{fyrsti} \dicPos{num} \dicFlx{ord}
\blspace

\subsection*{Slovesa}

Slovesa jsou uvedena v~infinitivu a~jsou označena zkratkou \textbf{v} (l. verbum).

\blspace
  \dicEntry \dicTerm{fara} \dicPos{v}
\blspace

Slovesa ve středním rodě jsou označena jako \textbf{refl} (angl. reflexive).

\blspace
  \dicEntry \dicTerm{nálgast} \dicPos{v} \dicFlx{refl}
\blspace

Slovesa, která jsou vždy neosobní, jsou označena \textbf{impers} (angl. impersonal).

\blspace
  \dicEntry \dicTerm{svima} \dicPos{v} \dicFlx{impers}
\blspace

U~sloves je dále uvedeno, s~jakým pádem se pojí. Zkratka \textbf{acc} označuje 4. pád, \textbf{dat} 3. pád, \textbf{gen} 2. pád a~\textbf{nom} 1. pád. V~případě, že se sloveso pojí s~různými pády, je použito lomítko.

\blspace
  \dicEntry \dicTerm{klóra} \dicPos{v} \dicFlx{acc / dat}
\blspace

V~případě, že se sloveso pojí s~více pády, je použito znaménko plus.

\blspace
  \dicEntry \dicTerm{gefa} \dicPos{v} \dicFlx{dat + acc}
\blspace

\subsection*{Příslovce}

Příslovce jsou označena zkratkou \textbf{adv} (l. adverbium). Pokud chybí informace o~stupni, jedná se o~1. stupeň stupňování (\textbf{pos}). 2. stupeň je vždy označen zkratkou \textbf{comp} a~3. stupeň zkratkou \textbf{sup}.

\blspace
  \dicEntry \dicTerm{þar} \dicPos{adv}
  \dicEntry \dicTerm{ofar} \dicPos{adv} \dicFlx{comp (pos uppi, sup efst)}
  \dicEntry \dicTerm{síðast} \dicPos{adv} \dicFlx{sup (comp síðar)}
\blspace

\subsection*{Předložky}

Předložky jsou označeny zkratkou \textbf{prep} (l. praepositio). Pokud se předložka pojí výhradně s~jedním pádem, jsou použity zkratky pro 4. pád (\textbf{acc}), 3. pád (\textbf{dat}), 2. pád (\textbf{gen}).

\blspace
  \dicEntry \dicTerm{af} \dicPos{prep / adv} \dicFlx{dat}
\blspace

V~případě, že se předložka pojí s~různými pády, je použito lomítko.

\blspace
  \dicEntry \dicTerm{fyrir} \dicPos{prep / adv} \dicFlx{acc / dat}
\blspace

\subsection*{Spojky}

Spojky jsou označeny zkratkou \textbf{conj} (l. conjunctio).

\subsection*{Částice}

Částice jsou označeny zkratkou \textbf{part} (l. particula).

\subsection*{Citoslovce}

Citoslovce jsou označeny zkratkou \textbf{inter} (l. interjectio).

\subsection*{Předpony}

Předpony jsou označeny zkratkou \textbf{predp}.

\subsection*{Zkratky}

Zkratky jsou označeny zkratkou \textbf{zkr}.

\section{Skloňování, časování a~stupňování}

Každé heslo obsahuje zápis morfologické informace, která přesně informuje, jak se heslo skloňuje / časuje nebo stupňuje.

V~slovníku jsou použity dva odlišné systémy zápisu morfologické informace, které se vzájemně doplňují:

\begin{enumerate}
\item tradiční systém - do závorky za slovním druhem se zapisují deklinační / konjugační koncovky. Počet koncovek závisí na slovním druhu (viz dále).
\item numerický systém (vychází ze systému Helgi Haraldssona (\cite {is6})) - číslice odkazující na morfologický klíč, kde jsou podrobně vypsány deklinační / konjugační třídy.
Další informace o~numerickém systému lze nalézt na straně \pageref{sec:dec}.
\end{enumerate}

Deklinační / konjugační tvary nebo koncovky podstatných jmen, přídavných jmen, zájmen, číslovek, sloves a~příslovcí jsou uvedeny v~závorkách za slovním druhem.

Hnízdovací šev \textbf{'|'} označuje místo ve slově, kde se deklinační koncovka připojuje ke slovu.

\blspace
  \dicEntry \dicTerm{heils|a{\textsuperscript{1}}} \dicPos{f} \dicFlx{(-u)}
\blspace

Pokud není hnízdovací šev \textbf{'|'} uveden, tak se koncovka připojuje k~samotnému heslu.

\blspace
  \dicEntry \dicTerm{mynd} \dicPos{f} \dicFlx{(-ar, -ir)}
\blspace

Pokud dochází v~ohýbání k~přehlásce nebo je ohýbání nepravidelné, pak jsou uvedeny celé tvary hesla.

\blspace
  \dicEntry \dicTerm{maður} \dicPos{m} \dicFlx{(manns, menn)}
\blspace

V~případě, že heslové slovo má více tvarů, jsou varianty koncovky odděleny lomítkem.

\blspace
  \dicEntry \dicTerm{beð|ur} \dicPos{m} \dicFlx{(-s / -jar, -ir)}
\blspace

\subsection*{Podstatná jména}

U~podstatných jmen je uveden tvar hesla pro 2.~pád jednotného čísla a~1.~pád množného čísla.

\blspace
  \dicEntry \dicTerm{hval|ur} \dicPos{m} \dicFlx{(-s, -ir)}
\blspace

kde '\textit{hvals}' je 2.~pád jednotného čísla a~'\textit{hvalir}' 1.~pád množného čísla.
Pokud je uveden pouze jeden tvar, jedná se o~2.~pád jednotného čísla a~znamená to rovněž, že heslové slovo se neohýbá v~množném čísle.

\blspace
  \dicEntry \dicTerm{heils|a{\textsuperscript{1}}} \dicPos{f} \dicFlx{(-u)}
\blspace

Pokud je použita zkratka pro množné číslo (\textbf{pl}), znamená to, že se heslo ohýbá pouze v~množném čísle.

\blspace
  \dicEntry \dicTerm{afar··kostir} \dicPos{m} \dicFlx{pl}
\blspace

Pokud je použita zkratka \textbf{indecl}, znamená to, že podstatné jméno je nesklonné.

\blspace
  \dicEntry \dicTerm{október} \dicPos{m} \dicFlx{indecl}
\blspace

\subsection*{Přídavná jména}

U~přídavných jmen jsou uvedeny v~závorkách tvary, které poukazují na nepravidelnosti ve skloňování.

Pokud přídavné jméno v~ženském rodě neztrácí koncové '\textit{-r}', je uveden tvar hesla v~ženském rodě.

\blspace
  \dicEntry \dicTerm{hýr} \dicPos{adj} \dicFlx{(f -)}
\blspace

Dále jsou uvedeny tvary ženského rodu, ve~kterých dochází k~přehlásce:

\blspace
  \dicEntry \dicTerm{reglu··|samur} \dicPos{adj} \dicFlx{(f -söm)}
\blspace

Pokud je přídavné jméno nesklonné, je uvedena zkratka \textbf{indecl}.

\blspace
  \dicEntry \dicTerm{tví··mála} \dicPos{adj} \dicFlx{indecl}
\blspace

Pokud má přídavné jméno nepravidelné tvary 2.~a~3.~stupně stupňování, jsou tyto tvary uvedeny v~závorce spolu se zkratkou pro 2.~stupeň (\textbf{comp}) a~3.~stupeň (\textbf{sup}).

\blspace
  \dicEntry \dicTerm{þunnur} \dicPos{adj} \dicFlx{(comp þynnri, sup þynnstur)}
\blspace

V~případě, že je uveden pouze jeden ze dvou zbývajících stupňů, znamená to, že heslo se stupňuje pouze v~uvedených stupních.

\blspace
  \dicEntry \dicTerm{syðri} \dicPos{adj} \dicFlx{comp (sup syðstur)}
\blspace

\subsection*{Slovesa}

Islandská slovesa se dělí na slabá, silná a~nepravidelná. Vzhledem k~tomuto dělení se liší počet slovních tvarů uvedených v~závorce za zkratkou \textbf{v}. Princip uvádění slovních tvarů spočívá v~tom, že uživatel
je schopen (více méně) z~těchto tvarů odvodit časování slovesa. V~tomto slovníku jsme
záměrně umístili některá slabá slovesa mezi silná (místo dvou slovesných tvarů je uvedeno pět tvarů) kvůli tomu, aby byl uživatel přesněji informován o~časování slovesa.

Slabá slovesa, která tvoří 1.~osobu jednotného čísla minulého času oznamovacího způsobu koncovkou '\textit{-aði}', mají v~závorce uvedenu jen jednu koncovku, jmenovitě '\textit{(-aði)}'.

\blspace
  \dicEntry \dicTerm{ætl|a} \dicPos{v} \dicFlx{(-aði)}
\blspace

U~zbývajících skupin slabých sloves se v~závorce nachází koncovka 1.~osoby jednotného čísla minulého času oznamovacího způsobu a~supinum.

\blspace
  \dicEntry \dicTerm{kenn|a} \dicPos{v} \dicFlx{(-di, -t)}
\blspace

Silná nebo nepravidelná slovesa mají v~závorce vždy pět tvarů, jmenovitě 1.~osobu jednotného čísla přítomného času oznamovacího způsobu, 1.~osobu jednotného čísla minulého času oznamovacího způsobu, 1.~osobu množného čísla minulého času oznamovacího způsobu, 1.~osobu jednotného čísla minulého času spojovacího způsobu a~supinum.

\blspace
  \dicEntry \dicTerm{grípa} \dicPos{v} \dicFlx{(gríp, greip, gripum, gripi, gripið)}
\blspace

Silná nebo nepravidelná slovesa vyskytující se pouze v~neosobním užití mají v~závorce 4 tvary, jmenovitě 3.~osobu jednotného čísla přítomného času oznamovacího způsobu, 3.~osobu jednotného čísla minulého času oznamovacího způsobu, 3.~osobu jednotného čísla minulého času spojovacího způsobu a~supinum.

\blspace
  \dicEntry \dicTerm{kross··|bregða} \dicPos{v} \dicFlx{(-bregður, -brá, -brygði, -brugðið)}
\blspace

U~některých významů v~heslovém slově se vyskytují gramatické informace, které popisují chování daného významu.

Zkratka \textbf{impers} označuje, že se daný význam slovesa vyskytuje v~neosobním užití. Pomocí slovního spojení je vyjádřeno užití slovesa.

\blspace
  \dicEntry \dicTerm{batn|a} \dicPos{v} \dicFlx{(-aði)} \dicTerm{e-m batnar} \dicFlx{impers} (kdo) se uzdravuje, (kdo) se zotavuje, (komu) je lépe
\blspace

Zkratka \textbf{refl} označuje mediopasivní tvary slovesa.

\blspace
  \dicEntry \dicTerm{að··lag|a} \dicPos{v} \dicFlx{(-aði)} \dicIdiom{aðlagast} {\dicPhraseIS{aðlagast e-u}} \dicFlx{refl} {přizpůsobit se (čemu)}
\blspace

\subsection*{Příslovce}

U~příslovcí se v~záhlaví nachází tvary stupňování. Tyto tvary jsou uvedeny v~závorce spolu se zkratkou pro 2.~stupeň (\textbf{comp}) a~3.~stupeň (\textbf{sup}).

\blspace
  \dicEntry \dicTerm{vel\textsuperscript{2}} \dicPos{adv} \dicFlx{(comp betur, sup best)}
\blspace

Nepravidelné tvary 2.~a~3.~stupně jsou uvedeny buď jako hesla odkazující na základní tvar v~1.~stupni (\textbf{pos}), anebo jako svébytná hesla. V~obou případech jsou uvedeny v~závorce zbývající stupně.

\blspace
  \dicEntry \dicTerm{síst} \dicPos{adv} \dicFlx{sup (pos varla, comp síður)}
\blspace

Pokud některý ze stupňů uveden není, znamená to, že pro dané slovo neexistuje.

\blspace
  \dicEntry \dicTerm{síðar} \dicPos{adv} \dicFlx{comp (sup síðast)}
\blspace

\section{Morfologická informace} \label{sec:dec}

V~záhlaví hesla se kromě deklinačních / konjugačních koncovek nachází také morfologické numerické informace, které odkazují na morfologický klíč (viz strana \pageref{sec:morpho}), kde jsou podrobně vypsány kompletní deklinační / konjugační tabulky.

\blspace
  \dicEntry \dicTerm{ap|i} \dicPos{m}[1] \dicFlx{(-a, -ar)}[8]
\blspace

Morfologická informace je zapsána podle následujícího schématu: \textbf{m{\textsubscript{1}} (-a, -ar){\textsubscript{8}}}

\blspace
\tableB
\blspace

Deklinační / konjugační třídy a~jejich počet nejsou v~islandštině ustálené a~slouží pouze pro rozlišení odlišných skloňování / časování.

Názvy podtříd slouží pro rozlišení odlišných skloňování / časování.

\subsection*{Práce s~morfologickým klíčem}

Pro vyhledání kompletních deklinačních / konjugačních tabulek slouží následující postup, který demonstrujeme na ukázce:

\blspace
  \dicEntry \dicTerm{kis|a} \dicPos{f}[1] \dicFlx{(-u, -ur)}[7]
\blspace

Nejdříve určíme slovní druh (v~případě podstatných jmen také rod), v~tomto případě podstatné jméno rodu ženského, a~navštívíme Morfologický klíč (str. \pageref{sec:morpho}) pro podstatná jména rodu ženského (str. \pageref{sec:morpho_f}).
Následně dohledáme třídu (v~tomto případě číslo {\textsubscript{\textbf{1}}}) a~následně dohledáme dodatečné numerické informace za závorkou (v~tomto případě číslo {\textsubscript{\textbf{7}}}).

\tableC

\section{Definice}

\subsection*{Základní tvar}

Základním tvarem definice je, že jednomu islandskému slovu odpovídá jeden nebo více českých významů.

\blspace
  \dicEntry \dicTerm{landa··fræð|i} \dicPos{f} \dicFlx{(-i)} \textCS{zeměpis, geografie}
\blspace

V~některých případech je v~české části použita závorka. V~prvním případě je závorka užívaná ke~zkrácení zápisu dvou slov ('\textit{pomoc}', '\textit{výpomoc}').

\blspace
  \dicEntry \dicTerm{að··stoð} \dicPos{f} \dicFlx{(-ar)} \textCS{(vý)pomoc, asistence}
\blspace

V~druhém případě se v~závorce nachází upřesňující informace, která zabraňuje dvojznačnosti.

\blspace
  \dicEntry \dicTerm{fíkj|a} \dicPos{f} \dicFlx{(-u, -ur)} \textCS{fík} {\footnotesize {\textCS{(plod)}}}
\blspace

V~mnoha případech je závorka v~české části použita pro příklad užití českého slova a~tím vymezení jeho významu.

\blspace
  \dicEntry \dicTerm{að·gengi··legur} \dicPos{adj} {přístupný, dostupný (informace ap.)}
\blspace

Lomítko \textbf{'/'} slouží v~české části významu k~oddělení více významů pojících se s~jedním nebo více slovy a~je použito v~české části následujícím způsobem:

\blspace
  \dicEntry \dicTerm{aðal··|gata} \dicPos{f} \dicFlx{(-götu, -götur)} \textCS{hlavní třída / ulice}
\blspace

Význam přečteme jako '\textit{hlavní třída, hlavní ulice}'.

\subsection*{Opis heslového slova}

V~případech, kdy islandské slovo nemá ekvivalent v~českém jazyce, jsou významy islandských slov opsány delším opisem a~zobrazeny menším písmem. Příkladem jsou slova z~lidových pověstí, kulinářské speciality nebo botanické pojmy.

\blspace
  \dicEntry \dicTerm{til··ber|i} \dicPos{m} \dicFlx{(-a, -ar)} {\footnotesize{pov.}} {\footnotesize{stvoření, které dojí krávy a~ovce hospodářů}}
\blspace

\subsection*{Slovní spojení}

V~definici hesla jsou kromě samotných významů uváděna také slovní spojení, která jsou sázena tučně. Po slovním spojení následuje překlad slovního spojení.

\blspace
  \dicEntry \dicTerm{lær|a} \dicPos{v} \dicFlx{(-ði, -t)} \textbf{læra e-ð utanbókar} {učit se (co) zpaměti, memorovat (co)}
\blspace

Závorky jsou použity na tu část slovního spojení, kterou je možné vypustit.

\blspace
  \dicEntry \dicTerm{ár\textsuperscript{1}} \dicPos{f} \dicFlx{(-ar, -ar)} \textbf{taka (of) djúpt í árinni} {\footnotesize{přen.}} {přehnat (co v~tvrzení)}
\blspace

Příklad můžeme přečíst jako '\textit{taka djúpt í árinni}' nebo '\textit{taka of djúpt í árinni}'.

Mezi slovní spojení v~tomto slovníku řadíme a) tvary heslového slova, b) příkladová spojení, c) ustálená slovní spojení a~d) rčení nebo přísloví.

\begin{enumerate}
\item Tvary heslového slova\\
Do této skupiny řadíme tvar množného čísla hesla (zkratka \textbf{pl}).
\blspace
\dicEntry \dicTerm{bót} \dicPos{f} \dicFlx{(-ar, bætur)} \textbf{bætur} \dicFlx{pl} {dávky, přídavky, příspěvky}\\

Dále do této skupiny řadíme tvar slovesa v~středním slovesném rodě (zkratka \textbf{refl}).\\

\dicEntry \dicTerm{finna} \dicPos{v} \dicFlx{(finn, fann, fundum, fyndi, fundið)} \dicIdiom{finnast} {\dicPhraseIS{finnast}} \dicFlx{refl} {být nalezen, najít se, nalézt se, nacházet se}
\blspace
\item Příkladová spojení\\
Příkladová spojení ilustrují komunikační chování heslových slov v~příslušných významech a~v~různých syntagmatech (např. sloveso s~určitou předložkou ap.). Jsou označena tučně modrou barvou a~odsazená pro snadnou orientaci jako svébytné heslo.
\blspace
\dicEntry \dicTerm{koma\textsuperscript{2}} \dicPos{v} \dicFlx{(kem, kom, komum, kæmi, komið)} \dicIdiom{koma}[að] {\dicPhraseIS{koma að e-m}} {přiblížit se ke (komu), přijít ke (komu)}
\blspace
\item Ustálená slovní spojení\\
Tento druh slovních spojení je charakterizován většinou přeneseným významem. Tato spojení jsou uváděna za posledním významem nebo příkladovým spojením a~jsou označená tučně modrou barvou a~odsazená pro snadnou orientaci jako svébytné heslo.
\blspace
\dicEntry \dicTerm{kast} \dicPos{n} \dicFlx{(-s, köst)} \dicIdiom{kast} {\dicPhraseIS{þegar til kastanna kemur}} {\footnotesize{přen.}} {když jde do tuhého}
\blspace
\item Rčení a~přísloví\\
Rčení a~přísloví se nacházejí vždy na konci definice a~jsou odlišená symbolem \dicsymProverb\ a~odsazená pro snadnou orientaci.
\blspace
\dicEntry \dicTerm{æfing} \dicPos{f} \dicFlx{(-ar, -ar)} \dicProverb{} \dicPhraseIS{Æfing skapar meistara.} {\footnotesize{přís.}} \textCS{Cvičení dělá mistra.}
\end{enumerate}

\section{Vícevýznamovost}

\subsection*{Vícevýznamová hesla}

Značné množství hesel má více než jeden význam. Významy jsou odlišeny arabskými číslicemi a~v~některých případech i~islandskými synonymy.

\blspace
  \dicEntry \dicTerm{land} \dicPos{n} \dicFlx{(-s, lönd)} {\textbf{1.}} \dicSynonym{þurrlendi} {souš, pevnina, země} {\textbf{2.}} \dicSynonym{árbakki} {břeh} {\textbf{3.}} \dicSynonym{ríki} {země, stát, vlast} {\textbf{4.}} \dicSynonym{landareign} {země, půda, pozemek}
\blspace

\subsection*{Vícevýznamová slovní spojení}

V~případě, že slovní spojení uvnitř definice má více významů, jsou jednotlivé významy odlišeny písmeny {\textbf{a., b., c.}} atd.

\blspace
  \dicEntry \dicTerm{drag} {\dicPos{n} \dicFlx{(-s, drög)} {\textbf{1.}} {mokřina, podmoklý terén;} {\textbf{2.}} \textbf{drög} {\footnotesize{pl}} {\textbf{a.}} \dicSynonym{uppsprettur} {prameny (řeky ap.)} {\textbf{b.}} \dicSynonym{undirbúningur} {náčrt, návrh.}}
\blspace

\section{Řazení významů a~slovních spojení v~definici}
\longsectionskip

V~definici jsou nejdříve seřazeny významy slova a~odlišeny arabskými číslicemi.

\subsection*{V rámci významu}

Slovní spojení jsou uváděna v~rámci jednotlivého významu, pokud se s~ním významově vážou.

\blspace
  \dicEntry \dicTerm{tím|i} \dicPos{m} \dicFlx{(-a, -ar)} {\textbf{1.}} \dicSynonym{tíð} {čas, doba}; \textbf{í þann tíma} \dicFlx{adv} \dicSynonym{þá\textsuperscript{2}} {pak} {\textbf{2.}} \dicSynonym{rétt stund} {(správný) čas / moment}
\blspace

\subsection*{Za posledním očíslovaným významem}

Pokud se slovní spojení neváže k~žádnému významu, je řazeno za posledním očíslovaným významem. Ustálená slovní spojení jsou uvedena symbolem \dicsymIdiom.

\blspace
  \dicEntry \dicTerm{finna} \dicPos{v} \dicFlx{(finn, fann, fundum, fyndi, fundið)} {\textbf{3.}} \dicSynonym{skynja} {cítit, vnímat} \dicIdiom{finna}[að] {\dicPhraseIS{finna að e-u}} {kritizovat (co), nacházet chyby na (čem)}
\blspace

\subsection*{Způsob řazení slovních spojení}

Slovní spojení, která se nachází za posledním očíslovaným významem, jsou řazena podle předložky / příslovce / zájmena, se kterou se pojí. Pro snadnou orientaci v~rozsáhlých definicích je uvedeno heslové slovo nebo tvar heslového slova plus předložka / příslovce / zájmeno a~barevně odlišeno modrou barvou.

\blspace
  \dicEntry \dicTerm{koma\textsuperscript{2}} \dicPos{v} \dicFlx{(kem, kom, komum, kæmi, komið)} \dicIdiom{koma}[út] {\dicPhraseIS{koma út}} {vyjít (kniha ap.), objevit se (na trhu ap.)}
\blspace

Poté následují slovní spojení v~slovesném středním rodě řazená opět podle předložky / příslovce / zájmena, se kterou se pojí:

\blspace
  \dicEntry \dicTerm{koma\textsuperscript{2}} \dicPos{v} \dicFlx{(kem, kom, komum, kæmi, komið)} \dicIdiom{komast} {\dicPhraseIS{komast}} \dicFlx{refl} {dostat se, dostávat se}
\blspace

Poté následují ustálená slovní spojení, která nebylo možné zařadit do předchozích položek:

\blspace
  \dicEntry \dicTerm{koma\textsuperscript{2}} \dicPos{v} \dicFlx{(kem, kom, komum, kæmi, komið)} \dicIdiom{koma} {\dicPhraseIS{(komdu) sæll (og blessaður)}} {ahoj, nazdar, dobrý den}
\blspace

Nakonec jsou uvedena přísloví a~rčení označená symbolem \dicsymProverb\:

\blspace
  \dicEntry \dicTerm{eik} \dicPos{n} \dicFlx{(-ar / -ur, -ur)} \dicProverb{} \dicPhraseIS{Ekki / eigi fellur eik við fyrsta högg.} {\footnotesize{přís.}} {Jednou ranou dub nespadne.}
\blspace

\section{Synonyma a~antonyma / hesla ke~srovnání}
\longsectionskip

\subsection*{Synonyma}

Synonyma nebo synonymní slovní spojení jsou uvedeny před opisem významu hesla. Jsou umístěna v~závorkách. Islandská synonyma slouží a) islandskému uživateli k~rozlišení českých významů b) českému uživateli jako dodatečná informace o~významu slova.

\blspace
  \dicEntry \dicTerm{staða} \dicPos{f} \dicFlx{(stöðu, stöður)} {\textbf{1.}} \dicSynonym{ástand} {situace, stav, poměry} {\textbf{2.}} \dicSynonym{embætti} {post, (pracovní) místo / pozice (ve firmě ap.)}
\blspace

V~případě, že synonymum lze nalézt jako heslo ve slovníku, je použita šipka '\dicsymSee\ '.

\blspace
  \dicEntry \dicTerm{bygging} {\dicPos{f} \dicFlx{(-ar, -ar)} {\textbf{1.}} \dicSynonym{hús} {budova, stavba, stavení} {\textbf{2.}} \textit{(það að byggja)} {(vý)stavba} }
\blspace

\subsection*{Antonyma / Hesla ke~srovnání}

Antonyma či hesla ke~srovnání jsou uvedena za opisem významu heslového slova a~jsou označena šipkou '$\shortuparrow$'.

\blspace
  \dicEntry \dicTerm{kaldur} \dicPos{adj} \dicFlx{(f köld)} {\textbf{1.}} {studený, chladný} \textit{($\shortuparrow$ heitur)}
\blspace

\section{Oborové a~stylové charakteristiky}

K~upřesnění významu heslového slova slouží oborové a~stylové charakteristiky.

\subsection*{Oborové charakteristiky}

Oborové charakteristiky zařazují islandské slovo do určitého oboru a~naznačují použití daného slova v~jazyce

\blspace
  \dicEntry \dicTerm{berg} \dicPos{n} \dicFlx{(-s, -)} {\textbf{1.}} {\footnotesize{geol.}} {hornina, kámen} {\textbf{2.}} \dicSynonym{klöpp\textsuperscript{1}} {skalní stěna, skalisko}
\blspace

a~vymezují význam českého slova.

\blspace
  \dicEntry \dicTerm{mús} \dicPos{f} \dicFlx{(-ar, mýs)} {\textbf{1.}} {\footnotesize{zool.}} {myš} (\textit{l. Mus}) {\textbf{2.}} {\footnotesize{poč.}} {myš}
\blspace

\subsection*{Stylové charakteristiky}

Stylové charakteristiky jsou uváděny v~případě, kdy význam heslového slova je možné použít pouze při určité příležitosti nebo je význam jazykově zabarven.

\blspace
  \dicEntry \dicTerm{gaur} \dicPos{m} \dicFlx{(-s, -ar)} {\textbf{1.}} {\footnotesize{hovor.}} {chlápek, maník, člověk}
\blspace

\section{Syntax}

Syntaktické informace popisují:

\begin{enumerate}
\item zdali je sloveso tranzitivní nebo intranzitivní a~s~jakým(i) pádem (pády) se pojí, pokud je tranzitivní,
\item zdali je podmět / předmět životný či neživotný,
\item zdali je sloveso osobní či neosobní a~v~jakém pádu je podmět, pokud je sloveso neosobní,
\item zdali je podmět gramatický,
\item zdali je sloveso použito ve středním rodě,
\item nejčastější příklady, které se pojí se slovesem,
\item se kterými příslovci nebo předložkami se sloveso pojí.
\end{enumerate}

Syntaktické informace se ve slovníku vyskytují ve třech formách: a) informace uvedené zkratkou, b) teoretické příklady a~c) praktické příklady.

\subsection*{Syntaktické informace uvedené zkratkou}

Zde máme na mysli informace o~pádě (\textbf{nom, acc, dat, gen}), slovesu v~středním slovesném rodě (\textbf{refl}), neosobním slovesu (\textbf{impers}), čísle (\textbf{sg, pl}), rozkazovacím způsobu (\textbf{imper}) a~spojovacím způsobu (\textbf{con}).

\subsection*{Teoretické příklady}

Teoretickým příkladem nazýváme příklad, kdy podmět nebo předmět je zastoupen zájmenem neurčitým v~pádě, se~kterým se slovo pojí a~sloveso je zpravidla v~infinitivu. V~případě, že podmět není v~1.~pádě (např. u~sloves neosobních), je sloveso uvedeno v~3.~osobě jednotného čísla přítomného času.
Ve slovních spojeních používáme zkratky pro tvary neurčitého zájmena \textbf{einhver} (někdo) a~\textbf{eitthvað} (něco). Tyto zkratky označují a) pád, b) číslo a~c) životnost nebo neživotnost podmětu / předmětu.\\
Následuje seznam zkratek neurčitých zájmen použitých ve slovníku: \\

\tableD
\blspace
Příklad použití zkratek s~neurčitými zájmeny:

\blspace
  \dicEntry \dicTerm{leit|a} \dicPos{v} \dicFlx{(-aði)} \textbf{leita að e-u / e-m} {hledat (co / koho)} \dicExampleIS{leita að lyklunum} \dicExampleCS{hledat klíče}
\blspace

Příklad tedy přečteme jako '\textit{leita að einhverju}' (\textit{hledat co}) a~'\textit{leita að einhverjum}' (\textit{hledat koho}). V~tomto případě se sloveso pojí s~3.~pádem a~předmět může být životný i~neživotný.

\subsection*{Praktické příklady}

Ve slovníku se vyskytuje velké množství praktických příkladů a~jejich překladů. Praktický příklad se liší od příkladu teoretického tím, že místo zájmen neurčitých jsou uvedena frekventovaná slova z~běžného projevu. Praktické příklady slouží k~ilustraci použití
heslového slova.\\
Praktickým příkladem může být a) (složené) slovo, b) slovní spojení nebo c) věta.

\begin{enumerate}
\item Slovo jako praktický příklad:
\blspace
\dicEntry \dicTerm{hett|a} \dicPos{f} \dicFlx{(-u, -ur)} \textbf{1.} {kapuce, kápě} \dicExampleIS{hettuúlpa} \dicExampleCS{bunda s~kapucí}
\blspace
\item Slovní spojení jako praktický příklad:
\blspace
\dicEntry \dicTerm{tím|i} \dicPos{m} \dicFlx{(-a, -ar)} \textbf{3.} \dicSynonym{klukkustund} {hodina} {\footnotesize{(šedesát minut)}} \dicExampleIS{tveggja tíma gangur} \dicExampleCS{dvouhodinový pochod}
\blspace
\item Věta jako praktický příklad:
\blspace
\dicEntry \dicTerm{mæt|a} \dicPos{v} \dicFlx{(-ti, -t)} \dicIdiom{mætast} \dicPhraseIS{mætast} \dicSynonym*{ná saman} \dicFlx{refl} {setkat se, potkat se} \dicExampleIS{Þau mættust á gangstéttinni.} \dicExampleCS{Potkali se na chodníku.}
\end{enumerate}

\section{Fotografie a~ilustrace}

Ve slovníku je publikováno 630 fotografií rostlin, živočichů a~fotografií, která dokumentují místopisně heslová slova (např. kulinářské speciality).
Fotografie jsou publikovány pod veřejnými licencemi. Jména autorů a~názvy licencí lze nalézt v~Seznamu autorů fotografií a~ilustrací na straně \pageref{sec:photo}. U~hesla a~konkrétně u~daného významu, ke kterému se fotografie vztahuje, je použit symbol fotoaparátu \dicsymPhoto. Fotografie je umístěna buď pod heslem, anebo v~horní části stránky, kde se heslo nachází, anebo na stránce (stránkách) následující(ch).

%\par\begin{center}\setlength\fboxsep{0pt}\setlength\fboxrule{0.5pt}
%\fbox{\includegraphics[width=6cm]{ds_image_posthus_0_2.jpg}}\end{center}
%\par\begin{center}\footnotesize {Pósthús}\end{center}

\end{multicols}

% ---------------------------------------- List of abbreviations and symbols -

\chapter{Seznam použitých zkratek a symbolů}
\begin{multicols}{2}
\begin{description}[leftmargin=!,labelwidth=1.5cm]
%[\breaklabel\setleftmargin{40pt}\setlabelstyle{\bfseries}
%\addtolength{\itemsep}{-0.5\baselineskip}]
\item[{abb}] {zkrácený}
\item[{acc}] {4. pád, akuzativ}
\item[{adj}] {přídavné jméno}
\item[{adv}] {příslovce}
\item[{akt}] {rod činný, aktivum}
\item[{anat.}] {anatomie}
\item[{angl.}] {anglicky}
\item[{astro.}] {astronomie, výzkum vesmíru}
\item[{ap.}] {a podobně}
\item[{básn.}] {básnický, poetický výraz}
\item[{biol.}] {biologie}
\item[{bot.}] {botanika}

\item[{comp}] {2. stupeň stupňování, komparativ}
\item[{con}] {spojovací způsob, konjunktiv}
\item[{conj}] {spojka}
\item[{dat}] {3. pád, dativ}
\item[{def}] {určitý tvar}
\item[{dem}] {zájmeno ukazovací}
\item[{dět.}] {dětsky}

\item[{e-að}] {co (sg nom neživotný) (\textit{eitthvað})}
\item[{e-ð}] {co (sg acc neživotný) (\textit{eitthvað})}
\item[{e-ir}] {kdo (pl nom životný) (\textit{einhverjir})}
\item[{e-ja}] {koho (pl acc životný) (\textit{einhverja})}
\item[{e-jum}] {komu (pl dat životný) (\textit{einhverjum})}
\item[{e-m}] {komu (sg dat životný) (\textit{einhverjum})}
\item[{e-n}] {koho (sg acc životný) (\textit{einhvern})}
\item[{e-r}] {kdo (sg nom životný) (\textit{einhver})}
\item[{e-rra}] {koho (pl gen životný) (\textit{einhverra})}
\item[{e-rs}] {koho (sg gen životný) (\textit{einhvers})}
\item[{e-s}] {čeho (sg gen neživotný) (\textit{einhvers})}
\item[{e-u}] {čemu (sg dat neživotný) (\textit{einhverju})}

\item[{ekon.}] {ekonomika, obchod}
\item[{elek.}] {elektřina}
\item[{f}] {rod ženský}
\item[{filos.}] {filosofie, logika}
\item[{form.}] {formálně}
\item[{fyz.}] {fyzika}
\item[{gen}] {2. pád, genitiv}
\item[{geog.}] {zeměpis, geografie}
\item[{geol.}] {geologie}
\item[{han.}] {hanlivě, pejorativně}
\item[{hist.}] {historie}
\item[{hovor.}] {hovorově}
\item[{hrub.}] {hrubě}
\item[{hud.}] {hudba}
\item[{imper}] {rozkazovací způsob, imperativ}
\item[{impers}] {sloveso neosobní}
\item[{ind}] {oznamovací způsob, indikativ}
\item[{indecl}] {nesklonný}
\item[{indef}] {neurčitý tvar, (pron +) zájmeno neurčité}
\item[{inf}] {infinitiv}
\item[{int}] {tázací zájmeno}
\item[{inter}] {citoslovce}
\item[{chem.}] {chemie}
\item[{jaz.}] {jazykověda}
\item[{kulin.}] {kulinářství, vaření}
\item[{l.}] {latinsky}
\item[{let.}] {letectví}
\item[{lit.}] {literatura, vydavatelství}
\item[{m}] {rod mužský}
\item[{mat.}] {matematika}
\item[{med. }] {lékařství}
\item[{meteo.}] {meteorologie}
\item[{myt.}] {mytologie}
\item[{n}] {rod střední}
\item[{náb.}] {náboženství}
\item[{nám.}] {námořnictví, rybolov}
\item[{nom}] {1. pád, nominativ}
\item[{num}] {číslovka}
\item[{ord}] {řadová číslovka}
\item[{p}] {osoba}
\item[{part}] {částice}
\item[{pers}] {(v +) osoba, (pron +) zájmeno osobní}
\item[{pf}] {minulý čas}
\item[{pl}] {množné číslo, plurál}
\item[{poč.}] {informatika}
\item[{pol.}] {politika, politologie}
\item[{poss}] {zájmeno přivlastňovací}
\item[{pos}] {1. stupeň stupňování, pozitiv}
\item[{pov.}] {lidové pověsti, folkloristika}
\item[{pp}] {příčestí minulé}
\item[{praes}] {přítomný čas}
\item[{práv.}] {právnictví, soudnictví}
\item[{predp}] {předpona}
\item[{prep}] {předložka}
\item[{presp}] {příčestí přítomné}
\item[{pron}] {zájmeno}
\item[{prop}] {vlastní jméno, proprium}
\item[{přen.}] {přeneseně}
\item[{přís.}] {přísloví}
\item[{psych.}] {psychologie}
\item[{refl}] {(v +) médium, (pron +) zvratné zájmeno}
\item[{rel}] {vztažné zájmeno}
\item[{sg}] {jednotné číslo, singulár}
\item[{slang.}] {slang}
\item[{sport.}] {sport}
\item[{stav.}] {stavebnictví, architektura}
\item[{subs}] {podstatné jméno (bez rodu)}
\item[{sup}] {3. stupeň stupňování, superlativ}
\item[{supin}] {supinum}
\item[{škol.}] {školství}
\item[{techn.}] {technika, mechanika}
\item[{v}] {sloveso}
\item[{voj.}] {vojenství}
\item[{zast.}] {zastarale}
\item[{zkr}] {zkratka}
\item[{zool.}] {zoologie}

\item[{;}] {středník uprostřed definice odděluje věty; středník na konci definice odděluje významy}

\item[{(...)}] {v kulatých závorkách jsou uvedeny doplňkové informace}
\item[{[...]}] {v hranatých závorkách je uveden fonetický zápis výslovnosti}\footnote{Vysvětlení fonetických symbolů je uvedeno v kapitole Fonetika \ref{sec:phon_phonems} na straně \pageref{sec:phon_phonems}}
\item[{|}] {naznačuje místo, kde se deklinační koncovka připojuje ke slovnímu základu}
\item[{·}] {vedlejší dělení složeného slova}
\item[{··}] {hlavní dělení složeného slova}
\item[{+}] {plus}
\item[{/}] {lomítko označuje více možností}
\item[{\dicsymSee}] {šipka odkazuje na heslo ve slovníku}
\item[{\dicsymCompare}] {šipka odkazuje na heslo ve slovníku ke srovnání}
\item[{\dicsymExampleIS}] {uvozuje islandský příklad}
\item[{\dicsymExampleCS}] {uvozuje český překlad příkladu}
\item[{\dicsymIdiom}] {uvozuje slovní spojení}
\item[{\dicsymFrequent}] {heslo patří mezi frekventované výrazy}

\item[{\dicsymPhoto}] {fotografie, ilustrace}
\item[{\dicsymProverb}] {přísloví, rčení}

\end{description}
\end{multicols}

\chapter{Skrá yfir skammstafanir og tákn}
\begin{multicols}{2}
\begin{description}[leftmargin=!,labelwidth=1.5cm]
%[\breaklabel\setleftmargin{40pt}\setlabelstyle{\bfseries}
%\addtolength{\itemsep}{-0.5\baselineskip}]
\item[{abb}] {stýfður}
\item[{acc}] {þolfall}
\item[{adj}] {lýsingarorð}
\item[{adv}] {atviksorð}
\item[{akt}] {germynd}
\item[{anat.}] {líffærafræði}
\item[{angl.}] {enska}
\item[{astro.}] {geimvísindi, stjörnufræði}
\item[{ap.}] {e.þ.h.}
\item[{básn.}] {ljóðrænn}
\item[{biol.}] {líffræði}
\item[{bot.}] {grasafræði}

\item[{comp}] {miðstig}
\item[{con}] {viðtengingarháttur}
\item[{conj}] {samtenging}
\item[{dat}] {þágufall}
\item[{def}] {með greini}
\item[{dem}] {ábendingarfornafn}
\item[{dět.}] {barnamál}

\item[{e-að}] {(\textit{eitthvað})}
\item[{e-ð}] {(\textit{eitthvað})}
\item[{e-ir}] {(\textit{einhverjir})}
\item[{e-ja}] {(\textit{einhverja})}
\item[{e-jum}] {(\textit{einhverjum})}
\item[{e-m}] {(\textit{einhverjum})}
\item[{e-n}] {(\textit{einhvern})}
\item[{e-r}] { (\textit{einhver})}
\item[{e-rra}] {(\textit{einhverra})}
\item[{e-rs}] {(\textit{einhvers})}
\item[{e-s}] {(\textit{einhvers})}
\item[{e-u}] { (\textit{einhverju})}

\item[{ekon.}] {hagfræði, viðskipti}
\item[{elek.}] {rafmagn}
\item[{f}] {kvenkyn}
\item[{filos.}] {heimspeki, rökfræði}
\item[{form.}] {formlegt mál}
\item[{fyz.}] {eðlisfræði}
\item[{gen}] {eignarfall}
\item[{geog.}] {landafræði}
\item[{geol.}] {jarðfræði}
\item[{han.}] {niðrandi}
\item[{hist.}] {sagnfræði}
\item[{hovor.}] {talmál}
\item[{hrub.}] {gróft}
\item[{hud.}] {tónlist}
\item[{imper}] {boðháttur}
\item[{impers}] {ópersónuleg sögn}
\item[{ind}] {framsöguháttur}
\item[{indecl}] {óbeygjanlegur}
\item[{indef}] {óákveðinn, (pron +) óákveðið fornafn}
\item[{inf}] {nafnháttur}
\item[{int}] {spurnigarfornafn}
\item[{inter}] {upphrópun}
\item[{chem.}] {efnafræði}
\item[{jaz.}] {málfræði, málvísindi}
\item[{kulin.}] {eldhús}
\item[{l.}] {latína}
\item[{let.}] {flug, flugmál}
\item[{lit.}] {skáldskapur}
\item[{m}] {nafnorð, karlkyn}
\item[{mat.}] {stærðfræði}
\item[{med. }] {læknisfræði}
\item[{meteo.}] {veðurfræði}
\item[{myt.}] {goðafræði}
\item[{n}] {nafnorð, hvorugkyn}
\item[{náb.}] {trúfræði}
\item[{nám.}] {sjómennska}
\item[{nom}] {nefnifall}
\item[{num}] {töluorð}
\item[{ord}] {raðtala}
\item[{p}] {persóna}
\item[{part}] {ögn}
\item[{pers}] {(v +) persóna, (pron +) persónufornafn}
\item[{pf}] {þátíð}
\item[{pl}] {fleirtala}
\item[{poč.}] {tölvufræði, tölvur}
\item[{pol.}] {stjórnmál, stjórnmálafræði}
\item[{poss}] {eignarfornafn}
\item[{pos}] {frumstig}
\item[{pov.}] {þjóðtrú}
\item[{pp}] {lýsingarháttur þátíðar}
\item[{praes}] {nútíð}
\item[{práv.}] {lögfræði, dómsmál}
\item[{predp}] {forskeyti}
\item[{prep}] {forsetning}
\item[{presp}] {lýsingarháttur nútíðar}
\item[{pron}] {fornafn}
\item[{prop}] {sérnafn}
\item[{přen.}] {afleidd merking}
\item[{přís.}] {málsháttur}
\item[{psych.}] {sálfræði}
\item[{refl}] {(v +) miðmynd, (pron +) afturbeygt fornafn}
\item[{rel}] {tilvísunarfornafn}
\item[{sg}] {eintala}
\item[{slang.}] {slangur}
\item[{sport.}] {íþróttir}
\item[{stav.}] {húsasmíði, byggingarlist}
\item[{subs}] {nafnorð (án kyns)}
\item[{sup}] {efstastig}
\item[{supin}] {sagnbót}
\item[{škol.}] {skólamál}
\item[{techn.}] {tækni, verkfræði}
\item[{v}] {sagnorð}
\item[{voj.}] {her}
\item[{zast.}] {úrelt mál}
\item[{zkr}] {skammstöfun}
\item[{zool.}] {dýrafræði}

\item[{;}] {semíkomma í miðri skýringu aðgreinir setningu; semíkomma í enda skýringar aðgreinir mismunandi merkingarsvið}

\item[{(...)}] {innan sviga eru viðbótarupplýsingar}
\item[{[...]}] {innan hornklofa er hljóðritaður framburður}
\item[{|}] {merkir stað, þar sem beygingarending bætist við stofn}
\item[{·}] {aukasamskeyti samsetts orðs}
\item[{··}] {aðalsamskeyti samsetts orðs}
\item[{+}] {plús}
\item[{/}] {skástrik táknar fleiri möguleika}
\item[{\dicsymSee}] {ör vísar til uppflettiorðs í orðabókinni}
\item[{\dicsymCompare}] {ör vísar til uppflettiorðs í bókinni til samanburðar}
\item[{\dicsymExampleIS}] {kemur á undan dæmi á íslensku}
\item[{\dicsymExampleCS}] {kemur á undan tékkneskri þýðingu á dæminu}
\item[{\dicsymIdiom}] {orð með háa tíðni}

\item[{\dicsymPhoto}] {ljósmynd}
\item[{\dicsymProverb}] {málsháttur}

\end{description}
\end{multicols}

% ------------------------------------------------- Headword part and alphabet

\chapter{Heslová část}
{\renewcommand{\arraystretch}{2}%
%\topskip0pt
\vstretch{1}
{\setlength{\tabcolsep}{1.5em}
\begin{table}[htpb]
\centering

\begin{tabular}{ccccccc} \hline
A a &  & G g &  & O o &  & V v \\
Á á  &  & H h &  & Ó ó &  & W w \\
B b &  & I i &  & P p &  & X x \\
C c &  & Í í &  & Q q &  & Y y \\
D d &  & J j &  & R r &  & Ý ý \\
Ð ð &  & K k &  & S s &  & Z z \\
E e &  & L l &  & T t &  & Þ þ \\
É é  &  & M m &  & U u &  & Æ æ \\
F f &  & N n &  & Ú ú &  & Ö ö \\
\hline
\end{tabular}

%\caption{\textbf{Islandská abeceda}}
\label{my-islandska-abeceda}
\end{table}
}
\vstretch{1}
{\renewcommand{\arraystretch}{1}%

\clearpage

\fi

% =========================================================== THE DICTIONARY =

\ifinputletters

\cleardoublepage

\dictionarygeometry
\pagestyle{myheadings}

\newcommand*{\alphabetF}{%
  a,aa,b,c,d,e,ee,f,g,h,i,ii,j,k,l,m,
  n,o,oo,p,r,s1,s2, t,u,uu,v,w,x,y,yy,th,ae,oe}

\ifPDFtest
\dicLetter{h}{letter10}
\dicEntry[H.Í.] \dicTerm{H.Í.} \dicPos{zkr} \dicPhraseIS{Háskóli Íslands} \dicDirectTranslationCS{Islandská univerzita}
\dicEntry[h.u.b.] \dicTerm{h.u.b.} \dicPos{zkr} \dicPhraseIS{hér um bil} \dicFlx{adv} \dicDirectTranslationCS{kolem, zhruba, přibližně}
\dicEntry[ha] \dicTerm{ha} \dicIPA{{h}{a}{\textlengthmark}} \dicPos{inter} \dicDirectTranslationCS{co?, cože?}
\dicEntry[ha ha] \dicTerm{ha ha} \dicIPA{{h}{a}{h}{a}} \dicPos{inter} \dicDirectTranslationCS{chacha, chichi}
\dicEntry[haf] \dicTerm{haf} \dicsymFrequent\  \dicIPA{{h}{a}{\textlengthmark}{f}} \dicPos{n}[2] \dicFlx{(‑s, höf)}[8] \dicDirectTranslationCS{moře, oceán} \dicExampleIS{í hafi} \dicExampleCS{na moři};  \dicPhraseIS{opið haf} \dicDirectTranslationCS{otevřené moře};  \dicPhraseIS{láta í haf} \dicSynonym{sigla\smash{\textsuperscript{2}}} \dicDirectTranslationCS{vyplout na moře}
\dicEntry[hafa] \dicTerm{hafa} \dicsymFrequent\  \dicIPA{{h}{a}{\textlengthmark}{v}{a}} \dicPos{v}[4] \dicFlx{(hef, hafði, höfðum, hefði, haft)}[18] \dicFlx{acc} \textbf{1.} \dicSynonym*{vera með} \dicDirectTranslationCS{mít, vlastnit} \dicExampleIS{hafa hníf í hendi} \dicExampleCS{mít nůž v~ruce}  \textbf{2.} \dicIndirectTranslationCS{pomocné sloveso -- se supinem vyjadřuje předpřítomný čas a~předminulý čas} \dicExampleIS{Hann hefur aldrei komið til Grænlands.} \dicExampleCS{Nikdy nebyl v~Grónsku.}  \textbf{3.} \dicPhraseIS{hafa sig} \dicDirectTranslationCS{odejít, ztratit se};  \dicIdiom{hafa}[að]{ \dicPhraseIS{hafa e‑ð að engu}} \dicDirectTranslationCS{nevšímat si (čeho), ignorovat (co)};  \dicIdiom{hafa}[af]{ \dicPhraseIS{hafa e‑ð af}} \dicDirectTranslationCS{udělat (co), podařit se (co)}; { \dicPhraseIS{hafa það af}} \dicDirectTranslationCS{přežít, vyváznout};  \dicIdiom{hafa}[á móti]{ \dicPhraseIS{hafa ekkert á móti e‑u}} \dicDirectTranslationCS{nemít nic proti (čemu)};  \dicIdiom{hafa}[eftir]{ \dicPhraseIS{hafa e‑ð eftir e‑m}} \dicDirectTranslationCS{mít (co) od (koho), opakovat (co) po (kom)};  \dicIdiom{hafa}[fram]{ \dicPhraseIS{hafa sig fram}} \dicDirectTranslationCS{prosadit svou (vůli)}; { \dicPhraseIS{hafa e‑ð fram}} \dicDirectTranslationCS{prosadit (co)};  \dicIdiom{hafa}[fram yfir]{ \dicPhraseIS{hafa e‑ð fram yfir e‑n}} \dicDirectTranslationCS{být lepší v~(čem) než (kdo)};  \dicIdiom{hafa}[fyrir]{ \dicPhraseIS{hafa fyrir e‑u}} \dicDirectTranslationCS{vynaložit na (co) úsilí}; { \dicPhraseIS{hafa e‑ð fyrir sér}} \dicDirectTranslationCS{opřít se o~(co) (o~vědecký základ ap.)};  \dicIdiom{hafa}[í]{ \dicPhraseIS{hafa í sig og á}} \dicDirectTranslationCS{mít na živobytí, uživit se} \dicIndirectTranslationCS{(používané zvláště se záporem)}; { \dicPhraseIS{hafa sig í (að hringja)}} \dicDirectTranslationCS{sebrat se (a~zavolat)};  \dicIdiom{hafa}[í frammi]{ \dicPhraseIS{hafa sig lítið í frammi}} \dicDirectTranslationCS{držet se zpátky};  \dicIdiom{hafa}[ofan af]{ \dicPhraseIS{hafa ofan af fyrir (sér) með e‑u}} \dicDirectTranslationCS{(za)bavit (se) (čím)};  \dicIdiom{hafa}[til]{ \dicPhraseIS{hafa sig til fyrir e‑ð}} \dicDirectTranslationCS{připravit se na (co), vystrojit se na (co)};  \dicIdiom{hafa}[undan]{ \dicPhraseIS{hafa undan e‑m}} \dicDirectTranslationCS{(u)držet krok s~(kým), stíhat (koho)};  \dicIdiom{hafa}[undir]{ \dicPhraseIS{hafa e‑ð undir höndum}} \dicDirectTranslationCS{mít (co) ve své péči}; { \dicPhraseIS{hafa e‑n undir}} \dicDirectTranslationCS{svrhnout (koho)};  \dicIdiom{hafa}[upp á sig]{ \dicPhraseIS{e‑að hefur lítið upp á sig}} \dicDirectTranslationCS{(co) se nevyplatí};  \dicIdiom{hafa}[uppi á]{ \dicPhraseIS{hafa uppi á e‑u\,/\addthin e‑m}} \dicSynonym{finna} \dicDirectTranslationCS{nalézt (co\,/\addthin koho)};  \dicIdiom{hafa}[upp úr]{ \dicPhraseIS{hafa e‑ð upp úr e‑u}} \dicDirectTranslationCS{mít (co) z~(čeho) (hodně peněz z~práce ap.)}; { \dicPhraseIS{hafa mikið upp úr sér}} \dicDirectTranslationCS{vydělat hodně peněz};  \dicIdiom{hafa}[út undan]{ \dicPhraseIS{hafa e‑n út undan}} \dicDirectTranslationCS{znevýhodnit (koho), držet (koho) stranou};  \dicIdiom{hafa}[út úr]{ \dicPhraseIS{hafa út úr e‑m e‑ð}} \dicDirectTranslationCS{vymámit (co) z~(koho)};  \dicIdiom{hafa}[við]{ \dicPhraseIS{hafa mikið við}} \dicDirectTranslationCS{udělat slavnostním}; { \dicPhraseIS{hafa við e‑m}} \dicDirectTranslationCS{mít na (koho), stačit (komu)};  \dicIdiom{hafa}[yfir]{ \dicPhraseIS{hafa yfir e‑ð}} \dicDirectTranslationCS{(od)recitovat (co)};  \dicIdiom{hafast}{ \dicPhraseIS{hafast}} \dicFlx{refl} \dicDirectTranslationCS{podařit se, povést se};  \dicIdiom{hafast}[að]{ \dicPhraseIS{hafast e‑ð að}} \dicFlx{refl} \dicDirectTranslationCS{provést (co), udělat (co)};  \dicIdiom{hafast}[við]{ \dicPhraseIS{e‑að hefst (vel) við}} \dicFlx{refl} \dicDirectTranslationCS{(co) se (dobře) hojí}; { \dicPhraseIS{hafast við e‑s staðar}} \dicFlx{refl} \dicDirectTranslationCS{zdržovat se (kde), přebývat (kde)};  \dicIdiom{hafa}{ \dicPhraseIS{hafðu það gott!}} \dicDirectTranslationCS{měj se dobře!}; { \dicPhraseIS{það verður að hafa það}} \dicDirectTranslationCS{nedá se s~tím nic dělat}
\dicEntry[hafbeit] \dicTerm{haf··beit} \dicIPA{{h}{a}{v}{\textsubring{b}}{ei}{\textsubring{d}}} \dicPos{f}[7] \dicFlx{(‑ar)}[3] \dicDirectTranslationCS{mořská akvakultura} \dicIndirectTranslationCS{(chov ryb na moři)}
\dicEntry[hafði] \dicTerm{hafði} \dicIPA{{h}{a}{v}{ð}{\textsci}} \dicPos{v} \dicFlx{ind pf sg 1 pers} \dicLink{hafa}
\dicEntry[haffræði] \dicTerm{haf··fræð|i} \dicIPA{{h}{a}{f}{\textlengthmark}{r}{a}{i}{ð}{\textsci}} \dicPos{f}[3] \dicFlx{(‑i)}[3] \dicDirectTranslationCS{oceánografie}
\dicEntry[haffræðingur] \dicTerm{haf·fræð··ing|ur} \dicIPA{{h}{a}{f}{\textlengthmark}{r}{a}{i}{ð}{i}{\ng}{\r{g}}{\textscy}{\textsubring{r}}} \dicPos{m}[6] \dicFlx{(‑s, ‑ar)}[8] \dicDirectTranslationCS{oceánograf(ka)}
\dicEntry[hafgola] \dicTerm{haf··gol|a} \dicIPA{{h}{a}{v}{\r{g}}{\textopeno}{l}{a}} \dicPos{f}[1] \dicFlx{(‑u, ‑ur)}[7] \dicDirectTranslationCS{mořský vánek}
\dicEntry[hafið] \dicTerm{hafið} \dicIPA{{h}{a}{\textlengthmark}{v}{\textsci}{\texttheta}} \dicPos{v} \dicFlx{supin} \dicLink{hefja}
\dicEntry[hafís] \dicTerm{haf··ís} \dicIPA{{h}{a}{\textlengthmark}{v}{i}{s}} \dicPos{m}[4] \dicFlx{(‑s, ‑ar)}[25] \dicDirectTranslationCS{plovoucí led}
\dicEntry[hafna] \dicTerm{hafn|a\smash{\textsuperscript{1}}} \dicsymFrequent\  \dicIPA{{h}{a}{\textsubring{b}}{n}{a}} \dicPos{v}[1] \dicFlx{(‑aði)}[50] \dicFlx{dat} \dicSynonym{neita} \dicDirectTranslationCS{odmítnout, odmítat, zamítnout, zamítat, neschválit, neschvalovat} \dicExampleIS{hafna tilboðinu} \dicExampleCS{odmítnout nabídku}
\dicEntry[hafna] \dicTerm{hafn|a\smash{\textsuperscript{2}}} \dicIPA{{h}{a}{\textsubring{b}}{n}{a}} \dicPos{v}[1] \dicFlx{(‑aði)}[50] \textbf{1.} \dicPhraseIS{hafna sig} \dicFieldCat{nám.} \dicDirectTranslationCS{přistát, zakotvit (loď ap.)}  \textbf{2.} \dicSynonym{lenda} \dicDirectTranslationCS{skončit, přistát (míč v~síti ap.)}
\dicEntry[hafnabolti] \dicTerm{hafna··bolt|i} \dicIPA{{h}{a}{\textsubring{b}}{n}{a}{\textsubring{b}}{\textopeno}{\textsubring{l}}{\textsubring{d}}{\textsci}} \dicPos{m}[1] \dicFlx{(‑a, ‑ar)}[1] \dicFieldCat{sport.} \dicDirectTranslationCS{baseball}
\dicEntry[hafnar] \dicTerm{hafnar} \dicIPA{{h}{a}{\textsubring{b}}{n}{a}{\textsubring{r}}} \dicPos{f} \dicFlx{sg gen} \dicLink{höfn}
\dicEntry[hafnarbakki] \dicTerm{hafnar··bakk|i} \dicIPA{{h}{a}{\textsubring{b}}{n}{a}{r}{\textsubring{b}}{a}{h}{\r{\textObardotlessj}}{\textsci}} \dicPos{m}[1] \dicFlx{(‑a, ‑ar)}[8] \dicFieldCat{nám.} \dicDirectTranslationCS{přístaviště, přístavní hráz}
\dicEntry[Hafnarfjörður] \dicTerm{Hafnar··|fjörður} \dicIPA{{h}{a}{\textsubring{b}}{n}{a}{\textsubring{r}}{f}{j}{\oe}{r}{ð}{\textscy}{\textsubring{r}}} \dicPos{m}[12] \dicFlx{(‑fjarðar)}[2] \dicFieldCat{geog.} \dicDirectTranslationCS{Hafnarfjörður} \dicIndirectTranslationCS{(druhé největší město Islandu nedaleko Reykjavíku)}
\dicEntry[hafnargarður] \dicTerm{hafnar··garð|ur} \dicIPA{{h}{a}{\textsubring{b}}{n}{a}{r}{\r{g}}{a}{r}{ð}{\textscy}{\textsubring{r}}} \dicPos{m}[6] \dicFlx{(‑s, ‑ar)}[5] \dicFieldCat{nám.} \dicDirectTranslationCS{molo}
\dicEntry[hafnarstjóri] \dicTerm{hafnar··stjór|i} \dicIPA{{h}{a}{\textsubring{b}}{n}{a}{\textsubring{r}}{s}{\textsubring{d}}{j}{ou}{r}{\textsci}} \dicPos{m}[1] \dicFlx{(‑a, ‑ar)}[1] \dicFieldCat{nám.} \dicDirectTranslationCS{přístavní kapitán(ka)}
\dicEntry[hafnarstjórn] \dicTerm{hafnar··stjórn} \dicIPA{{h}{a}{\textsubring{b}}{n}{a}{\textsubring{r}}{s}{\textsubring{d}}{j}{ou}{r}{\textsubring{d}}{\textsubring{n}}} \dicPos{f}[7] \dicFlx{(‑ar, ‑ir)}[1] \dicFieldCat{nám.} \dicDirectTranslationCS{přístavní správa\,/\addthin úřad}
\dicEntry[hafnarverkamaður] \dicTerm{hafnar··verka·|maður} \dicIPA{{h}{a}{\textsubring{b}}{n}{a}{r}{v}{\textepsilon}{\textsubring{r}}{\r{g}}{m}{a}{ð}{\textscy}{\textsubring{r}}} \dicPos{m}[13] \dicFlx{(‑manns, ‑menn)}[2] \dicFieldCat{nám.} \dicDirectTranslationCS{dokař(ka), přístavní dělník\,/\addthin dělnice}
\dicEntry[hafnir] \dicTerm{hafnir} \dicIPA{{h}{a}{\textsubring{b}}{n}{\textsci}{\textsubring{r}}} \dicPos{f} \dicFlx{pl nom} \dicLink{höfn}
\begin{xtolerant}{}{1pt}
\dicEntry[hafnsögumaður] \dicTerm{hafn·sögu··|maður}\addthinS\dicIPA{{h}{a}{\textsubring{b}}{\textsubring{n}}{s}{\oe}{\textbabygamma}{\textscy}{m}{a}{ð}{\textscy}{\textsubring{r}}}\addthinS\dicPos{m}[13]\addthinS\dicFlx{(‑manns, ‑menn)}[2] \dicFieldCat{nám.} \dicDirectTranslationCS{lodivod(ka)}
\end{xtolerant}
\dicEntry[hafragrautur] \dicTerm{hafra··graut|ur} \dicIPA{{h}{a}{v}{r}{a}{\r{g}}{r}{\oe i}{\textsubring{d}}{\textscy}{\textsubring{r}}} \dicPos{m}[6] \dicFlx{(‑ar, ‑ar)}[59] \dicDirectTranslationCS{ovesná kaše}
\dicEntry[hafragrjón] \dicTerm{hafra··grjón} \dicIPA{{h}{a}{v}{r}{a}{\r{g}}{r}{j}{ou}{\textsubring{n}}} \dicPos{n}[2] \dicFlx{pl}[1] \dicDirectTranslationCS{ovesné kroupy}
\dicEntry[haframjöl] \dicTerm{hafra··mjöl} \dicIPA{{h}{a}{v}{r}{a}{m}{j}{\oe}{\textsubring{l}}} \dicPos{n}[2] \dicFlx{(‑s)}[2] \dicDirectTranslationCS{ovesná mouka}
\dicEntry[hafrannsóknir] \dicTerm{haf··rann·sóknir} \dicIPA{{h}{a}{v}{r}{a}{n}{s}{ou}{h}{\r{g}}{n}{\textsci}{\textsubring{r}}} \dicPos{f}[7] \dicFlx{pl}[2] \dicDirectTranslationCS{mořský výzkum}
\dicEntry[hafrar] \dicTerm{hafrar} \dicIPA{{h}{a}{v}{r}{a}{\textsubring{r}}} \dicPos{m}[5] \dicFlx{pl}[5] \dicFieldCat{bot.} \dicDirectTranslationCS{oves, oves setý} \textit{(l.~{\textLA{Avena sativa}})}
\dicEntry[hafsbotn] \dicTerm{hafs··botn} \dicIPA{{h}{a}{f}{s}{\textsubring{b}}{\textopeno}{h}{\textsubring{d}}{\textsubring{n}}} \dicPos{m}[4] \dicFlx{(‑s, ‑ar)}[3] \dicDirectTranslationCS{mořské dno, dno oceánu}
\dicEntry[haft] \dicTerm{haft\smash{\textsuperscript{1}}} \dicIPA{{h}{a}{f}{\textsubring{d}}} \dicPos{n}[2] \dicFlx{(‑s, höft)}[8] \textbf{1.} \dicSynonym{fjötur} \dicDirectTranslationCS{pouto}  \textbf{2.} \dicSynonym{hömlun} \dicDirectTranslationCS{omezení, restrikce}
\dicEntry[haft] \dicTerm{haft\smash{\textsuperscript{2}}} \dicIPA{{h}{a}{f}{\textsubring{d}}} \dicPos{v} \dicFlx{supin} \dicLink{hafa}
\dicEntry[haftastefna] \dicTerm{hafta··stefn|a} \dicIPA{{h}{a}{f}{\textsubring{d}}{a}{s}{\textsubring{d}}{\textepsilon}{\textsubring{b}}{n}{a}} \dicPos{f}[1] \dicFlx{(‑u)}[5] \dicDirectTranslationCS{protekcionismus, ochranářství}
\dicEntry[haftyrðill] \dicTerm{haf··tyrð|ill} \dicIPA{{h}{a}{f}{t\smash{\textsuperscript{h}}}{\textsci}{r}{ð}{\textsci}{\textsubring{d}}{\textsubring{l}}} \dicPos{m}[6] \dicFlx{(‑ils, ‑lar)}[35] \dicFieldCat{zool.} \dicDirectTranslationCS{alkoun malý} \textit{(l.~{\textLA{Alle alle}})}  \dicsymPhoto\ 
\dicFigure{ds_image_haftyrdill_0_1.jpg}{Haftyrðill}{Haftyrðill - Michael Haferkamp, PD}
\dicEntry[hafur] \dicTerm{haf|ur} \dicIPA{{h}{a}{\textlengthmark}{v}{\textscy}{\textsubring{r}}} \dicPos{m}[5] \dicFlx{(‑urs, ‑rar)}[3] \dicDirectTranslationCS{kozel}
\dicEntry[hafurtask] \dicTerm{hafur··task} \dicIPA{{h}{a}{}{v}{\textscy}{\textsubring{r}}{t\smash{\textsuperscript{h}}}{a}{s}{\r{g}}} \dicPos{n}[2] \dicFlx{(‑s)}[2] \dicSynonym{farangur} \dicDirectTranslationCS{věci, majetek} \dicExampleIS{með allt sitt hafurtask} \dicExampleCS{s~veškerými (svými) věcmi}
\dicEntry[haförn] \dicTerm{haf··|örn} \dicIPA{{h}{a}{\textlengthmark}{v}{\oe}{r}{\textsubring{d}}{\textsubring{n}}} \dicPos{m}[11] \dicFlx{(‑arnar, ‑ernir)}[9] \dicFieldCat{zool.} \dicDirectTranslationCS{orel mořský} \textit{(l.~{\textLA{Haliaeetus albicilla}})}  \dicsymPhoto\ 
\dicFigure{ds_image_haforn_0_1.jpeg}{Haförn}{Haförn - Tobias Biehl, PD}
\dicEntry[haga] \dicTerm{hag|a} \dicsymFrequent\  \dicIPA{{h}{a}{\textlengthmark}{\textbabygamma}{a}} \dicPos{v}[1] \dicFlx{(‑aði)}[28] \dicFlx{dat} \textbf{1.} \dicPhraseIS{haga sér} \dicSynonym{hegða} \dicDirectTranslationCS{chovat se} \dicExampleIS{haga sér vel} \dicExampleCS{chovat se dobře}  \textbf{2.} \dicSynonym*{skipa fyrir} \dicDirectTranslationCS{uzpůsobit, uspořádat} \dicExampleIS{haga þessu þannig} \dicExampleCS{uspořádat to takto};  \dicPhraseIS{það hagar svo til} \dicDirectTranslationCS{situace je taková, že}
\dicEntry[hagamús] \dicTerm{haga··|mús} \dicIPA{{h}{a}{\textlengthmark}{\textbabygamma}{a}{m}{u}{s}} \dicPos{f}[9] \dicFlx{(‑músar, ‑mýs)}[6] \dicFieldCat{zool.} \dicDirectTranslationCS{myšice křovinná} \textit{(l.~{\textLA{Apodemus sylvaticus}})}  \dicsymPhoto\ 
\dicFigure{ds_image_hagamus_0_2.jpg}{Hagamús}{Hagamús - Hans Hillewaert, CC BY-SA 3.0}
\dicEntry[haganlegur] \dicTerm{hagan··legur} \dicIPA{{h}{a}{}{\textbabygamma}{a}{n}{l}{\textepsilon}{\textbabygamma}{\textscy}{\textsubring{r}}} \dicPos{adj}[1]\dicFlx{}[-8] \dicSynonym*{hagfelldur} \dicDirectTranslationCS{praktický, příhodný}
\dicEntry[hagfræði] \dicTerm{hag··fræð|i} \dicIPA{{h}{a}{x}{f}{r}{a}{i}{ð}{\textsci}} \dicPos{f}[3] \dicFlx{(‑i)}[3] \dicSynonym{viðskiptafræði} \dicDirectTranslationCS{ekonomie}
\dicEntry[hagfræðilegur] \dicTerm{hag·fræði··legur} \dicIPA{{h}{a}{x}{f}{r}{a}{i}{ð}{\textsci}{l}{\textepsilon}{\textbabygamma}{\textscy}{\textsubring{r}}} \dicPos{adj}[1]\dicFlx{}[-8] \dicDirectTranslationCS{ekonomický}
\dicEntry[hagfræðingur] \dicTerm{hag·fræð··ing|ur} \dicIPA{{h}{a}{x}{f}{r}{a}{i}{ð}{i}{\ng}{\r{g}}{\textscy}{\textsubring{r}}} \dicPos{m}[6] \dicFlx{(‑s, ‑ar)}[8] \dicDirectTranslationCS{ekonom(ka)}
\dicEntry[hagga] \dicTerm{hagg|a} \dicIPA{{h}{a}{\r{g}}{\textlengthmark}{a}} \dicPos{v}[1] \dicFlx{(‑aði)}[13] \dicFlx{dat} \dicSynonym{hreyfa} \dicDirectTranslationCS{(po)hnout} \dicIndirectTranslationCS{(často se záporem)} \dicExampleIS{hagga e‑u} \dicExampleCS{pohnout (s) (čím)};  \dicPhraseIS{e‑m verður ekki haggað} \dicFlx{impers} \dicLangCat{přen.} \dicDirectTranslationCS{s~(kým) se nedá hnout};  \dicIdiom{haggast}{ \dicPhraseIS{haggast}} \dicFlx{refl} {\textbf{a.}} \dicDirectTranslationCS{(po)hnout se,  hýbat se};  {\textbf{b.}} \dicDirectTranslationCS{hnout, vyvolat pohnutí} \dicIndirectTranslationCS{(se záporem)} \dicExampleIS{Hún haggaðist ekki.} \dicExampleCS{Nehnulo to s~ní.}
\dicEntry[hagi] \dicTerm{hag|i} \dicIPA{{h}{a}{i}{j}{\textlengthmark}{\textsci}} \dicPos{m}[1] \dicFlx{(‑a, ‑ar)}[8] \dicDirectTranslationCS{pastvina, louka}
\dicEntry[hagkerfi] \dicTerm{hag··kerfi} \dicIPA{{h}{a}{x}{c\smash{\textsuperscript{h}}}{\textepsilon}{r}{v}{\textsci}} \dicPos{n}[2] \dicFlx{(‑s, ‑)}[14] \dicDirectTranslationCS{ekonomický systém}
\dicEntry[hagkvæmni] \dicTerm{hag··kvæmn|i} \dicIPA{{h}{a}{x}{k\smash{\textsuperscript{h}}}{v}{a}{i}{m}{n}{\textsci}} \dicPos{f}[3] \dicFlx{(‑i)}[3] \textbf{1.} \dicFieldCat{ekon.} \dicSynonym*{fjárhagsleg hagkvæmni} \dicDirectTranslationCS{rentabilita, ziskovost}  \textbf{2.} \dicSynonym*{skilvirkni} \dicDirectTranslationCS{efektivita, efektivnost}  \textbf{3.} \dicSynonym{hentugleiki} \dicDirectTranslationCS{praktičnost, účelnost}
\dicEntry[hagkvæmur] \dicTerm{hag··kvæmur} \dicIPA{{h}{a}{x}{k\smash{\textsuperscript{h}}}{v}{a}{i}{m}{\textscy}{\textsubring{r}}} \dicPos{adj}[1]\dicFlx{}[-1] \dicDirectTranslationCS{efektivní, výnosný, rentabilní}
\dicEntry[hagl] \dicTerm{hagl} \dicIPA{{h}{a}{\r{g}}{\textsubring{l}}} \dicPos{n}[2] \dicFlx{(‑s, högl)}[8] \textbf{1.} \dicDirectTranslationCS{(ledová) kroupa}  \textbf{2.} \dicSynonym*{lítil blýkúla} \dicDirectTranslationCS{brok}
\dicEntry[haglabyssa] \dicTerm{hagla··byss|a} \dicIPA{{h}{a}{\r{g}}{l}{a}{\textsubring{b}}{\textsci}{s}{a}} \dicPos{f}[1] \dicFlx{(‑u, ‑ur)}[19] \dicDirectTranslationCS{brokovnice}
\dicEntry[haglegur] \dicTerm{hag··legur} \dicIPA{{h}{a}{\textbabygamma}{l}{\textepsilon}{\textbabygamma}{\textscy}{\textsubring{r}}} \dicPos{adj}[1]\dicFlx{}[-8] \dicDirectTranslationCS{mistrný, mistrovský (výrobek ap.)} \dicExampleIS{haglegt smíði} \dicExampleCS{mistrovský výrobek}
\dicEntry[hagleikur] \dicTerm{hag··leik|ur} \dicIPA{{h}{a}{\textbabygamma}{l}{ei}{\r{g}}{\textscy}{\textsubring{r}}} \dicPos{m}[6] \dicFlx{(‑s)}[17] \dicDirectTranslationCS{zručnost, dovednost}
\dicEntry[haglél] \dicTerm{hagl··él} \dicIPA{{h}{a}{\r{g}}{l}{j}{\textepsilon}{\textsubring{l}}} \dicPos{n}[2] \dicFlx{(‑s, ‑)}[13] \dicDirectTranslationCS{krupobití, kroupy}
\dicEntry[hagmælska] \dicTerm{hag··mælsk|a} \dicIPA{{h}{a}{\textbabygamma}{m}{a}{i}{l}{s}{\r{g}}{a}} \dicPos{f}[1] \dicFlx{(‑u)}[5] \dicDirectTranslationCS{básnické nadání}
\dicEntry[hagmæltur] \dicTerm{hag··mæltur} \dicIPA{{h}{a}{\textbabygamma}{m}{a}{i}{\textsubring{l}}{\textsubring{d}}{\textscy}{\textsubring{r}}} \dicPos{adj}[1]\dicFlx{}[-10] \dicDirectTranslationCS{básnicky nadaný}
\dicEntry[hagnaður] \dicTerm{hag··nað|ur} \dicsymFrequent\  \dicIPA{{h}{a}{\r{g}}{n}{a}{ð}{\textscy}{\textsubring{r}}} \dicPos{m}[10] \dicFlx{(‑ar)}[9] \dicFieldCat{ekon.} \dicDirectTranslationCS{zisk, profit} \dicExampleIS{hreinn hagnaður} \dicExampleCS{čistý zisk}
\dicEntry[hagnast] \dicTerm{hagn|ast} \dicIPA{{h}{a}{\r{g}}{n}{a}{s}{\textsubring{d}}} \dicPos{v}[1] \dicFlx{(‑aðist)}[101] \dicFlx{refl} \dicDirectTranslationCS{profitovat, mít zisk} \dicExampleIS{hagnast á e‑u} \dicExampleCS{profitovat z~(čeho)}
\dicEntry[hagnýta] \dicTerm{hag··nýt|a} \dicIPA{{h}{a}{\textbabygamma}{n}{i}{\textsubring{d}}{a}} \dicPos{v}[2] \dicFlx{(‑ti, ‑t)}[54] \dicFlx{acc} \dicDirectTranslationCS{upotřebit, využít, využívat, zužitkovat, zužitkovávat} \dicExampleIS{hagnýta auðlindir} \dicExampleCS{využívat přírodní bohatství}
\dicEntry[hagnýting] \dicTerm{hag··nýt·ing} \dicIPA{{h}{a}{\textbabygamma}{n}{i}{\textsubring{d}}{i}{\ng}{\r{g}}} \dicPos{f}[4] \dicFlx{(‑ar)}[7] \dicDirectTranslationCS{upotřebení, využití, zužitkování}
\dicEntry[hagnýtur] \dicTerm{hag··nýtur} \dicIPA{{h}{a}{\textbabygamma}{n}{i}{\textsubring{d}}{\textscy}{\textsubring{r}}} \dicPos{adj}[1]\dicFlx{}[-1] \textbf{1.} \dicSynonym{hentugur} \dicDirectTranslationCS{praktický, užitečný}  \textbf{2.} \dicDirectTranslationCS{užitý, aplikovaný} \dicExampleIS{hagnýt stærðfræði} \dicExampleCS{užitá matematika}
\dicEntry[hagræða] \dicTerm{hag··ræ|ða} \dicsymFrequent\  \dicIPA{{h}{a}{\textbabygamma}{r}{a}{i}{ð}{a}} \dicPos{v}[2] \dicFlx{(‑ddi, ‑tt)}[167] \dicFlx{dat} \textbf{1.} \dicPhraseIS{hagræða sér} \dicDirectTranslationCS{udělat si pohodlí, uvelebit se}  \textbf{2.} \dicSynonym{lagfæra} \dicDirectTranslationCS{upravit, upravovat, urovnat, urovnávat, srovnat, srovnávat} \dicExampleIS{hagræða blómunum} \dicExampleCS{upravit květiny}  \textbf{3.} \dicFieldCat{ekon.} \dicDirectTranslationCS{zefektivnit, (z)racionalizovat}  \textbf{4.} \dicDirectTranslationCS{zkreslit, zkreslovat, (z)deformovat (skutečnost ap.)} \dicExampleIS{hagræða sannleikanum} \dicExampleCS{zkreslovat pravdu}
\dicEntry[hagræði] \dicTerm{hag··ræði} \dicIPA{{h}{a}{\textbabygamma}{r}{a}{i}{ð}{\textsci}} \dicPos{n}[2] \dicFlx{(‑s)}[20] \dicSynonym{gagn\smash{\textsuperscript{1}}} \dicDirectTranslationCS{užitek, přínos, prospěch};  \dicPhraseIS{það er hagræði að e‑u} \dicDirectTranslationCS{v~(čem) je přínos}
\dicEntry[hagræðing] \dicTerm{hag··ræð·ing} \dicIPA{{h}{a}{\textbabygamma}{r}{a}{i}{ð}{i}{\ng}{\r{g}}} \dicPos{f}[4] \dicFlx{(‑ar)}[7] \dicDirectTranslationCS{racionalizace, zefektivnění}
\dicEntry[hagskýrsla] \dicTerm{hag··skýrsl|a} \dicIPA{{h}{a}{x}{s}{\r{\textObardotlessj}}{i}{\textsubring{r}}{s}{\textsubring{d}}{l}{a}} \dicPos{f}[1] \dicFlx{(‑u, ‑ur)}[13] \dicDirectTranslationCS{statistická zpráva}
\dicEntry[hagsmunaaðili] \dicTerm{hags·muna··aðil|i} \dicIPA{{h}{a}{x}{s}{m}{\textscy}{n}{a}{a}{ð}{\textsci}{l}{\textsci}} \dicPos{m}[1] \dicFlx{(‑a\,/\addthin ‑ja, ‑ar\,/\addthin ‑jar)}[15] \dicDirectTranslationCS{zainteresovaná strana}
\dicEntry[hagsmunaárekstur] \dicTerm{hags·muna··á·rekst|ur} \dicIPA{{h}\-{a}\-{x}\-{s}\-{m}\-{\textscy}\-{n}\-{a}\-{au}\-{r}\-{\textepsilon}\-{x}\-{s}\-{\textsubring{d}}\-{\textscy}\-{\textsubring{r}}\-} \dicPos{m}[5] \dicFlx{(‑rar\,/\addthin ‑urs, ‑rar)}[9] \dicFieldCat{práv.} \dicDirectTranslationCS{střet zájmů}
\dicEntry[hagsmunasamtök] \dicTerm{hags·muna··sam·tök} \dicIPA{{h}{a}{x}{s}{m}{\textscy}{n}{a}{s}{a}{m}{t\smash{\textsuperscript{h}}}{\oe}{\r{g}}} \dicPos{n}[2] \dicFlx{pl}[9] \dicFieldCat{ekon.} \dicDirectTranslationCS{zájmová skupina}
\dicEntry[hagsmunir] \dicTerm{hags··munir} \dicIPA{{h}{a}{x}{s}{m}{\textscy}{n}{\textsci}{\textsubring{r}}} \dicPos{m}[10] \dicFlx{pl}[1] \dicDirectTranslationCS{zájmy};  \dicPhraseIS{gæta hagsmuna e‑rs} \dicDirectTranslationCS{střežit\,/\addthin chránit (čí) zájmy};  \dicPhraseIS{e‑að þjónar hagsmunum e‑rs} \dicDirectTranslationCS{(co) slouží zájmům (koho)}
\dicEntry[hagstjórn] \dicTerm{hag··stjórn} \dicIPA{{h}{a}{x}{s}{\textsubring{d}}{j}{ou}{r}{\textsubring{d}}{\textsubring{n}}} \dicPos{f}[7] \dicFlx{(‑ar)}[3] \dicDirectTranslationCS{řízení ekonomiky}
\dicEntry[hagstofa] \dicTerm{hag··stof|a} \dicIPA{{h}{a}{x}{s}{\textsubring{d}}{\textopeno}{v}{a}} \dicPos{f}[1] \dicFlx{(‑u, ‑ur)}[7] \dicDirectTranslationCS{statistický úřad}
\dicEntry[hagstæður] \dicTerm{hag··stæður} \dicIPA{{h}{a}{x}{s}{\textsubring{d}}{a}{i}{ð}{\textscy}{\textsubring{r}}} \dicPos{adj}[2]\dicFlx{}[-6] \dicSynonym{heppilegur} \dicDirectTranslationCS{příhod\-ný, výhodný}
\dicEntry[hagsveifla] \dicTerm{hag··sveifl|a} \dicIPA{{h}{a}{x}{s}{v}{ei}{\textsubring{b}}{l}{a}} \dicPos{f}[1] \dicFlx{(‑u, ‑ur)}[19] \dicDirectTranslationCS{ekonomický cyklus}
\dicEntry[hagsýni] \dicTerm{hag··sýn|i} \dicIPA{{h}{a}{x}{s}{i}{n}{\textsci}} \dicPos{f}[3] \dicFlx{(‑i)}[3] \dicSynonym{sparsemi} \dicDirectTranslationCS{šetrnost, hospodárnost}
\dicEntry[hagsýnn] \dicTerm{hag··sýnn} \dicIPA{{h}{a}{x}{s}{i}{\textsubring{d}}{\textsubring{n}}} \dicPos{adj}[7]\dicFlx{}[-1] \dicSynonym{sparsamur} \dicDirectTranslationCS{šetrný, hospodárný}
\dicEntry[hagsæld] \dicTerm{hag··sæld} \dicIPA{{h}{a}{x}{s}{a}{i}{l}{\textsubring{d}}} \dicPos{f}[7] \dicFlx{(‑ar)}[3] \dicSynonym{velmegun} \dicDirectTranslationCS{blahobyt, prosperita}
\dicEntry[hagur] \dicTerm{hag|ur\smash{\textsuperscript{1}}} \dicsymFrequent\  \dicIPA{{h}{a}{\textlengthmark}{\textbabygamma}{\textscy}{\textsubring{r}}} \dicPos{m}[9] \dicFlx{(‑s, ‑ir)}[9] \textbf{1.} \dicSynonym{kjör} \dicDirectTranslationCS{(životní) podmínky, situace}  \textbf{2.} \dicDirectTranslationCS{přínos, užitek, prospěch} \dicExampleIS{Þetta er honum í hag.} \dicExampleCS{Přijde mu to vhod.};  \dicPhraseIS{e‑m gengur allt í haginn} \dicFlx{impers} \dicDirectTranslationCS{(komu) vše vychází, (komu) se vše daří} \dicExampleIS{Gangi þér allt í haginn!} \dicExampleCS{Ať se ti vše daří!};  \dicPhraseIS{sjá sér hag í e‑u} \dicDirectTranslationCS{vidět v~(čem) prospěch};  \dicPhraseIS{búa í haginn fyrir framtíðina} \dicDirectTranslationCS{investovat do budoucnosti}
\dicEntry[hagur] \dicTerm{hagur\smash{\textsuperscript{2}}} \dicIPA{{h}{a}{\textlengthmark}{\textbabygamma}{\textscy}{\textsubring{r}}} \dicPos{adj}[1] \dicFlx{(f hög)}[2] \dicDirectTranslationCS{šikovný, zručný}
\dicEntry[hagvöxtur] \dicTerm{hag··|vöxtur} \dicIPA{{h}{a}{\textbabygamma}{v}{\oe}{x}{s}{\textsubring{d}}{\textscy}{\textsubring{r}}} \dicPos{m}[11] \dicFlx{(‑vaxtar)}[6] \dicFieldCat{ekon.} \dicDirectTranslationCS{ekonomický růst}
\dicEntry[hagyrðingur] \dicTerm{hag··yrð·ing|ur} \dicIPA{{h}{a}{\textlengthmark}{\textbabygamma}{\textsci}{r}{ð}{i}{\ng}{\r{g}}{\textscy}{\textsubring{r}}} \dicPos{m}[6] \dicFlx{(‑s, ‑ar)}[8] \dicDirectTranslationCS{zručný básník, zručná básnířka}
\dicEntry[hagþróun] \dicTerm{hag··þró|un} \dicIPA{{h}{a}{x}{\texttheta}{r}{ou}{\textscy}{\textsubring{n}}} \dicPos{f}[7] \dicFlx{(‑unar)}[9] \dicFieldCat{ekon.} \dicDirectTranslationCS{ekonomický\,/\addthin hospodářský rozvoj}
\dicEntry[Haíti] \dicTerm{Haít|i}\dicTerm{, Haítibúi} \dicIPA{{h}\-{a}\-{i}\-{\textsubring{d}}\-{\textsci}\-} \dicPos{m}[1] \dicFlx{(‑a, ‑ar)}[1] \dicDirectTranslationCS{Haiťan(ka)}
\dicEntry[Haítibúi] \dicTerm{Haíti··bú|i} \dicIPA{{h}{a}{i}{\textsubring{d}}{\textsci}{\textsubring{b}}{u}{\textsci}} \dicPos{m}[1] \dicFlx{(‑a, ‑ar)}[1] \dicLink{Haíti}
\dicEntry[Haítí] \dicTerm{Haítí} \dicIPA{{h}{a}{i}{\textsubring{d}}{i}} \dicPos{n}[4] \dicFlx{indecl}[2] \dicFieldCat{geog.} \dicDirectTranslationCS{Haiti}
\dicEntry[haítískur] \dicTerm{haítískur} \dicIPA{{h}{a}{i}{\textsubring{d}}{i}{s}{\r{g}}{\textscy}{\textsubring{r}}} \dicPos{adj}[1]\dicFlx{}[-6] \dicDirectTranslationCS{haitský}
\dicEntry[hak] \dicTerm{hak} \dicIPA{{h}{a}{\textlengthmark}{\r{g}}} \dicPos{n}[2] \dicFlx{(‑s, hök)}[8] \textbf{1.} \dicDirectTranslationCS{zpětný háček, háček s~protihrotem}  \textbf{2.} \dicFieldCat{jaz.} \dicDirectTranslationCS{háček} \dicIndirectTranslationCS{(diakritické znaménko)}
\dicEntry[haka] \dicTerm{haka} \dicsymFrequent\  \dicIPA{{h}{a}{\textlengthmark}{\r{g}}{a}} \dicPos{f}[1] \dicFlx{(höku, hökur)}[8] \dicDirectTranslationCS{brada} \dicExampleIS{tvöföld haka} \dicExampleCS{dvojitá brada}
\dicEntry[hakakross] \dicTerm{haka··kross} \dicIPA{{h}{a}{\textlengthmark}{\r{g}}{a}{k\smash{\textsuperscript{h}}}{r}{\textopeno}{s}} \dicPos{m}[4] \dicFlx{(‑, ‑ar)}[28] \dicDirectTranslationCS{hákový kříž, svastika}
\dicEntry[haki] \dicTerm{hak|i} \dicIPA{{h}{a}{\textlengthmark}{\r{\textObardotlessj}}{\textsci}} \dicPos{m}[1] \dicFlx{(‑a, ‑ar)}[8] \textbf{1.} \dicDirectTranslationCS{krumpáč}  \textbf{2.} \dicSynonym{krókur} \dicDirectTranslationCS{skoba};  \dicPhraseIS{e‑að situr á hakanum} \dicLangCat{přen.} \dicDirectTranslationCS{(co) je zanedbávané\,/\addthin opomíjené (koníček ap.), (co) je pověšené na hřebík}
\dicEntry[hakk] \dicTerm{hakk} \dicIPA{{h}{a}{h}{\r{g}}} \dicPos{n}[2] \dicFlx{(‑s)}[2] \dicDirectTranslationCS{mleté (maso)}
\dicEntry[hakka] \dicTerm{hakk|a} \dicIPA{{h}{a}{h}{\r{g}}{a}} \dicPos{v}[1] \dicFlx{(‑aði)}[13] \dicFlx{acc} \textbf{1.} \dicDirectTranslationCS{mlít, umlít (maso ap.)} \dicExampleIS{hakka kjöt} \dicExampleCS{umlít maso}  \textbf{2.} \dicPhraseIS{hakka e‑ð í sig} \dicDirectTranslationCS{zhltnout (co)}
\dicEntry[hakkavél] \dicTerm{hakka··vél} \dicIPA{{h}{a}{h}{\r{g}}{a}{v}{j}{\textepsilon}{\textsubring{l}}} \dicPos{f}[4] \dicFlx{(‑ar, ‑ar)}[1] \dicDirectTranslationCS{mlýnek na maso}
\dicEntry[hala] \dicTerm{hal|a} \dicIPA{{h}{a}{\textlengthmark}{l}{a}} \dicPos{v}[1] \dicFlx{(‑aði)}[13] \dicFlx{acc} \textbf{1.} \dicSynonym{draga} \dicDirectTranslationCS{táhnout, vléct}  \textbf{2.} \dicPhraseIS{hala inn e‑u} \dicDirectTranslationCS{sklidit (co), získat (co) (hlasy ve volbách ap.)};  \dicIdiom{hala}[niður]{ \dicPhraseIS{hala niður e‑ð}} \dicFieldCat{poč.} \dicDirectTranslationCS{stáhnout\,/\addthin stahovat (co)}
\dicEntry[halarófa] \dicTerm{hala··róf|a} \dicIPA{{h}{a}{\textlengthmark}{l}{a}{r}{ou}{a}} \dicPos{f}[1] \dicFlx{(‑u, ‑ur)}[19] \dicDirectTranslationCS{šik, zástup};  \dicPhraseIS{ganga í halarófu} \dicDirectTranslationCS{jít v~řadě za sebou}
\dicEntry[halastjarna] \dicTerm{hala··|stjarna} \dicIPA{{h}{a}{\textlengthmark}{l}{a}{s}{\textsubring{d}}{j}{a}{r}{\textsubring{d}}{n}{a}} \dicPos{f}[1] \dicFlx{(‑stjörnu, ‑stjörnur)}[8] \dicFieldCat{astro.} \dicDirectTranslationCS{kometa, vlasatice}
\dicEntry[hald] \dicTerm{hald} \dicIPA{{h}{a}{l}{\textsubring{d}}} \dicPos{n}[2] \dicFlx{(‑s, höld)}[8] \textbf{1.} \dicDirectTranslationCS{držení, uchopení, chycení} \dicExampleIS{hafa hald á bandinu} \dicExampleCS{držet provaz}  \textbf{2.} \dicSynonym{hanki} \dicDirectTranslationCS{držadlo, držátko}  \textbf{3.} \dicSynonym{varðhald} \dicDirectTranslationCS{zadržení, vazba};  \dicPhraseIS{hafa e‑n í haldi} \dicDirectTranslationCS{zadržet (koho)};  \dicPhraseIS{leysa e‑n úr haldi} \dicDirectTranslationCS{propustit (koho) z~vazby}  \textbf{4.} \dicSynonym{gagn\smash{\textsuperscript{1}}} \dicDirectTranslationCS{užitek, prospěch};  \dicPhraseIS{koma að haldi} \dicDirectTranslationCS{být k~užitku, přijít vhod}  \textbf{5.} \dicSynonym{stoð} \dicDirectTranslationCS{podpora, opora};  \dicPhraseIS{eiga hald og traust í e‑m} \dicDirectTranslationCS{mít v~(kom) oporu}
\dicEntry[halda] \dicTerm{halda\smash{\textsuperscript{1}}} \dicIPA{{h}{a}{l}{\textsubring{d}}{a}} \dicPos{f}[1] \dicFlx{(höldu, höldur)}[20] \dicSynonym*{handarhald} \dicDirectTranslationCS{držátko, rukojeť}
\dicEntry[halda] \dicTerm{halda\smash{\textsuperscript{2}}} \dicsymFrequent\  \dicIPA{{h}{a}{l}{\textsubring{d}}{a}} \dicPos{v}[6] \dicFlx{(held, hélt, héldum, héldi, haldið)}[58] \dicFlx{acc\,/\addthin dat} \textbf{1.} \dicSynonym{telja} \dicDirectTranslationCS{domnívat se, myslet si} \dicExampleIS{Ég held að þessi maður sé einmana.} \dicExampleCS{Myslím si, že ten člověk je osamělý.}  \textbf{2.} \dicFlx{dat} \dicDirectTranslationCS{držet, mít v~rukou} \dicExampleIS{halda barninu} \dicExampleCS{držet dítě}  \textbf{3.} \dicDirectTranslationCS{držet (led ap.)} \dicExampleIS{Ísinn heldur ekki.} \dicExampleCS{Led nedrží.}  \textbf{4.} \dicFlx{dat} \dicDirectTranslationCS{pořádat, slavit, oslavovat} \dicExampleIS{halda hátíð} \dicExampleCS{slavit svátky}  \textbf{5.} \dicFlx{acc} \dicDirectTranslationCS{(do)držet (slib ap.)} \dicExampleIS{halda orð sín} \dicExampleCS{držet své slovo}  \textbf{6.} \dicDirectTranslationCS{jet, jít, cestovat} \dicExampleIS{halda burt} \dicExampleCS{odjet}  \textbf{7.} \dicPhraseIS{halda sig e‑s staðar} \dicDirectTranslationCS{držet se (kde), zdržovat se (kde)}  \textbf{8.} \dicPhraseIS{halda sér vel} \dicLangCat{přen.} \dicDirectTranslationCS{držet se dobře};  \dicIdiom{halda}[að]{ \dicPhraseIS{halda e‑m að verki}} \dicDirectTranslationCS{koncentrovat (koho) na práci} \dicExampleIS{halda sér að verki þrátt fyrir truflun} \dicExampleCS{soustředit se na práci i~přes vyrušování};  \dicIdiom{halda}[aftur af]{ \dicPhraseIS{halda aftur af sér}} \dicDirectTranslationCS{držet se zkrátka};  \dicIdiom{halda}[á]{ \dicPhraseIS{halda á e‑u}} \dicDirectTranslationCS{držet (co), nést (co)} \dicExampleIS{halda á bókinni í hendinni} \dicExampleCS{nést v~ruce knížku}; { \dicPhraseIS{þurfa á e‑u að halda}} \dicDirectTranslationCS{potřebovat (co)}; { \dicPhraseIS{halda vel á e‑u}} \dicDirectTranslationCS{pečovat o~(co) dobře (o~peníze ap.)}; { \dicPhraseIS{þegar\,/\addthin ef á þarf að halda}} \dicFlx{adv} \dicDirectTranslationCS{v~případě potřeby};  \dicIdiom{halda}[áfram]{ \dicPhraseIS{halda áfram}} \dicDirectTranslationCS{pokračovat, postupovat}; { \dicPhraseIS{halda e‑u áfram}} \dicDirectTranslationCS{pokračovat s~(čím), pokračovat v~(čem)}; { \dicPhraseIS{halda áfram með e‑ð}} \dicDirectTranslationCS{pokračovat s~(čím), pokračovat v~(čem)} \dicExampleIS{halda áfram með verkið} \dicExampleCS{pokračovat v~práci}; { \dicPhraseIS{e‑að heldur áfram}} \dicDirectTranslationCS{(co) pokračuje};  \dicIdiom{halda}[eftir]{ \dicPhraseIS{halda e‑u eftir}} \dicDirectTranslationCS{podržet (si) (co), uchovat (si) (co)} \dicExampleIS{halda eftir hluta af bókunum} \dicExampleCS{uchovat část knih};  \dicIdiom{halda}[fram]{ \dicPhraseIS{halda e‑u fram}} \dicDirectTranslationCS{tvrdit (co), prohlásit (co)}; { \dicPhraseIS{ef fram heldur sem horfir}} \dicDirectTranslationCS{pokud to bude pokračovat tímto způsobem};  \dicIdiom{halda}[framhjá]{ \dicPhraseIS{halda framhjá e‑m}} \dicDirectTranslationCS{být (komu) nevěrný, zahýbat (komu)};  \dicIdiom{halda}[frá]{ \dicPhraseIS{halda sig frá e‑u}} \dicDirectTranslationCS{vyhnout\,/\addthin vyhýbat se (čemu), zdržet se (čeho) (alkoholu ap.)};  \dicIdiom{halda}[fyrir]{ \dicPhraseIS{halda fyrir e‑ð}} \dicDirectTranslationCS{zakrýt si (co) (pusu ap.), zacpat si (co) (uši ap.)} \dicExampleIS{halda fyrir eyrun} \dicExampleCS{zacpat si uši}; { \dicPhraseIS{halda e‑u fyrir e‑m}} \dicDirectTranslationCS{skrývat (co) před (kým), tajit (co) před (kým)};  \dicIdiom{halda}[í]{ \dicPhraseIS{halda í e‑ð}} \dicDirectTranslationCS{držet se (čeho) (zábradlí ap.)}; { \dicPhraseIS{halda í vonina}} \dicLangCat{přen.} \dicDirectTranslationCS{uchovat si naději};  \dicIdiom{halda}[með]{ \dicPhraseIS{halda með e‑m}} \dicDirectTranslationCS{držet s~(kým), být na (čí) straně};  \dicIdiom{halda}[niðri]{ \dicPhraseIS{halda niðri í sér e‑u}} \dicDirectTranslationCS{zadržovat (co), potlačovat (co) (smích ap.)};  \dicIdiom{halda}[saman]{ \dicPhraseIS{halda saman e‑u}} \dicDirectTranslationCS{držet (co) pospolu, zachovávat (co) (staré účtenky ap.)};  \dicIdiom{halda}[til]{ \dicPhraseIS{halda til e‑s staðar}} \dicDirectTranslationCS{zdržovat se (kde), žít (kde)}; { \dicPhraseIS{halda sér til fyrir e‑m}} \dicDirectTranslationCS{vystrojit se kvůli (komu)};  \dicIdiom{halda}[upp á]{ \dicPhraseIS{halda upp á e‑ð}} \dicDirectTranslationCS{oslavovat (co), slavit (co)} \dicExampleIS{Heldur þú upp á afmælið þitt?} \dicExampleCS{Slavíš svoje narozeniny?};  \dicIdiom{halda}[uppi]{ \dicPhraseIS{halda uppi samræðum}} \dicDirectTranslationCS{konverzovat, udržovat rozhovor}; { \dicPhraseIS{halda uppi e‑u}} \dicDirectTranslationCS{udržovat (co), podporovat (co)}; { \dicPhraseIS{halda e‑m uppi}} \dicDirectTranslationCS{podpořit\,/\addthin podporovat (koho)};  \dicIdiom{halda}[utan um]{ \dicPhraseIS{halda utan um e‑n}} \dicDirectTranslationCS{držet (koho), objímat (koho)}; { \dicPhraseIS{halda utan um e‑ð}} \dicDirectTranslationCS{podržet (co), šetřit (co) (peníze ap.)};  \dicIdiom{halda}[út]{ \dicPhraseIS{halda e‑ð út}} \dicDirectTranslationCS{vydržet (co)};  \dicIdiom{halda}[úti]{ \dicPhraseIS{halda úti e‑u}} \dicDirectTranslationCS{publikovat (co)};  \dicIdiom{halda}[við]{ \dicPhraseIS{halda e‑u við}} \dicDirectTranslationCS{udržovat (co) (dům ap.)}; { \dicPhraseIS{halda við e‑n}} \dicDirectTranslationCS{mít aféru s~(kým)}; { \dicPhraseIS{halda við e‑ð}} \dicDirectTranslationCS{držet (co)}; { \dicPhraseIS{e‑m heldur við e‑u}} \dicFlx{impers} \dicDirectTranslationCS{(kdo) se začíná (co)} \dicExampleIS{Mér heldur við köfnun.} \dicExampleCS{Začínám se dusit.};  \dicIdiom{haldast}{ \dicPhraseIS{haldast}} \dicFlx{refl} \dicDirectTranslationCS{(se)trvat, udržet se} \dicExampleIS{Þetta hefur haldist fram á þennan dag.} \dicExampleCS{Trvá to až dodnes.};  \dicIdiom{haldast}[í]{ \dicPhraseIS{haldast í hendur}} \dicFlx{refl} {\textbf{a.}} \dicDirectTranslationCS{držet se za ruce} \dicExampleIS{Þau haldast í hendur.} \dicExampleCS{Drží se za ruce.};  {\textbf{b.}} \dicLangCat{přen.} \dicDirectTranslationCS{jít ruku v~ruce};  \dicIdiom{haldast}[við]{ \dicPhraseIS{haldast (ekki) við}} \dicFlx{refl} \dicDirectTranslationCS{(ne)vydržet (na nějakém místě ap.)};  \dicIdiom{halda}{ \dicPhraseIS{haltu kjafti!}} \dicFlx{imper} \dicDirectTranslationCS{drž hubu!}
\dicEntry[haldgóður] \dicTerm{hald··|góður} \dicIPA{{h}{a}{l}{\textsubring{d}}{\r{g}}{ou}{ð}{\textscy}{\textsubring{r}}} \dicPos{adj}[11] \dicFlx{(comp ‑betri, sup ‑bestur)}[1] \dicDirectTranslationCS{solidní, spolehlivý} \dicExampleIS{haldgóðar upplýsingar} \dicExampleCS{spolehlivé informace}
\dicEntry[haldið] \dicTerm{haldið} \dicIPA{{h}{a}{l}{\textsubring{d}}{\textsci}{\texttheta}} \dicPos{v} \dicFlx{supin} \dicLink{halda\smash{\textsuperscript{2}}}
\dicEntry[haldinn] \dicTerm{haldinn} \dicIPA{{h}{a}{l}{\textsubring{d}}{\textsci}{\textsubring{n}}} \dicPos{adj}[6]\dicFlx{}[-4] \dicPhraseIS{vera haldinn e‑u} \dicDirectTranslationCS{trpět (čím), být sužován (čím)} \dicExampleIS{vera haldinn sjúkdómi} \dicExampleCS{trpět nemocí};  \dicPhraseIS{vera haldinn illum anda} \dicDirectTranslationCS{být posedlý zlým duchem};  \dicPhraseIS{vera vel haldinn} \dicDirectTranslationCS{být v~dobrém stavu};  \dicPhraseIS{vera illa haldinn} \dicDirectTranslationCS{být ve špatném stavu}
\dicEntry[hali] \dicTerm{hal|i} \dicIPA{{h}{a}{\textlengthmark}{l}{\textsci}} \dicPos{m}[1] \dicFlx{(‑a, ‑ar)}[8] \dicSynonym{rófa} \dicDirectTranslationCS{ocas, oháňka, chvost};  \dicPhraseIS{leika lausum hala} \dicLangCat{přen.} \dicDirectTranslationCS{mít volnou ruku (člověk ap.), být na volno (pes ap.)}
\dicEntry[halla] \dicTerm{hall|a} \dicsymFrequent\  \dicIPA{{h}{a}{\textsubring{d}}{l}{a}} \dicPos{v}[1] \dicFlx{(‑aði)}[24] \dicFlx{dat} \textbf{1.} \dicDirectTranslationCS{naklonit, naklánět, nahnout, nahýbat, nachýlit, nachylovat} \dicExampleIS{halla glasinu} \dicExampleCS{naklonit skleničku}  \textbf{2.} \dicPhraseIS{halla sér} \dicDirectTranslationCS{zdřímnout si, schrupnout si} \dicExampleIS{halla sér á rúmið} \dicExampleCS{zdřímnout si na posteli}  \textbf{3.} \dicPhraseIS{e‑u hallar} \dicFlx{impers} \dicDirectTranslationCS{(co) se svažuje} \dicExampleIS{Landinu hallar niður að vatninu.} \dicExampleCS{Země se svažuje k~jezeru.}  \textbf{4.} \dicPhraseIS{e‑u hallar} \dicFlx{impers} \dicDirectTranslationCS{(co) se chýlí\,/\addthin schyluje ke konci} \dicExampleIS{Sumri hallar.} \dicExampleCS{Léto se chýlí ke konci.}  \textbf{5.} \dicPhraseIS{halla réttu máli} \dicDirectTranslationCS{lhát, neříkat pravdu};  \dicIdiom{halla}[að]{ \dicPhraseIS{halla sér að e‑m}} {\textbf{a.}} \dicDirectTranslationCS{opřít se o~(koho)} \dicIndirectTranslationCS{(podepřít se)};  {\textbf{b.}} \dicLangCat{přen.} \dicDirectTranslationCS{opřít se o~(koho)} \dicIndirectTranslationCS{(vyhledat pomoc)};  \dicIdiom{halla}[á]{ \dicPhraseIS{halla á e‑n}} \dicDirectTranslationCS{podvést (koho), ušít to na (koho)};  \dicIdiom{halla}[undan]{ \dicPhraseIS{það hallar undan fæti hjá e‑m}} \dicFlx{impers} \dicLangCat{přen.} \dicDirectTranslationCS{jde to z~kopce s~(kým)};  \dicIdiom{hallast}{ \dicPhraseIS{hallast}} \dicFlx{refl} \dicDirectTranslationCS{(na)klonit se, naklánět se, (na)chýlit se};  \dicIdiom{hallast}[að]{ \dicPhraseIS{hallast að e‑u}} \dicFlx{refl} \dicDirectTranslationCS{klonit\,/\addthin přiklánět se k~(čemu) (k~myšlence ap.)}
\dicEntry[hallalaus] \dicTerm{halla··laus} \dicIPA{{h}{a}{\textsubring{d}}{l}{a}{l}{\oe i}{s}} \dicPos{adj}[5]\dicFlx{}[-1] \textbf{1.} \dicSynonym*{taplaus} \dicDirectTranslationCS{vyrovnaný, vyvážený} \dicExampleIS{hallalaus fjárlög} \dicExampleCS{vyrovnaný rozpočet}  \textbf{2.} \dicSynonym{flatur} \dicDirectTranslationCS{rovný, plochý (povrch ap.)}
\dicEntry[hallamælir] \dicTerm{halla··mæl|ir} \dicIPA{{h}{a}{\textsubring{d}}{l}{a}{m}{a}{i}{l}{\textsci}{\textsubring{r}}} \dicPos{m}[7] \dicFlx{(‑is, ‑ar)}[1] \dicFieldCat{techn.} \dicDirectTranslationCS{libela, vodováha}
\dicEntry[hallar] \dicTerm{hallar} \dicIPA{{h}{a}{\textsubring{d}}{l}{a}{\textsubring{r}}} \dicPos{f} \dicFlx{sg gen} \dicLink{höll\smash{\textsuperscript{1}}}
\dicEntry[hallarekstur] \dicTerm{halla··rekst|ur} \dicIPA{{h}{a}{\textsubring{d}}{l}{a}{r}{\textepsilon}{x}{s}{\textsubring{d}}{\textscy}{\textsubring{r}}} \dicPos{m}[5] \dicFlx{(‑urs\,/\addthin ‑rar)}[10] \dicFieldCat{ekon.} \dicDirectTranslationCS{činnost vykazující deficit, prodělečná činnost} \dicExampleIS{Það er hallarekstur á fyrirtækinu.} \dicExampleCS{Firma vykazuje deficit.}
\dicEntry[halli] \dicTerm{hall|i} \dicsymFrequent\  \dicIPA{{h}{a}{\textsubring{d}}{l}{\textsci}} \dicPos{m}[1] \dicFlx{(‑a, ‑ar)}[8] \textbf{1.} \dicSynonym*{skái} \dicDirectTranslationCS{sklon, náklon, naklonění}  \textbf{2.} \dicSynonym{brekka} \dicDirectTranslationCS{stráň, svah}  \textbf{3.} \dicFieldCat{ekon.} \dicSynonym{tap} \dicDirectTranslationCS{ztráta, deficit} \dicExampleIS{halli á ríkissjóði} \dicExampleCS{deficit státní pokladny}
\dicEntry[hallir] \dicTerm{hallir} \dicIPA{{h}{a}{\textsubring{d}}{l}{\textsci}{\textsubring{r}}} \dicPos{f} \dicFlx{pl nom} \dicLink{höll\smash{\textsuperscript{1}}}
\dicEntry[halloka] \dicTerm{hall··oka} \dicIPA{{h}{a}{\textsubring{d}}{l}{\textopeno}{\r{g}}{a}} \dicPos{adj}[13] \dicFlx{indecl}[1] \dicPhraseIS{fara halloka (fyrir e‑m)} \dicDirectTranslationCS{prohrát (s~(kým)), utrpět porážku (s~(kým))}
\dicEntry[halló] \dicTerm{halló} \dicIPA{{h}{a}{l}{}{ou}} \dicPos{inter} \textbf{1.} \dicDirectTranslationCS{haló} \dicIndirectTranslationCS{(pozdrav na začátku telefonování)}  \textbf{2.} \dicLangCat{hovor.} \dicDirectTranslationCS{ahoj, čau}  \textbf{3.} \dicIndirectTranslationCS{zavolání při varování někoho ap.}
\dicEntry[hallur] \dicTerm{hallur} \dicIPA{{h}{a}{\textsubring{d}}{l}{\textscy}{\textsubring{r}}} \dicPos{adj}[1] \dicFlx{(f höll)}[2] \dicDirectTranslationCS{nakloněný, nachýlený, nahnutý};  \dicPhraseIS{standa höllum fæti} \dicLangCat{přen.} \dicDirectTranslationCS{stát na šikmé ploše};  \dicPhraseIS{vera hallur undir e‑n} \dicDirectTranslationCS{být nakloněn (komu)}
\dicEntry[hallæri] \dicTerm{hall··æri} \dicIPA{{h}{a}{\textsubring{d}}{l}{a}{i}{r}{\textsci}} \dicPos{n}[2] \dicFlx{(‑s, ‑)}[14] \dicDirectTranslationCS{hubená léta, neúroda, hladomor}
\dicEntry[hallærislegur] \dicTerm{hall·æris··legur} \dicIPA{{h}{a}{\textsubring{d}}{l}{a}{i}{r}{\textsci}{s}{l}{\textepsilon}{\textbabygamma}{\textscy}{\textsubring{r}}} \dicPos{adj}[1]\dicFlx{}[-8] \dicSynonym{aumur} \dicDirectTranslationCS{ubohý, bídný, mizerný}
\dicEntry[haltra] \dicTerm{haltr|a} \dicIPA{{h}{a}{\textsubring{l}}{\textsubring{d}}{r}{a}} \dicPos{v}[1] \dicFlx{(‑aði)}[50] \dicDirectTranslationCS{kulhat, pokulhávat}
\dicEntry[haltur] \dicTerm{haltur} \dicIPA{{h}{a}{\textsubring{l}}{\textsubring{d}}{\textscy}{\textsubring{r}}} \dicPos{adj}[1] \dicFlx{(f hölt)}[11] \dicDirectTranslationCS{kulhavý, kulhající}
\dicEntry[hamagangur] \dicTerm{hama··gang|ur} \dicIPA{{h}{a}{\textlengthmark}{m}{a}{\r{g}}{au}{\ng}{\r{g}}{\textscy}{\textsubring{r}}} \dicPos{m}[6] \dicFlx{(‑s)}[7] \dicDirectTranslationCS{povyk, rozruch, pozdvižení}
\dicEntry[hamar] \dicTerm{ham|ar} \dicsymFrequent\  \dicIPA{{h}{a}{\textlengthmark}{m}{a}{\textsubring{r}}} \dicPos{m}[5] \dicFlx{(‑ars, ‑rar)}[8] \textbf{1.} \dicDirectTranslationCS{kladivo, kladívko} \dicExampleIS{berja á naglann með hamri} \dicExampleCS{zatloukat hřebík kladívkem}  \textbf{2.} \dicSynonym{klettur} \dicDirectTranslationCS{skalisko, útes, sráz} \dicExampleIS{klífa upp hamarinn} \dicExampleCS{lézt na útes}  \textbf{3.} \dicFieldCat{anat.} \dicDirectTranslationCS{kladívko}
\dicEntry[hamast] \dicTerm{ham|ast} \dicsymFrequent\  \dicIPA{{h}{a}{\textlengthmark}{m}{a}{s}{\textsubring{d}}} \dicPos{v}[1] \dicFlx{(‑aðist)}[93] \dicFlx{refl} \dicSynonym*{vinna af kappi} \dicDirectTranslationCS{makat, dřít} \dicExampleIS{vera að hamast í heyskap} \dicExampleCS{dřít na senoseči}
\dicEntry[Hamborg] \dicTerm{Ham··borg} \dicIPA{{h}{a}{m}{\textsubring{b}}{\textopeno}{r}{\r{g}}} \dicPos{f}[7] \dicFlx{(‑ar)}[4] \dicFieldCat{geog.} \dicDirectTranslationCS{Hamburk} \dicIndirectTranslationCS{(město v~Německu)}
\dicEntry[hamborgarhryggur] \dicTerm{ham·borgar··hrygg|ur} \dicIPA{{h}{a}{m}{\textsubring{b}}{\textopeno}{r}{\r{g}}{a}{\textsubring{r}}{\textsci}{\r{g}}{\textscy}{\textsubring{r}}} \dicPos{m}[10] \dicFlx{(‑jar, ‑ir)}[23] \dicFieldCat{kulin.} \dicDirectTranslationCS{pečeně z~lehce uzeného vepřového krku}
\dicEntry[hamborgari] \dicTerm{ham·borg··ar|i} \dicIPA{{h}{a}{m}{\textsubring{b}}{\textopeno}{r}{\r{g}}{a}{r}{\textsci}} \dicPos{m}[1] \dicFlx{(‑a, ‑ar)}[13] \dicFieldCat{kulin.} \dicDirectTranslationCS{hamburger}
\dicEntry[hamdi] \dicTerm{hamdi} \dicIPA{{h}{a}{m}{\textsubring{d}}{\textsci}} \dicPos{v} \dicFlx{ind pf sg 1 pers} \dicLink{hemja\smash{\textsuperscript{2}}}
\dicEntry[hamfarir] \dicTerm{ham··farir} \dicIPA{{h}{a}{m}{f}{a}{r}{\textsci}{\textsubring{r}}} \dicPos{f}[7] \dicFlx{pl}[18] \dicDirectTranslationCS{katastrofa (přírodní ap.), pohroma, kalamita}
\dicEntry[hamið] \dicTerm{hamið} \dicIPA{{h}{a}{\textlengthmark}{m}{\textsci}{\texttheta}} \dicPos{v} \dicFlx{supin} \dicLink{hemja\smash{\textsuperscript{2}}}
\dicEntry[hamingja] \dicTerm{hamingj|a} \dicsymFrequent\  \dicIPA{{h}{a}{\textlengthmark}{m}{i}{\textltailn}{\r{\textObardotlessj}}{a}} \dicPos{f}[1] \dicFlx{(‑u)}[5] \dicSynonym{blessun} \dicDirectTranslationCS{štěstí, blaženost};  \dicPhraseIS{hamingjunni sé lof} \dicDirectTranslationCS{bohudíky};  \dicPhraseIS{óska e‑m til hamingju} \dicDirectTranslationCS{popřát (komu) štěstí};  \dicPhraseIS{til allrar hamingju} \dicFlx{adv} \dicDirectTranslationCS{naštěstí, šťastnou náhodou}
\dicEntry[hamingjusamur] \dicTerm{hamingju··|samur} \dicIPA{{h}{a}{\textlengthmark}{m}{i}{\textltailn}{\r{\textObardotlessj}}{\textscy}{s}{a}{m}{\textscy}{\textsubring{r}}} \dicPos{adj}[1] \dicFlx{(f ‑söm)}[2] \dicDirectTranslationCS{šťastný, spokojený}
\dicEntry[hamingjusemi] \dicTerm{hamingju··sem|i} \dicIPA{{h}{a}{\textlengthmark}{m}{i}{\textltailn}{\r{\textObardotlessj}}{\textscy}{s}{\textepsilon}{m}{\textsci}} \dicPos{f}[2] \dicFlx{(‑i)}[2] \dicDirectTranslationCS{štěstí, blaženost}
\dicEntry[hamla] \dicTerm{hamla\smash{\textsuperscript{1}}} \dicIPA{{h}{a}{m}{l}{a}} \dicPos{f}[1] \dicFlx{(hömlu, hömlur)}[8] \dicPhraseIS{hömlur} \dicFlx{pl} \dicSynonym*{takmarkanir} \dicDirectTranslationCS{zábrany, omezení};  \dicPhraseIS{draga e‑u úr hömlu} \dicDirectTranslationCS{protahovat (co) (rozhodnutí ap.)};  \dicPhraseIS{leggja hömlur á e‑u} \dicDirectTranslationCS{zavést omezení na (co)}
\dicEntry[hamla] \dicTerm{haml|a\smash{\textsuperscript{2}}} \dicIPA{{h}{a}{m}{l}{a}} \dicPos{v}[1] \dicFlx{(‑aði)}[13] \dicFlx{dat} \dicDirectTranslationCS{brzdit, blokovat} \dicExampleIS{hamla rannsókninni} \dicExampleCS{brzdit vyšetřování};  \dicIdiom{hamla}[á móti]{ \dicPhraseIS{hamla á móti e‑u}} \dicDirectTranslationCS{brzdit (co), zdržovat (co)};  \dicIdiom{hamla}[gegn]{ \dicPhraseIS{hamla gegn e‑u}} \dicDirectTranslationCS{bránit (čemu), předcházet (čemu)}
\dicEntry[hamlandi] \dicTerm{haml··andi} \dicIPA{{h}{a}{m}{l}{a}{n}{\textsubring{d}}{\textsci}} \dicPos{adj}[13] \dicFlx{indecl}[1] \dicDirectTranslationCS{brzdící, potlačující, inhibiční} \dicExampleIS{hamlandi áhrif} \dicExampleCS{brzdící účinek}
\dicEntry[hampa] \dicTerm{hamp|a} \dicIPA{{h}{a}{\textsubring{m}}{\textsubring{b}}{a}} \dicPos{v}[1] \dicFlx{(‑aði)}[50] \dicFlx{dat} \textbf{1.} \dicSynonym{dilla} \dicDirectTranslationCS{kolébat, houpat} \dicExampleIS{hampa barni} \dicExampleCS{houpat dítě}  \textbf{2.} \dicDirectTranslationCS{vychvalovat, opěvovat} \dicExampleIS{Bókinni var mikið hampað.} \dicExampleCS{Knížka byla velmi opěvována.}
\dicEntry[hampur] \dicTerm{hamp|ur} \dicIPA{{h}{a}{\textsubring{m}}{\textsubring{b}}{\textscy}{\textsubring{r}}} \dicPos{m}[6] \dicFlx{(‑s)}[7] \textbf{1.} \dicFieldCat{bot.} \dicDirectTranslationCS{konopí} \textit{(l.~{\textLA{Cannabis}})}  \dicsymPhoto\   \textbf{2.} \dicDirectTranslationCS{konopná vlákna} \dicExampleIS{hampreipi} \dicExampleCS{konopný provaz}  \textbf{3.} \dicLangCat{slang.} \dicSynonym{kannabis} \dicDirectTranslationCS{marihuana, hašiš}
\dicFigure{4766.jpg}{Hampur}{Hampur - Zicha Ondřej, Biolib, Copyright/CC-BY-NC}
\dicEntry[hamra] \dicTerm{hamr|a} \dicIPA{{h}{a}{m}{r}{a}} \dicPos{v}[1] \dicFlx{(‑aði)}[13] \dicFlx{acc} \textbf{1.} \dicDirectTranslationCS{bušit, tlouct, tepat};  \dicPhraseIS{hamra á e‑ð} \dicDirectTranslationCS{bušit do (čeho)} \dicExampleIS{hamra á lyklaborðið} \dicExampleCS{bušit do klávesnice}  \textbf{2.} \dicPhraseIS{hamra á e‑u} \dicDirectTranslationCS{omílat (co), mluvit stále dokola o~(čem)};  \dicIdiom{hamra}[í gegn]{ \dicPhraseIS{hamra e‑ð í gegn}} \dicDirectTranslationCS{protlačit (co) (nápad ap.)}
\begin{xtolerant}{}{1pt}
\dicEntry[hamsatólg] \dicTerm{hamsa··tólg} \dicIPA{{h}{a}{m}{s}{a}{t\smash{\textsuperscript{h}}}{ou}{l}{\r{g}}}\addthinS\dicPos{f}[4]\addthinS\dicFlx{(‑ar)}[3] \dicFieldCat{kulin.} \dicDirectTranslationCS{škvarky (z~jehněčího sádla)}
\end{xtolerant}
\dicEntry[hamskipti] \dicTerm{ham··skipti} \dicIPA{{h}{a}{m}{s}{\r{\textObardotlessj}}{\textsci}{f}{\textsubring{d}}{\textsci}} \dicPos{n}[2] \dicFlx{pl}[19] \dicDirectTranslationCS{metamorfóza, proměna}
\dicEntry[hamslaus] \dicTerm{hams··laus} \dicIPA{{h}{a}{m}{s}{l}{\oe i}{s}} \dicPos{adj}[5]\dicFlx{}[-1] \dicDirectTranslationCS{neovladatelný, nekontrolovatelný} \dicExampleIS{vera hamslaus af reiði} \dicExampleCS{být zuřivostí bez sebe}
\dicEntry[hamsleysi] \dicTerm{hams··leysi} \dicIPA{{h}{a}{m}{s}{l}{ei}{s}{\textsci}} \dicPos{n}[2] \dicFlx{(‑s)}[20] \dicDirectTranslationCS{nekontrolovatelnost, neovladatelnost}
\dicEntry[hamstola] \dicTerm{ham··stola} \dicIPA{{h}{a}{m}{s}{\textsubring{d}}{\textopeno}{l}{a}} \dicPos{adj}[13] \dicFlx{indecl}[1] \dicDirectTranslationCS{rozzuřený, rozzlobený}
\dicEntry[hamstra] \dicTerm{hamstr|a} \dicIPA{{h}{a}{m}{s}{\textsubring{d}}{r}{a}} \dicPos{v}[1] \dicFlx{(‑aði)}[43] \dicFlx{acc} \dicDirectTranslationCS{hromadit (nadměrné) zásoby, křečkovat}
\dicEntry[hamstur] \dicTerm{hamst|ur\smash{\textsuperscript{1}}} \dicIPA{{h}{a}{m}{s}{\textsubring{d}}{\textscy}{\textsubring{r}}} \dicPos{m}[5] \dicFlx{(‑urs, ‑rar)}[3] \dicFieldCat{zool.} \dicDirectTranslationCS{křeček} \textit{(l.~{\textLA{Cricetus}})}
\dicEntry[hamstur] \dicTerm{hamstur\smash{\textsuperscript{2}}} \dicIPA{{h}{a}{m}{s}{\textsubring{d}}{\textscy}{\textsubring{r}}} \dicPos{n}[2] \dicFlx{(‑s)}[28] \dicDirectTranslationCS{skupování (zboží ap.), hromadění zásob, křečkování}
\dicEntry[hamur] \dicTerm{ham|ur} \dicIPA{{h}{a}{\textlengthmark}{m}{\textscy}{\textsubring{r}}} \dicPos{m}[9] \dicFlx{(‑s, ‑ir)}[9] \textbf{1.} \dicSynonym{skinn} \dicDirectTranslationCS{kůže (ptáka ap.)}  \textbf{2.} \dicFieldCat{poč.} \dicDirectTranslationCS{režim, mód};  \dicPhraseIS{skipta hömum} \dicDirectTranslationCS{změnit podobu}
\dicEntry[hana] \dicTerm{hana} \dicIPA{{h}{a}{\textlengthmark}{n}{a}} \dicPos{pron}[17] \dicFlx{pers} \dicFlx{sg acc} \dicLink{hún}
\dicEntry[hanagal] \dicTerm{hana··gal} \dicIPA{{h}{a}{\textlengthmark}{n}{a}{\r{g}}{a}{\textsubring{l}}} \dicPos{n}[2] \dicFlx{(‑s)}[2] \dicDirectTranslationCS{(za)kokrhání, kuropění};  \dicPhraseIS{við fyrsta hanagal} \dicFlx{adv} \dicDirectTranslationCS{s~prvním zakokrháním\,/\addthin kuropěním}
\dicEntry[hanakambur] \dicTerm{hana··kamb|ur} \dicIPA{{h}{a}{\textlengthmark}{n}{a}{k\smash{\textsuperscript{h}}}{a}{m}{\textsubring{b}}{\textscy}{\textsubring{r}}} \dicPos{m}[6] \dicFlx{(‑s, ‑ar)}[11] \textbf{1.} \dicDirectTranslationCS{kohoutí hřebínek}  \textbf{2.} \dicDirectTranslationCS{číro} \dicIndirectTranslationCS{(účes)}  \textbf{3.} \dicFieldCat{bot.} \dicDirectTranslationCS{nevadlec hřebenitý} \textit{(l.~{\textLA{Celosia cristata}})}  \dicsymPhoto\ 
\dicFigure{ds_image_hanakambur_0_1.jpg}{Hanakambur}{Hanakambur - Earth, CC BY-SA 2.5}
\dicEntry[hanastél] \dicTerm{hana··stél} \dicIPA{{h}{a}{\textlengthmark}{n}{a}{s}{\textsubring{d}}{j}{\textepsilon}{\textsubring{l}}} \dicPos{n}[2] \dicFlx{(‑s, ‑)}[5] \dicSynonym*{vínblanda} \dicDirectTranslationCS{koktejl}
\dicEntry[hand] \dicTerm{hand-} \dicIPA{{h}{a}{n}{\textsubring{d}}} \dicPos{predp} \dicDirectTranslationCS{ruční}
\dicEntry[handa] \dicTerm{handa} \dicsymFrequent\  \dicIPA{{h}{a}{n}{\textsubring{d}}{a}} \dicPos{prep} \dicFlx{dat} \dicSynonym{fyrir} \dicDirectTranslationCS{pro} \dicExampleIS{Þessi bók er handa þér.} \dicExampleCS{Tato knížka je pro tebe.}
\dicEntry[handaband] \dicTerm{handa··|band} \dicIPA{{h}{a}{n}{\textsubring{d}}{a}{\textsubring{b}}{a}{n}{\textsubring{d}}} \dicPos{n}[2] \dicFlx{(‑bands, ‑bönd)}[8] \dicDirectTranslationCS{podání\,/\addthin stisk ruky, potřesení rukou}
\dicEntry[handahóf] \dicTerm{handa··hóf} \dicIPA{{h}{a}{n}{\textsubring{d}}{a}{h}{ou}{f}} \dicPos{n}[2] \dicFlx{(‑s)}[2] \dicDirectTranslationCS{nahodilost, náhodnost};  \dicPhraseIS{af handahófi} \dicFlx{adv} \dicDirectTranslationCS{nahodile, náhodně}
\begin{xtolerant}{}{1pt}
\dicEntry[handahófskenndur] \dicTerm{handa·hófs··kenndur}\addthinS\dicIPA{{h}{a}{n}{\textsubring{d}}{a}{h}{ou}{f}{s}{c\smash{\textsuperscript{h}}}{\textepsilon}{n}{\textsubring{d}}{\textscy}{\textsubring{r}}}\addthinS\dicPos{adj}[2]\dicFlx{}[-14]\addthinS\dicDirectTranslationCS{nahodilý, ná\-hodný}
\end{xtolerant}
\dicEntry[handalögmál] \dicTerm{handa··lög·mál} \dicIPA{{h}{a}{n}{\textsubring{d}}{a}{l}{\oe}{\textbabygamma}{m}{au}{\textsubring{l}}} \dicPos{n}[2] \dicFlx{(‑s, ‑)}[5] \dicSynonym{slagur} \dicDirectTranslationCS{rvačka, potyčka, pěstní souboj}
\dicEntry[handan] \dicTerm{handan} \dicsymFrequent\  \dicIPA{{h}{a}{n}{\textsubring{d}}{a}{\textsubring{n}}} \dicPos{prep\,/\addthin adv} \dicFlx{gen} \dicPhraseIS{handan e‑s} \dicFlx{prep} \dicSynonym*{hinum megin} \dicDirectTranslationCS{na druhé straně (čeho)} \dicExampleIS{handan árinnar} \dicExampleCS{na druhé straně řeky};  \dicIdiom{að}[handan]{ \dicPhraseIS{að handan}} \dicFlx{adv} {\textbf{a.}} \dicDirectTranslationCS{z~druhé strany};  {\textbf{b.}} \dicDirectTranslationCS{(vzkaz ap.) ze záhrobí, z~říše mrtvých};  \dicIdiom{handan}[við]{ \dicPhraseIS{handan við e‑ð}} \dicFlx{prep} \dicDirectTranslationCS{na druhé straně (čeho)}
\dicEntry[handapat] \dicTerm{handa··pat} \dicIPA{{h}{a}{n}{\textsubring{d}}{a}{p\smash{\textsuperscript{h}}}{a}{\textsubring{d}}} \dicPos{n}[2] \dicFlx{(‑s)}[2] \dicDirectTranslationCS{gesto, posunek}
\dicEntry[handar] \dicTerm{handar} \dicIPA{{h}{a}{n}{\textsubring{d}}{a}{\textsubring{r}}} \dicPos{f} \dicFlx{sg gen} \dicLink{hönd}
\dicEntry[handarbak] \dicTerm{handar··|bak} \dicIPA{{h}{a}{n}{\textsubring{d}}{a}{r}{\textsubring{b}}{a}{\r{g}}} \dicPos{n}[2] \dicFlx{(‑baks, ‑bök)}[8] \dicDirectTranslationCS{hřbet ruky}
\dicEntry[handargagn] \dicTerm{handar··|gagn} \dicIPA{{h}{a}{n}{\textsubring{d}}{a}{r}{\r{g}}{a}{\r{g}}{\textsubring{n}}} \dicPos{n}[2] \dicFlx{(‑gagns, ‑gögn)}[8] \dicPhraseIS{taka e‑ð til handargagns} \dicDirectTranslationCS{uložit\,/\addthin umístit (co) na dosah ruky}
\dicEntry[handarjaðar] \dicTerm{handar··jað|ar} \dicIPA{{h}{a}{n}{\textsubring{d}}{a}{r}{j}{a}{ð}{a}{\textsubring{r}}} \dicPos{m}[5] \dicFlx{(‑ars, ‑rar)}[8] \dicPhraseIS{vera undir handarjaðri e‑rs} \dicDirectTranslationCS{být pod (čí) ochranou}
\dicEntry[handarkriki] \dicTerm{handar··krik|i} \dicIPA{{h}{a}{n}{\textsubring{d}}{a}{\textsubring{r}}{k\smash{\textsuperscript{h}}}{r}{\textsci}{\r{\textObardotlessj}}{\textsci}} \dicPos{m}[1] \dicFlx{(‑a, ‑ar)}[1] \dicDirectTranslationCS{podpaží, podpaždí}
\dicEntry[handarvik] \dicTerm{handar··vik} \dicIPA{{h}{a}{n}{\textsubring{d}}{a}{r}{v}{\textsci}{\r{g}}} \dicPos{n}[2] \dicFlx{(‑s, ‑)}[5] \dicDirectTranslationCS{malá práce, prácička};  \dicPhraseIS{gera ekki handarvik} \dicLangCat{přen.} \dicDirectTranslationCS{nehnout ani prstem}
\dicEntry[handaskil] \dicTerm{handa··skil} \dicIPA{{h}{a}{n}{\textsubring{d}}{a}{s}{\r{\textObardotlessj}}{\textsci}{\textsubring{l}}} \dicPos{n}[2] \dicFlx{pl}[1] \dicPhraseIS{sjá ekki handaskil} \dicLangCat{přen.} \dicDirectTranslationCS{nevidět ani na krok}
\dicEntry[handaskol] \dicTerm{handa··skol} \dicIPA{{h}{a}{n}{\textsubring{d}}{a}{s}{\r{g}}{\textopeno}{\textsubring{l}}} \dicPos{n}[2] \dicFlx{pl}[1] \dicPhraseIS{e‑að fer í handaskolum hjá e‑m} \dicDirectTranslationCS{(co) se (komu) vymyká z~rukou}
\dicEntry[handavinna] \dicTerm{handa··vinn|a} \dicIPA{{h}{a}{n}{\textsubring{d}}{a}{v}{\textsci}{n}{a}} \dicPos{f}[1] \dicFlx{(‑u)}[5] \textbf{1.} \dicSynonym*{hannyrðir} \dicDirectTranslationCS{ruční práce (vyšívání, pletení ap.)}  \textbf{2.} \dicSynonym*{vinna með höndunum} \dicDirectTranslationCS{manuální\,/\addthin ruční práce}
\dicEntry[handbendi] \dicTerm{hand··bendi} \dicIPA{{h}{a}{n}{\textsubring{d}}{\textsubring{b}}{\textepsilon}{n}{\textsubring{d}}{\textsci}} \dicPos{n}[2] \dicFlx{(‑s, ‑)}[14] \dicDirectTranslationCS{nohsled(ka), přisluhovač(ka)}
\begin{xtolerant}{}{1pt}
\dicEntry[handboltamaður] \dicTerm{hand·bolta··|maður}\addthinSS\dicIPA{{h}{a}{n}{\textsubring{d}}{\textsubring{b}}{\textopeno}{\textsubring{l}}{\textsubring{d}}{a}{m}{a}{ð}{\textscy}{\textsubring{r}}}\addthinS\dicPos{m}[1\hspace{-0.06667em}3]\addthinS\dicFlx{(‑manns,\addthinS‑menn)}[2] \dicDirectTranslationCS{házenkář(ka)}
\end{xtolerant}
\dicEntry[handbolti] \dicTerm{hand··bolt|i} \dicIPA{{h}{a}{n}{\textsubring{d}}{\textsubring{b}}{\textopeno}{\textsubring{l}}{\textsubring{d}}{\textsci}} \dicPos{m}[1] \dicFlx{(‑a, ‑ar)}[1] \textbf{1.} \dicSynonym*{bolti til að kasta} \dicDirectTranslationCS{míč na házení}  \textbf{2.} \dicFieldCat{sport.} \dicSynonym{handknattleikur} \dicDirectTranslationCS{házená}
\dicEntry[handbók] \dicTerm{hand··|bók} \dicIPA{{h}{a}{n}{\textsubring{d}}{\textsubring{b}}{ou}{\r{g}}} \dicPos{f}[8] \dicFlx{(‑bókar, ‑bækur)}[5] \dicDirectTranslationCS{příručka, průvodce, manuál}
\dicEntry[handbragð] \dicTerm{hand··bragð} \dicIPA{{h}{a}{n}{\textsubring{d}}{\textsubring{b}}{r}{a}{\textbabygamma}{\texttheta}} \dicPos{n}[2] \dicFlx{(‑s)}[2] \textbf{1.} \dicSynonym{handverk} \dicDirectTranslationCS{(odborné) vykonání\,/\addthin provedení}  \textbf{2.} \dicSynonym{vinnubrögð} \dicDirectTranslationCS{pracovní metoda\,/\addthin technika}
\dicEntry[handbremsa] \dicTerm{hand··brems|a} \dicIPA{{h}{a}{n}{\textsubring{d}}{\textsubring{b}}{r}{\textepsilon}{m}{s}{a}} \dicPos{f}[1] \dicFlx{(‑u, ‑ur)}[7] \dicDirectTranslationCS{ruční brzda}
\dicEntry[handbrúða] \dicTerm{hand··brúð|a} \dicIPA{{h}{a}{n}{\textsubring{d}}{\textsubring{b}}{r}{u}{ð}{a}} \dicPos{f}[1] \dicFlx{(‑u, ‑ur)}[19] \dicDirectTranslationCS{maňásek}
\dicEntry[handbær] \dicTerm{hand··bær} \dicIPA{{h}{a}{n}{\textsubring{d}}{\textsubring{b}}{a}{i}{\textsubring{r}}} \dicPos{adj}[5] \dicFlx{(f ‑)}[8] \dicDirectTranslationCS{(jsoucí) při ruce, (jsoucí) k~dispozici}
\dicEntry[handbörur] \dicTerm{hand··börur} \dicIPA{{h}{a}{n}{\textsubring{d}}{\textsubring{b}}{\oe}{r}{\textscy}{\textsubring{r}}} \dicPos{f}[12] \dicFlx{pl}[4] \dicDirectTranslationCS{rudl}
\dicEntry[handfang] \dicTerm{hand··|fang} \dicIPA{{h}{a}{n}{\textsubring{d}}{f}{au}{\ng}{\r{g}}} \dicPos{n}[2] \dicFlx{(‑fangs, ‑föng)}[8] \dicDirectTranslationCS{rukojeť, držadlo}
\dicEntry[handfesta] \dicTerm{hand··fest|a} \dicIPA{{h}{a}{n}{\textsubring{d}}{f}{\textepsilon}{s}{\textsubring{d}}{a}} \dicPos{f}[1] \dicFlx{(‑u)}[5] \dicDirectTranslationCS{úchop (na ruku)}
\dicEntry[handfjalla] \dicTerm{hand··fjall|a} \dicIPA{{h}{a}{n}{\textsubring{d}}{f}{j}{a}{\textsubring{d}}{l}{a}} \dicPos{v}[1] \dicFlx{(‑aði)}[32] \dicFlx{acc} \dicDirectTranslationCS{manipulovat, zacházet}
\dicEntry[handfylli] \dicTerm{hand··fyll|i} \dicIPA{{h}{a}{n}{\textsubring{d}}{f}{\textsci}{\textsubring{d}}{l}{\textsci}} \dicPos{f}[5] \dicFlx{(‑i\,/\addthin ‑ar, ‑ir\,/\addthin ‑ar)}[9] \dicDirectTranslationCS{hrst} \dicExampleIS{handfylli af sandi} \dicExampleCS{hrst písku}
\dicEntry[handhafi] \dicTerm{hand··haf|i} \dicIPA{{h}{a}{n}{\textsubring{d}}{h}{a}{v}{\textsci}} \dicPos{m}[1] \dicFlx{(‑a, ‑ar)}[8] \dicDirectTranslationCS{držitel(ka)} \dicExampleIS{handhafar stúdentakortsins} \dicExampleCS{držitelé studentské karty}
\dicEntry[handhægur] \dicTerm{hand··hægur} \dicIPA{{h}{a}{n}{\textsubring{d}}{h}{a}{i}{\textbabygamma}{\textscy}{\textsubring{r}}} \dicPos{adj}[1]\dicFlx{}[-1] \dicDirectTranslationCS{praktický, užitečný, šikovný}
\dicEntry[handiðn] \dicTerm{hand··iðn} \dicIPA{{h}{a}{n}{\textsubring{d}}{\textsci}{ð}{\textsubring{n}}} \dicPos{f}[7] \dicFlx{(‑ar, ‑ir)}[1] \dicDirectTranslationCS{řemeslo, řemeslná výroba}
\dicEntry[handíð] \dicTerm{hand··íð} \dicIPA{{h}{a}{n}{\textsubring{d}}{i}{\texttheta}} \dicPos{f}[7] \dicFlx{(‑ar, ‑ir)}[1] \dicPhraseIS{handíðir} \dicFlx{pl} \dicDirectTranslationCS{řemeslná umění}
\dicEntry[handjárn] \dicTerm{hand··járn} \dicIPA{{h}{a}{n}{\textsubring{d}}{j}{au}{r}{\textsubring{d}}{\textsubring{n}}} \dicPos{n}[2] \dicFlx{(‑s, ‑)}[5] \dicDirectTranslationCS{pouta (na ruce)}
\dicEntry[handklæði] \dicTerm{hand··klæði} \dicsymFrequent\  \dicIPA{{h}{a}{n}{\textsubring{d}}{k\smash{\textsuperscript{h}}}{l}{a}{i}{ð}{\textsci}} \dicPos{n}[2] \dicFlx{(‑s, ‑)}[14] \dicDirectTranslationCS{ručník, osuška} \dicExampleIS{hreint handklæði} \dicExampleCS{čistý ručník}
\dicEntry[handknattleiksmaður] \dicTerm{hand·knatt·leiks··|maður} \dicIPA{{h}\-{a}\-{n}\-{\textsubring{d}}\-{k\smash{\textsuperscript{h}}}\-{n}\-{a}\-{h}\-{\textsubring{d}}\-{l}\-{ei}\-{\r{g}}\-{s}\-{m}\-{a}\-{ð}\-{\textscy}\-{\textsubring{r}}\-} \dicPos{m}[13] \dicFlx{(‑manns, ‑menn)}[2] \dicDirectTranslationCS{házenkář(ka)}
\dicEntry[handknattleikur] \dicTerm{hand··knatt·leik|ur} \dicIPA{{h}{a}{n}{\textsubring{d}}{k\smash{\textsuperscript{h}}}{n}{a}{h}{\textsubring{d}}{l}{ei}{\r{g}}{\textscy}{\textsubring{r}}} \dicPos{m}[9] \dicFlx{(‑s, ‑ir)}[15] \dicFieldCat{sport.} \dicSynonym{handbolti} \dicDirectTranslationCS{házená}
\begin{xtolerant}{}{1pt}
\dicEntry[handlaginn] \dicTerm{hand··laginn}\addthinS\dicIPA{{h}{a}{n}{\textsubring{d}}{l}{a}{i}{j}{\textsci}{\textsubring{n}}}\addthinS\dicPos{adj}[6]\dicFlx{}[-3]\addthinS\dicDirectTranslationCS{obratný, šikovný, zručný}
\end{xtolerant}
\dicEntry[handlagni] \dicTerm{hand··lagn|i} \dicIPA{{h}{a}{n}{\textsubring{d}}{l}{a}{\r{g}}{n}{\textsci}} \dicPos{f}[3] \dicFlx{(‑i)}[3] \dicDirectTranslationCS{obratnost, šikovnost, zručnost}
\dicEntry[handlaug] \dicTerm{hand··laug} \dicIPA{{h}{a}{n}{\textsubring{d}}{l}{\oe i}{x}} \dicPos{f}[4] \dicFlx{(‑ar, ‑ar)}[1] \dicDirectTranslationCS{umyvadlo (na ruce)}
\dicEntry[handleggur] \dicTerm{hand··legg|ur} \dicsymFrequent\  \dicIPA{{h}{a}{n}{\textsubring{d}}{l}{\textepsilon}{\r{g}}{\textscy}{\textsubring{r}}} \dicPos{m}[9] \dicFlx{(‑jar\,/\addthin ‑s, ‑ir)}[26] \dicFieldCat{anat.} \dicDirectTranslationCS{ruka, paže} \dicIndirectTranslationCS{(celá končetina)} \dicExampleIS{taka undir handlegg e‑rs} \dicExampleCS{vzít (koho) za ruku}
\dicEntry[handleiðsla] \dicTerm{hand··leiðsl|a} \dicIPA{{h}{a}{n}{\textsubring{d}}{l}{ei}{ð}{s}{\textsubring{d}}{l}{a}} \dicPos{f}[1] \dicFlx{(‑u)}[5] \dicDirectTranslationCS{vedení, dozor};  \dicPhraseIS{undir handleiðslu e‑rs} \dicDirectTranslationCS{pod (čím) vedením}
\dicEntry[handleik] \dicTerm{hand··leik} \dicIPA{{h}{a}{n}{\textsubring{d}}{l}{ei}{\r{g}}} \dicPos{v} \dicFlx{ind praes sg 1 pers} \dicLink{handleika}
\dicEntry[handleika] \dicTerm{hand··|leika} \dicIPA{{h}{a}{n}{\textsubring{d}}{l}{ei}{\r{g}}{a}} \dicPos{v}[6] \dicFlx{(‑leik, ‑lék, ‑lékum, ‑léki, ‑leikið)}[63] \dicFlx{acc} \dicDirectTranslationCS{manipulovat, zacházet} \dicExampleIS{handleika e‑ð} \dicExampleCS{zacházet s~(čím)}
\dicEntry[handleikið] \dicTerm{hand··leikið} \dicIPA{{h}{a}{n}{\textsubring{d}}{l}{ei}{\r{\textObardotlessj}}{\textsci}{\texttheta}} \dicPos{v} \dicFlx{supin} \dicLink{handleika}
\dicEntry[handlék] \dicTerm{hand··lék} \dicIPA{{h}{a}{n}{\textsubring{d}}{l}{j}{\textepsilon}{\r{g}}} \dicPos{v} \dicFlx{ind pf sg 1 pers} \dicLink{handleika}
\dicEntry[handléki] \dicTerm{hand··léki} \dicIPA{{h}{a}{n}{\textsubring{d}}{l}{j}{\textepsilon}{\r{\textObardotlessj}}{\textsci}} \dicPos{v} \dicFlx{con pf sg 1 pers} \dicLink{handleika}
\dicEntry[handlékum] \dicTerm{hand··lékum} \dicIPA{{h}{a}{n}{\textsubring{d}}{l}{j}{\textepsilon}{\r{g}}{\textscy}{\textsubring{m}}} \dicPos{v} \dicFlx{ind pf pl 1 pers} \dicLink{handleika}
\dicEntry[handónýtur] \dicTerm{hand··ó·nýtur} \dicIPA{{h}{a}{n}{\textsubring{d}}{ou}{n}{i}{\textsubring{d}}{\textscy}{\textsubring{r}}} \dicPos{adj}[1]\dicFlx{}[-6] \dicDirectTranslationCS{naprosto\,/\addthin zcela nepoužitelný}
\dicEntry[handrið] \dicTerm{hand··rið} \dicIPA{{h}{a}{n}{\textsubring{d}}{r}{\textsci}{\texttheta}} \dicPos{n}[2] \dicFlx{(‑s, ‑)}[5] \dicDirectTranslationCS{zábradlí, madlo, opěradlo}
\dicEntry[handrit] \dicTerm{hand··rit} \dicsymFrequent\  \dicIPA{{h}{a}{n}{\textsubring{d}}{r}{\textsci}{\textsubring{d}}} \dicPos{n}[2] \dicFlx{(‑s, ‑)}[5] \textbf{1.} \dicDirectTranslationCS{rukopis} \dicIndirectTranslationCS{(stará psaná literární památka)} \dicExampleIS{handrit að Egils sögu} \dicExampleCS{rukopis Ságy o~Egilovi}  \textbf{2.} \dicDirectTranslationCS{rukopis} \dicIndirectTranslationCS{(podoba slovesného díla jako podklad pro tisk)}  \textbf{3.} \dicDirectTranslationCS{scénář}
\dicEntry[handritshöfundur] \dicTerm{hand·rits··höf·und|ur} \dicIPA{{h}\-{a}\-{n}\-{\textsubring{d}}\-{r}\-{\textsci}\-{\textsubring{d}}\-{s}\-{h}\-{\oe}\-{v}\-{\textscy}\-{n}\-{\textsubring{d}}\-{\textscy}\-{\textsubring{r}}\-} \dicPos{m}[6] \dicFlx{(‑ar, ‑ar)}[56] \dicDirectTranslationCS{scenárista, scenáristka}
\dicEntry[handsama] \dicTerm{hand··sam|a} \dicIPA{{h}{a}{n}{\textsubring{d}}{s}{a}{m}{a}} \dicPos{v}[1] \dicFlx{(‑aði)}[13] \dicFlx{acc} \dicDirectTranslationCS{dopadnout, chytit, polapit} \dicExampleIS{handsama njósnara} \dicExampleCS{dopadnout špeha}
\dicEntry[handsápa] \dicTerm{hand··sáp|a} \dicIPA{{h}{a}{n}{\textsubring{d}}{s}{au}{\textsubring{b}}{a}} \dicPos{f}[1] \dicFlx{(‑u, ‑ur)}[19] \dicDirectTranslationCS{ruční mýdlo, mýdlo na ruce}
\dicEntry[handskrift] \dicTerm{hand··skrift} \dicIPA{{h}{a}{n}{\textsubring{d}}{s}{\r{g}}{r}{\textsci}{f}{\textsubring{d}}} \dicPos{f}[7] \dicFlx{(‑ar, ‑ir)}[1] \dicSynonym{rithönd} \dicDirectTranslationCS{rukopis, písmo} \dicIndirectTranslationCS{(způsob psaní charakteristický pro jedince)}
\dicEntry[handsnyrting] \dicTerm{hand··snyrt·ing} \dicIPA{{h}{a}{n}{\textsubring{d}}{s}{\textsubring{d}}{n}{\textsci}{\textsubring{r}}{\textsubring{d}}{i}{\ng}{\r{g}}} \dicPos{f}[4] \dicFlx{(‑ar)}[7] \dicDirectTranslationCS{manikúra} \dicIndirectTranslationCS{(odborná péče o~ruce a~nehty)}
\dicEntry[handsprengja] \dicTerm{hand··sprengj|a} \dicIPA{{h}{a}{n}{\textsubring{d}}{s}{\textsubring{b}}{r}{ei}{\textltailn}{\r{\textObardotlessj}}{a}} \dicPos{f}[1] \dicFlx{(‑u, ‑ur)}[25] \dicFieldCat{voj.} \dicDirectTranslationCS{(ruční) granát}
\dicEntry[handsög] \dicTerm{hand··|sög} \dicIPA{{h}{a}{n}{\textsubring{d}}{s}{\oe}{x}} \dicPos{f}[7] \dicFlx{(‑sagar, ‑sagir)}[16] \dicDirectTranslationCS{ruční pila}
\dicEntry[handtak] \dicTerm{hand··|tak} \dicIPA{{h}{a}{n}{t\smash{\textsuperscript{h}}}{a}{\r{g}}} \dicPos{n}[2] \dicFlx{(‑taks, ‑tök)}[8] \textbf{1.} \dicSynonym{handaband} \dicDirectTranslationCS{stisk (ruky)} \dicExampleIS{hlýtt handtak} \dicExampleCS{vřelý stisk}  \textbf{2.} \dicPhraseIS{handtök} \dicFlx{pl} \dicSynonym{aðferð} \dicDirectTranslationCS{(obratný) způsob, grif, kumšt};  \dicPhraseIS{kunna handtökin á e‑u} \dicDirectTranslationCS{znát na (co) ten správný grif};  \dicIdiom{handtak}{ \dicPhraseIS{gera ekki handtak}} \dicLangCat{přen.} \dicDirectTranslationCS{nehnout ani prstem}
\dicEntry[handtaka] \dicTerm{hand··|taka\smash{\textsuperscript{1}}} \dicIPA{{h}{a}{n}{t\smash{\textsuperscript{h}}}{a}{\r{g}}{a}} \dicPos{f}[1] \dicFlx{(‑töku, ‑tökur)}[20] \dicDirectTranslationCS{zatčení, zatknutí}
\dicEntry[handtaka] \dicTerm{hand··|taka\smash{\textsuperscript{2}}} \dicIPA{{h}{a}{n}{t\smash{\textsuperscript{h}}}{a}{\r{g}}{a}} \dicPos{v}[6] \dicFlx{(‑tek, ‑tók, ‑tókum, ‑tæki, ‑tekið)}[128] \dicFlx{acc} \dicDirectTranslationCS{zatknout} \dicExampleIS{handtaka innbrotsþjófinn} \dicExampleCS{zatknout lupiče}
\dicEntry[handtaska] \dicTerm{hand··|taska} \dicIPA{{h}{a}{n}{t\smash{\textsuperscript{h}}}{a}{s}{\r{g}}{a}} \dicPos{f}[1] \dicFlx{(‑tösku, ‑töskur)}[8] \dicDirectTranslationCS{kabelka}
\dicEntry[handtek] \dicTerm{hand··tek} \dicIPA{{h}{a}{n}{t\smash{\textsuperscript{h}}}{\textepsilon}{\r{g}}} \dicPos{v} \dicFlx{ind praes sg 1 pers} \dicLink{handtaka\smash{\textsuperscript{2}}}
\dicEntry[handtekið] \dicTerm{hand··tekið} \dicIPA{{h}{a}{n}{t\smash{\textsuperscript{h}}}{\textepsilon}{\r{\textObardotlessj}}{\textsci}{\texttheta}} \dicPos{v} \dicFlx{supin} \dicLink{handtaka\smash{\textsuperscript{2}}}
\dicEntry[handtók] \dicTerm{hand··tók} \dicIPA{{h}{a}{n}{t\smash{\textsuperscript{h}}}{ou}{\r{g}}} \dicPos{v} \dicFlx{ind pf sg 1 pers} \dicLink{handtaka\smash{\textsuperscript{2}}}
\dicEntry[handtókum] \dicTerm{hand··tókum} \dicIPA{{h}{a}{n}{t\smash{\textsuperscript{h}}}{ou}{\r{g}}{\textscy}{\textsubring{m}}} \dicPos{v} \dicFlx{ind pf pl 1 pers} \dicLink{handtaka\smash{\textsuperscript{2}}}
\dicEntry[handtæki] \dicTerm{hand··tæki} \dicIPA{{h}{a}{n}{t\smash{\textsuperscript{h}}}{a}{i}{\r{\textObardotlessj}}{\textsci}} \dicPos{v} \dicFlx{con pf sg 1 pers} \dicLink{handtaka\smash{\textsuperscript{2}}}
\dicEntry[handtölva] \dicTerm{hand··tölv|a} \dicIPA{{h}{a}{n}{t\smash{\textsuperscript{h}}}{\oe}{l}{v}{a}} \dicPos{f}[1] \dicFlx{(‑u, ‑ur)}[7] \dicFieldCat{poč.} \dicDirectTranslationCS{kapesní počítač, PDA}
\dicEntry[handveð] \dicTerm{hand··veð} \dicIPA{{h}{a}{n}{\textsubring{d}}{v}{\textepsilon}{\texttheta}} \dicPos{n}[2] \dicFlx{(‑s, ‑)}[5] \dicFieldCat{práv.} \dicDirectTranslationCS{zástava, zastavený předmět}
\dicEntry[handveðslán] \dicTerm{hand·veðs··lán} \dicIPA{{h}{a}{n}{\textsubring{d}}{v}{\textepsilon}{ð}{s}{l}{au}{\textsubring{n}}} \dicPos{n}[2] \dicFlx{(‑s, ‑)}[5] \dicFieldCat{ekon.} \dicDirectTranslationCS{kolaterální půjčka}
\dicEntry[handverk] \dicTerm{hand··verk} \dicIPA{{h}{a}{n}{\textsubring{d}}{v}{\textepsilon}{\textsubring{r}}{\r{g}}} \dicPos{n}[2] \dicFlx{(‑s)}[2] \dicSynonym{handiðn} \dicDirectTranslationCS{řemeslo, řemeslná výroba}
\begin{xtolerant}{}{1pt}
\dicEntry[handverksmaður] \dicTerm{hand··verks·|maður}\addthinSS\dicIPA{{h}{a}{n}{\textsubring{d}}{v}{\textepsilon}{\textsubring{r}}{\r{g}}{s}{m}{a}{ð}{\textscy}{\textsubring{r}}}\addthinSS\dicPos{m}[1\hspace{-0.06667em}3]\addthinSS\dicFlx{(‑manns,\addthinSS‑menn)}[2] \dicDirectTranslationCS{řemeslník, řemeslnice}
\end{xtolerant}
\dicEntry[handviss] \dicTerm{hand··viss} \dicIPA{{h}{a}{n}{\textsubring{d}}{v}{\textsci}{s}} \dicPos{adj}[5]\dicFlx{}[-4] \dicDirectTranslationCS{absolutně jistý, skálopevný}
\dicEntry[handvömm] \dicTerm{hand··|vömm} \dicIPA{{h}{a}{n}{\textsubring{d}}{v}{\oe}{\textsubring{m}}} \dicPos{f}[7] \dicFlx{(‑vammar)}[19] \dicDirectTranslationCS{ledabylost, nedbalost}
\dicEntry[hanga] \dicTerm{hanga} \dicsymFrequent\  \dicIPA{{h}{au}{\ng}{\r{g}}{a}} \dicPos{v}[5] \dicFlx{(hangi, hékk, héngum, héngi, hangið)}[11] \textbf{1.} \dicDirectTranslationCS{viset} \dicExampleIS{Spegillinn hangir á veggnum.} \dicExampleCS{Zrcadlo visí na zdi.}  \textbf{2.} \dicSynonym{slóra} \dicDirectTranslationCS{lenošit, flákat se} \dicExampleIS{hanga inni allan daginn} \dicExampleCS{flákat se doma celý den}  \textbf{3.} \dicDirectTranslationCS{viset (dítě na matce ap.), lnout}
\dicEntry[hangi] \dicTerm{hangi} \dicIPA{{h}{au}{\textltailn}{\r{\textObardotlessj}}{\textsci}} \dicPos{v} \dicFlx{ind praes sg 1 pers} \dicLink{hanga}
\dicEntry[hangið] \dicTerm{hangið} \dicIPA{{h}{au}{\textltailn}{\r{\textObardotlessj}}{\textsci}{\texttheta}} \dicPos{v} \dicFlx{supin} \dicLink{hanga}
\dicEntry[hangikjöt] \dicTerm{hangi··kjöt} \dicIPA{{h}{au}{\textltailn}{\r{\textObardotlessj}}{\textsci}{c\smash{\textsuperscript{h}}}{\oe}{\textsubring{d}}} \dicPos{n}[2] \dicFlx{(‑s)}[2] \dicDirectTranslationCS{uzené skopové (maso)}
\dicEntry[hangsa] \dicTerm{hangs|a} \dicIPA{{h}{au}{\ng}{s}{a}} \dicPos{v}[1] \dicFlx{(‑aði)}[23] \dicSynonym{slóra} \dicDirectTranslationCS{zahálet, lelkovat, zevlovat}
\dicEntry[hani] \dicTerm{han|i} \dicIPA{{h}{a}{\textlengthmark}{n}{\textsci}} \dicPos{m}[1] \dicFlx{(‑a, ‑ar)}[8] \textbf{1.} \dicDirectTranslationCS{kohout} \dicIndirectTranslationCS{(samec kura domácího)}  \textbf{2.} \dicSynonym{krani} \dicDirectTranslationCS{kohout, kohoutek, uzávěr}
\dicEntry[hankar] \dicTerm{hankar} \dicIPA{{h}{au}{\r{\ng}}{\r{g}}{a}{\textsubring{r}}} \dicPos{f} \dicFlx{sg gen} \dicLink{hönk}
\dicEntry[hanki] \dicTerm{hank|i} \dicIPA{{h}{au}{\r{\textltailn}}{\r{\textObardotlessj}}{\textsci}} \dicPos{m}[1] \dicFlx{(‑a, ‑ar)}[8] \textbf{1.} \dicSynonym{halda\smash{\textsuperscript{1}}} \dicDirectTranslationCS{držadlo, držátko, rukojeť}  \textbf{2.} \dicSynonym{lykkja} \dicDirectTranslationCS{poutko, ouško}
\dicEntry[hankir] \dicTerm{hankir} \dicIPA{{h}{au}{\r{\textltailn}}{\r{\textObardotlessj}}{\textsci}{\textsubring{r}}} \dicPos{f} \dicFlx{pl nom} \dicLink{hönk}
\dicEntry[hann] \dicTerm{hann} \dicsymFrequent\  \dicIPA{{h}{a}{n}{\textlengthmark}} \dicPos{pron}[17] \dicFlx{pers} \dicFlx{nom sg m} \textbf{1.} \dicDirectTranslationCS{on} \dicExampleIS{Hann lærir íslensku.} \dicExampleCS{Učí se islandsky.}  \textbf{2.} \dicIndirectTranslationCS{(užívané jako zpětný odkaz na maskulinum)}  \textbf{3.} \dicIndirectTranslationCS{(užívané před osobními jmény, vyjadřuje známost nebo blízké vztahy)}  \textbf{4.} \dicIndirectTranslationCS{(užívané podobně jako \clqq það\crqq  se slovesy opisujícími počasí)} \dicExampleIS{Hann kólnar.} \dicExampleCS{Ochlazuje se.}
\dicEntry[hanna] \dicTerm{hann|a} \dicIPA{{h}{a}{n}{\textlengthmark}{a}} \dicPos{v}[1] \dicFlx{(‑aði)}[13] \dicFlx{acc} \dicDirectTranslationCS{navrhovat (šaty ap.), konstruovat} \dicExampleIS{hanna húsgagn} \dicExampleCS{navrhovat nábytek}
\dicEntry[Hanoí] \dicTerm{Hanoí} \dicIPA{{h}{a}{\textlengthmark}{n}{\textopeno i}} \dicPos{subs} \dicFlx{indecl} \dicFieldCat{geog.} \dicDirectTranslationCS{Hanoj} \dicIndirectTranslationCS{(hlavní město Vietnamu)}
\dicEntry[hans] \dicTerm{hans\smash{\textsuperscript{1}}} \dicIPA{{h}{a}{n}{s}} \dicPos{pron}[17] \dicFlx{pers} \dicFlx{sg gen} \dicLink{hann}
\dicEntry[hans] \dicTerm{hans\smash{\textsuperscript{2}}} \dicIPA{{h}{a}{n}{s}} \dicPos{pron}[17] \dicFlx{poss m} \dicDirectTranslationCS{jeho} \dicExampleIS{faðir hans} \dicExampleCS{jeho otec}
\dicEntry[Hansakaupstaður] \dicTerm{Hansa··kaup·stað|ur} \dicIPA{{h}{a}{n}{s}{a}{k\smash{\textsuperscript{h}}}{\oe i}{\textsubring{b}}{s}{\textsubring{d}}{a}{ð}{\textscy}{\textsubring{r}}} \dicPos{m}[10] \dicFlx{(‑ar, ‑ir)}[14] \dicFieldCat{hist.} \dicDirectTranslationCS{hanzovní město}
\dicEntry[hanskahólf] \dicTerm{hanska··hólf} \dicIPA{{h}{a}{n}{s}{\r{g}}{a}{h}{ou}{l}{f}} \dicPos{n}[2] \dicFlx{(‑s, ‑)}[5] \dicDirectTranslationCS{palubní přihrádka, odkládací přihrádka u~spolujezdce} \dicIndirectTranslationCS{(v~autě)}
\dicEntry[hanski] \dicTerm{hansk|i} \dicIPA{{h}{a}{n}{s}{\r{\textObardotlessj}}{\textsci}} \dicPos{m}[1] \dicFlx{(‑a, ‑ar)}[8] \dicSynonym{vettlingur} \dicDirectTranslationCS{rukavice};  \dicPhraseIS{taka upp hanskann fyrir e‑n} \dicDirectTranslationCS{zastat\,/\addthin zastávat se (koho)};  \dicPhraseIS{kasta hanska í andlit e‑rs} \dicLangCat{přen.} \dicSynonym*{hefja árás} \dicDirectTranslationCS{hodit (komu) do tváře rukavici}
\dicEntry[happ] \dicTerm{happ} \dicIPA{{h}{a}{h}{\textsubring{b}}} \dicPos{n}[2] \dicFlx{(‑s, höpp)}[8] \dicDirectTranslationCS{štěstí, šťastná náhoda}
\dicEntry[happdrætti] \dicTerm{happ··drætti} \dicIPA{{h}{a}{h}{\textsubring{b}}{\textsubring{d}}{r}{a}{i}{h}{\textsubring{d}}{\textsci}} \dicPos{n}[2] \dicFlx{(‑s, ‑)}[14] \dicDirectTranslationCS{loterie} \dicExampleIS{vinna í happdrætti} \dicExampleCS{vyhrát v~loterii}
\dicEntry[harðahlaup] \dicTerm{harða··hlaup} \dicIPA{{h}{a}{r}{ð}{a}{\textsubring{l}}{\oe i}{\textsubring{b}}} \dicPos{n}[2] \dicFlx{pl}[1] \dicPhraseIS{koma á harðahlaupum} \dicDirectTranslationCS{sprintovat, běžet plnou rychlostí}
\dicEntry[harðbrjósta] \dicTerm{harð··brjósta} \dicIPA{{h}{a}{r}{ð}{\textsubring{b}}{r}{j}{ou}{s}{\textsubring{d}}{a}} \dicPos{adj}[13] \dicFlx{indecl}[1] \dicDirectTranslationCS{bezcitný, nelítostný, nemilosrdný}
\dicEntry[harðdiskur] \dicTerm{harð··disk|ur} \dicIPA{{h}{a}{r}{ð}{\textsubring{d}}{\textsci}{s}{\r{g}}{\textscy}{\textsubring{r}}} \dicPos{m}[6] \dicFlx{(‑s, ‑ar)}[10] \dicFieldCat{poč.} \dicDirectTranslationCS{harddisk, pevný disk}
\dicEntry[harðfengi] \dicTerm{harð··fengi\smash{\textsuperscript{1}}}\dicTerm{, harðfengi\smash{\textsuperscript{2}}} \dicIPA{{h}\-{a}\-{r}\-{\texttheta}\-{f}\-{ei}\-{\textltailn}\-{\r{\textObardotlessj}}\-{\textsci}\-} \dicPos{n}[2] \dicFlx{(‑s)}[20] \dicDirectTranslationCS{houževnatost, vytrvalost}
\dicEntry[harðfengi] \dicTerm{harð··feng|i\smash{\textsuperscript{2}}} \dicIPA{{h}{a}{r}{\texttheta}{f}{ei}{\textltailn}{\r{\textObardotlessj}}{\textsci}} \dicPos{f}[3] \dicFlx{(‑i)}[3] \dicLink{harðfengi\smash{\textsuperscript{1}}}
\dicEntry[harðfiskur] \dicTerm{harð··fisk|ur} \dicIPA{{h}{a}{r}{\texttheta}{f}{\textsci}{s}{\r{g}}{\textscy}{\textsubring{r}}} \dicPos{m}[6] \dicFlx{(‑s, ‑ar)}[10] \dicFieldCat{kulin.} \dicSynonym{skreið\smash{\textsuperscript{1}}} \dicDirectTranslationCS{sušená ryba}
\dicEntry[harðfylgi] \dicTerm{harð··fylgi} \dicIPA{{h}{a}{r}{\texttheta}{f}{\textsci}{l}{\r{\textObardotlessj}}{\textsci}} \dicPos{n}[2] \dicFlx{(‑s)}[20] \dicDirectTranslationCS{houževnatost, vytrvalost}
\dicEntry[harðfylginn] \dicTerm{harð··fylginn} \dicIPA{{h}{a}{r}{\texttheta}{f}{\textsci}{l}{\r{\textObardotlessj}}{\textsci}{\textsubring{n}}} \dicPos{adj}[6]\dicFlx{}[-2] \dicDirectTranslationCS{houževnatý, vytrvalý}
\dicEntry[harðhaus] \dicTerm{harð··haus} \dicIPA{{h}{a}{r}{\texttheta}{h}{\oe i}{s}} \dicPos{m}[4] \dicFlx{(‑s, ‑ar)}[14] \dicDirectTranslationCS{tvrďák, drsňák}
\dicEntry[harðindi] \dicTerm{harð··indi} \dicIPA{{h}{a}{r}{ð}{\textsci}{n}{\textsubring{d}}{\textsci}} \dicPos{n}[2] \dicFlx{pl}[19] \dicDirectTranslationCS{těžké časy}
\dicEntry[harðjaxl] \dicTerm{harð··jaxl} \dicIPA{{h}{a}{r}{ð}{j}{a}{x}{s}{\textsubring{d}}{\textsubring{l}}} \dicPos{m}[4] \dicFlx{(‑s, ‑ar)}[4] \dicDirectTranslationCS{tvrďák, drsňák}
\dicEntry[harðleikinn] \dicTerm{harð··leikinn} \dicIPA{{h}{a}{r}{ð}{l}{ei}{\r{\textObardotlessj}}{\textsci}{\textsubring{n}}} \dicPos{adj}[6]\dicFlx{}[-2] \dicDirectTranslationCS{tvrdý, drsný, brutální}
\dicEntry[harðleikni] \dicTerm{harð··leikn|i} \dicIPA{{h}{a}{r}{ð}{l}{ei}{h}{\r{g}}{n}{\textsci}} \dicPos{f}[3] \dicFlx{(‑i)}[3] \dicSynonym{fólska} \dicDirectTranslationCS{tvrdost, drsnost, brutalita}
\dicEntry[harðlífi] \dicTerm{harð··lífi} \dicIPA{{h}{a}{r}{ð}{l}{i}{v}{\textsci}} \dicPos{n}[2] \dicFlx{(‑s)}[20] \dicFieldCat{med.} \dicDirectTranslationCS{zácpa}
\dicEntry[harðlæstur] \dicTerm{harð··læstur} \dicIPA{{h}{a}{r}{ð}{l}{a}{i}{s}{\textsubring{d}}{\textscy}{\textsubring{r}}} \dicPos{adj}[1]\dicFlx{}[-13] \dicDirectTranslationCS{bezpečně (u)zamčený}
\dicEntry[harðmæli] \dicTerm{harð··mæli} \dicIPA{{h}{a}{r}{ð}{m}{a}{i}{l}{\textsci}} \dicPos{n}[2] \dicFlx{(‑s)}[20] \dicFieldCat{jaz.} \dicIndirectTranslationCS{druh výslovnosti -- používání aspirované plozivy ([p\smash{\textsuperscript{h}}], [t\smash{\textsuperscript{h}}], [k\smash{\textsuperscript{h}}], [c\smash{\textsuperscript{h}}]) uprostřed slova po dlouhé samohlásce (výslovnostní varianta používaná zejména na severovýchodě Islandu)} \dicAntonym{linmæli}
\dicEntry[harðna] \dicTerm{harðn|a} \dicIPA{{h}{a}{r}{\textsubring{d}}{n}{a}} \dicPos{v}[1] \dicFlx{(‑aði)}[42] \textbf{1.} \dicSynonym*{verða harður} \dicDirectTranslationCS{(z)tvrdnout, zatvrdnout} \dicExampleIS{Steypan harðnar.} \dicExampleCS{Beton tvrdne.}  \textbf{2.} \dicSynonym{versna} \dicDirectTranslationCS{přituhovat (zima ap.), zhoršovat se};  \dicIdiom{harðna}[að\,/\addthin á]{ \dicPhraseIS{það harðnar að\,/\addthin á}} \dicFlx{impers} \dicDirectTranslationCS{jde do tuhého, přituhuje (v~situaci ap.)};  \dicIdiom{harðna}[í]{ \dicPhraseIS{það harðnar í ári}} \dicFlx{impers} \dicDirectTranslationCS{zhoršuje se, horší se (situace ap.)}
\dicEntry[harðneita] \dicTerm{harð··neit|a} \dicIPA{{h}{a}{r}{ð}{n}{ei}{\textsubring{d}}{a}} \dicPos{v}[1] \dicFlx{(‑aði)}[44] \dicFlx{dat} \dicDirectTranslationCS{kategoricky\,/\addthin rozhodně odmítnout\,/\addthin odmítat}
\dicEntry[harðneskja] \dicTerm{harð··neskj|a} \dicIPA{{h}{a}{r}{ð}{n}{\textepsilon}{s}{\r{\textObardotlessj}}{a}} \dicPos{f}[1] \dicFlx{(‑u)}[5] \dicSynonym{harka\smash{\textsuperscript{1}}} \dicDirectTranslationCS{drsnost, tvrdost, přísnost}
\dicEntry[harðorður] \dicTerm{harð··orður} \dicIPA{{h}{a}{r}{ð}{\textopeno}{r}{ð}{\textscy}{\textsubring{r}}} \dicPos{adj}[2]\dicFlx{}[-1] \dicDirectTranslationCS{tvrdý, ostrý (výraz ap.)}
\dicEntry[harðskeyttur] \dicTerm{harð··skeyttur} \dicIPA{{h}{a}{r}{\texttheta}{s}{\r{\textObardotlessj}}{ei}{h}{\textsubring{d}}{\textscy}{\textsubring{r}}} \dicPos{adj}[1]\dicFlx{}[-10] \dicDirectTranslationCS{bojovný, urputný}
\dicEntry[harðsnúinn] \dicTerm{harð··snúinn} \dicIPA{{h}{a}{r}{\texttheta}{s}{\textsubring{d}}{n}{u}{\textsci}{\textsubring{n}}} \dicPos{adj}[6]\dicFlx{}[-2] \dicSynonym{hraustur} \dicDirectTranslationCS{houževnatý, tvrdý} \dicExampleIS{harðsnúið lið} \dicExampleCS{houževnatá družina}
\dicEntry[harðsoðinn] \dicTerm{harð··soðinn} \dicIPA{{h}{a}{r}{\texttheta}{s}{\textopeno}{ð}{\textsci}{\textsubring{n}}} \dicPos{adj}[6]\dicFlx{}[-2] \textbf{1.} \dicDirectTranslationCS{(uvařený) natvrdo} \dicExampleIS{harðsoðið egg} \dicExampleCS{vejce natvrdo}  \textbf{2.} \dicSynonym{svalur} \dicDirectTranslationCS{zocelený, ostřílený}
\dicEntry[harðsperrur] \dicTerm{harð··sperrur} \dicIPA{{h}{a}{r}{\texttheta}{s}{\textsubring{b}}{\textepsilon}{r}{\textscy}{\textsubring{r}}} \dicPos{f}[1] \dicFlx{pl}[12] \dicDirectTranslationCS{namožené svaly}
\dicEntry[harðstjóri] \dicTerm{harð··stjór|i} \dicIPA{{h}{a}{r}{\texttheta}{s}{\textsubring{d}}{j}{ou}{r}{\textsci}} \dicPos{m}[1] \dicFlx{(‑a, ‑ar)}[1] \dicDirectTranslationCS{tyran(ka), despota, despotka}
\dicEntry[harður] \dicTerm{harður} \dicsymFrequent\  \dicIPA{{h}{a}{r}{ð}{\textscy}{\textsubring{r}}} \dicPos{adj}[2] \dicFlx{(f hörð)}[2] \textbf{1.} \dicDirectTranslationCS{tvrdý, tuhý, pevný (materiál ap.)}  \textbf{2.} \dicDirectTranslationCS{tvrdý, drsný, nelítostný (boj ap.)}  \textbf{3.} \dicSynonym{strangur} \dicDirectTranslationCS{tvrdý, přísný (trest ap.)}  \textbf{4.} \dicSynonym{erfiður} \dicDirectTranslationCS{tvrdý, těžký (zima ap.)};  \dicIdiom{harður}{ \dicPhraseIS{láta hart mæta hörðu}} \dicLangCat{přen.} \dicDirectTranslationCS{zaujmout tvrdý postoj};  \dicPhraseIS{verður harður af sér} \dicDirectTranslationCS{být na sebe tvrdý, být statečný};  \dicPhraseIS{vera harður í horn að taka} \dicLangCat{přen.} \dicDirectTranslationCS{být osoba, se kterou se obtížně jedná};  \dicPhraseIS{vinna hörðum höndum að e‑u} \dicDirectTranslationCS{tvrdě na (čem) pracovat};  \dicPhraseIS{e‑m finnst það hart að (gera e‑ð)} \dicFlx{impers} \dicDirectTranslationCS{(komu) připadá těžké (udělat (co))}
\dicEntry[hark] \dicTerm{hark} \dicIPA{{h}{a}{\textsubring{r}}{\r{g}}} \dicPos{n}[2] \dicFlx{(‑s)}[2] \textbf{1.} \dicSynonym{hávaði} \dicDirectTranslationCS{hluk, rámus}  \textbf{2.} \dicDirectTranslationCS{pomocná práce}  \textbf{3.} \dicDirectTranslationCS{čára} \dicIndirectTranslationCS{(dětská hra s~mincemi)}
\dicEntry[harka] \dicTerm{harka\smash{\textsuperscript{1}}} \dicIPA{{h}{a}{\textsubring{r}}{\r{g}}{a}} \dicPos{f}[1] \dicFlx{(hörku, hörkur)}[20] \textbf{1.} \dicSynonym{harðneskja} \dicDirectTranslationCS{tvrdost, drsnost, surovost}  \textbf{2.} \dicFieldCat{geol.} \dicDirectTranslationCS{tvrdost}
\dicEntry[harka] \dicTerm{hark|a\smash{\textsuperscript{2}}} \dicIPA{{h}{a}{\textsubring{r}}{\r{g}}{a}} \dicPos{v}[1] \dicFlx{(‑aði)}[43] \dicPhraseIS{harka af sér} \dicLangCat{přen.} \dicDirectTranslationCS{zatnout zuby}
\dicEntry[harkalega] \dicTerm{harka··leg|a} \dicsymFrequent\  \dicIPA{{h}{a}{\textsubring{r}}{\r{g}}{a}{l}{\textepsilon}{\textbabygamma}{a}} \dicPos{adv} \dicFlx{(comp ‑ar, sup ‑ast)} \dicDirectTranslationCS{tvrdě, drsně, ostře, nelítostně} \dicExampleIS{gagnrýna e‑ð harkalega} \dicExampleCS{ostře (co) kritizovat}
\dicEntry[harkalegur] \dicTerm{harka··legur} \dicIPA{{h}{a}{\textsubring{r}}{\r{g}}{a}{l}{\textepsilon}{\textbabygamma}{\textscy}{\textsubring{r}}} \dicPos{adj}[1]\dicFlx{}[-8] \dicDirectTranslationCS{tvrdý, drsný, nelítostný}
\dicEntry[harla] \dicTerm{harla} \dicIPA{{h}{a}{r}{\textsubring{d}}{l}{a}} \dicPos{adv} \dicFlx{(comp fremur, sup fremst)} \textbf{1.} \dicSynonym{mjög} \dicDirectTranslationCS{velmi, hodně}  \textbf{2.} \dicSynonym{frekar} \dicDirectTranslationCS{docela, dost}
\dicEntry[harma] \dicTerm{harm|a} \dicIPA{{h}{a}{r}{m}{a}} \dicPos{v}[1] \dicFlx{(‑aði)}[43] \dicFlx{acc} \textbf{1.} \dicDirectTranslationCS{želet, truchlit, oplakávat};  \dicPhraseIS{harma e‑n} \dicDirectTranslationCS{truchlit nad (kým)}  \textbf{2.} \dicDirectTranslationCS{(hluboce) litovat};  \dicPhraseIS{harma e‑ð} \dicDirectTranslationCS{(hluboce) litovat (čeho)}
\dicEntry[harmdauði] \dicTerm{harm··dauði} \dicIPA{{h}{a}{r}{m}{\textsubring{d}}{\oe i}{ð}{\textsci}} \dicPos{adj}[13] \dicFlx{indecl}[1] \dicDirectTranslationCS{oplakávaný} \dicExampleIS{vera öllum harmdauði} \dicExampleCS{být všemi oplakávaný}
\dicEntry[harmkvæli] \dicTerm{harm··kvæli} \dicIPA{{h}{a}{r}{m}{k\smash{\textsuperscript{h}}}{v}{a}{i}{l}{\textsci}} \dicPos{n}[2] \dicFlx{pl}[19] \dicSynonym*{miklar þjáningar} \dicDirectTranslationCS{strast, útrapa};  \dicPhraseIS{(þetta tókst) með harmkvælum} \dicDirectTranslationCS{(povedlo se to) s~velkými obtížemi}
\dicEntry[harmleikur] \dicTerm{harm··leik|ur} \dicsymFrequent\  \dicIPA{{h}{a}{r}{m}{l}{ei}{\r{g}}{\textscy}{\textsubring{r}}} \dicPos{m}[9] \dicFlx{(‑s, ‑ir)}[15] \textbf{1.} \dicFieldCat{lit.} \dicSynonym*{sorgarleikrit} \dicDirectTranslationCS{tragédie}  \textbf{2.} \dicSynonym*{hryggilegur atburður} \dicDirectTranslationCS{tragédie, neštěstí} \dicExampleIS{persónulegur harmleikur} \dicExampleCS{osobní tragédie}
\dicEntry[harmoníka] \dicTerm{harmoník|a} \dicIPA{{h}{a}{r}{m}{\textopeno}{n}{i}{\r{g}}{a}} \dicPos{f}[1] \dicFlx{(‑u, ‑ur)}[7] \dicLink{harmónika}
\dicEntry[harmónika] \dicTerm{harmónik|a}\dicTerm{, harmoníka} \dicIPA{{h}\-{a}\-{r}\-{m}\-{ou}\-{n}\-{\textsci}\-{\r{g}}\-{a}\-} \dicPos{f}[1] \dicFlx{(‑u, ‑ur)}[7] \dicFieldCat{hud.} \dicDirectTranslationCS{(tahací) harmonika, akordeon}
\dicEntry[harmsaga] \dicTerm{harm··|saga} \dicIPA{{h}{a}{r}{m}{s}{a}{\textbabygamma}{a}} \dicPos{f}[1] \dicFlx{(‑sögu, ‑sögur)}[14] \dicDirectTranslationCS{tragický příběh, tragédie}
\dicEntry[harmur] \dicTerm{harm|ur} \dicIPA{{h}{a}{r}{m}{\textscy}{\textsubring{r}}} \dicPos{m}[6] \dicFlx{(‑s, ‑ar)}[5] \dicSynonym{sorg} \dicDirectTranslationCS{zármutek, hoře, žal};  \dicPhraseIS{hefna harma sinna} \dicDirectTranslationCS{(po)mstít se}
\dicEntry[harmþrunginn] \dicTerm{harm··þrunginn} \dicIPA{{h}{a}{r}{m}{\texttheta}{r}{u}{\textltailn}{\r{\textObardotlessj}}{\textsci}{\textsubring{n}}} \dicPos{adj}[6]\dicFlx{}[-2] \dicDirectTranslationCS{zarmoucený, nešťastný}
\dicEntry[harpa] \dicTerm{harpa} \dicIPA{{h}{a}{\textsubring{r}}{\textsubring{b}}{a}} \dicPos{f}[1] \dicFlx{(hörpu, hörpur)}[20] \dicFieldCat{hud.} \dicDirectTranslationCS{harfa}
\dicEntry[hart] \dicTerm{hart} \dicsymFrequent\  \dicIPA{{h}{a}{\textsubring{r}}{\textsubring{d}}} \dicPos{adv} \dicFlx{(comp harðar\,/\addthin harðara, sup harðast)} \dicDirectTranslationCS{tvrdě, drsně, nelítostně} \dicExampleIS{berjast hart} \dicExampleCS{nelítostně bojovat};  \dicPhraseIS{ef í hart fer} \dicFlx{adv} \dicDirectTranslationCS{v~nejhorším případě, v~případě nouze}
\dicEntry[hasar] \dicTerm{hasar} \dicIPA{{h}{a}{\textlengthmark}{s}{a}{\textsubring{r}}} \dicPos{m}[4] \dicFlx{(‑s)}[18] \dicDirectTranslationCS{napínavá situace\,/\addthin akce}
\dicEntry[hasarblað] \dicTerm{hasar··|blað} \dicIPA{{h}{a}{s}{a}{r}{\textsubring{b}}{l}{a}{\texttheta}} \dicPos{n}[2] \dicFlx{(‑blaðs, ‑blöð)}[8] \dicDirectTranslationCS{(akční) komiks}
\dicEntry[hasarmynd] \dicTerm{hasar··mynd} \dicIPA{{h}{a}{s}{a}{r}{m}{\textsci}{n}{\textsubring{d}}} \dicPos{f}[7] \dicFlx{(‑ar, ‑ir)}[1] \dicLangCat{hovor.} \dicSynonym{spennumynd} \dicDirectTranslationCS{akční film}
\dicEntry[hasla] \dicTerm{hasl|a} \dicIPA{{h}{a}{s}{\textsubring{d}}{l}{a}} \dicPos{v}[1] \dicFlx{(‑aði)}[43] \dicPhraseIS{hasla sér völl e‑s staðar} \dicDirectTranslationCS{etablovat se (kde)}
\dicEntry[hass] \dicTerm{hass} \dicIPA{{h}{a}{s}{\textlengthmark}} \dicPos{n}[2] \dicFlx{(‑)}[24] \dicDirectTranslationCS{hašiš}
\dicEntry[hassisti] \dicTerm{hass··ist|i} \dicIPA{{h}{a}{s}{\textlengthmark}{\textsci}{s}{\textsubring{d}}{\textsci}} \dicPos{m}[1] \dicFlx{(‑a, ‑ar)}[1] \dicDirectTranslationCS{uživatel(ka) hašiše}
\dicEntry[hast] \dicTerm{hast} \dicIPA{{h}{a}{s}{\textsubring{d}}} \dicPos{n}[2] \dicFlx{(‑s)}[2] \dicDirectTranslationCS{chvat, spěch};  \dicPhraseIS{gera e‑ð í hasti} \dicDirectTranslationCS{dělat (co) ve spěchu}
\dicEntry[hasta] \dicTerm{hast|a} \dicIPA{{h}{a}{s}{\textsubring{d}}{a}} \dicPos{v}[1] \dicFlx{(‑aði)}[50] \dicPhraseIS{hasta sér} \dicDirectTranslationCS{pospíchat};  \dicPhraseIS{hasta á e‑n} \dicDirectTranslationCS{napomenout (koho), udělit důtku (komu)}
\dicEntry[hastur] \dicTerm{hast|ur\smash{\textsuperscript{1}}} \dicIPA{{h}{a}{s}{\textsubring{d}}{\textscy}{\textsubring{r}}} \dicPos{m}[6] \dicFlx{(‑s)}[7] \dicSynonym{hast} \dicDirectTranslationCS{spěch, chvat}
\dicEntry[hastur] \dicTerm{hastur\smash{\textsuperscript{2}}} \dicIPA{{h}{a}{s}{\textsubring{d}}{\textscy}{\textsubring{r}}} \dicPos{adj}[1] \dicFlx{(f höst)}[11] \textbf{1.} \dicDirectTranslationCS{hrubý, drsný, surový}  \textbf{2.} \dicSynonym{hranalegur} \dicDirectTranslationCS{popudlivý, nevrlý}
\dicEntry[hata] \dicTerm{hat|a} \dicsymFrequent\  \dicIPA{{h}{a}{\textlengthmark}{\textsubring{d}}{a}} \dicPos{v}[1] \dicFlx{(‑aði)}[13] \dicFlx{acc} \dicDirectTranslationCS{nenávidět, nesnášet} \dicExampleIS{Þau hötuðu hann.} \dicExampleCS{Nenáviděli ho.};  \dicIdiom{hatast}[við]{ \dicPhraseIS{hatast við e‑n}} \dicFlx{refl} \dicDirectTranslationCS{nenávidět (koho), mít odpor ke (komu)}
\dicEntry[hattar] \dicTerm{hattar} \dicIPA{{h}{a}{h}{\textsubring{d}}{a}{\textsubring{r}}} \dicPos{m} \dicFlx{sg gen} \dicLink{höttur}
\dicEntry[hattur] \dicTerm{hatt|ur} \dicsymFrequent\  \dicIPA{{h}{a}{h}{\textsubring{d}}{\textscy}{\textsubring{r}}} \dicPos{m}[6] \dicFlx{(‑s, ‑ar)}[5] \textbf{1.} \dicDirectTranslationCS{klobouk, klobouček} \dicExampleIS{kvenhattur} \dicExampleCS{dámský klobouk};  \dicPhraseIS{fá skömm í hattinn} \dicLangCat{přen.} \dicDirectTranslationCS{dostat vynadáno};  \dicPhraseIS{góður fyrir sinn hatt} \dicDirectTranslationCS{(o~člověku) jistý sám sebou};  \dicPhraseIS{setja allt undir einn\,/\addthin sama hatt} \dicLangCat{přen.} \dicDirectTranslationCS{dávat vše do jednoho pytle};  \dicPhraseIS{vera á höttunum eftir e‑u} \dicSynonym*{leita eftir e‑u} \dicDirectTranslationCS{ohlížet se po (čem), pátrat po (čem)}  \textbf{2.} \dicFieldCat{bot.} \dicDirectTranslationCS{klobouk (houby)}
\dicEntry[hatur] \dicTerm{hatur} \dicIPA{{h}{a}{\textlengthmark}{\textsubring{d}}{\textscy}{\textsubring{r}}} \dicPos{n}[2] \dicFlx{(‑s)}[28] \dicDirectTranslationCS{nenávist};  \dicPhraseIS{leggja hatur á e‑n} \dicDirectTranslationCS{nenávidět (koho)}
\begin{xtolerant}{}{1pt}
\dicEntry[hatursfullur] \dicTerm{haturs··|fullur}\addthinS\dicIPA{{h}{a}{\textlengthmark}{\textsubring{d}}{\textscy}{\textsubring{r}}{s}{f}{\textscy}{\textsubring{d}}{l}{\textscy}{\textsubring{r}}}\addthinS\dicPos{adj}[10]\addthinS\dicFlx{(comp ‑fyllri, sup ‑fyllstur)}[7] \dicDirectTranslationCS{nenávistný, plný nenávisti, nenávidějící, záštiplný}
\end{xtolerant}
\dicEntry[haugarfi] \dicTerm{haug··arf|i} \dicIPA{{h}{\oe i}{\textlengthmark}{\textbabygamma}{a}{r}{v}{\textsci}} \dicPos{m}[1] \dicFlx{(‑a)}[3] \dicFieldCat{bot.} \dicDirectTranslationCS{ptačinec prostřední} \textit{(l.~{\textLA{Stellaria media}})}  \dicsymPhoto\ 
\dicFigure{ds_image_haugarfi_0_1.jpg}{Haugarfi}{Haugarfi - Kaldari, PD}
\dicEntry[haugur] \dicTerm{haug|ur} \dicIPA{{h}{\oe i}{\textlengthmark}{\textbabygamma}{\textscy}{\textsubring{r}}} \dicPos{m}[6] \dicFlx{(‑s, ‑ar)}[22] \textbf{1.} \dicSynonym*{hrúga} \dicDirectTranslationCS{hromada, halda, kupa} \dicExampleIS{haugur af rusli} \dicExampleCS{hromada smetí}  \textbf{2.} \dicSynonym{legstaður} \dicDirectTranslationCS{mohyla, hrobka}
\dicEntry[haukur] \dicTerm{hauk|ur} \dicIPA{{h}{\oe i}{\textlengthmark}{\r{g}}{\textscy}{\textsubring{r}}} \dicPos{m}[6] \dicFlx{(‑s, ‑ar)}[22] \dicFieldCat{zool.} \dicDirectTranslationCS{jestřáb, krahujec} \textit{(l.~{\textLA{Accipiter}})}  \dicsymPhoto\ ;  \dicPhraseIS{eiga hauk í horni} \dicLangCat{přen.} \dicDirectTranslationCS{mít eso v~rukávu}
\dicFigure{117602.jpg}{Haukur}{Haukur - Došlý Martin, Biolib, Copyright/PD}
\dicEntry[haus] \dicTerm{haus} \dicsymFrequent\  \dicIPA{{h}{\oe i}{\textlengthmark}{s}} \dicPos{m}[4] \dicFlx{(‑s, ‑ar)}[14] \textbf{1.} \dicSynonym{höfuð\smash{\textsuperscript{1}}} \dicDirectTranslationCS{hlava (lidská ap.)}  \textbf{2.} \dicDirectTranslationCS{hlava, hlavička (na hřebíku ap.)} \dicExampleIS{hausinn á naglanum} \dicExampleCS{hlavička hřebíku}  \textbf{3.} \dicSynonym{fyrirsögn} \dicDirectTranslationCS{hlavička, záhlaví};  \dicIdiom{haus}{ \dicPhraseIS{fara á hausinn}} \dicLangCat{přen.} \dicDirectTranslationCS{zbankrotovat, zkrachovat};  \dicPhraseIS{hengja haus} \dicLangCat{přen.} \dicDirectTranslationCS{věšet hlavu};  \dicPhraseIS{snúa öllu á haus} \dicLangCat{přen.} \dicDirectTranslationCS{postavit vše na hlavu};  \dicPhraseIS{standa á haus í e‑u} \dicLangCat{přen.} \dicDirectTranslationCS{být uprostřed (čeho) (činnosti ap.)};  \dicPhraseIS{vera á hausnum} \dicLangCat{přen.} \dicDirectTranslationCS{zkrachovat, zbankrotovat}
\dicEntry[hauskúpa] \dicTerm{haus··kúp|a} \dicIPA{{h}{\oe i}{\textlengthmark}{s}{k\smash{\textsuperscript{h}}}{u}{\textsubring{b}}{a}} \dicPos{f}[1] \dicFlx{(‑u, ‑ur)}[19] \dicFieldCat{anat.} \dicDirectTranslationCS{lebka}
\dicEntry[haust] \dicTerm{haust} \dicsymFrequent\  \dicIPA{{h}{\oe i}{s}{\textsubring{d}}} \dicPos{n}[2] \dicFlx{(‑s, ‑)}[5] \dicDirectTranslationCS{podzim} \dicExampleIS{á haustin} \dicExampleCS{na podzim}
\dicEntry[hausta] \dicTerm{haust|a} \dicIPA{{h}{\oe i}{s}{\textsubring{d}}{a}} \dicPos{v}[1] \dicFlx{(‑aði)}[83] \dicFlx{impers} \dicPhraseIS{það haustar} \dicDirectTranslationCS{přichází podzim}
\dicEntry[haustdagur] \dicTerm{haust··dag|ur} \dicIPA{{h}{\oe i}{s}{\textsubring{d}}{\textlengthmark}{a}{\textbabygamma}{\textscy}{\textsubring{r}}} \dicPos{m}[6] \dicFlx{(‑s, ‑ar)}[62] \dicDirectTranslationCS{podzimní den};  \dicPhraseIS{á haustdögum} \dicFlx{adv} \dicDirectTranslationCS{na podzim}
\dicEntry[haustjafndægur] \dicTerm{haust··jafn·dægur} \dicIPA{{h}{\oe i}{s}{\textsubring{d}}{j}{a}{\textsubring{b}}{\textsubring{n}}{\textsubring{d}}{a}{i}{\textbabygamma}{\textscy}{\textsubring{r}}} \dicPos{n}[2] \dicFlx{(‑s, ‑)}[25] \dicDirectTranslationCS{podzimní rovnodennost} \dicAntonym{vorjafndægur}
\dicEntry[haustlegur] \dicTerm{haust··legur} \dicIPA{{h}{\oe i}{s}{\textsubring{d}}{l}{\textepsilon}{\textbabygamma}{\textscy}{\textsubring{r}}} \dicPos{adj}[1]\dicFlx{}[-8] \dicDirectTranslationCS{podzimní, připomínající podzim}
\dicEntry[haustmisseri] \dicTerm{haust··misseri} \dicIPA{{h}{\oe i}{s}{\textsubring{d}}{m}{\textsci}{s}{\textepsilon}{r}{\textsci}} \dicPos{n}[2] \dicFlx{(‑s, ‑)}[14] \dicFieldCat{škol.} \dicDirectTranslationCS{podzimní\,/\addthin zimní semestr} \dicAntonym{vormisseri}
\dicEntry[haustönn] \dicTerm{haust··|önn} \dicIPA{{h}{\oe i}{s}{\textsubring{d}}{\oe}{\textsubring{n}}} \dicPos{f}[7] \dicFlx{(‑annar, ‑annir)}[16] \dicFieldCat{škol.} \dicDirectTranslationCS{podzimní\,/\addthin zimní semestr} \dicAntonym{vorönn}
\dicEntry[hausverkur] \dicTerm{haus··verk|ur} \dicIPA{{h}{\oe i}{\textlengthmark}{s}{v}{\textepsilon}{\textsubring{r}}{\r{g}}{\textscy}{\textsubring{r}}} \dicPos{m}[10] \dicFlx{(‑jar, ‑ir)}[23] \dicDirectTranslationCS{bolest hlavy}
\dicEntry[há] \dicTerm{há} \dicIPA{{h}{au}{\textlengthmark}} \dicPos{v}[2] \dicFlx{(‑ði, ‑ð)}[123] \dicFlx{dat} \dicDirectTranslationCS{znevýhodňovat} \dicExampleIS{e‑að háir e‑m} \dicExampleCS{(co koho) znevýhodňuje}
\dicEntry[háaloft] \dicTerm{háa··loft} \dicIPA{{h}{au}{\textlengthmark}{a}{l}{\textopeno}{f}{\textsubring{d}}} \dicPos{n}[2] \dicFlx{(‑s, ‑)}[5] \dicDirectTranslationCS{podkroví};  \dicPhraseIS{það er allt í háalofti} \dicLangCat{přen.} \dicDirectTranslationCS{všechno je vzhůru nohama}
\dicEntry[hádegi] \dicTerm{há··degi} \dicsymFrequent\  \dicIPA{{h}{au}{\textlengthmark}{\textsubring{d}}{ei}{\textsci}} \dicPos{n}[2] \dicFlx{(‑s, ‑)}[14] \dicDirectTranslationCS{poledne} \dicExampleIS{Borðum saman í hádeginu.} \dicExampleCS{Najezme se spolu v~poledne.};  \dicPhraseIS{eftir hádegi} \dicFlx{adv} \dicDirectTranslationCS{odpoledne};  \dicPhraseIS{fyrir hádegi} \dicFlx{adv} \dicDirectTranslationCS{dopoledne}
\dicEntry[hádegisbaugur] \dicTerm{há·degis··baug|ur} \dicIPA{{h}{au}{\textlengthmark}{\textsubring{d}}{ei}{\textsci}{s}{\textsubring{b}}{\oe i}{\textbabygamma}{\textscy}{\textsubring{r}}} \dicPos{m}[6] \dicFlx{(‑s, ‑ar)}[24] \dicFieldCat{geog.} \dicDirectTranslationCS{poledník}
\dicEntry[hádegismatur] \dicTerm{há·degis··mat|ur} \dicIPA{{h}{au}{\textlengthmark}{\textsubring{d}}{ei}{\textsci}{s}{m}{a}{\textsubring{d}}{\textscy}{\textsubring{r}}} \dicPos{m}[10] \dicFlx{(‑ar)}[15] \dicDirectTranslationCS{oběd}
\dicEntry[hádegisverður] \dicTerm{há·degis··verð|ur} \dicIPA{{h}{au}{\textlengthmark}{\textsubring{d}}{ei}{\textsci}{s}{v}{\textepsilon}{r}{ð}{\textscy}{\textsubring{r}}} \dicPos{m}[10] \dicFlx{(‑ar, ‑ir)}[17] \dicDirectTranslationCS{oběd}
\dicEntry[háð] \dicTerm{háð\smash{\textsuperscript{1}}} \dicIPA{{h}{au}{\textlengthmark}{\texttheta}} \dicPos{n}[2] \dicFlx{(‑s)}[2] \dicSynonym{spé} \dicDirectTranslationCS{ironie, posměch, výsměch} \dicExampleIS{segja e‑ð í háði} \dicExampleCS{říct (co) výsměšně}
\dicEntry[háð] \dicTerm{háð\smash{\textsuperscript{2}}} \dicIPA{{h}{au}{\textlengthmark}{\texttheta}} \dicPos{v} \dicFlx{supin} \dicLink{heyja\smash{\textsuperscript{2}}}
\dicEntry[háði] \dicTerm{háði} \dicIPA{{h}{au}{\textlengthmark}{ð}{\textsci}} \dicPos{v} \dicFlx{ind\,/\addthin con pf sg 1 pers} \dicLink{heyja\smash{\textsuperscript{2}}}
\dicEntry[háðslegur] \dicTerm{háðs··legur} \dicIPA{{h}{au}{ð}{s}{l}{\textepsilon}{\textbabygamma}{\textscy}{\textsubring{r}}} \dicPos{adj}[1]\dicFlx{}[-8] \dicDirectTranslationCS{posměšný, výsměšný}
\dicEntry[háðum] \dicTerm{háðum} \dicIPA{{h}{au}{\textlengthmark}{ð}{\textscy}{\textsubring{m}}} \dicPos{v} \dicFlx{ind pf pl 1 pers} \dicLink{heyja\smash{\textsuperscript{2}}}
\dicEntry[háðung] \dicTerm{háð··ung} \dicIPA{{h}{au}{\textlengthmark}{ð}{u}{\ng}{\r{g}}} \dicPos{f}[4] \dicFlx{(‑ar, ‑ar)}[1] \dicSynonym{skömm\smash{\textsuperscript{1}}} \dicDirectTranslationCS{hanba, ostuda, potupa}
\dicEntry[háður] \dicTerm{háður} \dicsymFrequent\  \dicIPA{{h}{au}{\textlengthmark}{ð}{\textscy}{\textsubring{r}}} \dicPos{adj}[2]\dicFlx{}[-12] \dicDirectTranslationCS{závislý, odkázaný};  \dicPhraseIS{vera háður e‑u} \dicDirectTranslationCS{být závislý na (čem)} \dicExampleIS{vera háður fíkniefnum} \dicExampleCS{být závislý na drogách}
\dicEntry[háfleygur] \dicTerm{há··fleygur} \dicIPA{{h}{au}{\textlengthmark}{f}{l}{ei}{\textbabygamma}{\textscy}{\textsubring{r}}} \dicPos{adj}[1]\dicFlx{}[-1] \textbf{1.} \dicSynonym*{sem flýgur hátt} \dicDirectTranslationCS{letící vysoko}  \textbf{2.} \dicSynonym*{rismikill} \dicDirectTranslationCS{pompézní, vzletný}
\dicEntry[háflæði] \dicTerm{há··flæði} \dicIPA{{h}{au}{\textlengthmark}{f}{l}{a}{i}{ð}{\textsci}} \dicPos{n}[2] \dicFlx{(‑s)}[20] \dicDirectTranslationCS{vysoký příliv}
\dicEntry[háfur] \dicTerm{háf|ur} \dicIPA{{h}{au}{v}{\textlengthmark}{\textscy}{\textsubring{r}}} \dicPos{m}[6] \dicFlx{(‑s, ‑ar)}[22] \textbf{1.} \dicDirectTranslationCS{podběrák}  \textbf{2.} \dicSynonym{vifta} \dicDirectTranslationCS{digestoř}  \textbf{3.} \dicFieldCat{zool.} \dicDirectTranslationCS{ostroun obecný} \textit{(l.~{\textLA{Squalus acanthias}})}  \dicsymPhoto\ 
\dicFigure{ds_image_hafur_0_1.jpg}{Háfur}{Háfur - Sanc1698 - Flickr - NOAA Photo Library, PD}
\dicEntry[háhestur] \dicTerm{há··hest|ur} \dicIPA{{h}{au}{\textlengthmark}{h}{\textepsilon}{s}{\textsubring{d}}{\textscy}{\textsubring{r}}} \dicPos{m}[6] \dicFlx{(‑s)}[7] \dicPhraseIS{bera e‑n á háhesti} \dicDirectTranslationCS{nést (koho) na koníčka, nést (koho) na ramenech}
\dicEntry[háhyrningur] \dicTerm{há··hyrn·ing|ur} \dicIPA{{h}{au}{\textlengthmark}{h}{\textsci}{r}{\textsubring{d}}{n}{i}{\ng}{\r{g}}{\textscy}{\textsubring{r}}} \dicPos{m}[6] \dicFlx{(‑s, ‑ar)}[8] \dicFieldCat{zool.} \dicDirectTranslationCS{kosatka dravá} \textit{(l.~{\textLA{Orcinus orca}})}  \dicsymPhoto\ 
\dicFigure{71471.jpg}{Háhyrningur}{Háhyrningur - ?, Biolib, PD}
\dicEntry[háhýsi] \dicTerm{há··hýsi} \dicIPA{{h}{au}{\textlengthmark}{h}{i}{s}{\textsci}} \dicPos{n}[2] \dicFlx{(‑s, ‑)}[14] \dicDirectTranslationCS{výšková budova, věžák, mrakodrap}
\dicEntry[hákarl] \dicTerm{há··karl} \dicIPA{{h}{au}{\textlengthmark}{k\smash{\textsuperscript{h}}}{a}{r}{\textsubring{d}}{\textsubring{l}}} \dicPos{m}[4] \dicFlx{(‑s, ‑ar)}[4] \dicFieldCat{zool.} \dicDirectTranslationCS{žralok, žralok malohlavý} \textit{(l.~{\textLA{Somniosus microcephalus}})}
\dicEntry[hákveða] \dicTerm{há··kveð|a} \dicIPA{{h}{au}{\textlengthmark}{k\smash{\textsuperscript{h}}}{v}{\textepsilon}{ð}{a}} \dicPos{f}[1] \dicFlx{(‑u, ‑ur)}[19] \dicFieldCat{lit.} \dicDirectTranslationCS{lichá stopa s~větším přízvukem} \dicAntonym{lágkveða}
\dicEntry[hálaunamaður] \dicTerm{há·launa··|maður} \dicIPA{{h}{au}{\textlengthmark}{l}{\oe i}{n}{a}{m}{a}{ð}{\textscy}{\textsubring{r}}} \dicPos{m}[13] \dicFlx{(‑manns, ‑menn)}[2] \dicDirectTranslationCS{člověk\,/\addthin osoba s~vysokým platem}
\dicEntry[háleitur] \dicTerm{há··leitur} \dicIPA{{h}{au}{\textlengthmark}{l}{ei}{\textsubring{d}}{\textscy}{\textsubring{r}}} \dicPos{adj}[1]\dicFlx{}[-1] \dicDirectTranslationCS{vzletný, vznosný} \dicExampleIS{háleitt kvæði} \dicExampleCS{vzletná báseň}
\dicEntry[hálendi] \dicTerm{há··lendi} \dicIPA{{h}{au}{\textlengthmark}{l}{\textepsilon}{n}{\textsubring{d}}{\textsci}} \dicPos{n}[2] \dicFlx{(‑s)}[20] \dicDirectTranslationCS{vysočina} \dicAntonym{láglendi}
\dicEntry[hálf] \dicTerm{hálf-} \dicIPA{{h}{au}{l}{f}} \dicPos{predp} \textbf{1.} \dicSynonym*{helmings-} \dicDirectTranslationCS{polo-, poloviční} \dicExampleIS{hálfsannleikur} \dicExampleCS{polopravda}  \textbf{2.} \dicSynonym{dálítið} \dicDirectTranslationCS{trochu, napolo} \dicExampleIS{hálfslappur} \dicExampleCS{trochu zesláblý}
\dicEntry[hálfa] \dicTerm{hálf|a} \dicsymFrequent\  \dicIPA{{h}{au}{l}{v}{a}} \dicPos{f}[1] \dicFlx{(‑u, ‑ur)}[19] \dicPhraseIS{af hálfu e‑rs} \dicFlx{prep} \dicDirectTranslationCS{jménem (koho), za (koho)};  \dicPhraseIS{af minni hálfu} \dicFlx{adv} \dicDirectTranslationCS{co se mě týče, z~mé strany}
\dicEntry[hálfbróðir] \dicTerm{hálf··|bróðir} \dicIPA{{h}{au}{l}{v}{\textsubring{b}}{r}{ou}{ð}{\textsci}{\textsubring{r}}} \dicPos{m}[13] \dicFlx{(‑bróður, ‑bræður)}[4] \dicDirectTranslationCS{nevlastní bratr}
\dicEntry[hálffimmtugur] \dicTerm{hálf··fimm·tugur} \dicIPA{{h}{au}{l}{f}{\textsci}{m}{t\smash{\textsuperscript{h}}}{\textscy}{\textbabygamma}{\textscy}{\textsubring{r}}} \dicPos{adj}[1]\dicFlx{}[-6] \dicDirectTranslationCS{čtyřicetipětiletý}
\dicEntry[hálffjarlægur] \dicTerm{hálf··fjar·lægur} \dicIPA{{h}{au}{l}{f}{j}{a}{\textsubring{r}}{l}{a}{i}{\textbabygamma}{\textscy}{\textsubring{r}}} \dicPos{adj}[1]\dicFlx{}[-1] \dicFieldCat{jaz.} \dicDirectTranslationCS{polootevřený (samohláska ap.)}
\dicEntry[hálfgerður] \dicTerm{hálf··gerður} \dicsymFrequent\  \dicIPA{{h}{au}{l}{v}{\r{\textObardotlessj}}{\textepsilon}{r}{ð}{\textscy}{\textsubring{r}}} \dicPos{adj}[2]\dicFlx{}[-4] \dicDirectTranslationCS{poloviční, částečný, polovičatý} \dicExampleIS{hálfgert flón} \dicExampleCS{poloviční hlupák}
\dicEntry[hálfkák] \dicTerm{hálf··kák} \dicIPA{{h}{au}{l}{v}{k\smash{\textsuperscript{h}}}{au}{\r{g}}} \dicPos{n}[2] \dicFlx{(‑s)}[2] \dicDirectTranslationCS{nedodělaná práce\,/\addthin věc}
\dicEntry[hálfkveðinn] \dicTerm{hálf··kveðinn} \dicIPA{{h}{au}{l}{v}{k\smash{\textsuperscript{h}}}{v}{\textepsilon}{ð}{\textsci}{\textsubring{n}}} \dicPos{adj}[6]\dicFlx{}[-6] \dicPhraseIS{tala í hálfkveðnum vísum} \dicDirectTranslationCS{mluvit v~náznacích\,/\addthin narážkách}
\dicEntry[hálfleikur] \dicTerm{hálf··leik|ur} \dicIPA{{h}{au}{l}{v}{l}{ei}{\r{g}}{\textscy}{\textsubring{r}}} \dicPos{m}[9] \dicFlx{(‑s, ‑ir)}[15] \dicFieldCat{sport.} \dicDirectTranslationCS{poločas (utkání ap.)} \dicExampleIS{í hálfleik} \dicExampleCS{v~poločase};  \dicPhraseIS{fyrri hálfleikur } \dicDirectTranslationCS{první poločas};  \dicPhraseIS{seinni\,/\addthin síðari hálfleikur} \dicDirectTranslationCS{druhý poločas}
\dicEntry[hálfmáni] \dicTerm{hálf··mán|i} \dicIPA{{h}{au}{l}{v}{m}{au}{n}{\textsci}} \dicPos{m}[1] \dicFlx{(‑a, ‑ar)}[1] \dicDirectTranslationCS{půlměsíc, srpek}
\dicEntry[hálfnaður] \dicTerm{hálf··n|aður} \dicIPA{{h}{au}{l}{n}{a}{ð}{\textscy}{\textsubring{r}}} \dicPos{adj}[3] \dicFlx{(f ‑uð)}[4] \dicDirectTranslationCS{(jsoucí) v~polovině, z~půlky hotový};  \dicProverb\  \dicPhraseIS{Hálfnað er verk sem hafið er.} \dicLangCat{přís.} \dicDirectTranslationCS{S~chutí do toho a~půl je hotovo.}
\dicEntry[hálfnálægur] \dicTerm{hálf··ná·lægur} \dicIPA{{h}{au}{l}{v}{n}{au}{l}{a}{i}{\textbabygamma}{\textscy}{\textsubring{r}}} \dicPos{adj}[1]\dicFlx{}[-1] \dicFieldCat{jaz.} \dicDirectTranslationCS{polozavřený (samohláska ap.)}
\dicEntry[hálfrím] \dicTerm{hálf··rím} \dicIPA{{h}{au}{l}{v}{r}{i}{\textsubring{m}}} \dicPos{n}[2] \dicFlx{(‑s)}[2] \dicFieldCat{lit.} \dicIndirectTranslationCS{vnitřní rým, ve kterém se rýmují buď pouze souhlásky, nebo pouze samohlásky (např. fugl : gull)} \dicAntonym{alrím}
\dicEntry[hálfskýjaður] \dicTerm{hálf··skýj·|aður} \dicIPA{{h}{au}{l}{v}{s}{\r{\textObardotlessj}}{i}{j}{a}{ð}{\textscy}{\textsubring{r}}} \dicPos{adj}[3] \dicFlx{(f ‑uð)}[4] \dicDirectTranslationCS{zatažený} \dicExampleIS{Það er hálfskýjað.} \dicExampleCS{Je zataženo.}
\dicEntry[hálfsystir] \dicTerm{hálf··syst|ir} \dicIPA{{h}{au}{l}{v}{s}{\textsci}{s}{\textsubring{d}}{\textsci}{\textsubring{r}}} \dicPos{f}[11] \dicFlx{(‑ur, ‑ur)}[3] \dicDirectTranslationCS{nevlastní sestra}
\dicEntry[hálftími] \dicTerm{hálf··tím|i} \dicsymFrequent\  \dicIPA{{h}{au}{l}{v}{t\smash{\textsuperscript{h}}}{i}{m}{\textsci}} \dicPos{m}[1] \dicFlx{(‑a, ‑ar)}[1] \dicDirectTranslationCS{půlhodina} \dicExampleIS{Hann kemur heim hálftíma áður en gestirnir koma.} \dicExampleCS{Přijde domů půlhodinu před příchodem hostů.}
\dicEntry[hálfur] \dicTerm{hálfur} \dicsymFrequent\  \dicIPA{{h}{au}{l}{v}{\textscy}{\textsubring{r}}} \dicPos{adj}[1]\dicFlx{}[-6] \textbf{1.} \dicDirectTranslationCS{poloviční} \dicExampleIS{Hann kemur eftir hálfan mánuð.} \dicExampleCS{Přijede za půl měsíce.}  \textbf{2.} \dicSynonym*{dálítið drukkinn} \dicDirectTranslationCS{mírně opilý};  \dicIdiom{hálfur}{ \dicPhraseIS{hálfur annar klukkutími}} \dicDirectTranslationCS{půldruhá hodina};  \dicPhraseIS{tala í hálfum hljóðum} \dicDirectTranslationCS{mluvit polohlasem};  \dicPhraseIS{til hálfs} \dicFlx{adv} \dicDirectTranslationCS{do půlky}
\dicEntry[hálfvirði] \dicTerm{hálf··virði} \dicIPA{{h}{au}{l}{v}{\textsci}{r}{ð}{\textsci}} \dicPos{n}[2] \dicFlx{(‑s)}[20] \dicDirectTranslationCS{poloviční cena} \dicExampleIS{á hálfvirði} \dicExampleCS{za poloviční cenu}
\dicEntry[hálfviti] \dicTerm{hálf··vit|i} \dicIPA{{h}{au}{l}{v}{\textsci}{\textsubring{d}}{\textsci}} \dicPos{m}[1] \dicFlx{(‑a, ‑ar)}[1] \dicLangCat{han.} \dicDirectTranslationCS{pitomec, idiot, kretén}
\dicEntry[hálfvolgur] \dicTerm{hálf··volgur} \dicIPA{{h}{au}{l}{v}{\textopeno}{l}{\r{g}}{\textscy}{\textsubring{r}}} \dicPos{adj}[1]\dicFlx{}[-6] \textbf{1.} \dicSynonym{hikandi} \dicDirectTranslationCS{váhavý, nerozhodný}  \textbf{2.} \dicSynonym*{ekki vel heitur} \dicDirectTranslationCS{vlažný, vlahý (voda ap.)}
\dicEntry[hálfþrítugur] \dicTerm{hálf··þrí·tugur} \dicIPA{{h}{au}{l}{v}{\texttheta}{r}{i}{\textsubring{d}}{\textscy}{\textbabygamma}{\textscy}{\textsubring{r}}} \dicPos{adj}[1]\dicFlx{}[-6] \dicDirectTranslationCS{pětadvacetiletý, dvacetipětiletý}
\dicEntry[hálka] \dicTerm{hálk|a} \dicIPA{{h}{au}{\textsubring{l}}{\r{g}}{a}} \dicPos{f}[1] \dicFlx{(‑u)}[5] \dicDirectTranslationCS{náledí, ledovka}
\dicEntry[háll] \dicTerm{háll} \dicIPA{{h}{au}{\textsubring{d}}{\textsubring{l}}} \dicPos{adj}[8]\dicFlx{}[-1] \dicDirectTranslationCS{kluzký, klouzavý}
\dicEntry[hálmstrá] \dicTerm{hálm··strá} \dicIPA{{h}{au}{l}{m}{s}{\textsubring{d}}{r}{au}} \dicPos{n}[2] \dicFlx{(‑s, ‑)}[5] \dicDirectTranslationCS{stéblo};  \dicPhraseIS{grípa í hálm\-stráið} \dicLangCat{přen.} \dicDirectTranslationCS{chytat se stébla}
\dicEntry[hálmur] \dicTerm{hálm|ur} \dicIPA{{h}{au}{l}{m}{\textscy}{\textsubring{r}}} \dicPos{m}[6] \dicFlx{(‑s)}[7] \dicDirectTranslationCS{sláma}
\dicEntry[háls] \dicTerm{háls} \dicsymFrequent\  \dicIPA{{h}{au}{l}{s}} \dicPos{m}[4] \dicFlx{(‑, ‑ar)}[28] \textbf{1.} \dicFieldCat{anat.} \dicDirectTranslationCS{krk} \dicExampleIS{Hjartað virtist slá uppi í hálsi mínum.} \dicExampleCS{Zdálo se, že mi srdce bije až v~krku.}  \textbf{2.} \dicDirectTranslationCS{hrdlo (láhve ap.)}  \textbf{3.} \dicSynonym{hæð} \dicDirectTranslationCS{vyvýšenina, vrch}  \textbf{4.} \dicFieldCat{hud.} \dicDirectTranslationCS{krk (kytary ap.)};  \dicIdiom{háls}{ \dicPhraseIS{góðir hálsar!}} \dicDirectTranslationCS{vážení!};  \dicPhraseIS{með öndina í hálsinum} \dicLangCat{přen.} \dicDirectTranslationCS{napnutý až k~prasknutí}; { \dicPhraseIS{liggja e‑m á hálsi fyrir e‑ð}} \dicDirectTranslationCS{jít (komu) po krku kvůli (čemu)};  \dicPhraseIS{syngja fullum hálsi} \dicDirectTranslationCS{zpívat z~plných plic};  \dicPhraseIS{vera búinn að fá upp í háls af e‑u} \dicLangCat{přen.} \dicDirectTranslationCS{mít (čeho) (až) po krk}
\dicEntry[hálsband] \dicTerm{háls··|band} \dicIPA{{h}{au}{l}{s}{\textsubring{b}}{a}{n}{\textsubring{d}}} \dicPos{n}[2] \dicFlx{(‑bands, ‑bönd)}[8] \dicDirectTranslationCS{obojek}
\dicEntry[hálsbindi] \dicTerm{háls··bindi} \dicIPA{{h}{au}{l}{s}{\textsubring{b}}{\textsci}{n}{\textsubring{d}}{\textsci}} \dicPos{n}[2] \dicFlx{(‑s, ‑)}[14] \dicDirectTranslationCS{vázanka, kravata}
\dicEntry[hálsbólga] \dicTerm{háls··bólg|a} \dicIPA{{h}{au}{l}{s}{\textsubring{b}}{ou}{l}{\r{g}}{a}} \dicPos{f}[1] \dicFlx{(‑u)}[5] \dicFieldCat{med.} \dicDirectTranslationCS{bolest v~krku, angína}
\dicEntry[hálsbraut] \dicTerm{háls··braut} \dicIPA{{h}{au}{l}{s}{\textsubring{b}}{r}{\oe i}{\textsubring{d}}} \dicPos{v} \dicFlx{ind pf sg 1 pers} \dicLink{hálsbrjóta}
\dicEntry[hálsbrjóta] \dicTerm{háls··|brjóta} \dicIPA{{h}{au}{l}{s}{\textsubring{b}}{r}{j}{ou}{\textsubring{d}}{a}} \dicPos{v}[6] \dicFlx{(‑brýt, ‑braut, ‑brutum, ‑bryti, ‑brotið)}[51] \dicFlx{acc} \dicDirectTranslationCS{zlomit vaz} \dicExampleIS{hálsbrjóta sig} \dicExampleCS{zlomit si vaz}
\dicEntry[hálsbrotið] \dicTerm{háls··brotið} \dicIPA{{h}{au}{l}{s}{\textsubring{b}}{r}{\textopeno}{\textsubring{d}}{\textsci}{\texttheta}} \dicPos{v} \dicFlx{supin} \dicLink{hálsbrjóta}
\dicEntry[hálsbrutum] \dicTerm{háls··brutum} \dicIPA{{h}{au}{l}{s}{\textsubring{b}}{r}{\textscy}{\textsubring{d}}{\textscy}{\textsubring{m}}} \dicPos{v} \dicFlx{ind pf pl 1 pers} \dicLink{hálsbrjóta}
\dicEntry[hálsbryti] \dicTerm{háls··bryti} \dicIPA{{h}{au}{l}{s}{\textsubring{b}}{r}{\textsci}{\textsubring{d}}{\textsci}} \dicPos{v} \dicFlx{con pf sg 1 pers} \dicLink{hálsbrjóta}
\dicEntry[hálsbrýt] \dicTerm{háls··brýt} \dicIPA{{h}{au}{l}{s}{\textsubring{b}}{r}{i}{\textsubring{d}}} \dicPos{v} \dicFlx{ind praes sg 1 pers} \dicLink{hálsbrjóta}
\dicEntry[hálsfesti] \dicTerm{háls··fest|i} \dicIPA{{h}{au}{l}{s}{f}{\textepsilon}{s}{\textsubring{d}}{\textsci}} \dicPos{f}[5] \dicFlx{(‑ar, ‑ar)}[3] \dicDirectTranslationCS{přívěsek (na řetízek), řetízek, náhrdelník}
\dicEntry[hálshegg] \dicTerm{háls··hegg} \dicIPA{{h}{au}{l}{s}{h}{\textepsilon}{\r{g}}} \dicPos{v} \dicFlx{ind praes sg 1 pers} \dicLink{hálshöggva}
\dicEntry[hálshjó] \dicTerm{háls··hjó} \dicIPA{{h}{au}{l}{s}{\c{c}}{ou}} \dicPos{v} \dicFlx{ind pf sg 1 pers} \dicLink{hálshöggva}
\dicEntry[hálshjuggum] \dicTerm{háls··hjuggum} \dicIPA{{h}{au}{l}{s}{\c{c}}{\textscy}{\r{g}}{\textscy}{\textsubring{m}}} \dicPos{v} \dicFlx{ind pf pl 1 pers} \dicLink{hálshöggva}
\dicEntry[hálshjyggi] \dicTerm{háls··hjyggi} \dicIPA{{h}{au}{l}{s}{\c{c}}{\textsci}{\r{\textObardotlessj}}{\textsci}} \dicPos{v} \dicFlx{con pf sg 1 pers} \dicLink{hálshöggva}
\dicEntry[hálshoggið] \dicTerm{háls··hoggið} \dicIPA{{h}{au}{l}{s}{h}{\textopeno}{\r{\textObardotlessj}}{\textsci}{\texttheta}} \dicPos{v} \dicFlx{supin} \dicLink{hálshöggva}
\dicEntry[hálshyggi] \dicTerm{háls··hyggi} \dicIPA{{h}{au}{l}{s}{h}{\textsci}{\r{\textObardotlessj}}{\textsci}} \dicPos{v} \dicFlx{con pf sg 1 pers} \dicLink{hálshöggva}
\dicEntry[hálshöggið] \dicTerm{háls··höggið} \dicIPA{{h}{au}{l}{s}{h}{\oe}{\r{\textObardotlessj}}{\textsci}{\texttheta}} \dicPos{v} \dicFlx{supin} \dicLink{hálshöggva}
\dicEntry[hálshöggva] \dicTerm{háls··|höggva} \dicIPA{{h}{au}{l}{s}{h}{\oe}{\r{g}}{v}{a}} \dicPos{v}[6] \dicFlx{(‑hegg, ‑hjó, ‑hjuggum, ‑hyggi\,/\addthin ‑hjyggi, ‑höggvið\,/\addthin ‑höggið\,/\addthin ‑hoggið)}[43] \dicFlx{acc} \dicDirectTranslationCS{stít\,/\addthin setnout hlavu}
\dicEntry[hálshöggvið] \dicTerm{háls··höggvið} \dicIPA{{h}{au}{l}{s}{h}{\oe}{\r{g}}{v}{\textsci}{\texttheta}} \dicPos{v} \dicFlx{supin} \dicLink{hálshöggva}
\dicEntry[hálskirtill] \dicTerm{háls··kirt|ill} \dicIPA{{h}{au}{l}{s}{c\smash{\textsuperscript{h}}}{\textsci}{\textsubring{r}}{\textsubring{d}}{\textsci}{\textsubring{d}}{\textsubring{l}}} \dicPos{m}[6] \dicFlx{(‑ils, ‑lar)}[35] \dicFieldCat{anat.} \dicDirectTranslationCS{(krční) mandle}
\dicEntry[hálsklútur] \dicTerm{háls··klút|ur} \dicIPA{{h}{au}{l}{s}{k\smash{\textsuperscript{h}}}{l}{u}{\textsubring{d}}{\textscy}{\textsubring{r}}} \dicPos{m}[6] \dicFlx{(‑s, ‑ar)}[22] \textbf{1.} \dicDirectTranslationCS{šátek, šála}  \textbf{2.} \dicSynonym{bindi} \dicDirectTranslationCS{kravata}
\dicEntry[hálsmál] \dicTerm{háls··mál} \dicIPA{{h}{au}{l}{s}{m}{au}{\textsubring{l}}} \dicPos{n}[2] \dicFlx{(‑s, ‑)}[5] \dicDirectTranslationCS{výstřih}
\dicEntry[hálsmen] \dicTerm{háls··men} \dicIPA{{h}{au}{l}{s}{m}{\textepsilon}{\textsubring{n}}} \dicPos{n}[2] \dicFlx{(‑s, ‑)}[5] \dicDirectTranslationCS{řetízek (na krk), náhrdelník}
\dicEntry[hálsól] \dicTerm{háls··ól} \dicIPA{{h}{au}{l}{s}{ou}{\textsubring{l}}} \dicPos{f}[4] \dicFlx{(‑ar, ‑ar)}[1] \textbf{1.} \dicDirectTranslationCS{řemínek (na krk) (fotoaparátu ap.)}  \textbf{2.} \dicDirectTranslationCS{obojek}
\dicEntry[hálstau] \dicTerm{háls··tau} \dicIPA{{h}{au}{l}{s}{t\smash{\textsuperscript{h}}}{\oe i}} \dicPos{n}[2] \dicFlx{(‑s, ‑)}[5] \dicDirectTranslationCS{vázanka, kravata, motýlek}
\dicEntry[háma] \dicTerm{hám|a} \dicIPA{{h}{au}{\textlengthmark}{m}{a}} \dicPos{v}[1] \dicFlx{(‑aði)}[44] \dicPhraseIS{háma í sig e‑ð} \dicDirectTranslationCS{hltat (co), cpát se (čím)}
\dicEntry[hámark] \dicTerm{há··|mark} \dicsymFrequent\  \dicIPA{{h}{au}{\textlengthmark}{m}{a}{\textsubring{r}}{\r{g}}} \dicPos{n}[2] \dicFlx{(‑marks, ‑mörk)}[8] \dicDirectTranslationCS{vrchol, kulminace, maximum} \dicExampleIS{hámark hamingjunnar} \dicExampleCS{vrchol štěstí} \dicAntonym{lágmark};  \dicPhraseIS{e‑að er í hámarki} \dicLangCat{přen.} \dicDirectTranslationCS{(co) je na vrcholu};  \dicPhraseIS{e‑að nær hámarki} \dicLangCat{přen.} \dicDirectTranslationCS{(co) dosahuje vrcholu, (co) kulminuje}
\dicEntry[hámarkshraði] \dicTerm{há·marks··hrað|i} \dicIPA{{h}{au}{\textlengthmark}{m}{a}{\textsubring{r}}{x}{s}{\textsubring{r}}{a}{ð}{\textsci}} \dicPos{m}[1] \dicFlx{(‑a)}[3] \textbf{1.} \dicDirectTranslationCS{nejvyšší (povolená) rychlost}  \textbf{2.} \dicDirectTranslationCS{maximální rychlost}
\dicEntry[hámæli] \dicTerm{há··mæli} \dicIPA{{h}{au}{\textlengthmark}{m}{a}{i}{l}{\textsci}} \dicPos{n}[2] \dicFlx{pl}[19] \dicPhraseIS{e‑að kemst í hámæli} \dicFlx{refl} \dicDirectTranslationCS{o~(čem) všichni mluví}
\dicEntry[hápunktur] \dicTerm{há··punkt|ur} \dicIPA{{h}{au}{\textlengthmark}{p\smash{\textsuperscript{h}}}{u}{\textsubring{n}}{\textsubring{d}}{\textscy}{\textsubring{r}}} \dicPos{m}[6] \dicFlx{(‑s, ‑ar)}[4] \dicDirectTranslationCS{vrchol, vyvrcholení} \dicExampleIS{hápunktur ævi} \dicExampleCS{vrchol života}
\dicEntry[hár] \dicTerm{hár\smash{\textsuperscript{1}}} \dicsymFrequent\  \dicIPA{{h}{au}{\textlengthmark}{\textsubring{r}}} \dicPos{n}[2] \dicFlx{(‑s, ‑)}[5] \textbf{1.} \dicDirectTranslationCS{vlas(y)} \dicExampleIS{þykkt sítt hár} \dicExampleCS{husté dlouhé vlasy}  \textbf{2.} \dicDirectTranslationCS{chlup}
\dicEntry[hár] \dicTerm{hár\smash{\textsuperscript{2}}} \dicsymFrequent\  \dicIPA{{h}{au}{\textlengthmark}{\textsubring{r}}} \dicPos{adj}[10] \dicFlx{(comp hærri, sup hæstur)}[14] \textbf{1.} \dicDirectTranslationCS{vysoký} \dicExampleIS{hátt hús} \dicExampleCS{vysoký dům}  \textbf{2.} \dicDirectTranslationCS{vysoký (plat ap.)} \dicExampleIS{háar tekjur} \dicExampleCS{vysoké příjmy}  \textbf{3.} \dicDirectTranslationCS{vysoký (pozice ap.)}  \textbf{4.} \dicDirectTranslationCS{hlasitý, hlučný} \dicExampleIS{tala hátt} \dicExampleCS{mluvit nahlas};  \dicPhraseIS{hafa hátt} \dicDirectTranslationCS{dělat hluk, hlučet}  \textbf{5.} \dicDirectTranslationCS{vysoký (tón ap.)}
\dicEntry[hárbursti] \dicTerm{hár··burst|i} \dicIPA{{h}{au}{r}{\textsubring{b}}{\textscy}{\textsubring{r}}{s}{\textsubring{d}}{\textsci}} \dicPos{m}[1] \dicFlx{(‑a, ‑ar)}[1] \dicDirectTranslationCS{kartáč na vlasy}
\dicEntry[háreysti] \dicTerm{há··reyst|i} \dicIPA{{h}{au}{\textlengthmark}{r}{ei}{s}{\textsubring{d}}{\textsci}} \dicPos{f}[3] \dicFlx{(‑i)}[3] \dicSynonym{hávaði} \dicDirectTranslationCS{hluk, halas, hlučení}
\dicEntry[hárgreiðsla] \dicTerm{hár··greiðsl|a} \dicIPA{{h}{au}{r}{\r{g}}{r}{ei}{ð}{s}{\textsubring{d}}{l}{a}} \dicPos{f}[1] \dicFlx{(‑u, ‑ur)}[13] \textbf{1.} \dicDirectTranslationCS{česání}  \textbf{2.} \dicDirectTranslationCS{účes}
\dicEntry[hárgreiðslukona] \dicTerm{hár·greiðslu··kon|a} \dicIPA{{h}{au}{r}{\r{g}}{r}{ei}{ð}{s}{\textsubring{d}}{l}{\textscy}{k\smash{\textsuperscript{h}}}{\textopeno}{n}{a}} \dicPos{f}[1] \dicFlx{(‑u, ‑ur)}[27] \dicSynonym*{hársnyrtir} \dicDirectTranslationCS{kadeřnice}
\dicEntry[hárkolla] \dicTerm{hár··koll|a} \dicIPA{{h}{au}{\textsubring{r}}{k\smash{\textsuperscript{h}}}{\textopeno}{\textsubring{d}}{l}{a}} \dicPos{f}[1] \dicFlx{(‑u, ‑ur)}[19] \dicDirectTranslationCS{paruka}
\dicEntry[hárlakk] \dicTerm{hár··|lakk} \dicIPA{{h}{au}{r}{l}{a}{h}{\r{g}}} \dicPos{n}[2] \dicFlx{(‑lakks, ‑lökk)}[8] \dicDirectTranslationCS{lak na vlasy}
\dicEntry[hárlaus] \dicTerm{hár··laus} \dicIPA{{h}{au}{r}{l}{\oe i}{s}} \dicPos{adj}[5]\dicFlx{}[-1] \dicDirectTranslationCS{bezvlasý}
\dicEntry[hárnál] \dicTerm{hár··nál} \dicIPA{{h}{au}{r}{n}{au}{\textsubring{l}}} \dicPos{f}[4] \dicFlx{(‑ar, ‑ar)}[1] \dicDirectTranslationCS{spona do vlasů, vlásenka, pinetka}
\dicEntry[háróma] \dicTerm{há··róma} \dicIPA{{h}{au}{\textlengthmark}{r}{ou}{m}{a}} \dicPos{adj}[13] \dicFlx{indecl}[1] \dicSynonym{hávær} \dicDirectTranslationCS{hlasitý, hlučný}
\dicEntry[hárpípa] \dicTerm{hár··píp|a} \dicIPA{{h}{au}{\textsubring{r}}{p\smash{\textsuperscript{h}}}{i}{\textsubring{b}}{a}} \dicPos{f}[1] \dicFlx{(‑u, ‑ur)}[19] \dicFieldCat{fyz.} \dicDirectTranslationCS{kapilára, vlásečnice} \dicIndirectTranslationCS{(úzká trubička s~nepatrnou světlostí)}
\dicEntry[hárrétt] \dicTerm{hár··rétt} \dicIPA{{h}{au}{r}{\textlengthmark}{j}{\textepsilon}{h}{\textsubring{d}}} \dicPos{adv} \dicDirectTranslationCS{naprosto\,/\addthin zcela správně};  \dicPhraseIS{hafa hárrétt fyrir sér} \dicDirectTranslationCS{mít naprostou pravdu}
\dicEntry[hárréttur] \dicTerm{hár··réttur}\addthinS\dicIPA{{h}{au}{r}{\textlengthmark}{j}{\textepsilon}{h}{\textsubring{d}}{\textscy}{\textsubring{r}}}\addthinS\dicPos{adj}[1]\dicFlx{}[-13] \dicDirectTranslationCS{zcela\,/\addthinS naprosto správ\-ný} \dicExampleIS{það er alveg hárrétt} \dicExampleCS{je to naprosto správně}
\dicEntry[hárskeri] \dicTerm{hár··sker|i} \dicIPA{{h}{au}{\textsubring{r}}{s}{\r{\textObardotlessj}}{\textepsilon}{r}{\textsci}} \dicPos{m}[1] \dicFlx{(‑a, ‑ar)}[1] \dicSynonym{rakari} \dicDirectTranslationCS{holič(ka), kadeřník, kadeřnice}
\dicEntry[hárspenna] \dicTerm{hár··spenn|a} \dicIPA{{h}{au}{\textsubring{r}}{s}{\textsubring{b}}{\textepsilon}{n}{a}} \dicPos{f}[1] \dicFlx{(‑u, ‑ur)}[7] \dicDirectTranslationCS{sponka, vlásenka, vlásnička}
\dicEntry[hártoga] \dicTerm{hár··tog|a} \dicIPA{{h}{au}{\textsubring{r}}{t\smash{\textsuperscript{h}}}{\textopeno}{\textbabygamma}{a}} \dicPos{v}[1] \dicFlx{(‑aði)}[1] \dicFlx{acc} \dicDirectTranslationCS{překroutit, překrucovat (slova ap.)}
\dicEntry[hártogun] \dicTerm{hár··tog|un} \dicIPA{{h}{au}{\textsubring{r}}{t\smash{\textsuperscript{h}}}{\textopeno}{\textbabygamma}{\textscy}{\textsubring{n}}} \dicPos{f}[7] \dicFlx{(‑unar, ‑anir)}[8] \dicDirectTranslationCS{puntičkářství, hnidopišství}
\dicEntry[hártoppur] \dicTerm{hár··topp|ur} \dicIPA{{h}{au}{\textsubring{r}}{t\smash{\textsuperscript{h}}}{\textopeno}{h}{\textsubring{b}}{\textscy}{\textsubring{r}}} \dicPos{m}[6] \dicFlx{(‑s, ‑ar)}[10] \dicDirectTranslationCS{příčesek, tupé}
\dicEntry[hárþurrka] \dicTerm{hár··þurrk|a} \dicIPA{{h}{au}{\textsubring{r}}{\texttheta}{\textscy}{\textsubring{r}}{\r{g}}{a}} \dicPos{f}[1] \dicFlx{(‑u, ‑ur)}[13] \dicDirectTranslationCS{fén} \dicIndirectTranslationCS{(elektrický vysoušeč vlasů)}
\dicEntry[háræð] \dicTerm{hár··æð} \dicIPA{{h}{au}{\textlengthmark}{r}{a}{i}{\texttheta}} \dicPos{f}[4] \dicFlx{(‑ar, ‑ar)}[1] \dicFieldCat{anat.} \dicDirectTranslationCS{kapilára, vlásečnice} \dicIndirectTranslationCS{(krevní céva nejmenšího průměru)}
\dicEntry[hás] \dicTerm{hás} \dicIPA{{h}{au}{\textlengthmark}{s}} \dicPos{adj}[5]\dicFlx{}[-1] \dicDirectTranslationCS{chraplavý, ochraptělý, chraptivý}
\dicEntry[háseti] \dicTerm{há··set|i} \dicIPA{{h}{au}{\textlengthmark}{s}{\textepsilon}{\textsubring{d}}{\textsci}} \dicPos{m}[1] \dicFlx{(‑a, ‑ar)}[1] \dicDirectTranslationCS{námořník, námořnice}
\dicEntry[háskalegur] \dicTerm{háska··legur} \dicIPA{{h}{au}{s}{\r{g}}{a}{l}{\textepsilon}{\textbabygamma}{\textscy}{\textsubring{r}}} \dicPos{adj}[1]\dicFlx{}[-8] \dicSynonym{hættulegur} \dicDirectTranslationCS{nebezpečný, krkolomný, riskantní}
\dicEntry[háskerpa] \dicTerm{há··skerp|a} \dicIPA{{h}{au}{\textlengthmark}{s}{\r{\textObardotlessj}}{\textepsilon}{\textsubring{r}}{\textsubring{b}}{a}} \dicPos{f}[1] \dicFlx{(‑u)}[5] \dicDirectTranslationCS{vysoké rozlišení} \dicExampleIS{háskerpusjónvarp} \dicExampleCS{televizor s~vysokým rozlišením}
\dicEntry[háski] \dicTerm{hásk|i} \dicIPA{{h}{au}{s}{\r{\textObardotlessj}}{\textsci}} \dicPos{m}[1] \dicFlx{(‑a, ‑ar)}[1] \dicSynonym{hætta\smash{\textsuperscript{1}}} \dicDirectTranslationCS{nebezpečí, riziko}
\dicEntry[háskólaborgari] \dicTerm{há·skóla··borg·ar|i} \dicIPA{{h}{au}{\textlengthmark}{s}{\r{g}}{ou}{l}{a}{\textsubring{b}}{\textopeno}{r}{\r{g}}{a}{r}{\textsci}} \dicPos{m}[1] \dicFlx{(‑a, ‑ar)}[13] \dicDirectTranslationCS{akademik, akademička, posluchač(ka) vysoké školy}
\dicEntry[háskólagráða] \dicTerm{há·skóla··gráð|a} \dicIPA{{h}{au}{\textlengthmark}{s}{\r{g}}{ou}{l}{a}{\r{g}}{r}{au}{ð}{a}} \dicPos{f}[1] \dicFlx{(‑u, ‑ur)}[19] \dicDirectTranslationCS{vysokoškolský titul}
\dicEntry[háskólakennari] \dicTerm{há·skóla··kenn·ar|i} \dicIPA{{h}{au}{\textlengthmark}{s}{\r{g}}{ou}{l}{a}{c\smash{\textsuperscript{h}}}{\textepsilon}{n}{a}{r}{\textsci}} \dicPos{m}[1] \dicFlx{(‑a, ‑ar)}[13] \dicFieldCat{škol.} \dicDirectTranslationCS{vysokoškolský pedagog\,/\addthin učitel, vysokoškolská pedagožka\,/\addthin učitelka}
\dicEntry[háskólamaður] \dicTerm{há·skóla··|maður} \dicIPA{{h}{au}{\textlengthmark}{s}{\r{g}}{ou}{l}{a}{m}{a}{ð}{\textscy}{\textsubring{r}}} \dicPos{m}[13] \dicFlx{(‑manns, ‑menn)}[2] \textbf{1.} \dicDirectTranslationCS{univerzitní zaměstnanec\,/\addthin zaměstnankyně}  \textbf{2.} \dicDirectTranslationCS{absolvent(ka) vysoké školy}
\dicEntry[háskólanemi] \dicTerm{há·skóla··nem|i} \dicIPA{{h}{au}{\textlengthmark}{s}{\r{g}}{ou}{l}{a}{n}{\textepsilon}{m}{\textsci}} \dicPos{m}[1] \dicFlx{(‑a, ‑ar)}[1] \dicDirectTranslationCS{vysokoškolský student, vysokoškolská studentka}
\dicEntry[háskólapróf] \dicTerm{há·skóla··próf} \dicIPA{{h}{au}{\textlengthmark}{s}{\r{g}}{ou}{l}{a}{p\smash{\textsuperscript{h}}}{r}{ou}{f}} \dicPos{n}[2] \dicFlx{(‑s, ‑)}[5] \dicDirectTranslationCS{(vysokoškolská) zkouška}
\dicEntry[háskólaráð] \dicTerm{há·skóla··ráð} \dicIPA{{h}{au}{\textlengthmark}{s}{\r{g}}{ou}{l}{a}{r}{au}{\texttheta}} \dicPos{n}[2] \dicFlx{(‑s, ‑)}[5] \dicFieldCat{škol.} \dicDirectTranslationCS{vysokoškolská rada}
\dicEntry[háskólastig] \dicTerm{há·skóla··stig} \dicIPA{{h}{au}{\textlengthmark}{s}{\r{g}}{ou}{l}{a}{s}{\textsubring{d}}{\textsci}{x}} \dicPos{n}[2] \dicFlx{(‑s)}[2] \dicFieldCat{škol.} \dicDirectTranslationCS{vysokoškolský stupeň}
\dicEntry[háskólastúdent] \dicTerm{há·skóla··stúdent} \dicIPA{{h}{au}{\textlengthmark}{s}{\r{g}}{ou}{l}{a}{s}{\textsubring{d}}{u}{\textsubring{d}}{\textepsilon}{\textsubring{n}}{\textsubring{d}}} \dicPos{m}[4] \dicFlx{(‑s, ‑ar)}[10] \dicDirectTranslationCS{vysokoškolský student, vysokoškolská studentka, vysokoškolák, vysokoškolačka}
\dicEntry[háskóli] \dicTerm{há··skól|i} \dicsymFrequent\  \dicIPA{{h}{au}{\textlengthmark}{s}{\r{g}}{ou}{l}{\textsci}} \dicPos{m}[1] \dicFlx{(‑a, ‑ar)}[1] \dicFieldCat{škol.} \dicDirectTranslationCS{univerzita, vysoká škola} \dicExampleIS{Háskóli Íslands var stofnaður 1911.} \dicExampleCS{Islandská univerzita byla založena v~roce 1911.}
\dicEntry[háslétta] \dicTerm{há··slétt|a} \dicIPA{{h}{au}{\textlengthmark}{s}{\textsubring{d}}{l}{j}{\textepsilon}{h}{\textsubring{d}}{a}} \dicPos{f}[1] \dicFlx{(‑u, ‑ur)}[19] \dicDirectTranslationCS{náhorní plošina\,/\addthin rovina}
\dicEntry[háspeki] \dicTerm{há··spek|i} \dicIPA{{h}{au}{\textlengthmark}{s}{\textsubring{b}}{\textepsilon}{\r{\textObardotlessj}}{\textsci}} \dicPos{f}[3] \dicFlx{(‑i)}[3] \dicFieldCat{filos.} \dicSynonym{frumspeki} \dicDirectTranslationCS{metafyzika}
\dicEntry[háspenna] \dicTerm{há··spenn|a} \dicIPA{{h}{au}{\textlengthmark}{s}{\textsubring{b}}{\textepsilon}{n}{a}} \dicPos{f}[1] \dicFlx{(‑u)}[5] \dicFieldCat{elek.} \dicDirectTranslationCS{vysoké napětí} \dicAntonym{lágspenna}
\dicEntry[hástafalæsihnappur] \dicTerm{há·stafa··læsi·hnapp|ur} \dicIPA{{h}\-{au}\-{\textlengthmark}\-{s}\-{\textsubring{d}}\-{a}\-{v}\-{a}\-{l}\-{a}\-{i}\-{s}\-{\textsci}\-{\textsubring{n}}\-{a}\-{h}\-{\textsubring{b}}\-{\textscy}\-{\textsubring{r}}\-} \dicPos{m}[6] \dicFlx{(‑s, ‑ar)}[25] \dicFieldCat{poč.} \dicDirectTranslationCS{klávesa Capslock}
\dicEntry[hástafur] \dicTerm{há··staf|ur} \dicIPA{{h}{au}{\textlengthmark}{s}{\textsubring{d}}{a}{v}{\textscy}{\textsubring{r}}} \dicPos{m}[9] \dicFlx{(‑s, ‑ir)}[9] \dicDirectTranslationCS{velké písmeno}
\dicEntry[hástemmdur] \dicTerm{há··stemmdur} \dicIPA{{h}{au}{\textlengthmark}{s}{\textsubring{d}}{\textepsilon}{m}{\textsubring{d}}{\textscy}{\textsubring{r}}} \dicPos{adj}[2]\dicFlx{}[-14] \dicDirectTranslationCS{velkohubý, nabubřelý}
\dicEntry[hástertur] \dicTerm{há··stert|ur} \dicIPA{{h}{au}{\textlengthmark}{s}{\textsubring{d}}{\textepsilon}{\textsubring{r}}{\textsubring{d}}{\textscy}{\textsubring{r}}} \dicPos{m}[6] \dicFlx{(‑s)}[3] \dicPhraseIS{hrósa e‑m upp í hástert} \dicLangCat{přen.} \dicDirectTranslationCS{pět ódy na (koho), vychvalovat (koho) až do nebes}
\dicEntry[hástétt] \dicTerm{há··stétt} \dicIPA{{h}{au}{\textlengthmark}{s}{\textsubring{d}}{j}{\textepsilon}{h}{\textsubring{d}}} \dicPos{f}[7] \dicFlx{(‑ar, ‑ir)}[1] \dicDirectTranslationCS{horní\,/\addthin vyšší třída (ve společnosti)} \dicAntonym{lágstétt}
\dicEntry[hástig] \dicTerm{há··stig} \dicIPA{{h}{au}{\textlengthmark}{s}{\textsubring{d}}{\textsci}{x}} \dicPos{n}[2] \dicFlx{(‑s, ‑)}[5] \textbf{1.} \dicSynonym{hámark} \dicDirectTranslationCS{vrchol, nejvyšší stupeň}  \textbf{2.} \dicFieldCat{jaz.} \dicSynonym{efstastig} \dicDirectTranslationCS{superlativ, 3. stupeň stupňování}
\dicEntry[hástöfum] \dicTerm{há··stöfum} \dicIPA{{h}{au}{\textlengthmark}{s}{\textsubring{d}}{\oe}{v}{\textscy}{\textsubring{m}}} \dicPos{adv} \dicDirectTranslationCS{hlasitě, nahlas}
\dicEntry[hástökk] \dicTerm{há··stökk} \dicIPA{{h}{au}{\textlengthmark}{s}{\textsubring{d}}{\oe}{h}{\r{g}}} \dicPos{n}[2] \dicFlx{(‑s, ‑)}[5] \dicFieldCat{sport.} \dicDirectTranslationCS{skok vysoký}
\dicEntry[hásuður] \dicTerm{há··suður} \dicIPA{{h}{au}{\textlengthmark}{s}{\textscy}{ð}{\textscy}{\textsubring{r}}} \dicPos{n}[2] \dicFlx{(‑s)}[28] \dicDirectTranslationCS{přímý jih}
\dicEntry[hásæti] \dicTerm{há··sæti} \dicIPA{{h}{au}{\textlengthmark}{s}{a}{i}{\textsubring{d}}{\textsci}} \dicPos{n}[2] \dicFlx{(‑s, ‑)}[14] \dicDirectTranslationCS{trůn, stolec}
\dicEntry[hátalarakerfi] \dicTerm{há·talara··kerfi} \dicIPA{{h}{au}{\textlengthmark}{t\smash{\textsuperscript{h}}}{a}{l}{a}{r}{a}{c\smash{\textsuperscript{h}}}{\textepsilon}{r}{v}{\textsci}} \dicPos{n}[2] \dicFlx{(‑s, ‑)}[14] \dicDirectTranslationCS{reprosoustava}
\dicEntry[hátalari] \dicTerm{há··tal·ar|i} \dicIPA{{h}{au}{\textlengthmark}{t\smash{\textsuperscript{h}}}{a}{l}{a}{r}{\textsci}} \dicPos{m}[1] \dicFlx{(‑a, ‑ar)}[10] \dicDirectTranslationCS{reproduktor}
\dicEntry[hátekjumaður] \dicTerm{há·tekju··|maður} \dicIPA{{h}{au}{\textlengthmark}{t\smash{\textsuperscript{h}}}{\textepsilon}{\r{\textObardotlessj}}{\textscy}{m}{a}{ð}{\textscy}{\textsubring{r}}} \dicPos{m}[13] \dicFlx{(‑manns, ‑menn)}[2] \dicDirectTranslationCS{osoba s~vysokým příjmem}
\dicEntry[hátign] \dicTerm{há··tign} \dicIPA{{h}{au}{\textlengthmark}{t\smash{\textsuperscript{h}}}{\textsci}{\r{g}}{\textsubring{n}}} \dicPos{f}[7] \dicFlx{(‑ar, ‑ir)}[1] \dicDirectTranslationCS{výsost, veličenstvo};  \dicPhraseIS{hans hátign} \dicDirectTranslationCS{jeho výsost};  \dicPhraseIS{hennar hátign} \dicDirectTranslationCS{její výsost};  \dicPhraseIS{yðar hátign} \dicDirectTranslationCS{Vaše výsost}
\dicEntry[hátindur] \dicTerm{há··tind|ur} \dicIPA{{h}{au}{\textlengthmark}{t\smash{\textsuperscript{h}}}{\textsci}{n}{\textsubring{d}}{\textscy}{\textsubring{r}}} \dicPos{m}[6] \dicFlx{(‑s, ‑ar)}[4] \textbf{1.} \dicDirectTranslationCS{vrchol, vrcholek (hory ap.)}  \textbf{2.} \dicSynonym{hápunktur} \dicDirectTranslationCS{vrchol, vyvrcholení (vyprávění ap.)}
\dicEntry[hátíð] \dicTerm{há··tíð} \dicIPA{{h}{au}{\textlengthmark}{t\smash{\textsuperscript{h}}}{i}{\texttheta}} \dicPos{f}[4] \dicFlx{(‑ar, ‑ir\,/\addthin ‑ar)}[30] \textbf{1.} \dicDirectTranslationCS{svátek, svátky (Vánoce ap.)} \dicExampleIS{Það voru allir heima yfir hátíðarnar.} \dicExampleCS{Na svátky byli všichni doma.}  \textbf{2.} \dicDirectTranslationCS{slavnost, festival};  \dicPhraseIS{gleðilega hátíð!} \dicDirectTranslationCS{šťastné a~veselé svátky!} \dicIndirectTranslationCS{(přání užívané o~Vánocích, Velikonocích ap.)};  \dicPhraseIS{halda hátíð} \dicDirectTranslationCS{slavit svátky};  \dicPhraseIS{um hátíðina} \dicFlx{adv} \dicDirectTranslationCS{o~svátcích};  \dicPhraseIS{um\,/\addthin yfir hátíðarnar\,/\addthin hátíðirnar} \dicFlx{adv} \dicDirectTranslationCS{o~Vánocích\,/\addthin Velikonocích}
\dicEntry[hátíðarlilja] \dicTerm{há·tíðar··lilj|a} \dicIPA{{h}{au}{\textlengthmark}{t\smash{\textsuperscript{h}}}{i}{ð}{a}{r}{l}{\textsci}{l}{j}{a}} \dicPos{f}[1] \dicFlx{(‑u, ‑ur)}[7] \dicFieldCat{bot.} \dicDirectTranslationCS{narcis} \textit{(l.~{\textLA{Narcissus}})}  \dicsymPhoto\ 
\dicFigure{6837.jpg}{Hátíðarlilja}{Hátíðarlilja - Antušek Ivo, Biolib, Copyright/CC-BY-NC}
\begin{xtolerant}{}{1pt}
\dicEntry[hátíðarskap] \dicTerm{há·tíðar··skap} \dicIPA{{h}{au}{\textlengthmark}{t\smash{\textsuperscript{h}}}{i}{ð}{a}{\textsubring{r}}{s}{\r{g}}{a}{\textsubring{b}}} \dicPos{n}[2] \dicFlx{(‑s)}[2] \dicDirectTranslationCS{slavnostní nálada}
\end{xtolerant}
\dicEntry[hátíðisdagur] \dicTerm{há·tíðis··dag|ur} \dicIPA{{h}{au}{\textlengthmark}{t\smash{\textsuperscript{h}}}{i}{ð}{\textsci}{s}{\textsubring{d}}{a}{\textbabygamma}{\textscy}{\textsubring{r}}} \dicPos{m}[6] \dicFlx{(‑s, ‑ar)}[62] \dicSynonym*{merkisdagur} \dicDirectTranslationCS{slavnostní den}
\dicEntry[hátíðlega] \dicTerm{há·tíð··lega} \dicIPA{{h}{au}{\textlengthmark}{t\smash{\textsuperscript{h}}}{i}{ð}{l}{\textepsilon}{\textbabygamma}{a}} \dicPos{adv} \dicDirectTranslationCS{slavnostně}
\dicEntry[hátíðlegur] \dicTerm{há·tíð··legur} \dicsymFrequent\  \dicIPA{{h}{au}{\textlengthmark}{t\smash{\textsuperscript{h}}}{i}{ð}{l}{\textepsilon}{\textbabygamma}{\textscy}{\textsubring{r}}} \dicPos{adj}[1]\dicFlx{}[-8] \textbf{1.} \dicDirectTranslationCS{sváteční}  \textbf{2.} \dicDirectTranslationCS{slavnostní, ceremoniální} \dicExampleIS{hátíðlegt loforð} \dicExampleCS{slavnostní přísaha}
\dicEntry[hátt] \dicTerm{hátt} \dicsymFrequent\  \dicIPA{{h}{au}{h}{\textsubring{d}}} \dicPos{adv} \dicFlx{(comp hærra, sup hæst)} \textbf{1.} \dicDirectTranslationCS{vysoko (o~pozici ap.)};  \dicPhraseIS{hátt og lágt} \dicFlx{adv} \dicDirectTranslationCS{všude (možně)} \dicExampleIS{leita hátt og lágt} \dicExampleCS{hledat všude (možně)}  \textbf{2.} \dicDirectTranslationCS{hlasitě (o~zvuku ap.)};  \dicPhraseIS{hafa hátt} \dicDirectTranslationCS{dělat hluk, hlučet};  \dicPhraseIS{hátt á e‑u} \dicDirectTranslationCS{téměř (co)} \dicExampleIS{hátt á þriðju viku} \dicExampleCS{téměř tři týdny};  \dicPhraseIS{hátt í e‑ð} \dicDirectTranslationCS{téměř (co)}
\dicEntry[hátta] \dicTerm{hátt|a} \dicsymFrequent\  \dicIPA{{h}{au}{h}{\textsubring{d}}{a}} \dicPos{v}[1] \dicFlx{(‑aði)}[63] \dicFlx{acc\,/\addthin dat} \textbf{1.} \dicFlx{acc} \dicSynonym*{ganga til hvílu} \dicDirectTranslationCS{jít do postele, jít spát} \dicExampleIS{Farið var að hátta kl. 11.} \dicExampleCS{V~11 hodin se šlo do postele.}  \textbf{2.} \dicFlx{dat} \dicSynonym{haga} \dicDirectTranslationCS{být uspořádán (určitým způsobem ap.)};  \dicPhraseIS{e‑u er þannig háttað} \dicFlx{impers} \dicDirectTranslationCS{(co) je tak uspořádáno} \dicExampleIS{skólanum er þannig háttað að} \dicExampleCS{škola je uspořádána tak, že}
\dicEntry[háttalag] \dicTerm{hátta··lag} \dicIPA{{h}{au}{h}{\textsubring{d}}{a}{l}{a}{x}} \dicPos{n}[2] \dicFlx{(‑s)}[2] \dicDirectTranslationCS{chování, způsoby}
\dicEntry[háttar] \dicTerm{háttar} \dicIPA{{h}{au}{h}{\textsubring{d}}{a}{\textsubring{r}}} \dicPos{m} \dicFlx{sg gen} \dicLink{háttur\smash{\textsuperscript{1}}}
\dicEntry[háttatími] \dicTerm{hátta··tím|i} \dicIPA{{h}{au}{h}{\textsubring{d}}{a}{t\smash{\textsuperscript{h}}}{i}{m}{\textsci}} \dicPos{m}[1] \dicFlx{(‑a, ‑ar)}[1] \dicDirectTranslationCS{čas jít spát, čas jít do postele}
\dicEntry[hátterni] \dicTerm{hátt··erni} \dicIPA{{h}{au}{h}{\textsubring{d}}{\textlengthmark}{\textepsilon}{r}{\textsubring{d}}{n}{\textsci}} \dicPos{n}[2] \dicFlx{(‑s)}[20] \dicDirectTranslationCS{chování, způsoby}
\dicEntry[háttprúður] \dicTerm{hátt··prúður} \dicIPA{{h}{au}{h}{\textsubring{d}}{p\smash{\textsuperscript{h}}}{r}{u}{ð}{\textscy}{\textsubring{r}}} \dicPos{adj}[2]\dicFlx{}[-6] \dicSynonym{kurteis} \dicDirectTranslationCS{zdvořilý, kultivovaný}
\dicEntry[háttprýði] \dicTerm{hátt··prýð|i} \dicIPA{{h}{au}{h}{\textsubring{d}}{p\smash{\textsuperscript{h}}}{r}{i}{ð}{\textsci}} \dicPos{f}[3] \dicFlx{(‑i)}[3] \dicDirectTranslationCS{zdvořilost, kultivovanost}
\dicEntry[háttsettur] \dicTerm{hátt··settur} \dicIPA{{h}{au}{h}{\textsubring{d}}{s}{\textepsilon}{h}{\textsubring{d}}{\textscy}{\textsubring{r}}} \dicPos{adj}[1]\dicFlx{}[-10] \dicDirectTranslationCS{vysoce postavený, prominentní} \dicExampleIS{háttsettur embættismaður} \dicExampleCS{vysoce postavený činitel}
\dicEntry[háttur] \dicTerm{háttur\smash{\textsuperscript{1}}} \dicsymFrequent\  \dicIPA{{h}{au}{h}{\textsubring{d}}{\textscy}{\textsubring{r}}} \dicPos{m}[12] \dicFlx{(háttar, hættir)}[7] \textbf{1.} \dicSynonym{venja\smash{\textsuperscript{1}}} \dicDirectTranslationCS{zvyk, obyčej} \dicExampleIS{Það er háttur hans að ganga út.} \dicExampleCS{Má ve zvyku chodit ven.}  \textbf{2.} \dicSynonym{máti} \dicDirectTranslationCS{způsob};  \dicPhraseIS{að hætti e‑rs} \dicDirectTranslationCS{podle (stejného) způsobu jako (kdo)};  \dicPhraseIS{á annan hátt} \dicFlx{adv} \dicDirectTranslationCS{jiným způsobem, jinak};  \dicPhraseIS{á engan hátt} \dicFlx{adv} \dicDirectTranslationCS{žádným způsobem, nijak};  \dicPhraseIS{á sama hátt} \dicFlx{adv} \dicDirectTranslationCS{stejným způsobem, stejně};  \dicPhraseIS{á ýmsan hátt} \dicFlx{adv} \dicDirectTranslationCS{v~mnoha ohledech};  \dicPhraseIS{með e‑u hætti} \dicDirectTranslationCS{(jakým) způsobem};  \dicPhraseIS{meiri háttar} \dicFlx{adj} \dicDirectTranslationCS{super, vynikající, výborný};  \dicPhraseIS{minni háttar} \dicFlx{adj} \dicDirectTranslationCS{menší, zanedbatelný};  \dicPhraseIS{þess háttar} \dicFlx{adj} \dicDirectTranslationCS{takový, takovýto, podobný};  \dicPhraseIS{þess háttar} \dicFlx{adv} \dicDirectTranslationCS{(tak) podobně}  \textbf{3.} \dicSynonym{útlit} \dicDirectTranslationCS{vzhled, zjev} \dicExampleIS{Hvernig er hann í hátt?} \dicExampleCS{Jak vypadá?}  \textbf{4.} \dicFieldCat{lit.} \dicSynonym{bragarháttur} \dicDirectTranslationCS{metrum}  \textbf{5.} \dicFieldCat{jaz.} \dicDirectTranslationCS{způsob} \dicExampleIS{boðháttur} \dicExampleCS{rozkazovací způsob}
\dicEntry[háttur] \dicTerm{háttur\smash{\textsuperscript{2}}} \dicIPA{{h}{au}{h}{\textsubring{d}}{\textscy}{\textsubring{r}}} \dicPos{m}[12] \dicFlx{(háttar, hættir)}[7] \dicPhraseIS{fara í háttinn} \dicDirectTranslationCS{jít do postele, jít si lehnout}
\dicEntry[háttvirtur] \dicTerm{hátt··virtur} \dicIPA{{h}{au}{h}{\textsubring{d}}{v}{\textsci}{\textsubring{r}}{\textsubring{d}}{\textscy}{\textsubring{r}}} \dicPos{adj}[1]\dicFlx{}[-13] \textbf{1.} \dicSynonym{virtur} \dicDirectTranslationCS{vážený, respektovaný}  \textbf{2.} \dicDirectTranslationCS{vážený} \dicIndirectTranslationCS{(zdvořilé oslovení poslanců)} \dicExampleIS{háttvirtur þingmaður} \dicExampleCS{vážený poslanec}
\begin{xtolerant}{}{1pt}
\dicEntry[háttvís] \dicTerm{hátt··vís} \dicIPA{{h}{au}{h}{\textsubring{d}}{v}{i}{s}}\addthinS\dicPos{adj}[5]\dicFlx{}[-1]\addthinS\dicSynonym{háttprúður} \dicDirectTranslationCS{korektní, taktní, diskrétní}
\end{xtolerant}
\dicEntry[háttvísi] \dicTerm{hátt··vís|i} \dicIPA{{h}{au}{h}{\textsubring{d}}{v}{i}{s}{\textsci}} \dicPos{f}[3] \dicFlx{(‑i)}[3] \dicSynonym{kurteisi} \dicDirectTranslationCS{takt, taktnost}
\dicEntry[hávaði] \dicTerm{há··vað|i} \dicsymFrequent\  \dicIPA{{h}{au}{\textlengthmark}{v}{a}{ð}{\textsci}} \dicPos{m}[1] \dicFlx{(‑a, ‑ar)}[8] \dicSynonym{hark} \dicDirectTranslationCS{hluk, halas, rámus};  \dicPhraseIS{hávaði í e‑u} \dicDirectTranslationCS{hluk (čeho)} \dicExampleIS{hávaði í vélinni} \dicExampleCS{hluk stroje};  \dicPhraseIS{hávaði frá e‑u} \dicDirectTranslationCS{hluk z~(čeho)}
\dicEntry[hávaxinn] \dicTerm{há··vaxinn} \dicsymFrequent\  \dicIPA{{h}{au}{\textlengthmark}{v}{a}{x}{s}{\textsci}{\textsubring{n}}} \dicPos{adj}[6]\dicFlx{}[-3] \dicDirectTranslationCS{vysoký, urostlý} \dicExampleIS{hávaxið gras} \dicExampleCS{vysoká tráva} \dicAntonym{smávaxinn}
\dicEntry[hávegir] \dicTerm{há··vegir} \dicIPA{{h}{au}{\textlengthmark}{v}{ei}{\textsci}{\textsubring{r}}} \dicPos{m}[10] \dicFlx{pl}[1] \dicPhraseIS{hafa e‑ð í hávegum} \dicDirectTranslationCS{chovat (co) v~úctě, cenit si (čeho)}
\dicEntry[hávella] \dicTerm{há··vell|a} \dicIPA{{h}{au}{\textlengthmark}{v}{\textepsilon}{\textsubring{d}}{l}{a}} \dicPos{f}[1] \dicFlx{(‑u, ‑ur)}[19] \dicFieldCat{zool.} \dicDirectTranslationCS{hoholka lední} \textit{(l.~{\textLA{Clangula hyemalis}})}  \dicsymPhoto\ 
\dicFigure{21083.jpg}{Hávella}{Hávella - Bowman, Tim, Biolib, PD}
\dicEntry[hávær] \dicTerm{há··vær} \dicsymFrequent\  \dicIPA{{h}{au}{\textlengthmark}{v}{a}{i}{\textsubring{r}}} \dicPos{adj}[5] \dicFlx{(f ‑)}[8] \dicDirectTranslationCS{hlučný, hlasitý} \dicExampleIS{hávær bjalla} \dicExampleCS{hlučný zvonek}
\dicEntry[háþróaður] \dicTerm{há··þró·|aður} \dicIPA{{h}{au}{\textlengthmark}{\texttheta}{r}{ou}{a}{ð}{\textscy}{\textsubring{r}}} \dicPos{adj}[3] \dicFlx{(f ‑uð)}[3] \textbf{1.} \dicSynonym*{á háu þróunarstigi} \dicDirectTranslationCS{pokročilý, vyspělý} \dicExampleIS{háþróuð menning} \dicExampleCS{vyspělá kultura}  \textbf{2.} \dicSynonym*{tæknilega fullkominn} \dicDirectTranslationCS{technologicky vyspělý} \dicExampleIS{háþróuð tækni} \dicExampleCS{vyspělá technologie}
\dicEntry[háþrýstingur] \dicTerm{há··þrýst·ing|ur} \dicIPA{{h}{au}{\textlengthmark}{\texttheta}{r}{i}{s}{\textsubring{d}}{i}{\ng}{\r{g}}{\textscy}{\textsubring{r}}} \dicPos{m}[6] \dicFlx{(‑s)}[9] \textbf{1.} \dicFieldCat{med.} \dicDirectTranslationCS{vysoký tlak}  \textbf{2.} \dicFieldCat{meteo.} \dicDirectTranslationCS{vysoký tlak}
\dicEntry[háþrýstisvæði] \dicTerm{há·þrýsti··svæði} \dicIPA{{h}{au}{\textlengthmark}{\texttheta}{r}{i}{s}{\textsubring{d}}{\textsci}{s}{v}{a}{i}{ð}{\textsci}} \dicPos{n}[2] \dicFlx{(‑s, ‑)}[14] \dicFieldCat{meteo.} \dicDirectTranslationCS{oblast vysokého tlaku} \dicAntonym{lágþrýstisvæði}
\dicEntry[háþýska] \dicTerm{há··þýsk|a} \dicIPA{{h}{au}{\textlengthmark}{\texttheta}{i}{s}{\r{g}}{a}} \dicPos{f}[1] \dicFlx{(‑u)}[5] \dicDirectTranslationCS{horní němčina} \dicIndirectTranslationCS{(dialekt němčiny používaný v~jižním Německu, Rakousku a~Švýcarsku)} \dicAntonym{lágþýska}
\dicEntry[he he] \dicTerm{he he} \dicIPA{{h}{\textepsilon}{h}{\textepsilon}} \dicPos{inter} \dicDirectTranslationCS{haha, chichi}
\dicEntry[hebreska] \dicTerm{hebresk|a} \dicIPA{{h}{\textepsilon}{\textsubring{b}}{r}{\textepsilon}{s}{\r{g}}{a}} \dicPos{f}[1] \dicFlx{(‑u)}[5] \dicDirectTranslationCS{hebrejština};  \dicPhraseIS{e‑að er hebreska fyrir e‑m} \dicLangCat{přen.} \dicDirectTranslationCS{(co) je pro (koho) španělská vesnice}
\dicEntry[hebreskur] \dicTerm{hebreskur} \dicIPA{{h}{\textepsilon}{\textsubring{b}}{r}{\textepsilon}{s}{\r{g}}{\textscy}{\textsubring{r}}} \dicPos{adj}[1]\dicFlx{}[-1] \dicDirectTranslationCS{hebrejský}
\dicEntry[hef] \dicTerm{hef} \dicIPA{{h}{\textepsilon}{\textlengthmark}{f}} \dicPos{v} \dicFlx{ind praes sg 1 pers} \textbf{1.} \dicLink{hafa}  \textbf{2.} \dicLink{hefja}
\dicEntry[hefa] \dicTerm{hef|a} \dicIPA{{h}{\textepsilon}{\textlengthmark}{v}{a}} \dicPos{v}[1] \dicFlx{(‑aði)}[1] \dicPhraseIS{hefa sig} \dicDirectTranslationCS{(na)kynout (těsto ap.)};  \dicIdiom{hefast}{ \dicPhraseIS{hefast}} \dicFlx{refl} \dicDirectTranslationCS{(na)kynout (těsto ap.)}
\dicEntry[hefð] \dicTerm{hefð} \dicIPA{{h}{\textepsilon}{v}{\texttheta}} \dicPos{f}[7] \dicFlx{(‑ar, ‑ir)}[1] \dicSynonym{siður} \dicDirectTranslationCS{tradice, obyčej, zvyk} \dicExampleIS{Það er orðin hefð hjá okkur.} \dicExampleCS{Už je to u~nás tradice.}
\dicEntry[hefðarfólk] \dicTerm{hefðar··fólk} \dicIPA{{h}{\textepsilon}{v}{ð}{a}{\textsubring{r}}{f}{ou}{\textsubring{l}}{\r{g}}} \dicPos{n}[2] \dicFlx{(‑s)}[2] \dicDirectTranslationCS{panstvo, urození lidé}
\dicEntry[hefðarkona] \dicTerm{hefðar··kon|a} \dicIPA{{h}{\textepsilon}{v}{ð}{a}{\textsubring{r}}{k\smash{\textsuperscript{h}}}{\textopeno}{n}{a}} \dicPos{f}[1] \dicFlx{(‑u, ‑ur)}[27] \dicDirectTranslationCS{urozená paní, paní z~lepší společnosti}
\dicEntry[hefðarmaður] \dicTerm{hefðar··|maður} \dicIPA{{h}{\textepsilon}{v}{ð}{a}{r}{m}{a}{ð}{\textscy}{\textsubring{r}}} \dicPos{m}[13] \dicFlx{(‑manns, ‑menn)}[2] \dicDirectTranslationCS{džentlmen, urozený pán, pán z~lepší společnosti}
\dicEntry[hefðbundinn] \dicTerm{hefð··bundinn} \dicsymFrequent\  \dicIPA{{h}{\textepsilon}{v}{ð}{\textsubring{b}}{\textscy}{n}{\textsubring{d}}{\textsci}{\textsubring{n}}} \dicPos{adj}[6]\dicFlx{}[-2] \dicDirectTranslationCS{tradiční} \dicExampleIS{á hefðbundinn hátt} \dicExampleCS{tradičním způsobem}
\dicEntry[hefði] \dicTerm{hefði} \dicIPA{{h}{\textepsilon}{v}{ð}{\textsci}} \dicPos{v} \dicFlx{con pf sg 1 pers} \dicLink{hafa}
\dicEntry[hefill] \dicTerm{hef|ill} \dicIPA{{h}{\textepsilon}{\textlengthmark}{v}{\textsci}{\textsubring{d}}{\textsubring{l}}} \dicPos{m}[6] \dicFlx{(‑ils, ‑lar)}[35] \textbf{1.} \dicDirectTranslationCS{hoblík}  \textbf{2.} \dicSynonym{veghefill} \dicDirectTranslationCS{srovnávač, grejdr} \dicIndirectTranslationCS{(stavební stroj)}
\dicEntry[hefja] \dicTerm{hefja} \dicsymFrequent\  \dicIPA{{h}{\textepsilon}{v}{j}{a}} \dicPos{v}[6] \dicFlx{(hef, hóf, hófum, hæfi, hafið)}[114] \dicFlx{acc} \textbf{1.} \dicSynonym{byrja} \dicDirectTranslationCS{začít, začínat, zahájit, zahajovat} \dicExampleIS{hefja verkið} \dicExampleCS{začít práci}  \textbf{2.} \dicSynonym{lyfta\smash{\textsuperscript{2}}} \dicDirectTranslationCS{zvednout, zvedat, pozvednout, pozvedat, vztyčit, vztyčovat} \dicExampleIS{hefja steininn á loft} \dicExampleCS{zvednout kámen do vzduchu};  \dicIdiom{hefja}[til]{ \dicPhraseIS{hefja sig til flugs}} \dicDirectTranslationCS{vzlétnout}; { \dicPhraseIS{hefja e‑n til skýjanna}} \dicLangCat{přen.} \dicDirectTranslationCS{vynášet (koho) do nebes};  \dicIdiom{hefja}[upp]{ \dicPhraseIS{hefja upp raust sína}} \dicDirectTranslationCS{zvýšit\,/\addthin zvednout hlas};  \dicIdiom{hefjast}{ \dicPhraseIS{hefjast}} \dicFlx{refl} \dicDirectTranslationCS{začít (se), začínat (se)} \dicExampleIS{Sagan hefst í Reykjavík.} \dicExampleCS{Příběh začíná v~Reykjavíku.}; { \dicPhraseIS{hefjast handa við e‑ð}} \dicFlx{refl} \dicDirectTranslationCS{pustit se do (čeho), dát se do (čeho)}
\dicEntry[hefla] \dicTerm{hefl|a} \dicIPA{{h}{\textepsilon}{\textsubring{b}}{l}{a}} \dicPos{v}[1] \dicFlx{(‑aði)}[1] \dicFlx{acc} \textbf{1.} \dicDirectTranslationCS{(o)hoblovat}  \textbf{2.} \dicDirectTranslationCS{srovnat (povrch ap.)}  \textbf{3.} \dicSynonym{lagfæra} \dicDirectTranslationCS{vytříbit, vypilovat, zdokonalit} \dicExampleIS{hefla orðalag} \dicExampleCS{vytříbit formulaci}
\dicEntry[heflaður] \dicTerm{hefl··|aður} \dicIPA{{h}{\textepsilon}{\textsubring{b}}{l}{a}{ð}{\textscy}{\textsubring{r}}} \dicPos{adj}[3] \dicFlx{(f ‑uð)}[3] \textbf{1.} \dicDirectTranslationCS{ohoblovaný}  \textbf{2.} \dicDirectTranslationCS{srovnaný (povrch ap.)}  \textbf{3.} \dicSynonym{fágaður} \dicDirectTranslationCS{vytříbený, kultivovaný, uhlazený} \dicExampleIS{heflaður í framkomu} \dicExampleCS{kultivovaný v~chování}
\dicEntry[hefna] \dicTerm{hefn|a} \dicIPA{{h}{\textepsilon}{\textsubring{b}}{n}{a}} \dicPos{v}[2] \dicFlx{(‑di, ‑t)}[153] \dicFlx{gen} \dicSynonym*{gjalda illt með illu} \dicDirectTranslationCS{(po)mstít};  \dicPhraseIS{hefna sín á e‑m fyrir e‑ð} \dicDirectTranslationCS{(po)mstít se (komu) za (co)};  \dicPhraseIS{eiga harma að hefna} \dicDirectTranslationCS{mít důvod k~odplatě};  \dicIdiom{hefnast}{ \dicPhraseIS{e‑m hefnist fyrir e‑ð}} \dicFlx{refl impers} \dicDirectTranslationCS{(co) se (komu) vymstí} \dicExampleIS{Henni hefndist fyrir mistökin.} \dicExampleCS{Ty chyby se jí vymstily.}
\dicEntry[hefnd] \dicTerm{hefnd} \dicIPA{{h}{\textepsilon}{m}{\textsubring{d}}} \dicPos{f}[7] \dicFlx{(‑ar, ‑ir)}[1] \dicDirectTranslationCS{(po)msta, odplata, odveta} \dicExampleIS{blóðhefnd} \dicExampleCS{krevní msta};  \dicPhraseIS{koma fram hefndum} \dicDirectTranslationCS{pomstít se};  \dicPhraseIS{leita hefnda} \dicDirectTranslationCS{vyhledat pomstu}
\dicEntry[hefndarþorsti] \dicTerm{hefndar··þorst|i} \dicIPA{{h}{\textepsilon}{m}{\textsubring{d}}{a}{\textsubring{r}}{\texttheta}{\textopeno}{\textsubring{r}}{s}{\textsubring{d}}{\textsci}} \dicPos{m}[1] \dicFlx{(‑a)}[3] \dicDirectTranslationCS{žízeň po pomstě}
\dicEntry[hefnigirni] \dicTerm{hefni··girn|i} \dicIPA{{h}{\textepsilon}{\textsubring{b}}{n}{\textsci}{\r{\textObardotlessj}}{\textsci}{r}{\textsubring{d}}{n}{\textsci}} \dicPos{f}[3] \dicFlx{(‑i)}[3] \dicDirectTranslationCS{pomstychtivost, mstivost}
\dicEntry[hefnigjarn] \dicTerm{hefni··|gjarn} \dicIPA{{h}{\textepsilon}{\textsubring{b}}{n}{\textsci}{\r{\textObardotlessj}}{a}{r}{\textsubring{d}}{\textsubring{n}}} \dicPos{adj}[5] \dicFlx{(f ‑gjörn)}[6] \dicSynonym*{hefndarfús} \dicDirectTranslationCS{pomstychtivý, mstivý}
\dicEntry[hefta] \dicTerm{heft|a\smash{\textsuperscript{1}}} \dicIPA{{h}{\textepsilon}{f}{\textsubring{d}}{a}} \dicPos{v}[1] \dicFlx{(‑aði)}[1] \dicFlx{acc} \dicDirectTranslationCS{svázat, sešít, secvaknout} \dicExampleIS{hefta bók} \dicExampleCS{svázat knížku}
\dicEntry[hefta] \dicTerm{heft|a\smash{\textsuperscript{2}}} \dicIPA{{h}{\textepsilon}{f}{\textsubring{d}}{a}} \dicPos{v}[2] \dicFlx{(‑i, ‑)}[12] \dicFlx{acc} \dicSynonym{hindra} \dicDirectTranslationCS{omezit, omezovat, ohraničit, (za)bránit} \dicExampleIS{hefta aðgang almennings að e‑u} \dicExampleCS{omezit přístup veřejnosti k~(čemu)}
\dicEntry[heftari] \dicTerm{heft··ar|i} \dicIPA{{h}{\textepsilon}{f}{\textsubring{d}}{a}{r}{\textsci}} \dicPos{m}[1] \dicFlx{(‑a, ‑ar)}[13] \dicDirectTranslationCS{sešívačka}
\dicEntry[hefti] \dicTerm{hefti} \dicsymFrequent\  \dicIPA{{h}{\textepsilon}{f}{\textsubring{d}}{\textsci}} \dicPos{n}[2] \dicFlx{(‑s, ‑)}[14] \textbf{1.} \dicDirectTranslationCS{sešit} \dicIndirectTranslationCS{(brožovaná tištěná publikace)} \dicExampleIS{Þetta hefti er einkum ætlað nemendum í 9. bekk.} \dicExampleCS{Tento sešit je určen zvláště žákům deváté třídy.}  \textbf{2.} \dicDirectTranslationCS{číslo} \dicIndirectTranslationCS{(výtisk publikace)}  \textbf{3.} \dicSynonym*{heftivír} \dicDirectTranslationCS{svorka (sešívací)}
\dicEntry[heftiplástur] \dicTerm{hefti··plást|ur} \dicIPA{{h}{\textepsilon}{f}{\textsubring{d}}{\textsci}{p\smash{\textsuperscript{h}}}{l}{au}{s}{\textsubring{d}}{\textscy}{\textsubring{r}}} \dicPos{m}[5] \dicFlx{(‑urs, ‑rar)}[1] \dicDirectTranslationCS{náplast, leukoplast}
\dicEntry[hegða] \dicTerm{hegð|a} \dicsymFrequent\  \dicIPA{{h}{\textepsilon}{\textbabygamma}{ð}{a}} \dicPos{v}[1] \dicFlx{(‑aði)}[34] \dicPhraseIS{hegða sér} \dicSynonym*{haga sér} \dicDirectTranslationCS{(za)chovat se} \dicExampleIS{hegða sér asnalega} \dicExampleCS{zachovat se hloupě}
\dicEntry[hegðun] \dicTerm{hegð|un} \dicsymFrequent\  \dicIPA{{h}{\textepsilon}{\textbabygamma}{ð}{\textscy}{\textsubring{n}}} \dicPos{f}[7] \dicFlx{(‑unar)}[9] \dicDirectTranslationCS{chování, jednání} \dicExampleIS{félagsleg hegðun} \dicExampleCS{společenské chování}
\dicEntry[hegg] \dicTerm{hegg} \dicIPA{{h}{\textepsilon}{\r{g}}{\textlengthmark}} \dicPos{v} \dicFlx{ind praes sg 1 pers} \dicLink{höggva}
\dicEntry[heggur] \dicTerm{hegg|ur} \dicIPA{{h}{\textepsilon}{\r{g}}{\textlengthmark}{\textscy}{\textsubring{r}}} \dicPos{m}[9] \dicFlx{(‑s, ‑ir)}[15] \dicFieldCat{bot.} \dicDirectTranslationCS{střemcha, střemcha obecná} \textit{(l.~{\textLA{Prunus padus}})}  \dicsymPhoto\ 
\dicFigure{ds_image_heggur_0_1.jpg}{Heggur}{Heggur - Udo Schröter, CC BY-SA 3.0}
\dicEntry[hegna] \dicTerm{hegn|a} \dicIPA{{h}{\textepsilon}{\r{g}}{n}{a}} \dicPos{v}[2] \dicFlx{(‑di, ‑t)}[149] \dicFlx{dat} \dicSynonym{refsa} \dicDirectTranslationCS{(po)trestat};  \dicPhraseIS{hegna e‑m fyrir e‑ð } \dicDirectTranslationCS{trestat (koho) za (co)}
\dicEntry[hegning] \dicTerm{hegn··ing} \dicIPA{{h}{\textepsilon}{\r{g}}{n}{i}{\ng}{\r{g}}} \dicPos{f}[4] \dicFlx{(‑ar, ‑ar)}[5] \dicSynonym{refsing} \dicDirectTranslationCS{(po)trestání} \dicExampleIS{fá hegningu} \dicExampleCS{být potrestán};  \dicPhraseIS{sæta hegningu} \dicDirectTranslationCS{být potrestán}
\dicEntry[hegningarhús] \dicTerm{hegningar··hús} \dicIPA{{h}{\textepsilon}{\r{g}}{n}{i}{\ng}{\r{g}}{a}{\textsubring{r}}{h}{u}{s}} \dicPos{n}[2] \dicFlx{(‑s, ‑)}[5] \dicDirectTranslationCS{káznice, trestnice}
\dicEntry[hegningarlög] \dicTerm{hegningar··lög} \dicIPA{{h}{\textepsilon}{\r{g}}{n}{i}{\ng}{\r{g}}{a}{r}{l}{\oe}{x}} \dicPos{n}[2] \dicFlx{pl}[9] \dicFieldCat{práv.} \dicDirectTranslationCS{trestní právo}
\dicEntry[hegri] \dicTerm{hegr|i} \dicIPA{{h}{\textepsilon}{\textbabygamma}{r}{\textsci}} \dicPos{m}[1] \dicFlx{(‑a, ‑ar)}[1] \textbf{1.} \dicFieldCat{techn.} \dicSynonym{krani} \dicDirectTranslationCS{jeřáb} \dicIndirectTranslationCS{(stroj)}  \textbf{2.} \dicFieldCat{zool.} \dicDirectTranslationCS{volavka} \textit{(l.~{\textLA{Ardea}})}  \dicsymPhoto\ 
\dicFigure{68087.jpg}{Hegri}{Hegri - Alena Andrej, Biolib, Copyright/CC-BY-NC}
\dicEntry[hei] \dicTerm{hei} \dicIPA{{h}{ei}{\textlengthmark}} \dicPos{inter} \dicDirectTranslationCS{hej} \dicIndirectTranslationCS{(vyjadřuje zavolání\,/\addthin upozornění)}
\dicEntry[heiðagæs] \dicTerm{heiða··gæs} \dicIPA{{h}{ei}{\textlengthmark}{ð}{a}{\r{\textObardotlessj}}{a}{i}{s}} \dicPos{f}[7] \dicFlx{(‑ar, ‑ir)}[1] \dicFieldCat{zool.} \dicDirectTranslationCS{husa krátkozobá} \textit{(l.~{\textLA{Anser brachyrhynchus}})}  \dicsymPhoto\ 
\dicFigure{ds_image_heidagaes_0_2.jpg}{Heiðagæs}{Heiðagæs - MPF, GFDL}
\begin{xtolerant}{}{1pt}
\dicEntry[heiðarlegur] \dicTerm{heiðar··legur}\addthinS\dicIPA{{h}{ei}{\textlengthmark}{ð}{a}{r}{l}{\textepsilon}{\textbabygamma}{\textscy}{\textsubring{r}}}\addthinS\dicPos{adj}[1]\dicFlx{}[-8]\addthinS\dicSynonym{áreiðanlegur} \dicDirectTranslationCS{čestný, poctivý} \dicExampleIS{vera heiðarlegur í viðskiptum} \dicExampleCS{poctivě obchodovat} \dicAntonym{óheiðarlegur}
\end{xtolerant}
\dicEntry[heiðarleiki] \dicTerm{heiðar··leik|i} \dicIPA{{h}{ei}{\textlengthmark}{ð}{a}{r}{l}{ei}{\r{\textObardotlessj}}{\textsci}} \dicPos{m}[1] \dicFlx{(‑a)}[3] \dicDirectTranslationCS{poctivost, čestnost}
\dicEntry[heiðblár] \dicTerm{heið··blár} \dicIPA{{h}{ei}{ð}{\textsubring{b}}{l}{au}{\textsubring{r}}} \dicPos{adj}[4]\dicFlx{}[-1] \dicDirectTranslationCS{azurový, blankytný}
\dicEntry[heiðhvolf] \dicTerm{heið··hvolf} \dicIPA{{h}{ei}{ð}{k\smash{\textsuperscript{h}}}{v}{\textopeno}{l}{f}} \dicPos{n}[2] \dicFlx{(‑s, ‑)}[5] \dicFieldCat{meteo.} \dicDirectTranslationCS{stratosféra}
\dicEntry[heiði] \dicTerm{heið|i\smash{\textsuperscript{1}}} \dicsymFrequent\  \dicIPA{{h}{ei}{\textlengthmark}{ð}{\textsci}} \dicPos{f}[5] \dicFlx{(‑ar, ‑ar)}[3] \textbf{1.} \dicDirectTranslationCS{vysočina, vrchovina} \dicExampleIS{fara upp á heiðina} \dicExampleCS{jet na vysočinu}  \textbf{2.} \dicDirectTranslationCS{vřesoviště}
\dicEntry[heiði] \dicTerm{heiði\smash{\textsuperscript{2}}} \dicIPA{{h}{ei}{\textlengthmark}{ð}{\textsci}} \dicPos{n}[2] \dicFlx{(‑s)}[20] \dicDirectTranslationCS{jasné nebe, jasná obloha}
\dicEntry[heiðingi] \dicTerm{heið··ing|i} \dicIPA{{h}{ei}{\textlengthmark}{ð}{i}{\textltailn}{\r{\textObardotlessj}}{\textsci}} \dicPos{m}[1] \dicFlx{(‑ja, ‑jar)}[14] \textbf{1.} \dicDirectTranslationCS{pohan(ka)} \dicIndirectTranslationCS{(člověk vyznávající jinou víru než křesťanství)}  \textbf{2.} \dicSynonym{trúleysingi} \dicDirectTranslationCS{ateista, ateistka}
\dicEntry[heiðinn] \dicTerm{heiðinn} \dicIPA{{h}{ei}{\textlengthmark}{ð}{\textsci}{\textsubring{n}}} \dicPos{adj}[6]\dicFlx{}[-2] \textbf{1.} \dicSynonym*{ókristinn} \dicDirectTranslationCS{pohanský}  \textbf{2.} \dicSynonym{trúlaus} \dicDirectTranslationCS{ateistický, nevěrecký}
\dicEntry[heiðlóa] \dicTerm{heið··ló|a} \dicIPA{{h}{ei}{ð}{l}{ou}{a}} \dicPos{f}[1] \dicFlx{(‑u, ‑ur)}[7] \dicFieldCat{zool.} \dicDirectTranslationCS{kulík zlatý} \textit{(l.~{\textLA{Pluvialis apricaria}})}  \dicsymPhoto\ 
\dicFigure{ds_image_heidloa_0_1.jpg}{Heiðlóa}{Heiðlóa - Quirin Herzog, CC BY-SA 3.0}
\dicEntry[heiðni] \dicTerm{heiðn|i} \dicIPA{{h}{ei}{ð}{n}{\textsci}} \dicPos{f}[3] \dicFlx{(‑i)}[3] \textbf{1.} \dicSynonym*{ekki kristni} \dicDirectTranslationCS{pohanství}  \textbf{2.} \dicFieldCat{náb.} \dicSynonym{ásatrú} \dicDirectTranslationCS{pohanská víra, ásatrú}
\dicEntry[heiðra] \dicTerm{heiðr|a} \dicIPA{{h}{ei}{ð}{r}{a}} \dicPos{v}[1] \dicFlx{(‑aði)}[34] \dicFlx{acc} \dicSynonym*{auðsýna heiður} \dicDirectTranslationCS{(po)ctít, uctívat, vzdát čest};  \dicPhraseIS{heiðra e‑n fyrir e‑ð} \dicDirectTranslationCS{vzdát čest (komu) za (co)} \dicExampleIS{heiðra hana fyrir merkilegt starf} \dicExampleCS{vzdát jí čest za vynikající práci}
\dicEntry[heiðskír] \dicTerm{heið··skír} \dicIPA{{h}{ei}{ð}{s}{\r{\textObardotlessj}}{i}{\textsubring{r}}} \dicPos{adj}[5] \dicFlx{(f ‑)}[8] \dicDirectTranslationCS{jasný, slunečný} \dicExampleIS{heiðskírt loft} \dicExampleCS{jasná obloha}
\dicEntry[heiður] \dicTerm{heið|ur\smash{\textsuperscript{1}}} \dicIPA{{h}{ei}{\textlengthmark}{ð}{\textscy}{\textsubring{r}}} \dicPos{m}[5] \dicFlx{(‑urs)}[2] \dicDirectTranslationCS{čest, pocta};  \dicPhraseIS{e‑m er heiður að e‑u} \dicFlx{impers} \dicDirectTranslationCS{(komu) je velkou ctí (co)};  \dicPhraseIS{eiga heiðurinn af e‑u} \dicDirectTranslationCS{zasloužit si poctu za (co)};  \dicPhraseIS{eiga heiður skilið\,/\addthin skilinn} \dicDirectTranslationCS{zasloužit si poctu};  \dicPhraseIS{halda e‑ð í heiðri} \dicDirectTranslationCS{mít (co) v~úctě, respektovat (co)}
\dicEntry[heiður] \dicTerm{heiður\smash{\textsuperscript{2}}} \dicsymFrequent\  \dicIPA{{h}{ei}{\textlengthmark}{ð}{\textscy}{\textsubring{r}}} \dicPos{adj}[2]\dicFlx{}[-6] \dicDirectTranslationCS{jasný, slunečný, bezmračný} \dicExampleIS{heiður himinn} \dicExampleCS{jasné nebe}
\dicEntry[heiðursborgari] \dicTerm{heiðurs··borg·ar|i} \dicIPA{{h}{ei}{\textlengthmark}{ð}{\textscy}{\textsubring{r}}{s}{\textsubring{b}}{\textopeno}{r}{\r{g}}{a}{r}{\textsci}} \dicPos{m}[1] \dicFlx{(‑a, ‑ar)}[13] \dicDirectTranslationCS{čestný občan, čestná občanka} \dicExampleIS{heiðursborgari Reykjavíkur} \dicExampleCS{čestný občan města Reykjavík}
\dicEntry[heiðursfélagi] \dicTerm{heiðurs··fé·lag|i} \dicIPA{{h}{ei}{\textlengthmark}{ð}{\textscy}{\textsubring{r}}{s}{f}{j}{\textepsilon}{l}{a}{i}{j}{\textsci}} \dicPos{m}[1] \dicFlx{(‑a, ‑ar)}[8] \dicDirectTranslationCS{čestný člen, čestná členka}
\dicEntry[heiðursgestur] \dicTerm{heiðurs··gest|ur} \dicIPA{{h}{ei}{\textlengthmark}{ð}{\textscy}{\textsubring{r}}{s}{\r{\textObardotlessj}}{\textepsilon}{s}{\textsubring{d}}{\textscy}{\textsubring{r}}} \dicPos{m}[9] \dicFlx{(‑s, ‑ir)}[4] \dicDirectTranslationCS{čestný host, čestná hostka}
\dicEntry[heiðursmerki] \dicTerm{heiðurs··merki} \dicIPA{{h}{ei}{\textlengthmark}{ð}{\textscy}{\textsubring{r}}{s}{m}{\textepsilon}{\textsubring{r}}{\r{\textObardotlessj}}{\textsci}} \dicPos{n}[2] \dicFlx{(‑s, ‑)}[16] \dicDirectTranslationCS{řád, vyznamenání}
\dicEntry[heiðvirður] \dicTerm{heið··virður} \dicIPA{{h}{ei}{ð}{v}{\textsci}{r}{ð}{\textscy}{\textsubring{r}}} \dicPos{adj}[2]\dicFlx{}[-1] \dicDirectTranslationCS{uznávaný, respektovaný, vážený}
\dicEntry[heift] \dicTerm{heift} \dicsymFrequent\  \dicIPA{{h}{ei}{f}{\textsubring{d}}} \dicPos{f}[7] \dicFlx{(‑ar, ‑ir)}[1] \dicSynonym{bræði} \dicDirectTranslationCS{hněv, vztek} \dicExampleIS{blanda af heift og ótta} \dicExampleCS{směs hněvu a~strachu}
\dicEntry[heiftrækinn] \dicTerm{heift··rækinn} \dicIPA{{h}{ei}{f}{\textsubring{d}}{r}{a}{i}{\r{\textObardotlessj}}{\textsci}{\textsubring{n}}} \dicPos{adj}[6]\dicFlx{}[-2] \dicDirectTranslationCS{pomstychtivý, mstivý}
\dicEntry[heiftrækni] \dicTerm{heift··rækn|i} \dicIPA{{h}{ei}{f}{\textsubring{d}}{r}{a}{i}{h}{\r{g}}{n}{\textsci}} \dicPos{f}[3] \dicFlx{(‑i)}[3] \dicDirectTranslationCS{pomstychtivost, mstivost}
\dicEntry[heiftúðugur] \dicTerm{heift·úð··ugur} \dicIPA{{h}{ei}{f}{\textsubring{d}}{u}{ð}{\textscy}{\textbabygamma}{\textscy}{\textsubring{r}}} \dicPos{adj}[1]\dicFlx{}[-8] \textbf{1.} \dicSynonym{hatursfullur} \dicDirectTranslationCS{záštiplný, nenávistný}  \textbf{2.} \dicSynonym{hefnigjarn} \dicDirectTranslationCS{pomstychtivý}
\dicEntry[heigull] \dicTerm{heig|ull} \dicIPA{{h}{ei}{\textlengthmark}{\textbabygamma}{\textscy}{\textsubring{d}}{\textsubring{l}}} \dicPos{m}[6] \dicFlx{(‑uls, ‑lar)}[31] \dicDirectTranslationCS{zbabělec, zbabělkyně}
\dicEntry[heilablóðfall] \dicTerm{heila··blóð·|fall} \dicIPA{{h}{ei}{\textlengthmark}{l}{a}{\textsubring{b}}{l}{ou}{ð}{f}{a}{\textsubring{d}}{\textsubring{l}}} \dicPos{n}[2] \dicFlx{(‑falls, ‑föll)}[8] \dicFieldCat{med.} \dicDirectTranslationCS{mrtvice}
\dicEntry[heilabrot] \dicTerm{heila··brot} \dicIPA{{h}{ei}{\textlengthmark}{l}{a}{\textsubring{b}}{r}{\textopeno}{\textsubring{d}}} \dicPos{n}[2] \dicFlx{(‑s, ‑)}[5] \dicDirectTranslationCS{hlavolam};  \dicPhraseIS{e‑að veldur e‑m heilabrotum} \dicLangCat{přen.} \dicDirectTranslationCS{(kdo) si láme hlavu s~(čím)}
\dicEntry[heilagleiki] \dicTerm{heilag··leik|i} \dicIPA{{h}{ei}{\textlengthmark}{l}{a}{\textbabygamma}{l}{ei}{\r{\textObardotlessj}}{\textsci}} \dicPos{m}[1] \dicFlx{(‑a)}[3] \dicDirectTranslationCS{svatost, posvátnost}
\dicEntry[heilagur] \dicTerm{heilagur} \dicsymFrequent\  \dicIPA{{h}{ei}{\textlengthmark}{l}{a}{\textbabygamma}{\textscy}{\textsubring{r}}} \dicPos{adj}[1] \dicFlx{(f heilög)}[2] \dicDirectTranslationCS{svatý} \dicExampleIS{heilagur sannleikur} \dicExampleCS{svatá pravda};  \dicPhraseIS{heilagur andi} \dicFieldCat{náb.} \dicDirectTranslationCS{Duch svatý}
\dicEntry[heilahimna] \dicTerm{heila··himn|a} \dicIPA{{h}{ei}{\textlengthmark}{l}{a}{h}{\textsci}{m}{n}{a}} \dicPos{f}[1] \dicFlx{(‑u, ‑ur)}[7] \dicFieldCat{anat.} \dicDirectTranslationCS{mozková blána}
\dicEntry[heilahimnubólga] \dicTerm{heila··himnu·bólg|a} \dicIPA{{h}{ei}{\textlengthmark}{l}{a}{h}{\textsci}{m}{n}{\textscy}{\textsubring{b}}{ou}{l}{\r{g}}{a}} \dicPos{f}[1] \dicFlx{(‑u)}[5] \dicFieldCat{med.} \dicDirectTranslationCS{zánět mozkových blan, meningitida}
\dicEntry[heilahristingur] \dicTerm{heila··hrist·ing|ur} \dicIPA{{h}{ei}{\textlengthmark}{l}{a}{\textsubring{r}}{\textsci}{s}{\textsubring{d}}{i}{\ng}{\r{g}}{\textscy}{\textsubring{r}}} \dicPos{m}[6] \dicFlx{(‑s)}[9] \dicFieldCat{med.} \dicDirectTranslationCS{otřes mozku}
\dicEntry[heilakveisa] \dicTerm{heila··kveis|a} \dicIPA{{h}{ei}{\textlengthmark}{l}{a}{k\smash{\textsuperscript{h}}}{v}{ei}{s}{a}} \dicPos{f}[1] \dicFlx{(‑u, ‑ur)}[7] \dicFieldCat{med.} \dicSynonym{mígreni} \dicDirectTranslationCS{migréna}
\dicEntry[heilalínurit] \dicTerm{heila··línu·rit} \dicIPA{{h}{ei}{\textlengthmark}{l}{a}{l}{i}{n}{\textscy}{r}{\textsci}{\textsubring{d}}} \dicPos{n}[2] \dicFlx{(‑s, ‑)}[5] \dicFieldCat{med.} \dicDirectTranslationCS{elektroencefalogram}
\dicEntry[heilarit] \dicTerm{heila··rit} \dicIPA{{h}{ei}{\textlengthmark}{l}{a}{r}{\textsci}{\textsubring{d}}} \dicPos{n}[2] \dicFlx{(‑s, ‑)}[5] \dicFieldCat{med.} \dicDirectTranslationCS{elektroencefalogram}
\dicEntry[heilaskemmd] \dicTerm{heila··skemmd} \dicIPA{{h}{ei}{\textlengthmark}{l}{a}{s}{\r{\textObardotlessj}}{\textepsilon}{m}{\textsubring{d}}} \dicPos{f}[7] \dicFlx{(‑ar, ‑ir)}[1] \dicFieldCat{med.} \dicDirectTranslationCS{poškození mozku}
\dicEntry[heilaskurður] \dicTerm{heila··skurð|ur} \dicIPA{{h}{ei}{\textlengthmark}{l}{a}{s}{\r{g}}{\textscy}{r}{ð}{\textscy}{\textsubring{r}}} \dicPos{m}[10] \dicFlx{(‑ar, ‑ir)}[4] \dicFieldCat{med.} \dicDirectTranslationCS{operace mozku}
\dicEntry[heilaspuni] \dicTerm{heila··spun|i} \dicIPA{{h}{ei}{\textlengthmark}{l}{a}{s}{\textsubring{b}}{\textscy}{n}{\textsci}} \dicPos{m}[1] \dicFlx{(‑a)}[3] \dicDirectTranslationCS{představa, výtvor\,/\addthin výplod fantazie}
\dicEntry[heilaþvottur] \dicTerm{heila··þvott|ur} \dicIPA{{h}{ei}{\textlengthmark}{l}{a}{\texttheta}{v}{\textopeno}{h}{\textsubring{d}}{\textscy}{\textsubring{r}}} \dicPos{m}[6] \dicFlx{(‑s\,/\addthin ‑ar, ‑ar)}[63] \dicLangCat{přen.} \dicDirectTranslationCS{vymývání mozku}
\dicEntry[heilaæxli] \dicTerm{heila··æxli} \dicIPA{{h}{ei}{\textlengthmark}{l}{a}{a}{i}{x}{s}{\textsubring{d}}{l}{\textsci}} \dicPos{n}[2] \dicFlx{(‑s, ‑)}[14] \dicFieldCat{med.} \dicDirectTranslationCS{nádor na mozku}
\dicEntry[heilbrigði] \dicTerm{heil··brigð|i\smash{\textsuperscript{1}}} \dicIPA{{h}{ei}{l}{\textsubring{b}}{r}{\textsci}{\textbabygamma}{ð}{\textsci}} \dicPos{f}[3] \dicFlx{(‑i)}[3] \dicLink{heilbrigði\smash{\textsuperscript{2}}}
\dicEntry[heilbrigði] \dicTerm{heil··brigði\smash{\textsuperscript{2}}}\dicTerm{, heilbrigði\smash{\textsuperscript{1}}} \dicIPA{{h}\-{ei}\-{l}\-{\textsubring{b}}\-{r}\-{\textsci}\-{\textbabygamma}\-{ð}\-{\textsci}\-} \dicPos{n}[2] \dicFlx{(‑s)}[20] \dicDirectTranslationCS{zdraví} \dicExampleIS{andleg(t) heilbrigði} \dicExampleCS{duševní zdraví}
\dicEntry[heilbrigðiseftirlit] \dicTerm{heil·brigðis··eftir·lit} \dicIPA{{h}\-{ei}\-{l}\-{\textsubring{b}}\-{r}\-{\textsci}\-{\textbabygamma}\-{ð}\-{\textsci}\-{s}\-{\textepsilon}\-{f}\-{\textsubring{d}}\-{\textsci}\-{r}\-{l}\-{\textsci}\-{\textsubring{d}}\-} \dicPos{n}[2] \dicFlx{(‑s)}[2] \dicDirectTranslationCS{zdravotní dozor}
\dicEntry[heilbrigðismál] \dicTerm{heil·brigðis··mál} \dicIPA{{h}{ei}{l}{\textsubring{b}}{r}{\textsci}{\textbabygamma}{ð}{\textsci}{s}{m}{au}{\textsubring{l}}} \dicPos{n}[2] \dicFlx{(‑s, ‑)}[5] \dicDirectTranslationCS{zdravotnictví, zdravotní péče}
\begin{xtolerant}{}{1pt}
\dicEntry[heilbrigðisráðherra] \dicTerm{heil·brigðis··ráð·herr|a} \dicIPA{{h}\-{ei}\-{l}\-{\textsubring{b}}\-{r}\-{\textsci}\-{\textbabygamma}\-{ð}\-{\textsci}\-{s}\-{r}\-{au}\-{\texttheta}\-{h}\-{\textepsilon}\-{r}\-{a}\-}\addthinS\dicPos{m}[1]\addthinS\dicFlx{(‑a, ‑ar)}[17] \dicDirectTranslationCS{ministr(yně) zdravotnictví}
\end{xtolerant}
\dicEntry[heilbrigðisráðuneyti] \dicTerm{heil·brigðis··ráðu·neyti} \dicIPA{{h}\-{ei}\-{l}\-{\textsubring{b}}\-{r}\-{\textsci}\-{\textbabygamma}\-{ð}\-{\textsci}\-{s}\-{r}\-{au}\-{ð}\-{\textscy}\-{n}\-{ei}\-{\textsubring{d}}\-{\textsci}\-} \dicPos{n}[2] \dicFlx{(‑s, ‑)}[14] \dicDirectTranslationCS{ministerstvo zdravotnictví}
\dicEntry[heilbrigðisstarfsfólk] \dicTerm{heil·brigðis··starfs·fólk} \dicIPA{{h}\-{ei}\-{l}\-{\textsubring{b}}\-{r}\-{\textsci}\-{\textbabygamma}\-{ð}\-{\textsci}\-{s}\-{\textsubring{d}}\-{a}\-{r}\-{v}\-{s}\-{f}\-{ou}\-{\textsubring{l}}\-{\r{g}}\-} \dicPos{n}[2] \dicFlx{(‑s)}[2] \dicDirectTranslationCS{zdravotnický personál}
\dicEntry[heilbrigðisstofnun] \dicTerm{heil·brigðis··stofn|un} \dicIPA{{h}\-{ei}\-{l}\-{\textsubring{b}}\-{r}\-{\textsci}\-{\textbabygamma}\-{ð}\-{\textsci}\-{s}\-{\textsubring{d}}\-{\textopeno}\-{\textsubring{b}}\-{n}\-{\textscy}\-{\textsubring{n}}\-} \dicPos{f}[7] \dicFlx{(‑unar, ‑anir)}[8] \dicDirectTranslationCS{zdravotnické zařízení}
\dicEntry[heilbrigðisvottorð] \dicTerm{heil·brigðis··vott·orð} \dicIPA{{h}\-{ei}\-{l}\-{\textsubring{b}}\-{r}\-{\textsci}\-{\textbabygamma}\-{ð}\-{\textsci}\-{s}\-{v}\-{\textopeno}\-{h}\-{\textsubring{d}}\-{\textopeno}\-{r}\-{\texttheta}\-} \dicPos{n}[2] \dicFlx{(‑s, ‑)}[5] \dicDirectTranslationCS{potvrzení o~zdravotním stavu, lékařské potvrzení}
\dicEntry[heilbrigðisþjónusta] \dicTerm{heil·brigðis··þjón·ust|a} \dicIPA{{h}\-{ei}\-{l}\-{\textsubring{b}}\-{r}\-{\textsci}\-{\textbabygamma}\-{ð}\-{\textsci}\-{s}\-{\texttheta}\-{j}\-{ou}\-{n}\-{\textscy}\-{s}\-{\textsubring{d}}\-{a}\-} \dicPos{f}[1] \dicFlx{(‑u)}[5] \dicDirectTranslationCS{zdravotnické služby}
\dicEntry[heilbrigður] \dicTerm{heil··brigður} \dicsymFrequent\  \dicIPA{{h}{ei}{l}{\textsubring{b}}{r}{\textsci}{\textbabygamma}{ð}{\textscy}{\textsubring{r}}} \dicPos{adj}[2]\dicFlx{}[-1] \dicSynonym{frískur} \dicDirectTranslationCS{zdravý} \dicExampleIS{heilbrigt líf} \dicExampleCS{zdravý život}
\dicEntry[heild] \dicTerm{heild} \dicsymFrequent\  \dicIPA{{h}{ei}{l}{\textsubring{d}}} \dicPos{f}[7] \dicFlx{(‑ar, ‑ir)}[1] \dicDirectTranslationCS{celek} \dicExampleIS{í heild} \dicExampleCS{v~celku}
\dicEntry[heilda] \dicTerm{heild|a} \dicIPA{{h}{ei}{l}{\textsubring{d}}{a}} \dicPos{v}[1] \dicFlx{(‑aði)}[1] \dicFieldCat{mat.} \dicSynonym{tegra} \dicDirectTranslationCS{integrovat}
\dicEntry[heildar] \dicTerm{heildar-} \dicIPA{{h}{ei}{l}{\textsubring{d}}{a}{r}} \dicPos{predp} \dicDirectTranslationCS{celkový, úhrnný}
\dicEntry[heildaráætlun] \dicTerm{heildar··á·ætl|un} \dicIPA{{h}{ei}{l}{\textsubring{d}}{a}{r}{au}{a}{i}{h}{\textsubring{d}}{l}{\textscy}{\textsubring{n}}} \dicPos{f}[7] \dicFlx{(‑unar, ‑anir)}[8] \dicDirectTranslationCS{celkový plán}
\dicEntry[heildarfjöldi] \dicTerm{heildar··fjöld|i} \dicIPA{{h}{ei}{l}{\textsubring{d}}{a}{\textsubring{r}}{f}{j}{\oe}{l}{\textsubring{d}}{\textsci}} \dicPos{m}[1] \dicFlx{(‑a)}[3] \dicDirectTranslationCS{celkový počet, celkové množství}
\dicEntry[heildarhyggja] \dicTerm{heildar··hyggj|a} \dicIPA{{h}{ei}{l}{\textsubring{d}}{a}{\textsubring{r}}{h}{\textsci}{\r{\textObardotlessj}}{a}} \dicPos{f}[1] \dicFlx{(‑u)}[5] \textbf{1.} \dicFieldCat{pol.} \dicDirectTranslationCS{kolektivismus}  \textbf{2.} \dicFieldCat{filos.} \dicDirectTranslationCS{holismus}
\dicEntry[heildarútgáfa] \dicTerm{heildar··út·gáf|a} \dicIPA{{h}{ei}{l}{\textsubring{d}}{a}{r}{u}{\textsubring{d}}{\r{g}}{au}{v}{a}} \dicPos{f}[1] \dicFlx{(‑u, ‑ur)}[13] \dicDirectTranslationCS{sebrané spisy}
\dicEntry[heildarútgjöld] \dicTerm{heildar··út·gjöld} \dicIPA{{h}{ei}{l}{\textsubring{d}}{a}{r}{u}{\textsubring{d}}{\r{\textObardotlessj}}{\oe}{l}{\textsubring{d}}} \dicPos{n}[2] \dicFlx{pl}[9] \dicDirectTranslationCS{celkové výdaje}
\dicEntry[heildi] \dicTerm{heildi} \dicIPA{{h}{ei}{l}{\textsubring{d}}{\textsci}} \dicPos{n}[2] \dicFlx{(‑s, ‑)}[14] \dicFieldCat{mat.} \dicSynonym{tegur} \dicDirectTranslationCS{integrál}
\dicEntry[heildrænn] \dicTerm{heild··rænn} \dicIPA{{h}{ei}{l}{\textsubring{d}}{r}{a}{i}{\textsubring{d}}{\textsubring{n}}} \dicPos{adj}[7]\dicFlx{}[-1] \dicDirectTranslationCS{komplexní, celostní, holistický}
\dicEntry[heildsala] \dicTerm{heild··|sala} \dicIPA{{h}{ei}{l}{\textsubring{d}}{s}{a}{l}{a}} \dicPos{f}[1] \dicFlx{(‑sölu, ‑sölur)}[20] \dicDirectTranslationCS{velkoobchod} \dicAntonym{smásala}
\dicEntry[heildsali] \dicTerm{heild··sal|i} \dicIPA{{h}{ei}{l}{\textsubring{d}}{s}{a}{l}{\textsci}} \dicPos{m}[1] \dicFlx{(‑a, ‑ar)}[8] \dicDirectTranslationCS{velkoobchodník, velkoobchodnice}
\dicEntry[heildsölufyrirtæki] \dicTerm{heild·sölu··fyrir·tæki} \dicIPA{{h}\-{ei}\-{l}\-{\textsubring{d}}\-{s}\-{\oe}\-{l}\-{\textscy}\-{f}\-{\textsci}\-{r}\-{\textsci}\-{\textsubring{r}}\-{t\smash{\textsuperscript{h}}}\-{a}\-{i}\-{\r{\textObardotlessj}}\-{\textsci}\-} \dicPos{n}[2] \dicFlx{(‑s, ‑)}[16] \dicDirectTranslationCS{velkoobchodní firma}
\dicEntry[heildsöluverð] \dicTerm{heild·sölu··verð} \dicIPA{{h}{ei}{l}{\textsubring{d}}{s}{\oe}{l}{\textscy}{v}{\textepsilon}{r}{\texttheta}} \dicPos{n}[2] \dicFlx{(‑s)}[2] \dicDirectTranslationCS{velkoobchodní cena}
\dicEntry[heildverslun] \dicTerm{heild··versl|un} \dicIPA{{h}{ei}{l}{\textsubring{d}}{v}{\textepsilon}{\textsubring{r}}{s}{\textsubring{d}}{l}{\textscy}{\textsubring{n}}} \dicPos{f}[7] \dicFlx{(‑unar, ‑anir)}[8] \dicDirectTranslationCS{velkoobchod}
\begin{xtolerant}{}{1pt}
\dicEntry[heilhveiti] \dicTerm{heil··hveiti}\addthinS\dicIPA{{h}{ei}{l}{k\smash{\textsuperscript{h}}}{v}{ei}{\textsubring{d}}{\textsci}}\addthinS\dicPos{n}[2]\addthinS\dicFlx{(‑s)}[20]\addthinS\dicDirectTranslationCS{hrubá\addthinS/\addthinSS celozrnná mouka}
\end{xtolerant}
\dicEntry[heili] \dicTerm{heil|i} \dicsymFrequent\  \dicIPA{{h}{ei}{\textlengthmark}{l}{\textsci}} \dicPos{m}[1] \dicFlx{(‑a, ‑ar)}[1] \dicFieldCat{anat.} \dicDirectTranslationCS{mozek};  \dicPhraseIS{brjóta heilann um e‑ð} \dicDirectTranslationCS{lámat si s~(čím) hlavu};  \dicPhraseIS{fá e‑ð á heilanum} \dicDirectTranslationCS{přijít (co) na mysl};  \dicPhraseIS{vera með e‑ð á heilanum} \dicDirectTranslationCS{mít (co) na mysli}
\dicEntry[heilindi] \dicTerm{heil··indi} \dicIPA{{h}{ei}{\textlengthmark}{l}{\textsci}{n}{\textsubring{d}}{\textsci}} \dicPos{n}[2] \dicFlx{(‑s, ‑)}[14] \dicSynonym{hreinskilni} \dicDirectTranslationCS{čestnost, poctivost}
\dicEntry[heilkenni] \dicTerm{heil··kenni} \dicIPA{{h}{ei}{l}{c\smash{\textsuperscript{h}}}{\textepsilon}{n}{\textsci}} \dicPos{n}[2] \dicFlx{(‑s, ‑)}[14] \dicFieldCat{med.} \dicDirectTranslationCS{syndrom} \dicExampleIS{Downs-heilkenni} \dicExampleCS{Downův syndrom}
\dicEntry[heill] \dicTerm{heill\smash{\textsuperscript{1}}} \dicIPA{{h}{ei}{\textsubring{d}}{\textsubring{l}}} \dicPos{f}[7] \dicFlx{(‑ar, ‑ir)}[1] \dicDirectTranslationCS{štěstí, zdar};  \dicPhraseIS{óska e‑m heilla} \dicDirectTranslationCS{přát (komu) štěstí};  \dicPhraseIS{vera heillum horfinn} \dicDirectTranslationCS{přestat mít štěstí};  \dicPhraseIS{góðu heilli} \dicFlx{adv} \dicDirectTranslationCS{naštěstí, šťastnou náhodou};  \dicPhraseIS{illu heilli} \dicFlx{adv} \dicDirectTranslationCS{bohužel, při vší smůle}
\dicEntry[heill] \dicTerm{heill\smash{\textsuperscript{2}}} \dicsymFrequent\  \dicIPA{{h}{ei}{\textsubring{d}}{\textsubring{l}}} \dicPos{adj}[8]\dicFlx{}[-1] \textbf{1.} \dicSynonym{allur} \dicDirectTranslationCS{celý, úplný, kompletní};  \dicPhraseIS{heilar tölur} \dicFieldCat{mat.} \dicDirectTranslationCS{celá čísla}  \textbf{2.} \dicSynonym{hreinskilinn} \dicDirectTranslationCS{poctivý, upřímný}  \textbf{3.} \dicSynonym{heilbrigður} \dicDirectTranslationCS{zdravý};  \dicPhraseIS{vera heill heilsu} \dicDirectTranslationCS{těšit se dobrému zdraví};  \dicPhraseIS{ganga ekki heill til skógar} \dicLangCat{přen.} \dicDirectTranslationCS{churavět, být nemocný}  \textbf{4.} \dicSynonym{óbrotinn} \dicDirectTranslationCS{celý, celistvý, neporušený};  \dicPhraseIS{heill á húfi\,/\addthin hófi} \dicDirectTranslationCS{celý, nezraněný, (jsoucí) bez zranění};  \dicIdiom{heill}{ \dicPhraseIS{illu heilli}} \dicFlx{adv} \dicDirectTranslationCS{bohužel, naneštěstí, želbohu}
\dicEntry[heilla] \dicTerm{heill|a} \dicIPA{{h}{ei}{\textsubring{d}}{l}{a}} \dicPos{v}[1] \dicFlx{(‑aði)}[1] \dicFlx{acc} \dicSynonym{hrífa\smash{\textsuperscript{2}}} \dicDirectTranslationCS{okouzlit, okouzlovat, uchvátit, uchvacovat, fascinovat};  \dicIdiom{heillast}{ \dicPhraseIS{heillast}} \dicFlx{refl} \dicDirectTranslationCS{být okouzlený\,/\addthin uchvácený} \dicExampleIS{Ég heillaðist af fegurð borgarinnar.} \dicExampleCS{Byl jsem uchvácený krásou města.}
\dicEntry[heillandi] \dicTerm{heill··andi} \dicsymFrequent\  \dicIPA{{h}{ei}{\textsubring{d}}{l}{a}{n}{\textsubring{d}}{\textsci}} \dicPos{adj}[13] \dicFlx{indecl}[1] \dicDirectTranslationCS{okouzlující, fascinující} \dicExampleIS{heillandi mynd} \dicExampleCS{okouzlující obrázek}
\dicEntry[heillaósk] \dicTerm{heilla··ósk} \dicIPA{{h}{ei}{\textsubring{d}}{l}{a}{ou}{s}{\r{g}}} \dicPos{f}[7] \dicFlx{(‑ar, ‑ir)}[1] \dicSynonym*{hamingjuósk} \dicDirectTranslationCS{přání štěstí\,/\addthin zdaru}
\dicEntry[heillaráð] \dicTerm{heilla··ráð} \dicIPA{{h}{ei}{\textsubring{d}}{l}{a}{r}{au}{\texttheta}} \dicPos{n}[2] \dicFlx{(‑s, ‑)}[5] \dicDirectTranslationCS{dobrá rada, rada nad zlato}
\dicEntry[heillegur] \dicTerm{heil··legur} \dicIPA{{h}{ei}{l}{\textlengthmark}{\textepsilon}{\textbabygamma}{\textscy}{\textsubring{r}}} \dicPos{adj}[1]\dicFlx{}[-8] \dicDirectTranslationCS{(zdánlivě) neporušený, celý}
\dicEntry[heilmikill] \dicTerm{heil··mikill} \dicIPA{{h}{ei}{l}{m}{\textsci}{\r{\textObardotlessj}}{\textsci}{\textsubring{d}}{\textsubring{l}}} \dicPos{adj}[11]\dicFlx{}[-13] \dicDirectTranslationCS{nemalý, značně velký}
\dicEntry[heilnæmi] \dicTerm{heil··næmi} \dicIPA{{h}{ei}{l}{n}{a}{i}{m}{\textsci}} \dicPos{n}[2] \dicFlx{(‑s)}[20] \dicDirectTranslationCS{zdravost, zdraví (obyvatel ap.)}
\begin{xtolerant}{}{1pt}
\dicEntry[heilnæmur] \dicTerm{heil··næmur} \dicIPA{{h}{ei}{l}{n}{a}{i}{m}{\textscy}{\textsubring{r}}}\addthinS\dicPos{adj}[1]\dicFlx{}[-1]\addthinS\dicSynonym{hollur} \dicDirectTranslationCS{zdravý, zdraví prospěšný}
\end{xtolerant}
\dicEntry[heilræði] \dicTerm{heil··ræði} \dicIPA{{h}{ei}{l}{r}{a}{i}{ð}{\textsci}} \dicPos{n}[2] \dicFlx{(‑s, ‑)}[14] \dicDirectTranslationCS{dobrá rada}
\dicEntry[heilsa] \dicTerm{heils|a\smash{\textsuperscript{1}}} \dicsymFrequent\  \dicIPA{{h}{ei}{l}{s}{a}} \dicPos{f}[1] \dicFlx{(‑u)}[5] \dicDirectTranslationCS{zdraví} \dicExampleIS{Heilsan fer batnandi.} \dicExampleCS{Zdraví se zlepšuje.};  \dicPhraseIS{komast til heilsu} \dicFlx{refl} \dicDirectTranslationCS{uzdravit se};  \dicPhraseIS{missa heilsuna} \dicDirectTranslationCS{onemocnět, ochořet};  \dicPhraseIS{ná heilsuna} \dicDirectTranslationCS{nabýt zdraví, uzdravit se};  \dicPhraseIS{vera við góðu heilsu} \dicDirectTranslationCS{těšit se dobrému zdraví}
\dicEntry[heilsa] \dicTerm{heils|a\smash{\textsuperscript{2}}} \dicsymFrequent\  \dicIPA{{h}{ei}{l}{s}{a}} \dicPos{v}[1] \dicFlx{(‑aði)}[58] \dicFlx{dat} \textbf{1.} \dicSynonym{ávarpa} \dicDirectTranslationCS{(po)zdravit} \dicExampleIS{heilsa honum með handabandi} \dicExampleCS{pozdravit ho podáním ruky}  \textbf{2.} \dicPhraseIS{biðja að heilsa e‑m} \dicDirectTranslationCS{pozdravovat (koho)} \dicIndirectTranslationCS{(nechat vyřídit pozdrav)} \dicExampleIS{Ég bið innilega að heilsa henni.} \dicExampleCS{Srdečně ji pozdravuji.};  \dicPhraseIS{heilsa upp á e‑n} \dicDirectTranslationCS{zajít pozdravit (koho)}  \textbf{3.} \dicPhraseIS{því er ekki að heilsa} \dicDirectTranslationCS{bohužel ne, tak tomu není (bohužel)};  \dicIdiom{heilsast}{ \dicPhraseIS{heilsast}} \dicFlx{refl} \dicDirectTranslationCS{pozdravit se, přivítat se}; { \dicPhraseIS{e‑m heilsast}} \dicFlx{refl impers} \dicDirectTranslationCS{(kdo) se uzdravuje}
\dicEntry[heilsíða] \dicTerm{heil··síð|a} \dicIPA{{h}{ei}{l}{s}{i}{ð}{a}} \dicPos{f}[1] \dicFlx{(‑u, ‑ur)}[13] \dicDirectTranslationCS{celá strana}
\dicEntry[heilsíðugrein] \dicTerm{heil·síðu··grein} \dicIPA{{h}{ei}{l}{s}{i}{ð}{\textscy}{\r{g}}{r}{ei}{\textsubring{n}}} \dicPos{f}[4] \dicFlx{(‑ar, ‑ar)}[1] \dicDirectTranslationCS{článek na celou stranu}
\dicEntry[heilsteyptur] \dicTerm{heil··steyptur} \dicIPA{{h}{ei}{l}{s}{\textsubring{d}}{ei}{f}{\textsubring{d}}{\textscy}{\textsubring{r}}} \dicPos{adj}[1]\dicFlx{}[-10] \dicSynonym{traustur} \dicDirectTranslationCS{solidní, spolehlivý, pevný (člověk ap.)}
\dicEntry[heilsubót] \dicTerm{heilsu··|bót} \dicIPA{{h}{ei}{l}{s}{\textscy}{\textsubring{b}}{ou}{\textsubring{d}}} \dicPos{f}[8] \dicFlx{(‑bótar, ‑bætur)}[5] \dicSynonym{bati} \dicDirectTranslationCS{zlepšení zdraví, rekreace};  \dicPhraseIS{ganga sér til heilsubótar} \dicDirectTranslationCS{chodit pro zdraví}
\dicEntry[heilsufar] \dicTerm{heilsu··far} \dicIPA{{h}{ei}{l}{s}{\textscy}{f}{a}{\textsubring{r}}} \dicPos{n}[2] \dicFlx{(‑s)}[2] \dicDirectTranslationCS{zdravotní stav}
\dicEntry[heilsufræði] \dicTerm{heilsu··fræð|i} \dicIPA{{h}{ei}{l}{s}{\textscy}{f}{r}{a}{i}{ð}{\textsci}} \dicPos{f}[3] \dicFlx{(‑i)}[3] \dicDirectTranslationCS{zdravotní výchova, zdravověda}
\dicEntry[heilsugóður] \dicTerm{heilsu··|góður} \dicIPA{{h}{ei}{l}{s}{\textscy}{\r{g}}{ou}{ð}{\textscy}{\textsubring{r}}} \dicPos{adj}[11] \dicFlx{(comp ‑betri, sup ‑bestur)}[1] \dicDirectTranslationCS{zdravý, (jsoucí) v~dobrém zdravotním stavu}
\dicEntry[heilsugæsla] \dicTerm{heilsu··gæsl|a} \dicIPA{{h}{ei}{l}{s}{\textscy}{\r{\textObardotlessj}}{a}{i}{s}{\textsubring{d}}{l}{a}} \dicPos{f}[1] \dicFlx{(‑u)}[5] \textbf{1.} \dicSynonym*{heilsufarseftirlit} \dicDirectTranslationCS{zdravotní péče}  \textbf{2.} \dicSynonym{heilsugæslustöð} \dicDirectTranslationCS{zdravotní středisko, poliklinika}
\dicEntry[heilsugæslulæknir] \dicTerm{heilsu·gæslu··lækn|ir} \dicIPA{{h}\-{ei}\-{l}\-{s}\-{\textscy}\-{\r{\textObardotlessj}}\-{a}\-{i}\-{s}\-{\textsubring{d}}\-{l}\-{\textscy}\-{l}\-{a}\-{i}\-{h}\-{\r{g}}\-{n}\-{\textsci}\-{\textsubring{r}}\-} \dicPos{m}[7] \dicFlx{(‑is, ‑ar)}[1] \dicDirectTranslationCS{lékař(ka)} \dicIndirectTranslationCS{(působící v~zdravotním středisku)}
\dicEntry[heilsugæslumál] \dicTerm{heilsu·gæslu··mál} \dicIPA{{h}{ei}{l}{s}{\textscy}{\r{\textObardotlessj}}{a}{i}{s}{\textsubring{d}}{l}{\textscy}{m}{au}{\textsubring{l}}} \dicPos{n}[2] \dicFlx{(‑s, ‑)}[5] \dicDirectTranslationCS{zdravotnictví, zdravotní péče}
\dicEntry[heilsugæslustöð] \dicTerm{heilsu·gæslu··stöð} \dicIPA{{h}{ei}{l}{s}{\textscy}{\r{\textObardotlessj}}{a}{i}{s}{\textsubring{d}}{l}{\textscy}{s}{\textsubring{d}}{\oe}{\texttheta}} \dicPos{f}[6] \dicFlx{(‑var, ‑var)}[1] \dicDirectTranslationCS{zdravotní středisko, poliklinika}
\dicEntry[heilsuhæli] \dicTerm{heilsu··hæli} \dicIPA{{h}{ei}{l}{s}{\textscy}{h}{a}{i}{l}{\textsci}} \dicPos{n}[2] \dicFlx{(‑s, ‑)}[14] \dicDirectTranslationCS{sanatorium, ozdravovna}
\dicEntry[heilsulaus] \dicTerm{heilsu··laus} \dicIPA{{h}{ei}{l}{s}{\textscy}{l}{\oe i}{s}} \dicPos{adj}[5]\dicFlx{}[-1] \dicDirectTranslationCS{churavý, neduživý, stonavý}
\dicEntry[heilsuleysi] \dicTerm{heilsu··leysi} \dicIPA{{h}{ei}{l}{s}{\textscy}{l}{ei}{s}{\textsci}} \dicPos{n}[2] \dicFlx{(‑s)}[20] \dicDirectTranslationCS{churavost, neduživost, nemocnost}
\dicEntry[heilsurækt] \dicTerm{heilsu··rækt} \dicIPA{{h}{ei}{l}{s}{\textscy}{r}{a}{i}{x}{\textsubring{d}}} \dicPos{f}[4] \dicFlx{(‑ar)}[3] \dicDirectTranslationCS{péče o~zdraví}
\dicEntry[heilsutjón] \dicTerm{heilsu··tjón} \dicIPA{{h}{ei}{l}{s}{\textscy}{t\smash{\textsuperscript{h}}}{j}{ou}{\textsubring{n}}} \dicPos{n}[2] \dicFlx{(‑s)}[2] \dicDirectTranslationCS{újma\,/\addthin škoda na zdraví}
\dicEntry[heilsutæpur] \dicTerm{heilsu··tæpur} \dicIPA{{h}{ei}{l}{s}{\textscy}{t\smash{\textsuperscript{h}}}{a}{i}{\textsubring{b}}{\textscy}{\textsubring{r}}} \dicPos{adj}[1]\dicFlx{}[-1] \dicDirectTranslationCS{stonavý, churavý}
\dicEntry[heilsuvernd] \dicTerm{heilsu··vernd} \dicIPA{{h}{ei}{l}{s}{\textscy}{v}{\textepsilon}{r}{n}{\textsubring{d}}} \dicPos{f}[4] \dicFlx{(‑ar)}[3] \dicDirectTranslationCS{ochrana zdraví}
\dicEntry[heilsuvörur] \dicTerm{heilsu··vörur} \dicIPA{{h}{ei}{l}{s}{\textscy}{v}{\oe}{r}{\textscy}{\textsubring{r}}} \dicPos{f}[1] \dicFlx{pl}[9] \dicDirectTranslationCS{bio produkty}
\dicEntry[heiltala] \dicTerm{heil··|tala} \dicIPA{{h}{ei}{l}{t\smash{\textsuperscript{h}}}{a}{l}{a}} \dicPos{f}[1] \dicFlx{(‑tölu, ‑tölur)}[14] \dicFieldCat{mat.} \dicDirectTranslationCS{celé číslo}
\dicEntry[heim] \dicTerm{heim} \dicsymFrequent\  \dicIPA{{h}{ei}{\textlengthmark}{\textsubring{m}}} \dicPos{adv} \dicDirectTranslationCS{domů} \dicExampleIS{bjóða e‑m heim} \dicExampleCS{pozvat (koho) domů} \dicAntonym{heiman};  \dicPhraseIS{fara heim} \dicDirectTranslationCS{jet\,/\addthin jít domů}
\dicEntry[heima] \dicTerm{heima} \dicsymFrequent\  \dicIPA{{h}{ei}{\textlengthmark}{m}{a}} \dicPos{adv} \dicDirectTranslationCS{doma};  \dicPhraseIS{eiga heima e‑s staðar} \dicDirectTranslationCS{bydlet (kde), mít domov (kde)} \dicExampleIS{Ég á heima í Reykjavík.} \dicExampleCS{Bydlím v~Reykjavíku.};  \dicPhraseIS{vera heima hjá sér} \dicDirectTranslationCS{být (u~sebe) doma};  \dicIdiom{heima}[fyrir]{ \dicPhraseIS{heima fyrir}} \dicFlx{adv} \dicDirectTranslationCS{(u~sebe) doma};  \dicIdiom{heima}[við]{ \dicPhraseIS{heima við}} \dicFlx{adv} \dicDirectTranslationCS{(u~sebe) doma}
\dicEntry[heimabakaður] \dicTerm{heima··|bak·aður} \dicIPA{{h}{ei}{\textlengthmark}{m}{a}{\textsubring{b}}{a}{\r{g}}{a}{ð}{\textscy}{\textsubring{r}}} \dicPos{adj}[3] \dicFlx{(f ‑bökuð)}[2] \dicDirectTranslationCS{doma pečený, po domácku pečený} \dicExampleIS{heimabakaðar kökur} \dicExampleCS{doma pečené koláče}
\dicEntry[heimabanki] \dicTerm{heima··bank|i} \dicIPA{{h}{ei}{\textlengthmark}{m}{a}{\textsubring{b}}{au}{\r{\textltailn}}{\r{\textObardotlessj}}{\textsci}} \dicPos{m}[1] \dicFlx{(‑a, ‑ar)}[8] \dicFieldCat{poč.} \dicDirectTranslationCS{internetové bankovnictví}
\dicEntry[heimabíó] \dicTerm{heima··bíó} \dicIPA{{h}{ei}{\textlengthmark}{m}{a}{\textsubring{b}}{i}{j}{ou}} \dicPos{n}[2] \dicFlx{(‑s, ‑)}[34] \dicDirectTranslationCS{domácí kino}
\dicEntry[heimaborg] \dicTerm{heima··borg} \dicIPA{{h}{ei}{\textlengthmark}{m}{a}{\textsubring{b}}{\textopeno}{r}{\r{g}}} \dicPos{f}[7] \dicFlx{(‑ar, ‑ir)}[1] \dicDirectTranslationCS{rodiště, rodné město}
\dicEntry[heimabrugg] \dicTerm{heima··brugg} \dicIPA{{h}{ei}{\textlengthmark}{m}{a}{\textsubring{b}}{r}{\textscy}{\r{g}}} \dicPos{n}[2] \dicFlx{(‑s)}[2] \dicDirectTranslationCS{domácí pálenka}
\dicEntry[heimagangur] \dicTerm{heima··gang|ur} \dicIPA{{h}{ei}{\textlengthmark}{m}{a}{\r{g}}{au}{\ng}{\r{g}}{\textscy}{\textsubring{r}}} \dicPos{m}[6] \dicFlx{(‑s, ‑ar)}[5] \dicDirectTranslationCS{častý návštěvník\,/\addthin host, častá návštěvnice\,/\addthin hostka}
\dicEntry[heimahjúkrun] \dicTerm{heima··hjúkr|un} \dicIPA{{h}{ei}{\textlengthmark}{m}{a}{\c{c}}{u}{\r{g}}{r}{\textscy}{\textsubring{n}}} \dicPos{f}[7] \dicFlx{(‑unar)}[9] \dicDirectTranslationCS{domácí (zdravotní) péče}
\dicEntry[heimakominn] \dicTerm{heima··kominn} \dicIPA{{h}{ei}{\textlengthmark}{m}{a}{k\smash{\textsuperscript{h}}}{\textopeno}{m}{\textsci}{\textsubring{n}}} \dicPos{adj}[6]\dicFlx{}[-2] \dicPhraseIS{gera sig heimakominn hjá e‑m} \dicDirectTranslationCS{cítit se jako doma u~(koho), udělat si pohodlí (jako doma) u~(koho)}
\dicEntry[heimaland] \dicTerm{heima··|land} \dicIPA{{h}{ei}{\textlengthmark}{m}{a}{l}{a}{n}{\textsubring{d}}} \dicPos{n}[2] \dicFlx{(‑lands, ‑lönd)}[8] \dicDirectTranslationCS{rodná země, vlast, domovina}
\dicEntry[heimamaður] \dicTerm{heima··|maður} \dicIPA{{h}{ei}{\textlengthmark}{m}{a}{m}{a}{ð}{\textscy}{\textsubring{r}}} \dicPos{m}[13] \dicFlx{(‑manns, ‑menn)}[2] \dicDirectTranslationCS{rodák, rodačka}
\dicEntry[heimamarkaður] \dicTerm{heima··mark·að|ur} \dicIPA{{h}{ei}{\textlengthmark}{m}{a}{m}{a}{\textsubring{r}}{\r{g}}{a}{ð}{\textscy}{\textsubring{r}}} \dicPos{m}[10] \dicFlx{(‑ar\,/\addthin ‑s, ‑ir)}[31] \dicFieldCat{ekon.} \dicDirectTranslationCS{domácí trh}
\dicEntry[heiman] \dicTerm{heiman} \dicsymFrequent\  \dicIPA{{h}{ei}{\textlengthmark}{m}{a}{\textsubring{n}}} \dicPos{adv} \dicDirectTranslationCS{z~domu\,/\addthin domova} \dicExampleIS{fara að heiman} \dicExampleCS{odejít z~domu} \dicAntonym{heim};  \dicIdiom{að}[heiman]{ \dicPhraseIS{að heiman}} \dicFlx{adv} {\textbf{a.}} \dicDirectTranslationCS{z~domu};  {\textbf{b.}} \dicDirectTranslationCS{mimo dům, (pryč) z~domu};  \dicIdiom{heiman}[að]{ \dicPhraseIS{heiman að}} \dicFlx{adv} \dicDirectTranslationCS{z~domu\,/\addthin domova}
\dicEntry[heimanfylgja] \dicTerm{heiman··fylgj|a} \dicIPA{{h}{ei}{m}{a}{n}{f}{\textsci}{l}{\r{\textObardotlessj}}{a}} \dicPos{f}[1] \dicFlx{(‑u, ‑ur)}[17] \dicLangCat{zast.} \dicDirectTranslationCS{věno}
\dicEntry[heimasíða] \dicTerm{heima··síð|a} \dicIPA{{h}{ei}{\textlengthmark}{m}{a}{s}{i}{ð}{a}} \dicPos{f}[1] \dicFlx{(‑u, ‑ur)}[13] \dicFieldCat{poč.} \dicDirectTranslationCS{domovská stránka} \dicIndirectTranslationCS{(na internetu)}
\dicEntry[heimasími] \dicTerm{heima··sím|i} \dicIPA{{h}{ei}{\textlengthmark}{m}{a}{s}{i}{m}{\textsci}} \dicPos{m}[1] \dicFlx{(‑a, ‑ar)}[1] \dicDirectTranslationCS{domácí telefon} \dicExampleIS{Hann gaf upp vinnusíma og heimasíma sinn.} \dicExampleCS{Uvedl telefon do práce a~telefon domů.}
\dicEntry[heimasæta] \dicTerm{heima··sæt|a} \dicIPA{{h}{ei}{\textlengthmark}{m}{a}{s}{a}{i}{\textsubring{d}}{a}} \dicPos{f}[1] \dicFlx{(‑u, ‑ur)}[19] \dicDirectTranslationCS{neprovdaná dívka žijící s~rodiči}
\dicEntry[heimatak] \dicTerm{heima··|tak} \dicIPA{{h}{ei}{\textlengthmark}{m}{a}{t\smash{\textsuperscript{h}}}{a}{\r{g}}} \dicPos{n}[2] \dicFlx{(‑taks, ‑tök)}[8] \dicDirectTranslationCS{hračka, lehká práce, snadná věc} \dicExampleIS{Það er ekkert heimatak.} \dicExampleCS{Není to hračka.}
\dicEntry[heimaverkefni] \dicTerm{heima··verk·efni} \dicIPA{{h}{ei}{\textlengthmark}{m}{a}{v}{\textepsilon}{\textsubring{r}}{\r{g}}{\textepsilon}{\textsubring{b}}{n}{\textsci}} \dicPos{n}[2] \dicFlx{(‑s, ‑)}[14] \dicDirectTranslationCS{domácí úkol} \dicExampleIS{skila heimaverkefni} \dicExampleCS{odevzdat domácí úkol}
\dicEntry[heimavinna] \dicTerm{heima··vinn|a} \dicIPA{{h}{ei}{\textlengthmark}{m}{a}{v}{\textsci}{n}{a}} \dicPos{f}[1] \dicFlx{(‑u)}[5] \dicDirectTranslationCS{domácí práce}
\dicEntry[heimavist] \dicTerm{heima··vist} \dicIPA{{h}{ei}{\textlengthmark}{m}{a}{v}{\textsci}{s}{\textsubring{d}}} \dicPos{f}[7] \dicFlx{(‑ar, ‑ir)}[1] \dicFieldCat{škol.} \dicDirectTranslationCS{internát};  \dicPhraseIS{búa í heimavist} \dicDirectTranslationCS{bydlet na internátu}
\dicEntry[heimavistarskóli] \dicTerm{heima·vistar··skól|i} \dicIPA{{h}{ei}{\textlengthmark}{m}{a}{v}{\textsci}{s}{\textsubring{d}}{a}{\textsubring{r}}{s}{\r{g}}{ou}{l}{\textsci}} \dicPos{m}[1] \dicFlx{(‑a, ‑ar)}[1] \dicFieldCat{škol.} \dicDirectTranslationCS{internátní škola}
\dicEntry[heimavöllur] \dicTerm{heima··|völlur} \dicIPA{{h}{ei}{\textlengthmark}{m}{a}{v}{\oe}{\textsubring{d}}{l}{\textscy}{\textsubring{r}}} \dicPos{m}[11] \dicFlx{(‑vallar, ‑vellir)}[5] \dicDirectTranslationCS{domácí stadión\,/\addthin hřiště} \dicExampleIS{keppa á heimavelli} \dicExampleCS{hrát na domácí půdě};  \dicPhraseIS{vera á heimavelli} \dicLangCat{přen.} \dicDirectTranslationCS{být ve svém živlu}
\dicEntry[heimboð] \dicTerm{heim··boð} \dicIPA{{h}{ei}{m}{\textsubring{b}}{\textopeno}{\texttheta}} \dicPos{n}[2] \dicFlx{(‑s, ‑)}[5] \dicDirectTranslationCS{pozvání domů (na návštěvu)}
\dicEntry[Heimdallur] \dicTerm{Heim··dall|ur} \dicIPA{{h}{ei}{m}{\textsubring{d}}{a}{\textsubring{d}}{l}{\textscy}{\textsubring{r}}} \dicPos{m}[10] \dicFlx{(‑ar\,/\addthin ‑s)}[28] \dicFieldCat{myt.} \dicDirectTranslationCS{Heimdall} \dicIndirectTranslationCS{(strážce Bifröstu a~Ásgardu)}
\dicEntry[heimferð] \dicTerm{heim··ferð} \dicIPA{{h}{ei}{m}{f}{\textepsilon}{r}{\texttheta}} \dicPos{f}[7] \dicFlx{(‑ar, ‑ir)}[1] \dicDirectTranslationCS{cesta domů}
\dicEntry[heimfæra] \dicTerm{heim··fær|a} \dicIPA{{h}{ei}{m}{f}{a}{i}{r}{a}} \dicPos{v}[2] \dicFlx{(‑ði, ‑t)}[99] \dicFlx{acc} \dicSynonym*{telja til} \dicDirectTranslationCS{přenést, přenášet} \dicIndirectTranslationCS{(učinit někoho jiného nositelem nějaké vlastnosti, práva ap.)};  \dicPhraseIS{heimfæra e‑ð upp á e‑ð} \dicDirectTranslationCS{přenášet (co) na (co) (okolnosti ap.)}
\dicEntry[heimför] \dicTerm{heim··|för} \dicIPA{{h}{ei}{m}{f}{\oe}{\textsubring{r}}} \dicPos{f}[7] \dicFlx{(‑farar)}[19] \dicSynonym{heimferð} \dicDirectTranslationCS{cesta domů}
\dicEntry[heimila] \dicTerm{heimil|a} \dicIPA{{h}{ei}{\textlengthmark}{m}{\textsci}{l}{a}} \dicPos{v}[1] \dicFlx{(‑aði)}[1] \dicFlx{dat + acc} \dicDirectTranslationCS{dovolit, povolit, umožnit, udělit povolení};  \dicPhraseIS{heimila e‑m e‑ð} \dicDirectTranslationCS{dovolit (komu co)}
\dicEntry[heimild] \dicTerm{heimild} \dicsymFrequent\  \dicIPA{{h}{ei}{\textlengthmark}{m}{\textsci}{l}{\textsubring{d}}} \dicPos{f}[7] \dicFlx{(‑ar, ‑ir)}[1] \textbf{1.} \dicSynonym{leyfi} \dicDirectTranslationCS{povolení, dovolení, svolení};  \dicPhraseIS{heimild fyrir e‑ð} \dicDirectTranslationCS{povolení na (co)};  \dicPhraseIS{heimild til e‑s} \dicDirectTranslationCS{povolení k~(čemu)} \dicExampleIS{heimild til framkvæmdar} \dicExampleCS{povolení k~realizaci};  \dicPhraseIS{veita e‑m heimild} \dicDirectTranslationCS{udělit (komu) povolení}  \textbf{2.} \dicDirectTranslationCS{zdroj\,/\addthin pramen (informací), evidence} \dicExampleIS{rannsaka heimildirnar} \dicExampleCS{prozkoumat zdroje};  \dicIdiom{heimild}{ \dicPhraseIS{villa á sér heimildir}} \dicDirectTranslationCS{vydávat se za někoho jiného}
\dicEntry[heimildarkvikmynd] \dicTerm{heimildar··kvik·mynd} \dicIPA{{h}{ei}{\textlengthmark}{m}{\textsci}{l}{\textsubring{d}}{a}{\textsubring{r}}{k\smash{\textsuperscript{h}}}{v}{\textsci}{\r{g}}{m}{\textsci}{n}{\textsubring{d}}} \dicPos{f}[7] \dicFlx{(‑ar, ‑ir)}[1] \dicDirectTranslationCS{dokument, dokumentární film}
\dicEntry[heimildarmaður] \dicTerm{heimildar··|maður} \dicIPA{{h}{ei}{\textlengthmark}{m}{\textsci}{l}{\textsubring{d}}{a}{r}{m}{a}{ð}{\textscy}{\textsubring{r}}} \dicPos{m}[13] \dicFlx{(‑manns, ‑menn)}[2] \dicDirectTranslationCS{zdroj (informací), informátor(ka)}
\dicEntry[heimildarmynd] \dicTerm{heimildar··mynd} \dicIPA{{h}{ei}{\textlengthmark}{m}{\textsci}{l}{\textsubring{d}}{a}{r}{m}{\textsci}{n}{\textsubring{d}}} \dicPos{f}[7] \dicFlx{(‑ar, ‑ir)}[1] \dicDirectTranslationCS{dokument, dokumentární film}
\dicEntry[heimildaskrá] \dicTerm{heimilda··skrá} \dicIPA{{h}{ei}{\textlengthmark}{m}{\textsci}{l}{\textsubring{d}}{a}{s}{\r{g}}{r}{au}} \dicPos{f}[4] \dicFlx{(‑r\,/\addthin ‑ar, ‑r)}[21] \dicDirectTranslationCS{seznam pramenů, bibliografie}
\dicEntry[heimili] \dicTerm{heimili} \dicsymFrequent\  \dicIPA{{h}{ei}{\textlengthmark}{m}{\textsci}{l}{\textsci}} \dicPos{n}[2] \dicFlx{(‑s, ‑)}[14] \textbf{1.} \dicDirectTranslationCS{domov, bydliště};  \dicPhraseIS{eiga heimili e‑s staðar} \dicDirectTranslationCS{bydlet (kde)};  \dicPhraseIS{vera til heimilis e‑s staðar} \dicDirectTranslationCS{bydlet (kde), mít bydliště (kde)}  \textbf{2.} \dicDirectTranslationCS{domácnost} \dicExampleIS{Það eru fimm menn í heimili.} \dicExampleCS{V~domácnosti je pět lidí.}
\dicEntry[heimilisfang] \dicTerm{heimilis··|fang} \dicsymFrequent\  \dicIPA{{h}{ei}{\textlengthmark}{m}{\textsci}{l}{\textsci}{s}{f}{au}{\ng}{\r{g}}} \dicPos{n}[2] \dicFlx{(‑fangs, ‑föng)}[8] \dicDirectTranslationCS{adresa (bydliště)} \dicExampleIS{biðja hana um heimilisfang hennar} \dicExampleCS{požádat ji o~její adresu}
\dicEntry[heimilisfólk] \dicTerm{heimilis··fólk} \dicIPA{{h}{ei}{\textlengthmark}{m}{\textsci}{l}{\textsci}{s}{f}{ou}{\textsubring{l}}{\r{g}}} \dicPos{n}[2] \dicFlx{(‑s)}[2] \dicDirectTranslationCS{rodina, domácnost, obyvatelé (domu)}
\dicEntry[heimilisfræði] \dicTerm{heimilis··fræð|i} \dicIPA{{h}{ei}{\textlengthmark}{m}{\textsci}{l}{\textsci}{s}{f}{r}{a}{i}{ð}{\textsci}} \dicPos{f}[3] \dicFlx{(‑i)}[3] \dicDirectTranslationCS{nauka o~vedení domácnosti}
\dicEntry[heimilishjálp] \dicTerm{heimilis··hjálp} \dicIPA{{h}{ei}{\textlengthmark}{m}{\textsci}{l}{\textsci}{s}{\c{c}}{au}{\textsubring{l}}{\textsubring{b}}} \dicPos{f}[7] \dicFlx{(‑ar)}[3] \textbf{1.} \dicSynonym*{aðstoð á heimili} \dicDirectTranslationCS{pomoc v~domácnosti} \dicIndirectTranslationCS{(sociální služba)}  \textbf{2.} \dicSynonym{hjálparmaður} \dicDirectTranslationCS{pomocník\,/\addthin pomocnice v~domácnosti}
\dicEntry[heimilisiðnaður] \dicTerm{heimilis··iðn·að|ur} \dicIPA{{h}{ei}{\textlengthmark}{m}{\textsci}{l}{\textsci}{s}{\textsci}{ð}{n}{a}{ð}{\textscy}{\textsubring{r}}} \dicPos{m}[10] \dicFlx{(‑ar)}[9] \dicDirectTranslationCS{domácký průmysl}
\dicEntry[heimilislaus] \dicTerm{heimilis··laus} \dicIPA{{h}{ei}{\textlengthmark}{m}{\textsci}{l}{\textsci}{s}{l}{\oe i}{s}} \dicPos{adj}[5]\dicFlx{}[-4] \dicDirectTranslationCS{(jsoucí) bez domova}
\dicEntry[heimilislæknir] \dicTerm{heimilis··lækn|ir} \dicIPA{{h}{ei}{\textlengthmark}{m}{\textsci}{l}{\textsci}{s}{l}{a}{i}{h}{\r{g}}{n}{\textsci}{\textsubring{r}}} \dicPos{m}[7] \dicFlx{(‑is, ‑ar)}[1] \textbf{1.} \dicSynonym*{læknir fjölskyldunnar} \dicDirectTranslationCS{rodinný lékař, rodinná lékařka}  \textbf{2.} \dicSynonym*{sérfræðingur í heimilislækningum} \dicDirectTranslationCS{praktický lékař, praktická lékařka}
\dicEntry[heimilistæki] \dicTerm{heimilis··tæki} \dicIPA{{h}{ei}{\textlengthmark}{m}{\textsci}{l}{\textsci}{s}{t\smash{\textsuperscript{h}}}{a}{i}{\r{\textObardotlessj}}{\textsci}} \dicPos{n}[2] \dicFlx{(‑s, ‑)}[16] \dicDirectTranslationCS{domácí spotřebič}
\dicEntry[heimilisuppbót] \dicTerm{heimilis··upp·|bót} \dicIPA{{h}{ei}{\textlengthmark}{m}{\textsci}{l}{\textsci}{s}{\textscy}{h}{\textsubring{b}}{ou}{\textsubring{d}}} \dicPos{f}[8] \dicFlx{(‑bótar, ‑bætur)}[5] \dicDirectTranslationCS{příspěvek na domácnost}
\dicEntry[heimill] \dicTerm{heimill} \dicIPA{{h}{ei}{\textlengthmark}{m}{\textsci}{\textsubring{d}}{\textsubring{l}}} \dicPos{adj}[8]\dicFlx{}[-3] \dicSynonym{leyfilegur} \dicDirectTranslationCS{povolený, dovolený, přípustný};  \dicPhraseIS{e‑m er e‑að heimilt} \dicDirectTranslationCS{(co) je (komu) povoleno} \dicExampleIS{Aðgangur er öllum heimill.} \dicExampleCS{Vstup je všem povolen.};  \dicPhraseIS{það er heimilt að (gera e‑ð)} \dicDirectTranslationCS{je povoleno (dělat (co))}
\dicEntry[heimkynni] \dicTerm{heim··kynni} \dicIPA{{h}{ei}{m}{c\smash{\textsuperscript{h}}}{\textsci}{n}{\textsci}} \dicPos{n}[2] \dicFlx{pl}[19] \dicDirectTranslationCS{domov, rodiště}
\dicEntry[heimleið] \dicTerm{heim··leið} \dicIPA{{h}{ei}{m}{l}{ei}{\texttheta}} \dicPos{f}[7] \dicFlx{(‑ar)}[3] \dicDirectTranslationCS{cesta domů, zpáteční cesta} \dicExampleIS{á heimleiðinni} \dicExampleCS{na cestě domů}
\dicEntry[heimsálfa] \dicTerm{heims··álf|a} \dicIPA{{h}{ei}{m}{s}{au}{l}{v}{a}} \dicPos{f}[1] \dicFlx{(‑u, ‑ur)}[7] \dicDirectTranslationCS{kontinent, světadíl}
\dicEntry[heimsborgari] \dicTerm{heims··borg·ar|i} \dicIPA{{h}{ei}{m}{s}{\textsubring{b}}{\textopeno}{r}{\r{g}}{a}{r}{\textsci}} \dicPos{m}[1] \dicFlx{(‑a, ‑ar)}[13] \dicDirectTranslationCS{světoobčan(ka), kosmopolita, kosmopolitka}
\dicEntry[heimsendir] \dicTerm{heims··end|ir} \dicIPA{{h}{ei}{m}{s}{\textepsilon}{n}{\textsubring{d}}{\textsci}{\textsubring{r}}} \dicPos{m}[7] \dicFlx{(‑is)}[2] \dicDirectTranslationCS{konec světa, soudný den, apokalypsa}
\dicEntry[heimsfaraldur] \dicTerm{heims··far·ald|ur} \dicIPA{{h}{ei}{m}{s}{f}{a}{r}{a}{l}{\textsubring{d}}{\textscy}{\textsubring{r}}} \dicPos{m}[5] \dicFlx{(‑urs, ‑rar)}[3] \dicFieldCat{med.} \dicDirectTranslationCS{pandemie}
\dicEntry[heimsfræði] \dicTerm{heims··fræð|i} \dicIPA{{h}{ei}{m}{s}{f}{r}{a}{i}{ð}{\textsci}} \dicPos{f}[3] \dicFlx{(‑i)}[3] \dicDirectTranslationCS{kosmologie}
\dicEntry[heimsfrægur] \dicTerm{heims··frægur} \dicIPA{{h}{ei}{m}{s}{f}{r}{a}{i}{\textbabygamma}{\textscy}{\textsubring{r}}} \dicPos{adj}[1]\dicFlx{}[-1] \dicDirectTranslationCS{světoznámý, světově proslulý}
\dicEntry[heimshlýnun] \dicTerm{heims··hlýn|un} \dicIPA{{h}{ei}{m}{s}{\textsubring{l}}{i}{n}{\textscy}{\textsubring{n}}} \dicPos{f}[7] \dicFlx{(‑unar)}[9] \dicFieldCat{meteo.} \dicDirectTranslationCS{globální oteplování}
\dicEntry[heimshorn] \dicTerm{heims··horn} \dicIPA{{h}{ei}{m}{s}{h}{\textopeno}{r}{\textsubring{d}}{\textsubring{n}}} \dicPos{n}[2] \dicFlx{(‑s, ‑)}[5] \dicDirectTranslationCS{kout světa};  \dicPhraseIS{frá öllum heimshornum} \dicFlx{adv} \dicDirectTranslationCS{ze všech koutů světa}
\dicEntry[heimska] \dicTerm{heimsk|a} \dicIPA{{h}{ei}{m}{s}{\r{g}}{a}} \dicPos{f}[1] \dicFlx{(‑u)}[5] \dicSynonym{asnaskapur} \dicDirectTranslationCS{hloupost, stupidita}
\dicEntry[heimskaut] \dicTerm{heim··skaut} \dicIPA{{h}{ei}{m}{s}{\r{g}}{\oe i}{\textsubring{d}}} \dicPos{n}[2] \dicFlx{(‑s, ‑)}[5] \dicFieldCat{geog.} \dicSynonym{póll} \dicDirectTranslationCS{pól}
\dicEntry[heimskautsbaugur] \dicTerm{heim·skauts··baug|ur} \dicIPA{{h}{ei}{m}{s}{\r{g}}{\oe i}{\textsubring{d}}{s}{\textsubring{b}}{\oe i}{\textbabygamma}{\textscy}{\textsubring{r}}} \dicPos{m}[6] \dicFlx{(‑s, ‑ar)}[24] \dicFieldCat{geog.} \dicDirectTranslationCS{polární kruh}
\dicEntry[heimskingi] \dicTerm{heimsk··ing|i} \dicIPA{{h}{ei}{m}{s}{\r{\textObardotlessj}}{i}{\textltailn}{\r{\textObardotlessj}}{\textsci}} \dicPos{m}[1] \dicFlx{(‑ja, ‑jar)}[14] \dicSynonym{bjáni} \dicDirectTranslationCS{hlupák, hlupačka, blbeček}
\dicEntry[heimskulegur] \dicTerm{heimsku··legur} \dicIPA{{h}{ei}{m}{s}{\r{g}}{\textscy}{l}{\textepsilon}{\textbabygamma}{\textscy}{\textsubring{r}}} \dicPos{adj}[1]\dicFlx{}[-8] \dicDirectTranslationCS{pitomý, hloupý, stupidní}
\dicEntry[heimskupar] \dicTerm{heimsku··|par} \dicIPA{{h}{ei}{m}{s}{\r{g}}{\textscy}{p\smash{\textsuperscript{h}}}{a}{\textsubring{r}}} \dicPos{n}[2] \dicFlx{(‑pars, ‑pör)}[8] \dicDirectTranslationCS{hloupost, pitomost} \dicIndirectTranslationCS{(zachovat se hloupě)}
\dicEntry[heimskur] \dicTerm{heimskur} \dicsymFrequent\  \dicIPA{{h}{ei}{m}{s}{\r{g}}{\textscy}{\textsubring{r}}} \dicPos{adj}[1]\dicFlx{}[-1] \dicDirectTranslationCS{hloupý, pitomý, blbý} \dicExampleIS{Hann er alltof heimskur til þess.} \dicExampleCS{Je na to příliš hloupý.}
\dicEntry[heimsmaður] \dicTerm{heims··|maður} \dicIPA{{h}{ei}{m}{s}{m}{a}{ð}{\textscy}{\textsubring{r}}} \dicPos{m}[13] \dicFlx{(‑manns, ‑menn)}[2] \dicDirectTranslationCS{světoobčan(ka), světoběžník, světoběžnice}
\dicEntry[heimsmarkaðsverð] \dicTerm{heims·markaðs··verð} \dicIPA{{h}{ei}{m}{s}{m}{a}{\textsubring{r}}{\r{g}}{a}{ð}{s}{v}{\textepsilon}{r}{\texttheta}} \dicPos{n}[2] \dicFlx{(‑s)}[2] \dicFieldCat{ekon.} \dicDirectTranslationCS{cena na světovém trhu}
\dicEntry[heimsmarkaður] \dicTerm{heims··mark·að|ur} \dicIPA{{h}{ei}{m}{s}{m}{a}{\textsubring{r}}{\r{g}}{a}{ð}{\textscy}{\textsubring{r}}} \dicPos{m}[10] \dicFlx{(‑ar\,/\addthin ‑s, ‑ir)}[31] \dicFieldCat{ekon.} \dicDirectTranslationCS{světový trh}
\dicEntry[heimsmeistari] \dicTerm{heims··meist·ar|i} \dicIPA{{h}{ei}{m}{s}{m}{ei}{s}{\textsubring{d}}{a}{r}{\textsci}} \dicPos{m}[1] \dicFlx{(‑a, ‑ar)}[13] \dicFieldCat{sport.} \dicDirectTranslationCS{mistr(yně) světa}
\dicEntry[heimsmet] \dicTerm{heims··met} \dicIPA{{h}{ei}{m}{s}{m}{\textepsilon}{\textsubring{d}}} \dicPos{n}[2] \dicFlx{(‑s, ‑)}[5] \dicFieldCat{sport.} \dicDirectTranslationCS{světový rekord};  \dicPhraseIS{setja heimsmet} \dicDirectTranslationCS{vytvořit světový rekord}
\dicEntry[heimsmynd] \dicTerm{heims··mynd} \dicIPA{{h}{ei}{m}{s}{m}{\textsci}{n}{\textsubring{d}}} \dicPos{f}[7] \dicFlx{(‑ar, ‑ir)}[1] \dicDirectTranslationCS{obraz světa} \dicIndirectTranslationCS{(idea)}
\dicEntry[heimsmælikvarði] \dicTerm{heims··mæli·kvarð|i} \dicIPA{{h}{ei}{m}{s}{m}{a}{i}{l}{\textsci}{k\smash{\textsuperscript{h}}}{v}{a}{r}{ð}{\textsci}} \dicPos{m}[1] \dicFlx{(‑a)}[3] \dicDirectTranslationCS{světové měřítko}
\dicEntry[heimsókn] \dicTerm{heim··sókn} \dicsymFrequent\  \dicIPA{{h}{ei}{m}{s}{ou}{h}{\r{g}}{\textsubring{n}}} \dicPos{f}[7] \dicFlx{(‑ar, ‑ir)}[1] \dicDirectTranslationCS{návštěva};  \dicPhraseIS{fara í heimsókn til e‑rs} \dicDirectTranslationCS{jít ke (komu) na návštěvu};  \dicPhraseIS{fá e‑n í heimsókn} \dicDirectTranslationCS{mít (koho) na návštěvě};  \dicPhraseIS{koma í heimsókn} \dicDirectTranslationCS{přijít na návštěvu}
\dicEntry[heimsóknartími] \dicTerm{heim·sóknar··tím|i} \dicIPA{{h}{ei}{m}{s}{ou}{h}{\r{g}}{n}{a}{\textsubring{r}}{t\smash{\textsuperscript{h}}}{i}{m}{\textsci}} \dicPos{m}[1] \dicFlx{(‑a, ‑ar)}[1] \dicDirectTranslationCS{návštěvní hodiny}
\dicEntry[heimsótt] \dicTerm{heim··sótt} \dicIPA{{h}{ei}{m}{s}{ou}{h}{\textsubring{d}}} \dicPos{v} \dicFlx{supin} \dicLink{heimsækja}
\dicEntry[heimsótti] \dicTerm{heim··sótti} \dicIPA{{h}{ei}{m}{s}{ou}{h}{\textsubring{d}}{\textsci}} \dicPos{v} \dicFlx{ind pf sg 1 pers} \dicLink{heimsækja}
\dicEntry[heimsóttum] \dicTerm{heim··sóttum} \dicIPA{{h}{ei}{m}{s}{ou}{h}{\textsubring{d}}{\textscy}{\textsubring{m}}} \dicPos{v} \dicFlx{ind pf pl 1 pers} \dicLink{heimsækja}
\dicEntry[heimspeki] \dicTerm{heim··spek|i} \dicIPA{{h}{ei}{m}{s}{\textsubring{b}}{\textepsilon}{\r{\textObardotlessj}}{\textsci}} \dicPos{f}[3] \dicFlx{(‑i)}[3] \dicFieldCat{filos.} \dicDirectTranslationCS{filosofie} \dicIndirectTranslationCS{(věda)}
\dicEntry[heimspekideild] \dicTerm{heim·speki··deild} \dicIPA{{h}{ei}{m}{s}{\textsubring{b}}{\textepsilon}{\r{\textObardotlessj}}{\textsci}{\textsubring{d}}{ei}{l}{\textsubring{d}}} \dicPos{f}[7] \dicFlx{(‑ar, ‑ir)}[1] \dicFieldCat{škol.} \dicDirectTranslationCS{filosofická fakulta}
\dicEntry[heimspekilegur] \dicTerm{heim·speki··legur} \dicIPA{{h}{ei}{m}{s}{\textsubring{b}}{\textepsilon}{\r{\textObardotlessj}}{\textsci}{l}{\textepsilon}{\textbabygamma}{\textscy}{\textsubring{r}}} \dicPos{adj}[1]\dicFlx{}[-8] \dicDirectTranslationCS{filosofický}
\dicEntry[heimsstyrjöld] \dicTerm{heims·styrj··|öld} \dicsymFrequent\  \dicIPA{{h}{ei}{m}{s}{\textsubring{d}}{\textsci}{r}{j}{\oe}{l}{\textsubring{d}}} \dicPos{f}[7] \dicFlx{(‑aldar, ‑aldir)}[16] \dicDirectTranslationCS{světová válka} \dicExampleIS{fyrri heimsstyrjöld} \dicExampleCS{první světová válka}
\dicEntry[heimsvaldasinni] \dicTerm{heims·valda··sinn|i} \dicIPA{{h}{ei}{m}{s}{v}{a}{l}{\textsubring{d}}{a}{s}{\textsci}{n}{\textsci}} \dicPos{m}[1] \dicFlx{(‑a, ‑ar)}[1] \dicDirectTranslationCS{imperialista, imperialistka}
\dicEntry[heimsvaldastefna] \dicTerm{heims·valda··stefn|a} \dicIPA{{h}{ei}{m}{s}{v}{a}{l}{\textsubring{d}}{a}{s}{\textsubring{d}}{\textepsilon}{\textsubring{b}}{n}{a}} \dicPos{f}[1] \dicFlx{(‑u)}[5] \dicFieldCat{pol.} \dicDirectTranslationCS{imperialismus}
\dicEntry[heimsveldi] \dicTerm{heims··veldi} \dicIPA{{h}{ei}{m}{s}{v}{\textepsilon}{l}{\textsubring{d}}{\textsci}} \dicPos{n}[2] \dicFlx{(‑s, ‑)}[14] \dicDirectTranslationCS{impérium, říše} \dicExampleIS{breska heimsveldið} \dicExampleCS{britské impérium}
\dicEntry[heimsvísa] \dicTerm{heims··vís|a} \dicIPA{{h}{ei}{m}{s}{v}{i}{s}{a}} \dicPos{f}[1] \dicFlx{(‑u)}[6] \dicPhraseIS{á heimsvísu} \dicFlx{adv} \dicDirectTranslationCS{globálně, všeobecně}
\dicEntry[heimsæki] \dicTerm{heim··sæki} \dicIPA{{h}{ei}{m}{s}{a}{i}{\r{\textObardotlessj}}{\textsci}} \dicPos{v} \dicFlx{ind praes sg 1 pers} \dicLink{heimsækja}
\dicEntry[heimsækja] \dicTerm{heim··|sækja} \dicsymFrequent\  \dicIPA{{h}{ei}{m}{s}{a}{i}{\r{\textObardotlessj}}{a}} \dicPos{v}[5] \dicFlx{(‑sæki, ‑sótti, ‑sóttum, ‑sækti, ‑sótt)}[7] \dicFlx{acc} \dicDirectTranslationCS{navštívit, navštěvovat} \dicExampleIS{heimsækja hann á spítalann} \dicExampleCS{navštívit ho v~nemocnici}
\dicEntry[heimsækti] \dicTerm{heim··sækti} \dicIPA{{h}{ei}{m}{s}{a}{i}{x}{\textsubring{d}}{\textsci}} \dicPos{v} \dicFlx{con pf sg 1 pers} \dicLink{heimsækja}
\dicEntry[heimta] \dicTerm{heimt|a\smash{\textsuperscript{1}}} \dicsymFrequent\  \dicIPA{{h}{ei}{\textsubring{m}}{\textsubring{d}}{a}} \dicPos{v}[1] \dicFlx{(‑aði)}[1] \dicFlx{acc} \dicDirectTranslationCS{požadovat, vyžadovat, vymáhat, dožadovat se} \dicExampleIS{heimta þjónustu af e‑m} \dicExampleCS{dožadovat se (čích) služeb}
\dicEntry[heimta] \dicTerm{heimt|a\smash{\textsuperscript{2}}} \dicIPA{{h}{ei}{\textsubring{m}}{\textsubring{d}}{a}} \dicPos{v}[2] \dicFlx{(‑i, ‑)}[12] \dicFlx{acc} \dicDirectTranslationCS{dostat zpět (půjčené peníze ap.)}
\dicEntry[heimting] \dicTerm{heimt··ing} \dicIPA{{h}{ei}{\textsubring{m}}{\textsubring{d}}{i}{\ng}{\r{g}}} \dicPos{f}[4] \dicFlx{(‑ar, ‑ar)}[5] \dicSynonym{krafa} \dicDirectTranslationCS{požadavek, nárok};  \dicPhraseIS{eiga heimting á e‑u} \dicDirectTranslationCS{dožadovat se (čeho)}
\dicEntry[heimtufrekja] \dicTerm{heimtu··frekj|a} \dicIPA{{h}{ei}{\textsubring{m}}{\textsubring{d}}{\textscy}{f}{r}{\textepsilon}{\r{\textObardotlessj}}{a}} \dicPos{f}[1] \dicFlx{(‑u)}[5] \dicDirectTranslationCS{nároky, požadavky}
\dicEntry[heimtufrekur] \dicTerm{heimtu··frekur} \dicIPA{{h}{ei}{\textsubring{m}}{\textsubring{d}}{\textscy}{f}{r}{\textepsilon}{\r{g}}{\textscy}{\textsubring{r}}} \dicPos{adj}[1]\dicFlx{}[-1] \dicDirectTranslationCS{náročný (mající velké nároky), (hodně) vyžadující}
\dicEntry[heimula] \dicTerm{heimul|a} \dicIPA{{h}{ei}{\textlengthmark}{m}{\textscy}{l}{a}} \dicPos{f}[1] \dicFlx{(‑u, ‑ur)}[19] \dicFieldCat{bot.} \dicSynonym{njóli} \dicDirectTranslationCS{šťovík, šťovík dlouholistý} \textit{(l.~{\textLA{Rumex longifolius}})}  \dicsymPhoto\ 
\dicFigure{ds_image_heimula_0_1.jpg}{Heimula}{Heimula - Sten Porse, GFDL}
\dicEntry[heimur] \dicTerm{heim|ur} \dicsymFrequent\  \dicIPA{{h}{ei}{\textlengthmark}{m}{\textscy}{\textsubring{r}}} \dicPos{m}[6] \dicFlx{(‑s, ‑ar)}[4] \textbf{1.} \dicSynonym*{alheimurinn} \dicDirectTranslationCS{svět, vesmír}  \textbf{2.} \dicSynonym*{jörðin} \dicDirectTranslationCS{svět} \dicIndirectTranslationCS{(planeta, na které žijeme)} \dicExampleIS{ferðast um heiminn} \dicExampleCS{cestovat po světě};  \dicPhraseIS{(fara) út í heim} \dicDirectTranslationCS{(vyrazit) do světa};  \dicPhraseIS{(ferðast) út um allan heim} \dicFlx{refl} \dicDirectTranslationCS{(cestovat) po celém světě};  \dicIdiom{heimur}{ \dicPhraseIS{koma í heiminn}} \dicDirectTranslationCS{přijít na svět};  \dicPhraseIS{leggja heiminn að fótum sér} \dicLangCat{přen.} \dicDirectTranslationCS{položit svět k~nohám, dobýt (celý) svět};  \dicPhraseIS{lifa í sínum eigin heimi} \dicLangCat{přen.} \dicDirectTranslationCS{žít ve svém vlastním světě};  \dicPhraseIS{sýna e‑m í tvo heimana} \dicLangCat{přen.} \dicDirectTranslationCS{dát (komu) nůž na krk};  \dicPhraseIS{vera í öðrum heimi} \dicLangCat{přen.} \dicDirectTranslationCS{žít\,/\addthin ocitnout se v~jiném světě}
\dicEntry[heimþrá] \dicTerm{heim··þrá} \dicIPA{{h}{ei}{m}{\texttheta}{r}{au}} \dicPos{f}[4] \dicFlx{(‑r)}[19] \dicDirectTranslationCS{touha\,/\addthin stesk po domově, nostalgie} \dicExampleIS{vera með heimþrá} \dicExampleCS{stýskat si po domově}
\dicEntry[heit] \dicTerm{heit} \dicIPA{{h}{ei}{\textlengthmark}{\textsubring{d}}} \dicPos{n}[2] \dicFlx{(‑s, ‑)}[5] \dicSynonym{loforð} \dicDirectTranslationCS{přísaha, slib};  \dicPhraseIS{strengja þess heit að (gera e‑ð)} \dicDirectTranslationCS{(slavnostně) slíbit (udělat (co))}
\dicEntry[heita] \dicTerm{heita} \dicsymFrequent\  \dicIPA{{h}{ei}{\textlengthmark}{\textsubring{d}}{a}} \dicPos{v}[6] \dicFlx{(heiti, hét, hétum, héti, heitið)}[57] \textbf{1.} \dicDirectTranslationCS{jmenovat se, nazývat se} \dicExampleIS{Ég heiti Jón.} \dicExampleCS{Jmenuji se Jón.}  \textbf{2.} \dicFlx{(dat +) dat} \dicSynonym{lofa} \dicDirectTranslationCS{slíbit, přísahat};  \dicPhraseIS{heita e‑m e‑u} \dicDirectTranslationCS{slíbit (komu co)} \dicExampleIS{heita honum trúnaði} \dicExampleCS{slíbit mu věrnost}  \textbf{3.} \dicPhraseIS{heita á e‑n} \dicDirectTranslationCS{vzývat (koho) (Boha ap.)}  \textbf{4.} \dicPhraseIS{ferðinni er heitið (til borgarinnar)} \dicFlx{impers} \dicDirectTranslationCS{výlet směřuje (do města)};  \dicIdiom{heita}{ \dicPhraseIS{láta gott heita}} \dicDirectTranslationCS{nechat toho (sportovní kariéry ap.)};  \dicPhraseIS{það á að heita svo} \dicDirectTranslationCS{dá se říct};  \dicPhraseIS{það má heita} \dicDirectTranslationCS{je téměř jisté};  \dicPhraseIS{það getur ekki heitið} \dicDirectTranslationCS{nedá se tvrdit};  \dicPhraseIS{það er ekkert sem heitir} \dicLangCat{přen.} \dicDirectTranslationCS{už žádné výmluvy}
\dicEntry[heiti] \dicTerm{heiti\smash{\textsuperscript{1}}} \dicIPA{{h}{ei}{\textlengthmark}{\textsubring{d}}{\textsci}} \dicPos{n}[2] \dicFlx{(‑s, ‑)}[14] \textbf{1.} \dicSynonym{nafn} \dicDirectTranslationCS{jméno, název, výraz, pojem, termín}  \textbf{2.} \dicFieldCat{lit.} \dicDirectTranslationCS{básnický výraz}
\dicEntry[heiti] \dicTerm{heiti\smash{\textsuperscript{2}}} \dicIPA{{h}{ei}{\textlengthmark}{\textsubring{d}}{\textsci}} \dicPos{v} \dicFlx{ind praes sg 1 pers} \dicLink{heita}
\dicEntry[heitið] \dicTerm{heitið} \dicIPA{{h}{ei}{\textlengthmark}{\textsubring{d}}{\textsci}{\texttheta}} \dicPos{v} \dicFlx{supin} \dicLink{heita}
\dicEntry[heitkona] \dicTerm{heit··kon|a} \dicIPA{{h}{ei}{\textlengthmark}{\textsubring{d}}{k\smash{\textsuperscript{h}}}{\textopeno}{n}{a}} \dicPos{f}[1] \dicFlx{(‑u, ‑ur)}[27] \dicDirectTranslationCS{snoubenka}
\dicEntry[heitmaður] \dicTerm{heit··|maður} \dicIPA{{h}{ei}{\textlengthmark}{\textsubring{d}}{m}{a}{ð}{\textscy}{\textsubring{r}}} \dicPos{m}[13] \dicFlx{(‑manns, ‑menn)}[2] \dicDirectTranslationCS{snoubenec}
\dicEntry[heitstrenging] \dicTerm{heit··streng·ing} \dicIPA{{h}{ei}{\textlengthmark}{\textsubring{d}}{s}{\textsubring{d}}{r}{ei}{\textltailn}{\r{\textObardotlessj}}{i}{\ng}{\r{g}}} \dicPos{f}[4] \dicFlx{(‑ar, ‑ar)}[5] \dicDirectTranslationCS{slavnostní přísaha\,/\addthin slib}
\dicEntry[heitur] \dicTerm{heitur} \dicsymFrequent\  \dicIPA{{h}{ei}{\textlengthmark}{\textsubring{d}}{\textscy}{\textsubring{r}}} \dicPos{adj}[1]\dicFlx{}[-1] \textbf{1.} \dicSynonym*{með miklum hita} \dicDirectTranslationCS{horký, teplý};  \dicPhraseIS{heitur matur} \dicDirectTranslationCS{teplé jídlo};  \dicPhraseIS{e‑m er heitt} \dicFlx{impers} \dicDirectTranslationCS{(komu) je teplo}  \textbf{2.} \dicSynonym{ákafur} \dicDirectTranslationCS{horoucí, vřelý, vroucí (city ap.)} \dicExampleIS{heit ást} \dicExampleCS{vroucí láska}  \textbf{3.} \dicSynonym*{umtalaður} \dicDirectTranslationCS{žhavý (novinka ap.), atraktivní} \dicExampleIS{einn heitasti staðurinn núna} \dicExampleCS{jedno z~nejatraktivnějších míst v~dnešní době}
\dicEntry[hekla] \dicTerm{hekl|a} \dicIPA{{h}{\textepsilon}{h}{\r{g}}{l}{a}} \dicPos{v}[1] \dicFlx{(‑aði)}[1] \dicFlx{acc} \dicDirectTranslationCS{háčkovat} \dicExampleIS{hekla peysu} \dicExampleCS{háčkovat svetr}
\dicEntry[heklunál] \dicTerm{heklu··nál} \dicIPA{{h}{\textepsilon}{h}{\r{g}}{l}{\textscy}{n}{au}{\textsubring{l}}} \dicPos{f}[4] \dicFlx{(‑ar, ‑ar)}[1] \dicDirectTranslationCS{háček (na háčkování ap.)}
\dicEntry[hektari] \dicTerm{hektar|i} \dicIPA{{h}{\textepsilon}{x}{\textsubring{d}}{a}{r}{\textsci}} \dicPos{m}[1] \dicFlx{(‑a, ‑ar)}[13] \dicDirectTranslationCS{hektar}
\dicEntry[hel] \dicTerm{hel} \dicIPA{{h}{\textepsilon}{\textlengthmark}{\textsubring{l}}} \dicPos{f}[4] \dicFlx{(‑jar)}[10] \textbf{1.} \dicSynonym*{ríki dauðra} \dicDirectTranslationCS{podsvětí, říše mrtvých}  \textbf{2.} \dicSynonym{bani} \dicDirectTranslationCS{záhuba, smrt};  \dicPhraseIS{heimta e‑n úr helju} \dicDirectTranslationCS{zachránit (koho) před smrtí};  \dicPhraseIS{milli heims og heljar} \dicFlx{adv} \dicDirectTranslationCS{na hranici života a~smrti}
\dicEntry[helber] \dicTerm{hel··ber} \dicIPA{{h}{\textepsilon}{l}{\textsubring{b}}{\textepsilon}{\textsubring{r}}} \dicPos{adj}[5] \dicFlx{(f ‑)}[10] \textbf{1.} \dicSynonym*{allsnakinn} \dicDirectTranslationCS{(úplně) holý\,/\addthin nahý}  \textbf{2.} \dicSynonym{einskær} \dicDirectTranslationCS{holý, zřetelný (lež ap.)} \dicExampleIS{helber lygi} \dicExampleCS{holá lež}
\dicEntry[held] \dicTerm{held} \dicIPA{{h}{\textepsilon}{l}{\textsubring{d}}} \dicPos{v} \dicFlx{ind praes sg 1 pers} \dicLink{halda\smash{\textsuperscript{2}}}
\dicEntry[heldri] \dicTerm{heldri} \dicsymFrequent\  \dicIPA{{h}{\textepsilon}{l}{\textsubring{d}}{r}{\textsci}} \dicPos{adj}[12] \dicFlx{comp (sup helstur)}[8] \dicDirectTranslationCS{lepší, urozenější} \dicExampleIS{heldri kona} \dicExampleCS{žena z~lepší společnosti}
\dicEntry[heldur] \dicTerm{heldur\smash{\textsuperscript{1}}} \dicIPA{{h}{\textepsilon}{l}{\textsubring{d}}{\textscy}{\textsubring{r}}} \dicPos{adj}[2]\dicFlx{}[-14] \textbf{1.} \dicSynonym{vatnsþéttur} \dicDirectTranslationCS{nepropustný, nepromokavý}  \textbf{2.} \dicSynonym*{manngengur} \dicDirectTranslationCS{držící (led ap.)} \dicExampleIS{Ísinn er orðinn heldur.} \dicExampleCS{Led už drží.}
\dicEntry[heldur] \dicTerm{heldur\smash{\textsuperscript{2}}} \dicsymFrequent\  \dicIPA{{h}{\textepsilon}{l}{\textsubring{d}}{\textscy}{\textsubring{r}}} \dicPos{adv} \dicFlx{comp (pos gjarna, sup helst)} \textbf{1.} \dicSynonym{frekar} \dicDirectTranslationCS{spíše, raději}  \textbf{2.} \dicSynonym{dálítið} \dicDirectTranslationCS{trochu, poněkud};  \dicIdiom{heldur}{ \dicPhraseIS{heldur betur}} \dicFlx{adv} \dicDirectTranslationCS{opravdu, skutečně};  \dicPhraseIS{hvort heldur -- eða} \dicFlx{conj} \dicDirectTranslationCS{ať již -- nebo};  \dicPhraseIS{vilja heldur} \dicDirectTranslationCS{chtít raději, dát\,/\addthin dávat přednost, preferovat} \dicExampleIS{Ég vil heldur kaffi en te.} \dicExampleCS{Raději bych si dal kafe než čaj.};  \dicPhraseIS{öllu heldur} \dicFlx{adv} \dicDirectTranslationCS{spíše, raději}
\dicEntry[heldur] \dicTerm{heldur\smash{\textsuperscript{3}}} \dicsymFrequent\  \dicIPA{{h}{\textepsilon}{l}{\textsubring{d}}{\textscy}{\textsubring{r}}} \dicPos{conj} \dicDirectTranslationCS{ale, ale spíše, nýbrž} \dicIndirectTranslationCS{(používané se záporem)} \dicExampleIS{Hann er ekki kennari heldur er hann læknir.} \dicExampleCS{Není učitel, ale doktor.};  \dicIdiom{heldur}[en]{ \dicPhraseIS{heldur en}} \dicFlx{conj} \dicDirectTranslationCS{(spíše) než}
\dicEntry[helför] \dicTerm{hel··|för} \dicIPA{{h}{\textepsilon}{l}{f}{\oe}{\textsubring{r}}} \dicPos{f}[7] \dicFlx{(‑farar)}[19] \dicDirectTranslationCS{vyhlazení, holokaust}
\dicEntry[helga] \dicTerm{helg|a} \dicIPA{{h}{\textepsilon}{l}{\r{g}}{a}} \dicPos{v}[1] \dicFlx{(‑aði)}[1] \textbf{1.} \dicFlx{acc + dat} \dicDirectTranslationCS{věnovat, obětovat, zasvětit, posvětit};  \dicPhraseIS{helga sig e‑u} \dicDirectTranslationCS{věnovat se (čemu)};  \dicPhraseIS{helga e‑ð e‑u} \dicDirectTranslationCS{zasvětit (čemu co)} \dicExampleIS{helga líf sitt starfi} \dicExampleCS{zasvětit svůj život práci}  \textbf{2.} \dicFlx{dat + acc} \dicDirectTranslationCS{přivlastnit si, přibrat si, vzít za své};  \dicPhraseIS{helga sér e‑ð} \dicDirectTranslationCS{přivlastnit si (co), vzít (co) za své}  \textbf{3.} \dicFlx{acc + dat} \dicFieldCat{náb.} \dicDirectTranslationCS{(za)světit};  \dicIdiom{helgast}{ \dicPhraseIS{e‑að helgast af e‑u}} \dicFlx{refl} \dicDirectTranslationCS{(co) vychází z~(čeho), (co) má své vysvětlení v~(čem)}
\dicEntry[helgarferð] \dicTerm{helgar··ferð} \dicIPA{{h}{\textepsilon}{l}{\r{g}}{a}{\textsubring{r}}{f}{\textepsilon}{r}{\texttheta}} \dicPos{f}[7] \dicFlx{(‑ar, ‑ir)}[1] \dicDirectTranslationCS{výlet na víkend, víkendový výlet}
\dicEntry[helgi] \dicTerm{helg|i\smash{\textsuperscript{1}}} \dicsymFrequent\  \dicIPA{{h}{\textepsilon}{l}{\r{\textObardotlessj}}{\textsci}} \dicPos{f}[5] \dicFlx{(‑ar, ‑ar)}[3] \textbf{1.} \dicSynonym{hátíðisdagur} \dicDirectTranslationCS{svátek, sváteční den}  \textbf{2.} \dicSynonym{frídagur} \dicDirectTranslationCS{víkend} \dicExampleIS{um helgina} \dicExampleCS{o~víkendu}
\dicEntry[helgi] \dicTerm{helg|i\smash{\textsuperscript{2}}} \dicIPA{{h}{\textepsilon}{l}{\r{\textObardotlessj}}{\textsci}} \dicPos{f}[3] \dicFlx{(‑i)}[3] \textbf{1.} \dicSynonym{heilagleiki} \dicDirectTranslationCS{svatost}  \textbf{2.} \dicSynonym{friðhelgi} \dicDirectTranslationCS{svatyně, svatostánek}
\dicEntry[helgidagur] \dicTerm{helgi··dag|ur} \dicIPA{{h}{\textepsilon}{l}{\r{\textObardotlessj}}{\textsci}{\textsubring{d}}{a}{\textbabygamma}{\textscy}{\textsubring{r}}} \dicPos{m}[6] \dicFlx{(‑s, ‑ar)}[62] \dicDirectTranslationCS{svátek} \dicIndirectTranslationCS{(volný den)}
\dicEntry[helgikvæði] \dicTerm{helgi··kvæði} \dicIPA{{h}{\textepsilon}{l}{\r{\textObardotlessj}}{\textsci}{k\smash{\textsuperscript{h}}}{v}{a}{i}{ð}{\textsci}} \dicPos{n}[2] \dicFlx{(‑s, ‑)}[14] \dicFieldCat{lit.} \dicDirectTranslationCS{hymnus}
\dicEntry[helgimynd] \dicTerm{helgi··mynd} \dicIPA{{h}{\textepsilon}{l}{\r{\textObardotlessj}}{\textsci}{m}{\textsci}{n}{\textsubring{d}}} \dicPos{f}[7] \dicFlx{(‑ar, ‑ir)}[1] \dicDirectTranslationCS{ikona, svatý obraz}
\dicEntry[helgisaga] \dicTerm{helgi··|saga} \dicIPA{{h}{\textepsilon}{l}{\r{\textObardotlessj}}{\textsci}{s}{a}{\textbabygamma}{a}} \dicPos{f}[1] \dicFlx{(‑sögu, ‑sögur)}[14] \dicFieldCat{lit.} \dicDirectTranslationCS{legenda} \dicIndirectTranslationCS{(epický útvar vyprávějící o~životě svatých)}
\dicEntry[helgisiður] \dicTerm{helgi··sið|ur} \dicIPA{{h}{\textepsilon}{l}{\r{\textObardotlessj}}{\textsci}{s}{\textsci}{ð}{\textscy}{\textsubring{r}}} \dicPos{m}[10] \dicFlx{(‑ar, ‑ir)}[16] \dicDirectTranslationCS{rituál}
\dicEntry[helgistaður] \dicTerm{helgi··stað|ur} \dicIPA{{h}{\textepsilon}{l}{\r{\textObardotlessj}}{\textsci}{s}{\textsubring{d}}{a}{ð}{\textscy}{\textsubring{r}}} \dicPos{m}[10] \dicFlx{(‑ar, ‑ir)}[14] \dicDirectTranslationCS{posvátné místo, svatostánek}
\dicEntry[helgur] \dicTerm{helgur} \dicsymFrequent\  \dicIPA{{h}{\textepsilon}{l}{\r{g}}{\textscy}{\textsubring{r}}} \dicPos{adj}[1]\dicFlx{}[-1] \dicDirectTranslationCS{svatý, posvátný} \dicExampleIS{helgur maður} \dicExampleCS{svatý muž}
\dicEntry[Helheimur] \dicTerm{Hel··heim|ur} \dicIPA{{h}{\textepsilon}{\textlengthmark}{l}{h}{ei}{m}{\textscy}{\textsubring{r}}} \dicPos{m}[6] \dicFlx{(‑s)}[3] \dicFieldCat{myt.} \dicDirectTranslationCS{Helheim} \dicIndirectTranslationCS{(říše mrtvých, podsvětí)}
\dicEntry[helín] \dicTerm{helín}\dicTerm{, helíum} \dicIPA{{h}\-{\textepsilon}\-{\textlengthmark}\-{l}\-{i}\-{\textsubring{n}}\-} \dicPos{n}[2] \dicFlx{(‑s)}[31] \dicFieldCat{chem.} \dicDirectTranslationCS{hélium} \textit{(l.~{\textLA{He, Helium}})}
\dicEntry[helíum] \dicTerm{helíum} \dicIPA{{h}{\textepsilon}{\textlengthmark}{l}{i}{j}{\textscy}{\textsubring{m}}} \dicPos{n}[2] \dicFlx{(‑s)}[31] \dicLink{helín}
\dicEntry[helja] \dicTerm{helj|a} \dicIPA{{h}{\textepsilon}{l}{j}{a}} \dicPos{f}[1] \dicFlx{(‑u, ‑ur)}[7] \dicDirectTranslationCS{říše mrtvých};  \dicPhraseIS{heimta e‑n úr helju} \dicDirectTranslationCS{zachránit (koho) před smrtí};  \dicPhraseIS{vera milli heims og helju} \dicLangCat{přen.} \dicDirectTranslationCS{být na hranici života a~smrti}
\dicEntry[heljarmenni] \dicTerm{heljar··menni} \dicIPA{{h}{\textepsilon}{l}{j}{a}{r}{m}{\textepsilon}{n}{\textsci}} \dicPos{n}[2] \dicFlx{(‑s, ‑)}[14] \dicDirectTranslationCS{svalovec, silák}
\dicEntry[heljarstökk] \dicTerm{heljar··stökk} \dicIPA{{h}{\textepsilon}{l}{j}{a}{\textsubring{r}}{s}{\textsubring{d}}{\oe}{h}{\r{g}}} \dicPos{n}[2] \dicFlx{(‑s, ‑)}[5] \dicDirectTranslationCS{salto, kotrmelec}
\dicEntry[hella] \dicTerm{hell|a\smash{\textsuperscript{1}}} \dicIPA{{h}{\textepsilon}{\textsubring{d}}{l}{a}} \dicPos{f}[1] \dicFlx{(‑u, ‑ur)}[19] \textbf{1.} \dicSynonym*{flatur steinn} \dicDirectTranslationCS{kamenná deska}  \textbf{2.} \dicSynonym*{eldavélarhella} \dicDirectTranslationCS{plotna, plotýnka};  \dicPhraseIS{fá hellu fyrir eyrun} \dicDirectTranslationCS{zalehnout v~uších};  \dicPhraseIS{e‑að er fyrir neðan allar hellur} \dicDirectTranslationCS{(co) je naprosto nepřijatelné (chování ap.)}
\dicEntry[hella] \dicTerm{hell|a\smash{\textsuperscript{2}}} \dicsymFrequent\  \dicIPA{{h}{\textepsilon}{\textsubring{d}}{l}{a}} \dicPos{v}[2] \dicFlx{(‑ti, ‑t)}[68] \dicFlx{dat} \dicDirectTranslationCS{(na)lít} \dicExampleIS{hella vatni úr flöskunni í glasið} \dicExampleCS{nalít vodu z~láhve do skleničky};  \dicIdiom{hella}[í]{ \dicPhraseIS{hella í sig}} \dicDirectTranslationCS{lít do sebe (alkohol ap.)};  \dicIdiom{hella}[niður]{ \dicPhraseIS{hella e‑u niður}} \dicDirectTranslationCS{rozlít, vylít} \dicExampleIS{hella mjólk niður} \dicExampleCS{rozlít mléko};  \dicIdiom{hella}[upp á]{ \dicPhraseIS{hella upp á}} \dicDirectTranslationCS{připravit kávu, dát na kávu}; { \dicPhraseIS{hella upp á e‑ð}} \dicDirectTranslationCS{připravit (co), uvařit (co) (kávu ap.)};  \dicIdiom{hella}[yfir]{ \dicPhraseIS{hella sér yfir e‑n}} \dicSynonym*{húðskamma e‑n} \dicDirectTranslationCS{vynadat (komu), vylít si na (kom) zlost};  \dicIdiom{hellast}{ \dicPhraseIS{hellast}} \dicFlx{refl} \dicDirectTranslationCS{vylít se, rozlít se};  \dicIdiom{hella}{ \dicPhraseIS{hella úr skálum reiði sinnar}} \dicLangCat{přen.} \dicDirectTranslationCS{vylít si zlost}
\dicEntry[hellingur] \dicTerm{hell··ing|ur} \dicIPA{{h}{\textepsilon}{\textsubring{d}}{l}{i}{\ng}{\r{g}}{\textscy}{\textsubring{r}}} \dicPos{m}[6] \dicFlx{(‑s)}[9] \dicSynonym*{mikið magn} \dicDirectTranslationCS{spousta, plno} \dicExampleIS{hellingur af appelsínusafa} \dicExampleCS{spousta pomerančového džusu}
\dicEntry[hellir] \dicTerm{hell|ir} \dicIPA{{h}{\textepsilon}{\textsubring{d}}{l}{\textsci}{\textsubring{r}}} \dicPos{m}[7] \dicFlx{(‑is, ‑ar)}[1] \dicDirectTranslationCS{jeskyně}
\dicEntry[hellirigning] \dicTerm{helli··rign·ing} \dicIPA{{h}{\textepsilon}{\textsubring{d}}{l}{\textsci}{r}{\textsci}{\r{g}}{n}{i}{\ng}{\r{g}}} \dicPos{f}[4] \dicFlx{(‑ar, ‑ar)}[5] \dicDirectTranslationCS{liják, lijavec, vydatný déšť}
\dicEntry[hellismunni] \dicTerm{hellis··munn|i} \dicIPA{{h}{\textepsilon}{\textsubring{d}}{l}{\textsci}{s}{m}{\textscy}{n}{\textsci}} \dicPos{m}[1] \dicFlx{(‑a, ‑ar)}[1] \dicDirectTranslationCS{ústí jeskyně}
\dicEntry[helluhnoðri] \dicTerm{hellu··hnoðr|i} \dicIPA{{h}{\textepsilon}{\textsubring{d}}{l}{\textscy}{\textsubring{n}}{\textopeno}{ð}{r}{\textsci}} \dicPos{m}[1] \dicFlx{(‑a, ‑ar)}[1] \dicFieldCat{bot.} \dicDirectTranslationCS{rozchodník ostrý} \textit{(l.~{\textLA{Sedum acre}})}  \dicsymPhoto\ 
\dicFigure{4002.jpg}{Helluhnoðri}{Helluhnoðri - Zicha Ondřej, Biolib, Copyright/CC-BY-NC}
\dicEntry[helluhraun] \dicTerm{hellu··hraun} \dicIPA{{h}{\textepsilon}{\textsubring{d}}{l}{\textscy}{\textsubring{r}}{\oe i}{\textsubring{n}}} \dicPos{n}[2] \dicFlx{(‑s, ‑)}[5] \dicFieldCat{geol.} \dicDirectTranslationCS{provazová láva}
\dicEntry[helmingaskipti] \dicTerm{helminga··skipti} \dicIPA{{h}{\textepsilon}{l}{m}{i}{\ng}{\r{g}}{a}{s}{\r{\textObardotlessj}}{\textsci}{f}{\textsubring{d}}{\textsci}} \dicPos{n}[2] \dicFlx{pl}[19] \dicDirectTranslationCS{(roz)půlení}
\dicEntry[helmingur] \dicTerm{helm··ing|ur} \dicsymFrequent\  \dicIPA{{h}{\textepsilon}{l}{m}{i}{\ng}{\r{g}}{\textscy}{\textsubring{r}}} \dicPos{m}[6] \dicFlx{(‑s, ‑ar)}[8] \dicDirectTranslationCS{půl, půlka, polovina};  \dicPhraseIS{betri helmingurinn} \dicLangCat{přen.} \dicDirectTranslationCS{lepší polovička} \dicIndirectTranslationCS{(manželka)};  \dicPhraseIS{skipta e‑u til helminga} \dicDirectTranslationCS{rozdělit (co) na půl};  \dicPhraseIS{um helming} \dicFlx{adv} \dicDirectTranslationCS{o~polovinu}
\dicEntry[helsi] \dicTerm{helsi} \dicIPA{{h}{\textepsilon}{l}{s}{\textsci}} \dicPos{n}[2] \dicFlx{(‑s, ‑)}[14] \textbf{1.} \dicSynonym{hálsband} \dicDirectTranslationCS{obojek}  \textbf{2.} \dicSynonym{fjötur} \dicDirectTranslationCS{okovy, pouta}  \textbf{3.} \dicSynonym{ánauð} \dicDirectTranslationCS{útisk, nesvoboda}
\dicEntry[helsingi] \dicTerm{hels··ing|i} \dicIPA{{h}{\textepsilon}{l}{s}{i}{\textltailn}{\r{\textObardotlessj}}{\textsci}} \dicPos{m}[1] \dicFlx{(‑ja, ‑jar)}[14] \dicFieldCat{zool.} \dicDirectTranslationCS{berneška bělolící} \textit{(l.~{\textLA{Branta leucopsis}})}  \dicsymPhoto\ 
\dicFigure{1198.jpg}{Helsingi}{Helsingi - Zicha Ondřej, Biolib, Copyright/CC-BY-NC}
\dicEntry[Helsingjabotn] \dicTerm{Helsingja··botn} \dicIPA{{h}{\textepsilon}{l}{s}{i}{\textltailn}{\r{\textObardotlessj}}{a}{\textsubring{b}}{\textopeno}{h}{\textsubring{d}}{\textsubring{n}}} \dicPos{m}[4] \dicFlx{(‑s)}[2] \dicFieldCat{geog.} \dicDirectTranslationCS{Botnický záliv}
\dicEntry[Helsingjaeyri] \dicTerm{Helsingja··eyr|i} \dicIPA{{h}{\textepsilon}{l}{s}{i}{\textltailn}{\r{\textObardotlessj}}{a}{ei}{r}{\textsci}} \dicPos{f}[5] \dicFlx{(‑ar)}[5] \dicFieldCat{geog.} \dicDirectTranslationCS{Helsingør} \dicIndirectTranslationCS{(přístavní dánské město)}
\dicEntry[Helsinki] \dicTerm{Helsinki} \dicIPA{{h}{\textepsilon}{l}{s}{i}{\r{\textltailn}}{\r{\textObardotlessj}}{\textsci}} \dicPos{subs} \dicFlx{indecl} \dicFieldCat{geog.} \dicDirectTranslationCS{Helsinky} \dicIndirectTranslationCS{(hlavní město Finska)}
\dicEntry[helst] \dicTerm{helst} \dicIPA{{h}{\textepsilon}{l}{s}{\textsubring{d}}} \dicPos{adv} \dicFlx{sup (pos gjarna, comp heldur)} \dicDirectTranslationCS{nejvíce, především};  \dicPhraseIS{helst til} \dicFlx{adv} \dicDirectTranslationCS{dost, značně, až příliš}
\dicEntry[helstur] \dicTerm{helstur} \dicIPA{{h}{\textepsilon}{l}{s}{\textsubring{d}}{\textscy}{\textsubring{r}}} \dicPos{adj}[12] \dicFlx{sup (comp heldri)}[8] \dicDirectTranslationCS{hlavní, nejdůležitější}
\dicEntry[helsærður] \dicTerm{hel··særður} \dicIPA{{h}{\textepsilon}{l}{s}{a}{i}{r}{ð}{\textscy}{\textsubring{r}}} \dicPos{adj}[2]\dicFlx{}[-4] \dicDirectTranslationCS{smrtelně zraněný}
\dicEntry[helvíti] \dicTerm{hel··víti\smash{\textsuperscript{1}}} \dicsymFrequent\  \dicIPA{{h}{\textepsilon}{l}{v}{i}{\textsubring{d}}{\textsci}} \dicPos{n}[2] \dicFlx{(‑s, ‑)}[14] \textbf{1.} \dicDirectTranslationCS{peklo}  \textbf{2.} \dicDirectTranslationCS{sakra, k~čertu} \dicIndirectTranslationCS{(nadávka)};  \dicPhraseIS{farið það í helvíti!} \dicDirectTranslationCS{k~čertu s~tím!};  \dicPhraseIS{helvítið á honum!} \dicDirectTranslationCS{ten mizera!};  \dicPhraseIS{helvíti!} \dicDirectTranslationCS{sakra!, k~čertu!};  \dicPhraseIS{helvítið þitt!} \dicDirectTranslationCS{ty mizero!, ty potvoro!}
\dicEntry[helvíti] \dicTerm{hel··víti\smash{\textsuperscript{2}}} \dicIPA{{h}{\textepsilon}{l}{v}{i}{\textsubring{d}}{\textsci}} \dicPos{adv} \dicDirectTranslationCS{sakra, zatraceně, zpropadeně} \dicExampleIS{helvíti góður læknir} \dicExampleCS{zatraceně dobrý lékař}
\dicEntry[helvítis] \dicTerm{hel··vítis} \dicIPA{{h}{\textepsilon}{l}{v}{i}{\textsubring{d}}{\textsci}{s}} \dicPos{adj}[13] \dicFlx{indecl}[1] \dicDirectTranslationCS{zpropadený, zatracený} \dicExampleIS{helvítis asni} \dicExampleCS{zatracený blbec}
\dicEntry[hem] \dicTerm{hem} \dicIPA{{h}{\textepsilon}{\textlengthmark}{\textsubring{m}}} \dicPos{v} \dicFlx{ind praes sg 1 pers} \dicLink{hemja\smash{\textsuperscript{2}}}
\dicEntry[hemdi] \dicTerm{hemdi} \dicIPA{{h}{\textepsilon}{m}{\textsubring{d}}{\textsci}} \dicPos{v} \dicFlx{con pf sg 1 pers} \dicLink{hemja\smash{\textsuperscript{2}}}
\dicEntry[hemill] \dicTerm{hem|ill} \dicIPA{{h}{\textepsilon}{\textlengthmark}{m}{\textsci}{\textsubring{d}}{\textsubring{l}}} \dicPos{m}[6] \dicFlx{(‑ils, ‑lar)}[35] \textbf{1.} \dicSynonym{hindrun} \dicDirectTranslationCS{brzda, zábrana, překážka};  \dicPhraseIS{e‑að er hemill á e‑u} \dicDirectTranslationCS{(co) brzdí (co), (co) je brzdou pro (co)};  \dicPhraseIS{hafa hemil á sér} \dicDirectTranslationCS{krotit se};  \dicPhraseIS{hafa hemil á e‑u\,/\addthin e‑m} \dicDirectTranslationCS{mít nad (čím\,/\addthin kým) kontrolu}  \textbf{2.} \dicSynonym{bremsa\smash{\textsuperscript{1}}} \dicDirectTranslationCS{brzda} \dicIndirectTranslationCS{(zařízení ke zpomalení nebo zastavení vozidla ap.)}
\dicEntry[hemja] \dicTerm{hemj|a\smash{\textsuperscript{1}}} \dicIPA{{h}{\textepsilon}{m}{j}{a}} \dicPos{f}[1] \dicFlx{(‑u)}[5] \dicDirectTranslationCS{hranice, míra};  \dicPhraseIS{e‑að er engin hemja} \dicLangCat{přen.} \dicDirectTranslationCS{(co) nezná hranic}
\dicEntry[hemja] \dicTerm{hemja\smash{\textsuperscript{2}}} \dicIPA{{h}{\textepsilon}{m}{j}{a}} \dicPos{v}[4] \dicFlx{(hem, hamdi, hömdum, hemdi, hamið)}[3] \dicFlx{acc} \dicDirectTranslationCS{(z)krotit, dostat pod kontrolu} \dicExampleIS{hemja tilfinningar sínar} \dicExampleCS{krotit své pocity}
\dicEntry[hemla] \dicTerm{heml|a} \dicIPA{{h}{\textepsilon}{m}{l}{a}} \dicPos{v}[1] \dicFlx{(‑aði)}[44] \dicFlx{acc} \dicSynonym*{stíga á hemil} \dicDirectTranslationCS{(za)brzdit} \dicExampleIS{hemla á rauðu ljósi} \dicExampleCS{zabrzdit na červenou}
\dicEntry[hempa] \dicTerm{hemp|a} \dicIPA{{h}{\textepsilon}{\textsubring{m}}{\textsubring{b}}{a}} \dicPos{f}[1] \dicFlx{(‑u, ‑ur)}[19] \dicFieldCat{náb.} \dicSynonym{kjóll} \dicDirectTranslationCS{sutana, klerika}
\dicEntry[henda] \dicTerm{hen|da} \dicsymFrequent\  \dicIPA{{h}{\textepsilon}{n}{\textsubring{d}}{a}} \dicPos{v}[2] \dicFlx{(‑ti, ‑t)}[42] \dicFlx{acc\,/\addthin dat} \textbf{1.} \dicFlx{dat} \dicSynonym{kasta} \dicDirectTranslationCS{hodit, vyhodit, odhodit} \dicExampleIS{henda ruslinu í tunnuna} \dicExampleCS{vyhodit smetí do popelnice}  \textbf{2.} \dicFlx{dat} \dicSynonym{fleygja} \dicDirectTranslationCS{vyhodit, zahodit, zbavit se}  \textbf{3.} \dicFlx{acc} \dicSynonym{grípa} \dicDirectTranslationCS{chytit, uchopit} \dicExampleIS{henda vopnið} \dicExampleCS{uchopit zbraň}  \textbf{4.} \dicSynonym*{gerast} \dicDirectTranslationCS{stát se, přihodit se} \dicExampleIS{Þetta hefur hent mig þrisvar á stuttum tíma.} \dicExampleCS{To se mi stalo třikrát v~krátkém čase.};  \dicIdiom{henda}[að]{ \dicPhraseIS{henda gaman að e‑m}} \dicDirectTranslationCS{dělat si z~(koho) legraci};  \dicIdiom{hendast}{ \dicPhraseIS{hendast}} \dicFlx{refl} \dicSynonym{þjóta} \dicDirectTranslationCS{hnát se, řítit se, uhánět} \dicExampleIS{hendast fram úr rúminu} \dicExampleCS{vyřítit se z~postele}
\dicEntry[hending] \dicTerm{hend··ing} \dicIPA{{h}{\textepsilon}{n}{\textsubring{d}}{i}{\ng}{\r{g}}} \dicPos{f}[4] \dicFlx{(‑ar, ‑ar)}[5] \textbf{1.} \dicSynonym{tilviljun} \dicDirectTranslationCS{náhoda} \dicExampleIS{af hendingu} \dicExampleCS{náhodou}  \textbf{2.} \dicFieldCat{lit.} \dicSynonym{vísuorð} \dicDirectTranslationCS{verš}  \textbf{3.} \dicFieldCat{lit.} \dicIndirectTranslationCS{slabika, která se rýmuje s~jinou slabikou ve stejném verši}
\dicEntry[hendur] \dicTerm{hendur} \dicIPA{{h}{\textepsilon}{n}{\textsubring{d}}{\textscy}{\textsubring{r}}} \dicPos{f} \dicFlx{pl nom} \dicLink{hönd}
\dicEntry[hengibrú] \dicTerm{hengi··|brú} \dicIPA{{h}{ei}{\textltailn}{\r{\textObardotlessj}}{\textsci}{\textsubring{b}}{r}{u}} \dicPos{f}[9] \dicFlx{(‑brúar, ‑brýr)}[4] \dicDirectTranslationCS{visutý most}
\dicEntry[hengiflug] \dicTerm{hengi··flug} \dicIPA{{h}{ei}{\textltailn}{\r{\textObardotlessj}}{\textsci}{f}{l}{\textscy}{x}} \dicPos{n}[2] \dicFlx{(‑s, ‑)}[5] \dicSynonym*{slútandi berg} \dicDirectTranslationCS{převis, strmá skála}
\dicEntry[hengilás] \dicTerm{hengi··lás} \dicIPA{{h}{ei}{\textltailn}{\r{\textObardotlessj}}{\textsci}{l}{au}{s}} \dicPos{m}[4] \dicFlx{(‑s, ‑ar)}[25] \dicDirectTranslationCS{visací zámek}
\dicEntry[henging] \dicTerm{heng··ing} \dicIPA{{h}{ei}{\textltailn}{\r{\textObardotlessj}}{i}{\ng}{\r{g}}} \dicPos{f}[4] \dicFlx{(‑ar, ‑ar)}[5] \textbf{1.} \dicDirectTranslationCS{oběšení, pověšení (na šibenici)}  \textbf{2.} \dicDirectTranslationCS{pověšení, zavěšení (obrazu ap.)}
\dicEntry[hengirúm] \dicTerm{hengi··rúm} \dicIPA{{h}{ei}{\textltailn}{\r{\textObardotlessj}}{\textsci}{r}{u}{\textsubring{m}}} \dicPos{n}[2] \dicFlx{(‑s, ‑)}[5] \dicDirectTranslationCS{houpací síť}
\dicEntry[hengja] \dicTerm{hengj|a\smash{\textsuperscript{1}}} \dicIPA{{h}{ei}{\textltailn}{\r{\textObardotlessj}}{a}} \dicPos{f}[1] \dicFlx{(‑u, ‑ur)}[25] \dicDirectTranslationCS{sněhový převis}
\dicEntry[hengja] \dicTerm{heng|ja\smash{\textsuperscript{2}}} \dicsymFrequent\  \dicIPA{{h}{ei}{\textltailn}{\r{\textObardotlessj}}{a}} \dicPos{v}[2] \dicFlx{(‑di, ‑t)}[137] \dicFlx{acc} \textbf{1.} \dicDirectTranslationCS{pověsit, zavěsit (obraz na zeď ap.)};  \dicPhraseIS{hengja e‑ð á e‑ð} \dicDirectTranslationCS{pověsit\,/\addthin zavěsit (co) na (co)} \dicExampleIS{hengja kápuna á snaga} \dicExampleCS{pověsit kabát na věšák};  \dicPhraseIS{hengja e‑ð upp} \dicDirectTranslationCS{pověsit (co)}  \textbf{2.} \dicDirectTranslationCS{pověsit, oběsit (na šibenici ap.)}
\dicEntry[hennar] \dicTerm{hennar\smash{\textsuperscript{1}}} \dicIPA{{h}{\textepsilon}{n}{\textlengthmark}{a}{\textsubring{r}}} \dicPos{pron}[17] \dicFlx{pers} \dicFlx{sg gen} \dicLink{hún}
\dicEntry[hennar] \dicTerm{hennar\smash{\textsuperscript{2}}} \dicIPA{{h}{\textepsilon}{n}{\textlengthmark}{a}{\textsubring{r}}} \dicPos{pron}[17] \dicFlx{poss f} \dicDirectTranslationCS{její} \dicExampleIS{tölvan hennar} \dicExampleCS{její počítač}
\dicEntry[henni] \dicTerm{henni} \dicIPA{{h}{\textepsilon}{n}{\textlengthmark}{\textsci}} \dicPos{pron}[17] \dicFlx{pers} \dicFlx{sg dat} \dicLink{hún}
\dicEntry[henta] \dicTerm{hent|a} \dicsymFrequent\  \dicIPA{{h}{\textepsilon}{\textsubring{n}}{\textsubring{d}}{a}} \dicPos{v}[1] \dicFlx{(‑aði)}[86] \dicDirectTranslationCS{vyhovovat, hodit se} \dicExampleIS{Þetta hentar mér illa.} \dicExampleCS{To mi moc nevyhovuje.}
\dicEntry[hentifáni] \dicTerm{henti··fán|i} \dicIPA{{h}{\textepsilon}{\textsubring{n}}{\textsubring{d}}{\textsci}{f}{au}{n}{\textsci}} \dicPos{m}[1] \dicFlx{(‑a, ‑ar)}[1] \dicFieldCat{nám.} \dicDirectTranslationCS{cizí vlajka};  \dicPhraseIS{sigla undir hentifána} \dicDirectTranslationCS{plavit se pod cizí vlajkou}
\dicEntry[hentisemi] \dicTerm{henti··sem|i} \dicIPA{{h}{\textepsilon}{\textsubring{n}}{\textsubring{d}}{\textsci}{s}{\textepsilon}{m}{\textsci}} \dicPos{f}[2] \dicFlx{(‑i)}[2] \dicDirectTranslationCS{rozmar, vrtoch};  \dicPhraseIS{hafa sína hentisemi} \dicDirectTranslationCS{dělat si podle libosti, co člověk chce}
\dicEntry[hentistefna] \dicTerm{henti··stefn|a} \dicIPA{{h}{\textepsilon}{\textsubring{n}}{\textsubring{d}}{\textsci}{s}{\textsubring{d}}{\textepsilon}{\textsubring{b}}{n}{a}} \dicPos{f}[1] \dicFlx{(‑u)}[5] \dicDirectTranslationCS{oportunismus, prospěchářství}
\dicEntry[hentugleiki] \dicTerm{hent·ug··leik|i} \dicIPA{{h}{\textepsilon}{\textsubring{n}}{\textsubring{d}}{\textscy}{\textbabygamma}{l}{ei}{\r{\textObardotlessj}}{\textsci}} \dicPos{m}[1] \dicFlx{(‑a, ‑ar)}[1] \dicDirectTranslationCS{vhodnost, příhodnost}
\dicEntry[hentugur] \dicTerm{hent··ugur} \dicIPA{{h}{\textepsilon}{\textsubring{n}}{\textsubring{d}}{\textscy}{\textbabygamma}{\textscy}{\textsubring{r}}} \dicPos{adj}[1]\dicFlx{}[-8] \dicSynonym{heppilegur} \dicDirectTranslationCS{vhodný, příhodný, případný} \dicExampleIS{hentugt byggingarefni} \dicExampleCS{vhodný stavební materiál}
\dicEntry[heppilega] \dicTerm{heppi··lega} \dicIPA{{h}{\textepsilon}{h}{\textsubring{b}}{\textsci}{l}{\textepsilon}{\textbabygamma}{a}} \dicPos{adv} \dicDirectTranslationCS{naštěstí, šťastnou náhodou}
\dicEntry[heppilegur] \dicTerm{heppi··legur} \dicsymFrequent\  \dicIPA{{h}{\textepsilon}{h}{\textsubring{b}}{\textsci}{l}{\textepsilon}{\textbabygamma}{\textscy}{\textsubring{r}}} \dicPos{adj}[1]\dicFlx{}[-8] \dicDirectTranslationCS{vhodný, pravý, příhodný} \dicExampleIS{heppilegur tími} \dicExampleCS{příhodný čas}
\dicEntry[heppinn] \dicTerm{heppinn} \dicsymFrequent\  \dicIPA{{h}{\textepsilon}{h}{\textsubring{b}}{\textsci}{\textsubring{n}}} \dicPos{adj}[6]\dicFlx{}[-2] \dicDirectTranslationCS{šťastný, mající štěstí} \dicExampleIS{Ég er heppinn að fá þetta verkefni.} \dicExampleCS{Mám štěstí, že jsem dostal tento úkol.};  \dicPhraseIS{vera heppinn með e‑ð} \dicDirectTranslationCS{mít štěstí s~(čím)}
\dicEntry[heppnaður] \dicTerm{heppn··|aður} \dicIPA{{h}{\textepsilon}{h}{\textsubring{b}}{n}{a}{ð}{\textscy}{\textsubring{r}}} \dicPos{adj}[3] \dicFlx{(f ‑uð)}[3] \dicPhraseIS{vel heppnaður} \dicDirectTranslationCS{podařený, vydařený}
\dicEntry[heppnast] \dicTerm{heppn|ast} \dicsymFrequent\  \dicIPA{{h}{\textepsilon}{h}{\textsubring{b}}{n}{a}{s}{\textsubring{d}}} \dicPos{v}[1] \dicFlx{(‑aðist)}[98] \dicFlx{refl} \dicPhraseIS{e‑að heppnast} \dicDirectTranslationCS{(co) se podaří\,/\addthin povede} \dicExampleIS{Tilraunin heppnaðist.} \dicExampleCS{Pokus se podařil.};  \dicPhraseIS{e‑m heppnast e‑að} \dicFlx{impers} \dicDirectTranslationCS{(co) se (komu) podaří\,/\addthin povede} \dicExampleIS{Mér heppnaðist áform mitt.} \dicExampleCS{Podařilo se mi, co jsem zamýšlel.}
\dicEntry[heppni] \dicTerm{heppn|i} \dicsymFrequent\  \dicIPA{{h}{\textepsilon}{h}{\textsubring{b}}{n}{\textsci}} \dicPos{f}[3] \dicFlx{(‑i)}[3] \dicDirectTranslationCS{štěstí, šťastná náhoda} \dicExampleIS{hafa heppnina með sér} \dicExampleCS{mít štěstí}
\dicEntry[her] \dicTerm{her} \dicsymFrequent\  \dicIPA{{h}{\textepsilon}{\textlengthmark}{\textsubring{r}}} \dicPos{m}[9] \dicFlx{(‑s, ‑ir)}[21] \textbf{1.} \dicSynonym*{hermenn} \dicDirectTranslationCS{armáda, vojsko};  \dicPhraseIS{ganga í herinn} \dicDirectTranslationCS{vstoupit do armády}  \textbf{2.} \dicSynonym{hópur} \dicDirectTranslationCS{skupina};  \dicPhraseIS{her manns\,/\addthin manna} \dicDirectTranslationCS{skupina lidí};  \dicIdiom{her}{ \dicPhraseIS{það er allt í hers höndum}} \dicLangCat{přen.} \dicDirectTranslationCS{všechno je vzhůru nohama}
\dicEntry[herafli] \dicTerm{her··afl|i} \dicIPA{{h}{\textepsilon}{\textlengthmark}{r}{a}{\textsubring{b}}{l}{\textsci}} \dicPos{m}[1] \dicFlx{(‑a)}[3] \dicDirectTranslationCS{vojenská\,/\addthin ozbrojená síla}
\dicEntry[herbergi] \dicTerm{her··bergi} \dicsymFrequent\  \dicIPA{{h}{\textepsilon}{r}{\textsubring{b}}{\textepsilon}{r}{\r{\textObardotlessj}}{\textsci}} \dicPos{n}[2] \dicFlx{(‑s, ‑)}[16] \dicDirectTranslationCS{místnost, pokoj} \dicExampleIS{í miðju herbergis} \dicExampleCS{uprostřed místnosti}
\dicEntry[herbergisfélagi] \dicTerm{her·bergis··fé·lag|i} \dicIPA{{h}{\textepsilon}{r}{\textsubring{b}}{\textepsilon}{r}{\r{\textObardotlessj}}{\textsci}{s}{f}{j}{\textepsilon}{l}{a}{i}{j}{\textsci}} \dicPos{m}[1] \dicFlx{(‑a, ‑ar)}[8] \dicDirectTranslationCS{spolubydlící}
\dicEntry[herdeild] \dicTerm{her··deild} \dicIPA{{h}{\textepsilon}{r}{\textsubring{d}}{ei}{l}{\textsubring{d}}} \dicPos{f}[7] \dicFlx{(‑ar, ‑ir)}[1] \dicFieldCat{voj.} \dicDirectTranslationCS{vojenská jednotka (pluk, divize ap.)}
\dicEntry[herða] \dicTerm{her|ða} \dicsymFrequent\  \dicIPA{{h}{\textepsilon}{r}{ð}{a}} \dicPos{v}[2] \dicFlx{(‑ti, ‑t)}[48] \dicFlx{acc} \textbf{1.} \dicSynonym*{gera harðan} \dicDirectTranslationCS{tvrdit, kalit} \dicExampleIS{herða málminn} \dicExampleCS{tvrdit kov}  \textbf{2.} \dicSynonym{efla} \dicDirectTranslationCS{přitvrdit, zesílit} \dicExampleIS{herða sóknina} \dicExampleCS{zesílit útok}  \textbf{3.} \dicPhraseIS{herða sig við e‑ð} \dicDirectTranslationCS{pohnout s~(čím), pospíšit si s~(čím)} \dicExampleIS{herða sig við verkið} \dicExampleCS{pospíšit si s~prací}  \textbf{4.} \dicPhraseIS{e‑ð herðir} \dicFlx{impers} \dicDirectTranslationCS{(co) sílí, (co) zesiluje} \dicExampleIS{Vindinn herti.} \dicExampleCS{Vítr zesílil.};  \dicIdiom{herða}[að]{ \dicPhraseIS{herða að e‑u}} \dicDirectTranslationCS{utáhnout (co), dotáhnout (co) (provaz ap.)};  \dicIdiom{herða}[á]{ \dicPhraseIS{herða á e‑u}} \dicDirectTranslationCS{utáhnout (co) (provaz ap.)} \dicExampleIS{herða á bandinu} \dicExampleCS{utáhnout provaz}; { \dicPhraseIS{það herðir á e‑u}} \dicFlx{impers} \dicDirectTranslationCS{(co) sílí, (co) zesiluje (mráz ap.)};  \dicIdiom{herða}[upp]{ \dicPhraseIS{herða sig upp}} \dicDirectTranslationCS{vyhecovat se}; { \dicPhraseIS{herða upp hugann}} \dicDirectTranslationCS{vzchopit se, sebrat se}
\dicEntry[herðablað] \dicTerm{herða··|blað} \dicIPA{{h}{\textepsilon}{r}{ð}{a}{\textsubring{b}}{l}{a}{\texttheta}} \dicPos{n}[2] \dicFlx{(‑blaðs, ‑blöð)}[8] \dicFieldCat{anat.} \dicDirectTranslationCS{lopatka}
\dicEntry[herðabreiður] \dicTerm{herða··breiður} \dicIPA{{h}{\textepsilon}{r}{ð}{a}{\textsubring{b}}{r}{ei}{ð}{\textscy}{\textsubring{r}}} \dicPos{adj}[2]\dicFlx{}[-6] \dicDirectTranslationCS{ramenatý}
\dicEntry[herðar] \dicTerm{herðar} \dicsymFrequent\  \dicIPA{{h}{\textepsilon}{r}{ð}{a}{\textsubring{r}}} \dicPos{f}[12] \dicFlx{pl}[2] \dicDirectTranslationCS{ramena};  \dicPhraseIS{taka e‑ð á sínar herðar} \dicLangCat{přen.} \dicDirectTranslationCS{vzít (co) na svá bedra};  \dicPhraseIS{e‑að hvílir á herðum e‑rs} \dicLangCat{přen.} \dicDirectTranslationCS{(co) leží (komu) na ramenou}
\dicEntry[herðatré] \dicTerm{herða··tré} \dicIPA{{h}{\textepsilon}{r}{ð}{a}{t\smash{\textsuperscript{h}}}{r}{j}{\textepsilon}} \dicPos{n}[2] \dicFlx{(‑s, ‑)}[36] \dicDirectTranslationCS{ramínko (na šaty)}
\dicEntry[herðing] \dicTerm{herð··ing} \dicIPA{{h}{\textepsilon}{r}{ð}{i}{\ng}{\r{g}}} \dicPos{f}[4] \dicFlx{(‑ar, ‑ar)}[5] \dicDirectTranslationCS{utvrzení, upevnění, zpevnění}
\dicEntry[herfa] \dicTerm{herf|a\smash{\textsuperscript{1}}} \dicIPA{{h}{\textepsilon}{r}{v}{a}} \dicPos{f}[1] \dicFlx{(‑u, ‑ur)}[19] \textbf{1.} \dicSynonym{afskræmi} \dicDirectTranslationCS{nestvůra, zrůda}  \textbf{2.} \dicSynonym{raggeit} \dicDirectTranslationCS{zbabělec, zbabělkyně}
\dicEntry[herfa] \dicTerm{herf|a\smash{\textsuperscript{2}}} \dicIPA{{h}{\textepsilon}{r}{v}{a}} \dicPos{v}[1] \dicFlx{(‑aði)}[34] \dicDirectTranslationCS{kypřit ornou půdu bránami}
\dicEntry[herfang] \dicTerm{her··|fang} \dicIPA{{h}{\textepsilon}{\textsubring{r}}{f}{au}{\ng}{\r{g}}} \dicPos{n}[2] \dicFlx{(‑fangs, ‑föng)}[8] \dicDirectTranslationCS{(válečná) kořist, plen} \dicExampleIS{taka herfang} \dicExampleCS{vzít válečnou kořist}
\dicEntry[herferð] \dicTerm{her··ferð} \dicIPA{{h}{\textepsilon}{\textsubring{r}}{f}{\textepsilon}{r}{\texttheta}} \dicPos{f}[7] \dicFlx{(‑ar, ‑ir)}[1] \textbf{1.} \dicSynonym*{herleiðangur} \dicDirectTranslationCS{(vojenské) tažení}  \textbf{2.} \dicSynonym{átak} \dicDirectTranslationCS{kampaň (reklamní ap.)}
\dicEntry[herfi] \dicTerm{herfi} \dicIPA{{h}{\textepsilon}{r}{v}{\textsci}} \dicPos{n}[2] \dicFlx{(‑s, ‑)}[14] \dicDirectTranslationCS{brány} \dicIndirectTranslationCS{(nářadí k~vláčení polí)}
\dicEntry[herfilegur] \dicTerm{herfi··legur} \dicIPA{{h}{\textepsilon}{r}{v}{\textsci}{l}{\textepsilon}{\textbabygamma}{\textscy}{\textsubring{r}}} \dicPos{adj}[1]\dicFlx{}[-8] \textbf{1.} \dicSynonym{ljótur} \dicDirectTranslationCS{ošklivý, hnusný}  \textbf{2.} \dicSynonym{slæmur} \dicDirectTranslationCS{příšerný, strašný} \dicExampleIS{herfileg mistök} \dicExampleCS{příšerná chyba}
\dicEntry[herflugvöllur] \dicTerm{her·flug··|völlur} \dicIPA{{h}{\textepsilon}{\textsubring{r}}{f}{l}{\textscy}{x}{v}{\oe}{\textsubring{d}}{l}{\textscy}{\textsubring{r}}} \dicPos{m}[11] \dicFlx{(‑vallar, ‑vellir)}[5] \dicFieldCat{voj.} \dicDirectTranslationCS{letecká základna}
\dicEntry[herflutningar] \dicTerm{her··flut·ningar} \dicIPA{{h}{\textepsilon}{\textsubring{r}}{f}{l}{\textscy}{h}{\textsubring{d}}{n}{i}{\ng}{\r{g}}{a}{\textsubring{r}}} \dicPos{m}[6] \dicFlx{pl}[2] \dicDirectTranslationCS{vojenská doprava}
\dicEntry[herforingi] \dicTerm{her··for·ing|i} \dicIPA{{h}{\textepsilon}{\textsubring{r}}{f}{\textopeno}{r}{i}{\textltailn}{\r{\textObardotlessj}}{\textsci}} \dicPos{m}[1] \dicFlx{(‑ja, ‑jar)}[14] \dicFieldCat{voj.} \dicDirectTranslationCS{důstojník, důstojnice}
\dicEntry[herforingjaráð] \dicTerm{her·foringja··ráð} \dicIPA{{h}{\textepsilon}{\textsubring{r}}{f}{\textopeno}{r}{i}{\textltailn}{\r{\textObardotlessj}}{a}{r}{au}{\texttheta}} \dicPos{n}[2] \dicFlx{(‑s, ‑)}[5] \dicFieldCat{voj.} \dicDirectTranslationCS{důstojnický štáb}
\dicEntry[herfugl] \dicTerm{her··fugl} \dicIPA{{h}{\textepsilon}{\textsubring{r}}{f}{\textscy}{\r{g}}{\textsubring{l}}} \dicPos{m}[4] \dicFlx{(‑s, ‑ar)}[3] \dicFieldCat{zool.} \dicDirectTranslationCS{dudek chocholatý} \textit{(l.~{\textLA{Upupa epops}})}  \dicsymPhoto\ 
\dicFigure{13223.jpg}{Herfugl}{Herfugl - Krejčík Stanislav, Biolib, Copyright/CC-BY}
\dicEntry[hergagnaverksmiðja] \dicTerm{her·gagna··verk·smiðj|a} \dicIPA{{h}\-{\textepsilon}\-{r}\-{\r{g}}\-{a}\-{\r{g}}\-{n}\-{a}\-{v}\-{\textepsilon}\-{\textsubring{r}}\-{\r{g}}\-{s}\-{m}\-{\textsci}\-{ð}\-{j}\-{a}\-} \dicPos{f}[1] \dicFlx{(‑u, ‑ur)}[7] \dicDirectTranslationCS{zbrojovka, závod na výrobu zbraní}
\dicEntry[herganga] \dicTerm{her··|ganga} \dicIPA{{h}{\textepsilon}{r}{\r{g}}{au}{\ng}{\r{g}}{a}} \dicPos{f}[1] \dicFlx{(‑göngu, ‑göngur)}[14] \dicDirectTranslationCS{(vojenský) pochod\,/\addthin marš}
\dicEntry[hergögn] \dicTerm{her··gögn} \dicIPA{{h}{\textepsilon}{r}{\r{g}}{\oe}{\r{g}}{\textsubring{n}}} \dicPos{n}[2] \dicFlx{pl}[9] \dicDirectTranslationCS{vojenský materiál (zbraně, munice)}
\dicEntry[hergöngulag] \dicTerm{her·göngu··|lag} \dicIPA{{h}{\textepsilon}{r}{\r{g}}{\oe i}{\ng}{\r{g}}{\textscy}{l}{a}{x}} \dicPos{n}[2] \dicFlx{(‑lags, ‑lög)}[8] \dicDirectTranslationCS{marš, pochodová skladba}
\dicEntry[herja] \dicTerm{herj|a} \dicIPA{{h}{\textepsilon}{r}{j}{a}} \dicPos{v}[1] \dicFlx{(‑aði)}[34] \dicFlx{acc} \textbf{1.} \dicDirectTranslationCS{pustošit, drancovat}  \textbf{2.} \dicPhraseIS{herja á e‑n} \dicDirectTranslationCS{sužovat (koho)}  \textbf{3.} \dicPhraseIS{herja e‑ð út úr e‑m} \dicDirectTranslationCS{vyždímat\,/\addthin vydolovat (co) z~(koho)}
\dicEntry[herkvaðning] \dicTerm{her··kvað·ning} \dicIPA{{h}{\textepsilon}{\textsubring{r}}{k\smash{\textsuperscript{h}}}{v}{a}{ð}{n}{i}{\ng}{\r{g}}} \dicPos{f}[4] \dicFlx{(‑ar, ‑ar)}[5] \dicFieldCat{voj.} \dicDirectTranslationCS{odvod}
\dicEntry[herkví] \dicTerm{her··kví} \dicIPA{{h}{\textepsilon}{\textsubring{r}}{k\smash{\textsuperscript{h}}}{v}{i}} \dicPos{f}[4] \dicFlx{(‑ar, ‑ar)}[5] \dicSynonym*{innilokun} \dicDirectTranslationCS{obležení, obléhání, blokáda} \dicExampleIS{vera í herkví} \dicExampleCS{být v~obležení}
\dicEntry[herlegur] \dicTerm{her··legur} \dicIPA{{h}{\textepsilon}{r}{l}{\textepsilon}{\textbabygamma}{\textscy}{\textsubring{r}}} \dicPos{adj}[1]\dicFlx{}[-8] \dicLangCat{hovor.} \dicSynonym{dásamlegur} \dicDirectTranslationCS{nádherný, skvělý}
\dicEntry[herlið] \dicTerm{her··lið} \dicIPA{{h}{\textepsilon}{r}{l}{\textsci}{\texttheta}} \dicPos{n}[2] \dicFlx{(‑s)}[2] \dicDirectTranslationCS{vojsko, vojáci}
\dicEntry[herlög] \dicTerm{her··lög} \dicIPA{{h}{\textepsilon}{r}{l}{\oe}{x}} \dicPos{n}[2] \dicFlx{pl}[9] \textbf{1.} \dicFieldCat{práv.} \dicDirectTranslationCS{vojenské právo}  \textbf{2.} \dicFieldCat{práv.} \dicDirectTranslationCS{stanné právo}
\dicEntry[herma] \dicTerm{herm|a} \dicIPA{{h}{\textepsilon}{r}{m}{a}} \dicPos{v}[2] \dicFlx{(‑di, ‑t)}[141] \textbf{1.} \dicDirectTranslationCS{napodobit, napodobovat};  \dicPhraseIS{herma eftir e‑m} \dicDirectTranslationCS{napodobit (koho)}  \textbf{2.} \dicDirectTranslationCS{sdělit, informovat};  \dicIdiom{herma}[upp á]{ \dicPhraseIS{herma e‑ð upp á e‑n}} \dicDirectTranslationCS{připomenout (komu co)}
\dicEntry[hermaður] \dicTerm{her··|maður} \dicsymFrequent\  \dicIPA{{h}{\textepsilon}{r}{m}{a}{ð}{\textscy}{\textsubring{r}}} \dicPos{m}[13] \dicFlx{(‑manns, ‑menn)}[2] \dicSynonym{bardagamaður} \dicDirectTranslationCS{voják(yně)} \dicExampleIS{vopnaðir hermenn} \dicExampleCS{ozbrojení vojáci};  \dicPhraseIS{fyrrum hermaður} \dicDirectTranslationCS{veterán(ka), vysloužilec};  \dicPhraseIS{óbreyttur hermaður} \dicDirectTranslationCS{vojín}
\dicEntry[hermannabraggi] \dicTerm{her·manna··bragg|i} \dicIPA{{h}{\textepsilon}{r}{m}{a}{n}{a}{\textsubring{b}}{r}{a}{\r{\textObardotlessj}}{\textsci}} \dicPos{m}[1] \dicFlx{(‑a, ‑ar)}[8] \dicDirectTranslationCS{kasárna}
\dicEntry[hermannabúningur] \dicTerm{her·manna··bú·ning|ur} \dicIPA{{h}\-{\textepsilon}\-{r}\-{m}\-{a}\-{n}\-{a}\-{\textsubring{b}}\-{u}\-{n}\-{i}\-{\ng}\-{\r{g}}\-{\textscy}\-{\textsubring{r}}\-} \dicPos{m}[6] \dicFlx{(‑s, ‑ar)}[8] \dicDirectTranslationCS{vojenská uniforma}
\dicEntry[hermdarverk] \dicTerm{hermdar··verk} \dicIPA{{h}{\textepsilon}{r}{m}{\textsubring{d}}{a}{r}{v}{\textepsilon}{\textsubring{r}}{\r{g}}} \dicPos{n}[2] \dicFlx{(‑s, ‑)}[5] \textbf{1.} \dicSynonym{ódæði} \dicDirectTranslationCS{teroristický čin}  \textbf{2.} \dicSynonym{spellvirki\smash{\textsuperscript{1}}} \dicDirectTranslationCS{sabotáž}
\begin{xtolerant}{}{1pt}
\dicEntry[hermdarverkamaður] \dicTerm{hermdar·verka··|maður}\addthinSS\dicIPA{{h}\-{\textepsilon}\-{r}\-{m}\-{\textsubring{d}}\-{a}\-{r}\-{v}\-{\textepsilon}\-{\textsubring{r}}\-{\r{g}}\-{a}\-{m}\-{a}\-{ð}\-{\textscy}\-{\textsubring{r}}}\addthinSS\dicPos{m}[1\hspace{-0.10667em}3]\addthinSS\dicFlx{(‑manns, ‑menn)}[2] \dicDirectTranslationCS{terorista, teroristka}
\end{xtolerant}
\dicEntry[hermennska] \dicTerm{her··mennsk|a} \dicIPA{{h}{\textepsilon}{r}{m}{\textepsilon}{n}{s}{\r{g}}{a}} \dicPos{f}[1] \dicFlx{(‑u)}[5] \dicDirectTranslationCS{vojenská služba}
\dicEntry[hermir] \dicTerm{herm|ir} \dicIPA{{h}{\textepsilon}{r}{m}{\textsci}{\textsubring{r}}} \dicPos{m}[7] \dicFlx{(‑is, ‑ar)}[1] \textbf{1.} \dicDirectTranslationCS{simulátor} \dicExampleIS{flughermir} \dicExampleCS{letecký simulátor}  \textbf{2.} \dicFieldCat{poč.} \dicDirectTranslationCS{emulátor}
\begin{xtolerant}{}{1pt}
\dicEntry[hernaðarbandalag] \dicTerm{hernaðar··banda·|lag}\addthinS\dicIPA{{h}{\textepsilon}{r}{\textsubring{d}}{n}{a}{ð}{a}{r}{\textsubring{b}}{a}{n}{\textsubring{d}}{a}{l}{a}{x}}\addthinS\dicPos{n}[2]\addthinS\dicFlx{(‑lags, ‑lög)}[8] \dicDirectTranslationCS{vojenská aliance}
\end{xtolerant}
\dicEntry[hernaðarfræði] \dicTerm{hernaðar··fræð|i} \dicIPA{{h}{\textepsilon}{r}{\textsubring{d}}{n}{a}{ð}{a}{\textsubring{r}}{f}{r}{a}{i}{ð}{\textsci}} \dicPos{f}[3] \dicFlx{(‑i)}[3] \dicDirectTranslationCS{vojenská věda, vojenství}
\dicEntry[hernaðarkapphlaup] \dicTerm{hernaðar··kapp·hlaup} \dicIPA{{h}{\textepsilon}{r}{\textsubring{d}}{n}{a}{ð}{a}{\textsubring{r}}{k\smash{\textsuperscript{h}}}{a}{h}{\textsubring{b}}{\textsubring{l}}{\oe i}{\textsubring{b}}} \dicPos{n}[2] \dicFlx{(‑s, ‑)}[5] \dicDirectTranslationCS{závody ve zbrojení}
\begin{xtolerant}{}{1pt}
\dicEntry[hernaðarsérfræðingur] \dicTerm{hernaðar··sér·fræð·ing|ur}\addthinS\dicIPA{{h}\-{\textepsilon}\-{r}\-{\textsubring{d}}\-{n}\-{a}\-{ð}\-{a}\-{\textsubring{r}}\-{s}\-{j}\-{\textepsilon}\-{\textsubring{r}}\-{f}\-{r}\-{a}\-{i}\-{ð}\-{i}\-{\ng}\-{\r{g}}\-{\textscy}\-{\textsubring{r}}}\addthinS\dicPos{m}[6]\addthinS\dicFlx{(‑s, ‑ar)}[8] \dicDirectTranslationCS{vojenský specialista, vojenská specialistka}
\end{xtolerant}
\dicEntry[hernaðarstefna] \dicTerm{hernaðar··stefn|a} \dicIPA{{h}{\textepsilon}{r}{\textsubring{d}}{n}{a}{ð}{a}{\textsubring{r}}{s}{\textsubring{d}}{\textepsilon}{\textsubring{b}}{n}{a}} \dicPos{f}[1] \dicFlx{(‑u)}[5] \dicDirectTranslationCS{militarismus}
\dicEntry[hernaðarveldi] \dicTerm{hernaðar··veldi} \dicIPA{{h}{\textepsilon}{r}{\textsubring{d}}{n}{a}{ð}{a}{r}{v}{\textepsilon}{l}{\textsubring{d}}{\textsci}} \dicPos{n}[2] \dicFlx{(‑s, ‑)}[14] \dicDirectTranslationCS{vojenská moc}
\dicEntry[hernaður] \dicTerm{her··nað|ur} \dicIPA{{h}{\textepsilon}{r}{\textsubring{d}}{n}{a}{ð}{\textscy}{\textsubring{r}}} \dicPos{m}[10] \dicFlx{(‑ar)}[9] \dicDirectTranslationCS{válčení, válka}
\dicEntry[hernám] \dicTerm{her··nám} \dicIPA{{h}{\textepsilon}{r}{n}{au}{\textsubring{m}}} \dicPos{n}[2] \dicFlx{(‑s)}[2] \dicDirectTranslationCS{okupace, obsazení, zábor}
\dicEntry[heróín] \dicTerm{heróín} \dicIPA{{h}{\textepsilon}{\textlengthmark}{r}{ou}{i}{\textsubring{n}}} \dicPos{n}[2] \dicFlx{(‑s)}[2] \dicDirectTranslationCS{heroin}
\dicEntry[herpa] \dicTerm{herp|a} \dicIPA{{h}{\textepsilon}{\textsubring{r}}{\textsubring{b}}{a}} \dicPos{v}[2] \dicFlx{(‑ti, ‑t)}[63] \dicFlx{acc} \dicDirectTranslationCS{stisknout, sevřít} \dicExampleIS{herpa saman varirnar} \dicExampleCS{sevřít rty}
\dicEntry[herra] \dicTerm{herr|a} \dicsymFrequent\  \dicIPA{{h}{\textepsilon}{r}{\textlengthmark}{a}} \dicPos{m}[1] \dicFlx{(‑a, ‑ar)}[17] \dicDirectTranslationCS{pán, pan} \dicExampleIS{herra forseti} \dicExampleCS{pan prezident};  \dicPhraseIS{herrar mínir og frúr} \dicDirectTranslationCS{dámy a~pánové};  \dicPhraseIS{dömur mínar og herrar} \dicDirectTranslationCS{dámy a~pánové};  \dicIdiom{herra}{ \dicPhraseIS{herra minn trúr!}} \dicDirectTranslationCS{ježíšmarjá!, ježíškote!}
\dicEntry[herraföt] \dicTerm{herra··föt} \dicIPA{{h}{\textepsilon}{r}{\textlengthmark}{a}{f}{\oe}{\textsubring{d}}} \dicPos{n}[2] \dicFlx{pl}[9] \dicDirectTranslationCS{pánské oblečení}
\dicEntry[herralegur] \dicTerm{herra··legur} \dicIPA{{h}{\textepsilon}{r}{\textlengthmark}{a}{l}{\textepsilon}{\textbabygamma}{\textscy}{\textsubring{r}}} \dicPos{adj}[1]\dicFlx{}[-8] \dicDirectTranslationCS{galantní, zdvořilý, džentlmenský}
\dicEntry[herramaður] \dicTerm{herra··|maður} \dicIPA{{h}{\textepsilon}{r}{\textlengthmark}{a}{m}{a}{ð}{\textscy}{\textsubring{r}}} \dicPos{m}[13] \dicFlx{(‑manns, ‑menn)}[2] \dicDirectTranslationCS{džentlmen, (vážený) pán}
\dicEntry[Hersegóvína] \dicTerm{Hersegóvín|a} \dicIPA{{h}{\textepsilon}{\textsubring{r}}{s}{\textepsilon}{\textbabygamma}{ou}{v}{i}{n}{a}} \dicPos{f}[1] \dicFlx{(‑u)}[6] \dicFieldCat{geog.} \dicDirectTranslationCS{Hercegovina}
\dicEntry[hershöfðingi] \dicTerm{hers··höfð·ing|i} \dicIPA{{h}{\textepsilon}{\textsubring{r}}{s}{h}{\oe}{v}{ð}{i}{\textltailn}{\r{\textObardotlessj}}{\textsci}} \dicPos{m}[1] \dicFlx{(‑ja, ‑jar)}[14] \dicFieldCat{voj.} \dicDirectTranslationCS{generál}
\dicEntry[herskár] \dicTerm{her··skár} \dicIPA{{h}{\textepsilon}{\textsubring{r}}{s}{\r{g}}{au}{\textsubring{r}}} \dicPos{adj}[4]\dicFlx{}[-1] \dicDirectTranslationCS{bojovný, výbojný, militantní}
\dicEntry[herskip] \dicTerm{her··skip} \dicIPA{{h}{\textepsilon}{\textsubring{r}}{s}{\r{\textObardotlessj}}{\textsci}{\textsubring{b}}} \dicPos{n}[2] \dicFlx{(‑s, ‑)}[5] \dicDirectTranslationCS{vojenská\,/\addthin válečná loď}
\dicEntry[herskylda] \dicTerm{her··skyld|a} \dicIPA{{h}{\textepsilon}{\textsubring{r}}{s}{\r{\textObardotlessj}}{\textsci}{l}{\textsubring{d}}{a}} \dicPos{f}[1] \dicFlx{(‑u)}[5] \dicFieldCat{voj.} \dicDirectTranslationCS{branná povinnost}
\dicEntry[hersla] \dicTerm{hersl|a} \dicIPA{{h}{\textepsilon}{\textsubring{r}}{s}{\textsubring{d}}{l}{a}} \dicPos{f}[1] \dicFlx{(‑u, ‑ur)}[13] \textbf{1.} \dicSynonym*{hersla málms} \dicDirectTranslationCS{tvrzení (železa ap.)}  \textbf{2.} \dicSynonym*{hersla fisks} \dicDirectTranslationCS{sušení (ryb)}
\dicEntry[herslumunur] \dicTerm{herslu··mun|ur} \dicIPA{{h}{\textepsilon}{\textsubring{r}}{s}{\textsubring{d}}{l}{\textscy}{m}{\textscy}{n}{\textscy}{\textsubring{r}}} \dicPos{m}[10] \dicFlx{(‑ar)}[12] \dicPhraseIS{það vantar herslumuninn} \dicDirectTranslationCS{zbývá už jen kousíček (k~dokončení něčeho ap.)}
\dicEntry[herspítali] \dicTerm{her··spítal|i} \dicIPA{{h}{\textepsilon}{\textsubring{r}}{s}{\textsubring{b}}{i}{\textsubring{d}}{a}{l}{\textsci}} \dicPos{m}[1] \dicFlx{(‑a, ‑ar)}[8] \dicDirectTranslationCS{vojenská nemocnice}
\dicEntry[herstjóri] \dicTerm{her··stjór|i} \dicIPA{{h}{\textepsilon}{\textsubring{r}}{s}{\textsubring{d}}{j}{ou}{r}{\textsci}} \dicPos{m}[1] \dicFlx{(‑a, ‑ar)}[1] \dicSynonym{herforingi} \dicDirectTranslationCS{důstojník, důstojnice}
\dicEntry[herstöð] \dicTerm{her··stöð} \dicIPA{{h}{\textepsilon}{\textsubring{r}}{s}{\textsubring{d}}{\oe}{\texttheta}} \dicPos{f}[6] \dicFlx{(‑var, ‑var)}[1] \dicDirectTranslationCS{vojenská základna}
\dicEntry[hersveit] \dicTerm{her··sveit} \dicIPA{{h}{\textepsilon}{\textsubring{r}}{s}{v}{ei}{\textsubring{d}}} \dicPos{f}[7] \dicFlx{(‑ar, ‑ir)}[1] \dicDirectTranslationCS{vojenská jednotka}
\dicEntry[hertaka] \dicTerm{her··|taka} \dicIPA{{h}{\textepsilon}{\textsubring{r}}{t\smash{\textsuperscript{h}}}{a}{\r{g}}{a}} \dicPos{v}[6] \dicFlx{(‑tek, ‑tók, ‑tókum, ‑tæki, ‑tekið)}[127] \dicFlx{acc} \dicDirectTranslationCS{obsadit, dobýt (území ap.)} \dicExampleIS{hertaka borgina} \dicExampleCS{dobýt město}
\dicEntry[hertek] \dicTerm{her··tek} \dicIPA{{h}{\textepsilon}{\textsubring{r}}{t\smash{\textsuperscript{h}}}{\textepsilon}{\r{g}}} \dicPos{v} \dicFlx{ind praes sg 1 pers} \dicLink{hertaka}
\dicEntry[hertekið] \dicTerm{her··tekið} \dicIPA{{h}{\textepsilon}{\textsubring{r}}{t\smash{\textsuperscript{h}}}{\textepsilon}{\r{\textObardotlessj}}{\textsci}{\texttheta}} \dicPos{v} \dicFlx{supin} \dicLink{hertaka}
\dicEntry[hertekinn] \dicTerm{her··tekinn} \dicIPA{{h}{\textepsilon}{\textsubring{r}}{t\smash{\textsuperscript{h}}}{\textepsilon}{\r{\textObardotlessj}}{\textsci}{\textsubring{n}}} \dicPos{adj}[6]\dicFlx{}[-6] \dicDirectTranslationCS{obsazený, dobytý} \dicExampleIS{hertekin borg} \dicExampleCS{dobyté město}
\dicEntry[hertogadæmi] \dicTerm{her·toga··dæmi} \dicIPA{{h}{\textepsilon}{\textsubring{r}}{t\smash{\textsuperscript{h}}}{\textopeno}{\textbabygamma}{a}{\textsubring{d}}{a}{i}{m}{\textsci}} \dicPos{n}[2] \dicFlx{(‑s, ‑)}[14] \dicDirectTranslationCS{vévodství}
\dicEntry[hertogaynja] \dicTerm{her·toga··ynj|a} \dicIPA{{h}{\textepsilon}{\textsubring{r}}{t\smash{\textsuperscript{h}}}{\textopeno}{\textbabygamma}{a}{\textsci}{n}{j}{a}} \dicPos{f}[1] \dicFlx{(‑u, ‑ur)}[7] \dicDirectTranslationCS{vévodkyně}
\dicEntry[hertogi] \dicTerm{her··tog|i} \dicIPA{{h}{\textepsilon}{\textsubring{r}}{t\smash{\textsuperscript{h}}}{\textopeno i}{j}{\textsci}} \dicPos{m}[1] \dicFlx{(‑a, ‑ar)}[1] \dicDirectTranslationCS{vévoda}
\dicEntry[hertók] \dicTerm{her··tók} \dicIPA{{h}{\textepsilon}{\textsubring{r}}{t\smash{\textsuperscript{h}}}{ou}{\r{g}}} \dicPos{v} \dicFlx{ind pf sg 1 pers} \dicLink{hertaka}
\dicEntry[hertókum] \dicTerm{her··tókum} \dicIPA{{h}{\textepsilon}{\textsubring{r}}{t\smash{\textsuperscript{h}}}{ou}{\r{g}}{\textscy}{\textsubring{m}}} \dicPos{v} \dicFlx{ind pf pl 1 pers} \dicLink{hertaka}
\dicEntry[hertæki] \dicTerm{her··tæki} \dicIPA{{h}{\textepsilon}{\textsubring{r}}{t\smash{\textsuperscript{h}}}{a}{i}{\r{\textObardotlessj}}{\textsci}} \dicPos{v} \dicFlx{con pf sg 1 pers} \dicLink{hertaka}
\dicEntry[heslihneta] \dicTerm{hesli··hnet|a} \dicIPA{{h}{\textepsilon}{s}{\textsubring{d}}{l}{\textsci}{\textsubring{n}}{\textepsilon}{\textsubring{d}}{a}} \dicPos{f}[1] \dicFlx{(‑u, ‑ur)}[7] \dicDirectTranslationCS{lískový ořech\,/\addthin oříšek} \dicIndirectTranslationCS{(plod)}
\dicEntry[hespa] \dicTerm{hesp|a\smash{\textsuperscript{1}}} \dicIPA{{h}{\textepsilon}{s}{\textsubring{b}}{a}} \dicPos{f}[1] \dicFlx{(‑u, ‑ur)}[19] \textbf{1.} \dicDirectTranslationCS{přadeno, klubko}  \textbf{2.} \dicSynonym{loka\smash{\textsuperscript{1}}} \dicDirectTranslationCS{petlice, západka}
\dicEntry[hespa] \dicTerm{hesp|a\smash{\textsuperscript{2}}} \dicIPA{{h}{\textepsilon}{s}{\textsubring{b}}{a}} \dicPos{v}[1] \dicFlx{(‑aði)}[1] \dicFlx{acc} \textbf{1.} \dicDirectTranslationCS{navíjet, navinout (nit ap.)} \dicExampleIS{hespa reipi} \dicExampleCS{navinout provaz}  \textbf{2.} \dicDirectTranslationCS{zavřít na petlici (dveře ap.)} \dicExampleIS{hespa aftur hurð} \dicExampleCS{zavřít dveře na petlici};  \dicIdiom{hespa}[af]{ \dicPhraseIS{hespa e‑ð af}} \dicDirectTranslationCS{rychle (co) vyřídit\,/\addthin dokončit}
\dicEntry[hestaat] \dicTerm{hesta··|at} \dicIPA{{h}{\textepsilon}{s}{\textsubring{d}}{a}{a}{\textsubring{d}}} \dicPos{n}[2] \dicFlx{(‑ats, ‑öt)}[8] \dicDirectTranslationCS{koňský souboj}
\dicEntry[hestafl] \dicTerm{hest··|afl} \dicIPA{{h}{\textepsilon}{s}{\textsubring{d}}{a}{\textsubring{b}}{\textsubring{l}}} \dicPos{n}[2] \dicFlx{(‑afls, ‑öfl)}[8] \dicDirectTranslationCS{koňská síla}
\dicEntry[hestakastanía] \dicTerm{hesta··kastaní|a} \dicIPA{{h}{\textepsilon}{s}{\textsubring{d}}{a}{k\smash{\textsuperscript{h}}}{a}{s}{\textsubring{d}}{a}{n}{i}{j}{a}} \dicPos{f}[1] \dicFlx{(‑u, ‑ur)}[7] \dicFieldCat{bot.} \dicDirectTranslationCS{jírovec maďal, koňský kaštan} \textit{(l.~{\textLA{Aesculus hippocastanum}})}  \dicsymPhoto\ 
\dicFigure{ds_image_hestakastania_0_1.jpg}{Hestakastanía}{Hestakastanía - Alvesgaspar, CC BY 2.5}
\dicEntry[hestaleiga] \dicTerm{hesta··leig|a} \dicIPA{{h}{\textepsilon}{s}{\textsubring{d}}{a}{l}{ei}{\textbabygamma}{a}} \dicPos{f}[1] \dicFlx{(‑u, ‑ur)}[19] \dicDirectTranslationCS{pronájem\,/\addthin půjčovna koní}
\dicEntry[hestamennska] \dicTerm{hesta··mennsk|a} \dicIPA{{h}{\textepsilon}{s}{\textsubring{d}}{a}{m}{\textepsilon}{n}{s}{\r{g}}{a}} \dicPos{f}[1] \dicFlx{(‑u)}[5] \dicDirectTranslationCS{jezdectví} \dicIndirectTranslationCS{(hobby a~sport)}
\dicEntry[hestaskál] \dicTerm{hesta··skál} \dicIPA{{h}{\textepsilon}{s}{\textsubring{d}}{a}{s}{\r{g}}{au}{\textsubring{l}}} \dicPos{f}[4] \dicFlx{(‑ar, ‑ar\,/\addthin ‑ir)}[30] \dicDirectTranslationCS{přípitek na zdraví} \dicIndirectTranslationCS{(před jízdou na koni)}
\dicEntry[hestbak] \dicTerm{hest··|bak} \dicIPA{{h}{\textepsilon}{s}{\textsubring{d}}{\textsubring{b}}{a}{\r{g}}} \dicPos{n}[2] \dicFlx{(‑baks, ‑bök)}[8] \dicDirectTranslationCS{koňský hřbet};  \dicPhraseIS{fara á hestbak} \dicDirectTranslationCS{jít jezdit (na koni), jet na vyjížďku (na koni)};  \dicPhraseIS{ferðast á hestbaki} \dicFlx{refl} \dicDirectTranslationCS{jet na koni}
\dicEntry[hesthús] \dicTerm{hest··hús} \dicIPA{{h}{\textepsilon}{s}{\textsubring{d}}{h}{u}{s}} \dicPos{n}[2] \dicFlx{(‑s, ‑)}[5] \dicDirectTranslationCS{stáj, konírna}
\dicEntry[hestur] \dicTerm{hest|ur} \dicsymFrequent\  \dicIPA{{h}{\textepsilon}{s}{\textsubring{d}}{\textscy}{\textsubring{r}}} \dicPos{m}[6] \dicFlx{(‑s, ‑ar)}[4] \textbf{1.} \dicFieldCat{zool.} \dicDirectTranslationCS{kůň} \textit{(l.~{\textLA{Equus}})}  \dicsymPhoto\  \dicExampleIS{Íslenski hesturinn er af mongólskum uppruna.} \dicExampleCS{Islandský kůň je mongolského původu.};  \dicPhraseIS{setja sig á háan hest} \dicLangCat{přen.} \dicSynonym*{gera sig merkilegan} \dicDirectTranslationCS{povyšovat se};  \dicPhraseIS{ríða ekki feitum hesti frá e‑u} \dicLangCat{přen.} \dicDirectTranslationCS{nezbohatnout z~(čeho)};  \dicPhraseIS{leiða saman hesta sína} \dicLangCat{přen.} \dicDirectTranslationCS{utkat se, střetnout se}  \textbf{2.} \dicFieldCat{sport.} \dicDirectTranslationCS{kůň} \dicIndirectTranslationCS{(sportovní náčiní v~gymnastice)}
\dicFigure{ds_image_hestur_0_1.jpg}{Hestur}{Hestur - Dorota Nierychlewska-Chejn, GPL}
\dicEntry[hestvagn] \dicTerm{hest··vagn} \dicIPA{{h}{\textepsilon}{s}{\textsubring{d}}{v}{a}{\r{g}}{\textsubring{n}}} \dicPos{m}[4] \dicFlx{(‑s, ‑ar)}[4] \dicDirectTranslationCS{kočár, drožka, dostavník}
\dicEntry[hetja] \dicTerm{hetj|a} \dicIPA{{h}{\textepsilon}{\textlengthmark}{\textsubring{d}}{j}{a}} \dicPos{f}[1] \dicFlx{(‑u, ‑ur)}[7] \dicDirectTranslationCS{hrdina, hrdinka}
\dicEntry[hetjudáð] \dicTerm{hetju··dáð} \dicIPA{{h}{\textepsilon}{\textlengthmark}{\textsubring{d}}{j}{\textscy}{\textsubring{d}}{au}{\texttheta}} \dicPos{f}[7] \dicFlx{(‑ar, ‑ir)}[1] \dicDirectTranslationCS{hrdinský čin\,/\addthin skutek}
\dicEntry[hetjukvæði] \dicTerm{hetju··kvæði} \dicIPA{{h}{\textepsilon}{\textlengthmark}{\textsubring{d}}{j}{\textscy}{k\smash{\textsuperscript{h}}}{v}{a}{i}{ð}{\textsci}} \dicPos{n}[2] \dicFlx{(‑s, ‑)}[14] \dicFieldCat{lit.} \dicDirectTranslationCS{hrdinská báseň}
\dicEntry[hetjulegur] \dicTerm{hetju··legur} \dicIPA{{h}{\textepsilon}{\textlengthmark}{\textsubring{d}}{j}{\textscy}{l}{\textepsilon}{\textbabygamma}{\textscy}{\textsubring{r}}} \dicPos{adj}[1]\dicFlx{}[-8] \dicDirectTranslationCS{hrdinský, heroický}
\dicEntry[hetjusaga] \dicTerm{hetju··|saga} \dicIPA{{h}{\textepsilon}{\textlengthmark}{\textsubring{d}}{j}{\textscy}{s}{a}{\textbabygamma}{a}} \dicPos{f}[1] \dicFlx{(‑sögu, ‑sögur)}[14] \dicFieldCat{lit.} \dicDirectTranslationCS{hrdinský příběh}
\dicEntry[hetjuskapur] \dicTerm{hetju··skap|ur} \dicIPA{{h}{\textepsilon}{\textlengthmark}{\textsubring{d}}{j}{\textscy}{s}{\r{g}}{a}{\textsubring{b}}{\textscy}{\textsubring{r}}} \dicPos{m}[10] \dicFlx{(‑ar)}[15] \dicDirectTranslationCS{hrdinství}
\dicEntry[hetta] \dicTerm{hett|a} \dicIPA{{h}{\textepsilon}{h}{\textsubring{d}}{a}} \dicPos{f}[1] \dicFlx{(‑u, ‑ur)}[7] \textbf{1.} \dicDirectTranslationCS{kapuce, kápě} \dicExampleIS{hettuúlpa} \dicExampleCS{bunda s~kapucí};  \dicPhraseIS{vera gamall í hettunni} \dicLangCat{přen.} \dicDirectTranslationCS{být mazák}  \textbf{2.} \dicDirectTranslationCS{pesar} \dicIndirectTranslationCS{(druh antikoncepčního prostředku)}
\dicEntry[hettir] \dicTerm{hettir} \dicIPA{{h}{\textepsilon}{h}{\textsubring{d}}{\textsci}{\textsubring{r}}} \dicPos{m} \dicFlx{pl nom} \dicLink{höttur}
\dicEntry[hettumávur] \dicTerm{hettu··máv|ur} \dicIPA{{h}{\textepsilon}{h}{\textsubring{d}}{\textscy}{m}{au}{v}{\textscy}{\textsubring{r}}} \dicPos{m}[6] \dicFlx{(‑s, ‑ar)}[19] \dicFieldCat{zool.} \dicDirectTranslationCS{racek chechtavý} \textit{(l.~{\textLA{Chroicocephalus ridibundus}})}  \dicsymPhoto\ 
\dicFigure{597.jpg}{Hettumávur}{Hettumávur - Zicha Ondřej, Biolib, Copyright/CC-BY-NC}
\dicEntry[hettuskrofa] \dicTerm{hettu··skrof|a} \dicIPA{{h}{\textepsilon}{h}{\textsubring{d}}{\textscy}{s}{\r{g}}{r}{\textopeno}{v}{a}} \dicPos{f}[1] \dicFlx{(‑u, ‑ur)}[19] \dicFieldCat{zool.} \dicDirectTranslationCS{buřňák velký} \textit{(l.~{\textLA{Puffinus gravis}})}  \dicsymPhoto\ 
\dicFigure{ds_image_hettuskrofa_0_2.jpg}{Hettuskrofa}{Hettuskrofa - Patrick Coin, cc-by-sa-2.5}
\dicEntry[hettusótt] \dicTerm{hettu··sótt} \dicIPA{{h}{\textepsilon}{h}{\textsubring{d}}{\textscy}{s}{ou}{h}{\textsubring{d}}} \dicPos{f}[7] \dicFlx{(‑ar)}[3] \dicFieldCat{med.} \dicDirectTranslationCS{příušnice}
\dicEntry[hettusöngvari] \dicTerm{hettu··söngv·ar|i} \dicIPA{{h}{\textepsilon}{h}{\textsubring{d}}{\textscy}{s}{\oe i}{\ng}{\r{g}}{v}{a}{r}{\textsci}} \dicPos{m}[1] \dicFlx{(‑a, ‑ar)}[13] \dicFieldCat{zool.} \dicDirectTranslationCS{pěnice černohlavá} \textit{(l.~{\textLA{Sylvia atricapilla}})}  \dicsymPhoto\ 
\dicFigure{ds_image_hettusongvari_0_1.jpg}{Hettusöngvari}{Hettusöngvari - Michael Apel, CC BY-SA 3.0}
\dicEntry[hettuúlpa] \dicTerm{hettu··úlp|a} \dicIPA{{h}{\textepsilon}{h}{\textsubring{d}}{\textscy}{u}{\textsubring{l}}{\textsubring{b}}{a}} \dicPos{f}[1] \dicFlx{(‑u, ‑ur)}[13] \dicDirectTranslationCS{bunda s~kapucí}
\dicEntry[hey] \dicTerm{hey\smash{\textsuperscript{1}}} \dicsymFrequent\  \dicIPA{{h}{ei}{\textlengthmark}} \dicPos{n}[2] \dicFlx{(‑s, ‑)}[13] \dicDirectTranslationCS{seno} \dicExampleIS{sækja hey handa kúnum} \dicExampleCS{dojít kravám pro seno};  \dicProverb\  \dicPhraseIS{Allt er hey í harðindum.} \dicLangCat{přís.} \dicDirectTranslationCS{V~nouzi přijde všechno vhod.}
\dicEntry[hey] \dicTerm{hey\smash{\textsuperscript{2}}} \dicIPA{{h}{ei}{\textlengthmark}} \dicPos{v} \dicFlx{ind praes sg 1 pers} \dicLink{heyja\smash{\textsuperscript{2}}}
\dicEntry[heybindivél] \dicTerm{hey·bindi··vél} \dicIPA{{h}{ei}{\textlengthmark}{\textsubring{b}}{\textsci}{n}{\textsubring{d}}{\textsci}{v}{j}{\textepsilon}{\textsubring{l}}} \dicPos{f}[4] \dicFlx{(‑ar, ‑ar)}[1] \dicDirectTranslationCS{lis} \dicIndirectTranslationCS{(stroj na lisování sena do balíků)}
\dicEntry[heygarðshorn] \dicTerm{hey·garðs··horn} \dicIPA{{h}{ei}{\textlengthmark}{\r{g}}{a}{r}{ð}{s}{h}{\textopeno}{r}{\textsubring{d}}{\textsubring{n}}} \dicPos{n}[2] \dicFlx{(‑s)}[2] \dicPhraseIS{vera við sama heygarðshornið} \dicLangCat{přen.} \dicDirectTranslationCS{pokračovat ve vyjetých kolejích}
\dicEntry[heyi] \dicTerm{heyi} \dicIPA{{h}{ei}{\textlengthmark}{\textsci}} \dicPos{v} \dicFlx{ind praes sg 1 pers} \dicLink{heyja\smash{\textsuperscript{2}}}
\dicEntry[heyja] \dicTerm{heyj|a\smash{\textsuperscript{1}}} \dicIPA{{h}{ei}{j}{\textlengthmark}{a}} \dicPos{v}[1] \dicFlx{(‑aði)}[30] \dicDirectTranslationCS{sklízet seno} \dicExampleIS{heyja túnið} \dicExampleCS{sklízet seno z~louky};  \dicIdiom{heyjast}{ \dicPhraseIS{það heyjast}} \dicFlx{refl impers} \dicDirectTranslationCS{(seno) se sklízí} \dicExampleIS{Það heyjast vel í dag.} \dicExampleCS{Dnes se dobře sklízí seno.}
\dicEntry[heyja] \dicTerm{heyja\smash{\textsuperscript{2}}} \dicIPA{{h}{ei}{j}{\textlengthmark}{a}} \dicPos{v}[5] \dicFlx{(hey\,/\addthin heyi, háði, háðum, háði, háð)}[12] \dicFlx{acc} \dicDirectTranslationCS{vést (kampaň ap.)} \dicExampleIS{heyja orustu} \dicExampleCS{vést bitvu}
\dicEntry[heykjast] \dicTerm{heyk|jast} \dicIPA{{h}{ei}{\textlengthmark}{\r{\textObardotlessj}}{a}{s}{\textsubring{d}}} \dicPos{v}[2] \dicFlx{(‑tist, ‑st)}[192] \dicFlx{refl} \dicPhraseIS{heykjast á e‑u} \dicSynonym{guggna} \dicDirectTranslationCS{vycouvat z~(čeho), ustoupit od (čeho) (kandidatury ap.)}
\dicEntry[heykvísl] \dicTerm{hey··kvísl} \dicIPA{{h}{ei}{\textlengthmark}{k\smash{\textsuperscript{h}}}{v}{i}{s}{\textsubring{d}}{\textsubring{l}}} \dicPos{f}[4] \dicFlx{(‑ar, ‑ar\,/\addthin ‑ir)}[30] \dicDirectTranslationCS{vidle}
\dicEntry[heyr] \dicTerm{heyr} \dicIPA{{h}{ei}{\textlengthmark}{\textsubring{r}}} \dicPos{inter} \dicDirectTranslationCS{bravo} \dicIndirectTranslationCS{(volat bravo při znamenitém proslovu ap.)}
\dicEntry[heyra] \dicTerm{heyr|a} \dicsymFrequent\  \dicIPA{{h}{ei}{\textlengthmark}{r}{a}} \dicPos{v}[2] \dicFlx{(‑ði, ‑t)}[104] \dicFlx{acc} \textbf{1.} \dicDirectTranslationCS{(u)slyšet, zaslechnout} \dicIndirectTranslationCS{(sluchem postihnout)} \dicExampleIS{heyra mannamál} \dicExampleCS{uslyšet lidské hlasy}  \textbf{2.} \dicDirectTranslationCS{(u)slyšet, zaslechnout, dovědět se z~doslechu};  \dicPhraseIS{hafa heyrt e‑u fleygt} \dicDirectTranslationCS{zaslechnout (co), slyšet zvěsti o~(čem)}  \textbf{3.} \dicSynonym{bænheyra} \dicDirectTranslationCS{vyslyšet (prosbu ap.)}  \textbf{4.} \dicPhraseIS{fá að heyra það} \dicLangCat{přen.} \dicDirectTranslationCS{říci to přímo, říci to bez obalu};  \dicPhraseIS{láta e‑n heyra það} \dicLangCat{přen.} \dicDirectTranslationCS{vynadat (komu), vyhubovat (komu)};  \dicIdiom{heyra}[til]{ \dicPhraseIS{það heyrir til}} \dicDirectTranslationCS{patří to k~(Vánocům ap.)}; { \dicPhraseIS{það heyrir til undantekninga}} \dicDirectTranslationCS{patří to k~výjimkám};  \dicIdiom{heyra}[undir]{ \dicPhraseIS{heyra undir e‑ð}} \dicDirectTranslationCS{náležet do (čeho), spadat pod (co)};  \dicIdiom{heyrast}{ \dicPhraseIS{heyrast}} \dicFlx{refl} \dicDirectTranslationCS{být slyšet}; { \dicPhraseIS{það heyrist (illa) í e‑u}} \dicFlx{refl impers} \dicDirectTranslationCS{(co) je (špatně) slyšet}; { \dicPhraseIS{það heyrist (vel) til (e‑rs)}} \dicFlx{refl impers} \dicDirectTranslationCS{(kdo) je (dobře) slyšet}; { \dicPhraseIS{það heyrðist ekki frá e‑m (árum saman)}} \dicFlx{refl impers} \dicDirectTranslationCS{nebylo o~(kom) (léta) slyšet}; { \dicPhraseIS{e‑m heyrist}} \dicFlx{refl impers} \dicDirectTranslationCS{(kdo) má za to (že slyšel)}
\dicEntry[heyrandi] \dicTerm{heyr··|andi} \dicIPA{{h}{ei}{\textlengthmark}{r}{a}{n}{\textsubring{d}}{\textsci}} \dicPos{m}[2] \dicFlx{(‑anda, ‑endur)}[1] \dicDirectTranslationCS{posluchač(ka)};  \dicPhraseIS{í heyranda hljóði} \dicFlx{adv} \dicDirectTranslationCS{veřejně, nahlas};  \dicProverb\  \dicPhraseIS{Oft er í holti heyrandi nær.} \dicLangCat{přís.} \dicDirectTranslationCS{I~stěny mají uši.}
\dicEntry[heyranlegur] \dicTerm{heyran··legur} \dicIPA{{h}{ei}{\textlengthmark}{r}{a}{n}{l}{\textepsilon}{\textbabygamma}{\textscy}{\textsubring{r}}} \dicPos{adj}[1]\dicFlx{}[-8] \dicDirectTranslationCS{slyšitelný}
\dicEntry[heyrn] \dicTerm{heyrn} \dicIPA{{h}{ei}{r}{\textsubring{d}}{\textsubring{n}}} \dicPos{f}[7] \dicFlx{(‑ar, ‑ir)}[1] \dicDirectTranslationCS{sluch} \dicExampleIS{hafa góða heyrn} \dicExampleCS{mít dobrý sluch}
\dicEntry[heyrnarlaus] \dicTerm{heyrnar··laus} \dicIPA{{h}{ei}{r}{\textsubring{d}}{n}{a}{r}{l}{\oe i}{s}} \dicPos{adj}[5]\dicFlx{}[-1] \dicDirectTranslationCS{hluchý, neslyšící}
\dicEntry[heyrnarleysi] \dicTerm{heyrnar··leysi} \dicIPA{{h}{ei}{r}{\textsubring{d}}{n}{a}{r}{l}{ei}{s}{\textsci}} \dicPos{n}[2] \dicFlx{(‑s)}[20] \dicDirectTranslationCS{hluchota}
\dicEntry[heyrnartap] \dicTerm{heyrnar··tap} \dicIPA{{h}{ei}{r}{\textsubring{d}}{n}{a}{\textsubring{r}}{t\smash{\textsuperscript{h}}}{a}{\textsubring{b}}} \dicPos{n}[2] \dicFlx{(‑s)}[2] \dicDirectTranslationCS{ztráta sluchu}
\dicEntry[heyrnartól] \dicTerm{heyrnar··tól} \dicIPA{{h}{ei}{r}{\textsubring{d}}{n}{a}{\textsubring{r}}{t\smash{\textsuperscript{h}}}{ou}{\textsubring{l}}} \dicPos{n}[2] \dicFlx{(‑s, ‑)}[5] \dicDirectTranslationCS{sluchátko} \dicIndirectTranslationCS{(většinou v~pl)}
\dicEntry[heyrnartæki] \dicTerm{heyrnar··tæki} \dicIPA{{h}{ei}{r}{\textsubring{d}}{n}{a}{\textsubring{r}}{t\smash{\textsuperscript{h}}}{a}{i}{\r{\textObardotlessj}}{\textsci}} \dicPos{n}[2] \dicFlx{(‑s, ‑)}[16] \dicDirectTranslationCS{naslouchátko, naslouchací přístroj}
\dicEntry[heyrnleysingi] \dicTerm{heyrn··leys·ing|i} \dicIPA{{h}{ei}{r}{\textsubring{d}}{\textsubring{n}}{l}{ei}{s}{i}{\textltailn}{\r{\textObardotlessj}}{\textsci}} \dicPos{m}[1] \dicFlx{(‑ja, ‑jar)}[14] \dicDirectTranslationCS{neslyšící} \dicIndirectTranslationCS{(osoba)}
\dicEntry[heyskapur] \dicTerm{hey··skap|ur} \dicIPA{{h}{ei}{\textlengthmark}{s}{\r{g}}{a}{\textsubring{b}}{\textscy}{\textsubring{r}}} \dicPos{m}[10] \dicFlx{(‑ar)}[15] \dicDirectTranslationCS{senoseč}
\dicEntry[heystakkur] \dicTerm{hey··stakk|ur} \dicIPA{{h}{ei}{\textlengthmark}{s}{\textsubring{d}}{a}{h}{\r{g}}{\textscy}{\textsubring{r}}} \dicPos{m}[6] \dicFlx{(‑s, ‑ar)}[25] \dicDirectTranslationCS{stoh sena}
\dicEntry[heyvagn] \dicTerm{hey··vagn} \dicIPA{{h}{ei}{\textlengthmark}{v}{a}{\r{g}}{\textsubring{n}}} \dicPos{m}[4] \dicFlx{(‑s, ‑ar)}[4] \dicDirectTranslationCS{vůz na seno}
\dicEntry[héðan] \dicTerm{héðan} \dicsymFrequent\  \dicIPA{{\c{c}}{\textepsilon}{\textlengthmark}{ð}{a}{\textsubring{n}}} \dicPos{adv} \dicDirectTranslationCS{odsud, odtud} \dicExampleIS{Við flugum héðan snemma að morgni.} \dicExampleCS{Odletěli jsme odtud brzy ráno.};  \dicIdiom{héðan}[af]{ \dicPhraseIS{héðan af}} \dicFlx{adv} \dicDirectTranslationCS{od teď, od této chvíle};  \dicIdiom{héðan}[í frá]{ \dicPhraseIS{héðan í frá}} \dicFlx{adv} \dicDirectTranslationCS{od teď, od této chvíle, od nynějška}
\dicEntry[hégómagirnd] \dicTerm{hé·góma··girnd} \dicIPA{{\c{c}}{\textepsilon}{\textlengthmark}{\r{g}}{ou}{m}{a}{\r{\textObardotlessj}}{\textsci}{r}{n}{\textsubring{d}}} \dicPos{f}[7] \dicFlx{(‑ar)}[3] \dicDirectTranslationCS{marnivost}
\dicEntry[hégómaskapur] \dicTerm{hé·góma··skap|ur} \dicIPA{{\c{c}}{\textepsilon}{\textlengthmark}{\r{g}}{ou}{m}{a}{s}{\r{g}}{a}{\textsubring{b}}{\textscy}{\textsubring{r}}} \dicPos{m}[10] \dicFlx{(‑ar)}[15] \dicDirectTranslationCS{marnivost}
\dicEntry[hégómi] \dicTerm{hé··góm|i} \dicIPA{{\c{c}}{\textepsilon}{\textlengthmark}{\r{g}}{ou}{m}{\textsci}} \dicPos{m}[1] \dicFlx{(‑a, ‑ar)}[1] \dicDirectTranslationCS{marnivost} \dicExampleIS{eintómur hégómi} \dicExampleCS{pouhopouhá marnivost};  \dicPhraseIS{leggja nafn guðs við hégóma} \dicDirectTranslationCS{brát jméno boží nadarmo}
\dicEntry[hégómlegur] \dicTerm{hé·góm··legur} \dicIPA{{\c{c}}{\textepsilon}{\textlengthmark}{\r{g}}{ou}{m}{l}{\textepsilon}{\textbabygamma}{\textscy}{\textsubring{r}}} \dicPos{adj}[1]\dicFlx{}[-8] \textbf{1.} \dicDirectTranslationCS{marnivý, pošetilý}  \textbf{2.} \dicSynonym{ómerkilegur} \dicDirectTranslationCS{okázalý, bezvýznamný}
\dicEntry[hégómleiki] \dicTerm{hé·góm··leik|i} \dicIPA{{\c{c}}{\textepsilon}{\textlengthmark}{\r{g}}{ou}{m}{l}{ei}{\r{\textObardotlessj}}{\textsci}} \dicPos{m}[1] \dicFlx{(‑a)}[3] \dicDirectTranslationCS{marnivost}
\dicEntry[hékk] \dicTerm{hékk} \dicIPA{{\c{c}}{\textepsilon}{h}{\r{g}}} \dicPos{v} \dicFlx{ind pf sg 1 pers} \dicLink{hanga}
\dicEntry[héla] \dicTerm{hél|a\smash{\textsuperscript{1}}} \dicIPA{{\c{c}}{\textepsilon}{\textlengthmark}{l}{a}} \dicPos{f}[1] \dicFlx{(‑u)}[5] \dicSynonym{hrím} \dicDirectTranslationCS{jinovatka} \dicExampleIS{héla á rúðum} \dicExampleCS{jinovatka na okenních tabulkách}
\dicEntry[héla] \dicTerm{hél|a\smash{\textsuperscript{2}}} \dicIPA{{\c{c}}{\textepsilon}{\textlengthmark}{l}{a}} \dicPos{v}[1] \dicFlx{(‑aði)}[80] \dicFlx{impers} \dicPhraseIS{e‑ð hélar} \dicDirectTranslationCS{(co) se pokrývá jinovatkou} \dicExampleIS{Gluggann hélaði.} \dicExampleCS{Okno se pokrylo jinovatkou.};  \dicPhraseIS{það hefur hélað (í nótt)} \dicDirectTranslationCS{padla jinovatka (v~noci)}
\dicEntry[héldi] \dicTerm{héldi} \dicIPA{{\c{c}}{\textepsilon}{l}{\textsubring{d}}{\textsci}} \dicPos{v} \dicFlx{con pf sg 1 pers} \dicLink{halda\smash{\textsuperscript{2}}}
\dicEntry[héldum] \dicTerm{héldum} \dicIPA{{\c{c}}{\textepsilon}{l}{\textsubring{d}}{\textscy}{\textsubring{m}}} \dicPos{v} \dicFlx{ind pf pl 1 pers} \dicLink{halda\smash{\textsuperscript{2}}}
\dicEntry[hélt] \dicTerm{hélt} \dicIPA{{\c{c}}{\textepsilon}{\textsubring{l}}{\textsubring{d}}} \dicPos{v} \dicFlx{ind pf sg 1 pers} \dicLink{halda\smash{\textsuperscript{2}}}
\dicEntry[héngi] \dicTerm{héngi} \dicIPA{{\c{c}}{\textepsilon}{\textltailn}{\r{\textObardotlessj}}{\textsci}} \dicPos{v} \dicFlx{con pf sg 1 pers} \dicLink{hanga}
\dicEntry[héngum] \dicTerm{héngum} \dicIPA{{\c{c}}{\textepsilon}{\ng}{\r{g}}{\textscy}{\textsubring{m}}} \dicPos{v} \dicFlx{ind pf pl 1 pers} \dicLink{hanga}
\dicEntry[hér] \dicTerm{hér} \dicsymFrequent\  \dicIPA{{\c{c}}{\textepsilon}{\textlengthmark}{\textsubring{r}}} \dicPos{adv} \textbf{1.} \dicSynonym{hérna} \dicDirectTranslationCS{zde, tady, tu} \dicExampleIS{Hann býr hér.} \dicExampleCS{Bydlí tady.}  \textbf{2.} \dicSynonym{hingað} \dicDirectTranslationCS{sem}  \textbf{3.} \dicSynonym*{þessi hér} \dicDirectTranslationCS{tady} \dicIndirectTranslationCS{(při ukazování na něco)}  \textbf{4.} \dicDirectTranslationCS{tady, prosím} \dicIndirectTranslationCS{(při předávání něčeho někomu)};  \dicIdiom{hér}{ \dicPhraseIS{hér að auki}} \dicFlx{adv} \dicDirectTranslationCS{navíc};  \dicPhraseIS{hér eftir} \dicFlx{adv} \dicDirectTranslationCS{od teď, od této chvíle};  \dicPhraseIS{hér með} \dicFlx{adv} \dicDirectTranslationCS{tímto};  \dicPhraseIS{hér og þar} \dicFlx{adv} \dicDirectTranslationCS{tady a~tam, na různých místech}; { \dicPhraseIS{hér um bil}} \dicFlx{adv} \dicSynonym{nærri\smash{\textsuperscript{1}}} \dicDirectTranslationCS{přibližně, asi tak, skoro}
\dicEntry[hérað] \dicTerm{hér|að} \dicsymFrequent\  \dicIPA{{\c{c}}{\textepsilon}{\textlengthmark}{r}{a}{\texttheta}} \dicPos{n}[3] \dicFlx{(‑aðs, ‑uð\,/\addthin ‑öð)}[4] \dicDirectTranslationCS{kraj, region, provincie} \dicExampleIS{ferðast um allt héraðið} \dicExampleCS{cestovat po celém kraji}
\dicEntry[héraðsdómari] \dicTerm{héraðs··dóm·ar|i} \dicIPA{{\c{c}}{\textepsilon}{\textlengthmark}{r}{a}{ð}{s}{\textsubring{d}}{ou}{m}{a}{r}{\textsci}} \dicPos{m}[1] \dicFlx{(‑a, ‑ar)}[13] \dicFieldCat{práv.} \dicDirectTranslationCS{krajský soudce, krajská soudkyně}
\dicEntry[héraðsdómur] \dicTerm{héraðs··dóm|ur} \dicIPA{{\c{c}}{\textepsilon}{\textlengthmark}{r}{a}{ð}{s}{\textsubring{d}}{ou}{m}{\textscy}{\textsubring{r}}} \dicPos{m}[6] \dicFlx{(‑s, ‑ar)}[10] \dicFieldCat{práv.} \dicDirectTranslationCS{krajský soud}
\dicEntry[héraðslæknir] \dicTerm{héraðs··lækn|ir} \dicIPA{{\c{c}}{\textepsilon}{\textlengthmark}{r}{a}{ð}{s}{l}{a}{i}{h}{\r{g}}{n}{\textsci}{\textsubring{r}}} \dicPos{m}[7] \dicFlx{(‑is, ‑ar)}[1] \dicDirectTranslationCS{krajský lékař, krajská lékařka}
\dicEntry[héri] \dicTerm{hér|i} \dicIPA{{\c{c}}{\textepsilon}{\textlengthmark}{r}{\textsci}} \dicPos{m}[1] \dicFlx{(‑a, ‑ar)}[1] \dicFieldCat{zool.} \dicDirectTranslationCS{zajíc} \textit{(l.~{\textLA{Lepus}})}  \dicsymPhoto\ 
\dicFigure{ds_image_heri_0_1.jpg}{Héri}{Héri - Fmickan, CC BY-SA 3.0}
\dicEntry[hérlendis] \dicTerm{hér··lendis} \dicsymFrequent\  \dicIPA{{\c{c}}{\textepsilon}{r}{l}{\textepsilon}{n}{\textsubring{d}}{\textsci}{s}} \dicPos{adv} \dicDirectTranslationCS{v~této zemi} \dicExampleIS{kvenréttindi hérlendis} \dicExampleCS{rovnoprávnost žen v~této zemi}
\dicEntry[hérna] \dicTerm{hérna} \dicsymFrequent\  \dicIPA{{\c{c}}{\textepsilon}{r}{\textsubring{d}}{n}{a}} \dicPos{adv} \textbf{1.} \dicSynonym{hér} \dicDirectTranslationCS{zde, tady, tadyhle, tu} \dicIndirectTranslationCS{(je možné na místo ukázat, určitější než příslovce \clqq hér\crqq )} \dicExampleIS{Býrðu hérna?} \dicExampleCS{Bydlíš tady?}  \textbf{2.} \dicSynonym{hingað} \dicDirectTranslationCS{sem}  \textbf{3.} \dicIndirectTranslationCS{(k~přidání důrazu)} \dicExampleIS{þessi hérna saga} \dicExampleCS{tady to vyprávění}  \textbf{4.} \dicIndirectTranslationCS{(jako slovní výplň)}  \textbf{5.} \dicDirectTranslationCS{tady, prosím} \dicExampleIS{Hérna, taktu peningana.} \dicExampleCS{Tady (prosím), vezmi si peníze.}
\dicEntry[héruð] \dicTerm{héruð} \dicIPA{{\c{c}}{\textepsilon}{\textlengthmark}{r}{\textscy}{\texttheta}} \dicPos{n} \dicFlx{pl nom} \dicLink{hérað}
\dicEntry[hét] \dicTerm{hét} \dicIPA{{\c{c}}{\textlengthmark}{\textepsilon}{\textsubring{d}}} \dicPos{v} \dicFlx{ind pf sg 1 pers} \dicLink{heita}
\dicEntry[héti] \dicTerm{héti} \dicIPA{{\c{c}}{\textlengthmark}{\textepsilon}{\textsubring{d}}{\textsci}} \dicPos{v} \dicFlx{con pf sg 1 pers} \dicLink{heita}
\dicEntry[hétum] \dicTerm{hétum} \dicIPA{{\c{c}}{\textlengthmark}{\textepsilon}{\textsubring{d}}{\textscy}{\textsubring{m}}} \dicPos{v} \dicFlx{ind pf pl 1 pers} \dicLink{heita}
\dicEntry[hf.] \dicTerm{hf.} \dicPos{zkr} \dicPhraseIS{hlutafélag} \dicDirectTranslationCS{akciová společnost}
\dicEntry[hið] \dicTerm{hið} \dicIPA{{h}{\textsci}{\textlengthmark}{\texttheta}} \dicPos{part} \dicLink{hinn\smash{\textsuperscript{2}}}
\dicEntry[hik] \dicTerm{hik} \dicIPA{{h}{\textsci}{\textlengthmark}{\r{g}}} \dicPos{n}[2] \dicFlx{(‑s)}[2] \dicDirectTranslationCS{(za)váhání, váhavost} \dicExampleIS{Það kom hik á mig.} \dicExampleCS{Zaváhal jsem.};  \dicPhraseIS{það er hik á e‑m} \dicDirectTranslationCS{(kdo) váhá}
\dicEntry[hika] \dicTerm{hik|a} \dicsymFrequent\  \dicIPA{{h}{\textsci}{\textlengthmark}{\r{g}}{a}} \dicPos{v}[1] \dicFlx{(‑aði)}[44] \dicSynonym*{draga við sig} \dicDirectTranslationCS{(za)váhat, rozpakovat se, kolísat} \dicExampleIS{Hún hikaði andartak.} \dicExampleCS{Na okamžik zaváhala.};  \dicPhraseIS{hika við e‑ð} \dicDirectTranslationCS{váhat s~(čím)};  \dicProverb\  \dicPhraseIS{Að hika er sama og tapa.} \dicLangCat{přís.} \dicDirectTranslationCS{Váhat je totéž co prohrát.}
\dicEntry[hikandi] \dicTerm{hik··andi} \dicIPA{{h}{\textsci}{\textlengthmark}{\r{g}}{a}{n}{\textsubring{d}}{\textsci}} \dicPos{adj}[13] \dicFlx{indecl}[1] \dicDirectTranslationCS{váhající, zdráhavý}
\dicEntry[hiklaus] \dicTerm{hik··laus} \dicIPA{{h}{\textsci}{\textlengthmark}{\r{g}}{l}{\oe i}{s}} \dicPos{adj}[5]\dicFlx{}[-1] \dicDirectTranslationCS{neváhající, rozhodný} \dicExampleIS{hiklaus svar} \dicExampleCS{odpověď bez váhání}
\dicEntry[hikorð] \dicTerm{hik··orð} \dicIPA{{h}{\textsci}{\textlengthmark}{\r{g}}{\textopeno}{r}{\texttheta}} \dicPos{n}[2] \dicFlx{(‑s, ‑)}[5] \dicFieldCat{jaz.} \dicDirectTranslationCS{slovní výplň\,/\addthin vata, vycpávka}
\dicEntry[hiksta] \dicTerm{hikst|a} \dicIPA{{h}{\textsci}{x}{s}{\textsubring{d}}{a}} \dicPos{v}[1] \dicFlx{(‑aði)}[22] \dicDirectTranslationCS{škytat, mít škytavku}
\dicEntry[hiksti] \dicTerm{hikst|i} \dicIPA{{h}{\textsci}{x}{s}{\textsubring{d}}{\textsci}} \dicPos{m}[1] \dicFlx{(‑a, ‑ar)}[1] \dicDirectTranslationCS{škytání, škytnutí, škytavka};  \dicPhraseIS{vera með hiksta} \dicDirectTranslationCS{mít škytavku}
\dicEntry[hilla] \dicTerm{hill|a\smash{\textsuperscript{1}}} \dicsymFrequent\  \dicIPA{{h}{\textsci}{\textsubring{d}}{l}{a}} \dicPos{f}[1] \dicFlx{(‑u, ‑ur)}[19] \dicDirectTranslationCS{police, polička} \dicExampleIS{Á veggnum voru hillur.} \dicExampleCS{Na stěně byly poličky.};  \dicPhraseIS{leggja e‑ð á hilluna} \dicLangCat{přen.} \dicDirectTranslationCS{pověsit (co) na hřebík};  \dicPhraseIS{vera á réttri\,/\addthin rangri hillu} \dicLangCat{přen.} \dicDirectTranslationCS{být na správném\,/\addthin nesprávném místě} \dicIndirectTranslationCS{(mít práci, která (ne)odpovídá schopnostem člověka)}
\dicEntry[hilla] \dicTerm{hill|a\smash{\textsuperscript{2}}} \dicIPA{{h}{\textsci}{\textsubring{d}}{l}{a}} \dicPos{v}[2] \dicFlx{(‑ti, ‑t)}[83] \dicFlx{impers} \dicPhraseIS{það hillir undir e‑u} \dicDirectTranslationCS{(co) je v~dosahu, (co) je na obzoru, (co) se rýsuje};  \dicPhraseIS{e‑ð hillir uppi} \dicDirectTranslationCS{(co) se rýsuje (na obzoru ap.)}
\dicEntry[hilling] \dicTerm{hill··ing} \dicIPA{{h}{\textsci}{\textsubring{d}}{l}{i}{\ng}{\r{g}}} \dicPos{f}[4] \dicFlx{(‑ar, ‑ar)}[5] \dicDirectTranslationCS{vidina, přelud, fata morgána}
\dicEntry[Himalæjafjöll] \dicTerm{Himalæja··fjöll} \dicIPA{{h}{\textsci}{\textlengthmark}{m}{a}{l}{a}{i}{j}{a}{f}{j}{\oe}{\textsubring{d}}{\textsubring{l}}} \dicPos{n}[2] \dicFlx{pl}[9] \dicFieldCat{geog.} \dicDirectTranslationCS{Himálaj(e)}
\dicEntry[himbrimi] \dicTerm{him··brim|i} \dicIPA{{h}{\textsci}{m}{\textsubring{b}}{r}{\textsci}{m}{\textsci}} \dicPos{m}[1] \dicFlx{(‑a, ‑ar)}[1] \dicFieldCat{zool.} \dicDirectTranslationCS{potáplice lední} \textit{(l.~{\textLA{Gavia immer}})}  \dicsymPhoto\ 
\dicFigure{20887.jpg}{Himbrimi}{Himbrimi - NWF, Biolib, PD}
\dicEntry[himinlifandi] \dicTerm{himin··lif·andi} \dicIPA{{h}{\textsci}{\textlengthmark}{m}{\textsci}{n}{l}{\textsci}{v}{a}{n}{\textsubring{d}}{\textsci}} \dicPos{adj}[13] \dicFlx{indecl}[1] \dicDirectTranslationCS{rozradostněný, (jsoucí) radostí bez sebe, uchvácený, nadšený}
\dicEntry[himinn] \dicTerm{him|inn} \dicsymFrequent\  \dicIPA{{h}{\textsci}{\textlengthmark}{m}{\textsci}{\textsubring{n}}} \dicPos{m}[6] \dicFlx{(‑ins, ‑nar)}[39] \textbf{1.} \dicDirectTranslationCS{obloha, nebe} \dicExampleIS{stjörnur á himni} \dicExampleCS{hvězdy na obloze}  \textbf{2.} \dicSynonym{himnaríki} \dicDirectTranslationCS{nebe, ráj} \dicExampleIS{guð á himnum} \dicExampleCS{Bůh na nebi};  \dicPhraseIS{vera í sjöunda himni} \dicLangCat{přen.} \dicDirectTranslationCS{být v~sedmém nebi}
\dicEntry[himna] \dicTerm{himn|a} \dicIPA{{h}{\textsci}{m}{n}{a}} \dicPos{f}[1] \dicFlx{(‑u, ‑ur)}[7] \textbf{1.} \dicDirectTranslationCS{film, povlak, blána}  \textbf{2.} \dicFieldCat{biol.} \dicDirectTranslationCS{membrána}
\dicEntry[himnaríki] \dicTerm{himna··ríki} \dicIPA{{h}{\textsci}{m}{n}{a}{r}{i}{\r{\textObardotlessj}}{\textsci}} \dicPos{n}[2] \dicFlx{(‑s, ‑)}[16] \dicDirectTranslationCS{nebe, ráj} \dicExampleIS{himnaríki og helvíti} \dicExampleCS{nebe a~peklo}
\dicEntry[himnastigi] \dicTerm{himna··stig|i} \dicIPA{{h}{\textsci}{m}{n}{a}{s}{\textsubring{d}}{i}{j}{\textsci}} \dicPos{m}[1] \dicFlx{(‑a, ‑ar)}[1] \dicDirectTranslationCS{schody do nebe}
\dicEntry[himneskur] \dicTerm{himn··eskur} \dicIPA{{h}{\textsci}{m}{n}{\textepsilon}{s}{\r{g}}{\textscy}{\textsubring{r}}} \dicPos{adj}[1]\dicFlx{}[-1] \textbf{1.} \dicDirectTranslationCS{nebeský}  \textbf{2.} \dicLangCat{přen.} \dicDirectTranslationCS{božský} \dicExampleIS{Veðrið er himneskt.} \dicExampleCS{Počasí je božské.}
\dicEntry[hind] \dicTerm{hind} \dicIPA{{h}{\textsci}{n}{\textsubring{d}}} \dicPos{f}[7] \dicFlx{(‑ar, ‑ir)}[1] \dicDirectTranslationCS{laň}
\dicEntry[hindarblóm] \dicTerm{hindar··blóm} \dicIPA{{h}{\textsci}{n}{\textsubring{d}}{a}{r}{\textsubring{b}}{l}{ou}{\textsubring{m}}} \dicPos{n}[2] \dicFlx{(‑s, ‑)}[5] \dicLink{hortensía}
\dicEntry[hindber] \dicTerm{hind··ber} \dicIPA{{h}{\textsci}{n}{\textsubring{d}}{\textsubring{b}}{\textepsilon}{\textsubring{r}}} \dicPos{n}[2] \dicFlx{(‑s, ‑)}[13] \dicDirectTranslationCS{malina} \dicIndirectTranslationCS{(plod)}
\dicEntry[hindberjarunni] \dicTerm{hind·berja··runn|i} \dicIPA{{h}{\textsci}{n}{\textsubring{d}}{\textsubring{b}}{\textepsilon}{r}{j}{a}{r}{\textscy}{n}{\textsci}} \dicPos{m}[1] \dicFlx{(‑a, ‑ar)}[1] \dicFieldCat{bot.} \dicDirectTranslationCS{maliník, maliník obecný, ostružiník maliník} \textit{(l.~{\textLA{Rubus idaeus}})}  \dicsymPhoto\ 
\dicFigure{ds_image_hindberjarunni_0_1.jpg}{Hindberjarunni}{Hindberjarunni - Karelj, PD}
\dicEntry[hindí] \dicTerm{hindí} \dicIPA{{h}{\textsci}{n}{\textsubring{d}}{i}} \dicPos{n}[4] \dicFlx{indecl}[2] \dicDirectTranslationCS{hindština}
\dicEntry[hindra] \dicTerm{hindr|a} \dicsymFrequent\  \dicIPA{{h}{\textsci}{n}{\textsubring{d}}{r}{a}} \dicPos{v}[1] \dicFlx{(‑aði)}[1] \dicFlx{acc} \dicDirectTranslationCS{(za)bránit, zamezit, překážet} \dicExampleIS{hindra framkvæmdir} \dicExampleCS{zabránit provedení}
\dicEntry[hindrun] \dicTerm{hindr|un} \dicIPA{{h}{\textsci}{n}{\textsubring{d}}{r}{\textscy}{\textsubring{n}}} \dicPos{f}[7] \dicFlx{(‑unar, ‑anir)}[8] \textbf{1.} \dicDirectTranslationCS{překážka, bariéra, zábrana}  \textbf{2.} \dicDirectTranslationCS{zabránění, zamezení}  \textbf{3.} \dicFieldCat{sport.} \dicDirectTranslationCS{překážka}
\dicEntry[hindrunarhlaup] \dicTerm{hindrunar··hlaup} \dicIPA{{h}{\textsci}{n}{\textsubring{d}}{r}{\textscy}{n}{a}{r}{\textsubring{l}}{\oe i}{\textsubring{b}}} \dicPos{n}[2] \dicFlx{(‑s, ‑)}[5] \dicFieldCat{sport.} \dicDirectTranslationCS{překážkový běh, steeplechase} \dicIndirectTranslationCS{(překážkový závod v~lehké atletice)}
\dicEntry[hindurvitni] \dicTerm{hindur··vitni} \dicIPA{{h}{\textsci}{n}{\textsubring{d}}{\textscy}{r}{v}{\textsci}{h}{\textsubring{d}}{n}{\textsci}} \dicPos{n}[2] \dicFlx{(‑s, ‑)}[14] \dicSynonym{hjátrú} \dicDirectTranslationCS{pověra}
\dicEntry[hindúi] \dicTerm{hindú|i} \dicIPA{{h}{\textsci}{n}{\textsubring{d}}{u}{\textsci}} \dicPos{m}[1] \dicFlx{(‑a, ‑ar)}[1] \dicDirectTranslationCS{Hind(ka)}
\dicEntry[hindúismi] \dicTerm{hindú··ism|i} \dicIPA{{h}{\textsci}{n}{\textsubring{d}}{u}{\textsci}{s}{m}{\textsci}} \dicPos{m}[1] \dicFlx{(‑a)}[3] \dicDirectTranslationCS{hinduismus}
\dicEntry[hingað] \dicTerm{hingað} \dicsymFrequent\  \dicIPA{{h}{i}{\ng}{\r{g}}{a}{\texttheta}} \dicPos{adv} \dicDirectTranslationCS{sem} \dicIndirectTranslationCS{(o~pohybu na místo zde)} \dicExampleIS{Komdu hingað til mín!} \dicExampleCS{Pojď sem ke mně!};  \dicPhraseIS{hingað og þangað} \dicFlx{adv} \dicDirectTranslationCS{sem a~tam};  \dicPhraseIS{hingað til} \dicFlx{adv} \dicDirectTranslationCS{dosud, doposud, doposavad}
\dicEntry[hinkra] \dicTerm{hinkr|a} \dicIPA{{h}{i}{\r{\ng}}{\r{g}}{r}{a}} \dicPos{v}[1] \dicFlx{(‑aði)}[1] \dicPhraseIS{hinkra við} \dicSynonym{doka} \dicDirectTranslationCS{posečkat, počkat chvíli\,/\addthin chvilku} \dicExampleIS{Nennir þú að hinkra aðeins?} \dicExampleCS{Mohl bys chvilku počkat?}
\dicEntry[hinn] \dicTerm{hinn\smash{\textsuperscript{1}}} \dicsymFrequent\  \dicIPA{{h}{\textsci}{n}{\textlengthmark}} \dicPos{pron}[15] \dicFlx{dem (n hitt)} \textbf{1.} \dicDirectTranslationCS{ten, onen, tamten, ten druhý} \dicIndirectTranslationCS{(s~podstatným jménem v~určitém tvaru)}  \textbf{2.} \dicDirectTranslationCS{ostatní, druhý} \dicIndirectTranslationCS{(stojící samostatně)} \dicExampleIS{Viltu kalla á hina?} \dicExampleCS{Můžeš zavolat na ostatní?}  \textbf{3.} \dicDirectTranslationCS{první} \dicIndirectTranslationCS{(stojící samostatně ve srovnání dvou věcí ap.)}  \textbf{4.} \dicDirectTranslationCS{pozítří} \dicIndirectTranslationCS{(stojící samostatně jako příslovce)} \dicExampleIS{ekki á morgun heldur hinn} \dicExampleCS{ne zítra, ale pozítří};  \dicPhraseIS{hinir og þessir} \dicDirectTranslationCS{různí, různí jiní};  \dicPhraseIS{þessi -- hinn} \dicDirectTranslationCS{tento -- ten druhý} \dicExampleIS{Þessi bók er ný en hin er gömul.} \dicExampleCS{Tato knížka je nová, ale ta druhá je stará.};  \dicPhraseIS{annar -- hinn} \dicDirectTranslationCS{první -- druhý} \dicIndirectTranslationCS{(o~dvou věcech ap.)} \dicExampleIS{Önnur myndin er falleg en hin er ljót.} \dicExampleCS{První obrázek je hezký, ale druhý je šeredný.}
\dicEntry[hinn] \dicTerm{hinn\smash{\textsuperscript{2}}} \dicsymFrequent\  \dicIPA{{h}{\textsci}{n}{\textlengthmark}} \dicPos{part} \dicFlx{(n hið)} \dicFieldCat{jaz.} \dicIndirectTranslationCS{člen určitý (vystupuje před přídavným jménem nebo číslovkou, používaný ve formální řeči)} \dicExampleIS{hinn góði} \dicExampleCS{ten dobrý}
\dicEntry[hins vegar] \dicTerm{hins vegar}\dicTerm{, hinsvegar} \dicIPA{{h}\-{\textsci}\-{n}\-{s}\-{v}\-{\textepsilon}\-{\textbabygamma}\-{a}\-{\textsubring{r}}\-} \dicPos{adv} \dicSynonym*{á hinn bóginn} \dicDirectTranslationCS{na druhou stranu, naproti tomu};  \dicPhraseIS{annars vegar -- hins vegar} \dicDirectTranslationCS{na jednu stranu -- na druhou stranu}
\dicEntry[hinstur] \dicTerm{hinstur} \dicIPA{{h}{\textsci}{n}{s}{\textsubring{d}}{\textscy}{\textsubring{r}}} \dicPos{adj}[12]\dicFlx{}[-12] \dicFlx{m sg nom sup} \dicSynonym{síðastur} \dicDirectTranslationCS{poslední} \dicExampleIS{til hinstu stundar} \dicExampleCS{do poslední chvíle}
\dicEntry[hinsvegar] \dicTerm{hins··vegar} \dicIPA{{h}{\textsci}{n}{s}{v}{\textepsilon}{\textbabygamma}{a}{\textsubring{r}}} \dicPos{adv} \dicLink{hins vegar}
\dicEntry[hinumegin] \dicTerm{hinu··megin} \dicIPA{{h}{\textsci}{\textlengthmark}{n}{\textscy}{m}{ei}{\textsci}{\textsubring{n}}} \dicPos{adv} \dicLink{megin}
\dicEntry[hippi] \dicTerm{hipp|i} \dicIPA{{h}{\textsci}{h}{\textsubring{b}}{\textsci}} \dicPos{m}[1] \dicFlx{(‑a, ‑ar)}[1] \dicDirectTranslationCS{hipík, hippie}
\dicEntry[hirð] \dicTerm{hirð} \dicIPA{{h}{\textsci}{r}{\texttheta}} \dicPos{f}[7] \dicFlx{(‑ar, ‑ir)}[1] \dicDirectTranslationCS{dvůr, dvořanstvo} \dicIndirectTranslationCS{(panovník s~dvořany)} \dicExampleIS{við hirðina} \dicExampleCS{u~dvora}
\dicEntry[hirða] \dicTerm{hirð|a\smash{\textsuperscript{1}}} \dicIPA{{h}{\textsci}{r}{ð}{a}} \dicPos{f}[1] \dicFlx{(‑u)}[5] \dicSynonym{hirðusemi} \dicDirectTranslationCS{péče, pečování, dbání, starání se}
\dicEntry[hirða] \dicTerm{hir|ða\smash{\textsuperscript{2}}} \dicsymFrequent\  \dicIPA{{h}{\textsci}{r}{ð}{a}} \dicPos{v}[2] \dicFlx{(‑ti, ‑t)}[47] \dicFlx{acc} \textbf{1.} \dicSynonym{geyma} \dicDirectTranslationCS{uschovat, uložit} \dicExampleIS{hirða umbúðirnar} \dicExampleCS{uschovat balicí materiál}  \textbf{2.} \dicDirectTranslationCS{pečovat, dbát} \dicExampleIS{hirða líkama sinn} \dicExampleCS{dbát o~své tělo}  \textbf{3.} \dicDirectTranslationCS{vybrat, sebrat} \dicExampleIS{hirða ruslið} \dicExampleCS{sebrat odpad}  \textbf{4.} \dicSynonym{gæta} \dicDirectTranslationCS{pečovat, starat se (o~ovce ap.)}  \textbf{5.} \dicPhraseIS{hirða um e‑ð} \dicSynonym*{annast} \dicDirectTranslationCS{starat se o~(co), věnovat (čemu) pozornost}
\dicEntry[hirðing] \dicTerm{hirð··ing} \dicIPA{{h}{\textsci}{r}{ð}{i}{\ng}{\r{g}}} \dicPos{f}[4] \dicFlx{(‑ar, ‑ar)}[5] \textbf{1.} \dicSynonym{umönnun} \dicDirectTranslationCS{péče, pečování}  \textbf{2.} \dicDirectTranslationCS{sběr, kolekce}
\dicEntry[hirðingi] \dicTerm{hirð··ing|i} \dicIPA{{h}{\textsci}{r}{ð}{i}{\textltailn}{\r{\textObardotlessj}}{\textsci}} \dicPos{m}[1] \dicFlx{(‑ja, ‑jar)}[14] \dicDirectTranslationCS{kočovník, kočovnice, nomád(ka)}
\dicEntry[hirðingjalíf] \dicTerm{hirð·ingja··líf} \dicIPA{{h}{\textsci}{r}{ð}{i}{\textltailn}{\r{\textObardotlessj}}{a}{l}{i}{f}} \dicPos{n}[2] \dicFlx{(‑s)}[2] \dicDirectTranslationCS{kočovný život}
\dicEntry[hirðir] \dicTerm{hirð|ir} \dicIPA{{h}{\textsci}{r}{ð}{\textsci}{\textsubring{r}}} \dicPos{m}[7] \dicFlx{(‑is, ‑ar)}[1] \textbf{1.} \dicSynonym{hjarðmaður} \dicDirectTranslationCS{pastýř(ka), pastevec, pastevkyně}  \textbf{2.} \dicSynonym{prestur} \dicDirectTranslationCS{pastýř (duchovní ap.)}
\dicEntry[hirðmaður] \dicTerm{hirð··|maður} \dicIPA{{h}{\textsci}{r}{ð}{m}{a}{ð}{\textscy}{\textsubring{r}}} \dicPos{m}[13] \dicFlx{(‑manns, ‑menn)}[2] \dicDirectTranslationCS{dvořan(ka)}
\dicEntry[hirðskáld] \dicTerm{hirð··skáld} \dicIPA{{h}{\textsci}{r}{ð}{s}{\r{g}}{au}{l}{\textsubring{d}}} \dicPos{n}[2] \dicFlx{(‑s, ‑)}[5] \dicDirectTranslationCS{dvorní básník\,/\addthin básnířka}
\dicEntry[hirðulaus] \dicTerm{hirðu··laus} \dicIPA{{h}{\textsci}{r}{ð}{\textscy}{l}{\oe i}{s}} \dicPos{adj}[5]\dicFlx{}[-1] \dicDirectTranslationCS{ledabylý, nedbalý, lehkovážný}
\dicEntry[hirðuleysi] \dicTerm{hirðu··leysi} \dicIPA{{h}{\textsci}{r}{ð}{\textscy}{l}{ei}{s}{\textsci}} \dicPos{n}[2] \dicFlx{(‑s)}[20] \dicDirectTranslationCS{nedbalost, ledabylost}
\dicEntry[hirðusamur] \dicTerm{hirðu··|samur} \dicIPA{{h}{\textsci}{r}{ð}{\textscy}{s}{a}{m}{\textscy}{\textsubring{r}}} \dicPos{adj}[1] \dicFlx{(f ‑söm)}[2] \dicDirectTranslationCS{spořádaný, svědomitý}
\dicEntry[hirðusemi] \dicTerm{hirðu··sem|i} \dicIPA{{h}{\textsci}{r}{ð}{\textscy}{s}{\textepsilon}{m}{\textsci}} \dicPos{f}[2] \dicFlx{(‑i)}[2] \dicDirectTranslationCS{spořádanost, svědomitost}
\dicEntry[hirsi] \dicTerm{hirsi} \dicIPA{{h}{\textsci}{\textsubring{r}}{s}{\textsci}} \dicPos{n}[2] \dicFlx{(‑s)}[20] \dicFieldCat{bot.} \dicDirectTranslationCS{proso seté} \textit{(l.~{\textLA{Panicum miliaceum}})}  \dicsymPhoto\ 
\dicFigure{ds_image_hirsi_0_2.jpg}{Hirsi}{Hirsi - Topjabot, GFDL}
\dicEntry[hirsla] \dicTerm{hirsl|a} \dicIPA{{h}{\textsci}{\textsubring{r}}{s}{\textsubring{d}}{l}{a}} \dicPos{f}[1] \dicFlx{(‑u, ‑ur)}[13] \dicDirectTranslationCS{skříň, skříňka} \dicExampleIS{hirsla undir peninga} \dicExampleCS{pokladnička}
\dicEntry[hirta] \dicTerm{hirt|a} \dicIPA{{h}{\textsci}{\textsubring{r}}{\textsubring{d}}{a}} \dicPos{v}[2] \dicFlx{(‑i, ‑)}[12] \dicFlx{acc} \dicSynonym{refsa} \dicDirectTranslationCS{(po)trestat, ztrestat};  \dicPhraseIS{hirta e‑n fyrir e‑ð} \dicDirectTranslationCS{potrestat (koho) za (co)}
\dicEntry[hirting] \dicTerm{hirt··ing} \dicIPA{{h}{\textsci}{\textsubring{r}}{\textsubring{d}}{i}{\ng}{\r{g}}} \dicPos{f}[4] \dicFlx{(‑ar, ‑ar)}[5] \dicSynonym{refsing} \dicDirectTranslationCS{(po)trestání, trest}
\dicEntry[hirtir] \dicTerm{hirtir} \dicIPA{{h}{\textsci}{\textsubring{r}}{\textsubring{d}}{\textsci}{\textsubring{r}}} \dicPos{m} \dicFlx{pl nom} \dicLink{hjörtur}
\dicEntry[hismi] \dicTerm{hismi} \dicIPA{{h}{\textsci}{s}{m}{\textsci}} \dicPos{n}[2] \dicFlx{(‑s)}[20] \dicDirectTranslationCS{pleva};  \dicPhraseIS{skilja hismið frá kjarnanum} \dicLangCat{přen.} \dicDirectTranslationCS{oddělit zrno od plev}
\dicEntry[hispurslaus] \dicTerm{hispurs··laus} \dicIPA{{h}{\textsci}{s}{\textsubring{b}}{\textscy}{\textsubring{r}}{s}{l}{\oe i}{s}} \dicPos{adj}[5]\dicFlx{}[-1] \dicSynonym{óþvingaður} \dicDirectTranslationCS{přímý, otevřený (v~projevu ap.), nevybíravý}
\dicEntry[hispursleysi] \dicTerm{hispurs··leysi} \dicIPA{{h}{\textsci}{s}{\textsubring{b}}{\textscy}{\textsubring{r}}{s}{l}{ei}{s}{\textsci}} \dicPos{n}[2] \dicFlx{(‑s)}[20] \dicDirectTranslationCS{přímost, otevřenost, nevybíravost}
\dicEntry[hissa] \dicTerm{hissa} \dicsymFrequent\  \dicIPA{{h}{\textsci}{s}{\textlengthmark}{a}} \dicPos{adj}[13] \dicFlx{indecl}[2] \dicSynonym{undrandi} \dicDirectTranslationCS{překvapený, udivený, užaslý} \dicExampleIS{Ég er ekki hissa á því.} \dicExampleCS{Nepřekvapuje mě to.};  \dicPhraseIS{verða hissa á e‑u} \dicDirectTranslationCS{být překvapený (čím), (po)divit se (čemu)}
\dicEntry[hita] \dicTerm{hit|a\smash{\textsuperscript{1}}} \dicIPA{{h}{\textsci}{\textlengthmark}{\textsubring{d}}{a}} \dicPos{f}[1] \dicFlx{(‑u, ‑ur)}[7] \dicLangCat{zast.} \dicDirectTranslationCS{ohřání, ohřátí, zahřátí (kafe ap.)} \dicExampleIS{ein hita af kaffi} \dicExampleCS{jedno ohřání kafe};  \dicPhraseIS{vera einn um hituna} \dicLangCat{přen.} \dicDirectTranslationCS{být na to sám (na práci ap.)}
\dicEntry[hita] \dicTerm{hit|a\smash{\textsuperscript{2}}} \dicsymFrequent\  \dicIPA{{h}{\textsci}{\textlengthmark}{\textsubring{d}}{a}} \dicPos{v}[1] \dicFlx{(‑aði)}[1] \dicFlx{acc\,/\addthin dat} \textbf{1.} \dicFlx{acc} \dicDirectTranslationCS{zahřát, ohřát, rozehřát} \dicIndirectTranslationCS{(učinit teplým)} \dicExampleIS{Hitið ofninn í 150 gráður.} \dicExampleCS{Zahřejte troubu na 150 stupňů.};  \dicPhraseIS{hita kaffi\,/\addthin te} \dicDirectTranslationCS{ohřát kávu\,/\addthin čaj}  \textbf{2.} \dicFlx{dat} \dicDirectTranslationCS{ohřát, zahřát} \dicIndirectTranslationCS{(zbavit pocitu chladu)} \dicExampleIS{hita e‑m á höndunum} \dicExampleCS{zahřát (komu) ruce};  \dicIdiom{hita}[upp]{ \dicPhraseIS{hita upp e‑ð}} {\textbf{a.}} \dicDirectTranslationCS{ohřívat (co), vyhřívat (co)} \dicExampleIS{hita upp hús} \dicExampleCS{vyhřívat dům};  {\textbf{b.}} \dicDirectTranslationCS{ohřát (co), zahřát (co)} \dicExampleIS{hita upp matinn} \dicExampleCS{ohřát jídlo};  \dicPhraseIS{hita upp} \dicDirectTranslationCS{rozehřát (tělo ap.)} \dicExampleIS{hita sig upp} \dicExampleCS{rozehřát se}
\dicEntry[hitabelti] \dicTerm{hita··belti} \dicIPA{{h}{\textsci}{\textlengthmark}{\textsubring{d}}{a}{\textsubring{b}}{\textepsilon}{\textsubring{l}}{\textsubring{d}}{\textsci}} \dicPos{n}[2] \dicFlx{(‑s, ‑)}[14] \dicFieldCat{meteo.} \dicDirectTranslationCS{tropický pás}
\dicEntry[hitabrúsi] \dicTerm{hita··brús|i} \dicIPA{{h}{\textsci}{\textlengthmark}{\textsubring{d}}{a}{\textsubring{b}}{r}{u}{s}{\textsci}} \dicPos{m}[1] \dicFlx{(‑a, ‑ar)}[1] \dicDirectTranslationCS{termoska}
\dicEntry[hitabylgja] \dicTerm{hita··bylgj|a} \dicIPA{{h}{\textsci}{\textlengthmark}{\textsubring{d}}{a}{\textsubring{b}}{\textsci}{l}{\r{\textObardotlessj}}{a}} \dicPos{f}[1] \dicFlx{(‑u, ‑ur)}[25] \dicFieldCat{meteo.} \dicDirectTranslationCS{horká vlna, vlna horka}
\dicEntry[hitaeining] \dicTerm{hita··ein·ing} \dicIPA{{h}{\textsci}{\textlengthmark}{\textsubring{d}}{a}{ei}{n}{i}{\ng}{\r{g}}} \dicPos{f}[4] \dicFlx{(‑ar, ‑ar)}[5] \dicSynonym{kaloría} \dicDirectTranslationCS{kalorie}
\dicEntry[hitakanna] \dicTerm{hita··|kanna} \dicIPA{{h}{\textsci}{\textlengthmark}{\textsubring{d}}{a}{k\smash{\textsuperscript{h}}}{a}{n}{a}} \dicPos{f}[1] \dicFlx{(‑könnu, ‑könnur)}[8] \dicDirectTranslationCS{termoska}
\dicEntry[hitamælir] \dicTerm{hita··mæl|ir} \dicIPA{{h}{\textsci}{\textlengthmark}{\textsubring{d}}{a}{m}{a}{i}{l}{\textsci}{\textsubring{r}}} \dicPos{m}[7] \dicFlx{(‑is, ‑ar)}[1] \dicDirectTranslationCS{teploměr}
\dicEntry[hitapoki] \dicTerm{hita··pok|i} \dicIPA{{h}{\textsci}{\textlengthmark}{\textsubring{d}}{a}{p\smash{\textsuperscript{h}}}{\textopeno}{\r{\textObardotlessj}}{\textsci}} \dicPos{m}[1] \dicFlx{(‑a, ‑ar)}[1] \dicDirectTranslationCS{termofor} \dicIndirectTranslationCS{(pryžový vak naplněný horkou vodou)}
\dicEntry[hitari] \dicTerm{hit··ar|i} \dicIPA{{h}{\textsci}{\textlengthmark}{\textsubring{d}}{a}{r}{\textsci}} \dicPos{m}[1] \dicFlx{(‑a, ‑ar)}[13] \dicSynonym*{hitatæki} \dicDirectTranslationCS{ohřívač}
\dicEntry[hitaskil] \dicTerm{hita··skil} \dicIPA{{h}{\textsci}{\textlengthmark}{\textsubring{d}}{a}{s}{\r{\textObardotlessj}}{\textsci}{\textsubring{l}}} \dicPos{n}[2] \dicFlx{pl}[1] \dicFieldCat{meteo.} \dicDirectTranslationCS{teplotní rozhraní}
\dicEntry[hitasótt] \dicTerm{hita··sótt} \dicIPA{{h}{\textsci}{\textlengthmark}{\textsubring{d}}{a}{s}{ou}{h}{\textsubring{d}}} \dicPos{f}[7] \dicFlx{(‑ar, ‑ir)}[1] \dicDirectTranslationCS{horečka, teplota}
\dicEntry[hitastig] \dicTerm{hita··stig} \dicIPA{{h}{\textsci}{\textlengthmark}{\textsubring{d}}{a}{s}{\textsubring{d}}{\textsci}{x}} \dicPos{n}[2] \dicFlx{(‑s, ‑)}[5] \dicDirectTranslationCS{teplota, teplotní stupeň}
\dicEntry[hitaveita] \dicTerm{hita··veit|a} \dicIPA{{h}{\textsci}{\textlengthmark}{\textsubring{d}}{a}{v}{ei}{\textsubring{d}}{a}} \dicPos{f}[1] \dicFlx{(‑u, ‑ur)}[13] \dicDirectTranslationCS{teplárna}
\dicEntry[hiti] \dicTerm{hit|i} \dicsymFrequent\  \dicIPA{{h}{\textsci}{\textlengthmark}{\textsubring{d}}{\textsci}} \dicPos{m}[1] \dicFlx{(‑a, ‑ar)}[1] \textbf{1.} \dicSynonym{ylur} \dicDirectTranslationCS{teplo, horko}  \textbf{2.} \dicSynonym{sótthiti} \dicDirectTranslationCS{teplota, mírná horečka} \dicExampleIS{vera með hita} \dicExampleCS{mít teplotu}  \textbf{3.} \dicSynonym{ákafi} \dicDirectTranslationCS{zápal, vášnivost};  \dicPhraseIS{það er hiti í e‑u} \dicLangCat{přen.} \dicDirectTranslationCS{(co) je vášnivé, (co) vyvolává vášně (volební kampaň ap.)}
\dicEntry[hitna] \dicTerm{hitn|a} \dicsymFrequent\  \dicIPA{{h}{\textsci}{h}{\textsubring{d}}{n}{a}} \dicPos{v}[1] \dicFlx{(‑aði)}[74] \dicDirectTranslationCS{zahřát se, ohřát se} \dicExampleIS{Ofninn hitnar fljótlega.} \dicExampleCS{Trouba se rychle rozehřeje.};  \dicPhraseIS{e‑m hitnar} \dicFlx{impers} \dicDirectTranslationCS{(kdo) se zahřívá\,/\addthin rozehřívá, (komu) se dělá teplo};  \dicPhraseIS{það hitnar} \dicFlx{impers} \dicDirectTranslationCS{dělá se teplo, otepluje se}
\dicEntry[hitta] \dicTerm{hitt|a} \dicsymFrequent\  \dicIPA{{h}{\textsci}{h}{\textsubring{d}}{a}} \dicPos{v}[2] \dicFlx{(‑i, ‑)}[1] \dicFlx{acc} \textbf{1.} \dicSynonym{mæta} \dicDirectTranslationCS{potkat, potkávat} \dicExampleIS{hitta e‑n af tilviljun} \dicExampleCS{potkat (koho) náhodou}  \textbf{2.} \dicSynonym{hæfa\smash{\textsuperscript{2}}} \dicDirectTranslationCS{trefit (se)} \dicExampleIS{hitta í markið} \dicExampleCS{trefit se do brány};  \dicIdiom{hitta}[á]{ \dicPhraseIS{hitta á e‑ð}} \dicDirectTranslationCS{trefit (co) (správné číslo ap.)}; { \dicPhraseIS{hitta á e‑n}} \dicDirectTranslationCS{(na)trefit na (koho)};  \dicIdiom{hittast}{ \dicPhraseIS{hittast}} \dicFlx{refl} \dicDirectTranslationCS{setkat se, sejít se, potkat se, potkávat se, scházet se} \dicExampleIS{Þau hittast oft í viku.} \dicExampleCS{Během týdne se často scházejí.};  \dicIdiom{hittast}[á]{ \dicPhraseIS{það hittist (þannig) á}} \dicFlx{refl} \dicDirectTranslationCS{tak se to sešlo\,/\addthin seběhlo}
\dicEntry[hittinn] \dicTerm{hittinn} \dicIPA{{h}{\textsci}{h}{\textsubring{d}}{\textsci}{\textsubring{n}}} \dicPos{adj}[6]\dicFlx{}[-2] \textbf{1.} \dicSynonym*{markvís} \dicDirectTranslationCS{přesný} \dicIndirectTranslationCS{(o~střelci)}  \textbf{2.} \dicSynonym{fundvís} \dicDirectTranslationCS{trefný (odpověď ap.)} \dicExampleIS{hittið svar} \dicExampleCS{trefná odpověď}
\dicEntry[hitun] \dicTerm{hit|un} \dicIPA{{h}{\textsci}{\textlengthmark}{\textsubring{d}}{\textscy}{\textsubring{n}}} \dicPos{f}[7] \dicFlx{(‑unar)}[9] \dicDirectTranslationCS{vytápění, vyhřívání} \dicExampleIS{hitun húsa} \dicExampleCS{vytápění domů}
\dicEntry[hí] \dicTerm{hí} \dicIPA{{h}{i}{\textlengthmark}} \dicPos{inter} \dicPhraseIS{hí á þig} \dicDirectTranslationCS{kiš kiš!} \dicIndirectTranslationCS{(vyjadřuje mírný výsměch)}
\dicEntry[hía] \dicTerm{hí|a} \dicIPA{{h}{i}{\textlengthmark}{j}{a}} \dicPos{v}[1] \dicFlx{(‑aði)}[1] \dicPhraseIS{hía á e‑n} \dicDirectTranslationCS{ukazovat si prstem na (koho), posmívat se (komu)}
\dicEntry[híbýli] \dicTerm{hí··býli}\dicTerm{, hýbýli} \dicIPA{{h}\-{i}\-{\textlengthmark}\-{\textsubring{b}}\-{i}\-{l}\-{\textsci}\-} \dicPos{n}[2] \dicFlx{pl}[19] \dicSynonym{hús} \dicDirectTranslationCS{obydlí, příbytek}
\dicEntry[híði] \dicTerm{híði} \dicIPA{{h}{i}{\textlengthmark}{ð}{\textsci}} \dicPos{n}[2] \dicFlx{(‑s, ‑)}[14] \dicDirectTranslationCS{doupě, brloh, pelech};  \dicPhraseIS{leggjast í híði} \dicFlx{refl} \dicDirectTranslationCS{přezimovat, zalehnout na zimu}
\dicEntry[híena] \dicTerm{híen|a} \dicIPA{{h}{i}{j}{\textepsilon}{n}{a}} \dicPos{f}[1] \dicFlx{(‑u, ‑ur)}[7] \dicFieldCat{zool.} \dicDirectTranslationCS{hyena} \textit{(l.~{\textLA{Hyaena}})}  \dicsymPhoto\ 
\dicFigure{ds_image_hiena_0_2.jpg}{Híena}{Híena - Emmanuel FAIVRE, PD}
\dicEntry[hífa] \dicTerm{híf|a} \dicIPA{{h}{i}{\textlengthmark}{v}{a}} \dicPos{v}[2] \dicFlx{(‑ði, ‑t)}[99] \dicFlx{acc} \dicSynonym{lyfta\smash{\textsuperscript{2}}} \dicDirectTranslationCS{zvednout, (vy)zdvihnout} \dicExampleIS{hífa e‑ð upp} \dicExampleCS{vyzdvihnout (co)}
\dicEntry[hífaður] \dicTerm{híf··|aður} \dicIPA{{h}{i}{\textlengthmark}{v}{a}{ð}{\textscy}{\textsubring{r}}} \dicPos{adj}[3] \dicFlx{(f ‑uð)}[3] \dicSynonym{drukkinn} \dicDirectTranslationCS{přiopilý, podnapilý}
\dicEntry[híma] \dicTerm{hím|a\smash{\textsuperscript{1}}} \dicIPA{{h}{i}{\textlengthmark}{m}{a}} \dicPos{f}[1] \dicFlx{(‑u, ‑ur)}[19] \textbf{1.} \dicSynonym{þerrir} \dicDirectTranslationCS{provázek, šňůrka}  \textbf{2.} \dicSynonym{blundur} \dicDirectTranslationCS{zdřímnutí}
\dicEntry[híma] \dicTerm{hím|a\smash{\textsuperscript{2}}} \dicIPA{{h}{i}{\textlengthmark}{m}{a}} \dicPos{v}[2] \dicFlx{(‑di, ‑t)}[160] \dicDirectTranslationCS{krčit se (před větrem ap.), choulit se}
\dicEntry[hírast] \dicTerm{hír|ast} \dicIPA{{h}{i}{\textlengthmark}{r}{a}{s}{\textsubring{d}}} \dicPos{v}[2] \dicFlx{(‑ðist, ‑st)}[202] \dicFlx{refl} \dicSynonym*{búa þröngt} \dicDirectTranslationCS{mačkat se, tísnit se}
\dicEntry[hjaðna] \dicTerm{hjaðn|a} \dicIPA{{\c{c}}{a}{ð}{n}{a}} \dicPos{v}[1] \dicFlx{(‑aði)}[13] \dicSynonym*{eyðast} \dicDirectTranslationCS{polevit, opadnout (protesty ap.)}
\dicEntry[hjakka] \dicTerm{hjakk|a} \dicIPA{{\c{c}}{a}{h}{\r{g}}{a}} \dicPos{v}[1] \dicFlx{(‑aði)}[13] \dicFlx{acc} \textbf{1.} \dicSynonym*{smáhöggva} \dicDirectTranslationCS{rozkrájet\,/\addthin rozřezat na malé kousky}  \textbf{2.} \dicDirectTranslationCS{opakovat se} \dicIndirectTranslationCS{(znovu totéž říkat, psát ap.)};  \dicPhraseIS{hjakka í sama farinu} \dicLangCat{přen.} \dicDirectTranslationCS{nehnout se z~místa}
\dicEntry[hjala] \dicTerm{hjal|a} \dicIPA{{\c{c}}{a}{\textlengthmark}{l}{a}} \dicPos{v}[1] \dicFlx{(‑aði)}[50] \textbf{1.} \dicSynonym*{ambra} \dicDirectTranslationCS{žvatlat (dítě ap.)}  \textbf{2.} \dicSynonym{skrafa} \dicDirectTranslationCS{klábosit, povídat si} \dicExampleIS{hjala um bækur og bíómyndir} \dicExampleCS{klábosit o~knížkách a~filmech}
\dicEntry[hjalli] \dicTerm{hjall|i} \dicIPA{{\c{c}}{a}{\textsubring{d}}{l}{\textsci}} \dicPos{m}[1] \dicFlx{(‑a, ‑ar)}[8] \dicSynonym*{bringur} \dicDirectTranslationCS{skalní římsa, terasa (ve svahu ap.)};  \dicPhraseIS{klífa örðugasta hjallann} \dicLangCat{přen.} \dicDirectTranslationCS{překonat nejtěžší překážku};  \dicPhraseIS{það er glatt á hjalla} \dicLangCat{přen.} \dicDirectTranslationCS{nálada je super (na večírku ap.)}
\dicEntry[hjallur] \dicTerm{hjall|ur} \dicIPA{{\c{c}}{a}{\textsubring{d}}{l}{\textscy}{\textsubring{r}}} \dicPos{m}[6] \dicFlx{(‑s, ‑ar)}[5] \textbf{1.} \dicSynonym*{rimlaskúr} \dicDirectTranslationCS{dřevěný přístřešek sloužící k~sušení (ryb ap.)} \dicExampleIS{fiskhjallur} \dicExampleCS{přístřešek na sušení ryb}  \textbf{2.} \dicSynonym*{lélegt hús} \dicDirectTranslationCS{kolna, bouda, chatrč}
\dicEntry[Hjaltland] \dicTerm{Hjalt··land} \dicIPA{{\c{c}}{a}{\textsubring{l}}{\textsubring{d}}{l}{a}{n}{\textsubring{d}}} \dicPos{n}[2] \dicFlx{(‑s)}[4] \dicFieldCat{geog.} \dicDirectTranslationCS{Shetlandy} \dicIndirectTranslationCS{(skotské souostroví)}
\dicEntry[hjara] \dicTerm{hjara\smash{\textsuperscript{1}}} \dicIPA{{\c{c}}{a}{\textlengthmark}{r}{a}} \dicPos{f}[1] \dicFlx{(hjöru, hjörur)}[8] \dicLangCat{zast.} \dicSynonym{hjarir} \dicDirectTranslationCS{pant} \dicIndirectTranslationCS{(dveřní\,/\addthin okenní závěs)}
\dicEntry[hjara] \dicTerm{hjar|a\smash{\textsuperscript{2}}} \dicIPA{{\c{c}}{a}{\textlengthmark}{r}{a}} \dicPos{v}[2] \dicFlx{(‑ði, ‑að)}[125] \dicSynonym{tóra} \dicDirectTranslationCS{přežívat, živořit, vegetovat} \dicExampleIS{hjara til vors} \dicExampleCS{živořit do jara}
\dicEntry[hjarðar] \dicTerm{hjarðar} \dicIPA{{\c{c}}{a}{r}{ð}{a}{\textsubring{r}}} \dicPos{f} \dicFlx{sg gen} \dicLink{hjörð}
\dicEntry[hjarðir] \dicTerm{hjarðir} \dicIPA{{\c{c}}{a}{r}{ð}{\textsci}{\textsubring{r}}} \dicPos{f} \dicFlx{pl nom} \dicLink{hjörð}
\dicEntry[hjarðljóð] \dicTerm{hjarð··ljóð} \dicIPA{{\c{c}}{a}{r}{ð}{l}{j}{ou}{\texttheta}} \dicPos{n}[2] \dicFlx{(‑s, ‑)}[5] \dicFieldCat{lit.} \dicDirectTranslationCS{pastýřská báseň, selanka}
\dicEntry[hjarðmaður] \dicTerm{hjarð··|maður} \dicIPA{{\c{c}}{a}{r}{ð}{m}{a}{ð}{\textscy}{\textsubring{r}}} \dicPos{m}[13] \dicFlx{(‑manns, ‑menn)}[2] \dicSynonym{hirðir} \dicDirectTranslationCS{pastýř(ka), pastevec, pastevkyně}
\dicEntry[hjarðskáldskapur] \dicTerm{hjarð··skáld·skap|ur} \dicIPA{{\c{c}}{a}{r}{ð}{s}{\r{g}}{au}{l}{\textsubring{d}}{s}{\r{g}}{a}{\textsubring{b}}{\textscy}{\textsubring{r}}} \dicPos{m}[10] \dicFlx{(‑ar)}[15] \dicFieldCat{lit.} \dicDirectTranslationCS{pastýřská poezie, bukolika}
\dicEntry[hjari] \dicTerm{hjar|i} \dicIPA{{\c{c}}{a}{\textlengthmark}{r}{\textsci}} \dicPos{m}[1] \dicFlx{(‑a, ‑ar)}[8] \dicSynonym{jaðar} \dicDirectTranslationCS{(o)kraj, hranice};  \dicPhraseIS{á hjara veraldar} \dicFlx{adv} \dicLangCat{přen.} \dicDirectTranslationCS{na kraji světa}
\dicEntry[hjarir] \dicTerm{hjarir} \dicIPA{{\c{c}}{a}{\textlengthmark}{r}{\textsci}{\textsubring{r}}} \dicPos{f}[12] \dicFlx{pl}[9] \dicSynonym*{lamir} \dicDirectTranslationCS{panty} \dicExampleIS{hjarir á hurðinni} \dicExampleCS{panty u~dveří}
\dicEntry[hjarn] \dicTerm{hjarn} \dicIPA{{\c{c}}{a}{r}{\textsubring{d}}{\textsubring{n}}} \dicPos{n}[2] \dicFlx{(‑s, hjörn)}[8] \dicDirectTranslationCS{firn, krusta} \dicIndirectTranslationCS{(hrubozrnná ledová hmota)}
\dicEntry[hjarna] \dicTerm{hjarn|a} \dicIPA{{\c{c}}{a}{r}{\textsubring{d}}{n}{a}} \dicPos{v}[1] \dicFlx{(‑aði)}[43] \dicPhraseIS{hjarna við} \dicSynonym*{rétta við} \dicDirectTranslationCS{zotavit se, uzdravit se}
\dicEntry[hjarta] \dicTerm{hjarta} \dicsymFrequent\  \dicIPA{{\c{c}}{a}{\textsubring{r}}{\textsubring{d}}{a}} \dicPos{n}[1] \dicFlx{(‑, hjörtu)}[6] \textbf{1.} \dicFieldCat{anat.} \dicDirectTranslationCS{srdce};  \dicPhraseIS{fá fyrir hjartað} \dicDirectTranslationCS{mít bolesti srdce}  \textbf{2.} \dicDirectTranslationCS{srdce} \dicIndirectTranslationCS{(citová oblast lidské psychiky)};  \dicPhraseIS{af öllu hjarta} \dicFlx{adv} \dicDirectTranslationCS{z~celého srdce, celým srdcem};  \dicPhraseIS{e‑m liggur e‑að á hjarta} \dicFlx{impers} \dicLangCat{přen.} \dicDirectTranslationCS{(komu co) leží na srdci};  \dicPhraseIS{hafa ekki hjarta í sér til að (gera e‑ð)} \dicLangCat{přen.} \dicDirectTranslationCS{nemít to srdce ((co) udělat)};  \dicPhraseIS{hafa gott hjarta} \dicLangCat{přen.} \dicDirectTranslationCS{mít dobré srdce};  \dicPhraseIS{hafa hjartað á réttum stað} \dicLangCat{přen.} \dicDirectTranslationCS{mít srdce na pravém místě};  \dicPhraseIS{hann er með hjartað í buxunum} \dicLangCat{přen.} \dicDirectTranslationCS{má srdce v~kalhotách};  \dicPhraseIS{hjartað mitt} \dicDirectTranslationCS{moje nejdražší, moje milovaná};  \dicPhraseIS{hugsa með hjartanu} \dicLangCat{přen.} \dicDirectTranslationCS{myslet srdcem}  \textbf{3.} \dicDirectTranslationCS{srdce} \dicIndirectTranslationCS{(barva v~kartách)}  \textbf{4.} \dicDirectTranslationCS{srdce, střed (země ap.)} \dicExampleIS{hjarta landsins} \dicExampleCS{srdce země}
\dicEntry[hjartaáfall] \dicTerm{hjarta··á·|fall} \dicIPA{{\c{c}}{a}{\textsubring{r}}{\textsubring{d}}{a}{au}{f}{a}{\textsubring{d}}{\textsubring{l}}} \dicPos{n}[2] \dicFlx{(‑falls, ‑föll)}[8] \dicFieldCat{med.} \dicDirectTranslationCS{infarkt}
\dicEntry[hjartagóður] \dicTerm{hjarta··|góður} \dicIPA{{\c{c}}{a}{\textsubring{r}}{\textsubring{d}}{a}{\r{g}}{ou}{ð}{\textscy}{\textsubring{r}}} \dicPos{adj}[11] \dicFlx{(comp ‑betri, sup ‑bestur)}[1] \dicSynonym*{góðviljaður} \dicDirectTranslationCS{dobrosrdečný}
\dicEntry[hjartaknosari] \dicTerm{hjarta··knos·ar|i} \dicIPA{{\c{c}}{a}{\textsubring{r}}{\textsubring{d}}{a}{k\smash{\textsuperscript{h}}}{n}{\textopeno}{s}{a}{r}{\textsci}} \dicPos{m}[1] \dicFlx{(‑a, ‑ar)}[13] \dicLink{hjartaknúsari}
\dicEntry[hjartaknúsari] \dicTerm{hjarta··knús·ar|i}\dicTerm{, hjartaknosari} \dicIPA{{\c{c}}\-{a}\-{\textsubring{r}}\-{\textsubring{d}}\-{a}\-{k\smash{\textsuperscript{h}}}\-{n}\-{u}\-{s}\-{a}\-{r}\-{\textsci}\-} \dicPos{m}[1] \dicFlx{(‑a, ‑ar)}[13] \dicDirectTranslationCS{lamač srdcí, miláček žen}
\dicEntry[hjartakveisa] \dicTerm{hjarta··kveis|a} \dicIPA{{\c{c}}{a}{\textsubring{r}}{\textsubring{d}}{a}{k\smash{\textsuperscript{h}}}{v}{ei}{s}{a}} \dicPos{f}[1] \dicFlx{(‑u, ‑ur)}[7] \dicFieldCat{med.} \dicDirectTranslationCS{angina pectoris}
\dicEntry[hjartalínurit] \dicTerm{hjarta··línu·rit} \dicIPA{{\c{c}}{a}{\textsubring{r}}{\textsubring{d}}{a}{l}{i}{n}{\textscy}{r}{\textsci}{\textsubring{d}}} \dicPos{n}[2] \dicFlx{(‑s, ‑)}[5] \dicFieldCat{med.} \dicDirectTranslationCS{(elektro)kardiogram, EKG}
\dicEntry[hjartalæknir] \dicTerm{hjarta··lækn|ir} \dicIPA{{\c{c}}{a}{\textsubring{r}}{\textsubring{d}}{a}{l}{a}{i}{h}{\r{g}}{n}{\textsci}{\textsubring{r}}} \dicPos{m}[7] \dicFlx{(‑is, ‑ar)}[1] \dicFieldCat{med.} \dicDirectTranslationCS{kardiolog, kardioložka}
\dicEntry[hjartar] \dicTerm{hjartar} \dicIPA{{\c{c}}{a}{\textsubring{r}}{\textsubring{d}}{a}{\textsubring{r}}} \dicPos{m} \dicFlx{sg gen} \dicLink{hjörtur}
\dicEntry[hjartarbjalla] \dicTerm{hjartar··|bjalla} \dicIPA{{\c{c}}{a}{\textsubring{r}}{\textsubring{d}}{a}{r}{\textsubring{b}}{j}{a}{\textsubring{d}}{l}{a}} \dicPos{f}[1] \dicFlx{(‑bjöllu, ‑bjöllur)}[20] \dicFieldCat{zool.} \dicDirectTranslationCS{roháč, roháč obecný} \textit{(l.~{\textLA{Lucanidae}})}  \dicsymPhoto\ 
\dicFigure{ds_image_hjartarbjalla_0_1.jpg}{Hjartarbjalla}{Hjartarbjalla - Fir0002/Flagstaffotos, GFDL}
\dicEntry[hjartarit] \dicTerm{hjarta··rit} \dicIPA{{\c{c}}{a}{\textsubring{r}}{\textsubring{d}}{a}{r}{\textsci}{\textsubring{d}}} \dicPos{n}[2] \dicFlx{(‑s, ‑)}[5] \dicFieldCat{med.} \dicDirectTranslationCS{(elektro)kardiogram, EKG}
\dicEntry[hjartariti] \dicTerm{hjarta··rit|i} \dicIPA{{\c{c}}{a}{\textsubring{r}}{\textsubring{d}}{a}{r}{\textsci}{\textsubring{d}}{\textsci}} \dicPos{m}[1] \dicFlx{(‑a, ‑ar)}[1] \dicFieldCat{med.} \dicDirectTranslationCS{elektrokardiograf, EKG}
\dicEntry[hjartaskel] \dicTerm{hjarta··skel} \dicIPA{{\c{c}}{a}{\textsubring{r}}{\textsubring{d}}{a}{s}{\r{\textObardotlessj}}{\textepsilon}{\textsubring{l}}} \dicPos{f}[4] \dicFlx{(‑jar, ‑jar)}[9] \dicFieldCat{zool.} \dicDirectTranslationCS{srdcovka jedlá} \textit{(l.~{\textLA{Cerastoderma edule}})}  \dicsymPhoto\ 
\dicFigure{ds_image_hjartaskel_0_1.jpg}{Hjartaskel}{Hjartaskel - Féron Benjamin, CC BY 2.0}
\dicEntry[hjartaslag] \dicTerm{hjarta··|slag} \dicIPA{{\c{c}}{a}{\textsubring{r}}{\textsubring{d}}{a}{s}{\textsubring{d}}{l}{a}{x}} \dicPos{n}[2] \dicFlx{(‑slags, ‑slög)}[8] \textbf{1.} \dicDirectTranslationCS{(srdeční) tep, puls} \dicExampleIS{80 hjartaslög á mínútu} \dicExampleCS{80 tepů za minutu}  \textbf{2.} \dicFieldCat{med.} \dicSynonym{hjartaáfall} \dicDirectTranslationCS{srdeční mrtvice, infarkt (myokardu)}
\dicEntry[hjartastuðtæki] \dicTerm{hjarta··stuð·tæki} \dicIPA{{\c{c}}{a}{\textsubring{r}}{\textsubring{d}}{a}{s}{\textsubring{d}}{\textscy}{ð}{t\smash{\textsuperscript{h}}}{a}{i}{\r{\textObardotlessj}}{\textsci}} \dicPos{n}[2] \dicFlx{(‑s, ‑)}[16] \dicFieldCat{med. } \dicDirectTranslationCS{defibrilátor}
\dicEntry[hjartnæmur] \dicTerm{hjart··næmur} \dicIPA{{\c{c}}{a}{\textsubring{r}}{\textsubring{d}}{n}{a}{i}{m}{\textscy}{\textsubring{r}}} \dicPos{adj}[1]\dicFlx{}[-1] \dicDirectTranslationCS{dojemný, srdcervoucí}
\dicEntry[hjartsláttur] \dicTerm{hjart··slátt|ur} \dicIPA{{\c{c}}{a}{\textsubring{r}}{\textsubring{d}}{s}{\textsubring{d}}{l}{au}{h}{\textsubring{d}}{\textscy}{\textsubring{r}}} \dicPos{m}[12] \dicFlx{(‑ar)}[8] \dicDirectTranslationCS{tlukot srdce, tep, puls} \dicExampleIS{hægur hjartsláttur} \dicExampleCS{pomalý tlukot srdce}
\dicEntry[hjartveiki] \dicTerm{hjart··veik|i} \dicIPA{{\c{c}}{a}{\textsubring{r}}{\textsubring{d}}{v}{ei}{\r{\textObardotlessj}}{\textsci}} \dicPos{f}[3] \dicFlx{(‑i)}[3] \dicDirectTranslationCS{srdeční slabost}
\dicEntry[hjartveikur] \dicTerm{hjart··veikur} \dicIPA{{\c{c}}{a}{\textsubring{r}}{\textsubring{d}}{v}{ei}{\r{g}}{\textscy}{\textsubring{r}}} \dicPos{adj}[1]\dicFlx{}[-1] \textbf{1.} \dicDirectTranslationCS{nemocný na srdce}  \textbf{2.} \dicSynonym{taugaóstyrkur} \dicDirectTranslationCS{nervózní, ustrašený}
\dicEntry[hjá] \dicTerm{hjá} \dicsymFrequent\  \dicIPA{{\c{c}}{au}{\textlengthmark}} \dicPos{prep} \dicFlx{dat} \textbf{1.} \dicSynonym*{við hliðina á} \dicDirectTranslationCS{u, při, vedle} \dicExampleIS{setjast hjá e‑m} \dicExampleCS{sednout si vedle (koho)};  \dicPhraseIS{rétt hjá e‑u} \dicDirectTranslationCS{v~blízkosti (čeho), blízko (čeho), hned vedle (čeho)}  \textbf{2.} \dicDirectTranslationCS{u} \dicIndirectTranslationCS{(o~pobytu někoho)} \dicExampleIS{Barnið er hjá ömmu sinni.} \dicExampleCS{Dítě je u~babičky.}  \textbf{3.} \dicSynonym{meðal\smash{\textsuperscript{2}}} \dicDirectTranslationCS{mezi, u};  \dicPhraseIS{heima hjá e‑m} \dicDirectTranslationCS{u~(koho) doma}  \textbf{4.} \dicDirectTranslationCS{vedle, ve srovnání s}  \textbf{5.} \dicFieldCat{poč.} \dicDirectTranslationCS{zavináč (@)} \dicIndirectTranslationCS{(znak užívaný v~e‑mailové adrese)}
\dicEntry[hjákátlegur] \dicTerm{hjá·kát··legur} \dicIPA{{\c{c}}{au}{\textlengthmark}{k\smash{\textsuperscript{h}}}{au}{\textsubring{d}}{l}{\textepsilon}{\textbabygamma}{\textscy}{\textsubring{r}}} \dicPos{adj}[1]\dicFlx{}[-8] \dicSynonym{undarlegur} \dicDirectTranslationCS{komický, směšný, bizarní}
\dicEntry[hjákona] \dicTerm{hjá··kon|a} \dicIPA{{\c{c}}{au}{\textlengthmark}{k\smash{\textsuperscript{h}}}{\textopeno}{n}{a}} \dicPos{f}[1] \dicFlx{(‑u, ‑ur)}[27] \dicDirectTranslationCS{milenka}
\dicEntry[hjálagður] \dicTerm{hjá··|lagður} \dicIPA{{\c{c}}{au}{\textlengthmark}{l}{a}{\textbabygamma}{ð}{\textscy}{\textsubring{r}}} \dicPos{adj}[2] \dicFlx{(f ‑lögð)}[3] \dicSynonym*{viðlagður} \dicDirectTranslationCS{přiložený} \dicExampleIS{sjá hjálagt bréf} \dicExampleCS{podívat se na přiložený dopis}
\dicEntry[hjálmur] \dicTerm{hjálm|ur} \dicIPA{{\c{c}}{au}{l}{m}{\textscy}{\textsubring{r}}} \dicPos{m}[6] \dicFlx{(‑s, ‑ar)}[4] \textbf{1.} \dicDirectTranslationCS{přilba, helma (cyklistická ap.)} \dicExampleIS{hafa hjálm á höfði} \dicExampleCS{mít helmu na hlavě}  \textbf{2.} \dicDirectTranslationCS{klobouk, stínítko (lampy ap.)}
\dicEntry[hjálp] \dicTerm{hjálp} \dicsymFrequent\  \dicIPA{{\c{c}}{au}{\textsubring{l}}{\textsubring{b}}} \dicPos{f}[7] \dicFlx{(‑ar, ‑ir)}[1] \textbf{1.} \dicSynonym{aðstoð} \dicDirectTranslationCS{(vý)pomoc, asistence} \dicExampleIS{biðja um hjálp} \dicExampleCS{požádat o~pomoc}  \textbf{2.} \dicSynonym{björgun} \dicDirectTranslationCS{pomoc, záchrana};  \dicPhraseIS{koma e‑m til hjálpar} \dicDirectTranslationCS{přijít (komu) na pomoc}
\dicEntry[hjálpa] \dicTerm{hjálp|a} \dicsymFrequent\  \dicIPA{{\c{c}}{au}{\textsubring{l}}{\textsubring{b}}{a}} \dicPos{v}[1] \dicFlx{(‑aði)}[22] \dicFlx{dat} \dicDirectTranslationCS{pomoct, pomoci, pomáhat, asistovat} \dicExampleIS{hjálpa e‑m við verkið} \dicExampleCS{pomoct (komu) s~prací};  \dicIdiom{hjálpa}[til]{ \dicPhraseIS{hjálpa til}} \dicDirectTranslationCS{vypomoct, vypomoci, vypomáhat};  \dicIdiom{hjálpast}[að]{ \dicPhraseIS{hjálpast að}} \dicFlx{refl} \dicDirectTranslationCS{vypomoci si, pomoci jeden druhému} \dicExampleIS{Þær hjálpuðust að við verkið.} \dicExampleCS{Jedna druhé pomáhaly s~prací.};  \dicIdiom{hjálpa}{ \dicPhraseIS{guð hjálpi okkur}} \dicDirectTranslationCS{nedej bože (aby se něco stalo ap.)}; { \dicPhraseIS{guð hjálpi þér!}} \dicDirectTranslationCS{na zdraví!, zdravíčko!} \dicIndirectTranslationCS{(používané při kýchnutí)}
\dicEntry[hjálpargagn] \dicTerm{hjálpar··|gagn} \dicIPA{{\c{c}}{au}{\textsubring{l}}{\textsubring{b}}{a}{r}{\r{g}}{a}{\r{g}}{\textsubring{n}}} \dicPos{n}[2] \dicFlx{(‑gagns, ‑gögn)}[8] \textbf{1.} \dicDirectTranslationCS{pomůcka} \dicExampleIS{mikilvægt hjálpargagn fyrir skóla} \dicExampleCS{důležitá školní pomůcka}  \textbf{2.} \dicPhraseIS{hjálpargögn} \dicFlx{pl} \dicDirectTranslationCS{potřeby, pomůcky} \dicExampleIS{flytja hjálpargögn til nauðstaddra} \dicExampleCS{dovézt potřeby lidem v~nouzi}
\dicEntry[hjálparlaus] \dicTerm{hjálpar··laus} \dicIPA{{\c{c}}{au}{\textsubring{l}}{\textsubring{b}}{a}{r}{l}{\oe i}{s}} \dicPos{adj}[5]\dicFlx{}[-1] \dicSynonym{hjálparvana} \dicDirectTranslationCS{bezradný, bezmocný}
\dicEntry[hjálparleysi] \dicTerm{hjálpar··leysi} \dicIPA{{\c{c}}{au}{\textsubring{l}}{\textsubring{b}}{a}{r}{l}{ei}{s}{\textsci}} \dicPos{n}[2] \dicFlx{(‑s)}[20] \dicDirectTranslationCS{bezmocnost, bezradnost}
\dicEntry[hjálparmaður] \dicTerm{hjálpar··|maður} \dicIPA{{\c{c}}{au}{\textsubring{l}}{\textsubring{b}}{a}{r}{m}{a}{ð}{\textscy}{\textsubring{r}}} \dicPos{m}[13] \dicFlx{(‑manns, ‑menn)}[2] \dicDirectTranslationCS{pomocník, pomocnice, asistent(ka)}
\dicEntry[hjálparsveit] \dicTerm{hjálpar··sveit} \dicIPA{{\c{c}}{au}{\textsubring{l}}{\textsubring{b}}{a}{\textsubring{r}}{s}{v}{ei}{\textsubring{d}}} \dicPos{f}[7] \dicFlx{(‑ar, ‑ir)}[1] \dicDirectTranslationCS{záchranný tým}
\dicEntry[hjálparsögn] \dicTerm{hjálpar··|sögn} \dicIPA{{\c{c}}{au}{\textsubring{l}}{\textsubring{b}}{a}{\textsubring{r}}{s}{\oe}{\r{g}}{\textsubring{n}}} \dicPos{f}[7] \dicFlx{(‑sagnar, ‑sagnir)}[16] \dicFieldCat{jaz.} \dicDirectTranslationCS{pomocné sloveso}
\dicEntry[hjálpartæki] \dicTerm{hjálpar··tæki} \dicIPA{{\c{c}}{au}{\textsubring{l}}{\textsubring{b}}{a}{\textsubring{r}}{t\smash{\textsuperscript{h}}}{a}{i}{\r{\textObardotlessj}}{\textsci}} \dicPos{n}[2] \dicFlx{(‑s, ‑)}[16] \dicDirectTranslationCS{pomůcka} \dicExampleIS{verslun með hjálpartæki ástarlífsins} \dicExampleCS{obchod s~erotickými pomůckami}
\dicEntry[hjálparvana] \dicTerm{hjálpar··vana} \dicIPA{{\c{c}}{au}{\textsubring{l}}{\textsubring{b}}{a}{r}{v}{a}{n}{a}} \dicPos{adj}[13] \dicFlx{indecl}[1] \dicSynonym{hjálparlaus} \dicDirectTranslationCS{bezmocný, bezbranný}
\dicEntry[hjálparþurfi] \dicTerm{hjálpar··þurfi} \dicIPA{{\c{c}}{au}{\textsubring{l}}{\textsubring{b}}{a}{\textsubring{r}}{\texttheta}{\textscy}{r}{v}{\textsci}} \dicPos{adj}[13] \dicFlx{indecl}[1] \dicDirectTranslationCS{potřebující pomoc}
\dicEntry[hjálpfús] \dicTerm{hjálp··fús} \dicIPA{{\c{c}}{au}{\textsubring{l}}{\textsubring{b}}{f}{u}{s}} \dicPos{adj}[5]\dicFlx{}[-1] \dicSynonym{greiðugur} \dicDirectTranslationCS{nápomocný, ochotný pomoci}
\dicEntry[hjálplegur] \dicTerm{hjálp··legur} \dicIPA{{\c{c}}{au}{\textsubring{l}}{\textsubring{b}}{l}{\textepsilon}{\textbabygamma}{\textscy}{\textsubring{r}}} \dicPos{adj}[1]\dicFlx{}[-8] \dicDirectTranslationCS{nápomocný, pomáhající};  \dicPhraseIS{vera e‑m hjálplegur við e‑ð} \dicDirectTranslationCS{být (komu) nápomocný v~(čem)}
\dicEntry[hjálpsamur] \dicTerm{hjálp··|samur} \dicIPA{{\c{c}}{au}{\textsubring{l}}{\textsubring{b}}{s}{a}{m}{\textscy}{\textsubring{r}}} \dicPos{adj}[1] \dicFlx{(f ‑söm)}[2] \dicSynonym{greiðvikinn} \dicDirectTranslationCS{nápomocný, ochotný pomoci}
\dicEntry[hjálpsemi] \dicTerm{hjálp··sem|i} \dicIPA{{\c{c}}{au}{\textsubring{l}}{\textsubring{b}}{s}{\textepsilon}{m}{\textsci}} \dicPos{f}[2] \dicFlx{(‑i)}[2] \dicDirectTranslationCS{nápomocnost, ochota pomoci}
\dicEntry[hjáróma] \dicTerm{hjá··róma} \dicIPA{{\c{c}}{au}{\textlengthmark}{r}{ou}{m}{a}} \dicPos{adj}[13] \dicFlx{indecl}[1] \dicSynonym{falskur} \dicDirectTranslationCS{falešný (tón ap.)}
\dicEntry[hjárænulegur] \dicTerm{hjá·rænu··legur} \dicIPA{{\c{c}}{au}{\textlengthmark}{r}{a}{i}{n}{\textscy}{l}{\textepsilon}{\textbabygamma}{\textscy}{\textsubring{r}}} \dicPos{adj}[1]\dicFlx{}[-8] \dicSynonym*{rolulegur} \dicDirectTranslationCS{zakřiknutý, nejistý}
\dicEntry[hjásvæfa] \dicTerm{hjá··svæf|a} \dicIPA{{\c{c}}{au}{\textlengthmark}{s}{v}{a}{i}{v}{a}} \dicPos{f}[1] \dicFlx{(‑u, ‑ur)}[19] \dicSynonym{rekkjunautur} \dicDirectTranslationCS{druh, družka} \dicIndirectTranslationCS{(osoba, s~níž má člověk sexuální vztah)}
\dicEntry[hjátrú] \dicTerm{hjá··trú} \dicIPA{{\c{c}}{au}{\textlengthmark}{t\smash{\textsuperscript{h}}}{r}{u}} \dicPos{f}[4] \dicFlx{(‑ar)}[27] \dicSynonym{hindurvitni} \dicDirectTranslationCS{pověra}
\dicEntry[hjátrúarfullur] \dicTerm{hjá·trúar··|fullur} \dicIPA{{\c{c}}{au}{\textlengthmark}{t\smash{\textsuperscript{h}}}{r}{u}{a}{\textsubring{r}}{f}{\textscy}{\textsubring{d}}{l}{\textscy}{\textsubring{r}}} \dicPos{adj}[10] \dicFlx{(comp ‑fyllri, sup ‑fyllstur)}[7] \dicDirectTranslationCS{pověrčivý}
\dicEntry[hjáverk] \dicTerm{hjá··verk} \dicIPA{{\c{c}}{au}{\textlengthmark}{v}{\textepsilon}{\textsubring{r}}{\r{g}}} \dicPos{n}[2] \dicFlx{(‑s, ‑)}[5] \dicSynonym*{aukaverk} \dicDirectTranslationCS{vedlejší práce};  \dicPhraseIS{gera e‑ð í hjáverkum} \dicDirectTranslationCS{dělat (co) bokem}
\dicEntry[hjó] \dicTerm{hjó} \dicIPA{{\c{c}}{ou}{\textlengthmark}} \dicPos{v} \dicFlx{ind pf sg 1 pers} \dicLink{höggva}
\dicEntry[hjól] \dicTerm{hjól} \dicsymFrequent\  \dicIPA{{\c{c}}{ou}{\textlengthmark}{\textsubring{l}}} \dicPos{n}[2] \dicFlx{(‑s, ‑)}[5] \textbf{1.} \dicDirectTranslationCS{kolo (auta ap.)}  \textbf{2.} \dicDirectTranslationCS{(jízdní) kolo, bicykl} \dicExampleIS{reiðhjól} \dicExampleCS{jízdní kolo}
\dicEntry[hjóla] \dicTerm{hjól|a} \dicsymFrequent\  \dicIPA{{\c{c}}{ou}{\textlengthmark}{l}{a}} \dicPos{v}[1] \dicFlx{(‑aði)}[44] \dicDirectTranslationCS{jet\,/\addthin jezdit na kole} \dicExampleIS{Ég hjóla alltaf í skólann á morgnanna.} \dicExampleCS{Ráno jezdím do školy vždycky na kole.}
\dicEntry[hjólabretti] \dicTerm{hjóla··bretti} \dicIPA{{\c{c}}{ou}{\textlengthmark}{l}{a}{\textsubring{b}}{r}{\textepsilon}{h}{\textsubring{d}}{\textsci}} \dicPos{n}[2] \dicFlx{(‑s, ‑)}[14] \dicFieldCat{sport.} \dicDirectTranslationCS{skateboard}
\dicEntry[hjólagrafa] \dicTerm{hjóla··|grafa} \dicIPA{{\c{c}}{ou}{\textlengthmark}{l}{a}{\r{g}}{r}{a}{v}{a}} \dicPos{f}[1] \dicFlx{(‑gröfu, ‑gröfur)}[8] \dicFieldCat{techn.} \dicDirectTranslationCS{bagr}
\dicEntry[hjólakeppni] \dicTerm{hjóla··keppn|i} \dicIPA{{\c{c}}{ou}{\textlengthmark}{l}{a}{c\smash{\textsuperscript{h}}}{\textepsilon}{h}{\textsubring{b}}{n}{\textsci}} \dicPos{f}[3] \dicFlx{(‑i, ‑ir)}[1] \dicFieldCat{sport.} \dicDirectTranslationCS{cyklistický závod}
\dicEntry[hjólaskautar] \dicTerm{hjóla··skautar} \dicIPA{{\c{c}}{ou}{\textlengthmark}{l}{a}{s}{\r{g}}{\oe i}{\textsubring{d}}{a}{\textsubring{r}}} \dicPos{m}[1] \dicFlx{pl}[2] \dicFieldCat{sport.} \dicDirectTranslationCS{kolečkové brusle}
\dicEntry[hjólaskip] \dicTerm{hjóla··skip} \dicIPA{{\c{c}}{ou}{\textlengthmark}{l}{a}{s}{\r{\textObardotlessj}}{\textsci}{\textsubring{b}}} \dicPos{n}[2] \dicFlx{(‑s, ‑)}[5] \dicFieldCat{nám.} \dicDirectTranslationCS{kolesový parník}
\dicEntry[hjólaskófla] \dicTerm{hjóla··skófl|a} \dicIPA{{\c{c}}{ou}{\textlengthmark}{l}{a}{s}{\r{g}}{ou}{\textsubring{b}}{l}{a}} \dicPos{f}[1] \dicFlx{(‑u, ‑ur)}[19] \dicFieldCat{techn.} \dicDirectTranslationCS{nakládač, nakladač} \dicIndirectTranslationCS{(stroj)}
\dicEntry[hjólastígur] \dicTerm{hjóla··stíg|ur} \dicIPA{{\c{c}}{ou}{\textlengthmark}{l}{a}{s}{\textsubring{d}}{i}{\textbabygamma}{\textscy}{\textsubring{r}}} \dicPos{m}[6] \dicFlx{(‑s, ‑ar\,/\addthin ‑ir)}[70] \dicDirectTranslationCS{cyklostezka}
\dicEntry[hjólastóll] \dicTerm{hjóla··stól|l} \dicIPA{{\c{c}}{ou}{\textlengthmark}{l}{a}{s}{\textsubring{d}}{ou}{\textsubring{d}}{\textsubring{l}}} \dicPos{m}[6] \dicFlx{(‑s, ‑ar)}[48] \dicDirectTranslationCS{invalidní vozík, vozíček} \dicExampleIS{vera í hjólastól} \dicExampleCS{být na vozíčku}
\dicEntry[hjólbarði] \dicTerm{hjól··barð|i} \dicIPA{{\c{c}}{ou}{l}{\textsubring{b}}{a}{r}{ð}{\textsci}} \dicPos{m}[1] \dicFlx{(‑a, ‑ar)}[8] \dicSynonym{dekk} \dicDirectTranslationCS{plášť} \dicIndirectTranslationCS{(obal duše)}
\dicEntry[hjólbeinóttur] \dicTerm{hjól··bein·óttur} \dicIPA{{\c{c}}{ou}{l}{\textsubring{b}}{ei}{n}{ou}{h}{\textsubring{d}}{\textscy}{\textsubring{r}}} \dicPos{adj}[1]\dicFlx{}[-10] \dicSynonym*{hjólfættur} \dicDirectTranslationCS{křivonohý, (jsoucí) s~křivýma nohama}
\dicEntry[hjólbörur] \dicTerm{hjól··börur} \dicIPA{{\c{c}}{ou}{l}{\textsubring{b}}{\oe}{r}{\textscy}{\textsubring{r}}} \dicPos{f}[12] \dicFlx{pl}[4] \dicDirectTranslationCS{kolečko} \dicIndirectTranslationCS{(nářadí s~jedním kolem)}
\dicEntry[hjólhlíf] \dicTerm{hjól··hlíf} \dicIPA{{\c{c}}{ou}{l}{\textsubring{l}}{i}{f}} \dicPos{f}[4] \dicFlx{(‑ar, ‑ar)}[1] \dicDirectTranslationCS{blatník}
\dicEntry[hjólhýsi] \dicTerm{hjól··hýsi} \dicIPA{{\c{c}}{ou}{l}{h}{i}{s}{\textsci}} \dicPos{n}[2] \dicFlx{(‑s, ‑)}[14] \dicDirectTranslationCS{karavan, obytný přívěs}
\dicEntry[hjólkeðja] \dicTerm{hjól··keðj|a} \dicIPA{{\c{c}}{ou}{l}{c\smash{\textsuperscript{h}}}{\textepsilon}{ð}{j}{a}} \dicPos{f}[1] \dicFlx{(‑u, ‑ur)}[7] \dicDirectTranslationCS{řetěz (u~kola)}
\dicEntry[hjólkoppur] \dicTerm{hjól··kopp|ur} \dicIPA{{\c{c}}{ou}{l}{k\smash{\textsuperscript{h}}}{\textopeno}{h}{\textsubring{b}}{\textscy}{\textsubring{r}}} \dicPos{m}[6] \dicFlx{(‑s, ‑ar)}[24] \dicDirectTranslationCS{poklice} \dicIndirectTranslationCS{(kryt kola u~auta)}
\dicEntry[hjólnöf] \dicTerm{hjól··|nöf} \dicIPA{{\c{c}}{ou}{l}{n}{\oe}{f}} \dicPos{f}[7] \dicFlx{(‑nafar, ‑nafir)}[16] \dicDirectTranslationCS{náboj} \dicIndirectTranslationCS{(část kola)}
\dicEntry[hjólpumpa] \dicTerm{hjól··pump|a} \dicIPA{{\c{c}}{ou}{l}{p\smash{\textsuperscript{h}}}{\textscy}{\textsubring{m}}{\textsubring{b}}{a}} \dicPos{f}[1] \dicFlx{(‑u, ‑ur)}[19] \dicDirectTranslationCS{pumpička, hustilka} \dicIndirectTranslationCS{(nástroj na nafukování duší u~kola ap.)}
\dicEntry[hjólreiðamaður] \dicTerm{hjól·reiða··|maður} \dicIPA{{\c{c}}{ou}{l}{r}{ei}{ð}{a}{m}{a}{ð}{\textscy}{\textsubring{r}}} \dicPos{m}[13] \dicFlx{(‑manns, ‑menn)}[2] \dicDirectTranslationCS{cyklista, cyklistka}
\dicEntry[hjólreiðar] \dicTerm{hjól··reiðar} \dicIPA{{\c{c}}{ou}{l}{r}{ei}{ð}{a}{\textsubring{r}}} \dicPos{f}[7] \dicFlx{pl}[2] \dicDirectTranslationCS{cyklistika}
\dicEntry[hjólteinn] \dicTerm{hjól··tein|n} \dicIPA{{\c{c}}{ou}{l}{t\smash{\textsuperscript{h}}}{ei}{\textsubring{d}}{\textsubring{n}}} \dicPos{m}[6] \dicFlx{(‑s, ‑ar)}[42] \dicDirectTranslationCS{výplet (kola)}
\dicEntry[hjóm] \dicTerm{hjóm} \dicIPA{{\c{c}}{ou}{\textlengthmark}{\textsubring{m}}} \dicPos{n}[2] \dicFlx{(‑s)}[2] \dicSynonym{fánýti} \dicDirectTranslationCS{nicotnost, bezcennost}
\dicEntry[hjón] \dicTerm{hjón} \dicsymFrequent\  \dicIPA{{\c{c}}{ou}{\textlengthmark}{\textsubring{n}}} \dicPos{n}[2] \dicFlx{pl}[1] \dicDirectTranslationCS{(manželský) pár, manželé} \dicExampleIS{Ung hjón með tvö börn keyptu húsið.} \dicExampleCS{Dům koupil mladý pár s~dvěma dětmi.}
\dicEntry[hjónaband] \dicTerm{hjóna··|band} \dicsymFrequent\  \dicIPA{{\c{c}}{ou}{\textlengthmark}{n}{a}{\textsubring{b}}{a}{n}{\textsubring{d}}} \dicPos{n}[2] \dicFlx{(‑bands, ‑bönd)}[8] \dicDirectTranslationCS{manželství} \dicExampleIS{tuttugu ára hjónaband} \dicExampleCS{dvacetileté manželství};  \dicPhraseIS{ganga í hjónaband} \dicDirectTranslationCS{oženit se}
\dicEntry[hjónagras] \dicTerm{hjóna··|gras} \dicIPA{{\c{c}}{ou}{\textlengthmark}{n}{a}{\r{g}}{r}{a}{s}} \dicPos{n}[2] \dicFlx{(‑grass, ‑grös)}[8] \dicFieldCat{bot.} \dicDirectTranslationCS{běloprstka bělavá} \textit{(l.~{\textLA{Pseudorchis albida}})}  \dicsymPhoto\ 
\dicFigure{ds_image_hjonagras_0_1.jpg}{Hjónagras}{Hjónagras - BerndH, CC BY-SA 2.5}
\dicEntry[hjónarúm] \dicTerm{hjóna··rúm} \dicIPA{{\c{c}}{ou}{\textlengthmark}{n}{a}{r}{u}{\textsubring{m}}} \dicPos{n}[2] \dicFlx{(‑s, ‑)}[5] \dicDirectTranslationCS{manželská postel}
\dicEntry[hjónaskilnaður] \dicTerm{hjóna··skil·nað|ur} \dicIPA{{\c{c}}{ou}{\textlengthmark}{n}{a}{s}{\r{\textObardotlessj}}{\textsci}{l}{n}{a}{ð}{\textscy}{\textsubring{r}}} \dicPos{m}[10] \dicFlx{(‑ar, ‑ir)}[6] \dicDirectTranslationCS{rozvod}
\dicEntry[hjuggum] \dicTerm{hjuggum} \dicIPA{{\c{c}}{\textscy}{\r{g}}{\textlengthmark}{\textscy}{\textsubring{m}}} \dicPos{v} \dicFlx{ind pf pl 1 pers} \dicLink{höggva}
\dicEntry[hjú] \dicTerm{hjú} \dicIPA{{\c{c}}{u}{\textlengthmark}} \dicPos{n}[2] \dicFlx{(‑s, ‑)}[5] \dicDirectTranslationCS{služebnictvo, sloužící}
\dicEntry[hjúfra] \dicTerm{hjúfr|a} \dicIPA{{\c{c}}{u}{v}{r}{a}} \dicPos{v}[1] \dicFlx{(‑aði)}[48] \textbf{1.} \dicPhraseIS{það hjúfrar} \dicFlx{impers} \dicDirectTranslationCS{(trochu) prší}  \textbf{2.} \dicPhraseIS{hjúfra sig upp að e‑m} \dicDirectTranslationCS{přitulit\,/\addthin přitisknout se ke (komu)}
\dicEntry[hjúkra] \dicTerm{hjúkr|a} \dicIPA{{\c{c}}{u}{\textlengthmark}{\r{g}}{r}{a}} \dicPos{v}[1] \dicFlx{(‑aði)}[44] \dicFlx{dat} \dicDirectTranslationCS{ošetřovat, pečovat, opatrovat} \dicExampleIS{hjúkra e‑m} \dicExampleCS{pečovat o~(koho)}
\dicEntry[hjúkrun] \dicTerm{hjúkr|un} \dicIPA{{\c{c}}{u}{\textlengthmark}{\r{g}}{r}{\textscy}{\textsubring{n}}} \dicPos{f}[7] \dicFlx{(‑unar)}[9] \dicDirectTranslationCS{ošetřování, pečování}
\dicEntry[hjúkrunarfræði] \dicTerm{hjúkrunar··fræð|i} \dicIPA{{\c{c}}{u}{\textlengthmark}{\r{g}}{r}{\textscy}{n}{a}{\textsubring{r}}{f}{r}{a}{i}{ð}{\textsci}} \dicPos{f}[3] \dicFlx{(‑i)}[3] \dicDirectTranslationCS{ošetřovatelství, pečovatelství} \dicIndirectTranslationCS{(vědní obor)}
\dicEntry[hjúkrunarfræðingur] \dicTerm{hjúkrunar·fræð··ing|ur} \dicIPA{{\c{c}}\-{u}\-{\textlengthmark}\-{\r{g}}\-{r}\-{\textscy}\-{n}\-{a}\-{\textsubring{r}}\-{f}\-{r}\-{a}\-{i}\-{ð}\-{i}\-{\ng}\-{\r{g}}\-{\textscy}\-{\textsubring{r}}\-} \dicPos{m}[6] \dicFlx{(‑s, ‑ar)}[8] \dicDirectTranslationCS{ošetřovatel(ka), zdravotní sestra}
\dicEntry[hjúkrunarheimili] \dicTerm{hjúkrunar··heimili} \dicIPA{{\c{c}}{u}{\textlengthmark}{\r{g}}{r}{\textscy}{n}{a}{\textsubring{r}}{h}{ei}{m}{\textsci}{l}{\textsci}} \dicPos{n}[2] \dicFlx{(‑s, ‑)}[14] \dicDirectTranslationCS{pečovatelský dům, dům s~pečovatelskou službou, sanatorium}
\dicEntry[hjúkrunarkona] \dicTerm{hjúkrunar··kon|a} \dicIPA{{\c{c}}{u}{\textlengthmark}{\r{g}}{r}{\textscy}{n}{a}{\textsubring{r}}{k\smash{\textsuperscript{h}}}{\textopeno}{n}{a}} \dicPos{f}[1] \dicFlx{(‑u, ‑ur)}[27] \dicDirectTranslationCS{sestřička, ošetřovatelka}
\dicEntry[hjúpa] \dicTerm{hjúp|a} \dicIPA{{\c{c}}{u}{\textlengthmark}{\textsubring{b}}{a}} \dicPos{v}[1] \dicFlx{(‑aði)}[1] \dicFlx{acc} \textbf{1.} \dicDirectTranslationCS{obalit, (hojně) pokrýt} \dicExampleIS{hjúpa e‑ð í e‑u} \dicExampleCS{obalit (co) v~(čem)}  \textbf{2.} \dicDirectTranslationCS{zahalit, zakrýt, přikrýt}
\dicEntry[hjúpur] \dicTerm{hjúp|ur} \dicIPA{{\c{c}}{u}{\textlengthmark}{\textsubring{b}}{\textscy}{\textsubring{r}}} \dicPos{m}[6] \dicFlx{(‑s, ‑ar)}[22] \dicSynonym{blæja} \dicDirectTranslationCS{povlak, obal, plášť}
\dicEntry[hjúskaparstétt] \dicTerm{hjú·skapar··stétt} \dicIPA{{\c{c}}{u}{\textlengthmark}{s}{\r{g}}{a}{\textsubring{b}}{a}{\textsubring{r}}{s}{\textsubring{d}}{j}{\textepsilon}{h}{\textsubring{d}}} \dicPos{f}[7] \dicFlx{(‑ar)}[3] \dicDirectTranslationCS{rodinný stav}
\dicEntry[hjúskapur] \dicTerm{hjú··skap|ur} \dicIPA{{\c{c}}{u}{\textlengthmark}{s}{\r{g}}{a}{\textsubring{b}}{\textscy}{\textsubring{r}}} \dicPos{m}[10] \dicFlx{(‑ar)}[15] \dicSynonym{hjónaband} \dicDirectTranslationCS{manželství}
\dicEntry[hjyggi] \dicTerm{hjyggi} \dicIPA{{\c{c}}{\textsci}{\r{\textObardotlessj}}{\textlengthmark}{\textsci}} \dicPos{v} \dicFlx{con pf sg 1 pers} \dicLink{höggva}
\dicEntry[hjörð] \dicTerm{hjörð} \dicIPA{{\c{c}}{\oe}{r}{\texttheta}} \dicPos{f}[7] \dicFlx{(hjarðar, hjarðir)}[16] \dicDirectTranslationCS{stádo} \dicExampleIS{hreindýrahjörð} \dicExampleCS{stádo sobů}
\dicEntry[hjörn] \dicTerm{hjörn} \dicIPA{{\c{c}}{\oe}{r}{\textsubring{d}}{\textsubring{n}}} \dicPos{n} \dicFlx{pl nom} \dicLink{hjarn}
\dicEntry[hjörtu] \dicTerm{hjörtu} \dicIPA{{\c{c}}{\oe}{\textsubring{r}}{\textsubring{d}}{\textscy}} \dicPos{n} \dicFlx{pl nom} \dicLink{hjarta}
\dicEntry[hjörtur] \dicTerm{hjörtur} \dicIPA{{\c{c}}{\oe}{\textsubring{r}}{\textsubring{d}}{\textscy}{\textsubring{r}}} \dicPos{m}[12] \dicFlx{(hjartar, hirtir)}[3] \dicFieldCat{zool.} \dicDirectTranslationCS{jelen} \textit{(l.~{\textLA{Cer-\linebreak vus}})} \dicsymPhoto\ 
\dicFigure{80409.jpg}{Hjörtur}{Hjörtur - U.S. Fish and Wildlife Service, Biolib, PD}
\dicEntry[hjöru] \dicTerm{hjöru} \dicIPA{{\c{c}}{\oe}{\textlengthmark}{r}{\textscy}} \dicPos{f} \dicFlx{sg gen} \dicLink{hjara\smash{\textsuperscript{1}}}
\dicEntry[hk.] \dicTerm{hk.}\dicTerm{, hvk.} \dicPos{zkr} \dicPhraseIS{hvorugkyn} \dicFieldCat{jaz.} \dicDirectTranslationCS{střední rod}
\dicEntry[hl.] \dicTerm{hl.} \dicPos{zkr} \dicPhraseIS{hluti} \dicDirectTranslationCS{část}
\dicEntry[hlað] \dicTerm{hlað} \dicsymFrequent\  \dicIPA{{\textsubring{l}}{a}{\textlengthmark}{\texttheta}} \dicPos{n}[2] \dicFlx{(‑s, hlöð)}[8] \dicDirectTranslationCS{dvůr (statku ap.)} \dicExampleIS{aka úr hlaði} \dicExampleCS{vyjet ze dvora};  \dicPhraseIS{fylgja e‑u úr hlaði} \dicLangCat{přen.} \dicDirectTranslationCS{připojit několik připomínek k~(čemu)}
\dicEntry[hlaða] \dicTerm{hlaða\smash{\textsuperscript{1}}} \dicIPA{{\textsubring{l}}{a}{\textlengthmark}{ð}{a}} \dicPos{f}[1] \dicFlx{(hlöðu, hlöður)}[20] \dicDirectTranslationCS{stodola} \dicExampleIS{safna heyinu í hlöðu} \dicExampleCS{sbírat seno do stodoly}
\dicEntry[hlaða] \dicTerm{hlaða\smash{\textsuperscript{2}}} \dicsymFrequent\  \dicIPA{{\textsubring{l}}{a}{\textlengthmark}{ð}{a}} \dicPos{v}[6] \dicFlx{(hleð, hlóð, hlóðum, hlæði, hlaðið)}[107] \dicFlx{acc\,/\addthin dat} \textbf{1.} \dicFlx{dat} \dicDirectTranslationCS{postavit, stavět (zeď ap.)} \dicExampleIS{hlaða vegg} \dicExampleCS{postavit zeď};  \dicPhraseIS{hlaða e‑u á e‑n} \dicDirectTranslationCS{nakupit (co) na (koho), naložit (co) na (koho)} \dicExampleIS{hlaða verkefnum á hana} \dicExampleCS{naložit na ni úkoly}  \textbf{2.} \dicFlx{acc} \dicSynonym{ferma} \dicDirectTranslationCS{naložit, nakládat} \dicExampleIS{hlaða bílinn varningi} \dicExampleCS{nakládat auto zbožím}  \textbf{3.} \dicFlx{acc} \dicDirectTranslationCS{nabít, nabíjet, dobít, dobíjet (baterii ap.)} \dicExampleIS{hlaða farsíma} \dicExampleCS{nabít mobil}  \textbf{4.} \dicFlx{acc} \dicDirectTranslationCS{nabít, nabíjet (zbraň ap.)}  \textbf{5.} \dicFlx{dat} \dicFieldCat{poč.} \dicDirectTranslationCS{nahrát, nahrávat, stáhnout, stahovat};  \dicIdiom{hlaða}[niður]{ \dicPhraseIS{hlaða niður e‑u}} \dicFieldCat{poč.} \dicSynonym{sækja} \dicDirectTranslationCS{stáhnout\,/\addthin stahovat (co)};  \dicIdiom{hlaða}[upp]{ \dicPhraseIS{hlaða upp e‑u}} \dicFieldCat{poč.} \dicDirectTranslationCS{nahrát\,/\addthin nahrávat (co)};  \dicIdiom{hlaðast}[á]{ \dicPhraseIS{það hlaðast á e‑n e‑að}} \dicFlx{impers refl} \dicDirectTranslationCS{(co) se valí na (koho) (úkoly ap.)};  \dicIdiom{hlaðast}[upp]{ \dicPhraseIS{hlaðast upp}} \dicFlx{refl} \dicDirectTranslationCS{nakupit se, nahromadit se}
\dicEntry[hlaðborð] \dicTerm{hlað··borð} \dicIPA{{\textsubring{l}}{a}{ð}{\textsubring{b}}{\textopeno}{r}{\texttheta}} \dicPos{n}[2] \dicFlx{(‑s, ‑)}[5] \dicDirectTranslationCS{bufet, švédský stůl}
\dicEntry[hlaði] \dicTerm{hlað|i} \dicIPA{{\textsubring{l}}{a}{\textlengthmark}{ð}{\textsci}} \dicPos{m}[1] \dicFlx{(‑a, ‑ar)}[8] \dicSynonym{stakkur} \dicDirectTranslationCS{hromada, hranice (dříví ap.)} \dicExampleIS{hlaði af bókum} \dicExampleCS{hromada knížek}
\dicEntry[hlaðið] \dicTerm{hlaðið} \dicIPA{{\textsubring{l}}{a}{\textlengthmark}{ð}{\textsci}{\texttheta}} \dicPos{v} \dicFlx{supin} \dicLink{hlaða\smash{\textsuperscript{2}}}
\dicEntry[hlaðinn] \dicTerm{hlaðinn} \dicsymFrequent\  \dicIPA{{\textsubring{l}}{a}{\textlengthmark}{ð}{\textsci}{\textsubring{n}}} \dicPos{adj}[6]\dicFlx{}[-3] \textbf{1.} \dicDirectTranslationCS{naložený, obtížený (úkoly ap.)} \dicExampleIS{hlaðinn e‑u} \dicExampleCS{obtížený (čím)}  \textbf{2.} \dicDirectTranslationCS{nabitý (baterie ap.)} \dicExampleIS{Fartölvan er ekki hlaðin.} \dicExampleCS{Notebook není nabitý.}  \textbf{3.} \dicDirectTranslationCS{nabitý (zbraň ap.)}  \textbf{4.} \dicDirectTranslationCS{nakupený, navršený} \dicExampleIS{hlaðinn veggur} \dicExampleCS{zděná zídka}
\dicEntry[hlaðkolla] \dicTerm{hlað··koll|a} \dicIPA{{\textsubring{l}}{a}{ð}{k\smash{\textsuperscript{h}}}{\textopeno}{\textsubring{d}}{l}{a}} \dicPos{f}[1] \dicFlx{(‑u, ‑ur)}[19] \dicFieldCat{bot.} \dicDirectTranslationCS{heřmánek terčovitý} \textit{(l.~{\textLA{Matricaria discoidea}})}  \dicsymPhoto\ 
\dicFigure{ds_image_hladkolla_0_2.jpg}{Hlaðkolla}{Hlaðkolla - Krzysztof Ziarnek, GFDL}
\dicEntry[hlakka] \dicTerm{hlakk|a} \dicsymFrequent\  \dicIPA{{\textsubring{l}}{a}{h}{\r{g}}{a}} \dicPos{v}[1] \dicFlx{(‑aði)}[26] \textbf{1.} \dicPhraseIS{hlakka til e‑s} \dicDirectTranslationCS{těšit se na (co)} \dicExampleIS{hlakka til jólanna} \dicExampleCS{těšit se na Vánoce}  \textbf{2.} \dicPhraseIS{hlakka yfir e‑u} \dicDirectTranslationCS{mít škodolibou radost (z~čeho), (škodolibě) se radovat z~(čeho)} \dicExampleIS{hlakka yfir óförum annarra} \dicExampleCS{radovat se z~nezdarů jiných};  \dicPhraseIS{það hlakkar í e‑n} \dicFlx{impers} \dicDirectTranslationCS{(kdo) se škodolibě těší\,/\addthin raduje}
\dicEntry[hlamma] \dicTerm{hlamm|a} \dicIPA{{\textsubring{l}}{a}{m}{\textlengthmark}{a}} \dicPos{v}[1] \dicFlx{(‑aði)}[28] \dicPhraseIS{hlamma sér niður} \dicDirectTranslationCS{praštit sebou, hlučně si sednout}
\dicEntry[hland] \dicTerm{hland} \dicIPA{{\textsubring{l}}{a}{n}{\textsubring{d}}} \dicPos{n}[2] \dicFlx{(‑s)}[2] \textbf{1.} \dicSynonym{þvag} \dicDirectTranslationCS{moč}  \textbf{2.} \dicSynonym*{lélegur drykkur} \dicDirectTranslationCS{břečka, slabý nápoj (kafe ap.)}
\dicEntry[hlass] \dicTerm{hlass} \dicIPA{{\textsubring{l}}{a}{s}{\textlengthmark}} \dicPos{n}[2] \dicFlx{(‑, hlöss)}[23] \dicDirectTranslationCS{náklad, fůra} \dicIndirectTranslationCS{(množství nákladu, které se vejde na vůz)}
\dicEntry[hlaup] \dicTerm{hlaup} \dicsymFrequent\  \dicIPA{{\textsubring{l}}{\oe i}{\textlengthmark}{\textsubring{b}}} \dicPos{n}[2] \dicFlx{(‑s, ‑)}[5] \textbf{1.} \dicSynonym{sprettur} \dicDirectTranslationCS{běh, běhání};  \dicPhraseIS{vera á hlaupum} \dicDirectTranslationCS{pobíhat} \dicIndirectTranslationCS{(spěchat)}  \textbf{2.} \dicSynonym{kapphlaup} \dicDirectTranslationCS{běh, běžecký závod}  \textbf{3.} \dicSynonym{jökulhlaup} \dicDirectTranslationCS{vzedmutí, příval (vody v~řece ap.)}  \textbf{4.} \dicFieldCat{kulin.} \dicDirectTranslationCS{aspik, rosol} \dicExampleIS{fiskur í hlaupi} \dicExampleCS{ryba v~aspiku}  \textbf{5.} \dicSynonym{byssuhlaup} \dicDirectTranslationCS{hlaveň (zbraně ap.)}  \textbf{6.} \dicFieldCat{kulin.} \dicDirectTranslationCS{želé} \dicIndirectTranslationCS{(druh bonbónů)}
\dicEntry[hlaupa] \dicTerm{hlaupa} \dicsymFrequent\  \dicIPA{{\textsubring{l}}{\oe i}{\textlengthmark}{\textsubring{b}}{a}} \dicPos{v}[6] \dicFlx{(hleyp, hljóp, hlupum, hlypi, hlaupið)}[90] \textbf{1.} \dicDirectTranslationCS{běžet, běhat, utíkat} \dicExampleIS{hlaupa upp á veg} \dicExampleCS{běžet na silnici};  \dicPhraseIS{hlaupa eins og fætur toga} \dicDirectTranslationCS{běžet jako o~život}  \textbf{2.} \dicSynonym*{dragast saman} \dicDirectTranslationCS{srazit se, scvrknout se (látka ap.)} \dicExampleIS{Jakkinn hefur hlaupið í þvotti.} \dicExampleCS{Bunda se srazila při praní.}  \textbf{3.} \dicDirectTranslationCS{vzedmout se (o~řece)};  \dicIdiom{hlaupa}[á]{ \dicPhraseIS{hlaupa á sig}} \dicSynonym*{gera skyssu} \dicDirectTranslationCS{přehmátnout se, zmýlit se};  \dicIdiom{hlaupa}[frá]{ \dicPhraseIS{hlaupa frá e‑u}} \dicDirectTranslationCS{upustit od (čeho), odbýt (co) (práci ap.)};  \dicIdiom{hlaupa}[í]{ \dicPhraseIS{það hleypur í e‑n e‑að}} \dicFlx{impers} \dicDirectTranslationCS{(co) do (koho) vchází (vztek ap.)};  \dicIdiom{hlaupa}[til]{ \dicPhraseIS{hlaupa til}} \dicDirectTranslationCS{přiběhnout, přiskočit};  \dicIdiom{hlaupa}[uppi]{ \dicPhraseIS{hlaupa e‑n uppi}} \dicDirectTranslationCS{dohonit (koho), dohnat (koho)};  \dicIdiom{hlaupa}[út undan]{ \dicPhraseIS{hlaupa út undan sér}} \dicLangCat{přen.} \dicDirectTranslationCS{šlápnout vedle};  \dicIdiom{hlaupa}[yfir]{ \dicPhraseIS{hlaupa yfir e‑ð}} \dicDirectTranslationCS{přeskočit (co), vynechat (co)} \dicExampleIS{hlaupa yfir einn kafla} \dicExampleCS{přeskočit jednu kapitolu};  \dicIdiom{hlaupast}{ \dicPhraseIS{hlaupast}} \dicFlx{refl} \dicDirectTranslationCS{utéct, uprchnout} \dicExampleIS{hlaupast undan ábyrgðinni} \dicExampleCS{utéct před zodpovědností}
\dicEntry[hlaupabóla] \dicTerm{hlaupa··ból|a} \dicIPA{{\textsubring{l}}{\oe i}{\textlengthmark}{\textsubring{b}}{a}{\textsubring{b}}{ou}{l}{a}} \dicPos{f}[1] \dicFlx{(‑u)}[5] \dicFieldCat{med.} \dicDirectTranslationCS{plané neštovice, varicela}
\dicEntry[hlaupabraut] \dicTerm{hlaupa··braut} \dicIPA{{\textsubring{l}}{\oe i}{\textlengthmark}{\textsubring{b}}{a}{\textsubring{b}}{r}{\oe i}{\textsubring{d}}} \dicPos{f}[7] \dicFlx{(‑ar, ‑ir)}[1] \dicFieldCat{sport.} \dicDirectTranslationCS{běžecká dráha}
\dicEntry[hlaupabretti] \dicTerm{hlaupa··bretti} \dicIPA{{\textsubring{l}}{\oe i}{\textlengthmark}{\textsubring{b}}{a}{\textsubring{b}}{r}{\textepsilon}{h}{\textsubring{d}}{\textsci}} \dicPos{n}[2] \dicFlx{(‑s, ‑)}[14] \dicDirectTranslationCS{běhací\,/\addthin běžecký pás}
\dicEntry[hlaupahjól] \dicTerm{hlaupa··hjól} \dicIPA{{\textsubring{l}}{\oe i}{\textlengthmark}{\textsubring{b}}{a}{\c{c}}{ou}{\textsubring{l}}} \dicPos{n}[2] \dicFlx{(‑s, ‑)}[5] \dicDirectTranslationCS{koloběžka}
\dicEntry[hlaupareikningur] \dicTerm{hlaupa··reikn·ing|ur} \dicIPA{{\textsubring{l}}{\oe i}{\textlengthmark}{\textsubring{b}}{a}{r}{ei}{h}{\r{g}}{n}{i}{\ng}{\r{g}}{\textscy}{\textsubring{r}}} \dicPos{m}[6] \dicFlx{(‑s, ‑ar)}[8] \dicDirectTranslationCS{přečerpatelný účet, kontokorent}
\dicEntry[hlaupari] \dicTerm{hlaup··ar|i} \dicIPA{{\textsubring{l}}{\oe i}{\textlengthmark}{\textsubring{b}}{a}{r}{\textsci}} \dicPos{m}[1] \dicFlx{(‑a, ‑ar)}[13] \dicDirectTranslationCS{běžec, běžkyně}
\dicEntry[hlaupaskór] \dicTerm{hlaupa··skó|r} \dicIPA{{\textsubring{l}}{\oe i}{\textlengthmark}{\textsubring{b}}{a}{s}{\r{g}}{ou}{\textsubring{r}}} \dicPos{m}[6] \dicFlx{(‑s, ‑r)}[52] \dicPhraseIS{hlaupaskór} \dicFlx{pl} \dicDirectTranslationCS{běžecké boty}
\dicEntry[hlaupár] \dicTerm{hlaup··ár} \dicIPA{{\textsubring{l}}{\oe i}{\textlengthmark}{\textsubring{b}}{au}{\textsubring{r}}} \dicPos{n}[2] \dicFlx{(‑s, ‑)}[5] \dicDirectTranslationCS{přestupný rok}
\dicEntry[hlaupársdagur] \dicTerm{hlaup·árs··dag|ur} \dicIPA{{\textsubring{l}}{\oe i}{\textlengthmark}{\textsubring{b}}{au}{\textsubring{r}}{s}{\textsubring{d}}{a}{\textbabygamma}{\textscy}{\textsubring{r}}} \dicPos{m}[6] \dicFlx{(‑s, ‑ar)}[62] \dicDirectTranslationCS{přestupný den}
\dicEntry[hlaupið] \dicTerm{hlaupið} \dicIPA{{\textsubring{l}}{\oe i}{\textlengthmark}{\textsubring{b}}{\textsci}{\texttheta}} \dicPos{v} \dicFlx{supin} \dicLink{hlaupa}
\dicEntry[hlaut] \dicTerm{hlaut} \dicIPA{{\textsubring{l}}{\oe i}{\textlengthmark}{\textsubring{d}}} \dicPos{v} \dicFlx{ind pf sg 1 pers} \dicLink{hljóta}
\dicEntry[hláka] \dicTerm{hlák|a} \dicIPA{{\textsubring{l}}{au}{\textlengthmark}{\r{g}}{a}} \dicPos{f}[1] \dicFlx{(‑u, ‑ur)}[19] \dicSynonym{bráðnun} \dicDirectTranslationCS{obleva, tání}
\dicEntry[hlálegur] \dicTerm{hlá··legur} \dicIPA{{\textsubring{l}}{au}{\textlengthmark}{l}{\textepsilon}{\textbabygamma}{\textscy}{\textsubring{r}}} \dicPos{adj}[1]\dicFlx{}[-8] \dicDirectTranslationCS{směšný, komický}
\dicEntry[hlána] \dicTerm{hlán|a} \dicIPA{{\textsubring{l}}{au}{\textlengthmark}{n}{a}} \dicPos{v}[1] \dicFlx{(‑aði)}[64] \dicFlx{impers} \dicPhraseIS{það\,/\addthin hann hlánar} \dicDirectTranslationCS{taje}
\dicEntry[hlátur] \dicTerm{hlát|ur} \dicsymFrequent\  \dicIPA{{\textsubring{l}}{au}{\textlengthmark}{\textsubring{d}}{\textscy}{\textsubring{r}}} \dicPos{m}[5] \dicFlx{(‑urs\,/\addthin ‑rar, ‑rar)}[9] \dicSynonym{athlægi} \dicDirectTranslationCS{smích} \dicExampleIS{Afslappaður hlátur hans feykti burt öllum áhyggjum.} \dicExampleCS{Jeho uvolněný smích zahnal všechny obavy.};  \dicPhraseIS{e‑m er ekki hlátur í hug} \dicFlx{impers} \dicDirectTranslationCS{(komu) není do smíchu};  \dicPhraseIS{reka upp hlátur} \dicLangCat{přen.} \dicDirectTranslationCS{vybuchnout smíchy};  \dicPhraseIS{veltast um af hlátri} \dicFlx{refl} \dicLangCat{přen.} \dicDirectTranslationCS{válet se smíchy}
\dicEntry[hláturmildur] \dicTerm{hlátur··mildur} \dicIPA{{\textsubring{l}}{au}{\textlengthmark}{\textsubring{d}}{\textscy}{r}{m}{\textsci}{l}{\textsubring{d}}{\textscy}{\textsubring{r}}} \dicPos{adj}[2]\dicFlx{}[-14] \dicDirectTranslationCS{usměvavý}
\dicEntry[hláturskast] \dicTerm{hláturs··|kast} \dicIPA{{\textsubring{l}}{au}{\textlengthmark}{\textsubring{d}}{\textscy}{\textsubring{r}}{s}{k\smash{\textsuperscript{h}}}{a}{s}{\textsubring{d}}} \dicPos{n}[2] \dicFlx{(‑kasts, ‑köst)}[8] \dicDirectTranslationCS{záchvat smíchu}
\dicEntry[hleð] \dicTerm{hleð} \dicIPA{{\textsubring{l}}{\textepsilon}{\textlengthmark}{\texttheta}} \dicPos{v} \dicFlx{ind praes sg 1 pers} \dicLink{hlaða\smash{\textsuperscript{2}}}
\dicEntry[hleðsla] \dicTerm{hleðsl|a} \dicIPA{{\textsubring{l}}{\textepsilon}{ð}{s}{\textsubring{d}}{l}{a}} \dicPos{f}[1] \dicFlx{(‑u, ‑ur)}[13] \textbf{1.} \dicSynonym*{veggjarhleðsla} \dicDirectTranslationCS{pokládání (zdi ap.)}  \textbf{2.} \dicSynonym*{rafmagnshleðsla} \dicDirectTranslationCS{nabíjení} \dicExampleIS{e‑að er í hleðslu} \dicExampleCS{(co) se nabíjí}  \textbf{3.} \dicSynonym*{hlaðinn veggur} \dicDirectTranslationCS{kamenná zeď}  \textbf{4.} \dicDirectTranslationCS{náboj, nábojnice}
\dicEntry[hlegið] \dicTerm{hlegið} \dicIPA{{\textsubring{l}}{ei}{\textsci}{\texttheta}} \dicPos{v} \dicFlx{supin} \dicLink{hlæja}
\dicEntry[hleifur] \dicTerm{hleif|ur} \dicIPA{{\textsubring{l}}{ei}{\textlengthmark}{v}{\textscy}{\textsubring{r}}} \dicPos{m}[6] \dicFlx{(‑s, ‑ar)}[22] \dicDirectTranslationCS{bochník, bochánek}
\dicEntry[hlekkja] \dicTerm{hlekkj|a} \dicIPA{{\textsubring{l}}{\textepsilon}{h}{\r{\textObardotlessj}}{a}} \dicPos{v}[1] \dicFlx{(‑aði)}[33] \dicFlx{acc} \dicDirectTranslationCS{spoutat\,/\addthin připoutat (řetězem)} \dicExampleIS{hlekkja hjólið sitt við grindverkið} \dicExampleCS{zamknout kolo do stojanu};  \dicIdiom{hlekkjast}[á]{ \dicPhraseIS{e‑u hlekkist á}} \dicFlx{refl impers} \dicDirectTranslationCS{(co) je poškozeno} \dicExampleIS{Skipinu hlekktist á.} \dicExampleCS{Loď byla poškozena.}
\dicEntry[hlekkur] \dicTerm{hlekk|ur} \dicIPA{{\textsubring{l}}{\textepsilon}{h}{\r{g}}{\textscy}{\textsubring{r}}} \dicPos{m}[9] \dicFlx{(‑s\,/\addthin ‑jar, ‑ir)}[26] \textbf{1.} \dicDirectTranslationCS{článek (řetězu ap.)};  \dicPhraseIS{veiki hlekkurinn í e‑u} \dicLangCat{přen.} \dicDirectTranslationCS{slabý článek (čeho)}  \textbf{2.} \dicPhraseIS{hlekkir} \dicFlx{pl} \dicSynonym*{fjötrar} \dicDirectTranslationCS{pouta}  \textbf{3.} \dicFieldCat{poč.} \dicSynonym{tengill} \dicDirectTranslationCS{(internetový) odkaz}
\dicEntry[hlemmur] \dicTerm{hlemm|ur} \dicIPA{{\textsubring{l}}{\textepsilon}{m}{\textlengthmark}{\textscy}{\textsubring{r}}} \dicPos{m}[6] \dicFlx{(‑s, ‑ar)}[4] \dicSynonym{hleri} \dicDirectTranslationCS{poklop, příklop}
\dicEntry[hler] \dicTerm{hler} \dicsymFrequent\  \dicIPA{{\textsubring{l}}{\textepsilon}{\textlengthmark}{\textsubring{r}}} \dicPos{n}[2] \dicFlx{(‑s)}[2] \dicDirectTranslationCS{odposlech, odposlouchání};  \dicPhraseIS{liggja\,/\addthin standa á hleri} \dicDirectTranslationCS{odposlouchávat, tajně poslouchat}
\dicEntry[hlera] \dicTerm{hler|a} \dicIPA{{\textsubring{l}}{\textepsilon}{\textlengthmark}{r}{a}} \dicPos{v}[1] \dicFlx{(‑aði)}[1] \dicFlx{acc} \dicDirectTranslationCS{odposlouchat, odposlouchávat} \dicExampleIS{hlera samtöl} \dicExampleCS{odposlouchávat rozhovory}
\dicEntry[hleri] \dicTerm{hler|i} \dicIPA{{\textsubring{l}}{\textepsilon}{\textlengthmark}{r}{\textsci}} \dicPos{m}[1] \dicFlx{(‑a, ‑ar)}[1] \dicDirectTranslationCS{poklop, padací dveře}
\dicEntry[hlerun] \dicTerm{hler|un} \dicIPA{{\textsubring{l}}{\textepsilon}{\textlengthmark}{r}{\textscy}{\textsubring{n}}} \dicPos{f}[7] \dicFlx{(‑unar, ‑anir)}[8] \dicDirectTranslationCS{odposlouchání, odposlech}
\dicEntry[hlessa] \dicTerm{hlessa} \dicIPA{{\textsubring{l}}{\textepsilon}{s}{\textlengthmark}{a}} \dicPos{adj}[13] \dicFlx{indecl}[2] \dicSynonym{hissa} \dicDirectTranslationCS{překvapený, zaskočený}
\dicEntry[hleyp] \dicTerm{hleyp} \dicIPA{{\textsubring{l}}{ei}{\textlengthmark}{\textsubring{b}}} \dicPos{v} \dicFlx{ind praes sg 1 pers} \dicLink{hlaupa}
\dicEntry[hleypa] \dicTerm{hleyp|a} \dicsymFrequent\  \dicIPA{{\textsubring{l}}{ei}{\textlengthmark}{\textsubring{b}}{a}} \dicPos{v}[2] \dicFlx{(‑ti, ‑t)}[66] \dicFlx{acc\,/\addthin dat} \textbf{1.} \dicFlx{dat} \dicSynonym{sleppa\smash{\textsuperscript{1}}} \dicDirectTranslationCS{(vy)pustit};  \dicPhraseIS{hleypa e‑m inn} \dicDirectTranslationCS{vpustit (do místnosti ap.)} \dicExampleIS{hleypa henni inn í húsið} \dicExampleCS{vpustit ji do domu};  \dicPhraseIS{hleypa e‑m út} \dicDirectTranslationCS{pustit (koho), nechat vyjít (koho) (z~místnosti ap.)}  \textbf{2.} \dicSynonym*{láta e‑n hlaupa} \dicDirectTranslationCS{pustit, hnát (koně ap.)} \dicExampleIS{hleypa hesti} \dicExampleCS{hnát koně}  \textbf{3.} \dicFlx{acc} \dicSynonym*{láta e‑ð þéttast} \dicDirectTranslationCS{srazit (mléko ap.), nechat ztuhnout} \dicExampleIS{hleypa skyrið} \dicExampleCS{srazit skyr}  \textbf{4.} \dicPhraseIS{hleypa í sig hörku} \dicDirectTranslationCS{vzmužit se};  \dicPhraseIS{hleypa í sig kjarki} \dicDirectTranslationCS{dodat si odvahu};  \dicIdiom{hleypa}[af]{ \dicPhraseIS{hleypa af (e‑u)}} \dicDirectTranslationCS{vystřelit (z~čeho), vypálit (z~čeho)} \dicExampleIS{hleypa af byssunni} \dicExampleCS{vystřelit z~pušky}; { \dicPhraseIS{hleypa e‑u af stokkunum}} \dicDirectTranslationCS{spustit (co), uvést (co) do provozu};  \dicIdiom{hleypa}[upp]{ \dicPhraseIS{hleypa e‑ð upp}} \dicDirectTranslationCS{vyhnat (co), vyšroubovat (co) (cenu ap.)}; { \dicPhraseIS{hleypa e‑m upp}} \dicDirectTranslationCS{vytočit (koho), rozčílit (koho)};  \dicIdiom{hleypa}[úr]{ \dicPhraseIS{hleypa úr e‑u}} \dicDirectTranslationCS{vypustit (co)}
\dicEntry[hleypidómafullur] \dicTerm{hleypi·dóma··|fullur} \dicIPA{{\textsubring{l}}{ei}{\textlengthmark}{\textsubring{b}}{\textsci}{\textsubring{d}}{ou}{m}{a}{f}{\textscy}{\textsubring{d}}{l}{\textscy}{\textsubring{r}}} \dicPos{adj}[10] \dicFlx{(comp ‑fyllri, sup ‑fyllstur)}[7] \dicDirectTranslationCS{plný předsudků}
\dicEntry[hleypidómar] \dicTerm{hleypi··dómar} \dicIPA{{\textsubring{l}}{ei}{\textlengthmark}{\textsubring{b}}{\textsci}{\textsubring{d}}{ou}{m}{a}{\textsubring{r}}} \dicPos{m}[6] \dicFlx{pl}[2] \dicDirectTranslationCS{předsudky}
\dicEntry[hlé] \dicTerm{hlé} \dicsymFrequent\  \dicIPA{{\textsubring{l}}{j}{\textepsilon}{\textlengthmark}} \dicPos{n}[2] \dicFlx{(‑s, ‑)}[5] \textbf{1.} \dicSynonym{skjól} \dicDirectTranslationCS{úkryt, přístřeší} \dicExampleIS{hlé fyrir stormi} \dicExampleCS{úkryt před bouřkou}  \textbf{2.} \dicSynonym{hvíld} \dicDirectTranslationCS{pauza, přestávka} \dicExampleIS{stutt hlé} \dicExampleCS{krátká přestávka}  \textbf{3.} \dicSynonym{stans} \dicDirectTranslationCS{přerušení, přestávka};  \dicPhraseIS{gera hlé í e‑u} \dicDirectTranslationCS{udělat přestávku v~(čem)};  \dicPhraseIS{draga sig í hlé} \dicLangCat{přen.} \dicDirectTranslationCS{stáhnout se do ústraní};  \dicPhraseIS{halda sér til hlés} \dicLangCat{přen.} \dicDirectTranslationCS{držet se zpátky}
\dicEntry[hlébarði] \dicTerm{hlé··barð|i} \dicIPA{{\textsubring{l}}{j}{\textepsilon}{\textlengthmark}{\textsubring{b}}{a}{r}{ð}{\textsci}} \dicPos{m}[1] \dicFlx{(‑a, ‑ar)}[8] \dicFieldCat{zool.} \dicDirectTranslationCS{levhart, levhart skvrnitý} \textit{(l.~{\textLA{Panthera pardus}})}  \dicsymPhoto\ 
\dicFigure{843.jpg}{Hlébarði}{Hlébarði - Zicha Ondřej, Biolib, Copyright/CC-BY-NC}
\dicEntry[hléborði] \dicTerm{hlé··borð|i} \dicIPA{{\textsubring{l}}{j}{\textepsilon}{\textlengthmark}{\textsubring{b}}{\textopeno}{r}{ð}{\textsci}} \dicPos{m}[1] \dicFlx{(‑a)}[3] \dicFieldCat{nám.} \dicDirectTranslationCS{závětrná strana} \dicAntonym{kulborði}
\dicEntry[hlédrægni] \dicTerm{hlé··drægn|i} \dicIPA{{\textsubring{l}}{j}{\textepsilon}{\textlengthmark}{\textsubring{d}}{r}{a}{i}{\r{g}}{n}{\textsci}} \dicPos{f}[3] \dicFlx{(‑i)}[3] \dicDirectTranslationCS{zdrženlivost, rezervovanost}
\dicEntry[hlédrægur] \dicTerm{hlé··drægur} \dicIPA{{\textsubring{l}}{j}{\textepsilon}{\textlengthmark}{\textsubring{d}}{r}{a}{i}{\textbabygamma}{\textscy}{\textsubring{r}}} \dicPos{adj}[1]\dicFlx{}[-1] \dicDirectTranslationCS{zdrženlivý, rezervovaný}
\dicEntry[hlið] \dicTerm{hlið\smash{\textsuperscript{1}}} \dicsymFrequent\  \dicIPA{{\textsubring{l}}{\textsci}{\textlengthmark}{\texttheta}} \dicPos{f}[4] \dicFlx{(‑ar, ‑ar)}[1] \textbf{1.} \dicDirectTranslationCS{strana (pravá ap.), bok} \dicExampleIS{hliðin á húsinu} \dicExampleCS{strana domu};  \dicPhraseIS{hlið við hlið} \dicFlx{adv} \dicDirectTranslationCS{bok po boku};  \dicPhraseIS{leggja til hliðar} \dicDirectTranslationCS{odkládat (peníze ap.), dávat stranou};  \dicPhraseIS{setja e‑ð til hliðar} \dicDirectTranslationCS{dát (co) na stranu};  \dicPhraseIS{sterk hlið} \dicLangCat{přen.} \dicDirectTranslationCS{silná stránka};  \dicPhraseIS{við hliðina á e‑u} \dicFlx{prep} \dicDirectTranslationCS{vedle (čeho)}  \textbf{2.} \dicDirectTranslationCS{úhel pohledu, hledisko} \dicExampleIS{Þetta er ný hlið á málinu.} \dicExampleCS{To je nový pohled na věc.}
\dicEntry[hlið] \dicTerm{hlið\smash{\textsuperscript{2}}} \dicsymFrequent\  \dicIPA{{\textsubring{l}}{\textsci}{\textlengthmark}{\texttheta}} \dicPos{n}[2] \dicFlx{(‑s, ‑)}[5] \dicDirectTranslationCS{brána, branka} \dicExampleIS{hlið að garðinum} \dicExampleCS{branka do zahrady}
\dicEntry[hliðargata] \dicTerm{hliðar··|gata} \dicIPA{{\textsubring{l}}{\textsci}{ð}{a}{r}{\r{g}}{a}{\textsubring{d}}{a}} \dicPos{f}[1] \dicFlx{(‑götu, ‑götur)}[14] \dicDirectTranslationCS{postranní\,/\addthin boční ulice}
\dicEntry[hliðarhljóð] \dicTerm{hliðar··hljóð} \dicIPA{{\textsubring{l}}{\textsci}{\textlengthmark}{ð}{a}{\textsubring{r}}{\textsubring{l}}{j}{ou}{\texttheta}} \dicPos{n}[2] \dicFlx{(‑s, ‑)}[5] \dicFieldCat{jaz.} \dicDirectTranslationCS{boková souhláska, laterála}
\dicEntry[hliðarlína] \dicTerm{hliðar··lín|a} \dicIPA{{\textsubring{l}}{\textsci}{\textlengthmark}{ð}{a}{r}{l}{i}{n}{a}} \dicPos{f}[1] \dicFlx{(‑u, ‑ur)}[7] \dicFieldCat{sport.} \dicDirectTranslationCS{postranní\,/\addthin pomezní čára}
\begin{xtolerant}{}{1pt}
\dicEntry[hliðarráðstöfun] \dicTerm{hliðar··ráð·|stöfun}\addthinSS\dicIPA{{\textsubring{l}}{\textsci}{\textlengthmark}{ð}{a}{r}{au}{ð}{s}{\textsubring{d}}{\oe}{v}{\textscy}{\textsubring{n}}}\addthinS\dicPos{f}[7]\addthinS\dicFlx{(‑stöfunar,\addthinS‑stafanir)}[11] \dicDirectTranslationCS{náhradní opatření}
\end{xtolerant}
\dicEntry[hliðarvindur] \dicTerm{hliðar··vind|ur} \dicIPA{{\textsubring{l}}{\textsci}{ð}{a}{r}{v}{\textsci}{n}{\textsubring{d}}{\textscy}{\textsubring{r}}} \dicPos{m}[6] \dicFlx{(‑s, ‑ar)}[4] \dicFieldCat{nám.} \dicDirectTranslationCS{boční vítr}
\dicEntry[hliðhollur] \dicTerm{hlið··hollur} \dicIPA{{\textsubring{l}}{\textsci}{\textlengthmark}{\texttheta}{h}{\textopeno}{\textsubring{d}}{l}{\textscy}{\textsubring{r}}} \dicPos{adj}[1]\dicFlx{}[-1] \dicDirectTranslationCS{přátelský, přátelsky nakloněný};  \dicPhraseIS{vera hliðhollur e‑m} \dicDirectTranslationCS{být přátelsky nakloněný (komu)}
\dicEntry[hliðmæltur] \dicTerm{hlið··mæltur} \dicIPA{{\textsubring{l}}{\textsci}{ð}{m}{a}{i}{\textsubring{l}}{\textsubring{d}}{\textscy}{\textsubring{r}}} \dicPos{adj}[1]\dicFlx{}[-13] \dicFieldCat{jaz.} \dicDirectTranslationCS{bokový, laterální}
\dicEntry[hliðra] \dicTerm{hliðr|a} \dicIPA{{\textsubring{l}}{\textsci}{ð}{r}{a}} \dicPos{v}[1] \dicFlx{(‑aði)}[1] \dicFlx{dat} \dicIdiom{hliðra}[hjá]{ \dicPhraseIS{hliðra sér hjá (því) að (gera e‑ð)}} \dicDirectTranslationCS{vyhýbat\,/\addthin stranit se (dělání (čeho))};  \dicIdiom{hliðra}[til]{ \dicPhraseIS{hliðra til fyrir e‑m}} \dicDirectTranslationCS{udělat (komu) místo};  \dicIdiom{hliðrast}{ \dicPhraseIS{hliðrast}} \dicFlx{refl} \dicDirectTranslationCS{posunout se, přesunout se}
\dicEntry[hliðrun] \dicTerm{hliðr|un} \dicIPA{{\textsubring{l}}{\textsci}{ð}{r}{\textscy}{\textsubring{n}}} \dicPos{f}[7] \dicFlx{(‑unar)}[9] \dicSynonym{tilfærsla} \dicDirectTranslationCS{posun, posunutí, přesun}
\dicEntry[hliðsjón] \dicTerm{hlið··sjón} \dicIPA{{\textsubring{l}}{\textsci}{ð}{s}{j}{ou}{\textsubring{n}}} \dicPos{f}[7] \dicFlx{(‑ar)}[3] \dicDirectTranslationCS{přihlédnutí};  \dicPhraseIS{hafa hliðsjón af e‑u} \dicDirectTranslationCS{přihlédnout k~(čemu), vzít v~úvahu (co)};  \dicPhraseIS{hafa e‑ð til hliðsjónar} \dicDirectTranslationCS{mít (co) k~přihlédnutí};  \dicPhraseIS{með hliðsjón af e‑u} \dicDirectTranslationCS{s~přihlédnutím k~(čemu)}
\dicEntry[hliðskipaður] \dicTerm{hlið··skip·|aður} \dicIPA{{\textsubring{l}}{\textsci}{ð}{s}{\r{\textObardotlessj}}{\textsci}{\textsubring{b}}{a}{ð}{\textscy}{\textsubring{r}}} \dicPos{adj}[3] \dicFlx{(f ‑uð)}[4] \dicFieldCat{jaz.} \dicDirectTranslationCS{souřadný}
\dicEntry[hliðskipun] \dicTerm{hlið··skip|un} \dicIPA{{\textsubring{l}}{\textsci}{ð}{s}{\r{\textObardotlessj}}{\textsci}{\textsubring{b}}{\textscy}{\textsubring{n}}} \dicPos{f}[7] \dicFlx{(‑unar)}[9] \dicFieldCat{jaz.} \dicDirectTranslationCS{souřadnost, parataxe} \dicAntonym{undirskipun}
\dicEntry[hliðstæða] \dicTerm{hlið··stæð|a} \dicIPA{{\textsubring{l}}{\textsci}{ð}{s}{\textsubring{d}}{a}{i}{ð}{a}} \dicPos{f}[1] \dicFlx{(‑u, ‑ur)}[13] \dicDirectTranslationCS{obdoba, analogie, paralela, podobnost}
\dicEntry[hliðstæður] \dicTerm{hlið··stæður} \dicIPA{{\textsubring{l}}{\textsci}{ð}{s}{\textsubring{d}}{a}{i}{ð}{\textscy}{\textsubring{r}}} \dicPos{adj}[2]\dicFlx{}[-6] \textbf{1.} \dicSynonym{sambærilegur} \dicDirectTranslationCS{obdobný, analogický}  \textbf{2.} \dicFieldCat{jaz.} \dicDirectTranslationCS{atributivní} \dicIndirectTranslationCS{(o~přídavných jménech, zájmenech a~číslovkách)} \dicAntonym{sérstæður}
\dicEntry[hlíð] \dicTerm{hlíð} \dicsymFrequent\  \dicIPA{{\textsubring{l}}{i}{\textlengthmark}{\texttheta}} \dicPos{f}[4] \dicFlx{(‑ar, ‑ar\,/\addthin ‑ir)}[30] \dicDirectTranslationCS{svah, stráň} \dicExampleIS{ganga niður hlíðina} \dicExampleCS{jít ze svahu}
\dicEntry[hlíf] \dicTerm{hlíf} \dicIPA{{\textsubring{l}}{i}{\textlengthmark}{f}} \dicPos{f}[4] \dicFlx{(‑ar, ‑ar)}[1] \dicSynonym{vernd} \dicDirectTranslationCS{ochrana} \dicIndirectTranslationCS{(ochranné zařízení\,/\addthin opatření ap.)}
\dicEntry[hlífa] \dicTerm{hlíf|a} \dicIPA{{\textsubring{l}}{i}{\textlengthmark}{v}{a}} \dicPos{v}[2] \dicFlx{(‑ði, ‑t)}[108] \dicFlx{dat} \textbf{1.} \dicSynonym{vægja} \dicDirectTranslationCS{(u)šetřit, mít ohled};  \dicPhraseIS{hlífa e‑m við e‑u} \dicDirectTranslationCS{ušetřit (koho čeho)};  \dicPhraseIS{hlífa sér} \dicDirectTranslationCS{šetřit se}  \textbf{2.} \dicSynonym{vernda} \dicDirectTranslationCS{(o)chránit, ochraňovat} \dicExampleIS{hlífa sér fyrir sterkri sól} \dicExampleCS{chránit se před silným sluncem};  \dicIdiom{hlífast}[við]{ \dicPhraseIS{hlífast við að (gera e‑ð)}} \dicFlx{refl} \dicDirectTranslationCS{snažit se vyhnout (dělání (čeho)), vyhýbat se (dělání (čeho))}
\dicEntry[hlífðarlaus] \dicTerm{hlífðar··laus} \dicIPA{{\textsubring{l}}{i}{v}{ð}{a}{r}{l}{\oe i}{s}} \dicPos{adj}[5]\dicFlx{}[-1] \dicDirectTranslationCS{nemilosrdný, nelítostný} \dicExampleIS{hlífðarlausir stormar} \dicExampleCS{nelítostné bouře}
\dicEntry[hlífiskjöldur] \dicTerm{hlífi··|skjöldur} \dicIPA{{\textsubring{l}}{i}{\textlengthmark}{v}{\textsci}{s}{\r{\textObardotlessj}}{\oe}{l}{\textsubring{d}}{\textscy}{\textsubring{r}}} \dicPos{m}[12] \dicFlx{(‑skjaldar)}[4] \dicPhraseIS{halda hlífiskildi yfir e‑m} \dicDirectTranslationCS{držet nad (kým) ochrannou ruku, chránit (koho)}
\dicEntry[hlít] \dicTerm{hlít} \dicIPA{{\textsubring{l}}{i}{\textlengthmark}{\textsubring{d}}} \dicPos{f}[4] \dicFlx{(‑ar)}[3] \dicPhraseIS{til hlítar} \dicFlx{adv} \dicSynonym{alveg} \dicDirectTranslationCS{zcela, pořádně}
\dicEntry[hlíta] \dicTerm{hlít|a} \dicIPA{{\textsubring{l}}{i}{\textlengthmark}{\textsubring{d}}{a}} \dicPos{v}[2] \dicFlx{(‑ti, ‑t)}[62] \dicFlx{dat} \dicDirectTranslationCS{řídit se, dodržovat} \dicExampleIS{hlíta reglum} \dicExampleCS{řídit se pravidly}
\dicEntry[hljóð] \dicTerm{hljóð} \dicsymFrequent\  \dicIPA{{\textsubring{l}}{j}{ou}{\textlengthmark}{\texttheta}} \dicPos{n}[2] \dicFlx{(‑s, ‑)}[5] \textbf{1.} \dicDirectTranslationCS{zvuk}  \textbf{2.} \dicSynonym{rödd} \dicDirectTranslationCS{hlas};  \dicPhraseIS{tala í hálfum hljóðum} \dicDirectTranslationCS{mluvit polohlasem};  \dicPhraseIS{í einu hljóði} \dicFlx{adv} \dicSynonym{einróma\smash{\textsuperscript{1}}} \dicDirectTranslationCS{jednohlasně, jedním hlasem}  \textbf{3.} \dicSynonym{óp} \dicDirectTranslationCS{výkřik, zakvílení};  \dicPhraseIS{reka upp hljóð} \dicDirectTranslationCS{vykřiknout, dát se do křiku}  \textbf{4.} \dicFieldCat{jaz.} \dicSynonym{málhljóð} \dicDirectTranslationCS{hláska, foném};  \dicIdiom{hljóð}{ \dicPhraseIS{biðja\,/\addthin kveðja sér hljóðs}} \dicDirectTranslationCS{požádat o~slovo}; { \dicPhraseIS{gefa hljóð}} \dicSynonym{þagna} \dicDirectTranslationCS{ztichnout, umlknout};  \dicPhraseIS{fá hljóð} \dicDirectTranslationCS{dostat slovo (při diskuzi ap.)};  \dicPhraseIS{í heyranda hljóði} \dicFlx{adv} \dicDirectTranslationCS{veřejně, nahlas};  \dicPhraseIS{í þegjanda hljóði} \dicFlx{adv} \dicDirectTranslationCS{potichu};  \dicPhraseIS{það er dauft í e‑m hljóðið} \dicDirectTranslationCS{(kdo) je přešlý\,/\addthin přepadlý};  \dicPhraseIS{þegja þunnu hljóði} \dicLangCat{přen.} \dicDirectTranslationCS{držet jazyk za zuby}
\dicEntry[hljóða] \dicTerm{hljóð|a} \dicsymFrequent\  \dicIPA{{\textsubring{l}}{j}{ou}{\textlengthmark}{ð}{a}} \dicPos{v}[1] \dicFlx{(‑aði)}[1] \textbf{1.} \dicSynonym{hljóma} \dicDirectTranslationCS{znít (věta ap.)} \dicIndirectTranslationCS{(mít slovní podobu)}  \textbf{2.} \dicSynonym{æpa} \dicDirectTranslationCS{vykřiknout, zavýt} \dicExampleIS{hljóða af sársauka} \dicExampleCS{vykřiknout bolestí}
\dicEntry[hljóðblöndun] \dicTerm{hljóð··blönd|un} \dicIPA{{\textsubring{l}}{j}{ou}{ð}{\textsubring{b}}{l}{\oe}{n}{\textsubring{d}}{\textscy}{\textsubring{n}}} \dicPos{f}[7] \dicFlx{(‑unar)}[12] \dicFieldCat{hud.} \dicDirectTranslationCS{mixování (zvuku)}
\dicEntry[hljóðbók] \dicTerm{hljóð··|bók} \dicIPA{{\textsubring{l}}{j}{ou}{ð}{\textsubring{b}}{ou}{\r{g}}} \dicPos{f}[8] \dicFlx{(‑bókar, ‑bækur)}[5] \dicDirectTranslationCS{audiokniha, audioknížka}
\dicEntry[hljóðbylgja] \dicTerm{hljóð··bylgj|a} \dicIPA{{\textsubring{l}}{j}{ou}{ð}{\textsubring{b}}{\textsci}{l}{\r{\textObardotlessj}}{a}} \dicPos{f}[1] \dicFlx{(‑u, ‑ur)}[25] \dicFieldCat{fyz.} \dicDirectTranslationCS{zvuková vlna}
\dicEntry[hljóðbær] \dicTerm{hljóð··bær} \dicIPA{{\textsubring{l}}{j}{ou}{ð}{\textsubring{b}}{a}{i}{\textsubring{r}}} \dicPos{adj}[5] \dicFlx{(f ‑)}[8] \dicPhraseIS{það er hljóðbært í e‑u} \dicDirectTranslationCS{je dobře slyšet v~(čem) (v~domě ap.)}
\begin{xtolerant}{}{1pt}
\dicEntry[hljóðeðlisfræði] \dicTerm{hljóð··eðlis·fræð|i}\addthinS\dicIPA{{\textsubring{l}}{j}{ou}{ð}{\textepsilon}{ð}{l}{\textsci}{s}{f}{r}{a}{i}{ð}{\textsci}}\addthinS\dicPos{f}[3]\addthinS\dicFlx{(‑i)}[3] \dicDirectTranslationCS{akustika} \dicIndirectTranslationCS{(nauka o~zvuku)}
\end{xtolerant}
\dicEntry[hljóðeinangraður] \dicTerm{hljóð··ein·angr·|aður} \dicIPA{{\textsubring{l}}\-{j}\-{ou}\-{ð}\-{ei}\-{n}\-{a}\-{n}\-{\r{g}}\-{r}\-{a}\-{ð}\-{\textscy}\-{\textsubring{r}}\-} \dicPos{adj}[3] \dicFlx{(f ‑uð)}[4] \dicDirectTranslationCS{zvukotěsný} \dicExampleIS{hljóðeinangrað herbergi} \dicExampleCS{zvukotěsná místnost}
\dicEntry[hljóðeinangrandi] \dicTerm{hljóð··ein·angr·andi} \dicIPA{{\textsubring{l}}{j}{ou}{ð}{ei}{n}{a}{n}{\r{g}}{r}{a}{n}{\textsubring{d}}{\textsci}} \dicPos{adj}[13] \dicFlx{indecl}[1] \dicDirectTranslationCS{zvukotěsný}
\dicEntry[hljóðeinangrun] \dicTerm{hljóð··ein·angr|un} \dicIPA{{\textsubring{l}}{j}{ou}{ð}{ei}{n}{au}{\ng}{\r{g}}{r}{\textscy}{\textsubring{n}}} \dicPos{f}[7] \dicFlx{(‑unar)}[9] \dicDirectTranslationCS{zvukotěsnost}
\dicEntry[hljóðfræði] \dicTerm{hljóð··fræð|i} \dicIPA{{\textsubring{l}}{j}{ou}{\texttheta}{f}{r}{a}{i}{ð}{\textsci}} \dicPos{f}[3] \dicFlx{(‑i)}[3] \dicFieldCat{jaz.} \dicDirectTranslationCS{fonetika}
\dicEntry[hljóðfræðingur] \dicTerm{hljóð·fræð··ing|ur} \dicIPA{{\textsubring{l}}{j}{ou}{\texttheta}{f}{r}{a}{i}{ð}{i}{\ng}{\r{g}}{\textscy}{\textsubring{r}}} \dicPos{m}[6] \dicFlx{(‑s, ‑ar)}[8] \dicDirectTranslationCS{fonetik, fonetička}
\dicEntry[hljóðfæraleikari] \dicTerm{hljóð·færa··leik·ar|i} \dicIPA{{\textsubring{l}}\-{j}\-{ou}\-{\texttheta}\-{f}\-{a}\-{i}\-{r}\-{a}\-{l}\-{ei}\-{\r{g}}\-{a}\-{r}\-{\textsci}\-} \dicPos{m}[1] \dicFlx{(‑a, ‑ar)}[13] \dicDirectTranslationCS{hráč(ka) na hudební nástroj}
\dicEntry[hljóðfæri] \dicTerm{hljóð··færi} \dicsymFrequent\  \dicIPA{{\textsubring{l}}{j}{ou}{\texttheta}{f}{a}{i}{r}{\textsci}} \dicPos{n}[2] \dicFlx{(‑s, ‑)}[14] \dicDirectTranslationCS{hudební nástroj} \dicExampleIS{læra á hljóðfæri} \dicExampleCS{učit se hrát na hudební nástroj}
\dicEntry[hljóðgervill] \dicTerm{hljóð··gerv|ill} \dicIPA{{\textsubring{l}}{j}{ou}{ð}{\r{\textObardotlessj}}{\textepsilon}{r}{v}{\textsci}{\textsubring{d}}{\textsubring{l}}} \dicPos{m}[6] \dicFlx{(‑ils, ‑lar)}[35] \dicFieldCat{hud.} \dicDirectTranslationCS{syntezátor, syntetizátor}
\dicEntry[hljóðgervingur] \dicTerm{hljóð··gerv·ing|ur} \dicIPA{{\textsubring{l}}{j}{ou}{ð}{\r{\textObardotlessj}}{\textepsilon}{r}{v}{i}{\ng}{\r{g}}{\textscy}{\textsubring{r}}} \dicPos{m}[6] \dicFlx{(‑s, ‑ar)}[8] \dicFieldCat{lit.\,/\addthin jaz.} \dicDirectTranslationCS{zvukomalba, onomatopoie}
\dicEntry[hljóðhimna] \dicTerm{hljóð··himn|a} \dicIPA{{\textsubring{l}}{j}{ou}{\textlengthmark}{\texttheta}{h}{\textsci}{m}{n}{a}} \dicPos{f}[1] \dicFlx{(‑u, ‑ur)}[7] \dicFieldCat{anat.} \dicDirectTranslationCS{(ušní) bubínek}
\dicEntry[hljóðkerfi] \dicTerm{hljóð··kerfi} \dicIPA{{\textsubring{l}}{j}{ou}{ð}{c\smash{\textsuperscript{h}}}{\textepsilon}{r}{v}{\textsci}} \dicPos{n}[2] \dicFlx{(‑s, ‑)}[14] \textbf{1.} \dicFieldCat{jaz.} \dicDirectTranslationCS{fonologie, fonologický systém}  \textbf{2.} \dicDirectTranslationCS{zvukový systém}
\dicEntry[hljóðkerfisfræði] \dicTerm{hljóð·kerfis··fræð|i} \dicIPA{{\textsubring{l}}{j}{ou}{ð}{c\smash{\textsuperscript{h}}}{\textepsilon}{r}{v}{\textsci}{s}{f}{r}{a}{i}{ð}{\textsci}} \dicPos{f}[3] \dicFlx{(‑i)}[3] \dicFieldCat{jaz.} \dicDirectTranslationCS{fonologie}
\dicEntry[hljóðkort] \dicTerm{hljóð··kort} \dicIPA{{\textsubring{l}}{j}{ou}{ð}{k\smash{\textsuperscript{h}}}{\textopeno}{\textsubring{r}}{\textsubring{d}}} \dicPos{n}[2] \dicFlx{(‑s, ‑)}[5] \dicFieldCat{poč.} \dicDirectTranslationCS{zvuková karta}
\dicEntry[hljóðkútur] \dicTerm{hljóð··kút|ur} \dicIPA{{\textsubring{l}}{j}{ou}{ð}{k\smash{\textsuperscript{h}}}{u}{\textsubring{d}}{\textscy}{\textsubring{r}}} \dicPos{m}[6] \dicFlx{(‑s, ‑ar)}[22] \dicFieldCat{techn.} \dicDirectTranslationCS{tlumič}
\dicEntry[hljóðlaus] \dicTerm{hljóð··laus} \dicIPA{{\textsubring{l}}{j}{ou}{ð}{l}{\oe i}{s}} \dicPos{adj}[5]\dicFlx{}[-1] \dicDirectTranslationCS{tichý, nehlučný}
\dicEntry[hljóðlátur] \dicTerm{hljóð··látur} \dicsymFrequent\  \dicIPA{{\textsubring{l}}{j}{ou}{ð}{l}{au}{\textsubring{d}}{\textscy}{\textsubring{r}}} \dicPos{adj}[1]\dicFlx{}[-1] \dicDirectTranslationCS{tichý, nehlučný} \dicExampleIS{hljóðlátt umhverfi} \dicExampleCS{tiché okolí}
\dicEntry[hljóðlega] \dicTerm{hljóð··lega} \dicsymFrequent\  \dicIPA{{\textsubring{l}}{j}{ou}{ð}{l}{\textepsilon}{\textbabygamma}{a}} \dicPos{adv} \dicDirectTranslationCS{tiše, nehlučně} \dicExampleIS{læðast hljóðlega að dyrunum} \dicExampleCS{tiše se krást ke dveřím}
\dicEntry[hljóðmerki] \dicTerm{hljóð··merki} \dicIPA{{\textsubring{l}}{j}{ou}{ð}{m}{\textepsilon}{\textsubring{r}}{\r{\textObardotlessj}}{\textsci}} \dicPos{n}[2] \dicFlx{(‑s, ‑)}[16] \dicDirectTranslationCS{zvukový signál}
\dicEntry[hljóðmúr] \dicTerm{hljóð··múr} \dicIPA{{\textsubring{l}}{j}{ou}{ð}{m}{u}{\textsubring{r}}} \dicPos{m}[4] \dicFlx{(‑s)}[18] \dicFieldCat{let.} \dicDirectTranslationCS{zvuková bariéra}
\dicEntry[hljóðna] \dicTerm{hljóðn|a} \dicIPA{{\textsubring{l}}{j}{ou}{ð}{n}{a}} \dicPos{v}[1] \dicFlx{(‑aði)}[7] \dicDirectTranslationCS{(u)tichnout, utichat, ztišit se, ztichnout} \dicExampleIS{Raddirnar hljóðnuðu.} \dicExampleCS{Hlasy utichly.};  \dicPhraseIS{það hljóðnar e‑s staðar} \dicFlx{impers} \dicDirectTranslationCS{(kde) to utichá, (kde) začíná být ticho}
\dicEntry[hljóðnemi] \dicTerm{hljóð··nem|i} \dicIPA{{\textsubring{l}}{j}{ou}{ð}{n}{\textepsilon}{m}{\textsci}} \dicPos{m}[1] \dicFlx{(‑a, ‑ar)}[1] \dicDirectTranslationCS{mikrofon}
\dicEntry[hljóðregla] \dicTerm{hljóð··regl|a} \dicIPA{{\textsubring{l}}{j}{ou}{ð}{r}{\textepsilon}{\r{g}}{l}{a}} \dicPos{f}[1] \dicFlx{(‑u, ‑ur)}[19] \dicFieldCat{jaz.} \dicDirectTranslationCS{fonetické pravidlo}
\dicEntry[hljóðrita] \dicTerm{hljóð··rit|a} \dicIPA{{\textsubring{l}}{j}{ou}{ð}{r}{\textsci}{\textsubring{d}}{a}} \dicPos{v}[1] \dicFlx{(‑aði)}[1] \dicFlx{acc} \textbf{1.} \dicDirectTranslationCS{nahrát, zaznamenat (zvukově)} \dicExampleIS{hljóðrita tónleikana} \dicExampleCS{zaznamenat koncert}  \textbf{2.} \dicFieldCat{jaz.} \dicDirectTranslationCS{zapsat\,/\addthin přepsat foneticky} \dicExampleIS{hljóðrita orð} \dicExampleCS{zapsat slovo foneticky}
\dicEntry[hljóðritun] \dicTerm{hljóð··rit|un} \dicIPA{{\textsubring{l}}{j}{ou}{ð}{r}{\textsci}{\textsubring{d}}{\textscy}{\textsubring{n}}} \dicPos{f}[7] \dicFlx{(‑unar)}[9] \textbf{1.} \dicDirectTranslationCS{nahrávka, (zvukový) záznam}  \textbf{2.} \dicFieldCat{jaz.} \dicDirectTranslationCS{fonetický zápis\,/\addthin přepis}
\dicEntry[hljóðritunartákn] \dicTerm{hljóð·ritunar··tákn} \dicIPA{{\textsubring{l}}{j}{ou}{ð}{\textsubring{r}}{\textsci}{\textsubring{d}}{\textscy}{n}{a}{\textsubring{r}}{t\smash{\textsuperscript{h}}}{au}{h}{\r{g}}{\textsubring{n}}} \dicPos{n}[2] \dicFlx{(‑s, ‑)}[5] \dicFieldCat{jaz.} \dicDirectTranslationCS{fonetický symbol}
\dicEntry[hljóðsjá] \dicTerm{hljóð··sjá} \dicIPA{{\textsubring{l}}{j}{ou}{ð}{s}{j}{au}} \dicPos{f}[4] \dicFlx{(‑r, ‑r)}[18] \dicFieldCat{techn.} \dicDirectTranslationCS{sonar}
\dicEntry[hljóðskipti] \dicTerm{hljóð··skipti} \dicIPA{{\textsubring{l}}{j}{ou}{ð}{s}{\r{\textObardotlessj}}{\textsci}{f}{\textsubring{d}}{\textsci}} \dicPos{n}[2] \dicFlx{(‑s, ‑)}[14] \dicFieldCat{jaz.} \dicDirectTranslationCS{ablaut, střída}
\dicEntry[hljóðstafur] \dicTerm{hljóð··staf|ur} \dicIPA{{\textsubring{l}}{j}{ou}{ð}{s}{\textsubring{d}}{a}{v}{\textscy}{\textsubring{r}}} \dicPos{m}[9] \dicFlx{(‑s, ‑ir)}[9] \dicPhraseIS{e‑að stendur í hljóðstöfum} \dicFieldCat{lit.} \dicDirectTranslationCS{(co) aliteruje, (co) obsahuje aliteraci}
\dicEntry[hljóðtákn] \dicTerm{hljóð··tákn} \dicIPA{{\textsubring{l}}{j}{ou}{ð}{t\smash{\textsuperscript{h}}}{au}{h}{\r{g}}{\textsubring{n}}} \dicPos{n}[2] \dicFlx{(‑s, ‑)}[5] \textbf{1.} \dicFieldCat{jaz.} \dicSynonym{bókstafur} \dicDirectTranslationCS{grafický\,/\addthin písemný znak, grafém}  \textbf{2.} \dicFieldCat{jaz.} \dicSynonym{hljóðritunartákn} \dicDirectTranslationCS{fonetický symbol}
\dicEntry[hljóðupptaka] \dicTerm{hljóð··upp·|taka} \dicIPA{{\textsubring{l}}{j}{ou}{ð}{\textscy}{h}{\textsubring{b}}{t\smash{\textsuperscript{h}}}{a}{\r{g}}{a}} \dicPos{f}[1] \dicFlx{(‑töku, ‑tökur)}[20] \dicDirectTranslationCS{nahrávka, záznam zvuku}
\dicEntry[hljóður] \dicTerm{hljóður} \dicsymFrequent\  \dicIPA{{\textsubring{l}}{j}{ou}{\textlengthmark}{ð}{\textscy}{\textsubring{r}}} \dicPos{adj}[2]\dicFlx{}[-6] \dicDirectTranslationCS{tichý, klidný} \dicExampleIS{Það er svo hljótt hérna.} \dicExampleCS{Je tady takové ticho.}
\dicEntry[hljóðvarp] \dicTerm{hljóð··|varp} \dicIPA{{\textsubring{l}}{j}{ou}{ð}{v}{a}{\textsubring{r}}{\textsubring{b}}} \dicPos{n}[2] \dicFlx{(‑varps, ‑vörp)}[8] \textbf{1.} \dicFieldCat{jaz.} \dicDirectTranslationCS{přehláska}  \textbf{2.} \dicSynonym{útvarp} \dicDirectTranslationCS{rádio, rozhlas} \dicIndirectTranslationCS{(instituce, vysílání)}
\dicEntry[hljóðver] \dicTerm{hljóð··ver} \dicIPA{{\textsubring{l}}{j}{ou}{ð}{v}{\textepsilon}{\textsubring{r}}} \dicPos{n}[2] \dicFlx{(‑s, ‑)}[5] \dicDirectTranslationCS{nahrávací studio}
\dicEntry[hljóma] \dicTerm{hljóm|a} \dicsymFrequent\  \dicIPA{{\textsubring{l}}{j}{ou}{\textlengthmark}{m}{a}} \dicPos{v}[1] \dicFlx{(‑aði)}[1] \textbf{1.} \dicSynonym{gjalla} \dicDirectTranslationCS{(za)znít} \dicExampleIS{Lagið hljómar vel.} \dicExampleCS{Písnička zní dobře.}  \textbf{2.} \dicDirectTranslationCS{znít} \dicIndirectTranslationCS{(působit nějakým dojmem)} \dicExampleIS{Það hljómar vel.} \dicExampleCS{To zní dobře.}
\dicEntry[hljómblær] \dicTerm{hljóm··blæ|r} \dicIPA{{\textsubring{l}}{j}{ou}{m}{\textsubring{b}}{l}{a}{i}{\textsubring{r}}} \dicPos{m}[9] \dicFlx{(‑jar\,/\addthin ‑s)}[30] \dicDirectTranslationCS{zabarvení (hlasu ap.), témbr}
\dicEntry[hljómborð] \dicTerm{hljóm··borð} \dicIPA{{\textsubring{l}}{j}{ou}{m}{\textsubring{b}}{\textopeno}{r}{\texttheta}} \dicPos{n}[2] \dicFlx{(‑s, ‑)}[5] \textbf{1.} \dicDirectTranslationCS{klaviatura, klávesnice}  \textbf{2.} \dicFieldCat{hud.} \dicDirectTranslationCS{klávesy}
\dicEntry[hljómbotn] \dicTerm{hljóm··botn} \dicIPA{{\textsubring{l}}{j}{ou}{m}{\textsubring{b}}{\textopeno}{h}{\textsubring{d}}{\textsubring{n}}} \dicPos{m}[4] \dicFlx{(‑s, ‑ar)}[3] \dicFieldCat{hud.} \dicDirectTranslationCS{rezonance, ozvuk}
\dicEntry[hljómburður] \dicTerm{hljóm··burð|ur} \dicIPA{{\textsubring{l}}{j}{ou}{m}{\textsubring{b}}{\textscy}{r}{ð}{\textscy}{\textsubring{r}}} \dicPos{m}[10] \dicFlx{(‑ar)}[7] \dicDirectTranslationCS{akustika, akustičnost} \dicIndirectTranslationCS{(zvukové vlastnosti nějakého prostoru)}
\dicEntry[hljómfall] \dicTerm{hljóm··fall} \dicIPA{{\textsubring{l}}{j}{ou}{m}{f}{a}{\textsubring{d}}{\textsubring{l}}} \dicPos{n}[2] \dicFlx{(‑s)}[2] \dicFieldCat{lit.\,/\addthin jaz.} \dicSynonym{tónfall} \dicDirectTranslationCS{rytmus}
\dicEntry[hljómgrunnur] \dicTerm{hljóm··grunn|ur} \dicIPA{{\textsubring{l}}{j}{ou}{m}{\r{g}}{r}{\textscy}{n}{\textscy}{\textsubring{r}}} \dicPos{m}[6] \dicFlx{(‑s)}[7] \dicPhraseIS{e‑að fær (góðan) hljómgrunn} \dicDirectTranslationCS{(co) má (dobrý) ohlas\,/\addthin odezvu}
\dicEntry[hljómkviða] \dicTerm{hljóm··kvið|a} \dicIPA{{\textsubring{l}}{j}{ou}{m}{k\smash{\textsuperscript{h}}}{v}{\textsci}{ð}{a}} \dicPos{f}[1] \dicFlx{(‑u, ‑ur)}[19] \dicFieldCat{hud.} \dicSynonym{sinfónía} \dicDirectTranslationCS{symfonie}
\dicEntry[hljómleikar] \dicTerm{hljóm··leikar} \dicIPA{{\textsubring{l}}{j}{ou}{m}{l}{ei}{\r{g}}{a}{\textsubring{r}}} \dicPos{m}[6] \dicFlx{pl}[2] \dicSynonym{tónleikar} \dicDirectTranslationCS{koncert}
\dicEntry[hljómlist] \dicTerm{hljóm··list} \dicIPA{{\textsubring{l}}{j}{ou}{m}{l}{\textsci}{s}{\textsubring{d}}} \dicPos{f}[7] \dicFlx{(‑ar)}[3] \dicSynonym{tónlist} \dicDirectTranslationCS{hudba}
\dicEntry[hljómlistarmaður] \dicTerm{hljóm·listar··|maður} \dicIPA{{\textsubring{l}}{j}{ou}{m}{l}{\textsci}{s}{\textsubring{d}}{a}{r}{m}{a}{ð}{\textscy}{\textsubring{r}}} \dicPos{m}[13] \dicFlx{(‑manns, ‑menn)}[2] \dicDirectTranslationCS{hudebník, hudebnice, muzikant(ka)}
\dicEntry[hljómplata] \dicTerm{hljóm··|plata} \dicIPA{{\textsubring{l}}{j}{ou}{m}{p\smash{\textsuperscript{h}}}{l}{a}{\textsubring{d}}{a}} \dicPos{f}[1] \dicFlx{(‑plötu, ‑plötur)}[14] \textbf{1.} \dicSynonym*{grammófónplata} \dicDirectTranslationCS{(gramofonová) deska}  \textbf{2.} \dicSynonym*{tónlistardiskur} \dicDirectTranslationCS{deska, album, hudební CD}
\dicEntry[hljómplötubúð] \dicTerm{hljóm·plötu··búð} \dicIPA{{\textsubring{l}}{j}{ou}{m}{p\smash{\textsuperscript{h}}}{l}{\oe}{\textsubring{d}}{\textscy}{\textsubring{b}}{u}{\texttheta}} \dicPos{f}[7] \dicFlx{(‑ar, ‑ir)}[1] \dicDirectTranslationCS{hudebniny, obchod s~hudebninami}
\dicEntry[hljómsveit] \dicTerm{hljóm··sveit} \dicsymFrequent\  \dicIPA{{\textsubring{l}}{j}{ou}{m}{s}{v}{ei}{\textsubring{d}}} \dicPos{f}[7] \dicFlx{(‑ar, ‑ir)}[1] \dicDirectTranslationCS{(hudební) skupina, kapela, orchestr} \dicExampleIS{stofna hljómsveit} \dicExampleCS{založit hudební skupinu}
\dicEntry[hljómsveitarstjóri] \dicTerm{hljóm·sveitar··stjór|i} \dicIPA{{\textsubring{l}}\-{j}\-{ou}\-{m}\-{s}\-{v}\-{ei}\-{\textsubring{d}}\-{a}\-{\textsubring{r}}\-{s}\-{\textsubring{d}}\-{j}\-{ou}\-{r}\-{\textsci}\-} \dicPos{m}[1] \dicFlx{(‑a, ‑ar)}[1] \dicDirectTranslationCS{kapelník, kapelnice, dirigent(ka)}
\dicEntry[hljómtæki] \dicTerm{hljóm··tæki} \dicIPA{{\textsubring{l}}{j}{ou}{m}{t\smash{\textsuperscript{h}}}{a}{i}{\r{\textObardotlessj}}{\textsci}} \dicPos{n}[2] \dicFlx{(‑s, ‑)}[16] \dicDirectTranslationCS{hi-fi věž, audio souprava}
\dicEntry[hljómur] \dicTerm{hljóm|ur} \dicIPA{{\textsubring{l}}{j}{ou}{\textlengthmark}{m}{\textscy}{\textsubring{r}}} \dicPos{m}[6] \dicFlx{(‑s, ‑ar)}[24] \textbf{1.} \dicSynonym{hljóð} \dicDirectTranslationCS{zvuk}  \textbf{2.} \dicFieldCat{hud.} \dicDirectTranslationCS{akord}
\dicEntry[hljóp] \dicTerm{hljóp} \dicIPA{{\textsubring{l}}{j}{ou}{\textlengthmark}{\textsubring{b}}} \dicPos{v} \dicFlx{ind pf sg 1 pers} \dicLink{hlaupa}
\dicEntry[hljóta] \dicTerm{hljóta} \dicsymFrequent\  \dicIPA{{\textsubring{l}}{j}{ou}{\textlengthmark}{\textsubring{d}}{a}} \dicPos{v}[6] \dicFlx{(hlýt, hlaut, hlutum, hlyti, hlotið)}[94] \dicFlx{acc} \textbf{1.} \dicSynonym{hlotnast} \dicDirectTranslationCS{obdržet, dostat} \dicExampleIS{hljóta viðurkenningu} \dicExampleCS{obdržet uznání}  \textbf{2.} \dicSynonym*{mega til} \dicDirectTranslationCS{muset} \dicIndirectTranslationCS{(modální sloveso vyjadřující nutnost)} \dicExampleIS{Ég hlýt að gera það.} \dicExampleCS{Musím to udělat.}  \textbf{3.} \dicDirectTranslationCS{muset} \dicIndirectTranslationCS{(modální sloveso vyjadřující pravděpodobnost)} \dicExampleIS{Það hlaut að vera!} \dicExampleCS{Muselo to tak být!};  \dicIdiom{hljótast}{ \dicPhraseIS{e‑að hlýst af e‑u}} \dicFlx{refl} \dicDirectTranslationCS{(co) vzniká z~důvodu (čeho)}
\dicEntry[hlotið] \dicTerm{hlotið} \dicIPA{{\textsubring{l}}{\textopeno}{\textlengthmark}{\textsubring{d}}{\textsci}{\texttheta}} \dicPos{v} \dicFlx{supin} \dicLink{hljóta}
\dicEntry[hlotnast] \dicTerm{hlotn|ast} \dicIPA{{\textsubring{l}}{\textopeno}{h}{\textsubring{d}}{n}{a}{s}{\textsubring{d}}} \dicPos{v}[1] \dicFlx{(‑aðist)}[102] \dicFlx{refl impers} \dicPhraseIS{e‑m hlotnast e‑að} \dicDirectTranslationCS{(kdo co) obdrží, (komu) se (co) dostává (sláva ap.)}
\dicEntry[hló] \dicTerm{hló} \dicIPA{{\textsubring{l}}{ou}{\textlengthmark}} \dicPos{v} \dicFlx{ind pf sg 1 pers} \dicLink{hlæja}
\dicEntry[hlóð] \dicTerm{hlóð} \dicIPA{{\textsubring{l}}{ou}{\textlengthmark}{\texttheta}} \dicPos{v} \dicFlx{ind pf sg 1 pers} \dicLink{hlaða\smash{\textsuperscript{2}}}
\dicEntry[hlóðir] \dicTerm{hlóðir} \dicIPA{{\textsubring{l}}{ou}{\textlengthmark}{ð}{\textsci}{\textsubring{r}}} \dicPos{f}[12] \dicFlx{pl}[10] \dicDirectTranslationCS{ohniště} \dicExampleIS{elda í hlóðum} \dicExampleCS{vařit na ohništi}
\dicEntry[hlóðum] \dicTerm{hlóðum} \dicIPA{{\textsubring{l}}{ou}{\textlengthmark}{ð}{\textscy}{\textsubring{m}}} \dicPos{v} \dicFlx{ind pf pl 1 pers} \dicLink{hlaða\smash{\textsuperscript{2}}}
\dicEntry[hlógum] \dicTerm{hlógum} \dicIPA{{\textsubring{l}}{ou}{\textlengthmark}{\textscy}{\textsubring{m}}} \dicPos{v} \dicFlx{ind pf pl 1 pers} \dicLink{hlæja}
\dicEntry[hlst.] \dicTerm{hlst.} \dicPos{zkr} \dicPhraseIS{hliðstæður} \dicDirectTranslationCS{podobný, analogický}
\dicEntry[hlunka] \dicTerm{hlunk|a} \dicIPA{{\textsubring{l}}{u}{\r{\ng}}{\r{g}}{a}} \dicPos{v}[1] \dicFlx{(‑aði)}[1] \dicDirectTranslationCS{žuchnout, bouchnout (o~zem ap.)} \dicExampleIS{hlunka e‑u niður} \dicExampleCS{bouchnout (čím) o~zem};  \dicIdiom{hlunkast}{ \dicPhraseIS{hlunkast}} \dicFlx{refl} \dicSynonym*{detta með dynk} \dicDirectTranslationCS{žuchnout, bouchnout (o~zem ap.)} \dicExampleIS{hlunkast á bakhlutann} \dicExampleCS{žuchnout sebou na zadek}
\dicEntry[hlunnindi] \dicTerm{hlunn··indi} \dicIPA{{\textsubring{l}}{\textscy}{n}{\textlengthmark}{\textsci}{n}{\textsubring{d}}{\textsci}} \dicPos{n}[2] \dicFlx{pl}[19] \dicDirectTranslationCS{výhody, požitky} \dicExampleIS{sérstök hlunnindi} \dicExampleCS{speciální výhody}
\dicEntry[hlunnur] \dicTerm{hlunn|ur} \dicIPA{{\textsubring{l}}{\textscy}{n}{\textlengthmark}{\textscy}{\textsubring{r}}} \dicPos{m}[6] \dicFlx{(‑s, ‑ar)}[4] \dicPhraseIS{vera kominn á fremsta hlunn með að (gera e‑ð)} \dicDirectTranslationCS{právě se chystat (co udělat)}
\dicEntry[hlupum] \dicTerm{hlupum} \dicIPA{{\textsubring{l}}{\textscy}{\textlengthmark}{\textsubring{b}}{\textscy}{\textsubring{m}}} \dicPos{v} \dicFlx{ind pf pl 1 pers} \dicLink{hlaupa}
\dicEntry[hlussa] \dicTerm{hluss|a} \dicIPA{{\textsubring{l}}{\textscy}{s}{\textlengthmark}{a}} \dicPos{f}[1] \dicFlx{(‑u, ‑ur)}[7] \dicSynonym*{þung kona} \dicDirectTranslationCS{postarší žena silnější postavy, matróna}
\dicEntry[hlust] \dicTerm{hlust} \dicIPA{{\textsubring{l}}{\textscy}{s}{\textsubring{d}}} \dicPos{f}[7] \dicFlx{(‑ar, ‑ir)}[1] \dicFieldCat{anat.} \dicDirectTranslationCS{zvukovod};  \dicPhraseIS{leggja við hlustir(nar)} \dicLangCat{přen.} \dicDirectTranslationCS{natáhnout\,/\addthin natahovat uši}
\dicEntry[hlusta] \dicTerm{hlust|a} \dicsymFrequent\  \dicIPA{{\textsubring{l}}{\textscy}{s}{\textsubring{d}}{a}} \dicPos{v}[1] \dicFlx{(‑aði)}[1] \dicDirectTranslationCS{poslouchat (rádio ap.), naslouchat};  \dicPhraseIS{hlusta á e‑ð} \dicDirectTranslationCS{poslouchat (co)} \dicExampleIS{hlusta með öðru eyranu á útvarpið} \dicExampleCS{poslouchat jedním uchem rádio};  \dicPhraseIS{hlusta á e‑n (syngja)} \dicDirectTranslationCS{poslouchat (koho) (zpívat)};  \dicPhraseIS{hlusta eftir e‑m} \dicDirectTranslationCS{podívat se na (koho), dohlédnout na (koho)}
\dicEntry[hlustandi] \dicTerm{hlust··|andi} \dicIPA{{\textsubring{l}}{\textscy}{s}{\textsubring{d}}{a}{n}{\textsubring{d}}{\textsci}} \dicPos{m}[2] \dicFlx{(‑anda, ‑endur)}[1] \dicDirectTranslationCS{posluchač(ka)} \dicExampleIS{Hlustendurnir mega hringja og biðja um lag.} \dicExampleCS{Posluchači mohou zavolat a~požádat o~písničku.}
\dicEntry[hlustun] \dicTerm{hlust|un} \dicIPA{{\textsubring{l}}{\textscy}{s}{\textsubring{d}}{\textscy}{\textsubring{n}}} \dicPos{f}[7] \dicFlx{(‑unar)}[9] \textbf{1.} \dicSynonym*{það að hlusta} \dicDirectTranslationCS{poslech, poslouchání, naslouchání}  \textbf{2.} \dicSynonym*{útvarpshlustun} \dicDirectTranslationCS{poslech, poslouchání (rádia ap.)}  \textbf{3.} \dicFieldCat{med.} \dicDirectTranslationCS{poslech, auskultace}
\dicEntry[hluta] \dicTerm{hlut|a} \dicIPA{{\textsubring{l}}{\textscy}{\textlengthmark}{\textsubring{d}}{a}} \dicPos{v}[1] \dicFlx{(‑aði)}[1] \dicFlx{acc} \dicSynonym{deila\smash{\textsuperscript{2}}} \dicDirectTranslationCS{rozdělit, rozporcovat} \dicExampleIS{hluta e‑ð sundur} \dicExampleCS{rozporcovat (co)};  \dicIdiom{hlutast}[til]{ \dicPhraseIS{hlutast til um e‑ð}} \dicFlx{refl} \dicDirectTranslationCS{zaplést se do (čeho), zapojit se do (čeho) (do konfliktu ap.)}
\dicEntry[hlutabréf] \dicTerm{hluta··bréf} \dicIPA{{\textsubring{l}}{\textscy}{\textlengthmark}{\textsubring{d}}{a}{\textsubring{b}}{r}{j}{\textepsilon}{f}} \dicPos{n}[2] \dicFlx{(‑s, ‑)}[5] \dicFieldCat{ekon.} \dicDirectTranslationCS{akcie}
\dicEntry[hlutabréfaeign] \dicTerm{hluta·bréfa··eign} \dicIPA{{\textsubring{l}}{\textscy}{\textlengthmark}{\textsubring{d}}{a}{\textsubring{b}}{r}{j}{\textepsilon}{v}{a}{ei}{\r{g}}{\textsubring{n}}} \dicPos{f}[7] \dicFlx{(‑ar)}[3] \dicFieldCat{ekon.} \dicDirectTranslationCS{držba akcií}
\dicEntry[hlutaðeigandi] \dicTerm{hlut··að·eig·|andi\smash{\textsuperscript{1}}} \dicIPA{{\textsubring{l}}{\textscy}{\textlengthmark}{\textsubring{d}}{a}{ð}{ei}{\textbabygamma}{a}{n}{\textsubring{d}}{\textsci}} \dicPos{m}[2] \dicFlx{(‑anda, ‑endur)}[1] \dicDirectTranslationCS{dotyčný, dotyčná} \dicIndirectTranslationCS{(dříve uvedená\,/\addthin zmíněná osoba)}
\dicEntry[hlutaðeigandi] \dicTerm{hlut··að·eig·andi\smash{\textsuperscript{2}}} \dicIPA{{\textsubring{l}}{\textscy}{\textlengthmark}{\textsubring{d}}{a}{ð}{ei}{\textbabygamma}{a}{n}{\textsubring{d}}{\textsci}} \dicPos{adj}[13] \dicFlx{indecl}[1] \dicDirectTranslationCS{dotyčný, příslušný}
\dicEntry[hlutafé] \dicTerm{hluta··|fé} \dicIPA{{\textsubring{l}}{\textscy}{\textlengthmark}{\textsubring{d}}{a}{f}{j}{\textepsilon}} \dicPos{n}[3] \dicFlx{(‑fjár)}[5] \dicFieldCat{ekon.} \dicDirectTranslationCS{akcie, akciový kapitál}
\dicEntry[hlutafélag] \dicTerm{hluta··fé·|lag} \dicsymFrequent\  \dicIPA{{\textsubring{l}}{\textscy}{\textlengthmark}{\textsubring{d}}{a}{f}{j}{\textepsilon}{l}{a}{x}} \dicPos{n}[2] \dicFlx{(‑lags, ‑lög)}[8] \dicFieldCat{ekon.} \dicDirectTranslationCS{akciová společnost} \dicExampleIS{framkvæmdastjóri hlutafélags} \dicExampleCS{výkonný ředitel akciové společnosti}
\dicEntry[hlutafjáreign] \dicTerm{hluta·fjár··eign} \dicIPA{{\textsubring{l}}{\textscy}{\textlengthmark}{\textsubring{d}}{a}{f}{j}{au}{r}{ei}{\r{g}}{\textsubring{n}}} \dicPos{f}[7] \dicFlx{(‑ar)}[3] \dicFieldCat{ekon.} \dicDirectTranslationCS{držba akcií}
\dicEntry[hlutastarf] \dicTerm{hluta··|starf} \dicIPA{{\textsubring{l}}{\textscy}{\textlengthmark}{\textsubring{d}}{a}{s}{\textsubring{d}}{a}{r}{f}} \dicPos{n}[2] \dicFlx{(‑starfs, ‑störf)}[8] \dicDirectTranslationCS{práce na částečný úvazek}
\dicEntry[hlutbundinn] \dicTerm{hlut··bundinn} \dicIPA{{\textsubring{l}}{\textscy}{\textlengthmark}{\textsubring{d}}{\textsubring{b}}{\textscy}{n}{\textsubring{d}}{\textsci}{\textsubring{n}}} \dicPos{adj}[6]\dicFlx{}[-2] \textbf{1.} \dicSynonym*{hlutkenndur} \dicDirectTranslationCS{figurativní} \dicExampleIS{hlutbundið málverk} \dicExampleCS{figurativní malba}  \textbf{2.} \dicDirectTranslationCS{poměrný, proporční, proporcionální} \dicExampleIS{hlutbundnar kosningar} \dicExampleCS{poměrný volební systém}
\dicEntry[hlutdeild] \dicTerm{hlut··deild} \dicIPA{{\textsubring{l}}{\textscy}{\textlengthmark}{\textsubring{d}}{ei}{l}{\textsubring{d}}} \dicPos{f}[7] \dicFlx{(‑ar)}[3] \dicSynonym{þátttaka} \dicDirectTranslationCS{podíl, (spolu)účast};  \dicPhraseIS{eiga hlutdeild í e‑u} \dicDirectTranslationCS{mít na (čem) podíl, účastnit se (čeho)}
\dicEntry[hlutdrægni] \dicTerm{hlut··drægn|i} \dicIPA{{\textsubring{l}}{\textscy}{\textlengthmark}{\textsubring{d}}{r}{a}{i}{\r{g}}{n}{\textsci}} \dicPos{f}[3] \dicFlx{(‑i)}[3] \dicDirectTranslationCS{zaujatost, stranictví}
\dicEntry[hlutdrægur] \dicTerm{hlut··drægur} \dicIPA{{\textsubring{l}}{\textscy}{\textlengthmark}{\textsubring{d}}{r}{a}{i}{\textbabygamma}{\textscy}{\textsubring{r}}} \dicPos{adj}[1]\dicFlx{}[-1] \dicDirectTranslationCS{zaujatý, stranící (jedné straně ap.)}
\dicEntry[hlutfall] \dicTerm{hlut··|fall} \dicIPA{{\textsubring{l}}{\textscy}{\textlengthmark}{\textsubring{d}}{f}{a}{\textsubring{d}}{\textsubring{l}}} \dicPos{n}[2] \dicFlx{(‑falls, ‑föll)}[8] \dicDirectTranslationCS{poměr, podíl, proporce} \dicExampleIS{hlutfallið 2:3} \dicExampleCS{poměr dva ku třem}
\dicEntry[hlutfallskosning] \dicTerm{hlut·falls··kos·ning} \dicIPA{{\textsubring{l}}{\textscy}{\textlengthmark}{\textsubring{d}}{f}{a}{l}{s}{k\smash{\textsuperscript{h}}}{\textopeno}{s}{\textsubring{d}}{n}{i}{\ng}{\r{g}}} \dicPos{f}[4] \dicFlx{(‑ar, ‑ar)}[5] \dicFieldCat{pol.} \dicDirectTranslationCS{proporční volby}
\dicEntry[hlutfallslegur] \dicTerm{hlut·falls··legur} \dicIPA{{\textsubring{l}}{\textscy}{\textlengthmark}{\textsubring{d}}{f}{a}{l}{s}{l}{\textepsilon}{\textbabygamma}{\textscy}{\textsubring{r}}} \dicPos{adj}[1]\dicFlx{}[-8] \dicDirectTranslationCS{poměrný, proporcionální}
\dicEntry[hlutgerving] \dicTerm{hlut··gerv·ing} \dicIPA{{\textsubring{l}}{\textscy}{\textlengthmark}{\textsubring{d}}{\r{\textObardotlessj}}{\textepsilon}{r}{v}{i}{\ng}{\r{g}}} \dicPos{f}[4] \dicFlx{(‑ar)}[7] \dicFieldCat{lit.} \dicDirectTranslationCS{zvěcnění, substanciace} \dicAntonym{persónugerving}
\dicEntry[hluthafafundur] \dicTerm{hlut·hafa··fund|ur} \dicIPA{{\textsubring{l}}{\textscy}{\textlengthmark}{\textsubring{d}}{h}{a}{v}{a}{f}{\textscy}{n}{\textsubring{d}}{\textscy}{\textsubring{r}}} \dicPos{m}[10] \dicFlx{(‑ar, ‑ir)}[4] \dicFieldCat{ekon.} \dicDirectTranslationCS{valná hromada}
\dicEntry[hluthafi] \dicTerm{hlut··haf|i} \dicIPA{{\textsubring{l}}{\textscy}{\textlengthmark}{\textsubring{d}}{h}{a}{v}{\textsci}} \dicPos{m}[1] \dicFlx{(‑a, ‑ar)}[8] \dicFieldCat{ekon.} \dicDirectTranslationCS{akcionář(ka)}
\dicEntry[hluthyggja] \dicTerm{hlut··hyggj|a} \dicIPA{{\textsubring{l}}{\textscy}{\textlengthmark}{\textsubring{d}}{h}{\textsci}{\r{\textObardotlessj}}{a}} \dicPos{f}[1] \dicFlx{(‑u)}[5] \dicFieldCat{filos.} \dicDirectTranslationCS{realismus}
\dicEntry[hluti] \dicTerm{hlut|i} \dicsymFrequent\  \dicIPA{{\textsubring{l}}{\textscy}{\textlengthmark}{\textsubring{d}}{\textsci}} \dicPos{m}[1] \dicFlx{(‑a, ‑ar)}[1] \dicDirectTranslationCS{část, díl} \dicExampleIS{síðari hluti verksins} \dicExampleCS{poslední část díla};  \dicPhraseIS{hluti af e‑u} \dicDirectTranslationCS{část (čeho)};  \dicPhraseIS{meiri hluti} \dicDirectTranslationCS{většina, větší část};  \dicPhraseIS{minni hluti} \dicDirectTranslationCS{menšina, menší část};  \dicPhraseIS{að hluta til} \dicFlx{adv} \dicDirectTranslationCS{zčásti, částečně}
\dicEntry[hlutkesti] \dicTerm{hlut··kesti} \dicIPA{{\textsubring{l}}{\textscy}{\textlengthmark}{\textsubring{d}}{c\smash{\textsuperscript{h}}}{\textepsilon}{s}{\textsubring{d}}{\textsci}} \dicPos{n}[2] \dicFlx{(‑s)}[20] \dicDirectTranslationCS{losování};  \dicPhraseIS{varpa hlutkesti um e‑ð} \dicDirectTranslationCS{losovat o~(co), hodit si o~(co) mincí}
\dicEntry[hlutlaus] \dicTerm{hlut··laus} \dicIPA{{\textsubring{l}}{\textscy}{\textlengthmark}{\textsubring{d}}{l}{\oe i}{s}} \dicPos{adj}[5]\dicFlx{}[-1] \dicSynonym{óhlutdrægur} \dicDirectTranslationCS{nezaujatý, nestranný, neutrální}
\dicEntry[hlutleysi] \dicTerm{hlut··leysi} \dicIPA{{\textsubring{l}}{\textscy}{\textlengthmark}{\textsubring{d}}{l}{ei}{s}{\textsci}} \dicPos{n}[2] \dicFlx{(‑s)}[20] \dicDirectTranslationCS{nezaujatost, nestrannost, neutralita}
\dicEntry[hlutlægni] \dicTerm{hlut··lægn|i} \dicIPA{{\textsubring{l}}{\textscy}{\textlengthmark}{\textsubring{d}}{l}{a}{i}{\r{g}}{n}{\textsci}} \dicPos{f}[3] \dicFlx{(‑i)}[3] \dicDirectTranslationCS{objektivnost, objektivita} \dicAntonym{huglægni}
\dicEntry[hlutlægur] \dicTerm{hlut··lægur} \dicIPA{{\textsubring{l}}{\textscy}{\textlengthmark}{\textsubring{d}}{l}{a}{i}{\textbabygamma}{\textscy}{\textsubring{r}}} \dicPos{adj}[1]\dicFlx{}[-1] \dicDirectTranslationCS{objektivní, nestranný, věcný} \dicAntonym{huglægur}
\dicEntry[hlutskarpur] \dicTerm{hlut··|skarpur} \dicIPA{{\textsubring{l}}{\textscy}{\textlengthmark}{\textsubring{d}}{s}{\r{g}}{a}{\textsubring{r}}{\textsubring{b}}{\textscy}{\textsubring{r}}} \dicPos{adj}[1] \dicFlx{(f ‑skörp)}[2] \dicDirectTranslationCS{vítězný};  \dicPhraseIS{verða hlutskarpari} \dicDirectTranslationCS{odnést vítězství}
\dicEntry[hlutskipti] \dicTerm{hlut··skipti} \dicIPA{{\textsubring{l}}{\textscy}{\textlengthmark}{\textsubring{d}}{s}{\r{\textObardotlessj}}{\textsci}{f}{\textsubring{d}}{\textsci}} \dicPos{n}[2] \dicFlx{(‑s)}[20] \dicDirectTranslationCS{úděl, osud} \dicExampleIS{vera ósáttur við hlutskipti sitt} \dicExampleCS{nebýt smířený se svým údělem}
\dicEntry[hlutstæður] \dicTerm{hlut··stæður} \dicIPA{{\textsubring{l}}{\textscy}{\textlengthmark}{\textsubring{d}}{s}{\textsubring{d}}{a}{i}{ð}{\textscy}{\textsubring{r}}} \dicPos{adj}[2]\dicFlx{}[-6] \dicSynonym{tiltekinn} \dicDirectTranslationCS{konkrétní}
\dicEntry[hluttekning] \dicTerm{hlut··tek·ning} \dicIPA{{\textsubring{l}}{\textscy}{\textlengthmark}{t\smash{\textsuperscript{h}}}{\textepsilon}{h}{\r{g}}{n}{i}{\ng}{\r{g}}} \dicPos{f}[4] \dicFlx{(‑ar)}[7] \dicDirectTranslationCS{soucítění, soustrast, empatie, účast}
\dicEntry[hlutum] \dicTerm{hlutum} \dicIPA{{\textsubring{l}}{\textscy}{\textlengthmark}{\textsubring{d}}{\textscy}{\textsubring{m}}} \dicPos{v} \dicFlx{ind pf pl 1 pers} \dicLink{hljóta}
\dicEntry[hlutur] \dicTerm{hlut|ur} \dicsymFrequent\  \dicIPA{{\textsubring{l}}{\textscy}{\textlengthmark}{\textsubring{d}}{\textscy}{\textsubring{r}}} \dicPos{m}[10] \dicFlx{(‑ar, ‑ir)}[13] \textbf{1.} \dicSynonym{gripur} \dicDirectTranslationCS{věc, předmět, objekt}  \textbf{2.} \dicSynonym{hluti} \dicDirectTranslationCS{(po)díl, část};  \dicPhraseIS{eiga hlut að e‑u} \dicDirectTranslationCS{mít na (čem) podíl (na hádce ap.)};  \dicPhraseIS{eiga hlut í e‑u} \dicDirectTranslationCS{mít podíl v~(čem)} \dicExampleIS{eiga hlut í fyrirtækinu} \dicExampleCS{mít podíl ve společnosti};  \dicIdiom{hlutur}{ \dicPhraseIS{bera\,/\addthin hafa hærra hlut}} \dicSynonym{sigra} \dicDirectTranslationCS{vítězit, vyhrávat, mít navrch};  \dicPhraseIS{bíða lægri hlut} \dicDirectTranslationCS{prohrávat, ztrácet};  \dicPhraseIS{taka e‑u eins og sjálfsögðum hlut} \dicDirectTranslationCS{brát (co) jako samozřejmost};  \dicPhraseIS{(þarna var) ekki nokkur skapaður hlutur} \dicDirectTranslationCS{(nebyla tam) ani jedna jediná věc}
\dicEntry[hlutverk] \dicTerm{hlut··verk} \dicIPA{{\textsubring{l}}{\textscy}{\textlengthmark}{\textsubring{d}}{v}{\textepsilon}{\textsubring{r}}{\r{g}}} \dicPos{n}[2] \dicFlx{(‑s, ‑)}[5] \textbf{1.} \dicSynonym{verkefni} \dicDirectTranslationCS{role, úloha} \dicExampleIS{fá e‑m sérstakt hlutverk} \dicExampleCS{přidělit (komu) speciální úlohu}  \textbf{2.} \dicSynonym*{í lekriti} \dicDirectTranslationCS{role, úloha} \dicIndirectTranslationCS{(herce v~divadelní hře nebo filmu)}
\dicEntry[hlutveruleiki] \dicTerm{hlut··veru·leik|i} \dicIPA{{\textsubring{l}}{\textscy}{\textlengthmark}{\textsubring{d}}{v}{\textepsilon}{r}{\textscy}{l}{ei}{\r{\textObardotlessj}}{\textsci}} \dicPos{m}[1] \dicFlx{(‑a)}[3] \dicDirectTranslationCS{objektivní realita}
\dicEntry[hlúa] \dicTerm{hlú|a} \dicIPA{{\textsubring{l}}{u}{\textlengthmark}{a}} \dicPos{v}[2] \dicFlx{(‑ði, ‑ð\,/\addthin ‑að)}[183] \dicPhraseIS{hlúa að e‑m} \dicSynonym{hlynna} \dicDirectTranslationCS{pečovat o~(koho), zahrnout (koho) péčí}
\dicEntry[hlykkjóttur] \dicTerm{hlykkj··óttur} \dicIPA{{\textsubring{l}}{\textsci}{h}{\r{\textObardotlessj}}{ou}{h}{\textsubring{d}}{\textscy}{\textsubring{r}}} \dicPos{adj}[1]\dicFlx{}[-10] \dicSynonym{bugðóttur} \dicDirectTranslationCS{klikatý, hadovitý}
\dicEntry[hlykkur] \dicTerm{hlykk|ur} \dicIPA{{\textsubring{l}}{\textsci}{h}{\r{g}}{\textscy}{\textsubring{r}}} \dicPos{m}[9] \dicFlx{(‑s, ‑ir)}[15] \dicSynonym{sveigja\smash{\textsuperscript{1}}} \dicDirectTranslationCS{ohyb, zákruta, klička}
\dicEntry[hlynna] \dicTerm{hlynn|a} \dicIPA{{\textsubring{l}}{\textsci}{n}{\textlengthmark}{a}} \dicPos{v}[2] \dicFlx{(‑ti, ‑t)}[75] \dicPhraseIS{hlynna að e‑m\,/\addthin e‑u} \dicDirectTranslationCS{pečovat o~(koho\,/\addthin co), poskytnout (komu\,/\addthin čemu) péči}
\dicEntry[hlynntur] \dicTerm{hlynntur} \dicIPA{{\textsubring{l}}{\textsci}{\textsubring{n}}{\textsubring{d}}{\textscy}{\textsubring{r}}} \dicPos{adj}[1]\dicFlx{}[-10] \dicPhraseIS{vera hlynntur e‑u} \dicDirectTranslationCS{být (čemu) nakloněný}
\dicEntry[hlynur] \dicTerm{hlyn|ur} \dicIPA{{\textsubring{l}}{\textsci}{\textlengthmark}{n}{\textscy}{\textsubring{r}}} \dicPos{m}[9] \dicFlx{(‑s, ‑ir)}[10] \dicFieldCat{bot.} \dicDirectTranslationCS{javor} \textit{(l.~{\textLA{Acer}})}  \dicsymPhoto\ 
\dicFigure{ds_image_hlynur_0_1.jpg}{Hlynur}{Hlynur - Bruce Marlin, CC BY-SA 2.5}
\dicEntry[hlypi] \dicTerm{hlypi} \dicIPA{{\textsubring{l}}{\textsci}{\textlengthmark}{\textsubring{b}}{\textsci}} \dicPos{v} \dicFlx{con pf sg 1 pers} \dicLink{hlaupa}
\dicEntry[hlyti] \dicTerm{hlyti} \dicIPA{{\textsubring{l}}{\textsci}{\textlengthmark}{\textsubring{d}}{\textsci}} \dicPos{v} \dicFlx{con pf sg 1 pers} \dicLink{hljóta}
\dicEntry[hlýða] \dicTerm{hlý|ða} \dicsymFrequent\  \dicIPA{{\textsubring{l}}{i}{\textlengthmark}{ð}{a}} \dicPos{v}[2] \dicFlx{(‑ddi, ‑tt)}[173] \dicFlx{dat} \textbf{1.} \dicDirectTranslationCS{(u)poslechnout, poslouchat (rodiče ap.)} \dicExampleIS{Hundurinn hlýðir mér.} \dicExampleCS{Pes mě poslouchá.}  \textbf{2.} \dicDirectTranslationCS{poslouchat (hudbu ap.)};  \dicPhraseIS{hlýða á e‑u} \dicDirectTranslationCS{poslouchat (co)} \dicExampleIS{hlýða á tónlistina} \dicExampleCS{poslouchat hudbu}  \textbf{3.} \dicDirectTranslationCS{hodit se, být vhodný} \dicExampleIS{Honum þykir ekki hlýða að þú reykir sígarettur hér.} \dicExampleCS{Připadá mu nevhodné, abys tu kouřil cigarety.};  \dicIdiom{hlýða}[yfir]{ \dicPhraseIS{hlýða e‑m yfir e‑ð}} \dicDirectTranslationCS{(ústně) vyzkoušet (koho) z~(čeho), přezkoušet (koho) z~(čeho)}
\dicEntry[hlýðinn] \dicTerm{hlýðinn} \dicIPA{{\textsubring{l}}{i}{\textlengthmark}{ð}{\textsci}{\textsubring{n}}} \dicPos{adj}[6]\dicFlx{}[-2] \dicDirectTranslationCS{poslušný};  \dicPhraseIS{vera hlýðinn við e‑n} \dicDirectTranslationCS{poslouchat (koho)} \dicIndirectTranslationCS{(dbát něčích rad a~napomenutí)}
\dicEntry[hlýðni] \dicTerm{hlýðn|i} \dicIPA{{\textsubring{l}}{i}{ð}{n}{\textsci}} \dicPos{f}[3] \dicFlx{(‑i)}[3] \dicDirectTranslationCS{poslušnost}
\dicEntry[hlýhugur] \dicTerm{hlý··hug|ur} \dicIPA{{\textsubring{l}}{i}{\textlengthmark}{h}{\textscy}{\textbabygamma}{\textscy}{\textsubring{r}}} \dicPos{m}[10] \dicFlx{(‑ar)}[12] \dicDirectTranslationCS{sympatie, laskavost, vřelost, srdečnost}
\dicEntry[hlýindi] \dicTerm{hlý··indi} \dicIPA{{\textsubring{l}}{i}{j}{\textsci}{n}{\textsubring{d}}{\textsci}} \dicPos{n}[2] \dicFlx{pl}[19] \dicDirectTranslationCS{teplo, teplé počasí}
\dicEntry[hlýja] \dicTerm{hlýj|a\smash{\textsuperscript{1}}} \dicsymFrequent\  \dicIPA{{\textsubring{l}}{i}{j}{\textlengthmark}{a}} \dicPos{f}[1] \dicFlx{(‑u)}[5] \textbf{1.} \dicSynonym{ylur} \dicDirectTranslationCS{teplo} \dicExampleIS{sólskinshlýja} \dicExampleCS{sluneční teplo}  \textbf{2.} \dicSynonym*{hlýtt viðmót} \dicDirectTranslationCS{vřelost, vroucnost, srdečnost}
\dicEntry[hlýja] \dicTerm{hlýj|a\smash{\textsuperscript{2}}} \dicIPA{{\textsubring{l}}{i}{j}{\textlengthmark}{a}} \dicPos{v}[1] \dicFlx{(‑aði)}[44] \dicFlx{dat} \dicDirectTranslationCS{(o)hřát, zahřát} \dicExampleIS{hlýja sér} \dicExampleCS{ohřát se};  \dicPhraseIS{e‑að hlýjar e‑m hjartarætur} \dicLangCat{přen.} \dicDirectTranslationCS{(co) hřeje (koho) u~srdce}
\dicEntry[hlýlegur] \dicTerm{hlý··legur} \dicIPA{{\textsubring{l}}{i}{\textlengthmark}{l}{\textepsilon}{\textbabygamma}{\textscy}{\textsubring{r}}} \dicPos{adj}[1]\dicFlx{}[-8] \textbf{1.} \dicSynonym{notalegur} \dicDirectTranslationCS{útulný, přítulný} \dicExampleIS{Það er hlýlegt inni.} \dicExampleCS{Uvnitř je útulně.}  \textbf{2.} \dicSynonym{vingjarnlegur} \dicDirectTranslationCS{srdečný, vřelý, přátelský}
\dicEntry[hlýna] \dicTerm{hlýn|a} \dicIPA{{\textsubring{l}}{i}{\textlengthmark}{n}{a}} \dicPos{v}[1] \dicFlx{(‑aði)}[74] \dicPhraseIS{e‑að hlýnar} \dicDirectTranslationCS{(co) se otepluje} \dicExampleIS{Veðráttan er farin að hlýna.} \dicExampleCS{Klima se otepluje.};  \dicPhraseIS{það\,/\addthin hann hlýnar} \dicFlx{impers} \dicDirectTranslationCS{otepluje se};  \dicPhraseIS{það hlýnar í veðri} \dicFlx{impers} \dicDirectTranslationCS{otepluje se};  \dicPhraseIS{e‑m hlýnar} \dicFlx{impers} \dicDirectTranslationCS{(komu) je tepleji};  \dicPhraseIS{e‑m hlýnar um hjartaræturnar} \dicFlx{impers} \dicDirectTranslationCS{(koho) hřeje u~srdce}
\dicEntry[hlýnun] \dicTerm{hlýn|un} \dicIPA{{\textsubring{l}}{i}{\textlengthmark}{n}{\textscy}{\textsubring{n}}} \dicPos{f}[7] \dicFlx{(‑unar)}[9] \dicDirectTranslationCS{oteplení, oteplování} \dicExampleIS{hlýnun jarðar} \dicExampleCS{oteplování Země}
\dicEntry[hlýr] \dicTerm{hlýr} \dicsymFrequent\  \dicIPA{{\textsubring{l}}{i}{\textlengthmark}{\textsubring{r}}} \dicPos{adj}[4]\dicFlx{}[-7] \textbf{1.} \dicDirectTranslationCS{teplý (počasí ap.)};  \dicPhraseIS{það er hlýtt (úti)} \dicDirectTranslationCS{je teplo (venku)}  \textbf{2.} \dicDirectTranslationCS{teplý, hřející (svetr ap.)} \dicExampleIS{hlý peysa} \dicExampleCS{teplý svetr}  \textbf{3.} \dicDirectTranslationCS{vřelý, přívětivý}
\dicEntry[hlýri] \dicTerm{hlýr|i} \dicIPA{{\textsubring{l}}{i}{\textlengthmark}{r}{\textsci}} \dicPos{m}[1] \dicFlx{(‑a, ‑ar)}[1] \textbf{1.} \dicFieldCat{zool.} \dicDirectTranslationCS{vlkouš skvrnitý} \textit{(l.~{\textLA{Anarhichas minor}})}  \textbf{2.} \dicDirectTranslationCS{ramínko, ramenní popruh}
\dicEntry[hlýt] \dicTerm{hlýt} \dicIPA{{\textsubring{l}}{i}{\textlengthmark}{\textsubring{d}}} \dicPos{v} \dicFlx{ind praes sg 1 pers} \dicLink{hljóta}
\dicEntry[hlæ] \dicTerm{hlæ} \dicIPA{{\textsubring{l}}{a}{i}{\textlengthmark}} \dicPos{v} \dicFlx{ind praes sg 1 pers} \dicLink{hlæja}
\dicEntry[hlæði] \dicTerm{hlæði} \dicIPA{{\textsubring{l}}{a}{i}{\textlengthmark}{ð}{\textsci}} \dicPos{v} \dicFlx{con pf sg 1 pers} \dicLink{hlaða\smash{\textsuperscript{2}}}
\dicEntry[hlægi] \dicTerm{hlægi} \dicIPA{{\textsubring{l}}{a}{i}{j}{\textlengthmark}{\textsci}} \dicPos{v} \dicFlx{con pf sg 1 pers} \dicLink{hlæja}
\dicEntry[hlægilegur] \dicTerm{hlægi··legur} \dicIPA{{\textsubring{l}}{a}{i}{j}{\textlengthmark}{\textsci}{l}{\textepsilon}{\textbabygamma}{\textscy}{\textsubring{r}}} \dicPos{adj}[1]\dicFlx{}[-8] \dicSynonym{afkáralegur} \dicDirectTranslationCS{směšný, legrační, úsměvný, zábavný}
\dicEntry[hlæja] \dicTerm{hlæja} \dicsymFrequent\  \dicIPA{{\textsubring{l}}{a}{i}{j}{\textlengthmark}{a}} \dicPos{v}[6] \dicFlx{(hlæ, hló, hlógum, hlægi, hlegið)}[138] \dicDirectTranslationCS{(za)smát se, (za)chechtat se} \dicExampleIS{hlæja af hjartans lyst} \dicExampleCS{smát se z~celého srdce};  \dicIdiom{hlæja}[að]{ \dicPhraseIS{hlæja að e‑m\,/\addthin e‑u}} \dicDirectTranslationCS{smát se (komu\,/\addthin čemu)};  \dicProverb\  \dicPhraseIS{Sá hlær best sem síðast hlær.} \dicLangCat{přís.} \dicDirectTranslationCS{Kdo se směje naposledy, ten se směje nejlépe.}
\dicEntry[hlæjandi] \dicTerm{hlæj··andi} \dicsymFrequent\  \dicIPA{{\textsubring{l}}{a}{i}{j}{\textlengthmark}{a}{n}{\textsubring{d}}{\textsci}} \dicPos{adj}[13] \dicFlx{indecl}[1] \dicDirectTranslationCS{smějící se, chechtající se} \dicExampleIS{hlæjandi börn} \dicExampleCS{smějící se děti}
\dicEntry[hlöð] \dicTerm{hlöð} \dicIPA{{\textsubring{l}}{\oe}{\textlengthmark}{\texttheta}} \dicPos{n} \dicFlx{pl nom} \dicLink{hlað}
\dicEntry[hlöðu] \dicTerm{hlöðu} \dicIPA{{\textsubring{l}}{\oe}{\textlengthmark}{ð}{\textscy}} \dicPos{f} \dicFlx{sg gen} \dicLink{hlaða\smash{\textsuperscript{1}}}
\dicEntry[hlöss] \dicTerm{hlöss} \dicIPA{{\textsubring{l}}{\oe}{s}{\textlengthmark}} \dicPos{n} \dicFlx{pl nom} \dicLink{hlass}
\dicEntry[hm] \dicTerm{hm} \dicIPA{{h}{\textsubring{m}}} \dicPos{inter} \dicLink{humm}
\dicEntry[hnakkabein] \dicTerm{hnakka··bein} \dicIPA{{\textsubring{n}}{a}{h}{\r{g}}{a}{\textsubring{b}}{ei}{\textsubring{n}}} \dicPos{n}[2] \dicFlx{(‑s, ‑)}[5] \dicFieldCat{anat.} \dicDirectTranslationCS{týlní kost}
\dicEntry[hnakkapúði] \dicTerm{hnakka··púð|i} \dicIPA{{\textsubring{n}}{a}{h}{\r{g}}{a}{p\smash{\textsuperscript{h}}}{u}{ð}{\textsci}} \dicPos{m}[1] \dicFlx{(‑a, ‑ar)}[1] \dicDirectTranslationCS{opěrka hlavy}
\dicEntry[hnakki] \dicTerm{hnakk|i} \dicsymFrequent\  \dicIPA{{\textsubring{n}}{a}{h}{\r{\textObardotlessj}}{\textsci}} \dicPos{m}[1] \dicFlx{(‑a, ‑ar)}[8] \dicDirectTranslationCS{šíje, vaz, týl} \dicExampleIS{finna fyrir sársauka aftan í hnakka} \dicExampleCS{cítit bolest vzadu na šíji}
\dicEntry[hnakkur] \dicTerm{hnakk|ur} \dicIPA{{\textsubring{n}}{a}{h}{\r{g}}{\textscy}{\textsubring{r}}} \dicPos{m}[6] \dicFlx{(‑s, ‑ar)}[25] \textbf{1.} \dicDirectTranslationCS{(jezdecké) sedlo} \dicExampleIS{leggja hnakk á hest} \dicExampleCS{osedlat koně}  \textbf{2.} \dicDirectTranslationCS{sedlo (kola, motocyklu ap.)}
\dicEntry[hnappa] \dicTerm{hnapp|a} \dicIPA{{\textsubring{n}}{a}{h}{\textsubring{b}}{a}} \dicPos{v}[1] \dicFlx{(‑aði)}[13] \dicPhraseIS{hnappa sig} \dicDirectTranslationCS{seskupit se, shromáždit se, vytvořit hlouček\,/\addthin stádo} \dicExampleIS{Ærnar hnappa sig.} \dicExampleCS{Ovce se shromažďují.};  \dicIdiom{hnappast}{ \dicPhraseIS{hnappast}} \dicFlx{refl} \dicDirectTranslationCS{seskupit se, shromáždit se}
\dicEntry[hnappagat] \dicTerm{hnappa··|gat} \dicIPA{{\textsubring{n}}{a}{h}{\textsubring{b}}{a}{\r{g}}{a}{\textsubring{d}}} \dicPos{n}[2] \dicFlx{(‑gats, ‑göt)}[8] \dicDirectTranslationCS{knoflíková dírka}
\dicEntry[hnappur] \dicTerm{hnapp|ur} \dicIPA{{\textsubring{n}}{a}{h}{\textsubring{b}}{\textscy}{\textsubring{r}}} \dicPos{m}[6] \dicFlx{(‑s, ‑ar)}[25] \textbf{1.} \dicSynonym{þyrping} \dicDirectTranslationCS{shluk, hlouček}  \textbf{2.} \dicSynonym*{tala} \dicDirectTranslationCS{(ouškový) knoflík} \dicExampleIS{hnappur á skyrtu} \dicExampleCS{knoflík na košili}  \textbf{3.} \dicSynonym{takki} \dicDirectTranslationCS{tlačítko, knoflík, čudlík};  \dicIdiom{hnappur}{ \dicPhraseIS{hafa öðrum hnöppum að hneppa}} \dicLangCat{přen.} \dicDirectTranslationCS{mít plné ruce práce, mít napilno}
\dicEntry[hnarreistur] \dicTerm{hnar··reistur} \dicIPA{{\textsubring{n}}{a}{r}{\textlengthmark}{ei}{s}{\textsubring{d}}{\textscy}{\textsubring{r}}} \dicPos{adj}[1]\dicFlx{}[-10] \dicDirectTranslationCS{vzpřímený, vztyčený}
\dicEntry[hnattar] \dicTerm{hnattar} \dicIPA{{\textsubring{n}}{a}{h}{\textsubring{d}}{a}{\textsubring{r}}} \dicPos{m} \dicFlx{sg gen} \dicLink{hnöttur}
\dicEntry[hnattferð] \dicTerm{hnatt··ferð} \dicIPA{{\textsubring{n}}{a}{h}{\textsubring{d}}{f}{\textepsilon}{r}{\texttheta}} \dicPos{f}[7] \dicFlx{(‑ar, ‑ir)}[1] \dicDirectTranslationCS{cesta kolem světa}
\dicEntry[hnattlíkan] \dicTerm{hnatt··lík|an} \dicIPA{{\textsubring{n}}{a}{h}{\textsubring{d}}{l}{i}{\r{g}}{a}{\textsubring{n}}} \dicPos{n}[2] \dicFlx{(‑ans, ‑ön)}[8] \dicDirectTranslationCS{glóbus}
\dicEntry[hnattrænn] \dicTerm{hnatt··rænn} \dicIPA{{\textsubring{n}}{a}{h}{\textsubring{d}}{r}{a}{i}{\textsubring{d}}{\textsubring{n}}} \dicPos{adj}[7]\dicFlx{}[-1] \dicDirectTranslationCS{globální} \dicExampleIS{hnattræn hlýnun} \dicExampleCS{globální oteplování}
\dicEntry[hnattstaða] \dicTerm{hnatt··|staða} \dicIPA{{\textsubring{n}}{a}{h}{\textsubring{d}}{s}{\textsubring{d}}{a}{ð}{a}} \dicPos{f}[1] \dicFlx{(‑stöðu)}[2] \dicDirectTranslationCS{zeměpisná pozice}
\dicEntry[hnaus] \dicTerm{hnaus} \dicIPA{{\textsubring{n}}{\oe i}{\textlengthmark}{s}} \dicPos{m}[4] \dicFlx{(‑s, ‑ar)}[25] \dicDirectTranslationCS{hrouda (hlíny), drn}
\dicEntry[hnausþykkur] \dicTerm{hnaus··þykkur} \dicIPA{{\textsubring{n}}{\oe i}{\textlengthmark}{s}{\texttheta}{\textsci}{h}{\r{g}}{\textscy}{\textsubring{r}}} \dicPos{adj}[1]\dicFlx{}[-6] \dicDirectTranslationCS{zahuštěný, velmi hustý}
\dicEntry[hnaut] \dicTerm{hnaut} \dicIPA{{\textsubring{n}}{\oe i}{\textlengthmark}{\textsubring{d}}} \dicPos{v} \dicFlx{ind pf sg 1 pers} \dicLink{hnjóta}
\dicEntry[hnefahögg] \dicTerm{hnefa··högg} \dicIPA{{\textsubring{n}}{\textepsilon}{\textlengthmark}{v}{a}{h}{\oe}{\r{g}}} \dicPos{n}[2] \dicFlx{(‑s, ‑)}[5] \dicDirectTranslationCS{rána\,/\addthin úder pěstí}
\dicEntry[hnefaleikar] \dicTerm{hnefa··leikar} \dicIPA{{\textsubring{n}}{\textepsilon}{\textlengthmark}{v}{a}{l}{ei}{\r{g}}{a}{\textsubring{r}}} \dicPos{m}[6] \dicFlx{pl}[2] \dicFieldCat{sport.} \dicDirectTranslationCS{box}
\dicEntry[hnefaleikari] \dicTerm{hnefa··leik·ar|i} \dicIPA{{\textsubring{n}}{\textepsilon}{\textlengthmark}{v}{a}{l}{ei}{\r{g}}{a}{r}{\textsci}} \dicPos{m}[1] \dicFlx{(‑a, ‑ar)}[13] \dicDirectTranslationCS{boxer(ka)}
\dicEntry[hnefi] \dicTerm{hnef|i} \dicsymFrequent\  \dicIPA{{\textsubring{n}}{\textepsilon}{\textlengthmark}{v}{\textsci}} \dicPos{m}[1] \dicFlx{(‑a, ‑ar)}[1] \textbf{1.} \dicSynonym*{kreppt hönd} \dicDirectTranslationCS{pěst};  \dicPhraseIS{ekki fyrr en í fulla hnefa} \dicLangCat{přen.} \dicDirectTranslationCS{dokud se nevyčerpají všechny možnosti};  \dicPhraseIS{setja hnefann í borðið} \dicLangCat{přen.} \dicDirectTranslationCS{udeřit pěstí do stolu} \dicIndirectTranslationCS{(říct jasně ne)};  \dicPhraseIS{steyta hnefann} \dicDirectTranslationCS{zatnout pěst}  \textbf{2.} \dicSynonym*{hnefafylli} \dicDirectTranslationCS{hrst} \dicExampleIS{hnefi af salti} \dicExampleCS{hrst soli}
\dicEntry[hnegg] \dicTerm{hnegg} \dicIPA{{\textsubring{n}}{\textepsilon}{\r{g}}{\textlengthmark}} \dicPos{n}[2] \dicFlx{(‑s)}[2] \dicDirectTranslationCS{ržání, řehtání}
\dicEntry[hneggja] \dicTerm{hneggj|a} \dicIPA{{\textsubring{n}}{\textepsilon}{\r{\textObardotlessj}}{\textlengthmark}{a}} \dicPos{v}[1] \dicFlx{(‑aði)}[44] \dicDirectTranslationCS{(za)ržát, (za)řehtat} \dicExampleIS{Hesturinn hneggjar.} \dicExampleCS{Kůň řehtá.}
\dicEntry[hneig] \dicTerm{hneig} \dicIPA{{\textsubring{n}}{ei}{\textlengthmark}{x}} \dicPos{v} \dicFlx{ind pf sg 1 pers} \dicLink{hníga}
\dicEntry[hneigð] \dicTerm{hneigð} \dicIPA{{\textsubring{n}}{ei}{\textbabygamma}{\texttheta}} \dicPos{f}[7] \dicFlx{(‑ar, ‑ir)}[1] \dicSynonym{tilhneiging} \dicDirectTranslationCS{sklon, inklinace, tíhnutí}
\dicEntry[hneigður] \dicTerm{hneigður} \dicIPA{{\textsubring{n}}{ei}{\textbabygamma}{ð}{\textscy}{\textsubring{r}}} \dicPos{adj}[2]\dicFlx{}[-1] \dicPhraseIS{vera hneigður fyrir e‑ð} \dicDirectTranslationCS{mít sklon k~(čemu)};  \dicPhraseIS{vera hneigður til e‑s} \dicDirectTranslationCS{mít pro (co) sklony}
\dicEntry[hneiging] \dicTerm{hneig··ing} \dicIPA{{\textsubring{n}}{ei}{j}{\textlengthmark}{i}{\ng}{\r{g}}} \dicPos{f}[4] \dicFlx{(‑ar, ‑ar)}[5] \dicSynonym{beyging} \dicDirectTranslationCS{poklona, úklona, klanění} \dicExampleIS{djúp hneiging} \dicExampleCS{hluboká poklona}
\dicEntry[hneigja] \dicTerm{hneig|ja} \dicsymFrequent\  \dicIPA{{\textsubring{n}}{ei}{j}{\textlengthmark}{a}} \dicPos{v}[2] \dicFlx{(‑ði, ‑t)}[89] \dicPhraseIS{hneigja sig} \dicSynonym{beygja\smash{\textsuperscript{2}}} \dicDirectTranslationCS{poklonit se, uklonit se, klanět se} \dicExampleIS{hneigja sig fyrir e‑m} \dicExampleCS{poklonit se (komu)};  \dicIdiom{hneigjast}[að]{ \dicPhraseIS{hneigjast að e‑u}} \dicFlx{refl} \dicSynonym*{hallast að e‑u} \dicDirectTranslationCS{tíhnout k~(čemu), mít tendenci k~(čemu)}
\dicEntry[hneisa] \dicTerm{hneis|a} \dicIPA{{\textsubring{n}}{ei}{\textlengthmark}{s}{a}} \dicPos{f}[1] \dicFlx{(‑u)}[5] \dicSynonym{smán} \dicDirectTranslationCS{potupa, hanba, ostuda};  \dicPhraseIS{gera e‑m hneisu} \dicDirectTranslationCS{způsobit (komu) hanbu, zahanbit (koho)}
\dicEntry[hnekkir] \dicTerm{hnekk|ir} \dicIPA{{\textsubring{n}}{\textepsilon}{h}{\r{\textObardotlessj}}{\textsci}{\textsubring{r}}} \dicPos{m}[7] \dicFlx{(‑is)}[2] \dicSynonym{áfall} \dicDirectTranslationCS{šok, rána, újma};  \dicPhraseIS{bíða hnekki} \dicDirectTranslationCS{přijít k~újmě}
\dicEntry[hnekkja] \dicTerm{hnekk|ja} \dicIPA{{\textsubring{n}}{\textepsilon}{h}{\r{\textObardotlessj}}{a}} \dicPos{v}[2] \dicFlx{(‑ti, ‑t)}[25] \dicFlx{dat} \textbf{1.} \dicSynonym{hrekja} \dicDirectTranslationCS{zvrátit, obrátit (výrok ap.)}  \textbf{2.} \dicSynonym{aftra} \dicDirectTranslationCS{zabránit, překazit}
\dicEntry[hneppa] \dicTerm{hnepp|a} \dicIPA{{\textsubring{n}}{\textepsilon}{h}{\textsubring{b}}{a}} \dicPos{v}[2] \dicFlx{(‑ti, ‑t)}[63] \dicFlx{acc\,/\addthin dat} \textbf{1.} \dicFlx{dat} \dicDirectTranslationCS{zapnout (na knoflík)} \dicExampleIS{hneppa hnappi} \dicExampleCS{zapnout knoflík};  \dicPhraseIS{hneppa e‑u að sér} \dicDirectTranslationCS{zapnout knoflíky u~(čeho)};  \dicPhraseIS{hneppa e‑u frá sér} \dicDirectTranslationCS{rozepnout knoflíky u~(čeho)} \dicExampleIS{hneppa frakkanum frá sér} \dicExampleCS{rozepnout si knoflíky u~kabátu}  \textbf{2.} \dicFlx{acc} \dicDirectTranslationCS{(u)zamknout, (u)zavřít};  \dicPhraseIS{hneppa e‑n í e‑ð} \dicDirectTranslationCS{zamknout (koho) do (čeho)} \dicExampleIS{hneppa hana í fangelsi} \dicExampleCS{zamknout ji do vězení};  \dicIdiom{hneppa}{ \dicPhraseIS{hafa öðrum hnöppum að hneppa}} \dicLangCat{přen.} \dicDirectTranslationCS{mít napilno, mít plné ruce práce}
\dicEntry[hnerra] \dicTerm{hnerr|a} \dicIPA{{\textsubring{n}}{\textepsilon}{r}{\textlengthmark}{a}} \dicPos{v}[1] \dicFlx{(‑aði)}[44] \dicDirectTranslationCS{kýchnout, kýchat} \dicExampleIS{hnerra hátt} \dicExampleCS{hlasitě kýchnout}
\dicEntry[hnerri] \dicTerm{hnerr|i} \dicIPA{{\textsubring{n}}{\textepsilon}{r}{\textlengthmark}{\textsci}} \dicPos{m}[1] \dicFlx{(‑a, ‑ar)}[1] \dicDirectTranslationCS{kýchnutí, kýchání}
\dicEntry[hneta] \dicTerm{hnet|a} \dicIPA{{\textsubring{n}}{\textepsilon}{\textlengthmark}{\textsubring{d}}{a}} \dicPos{f}[1] \dicFlx{(‑u, ‑ur)}[7] \dicDirectTranslationCS{ořech, oříšek}
\dicEntry[hnettir] \dicTerm{hnettir} \dicIPA{{\textsubring{n}}{\textepsilon}{h}{\textsubring{d}}{\textsci}{\textsubring{r}}} \dicPos{m} \dicFlx{pl nom} \dicLink{hnöttur}
\dicEntry[hnetubrjótur] \dicTerm{hnetu··brjót|ur} \dicIPA{{\textsubring{n}}{\textepsilon}{\textlengthmark}{\textsubring{d}}{\textscy}{\textsubring{b}}{r}{j}{ou}{\textsubring{d}}{\textscy}{\textsubring{r}}} \dicPos{m}[6] \dicFlx{(‑s, ‑ar)}[10] \dicDirectTranslationCS{louskáček}
\dicEntry[hneyksla] \dicTerm{hneyksl|a} \dicIPA{{\textsubring{n}}{ei}{x}{s}{\textsubring{d}}{l}{a}} \dicPos{v}[1] \dicFlx{(‑aði)}[1] \dicFlx{acc} \dicSynonym*{vekja hneykslun} \dicDirectTranslationCS{vzbudit pohoršení, pohoršit, pohoršovat} \dicExampleIS{hneyksla almenning} \dicExampleCS{pohoršit veřejnost};  \dicIdiom{hneykslast}{ \dicPhraseIS{hneykslast}} \dicFlx{refl} \dicDirectTranslationCS{být pohoršený, pohoršovat se};  \dicPhraseIS{hneykslast á e‑u} \dicFlx{refl} \dicDirectTranslationCS{pohoršovat se nad (čím)} \dicExampleIS{hneykslast á klámmyndum} \dicExampleCS{pohoršovat se nad pornofilmy}
\dicEntry[hneykslaður] \dicTerm{hneyksl··|aður} \dicsymFrequent\  \dicIPA{{\textsubring{n}}{ei}{x}{s}{\textsubring{d}}{l}{a}{ð}{\textscy}{\textsubring{r}}} \dicPos{adj}[3] \dicFlx{(f ‑uð)}[3] \dicDirectTranslationCS{pohoršený} \dicExampleIS{vera hneykslaður á e‑m} \dicExampleCS{být pohoršený (kým)}
\dicEntry[hneykslanlegur] \dicTerm{hneykslan··legur} \dicIPA{{\textsubring{n}}{ei}{x}{s}{\textsubring{d}}{l}{a}{n}{l}{\textepsilon}{\textbabygamma}{\textscy}{\textsubring{r}}} \dicPos{adj}[1]\dicFlx{}[-8] \dicDirectTranslationCS{skandální, šokující, pohoršující}
\dicEntry[hneyksli] \dicTerm{hneyksli} \dicIPA{{\textsubring{n}}{ei}{x}{s}{\textsubring{d}}{l}{\textsci}} \dicPos{n}[2] \dicFlx{(‑s, ‑)}[14] \dicDirectTranslationCS{skandál, veřejné pohoršení}
\dicEntry[hneykslun] \dicTerm{hneyksl|un} \dicIPA{{\textsubring{n}}{ei}{x}{s}{\textsubring{d}}{l}{\textscy}{\textsubring{n}}} \dicPos{f}[7] \dicFlx{(‑unar)}[9] \dicDirectTranslationCS{pohoršení, pobouření}
\dicEntry[hné] \dicTerm{hné}\dicTerm{, kné} \dicIPA{{\textsubring{n}}\-{j}\-{\textepsilon}\-{\textlengthmark}\-} \dicPos{n}[2] \dicFlx{(‑s, ‑)}[36] \dicDirectTranslationCS{koleno};  \dicPhraseIS{krjúpa á kné} \dicDirectTranslationCS{kleknout na kolena};  \dicPhraseIS{vera á báðum hnjánum, vera kominn á bæði hnén} \dicLangCat{přen.} \dicDirectTranslationCS{být na kolenou, být na mizině};  \dicPhraseIS{koma fyrir kné e‑m} \dicLangCat{přen.} \dicSynonym{knékrjúpa} \dicDirectTranslationCS{prosit (koho) na kolenou};  \dicPhraseIS{koma e‑m á kné} \dicLangCat{přen.} \dicDirectTranslationCS{dostat (koho) na kolena}
\dicEntry[hnébuxur] \dicTerm{hné··buxur} \dicIPA{{\textsubring{n}}{j}{\textepsilon}{\textlengthmark}{\textsubring{b}}{\textscy}{x}{s}{\textscy}{\textsubring{r}}} \dicPos{f}[12] \dicFlx{pl}[7] \dicDirectTranslationCS{krátké kalhoty, kraťasy}
\dicEntry[hnéskel] \dicTerm{hné··skel} \dicIPA{{\textsubring{n}}{j}{\textepsilon}{\textlengthmark}{s}{\r{\textObardotlessj}}{\textepsilon}{\textsubring{l}}} \dicPos{f}[4] \dicFlx{(‑jar, ‑jar)}[9] \dicFieldCat{anat.} \dicDirectTranslationCS{čéška}
\dicEntry[hnigi] \dicTerm{hnigi} \dicIPA{{\textsubring{n}}{i}{j}{\textlengthmark}{\textsci}} \dicPos{v} \dicFlx{con pf sg 1 pers} \dicLink{hníga}
\dicEntry[hnigið] \dicTerm{hnigið} \dicIPA{{\textsubring{n}}{i}{j}{\textlengthmark}{\textsci}{\texttheta}} \dicPos{v} \dicFlx{supin} \dicLink{hníga}
\dicEntry[hniginn] \dicTerm{hniginn} \dicIPA{{\textsubring{n}}{i}{j}{\textlengthmark}{\textsci}{\textsubring{n}}} \dicPos{adj}[6]\dicFlx{}[-6] \dicPhraseIS{hniginn að aldri} \dicDirectTranslationCS{starý (věkem), stárnoucí}
\dicEntry[hnigna] \dicTerm{hnign|a} \dicIPA{{\textsubring{n}}{\textsci}{\r{g}}{n}{a}} \dicPos{v}[1] \dicFlx{(‑aði)}[73] \dicFlx{impers} \dicPhraseIS{e‑u hnignar} \dicDirectTranslationCS{(co) upadá, s~(čím) to jde z~kopce}
\dicEntry[hnignun] \dicTerm{hnign|un} \dicIPA{{\textsubring{n}}{\textsci}{\r{g}}{n}{\textscy}{\textsubring{n}}} \dicPos{f}[7] \dicFlx{(‑unar)}[9] \dicDirectTranslationCS{úpadek, pokles}
\dicEntry[hnigum] \dicTerm{hnigum} \dicIPA{{\textsubring{n}}{\textsci}{\textlengthmark}{\textbabygamma}{\textscy}{\textsubring{m}}} \dicPos{v} \dicFlx{ind pf pl 1 pers} \dicLink{hníga}
\dicEntry[hnika] \dicTerm{hnik|a} \dicIPA{{\textsubring{n}}{\textsci}{\textlengthmark}{\r{g}}{a}} \dicPos{v}[1] \dicFlx{(‑aði)}[1] \dicFlx{dat} \dicSynonym{bifa} \dicDirectTranslationCS{posunout, posouvat, pohnout} \dicExampleIS{hnika borðinu} \dicExampleCS{posunout stůl};  \dicPhraseIS{e‑m verður ekki hnikað} \dicFlx{impers} \dicLangCat{přen.} \dicDirectTranslationCS{s~(kým) nejde pohnout, (kdo) se nedá přesvědčit};  \dicIdiom{hnika}[til]{ \dicPhraseIS{hnika e‑u til}} \dicDirectTranslationCS{upravit (co) (slovní obrat ap.)}
\dicEntry[hnippa] \dicTerm{hnipp|a} \dicIPA{{\textsubring{n}}{\textsci}{h}{\textsubring{b}}{a}} \dicPos{v}[2] \dicFlx{(‑ti, ‑t)}[78] \dicPhraseIS{hnippa í e‑n} \dicDirectTranslationCS{šťouchnout (koho)}
\dicEntry[hnipra] \dicTerm{hnipr|a} \dicIPA{{\textsubring{n}}{\textsci}{\textlengthmark}{\textsubring{b}}{r}{a}} \dicPos{v}[1] \dicFlx{(‑aði)}[1] \dicPhraseIS{hnipra sig saman} \dicDirectTranslationCS{schoulit se, stočit se do klubíčka}
\dicEntry[hnit] \dicTerm{hnit} \dicIPA{{\textsubring{n}}{\textsci}{\textlengthmark}{\textsubring{d}}} \dicPos{n}[2] \dicFlx{(‑s, ‑)}[5] \textbf{1.} \dicFieldCat{sport.} \dicSynonym{badminton} \dicDirectTranslationCS{badminton}  \textbf{2.} \dicPhraseIS{hnit} \dicFlx{pl} \dicFieldCat{mat.} \dicDirectTranslationCS{souřadnice}
\dicEntry[hnitakerfi] \dicTerm{hnita··kerfi} \dicIPA{{\textsubring{n}}{\textsci}{\textlengthmark}{\textsubring{d}}{a}{c\smash{\textsuperscript{h}}}{\textepsilon}{r}{v}{\textsci}} \dicPos{n}[2] \dicFlx{(‑s, ‑)}[14] \dicFieldCat{mat.} \dicDirectTranslationCS{souřadnicový systém}
\dicEntry[hnitmiða] \dicTerm{hnit··mið|a} \dicIPA{{\textsubring{n}}{\textsci}{\textlengthmark}{\textsubring{d}}{m}{\textsci}{ð}{a}} \dicPos{v}[1] \dicFlx{(‑aði)}[1] \dicFlx{acc} \dicDirectTranslationCS{upřesnit, zaměřit, vymezit} \dicExampleIS{hnitmiða orðalag} \dicExampleCS{upřesnit vyjádření}
\dicEntry[hnitmiðaður] \dicTerm{hnit··mið·|aður} \dicIPA{{\textsubring{n}}{\textsci}{\textlengthmark}{\textsubring{d}}{m}{\textsci}{ð}{a}{ð}{\textscy}{\textsubring{r}}} \dicPos{adj}[3] \dicFlx{(f ‑uð)}[3] \dicDirectTranslationCS{výstižný, přesný}
\dicEntry[hnífapar] \dicTerm{hnífa··|par} \dicIPA{{\textsubring{n}}{i}{\textlengthmark}{v}{a}{p\smash{\textsuperscript{h}}}{a}{\textsubring{r}}} \dicPos{n}[2] \dicFlx{(‑pars, ‑pör)}[8] \dicDirectTranslationCS{příbor}
\dicEntry[hnífur] \dicTerm{hníf|ur} \dicsymFrequent\  \dicIPA{{\textsubring{n}}{i}{\textlengthmark}{v}{\textscy}{\textsubring{r}}} \dicPos{m}[6] \dicFlx{(‑s, ‑ar)}[22] \dicSynonym{bredda} \dicDirectTranslationCS{nůž} \dicExampleIS{beittur hnífur} \dicExampleCS{ostrý nůž};  \dicPhraseIS{eiga\,/\addthin hafa ekki til hnífs og skeiðar} \dicDirectTranslationCS{nemít na jídlo ani na oblečení};  \dicPhraseIS{það gengur ekki hnífurinn á milli e‑rra} \dicLangCat{přen.} \dicDirectTranslationCS{(kdo) je zajedno}
\dicEntry[hníg] \dicTerm{hníg} \dicIPA{{\textsubring{n}}{i}{\textlengthmark}{x}} \dicPos{v} \dicFlx{ind praes sg 1 pers} \dicLink{hníga}
\dicEntry[hníga] \dicTerm{hníga} \dicIPA{{\textsubring{n}}{i}{\textlengthmark}{\textbabygamma}{a}} \dicPos{v}[6] \dicFlx{(hníg, hneig, hnigum, hnigi, hnigið)}[79] \dicSynonym{síga} \dicDirectTranslationCS{(po)klesnout, klesat} \dicExampleIS{Sólin hnígur.} \dicExampleCS{Slunce klesá.}
\begin{xtolerant}{}{1pt}
\dicEntry[hnípinn] \dicTerm{hnípinn}\addthinS\dicIPA{{\textsubring{n}}{i}{\textlengthmark}{\textsubring{b}}{\textsci}{\textsubring{n}}}\addthinS\dicPos{adj}[6]\dicFlx{}[-2]\addthinS\dicSynonym{dapur} \dicDirectTranslationCS{sklíčený, skleslý, schlíplý}
\end{xtolerant}
\dicEntry[hnísa] \dicTerm{hnís|a} \dicIPA{{\textsubring{n}}{i}{\textlengthmark}{s}{a}} \dicPos{f}[1] \dicFlx{(‑u, ‑ur)}[7] \dicFieldCat{zool.} \dicDirectTranslationCS{sviňucha, sviňucha obecná} \textit{(l.~{\textLA{Phocoena phocoena}})}  \dicsymPhoto\ 
\dicFigure{ds_image_hnisa_0_2.jpg}{Hnísa}{Hnísa - Malene Thyssen, GFDL, CC BY-SA 2.5}
\dicEntry[hnjask] \dicTerm{hnjask} \dicIPA{{\textsubring{n}}{j}{a}{s}{\r{g}}} \dicPos{n}[2] \dicFlx{(‑s)}[2] \dicDirectTranslationCS{špatné\,/\addthin hrubé zacházení};  \dicPhraseIS{e‑að verður fyrir hnjaski} \dicDirectTranslationCS{s~(čím) je špatně zacházeno}
\dicEntry[hnjóðsyrði] \dicTerm{hnjóðs··yrði} \dicIPA{{\textsubring{n}}{j}{ou}{ð}{s}{\textsci}{r}{ð}{\textsci}} \dicPos{n}[2] \dicFlx{(‑s, ‑)}[14] \dicSynonym*{ámælisorð} \dicDirectTranslationCS{posměšek, posměšná\,/\addthin výsměšná poznámka}
\dicEntry[hnjóta] \dicTerm{hnjóta} \dicIPA{{\textsubring{n}}{j}{ou}{\textlengthmark}{\textsubring{d}}{a}} \dicPos{v}[6] \dicFlx{(hnýt, hnaut, hnutum, hnyti, hnotið)}[103] \textbf{1.} \dicSynonym{hrasa} \dicDirectTranslationCS{zakopnout, klopýtnout, brknout, škobrtnout};  \dicPhraseIS{hnjóta um e‑ð} \dicDirectTranslationCS{klopýtnout o~(co)} \dicExampleIS{hnjóta um þúfu} \dicExampleCS{klopýtnout o~drn trávy}  \textbf{2.} \dicDirectTranslationCS{narazit (na článek ap.)};  \dicPhraseIS{hnjóta um e‑ð} \dicDirectTranslationCS{všimnout si (čeho), narazit na (co)} \dicExampleIS{hnjóta um smásögu í blaði} \dicExampleCS{narazit na povídku v~novinách}
\dicEntry[hnjúkur] \dicTerm{hnjúk|ur} \dicIPA{{\textsubring{n}}{j}{u}{\textlengthmark}{\r{g}}{\textscy}{\textsubring{r}}} \dicPos{m}[6] \dicFlx{(‑s, ‑ar)}[24] \dicDirectTranslationCS{vrcholek, štít (hory ap.)}
\dicEntry[hnoð] \dicTerm{hnoð} \dicIPA{{\textsubring{n}}{\textopeno}{\textlengthmark}{\texttheta}} \dicPos{n}[2] \dicFlx{(‑s, ‑)}[5] \textbf{1.} \dicSynonym*{það að hnoða} \dicDirectTranslationCS{hnětení (těsta ap.)}  \textbf{2.} \dicSynonym*{það að hnjóða} \dicDirectTranslationCS{nýtování}  \textbf{3.} \dicSynonym*{hnoðnagli} \dicDirectTranslationCS{nýt}  \textbf{4.} \dicSynonym{leirburður} \dicDirectTranslationCS{nekvalitní poezie}
\dicEntry[hnoða] \dicTerm{hnoð|a} \dicIPA{{\textsubring{n}}{\textopeno}{\textlengthmark}{ð}{a}} \dicPos{v}[1] \dicFlx{(‑aði)}[1] \dicFlx{acc\,/\addthin dat} \textbf{1.} \dicFlx{acc} \dicDirectTranslationCS{hníst, válet (těsto ap.)} \dicExampleIS{hnoða deig} \dicExampleCS{hníst těsto}  \textbf{2.} \dicFlx{acc} \dicDirectTranslationCS{nýtovat}  \textbf{3.} \dicFlx{dat} \dicSynonym*{hnuðla} \dicDirectTranslationCS{zmuchlat}  \textbf{4.} \dicPhraseIS{hnoða saman e‑u} \dicDirectTranslationCS{sesmolit (co) (báseň ap.)}
\dicEntry[hnoðri] \dicTerm{hnoðr|i} \dicIPA{{\textsubring{n}}{\textopeno}{ð}{r}{\textsci}} \dicPos{m}[1] \dicFlx{(‑a, ‑ar)}[1] \textbf{1.} \dicDirectTranslationCS{chomáček, obláček} \dicExampleIS{ullarhnoðri} \dicExampleCS{chomáček vlny}  \textbf{2.} \dicFieldCat{bot.} \dicDirectTranslationCS{rozchodník} \textit{(l.~{\textLA{Sedum}})}  \rlap{\dicsymPhoto\ }
\dicFigure{ds_image_hnodri_0_1.jpg}{Hnoðri}{Hnoðri - Stanislav Krejčík, CC BY-SA 3.0}
\dicEntry[hnossgæti] \dicTerm{hnoss··gæti} \dicIPA{{\textsubring{n}}{\textopeno}{s}{\r{\textObardotlessj}}{a}{i}{\textsubring{d}}{\textsci}} \dicPos{n}[2] \dicFlx{(‑s)}[20] \dicDirectTranslationCS{pochoutka, delikatesa}
\dicEntry[hnotið] \dicTerm{hnotið} \dicIPA{{\textsubring{n}}{\textopeno}{\textlengthmark}{\textsubring{d}}{\textsci}{\texttheta}} \dicPos{v} \dicFlx{supin} \dicLink{hnjóta}
\dicEntry[hnotskurn] \dicTerm{hnot··skurn} \dicIPA{{\textsubring{n}}{\textopeno}{\textlengthmark}{\textsubring{d}}{s}{\r{g}}{\textscy}{r}{\textsubring{d}}{\textsubring{n}}} \dicPos{f}[7] \dicFlx{(‑ar)}[4] \dicPhraseIS{í hnotskurn} \dicFlx{adv} \dicDirectTranslationCS{v~kostce, ve stručnosti}
\dicEntry[hnugginn] \dicTerm{hnugginn} \dicIPA{{\textsubring{n}}{\textscy}{\r{\textObardotlessj}}{\textlengthmark}{\textsci}{\textsubring{n}}} \dicPos{adj}[6]\dicFlx{}[-2] \dicDirectTranslationCS{sklíčený, zasmušilý, deprimovaný, skleslý}
\dicEntry[hnupla] \dicTerm{hnupl|a} \dicIPA{{\textsubring{n}}{\textscy}{h}{\textsubring{b}}{l}{a}} \dicPos{v}[1] \dicFlx{(‑aði)}[1] \dicFlx{dat} \dicDirectTranslationCS{(u)krást, ukradnout, šlohnout} \dicExampleIS{hnupla bókinni úr búðinni} \dicExampleCS{ukrást knížku z~obchodu}
\dicEntry[hnusa] \dicTerm{hnus|a} \dicIPA{{\textsubring{n}}{\textscy}{\textlengthmark}{s}{a}} \dicPos{v}[1] \dicFlx{(‑aði)}[1] \dicDirectTranslationCS{(vy)čenichat, (vy)čmuchat} \dicExampleIS{hnusa af e‑m} \dicExampleCS{čenichat ke (komu)}
\dicEntry[hnussa] \dicTerm{hnuss|a} \dicIPA{{\textsubring{n}}{\textscy}{s}{\textlengthmark}{a}} \dicPos{v}[1] \dicFlx{(‑aði)}[48] \dicDirectTranslationCS{funět, supět};  \dicPhraseIS{það hnussar í e‑m} \dicFlx{impers} \dicDirectTranslationCS{(kdo) funí}
\dicEntry[hnutum] \dicTerm{hnutum} \dicIPA{{\textsubring{n}}{\textscy}{\textlengthmark}{\textsubring{d}}{\textscy}{\textsubring{m}}} \dicPos{v} \dicFlx{ind pf pl 1 pers} \dicLink{hnjóta}
\dicEntry[hnúajárn] \dicTerm{hnúa··járn} \dicIPA{{\textsubring{n}}{u}{\textlengthmark}{a}{j}{au}{r}{\textsubring{d}}{\textsubring{n}}} \dicPos{n}[2] \dicFlx{(‑s, ‑)}[5] \dicDirectTranslationCS{boxer} \dicIndirectTranslationCS{(zbraň)}
\dicEntry[hnúðkál] \dicTerm{hnúð··kál} \dicIPA{{\textsubring{n}}{u}{ð}{k\smash{\textsuperscript{h}}}{au}{\textsubring{l}}} \dicPos{n}[2] \dicFlx{(‑s)}[2] \dicFieldCat{bot.} \dicDirectTranslationCS{kedluben} \textit{(l.~{\textLA{Brassica oleracea var. gongylodes}})}
\dicEntry[hnúðsellerí] \dicTerm{hnúð··sellerí} \dicIPA{{\textsubring{n}}{u}{ð}{s}{\textepsilon}{l}{\textepsilon}{r}{i}} \dicPos{n}[2] \dicFlx{(‑s)}[2] \dicFieldCat{bot.} \dicDirectTranslationCS{celer, miřík celer, celer bulvový} \textit{(l.~{\textLA{Apium graveolens var. rapaceum}})}
\dicEntry[hnúður] \dicTerm{hnúð|ur} \dicIPA{{\textsubring{n}}{u}{\textlengthmark}{ð}{\textscy}{\textsubring{r}}} \dicPos{m}[6] \dicFlx{(‑s, ‑ar)}[24] \textbf{1.} \dicSynonym*{kryppa á baki} \dicDirectTranslationCS{hrb} \dicIndirectTranslationCS{(vyboulení na zádech lidí nebo hřbetu zvířat)}  \textbf{2.} \dicFieldCat{bot.} \dicSynonym*{hnýði} \dicDirectTranslationCS{hlíza}
\dicEntry[hnúfubakur] \dicTerm{hnúfu··bak|ur} \dicIPA{{\textsubring{n}}{u}{\textlengthmark}{v}{\textscy}{\textsubring{b}}{a}{\r{g}}{\textscy}{\textsubring{r}}} \dicPos{m}[6] \dicFlx{(‑s, ‑ar)}[25] \dicFieldCat{zool.} \dicDirectTranslationCS{keporkak, plejtvák dlouhoploutvý} \textit{(l.~{\textLA{Megaptera novaeangliae}})}
\dicEntry[hnúi] \dicTerm{hnú|i} \dicIPA{{\textsubring{n}}{u}{\textlengthmark}{\textsci}} \dicPos{m}[1] \dicFlx{(‑a, ‑ar)}[1] \dicDirectTranslationCS{kloub} \dicIndirectTranslationCS{(na prstě)};  \dicPhraseIS{með hnúum og hnefum} \dicFlx{adv} \dicLangCat{přen.} \dicDirectTranslationCS{ze všech sil}
\dicEntry[hnúta] \dicTerm{hnút|a} \dicIPA{{\textsubring{n}}{u}{\textlengthmark}{\textsubring{d}}{a}} \dicPos{f}[1] \dicFlx{(‑u, ‑ur)}[7] \dicSynonym{sneið\smash{\textsuperscript{1}}} \dicDirectTranslationCS{špička, pichlavá narážka}
\dicEntry[hnútur] \dicTerm{hnút|ur} \dicsymFrequent\  \dicIPA{{\textsubring{n}}{u}{\textlengthmark}{\textsubring{d}}{\textscy}{\textsubring{r}}} \dicPos{m}[6] \dicFlx{(‑s, ‑ar)}[15] \textbf{1.} \dicDirectTranslationCS{uzel (na provázku ap.)};  \dicPhraseIS{hnýta\,/\addthin binda hnút} \dicDirectTranslationCS{uvázat uzel};  \dicPhraseIS{leysa hnútinn} \dicDirectTranslationCS{rozvázat uzel}  \textbf{2.} \dicDirectTranslationCS{keltský uzel}  \textbf{3.} \dicDirectTranslationCS{uzlina, bulka}  \textbf{4.} \dicFieldCat{nám.} \dicDirectTranslationCS{uzel} \dicIndirectTranslationCS{(jednotka rychlosti plavidla)};  \dicIdiom{hnútur}{ \dicPhraseIS{höggva á hnútinn}} \dicLangCat{přen.} \dicDirectTranslationCS{přetnout uzel};  \dicPhraseIS{e‑að er komið í hnút} \dicLangCat{přen.} \dicDirectTranslationCS{(co) se zamotalo\,/\addthin zkomplikovalo}
\dicEntry[hnykill] \dicTerm{hnyk|ill} \dicIPA{{\textsubring{n}}{\textsci}{\textlengthmark}{\r{\textObardotlessj}}{\textsci}{\textsubring{d}}{\textsubring{l}}} \dicPos{m}[6] \dicFlx{(‑ils, ‑lar)}[35] \textbf{1.} \dicDirectTranslationCS{klubko, klubíčko (vlny ap.)};  \dicPhraseIS{vinda band í hnykil} \dicDirectTranslationCS{svinout provázek do klubíčka}  \textbf{2.} \dicFieldCat{anat.} \dicDirectTranslationCS{malý mozek, mozeček}
\dicEntry[hnykkja] \dicTerm{hnykk|ja} \dicsymFrequent\  \dicIPA{{\textsubring{n}}{\textsci}{h}{\r{\textObardotlessj}}{\textlengthmark}{a}} \dicPos{v}[2] \dicFlx{(‑ti, ‑t)}[30] \dicFlx{dat} \textbf{1.} \dicSynonym{kippa\smash{\textsuperscript{2}}} \dicDirectTranslationCS{trhnout, škubnout, cuknout}  \textbf{2.} \dicPhraseIS{e‑m hnykkir við} \dicFlx{impers} \dicDirectTranslationCS{(kdo) se leká, (kdo) sebou trhá} \dicExampleIS{Öllum hnykkti við.} \dicExampleCS{Všichni se vylekali.};  \dicIdiom{hnykkja}[á]{ \dicPhraseIS{hnykkja á e‑u}} \dicDirectTranslationCS{připomenout (co), upozornit na (co)}
\dicEntry[hnykkur] \dicTerm{hnykk|ur} \dicIPA{{\textsubring{n}}{\textsci}{h}{\r{g}}{\textscy}{\textsubring{r}}} \dicPos{m}[9] \dicFlx{(‑s, ‑ir)}[15] \dicDirectTranslationCS{cuknutí, škubnutí, trhnutí}
\dicEntry[hnykla] \dicTerm{hnykl|a} \dicIPA{{\textsubring{n}}{\textsci}{h}{\r{g}}{l}{a}} \dicPos{v}[1] \dicFlx{(‑aði)}[1] \dicFlx{acc} \dicPhraseIS{hnykla brýnnar\,/\addthin brýrnar} \dicDirectTranslationCS{svraštit obočí};  \dicPhraseIS{hnykla vöðvana} \dicDirectTranslationCS{napnout svaly}
\dicEntry[hnyti] \dicTerm{hnyti} \dicIPA{{\textsubring{n}}{\textsci}{\textlengthmark}{\textsubring{d}}{\textsci}} \dicPos{v} \dicFlx{con pf sg 1 pers} \dicLink{hnjóta}
\dicEntry[hnyttilegur] \dicTerm{hnytti··legur} \dicIPA{{\textsubring{n}}{\textsci}{h}{\textsubring{d}}{\textsci}{l}{\textepsilon}{\textbabygamma}{\textscy}{\textsubring{r}}} \dicPos{adj}[1]\dicFlx{}[-8] \dicDirectTranslationCS{trefný, pádný (odpověď ap.)}
\dicEntry[hnyttinn] \dicTerm{hnyttinn} \dicIPA{{\textsubring{n}}{\textsci}{h}{\textsubring{d}}{\textsci}{\textsubring{n}}} \dicPos{adj}[6]\dicFlx{}[-2] \dicSynonym{fyndinn} \dicDirectTranslationCS{trefný, pohotový} \dicExampleIS{hnyttinn í tilsvörum} \dicExampleCS{pohotový v~odpovědích}
\dicEntry[hnyttni] \dicTerm{hnyttn|i} \dicIPA{{\textsubring{n}}{\textsci}{h}{\textsubring{d}}{n}{\textsci}} \dicPos{f}[3] \dicFlx{(‑i)}[3] \dicDirectTranslationCS{(slovní) pohotovost, pádnost (odpovědi ap.)}
\dicEntry[hnýsast] \dicTerm{hnýs|ast} \dicIPA{{\textsubring{n}}{i}{\textlengthmark}{s}{a}{s}{\textsubring{d}}} \dicPos{v}[2] \dicFlx{(‑tist, ‑t)}[193] \dicFlx{refl} \dicPhraseIS{hnýsast í e‑ð} \dicDirectTranslationCS{slídit v~(čem), strkat nos do (čeho)}
\dicEntry[hnýsilegur] \dicTerm{hnýsi··legur} \dicIPA{{\textsubring{n}}{i}{\textlengthmark}{s}{\textsci}{l}{\textepsilon}{\textbabygamma}{\textscy}{\textsubring{r}}} \dicPos{adj}[1]\dicFlx{}[-8] \dicSynonym{forvitnilegur} \dicDirectTranslationCS{pozoruhodný, zajímavý}
\dicEntry[hnýsinn] \dicTerm{hnýsinn} \dicIPA{{\textsubring{n}}{i}{\textlengthmark}{s}{\textsci}{\textsubring{n}}} \dicPos{adj}[6]\dicFlx{}[-2] \dicSynonym{forvitinn} \dicDirectTranslationCS{všetečný, (nadmíru) zvědavý}
\dicEntry[hnýsni] \dicTerm{hnýsn|i} \dicIPA{{\textsubring{n}}{i}{s}{\textsubring{d}}{n}{\textsci}} \dicPos{f}[3] \dicFlx{(‑i)}[3] \dicDirectTranslationCS{všetečnost, (nadměrná) zvědavost}
\dicEntry[hnýt] \dicTerm{hnýt} \dicIPA{{\textsubring{n}}{i}{\textlengthmark}{\textsubring{d}}} \dicPos{v} \dicFlx{ind praes sg 1 pers} \dicLink{hnjóta}
\dicEntry[hnýta] \dicTerm{hnýt|a} \dicsymFrequent\  \dicIPA{{\textsubring{n}}{i}{\textlengthmark}{\textsubring{d}}{a}} \dicPos{v}[2] \dicFlx{(‑ti, ‑t)}[54] \dicFlx{acc} \dicSynonym{binda} \dicDirectTranslationCS{(za)vázat, uvázat, přivázat} \dicExampleIS{hnýta slaufu á e‑ð} \dicExampleCS{uvázat mašli na (co)};  \dicIdiom{hnýta}[í]{ \dicPhraseIS{hnýta í e‑n}} \dicDirectTranslationCS{navážet se do (koho), dobírat si (koho)}
\dicEntry[hnökri] \dicTerm{hnökr|i} \dicIPA{{\textsubring{n}}{\oe}{\textlengthmark}{\r{g}}{r}{\textsci}} \dicPos{m}[1] \dicFlx{(‑a, ‑ar)}[1] \textbf{1.} \dicDirectTranslationCS{uzlík, žmolek (vlny ap.)} \dicExampleIS{Það eru hnökrar í ullinni.} \dicExampleCS{Ve vlně jsou uzlíky.}  \textbf{2.} \dicSynonym{galli\smash{\textsuperscript{2}}} \dicDirectTranslationCS{zádrhel, nedostatek (ve verši ap.)}
\dicEntry[hnöttóttur] \dicTerm{hnött··óttur} \dicIPA{{\textsubring{n}}{\oe}{h}{\textsubring{d}}{ou}{h}{\textsubring{d}}{\textscy}{\textsubring{r}}} \dicPos{adj}[1]\dicFlx{}[-10] \dicDirectTranslationCS{kulatý} \dicExampleIS{Jörðin er hnöttótt.} \dicExampleCS{Země je kulatá.}
\dicEntry[hnöttur] \dicTerm{hnöttur} \dicsymFrequent\  \dicIPA{{\textsubring{n}}{\oe}{h}{\textsubring{d}}{\textscy}{\textsubring{r}}} \dicPos{m}[11] \dicFlx{(hnattar, hnettir)}[5] \textbf{1.} \dicSynonym*{jarðarhnöttur} \dicDirectTranslationCS{zeměkoule, Země} \dicExampleIS{ferð umhverfis hnöttinn} \dicExampleCS{cesta kolem Země}  \textbf{2.} \dicSynonym*{himinhnöttur} \dicDirectTranslationCS{planeta} \dicExampleIS{ferð til annarra hnatta} \dicExampleCS{cesta na jiné planety}
\dicEntry[hobbí] \dicTerm{hobbí} \dicIPA{{h}{\textopeno}{\textsubring{b}}{\textlengthmark}{i}} \dicPos{n}[2] \dicFlx{(‑s, ‑)}[5] \dicLangCat{hovor.} \dicSynonym*{áhugaefni} \dicDirectTranslationCS{koníček, záliba, hobby}
\dicEntry[hof] \dicTerm{hof} \dicIPA{{h}{\textopeno}{\textlengthmark}{f}} \dicPos{n}[2] \dicFlx{(‑s, ‑)}[5] \dicDirectTranslationCS{chrám, svatyně}
\dicEntry[hoggið] \dicTerm{hoggið} \dicIPA{{h}{\textopeno}{\r{\textObardotlessj}}{\textlengthmark}{\textsci}{\texttheta}} \dicPos{v} \dicFlx{supin} \dicLink{höggva}
\dicEntry[hokinn] \dicTerm{hokinn} \dicIPA{{h}{\textopeno}{\textlengthmark}{\r{\textObardotlessj}}{\textsci}{\textsubring{n}}} \dicPos{adj}[6]\dicFlx{}[-2] \dicDirectTranslationCS{shrbený, nahrbený, sehnutý}
\dicEntry[hokkí] \dicTerm{hokkí} \dicIPA{{h}{\textopeno}{h}{\r{\textObardotlessj}}{i}} \dicPos{n}[2] \dicFlx{(‑s)}[33] \dicFieldCat{sport.} \dicDirectTranslationCS{hokej}
\dicEntry[hol] \dicTerm{hol} \dicIPA{{h}{\textopeno}{\textlengthmark}{\textsubring{l}}} \dicPos{n}[2] \dicFlx{(‑s, ‑)}[5] \textbf{1.} \dicSynonym{gap} \dicDirectTranslationCS{dutina}  \textbf{2.} \dicFieldCat{anat.} \dicSynonym*{holrúm líkamans} \dicDirectTranslationCS{(tělní) dutina}  \textbf{3.} \dicSynonym*{innri forstofa} \dicDirectTranslationCS{předsíň}
\dicEntry[hola] \dicTerm{hol|a\smash{\textsuperscript{1}}} \dicsymFrequent\  \dicIPA{{h}{\textopeno}{\textlengthmark}{l}{a}} \dicPos{f}[1] \dicFlx{(‑u, ‑ur)}[19] \textbf{1.} \dicSynonym{gryfja} \dicDirectTranslationCS{díra, jáma} \dicExampleIS{dýpka holu} \dicExampleCS{vyhloubit jámu}  \textbf{2.} \dicFieldCat{sport.} \dicDirectTranslationCS{(golfová) jamka};  \dicPhraseIS{fara\,/\addthin slá holu í höggi} \dicFieldCat{sport.} \dicDirectTranslationCS{zahrát hole-in-one}
\dicEntry[hola] \dicTerm{hol|a\smash{\textsuperscript{2}}} \dicIPA{{h}{\textopeno}{\textlengthmark}{l}{a}} \dicPos{v}[1] \dicFlx{(‑aði)}[1] \dicFlx{acc} \dicDirectTranslationCS{(vy)hloubit (příkop ap.)};  \dicIdiom{hola}[niður]{ \dicPhraseIS{hola sér niður}} \dicDirectTranslationCS{nacpat se, vtěsnat se}; { \dicPhraseIS{hola e‑u niður}} \dicDirectTranslationCS{zakopat (co)}
\dicEntry[hold] \dicTerm{hold} \dicsymFrequent\  \dicIPA{{h}{\textopeno}{l}{\textsubring{d}}} \dicPos{n}[2] \dicFlx{(‑s, ‑)}[5] \dicDirectTranslationCS{maso} \dicIndirectTranslationCS{(na těle)} \dicExampleIS{hold og blóð} \dicExampleCS{maso a~krev};  \dicPhraseIS{e‑m rís hold} \dicFlx{impers} \dicLangCat{přen.} \dicDirectTranslationCS{(komu) stojí, (kdo) má erekci};  \dicPhraseIS{holdi klæddur} \dicDirectTranslationCS{vtělený, ztělesněný};  \dicPhraseIS{safna holdum} \dicDirectTranslationCS{přibírat na váze};  \dicPhraseIS{vera í góðum holdum} \dicDirectTranslationCS{být při těle}
\dicEntry[holdakjúklingur] \dicTerm{holda··kjúkl·ing|ur} \dicIPA{{h}{\textopeno}{l}{\textsubring{d}}{a}{c\smash{\textsuperscript{h}}}{u}{h}{\r{g}}{l}{i}{\ng}{\r{g}}{\textscy}{\textsubring{r}}} \dicPos{m}[6] \dicFlx{(‑s, ‑ar)}[8] \dicDirectTranslationCS{brojler, chovné kuře}
\dicEntry[holdanaut] \dicTerm{holda··naut} \dicIPA{{h}{\textopeno}{l}{\textsubring{d}}{a}{n}{\oe i}{\textsubring{d}}} \dicPos{n}[2] \dicFlx{(‑s, ‑)}[5] \dicDirectTranslationCS{skot, hovězí dobytek}
\dicEntry[holdlegur] \dicTerm{hold··legur} \dicIPA{{h}{\textopeno}{l}{\textsubring{d}}{l}{\textepsilon}{\textbabygamma}{\textscy}{\textsubring{r}}} \dicPos{adj}[1]\dicFlx{}[-8] \dicDirectTranslationCS{tělesný, fyzický}
\dicEntry[holdsveiki] \dicTerm{holds··veik|i} \dicIPA{{h}{\textopeno}{l}{\textsubring{d}}{s}{v}{ei}{\r{\textObardotlessj}}{\textsci}} \dicPos{f}[3] \dicFlx{(‑i)}[3] \dicFieldCat{med.} \dicDirectTranslationCS{lepra, malomocenství}
\dicEntry[holdugur] \dicTerm{hold··ugur} \dicIPA{{h}{\textopeno}{l}{\textsubring{d}}{\textscy}{\textbabygamma}{\textscy}{\textsubring{r}}} \dicPos{adj}[1]\dicFlx{}[-8] \dicDirectTranslationCS{korpulentní, tělnatý, otylý}
\dicEntry[holdvotur] \dicTerm{hold··votur} \dicIPA{{h}{\textopeno}{l}{\textsubring{d}}{v}{\textopeno}{\textsubring{d}}{\textscy}{\textsubring{r}}} \dicPos{adj}[1]\dicFlx{}[-6] \dicDirectTranslationCS{zmáčený, promáčený, mokrý na kost}
\dicEntry[Holland] \dicTerm{Hol··land} \dicIPA{{h}{\textopeno}{l}{\textlengthmark}{a}{n}{\textsubring{d}}} \dicPos{n}[2] \dicFlx{(‑s)}[4] \dicFieldCat{geog.} \dicDirectTranslationCS{Holandsko, Nizozemsko}
\dicEntry[Hollendingur] \dicTerm{Hol·lend··ing|ur} \dicIPA{{h}{\textopeno}{l}{\textlengthmark}{\textepsilon}{n}{\textsubring{d}}{i}{\ng}{\r{g}}{\textscy}{\textsubring{r}}} \dicPos{m}[6] \dicFlx{(‑s, ‑ar)}[8] \dicDirectTranslationCS{Holanďan(ka), Nizozemec, Nizozemka}
\dicEntry[hollenska] \dicTerm{hol··lensk|a} \dicIPA{{h}{\textopeno}{l}{\textlengthmark}{\textepsilon}{n}{s}{\r{g}}{a}} \dicPos{f}[1] \dicFlx{(‑u)}[5] \dicDirectTranslationCS{holandština, nizozemština}
\dicEntry[hollenskur] \dicTerm{hol··lenskur} \dicIPA{{h}{\textopeno}{l}{\textlengthmark}{\textepsilon}{n}{s}{\r{g}}{\textscy}{\textsubring{r}}} \dicPos{adj}[1]\dicFlx{}[-1] \dicDirectTranslationCS{holandský, nizozemský}
\dicEntry[hollráður] \dicTerm{holl··ráður} \dicIPA{{h}{\textopeno}{\textsubring{d}}{l}{r}{au}{ð}{\textscy}{\textsubring{r}}} \dicPos{adj}[2]\dicFlx{}[-6] \dicDirectTranslationCS{dávající dobré rady, nápomocný}
\dicEntry[hollur] \dicTerm{hollur} \dicIPA{{h}{\textopeno}{\textsubring{d}}{l}{\textscy}{\textsubring{r}}} \dicPos{adj}[1]\dicFlx{}[-1] \textbf{1.} \dicSynonym*{heilsusamlegur} \dicDirectTranslationCS{zdravý, prospěšný, blahodárný (jídlo ap.)} \dicExampleIS{hollur matur} \dicExampleCS{zdravé jídlo};  \dicPhraseIS{það er holt að (gera e‑ð)} \dicDirectTranslationCS{je zdravé ((co) dělat)}  \textbf{2.} \dicSynonym{trúr} \dicDirectTranslationCS{věrný, oddaný, loajální} \dicExampleIS{vera hollur e‑m} \dicExampleCS{být (komu) oddaný}
\dicEntry[hollusta] \dicTerm{holl··ust|a} \dicIPA{{h}{\textopeno}{\textsubring{d}}{l}{\textscy}{s}{\textsubring{d}}{a}} \dicPos{f}[1] \dicFlx{(‑u)}[5] \textbf{1.} \dicSynonym{heilnæmi} \dicDirectTranslationCS{prospěšnost, blahodárnost, zdravost}  \textbf{2.} \dicSynonym{tryggð} \dicDirectTranslationCS{věrnost, oddanost, loajalita} \dicExampleIS{játa konungi hollustu sína} \dicExampleCS{odpřisáhnout králi věrnost}
\dicEntry[hollvinur] \dicTerm{holl··vin|ur} \dicIPA{{h}{\textopeno}{\textsubring{d}}{l}{v}{\textsci}{n}{\textscy}{\textsubring{r}}} \dicPos{m}[10] \dicFlx{(‑ar, ‑ir)}[17] \textbf{1.} \dicSynonym*{góður vinur} \dicDirectTranslationCS{dobrý kamarád\,/\addthin přítel, dobrá kamarádka\,/\addthin přítelkyně}  \textbf{2.} \dicSynonym*{stuðningsaðili} \dicDirectTranslationCS{sympatizant(ka), příznivec, příznivkyně}
\dicEntry[holóttur] \dicTerm{hol··óttur} \dicIPA{{h}{\textopeno}{\textlengthmark}{l}{ou}{h}{\textsubring{d}}{\textscy}{\textsubring{r}}} \dicPos{adj}[1]\dicFlx{}[-10] \dicDirectTranslationCS{děravý, plný děr (cesta ap.)} \dicExampleIS{holóttur vegur} \dicExampleCS{silnice plná děr}
\dicEntry[holræsi] \dicTerm{hol··ræsi} \dicIPA{{h}{\textopeno}{l}{r}{a}{i}{s}{\textsci}} \dicPos{n}[2] \dicFlx{(‑s, ‑)}[14] \dicSynonym{lokræsi} \dicDirectTranslationCS{kanál, kanalizace}
\dicEntry[holsjá] \dicTerm{hol··sjá} \dicIPA{{h}{\textopeno}{l}{s}{j}{au}} \dicPos{f}[4] \dicFlx{(‑r, ‑r)}[18] \dicFieldCat{med.} \dicDirectTranslationCS{endoskop}
\dicEntry[holskefla] \dicTerm{hol··skefl|a} \dicIPA{{h}{\textopeno}{l}{s}{\r{\textObardotlessj}}{\textepsilon}{\textsubring{b}}{l}{a}} \dicPos{f}[1] \dicFlx{(‑u, ‑ur)}[7] \textbf{1.} \dicSynonym{alda} \dicDirectTranslationCS{vysoká vlna}  \textbf{2.} \dicSynonym*{mikið magn} \dicDirectTranslationCS{značné množství, spousta}
\dicEntry[holt] \dicTerm{holt} \dicIPA{{h}{\textopeno}{\textsubring{l}}{\textsubring{d}}} \dicPos{n}[2] \dicFlx{(‑s, ‑)}[5] \dicDirectTranslationCS{pahorek};  \dicProverb\  \dicPhraseIS{Oft er í holti heyrandi nær.} \dicLangCat{přís.} \dicDirectTranslationCS{I~stěny mají uši.}
\dicEntry[holtasóley] \dicTerm{holta··sól·ey} \dicIPA{{h}{\textopeno}{\textsubring{l}}{\textsubring{d}}{a}{s}{ou}{l}{ei}} \dicPos{f}[4] \dicFlx{(‑jar, ‑jar)}[16] \dicFieldCat{bot.} \dicDirectTranslationCS{dryádka osmiplátečná} \textit{(l.~{\textLA{Dryas octopetala}})}  \dicsymPhoto\ 
\dicFigure{82307.jpg}{Holtasóley}{Holtasóley - Pleva František, Biolib, PD}
\dicEntry[holur] \dicTerm{holur} \dicIPA{{h}{\textopeno}{\textlengthmark}{l}{\textscy}{\textsubring{r}}} \dicPos{adj}[1]\dicFlx{}[-1] \dicDirectTranslationCS{dutý} \dicExampleIS{holur steinn} \dicExampleCS{dutý kámen}
\dicEntry[holurt] \dicTerm{hol··urt} \dicIPA{{h}{\textopeno}{\textlengthmark}{l}{\textscy}{\textsubring{r}}{\textsubring{d}}} \dicPos{f}[7] \dicFlx{(‑ar, ‑ir)}[1] \dicFieldCat{bot.} \dicDirectTranslationCS{silenka přímořská} \textit{(l.~{\textLA{Silene uniflora}})}  \dicsymPhoto\ 
\dicFigure{77782.jpg}{Holurt}{Holurt - Hroneš Michal, Biolib, Copyright/CC-BY-NC}
\dicEntry[hommi] \dicTerm{homm|i} \dicIPA{{h}{\textopeno}{m}{\textlengthmark}{\textsci}} \dicPos{m}[1] \dicFlx{(‑a, ‑ar)}[1] \dicLangCat{hovor.} \dicDirectTranslationCS{homosexuál, gay}
\dicEntry[Hondúras] \dicTerm{Hondúras} \dicIPA{{h}{\textopeno}{n}{\textsubring{d}}{u}{r}{a}{s}} \dicPos{n}[4] \dicFlx{indecl}[2] \dicFieldCat{geog.} \dicDirectTranslationCS{Honduras}
\dicEntry[Hondúri] \dicTerm{Hondúr|i} \dicIPA{{h}{\textopeno}{n}{\textsubring{d}}{u}{r}{\textsci}} \dicPos{m}[1] \dicFlx{(‑a, ‑ar)}[1] \dicDirectTranslationCS{Hondurasan(ka)}
\dicEntry[hondúrskur] \dicTerm{hondúrskur} \dicIPA{{h}{\textopeno}{n}{\textsubring{d}}{u}{\textsubring{r}}{s}{\r{g}}{\textscy}{\textsubring{r}}} \dicPos{adj}[1]\dicFlx{}[-6] \dicDirectTranslationCS{honduraský}
\dicEntry[Hong Kong] \dicTerm{Hong Kong} \dicIPA{{h}{\textopeno}{\ng}{\r{g}}{}{k\smash{\textsuperscript{h}}}{\textopeno}{\ng}{\r{g}}} \dicPos{subs} \dicFlx{indecl} \dicFieldCat{geog.} \dicDirectTranslationCS{Hongkong}
\dicEntry[hongkongskur] \dicTerm{hongkongskur} \dicIPA{{h}{\textopeno}{\ng}{\r{g}}{k\smash{\textsuperscript{h}}}{\textopeno}{\ng}{s}{\r{g}}{\textscy}{\textsubring{r}}} \dicPos{adj}[1]\dicFlx{}[-1] \dicDirectTranslationCS{hongkongský}
\dicEntry[honum] \dicTerm{honum} \dicIPA{{h}{\textopeno}{\textlengthmark}{n}{\textscy}{\textsubring{m}}} \dicPos{pron}[17] \dicFlx{pers} \dicFlx{sg dat} \dicLink{hann}
\dicEntry[hopa] \dicTerm{hop|a} \dicIPA{{h}{\textopeno}{\textlengthmark}{\textsubring{b}}{a}} \dicPos{v}[1] \dicFlx{(‑aði)}[44] \dicSynonym{hörfa} \dicDirectTranslationCS{ustoupit, ustupovat} \dicIndirectTranslationCS{(udělat krok zpět)};  \dicPhraseIS{hopa á hæl\,/\addthin hæli} \dicDirectTranslationCS{couvat}
\dicEntry[hopp] \dicTerm{hopp} \dicIPA{{h}{\textopeno}{h}{\textsubring{b}}} \dicPos{n}[2] \dicFlx{(‑s, ‑)}[5] \dicDirectTranslationCS{skákání, hopsání}
\dicEntry[hoppa] \dicTerm{hopp|a} \dicsymFrequent\  \dicIPA{{h}{\textopeno}{h}{\textsubring{b}}{a}} \dicPos{v}[1] \dicFlx{(‑aði)}[44] \dicDirectTranslationCS{skákat, poskakovat, hopsat} \dicExampleIS{hoppa yfir steininn} \dicExampleCS{skákat přes kámen}
\dicEntry[hor] \dicTerm{hor} \dicIPA{{h}{\textopeno}{\textlengthmark}{\textsubring{r}}} \dicPos{m}[4] \dicFlx{(‑s)}[18] \textbf{1.} \dicSynonym{megurð} \dicDirectTranslationCS{vyzáblost, vychrtlost};  \dicPhraseIS{deyja úr hor} \dicDirectTranslationCS{umřít hladem}  \textbf{2.} \dicSynonym{slím} \dicDirectTranslationCS{sopel, nudle}
\dicEntry[horaður] \dicTerm{hor··|aður} \dicIPA{{h}{\textopeno}{\textlengthmark}{r}{a}{ð}{\textscy}{\textsubring{r}}} \dicPos{adj}[3] \dicFlx{(f ‑uð)}[3] \dicDirectTranslationCS{vyzáblý, vychrtlý}
\dicEntry[horblaðka] \dicTerm{hor··|blaðka} \dicIPA{{h}{\textopeno}{r}{\textsubring{b}}{l}{a}{\texttheta}{\r{g}}{a}} \dicPos{f}[1] \dicFlx{(‑blöðku, ‑blöðkur)}[20] \dicFieldCat{bot.} \dicDirectTranslationCS{vachta trojlistá} \textit{(l.~{\textLA{Menyanthes trifoliata}})}  \dicsymPhoto\ 
\dicFigure{69605.jpg}{Horblaðka}{Horblaðka - Krejčík Stanislav, Biolib, Copyright/CC-BY}
\dicEntry[horf] \dicTerm{horf} \dicsymFrequent\  \dicIPA{{h}{\textopeno}{r}{f}} \dicPos{n}[2] \dicFlx{(‑s, ‑)}[5] \textbf{1.} \dicSynonym{stefna\smash{\textsuperscript{1}}} \dicDirectTranslationCS{směr, kurs} \dicExampleIS{halda í horfi} \dicExampleCS{držet směr}  \textbf{2.} \dicSynonym{ástand} \dicDirectTranslationCS{stav} \dicExampleIS{gott horf á e‑u} \dicExampleCS{dobrý stav (čeho)}  \textbf{3.} \dicSynonym{afstaða} \dicDirectTranslationCS{postoj, pohled} \dicExampleIS{horf við e‑u} \dicExampleCS{postoj k~(čemu)}  \textbf{4.} \dicFieldCat{jaz.} \dicDirectTranslationCS{(slovesný) vid, aspekt};  \dicPhraseIS{lokið horf} \dicFieldCat{jaz.} \dicDirectTranslationCS{dokonavý vid} \dicIndirectTranslationCS{(slovesný vid označující dokonavost děje vyjádřený v~islandštině dvěma způsoby: a) pomocné sloveso \clqq hafa\crqq  + supinum, b) sloveso \clqq vera\crqq  + \clqq búinn\crqq  + \clqq að\crqq  + infinitiv slovesa)};  \dicPhraseIS{ólokið horf} \dicFieldCat{jaz.} \dicDirectTranslationCS{nedokonavý vid} \dicIndirectTranslationCS{(slovesný vid označující nedokonavost děje)}
\dicEntry[horfa] \dicTerm{horf|a} \dicsymFrequent\  \dicIPA{{h}{\textopeno}{r}{v}{a}} \dicPos{v}[2] \dicFlx{(‑ði, ‑t)}[107] \textbf{1.} \dicSynonym{líta} \dicDirectTranslationCS{dívat se, sledovat, hledět} \dicExampleIS{horfa til himins} \dicExampleCS{dívat se do nebe};  \dicPhraseIS{horfa á e‑n} \dicDirectTranslationCS{sledovat (koho), dívat se na (koho)}  \textbf{2.} \dicSynonym{stefna\smash{\textsuperscript{2}}} \dicDirectTranslationCS{směřovat, být orientován, ležet} \dicExampleIS{Glugginn horfir í vestur.} \dicExampleCS{Okno směřuje na západ.}  \textbf{3.} \dicPhraseIS{það horfir vel\,/\addthin illa} \dicFlx{impers} \dicDirectTranslationCS{vypadá to dobře\,/\addthin špatně, jsou dobré\,/\addthin špatné vyhlídky};  \dicIdiom{horfa}[í]{ \dicPhraseIS{horfa í e‑ð}} \dicDirectTranslationCS{přihlédnout k~(čemu)};  \dicIdiom{horfa}[upp á]{ \dicPhraseIS{horfa upp á e‑ð}} \dicDirectTranslationCS{být svědkem (čeho), dívat se na (co)};  \dicIdiom{horfa}[við]{ \dicPhraseIS{e‑að horfir (þannig) við e‑m}} \dicDirectTranslationCS{(co) se (komu) jeví (nějak)}; { \dicPhraseIS{ef e‑m býður svo við að horfa}} \dicFlx{impers} \dicDirectTranslationCS{jestli se (komu) zachce, jestli se to (komu) hodí};  \dicIdiom{horfast}{ \dicPhraseIS{horfast}} \dicFlx{refl} \dicDirectTranslationCS{dívat se (na sebe), sledovat se}; { \dicPhraseIS{horfast í augu}} \dicFlx{refl} \dicDirectTranslationCS{dívat se do očí} \dicExampleIS{Þau horfðust lengi í augu.} \dicExampleCS{Dívali se dlouho do očí.}; { \dicPhraseIS{horfast í augu við e‑n}} \dicFlx{refl} \dicDirectTranslationCS{podívat se (komu) do očí}; { \dicPhraseIS{horfast í augu við e‑ð}} \dicFlx{refl} \dicLangCat{přen.} \dicDirectTranslationCS{podívat se (čemu) do očí (nebezpečí ap.)};  \dicIdiom{horfast}[á]{ \dicPhraseIS{það fór betur en á horfðist}} \dicFlx{impers} \dicDirectTranslationCS{šlo to lépe, než se očekávalo}
\dicEntry[horfið] \dicTerm{horfið} \dicIPA{{h}{\textopeno}{r}{v}{\textsci}{\texttheta}} \dicPos{v} \dicFlx{supin} \dicLink{hverfa\smash{\textsuperscript{2}}}
\dicEntry[horfinn] \dicTerm{horfinn} \dicIPA{{h}{\textopeno}{r}{v}{\textsci}{\textsubring{n}}} \dicPos{adj}[6]\dicFlx{}[-6] \dicDirectTranslationCS{zmizelý} \dicExampleIS{Hann er horfinn.} \dicExampleCS{Zmizel.}
\dicEntry[horfur] \dicTerm{horfur} \dicIPA{{h}{\textopeno}{r}{v}{\textscy}{\textsubring{r}}} \dicPos{f}[1] \dicFlx{pl}[12] \dicDirectTranslationCS{vyhlídky, výhled};  \dicPhraseIS{það er horfur á að (e‑að gerist)} \dicDirectTranslationCS{jsou vyhlídky na to, že (se (co) stane)}
\dicEntry[hormón] \dicTerm{hormón} \dicIPA{{h}{\textopeno}{r}{m}{ou}{\textsubring{n}}} \dicPos{n}[2] \dicFlx{(‑s, ‑)}[5] \dicFieldCat{med.} \dicDirectTranslationCS{hormon}
\dicEntry[horn] \dicTerm{horn} \dicIPA{{h}{\textopeno}{r}{\textsubring{d}}{\textsubring{n}}} \dicPos{n}[2] \dicFlx{(‑s, ‑)}[5] \textbf{1.} \dicSynonym{götuhorn} \dicDirectTranslationCS{roh (ulice ap.)} \dicExampleIS{götuhorn} \dicExampleCS{roh ulice}  \textbf{2.} \dicDirectTranslationCS{roh, kout} \dicExampleIS{sitja úti í horni} \dicExampleCS{sedět v~rohu}  \textbf{3.} \dicSynonym*{dýrahorn} \dicDirectTranslationCS{roh (zvířete ap.)}  \textbf{4.} \dicFieldCat{hud.} \dicDirectTranslationCS{roh}  \textbf{5.} \dicFieldCat{mat.} \dicDirectTranslationCS{úhel};  \dicPhraseIS{rétt horn} \dicFieldCat{mat.} \dicDirectTranslationCS{pravý úhel};  \dicPhraseIS{gleitt horn} \dicFieldCat{mat.} \dicDirectTranslationCS{tupý úhel};  \dicPhraseIS{hvasst horn} \dicFieldCat{mat.} \dicDirectTranslationCS{ostrý úhel}  \textbf{6.} \dicFieldCat{sport.} \dicDirectTranslationCS{roh};  \dicIdiom{horn}{ \dicPhraseIS{eiga hauk í horni}} \dicLangCat{přen.} \dicDirectTranslationCS{mít eso v~rukávu};  \dicPhraseIS{e‑r er harður í horn að taka} \dicLangCat{přen.} \dicDirectTranslationCS{s~(kým) je těžké pořízení};  \dicPhraseIS{hafa horn í síðu e‑rs} \dicLangCat{přen.} \dicDirectTranslationCS{mít výhrady vůči (komu)};  \dicPhraseIS{það skiptir í tvö horn} \dicLangCat{přen.} \dicDirectTranslationCS{existují dva tábory} \dicIndirectTranslationCS{(lidé jsou názorově rozdělení)}
\dicEntry[hornabolti] \dicTerm{horna··bolt|i} \dicIPA{{h}{\textopeno}{r}{\textsubring{d}}{n}{a}{\textsubring{b}}{\textopeno}{\textsubring{l}}{\textsubring{d}}{\textsci}} \dicPos{m}[1] \dicFlx{(‑a, ‑ar)}[1] \dicFieldCat{sport.} \dicSynonym{hafnabolti} \dicDirectTranslationCS{baseball}
\dicEntry[hornafall] \dicTerm{horna··fall} \dicIPA{{h}{\textopeno}{r}{\textsubring{d}}{n}{a}{f}{a}{\textsubring{d}}{\textsubring{l}}} \dicPos{n}[2] \dicFlx{(‑s)}[2] \dicFieldCat{mat.} \dicDirectTranslationCS{trigonometrická funkce}
\dicEntry[hornafræði] \dicTerm{horna··fræð|i} \dicIPA{{h}{\textopeno}{r}{\textsubring{d}}{n}{a}{f}{r}{a}{i}{ð}{\textsci}} \dicPos{f}[3] \dicFlx{(‑i)}[3] \dicFieldCat{mat.} \dicDirectTranslationCS{trigonometrie}
\dicEntry[hornauga] \dicTerm{horn··aug|a} \dicIPA{{h}{\textopeno}{r}{\textsubring{d}}{n}{\oe i}{\textbabygamma}{a}} \dicPos{n}[1] \dicFlx{(‑a)}[4] \dicPhraseIS{gefa e‑m hornauga} \dicLangCat{přen.} \dicDirectTranslationCS{uhnout\,/\addthin uhýbat (komu) pohledem};  \dicPhraseIS{líta e‑n hornauga} \dicDirectTranslationCS{(po)dívat se na (koho) škaredě}
\dicEntry[hornfiskur] \dicTerm{horn··fisk|ur} \dicIPA{{h}{\textopeno}{r}{\textsubring{d}}{\textsubring{n}}{f}{\textsci}{s}{\r{g}}{\textscy}{\textsubring{r}}} \dicPos{m}[6] \dicFlx{(‑s, ‑ar)}[10] \dicFieldCat{zool.} \dicDirectTranslationCS{jehlice rohozobá} \textit{(l.~{\textLA{Belone belone}})}
\dicEntry[hornhimna] \dicTerm{horn··himn|a} \dicIPA{{h}{\textopeno}{r}{\textsubring{d}}{\textsubring{n}}{h}{\textsci}{m}{n}{a}} \dicPos{f}[1] \dicFlx{(‑u, ‑ur)}[7] \dicFieldCat{anat.} \dicDirectTranslationCS{rohovka}
\dicEntry[Hornhöfði] \dicTerm{Horn··höfð|i} \dicIPA{{h}{\textopeno}{r}{\textsubring{d}}{\textsubring{n}}{h}{\oe}{v}{ð}{\textsci}} \dicPos{m}[1] \dicFlx{(‑a)}[5] \dicFieldCat{geog.} \dicDirectTranslationCS{mys Horn}
\dicEntry[hornklofi] \dicTerm{horn··klof|i} \dicIPA{{h}{\textopeno}{r}{\textsubring{d}}{\textsubring{n}}{k\smash{\textsuperscript{h}}}{l}{\textopeno}{v}{\textsci}} \dicPos{m}[1] \dicFlx{(‑a, ‑ar)}[1] \dicFieldCat{jaz.} \dicDirectTranslationCS{hranatá závorka} \dicIndirectTranslationCS{(interpunkční znaménko)}
\dicEntry[hornréttur] \dicTerm{horn··réttur} \dicIPA{{h}{\textopeno}{r}{\textsubring{d}}{n}{r}{j}{\textepsilon}{h}{\textsubring{d}}{\textscy}{\textsubring{r}}} \dicPos{adj}[1]\dicFlx{}[-10] \dicDirectTranslationCS{pravoúhlý}
\dicEntry[hornsíli] \dicTerm{horn··síli} \dicIPA{{h}{\textopeno}{r}{\textsubring{d}}{\textsubring{n}}{s}{i}{l}{\textsci}} \dicPos{n}[2] \dicFlx{(‑s, ‑)}[14] \dicFieldCat{zool.} \dicDirectTranslationCS{koljuška, koljuška tříostná} \textit{(l.~{\textLA{Gasterosteus aculeatus}})}  \dicsymPhoto\ 
\dicFigure{ds_image_hornsili_0_2.jpg}{Hornsíli}{Hornsíli - Ron Offermans, GFDL}
\dicEntry[hornsteinn] \dicTerm{horn··stein|n} \dicIPA{{h}{\textopeno}{r}{\textsubring{d}}{\textsubring{n}}{s}{\textsubring{d}}{ei}{\textsubring{d}}{\textsubring{n}}} \dicPos{m}[6] \dicFlx{(‑s, ‑ar)}[42] \dicSynonym{undirstaða} \dicDirectTranslationCS{základní kámen, základ};  \dicPhraseIS{leggja hornstein(inn) að e‑u} \dicLangCat{přen.} \dicDirectTranslationCS{položit základní kámen (čeho)}
\dicEntry[hortensía] \dicTerm{hortensí|a}\dicTerm{, hindarblóm} \dicIPA{{h}\-{\textopeno}\-{\textsubring{r}}\-{\textsubring{d}}\-{\textepsilon}\-{n}\-{s}\-{i}\-{j}\-{a}\-} \dicPos{f}[1] \dicFlx{(‑u, ‑ur)}[7] \dicFieldCat{bot.} \dicDirectTranslationCS{hortenzie} \textit{(l.~{\textLA{Hydrangea}})}  \dicsymPhoto\ 
\dicFigure{ds_image_hortensia_0_1.jpg}{Hortensía}{Hortensía - Nono vlf, CC BY-SA 3.0}
\dicEntry[hortittur] \dicTerm{hor··titt|ur} \dicIPA{{h}{\textopeno}{\textsubring{r}}{t\smash{\textsuperscript{h}}}{\textsci}{h}{\textsubring{d}}{\textscy}{\textsubring{r}}} \dicPos{m}[9] \dicFlx{(‑s, ‑ir)}[4] \textbf{1.} \dicFieldCat{techn.} \dicDirectTranslationCS{vycpávka}  \textbf{2.} \dicFieldCat{jaz.} \dicDirectTranslationCS{slovní vycpávka}
\dicEntry[hortugur] \dicTerm{hort··ugur} \dicIPA{{h}{\textopeno}{\textsubring{r}}{\textsubring{d}}{\textscy}{\textbabygamma}{\textscy}{\textsubring{r}}} \dicPos{adj}[1]\dicFlx{}[-8] \dicSynonym{ósvífinn} \dicDirectTranslationCS{drzý, troufalý}
\dicEntry[hossa] \dicTerm{hoss|a} \dicIPA{{h}{\textopeno}{s}{\textlengthmark}{a}} \dicPos{v}[1] \dicFlx{(‑aði)}[22] \dicFlx{dat} \dicDirectTranslationCS{kodrcat, nadhazovat (dítě na kolenou ap.)};  \dicIdiom{hossast}{ \dicPhraseIS{hossast}} \dicFlx{refl} \dicSynonym*{hristast} \dicDirectTranslationCS{nadskakovat, kodrcat se}
\dicEntry[hóa] \dicTerm{hó|a} \dicIPA{{h}{ou}{\textlengthmark}{a}} \dicPos{v}[1] \dicFlx{(‑aði)}[1] \dicDirectTranslationCS{houknout, zavolat};  \dicPhraseIS{hóa í e‑n} \dicDirectTranslationCS{houknout\,/\addthin zavolat na (koho)}
\dicEntry[hóf] \dicTerm{hóf\smash{\textsuperscript{1}}} \dicsymFrequent\  \dicIPA{{h}{ou}{\textlengthmark}{f}} \dicPos{n}[2] \dicFlx{(‑s, ‑)}[5] \textbf{1.} \dicSynonym{hófsemd} \dicDirectTranslationCS{míra, mez};  \dicPhraseIS{allt er best í hófi} \dicDirectTranslationCS{všeho s~mírou};  \dicPhraseIS{e‑að keyrir úr hófi} \dicDirectTranslationCS{(co) překračuje meze};  \dicPhraseIS{hafa hóf á e‑ð} \dicDirectTranslationCS{stanovit meze v~(čem)};  \dicPhraseIS{í hófi} \dicFlx{adv} \dicDirectTranslationCS{s~mírou,  v~mezích} \dicExampleIS{neyta áfengis í hófi} \dicExampleCS{konzumovat alkohol s~mírou};  \dicPhraseIS{kunna sér hóf} \dicDirectTranslationCS{znát svou mez}  \textbf{2.} \dicSynonym{samkvæmi} \dicDirectTranslationCS{sešlost, večírek}
\dicEntry[hóf] \dicTerm{hóf\smash{\textsuperscript{2}}} \dicIPA{{h}{ou}{\textlengthmark}{f}} \dicPos{v} \dicFlx{ind pf sg 1 pers} \dicLink{hefja}
\dicEntry[hófadynur] \dicTerm{hófa··dyn|ur} \dicIPA{{h}{ou}{\textlengthmark}{v}{a}{\textsubring{d}}{\textsci}{n}{\textscy}{\textsubring{r}}} \dicPos{m}[6] \dicFlx{(‑s)}[21] \dicDirectTranslationCS{dusot kopyt}
\dicEntry[hófdrykkja] \dicTerm{hóf··drykkj|a} \dicIPA{{h}{ou}{v}{\textsubring{d}}{r}{\textsci}{h}{\r{\textObardotlessj}}{a}} \dicPos{f}[1] \dicFlx{(‑u)}[5] \dicDirectTranslationCS{umírněné pití (alkoholu)}
\dicEntry[hófdýr] \dicTerm{hóf··dýr} \dicIPA{{h}{ou}{v}{\textsubring{d}}{i}{\textsubring{r}}} \dicPos{n}[2] \dicFlx{(‑s, ‑)}[5] \dicFieldCat{zool.} \dicDirectTranslationCS{lichokopytník} \textit{(l.~{\textLA{Perissodactyla}})}
\dicEntry[hóflaus] \dicTerm{hóf··laus} \dicIPA{{h}{ou}{v}{l}{\oe i}{s}} \dicPos{adj}[5]\dicFlx{}[-1] \dicDirectTranslationCS{nepřiměřený, nemírný, neúměrný} \dicExampleIS{hóflaus fjáreyðsla} \dicExampleCS{nepřiměřené utrácení}
\dicEntry[hóflegur] \dicTerm{hóf··legur} \dicIPA{{h}{ou}{v}{l}{\textepsilon}{\textbabygamma}{\textscy}{\textsubring{r}}} \dicPos{adj}[1]\dicFlx{}[-8] \dicDirectTranslationCS{střídmý, umírněný} \dicExampleIS{hófleg víndrykkja} \dicExampleCS{umírněné pití alkoholu}
\dicEntry[hófsamur] \dicTerm{hóf··|samur} \dicIPA{{h}{ou}{f}{s}{a}{m}{\textscy}{\textsubring{r}}} \dicPos{adj}[1] \dicFlx{(f ‑söm)}[2] \dicDirectTranslationCS{umírněný, přiměřený}
\dicEntry[hófsemd] \dicTerm{hóf··semd} \dicIPA{{h}{ou}{f}{s}{\textepsilon}{m}{\textsubring{d}}} \dicPos{f}[7] \dicFlx{(‑ar)}[3] \dicDirectTranslationCS{střídmost, umírněnost}
\dicEntry[hófsemi] \dicTerm{hóf··sem|i} \dicIPA{{h}{ou}{f}{s}{\textepsilon}{m}{\textsci}} \dicPos{f}[2] \dicFlx{(‑i)}[2] \dicDirectTranslationCS{střídmost, umírněnost}
\dicEntry[hófsóley] \dicTerm{hóf··sól·ey} \dicIPA{{h}{ou}{f}{s}{ou}{l}{ei}} \dicPos{f}[4] \dicFlx{(‑jar, ‑jar)}[16] \dicFieldCat{bot.} \dicDirectTranslationCS{blatouch bahenní} \textit{(l.~{\textLA{Caltha palustris}})}  \dicsymPhoto\ 
\dicFigure{922.jpg}{Hófsóley}{Hófsóley - Zicha Ondřej, Biolib, Copyright/CC-BY-NC}
\dicEntry[hófum] \dicTerm{hófum} \dicIPA{{h}{ou}{\textlengthmark}{v}{\textscy}{\textsubring{m}}} \dicPos{v} \dicFlx{ind pf pl 1 pers} \dicLink{hefja}
\dicEntry[hófur] \dicTerm{hóf|ur} \dicIPA{{h}{ou}{\textlengthmark}{v}{\textscy}{\textsubring{r}}} \dicPos{m}[6] \dicFlx{(‑s, ‑ar)}[8] \dicDirectTranslationCS{kopyto} \dicExampleIS{hófurinn á hestinum} \dicExampleCS{kopyto koně};  \dicPhraseIS{leita hófanna um e‑ð} \dicLangCat{přen.} \dicDirectTranslationCS{zkoumat možnosti (čeho)}
\dicEntry[hóglífi] \dicTerm{hóg··lífi} \dicIPA{{h}{ou}{\textbabygamma}{l}{i}{v}{\textsci}} \dicPos{n}[2] \dicFlx{(‑s)}[20] \dicDirectTranslationCS{poklidný život}
\dicEntry[hógvær] \dicTerm{hóg··vær} \dicIPA{{h}{ou}{\textbabygamma}{v}{a}{i}{\textsubring{r}}} \dicPos{adj}[5] \dicFlx{(f ‑)}[8] \dicSynonym{lítillátur} \dicDirectTranslationCS{skromný, pokorný}
\dicEntry[hógværð] \dicTerm{hóg··værð} \dicIPA{{h}{ou}{\textbabygamma}{v}{a}{i}{r}{\texttheta}} \dicPos{f}[7] \dicFlx{(‑ar)}[3] \dicDirectTranslationCS{skromnost, pokora}
\dicEntry[hól] \dicTerm{hól} \dicIPA{{h}{ou}{\textlengthmark}{\textsubring{l}}} \dicPos{n}[2] \dicFlx{(‑s)}[2] \dicSynonym{lof} \dicDirectTranslationCS{pochvala, lichotka}
\dicEntry[hólf] \dicTerm{hólf} \dicIPA{{h}{ou}{l}{f}} \dicPos{n}[2] \dicFlx{(‑s, ‑)}[5] \textbf{1.} \dicSynonym*{deild milli skilrúma} \dicDirectTranslationCS{oddělení, oddíl} \dicExampleIS{hólf í skáp} \dicExampleCS{oddělení ve skříni}  \textbf{2.} \dicSynonym{hirsla} \dicDirectTranslationCS{přihrádka, box} \dicExampleIS{pósthólf} \dicExampleCS{poštovní přihrádka}
\dicEntry[hólfa] \dicTerm{hólf|a} \dicIPA{{h}{ou}{l}{v}{a}} \dicPos{v}[1] \dicFlx{(‑aði)}[1] \dicFlx{acc} \dicDirectTranslationCS{oddělit (příčkou ap.), přepažit};  \dicIdiom{hólfa}[af]{ \dicPhraseIS{hólfa e‑ð af}} \dicDirectTranslationCS{oddělit (co) přepážkami};  \dicIdiom{hólfa}[niður]{ \dicPhraseIS{hólfa e‑ð niður\,/\addthin (í) sundur}} \dicDirectTranslationCS{rozdělit (co) (dům ap.)}
\dicEntry[hólkur] \dicTerm{hólk|ur} \dicIPA{{h}{ou}{\textsubring{l}}{\r{g}}{\textscy}{\textsubring{r}}} \dicPos{m}[6] \dicFlx{(‑s, ‑ar)}[24] \textbf{1.} \dicSynonym*{sívöl pípa} \dicDirectTranslationCS{(dutý) válec, cylindr}  \textbf{2.} \dicSynonym{átvagl} \dicDirectTranslationCS{žrout, jedlík}
\dicEntry[hóll] \dicTerm{hól|l} \dicsymFrequent\  \dicIPA{{h}{ou}{\textsubring{d}}{\textsubring{l}}} \dicPos{m}[6] \dicFlx{(‑s, ‑ar)}[51] \dicDirectTranslationCS{kopec, vrch} \dicExampleIS{grösugur hóll} \dicExampleCS{travnatý vrch}
\dicEntry[hólmganga] \dicTerm{hólm··|ganga} \dicIPA{{h}{ou}{l}{m}{\r{g}}{au}{\ng}{\r{g}}{a}} \dicPos{f}[1] \dicFlx{(‑göngu, ‑göngur)}[14] \dicDirectTranslationCS{duel, souboj}
\dicEntry[hólmi] \dicTerm{hólm|i} \dicIPA{{h}{ou}{l}{m}{\textsci}} \dicPos{m}[1] \dicFlx{(‑a, ‑ar)}[1] \dicDirectTranslationCS{ostrůvek, malý ostrov}
\dicEntry[hólmur] \dicTerm{hólm|ur} \dicsymFrequent\  \dicIPA{{h}{ou}{l}{m}{\textscy}{\textsubring{r}}} \dicPos{m}[6] \dicFlx{(‑s, ‑ar)}[4] \dicDirectTranslationCS{ostrůvek, malý ostrov};  \dicPhraseIS{fara með sigur af hólmi} \dicLangCat{přen.} \dicDirectTranslationCS{zvítězit, triumfovat};  \dicPhraseIS{leysa e‑n af hólmi} \dicDirectTranslationCS{nahradit (koho)} \dicExampleIS{Hún leysti hann af hólmi sem organisti í kirkjunni.} \dicExampleCS{Nahradila ho jako varhaník v~kostele.};  \dicPhraseIS{skora e‑n á hólm} \dicDirectTranslationCS{vyzvat (koho) na souboj}
\dicEntry[hólpinn] \dicTerm{hólpinn} \dicIPA{{h}{ou}{\textsubring{l}}{\textsubring{b}}{\textsci}{\textsubring{n}}} \dicPos{adj}[6]\dicFlx{}[-2] \dicDirectTranslationCS{zachráněný, spasený}
\dicEntry[hópferð] \dicTerm{hóp··ferð} \dicIPA{{h}{ou}{\textlengthmark}{\textsubring{b}}{f}{\textepsilon}{r}{\texttheta}} \dicPos{f}[7] \dicFlx{(‑ar, ‑ir)}[1] \dicDirectTranslationCS{skupinový zájezd\,/\addthin výlet}
\dicEntry[hópur] \dicTerm{hóp|ur} \dicsymFrequent\  \dicIPA{{h}{ou}{\textlengthmark}{\textsubring{b}}{\textscy}{\textsubring{r}}} \dicPos{m}[6] \dicFlx{(‑s, ‑ar)}[22] \textbf{1.} \dicSynonym{flokkur} \dicDirectTranslationCS{skupina, kolektiv, družina, parta} \dicExampleIS{Stór hópur neytir fíkniefna á Íslandi.} \dicExampleCS{Velká skupina lidí na Islandu užívá drogy.}  \textbf{2.} \dicDirectTranslationCS{stádo (ovcí ap.)};  \dicPhraseIS{e‑ir halda hópinn} \dicDirectTranslationCS{(kdo) drží spolu};  \dicPhraseIS{slást í hópinn} \dicFlx{refl} \dicDirectTranslationCS{začlenit se do skupiny};  \dicPhraseIS{upp til hópa} \dicFlx{adv} \dicDirectTranslationCS{co do jednoho}
\dicEntry[hópvinna] \dicTerm{hóp··vinn|a} \dicIPA{{h}{ou}{\textlengthmark}{\textsubring{b}}{v}{\textsci}{n}{a}} \dicPos{f}[1] \dicFlx{(‑u)}[5] \dicDirectTranslationCS{skupinová\,/\addthin týmová práce}
\dicEntry[hór] \dicTerm{hór} \dicIPA{{h}{ou}{\textlengthmark}{\textsubring{r}}} \dicPos{n}[2] \dicFlx{(‑s)}[2] \dicSynonym{hórdómur} \dicDirectTranslationCS{cizoložství, smilstvo};  \dicPhraseIS{drýgja hór} \dicDirectTranslationCS{cizoložit, smilnit}
\dicEntry[hóra] \dicTerm{hór|a} \dicIPA{{h}{ou}{\textlengthmark}{r}{a}} \dicPos{f}[1] \dicFlx{(‑u, ‑ur)}[7] \textbf{1.} \dicLangCat{hrub.} \dicDirectTranslationCS{děvka, kurva}  \textbf{2.} \dicLangCat{han.} \dicDirectTranslationCS{běhna, coura}
\dicEntry[hórdómur] \dicTerm{hór··dóm|ur} \dicIPA{{h}{ou}{r}{\textsubring{d}}{ou}{m}{\textscy}{\textsubring{r}}} \dicPos{m}[6] \dicFlx{(‑s)}[12] \dicSynonym{hór} \dicDirectTranslationCS{cizoložství, smilstvo}
\dicEntry[hórmangari] \dicTerm{hór··mang·ar|i} \dicIPA{{h}{ou}{r}{m}{au}{\ng}{\r{g}}{a}{r}{\textsci}} \dicPos{m}[1] \dicFlx{(‑a, ‑ar)}[10] \dicDirectTranslationCS{kuplíř(ka), pasák}
\dicEntry[hóruhús] \dicTerm{hóru··hús} \dicIPA{{h}{ou}{\textlengthmark}{r}{\textscy}{h}{u}{s}} \dicPos{n}[2] \dicFlx{(‑s, ‑)}[5] \dicDirectTranslationCS{nevěstinec, bordel}
\dicEntry[hósta] \dicTerm{hóst|a} \dicsymFrequent\  \dicIPA{{h}{ou}{s}{\textsubring{d}}{a}} \dicPos{v}[1] \dicFlx{(‑aði)}[22] \dicSynonym{kjölta} \dicDirectTranslationCS{(za)kašlat} \dicExampleIS{Í nótt hóstaði ég ekki neitt.} \dicExampleCS{V~noci jsem ani nezakašlal.}
\dicEntry[hóstakast] \dicTerm{hósta··|kast} \dicIPA{{h}{ou}{s}{\textsubring{d}}{a}{k\smash{\textsuperscript{h}}}{a}{s}{\textsubring{d}}} \dicPos{n}[2] \dicFlx{(‑kasts, ‑köst)}[8] \dicDirectTranslationCS{záchvat kašle}
\dicEntry[hóstamixtúra] \dicTerm{hósta··mixtúr|a} \dicIPA{{h}{ou}{s}{\textsubring{d}}{a}{m}{\textsci}{x}{s}{\textsubring{d}}{u}{r}{a}} \dicPos{f}[1] \dicFlx{(‑u, ‑ur)}[7] \dicDirectTranslationCS{sirup proti kašli}
\dicEntry[hósti] \dicTerm{hóst|i} \dicIPA{{h}{ou}{s}{\textsubring{d}}{\textsci}} \dicPos{m}[1] \dicFlx{(‑a, ‑ar)}[1] \dicDirectTranslationCS{kašel, (za)kašlání}
\dicEntry[hót] \dicTerm{hót} \dicIPA{{h}{ou}{\textlengthmark}{\textsubring{d}}} \dicPos{n}[2] \dicFlx{(‑s, ‑)}[5] \dicPhraseIS{ekki hót} \dicFlx{adv} \dicDirectTranslationCS{ani v~nejmenším};  \dicPhraseIS{ekki hætis hót} \dicFlx{adv} \dicDirectTranslationCS{absolutně nic, vůbec nic};  \dicPhraseIS{vera ekki hótinu betri en e‑n} \dicDirectTranslationCS{nebýt ani o~trochu lepší než (kdo)}
\dicEntry[hóta] \dicTerm{hót|a} \dicIPA{{h}{ou}{\textlengthmark}{\textsubring{d}}{a}} \dicPos{v}[1] \dicFlx{(‑aði)}[1] \dicFlx{dat + dat} \dicDirectTranslationCS{(po)hrozit, vyhrožovat};  \dicPhraseIS{hóta e‑m e‑u} \dicDirectTranslationCS{vyhrožovat (komu čím)};  \dicPhraseIS{hóta (því) að (gera e‑ð)} \dicDirectTranslationCS{hrozit ((co) udělat)} \dicExampleIS{Hann hótar að drepa hana.} \dicExampleCS{Vyhrožuje, že ji zabije.}
\dicEntry[hótel] \dicTerm{hótel} \dicsymFrequent\  \dicIPA{{h}{ou}{\textlengthmark}{t\smash{\textsuperscript{h}}}{\textepsilon}{\textsubring{l}}} \dicPos{n}[2] \dicFlx{(‑s, ‑)}[5] \dicDirectTranslationCS{hotel} \dicExampleIS{fjögurra stjörnu hótel} \dicExampleCS{čtyřhvězdičkový hotel}
\dicEntry[hótelherbergi] \dicTerm{hótel··her·bergi} \dicIPA{{h}{ou}{\textlengthmark}{\textsubring{d}}{\textepsilon}{l}{h}{\textepsilon}{r}{\textsubring{b}}{\textepsilon}{r}{\r{\textObardotlessj}}{\textsci}} \dicPos{n}[2] \dicFlx{(‑s, ‑)}[16] \dicDirectTranslationCS{hotelový pokoj}
\dicEntry[hótfyndinn] \dicTerm{hót··fyndinn} \dicIPA{{h}{ou}{\textlengthmark}{\textsubring{d}}{f}{\textsci}{n}{\textsubring{d}}{\textsci}{\textsubring{n}}} \dicPos{adj}[6]\dicFlx{}[-2] \dicDirectTranslationCS{sarkastický, jízlivý, uštěpačný}
\dicEntry[hótfyndni] \dicTerm{hót··fyndn|i} \dicIPA{{h}{ou}{\textlengthmark}{\textsubring{d}}{f}{\textsci}{n}{\textsubring{d}}{n}{\textsci}} \dicPos{f}[3] \dicFlx{(‑i)}[3] \dicDirectTranslationCS{sarkasmus, jízlivost, uštěpačnost}
\dicEntry[hótun] \dicTerm{hót|un} \dicIPA{{h}{ou}{\textlengthmark}{\textsubring{d}}{\textscy}{\textsubring{n}}} \dicPos{f}[7] \dicFlx{(‑unar, ‑anir)}[8] \dicDirectTranslationCS{hrozba, výhrůžka, pohrůžka} \dicExampleIS{hótun um árás} \dicExampleCS{hrozba napadení};  \dicPhraseIS{hafa í hótunum við e‑n} \dicDirectTranslationCS{vyhrožovat (komu)}
\dicEntry[hótunarbréf] \dicTerm{hótunar··bréf} \dicIPA{{h}{ou}{\textlengthmark}{\textsubring{d}}{\textscy}{n}{a}{r}{\textsubring{b}}{r}{j}{\textepsilon}{f}} \dicPos{n}[2] \dicFlx{(‑s, ‑)}[5] \dicDirectTranslationCS{výhružný dopis}
\dicEntry[hr.] \dicTerm{hr.} \dicPos{zkr} \dicPhraseIS{herra} \dicDirectTranslationCS{pán, pan}
\dicEntry[hrað] \dicTerm{hrað} \dicIPA{{\textsubring{r}}{a}{\textlengthmark}{\texttheta}} \dicPos{n}[2] \dicFlx{(‑s)}[2] \dicSynonym{flýtir} \dicDirectTranslationCS{spěch, kvap, chvat};  \dicPhraseIS{með hraði} \dicFlx{adv} \dicDirectTranslationCS{spěšně, kvapně, chvatně}
\dicEntry[hraða] \dicTerm{hrað|a} \dicsymFrequent\  \dicIPA{{\textsubring{r}}{a}{\textlengthmark}{ð}{a}} \dicPos{v}[1] \dicFlx{(‑aði)}[50] \dicFlx{dat} \dicDirectTranslationCS{zrychlit, urychlit, uspíšit} \dicExampleIS{hraða verkinu} \dicExampleCS{urychlit práci};  \dicPhraseIS{hraða sér} \dicDirectTranslationCS{pospíšit si, hnout sebou}
\dicEntry[hraðaaukning] \dicTerm{hraða··auk·ning} \dicIPA{{\textsubring{r}}{a}{\textlengthmark}{ð}{a}{\oe i}{h}{\r{g}}{n}{i}{\ng}{\r{g}}} \dicPos{f}[4] \dicFlx{(‑ar)}[7] \dicDirectTranslationCS{zrychlení, akcelerace, zvýšení rychlosti}
\dicEntry[hraðakstur] \dicTerm{hrað··akst|ur} \dicIPA{{\textsubring{r}}{a}{\textlengthmark}{ð}{a}{x}{s}{\textsubring{d}}{\textscy}{\textsubring{r}}} \dicPos{m}[5] \dicFlx{(‑urs)}[2] \dicDirectTranslationCS{rychlá jízda}
\dicEntry[hraðamælir] \dicTerm{hraða··mæl|ir} \dicIPA{{\textsubring{r}}{a}{\textlengthmark}{ð}{a}{m}{a}{i}{l}{\textsci}{\textsubring{r}}} \dicPos{m}[7] \dicFlx{(‑is, ‑ar)}[1] \dicDirectTranslationCS{rychloměr}
\dicEntry[hraðar] \dicTerm{hraðar} \dicIPA{{\textsubring{r}}{a}{\textlengthmark}{ð}{a}{\textsubring{r}}} \dicPos{adv} \dicFlx{comp} \dicLink{hratt}
\dicEntry[hraðara] \dicTerm{hraðara} \dicIPA{{\textsubring{r}}{a}{\textlengthmark}{ð}{a}{r}{a}} \dicPos{adv} \dicFlx{comp} \dicLink{hratt}
\dicEntry[hraðast] \dicTerm{hraðast} \dicIPA{{\textsubring{r}}{a}{\textlengthmark}{ð}{a}{s}{\textsubring{d}}} \dicPos{adv} \dicFlx{sup} \dicLink{hratt}
\dicEntry[hraðbanki] \dicTerm{hrað··bank|i} \dicIPA{{\textsubring{r}}{a}{ð}{\textsubring{b}}{au}{\r{\textltailn}}{\r{\textObardotlessj}}{\textsci}} \dicPos{m}[1] \dicFlx{(‑a, ‑ar)}[8] \dicDirectTranslationCS{bankomat, peněžní automat}
\dicEntry[hraðbátur] \dicTerm{hrað··bát|ur} \dicIPA{{\textsubring{r}}{a}{ð}{\textsubring{b}}{au}{\textsubring{d}}{\textscy}{\textsubring{r}}} \dicPos{m}[6] \dicFlx{(‑s, ‑ar)}[22] \dicFieldCat{nám.} \dicDirectTranslationCS{(rychlý) motorový člun}
\dicEntry[hraðberg] \dicTerm{hrað··berg} \dicIPA{{\textsubring{r}}{a}{ð}{\textsubring{b}}{\textepsilon}{r}{\r{g}}} \dicPos{n}[2] \dicFlx{(‑s)}[2] \dicPhraseIS{hafa e‑ð á hraðbergi} \dicLangCat{přen.} \dicDirectTranslationCS{mít (co) po ruce (trefnou odpověď ap.)}
\dicEntry[hraðbraut] \dicTerm{hrað··braut} \dicIPA{{\textsubring{r}}{a}{ð}{\textsubring{b}}{r}{\oe i}{\textsubring{d}}} \dicPos{f}[7] \dicFlx{(‑ar, ‑ir)}[1] \dicDirectTranslationCS{dálnice, rychlostní silnice}
\dicEntry[hraðbréf] \dicTerm{hrað··bréf} \dicIPA{{\textsubring{r}}{a}{ð}{\textsubring{b}}{r}{j}{\textepsilon}{f}} \dicPos{n}[2] \dicFlx{(‑s, ‑)}[5] \dicDirectTranslationCS{expresní dopis}
\dicEntry[hraðfrysta] \dicTerm{hrað··fryst|a} \dicIPA{{\textsubring{r}}{a}{ð}{f}{r}{\textsci}{s}{\textsubring{d}}{a}} \dicPos{v}[2] \dicFlx{(‑i, ‑)}[15] \dicFlx{acc} \dicDirectTranslationCS{(rychle) zamrazit}
\dicEntry[hraðfrystihús] \dicTerm{hrað·frysti··hús} \dicIPA{{\textsubring{r}}{a}{ð}{f}{r}{\textsci}{s}{\textsubring{d}}{\textsci}{h}{u}{s}} \dicPos{n}[2] \dicFlx{(‑s, ‑)}[5] \dicDirectTranslationCS{mrazírna}
\dicEntry[hraði] \dicTerm{hrað|i} \dicsymFrequent\  \dicIPA{{\textsubring{r}}{a}{\textlengthmark}{ð}{\textsci}} \dicPos{m}[1] \dicFlx{(‑a, ‑ar)}[8] \dicDirectTranslationCS{rychlost} \dicExampleIS{hraði nettengingar} \dicExampleCS{rychlost (internetového) připojení};  \dicPhraseIS{ólöglegur hraði} \dicDirectTranslationCS{nepovolená rychlost}
\dicEntry[hraðlest] \dicTerm{hrað··lest} \dicIPA{{\textsubring{r}}{a}{ð}{l}{\textepsilon}{s}{\textsubring{d}}} \dicPos{f}[7] \dicFlx{(‑ar, ‑ir)}[1] \dicDirectTranslationCS{rychlík} \dicIndirectTranslationCS{(druh vlaku)}
\dicEntry[hraðletur] \dicTerm{hrað··letur} \dicIPA{{\textsubring{r}}{a}{ð}{l}{\textepsilon}{\textsubring{d}}{\textscy}{\textsubring{r}}} \dicPos{n}[2] \dicFlx{(‑s, ‑)}[25] \dicDirectTranslationCS{těsnopisný znak}
\dicEntry[hraðrita] \dicTerm{hrað··rit|a} \dicIPA{{\textsubring{r}}{a}{ð}{r}{\textsci}{\textsubring{d}}{a}} \dicPos{v}[1] \dicFlx{(‑aði)}[1] \dicFlx{acc} \textbf{1.} \dicDirectTranslationCS{psát rychle}  \textbf{2.} \dicDirectTranslationCS{psát těsnopisem} \dicExampleIS{hraðrita ræðuna} \dicExampleCS{provést zápis proslovu těsnopisem}
\dicEntry[hraðritari] \dicTerm{hrað··rit·ar|i} \dicIPA{{\textsubring{r}}{a}{ð}{r}{\textsci}{\textsubring{d}}{a}{r}{\textsci}} \dicPos{m}[1] \dicFlx{(‑a, ‑ar)}[13] \dicDirectTranslationCS{těsnopisec, stenograf(ka)}
\dicEntry[hraðritun] \dicTerm{hrað··rit|un} \dicIPA{{\textsubring{r}}{a}{ð}{r}{\textsci}{\textsubring{d}}{\textscy}{\textsubring{n}}} \dicPos{f}[7] \dicFlx{(‑unar)}[9] \dicDirectTranslationCS{těsnopis, stenografie}
\dicEntry[hraðskreiður] \dicTerm{hrað··skreiður} \dicIPA{{\textsubring{r}}{a}{ð}{s}{\r{g}}{r}{ei}{ð}{\textscy}{\textsubring{r}}} \dicPos{adj}[2]\dicFlx{}[-6] \dicDirectTranslationCS{svižný, rychlý} \dicIndirectTranslationCS{(o~autě, letadle, lodi ap.)}
\dicEntry[hraðsuðuketill] \dicTerm{hrað·suðu··|ketill} \dicIPA{{\textsubring{r}}{a}{ð}{s}{\textscy}{ð}{\textscy}{c\smash{\textsuperscript{h}}}{\textepsilon}{\textsubring{d}}{\textsci}{\textsubring{d}}{\textsubring{l}}} \dicPos{m}[6] \dicFlx{(‑ketils, ‑katlar)}[79] \dicDirectTranslationCS{rychlovarná konvice}
\dicEntry[hraður] \dicTerm{hraður} \dicIPA{{\textsubring{r}}{a}{\textlengthmark}{ð}{\textscy}{\textsubring{r}}} \dicPos{adj}[2] \dicFlx{(f hröð)}[7] \dicSynonym{fljótur} \dicDirectTranslationCS{rychlý, spěšný} \dicExampleIS{hröð atburðarás} \dicExampleCS{rychlý sled událostí};  \dicPhraseIS{hafa hraðan á, hafa hraðar hendur} \dicDirectTranslationCS{pospíšit si} \dicIndirectTranslationCS{(vyvinout spěch při nějaké činnosti)}
\dicEntry[hrafl] \dicTerm{hrafl} \dicIPA{{\textsubring{r}}{a}{\textsubring{b}}{\textsubring{l}}} \dicPos{n}[2] \dicFlx{(‑s)}[2] \dicDirectTranslationCS{útržek, trocha};  \dicPhraseIS{kunna hrafl í e‑u} \dicDirectTranslationCS{umět trochu (co)} \dicExampleIS{Hann kann hrafl í ensku.} \dicExampleCS{Umí trochu anglicky.}
\dicEntry[hrafn] \dicTerm{hrafn} \dicIPA{{\textsubring{r}}{a}{\textsubring{b}}{\textsubring{n}}} \dicPos{m}[4] \dicFlx{(‑s, ‑ar)}[4] \dicFieldCat{zool.} \dicDirectTranslationCS{krkavec velký} \textit{(l.~{\textLA{Corvus corax}})}  \dicsymPhoto\ 
\dicFigure{6175.jpg}{Hrafn}{Hrafn - Deml Miroslav, Biolib, Copyright/CC-BY}
\dicEntry[hrafnafífa] \dicTerm{hrafna··fíf|a} \dicIPA{{\textsubring{r}}{a}{\textsubring{b}}{n}{a}{f}{i}{v}{a}} \dicPos{f}[1] \dicFlx{(‑u, ‑ur)}[19] \dicFieldCat{bot.} \dicDirectTranslationCS{suchopýr Scheuchzerův} \textit{(l.~{\textLA{Eriophorum scheuchzeri}})}  \dicsymPhoto\ 
\dicFigure{ds_image_hrafnafifa_0_2.jpg}{Hrafnafífa}{Hrafnafífa - Slaunger, GFDL}
\dicEntry[hrafnaklukka] \dicTerm{hrafna··klukk|a} \dicIPA{{\textsubring{r}}{a}{\textsubring{b}}{n}{a}{k\smash{\textsuperscript{h}}}{l}{\textscy}{h}{\r{g}}{a}} \dicPos{f}[1] \dicFlx{(‑u, ‑ur)}[13] \dicFieldCat{bot.} \dicIndirectTranslationCS{druh řeřišnice} \textit{(l.~{\textLA{Cardamine nymanii}})}
\dicEntry[hrafnreyður] \dicTerm{hrafn··reyð|ur} \dicIPA{{\textsubring{r}}{a}{\textsubring{b}}{n}{r}{ei}{ð}{\textscy}{\textsubring{r}}} \dicPos{f}[5] \dicFlx{(‑ar, ‑ar)}[1] \dicFieldCat{zool.} \dicSynonym{hrefna} \dicDirectTranslationCS{plejtvák malý} \textit{(l.~{\textLA{Balaenoptera acutorostrata}})}
\dicEntry[hrafnsönd] \dicTerm{hrafns··|önd} \dicIPA{{\textsubring{r}}{a}{\textsubring{b}}{\textsubring{n}}{s}{\oe}{n}{\textsubring{d}}} \dicPos{f}[8] \dicFlx{(‑andar, ‑endur)}[2] \dicFieldCat{zool.} \dicDirectTranslationCS{turpan černý} \textit{(l.~{\textLA{Melanitta nigra}})}
\dicEntry[hrafntinna] \dicTerm{hrafn··tinn|a} \dicIPA{{\textsubring{r}}{a}{\textsubring{b}}{\textsubring{n}}{t\smash{\textsuperscript{h}}}{\textsci}{n}{a}} \dicPos{f}[1] \dicFlx{(‑u, ‑ur)}[7] \dicFieldCat{geol.} \dicDirectTranslationCS{obsidián, sopečné sklo}
\dicEntry[hraka] \dicTerm{hrak|a} \dicIPA{{\textsubring{r}}{a}{\textlengthmark}{\r{g}}{a}} \dicPos{v}[1] \dicFlx{(‑aði)}[69] \dicFlx{impers} \dicPhraseIS{e‑m hrakar} \dicDirectTranslationCS{(komu) se daří hůř, (komu) se přitěžuje (zdraví ap.)};  \dicPhraseIS{e‑u hrakar} \dicDirectTranslationCS{(co) se zhoršuje}
\dicEntry[hrakfallabálkur] \dicTerm{hrak·falla··bálk|ur} \dicIPA{{\textsubring{r}}{a}{\textlengthmark}{\r{g}}{f}{a}{\textsubring{d}}{l}{a}{\textsubring{b}}{au}{\textsubring{l}}{\r{g}}{\textscy}{\textsubring{r}}} \dicPos{m}[6] \dicFlx{(‑s, ‑ar)}[24] \dicDirectTranslationCS{smolař(ka)}
\dicEntry[hrakhólar] \dicTerm{hrak··hólar} \dicIPA{{\textsubring{r}}{a}{\textlengthmark}{\r{g}}{h}{ou}{l}{a}{\textsubring{r}}} \dicPos{m}[6] \dicFlx{pl}[46] \dicPhraseIS{vera á hrakhólum} \dicLangCat{zast.} \dicDirectTranslationCS{mít potíže, být v~nesnázích}
\dicEntry[hrakið] \dicTerm{hrakið} \dicIPA{{\textsubring{r}}{a}{\textlengthmark}{\r{\textObardotlessj}}{\textsci}{\texttheta}} \dicPos{v} \dicFlx{supin} \dicLink{hrekja}
\dicEntry[hrakningur] \dicTerm{hrakn··ing|ur} \dicIPA{{\textsubring{r}}{a}{h}{\r{g}}{n}{i}{\ng}{\r{g}}{\textscy}{\textsubring{r}}} \dicPos{m}[6] \dicFlx{(‑s, ‑ar)}[8] \dicSynonym{flækingur} \dicDirectTranslationCS{tuláctví, potulování, potulka};  \dicPhraseIS{vera á hrakningi} \dicDirectTranslationCS{toulat se, potulovat se}
\dicEntry[hrakspá] \dicTerm{hrak··spá} \dicIPA{{\textsubring{r}}{a}{\textlengthmark}{\r{g}}{s}{\textsubring{b}}{au}} \dicPos{f}[4] \dicFlx{(‑r, ‑r)}[18] \dicDirectTranslationCS{katastrofická předpověď\,/\addthin prognóza}
\dicEntry[hrakti] \dicTerm{hrakti} \dicIPA{{\textsubring{r}}{a}{x}{\textsubring{d}}{\textsci}} \dicPos{v} \dicFlx{ind pf sg 1 pers} \dicLink{hrekja}
\dicEntry[hrakviðri] \dicTerm{hrak··viðri} \dicIPA{{\textsubring{r}}{a}{\textlengthmark}{\r{g}}{v}{\textsci}{ð}{r}{\textsci}} \dicPos{n}[2] \dicFlx{(‑s, ‑)}[14] \dicDirectTranslationCS{nečas, nepohoda, psota, slota}
\dicEntry[hrammur] \dicTerm{hramm|ur} \dicIPA{{\textsubring{r}}{a}{m}{\textlengthmark}{\textscy}{\textsubring{r}}} \dicPos{m}[6] \dicFlx{(‑s, ‑ar)}[5] \dicDirectTranslationCS{tlapa, tlapka} \dicExampleIS{hrammur á bjarndýri} \dicExampleCS{medvědí tlapa}
\dicEntry[hranalegur] \dicTerm{hrana··legur} \dicIPA{{\textsubring{r}}{a}{\textlengthmark}{n}{a}{l}{\textepsilon}{\textbabygamma}{\textscy}{\textsubring{r}}} \dicPos{adj}[1]\dicFlx{}[-8] \dicSynonym{hryssingslegur} \dicDirectTranslationCS{nevrlý, nepříjemný}
\dicEntry[hrannar] \dicTerm{hrannar} \dicIPA{{\textsubring{r}}{a}{n}{\textlengthmark}{a}{\textsubring{r}}} \dicPos{f} \dicFlx{sg gen} \dicLink{hrönn}
\dicEntry[hrannir] \dicTerm{hrannir} \dicIPA{{\textsubring{r}}{a}{n}{\textlengthmark}{\textsci}{\textsubring{r}}} \dicPos{f} \dicFlx{pl nom} \dicLink{hrönn}
\dicEntry[hrap] \dicTerm{hrap} \dicIPA{{\textsubring{r}}{a}{\textlengthmark}{\textsubring{b}}} \dicPos{n}[2] \dicFlx{(‑s, hröp)}[8] \dicSynonym{hrun} \dicDirectTranslationCS{pád, spadnutí, zřícení (na zem ap.)}
\dicEntry[hrapa] \dicTerm{hrap|a} \dicsymFrequent\  \dicIPA{{\textsubring{r}}{a}{\textlengthmark}{\textsubring{b}}{a}} \dicPos{v}[1] \dicFlx{(‑aði)}[13] \textbf{1.} \dicSynonym{falla} \dicDirectTranslationCS{spadnout, padat, sletět, zřítit se} \dicExampleIS{hrapa úr klettum} \dicExampleCS{spadnout ze skály}  \textbf{2.} \dicSynonym{lækka} \dicDirectTranslationCS{spadnout (ceny ap.)};  \dicIdiom{hrapa}[að]{ \dicPhraseIS{hrapa að e‑u}} \dicDirectTranslationCS{unáhlit se s~(čím)}
\dicEntry[hrapallegur] \dicTerm{hrapal··legur} \dicIPA{{\textsubring{r}}{a}{\textlengthmark}{\textsubring{b}}{a}{l}{\textepsilon}{\textbabygamma}{\textscy}{\textsubring{r}}} \dicPos{adj}[1]\dicFlx{}[-8] \dicSynonym{sviplegur} \dicDirectTranslationCS{fatální, katastrofální} \dicExampleIS{hrapallegt slys} \dicExampleCS{katastrofální nehoda}
\dicEntry[hrasa] \dicTerm{hras|a} \dicIPA{{\textsubring{r}}{a}{\textlengthmark}{s}{a}} \dicPos{v}[1] \dicFlx{(‑aði)}[13] \textbf{1.} \dicSynonym*{detta hálfvegis} \dicDirectTranslationCS{klopýtnout, zakopnout, škobrtnout} \dicExampleIS{hrasa um stein} \dicExampleCS{klopýtnout o~kámen}  \textbf{2.} \dicLangCat{přen.} \dicDirectTranslationCS{klopýtnout, škobrtnout} \dicIndirectTranslationCS{(dopustit se chyby)}
\dicEntry[hrat] \dicTerm{hrat} \dicIPA{{\textsubring{r}}{a}{\textlengthmark}{\textsubring{d}}} \dicPos{n}[2] \dicFlx{(‑s)}[2] \textbf{1.} \dicSynonym*{berjahrat} \dicDirectTranslationCS{(ovocný) odpad}  \textbf{2.} \dicSynonym*{fuglaskítur} \dicDirectTranslationCS{(ptačí) trus}
\dicEntry[hratt] \dicTerm{hratt\smash{\textsuperscript{1}}} \dicIPA{{\textsubring{r}}{a}{h}{\textsubring{d}}} \dicPos{v} \dicFlx{ind pf sg 1 pers} \dicLink{hrinda}
\dicEntry[hratt] \dicTerm{hratt\smash{\textsuperscript{2}}} \dicsymFrequent\  \dicIPA{{\textsubring{r}}{a}{h}{\textsubring{d}}} \dicPos{adv} \dicFlx{(comp hraðar\,/\addthin hraðara, sup hraðast)} \dicDirectTranslationCS{rychle, spěšně} \dicExampleIS{Upplýsingar úreldast hratt.} \dicExampleCS{Informace rychle zastarávají.}
\dicEntry[hraun] \dicTerm{hraun} \dicsymFrequent\  \dicIPA{{\textsubring{r}}{\oe i}{\textlengthmark}{\textsubring{n}}} \dicPos{n}[2] \dicFlx{(‑s, ‑)}[5] \textbf{1.} \dicFieldCat{geol.} \dicSynonym*{bráðið grjót} \dicDirectTranslationCS{láva, magma} \dicExampleIS{Hraun rennur úr gígnum.} \dicExampleCS{Z~kráteru vytéká láva.}  \textbf{2.} \dicSynonym*{hraunlendi} \dicDirectTranslationCS{láva, lávové pole}
\dicEntry[hraungjá] \dicTerm{hraun··gjá} \dicIPA{{\textsubring{r}}{\oe i}{n}{\r{\textObardotlessj}}{au}} \dicPos{f}[4] \dicFlx{(‑r, ‑r)}[18] \dicFieldCat{geol.} \dicDirectTranslationCS{lávová puklina}
\dicEntry[hraunkvika] \dicTerm{hraun··kvik|a} \dicIPA{{\textsubring{r}}{\oe i}{\ng}{k\smash{\textsuperscript{h}}}{v}{\textsci}{\r{g}}{a}} \dicPos{f}[1] \dicFlx{(‑u)}[5] \dicFieldCat{geol.} \dicDirectTranslationCS{magma}
\dicEntry[hraunrennsli] \dicTerm{hraun··rennsli} \dicIPA{{\textsubring{r}}{\oe i}{n}{r}{\textepsilon}{n}{s}{\textsubring{d}}{l}{\textsci}} \dicPos{n}[2] \dicFlx{(‑s, ‑)}[14] \dicFieldCat{geol.} \dicDirectTranslationCS{lávový proud}
\dicEntry[hraunstraumur] \dicTerm{hraun··straum|ur} \dicIPA{{\textsubring{r}}{\oe i}{n}{s}{\textsubring{d}}{r}{\oe i}{m}{\textscy}{\textsubring{r}}} \dicPos{m}[6] \dicFlx{(‑s, ‑ar)}[22] \dicFieldCat{geol.} \dicDirectTranslationCS{lávový proud}
\dicEntry[hrauntröð] \dicTerm{hraun··|tröð} \dicIPA{{\textsubring{r}}{\oe i}{n}{t\smash{\textsuperscript{h}}}{r}{\oe}{\texttheta}} \dicPos{f}[7] \dicFlx{(‑traðar, ‑traðir)}[16] \dicFieldCat{geol.} \dicDirectTranslationCS{lávový kanál}
\dicEntry[hraus] \dicTerm{hraus} \dicIPA{{\textsubring{r}}{\oe i}{\textlengthmark}{s}} \dicPos{v} \dicFlx{ind pf sg 1 pers} \dicLink{hrjósa}
\dicEntry[hraustur] \dicTerm{hraustur} \dicIPA{{\textsubring{r}}{\oe i}{s}{\textsubring{d}}{\textscy}{\textsubring{r}}} \dicPos{adj}[1]\dicFlx{}[-10] \textbf{1.} \dicSynonym{heilsugóður} \dicDirectTranslationCS{zdravý, vitální}  \textbf{2.} \dicSynonym{hugrakkur} \dicDirectTranslationCS{udatný, srdnatý, chrabrý}
\dicEntry[hraut] \dicTerm{hraut} \dicIPA{{\textsubring{r}}{\oe i}{\textlengthmark}{\textsubring{d}}} \dicPos{v} \dicFlx{ind pf sg 1 pers} \dicLink{hrjóta}
\dicEntry[hráefni] \dicTerm{hrá··efni} \dicIPA{{\textsubring{r}}{au}{\textlengthmark}{\textepsilon}{\textsubring{b}}{n}{\textsci}} \dicPos{n}[2] \dicFlx{(‑s, ‑)}[14] \dicDirectTranslationCS{surovina}
\dicEntry[hráki] \dicTerm{hrák|i} \dicIPA{{\textsubring{r}}{au}{\textlengthmark}{\r{\textObardotlessj}}{\textsci}} \dicPos{m}[1] \dicFlx{(‑a, ‑ar)}[1] \dicDirectTranslationCS{plivanec} \dicExampleIS{spýta hrákanum} \dicExampleCS{odplivnout si}
\dicEntry[hráolía] \dicTerm{hrá··olí|a} \dicIPA{{\textsubring{r}}{au}{\textlengthmark}{\textopeno}{l}{i}{j}{a}} \dicPos{f}[1] \dicFlx{(‑u)}[5] \dicDirectTranslationCS{ropa}
\dicEntry[hrár] \dicTerm{hrár} \dicIPA{{\textsubring{r}}{au}{\textlengthmark}{\textsubring{r}}} \dicPos{adj}[4]\dicFlx{}[-1] \textbf{1.} \dicSynonym{ósoðinn} \dicDirectTranslationCS{syrový} \dicIndirectTranslationCS{(zcela nebo částečně nezpracovaný, o~potravinách)} \dicExampleIS{hrár matur} \dicExampleCS{syrové jídlo}  \textbf{2.} \dicSynonym{ófullgerður} \dicDirectTranslationCS{hrubý, nedokončený}
\dicEntry[hráslagalegur] \dicTerm{hrá·slaga··legur} \dicIPA{{\textsubring{r}}{au}{s}{\textsubring{d}}{l}{a}{\textbabygamma}{a}{l}{\textepsilon}{\textbabygamma}{\textscy}{\textsubring{r}}} \dicPos{adj}[1]\dicFlx{}[-8] \dicDirectTranslationCS{syrový, sychravý} \dicExampleIS{hráslagalegur dagur} \dicExampleCS{sychravý den}
\dicEntry[hreðka] \dicTerm{hreðk|a} \dicIPA{{\textsubring{r}}{\textepsilon}{\texttheta}{\r{g}}{a}} \dicPos{f}[1] \dicFlx{(‑u, ‑ur)}[19] \dicFieldCat{bot.} \dicDirectTranslationCS{ředkvička, ředkev setá} \textit{(l.~{\textLA{Raphanus sativus}})}
\dicEntry[hrefna] \dicTerm{hrefn|a} \dicIPA{{\textsubring{r}}{\textepsilon}{\textsubring{b}}{n}{a}} \dicPos{f}[1] \dicFlx{(‑u, ‑ur)}[7] \textbf{1.} \dicFieldCat{zool.} \dicDirectTranslationCS{plejtvák malý} \textit{(l.~{\textLA{Balaenoptera acutorostrata}})}  \dicsymPhoto\   \textbf{2.} \dicSynonym*{kvenfugl hrafns} \dicDirectTranslationCS{samička krkavce}
\dicFigure{ds_image_hrefna_0_1.jpg}{Hrefna}{Hrefna - Mbz1, GFDL}
\dicEntry[hreiðra] \dicTerm{hreiðr|a} \dicIPA{{\textsubring{r}}{ei}{ð}{r}{a}} \dicPos{v}[1] \dicFlx{(‑aði)}[1] \dicPhraseIS{hreiðra sig} \dicDirectTranslationCS{uhnízdit se};  \dicPhraseIS{hreiðra um sig} \dicDirectTranslationCS{udělat si pohodlí, uvelebit se}
\dicEntry[hreiður] \dicTerm{hreiður} \dicIPA{{\textsubring{r}}{ei}{\textlengthmark}{ð}{\textscy}{\textsubring{r}}} \dicPos{n}[2] \dicFlx{(‑s, ‑)}[25] \dicSynonym{bæli} \dicDirectTranslationCS{hnízdo}
\dicEntry[hreif] \dicTerm{hreif} \dicIPA{{\textsubring{r}}{ei}{\textlengthmark}{f}} \dicPos{v} \dicFlx{ind pf sg 1 pers} \dicLink{hrífa\smash{\textsuperscript{2}}}
\dicEntry[hreifi] \dicTerm{hreif|i} \dicIPA{{\textsubring{r}}{ei}{\textlengthmark}{v}{\textsci}} \dicPos{m}[1] \dicFlx{(‑a, ‑ar)}[1] \dicDirectTranslationCS{zadní ploutev tuleně}
\dicEntry[hreifur] \dicTerm{hreifur} \dicIPA{{\textsubring{r}}{ei}{\textlengthmark}{v}{\textscy}{\textsubring{r}}} \dicPos{adj}[1]\dicFlx{}[-1] \textbf{1.} \dicSynonym{kátur} \dicDirectTranslationCS{veselý, rozjařený}  \textbf{2.} \dicSynonym{ölvaður} \dicDirectTranslationCS{veselý, mírně podnapilý} \dicExampleIS{vera hreifur af víni} \dicExampleCS{být lehce podnapilý}
\dicEntry[hreimur] \dicTerm{hreim|ur} \dicIPA{{\textsubring{r}}{ei}{\textlengthmark}{m}{\textscy}{\textsubring{r}}} \dicPos{m}[6] \dicFlx{(‑s, ‑ar)}[22] \dicSynonym{málrómur} \dicDirectTranslationCS{přízvuk, akcent} \dicExampleIS{tala með amerískum hreim} \dicExampleCS{mluvit s~americkým přízvukem}
\dicEntry[hrein] \dicTerm{hrein} \dicIPA{{\textsubring{r}}{ei}{\textlengthmark}{\textsubring{n}}} \dicPos{v} \dicFlx{ind pf sg 1 pers} \dicLink{hrína}
\dicEntry[hreindýr] \dicTerm{hrein··dýr} \dicIPA{{\textsubring{r}}{ei}{n}{\textsubring{d}}{i}{\textsubring{r}}} \dicPos{n}[2] \dicFlx{(‑s, ‑)}[5] \dicFieldCat{zool.} \dicDirectTranslationCS{sob, sob polární} \textit{(l.~{\textLA{Rangifer tarandus}})}  \dicsymPhoto\ 
\dicFigure{16462.jpg}{Hreindýr}{Hreindýr - Hanzlík Václav, Biolib, Copyright/CC-BY-NC}
\dicEntry[hreingerning] \dicTerm{hrein··ger·ning} \dicIPA{{\textsubring{r}}{ei}{\textltailn}{\r{\textObardotlessj}}{\textepsilon}{r}{\textsubring{d}}{n}{i}{\ng}{\r{g}}} \dicPos{f}[4] \dicFlx{(‑ar, ‑ar)}[5] \dicDirectTranslationCS{úklid, uklízení, čištění} \dicExampleIS{vinna við hreingerningar} \dicExampleCS{vydělávat si uklízením}
\dicEntry[hreinlega] \dicTerm{hrein··lega} \dicsymFrequent\  \dicIPA{{\textsubring{r}}{ei}{n}{l}{\textepsilon}{\textbabygamma}{a}} \dicPos{adv} \dicDirectTranslationCS{(zcela) zřejmě, jasně, očividně} \dicExampleIS{Þeir vilja hreinlega ekki tala við okkur.} \dicExampleCS{Očividně s~námi nechtějí mluvit.}
\dicEntry[hreinlegur] \dicTerm{hrein··legur} \dicIPA{{\textsubring{r}}{ei}{n}{l}{\textepsilon}{\textbabygamma}{\textscy}{\textsubring{r}}} \dicPos{adj}[1]\dicFlx{}[-8] \textbf{1.} \dicSynonym{hreinn} \dicDirectTranslationCS{čistý, čistotný}  \textbf{2.} \dicSynonym{ótvíræður} \dicDirectTranslationCS{očividný, zřejmý, zjevný, nepochybný}
\dicEntry[hreinleiki] \dicTerm{hrein··leik|i} \dicIPA{{\textsubring{r}}{ei}{n}{l}{ei}{\r{\textObardotlessj}}{\textsci}} \dicPos{m}[1] \dicFlx{(‑a)}[3] \textbf{1.} \dicSynonym*{það að e‑að er hreint} \dicDirectTranslationCS{čistota (oděvu ap.)}  \textbf{2.} \dicSynonym{sakleysi} \dicDirectTranslationCS{čistota, bezúhonnost}
\dicEntry[hreinlífi] \dicTerm{hrein··lífi} \dicIPA{{\textsubring{r}}{ei}{n}{l}{i}{v}{\textsci}} \dicPos{n}[2] \dicFlx{(‑s)}[20] \dicSynonym{skírlífi} \dicDirectTranslationCS{cudnost, mravní čistota, sexuální zdrženlivost}
\dicEntry[hreinlæti] \dicTerm{hrein··læti} \dicIPA{{\textsubring{r}}{ei}{n}{l}{a}{i}{\textsubring{d}}{\textsci}} \dicPos{n}[2] \dicFlx{(‑s)}[20] \dicDirectTranslationCS{čistota, hygiena} \dicExampleIS{gæta hreinlætis við e‑ð} \dicExampleCS{pečovat o~čistotu (čeho)}
\dicEntry[hreinn] \dicTerm{hreinn} \dicsymFrequent\  \dicIPA{{\textsubring{r}}{ei}{\textsubring{d}}{\textsubring{n}}} \dicPos{adj}[7]\dicFlx{}[-1] \textbf{1.} \dicSynonym*{óskítugur} \dicDirectTranslationCS{čistý, neposkvrněný}  \textbf{2.} \dicSynonym{bjartur} \dicDirectTranslationCS{čistý, jasný (obloha ap.)} \dicExampleIS{hreint veður} \dicExampleCS{jasno}  \textbf{3.} \dicSynonym{óblandaður} \dicDirectTranslationCS{čistý, ryzí} \dicIndirectTranslationCS{(neobsahující žádné příměsi)} \dicExampleIS{hreinn eplasafi} \dicExampleCS{100\%\,/\addthin stoprocentní jablečný džus};  \dicPhraseIS{hrein mey} \dicDirectTranslationCS{panna}  \textbf{4.} \dicFieldCat{ekon.} \dicDirectTranslationCS{čistý};  \dicPhraseIS{hreinar tekjur} \dicDirectTranslationCS{čisté zisky}  \textbf{5.} \dicSynonym{einber} \dicDirectTranslationCS{čistý, čirý, holý, úplný} \dicExampleIS{hreinn sannleikur} \dicExampleCS{čistá pravda}  \textbf{6.} \dicSynonym{greinilegur} \dicDirectTranslationCS{jasný, zřetelný} \dicExampleIS{hreint merki} \dicExampleCS{jasné znamení};  \dicIdiom{hreinn}{ \dicPhraseIS{gera hreint fyrir sínum dyrum}} \dicLangCat{přen.} \dicDirectTranslationCS{zamést si před vlastním prahem}; { \dicPhraseIS{vera hreinn og beinn}} \dicDirectTranslationCS{být otevřený\,/\addthin upřímný};  \dicPhraseIS{það er ekki á hreinu} \dicDirectTranslationCS{není jasné (kolik stojí jízdenka ap.)}
\dicEntry[hreinrit] \dicTerm{hrein··rit} \dicIPA{{\textsubring{r}}{ei}{n}{r}{\textsci}{\textsubring{d}}} \dicPos{n}[2] \dicFlx{(‑s, ‑)}[5] \dicDirectTranslationCS{čistopis}
\dicEntry[hreinræktaður] \dicTerm{hrein··rækt·|aður} \dicIPA{{\textsubring{r}}{ei}{n}{r}{a}{i}{x}{\textsubring{d}}{a}{ð}{\textscy}{\textsubring{r}}} \dicPos{adj}[3] \dicFlx{(f ‑uð)}[3] \textbf{1.} \dicDirectTranslationCS{plnokrevný, čistokrevný} \dicExampleIS{hreinræktaður hundur} \dicExampleCS{čistokrevný pes}  \textbf{2.} \dicSynonym{alger} \dicDirectTranslationCS{naprostý, úplný}
\dicEntry[hreinsa] \dicTerm{hreins|a} \dicsymFrequent\  \dicIPA{{\textsubring{r}}{ei}{n}{s}{a}} \dicPos{v}[1] \dicFlx{(‑aði)}[1] \dicFlx{acc} \dicDirectTranslationCS{(vy)čistit, očišťovat} \dicExampleIS{hreinsa ílátið} \dicExampleCS{vyčistit nádobu};  \dicPhraseIS{hreinsa e‑ð af e‑u} \dicDirectTranslationCS{očistit (co) od (čeho)};  \dicPhraseIS{láta hreinsa e‑ð} \dicDirectTranslationCS{nechat (co) vyčistit};  \dicIdiom{hreinsa}[til]{ \dicPhraseIS{hreinsa til}} \dicDirectTranslationCS{uklidit} \dicExampleIS{hreinsa til í skúrnum} \dicExampleCS{uklidit v~boudě};  \dicIdiom{hreinsast}{ \dicPhraseIS{hreinsast}} \dicFlx{refl} \dicDirectTranslationCS{očistit se, vyčistit se}
\dicEntry[hreinskilinn] \dicTerm{hrein··skilinn} \dicIPA{{\textsubring{r}}{ei}{n}{s}{\r{\textObardotlessj}}{\textsci}{l}{\textsci}{\textsubring{n}}} \dicPos{adj}[6]\dicFlx{}[-2] \dicDirectTranslationCS{upřímný, otevřený}
\dicEntry[hreinskilni] \dicTerm{hrein··skiln|i} \dicIPA{{\textsubring{r}}{ei}{n}{s}{\r{\textObardotlessj}}{\textsci}{l}{n}{\textsci}} \dicPos{f}[3] \dicFlx{(‑i)}[3] \dicDirectTranslationCS{upřímnost, otevřenost};  \dicPhraseIS{í hreinskilni sagt} \dicFlx{adv} \dicDirectTranslationCS{upřímně řečeno}
\dicEntry[hreinsun] \dicTerm{hreins|un} \dicIPA{{\textsubring{r}}{ei}{n}{s}{\textscy}{\textsubring{n}}} \dicPos{f}[7] \dicFlx{(‑unar, ‑anir)}[8] \textbf{1.} \dicDirectTranslationCS{(vy)čištění} \dicExampleIS{hreinsun á e‑u} \dicExampleCS{vyčištění (čeho)}  \textbf{2.} \dicDirectTranslationCS{rafinace}  \textbf{3.} \dicDirectTranslationCS{(politická) čistka} \dicExampleIS{pólitískar hreinsanir} \dicExampleCS{politické čistky}
\dicEntry[hreinsunareldur] \dicTerm{hreinsunar··eld|ur} \dicIPA{{\textsubring{r}}{ei}{n}{s}{\textscy}{n}{a}{r}{\textepsilon}{l}{\textsubring{d}}{\textscy}{\textsubring{r}}} \dicPos{m}[6] \dicFlx{(‑s)}[7] \dicFieldCat{náb.} \dicDirectTranslationCS{očistec, očistný oheň}
\dicEntry[hreint] \dicTerm{hrein|t} \dicsymFrequent\  \dicIPA{{\textsubring{r}}{ei}{\textsubring{n}}{\textsubring{d}}} \dicPos{adv} \dicFlx{(comp ‑na, sup ‑ast)} \dicSynonym{alveg} \dicDirectTranslationCS{dost(i), zcela, úplně, naprosto};  \dicPhraseIS{hreint ekki} \dicFlx{adv} \dicDirectTranslationCS{vůbec ne};  \dicPhraseIS{hreint og beint} \dicFlx{adv} \dicDirectTranslationCS{přímo řečeno, zkrátka, prostě};  \dicPhraseIS{hreint út sagt} \dicFlx{adv} \dicDirectTranslationCS{prostě, zkrátka, vysloveně, bez obalu} \dicExampleIS{Þetta var hreint út sagt hræðilegt.} \dicExampleCS{Bylo to prostě strašné.}
\dicEntry[hreintrúarstefna] \dicTerm{hrein·trúar··stefn|a} \dicIPA{{\textsubring{r}}{ei}{n}{t\smash{\textsuperscript{h}}}{\textsubring{r}}{u}{a}{\textsubring{r}}{s}{\textsubring{d}}{\textepsilon}{\textsubring{b}}{n}{a}} \dicPos{f}[1] \dicFlx{(‑u)}[5] \dicFieldCat{náb.} \dicDirectTranslationCS{puritánství}
\dicEntry[hreintungustefna] \dicTerm{hrein·tungu··stefn|a} \dicIPA{{\textsubring{r}}{ei}{n}{t\smash{\textsuperscript{h}}}{u}{\ng}{\r{g}}{\textscy}{s}{\textsubring{d}}{\textepsilon}{\textsubring{b}}{n}{a}} \dicPos{f}[1] \dicFlx{(‑u)}[5] \dicFieldCat{jaz.} \dicDirectTranslationCS{(jazykový) purismus}
\dicEntry[hreistur] \dicTerm{hreistur} \dicIPA{{\textsubring{r}}{ei}{s}{\textsubring{d}}{\textscy}{\textsubring{r}}} \dicPos{n}[2] \dicFlx{(‑s)}[28] \dicFieldCat{biol.} \dicDirectTranslationCS{šupina} \dicExampleIS{hreistrið á fiskinum} \dicExampleCS{rybí šupina}
\dicEntry[hrek] \dicTerm{hrek} \dicIPA{{\textsubring{r}}{\textepsilon}{\textlengthmark}{\r{g}}} \dicPos{v} \dicFlx{ind praes sg 1 pers} \dicLink{hrekja}
\dicEntry[hrekja] \dicTerm{hrekja} \dicsymFrequent\  \dicIPA{{\textsubring{r}}{\textepsilon}{\textlengthmark}{\r{\textObardotlessj}}{a}} \dicPos{v}[4] \dicFlx{(hrek, hrakti, hröktum, hrekti, hrakið)}[45] \dicFlx{acc} \textbf{1.} \dicSynonym{flæma} \dicDirectTranslationCS{zahnat, zahánět, odehnat, odhánět} \dicExampleIS{hrekja e‑n burt} \dicExampleCS{odehnat (koho) pryč}  \textbf{2.} \dicSynonym{afsanna} \dicDirectTranslationCS{vyvrátit, vyvracet} \dicIndirectTranslationCS{(dokázat nesprávnost)} \dicExampleIS{hrekja orð hennar} \dicExampleCS{vyvrátit její slova};  \dicIdiom{hrekjast}{ \dicPhraseIS{hrekjast}} \dicFlx{refl} {\textbf{a.}} \dicSynonym*{flæmast} \dicDirectTranslationCS{být vyhnán, muset uprchnout};  {\textbf{b.}} \dicDirectTranslationCS{trpět, strádat} \dicIndirectTranslationCS{(vlivem špatného počasí venku)}
\dicEntry[hrekk] \dicTerm{hrekk} \dicIPA{{\textsubring{r}}{\textepsilon}{h}{\r{g}}} \dicPos{v} \dicFlx{ind praes sg 1 pers} \dicLink{hrökkva}
\dicEntry[hrekkja] \dicTerm{hrekk|ja} \dicIPA{{\textsubring{r}}{\textepsilon}{h}{\r{\textObardotlessj}}{a}} \dicPos{v}[2] \dicFlx{(‑ti, ‑t)}[25] \dicFlx{acc} \dicSynonym{áreita} \dicDirectTranslationCS{(žertem) zlobit, škádlit, dobírat si} \dicExampleIS{hrekkja vinnufélaga} \dicExampleCS{zlobit kamarády}
\dicEntry[hrekkjalómur] \dicTerm{hrekkja··lóm|ur} \dicIPA{{\textsubring{r}}{\textepsilon}{h}{\r{\textObardotlessj}}{a}{l}{ou}{m}{\textscy}{\textsubring{r}}} \dicPos{m}[6] \dicFlx{(‑s, ‑ar)}[24] \dicDirectTranslationCS{rošťák, rošťačka, lump, uličník, uličnice}
\dicEntry[hrekkjóttur] \dicTerm{hrekkj··óttur} \dicIPA{{\textsubring{r}}{\textepsilon}{h}{\r{\textObardotlessj}}{ou}{h}{\textsubring{d}}{\textscy}{\textsubring{r}}} \dicPos{adj}[1]\dicFlx{}[-10] \dicDirectTranslationCS{nezbedný, rozpustilý, uličnický}
\dicEntry[hrekkur] \dicTerm{hrekk|ur} \dicIPA{{\textsubring{r}}{\textepsilon}{h}{\r{g}}{\textscy}{\textsubring{r}}} \dicPos{m}[9] \dicFlx{(‑s, ‑ir)}[15] \dicDirectTranslationCS{rošťárna, lumpárna, uličnictví}
\dicEntry[hrekti] \dicTerm{hrekti} \dicIPA{{\textsubring{r}}{\textepsilon}{x}{\textsubring{d}}{\textsci}} \dicPos{v} \dicFlx{con pf sg 1 pers} \dicLink{hrekja}
\dicEntry[hrella] \dicTerm{hrell|a} \dicIPA{{\textsubring{r}}{\textepsilon}{\textsubring{d}}{l}{a}} \dicPos{v}[2] \dicFlx{(‑di, ‑t)}[141] \dicFlx{acc} \dicSynonym*{valda leiða} \dicDirectTranslationCS{soužit, otravovat, obtěžovat}
\dicEntry[hremma] \dicTerm{hremm|a} \dicIPA{{\textsubring{r}}{\textepsilon}{m}{\textlengthmark}{a}} \dicPos{v}[2] \dicFlx{(‑di, ‑t)}[145] \dicFlx{acc} \dicSynonym{hrifsa} \dicDirectTranslationCS{popadnout, chňapnout} \dicExampleIS{hremma bráð} \dicExampleCS{popadnout kořist}
\dicEntry[hremming] \dicTerm{hremm··ing} \dicIPA{{\textsubring{r}}{\textepsilon}{m}{\textlengthmark}{i}{\ng}{\r{g}}} \dicPos{f}[4] \dicFlx{(‑ar, ‑ar)}[5] \textbf{1.} \dicSynonym*{það að hremma} \dicDirectTranslationCS{popadnutí, chňapnutí}  \textbf{2.} \dicSynonym*{slæm uppákoma} \dicDirectTranslationCS{těžkosti, nepříjemnosti}
\dicEntry[hreppa] \dicTerm{hrepp|a} \dicIPA{{\textsubring{r}}{\textepsilon}{h}{\textsubring{b}}{a}} \dicPos{v}[2] \dicFlx{(‑ti, ‑t)}[69] \dicFlx{acc} \textbf{1.} \dicDirectTranslationCS{chytit (špatné počasí)} \dicExampleIS{hreppa óveður í hafi} \dicExampleCS{chytit špatné počasí na moři}  \textbf{2.} \dicDirectTranslationCS{získat (cenu ap.)} \dicExampleIS{hreppa stóran vinning í happdrætti} \dicExampleCS{získat velkou výhru v~loterii}
\dicEntry[hreppsnefnd] \dicTerm{hrepps··nefnd} \dicIPA{{\textsubring{r}}{\textepsilon}{f}{s}{n}{\textepsilon}{m}{\textsubring{d}}} \dicPos{f}[7] \dicFlx{(‑ar, ‑ir)}[1] \dicDirectTranslationCS{okresní rada}
\dicEntry[hreppstjóri] \dicTerm{hrepp··stjór|i} \dicsymFrequent\  \dicIPA{{\textsubring{r}}{\textepsilon}{f}{s}{\textsubring{d}}{j}{ou}{r}{\textsci}} \dicPos{m}[1] \dicFlx{(‑a, ‑ar)}[1] \dicLangCat{zast.} \dicDirectTranslationCS{starosta\,/\addthin starostka (obce)} \dicExampleIS{Hreppstjórinn var dugnaðar maður og sinnti sínu starfi vel.} \dicExampleCS{Starosta byl schopný člověk a~svou práci vykonával dobře.}
\dicEntry[hreppur] \dicTerm{hrepp|ur} \dicIPA{{\textsubring{r}}{\textepsilon}{h}{\textsubring{b}}{\textscy}{\textsubring{r}}} \dicPos{m}[6] \dicFlx{(‑s, ‑ar)}[8] \dicDirectTranslationCS{okres, okrsek}
\dicEntry[hress] \dicTerm{hress} \dicsymFrequent\  \dicIPA{{\textsubring{r}}{\textepsilon}{s}{\textlengthmark}} \dicPos{adj}[5] \dicFlx{(f ‑)}[1] \textbf{1.} \dicSynonym{fjörlegur} \dicDirectTranslationCS{čilý, čiperný, svěží}  \textbf{2.} \dicSynonym{frískur} \dicDirectTranslationCS{(jsoucí) fit, (jsoucí) v~kondici} \dicExampleIS{Söngkonan er orðin hress á ný.} \dicExampleCS{Zpěvačka je znovu ve formě.}
\dicEntry[hressa] \dicTerm{hress|a} \dicIPA{{\textsubring{r}}{\textepsilon}{s}{\textlengthmark}{a}} \dicPos{v}[2] \dicFlx{(‑ti, ‑t)}[63] \dicDirectTranslationCS{osvěžit, osvěžovat, oživit, oživovat};  \dicPhraseIS{hressa sig á e‑u} \dicDirectTranslationCS{osvěžit se (čím)};  \dicIdiom{hressa}[upp á]{ \dicPhraseIS{hressa upp á e‑ð}} \dicDirectTranslationCS{osvěžit (co)};  \dicIdiom{hressa}[við]{ \dicPhraseIS{hressa e‑ð við}} \dicLangCat{přen.} \dicDirectTranslationCS{oživit (co)} \dicExampleIS{hressa íbúðina við} \dicExampleCS{oživit byt}; { \dicPhraseIS{hressa e‑n við}} \dicDirectTranslationCS{potěšit (koho), udělat (komu) radost};  \dicIdiom{hressast}{ \dicPhraseIS{hressast}} \dicFlx{refl} \dicDirectTranslationCS{zotavit se, vzpamatovat se, pookřát}
\dicEntry[hressandi] \dicTerm{hress··andi} \dicIPA{{\textsubring{r}}{\textepsilon}{s}{\textlengthmark}{a}{n}{\textsubring{d}}{\textsci}} \dicPos{adj}[13] \dicFlx{indecl}[1] \dicDirectTranslationCS{osvěžující, povzbuzující}
\dicEntry[hressilegur] \dicTerm{hressi··legur} \dicIPA{{\textsubring{r}}{\textepsilon}{s}{\textlengthmark}{\textsci}{l}{\textepsilon}{\textbabygamma}{\textscy}{\textsubring{r}}} \dicPos{adj}[1]\dicFlx{}[-8] \dicSynonym{hressandi} \dicDirectTranslationCS{čerstvý, svěží, osvěžující}
\dicEntry[hressing] \dicTerm{hress··ing} \dicIPA{{\textsubring{r}}{\textepsilon}{s}{\textlengthmark}{i}{\ng}{\r{g}}} \dicPos{f}[4] \dicFlx{(‑ar, ‑ar)}[5] \textbf{1.} \dicDirectTranslationCS{osvěžení, občerstvení};  \dicPhraseIS{gera e‑ð sér til hressingar} \dicDirectTranslationCS{udělat si (co) na občerstvení}  \textbf{2.} \dicSynonym{heilsubót} \dicDirectTranslationCS{rekreace, relaxace}
\dicEntry[hret] \dicTerm{hret} \dicIPA{{\textsubring{r}}{\textepsilon}{\textlengthmark}{\textsubring{d}}} \dicPos{n}[2] \dicFlx{(‑s, ‑)}[5] \dicSynonym{snjókoma} \dicDirectTranslationCS{plískanice}
\dicEntry[hreyfa] \dicTerm{hreyf|a} \dicsymFrequent\  \dicIPA{{\textsubring{r}}{ei}{\textlengthmark}{v}{a}} \dicPos{v}[2] \dicFlx{(‑ði, ‑t)}[102] \dicFlx{acc} \textbf{1.} \dicSynonym{færa} \dicDirectTranslationCS{(po)hnout, hýbat, pohybovat} \dicExampleIS{hreyfa höndina} \dicExampleCS{hýbat rukou};  \dicPhraseIS{hreyfa við e‑u} \dicDirectTranslationCS{hýbat s~(čím), pohnout s~(čím) (s~předmětem ap.)}  \textbf{2.} \dicPhraseIS{hreyfa sig} \dicDirectTranslationCS{(po)hnout se, hýbat se, pohybovat se}  \textbf{3.} \dicSynonym*{minnast á} \dicDirectTranslationCS{pohnout (s~tématem ap.), začít hovor} \dicExampleIS{hreyfa málinu} \dicExampleCS{pohnout se záležitostí}  \textbf{4.} \dicPhraseIS{það hreyfir ekki vind} \dicFlx{impers} \dicDirectTranslationCS{je bezvětří, nefouká vítr};  \dicIdiom{hreyfast}{ \dicPhraseIS{hreyfast}} \dicFlx{refl} \dicDirectTranslationCS{hýbat se, pohybovat se} \dicExampleIS{hreyfast hratt} \dicExampleCS{rychle se hýbat}
\dicEntry[hreyfanlegur] \dicTerm{hreyfan··legur} \dicIPA{{\textsubring{r}}{ei}{\textlengthmark}{v}{a}{n}{l}{\textepsilon}{\textbabygamma}{\textscy}{\textsubring{r}}} \dicPos{adj}[1]\dicFlx{}[-8] \dicDirectTranslationCS{pohyblivý, přenosný, mobilní}
\dicEntry[hreyfanleiki] \dicTerm{hreyfan··leik|i} \dicIPA{{\textsubring{r}}{ei}{\textlengthmark}{v}{a}{n}{l}{ei}{\r{\textObardotlessj}}{\textsci}} \dicPos{m}[1] \dicFlx{(‑a)}[3] \dicDirectTranslationCS{pohyblivost, mobilita}
\begin{xtolerant}{}{1pt}
\dicEntry[hreyfihamlaður] \dicTerm{hreyfi··|haml·aður}\addthinS\dicIPA{{\textsubring{r}}{ei}{\textlengthmark}{v}{\textsci}{h}{a}{m}{l}{a}{ð}{\textscy}{\textsubring{r}}} \dicPos{adj}[3] \dicFlx{(f ‑hömluð)}[1] \dicDirectTranslationCS{těles\-ně postižený, (pohybově) handicapovaný}
\end{xtolerant}
\dicEntry[hreyfill] \dicTerm{hreyf|ill} \dicIPA{{\textsubring{r}}{ei}{\textlengthmark}{v}{\textsci}{\textsubring{d}}{\textsubring{l}}} \dicPos{m}[6] \dicFlx{(‑ils, ‑lar)}[35] \dicDirectTranslationCS{motor, pohon}
\dicEntry[hreyfing] \dicTerm{hreyf··ing} \dicsymFrequent\  \dicIPA{{\textsubring{r}}{ei}{\textlengthmark}{v}{i}{\ng}{\r{g}}} \dicPos{f}[4] \dicFlx{(‑ar, ‑ar)}[5] \textbf{1.} \dicSynonym{hræring} \dicDirectTranslationCS{pohyb} \dicExampleIS{Það er hreyfing á bílnum.} \dicExampleCS{Auto se pohybuje.};  \dicPhraseIS{vera á hreyfingu} \dicDirectTranslationCS{být v~pohybu};  \dicPhraseIS{vera (mjúkur) í hreyfingum} \dicDirectTranslationCS{(elegantně) se pohybovat};  \dicPhraseIS{það er hreyfing á e‑u} \dicDirectTranslationCS{je pohyb na (čem) (na trhu práce ap.)}  \textbf{2.} \dicSynonym{samtök} \dicDirectTranslationCS{hnutí} \dicIndirectTranslationCS{(masové úsilí)}
\dicEntry[hreyfingarlaus] \dicTerm{hreyfingar··laus} \dicsymFrequent\  \dicIPA{{\textsubring{r}}{ei}{\textlengthmark}{v}{i}{\ng}{\r{g}}{a}{r}{l}{\oe i}{s}} \dicPos{adj}[5]\dicFlx{}[-1] \dicDirectTranslationCS{nehybný, strnulý} \dicExampleIS{liggja hreyfingarlaus} \dicExampleCS{ležet nehybně}
\dicEntry[hreyfingarleysi] \dicTerm{hreyfingar··leysi} \dicIPA{{\textsubring{r}}{ei}{\textlengthmark}{v}{i}{\ng}{\r{g}}{a}{r}{l}{ei}{s}{\textsci}} \dicPos{n}[2] \dicFlx{(‑s)}[20] \textbf{1.} \dicSynonym*{það að vera hreyfingarlaus} \dicDirectTranslationCS{nehybnost, strnulost}  \textbf{2.} \dicSynonym*{kyrrsetur} \dicDirectTranslationCS{fyzická nečinnost, nedostatek pohybu}
\dicEntry[hreyfiorka] \dicTerm{hreyfi··ork|a} \dicIPA{{\textsubring{r}}{ei}{\textlengthmark}{v}{\textsci}{\textopeno}{\textsubring{r}}{\r{g}}{a}} \dicPos{f}[1] \dicFlx{(‑u)}[5] \dicFieldCat{fyz.} \dicDirectTranslationCS{kinetická\,/\addthin pohybová energie} \dicAntonym{stöðuorka}
\dicEntry[hreykinn] \dicTerm{hreykinn} \dicsymFrequent\  \dicIPA{{\textsubring{r}}{ei}{\textlengthmark}{\r{\textObardotlessj}}{\textsci}{\textsubring{n}}} \dicPos{adj}[6]\dicFlx{}[-2] \dicSynonym*{upp með sér} \dicDirectTranslationCS{hrdý, pyšný} \dicExampleIS{hreykinn af e‑u\,/\addthin e‑m} \dicExampleCS{hrdý na (co\,/\addthin koho)}
\dicEntry[hreykja] \dicTerm{hreyk|ja} \dicIPA{{\textsubring{r}}{ei}{\textlengthmark}{\r{\textObardotlessj}}{a}} \dicPos{v}[2] \dicFlx{(‑ti, ‑t)}[29] \dicFlx{dat} \dicPhraseIS{hreykja sér} \dicDirectTranslationCS{chlubit se, chvástat se} \dicExampleIS{hreykja sér af e‑u} \dicExampleCS{chvástat se (čím)};  \dicPhraseIS{hreykja sér hátt} \dicDirectTranslationCS{naparovat se}
\dicEntry[hreysi] \dicTerm{hreysi} \dicIPA{{\textsubring{r}}{ei}{\textlengthmark}{s}{\textsci}} \dicPos{n}[2] \dicFlx{(‑s, ‑)}[14] \dicDirectTranslationCS{chatrč, bouda}
\dicEntry[hreysiköttur] \dicTerm{hreysi··|köttur} \dicIPA{{\textsubring{r}}{ei}{\textlengthmark}{s}{\textsci}{k\smash{\textsuperscript{h}}}{\oe}{h}{\textsubring{d}}{\textscy}{\textsubring{r}}} \dicPos{m}[11] \dicFlx{(‑kattar, ‑kettir)}[5] \dicFieldCat{zool.} \dicDirectTranslationCS{lasice hranostaj, hranostaj, hranostaj obecný} \textit{(l.~{\textLA{Mustela erminea}})}  \dicsymPhoto\ 
\dicFigure{ds_image_hreysikottur_0_2.jpg}{Hreysiköttur}{Hreysiköttur - Steve Hillebrand, USFWS, PD}
\begin{xtolerant}{}{1pt}
\dicEntry[hreysti] \dicTerm{hreyst|i}\addthinS\dicIPA{{\textsubring{r}}{ei}{s}{\textsubring{d}}{\textsci}}\addthinS\dicPos{f}[3]\addthinS\dicFlx{(‑i)}[3] \textbf{1.} \dicSynonym{hugrekki\smash{\textsuperscript{2}}} \dicDirectTranslationCS{udatnost, chrabrost}  \textbf{2.} \dicSynonym{heilbrigði\smash{\textsuperscript{2}}} \dicDirectTranslationCS{vitalita, dobré zdraví}  \textbf{3.} \dicFieldCat{sport.} \dicDirectTranslationCS{fitness}
\end{xtolerant}
\dicEntry[hreystiverk] \dicTerm{hreysti··verk} \dicIPA{{\textsubring{r}}{ei}{s}{\textsubring{d}}{\textsci}{v}{\textepsilon}{\textsubring{r}}{\r{g}}} \dicPos{n}[2] \dicFlx{(‑s, ‑)}[5] \dicDirectTranslationCS{udatný\,/\addthin chrabrý čin}
\dicEntry[hreyta] \dicTerm{hreyt|a} \dicIPA{{\textsubring{r}}{ei}{\textlengthmark}{\textsubring{d}}{a}} \dicPos{v}[2] \dicFlx{(‑ti, ‑t)}[54] \dicSynonym{kasta} \dicDirectTranslationCS{vrhnout, vrhat, mrštit} \dicExampleIS{hreyta grjóti} \dicExampleCS{mrštit kamením};  \dicIdiom{hreyta}[í]{ \dicPhraseIS{hreyta í e‑n ónotum}} \dicDirectTranslationCS{zasypat (koho) nadávkami};  \dicIdiom{hreyta}[út úr]{ \dicPhraseIS{hreyta e‑u út úr sér}} \dicDirectTranslationCS{vyhrknout (co) (poznámku ap.)}
\dicEntry[hrifi] \dicTerm{hrifi} \dicIPA{{\textsubring{r}}{\textsci}{\textlengthmark}{v}{\textsci}} \dicPos{v} \dicFlx{con pf sg 1 pers} \dicLink{hrífa\smash{\textsuperscript{2}}}
\dicEntry[hrifið] \dicTerm{hrifið} \dicIPA{{\textsubring{r}}{\textsci}{\textlengthmark}{v}{\textsci}{\texttheta}} \dicPos{v} \dicFlx{supin} \dicLink{hrífa\smash{\textsuperscript{2}}}
\dicEntry[hrifinn] \dicTerm{hrifinn} \dicsymFrequent\  \dicIPA{{\textsubring{r}}{\textsci}{\textlengthmark}{v}{\textsci}{\textsubring{n}}} \dicPos{adj}[6]\dicFlx{}[-2] \textbf{1.} \dicSynonym{hugfanginn} \dicDirectTranslationCS{nadšený, okouzlený, unesený, fascinovaný} \dicExampleIS{vera hrifinn af bókinni} \dicExampleCS{být nadšený knížkou}  \textbf{2.} \dicSynonym{ástfanginn} \dicDirectTranslationCS{zamilovaný} \dicExampleIS{vera hrifinn af e‑m} \dicExampleCS{být do (koho) zamilovaný}
\dicEntry[hrifning] \dicTerm{hrif··ning} \dicIPA{{\textsubring{r}}{\textsci}{\textsubring{b}}{n}{i}{\ng}{\r{g}}} \dicPos{f}[4] \dicFlx{(‑ar)}[7] \dicDirectTranslationCS{nadšení, okouzlení, fascinace, entuziasmus}
\dicEntry[hrifnæmi] \dicTerm{hrif··næmi} \dicIPA{{\textsubring{r}}{\textsci}{v}{n}{a}{i}{m}{\textsci}} \dicPos{n}[2] \dicFlx{(‑s)}[20] \dicDirectTranslationCS{vnímavost, citlivost}
\dicEntry[hrifnæmur] \dicTerm{hrif··næmur} \dicIPA{{\textsubring{r}}{\textsci}{v}{n}{a}{i}{m}{\textscy}{\textsubring{r}}} \dicPos{adj}[1]\dicFlx{}[-1] \dicDirectTranslationCS{vnímavý, citlivý}
\dicEntry[hrifsa] \dicTerm{hrifs|a} \dicIPA{{\textsubring{r}}{\textsci}{f}{s}{a}} \dicPos{v}[1] \dicFlx{(‑aði)}[1] \dicFlx{acc} \textbf{1.} \dicSynonym{grípa} \dicDirectTranslationCS{chňapnout, popadnout}  \textbf{2.} \dicSynonym{glefsa\smash{\textsuperscript{2}}} \dicDirectTranslationCS{chňapnout, kousnout} \dicExampleIS{Hundurinn hrifsaði í ærnar.} \dicExampleCS{Pes chňapal po ovcích.}
\dicEntry[hrifum] \dicTerm{hrifum} \dicIPA{{\textsubring{r}}{\textsci}{\textlengthmark}{v}{\textscy}{\textsubring{m}}} \dicPos{v} \dicFlx{ind pf pl 1 pers} \dicLink{hrífa\smash{\textsuperscript{2}}}
\dicEntry[hrikalegur] \dicTerm{hrika··legur} \dicsymFrequent\  \dicIPA{{\textsubring{r}}{\textsci}{\textlengthmark}{\r{g}}{a}{l}{\textepsilon}{\textbabygamma}{\textscy}{\textsubring{r}}} \dicPos{adj}[1]\dicFlx{}[-8] \textbf{1.} \dicSynonym{ægilegur} \dicDirectTranslationCS{strašlivý, děsivý, nahánějící hrůzu}  \textbf{2.} \dicDirectTranslationCS{strašný, ohromný} \dicIndirectTranslationCS{(k~zdůraznění)} \dicExampleIS{hrikaleg mistök} \dicExampleCS{strašná chyba}
\dicEntry[hrikta] \dicTerm{hrikt|a} \dicIPA{{\textsubring{r}}{\textsci}{x}{\textsubring{d}}{a}} \dicPos{v}[2] \dicFlx{(‑i, ‑)}[22] \dicFlx{impers} \dicPhraseIS{það hriktir í e‑u} \dicDirectTranslationCS{(co) vrže, (co) skřípá} \dicExampleIS{Það hriktir í hurðinni.} \dicExampleCS{Dveře skřípou.}
\dicEntry[hrina] \dicTerm{hrin|a} \dicIPA{{\textsubring{r}}{\textsci}{\textlengthmark}{n}{a}} \dicPos{f}[1] \dicFlx{(‑u, ‑ur)}[7] \textbf{1.} \dicDirectTranslationCS{sprcha (krupobití ap.), série (zemětřesení ap.)}  \textbf{2.} \dicDirectTranslationCS{vlna, série (útoků ap.)} \dicExampleIS{hrina sprengjuárása} \dicExampleCS{série bombových útoků}  \textbf{3.} \dicFieldCat{sport.} \dicDirectTranslationCS{set (v~tenise ap.)}
\dicEntry[hrinda] \dicTerm{hrin|da} \dicsymFrequent\  \dicIPA{{\textsubring{r}}{\textsci}{n}{\textsubring{d}}{a}} \dicPos{v}[7] \dicFlx{(‑ti, ‑t)\,/\addthin (hrindi, hratt, hrundum, hryndi, hrundið)}[1] \dicFlx{dat} \textbf{1.} \dicSynonym*{stjaka við} \dicDirectTranslationCS{(po)strčit, strkat} \dicExampleIS{hrinda e‑m niður} \dicExampleCS{strčit (koho) dolů}  \textbf{2.} \dicSynonym*{koma af stað} \dicDirectTranslationCS{začít, začínat, zahájit, zahajovat};  \dicPhraseIS{hrinda e‑u af stað} \dicDirectTranslationCS{zahájit (co), spustit (co)};  \dicPhraseIS{hrinda e‑u í framkvæmd} \dicDirectTranslationCS{zrealizovat (co), uskutečnit (co)} \dicExampleIS{hrinda verkinu í framkvæmd} \dicExampleCS{zrealizovat tu práci}  \textbf{3.} \dicDirectTranslationCS{odpuzovat (vodu ap.)}  \textbf{4.} \dicSynonym{hrekja} \dicDirectTranslationCS{odrazit, odrážet (útok ap.)};  \dicIdiom{hrinda}[upp]{ \dicPhraseIS{hrinda hurðinni upp}} \dicDirectTranslationCS{otevřít dveře}
\dicEntry[hrindi] \dicTerm{hrindi} \dicIPA{{\textsubring{r}}{\textsci}{n}{\textsubring{d}}{\textsci}} \dicPos{v} \dicFlx{ind praes sg 1 pers} \dicLink{hrinda}
\dicEntry[hringa] \dicTerm{hring|a} \dicIPA{{\textsubring{r}}{i}{\ng}{\r{g}}{a}} \dicPos{v}[1] \dicFlx{(‑aði)}[1] \dicFlx{acc} \dicDirectTranslationCS{svinout, svíjet, stočit, zkroutit} \dicExampleIS{hringa kaðal} \dicExampleCS{svinout provaz};  \dicPhraseIS{hringa sig} \dicDirectTranslationCS{stočit se, svinout se (do klubíčka ap.)};  \dicIdiom{hringast}{ \dicPhraseIS{hringast}} \dicFlx{refl} \dicDirectTranslationCS{stočit se, zkroutit se}
\dicEntry[hringbraut] \dicTerm{hring··braut} \dicIPA{{\textsubring{r}}{i}{\ng}{\r{g}}{\textsubring{b}}{r}{\oe i}{\textsubring{d}}} \dicPos{f}[7] \dicFlx{(‑ar, ‑ir)}[1] \dicDirectTranslationCS{okružní silnice}
\dicEntry[hringdúfa] \dicTerm{hring··dúf|a} \dicIPA{{\textsubring{r}}{i}{\ng}{\r{g}}{\textsubring{d}}{u}{v}{a}} \dicPos{f}[1] \dicFlx{(‑u, ‑ur)}[13] \dicFieldCat{zool.} \dicDirectTranslationCS{holub hřivnáč} \textit{(l.~{\textLA{Columba palumbus}})}
\dicEntry[hringekja] \dicTerm{hring··ekj|a} \dicIPA{{\textsubring{r}}{i}{\ng}{\r{g}}{\textepsilon}{\r{\textObardotlessj}}{a}} \dicPos{f}[1] \dicFlx{(‑u, ‑ur)}[7] \dicDirectTranslationCS{kolotoč}
\dicEntry[hringfari] \dicTerm{hring··far|i} \dicIPA{{\textsubring{r}}{i}{\ng}{\r{g}}{f}{a}{r}{\textsci}} \dicPos{m}[1] \dicFlx{(‑a, ‑ar)}[8] \dicSynonym{sirkill} \dicDirectTranslationCS{kružítko}
\dicEntry[hringferð] \dicTerm{hring··ferð} \dicIPA{{\textsubring{r}}{i}{\ng}{\r{g}}{f}{\textepsilon}{r}{\texttheta}} \dicPos{f}[7] \dicFlx{(‑ar, ‑ir)}[1] \dicDirectTranslationCS{okružní cesta\,/\addthin jízda} \dicExampleIS{hringferð um Ísland} \dicExampleCS{okružní cesta kolem Islandu}
\dicEntry[hringfingur] \dicTerm{hring··fing|ur} \dicIPA{{\textsubring{r}}{i}{\ng}{\r{g}}{f}{i}{\ng}{\r{g}}{\textscy}{\textsubring{r}}} \dicPos{m}[5] \dicFlx{(‑urs, ‑ur)}[13] \dicSynonym{baugfingur} \dicDirectTranslationCS{prsteník, prsteníček}
\dicEntry[hringgeiri] \dicTerm{hring··geir|i} \dicIPA{{\textsubring{r}}{i}{\ng}{\r{g}}{\r{\textObardotlessj}}{ei}{r}{\textsci}} \dicPos{m}[1] \dicFlx{(‑a, ‑ar)}[1] \dicFieldCat{mat.} \dicDirectTranslationCS{(kruhová) výseč}
\dicEntry[hringiða] \dicTerm{hring··ið|a} \dicIPA{{\textsubring{r}}{i}{\ng}{\r{g}}{\textsci}{ð}{a}} \dicPos{f}[1] \dicFlx{(‑u, ‑ur)}[19] \textbf{1.} \dicDirectTranslationCS{(vodní) vír, víření}  \textbf{2.} \dicLangCat{přen.} \dicDirectTranslationCS{vír (života ap.)}
\dicEntry[hringing] \dicTerm{hring··ing} \dicIPA{{\textsubring{r}}{i}{\textltailn}{\r{\textObardotlessj}}{i}{\ng}{\r{g}}} \dicPos{f}[4] \dicFlx{(‑ar, ‑ar)}[5] \textbf{1.} \dicSynonym{hringitónn} \dicDirectTranslationCS{vyzvánění, (za)zvonění} \dicExampleIS{hringing í vekjaraklukku} \dicExampleCS{zvonění budíku}  \textbf{2.} \dicSynonym*{það að hringja} \dicDirectTranslationCS{zazvonění, zatelefonování}
\dicEntry[hringitónn] \dicTerm{hringi··tón|n} \dicIPA{{\textsubring{r}}{i}{\textltailn}{\r{\textObardotlessj}}{\textsci}{t\smash{\textsuperscript{h}}}{ou}{\textsubring{d}}{\textsubring{n}}} \dicPos{m}[6] \dicFlx{(‑s, ‑ar)}[42] \dicDirectTranslationCS{vyzvánění, vyzváněcí tón}
\dicEntry[hringja] \dicTerm{hring|ja} \dicsymFrequent\  \dicIPA{{\textsubring{r}}{i}{\textltailn}{\r{\textObardotlessj}}{a}} \dicPos{v}[2] \dicFlx{(‑di, ‑t)}[135] \dicFlx{dat} \textbf{1.} \dicSynonym{klingja} \dicDirectTranslationCS{(za)zvonit, vyzvánět} \dicExampleIS{hringja bjöllu} \dicExampleCS{zazvonit na zvonek}  \textbf{2.} \dicDirectTranslationCS{(za)zvonit (budík ap.)} \dicExampleIS{Vekjaraklukkan hringir.} \dicExampleCS{Budík zvoní.}  \textbf{3.} \dicDirectTranslationCS{(za)zvonit, (za)volat (telefonem ap.), (za)telefonovat};  \dicPhraseIS{hringja í e‑n} \dicDirectTranslationCS{zavolat (komu) (telefonem ap.)}
\dicEntry[hringl] \dicTerm{hringl} \dicIPA{{\textsubring{r}}{i}{\ng}{\textsubring{l}}} \dicPos{n}[2] \dicFlx{(‑s)}[2] \textbf{1.} \dicDirectTranslationCS{cinkání, cinkot, řinkot}  \textbf{2.} \dicSynonym{óstöðugleiki} \dicDirectTranslationCS{nestálost, přelétavost}
\dicEntry[hringla] \dicTerm{hringl|a\smash{\textsuperscript{1}}} \dicIPA{{\textsubring{r}}{i}{\ng}{l}{a}} \dicPos{f}[1] \dicFlx{(‑u, ‑ur)}[19] \dicDirectTranslationCS{chrastítko, chřestítko}
\dicEntry[hringla] \dicTerm{hringl|a\smash{\textsuperscript{2}}} \dicIPA{{\textsubring{r}}{i}{\ng}{l}{a}} \dicPos{v}[1] \dicFlx{(‑aði)}[7] \dicFlx{dat} \textbf{1.} \dicPhraseIS{það hringlar í e‑u} \dicFlx{impers} \dicSynonym{glamra} \dicDirectTranslationCS{(co) chřestí, (co) řinčí}  \textbf{2.} \dicDirectTranslationCS{(poz)měnit, hýbat} \dicExampleIS{hringla með e‑ð} \dicExampleCS{měnit (co) ze dne na den}
\dicEntry[hringlaga] \dicTerm{hring··laga} \dicsymFrequent\  \dicIPA{{\textsubring{r}}{i}{\ng}{\r{g}}{l}{a}{\textbabygamma}{a}} \dicPos{adj}[13] \dicFlx{indecl}[1] \dicDirectTranslationCS{kulatý, kruhový} \dicExampleIS{Salurinn er fallegt hringlaga herbergi.} \dicExampleCS{Sál je pěkná kruhová místnost.}
\dicEntry[hringleikahús] \dicTerm{hring·leika··hús} \dicIPA{{\textsubring{r}}{i}{\ng}{\r{g}}{l}{ei}{\r{g}}{a}{h}{u}{s}} \dicPos{n}[2] \dicFlx{(‑s, ‑)}[5] \dicSynonym{fjölleikahús} \dicDirectTranslationCS{cirkus}
\dicEntry[hringrás] \dicTerm{hring··rás} \dicIPA{{\textsubring{r}}{i}{\ng}{\r{g}}{r}{au}{s}} \dicPos{f}[7] \dicFlx{(‑ar, ‑ir)}[1] \dicDirectTranslationCS{oběh, koloběh, cirkulace} \dicExampleIS{hringrás blóðsins} \dicExampleCS{krevní oběh}
\dicEntry[hringrif] \dicTerm{hring··rif} \dicIPA{{\textsubring{r}}{i}{\ng}{\r{g}}{r}{\textsci}{f}} \dicPos{n}[2] \dicFlx{(‑s, ‑)}[6] \dicDirectTranslationCS{atol}
\dicEntry[hringsjá] \dicTerm{hring··sjá} \dicIPA{{\textsubring{r}}{i}{\ng}{\r{g}}{s}{j}{au}} \dicPos{f}[4] \dicFlx{(‑r, ‑r)}[18] \textbf{1.} \dicFieldCat{nám.} \dicDirectTranslationCS{periskop}  \textbf{2.} \dicDirectTranslationCS{panoramatický průvodce} \dicIndirectTranslationCS{(ukazující pozici významných bodů v~krajině)}
\dicEntry[hringsneri] \dicTerm{hring··sneri} \dicIPA{{\textsubring{r}}{i}{\ng}{\r{g}}{s}{\textsubring{d}}{n}{\textepsilon}{r}{\textsci}} \dicPos{v} \dicFlx{ind\,/\addthin con pf sg 1 pers} \dicLink{hringsnúa}
\dicEntry[hringsnerum] \dicTerm{hring··snerum} \dicIPA{{\textsubring{r}}{i}{\ng}{\r{g}}{s}{\textsubring{d}}{n}{\textepsilon}{r}{\textscy}{\textsubring{m}}} \dicPos{v} \dicFlx{ind pf pl 1 pers} \dicLink{hringsnúa}
\dicEntry[hringsnéri] \dicTerm{hring··snéri} \dicIPA{{\textsubring{r}}{i}{\ng}{\r{g}}{s}{\textsubring{d}}{n}{j}{\textepsilon}{r}{\textsci}} \dicPos{v} \dicFlx{ind\,/\addthin con pf sg 1 pers} \dicLink{hringsnúa}
\dicEntry[hringsnérum] \dicTerm{hring··snérum} \dicIPA{{\textsubring{r}}{i}{\ng}{\r{g}}{s}{\textsubring{d}}{n}{j}{\textepsilon}{r}{\textscy}{\textsubring{m}}} \dicPos{v} \dicFlx{ind pf pl 1 pers} \dicLink{hringsnúa}
\dicEntry[hringsnúa] \dicTerm{hring··|snúa} \dicIPA{{\textsubring{r}}{i}{\ng}{\r{g}}{s}{\textsubring{d}}{n}{u}{a}} \dicPos{v}[5] \dicFlx{(‑sný, ‑sneri\,/\addthin ‑snéri, ‑snerum\,/\addthin ‑snérum, ‑sneri\,/\addthin ‑snéri, ‑snúið)}[5] \dicFlx{dat} \dicDirectTranslationCS{otočit, otáčet};  \dicIdiom{hringsnúast}{ \dicPhraseIS{hringsnúast}} \dicFlx{refl} \dicDirectTranslationCS{otočit se, otáčet se}
\dicEntry[hringsnúið] \dicTerm{hring··snúið} \dicIPA{{\textsubring{r}}{i}{\ng}{\r{g}}{s}{\textsubring{d}}{n}{u}{\textsci}{\texttheta}} \dicPos{v} \dicFlx{supin} \dicLink{hringsnúa}
\dicEntry[hringsný] \dicTerm{hring··sný} \dicIPA{{\textsubring{r}}{i}{\ng}{\r{g}}{s}{\textsubring{d}}{n}{i}} \dicPos{v} \dicFlx{ind praes sg 1 pers} \dicLink{hringsnúa}
\dicEntry[hringsóla] \dicTerm{hring··sól|a} \dicIPA{{\textsubring{r}}{i}{\ng}{\r{g}}{s}{ou}{l}{a}} \dicPos{v}[1] \dicFlx{(‑aði)}[1] \dicDirectTranslationCS{kroužit} \dicIndirectTranslationCS{(pohybovat se v~kruhu)}
\dicEntry[hringtorg] \dicTerm{hring··torg} \dicIPA{{\textsubring{r}}{i}{\ng}{\r{g}}{t\smash{\textsuperscript{h}}}{\textopeno}{r}{\r{g}}} \dicPos{n}[2] \dicFlx{(‑s, ‑)}[5] \dicDirectTranslationCS{kruhový objezd}
\dicEntry[hringur] \dicTerm{hring|ur} \dicsymFrequent\  \dicIPA{{\textsubring{r}}{i}{\ng}{\r{g}}{\textscy}{\textsubring{r}}} \dicPos{m}[6] \dicFlx{(‑s, ‑ir\,/\addthin ‑ar)}[67] \textbf{1.} \dicFlx{pl nom (-ir)} \dicDirectTranslationCS{kruh, obruč}  \textbf{2.} \dicFlx{pl nom (-ar)} \dicSynonym{baugur} \dicDirectTranslationCS{prsten, prstýnek} \dicExampleIS{gullhringur} \dicExampleCS{zlatý prsten}  \textbf{3.} \dicFieldCat{sport.} \dicDirectTranslationCS{ring}  \textbf{4.} \dicFieldCat{mat.} \dicDirectTranslationCS{kruh}  \textbf{5.} \dicSynonym{umferð} \dicDirectTranslationCS{kolo, runda (v~kartách ap.)}  \textbf{6.} \dicPhraseIS{hringir} \dicFlx{pl} \dicFieldCat{sport.} \dicDirectTranslationCS{kruhy}  \textbf{7.} \dicFieldCat{ekon.} \dicDirectTranslationCS{kartel}  \textbf{8.} \dicFieldCat{jaz.} \dicDirectTranslationCS{kroužek} \dicIndirectTranslationCS{(diakritické znaménko)}
\dicEntry[hringvegur] \dicTerm{hring··veg|ur} \dicIPA{{\textsubring{r}}{i}{\ng}{\r{g}}{v}{\textepsilon}{\textbabygamma}{\textscy}{\textsubring{r}}} \dicPos{m}[10] \dicFlx{(‑ar, ‑ir)}[4] \textbf{1.} \dicDirectTranslationCS{okružní silnice\,/\addthin cesta}  \textbf{2.} \dicDirectTranslationCS{okružní silnice} \dicIndirectTranslationCS{(silnice kolem pobřeží celého Islandu, nejčastěji s~určitým členem)} \dicExampleIS{hringvegurinn umhverfis Ísland} \dicExampleCS{okružní silnice kolem Islandu}
\dicEntry[hrini] \dicTerm{hrini} \dicIPA{{\textsubring{r}}{\textsci}{\textlengthmark}{n}{\textsci}} \dicPos{v} \dicFlx{con pf sg 1 pers} \dicLink{hrína}
\dicEntry[hrinið] \dicTerm{hrinið} \dicIPA{{\textsubring{r}}{\textsci}{\textlengthmark}{n}{\textsci}{\texttheta}} \dicPos{v} \dicFlx{supin} \dicLink{hrína}
\dicEntry[hrinum] \dicTerm{hrinum} \dicIPA{{\textsubring{r}}{\textsci}{\textlengthmark}{n}{\textscy}{\textsubring{m}}} \dicPos{v} \dicFlx{ind pf pl 1 pers} \dicLink{hrína}
\dicEntry[hripa] \dicTerm{hrip|a} \dicIPA{{\textsubring{r}}{\textsci}{\textlengthmark}{\textsubring{b}}{a}} \dicPos{v}[1] \dicFlx{(‑aði)}[1] \dicFlx{acc} \textbf{1.} \dicSynonym{leka} \dicDirectTranslationCS{vytékat, téct, kanout, crčet}  \textbf{2.} \dicDirectTranslationCS{načmárat, sesmolit (dopis ap.)} \dicExampleIS{hripa bréf} \dicExampleCS{načmárat dopis};  \dicIdiom{hripa}[niður]{ \dicPhraseIS{hripa niður e‑ð}} \dicDirectTranslationCS{zapsat si (co), poznamenat si (co) (zběžně)}
\dicEntry[hrista] \dicTerm{hrist|a} \dicsymFrequent\  \dicIPA{{\textsubring{r}}{\textsci}{s}{\textsubring{d}}{a}} \dicPos{v}[2] \dicFlx{(‑i, ‑)}[12] \dicFlx{acc} \dicDirectTranslationCS{(za)třást, (vy)třepat} \dicExampleIS{hrista teppið duglega} \dicExampleCS{vytřepat pořádně koberec};  \dicIdiom{hrista}[af]{ \dicPhraseIS{hrista e‑ð af sér}} \dicDirectTranslationCS{setřást ze sebe (co)} \dicExampleIS{hrista snjóinn af sér} \dicExampleCS{setřást ze sebe sníh};  \dicIdiom{hristast}{ \dicPhraseIS{hristast}} \dicFlx{refl} \dicDirectTranslationCS{(za)třást se, třepat se}
\dicEntry[hristari] \dicTerm{hrist··ar|i} \dicIPA{{\textsubring{r}}{\textsci}{s}{\textsubring{d}}{a}{r}{\textsci}} \dicPos{m}[1] \dicFlx{(‑a, ‑ar)}[13] \dicDirectTranslationCS{šejkr}
\dicEntry[hristingur] \dicTerm{hrist··ing|ur} \dicIPA{{\textsubring{r}}{\textsci}{s}{\textsubring{d}}{i}{\ng}{\r{g}}{\textscy}{\textsubring{r}}} \dicPos{m}[6] \dicFlx{(‑s, ‑ar)}[8] \dicSynonym*{það að hrista} \dicDirectTranslationCS{(za)třesení, (za)třepání}
\dicEntry[hríð] \dicTerm{hríð} \dicsymFrequent\  \dicIPA{{\textsubring{r}}{i}{\textlengthmark}{\texttheta}} \dicPos{f}[4] \dicFlx{(‑ar, ‑ir\,/\addthin ‑ar)}[30] \textbf{1.} \dicSynonym{snjókoma} \dicDirectTranslationCS{sněhová bouře, blizard}  \textbf{2.} \dicSynonym{lota} \dicDirectTranslationCS{chvíle} \dicExampleIS{koma aftur eftir skamma hríð} \dicExampleCS{vrátit se po krátké chvíli};  \dicPhraseIS{um hríð} \dicFlx{adv} \dicDirectTranslationCS{chvíli}
\dicEntry[hríðarbylur] \dicTerm{hríðar··byl|ur} \dicIPA{{\textsubring{r}}{i}{\textlengthmark}{ð}{a}{r}{\textsubring{b}}{\textsci}{l}{\textscy}{\textsubring{r}}} \dicPos{m}[9] \dicFlx{(‑s\,/\addthin ‑jar, ‑jir)}[32] \dicSynonym*{vindur með snjókomu} \dicDirectTranslationCS{sněhová bouře, vánice}
\dicEntry[hríðlækka] \dicTerm{hríð··lækk|a} \dicIPA{{\textsubring{r}}{i}{ð}{l}{a}{i}{h}{\r{g}}{a}} \dicPos{v}[1] \dicFlx{(‑aði)}[72] \dicFlx{acc} \dicDirectTranslationCS{prudce (po)klesnout, jít prudce dolů (o~cenách)} \dicExampleIS{Fiskverð hríðlækkaði.} \dicExampleCS{Cena ryb šla prudce dolů.}
\dicEntry[hríðskotabyssa] \dicTerm{hríð·skota··byss|a} \dicIPA{{\textsubring{r}}{i}{ð}{s}{\r{g}}{\textopeno}{\textsubring{d}}{a}{\textsubring{b}}{\textsci}{s}{a}} \dicPos{f}[1] \dicFlx{(‑u, ‑ur)}[19] \dicFieldCat{voj.} \dicDirectTranslationCS{samopal, rychlopalná zbraň}
\dicEntry[hríðversna] \dicTerm{hríð··versn|a} \dicIPA{{\textsubring{r}}{i}{ð}{v}{\textepsilon}{\textsubring{r}}{s}{\textsubring{d}}{n}{a}} \dicPos{v}[1] \dicFlx{(‑aði)}[73] \dicDirectTranslationCS{prudce\,/\addthin rychle se zhoršit} \dicExampleIS{Veðrið hríðversnaði.} \dicExampleCS{Počasí se rychle zhoršilo.};  \dicPhraseIS{e‑m hríðversnar} \dicFlx{impers} \dicDirectTranslationCS{(komu) se prudce přitěžuje} \dicExampleIS{Honum hríðversnaði dag frá degi.} \dicExampleCS{Přitížilo se mu prudce ze dne na den.}
\dicEntry[hrífa] \dicTerm{hríf|a\smash{\textsuperscript{1}}} \dicIPA{{\textsubring{r}}{i}{\textlengthmark}{v}{a}} \dicPos{f}[1] \dicFlx{(‑u, ‑ur)}[19] \dicDirectTranslationCS{hrábě}
\dicEntry[hrífa] \dicTerm{hrífa\smash{\textsuperscript{2}}} \dicsymFrequent\  \dicIPA{{\textsubring{r}}{i}{\textlengthmark}{v}{a}} \dicPos{v}[6] \dicFlx{(hríf, hreif, hrifum, hrifi, hrifið)}[69] \dicFlx{acc} \textbf{1.} \dicDirectTranslationCS{zapůsobit, uchvátit, nadchnout}  \textbf{2.} \dicDirectTranslationCS{působit, účinkovat, fungovat, mít efekt};  \dicPhraseIS{e‑að hrífur á e‑n} \dicDirectTranslationCS{(co) působí na (koho)};  \dicIdiom{hrífast}{ \dicPhraseIS{hrífast af e‑u}} \dicFlx{refl} \dicDirectTranslationCS{být unesen (čím), být uchvácen (čím), být okouzlen (čím)} \dicExampleIS{hrífast af útsýninu} \dicExampleCS{být uchvácen výhledem}
\dicEntry[hrífandi] \dicTerm{hríf··andi} \dicIPA{{\textsubring{r}}{i}{\textlengthmark}{v}{a}{n}{\textsubring{d}}{\textsci}} \dicPos{adj}[13] \dicFlx{indecl}[1] \dicSynonym{áhrifamikill} \dicDirectTranslationCS{působivý, úchvatný} \dicExampleIS{hrífandi tónlist} \dicExampleCS{působivá hudba}
\dicEntry[hrím] \dicTerm{hrím} \dicIPA{{\textsubring{r}}{i}{\textlengthmark}{\textsubring{m}}} \dicPos{n}[2] \dicFlx{(‑s)}[2] \textbf{1.} \dicSynonym{héla\smash{\textsuperscript{1}}} \dicDirectTranslationCS{jinovatka, jíní}  \textbf{2.} \dicSynonym{sót} \dicDirectTranslationCS{saze}
\dicEntry[hrín] \dicTerm{hrín} \dicIPA{{\textsubring{r}}{i}{\textlengthmark}{\textsubring{n}}} \dicPos{v} \dicFlx{ind praes sg 1 pers} \dicLink{hrína}
\dicEntry[hrína] \dicTerm{hrína} \dicIPA{{\textsubring{r}}{i}{\textlengthmark}{n}{a}} \dicPos{v}[6] \dicFlx{(hrín, hrein, hrinum, hrini, hrinið)}[83] \textbf{1.} \dicDirectTranslationCS{hýkat (osel ap.), kvičet (prase ap.)}  \textbf{2.} \dicDirectTranslationCS{řvát, ječet};  \dicIdiom{hrína}[á]{ \dicPhraseIS{það hrín ekki á e‑m}} \dicDirectTranslationCS{to na (koho) neplatí}
\dicEntry[hrís] \dicTerm{hrís\smash{\textsuperscript{1}}} \dicIPA{{\textsubring{r}}{i}{\textlengthmark}{s}} \dicPos{m}[4] \dicFlx{(‑s)}[26] \dicLink{hrís\smash{\textsuperscript{2}}}
\dicEntry[hrís] \dicTerm{hrís\smash{\textsuperscript{2}}}\dicTerm{, hrís\smash{\textsuperscript{1}}} \dicIPA{{\textsubring{r}}\-{i}\-{\textlengthmark}\-{s}\-} \dicPos{n}[2] \dicFlx{(‑s, ‑)}[5] \dicDirectTranslationCS{chrastí, chrundí}
\dicEntry[hrísgrjón] \dicTerm{hrís··grjón} \dicIPA{{\textsubring{r}}{i}{s}{\r{g}}{r}{j}{ou}{\textsubring{n}}} \dicPos{n}[2] \dicFlx{(‑s, ‑)}[5] \dicFieldCat{bot.} \dicDirectTranslationCS{rýže, rýže setá} \textit{(l.~{\textLA{Oryza sativa}})}  \dicsymPhoto\ 
\dicFigure{ds_image_hrisgrjon_0_2.jpg}{Hrísgrjón}{Hrísgrjón - IRRI Images, CC BY 2.0}
\dicEntry[hrísgrjónagrautur] \dicTerm{hrís·grjóna··graut|ur} \dicIPA{{\textsubring{r}}\-{i}\-{s}\-{\r{g}}\-{r}\-{j}\-{ou}\-{n}\-{a}\-{\r{g}}\-{r}\-{\oe i}\-{\textsubring{d}}\-{\textscy}\-{\textsubring{r}}\-} \dicPos{m}[6] \dicFlx{(‑ar, ‑ar)}[59] \dicFieldCat{kulin.} \dicDirectTranslationCS{rýžový nákyp, rýžová kaše}
\dicEntry[hrísla] \dicTerm{hrísl|a} \dicIPA{{\textsubring{r}}{i}{s}{\textsubring{d}}{l}{a}} \dicPos{f}[1] \dicFlx{(‑u, ‑ur)}[13] \textbf{1.} \dicSynonym{kvistur} \dicDirectTranslationCS{větvička}  \textbf{2.} \dicSynonym*{kræklótt tré} \dicDirectTranslationCS{stromek};  \dicPhraseIS{skjálfa eins og hrísla} \dicLangCat{přen.} \dicDirectTranslationCS{třást se jako osika}
\dicEntry[hrjósa] \dicTerm{hrjósa} \dicIPA{{\textsubring{r}}{j}{ou}{\textlengthmark}{s}{a}} \dicPos{v}[6] \dicFlx{(hrýs, hraus, hrysi, hrosið)}[53] \dicFlx{impers} \dicPhraseIS{e‑m hrýs hugur við e‑u} \dicDirectTranslationCS{(kdo) se třese při pomyšlení na (co)} \dicExampleIS{Okkur hraus hugur við fréttunum.} \dicExampleCS{Třásli jsme se při pomyšlení na ty zprávy.}
\dicEntry[hrjóstrugur] \dicTerm{hrjóstr··ugur} \dicIPA{{\textsubring{r}}{j}{ou}{s}{\textsubring{d}}{r}{\textscy}{\textbabygamma}{\textscy}{\textsubring{r}}} \dicPos{adj}[1]\dicFlx{}[-8] \dicDirectTranslationCS{neúrodný, neplodný} \dicExampleIS{hrjóstrugt land} \dicExampleCS{neúrodná půda}
\dicEntry[hrjóta] \dicTerm{hrjóta} \dicIPA{{\textsubring{r}}{j}{ou}{\textlengthmark}{\textsubring{d}}{a}} \dicPos{v}[6] \dicFlx{(hrýt, hraut, hrutum, hryti, hrotið)}[51] \textbf{1.} \dicDirectTranslationCS{chrápat} \dicExampleIS{hrjóta hátt} \dicExampleCS{hlasitě chrápat}  \textbf{2.} \dicSynonym{hrökkva} \dicDirectTranslationCS{vypadnout, vyklouznout, vyskočit} \dicExampleIS{Skífan hraut úr hendi hans.} \dicExampleCS{Podložka mu vyklouzla z~ruky.}  \textbf{3.} \dicDirectTranslationCS{vyklouznout, mimoděk uniknout (slovo ap.)}
\dicEntry[hrjúfur] \dicTerm{hrjúfur} \dicIPA{{\textsubring{r}}{j}{u}{\textlengthmark}{v}{\textscy}{\textsubring{r}}} \dicPos{adj}[1]\dicFlx{}[-1] \textbf{1.} \dicSynonym{ósléttur} \dicDirectTranslationCS{drsný, nerovný (povrch ap.)}  \textbf{2.} \dicDirectTranslationCS{hrubý, drsný (člověk ap.)} \dicExampleIS{vera hrjúfur í skapi} \dicExampleCS{být drsný od povahy}  \textbf{3.} \dicDirectTranslationCS{drsný, hrubý (hlas ap.)}
\dicEntry[hroðalegur] \dicTerm{hroða··legur} \dicIPA{{\textsubring{r}}{\textopeno}{\textlengthmark}{ð}{a}{l}{\textepsilon}{\textbabygamma}{\textscy}{\textsubring{r}}} \dicPos{adj}[1]\dicFlx{}[-8] \dicSynonym{óhugnanlegur} \dicDirectTranslationCS{děsivý, úděsný, hrůzný}
\dicEntry[hroðvirkni] \dicTerm{hroð··virkn|i} \dicIPA{{\textsubring{r}}{\textopeno}{ð}{v}{\textsci}{\textsubring{r}}{\r{g}}{n}{\textsci}} \dicPos{f}[3] \dicFlx{(‑i)}[3] \dicDirectTranslationCS{lajdáckost, nedbalost}
\dicEntry[hroðvirkur] \dicTerm{hroð··virkur} \dicIPA{{\textsubring{r}}{\textopeno}{ð}{v}{\textsci}{\textsubring{r}}{\r{g}}{\textscy}{\textsubring{r}}} \dicPos{adj}[1]\dicFlx{}[-1] \dicDirectTranslationCS{lajdácký, nedbalý}
\dicEntry[hrogn] \dicTerm{hrogn} \dicIPA{{\textsubring{r}}{\textopeno}{\r{g}}{\textsubring{n}}} \dicPos{n}[2] \dicFlx{(‑s, ‑)}[5] \dicFieldCat{biol.} \dicDirectTranslationCS{jikra}
\dicEntry[hrognamál] \dicTerm{hrogna··mál} \dicIPA{{\textsubring{r}}{\textopeno}{\r{g}}{n}{a}{m}{au}{\textsubring{l}}} \dicPos{n}[2] \dicFlx{(‑s, ‑)}[5] \dicDirectTranslationCS{hatmatilka, hantýrka}
\dicEntry[hrognkelsi] \dicTerm{hrogn··kelsi} \dicIPA{{\textsubring{r}}{\textopeno}{\r{g}}{\textsubring{n}}{c\smash{\textsuperscript{h}}}{\textepsilon}{l}{s}{\textsci}} \dicPos{n}[2] \dicFlx{(‑s, ‑)}[14] \dicFieldCat{zool.} \dicDirectTranslationCS{hranáč šedý} \textit{(l.~{\textLA{Cyclopterus lumpus}})}  \dicsymPhoto\ 
\dicFigure{26717.jpg}{Hrognkelsi}{Hrognkelsi - Dr. Aleksey Zuyev, Biolib, PD}
\dicEntry[hrokafullur] \dicTerm{hroka··|fullur} \dicIPA{{\textsubring{r}}{\textopeno}{\textlengthmark}{\r{g}}{a}{f}{\textscy}{\textsubring{d}}{l}{\textscy}{\textsubring{r}}} \dicPos{adj}[10] \dicFlx{(comp ‑fyllri, sup ‑fyllstur)}[7] \dicDirectTranslationCS{arogantní, domýšlivý}
\dicEntry[hroki] \dicTerm{hrok|i} \dicIPA{{\textsubring{r}}{\textopeno}{\textlengthmark}{\r{\textObardotlessj}}{\textsci}} \dicPos{m}[1] \dicFlx{(‑a)}[3] \dicDirectTranslationCS{arogance, domýšlivost}
\dicEntry[hrokkið] \dicTerm{hrokkið} \dicIPA{{\textsubring{r}}{\textopeno}{h}{\r{\textObardotlessj}}{\textsci}{\texttheta}} \dicPos{v} \dicFlx{supin} \dicLink{hrökkva}
\dicEntry[hrokkinn] \dicTerm{hrokkinn} \dicIPA{{\textsubring{r}}{\textopeno}{h}{\r{\textObardotlessj}}{\textsci}{\textsubring{n}}} \dicPos{adj}[6]\dicFlx{}[-2] \dicDirectTranslationCS{kudrnatý, kadeřavý, kučeravý}
\dicEntry[hrollur] \dicTerm{hroll|ur} \dicsymFrequent\  \dicIPA{{\textsubring{r}}{\textopeno}{\textsubring{d}}{l}{\textscy}{\textsubring{r}}} \dicPos{m}[6] \dicFlx{(‑s)}[7] \dicDirectTranslationCS{(za)chvění, (za)trnutí, zimnice} \dicExampleIS{Einkenni sjúkdómsins eru uppköst og hrollur.} \dicExampleCS{Příznaky choroby jsou zvracení a~zimnice.};  \dicPhraseIS{það er hrollur í e‑m} \dicFlx{impers} \dicDirectTranslationCS{(kdo) se třese (zimou ap.)};  \dicPhraseIS{það fer hrollur um e‑n} \dicFlx{impers} \dicDirectTranslationCS{(komu) běhá mráz po zádech}
\dicEntry[hrollvekja] \dicTerm{hroll··vekj|a} \dicIPA{{\textsubring{r}}{\textopeno}{\textsubring{d}}{l}{v}{\textepsilon}{\r{\textObardotlessj}}{a}} \dicPos{f}[1] \dicFlx{(‑u, ‑ur)}[25] \textbf{1.} \dicDirectTranslationCS{horor, hrůza}  \textbf{2.} \dicDirectTranslationCS{horor, thriller} \dicIndirectTranslationCS{(kniha, film ap.)}
\dicEntry[hrollvekjandi] \dicTerm{hroll··vekj·andi} \dicIPA{{\textsubring{r}}{\textopeno}{\textsubring{d}}{l}{v}{\textepsilon}{\r{\textObardotlessj}}{a}{n}{\textsubring{d}}{\textsci}} \dicPos{adj}[13] \dicFlx{indecl}[1] \dicDirectTranslationCS{děsivý, nahánějící hrůzu}
\dicEntry[hrosið] \dicTerm{hrosið} \dicIPA{{\textsubring{r}}{\textopeno}{\textlengthmark}{s}{\textsci}{\texttheta}} \dicPos{v} \dicFlx{supin} \dicLink{hrjósa}
\dicEntry[hross] \dicTerm{hross} \dicIPA{{\textsubring{r}}{\textopeno}{s}{\textlengthmark}} \dicPos{n}[2] \dicFlx{(‑, ‑)}[22] \dicDirectTranslationCS{kůň} \dicIndirectTranslationCS{(všeobecné označení bez rozdílu pohlaví)}
\dicEntry[hrossafluga] \dicTerm{hrossa··flug|a} \dicIPA{{\textsubring{r}}{\textopeno}{s}{\textlengthmark}{a}{f}{l}{\textscy}{\textbabygamma}{a}} \dicPos{f}[1] \dicFlx{(‑u, ‑ur)}[13] \dicFieldCat{zool.} \dicDirectTranslationCS{tiplice} \textit{(l.~{\textLA{Tipula}})}  \dicsymPhoto\ 
\dicFigure{ds_image_hrossafluga_0_1.jpg}{Hrossafluga}{Hrossafluga - James K. Lindsey, CC BY-SA 2.5}
\dicEntry[hrossagaukur] \dicTerm{hrossa··gauk|ur} \dicIPA{{\textsubring{r}}{\textopeno}{s}{\textlengthmark}{a}{\r{g}}{\oe i}{\r{g}}{\textscy}{\textsubring{r}}} \dicPos{m}[6] \dicFlx{(‑s, ‑ar)}[22] \dicFieldCat{zool.} \dicDirectTranslationCS{bekasina otavní} \textit{(l.~{\textLA{Gallinago gallinago}})}  \dicsymPhoto\ 
\dicFigure{20699.jpg}{Hrossagaukur}{Hrossagaukur - Ward, David, Biolib, PD}
\dicEntry[hrossakaup] \dicTerm{hrossa··kaup} \dicIPA{{\textsubring{r}}{\textopeno}{s}{\textlengthmark}{a}{k\smash{\textsuperscript{h}}}{\oe i}{\textsubring{b}}} \dicPos{n}[2] \dicFlx{pl}[1] \textbf{1.} \dicSynonym*{verslun með hesta} \dicDirectTranslationCS{koňský handl}  \textbf{2.} \dicSynonym*{siðlaus viðskipti} \dicDirectTranslationCS{handl, nepoctivý obchod}
\dicEntry[hrossalækning] \dicTerm{hrossa··lækn·ing} \dicIPA{{\textsubring{r}}{\textopeno}{s}{\textlengthmark}{a}{l}{a}{i}{h}{\r{g}}{n}{i}{\ng}{\r{g}}} \dicPos{f}[4] \dicFlx{(‑ar, ‑ar)}[5] \textbf{1.} \dicSynonym*{lækning hesta} \dicDirectTranslationCS{léčení koní}  \textbf{2.} \dicLangCat{přen.} \dicSynonym*{gróf læknisaðgerð} \dicDirectTranslationCS{drastická léčba}
\dicEntry[hrossanál] \dicTerm{hrossa··nál} \dicIPA{{\textsubring{r}}{\textopeno}{s}{\textlengthmark}{a}{n}{au}{\textsubring{l}}} \dicPos{f}[4] \dicFlx{(‑ar, ‑ar)}[1] \dicFieldCat{bot.} \dicDirectTranslationCS{sítina arktická} \textit{(l.~{\textLA{Juncus arcticus}})}
\dicEntry[hrossaþari] \dicTerm{hrossa··þar|i} \dicIPA{{\textsubring{r}}{\textopeno}{s}{\textlengthmark}{a}{\texttheta}{a}{r}{\textsci}} \dicPos{m}[1] \dicFlx{(‑a)}[3] \dicFieldCat{bot.} \dicDirectTranslationCS{čepelatka prstnatá} \textit{(l.~{\textLA{Laminaria digitata}})}  \dicsymPhoto\ 
\dicFigure{ds_image_hrossathari_0_2.jpg}{Hrossaþari}{Hrossaþari - Stemonitis, CC BY 2.5}
\dicEntry[hrota] \dicTerm{hrot|a} \dicIPA{{\textsubring{r}}{\textopeno}{\textlengthmark}{\textsubring{d}}{a}} \dicPos{f}[1] \dicFlx{(‑u, ‑ur)}[19] \dicPhraseIS{hrotur} \dicFlx{pl} \dicDirectTranslationCS{chrápání}
\dicEntry[hrotið] \dicTerm{hrotið} \dicIPA{{\textsubring{r}}{\textopeno}{\textlengthmark}{\textsubring{d}}{\textsci}{\texttheta}} \dicPos{v} \dicFlx{supin} \dicLink{hrjóta}
\dicEntry[hrottalegur] \dicTerm{hrotta··legur} \dicIPA{{\textsubring{r}}{\textopeno}{h}{\textsubring{d}}{a}{l}{\textepsilon}{\textbabygamma}{\textscy}{\textsubring{r}}} \dicPos{adj}[1]\dicFlx{}[-8] \dicDirectTranslationCS{brutální, drsný} \dicExampleIS{hrottaleg árás} \dicExampleCS{brutální útok}
\dicEntry[hrotti] \dicTerm{hrott|i} \dicIPA{{\textsubring{r}}{\textopeno}{h}{\textsubring{d}}{\textsci}} \dicPos{m}[1] \dicFlx{(‑a, ‑ar)}[1] \dicDirectTranslationCS{násilník, surovec, brutální člověk}
\dicEntry[Hróarskelda] \dicTerm{Hróars··keld|a} \dicIPA{{\textsubring{r}}{ou}{a}{\textsubring{r}}{s}{\r{\textObardotlessj}}{\textepsilon}{l}{\textsubring{d}}{a}} \dicPos{f}[1] \dicFlx{(‑u)}[6] \dicFieldCat{geog.} \dicDirectTranslationCS{Roskilde} \dicIndirectTranslationCS{(město v~Dánsku)}
\dicEntry[hróðugur] \dicTerm{hróð··ugur} \dicIPA{{\textsubring{r}}{ou}{\textlengthmark}{ð}{\textscy}{\textbabygamma}{\textscy}{\textsubring{r}}} \dicPos{adj}[1]\dicFlx{}[-8] \dicSynonym{hreykinn} \dicDirectTranslationCS{spokojený sám se sebou, samolibý}
\dicEntry[hróður] \dicTerm{hróð|ur} \dicIPA{{\textsubring{r}}{ou}{\textlengthmark}{ð}{\textscy}{\textsubring{r}}} \dicPos{m}[5] \dicFlx{(‑urs\,/\addthin ‑rar)}[10] \dicSynonym{frægð} \dicDirectTranslationCS{věhlas, reputace}
\dicEntry[hrófla] \dicTerm{hrófl|a} \dicIPA{{\textsubring{r}}{ou}{\textsubring{b}}{l}{a}} \dicPos{v}[1] \dicFlx{(‑aði)}[1] \textbf{1.} \dicPhraseIS{hrófla við e‑u} \dicDirectTranslationCS{hýbat s~(čím), pohnout s~(čím), pozměnit (co)}  \textbf{2.} \dicPhraseIS{hrófla upp e‑ð} \dicDirectTranslationCS{sbít\,/\addthin stlouct (co) (lajdácky)}
\dicEntry[hróka] \dicTerm{hrók|a} \dicIPA{{\textsubring{r}}{ou}{\textlengthmark}{\r{g}}{a}} \dicPos{v}[1] \dicFlx{(‑aði)}[1] \dicLink{hrókera}
\dicEntry[hrókera] \dicTerm{hróker|a}\dicTerm{, hróka} \dicIPA{{\textsubring{r}}\-{ou}\-{\textlengthmark}\-{\r{g}}\-{\textepsilon}\-{r}\-{a}\-} \dicPos{v}[1] \dicFlx{(‑aði)}[1] \textbf{1.} \dicDirectTranslationCS{udělat rošádu} \dicIndirectTranslationCS{(v~šachách)}  \textbf{2.} \dicLangCat{přen.} \dicDirectTranslationCS{udělat rošádu (na ministerstvu ap.)}
\dicEntry[hrókur] \dicTerm{hrók|ur} \dicIPA{{\textsubring{r}}{ou}{\textlengthmark}{\r{g}}{\textscy}{\textsubring{r}}} \dicPos{m}[6] \dicFlx{(‑s, ‑ar)}[22] \dicDirectTranslationCS{věž} \dicIndirectTranslationCS{(šachová figura)};  \dicPhraseIS{vera hrókur alls fagnaðar} \dicLangCat{přen.} \dicDirectTranslationCS{být duše a~srdce společnosti}
\dicEntry[hrókönd] \dicTerm{hrók··|önd} \dicIPA{{\textsubring{r}}{ou}{\textlengthmark}{\r{g}}{\oe}{n}{\textsubring{d}}} \dicPos{f}[8] \dicFlx{(‑andar, ‑endur)}[2] \dicFieldCat{zool.} \dicDirectTranslationCS{kachnice kaštanová} \textit{(l.~{\textLA{Oxyura jamaicensis}})}  \dicsymPhoto\ 
\dicFigure{20825.jpg}{Hrókönd}{Hrókönd - Hollingsworth, John and Karen, Biolib, PD}
\dicEntry[hróp] \dicTerm{hróp} \dicsymFrequent\  \dicIPA{{\textsubring{r}}{ou}{\textlengthmark}{\textsubring{b}}} \dicPos{n}[2] \dicFlx{(‑s, ‑)}[5] \dicDirectTranslationCS{(vý)křik, pokřik} \dicExampleIS{Á næsta augnabliki heyrðist hróp úr húsinu.} \dicExampleCS{V~následujícím okamžiku byl z~domu slyšet výkřik.};  \dicPhraseIS{gera hróp að e‑m} \dicDirectTranslationCS{pokřikovat na (koho)}
\dicEntry[hrópa] \dicTerm{hróp|a} \dicsymFrequent\  \dicIPA{{\textsubring{r}}{ou}{\textlengthmark}{\textsubring{b}}{a}} \dicPos{v}[1] \dicFlx{(‑aði)}[1] \dicDirectTranslationCS{(vy)křiknout, křičet, hulákat, zařvat, řvát} \dicExampleIS{hrópa á hjálp} \dicExampleCS{křičet o~pomoc};  \dicPhraseIS{hrópa upp yfir sig} \dicDirectTranslationCS{dát se do křiku}
\dicEntry[hróplegur] \dicTerm{hróp··legur} \dicIPA{{\textsubring{r}}{ou}{\textlengthmark}{\textsubring{b}}{l}{\textepsilon}{\textbabygamma}{\textscy}{\textsubring{r}}} \dicPos{adj}[1]\dicFlx{}[-8] \dicDirectTranslationCS{křiklavý, bijící do očí} \dicExampleIS{hróplegt ósamræmi} \dicExampleCS{křiklavý rozpor}
\dicEntry[hrós] \dicTerm{hrós} \dicIPA{{\textsubring{r}}{ou}{\textlengthmark}{s}} \dicPos{n}[2] \dicFlx{(‑s, ‑)}[5] \dicDirectTranslationCS{pochvala} \dicExampleIS{fá hrós fyrir e‑ð} \dicExampleCS{dostat pochvalu za (co)}
\dicEntry[hrósa] \dicTerm{hrós|a} \dicsymFrequent\  \dicIPA{{\textsubring{r}}{ou}{\textlengthmark}{s}{a}} \dicPos{v}[1] \dicFlx{(‑aði)}[1] \dicFlx{dat} \dicDirectTranslationCS{(po)chválit} \dicExampleIS{hrósa e‑m fyrir e‑ð} \dicExampleCS{pochválit (koho) za (co)};  \dicPhraseIS{hrósa sigri} \dicDirectTranslationCS{zvítězit, triumfovat}
\dicEntry[hrufl] \dicTerm{hrufl}\dicTerm{, hrufla\smash{\textsuperscript{1}}} \dicIPA{{\textsubring{r}}\-{\textscy}\-{\textsubring{b}}\-{\textsubring{l}}\-} \dicPos{n}[2] \dicFlx{(‑s)}[2] \dicSynonym{skeina\smash{\textsuperscript{1}}} \dicDirectTranslationCS{odření, odřenina}
\dicEntry[hrufla] \dicTerm{hrufl|a\smash{\textsuperscript{1}}} \dicIPA{{\textsubring{r}}{\textscy}{\textsubring{b}}{l}{a}} \dicPos{f}[1] \dicFlx{(‑u, ‑ur)}[19] \dicLink{hrufl}
\dicEntry[hrufla] \dicTerm{hrufl|a\smash{\textsuperscript{2}}} \dicIPA{{\textsubring{r}}{\textscy}{\textsubring{b}}{l}{a}} \dicPos{v}[1] \dicFlx{(‑aði)}[1] \dicFlx{acc} \dicDirectTranslationCS{odřít (si)} \dicExampleIS{hrufla bæði hnén} \dicExampleCS{odřít si obě kolena};  \dicPhraseIS{hrufla sig} \dicDirectTranslationCS{odřít se}
\dicEntry[hrufóttur] \dicTerm{hruf··óttur} \dicIPA{{\textsubring{r}}{\textscy}{\textlengthmark}{v}{ou}{h}{\textsubring{d}}{\textscy}{\textsubring{r}}} \dicPos{adj}[1]\dicFlx{}[-10] \dicSynonym{ósléttur} \dicDirectTranslationCS{zvrásněný, hrbolatý}
\dicEntry[hrukka] \dicTerm{hrukk|a\smash{\textsuperscript{1}}} \dicIPA{{\textsubring{r}}{\textscy}{h}{\r{g}}{a}} \dicPos{f}[1] \dicFlx{(‑u, ‑ur)}[19] \dicDirectTranslationCS{vráska, zvrásnění} \dicExampleIS{hrukka á enninu} \dicExampleCS{vráska na čele}
\dicEntry[hrukka] \dicTerm{hrukk|a\smash{\textsuperscript{2}}} \dicIPA{{\textsubring{r}}{\textscy}{h}{\r{g}}{a}} \dicPos{v}[1] \dicFlx{(‑aði)}[1] \dicDirectTranslationCS{zvrásnit, svraštit} \dicExampleIS{hrukka ennið} \dicExampleCS{svraštit čelo};  \dicIdiom{hrukkast}{ \dicPhraseIS{hrukkast}} \dicFlx{refl} \dicDirectTranslationCS{svraštit se, zvrásnit se} \dicExampleIS{Húðin hrukkast.} \dicExampleCS{Pleť se svrašťuje.}
\dicEntry[hrukkóttur] \dicTerm{hrukk··óttur} \dicIPA{{\textsubring{r}}{\textscy}{h}{\r{g}}{ou}{h}{\textsubring{d}}{\textscy}{\textsubring{r}}} \dicPos{adj}[1]\dicFlx{}[-10] \dicDirectTranslationCS{vrásčitý, zvrásněný}
\dicEntry[hrukkum] \dicTerm{hrukkum} \dicIPA{{\textsubring{r}}{\textscy}{h}{\r{g}}{\textscy}{\textsubring{m}}} \dicPos{v} \dicFlx{ind pf pl 1 pers} \dicLink{hrökkva}
\dicEntry[hrumur] \dicTerm{hrumur} \dicIPA{{\textsubring{r}}{\textscy}{\textlengthmark}{m}{\textscy}{\textsubring{r}}} \dicPos{adj}[1]\dicFlx{}[-1] \dicDirectTranslationCS{sešlý (věkem), vetchý}
\dicEntry[hrun] \dicTerm{hrun} \dicIPA{{\textsubring{r}}{\textscy}{\textlengthmark}{\textsubring{n}}} \dicPos{n}[2] \dicFlx{(‑s, ‑)}[5] \textbf{1.} \dicSynonym*{það að e‑að hrynur} \dicDirectTranslationCS{zhroucení, zřícení, zborcení}  \textbf{2.} \dicSynonym{verðfall} \dicDirectTranslationCS{(finanční) krach\,/\addthin kolaps}
\dicEntry[hrundi] \dicTerm{hrundi} \dicIPA{{\textsubring{r}}{\textscy}{n}{\textsubring{d}}{\textsci}} \dicPos{v} \dicFlx{ind pf sg 1 pers} \dicLink{hrynja}
\dicEntry[hrundið] \dicTerm{hrundið} \dicIPA{{\textsubring{r}}{\textscy}{n}{\textsubring{d}}{\textsci}{\texttheta}} \dicPos{v} \dicFlx{supin} \dicLink{hrinda}
\dicEntry[hrundum] \dicTerm{hrundum} \dicIPA{{\textsubring{r}}{\textscy}{n}{\textsubring{d}}{\textscy}{\textsubring{m}}} \dicPos{v} \dicFlx{ind pf pl 1 pers} \textbf{1.} \dicLink{hrinda}  \textbf{2.} \dicLink{hrynja}
\dicEntry[hrunið] \dicTerm{hrunið} \dicIPA{{\textsubring{r}}{\textscy}{\textlengthmark}{n}{\textsci}{\texttheta}} \dicPos{v} \dicFlx{supin} \dicLink{hrynja}
\dicEntry[hrutum] \dicTerm{hrutum} \dicIPA{{\textsubring{r}}{\textscy}{\textlengthmark}{\textsubring{d}}{\textscy}{\textsubring{m}}} \dicPos{v} \dicFlx{ind pf pl 1 pers} \dicLink{hrjóta}
\dicEntry[hrúður] \dicTerm{hrúður} \dicIPA{{\textsubring{r}}{u}{\textlengthmark}{ð}{\textscy}{\textsubring{r}}} \dicPos{n}[2] \dicFlx{(‑s, ‑)}[25] \dicDirectTranslationCS{strup} \dicExampleIS{hrúður á sári} \dicExampleCS{strup na ráně}
\dicEntry[hrúðurkarl] \dicTerm{hrúður··karl} \dicIPA{{\textsubring{r}}{u}{\textlengthmark}{ð}{\textscy}{\textsubring{r}}{k\smash{\textsuperscript{h}}}{a}{r}{\textsubring{d}}{\textsubring{l}}} \dicPos{m}[4] \dicFlx{(‑s, ‑ar)}[4] \dicFieldCat{zool.} \dicDirectTranslationCS{svijonožec} \textit{(l.~{\textLA{Balanus}})}
\dicEntry[hrúga] \dicTerm{hrúg|a\smash{\textsuperscript{1}}} \dicIPA{{\textsubring{r}}{u}{\textlengthmark}{a}} \dicPos{f}[1] \dicFlx{(‑u, ‑ur)}[19] \dicSynonym{dyngja} \dicDirectTranslationCS{hromádka, kupka} \dicExampleIS{fatahrúga} \dicExampleCS{hromádka prádla}
\dicEntry[hrúga] \dicTerm{hrúg|a\smash{\textsuperscript{2}}} \dicIPA{{\textsubring{r}}{u}{\textlengthmark}{a}} \dicPos{v}[1] \dicFlx{(‑aði)}[1] \dicFlx{dat} \dicDirectTranslationCS{(na)kupit, (na)hromadit} \dicExampleIS{hrúga e‑u upp} \dicExampleCS{nahromadit (co)};  \dicPhraseIS{hrúga e‑u saman} \dicDirectTranslationCS{(na)hromadit (co), (na)shromáždit (co)};  \dicIdiom{hrúgast}{ \dicPhraseIS{hrúgast}} \dicFlx{refl} \dicDirectTranslationCS{(na)kupit se, (na)shromáždit se (problémy ap.)}
\dicEntry[hrúgald] \dicTerm{hrúg··|ald} \dicIPA{{\textsubring{r}}{u}{\textlengthmark}{a}{l}{\textsubring{d}}} \dicPos{n}[2] \dicFlx{(‑alds, ‑öld)}[8] \dicDirectTranslationCS{(neurovnaná) hromada\,/\addthin kupa, změť}
\dicEntry[hrútaber] \dicTerm{hrúta··ber} \dicIPA{{\textsubring{r}}{u}{\textlengthmark}{\textsubring{d}}{a}{\textsubring{b}}{\textepsilon}{\textsubring{r}}} \dicPos{n}[2] \dicFlx{(‑s, ‑)}[13] \dicDirectTranslationCS{ostružina} \dicIndirectTranslationCS{(plod ostružiníku skalního)}
\dicEntry[hrútaberjaklungur] \dicTerm{hrúta·berja··klungur} \dicIPA{{\textsubring{r}}{u}{\textlengthmark}{\textsubring{d}}{a}{\textsubring{b}}{\textepsilon}{r}{j}{a}{k\smash{\textsuperscript{h}}}{l}{u}{\ng}{\r{g}}{\textscy}{\textsubring{r}}} \dicPos{n}[2] \dicFlx{(‑s, ‑)}[25] \dicLink{hrútaberjalyng}
\dicEntry[hrútaberjalyng] \dicTerm{hrúta·berja··lyng}\dicTerm{, hrútaberjaklungur} \dicIPA{{\textsubring{r}}\-{u}\-{\textlengthmark}\-{\textsubring{d}}\-{a}\-{\textsubring{b}}\-{\textepsilon}\-{r}\-{j}\-{a}\-{l}\-{i}\-{\ng}\-{\r{g}}\-} \dicPos{n}[2] \dicFlx{(‑s)}[2] \dicFieldCat{bot.} \dicDirectTranslationCS{ostružiník skalní} \textit{(l.~{\textLA{Rubus saxatilis}})}
\dicEntry[hrútur] \dicTerm{hrút|ur} \dicIPA{{\textsubring{r}}{u}{\textlengthmark}{\textsubring{d}}{\textscy}{\textsubring{r}}} \dicPos{m}[6] \dicFlx{(‑s, ‑ar)}[22] \textbf{1.} \dicDirectTranslationCS{beran}  \textbf{2.} \dicDirectTranslationCS{Beran} \dicIndirectTranslationCS{(znamení zvěrokruhu)} \textit{(l.~{\textLA{Aries}})}
\dicEntry[hryðja] \dicTerm{hryðj|a} \dicIPA{{\textsubring{r}}{\textsci}{ð}{j}{a}} \dicPos{f}[1] \dicFlx{(‑u, ‑ur)}[7] \dicDirectTranslationCS{smršť, průtrž (mračen)}
\dicEntry[hryðjuverk] \dicTerm{hryðju··verk} \dicIPA{{\textsubring{r}}{\textsci}{ð}{j}{\textscy}{v}{\textepsilon}{\textsubring{r}}{\r{g}}} \dicPos{n}[2] \dicFlx{(‑s, ‑)}[5] \dicDirectTranslationCS{teroristický útok\,/\addthin čin}
\dicEntry[hryðjuverkamaður] \dicTerm{hryðju·verka··|maður} \dicIPA{{\textsubring{r}}{\textsci}{ð}{j}{\textscy}{v}{\textepsilon}{\textsubring{r}}{\r{g}}{a}{m}{a}{ð}{\textscy}{\textsubring{r}}} \dicPos{m}[13] \dicFlx{(‑manns, ‑menn)}[2] \dicDirectTranslationCS{terorista, teroristka}
\dicEntry[hryðjuverkastarfsemi] \dicTerm{hryðju·verka··starf·sem|i} \dicIPA{{\textsubring{r}}\-{\textsci}\-{ð}\-{j}\-{\textscy}\-{v}\-{\textepsilon}\-{\textsubring{r}}\-{\r{g}}\-{a}\-{s}\-{\textsubring{d}}\-{a}\-{r}\-{v}\-{s}\-{\textepsilon}\-{m}\-{\textsci}\-} \dicPos{f}[2] \dicFlx{(‑i)}[2] \dicDirectTranslationCS{terorismus}
\dicEntry[hryggbraut] \dicTerm{hrygg··braut} \dicIPA{{\textsubring{r}}{\textsci}{\r{g}}{\textsubring{b}}{r}{\oe i}{\textsubring{d}}} \dicPos{v} \dicFlx{ind pf sg 1 pers} \dicLink{hryggbrjóta}
\dicEntry[hryggbrjóta] \dicTerm{hrygg··|brjóta} \dicIPA{{\textsubring{r}}{\textsci}{\r{g}}{\textsubring{b}}{r}{j}{ou}{\textsubring{d}}{a}} \dicPos{v}[6] \dicFlx{(‑brýt, ‑braut, ‑brutum, ‑bryti, ‑brotið)}[51] \dicFlx{acc} \textbf{1.} \dicSynonym*{brjóta hrygg} \dicDirectTranslationCS{zlomit vaz} \dicExampleIS{hryggbrjóta sig} \dicExampleCS{zlomit si vaz}  \textbf{2.} \dicSynonym*{hafna bónorði e‑rs} \dicDirectTranslationCS{odmítnout, odbýt (nápadníka ap.)} \dicExampleIS{hryggbrjóta e‑n} \dicExampleCS{odmítnout (koho)}
\dicEntry[hryggbrotið] \dicTerm{hrygg··brotið} \dicIPA{{\textsubring{r}}{\textsci}{\r{g}}{\textsubring{b}}{r}{\textopeno}{\textsubring{d}}{\textsci}{\texttheta}} \dicPos{v} \dicFlx{supin} \dicLink{hryggbrjóta}
\dicEntry[hryggbrutum] \dicTerm{hrygg··brutum} \dicIPA{{\textsubring{r}}{\textsci}{\r{g}}{\textsubring{b}}{r}{\textscy}{\textsubring{d}}{\textscy}{\textsubring{m}}} \dicPos{v} \dicFlx{ind pf pl 1 pers} \dicLink{hryggbrjóta}
\dicEntry[hryggbryti] \dicTerm{hrygg··bryti} \dicIPA{{\textsubring{r}}{\textsci}{\r{g}}{\textsubring{b}}{r}{\textsci}{\textsubring{d}}{\textsci}} \dicPos{v} \dicFlx{con pf sg 1 pers} \dicLink{hryggbrjóta}
\dicEntry[hryggbrýt] \dicTerm{hrygg··brýt} \dicIPA{{\textsubring{r}}{\textsci}{\r{g}}{\textsubring{b}}{r}{i}{\textsubring{d}}} \dicPos{v} \dicFlx{ind praes sg 1 pers} \dicLink{hryggbrjóta}
\dicEntry[hryggdýr] \dicTerm{hrygg··dýr} \dicIPA{{\textsubring{r}}{\textsci}{\r{g}}{\textsubring{d}}{i}{\textsubring{r}}} \dicPos{n}[2] \dicFlx{(‑s, ‑)}[5] \dicFieldCat{zool.} \dicDirectTranslationCS{obratlovec} \textit{(l.~{\textLA{Vertebrata}})}
\dicEntry[hryggð] \dicTerm{hryggð} \dicIPA{{\textsubring{r}}{\textsci}{\textbabygamma}{\texttheta}} \dicPos{f}[7] \dicFlx{(‑ar, ‑ir)}[1] \dicSynonym{sorg} \dicDirectTranslationCS{zármutek, žal}
\dicEntry[hryggilegur] \dicTerm{hryggi··legur} \dicIPA{{\textsubring{r}}{\textsci}{\r{\textObardotlessj}}{\textlengthmark}{\textsci}{l}{\textepsilon}{\textbabygamma}{\textscy}{\textsubring{r}}} \dicPos{adj}[1]\dicFlx{}[-8] \dicSynonym{dapurlegur} \dicDirectTranslationCS{žalostný, smutný}
\dicEntry[hryggja] \dicTerm{hrygg|ja} \dicIPA{{\textsubring{r}}{\textsci}{\r{\textObardotlessj}}{\textlengthmark}{a}} \dicPos{v}[2] \dicFlx{(‑ði, ‑t)}[89] \dicFlx{acc} \dicDirectTranslationCS{(za)rmoutit, způsobit smutek} \dicExampleIS{hryggja e‑n} \dicExampleCS{zarmoutit (koho)};  \dicPhraseIS{það hryggir e‑n að (heyra þetta)} \dicFlx{impers} \dicDirectTranslationCS{(koho) rmoutí (slyšet to)};  \dicIdiom{hryggjast}{ \dicPhraseIS{hryggjast}} \dicFlx{refl} \dicDirectTranslationCS{(za)rmoutit se}
\dicEntry[hryggleysingi] \dicTerm{hrygg··leys·ing|i} \dicIPA{{\textsubring{r}}{\textsci}{\r{g}}{l}{ei}{s}{i}{\textltailn}{\r{\textObardotlessj}}{\textsci}} \dicPos{m}[1] \dicFlx{(‑ja, ‑jar)}[14] \dicFieldCat{zool.} \dicDirectTranslationCS{bezobratlovec, bezobratlý} \textit{(l.~{\textLA{Invertebrata}})}
\dicEntry[hryggsúla] \dicTerm{hrygg··súl|a} \dicIPA{{\textsubring{r}}{\textsci}{\r{g}}{s}{u}{l}{a}} \dicPos{f}[1] \dicFlx{(‑u, ‑ur)}[13] \dicFieldCat{anat.} \dicDirectTranslationCS{páteř}
\dicEntry[hryggur] \dicTerm{hrygg|ur\smash{\textsuperscript{1}}} \dicsymFrequent\  \dicIPA{{\textsubring{r}}{\textsci}{\r{g}}{\textlengthmark}{\textscy}{\textsubring{r}}} \dicPos{m}[10] \dicFlx{(‑jar, ‑ir)}[23] \textbf{1.} \dicFieldCat{anat.} \dicSynonym{hryggsúla} \dicDirectTranslationCS{páteř} \dicExampleIS{slasast alvarlega á hrygg} \dicExampleCS{vážně si poranit páteř}  \textbf{2.} \dicSynonym{bak} \dicDirectTranslationCS{hřbet, záda};  \dicPhraseIS{fara á hrygginn} \dicDirectTranslationCS{spadnout na záda}  \textbf{3.} \dicSynonym*{fjallsás} \dicDirectTranslationCS{hřbet (hory ap.)}
\dicEntry[hryggur] \dicTerm{hryggur\smash{\textsuperscript{2}}} \dicIPA{{\textsubring{r}}{\textsci}{\r{g}}{\textlengthmark}{\textscy}{\textsubring{r}}} \dicPos{adj}[1]\dicFlx{}[-1] \dicDirectTranslationCS{zarmoucený, zasmušilý}
\dicEntry[hrygla] \dicTerm{hrygl|a\smash{\textsuperscript{1}}} \dicIPA{{\textsubring{r}}{\textsci}{\r{g}}{l}{a}} \dicPos{f}[1] \dicFlx{(‑u, ‑ur)}[19] \dicSynonym{korr} \dicDirectTranslationCS{chraptění, chrchlání}
\dicEntry[hrygla] \dicTerm{hrygl|a\smash{\textsuperscript{2}}} \dicIPA{{\textsubring{r}}{\textsci}{\r{g}}{l}{a}} \dicPos{v}[1] \dicFlx{(‑aði)}[75] \dicFlx{impers} \dicPhraseIS{það hryglar í e‑m} \dicDirectTranslationCS{(kdo) chrchlá\,/\addthin chraptí}
\dicEntry[hrygna] \dicTerm{hrygn|a} \dicIPA{{\textsubring{r}}{\textsci}{\r{g}}{n}{a}} \dicPos{v}[2] \dicFlx{(‑di, ‑t)}[155] \dicDirectTranslationCS{třít se} \dicIndirectTranslationCS{(vypouštět jikry nebo mlíčí k~oplodnění)}
\dicEntry[hrygning] \dicTerm{hrygn··ing} \dicIPA{{\textsubring{r}}{\textsci}{\r{g}}{n}{i}{\ng}{\r{g}}} \dicPos{f}[4] \dicFlx{(‑ar)}[7] \dicDirectTranslationCS{tření} \dicIndirectTranslationCS{(rozmnožování ryb)}
\dicEntry[hrykki] \dicTerm{hrykki} \dicIPA{{\textsubring{r}}{\textsci}{h}{\r{g}}{\textsci}} \dicPos{v} \dicFlx{con pf sg 1 pers} \dicLink{hrökkva}
\dicEntry[hrylla] \dicTerm{hryll|a} \dicIPA{{\textsubring{r}}{\textsci}{\textsubring{d}}{l}{a}} \dicPos{v}[2] \dicFlx{(‑ti, ‑t)}[79] \textbf{1.} \dicPhraseIS{e‑n hryllir við e‑u} \dicFlx{impers} \dicDirectTranslationCS{(komu) běhá mráz po zádech z~(čeho)}  \textbf{2.} \dicPhraseIS{hrylla sig} \dicDirectTranslationCS{třást se, klepat se (zimou ap.)}
\dicEntry[hryllilegur] \dicTerm{hrylli··legur} \dicIPA{{\textsubring{r}}{\textsci}{\textsubring{d}}{l}{\textsci}{l}{\textepsilon}{\textbabygamma}{\textscy}{\textsubring{r}}} \dicPos{adj}[1]\dicFlx{}[-8] \dicSynonym{voðalegur} \dicDirectTranslationCS{odporný, hrozný, hrůzný, otřesný} \dicExampleIS{hryllilegur glæpur} \dicExampleCS{odporný zločin}
\dicEntry[hryllingsmynd] \dicTerm{hryllings··mynd} \dicIPA{{\textsubring{r}}{\textsci}{\textsubring{d}}{l}{i}{\ng}{s}{m}{\textsci}{n}{\textsubring{d}}} \dicPos{f}[7] \dicFlx{(‑ar, ‑ir)}[1] \dicDirectTranslationCS{horor} \dicIndirectTranslationCS{(film)}
\dicEntry[hryllingssaga] \dicTerm{hryllings··|saga} \dicIPA{{\textsubring{r}}{\textsci}{\textsubring{d}}{l}{i}{\ng}{s}{a}{\textbabygamma}{a}} \dicPos{f}[1] \dicFlx{(‑sögu, ‑sögur)}[14] \dicDirectTranslationCS{strašidelný\,/\addthin hororový příběh, horor} \dicIndirectTranslationCS{(kniha)}
\dicEntry[hryn] \dicTerm{hryn} \dicIPA{{\textsubring{r}}{\textsci}{\textlengthmark}{\textsubring{n}}} \dicPos{v} \dicFlx{ind praes sg 1 pers} \dicLink{hrynja}
\dicEntry[hryndi] \dicTerm{hryndi} \dicIPA{{\textsubring{r}}{\textsci}{n}{\textsubring{d}}{\textsci}} \dicPos{v} \dicFlx{con pf sg 1 pers} \textbf{1.} \dicLink{hrinda}  \textbf{2.} \dicLink{hrynja}
\dicEntry[hrynja] \dicTerm{hrynja} \dicsymFrequent\  \dicIPA{{\textsubring{r}}{\textsci}{n}{j}{a}} \dicPos{v}[4] \dicFlx{(hryn, hrundi, hrundum, hryndi, hrunið)}[37] \textbf{1.} \dicDirectTranslationCS{zhroutit se, zbortit se} \dicExampleIS{Húsið hrundi eins og spilaborg.} \dicExampleCS{Dům se zbortil jako domeček z~karet.}  \textbf{2.} \dicDirectTranslationCS{zřítit se, spadnout};  \dicPhraseIS{það hrynur úr e‑u} \dicFlx{impers} \dicDirectTranslationCS{padají kameny z~(čeho)};  \dicIdiom{hrynja}[niður]{ \dicPhraseIS{hrynja niður}} \dicLangCat{přen.} \dicDirectTranslationCS{padat jako mouchy} \dicIndirectTranslationCS{(umírat ve velkém množství)}
\dicEntry[hrynjandi] \dicTerm{hrynj··|andi\smash{\textsuperscript{1}}} \dicIPA{{\textsubring{r}}{\textsci}{n}{j}{a}{n}{\textsubring{d}}{\textsci}} \dicPos{m}[2] \dicFlx{(‑anda)}[2] \dicLink{hrynjandi\smash{\textsuperscript{2}}}
\dicEntry[hrynjandi] \dicTerm{hrynj··and|i\smash{\textsuperscript{2}}}\dicTerm{, hrynjandi\smash{\textsuperscript{1}}} \dicIPA{{\textsubring{r}}\-{\textsci}\-{n}\-{j}\-{a}\-{n}\-{\textsubring{d}}\-{\textsci}\-} \dicPos{f}[3] \dicFlx{(‑i)}[3] \dicSynonym*{hljóðfall} \dicDirectTranslationCS{takt, rytmus, melodie} \dicExampleIS{hrynjandi íslenskrar tungu} \dicExampleCS{rytmus islandského jazyka}
\dicEntry[hrysi] \dicTerm{hrysi} \dicIPA{{\textsubring{r}}{\textsci}{\textlengthmark}{s}{\textsci}} \dicPos{v} \dicFlx{con pf sg 1 pers} \dicLink{hrjósa}
\dicEntry[hryssa] \dicTerm{hryss|a} \dicIPA{{\textsubring{r}}{\textsci}{s}{\textlengthmark}{a}} \dicPos{f}[1] \dicFlx{(‑u, ‑ur)}[19] \dicSynonym{meri} \dicDirectTranslationCS{klisna, kobyla}
\dicEntry[hryssingslegur] \dicTerm{hryssings··legur} \dicIPA{{\textsubring{r}}{\textsci}{s}{\textlengthmark}{i}{\ng}{s}{l}{\textepsilon}{\textbabygamma}{\textscy}{\textsubring{r}}} \dicPos{adj}[1]\dicFlx{}[-8] \dicDirectTranslationCS{strohý, příkrý}
\dicEntry[hryti] \dicTerm{hryti} \dicIPA{{\textsubring{r}}{\textsci}{\textlengthmark}{\textsubring{d}}{\textsci}} \dicPos{v} \dicFlx{con pf sg 1 pers} \dicLink{hrjóta}
\dicEntry[hrýs] \dicTerm{hrýs} \dicIPA{{\textsubring{r}}{i}{\textlengthmark}{s}} \dicPos{v} \dicFlx{ind praes sg 1 pers} \dicLink{hrjósa}
\dicEntry[hrýt] \dicTerm{hrýt} \dicIPA{{\textsubring{r}}{i}{\textlengthmark}{\textsubring{d}}} \dicPos{v} \dicFlx{ind praes sg 1 pers} \dicLink{hrjóta}
\dicEntry[hræ] \dicTerm{hræ} \dicIPA{{\textsubring{r}}{a}{i}{\textlengthmark}} \dicPos{n}[2] \dicFlx{(‑s, ‑)}[13] \textbf{1.} \dicSynonym*{dautt dýr} \dicDirectTranslationCS{zdechlina, mršina}  \textbf{2.} \dicSynonym*{bílhræ} \dicDirectTranslationCS{autovrak}
\dicEntry[hræddur] \dicTerm{hræddur} \dicsymFrequent\  \dicIPA{{\textsubring{r}}{a}{i}{\textsubring{d}}{\textlengthmark}{\textscy}{\textsubring{r}}} \dicPos{adj}[2]\dicFlx{}[-18] \dicSynonym{banginn} \dicDirectTranslationCS{mající strach, polekaný, vystrašený} \dicExampleIS{vera hræddur við dauðann} \dicExampleCS{bát se smrti};  \dicPhraseIS{vera hræddur um e‑ð} \dicDirectTranslationCS{bát se o~(co)}
\dicEntry[hræða] \dicTerm{hræð|a\smash{\textsuperscript{1}}} \dicIPA{{\textsubring{r}}{a}{i}{\textlengthmark}{ð}{a}} \dicPos{f}[1] \dicFlx{(‑u, ‑ur)}[19] \dicDirectTranslationCS{osoba, živáček} \dicIndirectTranslationCS{(nejčastěji se záporem)};  \dicPhraseIS{ekki nokkur hræða} \dicDirectTranslationCS{ani živá duše}
\dicEntry[hræða] \dicTerm{hræ|ða\smash{\textsuperscript{2}}} \dicIPA{{\textsubring{r}}{a}{i}{\textlengthmark}{ð}{a}} \dicPos{v}[2] \dicFlx{(‑ddi, ‑tt)}[167] \dicFlx{acc} \dicDirectTranslationCS{nahnat strach, (vy)lekat, (vy)strašit} \dicExampleIS{hræða e‑n} \dicExampleCS{nahnat (komu) strach};  \dicIdiom{hræðast}{ \dicPhraseIS{hræðast e‑ð}} \dicFlx{refl} \dicDirectTranslationCS{mít z~(čeho) strach, bát se (čeho), lekat se (čeho)} \dicExampleIS{hræðast myrkrið} \dicExampleCS{bát se tmy}
\dicEntry[hræðilegur] \dicTerm{hræði··legur} \dicsymFrequent\  \dicIPA{{\textsubring{r}}{a}{i}{\textlengthmark}{ð}{\textsci}{l}{\textepsilon}{\textbabygamma}{\textscy}{\textsubring{r}}} \dicPos{adj}[1]\dicFlx{}[-8] \dicDirectTranslationCS{hrozný, strašný, děsivý} \dicExampleIS{hræðilegt slys} \dicExampleCS{strašná nehoda}
\dicEntry[hræðsla] \dicTerm{hræðsl|a} \dicsymFrequent\  \dicIPA{{\textsubring{r}}{a}{i}{ð}{s}{\textsubring{d}}{l}{a}} \dicPos{f}[1] \dicFlx{(‑u)}[5] \dicSynonym{beygur} \dicDirectTranslationCS{strach} \dicExampleIS{hræðsla við myrkrið} \dicExampleCS{strach ze tmy};  \dicPhraseIS{finna fyrir hræðslu} \dicDirectTranslationCS{mít strach};  \dicPhraseIS{nötra\,/\addthin skjálfa\,/\addthin titra af hræðslu} \dicDirectTranslationCS{klepat\,/\addthin třást se strachy};  \dicPhraseIS{vera að deyja úr hræðslu} \dicLangCat{přen.} \dicDirectTranslationCS{umírat strachy};  \dicPhraseIS{það grípur e‑n hræðsla} \dicFlx{impers} \dicDirectTranslationCS{(kdo) začíná mít strach}
\dicEntry[hræfugl] \dicTerm{hræ··fugl} \dicIPA{{\textsubring{r}}{a}{i}{\textlengthmark}{f}{\textscy}{\r{g}}{\textsubring{l}}} \dicPos{m}[4] \dicFlx{(‑s, ‑ar)}[3] \dicDirectTranslationCS{mrchožrout} \dicIndirectTranslationCS{(pták)}
\dicEntry[hrægammur] \dicTerm{hræ··gamm|ur} \dicIPA{{\textsubring{r}}{a}{i}{\textlengthmark}{\r{g}}{a}{m}{\textscy}{\textsubring{r}}} \dicPos{m}[6] \dicFlx{(‑s, ‑ar)}[5] \dicFieldCat{zool.} \dicDirectTranslationCS{sup} \textit{(l.~{\textLA{Aegypius}})}  \dicsymPhoto\ 
\dicFigure{ds_image_hraegammur_0_2.jpg}{Hrægammur}{Hrægammur - BirdingInSpain, GFDL}
\dicEntry[hrækja] \dicTerm{hræk|ja} \dicIPA{{\textsubring{r}}{a}{i}{\textlengthmark}{\r{\textObardotlessj}}{a}} \dicPos{v}[2] \dicFlx{(‑ti, ‑t)}[28] \textbf{1.} \dicSynonym{spýta\smash{\textsuperscript{2}}} \dicDirectTranslationCS{(od)plivnout si, plivat} \dicExampleIS{hrækja á götuna} \dicExampleCS{odplivnout si na ulici}  \textbf{2.} \dicPhraseIS{hrækja út úr sér e‑u} \dicDirectTranslationCS{chrlit ze sebe (co) (nadávky ap.)}
\dicEntry[hræódýr] \dicTerm{hræ··ó·dýr} \dicIPA{{\textsubring{r}}{a}{i}{\textlengthmark}{ou}{\textsubring{d}}{i}{\textsubring{r}}} \dicPos{adj}[5] \dicFlx{(f ‑)}[10] \dicDirectTranslationCS{strašně\,/\addthin neuvěřitelně levný}
\dicEntry[hræra] \dicTerm{hrær|a\smash{\textsuperscript{1}}} \dicIPA{{\textsubring{r}}{a}{i}{\textlengthmark}{r}{a}} \dicPos{f}[1] \dicFlx{(‑u, ‑ur)}[7] \dicDirectTranslationCS{směs, míchanice}
\dicEntry[hræra] \dicTerm{hrær|a\smash{\textsuperscript{2}}} \dicsymFrequent\  \dicIPA{{\textsubring{r}}{a}{i}{\textlengthmark}{r}{a}} \dicPos{v}[2] \dicFlx{(‑ði, ‑t)}[99] \dicFlx{acc\,/\addthin dat} \textbf{1.} \dicFlx{acc} \dicSynonym{hreyfa} \dicDirectTranslationCS{(po)hnout, hýbat, pohybovat};  \dicPhraseIS{hræra hvorki hönd né fót} \dicLangCat{přen.} \dicDirectTranslationCS{nehnout ani prstem};  \dicPhraseIS{hræra hvorki legg né lið} \dicLangCat{přen.} \dicDirectTranslationCS{nemoct se pohnout}  \textbf{2.} \dicFlx{dat} \dicDirectTranslationCS{(pro)míchat, zamíchat (polévku ap.)} \dicExampleIS{hræra í pottinum} \dicExampleCS{míchat v~hrnci};  \dicPhraseIS{hræra e‑u saman við e‑ð} \dicDirectTranslationCS{promíchat (co) s~(čím)};  \dicPhraseIS{hræra saman e‑ð} \dicDirectTranslationCS{promíchat (co)}  \textbf{3.} \dicDirectTranslationCS{hýbat, pohnout} \dicIndirectTranslationCS{(změnit polohu)}  \textbf{4.} \dicDirectTranslationCS{pohnout, dojmout, zasáhnout} \dicIndirectTranslationCS{(vyvolat pohnutí)} \dicExampleIS{hræra hjörtu} \dicExampleCS{pohnout srdcem};  \dicPhraseIS{hræra e‑n til meðaumkunar} \dicDirectTranslationCS{pohnout (koho) k~lítosti};  \dicIdiom{hrærast}{ \dicPhraseIS{hrærast}} \dicFlx{refl} {\textbf{a.}} \dicDirectTranslationCS{(po)hnout se, hýbat se};  {\textbf{b.}} \dicDirectTranslationCS{(pro)míchat se (co) (těsto ap.)};  {\textbf{c.}} \dicDirectTranslationCS{míchat se (pocity ap.)}
\dicEntry[hrærigrautur] \dicTerm{hræri··graut|ur} \dicIPA{{\textsubring{r}}{a}{i}{\textlengthmark}{r}{\textsci}{\r{g}}{r}{\oe i}{\textsubring{d}}{\textscy}{\textsubring{r}}} \dicPos{m}[6] \dicFlx{(‑ar, ‑ar)}[59] \dicSynonym{benda\smash{\textsuperscript{1}}} \dicDirectTranslationCS{směsice, míchanice}
\dicEntry[hræring] \dicTerm{hrær··ing} \dicIPA{{\textsubring{r}}{a}{i}{\textlengthmark}{r}{i}{\ng}{\r{g}}} \dicPos{f}[4] \dicFlx{(‑ar, ‑ar)}[5] \textbf{1.} \dicSynonym{hreyfing} \dicDirectTranslationCS{(po)hnutí}  \textbf{2.} \dicSynonym{geðshræring} \dicDirectTranslationCS{pohnutí (citové)}  \textbf{3.} \dicPhraseIS{hræringar} \dicFlx{pl} \dicFieldCat{geol.} \dicDirectTranslationCS{(slabé) otřesy}
\dicEntry[hrærivél] \dicTerm{hræri··vél} \dicIPA{{\textsubring{r}}{a}{i}{\textlengthmark}{r}{\textsci}{v}{j}{\textepsilon}{\textsubring{l}}} \dicPos{f}[4] \dicFlx{(‑ar, ‑ar)}[1] \textbf{1.} \dicDirectTranslationCS{(elektrický) mixér}  \textbf{2.} \dicDirectTranslationCS{míchačka, míchací stroj}
\dicEntry[hræsna] \dicTerm{hræsn|a} \dicIPA{{\textsubring{r}}{a}{i}{s}{\textsubring{d}}{n}{a}} \dicPos{v}[1] \dicFlx{(‑aði)}[1] \dicDirectTranslationCS{předstírat, hrát};  \dicPhraseIS{hræsna fyrir e‑m} \dicDirectTranslationCS{předstírat před (kým)}
\dicEntry[hræsnari] \dicTerm{hræsn··ar|i} \dicIPA{{\textsubring{r}}{a}{i}{s}{\textsubring{d}}{n}{a}{r}{\textsci}} \dicPos{m}[1] \dicFlx{(‑a, ‑ar)}[13] \dicDirectTranslationCS{pokrytec, pokrytkyně}
\dicEntry[hræsni] \dicTerm{hræsn|i} \dicIPA{{\textsubring{r}}{a}{i}{s}{\textsubring{d}}{n}{\textsci}} \dicPos{f}[3] \dicFlx{(‑i)}[3] \dicDirectTranslationCS{pokrytectví, licoměrnost}
\dicEntry[hrævareldur] \dicTerm{hrævar··eld|ur} \dicIPA{{\textsubring{r}}{a}{i}{\textlengthmark}{v}{a}{r}{\textepsilon}{l}{\textsubring{d}}{\textscy}{\textsubring{r}}} \dicPos{m}[6] \dicFlx{(‑s, ‑ar)}[4] \dicDirectTranslationCS{bludička} \dicIndirectTranslationCS{(pohyblivé světélko nad bažinami)}
\dicEntry[hröðun] \dicTerm{hröð|un} \dicIPA{{\textsubring{r}}{\oe}{\textlengthmark}{ð}{\textscy}{\textsubring{n}}} \dicPos{f}[7] \dicFlx{(‑unar)}[12] \dicDirectTranslationCS{zrychlení, akcelerace}
\dicEntry[hrökk] \dicTerm{hrökk} \dicIPA{{\textsubring{r}}{\oe}{h}{\r{g}}} \dicPos{v} \dicFlx{ind pf sg 1 pers} \dicLink{hrökkva}
\dicEntry[hrökkáll] \dicTerm{hrökk··ál|l} \dicIPA{{\textsubring{r}}{\oe}{h}{\r{g}}{au}{\textsubring{d}}{\textsubring{l}}} \dicPos{m}[6] \dicFlx{(‑s, ‑ar)}[48] \dicFieldCat{zool.} \dicDirectTranslationCS{paúhoř elektrický} \textit{(l.~{\textLA{Electrophorus electricus}})}  \dicsymPhoto\ 
\dicFigure{ds_image_hrokkall_0_1.jpg}{Hrökkáll}{Hrökkáll - Steven G. Johnson, CC BY-SA 3.0}
\dicEntry[hrökkbrauð] \dicTerm{hrökk··brauð} \dicIPA{{\textsubring{r}}{\oe}{h}{\r{g}}{\textsubring{b}}{r}{\oe i}{\texttheta}} \dicPos{n}[2] \dicFlx{(‑s, ‑)}[5] \dicFieldCat{kulin.} \dicDirectTranslationCS{křehký žitný chléb}
\dicEntry[hrökklast] \dicTerm{hrökkl|ast} \dicIPA{{\textsubring{r}}{\oe}{h}{\r{g}}{l}{a}{s}{\textsubring{d}}} \dicPos{v}[1] \dicFlx{(‑aðist)}[91] \dicFlx{refl} \dicSynonym*{hrekjast} \dicDirectTranslationCS{být odvolán\,/\addthin zvolněn}
\dicEntry[hrökkva] \dicTerm{hrökkva} \dicsymFrequent\  \dicIPA{{\textsubring{r}}{\oe}{h}{\r{g}}{v}{a}} \dicPos{v}[6] \dicFlx{(hrekk, hrökk, hrukkum, hrykki, hrokkið)}[91] \textbf{1.} \dicDirectTranslationCS{cuknout, škubnout, trhnout};  \dicPhraseIS{hrökkva eða stökkva} \dicDirectTranslationCS{buď, anebo};  \dicPhraseIS{hrökkva í kút} \dicDirectTranslationCS{ucuknout (leknutím)}  \textbf{2.} \dicDirectTranslationCS{rozbít se, rozpadnout se};  \dicPhraseIS{hrökkva í sundur} \dicDirectTranslationCS{rozpadnout se, rozlomit se} \dicExampleIS{Plöturnar hrukku í sundur.} \dicExampleCS{Desky se rozpadly.};  \dicIdiom{hrökkva}[fyrir\,/\addthin til]{ \dicPhraseIS{hrökkva fyrir\,/\addthin til}} \dicDirectTranslationCS{(po)stačit, vystačit} \dicExampleIS{Peningarnir hrökkva ekki til.} \dicExampleCS{Peníze nestačí.};  \dicIdiom{hrökkva}[upp]{ \dicPhraseIS{hrökkva upp}} {\textbf{a.}} \dicDirectTranslationCS{rozletět\,/\addthin otevřít se dokořán (okno ap.)} \dicExampleIS{Glugginn hrökk upp.} \dicExampleCS{Okno se rozletělo dokořán.};  {\textbf{b.}} \dicDirectTranslationCS{probudit se náhle};  \dicIdiom{hrökkva}[upp af]{ \dicPhraseIS{hrökkva upp af}} \dicLangCat{hovor. přen.} \dicDirectTranslationCS{natáhnout bačkory};  \dicIdiom{hrökkva}[út úr]{ \dicPhraseIS{e‑að hrekkur út úr e‑m}} \dicDirectTranslationCS{(co komu) vypadne\,/\addthin vyklouzne (neopatrné slovo ap.)};  \dicIdiom{hrökkva}[við]{ \dicPhraseIS{hrökkva við}} \dicDirectTranslationCS{cuknout\,/\addthin trhnout sebou (leknutím)}
\dicEntry[hröktum] \dicTerm{hröktum} \dicIPA{{\textsubring{r}}{\oe}{x}{\textsubring{d}}{\textscy}{\textsubring{m}}} \dicPos{v} \dicFlx{ind pf pl 1 pers} \dicLink{hrekja}
\dicEntry[hrönn] \dicTerm{hrönn} \dicIPA{{\textsubring{r}}{\oe}{n}{\textlengthmark}} \dicPos{f}[7] \dicFlx{(hrannar, hrannir)}[16] \dicLangCat{básn.} \dicSynonym{alda} \dicDirectTranslationCS{vlna};  \dicPhraseIS{í hrönnum} \dicFlx{adv} \dicDirectTranslationCS{ve velkém množství}
\dicEntry[hröp] \dicTerm{hröp} \dicIPA{{\textsubring{r}}{\oe}{\textlengthmark}{\textsubring{b}}} \dicPos{n} \dicFlx{pl nom} \dicLink{hrap}
\dicEntry[hrörlegur] \dicTerm{hrör··legur} \dicIPA{{\textsubring{r}}{\oe}{r}{l}{\textepsilon}{\textbabygamma}{\textscy}{\textsubring{r}}} \dicPos{adj}[1]\dicFlx{}[-8] \textbf{1.} \dicSynonym*{illa kominn af elli} \dicDirectTranslationCS{vetchý, sešlý (věkem)}  \textbf{2.} \dicSynonym*{úr sér genginn} \dicDirectTranslationCS{zchátralý, chatrný, (jsoucí) na rozpadnutí (chata ap.)}
\dicEntry[hrörna] \dicTerm{hrörn|a} \dicIPA{{\textsubring{r}}{\oe}{r}{\textsubring{d}}{n}{a}} \dicPos{v}[1] \dicFlx{(‑aði)}[34] \dicDirectTranslationCS{rozpadat se, chátrat, scházet} \dicExampleIS{Bærinn var að hrörna niður.} \dicExampleCS{Usedlost se rozpadala.}
\dicEntry[hrörnun] \dicTerm{hrörn|un} \dicIPA{{\textsubring{r}}{\oe}{r}{\textsubring{d}}{n}{\textscy}{\textsubring{n}}} \dicPos{f}[7] \dicFlx{(‑unar)}[12] \dicDirectTranslationCS{úpadek, (z)chátrání}
\dicEntry[hrösun] \dicTerm{hrös|un} \dicIPA{{\textsubring{r}}{\oe}{\textlengthmark}{s}{\textscy}{\textsubring{n}}} \dicPos{f}[7] \dicFlx{(‑unar)}[12] \textbf{1.} \dicSynonym*{það að hrasa} \dicDirectTranslationCS{klopýtnutí, škobrtnutí}  \textbf{2.} \dicLangCat{přen.} \dicSynonym{yfirsjón} \dicDirectTranslationCS{klopýtnutí, přehmat}
\dicEntry[hst.] \dicTerm{hst.} \dicPos{zkr} \dicPhraseIS{hástig} \dicFieldCat{jaz.} \dicDirectTranslationCS{superlativ, 3. stupeň stupňování}
\dicEntry[Hudsonfljót] \dicTerm{Hudson··fljót} \dicIPA{{h}{\textscy}{\textsubring{d}}{s}{\textopeno}{n}{f}{l}{j}{ou}{\textsubring{d}}} \dicPos{n}[2] \dicFlx{(‑s)}[4] \dicFieldCat{geog.} \dicDirectTranslationCS{Hudson} \dicIndirectTranslationCS{(řeka v~USA)}
\dicEntry[huga] \dicTerm{hug|a} \dicsymFrequent\  \dicIPA{{h}{\textscy}{\textlengthmark}{\textbabygamma}{a}} \dicPos{v}[1] \dicFlx{(‑aði)}[1] \dicPhraseIS{huga að e‑u} \dicDirectTranslationCS{dát\,/\addthin dávat pozor na (co), myslet na (co)} \dicExampleIS{Alþjóðastofnunin hvetur stjórnvöld til að huga að framtíð landsins.} \dicExampleCS{Mezinárodní organizace vyzývá vládu, aby myslela na budoucnost země.}
\dicEntry[hugaður] \dicTerm{hug··|aður} \dicIPA{{h}{\textscy}{\textlengthmark}{\textbabygamma}{a}{ð}{\textscy}{\textsubring{r}}} \dicPos{adj}[3] \dicFlx{(f ‑uð)}[3] \dicSynonym{óhræddur} \dicDirectTranslationCS{odvážný, nebojácný}
\dicEntry[hugarástand] \dicTerm{hugar··á·stand} \dicIPA{{h}{\textscy}{\textlengthmark}{\textbabygamma}{a}{r}{au}{s}{\textsubring{d}}{a}{n}{\textsubring{d}}} \dicPos{n}[2] \dicFlx{(‑s)}[2] \dicDirectTranslationCS{duševní stav, rozpoložení}
\dicEntry[hugarburður] \dicTerm{hugar··burð|ur} \dicIPA{{h}{\textscy}{\textlengthmark}{\textbabygamma}{a}{r}{\textsubring{b}}{\textscy}{r}{ð}{\textscy}{\textsubring{r}}} \dicPos{m}[10] \dicFlx{(‑ar)}[7] \dicSynonym{ímyndun} \dicDirectTranslationCS{výmysl, smyšlenka, výplod\,/\addthin výtvor fantazie}
\dicEntry[hugarfar] \dicTerm{hugar··far} \dicIPA{{h}{\textscy}{\textlengthmark}{\textbabygamma}{a}{\textsubring{r}}{f}{a}{\textsubring{r}}} \dicPos{n}[2] \dicFlx{(‑s)}[2] \dicSynonym*{hugarþel} \dicDirectTranslationCS{postoj} \dicExampleIS{neikvætt hugarfar} \dicExampleCS{negativní postoj}
\dicEntry[hugarflug] \dicTerm{hugar··flug} \dicIPA{{h}{\textscy}{\textlengthmark}{\textbabygamma}{a}{\textsubring{r}}{f}{l}{\textscy}{x}} \dicPos{n}[2] \dicFlx{(‑s)}[2] \textbf{1.} \dicSynonym{ímyndunarafl} \dicDirectTranslationCS{obrazotvornost, představivost, fantazie}  \textbf{2.} \dicDirectTranslationCS{brainstorming}
\dicEntry[hugarlund] \dicTerm{hugar··lund} \dicsymFrequent\  \dicIPA{{h}{\textscy}{\textlengthmark}{\textbabygamma}{a}{r}{l}{\textscy}{n}{\textsubring{d}}} \dicPos{f}[7] \dicFlx{(‑ar)}[3] \dicPhraseIS{gera sér e‑ð í hugarlund} \dicSynonym*{ímynda sér e‑ð} \dicDirectTranslationCS{představit\,/\addthin představovat si (co)}
\dicEntry[hugarórar] \dicTerm{hugar··órar} \dicIPA{{h}{\textscy}{\textlengthmark}{\textbabygamma}{a}{r}{ou}{r}{a}{\textsubring{r}}} \dicPos{m}[1] \dicFlx{pl}[2] \dicSynonym{heilaspuni} \dicDirectTranslationCS{výtvor fantazie, smyšlenka}
\dicEntry[hugblær] \dicTerm{hug··blæ|r} \dicIPA{{h}{\textscy}{\textbabygamma}{\textsubring{b}}{l}{a}{i}{\textsubring{r}}} \dicPos{m}[9] \dicFlx{(‑s\,/\addthin ‑jar)}[30] \dicDirectTranslationCS{nálada, atmosféra} \dicExampleIS{hugblær augnabliksins} \dicExampleCS{atmosféra okamžiku}
\dicEntry[hugboð] \dicTerm{hug··boð} \dicIPA{{h}{\textscy}{\textbabygamma}{\textsubring{b}}{\textopeno}{\texttheta}} \dicPos{n}[2] \dicFlx{(‑s, ‑)}[5] \dicDirectTranslationCS{předtucha, tušení, pocit}
\dicEntry[hugbúnaður] \dicTerm{hug··bú·nað|ur} \dicIPA{{h}{\textscy}{\textbabygamma}{\textsubring{b}}{u}{n}{a}{ð}{\textscy}{\textsubring{r}}} \dicPos{m}[10] \dicFlx{(‑ar)}[9] \dicFieldCat{poč.} \dicDirectTranslationCS{software, programové vybavení}
\dicEntry[hugdetta] \dicTerm{hug··dett|a} \dicIPA{{h}{\textscy}{\textbabygamma}{\textsubring{d}}{\textepsilon}{h}{\textsubring{d}}{a}} \dicPos{f}[1] \dicFlx{(‑u, ‑ur)}[7] \dicSynonym{þanki} \dicDirectTranslationCS{nápad, (náhlá) myšlenka}
\dicEntry[hugdirfska] \dicTerm{hug··dirfsk|a} \dicIPA{{h}{\textscy}{\textbabygamma}{\textsubring{d}}{\textsci}{\textsubring{r}}{s}{\r{g}}{a}} \dicPos{f}[1] \dicFlx{(‑u)}[5] \dicSynonym{hugrekki\smash{\textsuperscript{2}}} \dicDirectTranslationCS{kuráž, odvaha}
\dicEntry[hugdjarfur] \dicTerm{hug··|djarfur} \dicIPA{{h}{\textscy}{\textbabygamma}{\textsubring{d}}{j}{a}{r}{v}{\textscy}{\textsubring{r}}} \dicPos{adj}[1] \dicFlx{(f ‑djörf)}[2] \dicDirectTranslationCS{odvážný, kurážný, smělý}
\dicEntry[hugðarefni] \dicTerm{hugðar··efni} \dicIPA{{h}{\textscy}{\textbabygamma}{ð}{a}{r}{\textepsilon}{\textsubring{b}}{n}{\textsci}} \dicPos{n}[2] \dicFlx{(‑s, ‑)}[14] \dicSynonym{áhugamál} \dicDirectTranslationCS{záliba, koníček}
\dicEntry[hugði] \dicTerm{hugði} \dicIPA{{h}{\textscy}{\textbabygamma}{ð}{\textsci}} \dicPos{v} \dicFlx{ind pf sg 1 pers} \dicLink{hyggja\smash{\textsuperscript{2}}}
\dicEntry[hugðum] \dicTerm{hugðum} \dicIPA{{h}{\textscy}{\textbabygamma}{ð}{\textscy}{\textsubring{m}}} \dicPos{v} \dicFlx{ind pf pl 1 pers} \dicLink{hyggja\smash{\textsuperscript{2}}}
\dicEntry[hugfanginn] \dicTerm{hug··fanginn} \dicIPA{{h}{\textscy}{x}{f}{au}{\textltailn}{\r{\textObardotlessj}}{\textsci}{\textsubring{n}}} \dicPos{adj}[6]\dicFlx{}[-3] \dicSynonym{hrifinn} \dicDirectTranslationCS{fascinovaný, okouzlený};  \dicPhraseIS{vera hugfanginn af e‑u} \dicDirectTranslationCS{být okouzlený (čím)}
\dicEntry[hugfastur] \dicTerm{hug··|fastur} \dicIPA{{h}{\textscy}{x}{f}{a}{s}{\textsubring{d}}{\textscy}{\textsubring{r}}} \dicPos{adj}[1] \dicFlx{(f ‑föst)}[11] \dicPhraseIS{hafa e‑ð hugfast} \dicDirectTranslationCS{(dobře) si zapamatovat (co), zapsat si (co) za uši}
\dicEntry[hugga] \dicTerm{hugg|a} \dicIPA{{h}{\textscy}{\r{g}}{\textlengthmark}{a}} \dicPos{v}[1] \dicFlx{(‑aði)}[1] \dicFlx{acc} \dicDirectTranslationCS{utěšit, utěšovat, (u)konejšit} \dicExampleIS{hugga barn} \dicExampleCS{ukonejšit dítě};  \dicPhraseIS{hugga sig við e‑ð} \dicDirectTranslationCS{utěšovat se (čím)}
\dicEntry[huggulegur] \dicTerm{huggu··legur} \dicIPA{{h}{\textscy}{\r{g}}{\textlengthmark}{\textscy}{l}{\textepsilon}{\textbabygamma}{\textscy}{\textsubring{r}}} \dicPos{adj}[1]\dicFlx{}[-8] \textbf{1.} \dicSynonym{notalegur} \dicDirectTranslationCS{útulný, příjemný}  \textbf{2.} \dicDirectTranslationCS{pohledný, atraktivní (dívka ap.)}
\dicEntry[huggun] \dicTerm{hugg|un} \dicIPA{{h}{\textscy}{\r{g}}{\textlengthmark}{\textscy}{\textsubring{n}}} \dicPos{f}[7] \dicFlx{(‑unar, ‑anir)}[8] \dicDirectTranslationCS{útěcha, utěšování} \dicExampleIS{finna huggun í trúnni} \dicExampleCS{nalézt útěchu ve víře};  \dicPhraseIS{e‑m er huggun í e‑u} \dicFlx{impers} \dicDirectTranslationCS{(kdo) nachází útěchu v~(čem)}
\dicEntry[huggunarorð] \dicTerm{huggunar··orð} \dicIPA{{h}{\textscy}{\r{g}}{\textlengthmark}{\textscy}{n}{a}{r}{\textopeno}{r}{\texttheta}} \dicPos{n}[2] \dicFlx{(‑s, ‑)}[5] \dicDirectTranslationCS{slovo útěchy}
\dicEntry[hugheill] \dicTerm{hug··heill} \dicIPA{{h}{\textscy}{\textlengthmark}{x}{h}{ei}{\textsubring{d}}{\textsubring{l}}} \dicPos{adj}[8]\dicFlx{}[-1] \dicSynonym{innilegur} \dicDirectTranslationCS{srdečný, vřelý, vroucí} \dicExampleIS{hugheilar kveðjur} \dicExampleCS{srdečné pozdravy}
\dicEntry[hughraustur] \dicTerm{hug··hraustur} \dicIPA{{h}{\textscy}{x}{\textsubring{r}}{\oe i}{s}{\textsubring{d}}{\textscy}{\textsubring{r}}} \dicPos{adj}[1]\dicFlx{}[-10] \dicSynonym{djarfur} \dicDirectTranslationCS{smělý, chrabrý}
\dicEntry[hughreysta] \dicTerm{hug··hreyst|a} \dicIPA{{h}{\textscy}{x}{\textsubring{r}}{ei}{s}{\textsubring{d}}{a}} \dicPos{v}[2] \dicFlx{(‑i, ‑)}[12] \dicFlx{acc} \dicDirectTranslationCS{povzbudit, povzbuzovat} \dicExampleIS{hughreysta e‑n með e‑u} \dicExampleCS{povzbudit (koho čím)}
\begin{xtolerant}{}{1pt}
\dicEntry[hughreysting] \dicTerm{hug··hreyst·ing}\addthinS\dicIPA{{h}{\textscy}{x}{\textsubring{r}}{ei}{s}{\textsubring{d}}{i}{\ng}{\r{g}}}\addthinS\dicPos{f}[4]\addthinS\dicFlx{(‑ar)}[7]\addthinS\dicDirectTranslationCS{povzbuzení, vzpru\-ha}
\end{xtolerant}
\dicEntry[hughrif] \dicTerm{hug··hrif} \dicIPA{{h}{\textscy}{x}{\textsubring{r}}{\textsci}{f}} \dicPos{n}[2] \dicFlx{pl}[1] \dicDirectTranslationCS{dojem, imprese} \dicExampleIS{skapa hughrif} \dicExampleCS{udělat dojem}
\dicEntry[hughvarf] \dicTerm{hug··hvarf} \dicIPA{{h}{\textscy}{\textlengthmark}{k\smash{\textsuperscript{h}}}{v}{a}{r}{f}} \dicPos{n}[2] \dicFlx{(‑s)}[2] \dicPhraseIS{telja e‑m hughvarf} \dicDirectTranslationCS{přimět (koho) změnit názor, přemluvit (koho)}
\dicEntry[hughyggja] \dicTerm{hug··hyggj|a} \dicIPA{{h}{\textscy}{\textlengthmark}{x}{h}{\textsci}{\r{\textObardotlessj}}{a}} \dicPos{f}[1] \dicFlx{(‑u)}[5] \dicFieldCat{filos.} \dicDirectTranslationCS{idealismus}
\dicEntry[hugi] \dicTerm{hug|i} \dicsymFrequent\  \dicIPA{{h}{\textscy}{i}{j}{\textlengthmark}{\textsci}} \dicPos{m}[1] \dicFlx{(‑a, ‑ar)}[1] \dicDirectTranslationCS{mysl};  \dicPhraseIS{af heilum huga} \dicFlx{adv} \dicDirectTranslationCS{z~celého srdce};  \dicPhraseIS{dreifa huganum} \dicDirectTranslationCS{rozptýlit mysl};  \dicPhraseIS{hafa e‑ð í huga} \dicDirectTranslationCS{mít (co) na mysli};  \dicPhraseIS{láta hugann reika} \dicDirectTranslationCS{nechat (volně) plynout myšlenky};  \dicPhraseIS{leiða hugann að e‑u} \dicDirectTranslationCS{přemýšlet o~(čem)};  \dicPhraseIS{með hálfum huga} \dicFlx{adv} \dicDirectTranslationCS{vlažně, laxně, polovičatě};  \dicPhraseIS{með opnum huga} \dicFlx{adv} \dicDirectTranslationCS{s~otevřenou myslí};  \dicPhraseIS{vera með hugann við e‑ð} \dicDirectTranslationCS{soustředit\,/\addthin koncentrovat se na (co)}
\dicEntry[hugkvæmni] \dicTerm{hug··kvæmn|i} \dicIPA{{h}{\textscy}{x}{k\smash{\textsuperscript{h}}}{v}{a}{i}{m}{n}{\textsci}} \dicPos{f}[3] \dicFlx{(‑i)}[3] \dicDirectTranslationCS{vynalézavost, tvořivost, nápaditost}
\dicEntry[hugkvæmur] \dicTerm{hug··kvæmur} \dicIPA{{h}{\textscy}{x}{k\smash{\textsuperscript{h}}}{v}{a}{i}{m}{\textscy}{\textsubring{r}}} \dicPos{adj}[1]\dicFlx{}[-1] \dicDirectTranslationCS{tvůrčí, tvořivý, vynalézavý, nápaditý}
\dicEntry[huglaus] \dicTerm{hug··laus} \dicIPA{{h}{\textscy}{\textbabygamma}{l}{\oe i}{s}} \dicPos{adj}[5]\dicFlx{}[-1] \dicDirectTranslationCS{zbabělý, bojácný}
\dicEntry[hugleiða] \dicTerm{hug··lei|ða} \dicsymFrequent\  \dicIPA{{h}{\textscy}{\textbabygamma}{l}{ei}{ð}{a}} \dicPos{v}[2] \dicFlx{(‑ddi, ‑tt)}[167] \dicFlx{acc} \textbf{1.} \dicSynonym*{velta fyrir sér} \dicDirectTranslationCS{přemýšlet, uvažovat, zvažovat, hloubat} \dicExampleIS{Landið hugleiðir að biðja um aðstoð.} \dicExampleCS{Země uvažuje, že požádá o~pomoc.}  \textbf{2.} \dicDirectTranslationCS{meditovat}
\dicEntry[hugleiðing] \dicTerm{hug··leið·ing} \dicIPA{{h}{\textscy}{\textbabygamma}{l}{ei}{ð}{i}{\ng}{\r{g}}} \dicPos{f}[4] \dicFlx{(‑ar, ‑ar)}[5] \dicDirectTranslationCS{uvažování, přemýšlení, hloubání}
\dicEntry[hugleiðsla] \dicTerm{hug··leiðsl|a} \dicIPA{{h}{\textscy}{\textbabygamma}{l}{ei}{ð}{s}{\textsubring{d}}{l}{a}} \dicPos{f}[1] \dicFlx{(‑u)}[5] \dicDirectTranslationCS{meditace}
\dicEntry[hugleysi] \dicTerm{hug··leysi} \dicIPA{{h}{\textscy}{\textbabygamma}{l}{ei}{s}{\textsci}} \dicPos{n}[2] \dicFlx{(‑s)}[20] \dicDirectTranslationCS{zbabělost, bojácnost}
\dicEntry[hugljómun] \dicTerm{hug··ljóm|un} \dicIPA{{h}{\textscy}{\textbabygamma}{l}{j}{ou}{m}{\textscy}{\textsubring{n}}} \dicPos{f}[7] \dicFlx{(‑unar)}[9] \dicDirectTranslationCS{(náhlý) vhled, inspirace, vnuknutí}
\dicEntry[hugljúfi] \dicTerm{hug··ljúf|i} \dicIPA{{h}{\textscy}{\textbabygamma}{l}{j}{u}{v}{\textsci}} \dicPos{m}[1] \dicFlx{(‑a)}[5] \dicDirectTranslationCS{miláček, oblíbenec, oblíbenkyně, oblíbený člověk};  \dicPhraseIS{vera hvers manns hugljúfi} \dicDirectTranslationCS{být miláčkem všech}
\dicEntry[hugljúfur] \dicTerm{hug··ljúfur} \dicIPA{{h}{\textscy}{\textbabygamma}{l}{j}{u}{v}{\textscy}{\textsubring{r}}} \dicPos{adj}[1]\dicFlx{}[-1] \dicSynonym{geðþekkur} \dicDirectTranslationCS{milý, líbezný}
\dicEntry[huglægni] \dicTerm{hug··lægn|i} \dicIPA{{h}{\textscy}{\textbabygamma}{l}{a}{i}{\r{g}}{n}{\textsci}} \dicPos{f}[3] \dicFlx{(‑i)}[3] \dicDirectTranslationCS{subjektivita, subjektivnost} \dicAntonym{hlutlægni}
\dicEntry[huglægur] \dicTerm{hug··lægur} \dicIPA{{h}{\textscy}{\textbabygamma}{l}{a}{i}{\textbabygamma}{\textscy}{\textsubring{r}}} \dicPos{adj}[1]\dicFlx{}[-1] \dicDirectTranslationCS{subjektivní} \dicExampleIS{huglæg frásögn} \dicExampleCS{subjektivní vyprávění} \dicAntonym{hlutlægur}
\dicEntry[hugmynd] \dicTerm{hug··mynd} \dicsymFrequent\  \dicIPA{{h}{\textscy}{\textbabygamma}{m}{\textsci}{n}{\textsubring{d}}} \dicPos{f}[7] \dicFlx{(‑ar, ‑ir)}[1] \dicDirectTranslationCS{myšlenka, nápad, idea} \dicExampleIS{Þú hefur góða hugmynd.} \dicExampleCS{Máš dobrý nápad.};  \dicPhraseIS{fá hugmynd} \dicDirectTranslationCS{dostat nápad};  \dicPhraseIS{hafa ekki hugmynd um e‑ð} \dicDirectTranslationCS{nemít o~(čem) ponětí, netušit (co)}
\dicEntry[hugmyndaflug] \dicTerm{hug·mynda··flug} \dicIPA{{h}{\textscy}{\textbabygamma}{m}{\textsci}{n}{\textsubring{d}}{a}{f}{l}{\textscy}{x}} \dicPos{n}[2] \dicFlx{(‑s)}[2] \dicDirectTranslationCS{obrazotvornost, představivost, fantazie}
\dicEntry[hugmyndafræði] \dicTerm{hug·mynda··fræð|i} \dicIPA{{h}{\textscy}{\textbabygamma}{m}{\textsci}{n}{\textsubring{d}}{a}{f}{r}{a}{i}{ð}{\textsci}} \dicPos{f}[3] \dicFlx{(‑i)}[3] \dicDirectTranslationCS{ideologie}
\dicEntry[hugnæmur] \dicTerm{hug··næmur} \dicIPA{{h}{\textscy}{\textbabygamma}{n}{a}{i}{m}{\textscy}{\textsubring{r}}} \dicPos{adj}[1]\dicFlx{}[-1] \dicSynonym{áhrifamikill} \dicDirectTranslationCS{dojemný, jímavý, tklivý}
\dicEntry[hugprúður] \dicTerm{hug··prúður} \dicIPA{{h}{\textscy}{x}{p\smash{\textsuperscript{h}}}{r}{u}{ð}{\textscy}{\textsubring{r}}} \dicPos{adj}[2]\dicFlx{}[-6] \dicLangCat{kniž.} \dicSynonym{kjarkmikill} \dicDirectTranslationCS{udatný, chrabrý}
\dicEntry[hugrakkur] \dicTerm{hug··|rakkur} \dicIPA{{h}{\textscy}{\textbabygamma}{r}{a}{h}{\r{g}}{\textscy}{\textsubring{r}}} \dicPos{adj}[1] \dicFlx{(f ‑rökk)}[2] \dicSynonym{djarfur} \dicDirectTranslationCS{statečný, udatný, neohrožený}
\dicEntry[hugrekki] \dicTerm{hug··rekk|i\smash{\textsuperscript{1}}} \dicIPA{{h}{\textscy}{\textbabygamma}{r}{\textepsilon}{h}{\r{\textObardotlessj}}{\textsci}} \dicPos{f}[3] \dicFlx{(‑i)}[3] \dicLink{hugrekki\smash{\textsuperscript{2}}}
\dicEntry[hugrekki] \dicTerm{hug··rekki\smash{\textsuperscript{2}}}\dicTerm{, hugrekki\smash{\textsuperscript{1}}} \dicIPA{{h}\-{\textscy}\-{\textbabygamma}\-{r}\-{\textepsilon}\-{h}\-{\r{\textObardotlessj}}\-{\textsci}\-} \dicPos{n}[2] \dicFlx{(‑s)}[20] \dicSynonym{dirfska} \dicDirectTranslationCS{odvaha, statečnost}
\dicEntry[hugrenning] \dicTerm{hug··renn·ing} \dicIPA{{h}{\textscy}{\textbabygamma}{r}{\textepsilon}{n}{i}{\ng}{\r{g}}} \dicPos{f}[4] \dicFlx{(‑ar, ‑ar)}[5] \dicPhraseIS{hugrenningar} \dicFlx{pl} \dicSynonym{hugsun} \dicDirectTranslationCS{zamyšlení}
\dicEntry[hugrenningatengsl] \dicTerm{hug·renninga··tengsl} \dicIPA{{h}{\textscy}{\textbabygamma}{r}{\textepsilon}{n}{i}{\ng}{\r{g}}{a}{t\smash{\textsuperscript{h}}}{\textepsilon}{\ng}{s}{\textsubring{d}}{\textsubring{l}}} \dicPos{n}[2] \dicFlx{pl}[1] \dicDirectTranslationCS{(myšlenková) asociace, (myšlenkové) spojení}
\dicEntry[hugsa] \dicTerm{hugs|a} \dicsymFrequent\  \dicIPA{{h}{\textscy}{x}{s}{a}} \dicPos{v}[1] \dicFlx{(‑aði)}[15] \dicFlx{acc} \textbf{1.} \dicDirectTranslationCS{myslet, myslit, přemýšlet, zamyslet se};  \dicPhraseIS{hugsa sig um} \dicDirectTranslationCS{rozmyslet se};  \dicPhraseIS{hugsa um e‑ð} \dicDirectTranslationCS{přemýšlet o~(čem), myslet na (co)} \dicExampleIS{Hann hugsaði um það sem hún sagði.} \dicExampleCS{Přemýšlel o~tom, co řekla.};  \dicPhraseIS{hugsa um e‑n} \dicDirectTranslationCS{myslet na (koho)}  \textbf{2.} \dicDirectTranslationCS{domnívat se, myslet si, myslit si} \dicExampleIS{Sumir hugsa að það þýði ekkert.} \dicExampleCS{Někteří si myslí, že to nic neznamená.}  \textbf{3.} \dicPhraseIS{geta hugsað sér e‑ð} \dicDirectTranslationCS{umět si představit (co)}  \textbf{4.} \dicPhraseIS{að hugsa sér} \dicDirectTranslationCS{představit si};  \dicPhraseIS{hugsaðu þér} \dicDirectTranslationCS{zamysli se, podívej se (jak je tu hezky ap.)};  \dicIdiom{hugsa}[fyrir]{ \dicPhraseIS{hugsa fyrir e‑u}} \dicDirectTranslationCS{promyslet (co), rozmyslet si (co)};  \dicIdiom{hugsa}[til]{ \dicPhraseIS{hugsa til e‑rs}} \dicDirectTranslationCS{myslet na (koho)}; { \dicPhraseIS{geta ekki til þess hugsað}} \dicDirectTranslationCS{nemoct na to ani pomyslet};  \dicIdiom{hugsa}[um]{ \dicPhraseIS{hugsa um e‑ð}} \dicDirectTranslationCS{starat se o~(co), pečovat o~(co)} \dicExampleIS{hugsa um garðinn} \dicExampleCS{pečovat o~zahrádku};  \dicIdiom{hugsa}[upp]{ \dicPhraseIS{hugsa e‑ð upp}} \dicDirectTranslationCS{vymyslet\,/\addthin vymýšlet (co)};  \dicIdiom{hugsa}[út í]{ \dicPhraseIS{hugsa út í e‑ð}} \dicDirectTranslationCS{zamyslet se nad (čím)};  \dicIdiom{hugsast}{ \dicPhraseIS{það getur hugsast}} \dicFlx{refl} \dicDirectTranslationCS{je to možné, může být}
\dicEntry[hugsanaflutningur] \dicTerm{hugsana··flut·ning|ur} \dicIPA{{h}\-{\textscy}\-{x}\-{s}\-{a}\-{n}\-{a}\-{f}\-{l}\-{\textscy}\-{h}\-{\textsubring{d}}\-{n}\-{i}\-{\ng}\-{\r{g}}\-{\textscy}\-{\textsubring{r}}\-} \dicPos{m}[6] \dicFlx{(‑s)}[9] \dicFieldCat{psych.} \dicDirectTranslationCS{telepatie}
\dicEntry[hugsanagangur] \dicTerm{hugsana··gang|ur} \dicIPA{{h}{\textscy}{x}{s}{a}{n}{a}{\r{g}}{au}{\ng}{\r{g}}{\textscy}{\textsubring{r}}} \dicPos{m}[6] \dicFlx{(‑s)}[7] \dicDirectTranslationCS{myšlenkový pochod, uvažování}
\dicEntry[hugsandi] \dicTerm{hugs··andi} \dicsymFrequent\  \dicIPA{{h}{\textscy}{x}{s}{a}{n}{\textsubring{d}}{\textsci}} \dicPos{adj}[13] \dicFlx{indecl}[1] \textbf{1.} \dicSynonym{íhugull} \dicDirectTranslationCS{přemýšlející, zamyšlený, hloubající} \dicExampleIS{Kennarinn var mjög hugsandi á svip.} \dicExampleCS{Učitel se tvářil velmi zamyšleně.}  \textbf{2.} \dicSynonym{hugsanlegur} \dicDirectTranslationCS{možný, představitelný};  \dicPhraseIS{það er hugsandi að (gera e‑ð)} \dicDirectTranslationCS{přichází v~úvahu (udělat (co))}
\dicEntry[hugsanlegur] \dicTerm{hugsan··legur} \dicsymFrequent\  \dicIPA{{h}{\textscy}{x}{s}{a}{n}{l}{\textepsilon}{\textbabygamma}{\textscy}{\textsubring{r}}} \dicPos{adj}[1]\dicFlx{}[-8] \dicSynonym{mögulegur} \dicDirectTranslationCS{myslitelný, možný, potencionální} \dicExampleIS{Við leituðum á öllum hugsanlegum stöðum.} \dicExampleCS{Hledali jsme na všech možných místech.}
\dicEntry[hugsi] \dicTerm{hugsi} \dicsymFrequent\  \dicIPA{{h}{\textscy}{x}{s}{\textsci}} \dicPos{adj}[13] \dicFlx{indecl}[2] \dicSynonym*{í þungum þönkum} \dicDirectTranslationCS{zamyšlený, zadumaný} \dicExampleIS{Presturinn var svo hugsi að hann tók ekki eftir þeim.} \dicExampleCS{Kněz byl tak zamyšlený, že si jich nevšiml.};  \dicPhraseIS{vera hugsi út af e‑u} \dicDirectTranslationCS{spekulovat nad (čím), hloubat o~(čem)}
\dicEntry[hugsjón] \dicTerm{hug··sjón} \dicIPA{{h}{\textscy}{x}{s}{j}{ou}{\textsubring{n}}} \dicPos{f}[7] \dicFlx{(‑ar, ‑ir)}[1] \dicDirectTranslationCS{ideál} \dicExampleIS{berjast fyrir hugsjónum sínum} \dicExampleCS{bojovat za své ideály}
\begin{xtolerant}{}{1pt}
\dicEntry[hugsjónamaður] \dicTerm{hug·sjóna··|maður}\addthinS\dicIPA{{h}{\textscy}{x}{s}{j}{ou}{n}{a}{m}{a}{ð}{\textscy}{\textsubring{r}}}\addthinS\dicPos{m}[13]\addthinS\dicFlx{(‑manns, ‑menn)}[2] \dicDirectTranslationCS{idealista, idealistka}
\end{xtolerant}
\dicEntry[hugsjúkur] \dicTerm{hug··sjúkur} \dicIPA{{h}{\textscy}{x}{s}{j}{u}{\r{g}}{\textscy}{\textsubring{r}}} \dicPos{adj}[1]\dicFlx{}[-1] \textbf{1.} \dicSynonym{þunglyndur} \dicDirectTranslationCS{melancholický, deprimovaný, sklíčený}  \textbf{2.} \dicFieldCat{psych.} \dicDirectTranslationCS{neurotický}
\dicEntry[hugskot] \dicTerm{hug··skot} \dicIPA{{h}{\textscy}{x}{s}{\r{g}}{\textopeno}{\textsubring{d}}} \dicPos{n}[2] \dicFlx{(‑s, ‑)}[5] \dicSynonym{hugur} \dicDirectTranslationCS{mysl, vědomí} \dicExampleIS{í innsta hugskoti manns} \dicExampleCS{v~nejzazším vědomí člověka}
\dicEntry[hugskotssjónir] \dicTerm{hug·skots··sjónir} \dicIPA{{h}{\textscy}{x}{s}{\r{g}}{\textopeno}{\textsubring{d}}{s}{j}{ou}{n}{\textsci}{\textsubring{r}}} \dicPos{f}[7] \dicFlx{pl}[2] \dicPhraseIS{e‑að stendur e‑m fyrir hugskotssjónum} \dicDirectTranslationCS{(kdo) vidí (co), jako by to bylo před (kým); (kdo) má (co) živě před očima}
\dicEntry[hugsuður] \dicTerm{hugs··uð|ur} \dicIPA{{h}{\textscy}{x}{s}{\textscy}{ð}{\textscy}{\textsubring{r}}} \dicPos{m}[10] \dicFlx{(‑ar, ‑ir)}[4] \dicDirectTranslationCS{myslitel(ka)}
\dicEntry[hugsun] \dicTerm{hugs|un} \dicsymFrequent\  \dicIPA{{h}{\textscy}{x}{s}{\textscy}{\textsubring{n}}} \dicPos{f}[7] \dicFlx{(‑unar, ‑anir)}[8] \dicDirectTranslationCS{myšlení, přemýšlení, uvažování} \dicExampleIS{jákvæð hugsun} \dicExampleCS{pozitivní myšlení};  \dicPhraseIS{hafa hugsun á að (gera e‑ð)} \dicDirectTranslationCS{myslet na (udělání (čeho))}
\dicEntry[hugsunarháttur] \dicTerm{hugsunar··|háttur} \dicIPA{{h}{\textscy}{x}{s}{\textscy}{n}{a}{\textsubring{r}}{h}{au}{h}{\textsubring{d}}{\textscy}{\textsubring{r}}} \dicPos{m}[12] \dicFlx{(‑háttar)}[8] \dicDirectTranslationCS{způsob myšlení, mentalita}
\dicEntry[hugsunarlaus] \dicTerm{hugsunar··laus} \dicIPA{{h}{\textscy}{x}{s}{\textscy}{n}{a}{r}{l}{\oe i}{s}} \dicPos{adj}[5]\dicFlx{}[-1] \dicDirectTranslationCS{(duchem) nepřítomný, roztržitý}
\dicEntry[hugsunarleysi] \dicTerm{hugsunar··leysi} \dicIPA{{h}{\textscy}{x}{s}{\textscy}{n}{a}{r}{l}{ei}{s}{\textsci}} \dicPos{n}[2] \dicFlx{(‑s)}[20] \dicDirectTranslationCS{nepřítomnost (duchem), roztržitost}
\dicEntry[hugtak] \dicTerm{hug··|tak} \dicsymFrequent\  \dicIPA{{h}{\textscy}{x}{t\smash{\textsuperscript{h}}}{a}{\r{g}}} \dicPos{n}[2] \dicFlx{(‑taks, ‑tök)}[8] \dicDirectTranslationCS{pojem, pojetí, koncept} \dicExampleIS{skilgreina hugtakið} \dicExampleCS{definovat pojem}
\dicEntry[hugulsamur] \dicTerm{hugul··|samur} \dicIPA{{h}{\textscy}{\textlengthmark}{\textbabygamma}{\textscy}{l}{s}{a}{m}{\textscy}{\textsubring{r}}} \dicPos{adj}[1] \dicFlx{(f ‑söm)}[2] \dicSynonym{nærgætinn} \dicDirectTranslationCS{rozvážný, pozorný}
\dicEntry[hugulsemi] \dicTerm{hugul··sem|i} \dicIPA{{h}{\textscy}{\textlengthmark}{\textbabygamma}{\textscy}{l}{s}{\textepsilon}{m}{\textsci}} \dicPos{f}[2] \dicFlx{(‑i)}[2] \dicDirectTranslationCS{rozvážnost, pozornost}
\dicEntry[hugur] \dicTerm{hug|ur} \dicsymFrequent\  \dicIPA{{h}{\textscy}{\textlengthmark}{\textbabygamma}{\textscy}{\textsubring{r}}} \dicPos{m}[10] \dicFlx{(‑ar, ‑ir)}[10] \textbf{1.} \dicSynonym{hugsun} \dicDirectTranslationCS{mysl};  \dicPhraseIS{e‑m kemur e‑að í hug} \dicFlx{impers} \dicDirectTranslationCS{(co koho) napadá} \dicExampleIS{Það fyrsta sem mér kom í hug.} \dicExampleCS{To první, co mě napadlo.}  \textbf{2.} \dicSynonym{áhugi} \dicDirectTranslationCS{zájem, zaujetí};  \dicPhraseIS{hafa hug á e‑ð} \dicDirectTranslationCS{mít zájem o~(co)};  \dicIdiom{hugur}{ \dicPhraseIS{e‑ir fella hugi saman}} \dicDirectTranslationCS{(kdo) si padne do oka}; { \dicPhraseIS{e‑m er ekki hlátur í hug}} \dicFlx{impers} \dicDirectTranslationCS{(komu) není do smíchu};  \dicPhraseIS{vera alls hugar feginn} \dicDirectTranslationCS{být štěstím bez sebe};  \dicPhraseIS{vera annars hugar} \dicDirectTranslationCS{být duchem nepřítomný, být myšlenkami někde jinde};  \dicPhraseIS{það er hugur í e‑m} \dicLangCat{přen.} \dicDirectTranslationCS{(kdo) má filipa}
\dicEntry[hugvitsamlegur] \dicTerm{hug·vit·sam··legur} \dicIPA{{h}{\textscy}{\textbabygamma}{v}{\textsci}{\textsubring{d}}{s}{a}{m}{l}{\textepsilon}{\textbabygamma}{\textscy}{\textsubring{r}}} \dicPos{adj}[1]\dicFlx{}[-8] \dicDirectTranslationCS{vynalézavý, tvůrčí, kreativní}
\dicEntry[hugvitsmaður] \dicTerm{hug·vits··|maður} \dicIPA{{h}{\textscy}{\textbabygamma}{v}{\textsci}{\textsubring{d}}{s}{m}{a}{ð}{\textscy}{\textsubring{r}}} \dicPos{m}[13] \dicFlx{(‑manns, ‑menn)}[2] \dicDirectTranslationCS{vynálezce, vynálezkyně}
\dicEntry[hugvísindadeild] \dicTerm{hug·vísinda··deild} \dicIPA{{h}{\textscy}{\textbabygamma}{v}{i}{s}{\textsci}{n}{\textsubring{d}}{a}{\textsubring{d}}{ei}{l}{\textsubring{d}}} \dicPos{f}[7] \dicFlx{(‑ar, ‑ir)}[1] \dicFieldCat{škol.} \dicDirectTranslationCS{fakulta humanitních studií}
\dicEntry[hugvísindi] \dicTerm{hug··vís·indi} \dicIPA{{h}{\textscy}{\textbabygamma}{v}{i}{s}{\textsci}{n}{\textsubring{d}}{\textsci}} \dicPos{n}[2] \dicFlx{pl}[19] \dicDirectTranslationCS{humanitní vědy\,/\addthin studia}
\dicEntry[hugþekkur] \dicTerm{hug··þekkur} \dicIPA{{h}{\textscy}{\textbabygamma}{\texttheta}{\textepsilon}{h}{\r{g}}{\textscy}{\textsubring{r}}} \dicPos{adj}[1]\dicFlx{}[-1] \dicSynonym{geðfelldur} \dicDirectTranslationCS{sympatický, příjemný}
\dicEntry[hula] \dicTerm{hul|a} \dicIPA{{h}{\textscy}{\textlengthmark}{l}{a}} \dicPos{f}[1] \dicFlx{(‑u)}[5] \textbf{1.} \dicSynonym{slæða\smash{\textsuperscript{1}}} \dicDirectTranslationCS{závoj, zahalení};  \dicPhraseIS{svipta hulunni af e‑u} \dicDirectTranslationCS{odhalit (co)} \dicExampleIS{svipta hulunni af leyndarmáli} \dicExampleCS{odhalit tajemství}  \textbf{2.} \dicSynonym{móða} \dicDirectTranslationCS{opar, mlha}
\dicEntry[hulda] \dicTerm{huld|a} \dicIPA{{h}{\textscy}{l}{\textsubring{d}}{a}} \dicPos{f}[1] \dicFlx{(‑u, ‑ur)}[13] \dicSynonym{leynd} \dicDirectTranslationCS{tajemství, tajemnost};  \dicPhraseIS{e‑að er á huldu} \dicDirectTranslationCS{(co) je zahaleno tajemstvím, (co) je nevyjasněné}
\dicEntry[huldi] \dicTerm{huldi} \dicIPA{{h}{\textscy}{l}{\textsubring{d}}{\textsci}} \dicPos{v} \dicFlx{ind pf sg 1 pers} \dicLink{hylja}
\dicEntry[huldufólk] \dicTerm{huldu··fólk} \dicIPA{{h}{\textscy}{l}{\textsubring{d}}{\textscy}{f}{ou}{\textsubring{l}}{\r{g}}} \dicPos{n}[2] \dicFlx{(‑s)}[2] \dicFieldCat{pov.} \dicSynonym*{álfar} \dicDirectTranslationCS{elfové}
\dicEntry[huldukona] \dicTerm{huldu··kon|a} \dicIPA{{h}{\textscy}{l}{\textsubring{d}}{\textscy}{k\smash{\textsuperscript{h}}}{\textopeno}{n}{a}} \dicPos{f}[1] \dicFlx{(‑u, ‑ur)}[27] \dicFieldCat{pov.} \dicSynonym{álfkona} \dicDirectTranslationCS{elfka, divoženka, víla} \dicIndirectTranslationCS{(nadpřirozená bytost ženského pohlaví podobná lidem zjevem i~způsoby)}
\dicEntry[huldum] \dicTerm{huldum} \dicIPA{{h}{\textscy}{l}{\textsubring{d}}{\textscy}{\textsubring{m}}} \dicPos{v} \dicFlx{ind pf pl 1 pers} \dicLink{hylja}
\dicEntry[huldumaður] \dicTerm{huldu··|maður} \dicIPA{{h}{\textscy}{l}{\textsubring{d}}{\textscy}{m}{a}{ð}{\textscy}{\textsubring{r}}} \dicPos{m}[13] \dicFlx{(‑manns, ‑menn)}[2] \textbf{1.} \dicFieldCat{pov.} \dicSynonym{álfur} \dicDirectTranslationCS{skřítek, elf}  \textbf{2.} \dicLangCat{přen.} \dicDirectTranslationCS{muž v~pozadí}
\dicEntry[hulið] \dicTerm{hulið} \dicIPA{{h}{\textscy}{\textlengthmark}{l}{\textsci}{\texttheta}} \dicPos{v} \dicFlx{supin} \dicLink{hylja}
\dicEntry[hulinn] \dicTerm{hulinn} \dicsymFrequent\  \dicIPA{{h}{\textscy}{\textlengthmark}{l}{\textsci}{\textsubring{n}}} \dicPos{adj}[6]\dicFlx{}[-7] \dicDirectTranslationCS{zahalený, zakrytý, skrytý};  \dicPhraseIS{vera hulinn e‑u} \dicDirectTranslationCS{být (čím) zakrytý} \dicExampleIS{Tindurinn er oft hulinn þoku.} \dicExampleCS{Vrchol je často ukrytý v~mlze.}
\dicEntry[hulstur] \dicTerm{hulstur} \dicIPA{{h}{\textscy}{l}{s}{\textsubring{d}}{\textscy}{\textsubring{r}}} \dicPos{n}[2] \dicFlx{(‑s, ‑)}[25] \dicSynonym{hylki} \dicDirectTranslationCS{pouzdro, kryt};  \dicPhraseIS{hulstur utan um e‑ð} \dicDirectTranslationCS{pouzdro na (co)} \dicExampleIS{hulstur utan um gleraugun} \dicExampleCS{pouzdro na brýle}
\dicEntry[humall] \dicTerm{hum|all} \dicIPA{{h}{\textscy}{\textlengthmark}{m}{a}{\textsubring{d}}{\textsubring{l}}} \dicPos{m}[6] \dicFlx{(‑als, ‑lar)}[27] \dicFieldCat{bot.} \dicDirectTranslationCS{chmel, chmel otáčivý} \textit{(l.~{\textLA{Humulus lupulus}})}  \dicsymPhoto\ 
\dicFigure{599.jpg}{Humall}{Humall - Zicha Ondřej, Biolib, Copyright/CC-BY-NC}
\dicEntry[humar] \dicTerm{hum|ar} \dicIPA{{h}{\textscy}{\textlengthmark}{m}{a}{\textsubring{r}}} \dicPos{m}[5] \dicFlx{(‑ars, ‑rar)}[7] \dicFieldCat{zool.} \dicDirectTranslationCS{humr, humr norský} \textit{(l.~{\textLA{Nephrops norvegicus}})}  \dicsymPhoto\ 
\dicFigure{ds_image_humar_0_1.jpg}{Humar}{Humar - Hans Hillewaert, CC BY-SA 3.0}
\dicEntry[humátt] \dicTerm{hum··átt} \dicIPA{{h}{\textscy}{\textlengthmark}{m}{au}{h}{\textsubring{d}}} \dicPos{f}[7] \dicFlx{(‑ar)}[3] \dicPhraseIS{fara í humátt á eftir e‑m} \dicDirectTranslationCS{jít za (kým) v~závěsu}
\dicEntry[humla] \dicTerm{huml|a} \dicIPA{{h}{\textscy}{m}{l}{a}} \dicPos{f}[1] \dicFlx{(‑u, ‑ur)}[7] \dicLink{hunangsfluga}
\dicEntry[humm] \dicTerm{humm}\dicTerm{, hm} \dicIPA{{h}\-{\textscy}\-{m}\-{\textlengthmark}\-} \dicPos{inter} \dicDirectTranslationCS{hm} \dicIndirectTranslationCS{(citoslovce vyjadřující váhání nebo výhradu)}
\dicEntry[humma] \dicTerm{humm|a} \dicIPA{{h}{\textscy}{m}{\textlengthmark}{a}} \dicPos{v}[1] \dicFlx{(‑aði)}[1] \dicSynonym{ræskja} \dicDirectTranslationCS{odkašlat si (v~rozpacích ap.)};  \dicPhraseIS{humma e‑ð fram af sér} \dicDirectTranslationCS{kličkovat před (čím), vyhýbat se (čemu)}
\dicEntry[hunang] \dicTerm{hunang} \dicIPA{{h}{\textscy}{\textlengthmark}{n}{au}{\ng}{\r{g}}} \dicPos{n}[2] \dicFlx{(‑s)}[2] \dicDirectTranslationCS{med}
\dicEntry[hunangsfluga] \dicTerm{hunangs··flug|a}\dicTerm{, humla} \dicIPA{{h}\-{\textscy}\-{\textlengthmark}\-{n}\-{au}\-{\ng}\-{s}\-{f}\-{l}\-{\textscy}\-{\textbabygamma}\-{a}\-} \dicPos{f}[1] \dicFlx{(‑u, ‑ur)}[13] \dicFieldCat{zool.} \dicDirectTranslationCS{čmelák} \textit{(l.~{\textLA{Bombus}})}  \dicsymPhoto\ 
\dicFigure{125.jpg}{Hunangsfluga}{Hunangsfluga - Zicha Ondřej, Biolib, Copyright/CC-BY-NC}
\dicEntry[hundadagar] \dicTerm{hunda··dagar} \dicIPA{{h}{\textscy}{n}{\textsubring{d}}{a}{\textsubring{d}}{a}{\textbabygamma}{a}{\textsubring{r}}} \dicPos{m}[6] \dicFlx{pl}[60] \dicLangCat{zast.} \dicDirectTranslationCS{psí dny, dny pod psa}
\dicEntry[hundaheppni] \dicTerm{hunda··heppn|i} \dicIPA{{h}{\textscy}{n}{\textsubring{d}}{a}{h}{\textepsilon}{h}{\textsubring{b}}{n}{\textsci}} \dicPos{f}[3] \dicFlx{(‑i)}[3] \dicDirectTranslationCS{klika, pořádná náhoda}
\dicEntry[hundakofi] \dicTerm{hunda··kof|i} \dicIPA{{h}{\textscy}{n}{\textsubring{d}}{a}{k\smash{\textsuperscript{h}}}{\textopeno}{v}{\textsci}} \dicPos{m}[1] \dicFlx{(‑a, ‑ar)}[1] \dicDirectTranslationCS{psí bouda}
\dicEntry[hundalíf] \dicTerm{hunda··líf} \dicIPA{{h}{\textscy}{n}{\textsubring{d}}{a}{l}{i}{f}} \dicPos{n}[2] \dicFlx{(‑s)}[2] \dicLangCat{přen.} \dicDirectTranslationCS{psí život}
\dicEntry[hundamatur] \dicTerm{hunda··mat|ur} \dicIPA{{h}{\textscy}{n}{\textsubring{d}}{a}{m}{a}{\textsubring{d}}{\textscy}{\textsubring{r}}} \dicPos{m}[10] \dicFlx{(‑ar)}[15] \dicDirectTranslationCS{jídlo pro psy}
\dicEntry[hundasúra] \dicTerm{hunda··súr|a} \dicIPA{{h}{\textscy}{n}{\textsubring{d}}{a}{s}{u}{r}{a}} \dicPos{f}[1] \dicFlx{(‑u, ‑ur)}[7] \dicFieldCat{bot.} \dicDirectTranslationCS{šťovík menší} \textit{(l.~{\textLA{Rumex acetosella}})}  \dicsymPhoto\ 
\dicFigure{ds_image_hundasura_0_1.jpg}{Hundasúra}{Hundasúra - Alterego(fin), CC-BY-SA 3.0}
\dicEntry[hundaæði] \dicTerm{hunda··æði} \dicIPA{{h}{\textscy}{n}{\textsubring{d}}{a}{a}{i}{ð}{\textsci}} \dicPos{n}[2] \dicFlx{(‑s)}[20] \dicFieldCat{med.} \dicDirectTranslationCS{vzteklina, rabies}
\dicEntry[hundflatur] \dicTerm{hund··|flatur} \dicIPA{{h}{\textscy}{n}{\textsubring{d}}{f}{l}{a}{\textsubring{d}}{\textscy}{\textsubring{r}}} \dicPos{adj}[1] \dicFlx{(f ‑flöt)}[3] \dicPhraseIS{liggja hundflatur fyrir e‑m} \dicLangCat{přen.} \dicDirectTranslationCS{plazit se před (kým), podlézat (komu)}
\dicEntry[hundingjaháttur] \dicTerm{hundingja··hátt|ur} \dicIPA{{h}{\textscy}{n}{\textsubring{d}}{i}{\textltailn}{\r{\textObardotlessj}}{a}{h}{au}{h}{\textsubring{d}}{\textscy}{\textsubring{r}}} \dicPos{m}[12] \dicFlx{(‑ar)}[8] \dicDirectTranslationCS{cynismus}
\dicEntry[hundleiðinlegur] \dicTerm{hund··leiðin·legur} \dicIPA{{h}{\textscy}{n}{\textsubring{d}}{l}{ei}{ð}{\textsci}{n}{l}{\textepsilon}{\textbabygamma}{\textscy}{\textsubring{r}}} \dicPos{adj}[1]\dicFlx{}[-8] \dicDirectTranslationCS{strašně\,/\addthin příšerně nudný}
\dicEntry[hundrað] \dicTerm{hundr|að\smash{\textsuperscript{1}}} \dicsymFrequent\  \dicIPA{{h}{\textscy}{n}{\textsubring{d}}{r}{a}{\texttheta}} \dicPos{n}[3] \dicFlx{(‑aðs, ‑uð)}[3] \dicDirectTranslationCS{sto, stovka} \dicExampleIS{Þeir ætla að kaupa hundruð nýja bíla.} \dicExampleCS{Chtějí nakoupit stovky nových aut.}
\dicEntry[hundrað] \dicTerm{hundrað\smash{\textsuperscript{2}}} \dicIPA{{h}{\textscy}{n}{\textsubring{d}}{r}{a}{\texttheta}} \dicPos{num} \dicFlx{indecl} \dicDirectTranslationCS{100, sto}
\dicEntry[hundraðasti] \dicTerm{hundrað··asti} \dicIPA{{h}{\textscy}{n}{\textsubring{d}}{r}{a}{ð}{a}{s}{\textsubring{d}}{\textsci}} \dicPos{num}[8] \dicFlx{ord} \dicDirectTranslationCS{stý}
\dicEntry[hundraðfaldur] \dicTerm{hundrað··|faldur} \dicIPA{{h}{\textscy}{n}{\textsubring{d}}{r}{a}{ð}{f}{a}{l}{\textsubring{d}}{\textscy}{\textsubring{r}}} \dicPos{adj}[2] \dicFlx{(f ‑föld)}[16] \textbf{1.} \dicDirectTranslationCS{stonásobný}  \textbf{2.} \dicDirectTranslationCS{stodílný}
\dicEntry[hundraðkall] \dicTerm{hundrað··kall} \dicIPA{{h}{\textscy}{n}{\textsubring{d}}{r}{a}{ð}{k\smash{\textsuperscript{h}}}{a}{\textsubring{d}}{\textsubring{l}}} \dicPos{m}[4] \dicFlx{(‑s, ‑ar)}[4] \textbf{1.} \dicLangCat{hovor.} \dicDirectTranslationCS{sto (korun)} \dicIndirectTranslationCS{(částka)}  \textbf{2.} \dicLangCat{hovor.} \dicDirectTranslationCS{stokoruna, stovka} \dicIndirectTranslationCS{(mince)}
\dicEntry[hundraðshluti] \dicTerm{hundraðs··hlut|i} \dicIPA{{h}{\textscy}{n}{\textsubring{d}}{r}{a}{ð}{s}{\textsubring{l}}{\textscy}{\textsubring{d}}{\textsci}} \dicPos{m}[1] \dicFlx{(‑a, ‑ar)}[1] \dicSynonym{prósent} \dicDirectTranslationCS{procento, procentní díl}
\dicEntry[hundruð] \dicTerm{hundruð} \dicIPA{{h}{\textscy}{n}{\textsubring{d}}{r}{\textscy}{\texttheta}} \dicPos{n} \dicFlx{pl nom} \dicLink{hundrað\smash{\textsuperscript{1}}}
\dicEntry[hundshaus] \dicTerm{hunds··haus} \dicIPA{{h}{\textscy}{n}{\textsubring{d}}{s}{h}{\oe i}{s}} \dicPos{m}[4] \dicFlx{(‑s, ‑ar)}[14] \dicPhraseIS{setja upp hundshaus} \dicDirectTranslationCS{zatvářit se mrzutě}
\dicEntry[hundskamma] \dicTerm{hund··skamm|a} \dicIPA{{h}{\textscy}{n}{\textsubring{d}}{s}{\r{g}}{a}{m}{a}} \dicPos{v}[1] \dicFlx{(‑aði)}[13] \dicFlx{acc} \dicDirectTranslationCS{vynadat, nadávat, (vy)hubovat, (vy)plísnit} \dicExampleIS{hundskamma e‑n fyrir e‑ð} \dicExampleCS{vyplísnit (koho) za (co)}
\dicEntry[hundspott] \dicTerm{hund··spott} \dicIPA{{h}{\textscy}{n}{\textsubring{d}}{s}{\textsubring{b}}{\textopeno}{h}{\textsubring{d}}} \dicPos{n}[2] \dicFlx{(‑s)}[2] \dicLangCat{han.} \dicSynonym{úrþvætti} \dicDirectTranslationCS{šmejd, hovado}
\dicEntry[hundur] \dicTerm{hund|ur} \dicsymFrequent\  \dicIPA{{h}{\textscy}{n}{\textsubring{d}}{\textscy}{\textsubring{r}}} \dicPos{m}[6] \dicFlx{(‑s, ‑ar)}[4] \dicFieldCat{zool.} \dicDirectTranslationCS{pes} \textit{(l.~{\textLA{Canis}})}  \dicsymPhoto\ ;  \dicPhraseIS{rauðir hundar} \dicFieldCat{med.} \dicDirectTranslationCS{zarděnky};  \dicPhraseIS{þar liggur hundurinn grafinn} \dicLangCat{přen.} \dicDirectTranslationCS{v~tom je zakopaný pes};  \dicPhraseIS{bjóða hundi heila köku} \dicLangCat{přen.} \dicDirectTranslationCS{házet perly sviním};  \dicPhraseIS{éttu hund á helginni og mús (hrafn) í miðri viku} \dicLangCat{přen.} \dicDirectTranslationCS{měl by ses stydět};  \dicPhraseIS{fara í hundana (í hund og kött\,/\addthin í hund og hrafn)} \dicLangCat{přen.} \dicDirectTranslationCS{mařit svůj život, ničit si život};  \dicPhraseIS{standa á e‑u eins og hundur á roði} \dicLangCat{přen.} \dicSynonym*{halda fast við e‑ð} \dicDirectTranslationCS{držet se (čeho) jako klíště};  \dicPhraseIS{áður en köttinn dagar og hundinum er setuljóst } \dicFlx{adv} \dicSynonym*{eldsnemma dags} \dicDirectTranslationCS{za kuropění}
\dicFigure{ds_image_hundur_0_2.jpg}{Hundur}{Hundur - hvalur.org, CC BY 3.0}
\dicEntry[hungra] \dicTerm{hungr|a} \dicIPA{{h}{\textscy}{\ng}{\r{g}}{r}{a}} \dicPos{v}[1] \dicFlx{(‑aði)}[66] \dicFlx{impers} \dicPhraseIS{e‑n hungrar} \dicDirectTranslationCS{(kdo) hladoví};  \dicPhraseIS{e‑n hungrar eftir e‑u} \dicDirectTranslationCS{(kdo) po (čem) prahne}
\dicEntry[hungraður] \dicTerm{hungr··|aður} \dicIPA{{h}{u}{\ng}{\r{g}}{r}{a}{ð}{\textscy}{\textsubring{r}}} \dicPos{adj}[3] \dicFlx{(f ‑uð)}[3] \dicSynonym{soltinn} \dicDirectTranslationCS{hladový, hladovějící, mající hlad}
\dicEntry[hungur] \dicTerm{hungur} \dicsymFrequent\  \dicIPA{{h}{u}{\ng}{\r{g}}{\textscy}{\textsubring{r}}} \dicPos{n}[2] \dicFlx{(‑s)}[28] \dicSynonym{sultur} \dicDirectTranslationCS{hlad, hladovění} \dicExampleIS{eyða hungri í heiminum} \dicExampleCS{zastavit hladovění ve světě}
\dicEntry[hungursneyð] \dicTerm{hungurs··neyð} \dicIPA{{h}{u}{\ng}{\r{g}}{\textscy}{\textsubring{r}}{s}{\textsubring{d}}{n}{ei}{\texttheta}} \dicPos{f}[7] \dicFlx{(‑ar, ‑ir)}[1] \dicDirectTranslationCS{hladomor}
\dicEntry[hungurverkfall] \dicTerm{hungur··verk·|fall} \dicIPA{{h}{u}{\ng}{\r{g}}{\textscy}{r}{v}{\textepsilon}{\textsubring{r}}{\r{g}}{f}{a}{\textsubring{d}}{\textsubring{l}}} \dicPos{n}[2] \dicFlx{(‑falls, ‑föll)}[8] \dicDirectTranslationCS{hladovka}
\dicEntry[hunsa] \dicTerm{huns|a} \dicIPA{{h}{\textscy}{n}{s}{a}} \dicPos{v}[1] \dicFlx{(‑aði)}[1] \dicFlx{acc} \dicDirectTranslationCS{ignorovat, přehlížet} \dicExampleIS{hunsa fyrirmæli læknisins} \dicExampleCS{ignorovat pokyn lékaře}
\dicEntry[hurð] \dicTerm{hurð} \dicsymFrequent\  \dicIPA{{h}{\textscy}{r}{\texttheta}} \dicPos{f}[7] \dicFlx{(‑ar, ‑ir)}[1] \dicDirectTranslationCS{dveře} \dicExampleIS{hurðin á bílnum} \dicExampleCS{dveře u~auta}
\dicEntry[Hurðaskellir] \dicTerm{Hurða··skell|ir}\dicTerm{, Faldafeykir} \dicIPA{{h}\-{\textscy}\-{r}\-{ð}\-{a}\-{s}\-{\r{\textObardotlessj}}\-{\textepsilon}\-{\textsubring{d}}\-{l}\-{\textsci}\-{\textsubring{r}}\-} \dicPos{m}[7] \dicFlx{(‑is)}[3] \dicFlx{prop} \dicFieldCat{pov.} \dicDirectTranslationCS{Hurðaskellir} \dicIndirectTranslationCS{(7. vánoční skřítek, který rád bouchá dveřmi, především v~noci; přichází 18. prosince)}
\dicEntry[hurfum] \dicTerm{hurfum} \dicIPA{{h}{\textscy}{r}{v}{\textscy}{\textsubring{m}}} \dicPos{v} \dicFlx{ind pf pl 1 pers} \dicLink{hverfa\smash{\textsuperscript{2}}}
\dicEntry[húdd] \dicTerm{húdd} \dicIPA{{h}{u}{\textsubring{d}}{\textlengthmark}} \dicPos{n}[2] \dicFlx{(‑s, ‑)}[5] \dicLangCat{hovor.} \dicDirectTranslationCS{kapota}
\dicEntry[húð] \dicTerm{húð} \dicsymFrequent\  \dicIPA{{h}{u}{\textlengthmark}{\texttheta}} \dicPos{f}[7] \dicFlx{(‑ar, ‑ir)}[1] \textbf{1.} \dicSynonym{skinn} \dicDirectTranslationCS{pleť, kůže} \dicExampleIS{vera með dökka húð} \dicExampleCS{mít tmavou pleť}  \textbf{2.} \dicSynonym*{þunnt lag} \dicDirectTranslationCS{povlak, potah};  \dicPhraseIS{í húð og hár} \dicFlx{adv} \dicDirectTranslationCS{tělem i~duší}
\dicEntry[húða] \dicTerm{húð|a} \dicIPA{{h}{u}{\textlengthmark}{ð}{a}} \dicPos{v}[1] \dicFlx{(‑aði)}[1] \dicFlx{acc} \dicDirectTranslationCS{nanést, nanášet (vrstvu ap.)}
\dicEntry[húðflúr] \dicTerm{húð··flúr} \dicIPA{{h}{u}{ð}{f}{l}{u}{\textsubring{r}}} \dicPos{n}[2] \dicFlx{(‑s)}[2] \dicDirectTranslationCS{tetování}
\dicEntry[húðflúra] \dicTerm{húð··flúr|a} \dicIPA{{h}{u}{ð}{f}{l}{u}{r}{a}} \dicPos{v}[1] \dicFlx{(‑aði)}[1] \dicFlx{acc} \dicDirectTranslationCS{(vy)tetovat}
\dicEntry[húðkeipur] \dicTerm{húð··keip|ur} \dicIPA{{h}{u}{ð}{c\smash{\textsuperscript{h}}}{ei}{\textsubring{b}}{\textscy}{\textsubring{r}}} \dicPos{m}[6] \dicFlx{(‑s, ‑ar)}[22] \dicFieldCat{nám.} \dicDirectTranslationCS{kajak}
\dicEntry[húðsjúkdómafræði] \dicTerm{húð·sjúk·dóma··fræð|i} \dicIPA{{h}\-{u}\-{ð}\-{s}\-{j}\-{u}\-{\r{g}}\-{\textsubring{d}}\-{ou}\-{m}\-{a}\-{f}\-{r}\-{a}\-{i}\-{ð}\-{\textsci}\-} \dicPos{f}[3] \dicFlx{(‑i)}[3] \dicFieldCat{med.} \dicDirectTranslationCS{dermatologie}
\dicEntry[húðsjúkdómalæknir] \dicTerm{húð·sjúk·dóma··lækn|ir} \dicIPA{{h}\-{u}\-{ð}\-{s}\-{j}\-{u}\-{\r{g}}\-{\textsubring{d}}\-{ou}\-{m}\-{a}\-{l}\-{a}\-{i}\-{h}\-{\r{g}}\-{n}\-{\textsci}\-{\textsubring{r}}\-} \dicPos{m}[7] \dicFlx{(‑is, ‑ar)}[1] \dicDirectTranslationCS{dermatolog, dermatoložka}
\dicEntry[húðskamma] \dicTerm{húð··skamm|a} \dicIPA{{h}{u}{ð}{s}{\r{g}}{a}{m}{a}} \dicPos{v}[1] \dicFlx{(‑aði)}[13] \dicFlx{acc} \dicSynonym{hundskamma} \dicDirectTranslationCS{peskovat, (vy)plísnit, (vy)hubovat} \dicExampleIS{húðskamma e‑n fyrir e‑ð} \dicExampleCS{vyplísnit (koho) za (co)}
\dicEntry[húðstrýking] \dicTerm{húð··strýk·ing} \dicIPA{{h}{u}{ð}{s}{\textsubring{d}}{r}{i}{\r{\textObardotlessj}}{i}{\ng}{\r{g}}} \dicPos{f}[4] \dicFlx{(‑ar, ‑ar)}[5] \dicDirectTranslationCS{mrskání, bičování}
\dicEntry[húðstrýkja] \dicTerm{húð··strýk|ja} \dicIPA{{h}{u}{ð}{s}{\textsubring{d}}{r}{i}{\r{\textObardotlessj}}{a}} \dicPos{v}[2] \dicFlx{(‑ti, ‑t)}[25] \dicFlx{acc} \dicDirectTranslationCS{mrskat, bičovat}
\dicEntry[húfa] \dicTerm{húf|a} \dicsymFrequent\  \dicIPA{{h}{u}{\textlengthmark}{a}} \dicPos{f}[1] \dicFlx{(‑u, ‑ur)}[19] \dicDirectTranslationCS{čepice, kulich, čapka, čepec} \dicExampleIS{Þær sauma húfur fyrir lítil börn.} \dicExampleCS{Šijí čepice pro malé děti.};  \dicPhraseIS{setja upp húfu(na)} \dicDirectTranslationCS{nasadit si čepici};  \dicPhraseIS{taka ofan húfu} \dicDirectTranslationCS{sundat si čepici}
\dicEntry[húftrygging] \dicTerm{húf··trygg·ing} \dicIPA{{h}{u}{v}{t\smash{\textsuperscript{h}}}{r}{\textsci}{\r{\textObardotlessj}}{i}{\ng}{\r{g}}} \dicPos{f}[4] \dicFlx{(‑ar, ‑ar)}[5] \dicFieldCat{práv.} \dicDirectTranslationCS{kasko pojištění}
\dicEntry[húfur] \dicTerm{húf|ur} \dicIPA{{h}{u}{\textlengthmark}{v}{\textscy}{\textsubring{r}}} \dicPos{m}[6] \dicFlx{(‑s, ‑ar)}[8] \dicPhraseIS{heill á húfi} \dicFlx{adj} \dicDirectTranslationCS{živ a~zdráv, nezraněný};  \dicPhraseIS{(það er) mikið í húfi} \dicLangCat{přen.} \dicDirectTranslationCS{(je) hodně v~sázce}
\dicEntry[húka] \dicTerm{húk|a} \dicIPA{{h}{u}{\textlengthmark}{\r{g}}{a}} \dicPos{v}[2] \dicFlx{(‑ti, ‑t)}[75] \dicDirectTranslationCS{dřepět,  vysedávat}
\dicEntry[húkka] \dicTerm{húkk|a} \dicIPA{{h}{u}{h}{\r{g}}{a}} \dicPos{v}[1] \dicFlx{(‑aði)}[1] \dicFlx{acc} \textbf{1.} \dicSynonym*{krækja í} \dicDirectTranslationCS{zaháknout (rybu na háček ap.)}  \textbf{2.} \dicSynonym*{ná í} \dicDirectTranslationCS{dostat, zajistit si (odvoz ap.)};  \dicIdiom{húkkast}{ \dicPhraseIS{húkkast}} \dicFlx{refl} \dicDirectTranslationCS{stát se závislý (na počítačové hře ap.)}
\dicEntry[húm] \dicTerm{húm} \dicIPA{{h}{u}{\textlengthmark}{\textsubring{m}}} \dicPos{n}[2] \dicFlx{(‑s)}[2] \dicSynonym{rökkur} \dicDirectTranslationCS{soumrak, šero}
\dicEntry[húma] \dicTerm{húm|a} \dicIPA{{h}{u}{\textlengthmark}{m}{a}} \dicPos{v}[1] \dicFlx{(‑aði)}[75] \dicPhraseIS{það húmar (að)} \dicFlx{impers} \dicDirectTranslationCS{smráká se, padá soumrak}
\dicEntry[húmanismi] \dicTerm{húman··ism|i} \dicIPA{{h}{u}{\textlengthmark}{m}{a}{n}{\textsci}{s}{m}{\textsci}} \dicPos{m}[1] \dicFlx{(‑a)}[3] \dicDirectTranslationCS{humanismus}
\dicEntry[húmanisti] \dicTerm{húman··ist|i} \dicIPA{{h}{u}{\textlengthmark}{m}{a}{n}{\textsci}{s}{\textsubring{d}}{\textsci}} \dicPos{m}[1] \dicFlx{(‑a, ‑ar)}[1] \textbf{1.} \dicSynonym{mannvinur} \dicDirectTranslationCS{humanista, humanistka} \dicIndirectTranslationCS{(humánní člověk)}  \textbf{2.} \dicSynonym*{hugvísindamaður} \dicDirectTranslationCS{humanista, humanistka} \dicIndirectTranslationCS{(odborník\,/\addthin student humanitních věd)}
\dicEntry[húmanískur] \dicTerm{húmanískur} \dicIPA{{h}{u}{\textlengthmark}{m}{a}{n}{i}{s}{\r{g}}{\textscy}{\textsubring{r}}} \dicPos{adj}[1]\dicFlx{}[-1] \dicDirectTranslationCS{humanistický}
\dicEntry[húmor] \dicTerm{húmor} \dicIPA{{h}{u}{\textlengthmark}{m}{\textopeno}{\textsubring{r}}} \dicPos{m}[4] \dicFlx{(‑s)}[18] \dicSynonym{spaug} \dicDirectTranslationCS{humor} \dicExampleIS{svartur húmor} \dicExampleCS{černý humor}
\dicEntry[hún] \dicTerm{hún} \dicsymFrequent\  \dicIPA{{h}{u}{\textlengthmark}{\textsubring{n}}} \dicPos{pron}[17] \dicFlx{pers} \dicFlx{nom sg f} \dicDirectTranslationCS{ona} \dicExampleIS{Hún hefur gefið út sína fyrstu bók.} \dicExampleCS{Vydala svou první knihu.}
\dicEntry[húnn] \dicTerm{hún|n} \dicIPA{{h}{u}{\textsubring{d}}{\textsubring{n}}} \dicPos{m}[6] \dicFlx{(‑s, ‑ar)}[42] \textbf{1.} \dicDirectTranslationCS{koule (na zakončení stožáru ap.)}  \textbf{2.} \dicSynonym{handfang} \dicDirectTranslationCS{kulatá klika (u~dveří ap.)}  \textbf{3.} \dicSynonym*{bjarndýrsungi} \dicDirectTranslationCS{medvídě, mládě medvěda}
\dicEntry[húrra] \dicTerm{húrr|a\smash{\textsuperscript{1}}} \dicIPA{{h}{u}{r}{\textlengthmark}{a}} \dicPos{v}[1] \dicFlx{(‑aði)}[1] \dicSynonym{renna\smash{\textsuperscript{3}}} \dicDirectTranslationCS{pádit, uhánět, letět} \dicExampleIS{Ég húrraði niður alla brekkuna.} \dicExampleCS{Pádil jsem dolů z~kopce.}
\dicEntry[húrra] \dicTerm{húrra\smash{\textsuperscript{2}}} \dicIPA{{h}{u}{r}{\textlengthmark}{a}} \dicPos{inter} \dicDirectTranslationCS{hurá}
\dicEntry[hús] \dicTerm{hús} \dicsymFrequent\  \dicIPA{{h}{u}{\textlengthmark}{s}} \dicPos{n}[2] \dicFlx{(‑s, ‑)}[5] \dicSynonym{bygging} \dicDirectTranslationCS{dům, stavení} \dicExampleIS{Hjón eru að leita að nýju húsi.} \dicExampleCS{Manželé hledají nový dům.}
\dicEntry[húsagarður] \dicTerm{húsa··garð|ur} \dicIPA{{h}{u}{\textlengthmark}{s}{a}{\r{g}}{a}{r}{ð}{\textscy}{\textsubring{r}}} \dicPos{m}[6] \dicFlx{(‑s, ‑ar)}[5] \dicDirectTranslationCS{dvůr, dvorek} \dicIndirectTranslationCS{(nekryté uzavřené prostranství u~domu)}
\dicEntry[húsageitungur] \dicTerm{húsa··geit·ung|ur} \dicIPA{{h}{u}{\textlengthmark}{s}{a}{\r{\textObardotlessj}}{ei}{\textsubring{d}}{u}{\ng}{\r{g}}{\textscy}{\textsubring{r}}} \dicPos{m}[6] \dicFlx{(‑s, ‑ar)}[8] \dicFieldCat{zool.} \dicDirectTranslationCS{vosa útočná} \textit{(l.~{\textLA{Vespula germanica}})}
\dicEntry[húsagerðarlist] \dicTerm{húsa·gerðar··list} \dicIPA{{h}{u}{\textlengthmark}{s}{a}{\r{\textObardotlessj}}{\textepsilon}{r}{ð}{a}{r}{l}{\textsci}{s}{\textsubring{d}}} \dicPos{f}[7] \dicFlx{(‑ar)}[3] \dicSynonym{byggingarlist} \dicDirectTranslationCS{stavitelství, architektura}
\dicEntry[húsakynni] \dicTerm{húsa··kynni} \dicIPA{{h}{u}{\textlengthmark}{s}{a}{c\smash{\textsuperscript{h}}}{\textsci}{n}{\textsci}} \dicPos{n}[2] \dicFlx{pl}[19] \dicDirectTranslationCS{prostory (domu ap.)}
\dicEntry[húsaleiga] \dicTerm{húsa··leig|a} \dicIPA{{h}{u}{\textlengthmark}{s}{a}{l}{ei}{\textbabygamma}{a}} \dicPos{f}[1] \dicFlx{(‑u)}[5] \dicDirectTranslationCS{nájemné, nájem, činže}
\dicEntry[húsamálari] \dicTerm{húsa··mál·ar|i} \dicIPA{{h}{u}{\textlengthmark}{s}{a}{m}{au}{l}{a}{r}{\textsci}} \dicPos{m}[1] \dicFlx{(‑a, ‑ar)}[13] \dicDirectTranslationCS{malíř(ka)} \dicIndirectTranslationCS{(řemeslník zabývající se malováním pokojů ap.)}
\dicEntry[húsameistari] \dicTerm{húsa··meist·ar|i} \dicIPA{{h}{u}{\textlengthmark}{s}{a}{m}{ei}{s}{\textsubring{d}}{a}{r}{\textsci}} \dicPos{m}[1] \dicFlx{(‑a, ‑ar)}[13] \dicDirectTranslationCS{stavitel(ka), architekt(ka)}
\dicEntry[húsasmiður] \dicTerm{húsa··smið|ur} \dicIPA{{h}{u}{\textlengthmark}{s}{a}{s}{m}{\textsci}{ð}{\textscy}{\textsubring{r}}} \dicPos{m}[9] \dicFlx{(‑s, ‑ir)}[8] \dicDirectTranslationCS{tesař(ka)}
\dicEntry[húsbátur] \dicTerm{hús··bát|ur} \dicIPA{{h}{u}{\textlengthmark}{s}{\textsubring{b}}{au}{\textsubring{d}}{\textscy}{\textsubring{r}}} \dicPos{m}[6] \dicFlx{(‑s, ‑ar)}[22] \dicFieldCat{nám.} \dicDirectTranslationCS{hausbót}
\dicEntry[húsbíll] \dicTerm{hús··bíl|l} \dicIPA{{h}{u}{\textlengthmark}{s}{\textsubring{b}}{i}{\textsubring{d}}{\textsubring{l}}} \dicPos{m}[6] \dicFlx{(‑s, ‑ar)}[48] \dicDirectTranslationCS{karavan}
\dicEntry[húsbóndi] \dicTerm{hús··|bóndi} \dicsymFrequent\  \dicIPA{{h}{u}{\textlengthmark}{s}{\textsubring{b}}{ou}{n}{\textsubring{d}}{\textsci}} \dicPos{m}[2] \dicFlx{(‑bónda, ‑bændur)}[4] \dicDirectTranslationCS{pan domácí, hospodář} \dicExampleIS{harður húsbóndi} \dicExampleCS{přísný pán domu}
\dicEntry[húsbúnaður] \dicTerm{hús··bú·nað|ur} \dicIPA{{h}{u}{\textlengthmark}{s}{\textsubring{b}}{u}{n}{a}{ð}{\textscy}{\textsubring{r}}} \dicPos{m}[10] \dicFlx{(‑ar)}[9] \dicDirectTranslationCS{vybavení domu}
\dicEntry[húsdúfa] \dicTerm{hús··dúf|a} \dicIPA{{h}{u}{\textlengthmark}{s}{\textsubring{d}}{u}{v}{a}} \dicPos{f}[1] \dicFlx{(‑u, ‑ur)}[13] \dicFieldCat{zool.} \dicDirectTranslationCS{holub, holub domácí} \textit{(l.~{\textLA{Columba livia f. domestica}})}
\dicEntry[húsdýr] \dicTerm{hús··dýr} \dicIPA{{h}{u}{\textlengthmark}{s}{\textsubring{d}}{i}{\textsubring{r}}} \dicPos{n}[2] \dicFlx{(‑s, ‑)}[5] \dicDirectTranslationCS{domácí zvíře}
\dicEntry[húseigandi] \dicTerm{hús··eig·|andi} \dicIPA{{h}{u}{\textlengthmark}{s}{ei}{\textbabygamma}{a}{n}{\textsubring{d}}{\textsci}} \dicPos{m}[2] \dicFlx{(‑anda, ‑endur)}[1] \dicDirectTranslationCS{bytný, bytná, majitel(ka) domu}
\dicEntry[húsfreyja] \dicTerm{hús··freyj|a} \dicIPA{{h}{u}{\textlengthmark}{s}{f}{r}{ei}{j}{a}} \dicPos{f}[1] \dicFlx{(‑u, ‑ur)}[7] \dicDirectTranslationCS{hospodyně}
\dicEntry[húsgagn] \dicTerm{hús··|gagn} \dicsymFrequent\  \dicIPA{{h}{u}{\textlengthmark}{s}{\r{g}}{a}{\r{g}}{\textsubring{n}}} \dicPos{n}[2] \dicFlx{(‑gagns, ‑gögn)}[8] \dicDirectTranslationCS{nábytek} \dicIndirectTranslationCS{(zvláště v~pl)} \dicExampleIS{setustofa með gömlum húsgögnum} \dicExampleCS{obývací pokoj se starým nábytkem}
\dicEntry[húsgagnasmiður] \dicTerm{hús·gagna··smið|ur} \dicIPA{{h}{u}{\textlengthmark}{s}{\r{g}}{a}{\r{g}}{n}{a}{s}{m}{\textsci}{ð}{\textscy}{\textsubring{r}}} \dicPos{m}[9] \dicFlx{(‑s, ‑ir)}[8] \dicDirectTranslationCS{nábytkář(ka), truhlář(ka)}
\dicEntry[húsgögn] \dicTerm{hús··gögn} \dicIPA{{h}{u}{\textlengthmark}{s}{\r{g}}{\oe}{\r{g}}{\textsubring{n}}} \dicPos{n}[2] \dicFlx{pl}[9] \dicLink{húsgagn}
\dicEntry[húshitun] \dicTerm{hús··hit|un} \dicIPA{{h}{u}{\textlengthmark}{s}{h}{\textsci}{\textsubring{d}}{\textscy}{\textsubring{n}}} \dicPos{f}[7] \dicFlx{(‑unar)}[9] \dicDirectTranslationCS{vytápění (domu)}
\dicEntry[húsleit] \dicTerm{hús··leit} \dicIPA{{h}{u}{\textlengthmark}{s}{l}{ei}{\textsubring{d}}} \dicPos{f}[7] \dicFlx{(‑ar, ‑ir)}[1] \dicDirectTranslationCS{domovní prohlídka}
\dicEntry[húsleitarheimild] \dicTerm{hús·leitar··heimild} \dicIPA{{h}{u}{\textlengthmark}{s}{l}{ei}{\textsubring{d}}{a}{\textsubring{r}}{h}{ei}{m}{\textsci}{l}{\textsubring{d}}} \dicPos{f}[7] \dicFlx{(‑ar, ‑ir)}[1] \dicDirectTranslationCS{povolení k~domovní prohlídce}
\dicEntry[húsmóðir] \dicTerm{hús··|móðir} \dicsymFrequent\  \dicIPA{{h}{u}{\textlengthmark}{s}{m}{ou}{ð}{\textsci}{\textsubring{r}}} \dicPos{f}[11] \dicFlx{(‑móður, ‑mæður)}[1] \dicDirectTranslationCS{paní domácí, hospodyně} \dicExampleIS{hagsýn húsmóðir} \dicExampleCS{šetrná hospodyně}
\dicEntry[húsmunir] \dicTerm{hús··munir} \dicIPA{{h}{u}{\textlengthmark}{s}{m}{\textscy}{n}{\textsci}{\textsubring{r}}} \dicPos{m}[10] \dicFlx{pl}[1] \dicDirectTranslationCS{nábytek, bytové zařízení}
\dicEntry[húsnæði] \dicTerm{hús··næði} \dicsymFrequent\  \dicIPA{{h}{u}{\textlengthmark}{s}{n}{a}{i}{ð}{\textsci}} \dicPos{n}[2] \dicFlx{(‑s)}[20] \dicDirectTranslationCS{bydlení, ubytování} \dicExampleIS{bjóða húsnæði til leigu} \dicExampleCS{nabízet bydlení k~pronájmu};  \dicPhraseIS{fæði og húsnæði} \dicDirectTranslationCS{strava a~ubytování}
\dicEntry[húsnæðislaus] \dicTerm{hús·næðis··laus} \dicIPA{{h}{u}{\textlengthmark}{s}{n}{a}{i}{ð}{\textsci}{s}{l}{\oe i}{s}} \dicPos{adj}[5]\dicFlx{}[-4] \dicDirectTranslationCS{nemající\,/\addthin postrádající bydlení}
\dicEntry[húsnæðislán] \dicTerm{hús·næðis··lán} \dicIPA{{h}{u}{\textlengthmark}{s}{n}{a}{i}{ð}{\textsci}{s}{l}{au}{\textsubring{n}}} \dicPos{n}[2] \dicFlx{(‑s, ‑)}[5] \dicFieldCat{ekon.} \dicDirectTranslationCS{úvěr na bydlení}
\dicEntry[húsnæðismál] \dicTerm{hús·næðis··mál} \dicIPA{{h}{u}{\textlengthmark}{s}{n}{a}{i}{ð}{\textsci}{s}{m}{au}{\textsubring{l}}} \dicPos{n}[2] \dicFlx{(‑s, ‑)}[5] \dicDirectTranslationCS{bytová politika, otázka bydlení}
\dicEntry[húsnæðisskortur] \dicTerm{hús·næðis··skort|ur} \dicIPA{{h}{u}{\textlengthmark}{s}{n}{a}{i}{ð}{\textsci}{s}{\r{g}}{\textopeno}{\textsubring{r}}{\textsubring{d}}{\textscy}{\textsubring{r}}} \dicPos{m}[6] \dicFlx{(‑s)}[7] \dicDirectTranslationCS{bytová nouze, nedostatek bytů}
\dicEntry[húsverk] \dicTerm{hús··verk} \dicIPA{{h}{u}{\textlengthmark}{s}{v}{\textepsilon}{\textsubring{r}}{\r{g}}} \dicPos{n}[2] \dicFlx{(‑s, ‑)}[5] \dicDirectTranslationCS{domácí práce, práce v~domácnosti}
\dicEntry[húsvörður] \dicTerm{hús··|vörður} \dicIPA{{h}{u}{\textlengthmark}{s}{v}{\oe}{r}{ð}{\textscy}{\textsubring{r}}} \dicPos{m}[11] \dicFlx{(‑varðar, ‑verðir)}[5] \dicDirectTranslationCS{správce, správkyně, domovník, domovnice, školník, školnice}
\dicEntry[húsönd] \dicTerm{hús··|önd} \dicIPA{{h}{u}{\textlengthmark}{s}{\oe}{n}{\textsubring{d}}} \dicPos{f}[8] \dicFlx{(‑andar, ‑endur)}[2] \dicFieldCat{zool.} \dicDirectTranslationCS{hohol islandský} \textit{(l.~{\textLA{Bucephala islandica}})}  \dicsymPhoto\ 
\dicFigure{20961.jpg}{Húsönd}{Húsönd - Dewhurst, Donna, Biolib, PD}
\dicEntry[hvað] \dicTerm{hvað} \dicIPA{{k\smash{\textsuperscript{h}}}{v}{a}{\textlengthmark}{\texttheta}} \dicPos{pron} \dicFlx{indecl} \textbf{1.} \dicDirectTranslationCS{co, copak} \dicExampleIS{Hvað ertu að gera?} \dicExampleCS{Co děláš?}  \textbf{2.} \dicDirectTranslationCS{jak} \dicIndirectTranslationCS{(užívané s~příd. jménem nebo příslovcem)} \dicExampleIS{Hvað gamall?} \dicExampleCS{Jak starý?}  \textbf{3.} \dicDirectTranslationCS{kolik} \dicExampleIS{Hvað er klukkan?} \dicExampleCS{Kolik je hodin?};  \dicIdiom{hvað}{ \dicPhraseIS{hvað e‑ð snertir\,/\addthin varðar}} \dicDirectTranslationCS{co se (čeho) týče};  \dicPhraseIS{hvað þá} \dicFlx{adv} {\textbf{a.}} \dicDirectTranslationCS{natož, tím spíše, co teprve};  {\textbf{b.}} \dicDirectTranslationCS{co pak\,/\addthin potom};  \dicPhraseIS{sitt á hvað} \dicFlx{adv} \dicDirectTranslationCS{střídavě, ze strany na stranu}
\dicEntry[hvaða] \dicTerm{hvaða} \dicsymFrequent\  \dicIPA{{k\smash{\textsuperscript{h}}}{v}{a}{\textlengthmark}{ð}{a}} \dicPos{pron} \dicFlx{indecl} \dicFlx{int\,/\addthin indef\,/\addthin dem} \textbf{1.} \dicFlx{int} \dicSynonym{hverskonar} \dicDirectTranslationCS{jaký, který, co} \dicExampleIS{Hvaða mann talaðir þú við?} \dicExampleCS{S~kým si hovořil?}  \textbf{2.} \dicFlx{indef} \dicDirectTranslationCS{jakýkoli(v)};  \dicPhraseIS{hvaða (atburður) sem er} \dicDirectTranslationCS{jakoukoliv (událost)}  \textbf{3.} \dicFlx{dem} \dicDirectTranslationCS{co, takový};  \dicPhraseIS{hvaða vitleysa!} \dicDirectTranslationCS{takový nesmysl!}
\dicEntry[hvaðan] \dicTerm{hvaðan} \dicsymFrequent\  \dicIPA{{k\smash{\textsuperscript{h}}}{v}{a}{\textlengthmark}{ð}{a}{\textsubring{n}}} \dicPos{adv} \textbf{1.} \dicDirectTranslationCS{odkud} \dicIndirectTranslationCS{(užívané v~přímé otázce)} \dicExampleIS{Hvaðan ertu?} \dicExampleCS{Odkud jsi?}  \textbf{2.} \dicDirectTranslationCS{odkud} \dicIndirectTranslationCS{(užívané ve vedlejší větě)};  \dicIdiom{hvaðan}{ \dicPhraseIS{hvaðan sem}} \dicFlx{conj} \dicDirectTranslationCS{odkudkoli(v)};  \dicPhraseIS{hvaðan sem er} \dicFlx{adv} \dicDirectTranslationCS{odkudkoli(v)}
\dicEntry[hvalasafn] \dicTerm{hvala··|safn} \dicIPA{{k\smash{\textsuperscript{h}}}{v}{a}{\textlengthmark}{l}{a}{s}{a}{\textsubring{b}}{\textsubring{n}}} \dicPos{n}[2] \dicFlx{(‑safns, ‑söfn)}[8] \dicDirectTranslationCS{muzeum velryb}
\dicEntry[hvalaskoðun] \dicTerm{hvala··skoð|un} \dicIPA{{k\smash{\textsuperscript{h}}}{v}{a}{\textlengthmark}{l}{a}{s}{\r{g}}{\textopeno}{ð}{\textscy}{\textsubring{n}}} \dicPos{f}[7] \dicFlx{(‑unar)}[9] \dicDirectTranslationCS{pozorování velryb (turistická atrakce)}
\dicEntry[hvalreki] \dicTerm{hval··rek|i} \dicIPA{{k\smash{\textsuperscript{h}}}{v}{a}{l}{r}{\textepsilon}{\r{\textObardotlessj}}{\textsci}} \dicPos{m}[1] \dicFlx{(‑a, ‑ar)}[1] \textbf{1.} \dicSynonym*{reki hvals} \dicDirectTranslationCS{uváznutí velryby na mělčině}  \textbf{2.} \dicSynonym{happ} \dicDirectTranslationCS{klika, nečekané štěstí}
\dicEntry[hvalstöð] \dicTerm{hval··stöð} \dicIPA{{k\smash{\textsuperscript{h}}}{v}{a}{l}{s}{\textsubring{d}}{\oe}{\texttheta}} \dicPos{f}[6] \dicFlx{(‑var, ‑var)}[1] \dicDirectTranslationCS{velrybářská stanice} \dicIndirectTranslationCS{(závod na zpracování velryb)}
\dicEntry[hvalur] \dicTerm{hval|ur} \dicIPA{{k\smash{\textsuperscript{h}}}{v}{a}{\textlengthmark}{l}{\textscy}{\textsubring{r}}} \dicPos{m}[9] \dicFlx{(‑s, ‑ir)}[9] \textbf{1.} \dicDirectTranslationCS{velryba} \dicExampleIS{Hvalurinn blæs.} \dicExampleCS{Velryba vydechuje.}  \textbf{2.} \dicPhraseIS{hvalir} \dicFlx{pl} \dicFieldCat{zool.} \dicDirectTranslationCS{kytovci} \textit{(l.~{\textLA{Cetacea}})}  \dicsymPhoto\ 
\dicFigure{ds_image_hvalur_0_1.jpg}{Hvalur}{Hvalur - Zorankovacevic, CC BY-SA 2.5}
\dicEntry[hvalveiði] \dicTerm{hval··veið|i} \dicIPA{{k\smash{\textsuperscript{h}}}{v}{a}{l}{v}{ei}{ð}{\textsci}} \dicPos{f}[2] \dicFlx{(‑i\,/\addthin ‑ar, ‑ar)}[3] \dicDirectTranslationCS{velrybaření}
\dicEntry[hvalveiðibátur] \dicTerm{hval·veiði··bát|ur} \dicIPA{{k\smash{\textsuperscript{h}}}{v}{a}{l}{v}{ei}{ð}{\textsci}{\textsubring{b}}{au}{\textsubring{d}}{\textscy}{\textsubring{r}}} \dicPos{m}[6] \dicFlx{(‑s, ‑ar)}[22] \dicFieldCat{nám.} \dicDirectTranslationCS{velrybářský člun}
\dicEntry[hvalveiðiskip] \dicTerm{hval·veiði··skip} \dicIPA{{k\smash{\textsuperscript{h}}}{v}{a}{l}{v}{ei}{ð}{\textsci}{s}{\r{\textObardotlessj}}{\textsci}{\textsubring{b}}} \dicPos{n}[2] \dicFlx{(‑s, ‑)}[5] \dicFieldCat{nám.} \dicDirectTranslationCS{velrybářská loď}
\dicEntry[hvammur] \dicTerm{hvamm|ur} \dicIPA{{k\smash{\textsuperscript{h}}}{v}{a}{m}{\textlengthmark}{\textscy}{\textsubring{r}}} \dicPos{m}[6] \dicFlx{(‑s, ‑ar)}[5] \dicDirectTranslationCS{(travnatá) dolina, (zarostlý) žleb}
\dicEntry[Hvannadalshnúkur] \dicTerm{Hvanna·dals··hnúk|ur} \dicIPA{{k\smash{\textsuperscript{h}}}{v}{a}{n}{\textlengthmark}{a}{\textsubring{d}}{a}{l}{s}{\textsubring{n}}{u}{\r{g}}{\textscy}{\textsubring{r}}} \dicPos{m}[6] \dicFlx{(‑s)}[18] \dicFieldCat{geog.} \dicDirectTranslationCS{Hvannadalshnúkur} \dicIndirectTranslationCS{(nejvyšší hora Islandu)}
\dicEntry[hvannar] \dicTerm{hvannar} \dicIPA{{k\smash{\textsuperscript{h}}}{v}{a}{n}{\textlengthmark}{a}{\textsubring{r}}} \dicPos{f} \dicFlx{sg gen} \dicLink{hvönn}
\dicEntry[Hvanneyrarveiki] \dicTerm{Hvann·eyrar··veik|i}\dicTerm{, votheysveiki} \dicIPA{{k\smash{\textsuperscript{h}}}\-{v}\-{a}\-{n}\-{ei}\-{r}\-{a}\-{r}\-{v}\-{ei}\-{\r{\textObardotlessj}}\-{\textsci}\-} \dicPos{f}[3] \dicFlx{(‑i)}[3] \dicFieldCat{med.} \dicDirectTranslationCS{listerióza} \dicIndirectTranslationCS{(infekční onemocnění způsobené grampozitivní bakterií Listeria monocytogenes)}
\dicEntry[hvannir] \dicTerm{hvannir} \dicIPA{{k\smash{\textsuperscript{h}}}{v}{a}{n}{\textlengthmark}{\textsci}{\textsubring{r}}} \dicPos{f} \dicFlx{pl nom} \dicLink{hvönn}
\dicEntry[hvar] \dicTerm{hvar} \dicsymFrequent\  \dicIPA{{k\smash{\textsuperscript{h}}}{v}{a}{\textlengthmark}{\textsubring{r}}} \dicPos{adv} \textbf{1.} \dicDirectTranslationCS{kde, kdepak} \dicIndirectTranslationCS{(užívané v~přímé otázce)} \dicExampleIS{Hvar er hann?} \dicExampleCS{Kde je (on)?}  \textbf{2.} \dicDirectTranslationCS{kde, kdepak} \dicIndirectTranslationCS{(užívané ve vedlejší větě)};  \dicIdiom{hvar}{ \dicPhraseIS{hér og hvar}} \dicFlx{adv} \dicDirectTranslationCS{tu a~tam};  \dicPhraseIS{hvar sem} \dicFlx{conj} \dicDirectTranslationCS{kdekoli(v)};  \dicPhraseIS{hvar sem er} \dicFlx{adv} \dicDirectTranslationCS{kdekoli(v)};  \dicPhraseIS{víðast hvar} \dicFlx{adv} \dicDirectTranslationCS{téměř všude}
\dicEntry[hvarf] \dicTerm{hvarf\smash{\textsuperscript{1}}} \dicIPA{{k\smash{\textsuperscript{h}}}{v}{a}{r}{f}} \dicPos{n}[2] \dicFlx{(‑s, hvörf)}[8] \textbf{1.} \dicSynonym*{það að hverfa} \dicDirectTranslationCS{zmizení, ztracení (se)};  \dicPhraseIS{fara í\,/\addthin á hvarf} \dicDirectTranslationCS{ztratit se z~dohledu}  \textbf{2.} \dicSynonym{hæð} \dicDirectTranslationCS{vyvýšenina (bránící výhledu)}  \textbf{3.} \dicSynonym{hola\smash{\textsuperscript{1}}} \dicDirectTranslationCS{díra, prohlubeň (na cestě ap.)} \dicIndirectTranslationCS{(užívané zvláště v~pl)}  \textbf{4.} \dicFieldCat{chem.} \dicSynonym{efnahvarf} \dicDirectTranslationCS{(chemická) reakce}
\dicEntry[hvarf] \dicTerm{hvarf\smash{\textsuperscript{2}}} \dicIPA{{k\smash{\textsuperscript{h}}}{v}{a}{r}{f}} \dicPos{v} \dicFlx{ind pf sg 1 pers} \dicLink{hverfa\smash{\textsuperscript{2}}}
\dicEntry[hvarfbaugur] \dicTerm{hvarf··baug|ur} \dicIPA{{k\smash{\textsuperscript{h}}}{v}{a}{r}{v}{\textsubring{b}}{\oe i}{\textbabygamma}{\textscy}{\textsubring{r}}} \dicPos{m}[6] \dicFlx{(‑s, ‑ar)}[24] \dicFieldCat{geog.} \dicDirectTranslationCS{obratník};  \dicPhraseIS{nyrðri hvarfbaugur, hvarfbaugur krabbans} \dicFieldCat{geog.} \dicDirectTranslationCS{obratník Raka};  \dicPhraseIS{syðri hvarfbaugur, hvarfbaugur steingeitarinnar} \dicFieldCat{geog.} \dicDirectTranslationCS{obratník Kozoroha}
\dicEntry[hvarfla] \dicTerm{hvarfl|a} \dicsymFrequent\  \dicIPA{{k\smash{\textsuperscript{h}}}{v}{a}{r}{\textsubring{d}}{l}{a}} \dicPos{v}[1] \dicFlx{(‑aði)}[50] \dicFlx{dat} \textbf{1.} \dicDirectTranslationCS{bloudit, těkat (o~očích ap.)} \dicExampleIS{Augun hvarfla til og frá.} \dicExampleCS{Oči bloudí sem a~tam.}  \textbf{2.} \dicPhraseIS{það hvarflar að e‑m að (fara)} \dicFlx{impers} \dicDirectTranslationCS{(koho) napadá (odejít)}
\dicEntry[hvarmur] \dicTerm{hvarm|ur} \dicIPA{{k\smash{\textsuperscript{h}}}{v}{a}{r}{m}{\textscy}{\textsubring{r}}} \dicPos{m}[6] \dicFlx{(‑s, ‑ar)}[5] \dicSynonym{augnalok} \dicDirectTranslationCS{(oční) víčko} \dicExampleIS{kipra hvarmana} \dicExampleCS{přimhouřit oči}
\dicEntry[hvass] \dicTerm{hvass} \dicsymFrequent\  \dicIPA{{k\smash{\textsuperscript{h}}}{v}{a}{s}{\textlengthmark}} \dicPos{adj}[5] \dicFlx{(f hvöss)}[2] \textbf{1.} \dicDirectTranslationCS{ostrý (nůž ap.)} \dicExampleIS{hvass hnífur} \dicExampleCS{ostrý nůž}  \textbf{2.} \dicSynonym{napur} \dicDirectTranslationCS{ostrý, jízlivý} \dicExampleIS{hvass í orðum} \dicExampleCS{používající ostrá slova}  \textbf{3.} \dicSynonym{skarpur} \dicDirectTranslationCS{ostrý, pronikavý, bystrý (pohled ap.)} \dicExampleIS{hvöss sjón} \dicExampleCS{pronikavý pohled}  \textbf{4.} \dicSynonym{harður} \dicDirectTranslationCS{ostrý, prudký (boj ap.)} \dicExampleIS{hvöss árás} \dicExampleCS{prudký boj}  \textbf{5.} \dicDirectTranslationCS{prudký, silný (vítr ap.)};  \dicPhraseIS{það er hvasst} \dicDirectTranslationCS{je větrno}
\dicEntry[hvasst] \dicTerm{hvass|t} \dicsymFrequent\  \dicIPA{{k\smash{\textsuperscript{h}}}{v}{a}{s}{\textsubring{d}}} \dicPos{adv} \dicFlx{(comp ‑ar, sup ‑ast)} \dicDirectTranslationCS{ostře, prudce} \dicExampleIS{snúa sér hvasst að e‑m} \dicExampleCS{prudce se obrátit ke (komu)}
\dicEntry[hvassviðri] \dicTerm{hvass··viðri} \dicIPA{{k\smash{\textsuperscript{h}}}{v}{a}{s}{\textlengthmark}{v}{\textsci}{ð}{r}{\textsci}} \dicPos{n}[2] \dicFlx{(‑s, ‑)}[14] \dicDirectTranslationCS{bouřlivý vítr} \dicIndirectTranslationCS{(Beaufort 8) (17,2--20,7~m/s)}
\dicEntry[hvassyrtur] \dicTerm{hvass··yrtur} \dicIPA{{k\smash{\textsuperscript{h}}}{v}{a}{s}{\textlengthmark}{\textsci}{\textsubring{r}}{\textsubring{d}}{\textscy}{\textsubring{r}}} \dicPos{adj}[1]\dicFlx{}[-10] \dicLangCat{přen.} \dicDirectTranslationCS{mající ostrý jazyk}
\dicEntry[hvatamaður] \dicTerm{hvata··|maður} \dicIPA{{k\smash{\textsuperscript{h}}}{v}{a}{\textlengthmark}{\textsubring{d}}{a}{m}{a}{ð}{\textscy}{\textsubring{r}}} \dicPos{m}[13] \dicFlx{(‑manns, ‑menn)}[2] \dicDirectTranslationCS{podněcovatel(ka), hybatel(ka), iniciátor(ka)}
\dicEntry[hvatar] \dicTerm{hvatar} \dicIPA{{k\smash{\textsuperscript{h}}}{v}{a}{\textlengthmark}{\textsubring{d}}{a}{\textsubring{r}}} \dicPos{f} \dicFlx{sg gen} \dicLink{hvöt}
\dicEntry[hvati] \dicTerm{hvat|i} \dicIPA{{k\smash{\textsuperscript{h}}}{v}{a}{\textlengthmark}{\textsubring{d}}{\textsci}} \dicPos{m}[1] \dicFlx{(‑a, ‑ar)}[8] \textbf{1.} \dicSynonym{hvatning} \dicDirectTranslationCS{podnět, stimul, hnací síla}  \textbf{2.} \dicFieldCat{chem.} \dicDirectTranslationCS{katalyzátor}
\dicEntry[hvatir] \dicTerm{hvatir} \dicIPA{{k\smash{\textsuperscript{h}}}{v}{a}{\textlengthmark}{\textsubring{d}}{\textsci}{\textsubring{r}}} \dicPos{f} \dicFlx{pl nom} \dicLink{hvöt}
\dicEntry[hvatlegur] \dicTerm{hvat··legur} \dicIPA{{k\smash{\textsuperscript{h}}}{v}{a}{\textlengthmark}{\textsubring{d}}{l}{\textepsilon}{\textbabygamma}{\textscy}{\textsubring{r}}} \dicPos{adj}[1]\dicFlx{}[-8] \dicSynonym{fjörlegur} \dicDirectTranslationCS{energický, rázný}
\dicEntry[hvatning] \dicTerm{hvat··ning} \dicIPA{{k\smash{\textsuperscript{h}}}{v}{a}{h}{\textsubring{d}}{n}{i}{\ng}{\r{g}}} \dicPos{f}[4] \dicFlx{(‑ar, ‑ar)}[5] \dicDirectTranslationCS{pobídka, pobídnutí, impuls, výzva} \dicExampleIS{fá hvatningu frá e‑m} \dicExampleCS{dostat impuls od (koho)}
\dicEntry[hvatt] \dicTerm{hvatt} \dicIPA{{k\smash{\textsuperscript{h}}}{v}{a}{h}{\textsubring{d}}} \dicPos{v} \dicFlx{supin} \dicLink{hvetja}
\dicEntry[hvatti] \dicTerm{hvatti} \dicIPA{{k\smash{\textsuperscript{h}}}{v}{a}{h}{\textsubring{d}}{\textsci}} \dicPos{v} \dicFlx{ind pf sg 1 pers} \dicLink{hvetja}
\dicEntry[hvatur] \dicTerm{hvatur} \dicIPA{{k\smash{\textsuperscript{h}}}{v}{a}{\textlengthmark}{\textsubring{d}}{\textscy}{\textsubring{r}}} \dicPos{adj}[1] \dicFlx{(f hvöt)}[2] \dicSynonym{skjótur} \dicDirectTranslationCS{čilý, svižný};  \dicPhraseIS{hvatur í spori} \dicDirectTranslationCS{(jsoucí) se svižným krokem};  \dicPhraseIS{hvatur köttur} \dicSynonym{fress} \dicDirectTranslationCS{kocour}
\dicEntry[hvatvís] \dicTerm{hvat··vís} \dicIPA{{k\smash{\textsuperscript{h}}}{v}{a}{\textlengthmark}{\textsubring{d}}{v}{i}{s}} \dicPos{adj}[5]\dicFlx{}[-1] \dicDirectTranslationCS{zbrklý, impulzivní, spontánní}
\dicEntry[hvatvísi] \dicTerm{hvat··vís|i} \dicIPA{{k\smash{\textsuperscript{h}}}{v}{a}{\textlengthmark}{\textsubring{d}}{v}{i}{s}{\textsci}} \dicPos{f}[3] \dicFlx{(‑i)}[3] \dicDirectTranslationCS{zbrklost, impulzivnost, spontánnost}
\dicEntry[hvá] \dicTerm{hvá} \dicIPA{{k\smash{\textsuperscript{h}}}{v}{au}{\textlengthmark}} \dicPos{v}[2] \dicFlx{(‑ði, ‑ð)}[120] \dicDirectTranslationCS{přeptat se (při přeslechnutí)};  \dicPhraseIS{hvá eftir e‑u} \dicDirectTranslationCS{přeptat se na (co)}
\dicEntry[hve] \dicTerm{hve} \dicsymFrequent\  \dicIPA{{k\smash{\textsuperscript{h}}}{v}{\textepsilon}{\textlengthmark}} \dicPos{adv} \dicSynonym{hversu} \dicDirectTranslationCS{jak, jako} \dicExampleIS{hve mikið} \dicExampleCS{jak moc}
\dicEntry[hvein] \dicTerm{hvein} \dicIPA{{k\smash{\textsuperscript{h}}}{v}{ei}{\textlengthmark}{\textsubring{n}}} \dicPos{v} \dicFlx{ind pf sg 1 pers} \dicLink{hvína}
\dicEntry[hveiti] \dicTerm{hveiti} \dicIPA{{k\smash{\textsuperscript{h}}}{v}{ei}{\textlengthmark}{\textsubring{d}}{\textsci}} \dicPos{n}[2] \dicFlx{(‑s)}[20] \textbf{1.} \dicDirectTranslationCS{mouka}  \textbf{2.} \dicFieldCat{bot.} \dicDirectTranslationCS{pšenice, pšenice setá} \textit{(l.~{\textLA{Triticum aestivum}})}  \dicsymPhoto\ 
\dicFigure{ds_image_hveiti_0_1.jpg}{Hveiti}{Hveiti - Hans Hillewaert, CC BY-SA 3.0}
\dicEntry[hveitibrauð] \dicTerm{hveiti··brauð} \dicIPA{{k\smash{\textsuperscript{h}}}{v}{ei}{\textlengthmark}{\textsubring{d}}{\textsci}{\textsubring{b}}{r}{\oe i}{\texttheta}} \dicPos{n}[2] \dicFlx{(‑s, ‑)}[5] \dicDirectTranslationCS{pšeničný chléb}
\begin{xtolerant}{}{1pt}
\dicEntry[hveitibrauðsdagar] \dicTerm{hveiti·brauðs··dagar}\addthinSS\dicIPA{{k\smash{\textsuperscript{h}}}{v}{ei}{\textlengthmark}{\textsubring{d}}{\textsci}{\textsubring{b}}{r}{\oe i}{ð}{s}{\textsubring{d}}{a}{\textbabygamma}{a}{\textsubring{r}}}\addthinS\dicPos{m}[6]\addthinS\dicFlx{pl}[60]\addthinS\dicDirectTranslationCS{líbán\-ky}
\end{xtolerant}
\dicEntry[hveitimjöl] \dicTerm{hveiti··mjöl} \dicIPA{{k\smash{\textsuperscript{h}}}{v}{ei}{\textlengthmark}{\textsubring{d}}{\textsci}{m}{j}{\oe}{\textsubring{l}}} \dicPos{n}[2] \dicFlx{(‑s)}[2] \dicDirectTranslationCS{pšeničná mouka}
\dicEntry[hveitipípa] \dicTerm{hveiti··píp|a} \dicIPA{{k\smash{\textsuperscript{h}}}{v}{ei}{\textlengthmark}{\textsubring{d}}{\textsci}{p\smash{\textsuperscript{h}}}{i}{\textsubring{b}}{a}} \dicPos{f}[1] \dicFlx{(‑u, ‑ur)}[19] \dicSynonym{makkaróna} \dicDirectTranslationCS{makarón}
\dicEntry[hvekkja] \dicTerm{hvekk|ja} \dicIPA{{k\smash{\textsuperscript{h}}}{v}{\textepsilon}{h}{\r{\textObardotlessj}}{a}} \dicPos{v}[2] \dicFlx{(‑ti, ‑t)}[25] \dicFlx{acc} \dicDirectTranslationCS{štvát, vytáčet, lézt na nervy};  \dicIdiom{hvekkjast}[á]{ \dicPhraseIS{hvekkjast á e‑u}} \dicFlx{refl} \dicDirectTranslationCS{mít dost (čeho), být naštvaný kvůli (čemu)}
\dicEntry[hvel] \dicTerm{hvel} \dicIPA{{k\smash{\textsuperscript{h}}}{v}{\textepsilon}{\textlengthmark}{\textsubring{l}}} \dicPos{n}[2] \dicFlx{(‑s, ‑)}[5] \textbf{1.} \dicSynonym*{jarðarhvel} \dicDirectTranslationCS{polokoule}  \textbf{2.} \dicFieldCat{anat.} \dicDirectTranslationCS{hemisféra}
\dicEntry[hvelaheili] \dicTerm{hvela··heil|i} \dicIPA{{k\smash{\textsuperscript{h}}}{v}{\textepsilon}{\textlengthmark}{l}{a}{h}{ei}{l}{\textsci}} \dicPos{m}[1] \dicFlx{(‑a, ‑ar)}[1] \dicFieldCat{anat.} \dicDirectTranslationCS{velký mozek, cerebrum}
\dicEntry[hvelfing] \dicTerm{hvelf··ing} \dicIPA{{k\smash{\textsuperscript{h}}}{v}{\textepsilon}{l}{v}{i}{\ng}{\r{g}}} \dicPos{f}[4] \dicFlx{(‑ar, ‑ar)}[5] \dicDirectTranslationCS{klenba, klenutí} \dicExampleIS{hvelfing kirkjunnar} \dicExampleCS{kostelní klenba}
\dicEntry[hvelja] \dicTerm{hvelj|a} \dicIPA{{k\smash{\textsuperscript{h}}}{v}{\textepsilon}{l}{j}{a}} \dicPos{f}[1] \dicFlx{(‑u, ‑ur)}[7] \textbf{1.} \dicFieldCat{zool.} \dicDirectTranslationCS{polypovec} \textit{(l.~{\textLA{Hydrozoa}})}  \dicsymPhoto\   \textbf{2.} \dicSynonym*{hvalshúð} \dicDirectTranslationCS{velrybí kůže}
\dicFigure{43488.jpg}{Hvelja}{Hvelja - Dr. Anthony R. Picciolo, NOAA NODC, Biolib, PD}
\begin{xtolerant}{}{1pt}
\dicEntry[hvellur] \dicTerm{hvell|ur} \dicIPA{{k\smash{\textsuperscript{h}}}{v}{\textepsilon}{\textsubring{d}}{l}{\textscy}{\textsubring{r}}} \dicPos{m}[9] \dicFlx{(‑s, ‑ir)}[4] \textbf{1.} \dicSynonym{smellur} \dicDirectTranslationCS{rána, třesknutí, zaburácení}  \textbf{2.} \dicSynonym{óveður} \dicDirectTranslationCS{nečas, slota}  \textbf{3.} \dicSynonym*{uppistand} \dicDirectTranslationCS{pozdvižení};  \dicIdiom{hvellur}{ \dicPhraseIS{í\addthinS(einum\,/\addthinSS grænum\,/\addthinSS hvínandi\,/\addthinSS logandi)\addthinS hvelli}} \dicFlx{adv} \dicSynonym{undireins} \dicDirectTranslationCS{okamžitě, ihned}
\end{xtolerant}
\dicEntry[hvenær] \dicTerm{hvenær} \dicsymFrequent\  \dicIPA{{k\smash{\textsuperscript{h}}}{v}{\textepsilon}{\textlengthmark}{n}{a}{i}{\textsubring{r}}} \dicPos{adv} \textbf{1.} \dicDirectTranslationCS{kdy} \dicIndirectTranslationCS{(užívané v~přímé otázce)} \dicExampleIS{Hvenær á ég að byrja?} \dicExampleCS{Kdy mám začít?}  \textbf{2.} \dicDirectTranslationCS{kdy} \dicIndirectTranslationCS{(užívané ve vedlejší větě)} \dicExampleIS{Við vitum ekki hvenær þau koma aftur.} \dicExampleCS{Nevíme, kdy se vrátí.};  \dicPhraseIS{hvenær sem} \dicFlx{conj} \dicDirectTranslationCS{kdykoli(v)};  \dicPhraseIS{hvenær sem er} \dicFlx{adv} \dicDirectTranslationCS{kdykoli(v)}
\dicEntry[hver] \dicTerm{hver\smash{\textsuperscript{1}}} \dicIPA{{k\smash{\textsuperscript{h}}}{v}{\textepsilon}{\textlengthmark}{\textsubring{r}}} \dicPos{m}[6] \dicFlx{(‑s, ‑ir\,/\addthin ‑ar)}[70] \dicFieldCat{geol.} \dicDirectTranslationCS{gejzír, vřídlo, horký pramen}
\dicEntry[hver] \dicTerm{hver\smash{\textsuperscript{2}}} \dicsymFrequent\  \dicIPA{{k\smash{\textsuperscript{h}}}{v}{\textepsilon}{\textlengthmark}{\textsubring{r}}} \dicPos{pron}[1] \dicFlx{int} \textbf{1.} \dicDirectTranslationCS{kdo} \dicIndirectTranslationCS{(užívané jako zájmeno tázací)} \dicExampleIS{Hver er hann?} \dicExampleCS{Kdo to je?}  \textbf{2.} \dicDirectTranslationCS{který} \dicExampleIS{hver systranna} \dicExampleCS{která ze sester}  \textbf{3.} \dicDirectTranslationCS{jaký} \dicIndirectTranslationCS{(stojící samostatně)} \dicExampleIS{hver er munurinn} \dicExampleCS{jaký je rozdíl}  \textbf{4.} \dicDirectTranslationCS{jaký} \dicIndirectTranslationCS{(stojící jako atribut)} \dicExampleIS{hver áhrif} \dicExampleCS{jaký vliv};  \dicPhraseIS{hver eðlis} \dicDirectTranslationCS{jaké povahy};  \dicPhraseIS{hver lags} \dicDirectTranslationCS{jakého typu, jaké povahy}  \textbf{5.} \dicDirectTranslationCS{jaký} \dicIndirectTranslationCS{(užívané při zvolání s~vulgárním výrazem)} \dicExampleIS{Hver andskotinn gengur þarna?} \dicExampleCS{Kdo to tam sakra jde?}  \textbf{6.} \dicSynonym{sérhver} \dicDirectTranslationCS{každý} \dicIndirectTranslationCS{(užívané jako zájmeno neurčité)}  \textbf{7.} \dicDirectTranslationCS{každý (druhý den ap.)} \dicIndirectTranslationCS{(užívané s~číslovkami řadovými)} \dicExampleIS{annan hvern dag} \dicExampleCS{každý druhý den}  \textbf{8.} \dicDirectTranslationCS{navzájem} \dicIndirectTranslationCS{(užívané jako zájmeno reciproční)} \dicExampleIS{Þeir hjálpuðu hver öðrum.} \dicExampleCS{Pomáhali si navzájem.}  \textbf{9.} \dicDirectTranslationCS{každý (svůj míč ap.)} \dicIndirectTranslationCS{(užívané se vztažným přivlastňovacím zájmenem sinn)} \dicExampleIS{Þeir tóku hver sinn bolta.} \dicExampleCS{Každý si vzal svůj míč.};  \dicIdiom{hver}[annar]{ \dicPhraseIS{hver annar}} \dicFlx{pron} {\textbf{a.}} \dicDirectTranslationCS{(sebe) navzájem} \dicExampleIS{Þeir leiðast til að týna ekki hver öðrum.} \dicExampleCS{Drží se za ruce, aby se navzájem neztratili.};  {\textbf{b.}} \dicDirectTranslationCS{sebe} \dicExampleIS{lita hver á annan} \dicExampleCS{dívat se na sebe};  {\textbf{c.}} \dicDirectTranslationCS{každý jiný} \dicExampleIS{eins og hver annar} \dicExampleCS{jako každý jiný};  {\textbf{d.}} \dicDirectTranslationCS{jeden -- než druhý} \dicExampleIS{hver tölvan annarri betri} \dicExampleCS{jeden počítač lepší než druhý};  \dicIdiom{hver}[og]{ \dicPhraseIS{hver og einn}} \dicSynonym{allir} \dicDirectTranslationCS{všichni, každý};  \dicIdiom{hver}[sem]{ \dicPhraseIS{hver sem er}} \dicFlx{pron} \dicDirectTranslationCS{kdokoli(v)};  \dicIdiom{hver}[sinn]{ \dicPhraseIS{hver sinn}} \dicFlx{pron} \dicDirectTranslationCS{každý svůj} \dicExampleIS{hver með sínum hætti} \dicExampleCS{každý svým způsobem};  \dicIdiom{hver}[um]{ \dicPhraseIS{hver um sig}} \dicDirectTranslationCS{každý (sám o~sobě)};  \dicIdiom{hver}[öðrum]{ \dicPhraseIS{hver af öðrum}} \dicFlx{adv} \dicDirectTranslationCS{jeden po druhém};  \dicPhraseIS{hver (á) eftir öðrum} \dicFlx{adv} \dicDirectTranslationCS{jeden po druhém};  \dicPhraseIS{hver á fætur öðrum} \dicFlx{adv} \dicDirectTranslationCS{jeden za druhým, jeden po druhém};  \dicIdiom{sinn}[hver]{ \dicPhraseIS{sinn hver}} \dicFlx{pron} \dicDirectTranslationCS{každý jiný};  \dicPhraseIS{sitt hvað} {\textbf{a.}} \dicDirectTranslationCS{trochu od všeho, různé věci};  {\textbf{b.}} \dicDirectTranslationCS{odlišné věci} \dicExampleIS{Það er sitt hvað, e‑að og e‑að.} \dicExampleCS{(co) a~(co) jsou dvě odlišné věci.};  \dicPhraseIS{sitt (lítið) af hverju} \dicDirectTranslationCS{trochu od všeho, různé věci};  \dicPhraseIS{sitt af hverju tagi} \dicDirectTranslationCS{trochu od všeho, různé věci};  \dicPhraseIS{sitt sýnist hverjum um e‑ð} \dicFlx{refl} \dicDirectTranslationCS{existují různé názory na (co)}
\dicEntry[hverasvæði] \dicTerm{hvera··svæði} \dicIPA{{k\smash{\textsuperscript{h}}}{v}{\textepsilon}{\textlengthmark}{r}{a}{s}{v}{a}{i}{ð}{\textsci}} \dicPos{n}[2] \dicFlx{(‑s, ‑)}[14] \dicDirectTranslationCS{geotermální oblast}
\dicEntry[hverf] \dicTerm{hverf} \dicIPA{{k\smash{\textsuperscript{h}}}{v}{\textepsilon}{r}{f}} \dicPos{v} \dicFlx{ind praes sg 1 pers} \dicLink{hverfa\smash{\textsuperscript{2}}}
\dicEntry[hverfa] \dicTerm{hverf|a\smash{\textsuperscript{1}}} \dicsymFrequent\  \dicIPA{{k\smash{\textsuperscript{h}}}{v}{\textepsilon}{r}{v}{a}} \dicPos{v}[2] \dicFlx{(‑ði, ‑t)}[99] \textbf{1.} \dicSynonym*{láta snúast} \dicDirectTranslationCS{(o)točit, otáčet} \dicExampleIS{hverfa e‑u} \dicExampleCS{otáčet (co)}  \textbf{2.} \dicSynonym{breyta\smash{\textsuperscript{2}}} \dicDirectTranslationCS{(z)měnit} \dicExampleIS{hverfa e‑u í e‑ð} \dicExampleCS{změnit (co) v~(co)};  \dicPhraseIS{hverfa huga e‑rs} \dicDirectTranslationCS{změnit (čí) názor};  \dicIdiom{hverfast}{ \dicPhraseIS{hverfast}} \dicFlx{refl} \dicSynonym*{snúast} \dicDirectTranslationCS{točit se, týkat se} \dicExampleIS{Allt málið hverfist um einn depil.} \dicExampleCS{Všechno se točí kolem jednoho bodu.}
\dicEntry[hverfa] \dicTerm{hverfa\smash{\textsuperscript{2}}} \dicIPA{{k\smash{\textsuperscript{h}}}{v}{\textepsilon}{r}{v}{a}} \dicPos{v}[6] \dicFlx{(hverf, hvarf, hurfum, hyrfi, horfið)}[35] \dicDirectTranslationCS{(z)mizet, ztratit se, ztrácet se, vytratit se, vytrácet se};  \dicIdiom{hverfa}[að]{ \dicPhraseIS{hverfa að e‑u}} \dicDirectTranslationCS{uchýlit se k~(čemu)};  \dicIdiom{hverfa}[af]{ \dicPhraseIS{hverfa af sjónarsviðinu}} \dicDirectTranslationCS{vyhynout, vymizet};  \dicIdiom{hverfa}[frá]{ \dicPhraseIS{hverfa frá e‑u}} \dicDirectTranslationCS{upustit od (čeho), odklonit se od (čeho)} \dicExampleIS{hverfa frá þessari fyrirætlun} \dicExampleCS{upustit od toho záměru};  \dicIdiom{hverfa}[til]{ \dicPhraseIS{hverfa til e‑s}} \dicDirectTranslationCS{vrátit\,/\addthin vracet se k~(čemu) (ke starému způsobu ap.)};  \dicIdiom{hverfa}[úr]{ \dicPhraseIS{hverfa úr augsýn}} \dicDirectTranslationCS{zmizet z~dohledu};  \dicIdiom{hverfa}{ \dicPhraseIS{hverfa eins og dögg fyrir sólu, hverfa eins og hendi væri veifað}} \dicLangCat{přen.} \dicDirectTranslationCS{zmizet jako mávnutím (kouzelného) proutku}
\dicEntry[hverfi] \dicTerm{hverfi} \dicsymFrequent\  \dicIPA{{k\smash{\textsuperscript{h}}}{v}{\textepsilon}{r}{v}{\textsci}} \dicPos{n}[2] \dicFlx{(‑s, ‑)}[14] \dicDirectTranslationCS{čtvrť, obvod} \dicExampleIS{búa í rólegu hverfi} \dicExampleCS{bydlet v~klidné čtvrti}
\dicEntry[hverfisskóli] \dicTerm{hverfis··skól|i} \dicIPA{{k\smash{\textsuperscript{h}}}{v}{\textepsilon}{r}{v}{\textsci}{s}{\r{g}}{ou}{l}{\textsci}} \dicPos{m}[1] \dicFlx{(‑a, ‑ar)}[1] \dicFieldCat{škol.} \dicDirectTranslationCS{spádová škola}
\dicEntry[hverfisteinn] \dicTerm{hverfi··stein|n} \dicIPA{{k\smash{\textsuperscript{h}}}{v}{\textepsilon}{r}{v}{\textsci}{s}{\textsubring{d}}{ei}{\textsubring{d}}{\textsubring{n}}} \dicPos{m}[6] \dicFlx{(‑s, ‑ar)}[42] \dicDirectTranslationCS{brusný kotouč}
\dicEntry[hverflyndur] \dicTerm{hverf··lyndur} \dicIPA{{k\smash{\textsuperscript{h}}}{v}{\textepsilon}{r}{v}{l}{\textsci}{n}{\textsubring{d}}{\textscy}{\textsubring{r}}} \dicPos{adj}[2]\dicFlx{}[-14] \dicDirectTranslationCS{nestálý, vrtkavý} \dicExampleIS{hverflyndur í ástum} \dicExampleCS{nestálý v~lásce}
\dicEntry[hverfull] \dicTerm{hverfull} \dicIPA{{k\smash{\textsuperscript{h}}}{v}{\textepsilon}{r}{v}{\textscy}{\textsubring{d}}{\textsubring{l}}} \dicPos{adj}[8]\dicFlx{}[-4] \dicDirectTranslationCS{pomíjivý, prchavý} \dicExampleIS{hverfult líf} \dicExampleCS{pomíjivý život}
\dicEntry[hverfulleiki] \dicTerm{hverful··leik|i} \dicIPA{{k\smash{\textsuperscript{h}}}{v}{\textepsilon}{r}{v}{\textscy}{l}{ei}{\r{\textObardotlessj}}{\textsci}} \dicPos{m}[1] \dicFlx{(‑a)}[3] \dicDirectTranslationCS{pomíjivost, prchavost}
\dicEntry[hvergi] \dicTerm{hvergi} \dicsymFrequent\  \dicIPA{{k\smash{\textsuperscript{h}}}{v}{\textepsilon}{r}{\r{\textObardotlessj}}{\textsci}} \dicPos{adv} \dicSynonym*{á engum stað} \dicDirectTranslationCS{nikde} \dicExampleIS{Tónlist heyrist hvergi.} \dicExampleCS{Nikde není slyšet hudba.} \dicAntonym{allsstaðar};  \dicPhraseIS{hvergi nærri} \dicFlx{adv} \dicDirectTranslationCS{ani zdaleka, vůbec ne}
\dicEntry[hvernig] \dicTerm{hvernig\smash{\textsuperscript{1}}} \dicIPA{{k\smash{\textsuperscript{h}}}{v}{\textepsilon}{r}{\textsubring{d}}{n}{\textsci}{x}} \dicPos{adj}[13] \dicFlx{indecl}[1] \textbf{1.} \dicDirectTranslationCS{jaký, jakého druhu} \dicIndirectTranslationCS{(užívané v~přímé otázce)} \dicExampleIS{Hvernig er maturinn?} \dicExampleCS{Jaké je jídlo?}  \textbf{2.} \dicDirectTranslationCS{jaký, jakého druhu} \dicIndirectTranslationCS{(užívané ve vedlejší větě)} \dicExampleIS{Flestum er ljóst hvernig staðan er.} \dicExampleCS{Většině je jasné, jaká je situace.}
\dicEntry[hvernig] \dicTerm{hvernig\smash{\textsuperscript{2}}} \dicsymFrequent\  \dicIPA{{k\smash{\textsuperscript{h}}}{v}{\textepsilon}{r}{\textsubring{d}}{n}{\textsci}{x}} \dicPos{adv} \textbf{1.} \dicDirectTranslationCS{jak, jakpak} \dicIndirectTranslationCS{(užívané v~přímé otázce)} \dicExampleIS{Hvernig kynntust þið?} \dicExampleCS{Jak jste se poznali?}  \textbf{2.} \dicDirectTranslationCS{jak, jakpak} \dicIndirectTranslationCS{(užívané ve vedlejší větě)} \dicExampleIS{Þau vita hvernig myndin endar.} \dicExampleCS{Vědí, jak film dopadne.};  \dicPhraseIS{hvernig sem} \dicFlx{conj} \dicDirectTranslationCS{jakkoli(v)}
\dicEntry[hvers kyns] \dicTerm{hvers kyns}\dicTerm{, hverskyns} \dicIPA{{k\smash{\textsuperscript{h}}}\-{v}\-{\textepsilon}\-{\textsubring{r}}\-{s}\-{\r{\textObardotlessj}}\-{\textsci}\-{n}\-{s}\-} \dicPos{adj}[13] \dicFlx{indecl}[1] \textbf{1.} \dicDirectTranslationCS{jakého druhu, jaký}  \textbf{2.} \dicSynonym*{alls konar} \dicDirectTranslationCS{všelijaký, všemožný}  \textbf{3.} \dicDirectTranslationCS{jaký} \dicIndirectTranslationCS{(vyjadřuje údiv)}
\dicEntry[hversdags] \dicTerm{hvers··dags\smash{\textsuperscript{1}}} \dicIPA{{k\smash{\textsuperscript{h}}}{v}{\textepsilon}{\textsubring{r}}{s}{\textsubring{d}}{a}{x}{s}} \dicPos{adv} \dicDirectTranslationCS{každý den, každodenně}
\dicEntry[hversdags] \dicTerm{hvers··dags-\smash{\textsuperscript{2}}} \dicIPA{{k\smash{\textsuperscript{h}}}{v}{\textepsilon}{\textsubring{r}}{s}{\textsubring{d}}{a}{x}{s}} \dicPos{predp} \dicDirectTranslationCS{všední, každodenní} \dicAntonym{spari\smash{\textsuperscript{2}}}
\dicEntry[hversdagslegur] \dicTerm{hvers·dags··legur} \dicsymFrequent\  \dicIPA{{k\smash{\textsuperscript{h}}}{v}{\textepsilon}{\textsubring{r}}{s}{\textsubring{d}}{a}{x}{s}{l}{\textepsilon}{\textbabygamma}{\textscy}{\textsubring{r}}} \dicPos{adj}[1]\dicFlx{}[-8] \dicDirectTranslationCS{každodenní, všední} \dicExampleIS{hversdagslegt líf} \dicExampleCS{každodenní život}
\dicEntry[hversdagsleiki] \dicTerm{hvers·dags··leik|i} \dicIPA{{k\smash{\textsuperscript{h}}}{v}{\textepsilon}{\textsubring{r}}{s}{\textsubring{d}}{a}{x}{s}{l}{ei}{\r{\textObardotlessj}}{\textsci}} \dicPos{m}[1] \dicFlx{(‑a)}[3] \dicDirectTranslationCS{každodennost, všednost}
\dicEntry[hverskonar] \dicTerm{hvers··konar} \dicIPA{{k\smash{\textsuperscript{h}}}{v}{\textepsilon}{\textsubring{r}}{s}{k\smash{\textsuperscript{h}}}{\textopeno}{n}{a}{\textsubring{r}}} \dicPos{adj}[4]\dicFlx{}[-5] \dicLink{konar}
\dicEntry[hverskyns] \dicTerm{hvers··kyns} \dicIPA{{k\smash{\textsuperscript{h}}}{v}{\textepsilon}{\textsubring{r}}{s}{\r{\textObardotlessj}}{\textsci}{n}{s}} \dicPos{adj}[13] \dicFlx{indecl}[1] \dicLink{hvers kyns}
\dicEntry[hversu] \dicTerm{hversu} \dicsymFrequent\  \dicIPA{{k\smash{\textsuperscript{h}}}{v}{\textepsilon}{\textsubring{r}}{s}{\textscy}} \dicPos{adv} \dicSynonym{hve} \dicDirectTranslationCS{jak} \dicExampleIS{hversu mikið} \dicExampleCS{jak moc}
\dicEntry[hvert] \dicTerm{hvert\smash{\textsuperscript{1}}} \dicIPA{{k\smash{\textsuperscript{h}}}{v}{\textepsilon}{\textsubring{r}}{\textsubring{d}}} \dicPos{pron}[1] \dicFlx{n} \dicLink{hver\smash{\textsuperscript{3}}}
\dicEntry[hvert] \dicTerm{hvert\smash{\textsuperscript{2}}} \dicsymFrequent\  \dicIPA{{k\smash{\textsuperscript{h}}}{v}{\textepsilon}{\textsubring{r}}{\textsubring{d}}} \dicPos{adv} \textbf{1.} \dicDirectTranslationCS{kam} \dicIndirectTranslationCS{(užívané v~přímé otázce)} \dicExampleIS{Hvert ertu að fara?} \dicExampleCS{Kam jdeš?}  \textbf{2.} \dicDirectTranslationCS{kam} \dicIndirectTranslationCS{(užívané ve vedlejší větě)};  \dicPhraseIS{hvert sem} \dicFlx{conj} \dicDirectTranslationCS{kamkoli(v)};  \dicPhraseIS{hvert sem er} \dicFlx{adv} \dicDirectTranslationCS{kamkoli(v)}
\dicEntry[hvessa] \dicTerm{hvess|a} \dicsymFrequent\  \dicIPA{{k\smash{\textsuperscript{h}}}{v}{\textepsilon}{s}{\textlengthmark}{a}} \dicPos{v}[2] \dicFlx{(‑ti, ‑t)}[67] \dicFlx{acc} \textbf{1.} \dicSynonym{brýna} \dicDirectTranslationCS{(na)ostřit, (na)brousit} \dicExampleIS{hvessa hnífinn} \dicExampleCS{naostřit nůž}  \textbf{2.} \dicPhraseIS{það\,/\addthin hann hvessir} \dicFlx{impers} \dicDirectTranslationCS{zvedá\,/\addthin žene se vítr}  \textbf{3.} \dicPhraseIS{hvessa augun á e‑n\,/\addthin e‑ð} \dicDirectTranslationCS{zaostřit na (koho\,/\addthin co)} \dicExampleIS{hvessa augun á myndina} \dicExampleCS{zaostřit na obrázek}
\dicEntry[hvet] \dicTerm{hvet} \dicIPA{{k\smash{\textsuperscript{h}}}{v}{\textepsilon}{\textlengthmark}{\textsubring{d}}} \dicPos{v} \dicFlx{ind praes sg 1 pers} \dicLink{hvetja}
\dicEntry[hvetja] \dicTerm{hvetja} \dicsymFrequent\  \dicIPA{{k\smash{\textsuperscript{h}}}{v}{\textepsilon}{\textlengthmark}{\textsubring{d}}{j}{a}} \dicPos{v}[4] \dicFlx{(hvet, hvatti, hvöttum, hvetti, hvatt)}[49] \dicFlx{acc} \textbf{1.} \dicSynonym{örva} \dicDirectTranslationCS{pobízet, pobídnout, vybízet, vybídnout, vyzvat, vyzývat, (vy)burcovat} \dicExampleIS{hvetja e‑n til aðgerða} \dicExampleCS{pobídnout (koho) k~akci}  \textbf{2.} \dicSynonym{brýna} \dicDirectTranslationCS{(na)ostřit, (na)brousit} \dicExampleIS{hvetja hnífinn} \dicExampleCS{naostřit nůž}
\dicEntry[hvetjandi] \dicTerm{hvetj··andi} \dicIPA{{k\smash{\textsuperscript{h}}}{v}{\textepsilon}{\textlengthmark}{\textsubring{d}}{j}{a}{n}{\textsubring{d}}{\textsci}} \dicPos{adj}[13] \dicFlx{indecl}[1] \dicDirectTranslationCS{povzbuzující, podněcující, stimulující} \dicExampleIS{hvetjandi áhrif} \dicExampleCS{povzbuzující vliv}
\dicEntry[hvetti] \dicTerm{hvetti} \dicIPA{{k\smash{\textsuperscript{h}}}{v}{\textepsilon}{h}{\textsubring{d}}{\textsci}} \dicPos{v} \dicFlx{con pf sg 1 pers} \dicLink{hvetja}
\dicEntry[hviða] \dicTerm{hvið|a} \dicIPA{{k\smash{\textsuperscript{h}}}{v}{\textsci}{\textlengthmark}{ð}{a}} \dicPos{f}[1] \dicFlx{(‑u, ‑ur)}[19] \dicDirectTranslationCS{poryv (větru), závan, zadutí}
\dicEntry[hvika] \dicTerm{hvik|a} \dicIPA{{k\smash{\textsuperscript{h}}}{v}{\textsci}{\textlengthmark}{\r{g}}{a}} \dicPos{v}[1] \dicFlx{(‑aði)}[44] \dicSynonym{víkja} \dicDirectTranslationCS{ustoupit, ustupovat, couvnout, couvat} \dicIndirectTranslationCS{(nejčastěji se záporem)} \dicExampleIS{hvika ekki frá ákvörðun sinni} \dicExampleCS{neustoupit ze svého rozhodnutí}
\dicEntry[hviklyndur] \dicTerm{hvik··lyndur} \dicIPA{{k\smash{\textsuperscript{h}}}{v}{\textsci}{\textlengthmark}{\r{g}}{l}{\textsci}{n}{\textsubring{d}}{\textscy}{\textsubring{r}}} \dicPos{adj}[2]\dicFlx{}[-14] \dicDirectTranslationCS{nestálý, nestabilní (v~lásce ap.)}
\dicEntry[hviksaga] \dicTerm{hvik··|saga} \dicIPA{{k\smash{\textsuperscript{h}}}{v}{\textsci}{\textlengthmark}{\r{g}}{s}{a}{\textbabygamma}{a}} \dicPos{f}[1] \dicFlx{(‑sögu, ‑sögur)}[14] \dicDirectTranslationCS{fáma, pověst}
\dicEntry[hvikull] \dicTerm{hvikull} \dicIPA{{k\smash{\textsuperscript{h}}}{v}{\textsci}{\textlengthmark}{\r{g}}{\textscy}{\textsubring{d}}{\textsubring{l}}} \dicPos{adj}[8]\dicFlx{}[-4] \dicDirectTranslationCS{nestálý, proměnlivý, prchavý}
\dicEntry[hvilft] \dicTerm{hvilft} \dicIPA{{k\smash{\textsuperscript{h}}}{v}{\textsci}{\textsubring{l}}{\textsubring{d}}} \dicPos{f}[7] \dicFlx{(‑ar, ‑ir)}[1] \dicDirectTranslationCS{dolík, proláklina}
\dicEntry[hvimleiður] \dicTerm{hvim··leiður} \dicIPA{{k\smash{\textsuperscript{h}}}{v}{\textsci}{m}{l}{ei}{ð}{\textscy}{\textsubring{r}}} \dicPos{adj}[2]\dicFlx{}[-6] \dicSynonym{ógeðfelldur} \dicDirectTranslationCS{nepříjemný, obtěžující, iritující}
\dicEntry[hvini] \dicTerm{hvini} \dicIPA{{k\smash{\textsuperscript{h}}}{v}{\textsci}{\textlengthmark}{n}{\textsci}} \dicPos{v} \dicFlx{con pf sg 1 pers} \dicLink{hvína}
\dicEntry[hvinið] \dicTerm{hvinið} \dicIPA{{k\smash{\textsuperscript{h}}}{v}{\textsci}{\textlengthmark}{n}{\textsci}{\texttheta}} \dicPos{v} \dicFlx{supin} \dicLink{hvína}
\dicEntry[hvinum] \dicTerm{hvinum} \dicIPA{{k\smash{\textsuperscript{h}}}{v}{\textsci}{\textlengthmark}{n}{\textscy}{\textsubring{m}}} \dicPos{v} \dicFlx{ind pf pl 1 pers} \dicLink{hvína}
\dicEntry[hvinur] \dicTerm{hvin|ur} \dicIPA{{k\smash{\textsuperscript{h}}}{v}{\textsci}{\textlengthmark}{n}{\textscy}{\textsubring{r}}} \dicPos{m}[6] \dicFlx{(‑s)}[21] \dicDirectTranslationCS{hvizd, hvízdot, hvízdání} \dicExampleIS{hvinur stormsins} \dicExampleCS{hvízdání bouře}
\dicEntry[hvinönd] \dicTerm{hvin··|önd} \dicIPA{{k\smash{\textsuperscript{h}}}{v}{\textsci}{\textlengthmark}{n}{\oe}{n}{\textsubring{d}}} \dicPos{f}[8] \dicFlx{(‑andar, ‑endur)}[2] \dicFieldCat{zool.} \dicDirectTranslationCS{hohol severní} \textit{(l.~{\textLA{Bucephala clangula}})}  \dicsymPhoto\ 
\dicFigure{20962.jpg}{Hvinönd}{Hvinönd - Fenner, M., Biolib, PD}
\dicEntry[hvirfilbein] \dicTerm{hvirfil··bein} \dicIPA{{k\smash{\textsuperscript{h}}}{v}{\textsci}{r}{v}{\textsci}{l}{\textsubring{b}}{ei}{\textsubring{n}}} \dicPos{n}[2] \dicFlx{(‑s, ‑)}[5] \dicFieldCat{anat.} \dicDirectTranslationCS{temenní kost}
\dicEntry[hvirfilbylur] \dicTerm{hvirfil··byl|ur} \dicIPA{{k\smash{\textsuperscript{h}}}{v}{\textsci}{r}{v}{\textsci}{l}{\textsubring{b}}{\textsci}{l}{\textscy}{\textsubring{r}}} \dicPos{m}[9] \dicFlx{(‑s\,/\addthin ‑jar, ‑jir)}[32] \dicFieldCat{meteo.} \dicDirectTranslationCS{tornádo}
\dicEntry[hvirfill] \dicTerm{hvirf|ill} \dicIPA{{k\smash{\textsuperscript{h}}}{v}{\textsci}{r}{v}{\textsci}{\textsubring{d}}{\textsubring{l}}} \dicPos{m}[6] \dicFlx{(‑ils, ‑lar)}[35] \textbf{1.} \dicSynonym{kollur} \dicDirectTranslationCS{temeno}  \textbf{2.} \dicDirectTranslationCS{zenit, vrcholný bod};  \dicPhraseIS{frá hvirfli til ilja} \dicFlx{adv} \dicLangCat{přen.} \dicDirectTranslationCS{od hlavy (až) k~patě}
\dicEntry[hvirfilpunktur] \dicTerm{hvirfil··punkt|ur} \dicIPA{{k\smash{\textsuperscript{h}}}{v}{\textsci}{r}{v}{\textsci}{l}{p\smash{\textsuperscript{h}}}{u}{\textsubring{n}}{\textsubring{d}}{\textscy}{\textsubring{r}}} \dicPos{m}[6] \dicFlx{(‑s, ‑ar)}[4] \dicFieldCat{astro.} \dicDirectTranslationCS{zenit, nadhlavník}
\dicEntry[hvirfilvindur] \dicTerm{hvirfil··vind|ur} \dicIPA{{k\smash{\textsuperscript{h}}}{v}{\textsci}{r}{v}{\textsci}{l}{v}{\textsci}{n}{\textsubring{d}}{\textscy}{\textsubring{r}}} \dicPos{m}[6] \dicFlx{(‑s, ‑ar)}[4] \dicFieldCat{meteo.} \dicDirectTranslationCS{tornádo}
\dicEntry[hvirfing] \dicTerm{hvirf··ing} \dicIPA{{k\smash{\textsuperscript{h}}}{v}{\textsci}{r}{v}{i}{\ng}{\r{g}}} \dicPos{f}[4] \dicFlx{(‑ar, ‑ar)}[5] \textbf{1.} \dicSynonym{hringur} \dicDirectTranslationCS{kruh} \dicExampleIS{sitja í hvirfingu um eldinn} \dicExampleCS{sedět v~kruhu kolem ohně}  \textbf{2.} \dicSynonym{þyrping} \dicDirectTranslationCS{skupina, shluk, hlouček} \dicExampleIS{mörg hús í hvirfingu} \dicExampleCS{mnoho domů v~hloučku}
\dicEntry[hvissa] \dicTerm{hviss|a} \dicIPA{{k\smash{\textsuperscript{h}}}{v}{\textsci}{s}{\textlengthmark}{a}} \dicPos{v}[1] \dicFlx{(‑aði)}[34] \textbf{1.} \dicSynonym*{fara hratt} \dicDirectTranslationCS{svištět, řítit se (řeka ap.)}  \textbf{2.} \dicSynonym{hvína} \dicDirectTranslationCS{syčet, svištět (kosa ap.)}
\dicEntry[hví] \dicTerm{hví} \dicIPA{{k\smash{\textsuperscript{h}}}{v}{i}{\textlengthmark}} \dicPos{adv} \dicDirectTranslationCS{proč}
\dicEntry[hvíla] \dicTerm{hvíl|a\smash{\textsuperscript{1}}} \dicIPA{{k\smash{\textsuperscript{h}}}{v}{i}{\textlengthmark}{l}{a}} \dicPos{f}[1] \dicFlx{(‑u, ‑ur)}[19] \dicSynonym{bæli} \dicDirectTranslationCS{pohovka, postel, lože};  \dicPhraseIS{ganga til hvílu} \dicDirectTranslationCS{jít do postele, jít si lehnout}
\dicEntry[hvíla] \dicTerm{hvíl|a\smash{\textsuperscript{2}}} \dicsymFrequent\  \dicIPA{{k\smash{\textsuperscript{h}}}{v}{i}{\textlengthmark}{l}{a}} \dicPos{v}[2] \dicFlx{(‑di, ‑t)}[141] \dicFlx{acc} \textbf{1.} \dicSynonym*{gera afþreyttan} \dicDirectTranslationCS{(nechat) odpočinout\,/\addthin odpočívat} \dicExampleIS{hvíla hestinn} \dicExampleCS{nechat odpočinout koně};  \dicPhraseIS{hvíla sig} \dicDirectTranslationCS{odpočinout si, odpočívat} \dicExampleIS{hvíla sig eftir langan dag} \dicExampleCS{odpočinout si po dlouhém dni}  \textbf{2.} \dicSynonym{liggja} \dicDirectTranslationCS{ležet, spočívat} \dicExampleIS{hvíla í faðmi hennar} \dicExampleCS{spočívat v~jejím náručí}  \textbf{3.} \dicDirectTranslationCS{odpočívat} \dicIndirectTranslationCS{(být pohřben)} \dicExampleIS{hvíla í friði} \dicExampleCS{odpočívat v~pokoji};  \dicPhraseIS{hér hvílir} \dicDirectTranslationCS{zde odpočívá (vytesané na náhrobku ap.)};  \dicIdiom{hvíla}[á]{ \dicPhraseIS{hvíla á e‑m}} \dicDirectTranslationCS{spočívat na (kom), ležet na (kom)} \dicExampleIS{Allt starfið hvílir á honum.} \dicExampleCS{Všechna práce leží na něm.}; { \dicPhraseIS{það hvílir e‑að á e‑u}} \dicFlx{impers} \dicDirectTranslationCS{(co) spočívá na (čem) (kletba na domě ap.)} \dicExampleIS{Það hvílir mikil ábyrgð á okkur.} \dicExampleCS{Spočívá na nás velká zodpovědnost.};  \dicIdiom{hvílast}{ \dicPhraseIS{hvílast}} \dicFlx{refl} \dicDirectTranslationCS{odpočinout si, odpočívat, vydechnout si}
\dicEntry[hvíld] \dicTerm{hvíld} \dicsymFrequent\  \dicIPA{{k\smash{\textsuperscript{h}}}{v}{i}{l}{\textsubring{d}}} \dicPos{f}[7] \dicFlx{(‑ar, ‑ir)}[1] \dicDirectTranslationCS{odpočinek, odpočinutí};  \dicPhraseIS{leggjast til hvíldar} \dicFlx{refl} \dicDirectTranslationCS{lehnout si (ke spánku)};  \dicPhraseIS{taka sér hvíld} \dicDirectTranslationCS{vydechnout si, odpočinout si};  \dicPhraseIS{unna sér ekki hvíldar} \dicDirectTranslationCS{nedopřát si odpočinku}
\dicEntry[hvíldarlaus] \dicTerm{hvíldar··laus} \dicIPA{{k\smash{\textsuperscript{h}}}{v}{i}{l}{\textsubring{d}}{a}{r}{l}{\oe i}{s}} \dicPos{adj}[5]\dicFlx{}[-1] \dicDirectTranslationCS{nepřetržitý, kontinuální}
\dicEntry[hvílíkur] \dicTerm{hví··líkur} \dicsymFrequent\  \dicIPA{{k\smash{\textsuperscript{h}}}{v}{i}{\textlengthmark}{l}{i}{\r{g}}{\textscy}{\textsubring{r}}} \dicPos{pron}[13] \dicFlx{int} \textbf{1.} \dicDirectTranslationCS{jaký}  \textbf{2.} \dicDirectTranslationCS{takový, co za} \dicIndirectTranslationCS{(používané ve zvoláních)} \dicExampleIS{hvílík heppni!} \dicExampleCS{takové štěstí!}
\dicEntry[hvílubrögð] \dicTerm{hvílu··brögð} \dicIPA{{k\smash{\textsuperscript{h}}}{v}{i}{\textlengthmark}{l}{\textscy}{\textsubring{b}}{r}{\oe}{\textbabygamma}{\texttheta}} \dicPos{n}[2] \dicFlx{pl}[9] \dicDirectTranslationCS{soulož, souložení}
\dicEntry[hvín] \dicTerm{hvín} \dicIPA{{k\smash{\textsuperscript{h}}}{v}{i}{\textlengthmark}{\textsubring{n}}} \dicPos{v} \dicFlx{ind praes sg 1 pers} \dicLink{hvína}
\dicEntry[hvína] \dicTerm{hvína} \dicIPA{{k\smash{\textsuperscript{h}}}{v}{i}{\textlengthmark}{n}{a}} \dicPos{v}[6] \dicFlx{(hvín, hvein, hvinum, hvini, hvinið)}[82] \dicDirectTranslationCS{svištět, hvízdat (vítr ap.)};  \dicPhraseIS{það hvín í e‑u} \dicFlx{impers} \dicDirectTranslationCS{(co) hvízdá (vítr ap.)}
\dicEntry[hvínandi] \dicTerm{hvín··andi} \dicIPA{{k\smash{\textsuperscript{h}}}{v}{i}{\textlengthmark}{n}{a}{n}{\textsubring{d}}{\textsci}} \dicPos{adj}[13] \dicFlx{indecl}[1] \dicDirectTranslationCS{hvízdající, svištící (vichřice ap.)}
\dicEntry[hvísl] \dicTerm{hvísl} \dicIPA{{k\smash{\textsuperscript{h}}}{v}{i}{s}{\textsubring{d}}{\textsubring{l}}} \dicPos{n}[2] \dicFlx{(‑s)}[2] \dicDirectTranslationCS{šepot, šeptání}
\dicEntry[hvísla] \dicTerm{hvísl|a} \dicsymFrequent\  \dicIPA{{k\smash{\textsuperscript{h}}}{v}{i}{s}{\textsubring{d}}{l}{a}} \dicPos{v}[1] \dicFlx{(‑aði)}[1] \dicFlx{dat} \dicDirectTranslationCS{(za)šeptat, pošeptat} \dicExampleIS{Hún hvíslaði einhverjum óskiljanlegum orðum.} \dicExampleCS{Zašeptala nějaká nesrozumitelná slova.};  \dicPhraseIS{hvísla e‑u að e‑m} \dicDirectTranslationCS{pošeptat (co komu)};  \dicIdiom{hvíslast}[á]{ \dicPhraseIS{hvíslast á}} \dicFlx{refl} \dicDirectTranslationCS{šeptat si}
\dicEntry[hvíslari] \dicTerm{hvísl··ar|i} \dicIPA{{k\smash{\textsuperscript{h}}}{v}{i}{s}{\textsubring{d}}{l}{a}{r}{\textsci}} \dicPos{m}[1] \dicFlx{(‑a, ‑ar)}[13] \dicDirectTranslationCS{nápověda, nápovědka} \dicIndirectTranslationCS{(osoba napovídající hercům v~divadle)}
\dicEntry[hvísltúlkun] \dicTerm{hvísl··túlk|un} \dicIPA{{k\smash{\textsuperscript{h}}}{v}{i}{s}{\textsubring{d}}{\textsubring{l}}{t\smash{\textsuperscript{h}}}{u}{\textsubring{l}}{\r{g}}{\textscy}{\textsubring{n}}} \dicPos{f}[7] \dicFlx{(‑unar, ‑anir)}[8] \dicDirectTranslationCS{šušotáž, simultánní tlumočení šeptem}
\dicEntry[hvíta] \dicTerm{hvít|a} \dicIPA{{k\smash{\textsuperscript{h}}}{v}{i}{\textlengthmark}{\textsubring{d}}{a}} \dicPos{f}[1] \dicFlx{(‑u, ‑ur)}[7] \textbf{1.} \dicSynonym*{hvítleiki} \dicDirectTranslationCS{bělost, běloba}  \textbf{2.} \dicFieldCat{biol.} \dicSynonym{prótín} \dicDirectTranslationCS{bílkovina, protein}  \textbf{3.} \dicSynonym{eggjahvíta} \dicDirectTranslationCS{bílek}  \textbf{4.} \dicFieldCat{anat.} \dicSynonym*{augnhvíta} \dicDirectTranslationCS{bělmo} \dicIndirectTranslationCS{(bílá zevní část oka)}
\dicEntry[Hvíta-Rússland] \dicTerm{Hvíta-Rúss··land} \dicIPA{{k\smash{\textsuperscript{h}}}{v}{i}{\textlengthmark}{\textsubring{d}}{a}{r}{u}{s}{\textlengthmark}{l}{a}{n}{\textsubring{d}}} \dicPos{n}[2] \dicFlx{(‑s)}[4] \dicFieldCat{geog.} \dicDirectTranslationCS{Bělorusko}
\dicEntry[hvítabjörn] \dicTerm{hvíta··|björn} \dicIPA{{k\smash{\textsuperscript{h}}}{v}{i}{\textlengthmark}{\textsubring{d}}{a}{\textsubring{b}}{j}{\oe}{r}{\textsubring{d}}{\textsubring{n}}} \dicPos{m}[12] \dicFlx{(‑bjarnar, ‑birnir)}[5] \dicFieldCat{zool.} \dicDirectTranslationCS{lední medvěd} \textit{(l.~{\textLA{Ursus maritimus}})}  \dicsymPhoto\ 
\dicFigure{21417.jpg}{Hvítabjörn}{Hvítabjörn - Olsen, Dave, Biolib, PD}
\dicEntry[hvítagaldur] \dicTerm{hvíta··gald|ur} \dicIPA{{k\smash{\textsuperscript{h}}}{v}{i}{\textlengthmark}{\textsubring{d}}{a}{\r{g}}{a}{l}{\textsubring{d}}{\textscy}{\textsubring{r}}} \dicPos{m}[5] \dicFlx{(‑urs)}[2] \dicDirectTranslationCS{bílá magie}
\dicEntry[hvítasunna] \dicTerm{hvíta··sunn|a} \dicIPA{{k\smash{\textsuperscript{h}}}{v}{i}{\textlengthmark}{\textsubring{d}}{a}{s}{\textscy}{n}{a}} \dicPos{f}[1] \dicFlx{(‑u, ‑ur)}[7] \dicDirectTranslationCS{svatodušní svátky, letnice}
\dicEntry[hvítasunnudagur] \dicTerm{hvíta·sunnu··dag|ur} \dicIPA{{k\smash{\textsuperscript{h}}}{v}{i}{\textlengthmark}{\textsubring{d}}{a}{s}{\textscy}{n}{\textscy}{\textsubring{d}}{a}{\textbabygamma}{\textscy}{\textsubring{r}}} \dicPos{m}[6] \dicFlx{(‑s, ‑ar)}[62] \dicDirectTranslationCS{svatodušní neděle}
\dicEntry[hvítblæði] \dicTerm{hvít··blæði} \dicIPA{{k\smash{\textsuperscript{h}}}{v}{i}{\textlengthmark}{\textsubring{d}}{\textsubring{b}}{l}{a}{i}{ð}{\textsci}} \dicPos{n}[2] \dicFlx{(‑s)}[20] \dicFieldCat{med.} \dicDirectTranslationCS{leukémie}
\dicEntry[hvíthærður] \dicTerm{hvít··hærður} \dicIPA{{k\smash{\textsuperscript{h}}}{v}{i}{\textlengthmark}{\textsubring{d}}{h}{a}{i}{r}{ð}{\textscy}{\textsubring{r}}} \dicPos{adj}[2]\dicFlx{}[-1] \dicDirectTranslationCS{bělovlasý}
\dicEntry[hvítingi] \dicTerm{hvít··ing|i} \dicIPA{{k\smash{\textsuperscript{h}}}{v}{i}{\textlengthmark}{\textsubring{d}}{i}{\textltailn}{\r{\textObardotlessj}}{\textsci}} \dicPos{m}[1] \dicFlx{(‑ja, ‑jar)}[14] \dicDirectTranslationCS{albín(ka)}
\dicEntry[hvítkál] \dicTerm{hvít··kál} \dicIPA{{k\smash{\textsuperscript{h}}}{v}{i}{\textlengthmark}{\textsubring{d}}{k\smash{\textsuperscript{h}}}{au}{\textsubring{l}}} \dicPos{n}[2] \dicFlx{(‑s)}[2] \dicFieldCat{bot.} \dicDirectTranslationCS{zelí, zelí hlávkové} \textit{(l.~{\textLA{Brassica oleracea var. capitata}})}
\dicEntry[hvítkorn] \dicTerm{hvít··korn} \dicIPA{{k\smash{\textsuperscript{h}}}{v}{i}{\textlengthmark}{\textsubring{d}}{k\smash{\textsuperscript{h}}}{\textopeno}{r}{\textsubring{d}}{\textsubring{n}}} \dicPos{n}[2] \dicFlx{(‑s, ‑)}[5] \dicFieldCat{biol.} \dicDirectTranslationCS{bílá krvinka, leukocyt}
\dicEntry[hvítlaukspressa] \dicTerm{hvít·lauks··press|a} \dicIPA{{k\smash{\textsuperscript{h}}}{v}{i}{\textlengthmark}{\textsubring{d}}{l}{\oe i}{\r{g}}{s}{p\smash{\textsuperscript{h}}}{r}{\textepsilon}{s}{a}} \dicPos{f}[1] \dicFlx{(‑u, ‑ur)}[7] \dicDirectTranslationCS{lis na česnek}
\dicEntry[hvítlaukur] \dicTerm{hvít··lauk|ur} \dicIPA{{k\smash{\textsuperscript{h}}}{v}{i}{\textlengthmark}{\textsubring{d}}{l}{\oe i}{\r{g}}{\textscy}{\textsubring{r}}} \dicPos{m}[6] \dicFlx{(‑s, ‑ar)}[22] \dicFieldCat{bot.} \dicDirectTranslationCS{česnek, česnek kuchyňský} \textit{(l.~{\textLA{Allium sativum}})}  \dicsymPhoto\ 
\dicFigure{4681.jpg}{Hvítlaukur}{Hvítlaukur - Zicha Ondřej, Biolib, Copyright/CC-BY-NC}
\dicEntry[hvítmaðra] \dicTerm{hvít··|maðra} \dicIPA{{k\smash{\textsuperscript{h}}}{v}{i}{\textlengthmark}{\textsubring{d}}{m}{a}{ð}{r}{a}} \dicPos{f}[1] \dicFlx{(‑möðru, ‑möðrur)}[8] \dicFieldCat{bot.} \dicIndirectTranslationCS{rostlina z~čeledi mořenovitých} \textit{(l.~{\textLA{Galium normanii}})}
\dicEntry[hvítmávur] \dicTerm{hvít··máv|ur} \dicIPA{{k\smash{\textsuperscript{h}}}{v}{i}{\textlengthmark}{\textsubring{d}}{m}{au}{v}{\textscy}{\textsubring{r}}} \dicPos{m}[6] \dicFlx{(‑s, ‑ar)}[19] \dicFieldCat{zool.} \dicDirectTranslationCS{racek šedý} \textit{(l.~{\textLA{Larus hyperboreus}})}  \dicsymPhoto\ 
\dicFigure{20872.jpg}{Hvítmávur}{Hvítmávur - Sowls, Art, Biolib, PD}
\dicEntry[hvítna] \dicTerm{hvítn|a} \dicIPA{{k\smash{\textsuperscript{h}}}{v}{i}{h}{\textsubring{d}}{n}{a}} \dicPos{v}[1] \dicFlx{(‑aði)}[44] \dicDirectTranslationCS{(z)bělat} \dicExampleIS{Hár hvítnaði.} \dicExampleCS{Vlasy zbělaly.}
\dicEntry[Hvítrússi] \dicTerm{Hvít··rúss|i} \dicIPA{{k\smash{\textsuperscript{h}}}{v}{i}{\textlengthmark}{\textsubring{d}}{r}{u}{s}{\textsci}} \dicPos{m}[1] \dicFlx{(‑a, ‑ar)}[1] \dicDirectTranslationCS{Bělorus(ka)}
\dicEntry[hvítrússneska] \dicTerm{hvít··rúss·nesk|a} \dicIPA{{k\smash{\textsuperscript{h}}}{v}{i}{\textlengthmark}{\textsubring{d}}{r}{u}{s}{\textsubring{d}}{n}{\textepsilon}{s}{\r{g}}{a}} \dicPos{f}[1] \dicFlx{(‑u)}[5] \dicDirectTranslationCS{běloruština}
\dicEntry[hvítrússneskur] \dicTerm{hvít··rúss·neskur} \dicIPA{{k\smash{\textsuperscript{h}}}{v}{i}{\textlengthmark}{\textsubring{d}}{r}{u}{s}{\textsubring{d}}{n}{\textepsilon}{s}{\r{g}}{\textscy}{\textsubring{r}}} \dicPos{adj}[1]\dicFlx{}[-6] \dicDirectTranslationCS{běloruský}
\dicEntry[hvítstorkur] \dicTerm{hvít··stork|ur} \dicIPA{{k\smash{\textsuperscript{h}}}{v}{i}{\textlengthmark}{\textsubring{d}}{s}{\textsubring{d}}{\textopeno}{\textsubring{r}}{\r{g}}{\textscy}{\textsubring{r}}} \dicPos{m}[6] \dicFlx{(‑s, ‑ar)}[10] \dicFieldCat{zool.} \dicDirectTranslationCS{čáp, čáp bílý} \textit{(l.~{\textLA{Ciconia ciconia}})}  \dicsymPhoto\ 
\dicFigure{116.jpg}{Hvítstorkur}{Hvítstorkur - Zicha Ondřej, Biolib, Copyright/CC-BY-NC}
\dicEntry[hvítur] \dicTerm{hvít|ur\smash{\textsuperscript{1}}} \dicIPA{{k\smash{\textsuperscript{h}}}{v}{i}{\textlengthmark}{\textsubring{d}}{\textscy}{\textsubring{r}}} \dicPos{m}[6] \dicFlx{(‑s)}[3] \dicDirectTranslationCS{bílý} \dicIndirectTranslationCS{(hráč mající bílé figurky)}
\dicEntry[hvítur] \dicTerm{hvítur\smash{\textsuperscript{2}}} \dicsymFrequent\  \dicIPA{{k\smash{\textsuperscript{h}}}{v}{i}{\textlengthmark}{\textsubring{d}}{\textscy}{\textsubring{r}}} \dicPos{adj}[1]\dicFlx{}[-1] \dicDirectTranslationCS{bílý} \dicExampleIS{hvítt duft} \dicExampleCS{bílý prášek};  \dicPhraseIS{hvítur fyrir hærum} \dicDirectTranslationCS{mající šediny}
\dicEntry[hvítvín] \dicTerm{hvít··vín} \dicIPA{{k\smash{\textsuperscript{h}}}{v}{i}{\textlengthmark}{\textsubring{d}}{v}{i}{\textsubring{n}}} \dicPos{n}[2] \dicFlx{(‑s, ‑)}[5] \dicDirectTranslationCS{bílé víno}
\dicEntry[hvk.] \dicTerm{hvk.} \dicPos{zkr} \dicLink{hk.}
\dicEntry[hvoftur] \dicTerm{hvoft|ur} \dicIPA{{k\smash{\textsuperscript{h}}}{v}{\textopeno}{f}{\textsubring{d}}{\textscy}{\textsubring{r}}} \dicPos{m}[6] \dicFlx{(‑s, ‑ar)}[4] \dicSynonym{gin\smash{\textsuperscript{1}}} \dicDirectTranslationCS{tlama}
\dicEntry[hvolf] \dicTerm{hvolf} \dicIPA{{k\smash{\textsuperscript{h}}}{v}{\textopeno}{l}{f}} \dicPos{n}[2] \dicFlx{(‑s, ‑)}[5] \textbf{1.} \dicSynonym*{himinhvolf} \dicDirectTranslationCS{klenba, klenutí}  \textbf{2.} \dicFieldCat{anat.} \dicSynonym*{hjartahvolf} \dicDirectTranslationCS{srdeční komora};  \dicIdiom{hvolf}{ \dicPhraseIS{e‑að liggur á hvolfi}} \dicDirectTranslationCS{(co) leží vzhůru nohama};  \dicPhraseIS{vera á hvolfi} \dicDirectTranslationCS{cítit se mizerně\,/\addthin bídně (z~opilosti nebo jiných příčin)};  \dicPhraseIS{það er allt á hvolfi} \dicLangCat{přen.} \dicDirectTranslationCS{všechno je vzhůru nohama}
\dicEntry[hvolfa] \dicTerm{hvolf|a} \dicsymFrequent\  \dicIPA{{k\smash{\textsuperscript{h}}}{v}{\textopeno}{l}{v}{a}} \dicPos{v}[2] \dicFlx{(‑di, ‑t)}[161] \dicFlx{dat} \textbf{1.} \dicDirectTranslationCS{překlopit, překlápět, převrátit, převracet} \dicExampleIS{hvolfa ílátinu} \dicExampleCS{překlopit nádobu}  \textbf{2.} \dicPhraseIS{e‑u hvolfir} \dicFlx{impers} \dicDirectTranslationCS{(co) se překlopí (loď ap.)};  \dicIdiom{hvolfa}[í]{ \dicPhraseIS{hvolfa e‑u í sig}} \dicDirectTranslationCS{hodit (co) do sebe, vypít (co) (naráz)} \dicExampleIS{hvolfa bjórnum í sig} \dicExampleCS{hodit do sebe pivo};  \dicIdiom{hvolfast}{ \dicPhraseIS{hvolfast}} \dicFlx{refl} \dicDirectTranslationCS{klenout se, rozkládat se (tma nad městem ap.)}
\dicEntry[hvolfþak] \dicTerm{hvolf··|þak} \dicIPA{{k\smash{\textsuperscript{h}}}{v}{\textopeno}{l}{v}{\texttheta}{a}{\r{g}}} \dicPos{n}[2] \dicFlx{(‑þaks, ‑þök)}[8] \dicDirectTranslationCS{kopule, klenutá střecha}
\dicEntry[hvolpur] \dicTerm{hvolp|ur} \dicIPA{{k\smash{\textsuperscript{h}}}{v}{\textopeno}{\textsubring{l}}{\textsubring{b}}{\textscy}{\textsubring{r}}} \dicPos{m}[6] \dicFlx{(‑s, ‑ar)}[10] \textbf{1.} \dicSynonym*{ungur hundur} \dicDirectTranslationCS{štěně, štěňátko}  \textbf{2.} \dicSynonym*{ungur maður} \dicDirectTranslationCS{štěně} \dicIndirectTranslationCS{(mladý nezkušený člověk)}
\dicEntry[hvor] \dicTerm{hvor} \dicsymFrequent\  \dicIPA{{k\smash{\textsuperscript{h}}}{v}{\textopeno}{\textlengthmark}{\textsubring{r}}} \dicPos{pron}[2] \textbf{1.} \dicFlx{int} \dicDirectTranslationCS{který (ze dvou)} \dicExampleIS{Hvorir unnu, Frakkar eða Króatar?} \dicExampleCS{Kdo vyhrál? Francouzi, nebo Chorvati?}  \textbf{2.} \dicFlx{indef} \dicDirectTranslationCS{každý (ze dvou)} \dicExampleIS{vera með tvo hringa á hvorri hendi} \dicExampleCS{mít dva prsteny na každé ruce};  \dicIdiom{hvor}[annar]{ \dicPhraseIS{hvor annar}} \dicFlx{pron} \dicDirectTranslationCS{vzájemně, navzájem} \dicIndirectTranslationCS{(o~dvou)};  \dicIdiom{hvor}[heldur]{ \dicPhraseIS{hvor heldur (sem)}} \dicFlx{conj} \dicDirectTranslationCS{tak či onak, tak jako tak};  \dicIdiom{hvor}[sinn]{ \dicPhraseIS{hvor sinn}} \dicFlx{pron} \dicDirectTranslationCS{každý svůj} \dicIndirectTranslationCS{(o~dvou)} \dicExampleIS{Við fórum hvort í sína áttina.} \dicExampleCS{Každý jsme odešli ve svém směru.};  \dicIdiom{hvor}[tveggja]{ \dicPhraseIS{hvor tveggja}} \dicFlx{pron} \dicDirectTranslationCS{oba, obě} \dicIndirectTranslationCS{(o~dvou)} \dicExampleIS{Mig vantar hveiti og sykur. Kauptu fyrir mig hvort tveggja.} \dicExampleCS{Chybí mi mouka a~cukr. Kup mi oboje.};  \dicIdiom{sinn}[hvor]{ \dicPhraseIS{sinn hvor}} \dicFlx{pron} \dicDirectTranslationCS{každý jiný} \dicIndirectTranslationCS{(o~dvou)};  \dicPhraseIS{gefa\,/\addthin reka e‑m sitt undir hvorn} \dicDirectTranslationCS{dát (komu) pohlavek\,/\addthin políček};  \dicPhraseIS{sitt af hvoru} \dicDirectTranslationCS{trochu od všeho\,/\addthin každého};  \dicPhraseIS{sitt af hvoru tagi} \dicDirectTranslationCS{každý jiného druhu} \dicIndirectTranslationCS{(o~dvou)}
\dicEntry[hvorki] \dicTerm{hvorki} \dicsymFrequent\  \dicIPA{{k\smash{\textsuperscript{h}}}{v}{\textopeno}{\textsubring{r}}{\r{\textObardotlessj}}{\textsci}} \dicPos{conj} \dicPhraseIS{hvorki -- né} \dicDirectTranslationCS{(ani) -- ani} \dicExampleIS{Hún á hvorki eiginmann né börn.} \dicExampleCS{Nemá ani manžela, ani děti.}
\dicEntry[hvort] \dicTerm{hvort} \dicsymFrequent\  \dicIPA{{k\smash{\textsuperscript{h}}}{v}{\textopeno}{\textsubring{r}}{\textsubring{d}}} \dicPos{conj} \dicDirectTranslationCS{jestli, zda, zdali, zdalipak} \dicExampleIS{Hann spurði hvort ég kæmi aftur á morgun.} \dicExampleCS{Zeptal se, zdali přijdu zase zítra.};  \dicPhraseIS{hvort -- eða} \dicFlx{conj} \dicDirectTranslationCS{(to první ap.) -- (a)nebo (to druhé ap.)} \dicIndirectTranslationCS{(používané v~otázce)} \dicExampleIS{Hvort viltu te eða kaffi?} \dicExampleCS{Chceš čaj, nebo kávu?};  \dicPhraseIS{hvort eð er} \dicFlx{adv} \dicDirectTranslationCS{tak či onak, tak jako tak, stejně};  \dicPhraseIS{hvort heldur -- eða} \dicFlx{conj} \dicDirectTranslationCS{ať již -- nebo\,/\addthin či};  \dicPhraseIS{hvort sem er -- eða} \dicFlx{conj} \dicDirectTranslationCS{ať -- či, buď -- (a)nebo};  \dicPhraseIS{hvort sem er} \dicFlx{adv} \dicSynonym*{í öllu falli} \dicDirectTranslationCS{tak jako tak, stejně};  \dicPhraseIS{hvort sem -- eða} \dicFlx{conj} \dicDirectTranslationCS{ať -- či\,/\addthin nebo}
\dicEntry[hvorugkyn] \dicTerm{hvorug··kyn} \dicIPA{{k\smash{\textsuperscript{h}}}{v}{\textopeno}{\textlengthmark}{r}{\textscy}{\textbabygamma}{c\smash{\textsuperscript{h}}}{\textsci}{\textsubring{n}}} \dicPos{n}[2] \dicFlx{(‑s, ‑)}[13] \dicFieldCat{jaz.} \dicDirectTranslationCS{střední rod, neutrum}
\dicEntry[hvorugkyns] \dicTerm{hvorug··kyns} \dicIPA{{k\smash{\textsuperscript{h}}}{v}{\textopeno}{\textlengthmark}{r}{\textscy}{\textbabygamma}{c\smash{\textsuperscript{h}}}{\textsci}{n}{s}} \dicPos{adj}[13] \dicFlx{indecl}[1] \dicFieldCat{jaz.} \dicDirectTranslationCS{(jsoucí) středního rodu}
\dicEntry[hvorugur] \dicTerm{hvor··ugur} \dicsymFrequent\  \dicIPA{{k\smash{\textsuperscript{h}}}{v}{\textopeno}{\textlengthmark}{r}{\textscy}{\textbabygamma}{\textscy}{\textsubring{r}}} \dicPos{pron}[13] \dicFlx{indef} \dicDirectTranslationCS{žádný (ze dvou), ani jeden (ze dvou)} \dicExampleIS{Hvorugur strákurinn er ljóshærður.} \dicExampleCS{Ani jeden z~kluků není blonďatý.}
\dicEntry[hvorumegin] \dicTerm{hvoru··megin} \dicIPA{{k\smash{\textsuperscript{h}}}{v}{\textopeno}{\textlengthmark}{r}{\textscy}{m}{ei}{\textsci}{\textsubring{n}}} \dicPos{adv} \dicLink{megin}
\dicEntry[hvumpinn] \dicTerm{hvumpinn} \dicIPA{{k\smash{\textsuperscript{h}}}{v}{\textscy}{\textsubring{m}}{\textsubring{b}}{\textsci}{\textsubring{n}}} \dicPos{adj}[6]\dicFlx{}[-2] \dicDirectTranslationCS{podrážděný, nevrlý}
\dicEntry[hvumsa] \dicTerm{hvumsa} \dicIPA{{k\smash{\textsuperscript{h}}}{v}{\textscy}{m}{s}{a}} \dicPos{adj}[13] \dicFlx{indecl}[2] \dicDirectTranslationCS{(jsoucí) beze slova, oněmělý} \dicExampleIS{verða hvumsa (við e‑ð)} \dicExampleCS{oněmět (z~(čeho))}
\dicEntry[hvæs] \dicTerm{hvæs} \dicIPA{{k\smash{\textsuperscript{h}}}{v}{a}{i}{\textlengthmark}{s}} \dicPos{n}[2] \dicFlx{(‑s, ‑)}[5] \dicDirectTranslationCS{syčení, prskání (kočky ap.)}
\dicEntry[hvæsa] \dicTerm{hvæs|a} \dicsymFrequent\  \dicIPA{{k\smash{\textsuperscript{h}}}{v}{a}{i}{\textlengthmark}{s}{a}} \dicPos{v}[2] \dicFlx{(‑ti, ‑t)}[75] \textbf{1.} \dicDirectTranslationCS{(za)syčet, (za)prskat (kočka ap.)} \dicExampleIS{Kötturinn hvæsir á hana.} \dicExampleCS{Kočka na ni syčí.}  \textbf{2.} \dicDirectTranslationCS{(za)syčet} \dicIndirectTranslationCS{(zlostně mluvit)}
\dicEntry[hvönn] \dicTerm{hvönn} \dicIPA{{k\smash{\textsuperscript{h}}}{v}{\oe}{n}{\textlengthmark}} \dicPos{f}[7] \dicFlx{(hvannar, hvannir)}[16] \dicFieldCat{bot.} \dicDirectTranslationCS{andělika lékařská} \textit{(l.~{\textLA{Angelica archangelica}})}  \dicsymPhoto\ 
\dicFigure{74973.jpg}{Hvönn}{Hvönn - Kesl Michael, Biolib, Copyright/CC-BY-NC}
\dicEntry[hvörf] \dicTerm{hvörf} \dicIPA{{k\smash{\textsuperscript{h}}}{v}{\oe}{r}{f}} \dicPos{n} \dicFlx{pl nom} \dicLink{hvarf\smash{\textsuperscript{1}}}
\dicEntry[hvöt] \dicTerm{hvöt} \dicIPA{{k\smash{\textsuperscript{h}}}{v}{\oe}{\textlengthmark}{\textsubring{d}}} \dicPos{f}[7] \dicFlx{(hvatar, hvatir)}[16] \textbf{1.} \dicSynonym{hvatning} \dicDirectTranslationCS{popud, pohnutka, potřeba, nutkání}  \textbf{2.} \dicFieldCat{psych.} \dicDirectTranslationCS{pud, instinkt} \dicExampleIS{kynhvöt} \dicExampleCS{pohlavní pud}
\dicEntry[hvöttum] \dicTerm{hvöttum} \dicIPA{{k\smash{\textsuperscript{h}}}{v}{\oe}{h}{\textsubring{d}}{\textscy}{\textsubring{m}}} \dicPos{v} \dicFlx{ind pf pl 1 pers} \dicLink{hvetja}
\dicEntry[hygði] \dicTerm{hygði} \dicIPA{{h}{\textsci}{\textbabygamma}{ð}{\textsci}} \dicPos{v} \dicFlx{con pf sg 1 pers} \dicLink{hyggja\smash{\textsuperscript{2}}}
\dicEntry[hygg] \dicTerm{hygg} \dicIPA{{h}{\textsci}{\r{g}}{\textlengthmark}} \dicPos{v} \dicFlx{ind praes sg 1 pers} \dicLink{hyggja\smash{\textsuperscript{2}}}
\dicEntry[hyggi] \dicTerm{hyggi} \dicIPA{{h}{\textsci}{\r{\textObardotlessj}}{\textlengthmark}{\textsci}} \dicPos{v} \dicFlx{con pf sg 1 pers} \dicLink{höggva}
\dicEntry[hyggilegur] \dicTerm{hyggi··legur} \dicIPA{{h}{\textsci}{\r{\textObardotlessj}}{\textlengthmark}{\textsci}{l}{\textepsilon}{\textbabygamma}{\textscy}{\textsubring{r}}} \dicPos{adj}[1]\dicFlx{}[-8] \dicDirectTranslationCS{rozvážný, uvážlivý, rozumný} \dicExampleIS{hyggilegt ráð} \dicExampleCS{rozvážná rada}
\dicEntry[hyggindi] \dicTerm{hygg··indi} \dicIPA{{h}{\textsci}{\r{\textObardotlessj}}{\textlengthmark}{\textsci}{n}{\textsubring{d}}{\textsci}} \dicPos{n}[2] \dicFlx{pl}[19] \dicSynonym{klókindi} \dicDirectTranslationCS{rozvážnost, uvážlivost}
\dicEntry[hygginn] \dicTerm{hygginn} \dicIPA{{h}{\textsci}{\r{\textObardotlessj}}{\textlengthmark}{\textsci}{\textsubring{n}}} \dicPos{adj}[6]\dicFlx{}[-2] \dicSynonym*{skynugur} \dicDirectTranslationCS{rozvážný, uvážlivý}
\dicEntry[hyggja] \dicTerm{hyggj|a\smash{\textsuperscript{1}}} \dicIPA{{h}{\textsci}{\r{\textObardotlessj}}{\textlengthmark}{a}} \dicPos{f}[1] \dicFlx{(‑u, ‑ur)}[25] \textbf{1.} \dicSynonym{hugur} \dicDirectTranslationCS{mysl, chápání};  \dicPhraseIS{að e‑rs hyggju, að hyggju e‑rs} \dicDirectTranslationCS{podle chápání (koho)} \dicExampleIS{að minni hyggju} \dicExampleCS{podle mého chápání};  \dicPhraseIS{hafa e‑ð í hyggju} \dicDirectTranslationCS{mít (co) v~plánu, hodlat (co)}  \textbf{2.} \dicIndirectTranslationCS{přípona složených slov vyjadřujících pocity a~vědecké a~názorové směry} \dicExampleIS{rökhyggja} \dicExampleCS{racionalismus}
\dicEntry[hyggja] \dicTerm{hyggja\smash{\textsuperscript{2}}} \dicsymFrequent\  \dicIPA{{h}{\textsci}{\r{\textObardotlessj}}{\textlengthmark}{a}} \dicPos{v}[4] \dicFlx{(hygg, hugði, hugðum, hygði, hugað)}[50] \dicFlx{acc} \dicSynonym{álíta} \dicDirectTranslationCS{mít za to, domnívat se} \dicExampleIS{Ég hugði hann vera sofandi.} \dicExampleCS{Měl jsem za to, že spí.};  \dicIdiom{hyggja}[að]{ \dicPhraseIS{hyggja að e‑u}} \dicDirectTranslationCS{obrátit pozornost na (co), prověřit (co)} \dicExampleIS{hyggja að afleiðingunum} \dicExampleCS{obrátit pozornost na důsledky};  \dicIdiom{hyggja}[á]{ \dicPhraseIS{hyggja á e‑ð}} \dicDirectTranslationCS{zamýšlet (co), strojit (co), plánovat (co)} \dicExampleIS{hyggja á árás} \dicExampleCS{zamýšlet útok};  \dicPhraseIS{eftir á að hyggja} \dicFlx{adv} \dicDirectTranslationCS{po úvaze};  \dicIdiom{hyggjast}{ \dicPhraseIS{hyggjast}} \dicFlx{refl} \dicDirectTranslationCS{plánovat, zamýšlet, hodlat} \dicExampleIS{Hann hyggst rannsaka þetta nánar.} \dicExampleCS{Plánuje, že to blíže prozkoumá.}
\dicEntry[hyggjuvit] \dicTerm{hyggju··vit} \dicIPA{{h}{\textsci}{\r{\textObardotlessj}}{\textlengthmark}{\textscy}{v}{\textsci}{\textsubring{d}}} \dicPos{n}[2] \dicFlx{(‑s)}[2] \dicDirectTranslationCS{rozum, zdravý úsudek, soudnost}
\dicEntry[hygla] \dicTerm{hygl|a} \dicIPA{{h}{\textsci}{\r{g}}{l}{a}} \dicPos{v}[1] \dicFlx{(‑aði)}[34] \dicFlx{dat} \dicSynonym{mismuna} \dicDirectTranslationCS{protežovat, favorizovat} \dicExampleIS{hygla vinum sínum} \dicExampleCS{protežovat své přátele}
\dicEntry[hyl] \dicTerm{hyl} \dicIPA{{h}{\textsci}{\textlengthmark}{\textsubring{l}}} \dicPos{v} \dicFlx{ind praes sg 1 pers} \dicLink{hylja}
\dicEntry[hyldi] \dicTerm{hyldi} \dicIPA{{h}{\textsci}{l}{\textsubring{d}}{\textsci}} \dicPos{v} \dicFlx{con pf sg 1 pers} \dicLink{hylja}
\dicEntry[hyldýpi] \dicTerm{hyl··dýpi} \dicIPA{{h}{\textsci}{l}{\textsubring{d}}{i}{\textsubring{b}}{\textsci}} \dicPos{n}[2] \dicFlx{(‑s)}[20] \dicDirectTranslationCS{(velká) hloubka (ve vodě ap.), hlubina}
\dicEntry[hylja] \dicTerm{hylja} \dicsymFrequent\  \dicIPA{{h}{\textsci}{l}{j}{a}} \dicPos{v}[4] \dicFlx{(hyl, huldi, huldum, hyldi, hulið)}[32] \dicFlx{acc} \dicSynonym{fela\smash{\textsuperscript{1}}} \dicDirectTranslationCS{(za)krýt, zakrývat, (za)halit, zahalovat, zastřít, zastírat} \dicExampleIS{hylja andlit sitt} \dicExampleCS{zakrýt si tvář};  \dicPhraseIS{e‑að hylur e‑m sýn} \dicDirectTranslationCS{(co komu) zakrývá výhled}
\dicEntry[hylki] \dicTerm{hylki} \dicIPA{{h}{\textsci}{\textsubring{l}}{\r{\textObardotlessj}}{\textsci}} \dicPos{n}[2] \dicFlx{(‑s, ‑)}[16] \textbf{1.} \dicSynonym{hirsla} \dicDirectTranslationCS{kazeta, obal, pouzdro};  \dicPhraseIS{hylki utan um e‑ð} \dicDirectTranslationCS{obal na (co)}  \textbf{2.} \dicSynonym*{lyfjahylki} \dicDirectTranslationCS{kapsle} \dicIndirectTranslationCS{(obal na léčiva)}
\dicEntry[hylla] \dicTerm{hyll|a} \dicIPA{{h}{\textsci}{\textsubring{d}}{l}{a}} \dicPos{v}[2] \dicFlx{(‑ti, ‑t)}[63] \dicFlx{acc} \dicDirectTranslationCS{projevit přízeň, uctít} \dicExampleIS{hylla konung} \dicExampleCS{uctít krále}
\dicEntry[hylli] \dicTerm{hyll|i} \dicIPA{{h}{\textsci}{\textsubring{d}}{l}{\textsci}} \dicPos{f}[3] \dicFlx{(‑i)}[3] \dicSynonym{náð\smash{\textsuperscript{1}}} \dicDirectTranslationCS{přízeň, obliba};  \dicPhraseIS{ná hylli e‑rs} \dicDirectTranslationCS{získat (čí) přízeň\,/\addthin oblibu}
\dicEntry[hylling] \dicTerm{hyll··ing} \dicIPA{{h}{\textsci}{\textsubring{d}}{l}{i}{\ng}{\r{g}}} \dicPos{f}[4] \dicFlx{(‑ar, ‑ar)}[5] \dicDirectTranslationCS{pocta, poctění}
\dicEntry[hylma] \dicTerm{hylm|a} \dicIPA{{h}{\textsci}{l}{m}{a}} \dicPos{v}[3] \dicFlx{(‑aði)\,/\addthin (‑di, ‑t)}[8] \dicPhraseIS{hylma yfir e‑ð} \dicDirectTranslationCS{(za)tajit (co), zamlčet\,/\addthin zamlčovat (co) (zločin ap.)};  \dicPhraseIS{hylma yfir með e‑m} \dicDirectTranslationCS{krýt (koho) (ve výpovědi ap.)} \dicExampleIS{hylma yfir með þjófi} \dicExampleCS{krýt zloděje}
\dicEntry[hylur] \dicTerm{hyl|ur} \dicIPA{{h}{\textsci}{\textlengthmark}{l}{\textscy}{\textsubring{r}}} \dicPos{m}[9] \dicFlx{(‑s\,/\addthin ‑jar, ‑jir\,/\addthin ‑ir)}[33] \dicDirectTranslationCS{tůň} \dicExampleIS{Það veiðist stundum lax í hylnum.} \dicExampleCS{V~té tůni se dá občas chytit losos.}
\dicEntry[hypja] \dicTerm{hypj|a} \dicIPA{{h}{\textsci}{\textlengthmark}{\textsubring{b}}{j}{a}} \dicPos{v}[1] \dicFlx{(‑aði)}[34] \dicPhraseIS{hypja sig} \dicLangCat{hovor.} \dicDirectTranslationCS{vypadnout} \dicIndirectTranslationCS{(zmizet)};  \dicPhraseIS{hypja sig í föt} \dicLangCat{hovor.} \dicDirectTranslationCS{obléct se (rychle), hodit na sebe};  \dicPhraseIS{hypjaðu þig!} \dicFlx{imper} \dicLangCat{hovor.} \dicDirectTranslationCS{dej si odchod!, vypadni!}
\dicEntry[hyrfi] \dicTerm{hyrfi} \dicIPA{{h}{\textsci}{r}{v}{\textsci}} \dicPos{v} \dicFlx{con pf sg 1 pers} \dicLink{hverfa\smash{\textsuperscript{2}}}
\dicEntry[hyrna] \dicTerm{hyrn|a} \dicIPA{{h}{\textsci}{r}{\textsubring{d}}{n}{a}} \dicPos{f}[1] \dicFlx{(‑u, ‑ur)}[7] \textbf{1.} \dicDirectTranslationCS{špička, štít (hory ap.)}  \textbf{2.} \dicSynonym*{þríhyrna} \dicDirectTranslationCS{trojcípý šátek}
\dicEntry[hyrndur] \dicTerm{hyrndur} \dicIPA{{h}{\textsci}{r}{n}{\textsubring{d}}{\textscy}{\textsubring{r}}} \dicPos{adj}[2]\dicFlx{}[-14] \dicDirectTranslationCS{rohatý, mající rohy} \dicExampleIS{hyrndur hrútur} \dicExampleCS{rohatý beran}
\dicEntry[hysja] \dicTerm{hysj|a} \dicIPA{{h}{\textsci}{\textlengthmark}{s}{j}{a}} \dicPos{v}[1] \dicFlx{(‑aði)}[34] \dicFlx{acc} \dicDirectTranslationCS{(po)vytáhnout} \dicExampleIS{hysja upp um sig buxurnar} \dicExampleCS{povytáhnout si kalhoty}
\dicEntry[hyski] \dicTerm{hyski} \dicIPA{{h}{\textsci}{s}{\r{\textObardotlessj}}{\textsci}} \dicPos{n}[2] \dicFlx{(‑s)}[20] \textbf{1.} \dicSynonym{fjölskylda} \dicDirectTranslationCS{rodinní příslušníci, rodina, blízcí}  \textbf{2.} \dicSynonym{illþýði} \dicDirectTranslationCS{sebranka, chátra}
\dicEntry[hyskinn] \dicTerm{hyskinn} \dicIPA{{h}{\textsci}{s}{\r{\textObardotlessj}}{\textsci}{\textsubring{n}}} \dicPos{adj}[6]\dicFlx{}[-2] \dicSynonym{latur} \dicDirectTranslationCS{lenivý, líný}
\dicEntry[hysterí] \dicTerm{hysterí} \dicIPA{{h}{\textsci}{s}{\textsubring{d}}{\textepsilon}{r}{i}} \dicPos{n}[2] \dicFlx{(‑s)}[2] \dicLink{hystería}
\dicEntry[hystería] \dicTerm{hysterí|a}\dicTerm{, hysterí} \dicIPA{{h}\-{\textsci}\-{s}\-{\textsubring{d}}\-{\textepsilon}\-{r}\-{i}\-{j}\-{a}\-} \dicPos{f}[1] \dicFlx{(‑u)}[5] \dicDirectTranslationCS{hysterie}
\dicEntry[hysterískur] \dicTerm{hysterískur} \dicIPA{{h}{\textsci}{s}{\textsubring{d}}{\textepsilon}{r}{i}{s}{\r{g}}{\textscy}{\textsubring{r}}} \dicPos{adj}[1]\dicFlx{}[-1] \dicDirectTranslationCS{hysterický}
\dicEntry[hýasinta] \dicTerm{hýasint|a} \dicIPA{{h}{i}{\textlengthmark}{a}{s}{\textsci}{\textsubring{n}}{\textsubring{d}}{a}} \dicPos{f}[1] \dicFlx{(‑u, ‑ur)}[7] \dicFieldCat{bot.} \dicDirectTranslationCS{hyacint, hyacint východní} \textit{(l.~{\textLA{Hyacinthus orientalis}})}  \dicsymPhoto\ 
\dicFigure{ds_image_hyasinta_0_1.jpg}{Hýasinta}{Hýasinta - John O'Neill, CC BY-SA 3.0}
\dicEntry[hýbýli] \dicTerm{hý··býli} \dicIPA{{h}{i}{\textlengthmark}{\textsubring{b}}{i}{l}{\textsci}} \dicPos{n}[2] \dicFlx{pl}[19] \dicLink{híbýli}
\dicEntry[hýða] \dicTerm{hý|ða} \dicIPA{{h}{i}{\textlengthmark}{ð}{a}} \dicPos{v}[2] \dicFlx{(‑ddi, ‑tt)}[167] \dicFlx{acc} \textbf{1.} \dicDirectTranslationCS{(o)loupat} \dicExampleIS{hýða kartöflur} \dicExampleCS{oloupat brambory}  \textbf{2.} \dicSynonym{flengja} \dicDirectTranslationCS{(z)mrskat, (z)bičovat}
\dicEntry[hýði] \dicTerm{hýði} \dicIPA{{h}{i}{\textlengthmark}{ð}{\textsci}} \dicPos{n}[2] \dicFlx{(‑s, ‑)}[14] \dicDirectTranslationCS{slupka, šlupka (brambory ap.), kůra (pomeranče ap.)} \dicExampleIS{taka hýði af kartöflum} \dicExampleCS{oloupat brambory}
\dicEntry[hýðing] \dicTerm{hýð··ing} \dicIPA{{h}{i}{\textlengthmark}{ð}{i}{\ng}{\r{g}}} \dicPos{f}[4] \dicFlx{(‑ar, ‑ar)}[5] \dicDirectTranslationCS{mrskání, bičování}
\dicEntry[hýena] \dicTerm{hýen|a} \dicIPA{{h}{i}{j}{\textepsilon}{n}{a}} \dicPos{f}[1] \dicFlx{(‑u, ‑ur)}[7] \dicLink{híena}
\dicEntry[hýjungur] \dicTerm{hýj··ung|ur} \dicIPA{{h}{i}{j}{\textlengthmark}{u}{\ng}{\r{g}}{\textscy}{\textsubring{r}}} \dicPos{m}[6] \dicFlx{(‑s)}[9] \dicDirectTranslationCS{chmýří, chloupky} \dicIndirectTranslationCS{(jemné vousy)}
\dicEntry[hýr] \dicTerm{hýr} \dicIPA{{h}{i}{\textlengthmark}{\textsubring{r}}} \dicPos{adj}[5] \dicFlx{(f ‑)}[8] \textbf{1.} \dicSynonym{glaðlegur} \dicDirectTranslationCS{veselý, rozradostněný}  \textbf{2.} \dicSynonym{blíður} \dicDirectTranslationCS{líbezný, vlídný};  \dicPhraseIS{líta e‑n hýru auga} \dicDirectTranslationCS{dívat se na (koho) vlídnýma očima}  \textbf{3.} \dicLangCat{slang.} \dicSynonym{samkynhneigður} \dicDirectTranslationCS{homosexuální}
\dicEntry[hýra] \dicTerm{hýr|a} \dicIPA{{h}{i}{\textlengthmark}{r}{a}} \dicPos{f}[1] \dicFlx{(‑u, ‑ur)}[7] \textbf{1.} \dicLangCat{zast.} \dicSynonym{lögg} \dicDirectTranslationCS{kapka, doušek}  \textbf{2.} \dicLangCat{zast.} \dicSynonym{hlýja\smash{\textsuperscript{1}}} \dicDirectTranslationCS{vlídnost, líbeznost}  \textbf{3.} \dicLangCat{zast.} \dicSynonym{laun\smash{\textsuperscript{2}}} \dicDirectTranslationCS{honorář, mzda}
\dicEntry[hýrga] \dicTerm{hýrg|a} \dicIPA{{h}{i}{r}{\r{g}}{a}} \dicPos{v}[1] \dicFlx{(‑aði)}[1] \dicFlx{acc} \dicSynonym{gleðja} \dicDirectTranslationCS{rozveselit, rozjitřit} \dicExampleIS{hýrga sig á brennivíni} \dicExampleCS{rozveselit se pálenkou}
\dicEntry[hýrlegur] \dicTerm{hýr··legur} \dicIPA{{h}{i}{r}{l}{\textepsilon}{\textbabygamma}{\textscy}{\textsubring{r}}} \dicPos{adj}[1]\dicFlx{}[-8] \dicSynonym*{brosleitur} \dicDirectTranslationCS{usměvavý, veselý, radostný}
\dicEntry[hýrna] \dicTerm{hýrn|a} \dicIPA{{h}{i}{r}{\textsubring{d}}{n}{a}} \dicPos{v}[1] \dicFlx{(‑aði)}[47] \dicDirectTranslationCS{rozveselit se, rozjitřit se} \dicExampleIS{Hann hýrnaði fljótt.} \dicExampleCS{Rychle se rozveselil.};  \dicPhraseIS{það hýrnar yfir e‑m} \dicFlx{impers} \dicDirectTranslationCS{(kdo) se rozveseluje}
\dicEntry[hýsa] \dicTerm{hýs|a} \dicIPA{{h}{i}{\textlengthmark}{s}{a}} \dicPos{v}[2] \dicFlx{(‑ti, ‑t)}[63] \dicFlx{acc} \textbf{1.} \dicDirectTranslationCS{ubytovat, poskytnout přístřeší} \dicExampleIS{hýsa gesti} \dicExampleCS{ubytovat hosty}  \textbf{2.} \dicDirectTranslationCS{ustájit} \dicExampleIS{hýsa hesta} \dicExampleCS{ustájit koně}  \textbf{3.} \dicSynonym{rúma} \dicDirectTranslationCS{pojmout, obsáhnout} \dicExampleIS{Stofan hýsir 6 manns.} \dicExampleCS{Místnost pojme 6 lidí.}  \textbf{4.} \dicFieldCat{poč.} \dicDirectTranslationCS{hostovat} \dicExampleIS{hýsa vefsíður} \dicExampleCS{hostovat webové stránky}
\dicEntry[hýsing] \dicTerm{hýs··ing} \dicIPA{{h}{i}{\textlengthmark}{s}{i}{\ng}{\r{g}}} \dicPos{f}[4] \dicFlx{(‑ar, ‑ar)}[5] \textbf{1.} \dicDirectTranslationCS{ubytování}  \textbf{2.} \dicDirectTranslationCS{ustájení} \dicExampleIS{hýsing á hestum} \dicExampleCS{ustájení koní}  \textbf{3.} \dicFieldCat{poč.} \dicDirectTranslationCS{hosting, webhosting}
\dicEntry[hæ] \dicTerm{hæ} \dicIPA{{h}{a}{i}{\textlengthmark}} \dicPos{inter} \dicDirectTranslationCS{ahoj!, čau!, nazdar!}
\dicEntry[hæð] \dicTerm{hæð} \dicsymFrequent\  \dicIPA{{h}{a}{i}{\textlengthmark}{\texttheta}} \dicPos{f}[7] \dicFlx{(‑ar, ‑ir)}[1] \textbf{1.} \dicSynonym{hóll} \dicDirectTranslationCS{kopec, vyvýšenina}  \textbf{2.} \dicDirectTranslationCS{poschodí, patro, podlaží} \dicExampleIS{á fyrstu hæð og á annarri hæð} \dicExampleCS{v~přízemí a~v~prvním patře}  \textbf{3.} \dicSynonym*{það að vera hár} \dicDirectTranslationCS{výše, výška} \dicExampleIS{hæð fjalls} \dicExampleCS{výška hory}  \textbf{4.} \dicFieldCat{meteo.} \dicSynonym{háþrýstisvæði} \dicDirectTranslationCS{tlaková výše} \dicAntonym{lægð}
\dicEntry[hæða] \dicTerm{hæ|ða} \dicIPA{{h}{a}{i}{\textlengthmark}{ð}{a}} \dicPos{v}[2] \dicFlx{(‑ddi, ‑tt)}[167] \dicFlx{acc} \dicSynonym*{hafa að athlægi} \dicDirectTranslationCS{vysmívat se, posmívat se, zesměšňovat} \dicExampleIS{hæða e‑n fyrir e‑ð} \dicExampleCS{posmívat se (komu) za (co)};  \dicIdiom{hæðast}[að]{ \dicPhraseIS{hæðast að e‑m}} \dicFlx{refl} \dicDirectTranslationCS{posmívat se (komu), vysmívat se (komu)}
\dicEntry[hæðardrag] \dicTerm{hæðar··|drag} \dicIPA{{h}{a}{i}{ð}{a}{r}{\textsubring{d}}{r}{a}{x}} \dicPos{n}[2] \dicFlx{(‑drags, ‑drög)}[8] \dicDirectTranslationCS{předhůří}
\dicEntry[hæðarlína] \dicTerm{hæðar··lín|a} \dicIPA{{h}{a}{i}{ð}{a}{r}{l}{i}{n}{a}} \dicPos{f}[1] \dicFlx{(‑u, ‑ur)}[7] \dicFieldCat{geog.} \dicDirectTranslationCS{vrstevnice}
\dicEntry[hæðarmiðja] \dicTerm{hæðar··miðj|a} \dicIPA{{h}{a}{i}{}{ð}{a}{r}{m}{\textsci}{ð}{j}{a}} \dicPos{f}[1] \dicFlx{(‑u, ‑ur)}[7] \dicFieldCat{meteo.} \dicDirectTranslationCS{anticyklóna, oblast vysokého tlaku}
\dicEntry[hæðarmælir] \dicTerm{hæðar··mæl|ir} \dicIPA{{h}{a}{i}{}{ð}{a}{r}{m}{a}{i}{l}{\textsci}{\textsubring{r}}} \dicPos{m}[7] \dicFlx{(‑is, ‑ar)}[1] \dicDirectTranslationCS{výškoměr}
\dicEntry[hæðinn] \dicTerm{hæðinn} \dicIPA{{h}{a}{i}{}{ð}{\textsci}{\textsubring{n}}} \dicPos{adj}[6]\dicFlx{}[-2] \dicDirectTranslationCS{posměvačný, výsměšný, ironický}
\dicEntry[hæðni] \dicTerm{hæðn|i} \dicIPA{{h}{a}{i}{ð}{n}{\textsci}} \dicPos{f}[3] \dicFlx{(‑i)}[3] \dicDirectTranslationCS{posměch, výsměch, ironie}
\dicEntry[hæðnislega] \dicTerm{hæðnis··lega} \dicIPA{{h}{a}{i}{ð}{n}{\textsci}{s}{l}{\textepsilon}{\textbabygamma}{a}} \dicPos{adv} \dicDirectTranslationCS{posměšně, výsměšně, ironicky}
\dicEntry[hæðnislegur] \dicTerm{hæðnis··legur} \dicIPA{{h}{a}{i}{ð}{n}{\textsci}{s}{l}{\textepsilon}{\textbabygamma}{\textscy}{\textsubring{r}}} \dicPos{adj}[1]\dicFlx{}[-8] \dicDirectTranslationCS{posměvačný, výsměšný, ironický}
\dicEntry[hæðóttur] \dicTerm{hæð··óttur} \dicIPA{{h}{a}{i}{\textlengthmark}{ð}{ou}{h}{\textsubring{d}}{\textscy}{\textsubring{r}}} \dicPos{adj}[1]\dicFlx{}[-10] \dicDirectTranslationCS{kopcovitý, zvlněný (krajina ap.)}
\dicEntry[hæfa] \dicTerm{hæf|a\smash{\textsuperscript{1}}} \dicIPA{{h}{a}{i}{\textlengthmark}{v}{a}} \dicPos{f}[1] \dicFlx{(‑u)}[5] \dicPhraseIS{það er engin hæfa í e‑ð} \dicDirectTranslationCS{není ani trochu pravdy na (čem)}
\dicEntry[hæfa] \dicTerm{hæf|a\smash{\textsuperscript{2}}} \dicIPA{{h}{a}{i}{\textlengthmark}{v}{a}} \dicPos{v}[2] \dicFlx{(‑ði, ‑t)}[111] \dicFlx{acc\,/\addthin dat} \textbf{1.} \dicFlx{acc} \dicSynonym{hitta} \dicDirectTranslationCS{trefit, zasáhnout} \dicExampleIS{hæfa mark} \dicExampleCS{trefit branku}  \textbf{2.} \dicFlx{dat} \dicDirectTranslationCS{být vhodný, hodit se, vyhovovat} \dicExampleIS{Fötin hæfa honum vel.} \dicExampleCS{Šaty se k~němu hodí.}
\dicEntry[hæfi] \dicTerm{hæfi\smash{\textsuperscript{1}}} \dicsymFrequent\  \dicIPA{{h}{a}{i}{\textlengthmark}{v}{\textsci}} \dicPos{n}[2] \dicFlx{(‑s)}[20] \dicDirectTranslationCS{vhodnost, náležitost};  \dicPhraseIS{e‑að er við hæfi e‑rs} \dicDirectTranslationCS{(co) je vhodné pro (koho)} \dicExampleIS{Þessi bíómynd er ekki við hæfi barna.} \dicExampleCS{Tenhle film není vhodný pro děti.};  \dicPhraseIS{gera e‑m til hæfis} \dicDirectTranslationCS{vyhovět (komu)}
\dicEntry[hæfi] \dicTerm{hæfi\smash{\textsuperscript{2}}} \dicIPA{{h}{a}{i}{\textlengthmark}{v}{\textsci}} \dicPos{v} \dicFlx{con pf sg 1 pers} \dicLink{hefja}
\dicEntry[hæfilegur] \dicTerm{hæfi··legur} \dicIPA{{h}{a}{i}{\textlengthmark}{v}{\textsci}{l}{\textepsilon}{\textbabygamma}{\textscy}{\textsubring{r}}} \dicPos{adj}[1]\dicFlx{}[-8] \dicDirectTranslationCS{vhodný, vyhovující, přijatelný} \dicExampleIS{hæfileg laun} \dicExampleCS{vyhovující mzda}
\dicEntry[hæfileikaríkur] \dicTerm{hæfi·leika··ríkur} \dicIPA{{h}{a}{i}{\textlengthmark}{v}{\textsci}{l}{ei}{\r{g}}{a}{r}{i}{\r{g}}{\textscy}{\textsubring{r}}} \dicPos{adj}[1]\dicFlx{}[-1] \dicDirectTranslationCS{nadaný, talentovaný}
\dicEntry[hæfileiki] \dicTerm{hæfi··leik|i} \dicsymFrequent\  \dicIPA{{h}{a}{i}{\textlengthmark}{v}{\textsci}{l}{ei}{\r{\textObardotlessj}}{\textsci}} \dicPos{m}[1] \dicFlx{(‑a, ‑ar)}[1] \textbf{1.} \dicDirectTranslationCS{schopnost} \dicExampleIS{búa yfir óvenjulegum hæfileikum} \dicExampleCS{mít neobvyklé schopnosti}  \textbf{2.} \dicDirectTranslationCS{talent, nadání}
\dicEntry[hæfni] \dicTerm{hæfn|i} \dicIPA{{h}{a}{i}{\textsubring{b}}{n}{\textsci}} \dicPos{f}[3] \dicFlx{(‑i)}[3] \dicDirectTranslationCS{schopnost, způsobilost}
\dicEntry[hæfnispróf] \dicTerm{hæfnis··próf} \dicIPA{{h}{a}{i}{\textsubring{b}}{n}{\textsci}{s}{p\smash{\textsuperscript{h}}}{r}{ou}{f}} \dicPos{n}[2] \dicFlx{(‑s, ‑)}[5] \dicDirectTranslationCS{zkouška způsobilosti\,/\addthin kompetence, talentová zkouška}
\dicEntry[hæfur] \dicTerm{hæfur} \dicIPA{{h}{a}{i}{\textlengthmark}{v}{\textscy}{\textsubring{r}}} \dicPos{adj}[1]\dicFlx{}[-1] \textbf{1.} \dicDirectTranslationCS{způsobilý, kompetentní, kvalifikovaný};  \dicPhraseIS{vera hæfur til e‑s} \dicDirectTranslationCS{být způsobilý\,/\addthin kompetentní k~(čemu)}  \textbf{2.} \dicSynonym{nothæfur} \dicDirectTranslationCS{vhodný, použitelný}  \textbf{3.} \dicPhraseIS{hæft} \dicFlx{n} \dicSynonym*{raunsannur} \dicDirectTranslationCS{pravdivý, nevymyšlený} \dicIndirectTranslationCS{(používané většinou v~záporu)} \dicExampleIS{Er eitthvað hæft í þessu?} \dicExampleCS{Je na tom něco pravdy?}
\dicEntry[hægagangur] \dicTerm{hæga··gang|ur} \dicIPA{{h}{a}{i}{\textlengthmark}{\textbabygamma}{a}{\r{g}}{au}{\ng}{\r{g}}{\textscy}{\textsubring{r}}} \dicPos{m}[6] \dicFlx{(‑s)}[7] \textbf{1.} \dicLangCat{přen.} \dicSynonym*{hæg ferð} \dicDirectTranslationCS{hlemýždí tempo}  \textbf{2.} \dicDirectTranslationCS{volnoběh}
\dicEntry[hægð] \dicTerm{hægð} \dicsymFrequent\  \dicIPA{{h}{a}{i}{\textbabygamma}{\texttheta}} \dicPos{f}[7] \dicFlx{(‑ar, ‑ir)}[1] \textbf{1.} \dicSynonym*{kyrrlæti} \dicDirectTranslationCS{klid, pohoda} \dicExampleIS{svara með hægð} \dicExampleCS{odpovědět s~klidem};  \dicPhraseIS{gera e‑ð í hægðum sínum} \dicDirectTranslationCS{dělat (co) zvolna, dělat (co) vlastním tempem}  \textbf{2.} \dicSynonym*{auðveldni} \dicDirectTranslationCS{snadnost, lehkost}  \textbf{3.} \dicPhraseIS{hægðir} \dicFlx{pl} \dicSynonym{saur} \dicDirectTranslationCS{výkaly, lejna};  \dicPhraseIS{hafa hægðir} \dicDirectTranslationCS{vyprázdnit se, vykonat potřebu}
\dicEntry[hægðalyf] \dicTerm{hægða··lyf} \dicIPA{{h}{a}{i}{\textbabygamma}{ð}{a}{l}{\textsci}{f}} \dicPos{n}[2] \dicFlx{(‑s, ‑)}[13] \dicFieldCat{med.} \dicDirectTranslationCS{projímadlo, laxativum}
\dicEntry[hægðarauki] \dicTerm{hægðar··auk|i} \dicIPA{{h}{a}{i}{\textbabygamma}{ð}{a}{r}{\oe i}{\r{\textObardotlessj}}{\textsci}} \dicPos{m}[1] \dicFlx{(‑a)}[3] \dicDirectTranslationCS{usnadnění, ulehčení, větší pohodlí};  \dicPhraseIS{til hægðarauka} \dicFlx{adv} \dicDirectTranslationCS{pro usnadnění, pro větší pohodlí}
\begin{xtolerant}{}{1pt}
\dicEntry[hægðarleikur] \dicTerm{hægðar··leik|ur}\addthinS\dicIPA{{h}{a}{i}{\textbabygamma}{ð}{a}{r}{l}{ei}{\r{g}}{\textscy}{\textsubring{r}}}\addthinS\dicPos{m}[6]\addthinS\dicFlx{(‑s)}[17] \dicDirectTranslationCS{lehká věc (k~udělání ap.), hračka};  \dicPhraseIS{það er hægðarleikur að (gera e‑ð)} \dicDirectTranslationCS{je hračka ((co) udělat)}
\end{xtolerant}
\dicEntry[hægðatregða] \dicTerm{hægða··tregð|a} \dicIPA{{h}{a}{i}{\textbabygamma}{ð}{a}{t\smash{\textsuperscript{h}}}{r}{\textepsilon}{\textbabygamma}{ð}{a}} \dicPos{f}[1] \dicFlx{(‑u)}[5] \dicFieldCat{med.} \dicDirectTranslationCS{zácpa}
\dicEntry[hægfara] \dicTerm{hæg··fara} \dicIPA{{h}{a}{i}{x}{f}{a}{r}{a}} \dicPos{adj}[13] \dicFlx{indecl}[1] \dicSynonym{seinn} \dicDirectTranslationCS{pozvolný, postupný}
\dicEntry[hægindastóll] \dicTerm{hæginda··stól|l} \dicIPA{{h}{a}{i}{j}{\textlengthmark}{\textsci}{n}{\textsubring{d}}{a}{s}{\textsubring{d}}{ou}{\textsubring{d}}{\textsubring{l}}} \dicPos{m}[6] \dicFlx{(‑s, ‑ar)}[48] \dicDirectTranslationCS{lenoška, (pohodlné) křeslo}
\dicEntry[hægindi] \dicTerm{hæg··indi} \dicIPA{{h}{a}{i}{j}{\textlengthmark}{\textsci}{n}{\textsubring{d}}{\textsci}} \dicPos{n}[2] \dicFlx{(‑s, ‑)}[14] \textbf{1.} \dicSynonym{hægindastóll} \dicDirectTranslationCS{lenoška, (pohodlné) křeslo}  \textbf{2.} \dicSynonym{þægindi} \dicDirectTranslationCS{pohodlí}
\dicEntry[hægja] \dicTerm{hæg|ja} \dicsymFrequent\  \dicIPA{{h}{a}{i}{j}{\textlengthmark}{a}} \dicPos{v}[2] \dicFlx{(‑ði, ‑t)}[93] \dicFlx{acc} \textbf{1.} \dicPhraseIS{hægja ferðina} \dicDirectTranslationCS{zpomalit (jízdu)} \dicExampleIS{hægja ferðina í beygjunni} \dicExampleCS{zpomalit v~zatáčce};  \dicPhraseIS{hægja á sér} \dicDirectTranslationCS{zpomalit} \dicExampleIS{Báturinn hægði á sér.} \dicExampleCS{Člun zpomalil.}  \textbf{2.} \dicPhraseIS{hægja sér} \dicSynonym{kúka} \dicDirectTranslationCS{vykakat se, vyprázdnit se}  \textbf{3.} \dicPhraseIS{það hægir} \dicFlx{impers} \dicDirectTranslationCS{vítr se utišuje}  \textbf{4.} \dicPhraseIS{e‑m hægir við e‑ð} \dicFlx{impers} \dicDirectTranslationCS{(kdo) se uklidňuje (čím)};  \dicIdiom{hægjast}{ \dicPhraseIS{hægjast}} \dicFlx{refl} \dicDirectTranslationCS{uklidnit se, zklidnit se, utišit se};  \dicPhraseIS{það hægist um} \dicFlx{refl impers} \dicDirectTranslationCS{zklidňuje se (období v~životě ap.)};  \dicPhraseIS{það hægist á e‑u} \dicFlx{refl impers} \dicDirectTranslationCS{padá klid na (co), (co) se uklidňuje}
\dicEntry[hæglátur] \dicTerm{hæg··látur} \dicIPA{{h}{a}{i}{\textbabygamma}{l}{au}{\textsubring{d}}{\textscy}{\textsubring{r}}} \dicPos{adj}[1]\dicFlx{}[-1] \dicSynonym{rólegur} \dicDirectTranslationCS{klidný, pokojný, rozvážný} \dicExampleIS{hæglátur stúdent} \dicExampleCS{rozvážný student}
\dicEntry[hægri] \dicTerm{hægr|i} \dicsymFrequent\  \dicIPA{{h}{a}{i}{\textbabygamma}{r}{\textsci}} \dicPos{adj}[12] \dicFlx{(n ‑a)}[3] \textbf{1.} \dicDirectTranslationCS{pravý} \dicIndirectTranslationCS{(jsoucí na druhé straně těla než je srdce; opak levý)} \dicAntonym{vinstri};  \dicPhraseIS{til hægri} \dicFlx{adv} \dicDirectTranslationCS{doprava, napravo}  \textbf{2.} \dicDirectTranslationCS{pravicový};  \dicPhraseIS{hægri maður} \dicDirectTranslationCS{stoupenec\,/\addthin stoupenkyně pravice};  \dicPhraseIS{vera hægri hönd e‑rs} \dicLangCat{přen.} \dicDirectTranslationCS{být pravou rukou (koho)}
\dicEntry[hægriflokkur] \dicTerm{hægri··flokk|ur} \dicIPA{{h}{a}{i}{\textbabygamma}{r}{\textsci}{f}{l}{\textopeno}{h}{\r{g}}{\textscy}{\textsubring{r}}} \dicPos{m}[6] \dicFlx{(‑s, ‑ar)}[8] \dicDirectTranslationCS{pravicová strana} \dicAntonym{vinstriflokkur}
\dicEntry[hægrimaður] \dicTerm{hægri··|maður} \dicIPA{{h}{a}{i}{\textbabygamma}{r}{\textsci}{m}{a}{ð}{\textscy}{\textsubring{r}}} \dicPos{m}[13] \dicFlx{(‑manns, ‑menn)}[2] \dicDirectTranslationCS{pravicově orientovaný člověk, stoupenec\,/\addthin stoupenkyně pravice, pravičák, pravičačka} \dicAntonym{vinstrimaður}
\dicEntry[hægrisinnaður] \dicTerm{hægri··sinn·|aður} \dicIPA{{h}{a}{i}{\textbabygamma}{r}{\textsci}{s}{\textsci}{n}{a}{ð}{\textscy}{\textsubring{r}}} \dicPos{adj}[3] \dicFlx{(f ‑uð)}[3] \dicDirectTranslationCS{pravicový, pravicově smýšlející\,/\addthin orientovaný} \dicAntonym{vinstrisinnaður}
\dicEntry[hægrisinni] \dicTerm{hægri··sinn|i} \dicIPA{{h}{a}{i}{\textbabygamma}{r}{\textsci}{s}{\textsci}{n}{\textsci}} \dicPos{m}[1] \dicFlx{(‑a, ‑ar)}[1] \dicDirectTranslationCS{pravicově orientovaný člověk, stoupenec\,/\addthin stoupenkyně pravice} \dicAntonym{vinstrisinni}
\dicEntry[hægristjórn] \dicTerm{hægri··stjórn}\dicTerm{, hægri stjórn} \dicIPA{{h}\-{a}\-{i}\-{\textbabygamma}\-{r}\-{\textsci}\-{s}\-{\textsubring{d}}\-{j}\-{ou}\-{r}\-{\textsubring{d}}\-{\textsubring{n}}\-} \dicPos{f}[7] \dicFlx{(‑ar, ‑ir)}[1] \dicDirectTranslationCS{pravicová vláda} \dicAntonym{vinstristjórn}
\dicEntry[hægt] \dicTerm{hæg|t} \dicsymFrequent\  \dicIPA{{h}{a}{i}{x}{\textsubring{d}}} \dicPos{adv} \dicFlx{(comp ‑ar\,/\addthin ‑ara, sup ‑ast)} \dicDirectTranslationCS{pomalu, zvolna};  \dicPhraseIS{hægt og rólega} \dicFlx{adv} \dicDirectTranslationCS{pomalu a~klidně, krůček po krůčku};  \dicPhraseIS{ef hægt er} \dicDirectTranslationCS{pokud je to možné};  \dicPhraseIS{það er hægt að (gera e‑ð)} \dicDirectTranslationCS{je možné ((co) udělat), lze ((co) udělat)};  \dicPhraseIS{það er hægara sagt en gert} \dicDirectTranslationCS{to se snadněji řekne, než udělá}
\dicEntry[hægur] \dicTerm{hægur} \dicsymFrequent\  \dicIPA{{h}{a}{i}{\textlengthmark}{\textbabygamma}{\textscy}{\textsubring{r}}} \dicPos{adj}[1]\dicFlx{}[-1] \textbf{1.} \dicDirectTranslationCS{slabý, mírný (vítr ap.)}  \textbf{2.} \dicDirectTranslationCS{pomalý, volný (vývoj ap.)}  \textbf{3.} \dicDirectTranslationCS{pomalý, poklidný (jízda ap.)} \dicExampleIS{hægur akstur ökumanna} \dicExampleCS{pomalá jízda řidičů}  \textbf{4.} \dicSynonym{rólegur} \dicDirectTranslationCS{klidný, rozvážný};  \dicPhraseIS{hafa hægt um sig} \dicDirectTranslationCS{držet se zpátky}
\dicEntry[hægviðri] \dicTerm{hæg··viðri} \dicIPA{{h}{a}{i}{\textbabygamma}{v}{\textsci}{ð}{r}{\textsci}} \dicPos{n}[2] \dicFlx{(‑s, ‑)}[14] \dicDirectTranslationCS{klidné počasí}
\dicEntry[hækja] \dicTerm{hækj|a} \dicIPA{{h}{a}{i}{\textlengthmark}{\r{\textObardotlessj}}{a}} \dicPos{f}[1] \dicFlx{(‑u, ‑ur)}[7] \dicDirectTranslationCS{berla};  \dicPhraseIS{sitja á hækjum sér\,/\addthin sínum} \dicDirectTranslationCS{dřepět}
\dicEntry[hækka] \dicTerm{hækk|a} \dicsymFrequent\  \dicIPA{{h}{a}{i}{h}{\r{g}}{a}} \dicPos{v}[1] \dicFlx{(‑aði)}[44] \dicFlx{acc} \textbf{1.} \dicDirectTranslationCS{(vy)růst (ceny ap.), stoupat, stoupnout};  \dicPhraseIS{e‑að hækkar í verði} \dicDirectTranslationCS{(co) stoupá v~ceně};  \dicPhraseIS{það hækkar í e‑u} \dicFlx{impers} \dicDirectTranslationCS{(co) stoupá (voda v~řece ap.)}  \textbf{2.} \dicDirectTranslationCS{zvýšit, zvyšovat} \dicExampleIS{hækka laun} \dicExampleCS{zvýšit plat};  \dicPhraseIS{hækka röddina\,/\addthin róminn} \dicDirectTranslationCS{zvýšit hlas};  \dicPhraseIS{hækka í e‑u} \dicDirectTranslationCS{dát (víc) nahlas (co), zvýšit hlasitost (čeho)} \dicExampleIS{hækka í útvarpinu} \dicExampleCS{dát rádio (víc) nahlas};  \dicPhraseIS{hækka upp} \dicFieldCat{mat.} \dicDirectTranslationCS{zaokrouhlit}
\dicEntry[hækkun] \dicTerm{hækk|un} \dicIPA{{h}{a}{i}{h}{\r{g}}{\textscy}{\textsubring{n}}} \dicPos{f}[7] \dicFlx{(‑unar, ‑anir)}[8] \textbf{1.} \dicSynonym{aukning} \dicDirectTranslationCS{zvýšení, stoupání, nárůst}  \textbf{2.} \dicSynonym{verðhækkun} \dicDirectTranslationCS{zdražení, zdražování, zvýšení\,/\addthin zvyšování ceny}
\dicEntry[hæla] \dicTerm{hæl|a} \dicIPA{{h}{a}{i}{\textlengthmark}{l}{a}} \dicPos{v}[2] \dicFlx{(‑di, ‑t)}[149] \dicFlx{dat} \dicSynonym{lofa} \dicDirectTranslationCS{(po)chválit};  \dicPhraseIS{hæla e‑m á hvert reipi} \dicLangCat{přen.} \dicDirectTranslationCS{vychvalovat (koho) do nebes};  \dicIdiom{hælast}[um]{ \dicPhraseIS{hælast um af e‑u}} \dicFlx{refl} \dicSynonym{gorta} \dicDirectTranslationCS{chlubit se (čím), chvástat se (čím)}
\dicEntry[hælahár] \dicTerm{hæla··|hár} \dicIPA{{h}{a}{i}{\textlengthmark}{l}{a}{h}{au}{\textsubring{r}}} \dicPos{adj}[10] \dicFlx{(comp ‑hærri, sup ‑hæstur)}[14] \dicDirectTranslationCS{(jsoucí) s~vysokými podpatky} \dicExampleIS{hælaháir skór} \dicExampleCS{boty s~vysokými podpatky}
\dicEntry[hælbein] \dicTerm{hæl··bein} \dicIPA{{h}{a}{i}{l}{\textsubring{b}}{ei}{\textsubring{n}}} \dicPos{n}[2] \dicFlx{(‑s, ‑)}[5] \dicFieldCat{anat.} \dicDirectTranslationCS{kost patní} \textit{(l.~{\textLA{calcaneus}})}
\dicEntry[hæli] \dicTerm{hæli} \dicIPA{{h}{a}{i}{\textlengthmark}{l}{\textsci}} \dicPos{n}[2] \dicFlx{(‑s, ‑)}[14] \textbf{1.} \dicSynonym{skjól} \dicDirectTranslationCS{útočiště, azyl} \dicExampleIS{pólitískt hæli} \dicExampleCS{politický azyl};  \dicPhraseIS{leita hælis e‑s staðar} \dicDirectTranslationCS{hledat (kde) útočiště}  \textbf{2.} \dicSynonym{sjúkrahús} \dicDirectTranslationCS{sanatorium, institut pro léčbu pacientů} \dicExampleIS{berklahæli} \dicExampleCS{institut pro léčbu tuberkulózy}
\dicEntry[hæll] \dicTerm{hæl|l} \dicsymFrequent\  \dicIPA{{h}{a}{i}{\textsubring{d}}{\textsubring{l}}} \dicPos{m}[6] \dicFlx{(‑s, ‑ar)}[51] \textbf{1.} \dicFieldCat{anat.} \dicDirectTranslationCS{pata} \dicExampleIS{hællinn á mér} \dicExampleCS{moje pata};  \dicPhraseIS{vera á hælunum á e‑m} \dicDirectTranslationCS{být (komu) v~patách};  \dicPhraseIS{það er undir hælinn lagt} \dicDirectTranslationCS{není to jisté};  \dicPhraseIS{um hæl} \dicFlx{adv} \dicDirectTranslationCS{obratem, okamžitě} \dicExampleIS{svara um hæl} \dicExampleCS{odpovědět obratem}  \textbf{2.} \dicDirectTranslationCS{pata (u~boty ap.)}  \textbf{3.} \dicDirectTranslationCS{podpatek} \dicExampleIS{háir hælar} \dicExampleCS{vysoké podpatky};  \dicPhraseIS{snúast á hæli} \dicFlx{refl} \dicDirectTranslationCS{otočit se na podpatku}  \textbf{4.} \dicSynonym*{stautur} \dicDirectTranslationCS{kolík, kolíček} \dicExampleIS{tjaldhæll} \dicExampleCS{stanový kolík}
\dicEntry[hæna] \dicTerm{hæn|a\smash{\textsuperscript{1}}} \dicIPA{{h}{a}{i}{\textlengthmark}{n}{a}} \dicPos{f}[1] \dicFlx{(‑u, ‑ur)}[7] \dicDirectTranslationCS{slepice} \dicIndirectTranslationCS{(samice kura domácího)} \dicExampleIS{hani og hæna} \dicExampleCS{kohout a~slepice}
\dicEntry[hæna] \dicTerm{hæn|a\smash{\textsuperscript{2}}} \dicIPA{{h}{a}{i}{\textlengthmark}{n}{a}} \dicPos{v}[2] \dicFlx{(‑di, ‑t)}[141] \dicFlx{acc} \dicSynonym{lokka} \dicDirectTranslationCS{(při)lákat, přitahovat, (při)vábit, nalákat};  \dicPhraseIS{hæna e‑n að sér} \dicDirectTranslationCS{přilákat (koho) k~sobě};  \dicIdiom{hænast}[að]{ \dicPhraseIS{hænast að e‑m}} \dicFlx{refl} \dicDirectTranslationCS{být přitahován\,/\addthin lákán ke (komu)}
\dicEntry[hændur] \dicTerm{hændur} \dicIPA{{h}{a}{i}{n}{\textsubring{d}}{\textscy}{\textsubring{r}}} \dicPos{adj}[2]\dicFlx{}[-14] \dicPhraseIS{vera hændur að e‑m} \dicDirectTranslationCS{být přitahován ke (komu)}
\dicEntry[hængur] \dicTerm{hæng|ur} \dicIPA{{h}{a}{i}{\ng}{\r{g}}{\textscy}{\textsubring{r}}} \dicPos{m}[6] \dicFlx{(‑s, ‑ir\,/\addthin ‑ar)}[70] \textbf{1.} \dicDirectTranslationCS{mlíčák, mlíčník, rybí samec}  \textbf{2.} \dicSynonym{galli\smash{\textsuperscript{2}}} \dicDirectTranslationCS{zádrhel, háček} \dicExampleIS{Það er hængur á því.} \dicExampleCS{Je v~tom háček.}
\dicEntry[hænsn] \dicTerm{hænsn} \dicIPA{{h}{a}{i}{n}{s}{\textsubring{d}}{\textsubring{n}}} \dicPos{n}[2] \dicFlx{(‑s, ‑)}[5] \dicLink{hænsni}
\dicEntry[hænsnabú] \dicTerm{hænsna··bú} \dicIPA{{h}{a}{i}{n}{s}{\textsubring{d}}{n}{a}{\textsubring{b}}{u}} \dicPos{n}[2] \dicFlx{(‑s, ‑)}[5] \dicDirectTranslationCS{slepičárna}
\dicEntry[hænsnahús] \dicTerm{hænsna··hús} \dicIPA{{h}{a}{i}{n}{s}{\textsubring{d}}{n}{a}{h}{u}{s}} \dicPos{n}[2] \dicFlx{(‑s, ‑)}[5] \dicDirectTranslationCS{drůbežárna}
\dicEntry[hænsnakofi] \dicTerm{hænsna··kof|i} \dicIPA{{h}{a}{i}{n}{s}{\textsubring{d}}{n}{a}{k\smash{\textsuperscript{h}}}{\textopeno}{v}{\textsci}} \dicPos{m}[1] \dicFlx{(‑a, ‑ar)}[1] \dicDirectTranslationCS{kurník, slepičárna}
\dicEntry[hænsni] \dicTerm{hænsni}\dicTerm{, hænsn} \dicIPA{{h}\-{a}\-{i}\-{n}\-{s}\-{\textsubring{d}}\-{n}\-{\textsci}\-} \dicPos{n}[2] \dicFlx{(‑s, ‑)}[14] \dicFieldCat{zool.} \dicDirectTranslationCS{kur, kur domácí} \textit{(l.~{\textLA{Gallus gallus f. domestica}})}  \dicsymPhoto\ 
\dicFigure{29593.jpg}{Hænsni}{Hænsni - Velechovský Marek, Biolib, Copyright/CC-BY-NC}
\dicEntry[hænuegg] \dicTerm{hænu··egg} \dicIPA{{h}{a}{i}{\textlengthmark}{n}{\textscy}{\textepsilon}{\r{g}}} \dicPos{n}[2] \dicFlx{(‑s, ‑)}[13] \dicDirectTranslationCS{slepičí vejce}
\dicEntry[hænufet] \dicTerm{hænu··fet} \dicIPA{{h}{a}{i}{\textlengthmark}{n}{\textscy}{f}{\textepsilon}{\textsubring{d}}} \dicPos{n}[2] \dicFlx{(‑s, ‑)}[5] \dicLangCat{přen.} \dicDirectTranslationCS{slepičí krok} \dicIndirectTranslationCS{(velmi krátký krok)};  \dicPhraseIS{komast\,/\addthin fara ekki hænufet} \dicLangCat{přen.} \dicDirectTranslationCS{nehnout se z~místa}
\dicEntry[hænuhaus] \dicTerm{hænu··haus} \dicIPA{{h}{a}{i}{\textlengthmark}{n}{\textscy}{h}{\oe i}{s}} \dicPos{m}[4] \dicFlx{(‑s, ‑ar)}[14] \textbf{1.} \dicDirectTranslationCS{slepičí hlava}  \textbf{2.} \dicIndirectTranslationCS{člověk, který se opije po malém množství alkoholu}
\dicEntry[hænuungi] \dicTerm{hænu··ung|i} \dicIPA{{h}{a}{i}{\textlengthmark}{n}{\textscy}{u}{\textltailn}{\r{\textObardotlessj}}{\textsci}} \dicPos{m}[1] \dicFlx{(‑a, ‑ar)}[1] \dicSynonym{kjúklingur} \dicDirectTranslationCS{kuře}
\dicEntry[hæpinn] \dicTerm{hæpinn} \dicIPA{{h}{a}{i}{\textlengthmark}{\textsubring{b}}{\textsci}{\textsubring{n}}} \dicPos{adj}[6]\dicFlx{}[-2] \dicSynonym{vafasamur} \dicDirectTranslationCS{pochybný, problematický, sporný}
\dicEntry[hæra] \dicTerm{hær|a} \dicIPA{{h}{a}{i}{\textlengthmark}{r}{a}} \dicPos{f}[1] \dicFlx{(‑u, ‑ur)}[7] \textbf{1.} \dicDirectTranslationCS{šedivé vlasy, šediny};  \dicPhraseIS{hvítur fyrir hærum} \dicDirectTranslationCS{(jsoucí) s~prošedivělými vlasy};  \dicPhraseIS{kemba ekki hærurnar} \dicLangCat{přen.} \dicDirectTranslationCS{nestačit zešedivět} \dicIndirectTranslationCS{(umřít mladý)}  \textbf{2.} \dicFieldCat{bot.} \dicDirectTranslationCS{bika} \textit{(l.~{\textLA{Luzula}})}  \dicsymPhoto\ 
\dicFigure{ds_image_haera_0_1.jpg}{Hæra}{Hæra - Sten Porse, CC BY-SA 3.0}
\dicEntry[hærast] \dicTerm{hær|ast} \dicIPA{{h}{a}{i}{\textlengthmark}{r}{a}{s}{\textsubring{d}}} \dicPos{v}[2] \dicFlx{(‑ðist, ‑st)}[200] \dicFlx{refl} \dicSynonym*{verða gráhærður} \dicDirectTranslationCS{zešedivět} \dicIndirectTranslationCS{(o~vlasech)}
\dicEntry[hærður] \dicTerm{hærður} \dicIPA{{h}{a}{i}{r}{ð}{\textscy}{\textsubring{r}}} \dicPos{adj}[2]\dicFlx{}[-1] \dicDirectTranslationCS{vlasatý, srstnatý, chlupatý}
\dicEntry[hærra] \dicTerm{hærra} \dicIPA{{h}{a}{i}{r}{\textlengthmark}{a}} \dicPos{adv} \dicFlx{comp} \dicLink{hátt}
\dicEntry[hærri] \dicTerm{hærri} \dicIPA{{h}{a}{i}{r}{\textlengthmark}{\textsci}} \dicPos{adj} \dicFlx{comp m} \dicLink{hár\smash{\textsuperscript{2}}}
\dicEntry[hæsi] \dicTerm{hæs|i} \dicIPA{{h}{a}{i}{\textlengthmark}{s}{\textsci}} \dicPos{f}[3] \dicFlx{(‑i)}[3] \dicDirectTranslationCS{chrapot, chraptivost, ochraptělost}
\dicEntry[hæst] \dicTerm{hæst} \dicIPA{{h}{a}{i}{s}{\textsubring{d}}} \dicPos{adv} \dicFlx{sup} \dicLink{hátt}
\dicEntry[hæstaréttardómari] \dicTerm{hæsta·réttar··dóm·ar|i} \dicIPA{{h}\-{a}\-{i}\-{s}\-{\textsubring{d}}\-{a}\-{r}\-{j}\-{\textepsilon}\-{h}\-{\textsubring{d}}\-{a}\-{r}\-{\textsubring{d}}\-{ou}\-{m}\-{a}\-{r}\-{\textsci}\-} \dicPos{m}[1] \dicFlx{(‑a, ‑ar)}[13] \dicFieldCat{práv.} \dicDirectTranslationCS{soudce\,/\addthin soudkyně Nejvyššího soudu}
\begin{xtolerant}{}{1pt}
\dicEntry[hæstaréttardómur] \dicTerm{hæsta·réttar··dóm|ur}\addthinS\dicIPA{{h}{a}{i}{s}{\textsubring{d}}{a}{r}{j}{\textepsilon}{h}{\textsubring{d}}{a}{r}{\textsubring{d}}{ou}{m}{\textscy}{\textsubring{r}}}\addthinS\dicPos{m}[6]\addthinS\dicFlx{(‑s,\addthinSS‑ar)}[10] \dicFieldCat{práv.} \dicDirectTranslationCS{výrok Nejvyššího soudu}
\end{xtolerant}
\begin{xtolerant}{}{1pt}
\dicEntry[hæstaréttarlögmaður] \dicTerm{hæsta·réttar··lög·|maður}\addthinSS\dicIPA{{h}\-{a}\-{i}\-{s}\-{\textsubring{d}}\-{a}\-{r}\-{j}\-{\textepsilon}\-{h}\-{\textsubring{d}}\-{a}\-{r}\-{l}\-{\oe}\-{\textbabygamma}\-{m}\-{a}\-{ð}\-{\textscy}\-{\textsubring{r}}\-}\addthinS\dicPos{m}[13]\addthinS\dicFlx{(‑manns, ‑menn)}[2] \dicFieldCat{práv.} \dicDirectTranslationCS{právník\,/\addthinSS právnička Nejvyššího sou\-du}
\end{xtolerant}
\dicEntry[hæstbjóðandi] \dicTerm{hæst··bjóð·|andi} \dicIPA{{h}{a}{i}{s}{\textsubring{d}}{\textsubring{b}}{j}{ou}{ð}{a}{n}{\textsubring{d}}{\textsci}} \dicPos{m}[2] \dicFlx{(‑anda, ‑endur)}[1] \dicDirectTranslationCS{dražitel(ka) s~nejvyšší nabídkou}
\dicEntry[hæstiréttur] \dicTerm{hæsti··réttur} \dicIPA{{h}{a}{i}{s}{\textsubring{d}}{\textsci}{r}{j}{\textepsilon}{h}{\textsubring{d}}{\textscy}{\textsubring{r}}} \dicPos{m}[10] \dicFlx{(hæstaréttar, hæsturéttir)}[5] \dicFieldCat{práv.} \dicDirectTranslationCS{Nejvyšší soud}
\dicEntry[hæstur] \dicTerm{hæstur} \dicIPA{{h}{a}{i}{s}{\textsubring{d}}{\textscy}{\textsubring{r}}} \dicPos{adj} \dicFlx{m sg nom sup} \dicLink{hár\smash{\textsuperscript{2}}}
\dicEntry[hæstvirtur] \dicTerm{hæst··virtur} \dicIPA{{h}{a}{i}{s}{\textsubring{d}}{v}{\textsci}{\textsubring{r}}{\textsubring{d}}{\textscy}{\textsubring{r}}} \dicPos{adj}[1]\dicFlx{}[-13] \dicDirectTranslationCS{vážený} \dicIndirectTranslationCS{(zdvořilostní oslovení)} \dicExampleIS{hæstvirtur forsætisráðherra} \dicExampleCS{vážený pane premiére}
\dicEntry[hætta] \dicTerm{hætt|a\smash{\textsuperscript{1}}} \dicsymFrequent\  \dicIPA{{h}{a}{i}{h}{\textsubring{d}}{a}} \dicPos{f}[1] \dicFlx{(‑u, ‑ur)}[7] \dicSynonym{háski} \dicDirectTranslationCS{nebezpečí, riziko} \dicExampleIS{hætta á e‑u} \dicExampleCS{nebezpečí (čeho)};  \dicPhraseIS{bjóða hættunni heim} \dicLangCat{přen.} \dicDirectTranslationCS{pokoušet osud};  \dicPhraseIS{e‑að dregur úr hættu á e‑u} \dicDirectTranslationCS{(co) minimalizuje riziko (čeho)};  \dicPhraseIS{tefla á tvær hættur} \dicLangCat{přen.} \dicDirectTranslationCS{riskovat, hazardovat};  \dicPhraseIS{vera í hættu} \dicDirectTranslationCS{být v~nebezpečí};  \dicPhraseIS{vera úr hættu} \dicDirectTranslationCS{být mimo nebezpečí};  \dicPhraseIS{það er hætta á ferðum} \dicDirectTranslationCS{na cestách se vyskytuje nebezpečí};  \dicPhraseIS{það er hætta á því að} \dicDirectTranslationCS{je nebezpečí, že}
\dicEntry[hætta] \dicTerm{hætt|a\smash{\textsuperscript{2}}} \dicsymFrequent\  \dicIPA{{h}{a}{i}{h}{\textsubring{d}}{a}} \dicPos{v}[2] \dicFlx{(‑i, ‑)}[7] \dicFlx{dat} \textbf{1.} \dicSynonym{linna} \dicDirectTranslationCS{(s)končit, ukončit, přestat, přestávat};  \dicPhraseIS{hætta e‑u} \dicDirectTranslationCS{přestat s~(čím), skončit s~(čím)} \dicExampleIS{hætta reykingum} \dicExampleCS{přestat s~kouřením}  \textbf{2.} \dicDirectTranslationCS{riskovat, dávat v~sázku};  \dicIdiom{hætta}[á]{ \dicPhraseIS{hætta á e‑ð}} \dicDirectTranslationCS{riskovat (co)};  \dicIdiom{hætta}[til]{ \dicPhraseIS{e‑m hættir til að (detta)}} \dicFlx{impers} \dicDirectTranslationCS{(kdo) má tendenci k~(padání)};  \dicIdiom{hætta}[við]{ \dicPhraseIS{hætta við e‑ð}} \dicDirectTranslationCS{ukončit (co), zanechat (čeho)}
\dicEntry[hættir] \dicTerm{hættir} \dicIPA{{h}{a}{i}{h}{\textsubring{d}}{\textsci}{\textsubring{r}}} \dicPos{m} \dicFlx{pl nom} \dicLink{háttur\smash{\textsuperscript{1}}}
\dicEntry[hættulaus] \dicTerm{hættu··laus} \dicIPA{{h}{a}{i}{h}{\textsubring{d}}{\textscy}{l}{\oe i}{s}} \dicPos{adj}[5]\dicFlx{}[-1] \dicSynonym*{háskalaus} \dicDirectTranslationCS{bezpečný, (jsoucí) bez nebezpečí\,/\addthin rizika}
\dicEntry[hættulegur] \dicTerm{hættu··legur} \dicsymFrequent\  \dicIPA{{h}{a}{i}{h}{\textsubring{d}}{\textscy}{l}{\textepsilon}{\textbabygamma}{\textscy}{\textsubring{r}}} \dicPos{adj}[1]\dicFlx{}[-8] \dicSynonym{glæfralegur} \dicDirectTranslationCS{nebezpečný, riskantní} \dicExampleIS{hættulegur maður} \dicExampleCS{nebezpečný člověk}
\dicEntry[hættur] \dicTerm{hættur} \dicIPA{{h}{a}{i}{h}{\textsubring{d}}{\textscy}{\textsubring{r}}} \dicPos{adj}[1]\dicFlx{}[-10] \dicDirectTranslationCS{skončený, ukončený} \dicExampleIS{vera hættur e‑u} \dicExampleCS{skončit s~(čím)};  \dicPhraseIS{e‑r er hættur að (gera e‑ð)} \dicDirectTranslationCS{(kdo) přestal (dělat (co))} \dicExampleIS{Hún er hætt að reykja.} \dicExampleCS{Přestala kouřit.}
\dicEntry[hættuspil] \dicTerm{hættu··spil} \dicIPA{{h}{a}{i}{h}{\textsubring{d}}{\textscy}{s}{\textsubring{b}}{\textsci}{\textsubring{l}}} \dicPos{n}[2] \dicFlx{(‑s, ‑)}[5] \dicDirectTranslationCS{hazard, hazardní hra}
\dicEntry[hættusvæði] \dicTerm{hættu··svæði} \dicIPA{{h}{a}{i}{h}{\textsubring{d}}{\textscy}{s}{v}{a}{i}{ð}{\textsci}} \dicPos{n}[2] \dicFlx{(‑s, ‑)}[14] \dicDirectTranslationCS{nebezpečná oblast}
\dicEntry[hæverska] \dicTerm{hæ··versk|a} \dicIPA{{h}{a}{i}{\textlengthmark}{v}{\textepsilon}{\textsubring{r}}{s}{\r{g}}{a}} \dicPos{f}[1] \dicFlx{(‑u)}[5] \dicDirectTranslationCS{uctivost, zdvořilost, slušnost}
\dicEntry[hæverskur] \dicTerm{hæ··verskur} \dicIPA{{h}{a}{i}{\textlengthmark}{v}{\textepsilon}{\textsubring{r}}{s}{\r{g}}{\textscy}{\textsubring{r}}} \dicPos{adj}[1]\dicFlx{}[-1] \dicSynonym{kurteis} \dicDirectTranslationCS{uctivý, zdvořilý, slušný}
\dicEntry[höf] \dicTerm{höf} \dicIPA{{h}{\oe}{\textlengthmark}{f}} \dicPos{n} \dicFlx{pl nom} \dicLink{haf}
\dicEntry[höf.] \dicTerm{höf.} \dicPos{zkr} \dicPhraseIS{höfundur} \dicDirectTranslationCS{autor(ka)}
\dicEntry[höfða] \dicTerm{höfð|a} \dicIPA{{h}{\oe}{v}{ð}{a}} \dicPos{v}[1] \dicFlx{(‑aði)}[1] \dicPhraseIS{höfða mál} \dicFieldCat{práv.} \dicDirectTranslationCS{podat žalobu, zažalovat} \dicExampleIS{höfða mál gegn e‑m} \dicExampleCS{zažalovat (koho)};  \dicPhraseIS{höfða til e‑rs} \dicDirectTranslationCS{oslovit\,/\addthin oslovovat (koho), obrátit\,/\addthin obracet se na (koho), apelovat na (koho)} \dicExampleIS{Bókin höfðar til yngri lesenda.} \dicExampleCS{Knížka se obrací na mladé čtenáře.}
\dicEntry[Höfðaborg] \dicTerm{Höfða··borg} \dicIPA{{h}{\oe}{v}{ð}{a}{\textsubring{b}}{\textopeno}{r}{\r{g}}} \dicPos{f}[7] \dicFlx{(‑ar)}[4] \dicFieldCat{geog.} \dicDirectTranslationCS{Kapské Město} \dicIndirectTranslationCS{(město v~Jihoafrické republice)}
\dicEntry[höfðalag] \dicTerm{höfða··|lag} \dicIPA{{h}{\oe}{v}{ð}{a}{l}{a}{x}} \dicPos{n}[2] \dicFlx{(‑lags, ‑lög)}[8] \dicDirectTranslationCS{čelo  postele}
\dicEntry[höfðatala] \dicTerm{höfða··|tala} \dicIPA{{h}{\oe}{v}{ð}{a}{t\smash{\textsuperscript{h}}}{a}{l}{a}} \dicPos{f}[1] \dicFlx{(‑tölu)}[2] \dicDirectTranslationCS{populace, počet obyvatel};  \dicPhraseIS{miðað við höfðatölu} \dicFlx{adv} \dicDirectTranslationCS{v~poměru k~počtu obyvatel}
\dicEntry[höfði] \dicTerm{höfð|i} \dicIPA{{h}{\oe}{v}{ð}{\textsci}} \dicPos{m}[1] \dicFlx{(‑a, ‑ar)}[1] \dicDirectTranslationCS{mys}
\dicEntry[höfðingi] \dicTerm{höfð··ing|i} \dicIPA{{h}{\oe}{v}{ð}{i}{\textltailn}{\r{\textObardotlessj}}{\textsci}} \dicPos{m}[1] \dicFlx{(‑ja, ‑jar)}[14] \textbf{1.} \dicSynonym{foringi} \dicDirectTranslationCS{vůdce, vůdkyně, náčelník, náčelnice}  \textbf{2.} \dicSynonym{valdamaður} \dicDirectTranslationCS{panovník, panovnice, vladař(ka), mocnář(ka)}  \textbf{3.} \dicSynonym*{örlátur maður} \dicDirectTranslationCS{velkorysý člověk}
\dicEntry[höfðingjasleikja] \dicTerm{höfðingja··sleikj|a} \dicIPA{{h}{\oe}{v}{ð}{i}{\textltailn}{\r{\textObardotlessj}}{a}{s}{\textsubring{d}}{l}{ei}{\r{\textObardotlessj}}{a}} \dicPos{f}[1] \dicFlx{(‑u, ‑ur)}[7] \dicDirectTranslationCS{patolízal(ka), pochlebník, pochlebnice}
\dicEntry[höfðinglegur] \dicTerm{höfðing··legur} \dicIPA{{h}{\oe}{v}{ð}{i}{\ng}{\r{g}}{l}{\textepsilon}{\textbabygamma}{\textscy}{\textsubring{r}}} \dicPos{adj}[1]\dicFlx{}[-8] \textbf{1.} \dicSynonym{göfugur} \dicDirectTranslationCS{velkodušný, velkorysý}  \textbf{2.} \dicSynonym{rausnarlegur} \dicDirectTranslationCS{královský, skvostný}
\dicEntry[höfðingsskapur] \dicTerm{höfðings··skap|ur} \dicIPA{{h}{\oe}{v}{ð}{i}{\ng}{s}{\r{g}}{a}{\textsubring{b}}{\textscy}{\textsubring{r}}} \dicPos{m}[10] \dicFlx{(‑ar)}[15] \dicSynonym{göfuglyndi} \dicDirectTranslationCS{velkorysost, velkodušnost}
\dicEntry[höfðum] \dicTerm{höfðum} \dicIPA{{h}{\oe}{v}{ð}{\textscy}{\textsubring{m}}} \dicPos{v} \dicFlx{ind pf pl 1 pers} \dicLink{hafa}
\dicEntry[höfgi] \dicTerm{höfg|i} \dicIPA{{h}{\oe}{v}{\r{\textObardotlessj}}{\textsci}} \dicPos{m}[1] \dicFlx{(‑a)}[3] \textbf{1.} \dicSynonym*{svefnmók} \dicDirectTranslationCS{ospalost} \dicExampleIS{Það seig á hann höfgi.} \dicExampleCS{Přepadla ho ospalost.}  \textbf{2.} \dicSynonym{þungi} \dicDirectTranslationCS{tíha, tíže}
\dicEntry[höfn] \dicTerm{höfn} \dicsymFrequent\  \dicIPA{{h}{\oe}{\textsubring{b}}{\textsubring{n}}} \dicPos{f}[7] \dicFlx{(hafnar, hafnir)}[16] \dicDirectTranslationCS{přístav, přístaviště} \dicExampleIS{starfsmenn hafnarinnar} \dicExampleCS{pracovníci přístavu};  \dicIdiom{höfn}{ \dicPhraseIS{e‑að er í höfn}} \dicLangCat{přen.} \dicDirectTranslationCS{(co) je doma} \dicIndirectTranslationCS{(co je zdárně ukončeno)};  \dicPhraseIS{koma e‑u í höfn} \dicLangCat{přen.} \dicDirectTranslationCS{(zdárně) (co) ukončit};  \dicPhraseIS{leysast höfn} \dicFlx{refl} \dicDirectTranslationCS{potratit}
\dicEntry[höfnun] \dicTerm{höfnun} \dicIPA{{h}{\oe}{\textsubring{b}}{n}{\textscy}{\textsubring{n}}} \dicPos{f}[7] \dicFlx{(‑ar, hafnanir)}[16] \dicDirectTranslationCS{odmítnutí, zamítnutí, odřeknutí}
\dicEntry[höfrungur] \dicTerm{höfr··ung|ur} \dicIPA{{h}{\oe}{v}{r}{u}{\ng}{\r{g}}{\textscy}{\textsubring{r}}} \dicPos{m}[6] \dicFlx{(‑s, ‑ar)}[8] \dicFieldCat{zool.} \dicSynonym{léttir} \dicDirectTranslationCS{delfín, delfín obecný} \textit{(l.~{\textLA{Delphinus delphis}})}
\dicEntry[höft] \dicTerm{höft} \dicIPA{{h}{\oe}{f}{\textsubring{d}}} \dicPos{n} \dicFlx{pl nom} \dicLink{haft\smash{\textsuperscript{1}}}
\dicEntry[höfuð] \dicTerm{höfuð\smash{\textsuperscript{1}}} \dicIPA{{h}{\oe}{\textlengthmark}{v}{\textscy}{\texttheta}} \dicPos{n}[2] \dicFlx{(‑s, ‑)}[37] \dicSynonym{kollur} \dicDirectTranslationCS{hlava};  \dicPhraseIS{bera hönd fyrir höfuð sér} \dicDirectTranslationCS{bránit se, hájit se};  \dicPhraseIS{bíta höfuðið af skömminni} \dicLangCat{přen.} \dicDirectTranslationCS{mít z~ostudy kabát};  \dicPhraseIS{berja höfðinu við steininn} \dicLangCat{přen.} \dicDirectTranslationCS{jít hlavou proti zdi, bojovat s~větrnými mlýny};  \dicPhraseIS{e‑að vex e‑m yfir höfuð} \dicLangCat{přen.} \dicDirectTranslationCS{(co komu) přerůstá přes hlavu};  \dicPhraseIS{eiga hvergi höfði sínu að halla} \dicLangCat{přen.} \dicDirectTranslationCS{nemít kde složit hlavu};  \dicPhraseIS{fara á höfuðið} \dicDirectTranslationCS{zbankrotovat, přijít na mizinu};  \dicPhraseIS{fara huldu höfði} \dicDirectTranslationCS{zamaskovat se};  \dicPhraseIS{hitta naglann á höfuðið} \dicLangCat{přen.} \dicDirectTranslationCS{uhodit hřebíček na hlavičku};  \dicPhraseIS{láta undir höfuð leggjast að (gera e‑ð)} \dicLangCat{přen.} \dicDirectTranslationCS{pustit z~hlavy (dělání (čeho))};  \dicPhraseIS{leggja fé til höfuðs e‑m} \dicDirectTranslationCS{vypsat odměnu za hlavu (koho)};  \dicPhraseIS{stinga höfðinu í sandinn} \dicLangCat{přen.} \dicDirectTranslationCS{strkat hlavu do písku};  \dicPhraseIS{það stígur e‑m til höfuðs} \dicFlx{impers} \dicLangCat{přen.} \dicDirectTranslationCS{stoupá to (komu) do hlavy}
\dicEntry[höfuð] \dicTerm{höfuð-\smash{\textsuperscript{2}}} \dicIPA{{h}{\oe}{\textlengthmark}{v}{\textscy}{\texttheta}} \dicPos{predp} \dicDirectTranslationCS{hlavní}
\dicEntry[höfuðatriði] \dicTerm{höfuð··at·riði} \dicIPA{{h}{\oe}{\textlengthmark}{v}{\textscy}{ð}{a}{\textsubring{d}}{r}{\textsci}{ð}{\textsci}} \dicPos{n}[2] \dicFlx{(‑s, ‑)}[14] \dicDirectTranslationCS{hlavní věc\,/\addthin bod}
\dicEntry[höfuðborg] \dicTerm{höfuð··borg} \dicsymFrequent\  \dicIPA{{h}{\oe}{\textlengthmark}{v}{\textscy}{ð}{\textsubring{b}}{\textopeno}{r}{\r{g}}} \dicPos{f}[7] \dicFlx{(‑ar, ‑ir)}[1] \dicDirectTranslationCS{hlavní město} \dicExampleIS{fara til höfuðborga nágrannalanda} \dicExampleCS{cestovat do hlavních měst sousedních zemí}
\dicEntry[höfuðborgarsvæði] \dicTerm{höfuð·borgar··svæði} \dicIPA{{h}{\oe}{\textlengthmark}{v}{\textscy}{ð}{\textsubring{b}}{\textopeno}{r}{\r{g}}{a}{\textsubring{r}}{s}{v}{a}{i}{ð}{\textsci}} \dicPos{n}[2] \dicFlx{(‑s, ‑)}[14] \dicDirectTranslationCS{oblast hlavního města};  \dicPhraseIS{höfuðborgarsvæðið} \dicFlx{def} \dicDirectTranslationCS{oblast hlavního města Reykjavíku}
\dicEntry[höfuðbók] \dicTerm{höfuð··|bók} \dicIPA{{h}{\oe}{\textlengthmark}{v}{\textscy}{ð}{\textsubring{b}}{ou}{\r{g}}} \dicPos{f}[8] \dicFlx{(‑bókar, ‑bækur)}[5] \dicFieldCat{ekon.} \dicDirectTranslationCS{účetní kniha}
\dicEntry[höfuðfat] \dicTerm{höfuð··|fat} \dicIPA{{h}{\oe}{\textlengthmark}{v}{\textscy}{ð}{f}{a}{\textsubring{d}}} \dicPos{n}[2] \dicFlx{(‑fats, ‑föt)}[8] \dicDirectTranslationCS{pokrývka hlavy}
\dicEntry[höfuðkúpa] \dicTerm{höfuð··kúp|a} \dicIPA{{h}{\oe}{\textlengthmark}{v}{\textscy}{ð}{k\smash{\textsuperscript{h}}}{u}{\textsubring{b}}{a}} \dicPos{f}[1] \dicFlx{(‑u, ‑ur)}[19] \dicFieldCat{anat.} \dicDirectTranslationCS{lebka}
\dicEntry[höfuðleður] \dicTerm{höfuð··leður} \dicIPA{{h}{\oe}{\textlengthmark}{v}{\textscy}{ð}{l}{\textepsilon}{ð}{\textscy}{\textsubring{r}}} \dicPos{n}[2] \dicFlx{(‑s, ‑)}[25] \textbf{1.} \dicSynonym*{höfuðhúð} \dicDirectTranslationCS{skalp}  \textbf{2.} \dicSynonym*{beislishluti} \dicDirectTranslationCS{ohlávka}
\dicEntry[höfuðlús] \dicTerm{höfuð··|lús} \dicIPA{{h}{\oe}{\textlengthmark}{v}{\textscy}{ð}{l}{u}{s}} \dicPos{f}[9] \dicFlx{(‑lúsar, ‑lýs)}[6] \dicFieldCat{zool.} \dicDirectTranslationCS{veš, veš dětská} \textit{(l.~{\textLA{Pediculus capitis}})}
\dicEntry[höfuðmál] \dicTerm{höfuð··mál} \dicIPA{{h}{\oe}{\textlengthmark}{v}{\textscy}{ð}{m}{au}{\textsubring{l}}} \dicPos{n}[2] \dicFlx{(‑s, ‑)}[5] \dicPhraseIS{skipta ekki höfuðmáli} \dicLangCat{přen.} \dicDirectTranslationCS{nehrát velkou roli}
\dicEntry[höfuðpaur] \dicTerm{höfuð··paur} \dicIPA{{h}{\oe}{\textlengthmark}{v}{\textscy}{ð}{p\smash{\textsuperscript{h}}}{\oe i}{\textsubring{r}}} \dicPos{m}[4] \dicFlx{(‑s, ‑ar)}[14] \dicLangCat{han.} \dicDirectTranslationCS{boss, kápo}
\dicEntry[höfuðsmaður] \dicTerm{höfuðs··|maður} \dicIPA{{h}{\oe}{\textlengthmark}{v}{\textscy}{ð}{s}{m}{a}{ð}{\textscy}{\textsubring{r}}} \dicPos{m}[13] \dicFlx{(‑manns, ‑menn)}[2] \textbf{1.} \dicFieldCat{hist.} \dicDirectTranslationCS{místodržící, místokrál}  \textbf{2.} \dicFieldCat{voj.} \dicDirectTranslationCS{kapitán}
\dicEntry[höfuðstaður] \dicTerm{höfuð··stað|ur} \dicIPA{{h}{\oe}{\textlengthmark}{v}{\textscy}{ð}{s}{\textsubring{d}}{a}{ð}{\textscy}{\textsubring{r}}} \dicPos{m}[10] \dicFlx{(‑ar, ‑ir)}[14] \dicDirectTranslationCS{hlavní město}
\dicEntry[höfuðstafur] \dicTerm{höfuð··staf|ur} \dicIPA{{h}{\oe}{\textlengthmark}{v}{\textscy}{ð}{s}{\textsubring{d}}{a}{v}{\textscy}{\textsubring{r}}} \dicPos{m}[9] \dicFlx{(‑s, ‑ir)}[9] \dicFieldCat{lit.} \dicIndirectTranslationCS{aliterace na začátku  sudého verše}
\dicEntry[höfuðstóll] \dicTerm{höfuð··stól|l} \dicIPA{{h}{\oe}{\textlengthmark}{v}{\textscy}{ð}{s}{\textsubring{d}}{ou}{\textsubring{d}}{\textsubring{l}}} \dicPos{m}[6] \dicFlx{(‑s)}[50] \dicFieldCat{ekon.} \dicDirectTranslationCS{kapitál, jistina} \dicExampleIS{vextir af höfuðstól} \dicExampleCS{úroky z~jistiny}
\dicEntry[höfuðstöðvar] \dicTerm{höfuð··stöðvar} \dicIPA{{h}{\oe}{\textlengthmark}{v}{\textscy}{ð}{s}{\textsubring{d}}{\oe}{ð}{v}{a}{\textsubring{r}}} \dicPos{f}[6] \dicFlx{pl}[2] \dicDirectTranslationCS{ústředí, centrála, ředitelství}
\dicEntry[höfuðsynd] \dicTerm{höfuð··synd} \dicIPA{{h}{\oe}{\textlengthmark}{v}{\textscy}{ð}{s}{\textsci}{n}{\textsubring{d}}} \dicPos{f}[7] \dicFlx{(‑ar, ‑ir)}[1] \dicDirectTranslationCS{smrtelný hřích}
\dicEntry[höfuðverkur] \dicTerm{höfuð··verk|ur} \dicsymFrequent\  \dicIPA{{h}{\oe}{\textlengthmark}{v}{\textscy}{ð}{v}{\textepsilon}{\textsubring{r}}{\r{g}}{\textscy}{\textsubring{r}}} \dicPos{m}[10] \dicFlx{(‑jar, ‑ir)}[23] \dicDirectTranslationCS{bolest\,/\addthin bolení hlavy, migréna} \dicExampleIS{taka töflu við höfuðverknum} \dicExampleCS{vzít si prášek na bolest hlavy}
\dicEntry[höfundarlaun] \dicTerm{höfundar··laun} \dicIPA{{h}{\oe}{\textlengthmark}{v}{\textscy}{n}{\textsubring{d}}{a}{r}{l}{\oe i}{\textsubring{n}}} \dicPos{n}[2] \dicFlx{pl}[1] \dicDirectTranslationCS{honorář}
\dicEntry[höfundur] \dicTerm{höf··und|ur} \dicsymFrequent\  \dicIPA{{h}{\oe}{\textlengthmark}{v}{\textscy}{n}{\textsubring{d}}{\textscy}{\textsubring{r}}} \dicPos{m}[6] \dicFlx{(‑ar, ‑ar)}[56] \dicDirectTranslationCS{autor(ka), spisovatel(ka)} \dicExampleIS{höfundur að verkinu} \dicExampleCS{autor díla}
\dicEntry[hög] \dicTerm{hög} \dicIPA{{h}{\oe}{\textlengthmark}{x}} \dicPos{adj} \dicFlx{f sg nom pos} \dicLink{hagur\smash{\textsuperscript{2}}}
\dicEntry[högg] \dicTerm{högg} \dicsymFrequent\  \dicIPA{{h}{\oe}{\r{g}}{\textlengthmark}} \dicPos{n}[2] \dicFlx{(‑s, ‑)}[5] \dicSynonym{barsmíð} \dicDirectTranslationCS{rána, úder, uhození, bouchnutí} \dicExampleIS{Hann fékk högg í andlitið.} \dicExampleCS{Dostal úder do obličeje.};  \dicPhraseIS{eiga í höggi við e‑n} \dicDirectTranslationCS{bojovat s~(kým), být v~konfliktu s~(kým)};  \dicPhraseIS{eiga undir högg að sækja} \dicLangCat{přen.} \dicDirectTranslationCS{nacházet se v~prekérní situaci};  \dicPhraseIS{högg undir beltisstað} \dicLangCat{přen.} \dicDirectTranslationCS{rána pod pás};  \dicPhraseIS{liggja vel við höggi} \dicLangCat{přen.} \dicDirectTranslationCS{být\,/\addthin stát na ráně};  \dicPhraseIS{slá tvær flugur í einu höggi} \dicLangCat{přen.} \dicDirectTranslationCS{zabít dvě mouchy jednou ranou}
\dicEntry[höggdeyfir] \dicTerm{högg··deyf|ir} \dicIPA{{h}{\oe}{\r{g}}{\textsubring{d}}{ei}{v}{\textsci}{\textsubring{r}}} \dicPos{m}[7] \dicFlx{(‑is, ‑ar)}[1] \dicDirectTranslationCS{tlumič nárazů}
\dicEntry[höggdofa] \dicTerm{högg··dofa} \dicIPA{{h}{\oe}{\r{g}}{\textsubring{d}}{\textopeno}{v}{a}} \dicPos{adj}[13] \dicFlx{indecl}[1] \dicSynonym{undrandi} \dicDirectTranslationCS{udivený, ohromený, užaslý}
\dicEntry[höggið] \dicTerm{höggið} \dicIPA{{h}{\oe}{\r{\textObardotlessj}}{\textlengthmark}{\textsci}{\texttheta}} \dicPos{v} \dicFlx{supin} \dicLink{höggva}
\dicEntry[höggmynd] \dicTerm{högg··mynd} \dicIPA{{h}{\oe}{\r{g}}{m}{\textsci}{n}{\textsubring{d}}} \dicPos{f}[7] \dicFlx{(‑ar, ‑ir)}[1] \dicDirectTranslationCS{socha} \dicExampleIS{höggmynd af forsetanum} \dicExampleCS{socha prezidenta}
\dicEntry[höggmyndalist] \dicTerm{högg·mynda··list} \dicIPA{{h}{\oe}{\r{g}}{m}{\textsci}{n}{\textsubring{d}}{a}{l}{\textsci}{s}{\textsubring{d}}} \dicPos{f}[7] \dicFlx{(‑ar)}[3] \dicDirectTranslationCS{sochařství}
\dicEntry[höggormur] \dicTerm{högg··orm|ur} \dicIPA{{h}{\oe}{\r{g}}{\textlengthmark}{\textopeno}{r}{m}{\textscy}{\textsubring{r}}} \dicPos{m}[6] \dicFlx{(‑s, ‑ar)}[4] \dicFieldCat{zool.} \dicDirectTranslationCS{zmije, zmije obecná} \textit{(l.~{\textLA{Vipera berus}})}  \dicsymPhoto\ 
\dicFigure{ds_image_hoggormur_0_1.jpg}{Höggormur}{Höggormur - Zdeněk Fric, CC BY-SA 3.0}
\dicEntry[höggstaður] \dicTerm{högg··stað|ur} \dicIPA{{h}{\oe}{\r{g}}{s}{\textsubring{d}}{a}{ð}{\textscy}{\textsubring{r}}} \dicPos{m}[10] \dicFlx{(‑ar, ‑ir)}[14] \dicLangCat{přen.} \dicDirectTranslationCS{slabé místo};  \dicPhraseIS{gefa höggstað á sér} \dicLangCat{přen.} \dicDirectTranslationCS{ukázat své slabé místo};  \dicPhraseIS{leita höggstaðar á e‑m} \dicLangCat{přen.} \dicDirectTranslationCS{hledat (čí) slabé místo}
\dicEntry[höggva] \dicTerm{höggva} \dicsymFrequent\  \dicIPA{{h}{\oe}{\r{g}}{v}{a}} \dicPos{v}[6] \dicFlx{(hegg, hjó, hjuggum, hyggi\,/\addthin hjyggi, höggvið\,/\addthin höggið\,/\addthin hoggið)}[43] \dicFlx{acc} \textbf{1.} \dicDirectTranslationCS{(na)štípat, (na)sekat, tnout (dříví ap.)} \dicExampleIS{höggva tré með öxi} \dicExampleCS{štípat dříví sekyrou}  \textbf{2.} \dicDirectTranslationCS{prosekat, proklestit (cestu ap.)};  \dicIdiom{höggva}[eftir]{ \dicPhraseIS{höggva eftir e‑u}} \dicSynonym*{hvá eftir e‑u} \dicDirectTranslationCS{všimnout\,/\addthin všímat si (čeho)} \dicExampleIS{Ég hjó eftir því að þú elskir osta.} \dicExampleCS{Všiml jsem si, že miluješ sýry.};  \dicIdiom{höggvast}{ \dicPhraseIS{höggvast}} \dicFlx{refl} \dicDirectTranslationCS{sekat se, bít se}
\dicEntry[höggvopn] \dicTerm{högg··vopn} \dicIPA{{h}{\oe}{\r{g}}{v}{\textopeno}{h}{\textsubring{b}}{\textsubring{n}}} \dicPos{n}[2] \dicFlx{(‑s, ‑)}[5] \dicDirectTranslationCS{sečná zbraň}
\dicEntry[högl] \dicTerm{högl} \dicIPA{{h}{\oe}{\r{g}}{\textsubring{l}}} \dicPos{n} \dicFlx{pl nom} \dicLink{hagl}
\dicEntry[högni] \dicTerm{högn|i} \dicIPA{{h}{\oe}{\r{g}}{n}{\textsci}} \dicPos{m}[1] \dicFlx{(‑a, ‑ar)}[1] \dicSynonym{fress} \dicDirectTranslationCS{kocour}
\dicEntry[hökta] \dicTerm{hökt|a} \dicIPA{{h}{\oe}{x}{\textsubring{d}}{a}} \dicPos{v}[2] \dicFlx{(‑i, ‑)}[17] \textbf{1.} \dicSynonym{haltra} \dicDirectTranslationCS{kulhat, pokulhávat, pajdat}  \textbf{2.} \dicDirectTranslationCS{škytnout, škytat} \dicIndirectTranslationCS{(o~autě ap.)}
\dicEntry[höku] \dicTerm{höku} \dicIPA{{h}{\oe}{\textlengthmark}{\r{g}}{\textscy}} \dicPos{f} \dicFlx{sg gen} \dicLink{haka}
\dicEntry[hökuskarð] \dicTerm{höku··|skarð} \dicIPA{{h}{\oe}{\textlengthmark}{\r{g}}{\textscy}{s}{\r{g}}{a}{r}{\texttheta}} \dicPos{n}[2] \dicFlx{(‑skarðs, ‑skörð)}[8] \dicDirectTranslationCS{důlek} \dicIndirectTranslationCS{(na bradě)}
\dicEntry[hökutoppur] \dicTerm{höku··topp|ur} \dicIPA{{h}{\oe}{\textlengthmark}{\r{g}}{\textscy}{t\smash{\textsuperscript{h}}}{\textopeno}{h}{\textsubring{b}}{\textscy}{\textsubring{r}}} \dicPos{m}[6] \dicFlx{(‑s, ‑ar)}[10] \dicDirectTranslationCS{kozí bradka} \dicIndirectTranslationCS{(druh úpravy vousů)}
\dicEntry[höld] \dicTerm{höld} \dicIPA{{h}{\oe}{l}{\textsubring{d}}} \dicPos{n} \dicFlx{pl nom} \dicLink{hald}
\dicEntry[höldu] \dicTerm{höldu} \dicIPA{{h}{\oe}{l}{\textsubring{d}}{\textscy}} \dicPos{f} \dicFlx{sg gen} \dicLink{halda\smash{\textsuperscript{1}}}
\dicEntry[höll] \dicTerm{höll\smash{\textsuperscript{1}}} \dicIPA{{h}{\oe}{\textsubring{d}}{\textsubring{l}}} \dicPos{f}[7] \dicFlx{(hallar, hallir)}[22] \textbf{1.} \dicSynonym*{konungshöll} \dicDirectTranslationCS{zámek, palác}  \textbf{2.} \dicLangCat{hovor.} \dicSynonym*{lúxusvilla} \dicDirectTranslationCS{přepychová vila, palác}
\dicEntry[höll] \dicTerm{höll\smash{\textsuperscript{2}}} \dicIPA{{h}{\oe}{\textsubring{d}}{\textsubring{l}}} \dicPos{adj} \dicFlx{f sg nom pos} \dicLink{hallur}
\dicEntry[hölt] \dicTerm{hölt} \dicIPA{{h}{\oe}{\textsubring{l}}{\textsubring{d}}} \dicPos{adj} \dicFlx{f sg nom pos} \dicLink{haltur}
\dicEntry[hömdum] \dicTerm{hömdum} \dicIPA{{h}{\oe}{m}{\textsubring{d}}{\textscy}{\textsubring{m}}} \dicPos{v} \dicFlx{ind pf pl 1 pers} \dicLink{hemja\smash{\textsuperscript{2}}}
\dicEntry[hömlu] \dicTerm{hömlu} \dicIPA{{h}{\oe}{m}{l}{\textscy}} \dicPos{f} \dicFlx{sg gen} \dicLink{hamla\smash{\textsuperscript{1}}}
\dicEntry[hömlulaus] \dicTerm{hömlu··laus} \dicIPA{{h}{\oe}{m}{l}{\textscy}{l}{\oe i}{s}} \dicPos{adj}[5]\dicFlx{}[-1] \dicSynonym{gegndarlaus} \dicDirectTranslationCS{nespoutaný, nevázaný, (jsoucí) bez zábran}
\dicEntry[hömlun] \dicTerm{hömlun} \dicIPA{{h}{\oe}{m}{l}{\textscy}{\textsubring{n}}} \dicPos{f}[7] \dicFlx{(‑ar, hamlanir)}[16] \dicSynonym*{afstýring} \dicDirectTranslationCS{zábrana, inhibice}
\dicEntry[hönd] \dicTerm{hönd} \dicsymFrequent\  \dicIPA{{h}{\oe}{n}{\textsubring{d}}} \dicPos{f}[8] \dicFlx{(handar, hendur)}[4] \textbf{1.} \dicFieldCat{anat.} \dicDirectTranslationCS{ruka} \dicIndirectTranslationCS{(dolní část horní končetiny)};  \dicPhraseIS{e‑ir haldast í hendur} \dicFlx{refl} \dicDirectTranslationCS{(kdo) se drží za ruce};  \dicPhraseIS{rétta e‑m höndina} \dicDirectTranslationCS{podat (komu) ruku};  \dicPhraseIS{taka í höndina á e‑m} \dicDirectTranslationCS{stisknout (čí) ruku}  \textbf{2.} \dicSynonym{handleggur} \dicDirectTranslationCS{ruka, paže} \dicIndirectTranslationCS{(celá horní končetina)};  \dicPhraseIS{baða út höndunum} \dicDirectTranslationCS{mávat rukama};  \dicPhraseIS{rétta upp hönd\,/\addthin höndina} \dicDirectTranslationCS{natáhnout ruku}  \textbf{3.} \dicSynonym{rithönd} \dicDirectTranslationCS{rukopis, písmo} \dicIndirectTranslationCS{(způsob psaní charakteristický pro jedince)};  \dicPhraseIS{e‑r hefur fallega hönd} \dicDirectTranslationCS{(kdo) má pěkný rukopis}  \textbf{4.} \dicSynonym*{hlið} \dicDirectTranslationCS{strana, ruka};  \dicPhraseIS{á hægri\,/\addthin vinstri hönd} \dicFlx{adv} \dicDirectTranslationCS{po pravé\,/\addthin levé ruce, na\,/\addthin po pravé\,/\addthin levé straně};  \dicIdiom{hönd}{ \dicPhraseIS{bera e‑n á höndum sér}} \dicLangCat{přen.} \dicDirectTranslationCS{nosit (koho) na rukou}; { \dicPhraseIS{bera hönd fyrir höfuð sér}} \dicDirectTranslationCS{bránit se, hájit se};  \dicPhraseIS{binda hendur e‑rs} \dicLangCat{přen.} \dicDirectTranslationCS{svázat (komu) ruce};  \dicPhraseIS{e‑ð ber að höndum} \dicFlx{impers} \dicDirectTranslationCS{(co) vyvstává (problém ap.)};  \dicPhraseIS{e‑að fer í hönd} \dicLangCat{přen.} \dicDirectTranslationCS{(co) je za dveřmi (Vánoce ap.)};  \dicPhraseIS{e‑að helst í hendur} \dicFlx{refl} \dicLangCat{přen.} \dicDirectTranslationCS{(co) jde ruku v~ruce};  \dicPhraseIS{e‑að snýst í höndunum á e‑m} \dicFlx{refl} \dicLangCat{přen.} \dicDirectTranslationCS{(co) se obrací proti (komu)};  \dicPhraseIS{eiga e‑ð fyrir höndum} \dicLangCat{přen.} \dicDirectTranslationCS{mít (co) před sebou (dlouhou cestu ap.)};  \dicPhraseIS{e‑ir láta hendur skipta} \dicLangCat{přen.} \dicDirectTranslationCS{(kdo) nechá promluvit ruce};  \dicPhraseIS{e‑m ferst e‑að vel úr hendi} \dicFlx{impers} \dicLangCat{přen.} \dicDirectTranslationCS{(komu) jde (co) pěkně od ruky};  \dicPhraseIS{fá e‑ð í hendur} \dicDirectTranslationCS{dostat (co) do rukou (dokumenty ap.)};  \dicPhraseIS{fegins hendi} \dicFlx{adv} \dicDirectTranslationCS{s~otevřenou náručí, s~radostí};  \dicPhraseIS{geta ekki\,/\addthin hvorki hreyft hönd né fót} \dicDirectTranslationCS{nemoci se ani trochu pohnout};  \dicPhraseIS{glopra e‑u út úr höndum sér} \dicLangCat{přen.} \dicDirectTranslationCS{nechat si (co) proklouznout mezi prsty};  \dicPhraseIS{hafa alla þræði í hendi sér} \dicLangCat{přen.} \dicDirectTranslationCS{mít vše pod kontrolou, tahat za všechny nitky};  \dicPhraseIS{hafa e‑ð með höndum} \dicDirectTranslationCS{dohlížet na (co)};  \dicPhraseIS{hafa e‑ð undir höndum} \dicDirectTranslationCS{mít (co) v~opatrování};  \dicPhraseIS{hafa e‑ð við höndina} \dicDirectTranslationCS{mít (co) při ruce};  \dicPhraseIS{hafa frjálsar hendur} \dicLangCat{přen.} \dicDirectTranslationCS{mít volné ruce};  \dicPhraseIS{hafa hendur í hári e‑rs} \dicDirectTranslationCS{dostat se (komu) na kobylku, chytit (koho) (zloděje ap.)};  \dicPhraseIS{hafa hraðar hendur} \dicDirectTranslationCS{pospíšit si};  \dicPhraseIS{hefjast handa} \dicFlx{refl} \dicDirectTranslationCS{začít};  \dicPhraseIS{kasta höndunum til e‑s} \dicDirectTranslationCS{vrhnout se na (co)};  \dicPhraseIS{láta e‑ð af hendi} \dicDirectTranslationCS{dát (co) z~ruky};  \dicPhraseIS{láta hendur standa fram úr ermum} \dicLangCat{přen.} \dicDirectTranslationCS{vykasat si rukávy} \dicIndirectTranslationCS{(pustit se do něčeho rázně)};  \dicPhraseIS{leggja hendur á e‑n} \dicDirectTranslationCS{vztáhnout na (koho) ruku};  \dicPhraseIS{leggja hönd á plóginn} \dicLangCat{přen.} \dicDirectTranslationCS{přiložit ruku k~dílu};  \dicPhraseIS{leysa e‑ð af hendi} \dicDirectTranslationCS{hravě (co) vyřešit, vyřešit (co) raz dva};  \dicPhraseIS{segja e‑m stríð á hendur} \dicDirectTranslationCS{vyhlásit (komu) válku};  \dicPhraseIS{sitja auðum höndum} \dicLangCat{přen.} \dicDirectTranslationCS{sedět s~prázdnýma rukama};  \dicPhraseIS{sjá sig um hönd} \dicDirectTranslationCS{rozmyslet se, změnit názor};  \dicPhraseIS{(skrifa undir) fyrir hönd e‑rs} \dicDirectTranslationCS{(podepsat se) jménem (koho)};  \dicPhraseIS{(stjórna) með harðri hendi} \dicDirectTranslationCS{(vládnout) tvrdou rukou};  \dicPhraseIS{taka e‑ð í sínar hendur} \dicLangCat{přen.} \dicDirectTranslationCS{vzít (co) do svých rukou};  \dicPhraseIS{taka e‑ð ófrjálsri hendi} \dicDirectTranslationCS{ukrást (co)};  \dicPhraseIS{taka e‑m tveim höndum} \dicDirectTranslationCS{přijmout (koho) s~otevřenou náručí};  \dicPhraseIS{taka fram fyrir hendurnar á e‑m} \dicLangCat{přen.} \dicDirectTranslationCS{rozhodovat za (koho), lézt (komu) do zelí};  \dicPhraseIS{taka sér e‑ð fyrir hendur} \dicDirectTranslationCS{nalézt si (co) na práci};  \dicPhraseIS{taka til hendi\,/\addthin hendinni} \dicDirectTranslationCS{pustit se do práce};  \dicPhraseIS{vera e‑m innan handar} \dicDirectTranslationCS{být (komu) po ruce, být (komu) k~dispozici};  \dicPhraseIS{vera hægri hönd e‑rs} \dicLangCat{přen.} \dicDirectTranslationCS{být pravou rukou (koho)};  \dicPhraseIS{það er allt í hers höndum} \dicLangCat{přen.} \dicDirectTranslationCS{všechno je vzhůru nohama};  \dicPhraseIS{það eru (tveir kostir) fyrir hendi} \dicDirectTranslationCS{na výběr jsou (dvě možnosti)};  \dicPhraseIS{það leikur allt í höndunum á e‑m} \dicLangCat{přen.} \dicDirectTranslationCS{všechno (komu) hraje do karet};  \dicPhraseIS{þvo hendur sínar} \dicLangCat{přen.} \dicDirectTranslationCS{umýt si ruce}
\dicEntry[höndla] \dicTerm{höndl|a} \dicIPA{{h}{\oe}{n}{\textsubring{d}}{l}{a}} \dicPos{v}[1] \dicFlx{(‑aði)}[1] \dicFlx{acc} \textbf{1.} \dicDirectTranslationCS{obchodovat, handlovat}  \textbf{2.} \dicDirectTranslationCS{lapit, popadnout} \dicExampleIS{höndla hamingjuna} \dicExampleCS{popadnout štěstí}  \textbf{3.} \dicLangCat{hovor.} \dicDirectTranslationCS{zvládnout, poradit si}
\dicEntry[hönk] \dicTerm{hönk} \dicIPA{{h}{\oe i}{\r{\ng}}{\r{g}}} \dicPos{f}[7] \dicFlx{(hankar, hankir)}[16] \dicSynonym{lykkja} \dicDirectTranslationCS{smyčka};  \dicPhraseIS{það er allt í hönk} \dicLangCat{přen.} \dicDirectTranslationCS{je to blázinec, všechno je vzhůru nohama};  \dicPhraseIS{eiga hönk upp í bakið á e‑m} \dicLangCat{přen.} \dicDirectTranslationCS{mít to u~(koho) (protislužbu ap.)}
\dicEntry[hönnuður] \dicTerm{hönn··uð|ur} \dicIPA{{h}{\oe}{n}{\textlengthmark}{\textscy}{ð}{\textscy}{\textsubring{r}}} \dicPos{m}[10] \dicFlx{(‑ar, ‑ir)}[4] \dicDirectTranslationCS{výtvarník, výtvarnice, designér(ka), návrhář(ka)}
\dicEntry[hönnun] \dicTerm{hönn|un} \dicIPA{{h}{\oe}{n}{\textlengthmark}{\textscy}{\textsubring{n}}} \dicPos{f}[7] \dicFlx{(‑unar)}[12] \textbf{1.} \dicDirectTranslationCS{výtvarný návrh\,/\addthin projekt, design} \dicExampleIS{hönnun byggingarinnar} \dicExampleCS{výtvarný návrh budovy}  \textbf{2.} \dicDirectTranslationCS{návrhářství, design} \dicIndirectTranslationCS{(předmět ve škole)}
\dicEntry[höpp] \dicTerm{höpp} \dicIPA{{h}{\oe}{h}{\textsubring{b}}} \dicPos{n} \dicFlx{pl nom} \dicLink{happ}
\dicEntry[hör] \dicTerm{hör} \dicIPA{{h}{\oe}{\textlengthmark}{\textsubring{r}}} \dicPos{m}[4] \dicFlx{(‑s)}[18] \dicFieldCat{bot.} \dicDirectTranslationCS{len, len setý} \textit{(l.~{\textLA{Linum usitatissimum}})}
\dicEntry[hörð] \dicTerm{hörð} \dicIPA{{h}{\oe}{r}{\texttheta}} \dicPos{adj} \dicFlx{f sg nom pos} \dicLink{harður}
\dicEntry[hörðnun] \dicTerm{hörðn|un} \dicIPA{{h}{\oe}{r}{\textsubring{d}}{n}{\textscy}{\textsubring{n}}} \dicPos{f}[7] \dicFlx{(‑unar)}[12] \dicDirectTranslationCS{(z)tvrdnutí}
\dicEntry[hörfa] \dicTerm{hörf|a} \dicsymFrequent\  \dicIPA{{h}{\oe}{r}{v}{a}} \dicPos{v}[1] \dicFlx{(‑aði)}[34] \dicDirectTranslationCS{ustoupit, ustupovat, couvnout, couvat, stáhnout se} \dicExampleIS{hörfa aftur á bak} \dicExampleCS{ustoupit zpátky dozadu}
\dicEntry[hörgull] \dicTerm{hörg|ull} \dicIPA{{h}{\oe}{r}{\r{g}}{\textscy}{\textsubring{d}}{\textsubring{l}}} \dicPos{m}[6] \dicFlx{(‑uls)}[32] \dicSynonym{skortur} \dicDirectTranslationCS{nedostatek, nouze};  \dicPhraseIS{hörgull á e‑u} \dicDirectTranslationCS{nouze o~(co)};  \dicPhraseIS{út í hörgul} \dicFlx{adv} \dicDirectTranslationCS{důkladně, do nejmenšího detailu}
\dicEntry[hörku] \dicTerm{hörku\smash{\textsuperscript{1}}} \dicIPA{{h}{\oe}{\textsubring{r}}{\r{g}}{\textscy}} \dicPos{f} \dicFlx{sg gen} \dicLink{harka\smash{\textsuperscript{1}}}
\dicEntry[hörku] \dicTerm{hörku-\smash{\textsuperscript{2}}} \dicIPA{{h}{\oe}{\textsubring{r}}{\r{g}}{\textscy}} \dicPos{predp} \dicDirectTranslationCS{pořádný, drsný} \dicExampleIS{hörkufrost} \dicExampleCS{pořádný mráz}
\dicEntry[hörkulegur] \dicTerm{hörku··legur} \dicIPA{{h}{\oe}{\textsubring{r}}{\r{g}}{\textscy}{l}{\textepsilon}{\textbabygamma}{\textscy}{\textsubring{r}}} \dicPos{adj}[1]\dicFlx{}[-8] \dicDirectTranslationCS{drsný, hrozný}
\dicEntry[hörkutól] \dicTerm{hörku··tól} \dicIPA{{h}{\oe}{\textsubring{r}}{\r{g}}{\textscy}{t\smash{\textsuperscript{h}}}{ou}{\textsubring{l}}} \dicPos{n}[2] \dicFlx{(‑s, ‑)}[5] \dicDirectTranslationCS{drsňák, pořízek}
\dicEntry[hörkuvinna] \dicTerm{hörku··vinn|a} \dicIPA{{h}{\oe}{\textsubring{r}}{\r{g}}{\textscy}{v}{\textsci}{n}{a}} \dicPos{f}[1] \dicFlx{(‑u)}[5] \dicDirectTranslationCS{pořádná práce\,/\addthin dřina}
\dicEntry[hörmulegur] \dicTerm{hörmu··legur} \dicIPA{{h}{\oe}{r}{m}{\textscy}{l}{\textepsilon}{\textbabygamma}{\textscy}{\textsubring{r}}} \dicPos{adj}[1]\dicFlx{}[-8] \dicSynonym{sorglegur} \dicDirectTranslationCS{tragický, žalostný} \dicExampleIS{hörmulegt slys} \dicExampleCS{tragická nehoda}
\dicEntry[hörmung] \dicTerm{hörm··ung} \dicIPA{{h}{\oe}{r}{m}{u}{\ng}{\r{g}}} \dicPos{f}[4] \dicFlx{(‑ar, ‑ar)}[1] \textbf{1.} \dicSynonym*{lélegt verk} \dicDirectTranslationCS{fušeřina, fušerská práce}  \textbf{2.} \dicPhraseIS{hörmungar} \dicFlx{pl} \dicSynonym*{skelfilegt ástand} \dicDirectTranslationCS{pohroma, neštěstí};  \dicPhraseIS{hörmung er að sjá e‑ð} \dicDirectTranslationCS{je to smutný pohled vidět (co)}
\dicEntry[hörpu] \dicTerm{hörpu} \dicIPA{{h}{\oe}{\textsubring{r}}{\textsubring{b}}{\textscy}} \dicPos{f} \dicFlx{sg gen} \dicLink{harpa}
\dicEntry[hörpudiskur] \dicTerm{hörpu··disk|ur}\dicTerm{, hörpuskel} \dicIPA{{h}\-{\oe}\-{\textsubring{r}}\-{\textsubring{b}}\-{\textscy}\-{\textsubring{d}}\-{\textsci}\-{s}\-{\r{g}}\-{\textscy}\-{\textsubring{r}}\-} \dicPos{m}[6] \dicFlx{(‑s, ‑ar)}[10] \dicFieldCat{zool.} \dicDirectTranslationCS{hřebenatka islandská} \textit{(l.~{\textLA{Chlamys islandica}})}
\dicEntry[hörpuleikari] \dicTerm{hörpu··leik·ar|i} \dicIPA{{h}{\oe}{\textsubring{r}}{\textsubring{b}}{\textscy}{l}{ei}{\r{g}}{a}{r}{\textsci}} \dicPos{m}[1] \dicFlx{(‑a, ‑ar)}[13] \dicDirectTranslationCS{harfenista, harfenistka}
\dicEntry[hörpuskel] \dicTerm{hörpu··skel} \dicIPA{{h}{\oe}{\textsubring{r}}{\textsubring{b}}{\textscy}{s}{\r{\textObardotlessj}}{\textepsilon}{\textsubring{l}}} \dicPos{f}[4] \dicFlx{(‑jar, ‑jar)}[9] \dicLink{hörpudiskur}
\dicEntry[hörund] \dicTerm{hörund} \dicsymFrequent\  \dicIPA{{h}{\oe}{\textlengthmark}{r}{\textscy}{n}{\textsubring{d}}} \dicPos{n}[2] \dicFlx{(‑s)}[2] \dicDirectTranslationCS{pleť, pokožka} \dicExampleIS{vera ljós á hörund} \dicExampleCS{mít světlou pleť}
\dicEntry[hörundsár] \dicTerm{hörund··sár} \dicIPA{{h}{\oe}{\textlengthmark}{r}{\textscy}{n}{\textsubring{d}}{s}{au}{\textsubring{r}}} \dicPos{adj}[5] \dicFlx{(f ‑)}[8] \dicDirectTranslationCS{citlivý, choulostivý}
\dicEntry[hörundslitur] \dicTerm{hörunds··lit|ur} \dicIPA{{h}{\oe}{\textlengthmark}{r}{\textscy}{n}{\textsubring{d}}{s}{l}{\textsci}{\textsubring{d}}{\textscy}{\textsubring{r}}} \dicPos{m}[10] \dicFlx{(‑ar, ‑ir)}[13] \dicDirectTranslationCS{barva pleti}
\dicEntry[höst] \dicTerm{höst} \dicIPA{{h}{\oe}{s}{\textsubring{d}}} \dicPos{adj} \dicFlx{f sg nom pos} \dicLink{hastur\smash{\textsuperscript{2}}}
\begin{xtolerant}{}{1pt}
\dicEntry[höstugur] \dicTerm{höst··ugur}\addthinS\dicIPA{{h}{\oe}{s}{\textsubring{d}}{\textscy}{\textbabygamma}{\textscy}{\textsubring{r}}}\addthinS\dicPos{adj}[1]\dicFlx{}[-8]\addthinS\textbf{1.} \dicSynonym{óvingjarnlegur} \dicDirectTranslationCS{strohý, příkrý}  \textbf{2.} \dicSynonym{slæmur} \dicDirectTranslationCS{zlý, závažný (nemoc ap.)}
\end{xtolerant}
\dicEntry[höttur] \dicTerm{höttur} \dicIPA{{h}{\oe}{h}{\textsubring{d}}{\textscy}{\textsubring{r}}} \dicPos{m}[11] \dicFlx{(hattar, hettir)}[5] \dicPhraseIS{vera á höttunum eftir e‑u} \dicDirectTranslationCS{ohlížet se po (čem), hledat (co)};  \dicPhraseIS{Hrói Höttur} \dicFlx{prop} \dicDirectTranslationCS{Robin Hood};  \dicPhraseIS{út í hött} \dicFlx{adj\,/\addthin adv} \dicDirectTranslationCS{(jsoucí) mimo, absurdní}

\else
\foreach \l in \alphabetF{
  \dicLetter{l}{letter15}
\dicEntry[l] \dicTerm{l} \dicPos{zkr} \dicPhraseIS{lítri} \dicDirectTranslationCS{litr}
\dicEntry[l.] \dicTerm{l.} \dicPos{zkr} \dicPhraseIS{lýsingarorð} \dicFieldCat{jaz.} \dicDirectTranslationCS{přídavné jméno}
\dicEntry[labb] \dicTerm{labb} \dicIPA{{l}{a}{\textsubring{b}}{\textlengthmark}} \dicPos{n}[2] \dicFlx{(‑s)}[2] \dicDirectTranslationCS{procházka, špacír}
\dicEntry[labba] \dicTerm{labb|a} \dicsymFrequent\  \dicIPA{{l}{a}{\textsubring{b}}{\textlengthmark}{a}} \dicPos{v}[1] \dicFlx{(‑aði)}[13] \dicDirectTranslationCS{projít se, procházet se, jít na procházku, špacírovat} \dicExampleIS{Ég labbaði í bæinn.} \dicExampleCS{Prošel jsem se do města.}
\dicEntry[labbakútur] \dicTerm{labba··kút|ur} \dicIPA{{l}{a}{\textsubring{b}}{\textlengthmark}{a}{k\smash{\textsuperscript{h}}}{u}{\textsubring{d}}{\textscy}{\textsubring{r}}} \dicPos{m}[6] \dicFlx{(‑s, ‑ar)}[22] \textbf{1.} \dicSynonym{barn} \dicDirectTranslationCS{klouček, prcek}  \textbf{2.} \dicSynonym*{mannleysa} \dicDirectTranslationCS{trdlo, janek}
\dicEntry[laða] \dicTerm{lað|a} \dicIPA{{l}{a}{\textlengthmark}{ð}{a}} \dicPos{v}[1] \dicFlx{(‑aði)}[13] \dicFlx{acc} \dicDirectTranslationCS{přitahovat, (při)lákat, (při)vábit} \dicExampleIS{laða e‑n að sér} \dicExampleCS{přilákat (koho) k~sobě};  \dicIdiom{laða}[fram]{ \dicPhraseIS{laða e‑u fram}} \dicDirectTranslationCS{vyvolat (co), přivolat (co), přivodit (co)};  \dicIdiom{laðast}{ \dicPhraseIS{laðast}} \dicFlx{refl} \dicDirectTranslationCS{být přitahován} \dicExampleIS{laðast að e‑u} \dicExampleCS{být přitahován (čím)}
\dicEntry[lafa] \dicTerm{laf|a} \dicIPA{{l}{a}{\textlengthmark}{v}{a}} \dicPos{v}[2] \dicFlx{(‑ði, ‑að)}[125] \textbf{1.} \dicSynonym{hanga} \dicDirectTranslationCS{viset} \dicExampleIS{láta fæturna lafa} \dicExampleCS{nechat nohy viset}  \textbf{2.} \dicDirectTranslationCS{prodlévat, zdržovat se} \dicExampleIS{lafa í þorpinu} \dicExampleCS{zdržovat se na vsi}
\dicEntry[lag] \dicTerm{lag} \dicsymFrequent\  \dicIPA{{l}{a}{\textlengthmark}{x}} \dicPos{n}[2] \dicFlx{(‑s, lög)}[8] \textbf{1.} \dicSynonym{samfella} \dicDirectTranslationCS{vrstva} \dicExampleIS{lag af kolum} \dicExampleCS{vrstva uhlí}  \textbf{2.} \dicDirectTranslationCS{(pro)bodnutí (zbraní ap.)} \dicExampleIS{Lagið nam í hjartastað.} \dicExampleCS{Bodnutí zasáhlo srdce.}  \textbf{3.} \dicSynonym{aðferð} \dicDirectTranslationCS{způsob, metoda};  \dicPhraseIS{finna lagið á e‑u} \dicDirectTranslationCS{přijít na to, jak se (co) dělá};  \dicPhraseIS{hafa lag á e‑u} \dicDirectTranslationCS{mít talent na (co)};  \dicPhraseIS{koma e‑m (upp) á lagið með e‑ð} \dicDirectTranslationCS{naučit (koho), jak se (co) dělá};  \dicPhraseIS{með þessu lagi} \dicFlx{adv} \dicDirectTranslationCS{takto, tímto způsobem}  \textbf{4.} \dicSynonym{magn} \dicDirectTranslationCS{množství};  \dicPhraseIS{í einu lagi} \dicFlx{adv} \dicSynonym*{í heild} \dicDirectTranslationCS{najednou, naráz};  \dicPhraseIS{í fyrsta lagi} \dicFlx{adv} {\textbf{a.}} \dicDirectTranslationCS{ne dříve než} \dicExampleIS{Ég kem í fyrsta lagi kl. þrjú.} \dicExampleCS{Nepřijdu dřív než ve tři.};  {\textbf{b.}} \dicDirectTranslationCS{za prvé, v~první řadě};  \dicPhraseIS{í hæsta lagi} \dicFlx{adv} \dicDirectTranslationCS{nanejvýš, maximálně, nejvýše} \dicExampleIS{Þetta tekur í hæsta lagi fimm mínútur.} \dicExampleCS{Zabere to nejvýše pět minut.};  \dicPhraseIS{í meira lagi} \dicFlx{adv} \dicDirectTranslationCS{dosti, značně};  \dicPhraseIS{í mesta lagi} \dicFlx{adv} \dicDirectTranslationCS{ne více než};  \dicPhraseIS{í öðru lagi} \dicFlx{adv} \dicDirectTranslationCS{za druhé}  \textbf{5.} \dicSynonym{regla} \dicDirectTranslationCS{řád, pořádek};  \dicPhraseIS{allt í lagi!} \dicDirectTranslationCS{ok!, všechno v~pořádku!};  \dicPhraseIS{e‑að er fjarri lagi} \dicDirectTranslationCS{(co) není zdaleka v~pořádku};  \dicPhraseIS{koma e‑u í lag} \dicSynonym{lagfæra} \dicDirectTranslationCS{dát (co) do pořádku};  \dicPhraseIS{nærri lagi} \dicFlx{adv} \dicDirectTranslationCS{téměř v~pořádku}  \textbf{6.} \dicSynonym{tónverk} \dicDirectTranslationCS{píseň, písnička, melodie} \dicExampleIS{lag eftir e‑n} \dicExampleCS{(čí) písnička};  \dicPhraseIS{fara út af laginu} \dicDirectTranslationCS{vypadnout z~rytmu (v~písni), zazpívat falešný tón};  \dicPhraseIS{taka lagið} \dicDirectTranslationCS{(začít) zpívat}  \textbf{7.} \dicSynonym{verðlag} \dicDirectTranslationCS{cena, hodnota};  \dicPhraseIS{leggja lag á e‑ð} \dicDirectTranslationCS{ocenit (co), ohodnotit (co), určit cenu (čeho)}  \textbf{8.} \dicSynonym{samneyti} \dicDirectTranslationCS{společnost, přítomnost} \dicIndirectTranslationCS{(společenský styk)};  \dicPhraseIS{leggja lag sitt við e‑n} \dicDirectTranslationCS{dělat (komu) společnost}  \textbf{9.} \dicPhraseIS{lög} \dicFlx{pl} \dicLink{lög};  \dicIdiom{lag}{ \dicPhraseIS{slá e‑n út af laginu}} \dicDirectTranslationCS{vyvést (koho) z~míry, vytočit (koho)}; { \dicPhraseIS{sæta lagi}} \dicDirectTranslationCS{počkat si (na příležitost)}
\dicEntry[laga] \dicTerm{lag|a} \dicsymFrequent\  \dicIPA{{l}{a}{\textlengthmark}{\textbabygamma}{a}} \dicPos{v}[1] \dicFlx{(‑aði)}[13] \dicFlx{acc} \textbf{1.} \dicSynonym{prýða} \dicDirectTranslationCS{upravit, upravovat, urovnat, urovnávat} \dicExampleIS{laga á sér hárið} \dicExampleCS{upravit si vlasy}  \textbf{2.} \dicSynonym*{koma e‑u í lag} \dicDirectTranslationCS{opravit, opravovat, spravit, spravovat, dát\,/\addthin dávat do pořádku} \dicExampleIS{laga vélina} \dicExampleCS{opravit stroj}  \textbf{3.} \dicDirectTranslationCS{připravit, připravovat (kávu ap.)} \dicExampleIS{laga kaffi og te} \dicExampleCS{připravit kávu a~čaj}  \textbf{4.} \dicSynonym{streyma} \dicDirectTranslationCS{proudit, prýštit (krev z~rány ap.)};  \dicIdiom{laga}[að]{ \dicPhraseIS{laga sig að e‑u}} \dicSynonym{aðhæfa} \dicDirectTranslationCS{přizpůsobit se (čemu)};  \dicIdiom{laga}[til]{ \dicPhraseIS{laga til}} \dicDirectTranslationCS{uklízet, poklízet, uklidit} \dicExampleIS{laga til í íbúðinni} \dicExampleCS{uklidit v~bytě}; { \dicPhraseIS{laga sig til}} \dicDirectTranslationCS{poupravit se};  \dicIdiom{lagast}{ \dicPhraseIS{lagast}} \dicFlx{refl} {\textbf{a.}} \dicDirectTranslationCS{(z)lepšit se, vylepšit se} \dicExampleIS{Ástandið lagast af sjálfu sér.} \dicExampleCS{Situace se zlepšuje sama od sebe.};  {\textbf{b.}} \dicDirectTranslationCS{ustoupit, povolit (bolest ap.)}
\dicEntry[lagaákvæði] \dicTerm{laga··á·kvæði} \dicIPA{{l}{a}{\textlengthmark}{\textbabygamma}{a}{au}{k\smash{\textsuperscript{h}}}{v}{a}{i}{ð}{\textsci}} \dicPos{n}[2] \dicFlx{(‑s, ‑)}[14] \dicFieldCat{práv.} \dicDirectTranslationCS{právní předpis\,/\addthin ustanovení}
\dicEntry[lagabálkur] \dicTerm{laga··bálk|ur} \dicIPA{{l}{a}{\textlengthmark}{\textbabygamma}{a}{\textsubring{b}}{au}{\textsubring{l}}{\r{g}}{\textscy}{\textsubring{r}}} \dicPos{m}[6] \dicFlx{(‑s, ‑ar)}[24] \dicFieldCat{práv.} \dicDirectTranslationCS{kodex, zákoník}
\dicEntry[lagabreyting] \dicTerm{laga··breyt·ing} \dicIPA{{l}{a}{\textlengthmark}{\textbabygamma}{a}{\textsubring{b}}{r}{ei}{\textsubring{d}}{i}{\ng}{\r{g}}} \dicPos{f}[4] \dicFlx{(‑ar, ‑ar)}[5] \dicFieldCat{práv.} \dicDirectTranslationCS{novela zákona}
\dicEntry[lagabrot] \dicTerm{laga··brot} \dicIPA{{l}{a}{\textlengthmark}{\textbabygamma}{a}{\textsubring{b}}{r}{\textopeno}{\textsubring{d}}} \dicPos{n}[2] \dicFlx{(‑s, ‑)}[5] \dicDirectTranslationCS{porušení zákona}
\dicEntry[lagadeild] \dicTerm{laga··deild} \dicIPA{{l}{a}{\textlengthmark}{\textbabygamma}{a}{\textsubring{d}}{ei}{l}{\textsubring{d}}} \dicPos{f}[7] \dicFlx{(‑ar, ‑ir)}[1] \dicDirectTranslationCS{právnická fakulta}
\dicEntry[lagafrumvarp] \dicTerm{laga··frum·|varp} \dicIPA{{l}{a}{\textlengthmark}{\textbabygamma}{a}{f}{r}{\textscy}{m}{v}{a}{\textsubring{r}}{\textsubring{b}}} \dicPos{n}[2] \dicFlx{(‑varps, ‑vörp)}[8] \dicFieldCat{práv.} \dicDirectTranslationCS{návrh zákona}
\dicEntry[lagagrein] \dicTerm{laga··grein} \dicIPA{{l}{a}{\textlengthmark}{\textbabygamma}{a}{\r{g}}{r}{ei}{\textsubring{n}}} \dicPos{f}[4] \dicFlx{(‑ar, ‑ar)}[1] \dicFieldCat{práv.} \dicSynonym{grein} \dicDirectTranslationCS{odstavec zákona, paragraf}
\dicEntry[lagaheimild] \dicTerm{laga··heimild} \dicIPA{{l}{a}{\textlengthmark}{\textbabygamma}{a}{h}{ei}{m}{\textsci}{l}{\textsubring{d}}} \dicPos{f}[7] \dicFlx{(‑ar, ‑ir)}[1] \dicDirectTranslationCS{moc zákona, zákonnost}
\dicEntry[lagahöfundur] \dicTerm{laga··höf·und|ur} \dicIPA{{l}{a}{\textlengthmark}{\textbabygamma}{a}{h}{\oe}{v}{\textscy}{n}{\textsubring{d}}{\textscy}{\textsubring{r}}} \dicPos{m}[6] \dicFlx{(‑ar, ‑ar)}[56] \dicDirectTranslationCS{písničkář(ka), skladatel(ka)}
\dicEntry[lagalegur] \dicTerm{laga··legur} \dicIPA{{l}{a}{\textlengthmark}{\textbabygamma}{a}{l}{\textepsilon}{\textbabygamma}{\textscy}{\textsubring{r}}} \dicPos{adj}[1]\dicFlx{}[-8] \dicDirectTranslationCS{zákonný, právní}
\dicEntry[laganir] \dicTerm{laganir} \dicIPA{{l}{a}{\textlengthmark}{\textbabygamma}{a}{n}{\textsci}{\textsubring{r}}} \dicPos{f} \dicFlx{pl nom} \dicLink{lögun}
\dicEntry[lagar] \dicTerm{lagar} \dicIPA{{l}{a}{\textlengthmark}{\textbabygamma}{a}{\textsubring{r}}} \dicPos{m} \dicFlx{sg gen} \dicLink{lögur}
\dicEntry[lagaregla] \dicTerm{laga··regl|a} \dicIPA{{l}{a}{\textlengthmark}{\textbabygamma}{a}{r}{\textepsilon}{\r{g}}{l}{a}} \dicPos{f}[1] \dicFlx{(‑u, ‑ur)}[19] \dicFieldCat{práv.} \dicDirectTranslationCS{právní předpis}
\dicEntry[lagasafn] \dicTerm{laga··|safn} \dicIPA{{l}{a}{\textlengthmark}{\textbabygamma}{a}{s}{a}{\textsubring{b}}{\textsubring{n}}} \dicPos{n}[2] \dicFlx{(‑safns, ‑söfn)}[8] \dicFieldCat{práv.} \dicDirectTranslationCS{sbírka zákonů}
\dicEntry[lager] \dicTerm{lager} \dicIPA{{l}{a}{\textlengthmark}{\textbabygamma}{\textepsilon}{\textsubring{r}}} \dicPos{m}[4] \dicFlx{(‑s, ‑ar)}[14] \textbf{1.} \dicSynonym{vörubirgðir} \dicDirectTranslationCS{sklad, zásoba} \dicExampleIS{Vörur eru sendar samdægurs séu þær til á lager.} \dicExampleCS{Zboží je odesíláno ve stejný den, pokud je na skladě.};  \dicPhraseIS{hafa e‑ð á lager} \dicDirectTranslationCS{mít (co) na skladě}  \textbf{2.} \dicSynonym{vörugeymsla} \dicDirectTranslationCS{sklad, skladiště} \dicExampleIS{vinna á lager} \dicExampleCS{pracovat ve skladu}
\dicEntry[lagermaður] \dicTerm{lager··|maður} \dicIPA{{l}{a}{\textlengthmark}{\textbabygamma}{\textepsilon}{r}{m}{a}{ð}{\textscy}{\textsubring{r}}} \dicPos{m}[13] \dicFlx{(‑manns, ‑menn)}[2] \dicDirectTranslationCS{skladník, skladnice}
\dicEntry[lagfæra] \dicTerm{lag··fær|a} \dicIPA{{l}{a}{x}{f}{a}{i}{r}{a}} \dicPos{v}[2] \dicFlx{(‑ði, ‑t)}[105] \dicFlx{acc} \textbf{1.} \dicDirectTranslationCS{dát\,/\addthin uvést do pořádku, opravit} \dicExampleIS{lagfæra skemmdir á húsi} \dicExampleCS{dát do pořádku poškození na domě}  \textbf{2.} \dicSynonym{endurbæta} \dicDirectTranslationCS{zlepšit, vylepšit, uzpůsobit}
\dicEntry[lagfæring] \dicTerm{lag··fær·ing} \dicIPA{{l}{a}{x}{f}{a}{i}{r}{i}{\ng}{\r{g}}} \dicPos{f}[4] \dicFlx{(‑ar, ‑ar)}[5] \dicSynonym{endurbót} \dicDirectTranslationCS{oprava, náprava, korekce}
\dicEntry[laggar] \dicTerm{laggar} \dicIPA{{l}{a}{\r{g}}{\textlengthmark}{a}{\textsubring{r}}} \dicPos{f} \dicFlx{sg gen} \dicLink{lögg}
\dicEntry[laggir] \dicTerm{laggir} \dicIPA{{l}{a}{\r{\textObardotlessj}}{\textlengthmark}{\textsci}{\textsubring{r}}} \dicPos{f} \dicFlx{pl nom} \dicLink{lögg}
\dicEntry[laggóður] \textls[15]{\dicTerm{lag··|góður} \dicIPA{{l}{a}{\textbabygamma}{\r{g}}{ou}{ð}{\textscy}{\textsubring{r}}} \dicPos{adj}[11] \dicFlx{(comp ‑betri, sup ‑bestur)}[1] \dicSynonym{hnyttinn} \dicDirectTranslationCS{výstižný, trefný} \dicExampleIS{stutt og laggott svar} \dicExampleCS{krátká a~výstižná odpověď}}
\dicEntry[laghendur] \dicTerm{lag··hendur}\dicTerm{, laghentur} \dicIPA{{l}\-{a}\-{\textlengthmark}\-{x}\-{h}\-{\textepsilon}\-{n}\-{\textsubring{d}}\-{\textscy}\-{\textsubring{r}}\-} \dicPos{adj}[2]\dicFlx{}[-14] \dicSynonym{handlaginn} \dicDirectTranslationCS{obratný, zručný}
\dicEntry[laghentur] \dicTerm{lag··hentur} \dicIPA{{l}{a}{\textlengthmark}{x}{h}{\textepsilon}{\textsubring{n}}{\textsubring{d}}{\textscy}{\textsubring{r}}} \dicPos{adj}[1]\dicFlx{}[-10] \dicLink{laghendur}
\dicEntry[laginn] \dicTerm{laginn} \dicIPA{{l}{a}{i}{j}{\textlengthmark}{\textsci}{\textsubring{n}}} \dicPos{adj}[6]\dicFlx{}[-3] \dicDirectTranslationCS{zručný, šikovný};  \dicPhraseIS{e‑r er laginn í höndunum} \dicDirectTranslationCS{(kdo) má šikovné ruce};  \dicPhraseIS{e‑r er laginn við e‑n} \dicDirectTranslationCS{(kdo) má dobrý přístup ke (komu) (zvířatům ap.)}
\dicEntry[lagkaka] \dicTerm{lag··|kaka} \dicIPA{{l}{a}{\textlengthmark}{\textbabygamma}{k\smash{\textsuperscript{h}}}{a}{\r{g}}{a}} \dicPos{f}[1] \dicFlx{(‑köku, ‑kökur)}[20] \dicFieldCat{kulin.} \dicDirectTranslationCS{řezy (čokoládové ap.), vrstvený koláč, dortík, zákusek}
\dicEntry[laglaus] \dicTerm{lag··laus} \dicIPA{{l}{a}{\textbabygamma}{l}{\oe i}{s}} \dicPos{adj}[5]\dicFlx{}[-1] \dicDirectTranslationCS{(jsoucí) bez hudebního sluchu, hluchý (k~hudbě ap.)}
\dicEntry[laglegur] \dicTerm{lag··legur} \dicsymFrequent\  \dicIPA{{l}{a}{\textbabygamma}{l}{\textepsilon}{\textbabygamma}{\textscy}{\textsubring{r}}} \dicPos{adj}[1]\dicFlx{}[-8] \textbf{1.} \dicSynonym{fríður} \dicDirectTranslationCS{pohledný, půvabný} \dicExampleIS{lagleg stúlka} \dicExampleCS{pohledná dívka}  \textbf{2.} \dicSynonym{snotur} \dicDirectTranslationCS{krásný, hezký}  \textbf{3.} \dicDirectTranslationCS{pěkný, pořádný} \dicIndirectTranslationCS{(k~zdůraznění)}
\dicEntry[laglína] \dicTerm{lag··lín|a} \dicIPA{{l}{a}{\textbabygamma}{l}{i}{n}{a}} \dicPos{f}[1] \dicFlx{(‑u, ‑ur)}[7] \dicDirectTranslationCS{melodie}
\dicEntry[lagnar] \dicTerm{lagnar} \dicIPA{{l}{a}{\r{g}}{n}{a}{\textsubring{r}}} \dicPos{f} \dicFlx{sg gen} \dicLink{lögn}
\dicEntry[lagni] \dicTerm{lagn|i} \dicIPA{{l}{a}{\r{g}}{n}{\textsci}} \dicPos{f}[3] \dicFlx{(‑i)}[3] \dicSynonym{leikni} \dicDirectTranslationCS{šikovnost, zručnost, obratnost}
\dicEntry[lagning] \dicTerm{lag··ning} \dicIPA{{l}{a}{\r{g}}{n}{i}{\ng}{\r{g}}} \dicPos{f}[4] \dicFlx{(‑ar, ‑ar)}[5] \textbf{1.} \dicDirectTranslationCS{kladení, pokládání, instalace} \dicExampleIS{lagning vegar} \dicExampleCS{kladení cesty}  \textbf{2.} \dicSynonym*{hárlagning} \dicDirectTranslationCS{úprava\,/\addthin upravení vlasů}
\dicEntry[lagnir] \dicTerm{lagnir} \dicIPA{{l}{a}{\r{g}}{n}{\textsci}{\textsubring{r}}} \dicPos{f} \dicFlx{pl nom} \dicLink{lögn}
\dicEntry[lagskipting] \dicTerm{lag··skipt·ing} \dicIPA{{l}{a}{x}{s}{\r{\textObardotlessj}}{\textsci}{f}{\textsubring{d}}{i}{\ng}{\r{g}}} \dicPos{f}[4] \dicFlx{(‑ar)}[7] \dicDirectTranslationCS{(roz)vrstvení}
\dicEntry[lagsmaður] \dicTerm{lags··|maður} \dicIPA{{l}{a}{x}{s}{m}{a}{ð}{\textscy}{\textsubring{r}}} \dicPos{m}[13] \dicFlx{(‑manns, ‑menn)}[2] \dicSynonym{félagi} \dicDirectTranslationCS{kamarád(ka), přítel(kyně), společník, společnice}
\dicEntry[lagviss] \dicTerm{lag··viss} \dicIPA{{l}{a}{\textbabygamma}{v}{\textsci}{s}} \dicPos{adj}[5]\dicFlx{}[-1] \dicDirectTranslationCS{zpívající čistě, muzikální}
\dicEntry[lak] \dicTerm{lak\smash{\textsuperscript{1}}} \dicIPA{{l}{a}{\textlengthmark}{\r{g}}} \dicPos{n}[2] \dicFlx{(‑s, lök)}[8] \dicDirectTranslationCS{prostěradlo}
\dicEntry[lak] \dicTerm{lak\smash{\textsuperscript{2}}} \dicIPA{{l}{a}{\textlengthmark}{\r{g}}} \dicPos{v} \dicFlx{ind pf sg 1 pers} \dicLink{leka}
\dicEntry[lakk] \dicTerm{lakk} \dicIPA{{l}{a}{h}{\r{g}}} \dicPos{n}[2] \dicFlx{(‑s, lökk)}[8] \textbf{1.} \dicSynonym*{lakkmálning} \dicDirectTranslationCS{lak, nátěr}  \textbf{2.} \dicSynonym{innsigli} \dicDirectTranslationCS{pečeť}  \textbf{3.} \dicSynonym*{innsiglislakk} \dicDirectTranslationCS{pečetní vosk}
\dicEntry[lakka] \dicTerm{lakk|a} \dicIPA{{l}{a}{h}{\r{g}}{a}} \dicPos{v}[1] \dicFlx{(‑aði)}[13] \dicFlx{acc} \textbf{1.} \dicDirectTranslationCS{(na)lakovat} \dicExampleIS{lakka gluggana} \dicExampleCS{nalakovat okna}  \textbf{2.} \dicDirectTranslationCS{zapečetit} \dicExampleIS{lakka bréf} \dicExampleCS{zapečetit dopis}
\dicEntry[lakkrís] \dicTerm{lakk··rís} \dicIPA{{l}{a}{h}{\r{g}}{r}{i}{s}} \dicPos{m}[4] \dicFlx{(‑s)}[18] \dicDirectTranslationCS{lékořice} \dicExampleIS{fá einn poka af lakkrís} \dicExampleCS{dostat jeden pytlík lékořice}
\dicEntry[lakkrísrót] \dicTerm{lakk·rís··|rót} \dicIPA{{l}{a}{h}{\r{g}}{r}{i}{s}{r}{ou}{\textsubring{d}}} \dicPos{f}[8] \dicFlx{(‑rótar, ‑rætur)}[5] \dicFieldCat{bot.} \dicDirectTranslationCS{kořen lékořice lysé} \textit{(l.~{\textLA{Glycyrrhiza glabra}})}
\dicEntry[laktósi] \dicTerm{laktós|i} \dicIPA{{l}{a}{x}{\textsubring{d}}{ou}{s}{\textsci}} \dicPos{m}[1] \dicFlx{(‑a)}[3] \dicFieldCat{chem.} \dicSynonym*{mjólkursykur} \dicDirectTranslationCS{laktóza, mléčný cukr}
\dicEntry[lakur] \dicTerm{lakur} \dicIPA{{l}{a}{\textlengthmark}{\r{g}}{\textscy}{\textsubring{r}}} \dicPos{adj}[1] \dicFlx{(f lök)}[2] \textbf{1.} \dicSynonym{lélegur} \dicDirectTranslationCS{chabý, mizerný (výsledek ap.)} \dicExampleIS{lakur árangur} \dicExampleCS{mizerný výsledek}  \textbf{2.} \dicSynonym{lasinn} \dicDirectTranslationCS{chorý, neduživý}  \textbf{3.} \dicSynonym{tæpur} \dicDirectTranslationCS{necelý, pouhý}
\dicEntry[lama] \dicTerm{lam|a} \dicIPA{{l}{a}{\textlengthmark}{m}{a}} \dicPos{v}[1] \dicFlx{(‑aði)}[13] \dicDirectTranslationCS{ochromit, paralyzovat} \dicExampleIS{Sjúkdómur lamar líkamann.} \dicExampleCS{Nemoc ochromuje tělo.};  \dicIdiom{lamast}{ \dicPhraseIS{lamast}} \dicFlx{refl} {\textbf{a.}} \dicDirectTranslationCS{ochrnout (o~člověku ap.)} \dicExampleIS{lamast við slysið} \dicExampleCS{ochrnout při nehodě};  {\textbf{b.}} \dicDirectTranslationCS{být paralyzován\,/\addthin ochromen (dopravní ruch ap.)}
\dicEntry[lamadýr] \dicTerm{lama··dýr} \dicIPA{{l}{a}{\textlengthmark}{m}{a}{\textsubring{d}}{i}{\textsubring{r}}} \dicPos{n}[2] \dicFlx{(‑s, ‑)}[5] \dicFieldCat{zool.} \dicDirectTranslationCS{lama} \textit{(l.~{\textLA{Lama}})}  \dicsymPhoto\ 
\dicFigure{ds_image_lamadyr_0_1.jpg}{Lamadýr}{Lamadýr - Luca Galuzzi, CC BY-SA 2.5}
\dicEntry[lamaður] \dicTerm{lam··aður} \dicIPA{{l}{a}{\textlengthmark}{m}{a}{ð}{\textscy}{\textsubring{r}}} \dicPos{adj}[3] \dicFlx{(f lömuð)}[1] \dicDirectTranslationCS{chromý, ochrnutý}
\dicEntry[lamasess] \dicTerm{lama··sess} \dicIPA{{l}{a}{\textlengthmark}{m}{a}{s}{\textepsilon}{s}} \dicPos{m}[4] \dicFlx{(‑)}[34] \dicPhraseIS{liggja\,/\addthin vera í lamasessi} \dicDirectTranslationCS{být mimo provoz, nefungovat (o~věcech)}
\dicEntry[lamb] \dicTerm{lamb} \dicsymFrequent\  \dicIPA{{l}{a}{m}{\textsubring{b}}} \dicPos{n}[2] \dicFlx{(‑s, lömb)}[8] \textbf{1.} \dicDirectTranslationCS{jehně} \dicExampleIS{nýfædd lömb} \dicExampleCS{novorozená jehňata}  \textbf{2.} \dicDirectTranslationCS{drahoušek, miláček};  \dicPhraseIS{launa e‑m lambið gráa} \dicLangCat{přen.} \dicDirectTranslationCS{oplatit (komu) stejnou mincí}
\dicEntry[lambagras] \dicTerm{lamba··|gras} \dicIPA{{l}{a}{m}{\textsubring{b}}{a}{\r{g}}{r}{a}{s}} \dicPos{n}[2] \dicFlx{(‑grass, ‑grös)}[8] \dicFieldCat{bot.} \dicDirectTranslationCS{silenka bezlodyžná} \textit{(l.~{\textLA{Silene acaulis}})}  \dicsymPhoto\ 
\dicFigure{9260.jpg}{Lambagras}{Lambagras - Daněk Pavel, Biolib, Copyright/CC-BY-SA}
\dicEntry[lambakjöt] \dicTerm{lamba··kjöt} \dicIPA{{l}{a}{m}{\textsubring{b}}{a}{c\smash{\textsuperscript{h}}}{\oe}{\textsubring{d}}} \dicPos{n}[2] \dicFlx{(‑s)}[2] \dicDirectTranslationCS{jehněčí (maso)}
\dicEntry[lambalæri] \dicTerm{lamba··læri} \dicIPA{{l}{a}{m}{\textsubring{b}}{a}{l}{a}{i}{r}{\textsci}} \dicPos{n}[2] \dicFlx{(‑s, ‑)}[14] \dicDirectTranslationCS{jehněčí kýta}
\dicEntry[lamdi] \dicTerm{lamdi} \dicIPA{{l}{a}{m}{\textsubring{d}}{\textsci}} \dicPos{v} \dicFlx{ind pf sg 1 pers} \dicLink{lemja}
\dicEntry[lamið] \dicTerm{lamið} \dicIPA{{l}{a}{\textlengthmark}{m}{\textsci}{\texttheta}} \dicPos{v} \dicFlx{supin} \dicLink{lemja}
\dicEntry[lampi] \dicTerm{lamp|i} \dicsymFrequent\  \dicIPA{{l}{a}{\textsubring{m}}{\textsubring{b}}{\textsci}} \dicPos{m}[1] \dicFlx{(‑a, ‑ar)}[8] \textbf{1.} \dicDirectTranslationCS{lampa, svítilna} \dicExampleIS{kveikja á lampa} \dicExampleCS{zapálit lampu}  \textbf{2.} \dicFieldCat{elek.} \dicDirectTranslationCS{výbojka}
\dicEntry[land] \dicTerm{land} \dicsymFrequent\  \dicIPA{{l}{a}{n}{\textsubring{d}}} \dicPos{n}[2] \dicFlx{(‑s, lönd)}[8] \textbf{1.} \dicSynonym{þurrlendi} \dicDirectTranslationCS{souš, pevnina, země} \dicExampleIS{á sjó og landi} \dicExampleCS{na moři a~na pevnině}  \textbf{2.} \dicSynonym{árbakki} \dicDirectTranslationCS{břeh} \dicExampleIS{á öðru landinu} \dicExampleCS{na druhém břehu}  \textbf{3.} \dicSynonym{ríki} \dicDirectTranslationCS{země, stát, vlast} \dicExampleIS{höfuðborg landsins} \dicExampleCS{hlavní město země}  \textbf{4.} \dicSynonym{landareign} \dicDirectTranslationCS{země, půda, pozemek};  \dicIdiom{land}{ \dicPhraseIS{e‑að á langt í land}} \dicLangCat{přen.} \dicDirectTranslationCS{(co) je ještě daleko (cíl ap.)}; { \dicPhraseIS{eiga langt í land}} \dicLangCat{přen.} \dicDirectTranslationCS{mít ještě hodně před sebou (větší část práce ap.)}; { \dicPhraseIS{draga í land}} \dicLangCat{přen.} \dicDirectTranslationCS{ustoupit (ze zásad ap.), rezignovat, upustit (od plánů ap.)}; { \dicPhraseIS{draga e‑n að landi}} \dicDirectTranslationCS{dojíst po (kom)}; { \dicPhraseIS{fara með löndum}} \dicLangCat{přen.} \dicDirectTranslationCS{postupovat\,/\addthin jít opatrně}; { \dicPhraseIS{láta e‑ð lönd og leið}} \dicLangCat{přen.} \dicDirectTranslationCS{nestarat se o~(co), nechat (co) být}; { \dicPhraseIS{leggja land undir fót}} \dicDirectTranslationCS{jít pěšky, pochodovat, kráčet}
\dicEntry[landa] \dicTerm{land|a} \dicIPA{{l}{a}{n}{\textsubring{d}}{a}} \dicPos{v}[1] \dicFlx{(‑aði)}[13] \dicFlx{dat} \dicDirectTranslationCS{vyložit, vylodit} \dicExampleIS{landa aflanum} \dicExampleCS{vyložit úlovek}
\dicEntry[landabréf] \dicTerm{landa··bréf} \dicIPA{{l}{a}{n}{\textsubring{d}}{a}{\textsubring{b}}{r}{j}{\textepsilon}{f}} \dicPos{n}[2] \dicFlx{(‑s, ‑)}[5] \dicSynonym{landakort} \dicDirectTranslationCS{mapa (zeměpisná ap.)}
\dicEntry[landafræði] \dicTerm{landa··fræð|i} \dicIPA{{l}{a}{n}{\textsubring{d}}{a}{f}{r}{a}{i}{ð}{\textsci}} \dicPos{f}[3] \dicFlx{(‑i)}[3] \dicDirectTranslationCS{zeměpis, geografie}
\dicEntry[landakort] \dicTerm{landa··kort} \dicIPA{{l}{a}{n}{\textsubring{d}}{a}{k\smash{\textsuperscript{h}}}{\textopeno}{\textsubring{r}}{\textsubring{d}}} \dicPos{n}[2] \dicFlx{(‑s, ‑)}[5] \dicDirectTranslationCS{mapa}
\dicEntry[landakortabók] \dicTerm{landa·korta··|bók} \dicIPA{{l}{a}{n}{\textsubring{d}}{a}{k\smash{\textsuperscript{h}}}{\textopeno}{\textsubring{r}}{\textsubring{d}}{a}{\textsubring{b}}{ou}{\r{g}}} \dicPos{f}[8] \dicFlx{(‑bókar, ‑bækur)}[5] \dicSynonym{atlas} \dicDirectTranslationCS{atlas}
\dicEntry[landalda] \dicTerm{land··|alda} \dicIPA{{l}{a}{n}{\textsubring{d}}{a}{l}{\textsubring{d}}{a}} \dicPos{f}[1] \dicFlx{(‑öldu, ‑öldur)}[20] \dicDirectTranslationCS{pobřežní vlna}
\dicEntry[landamerki] \dicTerm{landa··merki} \dicIPA{{l}{a}{n}{\textsubring{d}}{a}{m}{\textepsilon}{\textsubring{r}}{\r{\textObardotlessj}}{\textsci}} \dicPos{n}[2] \dicFlx{(‑s, ‑)}[16] \dicSynonym{landamæri} \dicDirectTranslationCS{pomezí, hranice (státní ap.)}
\dicEntry[landamæri] \dicTerm{landa··mæri} \dicIPA{{l}{a}{n}{\textsubring{d}}{a}{m}{a}{i}{r}{\textsci}} \dicPos{n}[2] \dicFlx{pl}[19] \dicDirectTranslationCS{(státní) hranice} \dicExampleIS{vera kominn yfir landamærin} \dicExampleCS{překročit státní hranici}
\dicEntry[landareign] \dicTerm{landar··eign} \dicIPA{{l}{a}{n}{\textsubring{d}}{a}{r}{ei}{\r{g}}{\textsubring{n}}} \dicPos{f}[7] \dicFlx{(‑ar, ‑ir)}[1] \dicDirectTranslationCS{pozemek, půda, pozemkové vlastnictví}
\dicEntry[landbúnaðarframleiðsla] \dicTerm{land·búnaðar··fram·leiðsl|a} \dicIPA{{l}\-{a}\-{n}\-{\textsubring{d}}\-{\textsubring{b}}\-{u}\-{n}\-{a}\-{ð}\-{a}\-{\textsubring{r}}\-{f}\-{r}\-{a}\-{m}\-{l}\-{ei}\-{ð}\-{s}\-{\textsubring{d}}\-{l}\-{a}\-} \dicPos{f}[1] \dicFlx{(‑u)}[5] \dicDirectTranslationCS{zemědělská produkce\,/\addthin výroba}
\dicEntry[landbúnaðarráðherra] \dicTerm{land·búnaðar··ráð·herr|a} \dicIPA{{l}\-{a}\-{n}\-{\textsubring{d}}\-{\textsubring{b}}\-{u}\-{n}\-{a}\-{ð}\-{a}\-{r}\-{au}\-{\texttheta}\-{h}\-{\textepsilon}\-{r}\-{a}\-} \dicPos{m}[1] \dicFlx{(‑a, ‑ar)}[17] \dicDirectTranslationCS{ministr(yně) zemědělství}
\dicEntry[landbúnaðarráðuneyti] \textls[15]{\dicTerm{land·búnaðar··ráðu·neyti} \dicIPA{{l}\-{a}\-{n}\-{\textsubring{d}}\-{\textsubring{b}}\-{u}\-{n}\-{a}\-{ð}\-{a}\-{r}\-{au}\-{ð}\-{\textscy}\-{n}\-{ei}\-{\textsubring{d}}\-{\textsci}}} \dicPos{n}[2] \dicFlx{(‑s, ‑)}[14] \dicDirectTranslationCS{ministerstvo zemědělství}
\dicEntry[landbúnaður] \dicTerm{land··bú·nað|ur} \dicIPA{{l}{a}{n}{\textsubring{d}}{\textsubring{b}}{u}{n}{a}{ð}{\textscy}{\textsubring{r}}} \dicPos{m}[10] \dicFlx{(‑ar)}[9] \dicDirectTranslationCS{zemědělství}
\dicEntry[landeigandi] \dicTerm{land··eig·|andi} \dicIPA{{l}{a}{n}{\textsubring{d}}{ei}{\textbabygamma}{a}{n}{\textsubring{d}}{\textsci}} \dicPos{m}[2] \dicFlx{(‑anda, ‑endur)}[1] \dicDirectTranslationCS{majitel(ka) půdy\,/\addthin pozemku}
\dicEntry[landfastur] \dicTerm{land··|fastur} \dicIPA{{l}{a}{n}{\textsubring{d}}{f}{a}{s}{\textsubring{d}}{\textscy}{\textsubring{r}}} \dicPos{adj}[1] \dicFlx{(f ‑föst)}[11] \dicDirectTranslationCS{(jsoucí) při pobřeží (led ap.)}
\dicEntry[landflótta] \dicTerm{land··flótta} \dicIPA{{l}{a}{n}{\textsubring{d}}{f}{l}{ou}{h}{\textsubring{d}}{a}} \dicPos{adj}[13] \dicFlx{indecl}[1] \dicDirectTranslationCS{(jsoucí) v~exilu, (jsoucí) ve vyhnanství}
\dicEntry[landfræðilegur] \dicTerm{land·fræði··legur} \dicIPA{{l}{a}{n}{\textsubring{d}}{f}{r}{a}{i}{ð}{\textsci}{l}{\textepsilon}{\textbabygamma}{\textscy}{\textsubring{r}}} \dicPos{adj}[1]\dicFlx{}[-8] \dicDirectTranslationCS{geografický}
\dicEntry[landfræðingur] \dicTerm{land·fræð··ing|ur} \dicIPA{{l}{a}{n}{\textsubring{d}}{f}{r}{a}{i}{ð}{i}{\ng}{\r{g}}{\textscy}{\textsubring{r}}} \dicPos{m}[6] \dicFlx{(‑s, ‑ar)}[8] \dicDirectTranslationCS{geograf(ka)}
\dicEntry[landgrunn] \dicTerm{land··grunn} \dicIPA{{l}{a}{n}{\textsubring{d}}{\r{g}}{r}{\textscy}{\textsubring{n}}} \dicPos{n}[2] \dicFlx{(‑s, ‑)}[5] \dicFieldCat{geol.} \dicDirectTranslationCS{kontinentální\,/\addthin pevninský šelf, pevninský práh}
\dicEntry[landgræðsla] \dicTerm{land··græðsl|a} \dicIPA{{l}{a}{n}{\textsubring{d}}{\r{g}}{r}{a}{i}{ð}{s}{\textsubring{d}}{l}{a}} \dicPos{f}[1] \dicFlx{(‑u)}[5] \dicDirectTranslationCS{meliorace}
\dicEntry[landgöngubrú] \dicTerm{land··göngu·|brú} \dicIPA{{l}{a}{n}{\textsubring{d}}{\r{g}}{\oe i}{\ng}{\r{g}}{\textscy}{\textsubring{b}}{r}{u}} \dicPos{f}[9] \dicFlx{(‑brúar, ‑brýr)}[4] \dicDirectTranslationCS{lodní lávka}
\dicEntry[landhelgi] \dicTerm{land··helg|i} \dicsymFrequent\  \dicIPA{{l}{a}{n}{\textsubring{d}}{h}{\textepsilon}{l}{\r{\textObardotlessj}}{\textsci}} \dicPos{f}[3] \dicFlx{(‑i)}[3] \dicDirectTranslationCS{výsostné\,/\addthin teritoriální vody} \dicExampleIS{í landhelgi Íslands} \dicExampleCS{v~islandských výsostných vodách}
\dicEntry[landhelgisgæsla] \dicTerm{land·helgis··gæsl|a} \dicIPA{{l}{a}{n}{\textsubring{d}}{h}{\textepsilon}{l}{\r{\textObardotlessj}}{\textsci}{s}{\r{\textObardotlessj}}{a}{i}{s}{\textsubring{d}}{l}{a}} \dicPos{f}[1] \dicFlx{(‑u)}[5] \dicDirectTranslationCS{pobřežní stráž\,/\addthin hlídka}
\dicEntry[landher] \dicTerm{land··her} \dicIPA{{l}{a}{n}{\textsubring{d}}{h}{\textepsilon}{\textsubring{r}}} \dicPos{m}[9] \dicFlx{(‑s, ‑ir)}[21] \dicDirectTranslationCS{pozemní vojsko}
\dicEntry[landhreinsun] \dicTerm{land··hreins|un} \dicIPA{{l}{a}{n}{\textsubring{d}}{\textsubring{r}}{ei}{n}{s}{\textscy}{\textsubring{n}}} \dicPos{f}[7] \dicFlx{(‑unar)}[9] \dicPhraseIS{það er landhreinsun að e‑m} \dicLangCat{přen.} \dicDirectTranslationCS{no sláva, že už je (kdo) pryč}
\dicEntry[landi] \dicTerm{land|i} \dicIPA{{l}{a}{n}{\textsubring{d}}{\textsci}} \dicPos{m}[1] \dicFlx{(‑a, ‑ar)}[8] \textbf{1.} \dicSynonym{samlandi} \dicDirectTranslationCS{krajan(ka)}  \textbf{2.} \dicSynonym{brugg} \dicDirectTranslationCS{(domácí) pálenka}
\dicEntry[landkrabbi] \dicTerm{land··krabb|i} \dicIPA{{l}{a}{n}{\textsubring{d}}{k\smash{\textsuperscript{h}}}{r}{a}{\textsubring{b}}{\textsci}} \dicPos{m}[1] \dicFlx{(‑a, ‑ar)}[8] \textbf{1.} \dicFieldCat{zool.} \dicDirectTranslationCS{krab} \textit{(l.~{\textLA{Gecarcinus}})}  \dicsymPhoto\   \textbf{2.} \dicDirectTranslationCS{suchozemská krysa} \dicIndirectTranslationCS{(o~člověku, který ví málo o~práci na moři)}
\dicFigure{ds_image_landkrabbi_0_1.jpg}{Landkrabbi}{Landkrabbi - Dozenist, CC BY-SA 3.0}
\dicEntry[landkynning] \dicTerm{land··kynn·ing} \dicIPA{{l}{a}{n}{\textsubring{d}}{c\smash{\textsuperscript{h}}}{\textsci}{n}{i}{\ng}{\r{g}}} \dicPos{f}[4] \dicFlx{(‑ar, ‑ar)}[5] \dicDirectTranslationCS{propagace země}
\dicEntry[landkönnuður] \dicTerm{land··könn·uð|ur} \dicIPA{{l}{a}{n}{\textsubring{d}}{k\smash{\textsuperscript{h}}}{\oe}{n}{\textscy}{ð}{\textscy}{\textsubring{r}}} \dicPos{m}[10] \dicFlx{(‑ar, ‑ir)}[4] \dicDirectTranslationCS{cestovatel(ka), průzkumník, průzkumnice, objevitel(ka) (neznámých končin ap.)}
\dicEntry[landluktur] \dicTerm{land··luktur} \dicIPA{{l}{a}{n}{\textsubring{d}}{l}{\textscy}{x}{\textsubring{d}}{\textscy}{\textsubring{r}}} \dicPos{adj}[1]\dicFlx{}[-10] \dicDirectTranslationCS{vnitrozemský}
\dicEntry[landlægur] \dicTerm{land··lægur} \dicIPA{{l}{a}{n}{\textsubring{d}}{l}{a}{i}{\textbabygamma}{\textscy}{\textsubring{r}}} \dicPos{adj}[1]\dicFlx{}[-1] \textbf{1.} \dicFieldCat{med.} \dicDirectTranslationCS{endemický} \dicExampleIS{landlægur sjúkdómur} \dicExampleCS{endemické onemocnění}  \textbf{2.} \dicSynonym{fastur} \dicDirectTranslationCS{zakořeněný} \dicExampleIS{landlægur siður} \dicExampleCS{zakořeněný zvyk}
\dicEntry[landlæknir] \dicTerm{land··lækn|ir} \dicIPA{{l}{a}{n}{\textsubring{d}}{l}{a}{i}{h}{\r{g}}{n}{\textsci}{\textsubring{r}}} \dicPos{m}[7] \dicFlx{(‑is, ‑ar)}[1] \dicDirectTranslationCS{ředitel(ka) zdravotnictví}
\dicEntry[Landmannalaugar] \dicTerm{Land·manna··laugar} \dicIPA{{l}{a}{n}{\textsubring{d}}{m}{a}{n}{a}{l}{\oe i}{\textbabygamma}{a}{\textsubring{r}}} \dicPos{f}[4] \dicFlx{pl}[2] \dicFieldCat{geog.} \dicDirectTranslationCS{Landmannalaugar} \dicIndirectTranslationCS{(oblast s~horkými prameny a~barevnými skalami na jihu Islandu, oblíbený turistický cíl)}
\dicEntry[landmótunarfræði] \dicTerm{land·mótunar··fræð|i} \dicIPA{{l}{a}{n}{\textsubring{d}}{m}{ou}{\textsubring{d}}{\textscy}{n}{a}{\textsubring{r}}{f}{r}{a}{i}{ð}{\textsci}} \dicPos{f}[3] \dicFlx{(‑i)}[3] \dicDirectTranslationCS{geomorfologie}
\dicEntry[landmælingafræði] \dicTerm{land·mælinga··fræð|i} \dicIPA{{l}{a}{n}{\textsubring{d}}{m}{a}{i}{l}{i}{\ng}{\r{g}}{a}{f}{r}{a}{i}{ð}{\textsci}} \dicPos{f}[3] \dicFlx{(‑i)}[3] \dicDirectTranslationCS{geodézie, zeměměřičství}
\dicEntry[landmælingar] \dicTerm{land··mæl·ingar} \dicIPA{{l}{a}{n}{\textsubring{d}}{m}{a}{i}{l}{i}{\ng}{\r{g}}{a}{\textsubring{r}}} \dicPos{f}[4] \dicFlx{pl}[6] \dicDirectTranslationCS{geodetické měření}
\dicEntry[landnám] \dicTerm{land··nám} \dicIPA{{l}{a}{n}{\textsubring{d}}{n}{au}{\textsubring{m}}} \dicPos{n}[2] \dicFlx{(‑s)}[2] \dicDirectTranslationCS{osídlení, osídlování, kolonizace, kolonizování}

\dicEntry[landnámsmaður] \textls[19]{\dicTerm{land·náms··|maður} \dicIPA{{l}{a}{n}{\textsubring{d}}{n}{au}{m}{s}{m}{a}{ð}{\textscy}{\textsubring{r}}} \dicPos{m}[13] \dicFlx{(‑manns, ‑menn)}[2]} \dicDirectTranslationCS{osadník, osadnice, kolonista, kolonistka}

\dicEntry[landnámsöld] \dicTerm{land·náms··|öld} \dicIPA{{l}{a}{n}{\textsubring{d}}{n}{au}{m}{s}{\oe}{l}{\textsubring{d}}} \dicPos{f}[7] \dicFlx{(‑aldar)}[19] \dicFieldCat{hist.} \dicDirectTranslationCS{doba osídlování}
\dicEntry[landnemi] \dicTerm{land··nem|i} \dicIPA{{l}{a}{n}{\textsubring{d}}{n}{\textepsilon}{m}{\textsci}} \dicPos{m}[1] \dicFlx{(‑a, ‑ar)}[1] \dicDirectTranslationCS{osadník, osadnice, kolonista, kolonistka} \dicExampleIS{íslensku landnemarnir í Vesturheimi} \dicExampleCS{islandští osadníci v~Severní Americe}
\dicEntry[landráð] \dicTerm{land··ráð} \dicIPA{{l}{a}{n}{\textsubring{d}}{r}{au}{\texttheta}} \dicPos{n}[2] \dicFlx{pl}[1] \dicFieldCat{práv.} \dicDirectTranslationCS{velezrada, vlastizrada}
\dicEntry[landrek] \dicTerm{land··rek} \dicIPA{{l}{a}{n}{\textsubring{d}}{r}{\textepsilon}{\r{g}}} \dicPos{n}[2] \dicFlx{(‑s)}[2] \dicFieldCat{geol.} \dicDirectTranslationCS{kontinentální drift}
\dicEntry[landsbanki] \dicTerm{lands··bank|i} \dicIPA{{l}{a}{n}{\textsubring{d}}{s}{\textsubring{b}}{au}{\r{\textltailn}}{\r{\textObardotlessj}}{\textsci}} \dicPos{m}[1] \dicFlx{(‑a, ‑ar)}[8] \dicDirectTranslationCS{národní banka}
\dicEntry[landsbyggð] \dicTerm{lands··byggð} \dicIPA{{l}{a}{n}{\textsubring{d}}{s}{\textsubring{b}}{\textsci}{\textbabygamma}{\texttheta}} \dicPos{f}[7] \dicFlx{(‑ar)}[3] \dicDirectTranslationCS{venkov, venkovský kraj} \dicExampleIS{úti á landsbyggðinni} \dicExampleCS{na venkově}
\dicEntry[landsfrægur] \dicTerm{lands··frægur} \dicIPA{{l}{a}{n}{\textsubring{d}}{s}{f}{r}{a}{i}{\textbabygamma}{\textscy}{\textsubring{r}}} \dicPos{adj}[1]\dicFlx{}[-1] \dicDirectTranslationCS{známý\,/\addthin proslavený po celé zemi}
\dicEntry[landshluti] \dicTerm{lands··hlut|i} \dicIPA{{l}{a}{n}{\textsubring{d}}{s}{\textsubring{l}}{\textscy}{\textsubring{d}}{\textsci}} \dicPos{m}[1] \dicFlx{(‑a, ‑ar)}[1] \dicDirectTranslationCS{část země, region}
\dicEntry[landsími] \dicTerm{land··sím|i} \dicIPA{{l}{a}{n}{\textsubring{d}}{s}{i}{m}{\textsci}} \dicPos{m}[1] \dicFlx{(‑a, ‑ar)}[1] \dicFieldCat{hist.} \dicDirectTranslationCS{národní telefonní společnost}
\dicEntry[landskjálfti] \dicTerm{land··skjálft|i} \dicIPA{{l}{a}{n}{\textsubring{d}}{s}{\r{\textObardotlessj}}{au}{\textsubring{l}}{\textsubring{d}}{\textsci}} \dicPos{m}[1] \dicFlx{(‑a, ‑ar)}[1] \dicFieldCat{geol.} \dicLangCat{zast.} \dicSynonym{jarðskjálfti} \dicDirectTranslationCS{zemětřesení}
\dicEntry[landslag] \dicTerm{lands··lag} \dicsymFrequent\  \dicIPA{{l}{a}{n}{\textsubring{d}}{s}{l}{a}{x}} \dicPos{n}[2] \dicFlx{(‑s)}[2] \dicDirectTranslationCS{krajina, terén} \dicExampleIS{skoða landslag í firðinum} \dicExampleCS{zkoumat krajinu ve fjordu}
\dicEntry[landslagsarkitekt] \dicTerm{lands·lags··arkitekt} \dicIPA{{l}{a}{n}{\textsubring{d}}{s}{l}{a}{x}{s}{a}{\textsubring{r}}{\r{\textObardotlessj}}{\textsci}{\textsubring{d}}{\textepsilon}{x}{\textsubring{d}}} \dicPos{m}[4] \dicFlx{(‑s, ‑ar)}[10] \dicDirectTranslationCS{krajinný architekt, krajinná architektka}
\dicEntry[landslagsmynd] \dicTerm{lands·lags··mynd} \dicIPA{{l}{a}{n}{\textsubring{d}}{s}{l}{a}{x}{s}{m}{\textsci}{n}{\textsubring{d}}} \dicPos{f}[7] \dicFlx{(‑ar, ‑ir)}[1] \dicDirectTranslationCS{krajina (obraz), krajinomalba}
\dicEntry[landsleikur] \dicTerm{lands··leik|ur} \dicIPA{{l}{a}{n}{\textsubring{d}}{s}{l}{ei}{\r{g}}{\textscy}{\textsubring{r}}} \dicPos{m}[9] \dicFlx{(‑s, ‑ir)}[15] \dicDirectTranslationCS{mezistátní\,/\addthin mezinárodní zápas\,/\addthin utkání (ve fotbale ap.)}
\dicEntry[landslið] \dicTerm{lands··lið} \dicIPA{{l}{a}{n}{\textsubring{d}}{s}{l}{\textsci}{\texttheta}} \dicPos{n}[2] \dicFlx{(‑s, ‑)}[5] \dicDirectTranslationCS{národní tým\,/\addthin mužstvo}
\dicEntry[landsliðsþjálfari] \dicTerm{lands·liðs··þjálf·ar|i} \dicIPA{{l}\-{a}\-{n}\-{\textsubring{d}}\-{s}\-{l}\-{\textsci}\-{ð}\-{s}\-{\texttheta}\-{j}\-{au}\-{l}\-{v}\-{a}\-{r}\-{\textsci}\-} \dicPos{m}[1] \dicFlx{(‑a, ‑ar)}[13] \dicDirectTranslationCS{trenér(ka) národního týmu}
\dicEntry[landslýður] \dicTerm{lands··lýð|ur} \dicIPA{{l}{a}{n}{\textsubring{d}}{s}{l}{i}{ð}{\textscy}{\textsubring{r}}} \dicPos{m}[9] \dicFlx{(‑s, ‑ir)}[8] \dicDirectTranslationCS{obyvatelstvo}
\dicEntry[landsmaður] \dicTerm{lands··|maður} \dicsymFrequent\  \dicIPA{{l}{a}{n}{\textsubring{d}}{s}{m}{a}{ð}{\textscy}{\textsubring{r}}} \dicPos{m}[13] \dicFlx{(‑manns, ‑menn)}[2] \dicDirectTranslationCS{obyvatel(ka)} \dicExampleIS{allir landsmenn undir fimmtugu} \dicExampleCS{všichni obyvatelé pod padesát let}
\dicEntry[landsmót] \dicTerm{lands··mót} \dicIPA{{l}{a}{n}{\textsubring{d}}{s}{m}{ou}{\textsubring{d}}} \dicPos{n}[2] \dicFlx{(‑s, ‑)}[5] \textbf{1.} \dicDirectTranslationCS{celorepublikový\,/\addthin národní sjezd\,/\addthin setkání}  \textbf{2.} \dicDirectTranslationCS{celorepublikový\,/\addthin národní turnaj (v~házené ap.)}
\dicEntry[landspítali] \dicTerm{land··spítal|i} \dicIPA{{l}{a}{n}{\textsubring{d}}{s}{\textsubring{b}}{i}{t\smash{\textsuperscript{h}}}{a}{l}{\textsci}} \dicPos{m}[1] \dicFlx{(‑a, ‑ar)}[8] \dicDirectTranslationCS{národní univerzitní nemocnice}
\begin{xtolerant}{}{1pt}
\dicEntry[landssamband] \dicTerm{lands··sam·|band} \dicIPA{{l}{a}{n}{\textsubring{d}}{s}{a}{m}{\textsubring{b}}{a}{n}{\textsubring{d}}} \dicPos{n}[2] \dicFlx{(‑bands, ‑bönd)}[8] \dicDirectTranslationCS{národní\,/\addthin celorepublikový svaz}
\end{xtolerant}
\dicEntry[landstjóri] \dicTerm{land··stjór|i} \dicIPA{{l}{a}{n}{\textsubring{d}}{s}{\textsubring{d}}{j}{ou}{r}{\textsci}} \dicPos{m}[1] \dicFlx{(‑a, ‑ar)}[1] \dicDirectTranslationCS{místodržitel(ka), guvernér(ka), zemský správce, zemská správkyně}
\dicEntry[landsvala] \dicTerm{land··|svala} \dicIPA{{l}{a}{n}{\textsubring{d}}{s}{v}{a}{l}{a}} \dicPos{f}[1] \dicFlx{(‑svölu, ‑svölur)}[20] \dicFieldCat{zool.} \dicDirectTranslationCS{vlaštovka, vlaštovka obecná} \textit{(l.~{\textLA{Hirundo rustica}})}  \dicsymPhoto\ 
\dicFigure{20731.jpg}{Landsvala}{Landsvala - Menke, Dave, Biolib, PD}
\dicEntry[Landsvirkjun] \dicTerm{Lands··virkj|un} \dicIPA{{l}{a}{n}{\textsubring{d}}{s}{v}{\textsci}{\textsubring{r}}{\r{\textObardotlessj}}{\textscy}{\textsubring{n}}} \dicPos{f}[7] \dicFlx{(‑unar)}[10] \dicIndirectTranslationCS{islandská společnost provozující elektrárny}
\dicEntry[landsvæði] \dicTerm{land··svæði} \dicIPA{{l}{a}{n}{\textsubring{d}}{s}{v}{a}{i}{ð}{\textsci}} \dicPos{n}[2] \dicFlx{(‑s, ‑)}[14] \dicDirectTranslationCS{území, oblast, teritorium}
\dicEntry[langa] \dicTerm{langa\smash{\textsuperscript{1}}} \dicIPA{{l}{au}{\ng}{\r{g}}{a}} \dicPos{f}[1] \dicFlx{(löngu, löngur)}[20] \dicFieldCat{zool.} \dicDirectTranslationCS{mník mořský} \textit{(l.~{\textLA{Molva molva}})}
\dicEntry[langa] \dicTerm{lang|a\smash{\textsuperscript{2}}} \dicsymFrequent\  \dicIPA{{l}{au}{\ng}{\r{g}}{a}} \dicPos{v}[1] \dicFlx{(‑aði)}[79] \dicFlx{impers} \dicPhraseIS{e‑n langar} \dicDirectTranslationCS{(kdo) touží, (kdo) chce, (kdo) má chuť} \dicExampleIS{Mig langar að fara.} \dicExampleCS{Chce se mi jít.};  \dicPhraseIS{e‑n langar í e‑ð} \dicDirectTranslationCS{(kdo) touží po (čem), (kdo) chce (co), (kdo) má chuť na (co)} \dicExampleIS{Mig langar í sælgæti.} \dicExampleCS{Mám chuť na něco sladkého.};  \dicPhraseIS{e‑n langar til e‑s} \dicDirectTranslationCS{(kdo) touží po (čem), (kdo) chce (co)} \dicExampleIS{Mig langar til að lesa bókina.} \dicExampleCS{Mám chuť si číst knížku.};  \dicPhraseIS{e‑n langar eftir e‑u} \dicDirectTranslationCS{(kdo) touží po (čem), (kdo) prahne po (čem)} \dicExampleIS{Hana langaði eftir friði.} \dicExampleCS{Toužila po klidu.}
\dicEntry[langafasta] \dicTerm{langa··fasta} \dicIPA{{l}{au}{\ng}{\r{g}}{a}{f}{a}{s}{\textsubring{d}}{a}} \dicPos{f}[1] \dicFlx{(lönguföstu, lönguföstur)}[21] \dicFieldCat{náb.} \dicDirectTranslationCS{půst} \dicIndirectTranslationCS{(trvající 7 týdnů od Popeleční středy do Velikonoc)}
\dicEntry[langafi] \dicTerm{lang··af|i} \dicIPA{{l}{au}{\ng}{\r{g}}{a}{v}{\textsci}} \dicPos{m}[1] \dicFlx{(‑a, ‑ar)}[8] \dicDirectTranslationCS{pradědeček} \dicIndirectTranslationCS{(otec dědečka nebo babičky)}
\dicEntry[langamma] \dicTerm{lang··|amma} \dicIPA{{l}{au}{\ng}{\r{g}}{a}{m}{a}} \dicPos{f}[1] \dicFlx{(‑ömmu, ‑ömmur)}[8] \dicDirectTranslationCS{prababička} \dicIndirectTranslationCS{(matka dědečka nebo babičky)}
\dicEntry[langanir] \dicTerm{langanir} \dicIPA{{l}{au}{\ng}{\r{g}}{a}{n}{\textsci}{\textsubring{r}}} \dicPos{f} \dicFlx{pl nom} \dicLink{löngun}
\dicEntry[langatöng] \dicTerm{langa··töng}\dicTerm{, löngutöng} \dicIPA{{l}\-{au}\-{\ng}\-{\r{g}}\-{a}\-{t\smash{\textsuperscript{h}}}\-{\oe i}\-{\ng}\-{\r{g}}\-} \dicPos{f}[8] \dicFlx{(löngutangar, löngutengur\,/\addthin löngutangir)}[10] \dicDirectTranslationCS{prostředník (prst), prostředníček}
\dicEntry[langbestur] \dicTerm{lang··bestur} \dicIPA{{l}{au}{\ng}{\r{g}}{\textsubring{b}}{\textepsilon}{s}{\textsubring{d}}{\textscy}{\textsubring{r}}} \dicPos{adj}[11]\dicFlx{}[-2] \dicFlx{m sg nom sup} \dicDirectTranslationCS{úplně\,/\addthin zdaleka nejlepší}
\dicEntry[langbylgja] \dicTerm{lang··bylgj|a} \dicIPA{{l}{au}{\ng}{\r{g}}{\textsubring{b}}{\textsci}{l}{\r{\textObardotlessj}}{a}} \dicPos{f}[1] \dicFlx{(‑u, ‑ur)}[25] \dicFieldCat{fyz.} \dicDirectTranslationCS{dlouhá vlna} \dicAntonym{stuttbylgja}
\dicEntry[langdreginn] \dicTerm{lang··dreginn} \dicIPA{{l}{au}{\ng}{\r{g}}{\textsubring{d}}{r}{ei}{\textsci}{\textsubring{n}}} \dicPos{adj}[6]\dicFlx{}[-2] \textbf{1.} \dicDirectTranslationCS{zdlouhavý, rozvláč\-ný} \dicExampleIS{langdreginn ræða} \dicExampleCS{zdlouhavý proslov}  \textbf{2.} \dicDirectTranslationCS{dlouhotrvající} \dicExampleIS{langdreginn sjúkdómur} \dicExampleCS{dlouhotrvající onemocnění}
\dicEntry[langdvöl] \dicTerm{lang··|dvöl} \dicIPA{{l}{au}{\ng}{\r{g}}{\textsubring{d}}{v}{\oe}{\textsubring{l}}} \dicPos{f}[7] \dicFlx{(‑dvalar, ‑dvalir)}[16] \dicDirectTranslationCS{dlouhotrvající pobyt};  \dicPhraseIS{dvelja\,/\addthin vera e‑s staðar langdvölum} \dicDirectTranslationCS{pobývat (kde) dlouho}
\dicEntry[langeygur] \dicTerm{lang··eygur} \dicIPA{{l}{au}{\ng}{\r{g}}{ei}{\textbabygamma}{\textscy}{\textsubring{r}}} \dicPos{adj}[1]\dicFlx{}[-1] \dicPhraseIS{vera langeygur eftir e‑u} \dicDirectTranslationCS{být nedočkavý (čeho), nemoct se dočkat (čeho)}
\dicEntry[langframi] \dicTerm{lang··fram|i} \dicIPA{{l}{au}{\ng}{\r{g}}{f}{r}{a}{m}{\textsci}} \dicPos{m}[1] \dicFlx{(‑a)}[3] \dicPhraseIS{til langframa} \dicFlx{adv} \dicDirectTranslationCS{dlouhodobě}
\dicEntry[langhlaup] \dicTerm{lang··hlaup} \dicIPA{{l}{au}{\ng}{\r{g}}{\textsubring{l}}{\oe i}{\textsubring{b}}} \dicPos{n}[2] \dicFlx{(‑s, ‑)}[5] \dicFieldCat{sport.} \dicDirectTranslationCS{dálkový běh}
\dicEntry[langlífi] \dicTerm{lang··lífi} \dicIPA{{l}{au}{\ng}{\r{g}}{l}{i}{v}{\textsci}} \dicPos{n}[2] \dicFlx{(‑s)}[20] \dicDirectTranslationCS{dlouhověkost}
\dicEntry[langlífur] \dicTerm{lang··lífur} \dicIPA{{l}{au}{\ng}{\r{g}}{l}{i}{v}{\textscy}{\textsubring{r}}} \dicPos{adj}[1]\dicFlx{}[-1] \dicDirectTranslationCS{dlouhověký}
\dicEntry[langloka] \dicTerm{lang··lok|a} \dicIPA{{l}{au}{\ng}{\r{g}}{l}{\textopeno}{\r{g}}{a}} \dicPos{f}[1] \dicFlx{(‑u, ‑ur)}[19] \textbf{1.} \dicSynonym{romsa\smash{\textsuperscript{1}}} \dicDirectTranslationCS{litanie, nudná a~zdlouhavá řeč, nudný a~zdlouhavý text}  \textbf{2.} \dicSynonym{samloka} \dicDirectTranslationCS{bageta}
\dicEntry[langorður] \dicTerm{lang··orður} \dicIPA{{l}{au}{\ng}{\r{g}}{\textopeno}{r}{ð}{\textscy}{\textsubring{r}}} \dicPos{adj}[2]\dicFlx{}[-1] \dicDirectTranslationCS{rozvláčný, upovídaný}
\dicEntry[langreyður] \dicTerm{lang··reyð|ur} \dicIPA{{l}{au}{\ng}{\r{g}}{r}{ei}{ð}{\textscy}{\textsubring{r}}} \dicPos{f}[5] \dicFlx{(‑ar, ‑ar)}[1] \dicFieldCat{zool.} \dicDirectTranslationCS{plejtvák myšok} \textit{(l.~{\textLA{Balaenoptera physalus}})}
\dicEntry[langrækinn] \dicTerm{lang··rækinn} \dicIPA{{l}{au}{\ng}{\r{g}}{r}{a}{i}{\r{\textObardotlessj}}{\textsci}{\textsubring{n}}} \dicPos{adj}[6]\dicFlx{}[-2] \dicDirectTranslationCS{neodpouštějící, nesmiřitelný}
\dicEntry[langskip] \dicTerm{lang··skip} \dicIPA{{l}{au}{\ng}{\r{g}}{s}{\r{\textObardotlessj}}{\textsci}{\textsubring{b}}} \dicPos{n}[2] \dicFlx{(‑s, ‑)}[5] \dicFieldCat{nám.} \dicDirectTranslationCS{vikinská válečná loď}
\dicEntry[langsnið] \dicTerm{lang··snið} \dicIPA{{l}{au}{\ng}{\r{g}}{s}{\textsubring{d}}{n}{\textsci}{\texttheta}} \dicPos{n}[2] \dicFlx{(‑s, ‑)}[5] \dicFieldCat{poč.} \dicDirectTranslationCS{formát na délku}
\dicEntry[langsóttur] \dicTerm{lang··sóttur} \dicIPA{{l}{au}{\ng}{\r{g}}{s}{ou}{h}{\textsubring{d}}{\textscy}{\textsubring{r}}} \dicPos{adj}[1]\dicFlx{}[-10] \dicDirectTranslationCS{nevěrohodný, přitažený za vlasy} \dicExampleIS{langsótt sönnun} \dicExampleCS{nevěrohodný důkaz}
\dicEntry[langspil] \dicTerm{lang··spil} \dicIPA{{l}{au}{\ng}{\r{g}}{s}{\textsubring{b}}{\textsci}{\textsubring{l}}} \dicPos{n}[2] \dicFlx{(‑s, ‑)}[5] \dicFieldCat{hud.} \dicIndirectTranslationCS{islandský strunný hudební nástroj}
\dicEntry[langstökk] \dicTerm{lang··stökk} \dicIPA{{l}{au}{\ng}{\r{g}}{s}{\textsubring{d}}{\oe}{h}{\r{g}}} \dicPos{n}[2] \dicFlx{(‑s, ‑)}[5] \dicFieldCat{sport.} \dicDirectTranslationCS{skok do dálky, skok daleký}
\dicEntry[langt] \dicTerm{langt} \dicsymFrequent\  \dicIPA{{l}{au}{\r{\ng}}{\textsubring{d}}} \dicPos{adv} \dicFlx{(comp lengra, sup lengst)} \dicDirectTranslationCS{daleko} \dicExampleIS{Þetta liggur langt inni í fortíðinni.} \dicExampleCS{Má to původ v~daleké minulosti.};  \dicPhraseIS{langt í burtu} \dicFlx{adv} \dicDirectTranslationCS{daleko};  \dicPhraseIS{langt síðan} \dicFlx{adv} \dicDirectTranslationCS{dávno}
\dicEntry[langtíma] \dicTerm{lang··tíma-} \dicIPA{{l}{au}{\ng}{\r{g}}{t\smash{\textsuperscript{h}}}{i}{m}{a}} \dicPos{predp} \dicDirectTranslationCS{dlouhodobý}
\dicEntry[langtímalán] \dicTerm{lang·tíma··lán} \dicIPA{{l}{au}{\ng}{\r{g}}{t\smash{\textsuperscript{h}}}{i}{m}{a}{l}{au}{\textsubring{n}}} \dicPos{n}[2] \dicFlx{(‑s, ‑)}[5] \dicFieldCat{ekon.} \dicDirectTranslationCS{dlouhodobá půjčka}
\dicEntry[langtímasamningur] \dicTerm{lang·tíma··samn·ing|ur} \dicIPA{{l}\-{au}\-{\ng}\-{\r{g}}\-{t\smash{\textsuperscript{h}}}\-{i}\-{m}\-{a}\-{s}\-{a}\-{m}\-{n}\-{i}\-{\ng}\-{\r{g}}\-{\textscy}\-{\textsubring{r}}\-} \dicPos{m}[6] \dicFlx{(‑s, ‑ar)}[8] \dicDirectTranslationCS{dlouhodobá smlouva}
\dicEntry[langur] \dicTerm{lang|ur\smash{\textsuperscript{1}}} \dicIPA{{l}{au}{\ng}{\r{g}}{\textscy}{\textsubring{r}}} \dicPos{m}[6] \dicFlx{(‑s, ‑ar)}[16] \dicPhraseIS{draga e‑ð á langinn} \dicDirectTranslationCS{odložit (co), odkládat (co)}
\dicEntry[langur] \dicTerm{langur\smash{\textsuperscript{2}}} \dicsymFrequent\  \dicIPA{{l}{au}{\ng}{\r{g}}{\textscy}{\textsubring{r}}} \dicPos{adj}[10] \dicFlx{(f löng, comp lengri, sup lengstur)}[2] \mbox{\textbf{1.} \dicDirectTranslationCS{dlouhý}} \dicIndirectTranslationCS{(o~vzdálenosti)} \dicExampleIS{löng leið} \dicExampleCS{dlouhá cesta} \dicAntonym{stuttur}  \textbf{2.} \dicDirectTranslationCS{vysoký} \dicExampleIS{langur maður} \dicExampleCS{vysoký člověk}  \textbf{3.} \dicDirectTranslationCS{dlouhý} \dicIndirectTranslationCS{(o~čase)} \dicExampleIS{langur tími} \dicExampleCS{dlouhý čas};  \dicPhraseIS{fyrir löngu} \dicFlx{adv} \dicDirectTranslationCS{dávno, před dlouhou dobou};  \dicPhraseIS{löngum stundum} \dicFlx{adv} \dicDirectTranslationCS{dlouho, dlouhé hodiny};  \dicPhraseIS{það er langt til e‑s} \dicDirectTranslationCS{je daleko do (čeho) (Vánoc ap.)}
\dicEntry[langvarandi] \dicTerm{lang··var·andi} \dicsymFrequent\  \dicIPA{{l}{au}{\ng}{\r{g}}{v}{a}{r}{a}{n}{\textsubring{d}}{\textsci}} \dicPos{adj}[13] \dicFlx{indecl}[1] \dicDirectTranslationCS{dlouhotrvající, prodlužovaný} \dicExampleIS{langvarandi erfiðleikar} \dicExampleCS{dlouhotrvající potíže}
\dicEntry[langvinnur] \textls[15]{\dicTerm{lang··vinnur} \dicIPA{{l}{au}{\ng}{\r{g}}{v}{\textsci}{n}{\textscy}{\textsubring{r}}} \dicPos{adj}[1]\dicFlx{}[-1] \dicDirectTranslationCS{přetrvávající, táhlý, dlouhotrvající} \dicExampleIS{langvinnur sjúkdómur} \dicExampleCS{přetrvávající onemocnění}}
\dicEntry[langvía] \dicTerm{lang··ví|a} \dicIPA{{l}{au}{\ng}{\r{g}}{v}{i}{j}{a}} \dicPos{f}[1] \dicFlx{(‑u, ‑ur)}[7] \dicFieldCat{zool.} \dicDirectTranslationCS{alkoun úzkozobý} \textit{(l.~{\textLA{Uria aalge}})}  \dicsymPhoto\ 
\dicFigure{20877.jpg}{Langvía}{Langvía - Sowls, Art, Biolib, PD}
\dicEntry[langþráður] \dicTerm{lang··þráður} \dicIPA{{l}{au}{\ng}{\r{g}}{\texttheta}{r}{au}{ð}{\textscy}{\textsubring{r}}} \dicPos{adj}[2]\dicFlx{}[-12] \dicDirectTranslationCS{vytoužený, (dlouho) očekávaný}
\dicEntry[langþreyttur] \dicTerm{lang··þreyttur} \dicIPA{{l}{au}{\ng}{\r{g}}{\texttheta}{r}{ei}{h}{\textsubring{d}}{\textscy}{\textsubring{r}}} \dicPos{adj}[1]\dicFlx{}[-13] \dicDirectTranslationCS{znavený, (velmi) unavený}
\dicEntry[lantan] \dicTerm{lantan} \dicIPA{{l}{a}{\textsubring{n}}{\textsubring{d}}{a}{\textsubring{n}}} \dicPos{n}[2] \dicFlx{(‑s)}[2] \dicFieldCat{chem.} \dicDirectTranslationCS{lanthan} \textit{(l.~{\textLA{La, Lanthanum}})}
\dicEntry[Laos] \dicTerm{Laos} \dicIPA{{l}{a}{\textlengthmark}{\textopeno}{s}} \dicPos{n}[4] \dicFlx{indecl}[2] \dicFieldCat{geog.} \dicDirectTranslationCS{Laos}
\dicEntry[Laosi] \dicTerm{Laos|i} \dicIPA{{l}{a}{\textlengthmark}{\textopeno}{s}{\textsci}} \dicPos{m}[1] \dicFlx{(‑a, ‑ar)}[1] \dicDirectTranslationCS{Laosan(ka)}
\dicEntry[laoska] \dicTerm{laosk|a} \dicIPA{{l}{a}{\textlengthmark}{\textopeno}{s}{\r{g}}{a}} \dicPos{f}[1] \dicFlx{(‑u)}[5] \dicDirectTranslationCS{laoština}
\dicEntry[laoskur] \dicTerm{laoskur} \dicIPA{{l}{a}{\textlengthmark}{\textopeno}{s}{\r{g}}{\textscy}{\textsubring{r}}} \dicPos{adj}[1]\dicFlx{}[-6] \dicDirectTranslationCS{laoský}
\dicEntry[lapið] \dicTerm{lapið} \dicIPA{{l}{a}{\textlengthmark}{\textsubring{b}}{\textsci}{\texttheta}} \dicPos{v} \dicFlx{supin} \dicLink{lepja}
\dicEntry[lappa] \dicTerm{lapp|a} \dicIPA{{l}{a}{h}{\textsubring{b}}{a}} \dicPos{v}[1] \dicFlx{(‑aði)}[13] \dicDirectTranslationCS{ochodit, prochodit (boty ap.)} \dicExampleIS{lappa skóna} \dicExampleCS{prochodit boty};  \dicPhraseIS{e‑að lappar e‑n} \dicDirectTranslationCS{(co) je (komu) volné (boty ap.)} \dicExampleIS{Skórinn lappar mig.} \dicExampleCS{Boty mi jsou volné.};  \dicIdiom{lappa}[upp á]{ \dicPhraseIS{lappa upp á e‑ð}} \dicDirectTranslationCS{osvěžit (co), vyspravit (co)}
\dicEntry[lappajaðrakan] \dicTerm{lappa··jað·rakan} \dicIPA{{l}{a}{h}{\textsubring{b}}{a}{j}{a}{ð}{r}{a}{\r{g}}{a}{\textsubring{n}}} \dicPos{m}[4] \dicFlx{(‑s, ‑ar)}[12] \dicFieldCat{zool.} \dicDirectTranslationCS{břehouš rudý} \textit{(l.~{\textLA{Limosa lapponica}})}  \dicsymPhoto\ 
\dicFigure{ds_image_lappajadrakan_0_1.jpg}{Lappajaðrakan}{Lappajaðrakan - Andreas Trepte, CC BY-SA 2.5}
\dicEntry[lappalaus] \dicTerm{lappa··laus} \dicIPA{{l}{a}{h}{\textsubring{b}}{a}{l}{\oe i}{s}} \dicPos{adj}[5]\dicFlx{}[-1] \dicPhraseIS{vera orðinn lappalaus} \dicLangCat{přen.} \dicDirectTranslationCS{být k~smrti unavený}
\dicEntry[lappar] \dicTerm{lappar} \dicIPA{{l}{a}{h}{\textsubring{b}}{a}{\textsubring{r}}} \dicPos{f} \dicFlx{sg gen} \dicLink{löpp}
\dicEntry[Lappi] \dicTerm{Lapp|i} \dicIPA{{l}{a}{h}{\textsubring{b}}{\textsci}} \dicPos{m}[1] \dicFlx{(‑a, ‑ar)}[8] \dicDirectTranslationCS{Laponec, Laponka}
\dicEntry[lappir] \dicTerm{lappir} \dicIPA{{l}{a}{h}{\textsubring{b}}{\textsci}{\textsubring{r}}} \dicPos{f} \dicFlx{pl nom} \dicLink{löpp}
\dicEntry[Lappland] \dicTerm{Lapp··land} \dicIPA{{l}{a}{h}{\textsubring{b}}{l}{a}{n}{\textsubring{d}}} \dicPos{n}[2] \dicFlx{(‑s)}[4] \dicFieldCat{geog.} \dicSynonym{Samaland} \dicDirectTranslationCS{Laponsko}
\dicEntry[lappneska] \dicTerm{lapp··nesk|a} \dicIPA{{l}{a}{h}{\textsubring{b}}{n}{\textepsilon}{s}{\r{g}}{a}} \dicPos{f}[1] \dicFlx{(‑u)}[5] \dicDirectTranslationCS{laponština}
\dicEntry[lappneskur] \dicTerm{lapp··neskur} \dicIPA{{l}{a}{h}{\textsubring{b}}{n}{\textepsilon}{s}{\r{g}}{\textscy}{\textsubring{r}}} \dicPos{adj}[1]\dicFlx{}[-1] \dicDirectTranslationCS{laponský}
\dicEntry[lapti] \dicTerm{lapti} \dicIPA{{l}{a}{f}{\textsubring{d}}{\textsci}} \dicPos{v} \dicFlx{ind pf sg 1 pers} \dicLink{lepja}
\dicEntry[las] \dicTerm{las} \dicIPA{{l}{a}{\textlengthmark}{s}} \dicPos{v} \dicFlx{ind pf sg 1 pers} \dicLink{lesa}
\dicEntry[lasanja] \dicTerm{lasanja} \dicIPA{{l}{a}{\textlengthmark}{s}{a}{n}{j}{a}} \dicPos{n}[4] \dicFlx{indecl}[1] \dicFieldCat{kulin.} \dicDirectTranslationCS{lasagne}
\dicEntry[lasburða] \dicTerm{las··burða} \dicIPA{{l}{a}{\textlengthmark}{s}{\textsubring{b}}{\textscy}{r}{ð}{a}} \dicPos{adj}[13] \dicFlx{indecl}[1] \dicDirectTranslationCS{churavý, neduživý}
\dicEntry[lasinn] \dicTerm{lasinn} \dicIPA{{l}{a}{\textlengthmark}{s}{\textsci}{\textsubring{n}}} \dicPos{adj}[6]\dicFlx{}[-3] \textbf{1.} \dicSynonym{veikur} \dicDirectTranslationCS{nemocný, indisponovaný}  \textbf{2.} \dicSynonym{lélegur} \dicDirectTranslationCS{opotřebovaný}
\dicEntry[laska] \dicTerm{lask|a} \dicIPA{{l}{a}{s}{\r{g}}{a}} \dicPos{v}[1] \dicFlx{(‑aði)}[13] \dicFlx{acc} \dicSynonym{skadda} \dicDirectTranslationCS{poškodit, způsobit škodu}
\dicEntry[lasleiki] \dicTerm{las··leik|i} \dicIPA{{l}{a}{\textlengthmark}{s}{l}{ei}{\r{\textObardotlessj}}{\textsci}} \dicPos{m}[1] \dicFlx{(‑a)}[3] \dicDirectTranslationCS{indispozice, churavost}
\dicEntry[last] \dicTerm{last} \dicIPA{{l}{a}{s}{\textsubring{d}}} \dicPos{n}[2] \dicFlx{(‑s)}[2] \dicSynonym*{baktal} \dicDirectTranslationCS{pomluva, osočení};  \dicPhraseIS{leggja e‑m e‑ð til lasts} \dicDirectTranslationCS{dávat (co komu) za vinu}
\dicEntry[lasta] \dicTerm{last|a} \dicIPA{{l}{a}{s}{\textsubring{d}}{a}} \dicPos{v}[1] \dicFlx{(‑aði)}[13] \dicFlx{acc} \dicSynonym{gagnrýna} \dicDirectTranslationCS{očernit, očerňovat, (po)hanit, pohanět} \dicExampleIS{lasta þýðingu á bók} \dicExampleCS{očerňovat překlad knížky}
\dicEntry[lastar] \dicTerm{lastar} \dicIPA{{l}{a}{s}{\textsubring{d}}{a}{\textsubring{r}}} \dicPos{m} \dicFlx{sg gen} \dicLink{löstur}
\dicEntry[lastyrði] \dicTerm{last··yrði} \dicIPA{{l}{a}{s}{\textsubring{d}}{\textsci}{r}{ð}{\textsci}} \dicPos{n}[2] \dicFlx{(‑s, ‑)}[14] \dicSynonym{last} \dicDirectTranslationCS{hana, hanobení}
\dicEntry[latína] \dicTerm{latín|a} \dicIPA{{l}{a}{\textlengthmark}{t\smash{\textsuperscript{h}}}{i}{n}{a}} \dicPos{f}[1] \dicFlx{(‑u)}[5] \dicDirectTranslationCS{latina}
\dicEntry[latínusegl] \dicTerm{latínu··segl} \dicIPA{{l}{a}{\textlengthmark}{t\smash{\textsuperscript{h}}}{i}{n}{\textscy}{s}{\textepsilon}{\r{g}}{\textsubring{l}}} \dicPos{n}[2] \dicFlx{(‑s, ‑)}[5] \dicFieldCat{nám.} \dicDirectTranslationCS{latinská plachta}
\dicEntry[latneskur] \dicTerm{latn··eskur} \dicIPA{{l}{a}{h}{\textsubring{d}}{n}{\textepsilon}{s}{\r{g}}{\textscy}{\textsubring{r}}} \dicPos{adj}[1]\dicFlx{}[-1] \dicDirectTranslationCS{latinský}
\dicEntry[latt] \dicTerm{latt} \dicIPA{{l}{a}{h}{\textsubring{d}}} \dicPos{v} \dicFlx{supin} \dicLink{letja}
\dicEntry[latti] \dicTerm{latti} \dicIPA{{l}{a}{h}{\textsubring{d}}{\textsci}} \dicPos{v} \dicFlx{ind pf sg 1 pers} \dicLink{letja}
\dicEntry[latur] \dicTerm{latur} \dicIPA{{l}{a}{\textlengthmark}{\textsubring{d}}{\textscy}{\textsubring{r}}} \dicPos{adj}[1] \dicFlx{(f löt)}[2] \dicDirectTranslationCS{líný, lenivý}
\dicEntry[lauf] \dicTerm{lauf} \dicsymFrequent\  \dicIPA{{l}{\oe i}{\textlengthmark}{f}} \dicPos{n}[2] \dicFlx{(‑s, ‑)}[5] \textbf{1.} \dicSynonym{laufblað} \dicDirectTranslationCS{list (stromu ap.)} \dicExampleIS{gulnuð lauf að hausti} \dicExampleCS{zežloutlé listy na podzim}  \textbf{2.} \dicSynonym*{litur í spilum} \dicDirectTranslationCS{trefy, kříže} \dicIndirectTranslationCS{(barva v~kartách)}
\dicEntry[laufabrauð] \dicTerm{laufa··brauð} \dicIPA{{l}{\oe i}{\textlengthmark}{v}{a}{\textsubring{b}}{r}{\oe i}{\texttheta}} \dicPos{n}[2] \dicFlx{(‑s, ‑)}[5] \dicFieldCat{kulin.} \dicIndirectTranslationCS{islandský vánoční chléb ručně zdobený různými motivy} \dicsymPhoto\ 
\dicFigure{ds_image_laufabraud_0_2.jpg}{Laufabrauð}{Laufabrauð - Fany Larota Catunta, PD}
\dicEntry[laufblað] \dicTerm{lauf··|blað} \dicIPA{{l}{\oe i}{v}{\textsubring{b}}{l}{a}{\texttheta}} \dicPos{n}[2] \dicFlx{(‑blaðs, ‑blöð)}[8] \dicSynonym{lauf} \dicDirectTranslationCS{list (stromu ap.)}
\dicEntry[laufgast] \dicTerm{laufg|ast} \dicIPA{{l}{\oe i}{v}{\r{g}}{a}{s}{\textsubring{d}}} \dicPos{v}[1] \dicFlx{(‑aðist)}[95] \dicFlx{refl} \dicDirectTranslationCS{nasadit listy, zazelenat se}
\dicEntry[laufglói] \dicTerm{lauf··gló|i} \dicIPA{{l}{\oe i}{v}{\r{g}}{l}{ou}{\textsci}} \dicPos{m}[1] \dicFlx{(‑a, ‑ar)}[1] \dicFieldCat{zool.} \dicSynonym{gullþröstur} \dicDirectTranslationCS{žluva, žluva hajní} \textit{(l.~{\textLA{Oriolus oriolus}})}  \dicsymPhoto\ 
\dicFigure{ds_image_laufgloi_0_1.jpg}{Laufglói}{Laufglói - Dixi, GFDL}
\dicEntry[laufgræna] \dicTerm{lauf··græn|a} \dicIPA{{l}{\oe i}{v}{\r{g}}{r}{a}{i}{n}{a}} \dicPos{f}[1] \dicFlx{(‑u)}[5] \dicFieldCat{bot.} \dicSynonym{blaðgræna} \dicDirectTranslationCS{chlorofyl}
\dicEntry[laufskógur] \dicTerm{lauf··skóg|ur} \dicIPA{{l}{\oe i}{f}{s}{\r{g}}{ou}{\textscy}{\textsubring{r}}} \dicPos{m}[6] \dicFlx{(‑ar, ‑ar)}[56] \dicDirectTranslationCS{listnatý les}
\dicEntry[laufsöngvari] \dicTerm{lauf··söngv·ar|i} \dicIPA{{l}{\oe i}{f}{s}{\oe i}{\ng}{\r{g}}{v}{a}{r}{\textsci}} \dicPos{m}[1] \dicFlx{(‑a, ‑ar)}[13] \dicFieldCat{zool.} \dicDirectTranslationCS{budníček větší} \textit{(l.~{\textLA{Phylloscopus trochilus}})}  \dicsymPhoto\ 
\dicFigure{ds_image_laufsongvari_0_1.jpg}{Laufsöngvari}{Laufsöngvari - Vogelartinfo, CC BY-SA 3.0}
\dicEntry[lauftré] \dicTerm{lauf··tré} \dicIPA{{l}{\oe i}{v}{t\smash{\textsuperscript{h}}}{r}{j}{\textepsilon}} \dicPos{n}[2] \dicFlx{(‑s, ‑)}[36] \dicDirectTranslationCS{listnatý strom, listnáč}
\dicEntry[laug] \dicTerm{laug\smash{\textsuperscript{1}}} \dicIPA{{l}{\oe i}{\textlengthmark}{x}} \dicPos{f}[4] \dicFlx{(‑ar, ‑ar)}[1] \textbf{1.} \dicSynonym{sundlaug} \dicDirectTranslationCS{(plavecký) bazén, koupaliště, plovárna}  \textbf{2.} \dicSynonym*{heit lind} \dicDirectTranslationCS{vřídlo, horký pramen}
\dicEntry[laug] \dicTerm{laug\smash{\textsuperscript{2}}} \dicIPA{{l}{\oe i}{\textlengthmark}{x}} \dicPos{v} \dicFlx{ind pf sg 1 pers} \dicLink{ljúga}
\dicEntry[lauga] \dicTerm{laug|a} \dicIPA{{l}{\oe i}{\textlengthmark}{\textbabygamma}{a}} \dicPos{v}[1] \dicFlx{(‑aði)}[1] \dicFlx{acc} \dicLangCat{zast.} \dicSynonym{þvo} \dicDirectTranslationCS{(vy)koupat} \dicExampleIS{lauga barn} \dicExampleCS{umýt dítě}
\dicEntry[laugardagskvöld] \dicTerm{laugar·dags··kvöld} \dicIPA{{l}{\oe i}{\textlengthmark}{\textbabygamma}{a}{r}{\textsubring{d}}{a}{x}{s}{k\smash{\textsuperscript{h}}}{v}{\oe}{l}{\textsubring{d}}} \dicPos{n}[2] \dicFlx{(‑s, ‑)}[5] \dicDirectTranslationCS{sobotní večer, sobota večer}
\dicEntry[laugardagur] \dicTerm{laugar··dag|ur} \dicIPA{{l}{\oe i}{\textlengthmark}{\textbabygamma}{a}{r}{\textsubring{d}}{a}{\textbabygamma}{\textscy}{\textsubring{r}}} \dicPos{m}[6] \dicFlx{(‑s, ‑ar)}[62] \dicDirectTranslationCS{sobota}
\dicEntry[Laugavegur] \dicTerm{Lauga··veg|ur} \dicIPA{{l}{\oe i}{\textlengthmark}{\textbabygamma}{a}{v}{\textepsilon}{\textbabygamma}{\textscy}{\textsubring{r}}} \dicPos{m}[10] \dicFlx{(‑ar\,/\addthin ‑s)}[28] \dicFlx{prop} \dicDirectTranslationCS{Laugavegur} \dicIndirectTranslationCS{(hlavní nákupní třída v~Reykjavíku)}
\dicEntry[lauk] \dicTerm{lauk} \dicIPA{{l}{\oe i}{\textlengthmark}{\r{g}}} \dicPos{v} \dicFlx{ind pf sg 1 pers} \dicLink{ljúka}
\dicEntry[laukur] \dicTerm{lauk|ur} \dicIPA{{l}{\oe i}{\textlengthmark}{\r{g}}{\textscy}{\textsubring{r}}} \dicPos{m}[6] \dicFlx{(‑s, ‑ar)}[22] \dicFieldCat{bot.} \dicDirectTranslationCS{cibule, cibule kuchyňská} \textit{(l.~{\textLA{Allium cepa}})};  \dicPhraseIS{laukur ættar} \dicLangCat{přen.} \dicDirectTranslationCS{pýcha rodiny}
\dicEntry[laum] \dicTerm{laum} \dicIPA{{l}{\oe i}{\textlengthmark}{\textsubring{m}}} \dicPos{n}[2] \dicFlx{(‑s)}[2] \dicPhraseIS{í laumi} \dicFlx{adv} \dicDirectTranslationCS{potají, v~tajnosti, tajně}
\dicEntry[lauma] \dicTerm{laum|a} \dicsymFrequent\  \dicIPA{{l}{\oe i}{\textlengthmark}{m}{a}} \dicPos{v}[1] \dicFlx{(‑aði)}[1] \dicFlx{dat} \dicDirectTranslationCS{(nenápadně) dát, vsunout, podstrčit, propašovat} \dicExampleIS{lauma bréfi að e‑m} \dicExampleCS{nenápadně (komu) podstrčit dopis};  \dicIdiom{laumast}{ \dicPhraseIS{laumast}} \dicFlx{refl} \dicSynonym*{læðast} \dicDirectTranslationCS{(v)plížit se, (v)krást se} \dicExampleIS{laumast úr bænum} \dicExampleCS{vykrást se z~města}
\dicEntry[laumufarþegi] \dicTerm{laumu··far·þeg|i} \dicIPA{{l}{\oe i}{\textlengthmark}{m}{\textscy}{f}{a}{\textsubring{r}}{\texttheta}{ei}{\textsci}} \dicPos{m}[1] \dicFlx{(‑a, ‑ar)}[1] \dicDirectTranslationCS{černý pasažér, černá pasažérka}
\dicEntry[laumulegur] \dicTerm{laumu··legur} \dicIPA{{l}{\oe i}{\textlengthmark}{m}{\textscy}{l}{\textepsilon}{\textbabygamma}{\textscy}{\textsubring{r}}} \dicPos{adj}[1]\dicFlx{}[-8] \dicDirectTranslationCS{tajný, skrytý, kradmý}
\dicEntry[laumuspil] \dicTerm{laumu··spil} \dicIPA{{l}{\oe i}{\textlengthmark}{m}{\textscy}{s}{\textsubring{b}}{\textsci}{\textsubring{l}}} \dicPos{n}[2] \dicFlx{(‑s, ‑)}[5] \dicDirectTranslationCS{tajnosti, tajnůstkaření}
\dicEntry[laun] \dicTerm{laun\smash{\textsuperscript{1}}} \dicIPA{{l}{\oe i}{\textlengthmark}{\textsubring{n}}} \dicPos{f}[4] \dicFlx{(‑ar)}[3] \dicSynonym{leynd} \dicDirectTranslationCS{tajemství};  \dicPhraseIS{á laun} \dicFlx{adv} \dicDirectTranslationCS{potají, v~tajnosti}
\dicEntry[laun] \dicTerm{laun\smash{\textsuperscript{2}}} \dicsymFrequent\  \dicIPA{{l}{\oe i}{\textlengthmark}{\textsubring{n}}} \dicPos{n}[2] \dicFlx{pl}[1] \textbf{1.} \dicSynonym{kaup} \dicDirectTranslationCS{plat, výplata, mzda} \dicExampleIS{há laun} \dicExampleCS{vysoký plat}  \textbf{2.} \dicSynonym{endurgjald} \dicDirectTranslationCS{vděk, odměna}  \textbf{3.} \dicSynonym{borgun} \dicDirectTranslationCS{(peněžitá) odměna, honorář}
\dicEntry[launa] \dicTerm{laun|a} \dicIPA{{l}{\oe i}{\textlengthmark}{n}{a}} \dicPos{v}[1] \dicFlx{(‑aði)}[1] \dicFlx{dat + acc} \textbf{1.} \dicSynonym{endurgjalda} \dicDirectTranslationCS{odplatit, oplatit, odvděčit se};  \dicPhraseIS{launa e‑m e‑ð} \dicDirectTranslationCS{odvděčit se (komu čím), oplatit (komu čím)}  \textbf{2.} \dicDirectTranslationCS{ohodnotit, honorovat, platit} \dicExampleIS{illa launað starf} \dicExampleCS{špatně honorovaná práce}
\dicEntry[launaauki] \dicTerm{launa··auk|i} \dicIPA{{l}{\oe i}{\textlengthmark}{n}{a}{\oe i}{\r{\textObardotlessj}}{\textsci}} \dicPos{m}[1] \dicFlx{(‑a)}[3] \dicDirectTranslationCS{příplatek\,/\addthin přídavek k~platu, prémie, bonus}
\dicEntry[launadeila] \dicTerm{launa··deil|a} \dicIPA{{l}{\oe i}{\textlengthmark}{n}{a}{\textsubring{d}}{ei}{l}{a}} \dicPos{f}[1] \dicFlx{(‑u, ‑ur)}[13] \dicDirectTranslationCS{mzdový spor}
\dicEntry[launaflokkur] \dicTerm{launa··flokk|ur} \dicIPA{{l}{\oe i}{\textlengthmark}{n}{a}{f}{l}{\textopeno}{h}{\r{g}}{\textscy}{\textsubring{r}}} \dicPos{m}[6] \dicFlx{(‑s, ‑ar)}[8] \dicDirectTranslationCS{platové zařazení, platová skupina\,/\addthin třída}
\dicEntry[launagreiðsla] \dicTerm{launa··greiðsl|a} \dicIPA{{l}{\oe i}{\textlengthmark}{n}{a}{\r{g}}{r}{ei}{ð}{s}{\textsubring{d}}{l}{a}} \dicPos{f}[1] \dicFlx{(‑u, ‑ur)}[13] \dicDirectTranslationCS{výplata mzdy}
\dicEntry[launahækkun] \dicTerm{launa··hækk|un} \dicIPA{{l}{\oe i}{\textlengthmark}{n}{a}{h}{a}{i}{h}{\r{g}}{\textscy}{\textsubring{n}}} \dicPos{f}[7] \dicFlx{(‑unar, ‑anir)}[8] \dicDirectTranslationCS{zvýšení platu, mzdový růst}
\dicEntry[launakjör] \dicTerm{launa··kjör} \dicIPA{{l}{\oe i}{\textlengthmark}{n}{a}{c\smash{\textsuperscript{h}}}{\oe}{\textsubring{r}}} \dicPos{n}[2] \dicFlx{pl}[9] \dicDirectTranslationCS{platové\,/\addthin mzdové podmínky}
\begin{xtolerant}{}{1pt}
\dicEntry[launakostnaður] \dicTerm{launa··kost·nað|ur}\addthinS\dicIPA{{l}{\oe i}{\textlengthmark}{n}{a}{k\smash{\textsuperscript{h}}}{\textopeno}{s}{\textsubring{d}}{n}{a}{ð}{\textscy}{\textsubring{r}}}\addthinS\dicPos{m}[10]\addthinS\dicFlx{(‑ar)}[9] \dicDirectTranslationCS{mzdové náklady}
\end{xtolerant}
\dicEntry[launalækkun] \dicTerm{launa··lækk|un} \dicIPA{{l}{\oe i}{\textlengthmark}{n}{a}{l}{a}{i}{h}{\r{g}}{\textscy}{\textsubring{n}}} \dicPos{f}[7] \dicFlx{(‑unar)}[9] \dicDirectTranslationCS{snížení platu, mzdový pokles}
\dicEntry[launamaður] \dicTerm{launa··|maður} \dicIPA{{l}{\oe i}{\textlengthmark}{n}{a}{m}{a}{ð}{\textscy}{\textsubring{r}}} \dicPos{m}[13] \dicFlx{(‑manns, ‑menn)}[2] \dicDirectTranslationCS{zaměstnanec, zaměstnankyně}
\dicEntry[launaskrá] \dicTerm{launa··skrá} \dicIPA{{l}{\oe i}{\textlengthmark}{n}{a}{s}{\r{g}}{r}{au}} \dicPos{f}[4] \dicFlx{(‑r\,/\addthin ‑ar, ‑r)}[21] \dicDirectTranslationCS{výplatní listina}
\dicEntry[launaskrið] \dicTerm{launa··skrið} \dicIPA{{l}{\oe i}{\textlengthmark}{n}{a}{s}{\r{g}}{r}{\textsci}{\texttheta}} \dicPos{n}[2] \dicFlx{(‑s)}[2] \dicFieldCat{ekon.} \dicDirectTranslationCS{stoupání mezd}
\dicEntry[launastigi] \dicTerm{launa··stig|i} \dicIPA{{l}{\oe i}{\textlengthmark}{n}{a}{s}{\textsubring{d}}{i}{j}{\textsci}} \dicPos{m}[1] \dicFlx{(‑a)}[3] \dicDirectTranslationCS{mzdový tarif}
\dicEntry[launatengdur] \dicTerm{launa··tengdur} \dicIPA{{l}{\oe i}{\textlengthmark}{n}{a}{t\smash{\textsuperscript{h}}}{\textepsilon}{\ng}{\textsubring{d}}{\textscy}{\textsubring{r}}} \dicPos{adj}[2]\dicFlx{}[-17] \dicDirectTranslationCS{mzdový, spojený se mzdou (náklady ap.)};  \dicPhraseIS{launatengd gjöld} \dicDirectTranslationCS{příspěvky zaměstnavatele}
\dicEntry[launauppbót] \dicTerm{launa··upp·|bót} \dicIPA{{l}{\oe i}{\textlengthmark}{n}{a}{\textscy}{h}{\textsubring{b}}{ou}{\textsubring{d}}} \dicPos{f}[8] \dicFlx{(‑bótar, ‑bætur)}[5] \dicDirectTranslationCS{platový bonus}
\dicEntry[launráð] \dicTerm{laun··ráð} \dicIPA{{l}{\oe i}{n}{r}{au}{\texttheta}} \dicPos{n}[2] \dicFlx{pl}[1] \dicSynonym{ráðabrugg} \dicDirectTranslationCS{pikle, intriky, komplot}
\dicEntry[launsátur] \dicTerm{laun··sátur} \dicIPA{{l}{\oe i}{n}{s}{au}{\textsubring{d}}{\textscy}{\textsubring{r}}} \dicPos{n}[2] \dicFlx{(‑s, ‑)}[25] \dicSynonym{fyrirsát} \dicDirectTranslationCS{léčka, nástraha}
\dicEntry[launung] \dicTerm{laun··ung} \dicIPA{{l}{\oe i}{\textlengthmark}{n}{u}{\ng}{\r{g}}} \dicPos{f}[4] \dicFlx{(‑ar)}[3] \dicSynonym{leynd} \dicDirectTranslationCS{tajemství, tajnost};  \dicPhraseIS{með launung} \dicFlx{adv} \dicSynonym*{leynilega} \dicDirectTranslationCS{tajně};  \dicPhraseIS{e‑m er engin launung á e‑u} \dicFlx{impers} \dicDirectTranslationCS{(kdo) se netají (čím)}
\dicEntry[launþegi] \dicTerm{laun··þeg|i} \dicIPA{{l}{\oe i}{n}{\texttheta}{ei}{\textsci}} \dicPos{m}[1] \dicFlx{(‑a, ‑ar)}[1] \dicDirectTranslationCS{zaměstnanec, zaměstnankyně, pracovník, pracovnice}
\dicEntry[laus] \dicTerm{laus} \dicsymFrequent\  \dicIPA{{l}{\oe i}{\textlengthmark}{s}} \dicPos{adj}[5] \dicFlx{(f ‑)}[1] \textbf{1.} \dicSynonym{óbundinn} \dicDirectTranslationCS{volný, uvolněný, povolený} \dicExampleIS{Steinninn er orðinn laus frá veggnum.} \dicExampleCS{Kámen se už uvolnil ze zdi.}  \textbf{2.} \dicSynonym{frjáls} \dicDirectTranslationCS{volný, svobodný} \dicExampleIS{Ertu laus í kvöld?} \dicExampleCS{Máš večer čas?};  \dicPhraseIS{ganga laus} \dicDirectTranslationCS{jít volně (pes ap.)};  \dicPhraseIS{láta e‑n lausan} \dicDirectTranslationCS{pustit (koho) na svobodu};  \dicPhraseIS{vera laus við e‑ð} \dicDirectTranslationCS{zbavit se (čeho), oprostit se od (čeho)}  \textbf{3.} \dicSynonym*{auður} \dicDirectTranslationCS{volný, prázdný, neobsazený} \dicExampleIS{Er þetta sæti laust?} \dicExampleCS{Je toto místo volné?} \dicAntonym{upptekinn}  \textbf{4.} \dicSynonym{léttur} \dicDirectTranslationCS{lehký} \dicExampleIS{laust högg} \dicExampleCS{lehká rána}  \textbf{5.} \dicSynonym{óáreiðanlegur} \dicDirectTranslationCS{nespolehlivý, nestabilní};  \dicPhraseIS{vera á lausu} \dicDirectTranslationCS{být nezadaný};  \dicPhraseIS{vera laus á kostunum} \dicSynonym{lauslátur} \dicDirectTranslationCS{mít volné mravy} \dicIndirectTranslationCS{(zvláště o~ženě)}  \textbf{6.} \dicSynonym{örlátur} \dicDirectTranslationCS{štědrý};  \dicPhraseIS{vera laus á fé} \dicDirectTranslationCS{být štědrý}
\dicEntry[lausafé] \dicTerm{lausa··|fé} \dicIPA{{l}{\oe i}{\textlengthmark}{s}{a}{f}{j}{\textepsilon}} \dicPos{n}[3] \dicFlx{(‑fjár)}[5] \dicDirectTranslationCS{hotovost}
\dicEntry[lausafjárstaða] \dicTerm{lausa··fjár·|staða} \dicIPA{{l}{\oe i}{\textlengthmark}{s}{a}{f}{j}{au}{\textsubring{r}}{s}{\textsubring{d}}{a}{ð}{a}} \dicPos{f}[1] \dicFlx{(‑stöðu)}[2] \dicFieldCat{ekon.} \dicDirectTranslationCS{likvidita}
\dicEntry[lausaleikur] \dicTerm{lausa··leik|ur} \dicIPA{{l}{\oe i}{\textlengthmark}{s}{a}{l}{ei}{\r{g}}{\textscy}{\textsubring{r}}} \dicPos{m}[6] \dicFlx{(‑s)}[17] \dicPhraseIS{eiga\,/\addthin eignast barn í lausaleik} \dicDirectTranslationCS{mít nemanželské dítě}
\dicEntry[lausamennska] \dicTerm{lausa··mennsk|a} \dicIPA{{l}{\oe i}{\textlengthmark}{s}{a}{m}{\textepsilon}{n}{s}{\r{g}}{a}} \dicPos{f}[1] \dicFlx{(‑u)}[5] \dicDirectTranslationCS{práce na volné noze}
\dicEntry[lausamjöll] \dicTerm{lausa··|mjöll} \dicIPA{{l}{\oe i}{\textlengthmark}{s}{a}{m}{j}{\oe}{\textsubring{d}}{\textsubring{l}}} \dicPos{f}[7] \dicFlx{(‑mjallar)}[19] \dicDirectTranslationCS{čerstvě napadlý sníh, prašan}
\dicEntry[lausasala] \dicTerm{lausa··|sala} \dicIPA{{l}{\oe i}{\textlengthmark}{s}{a}{s}{a}{l}{a}} \dicPos{f}[1] \dicFlx{(‑sölu, ‑sölur)}[20] \dicDirectTranslationCS{volný prodej}
\dicEntry[lausavísa] \dicTerm{lausa··vís|a} \dicIPA{{l}{\oe i}{\textlengthmark}{s}{a}{v}{i}{s}{a}} \dicPos{f}[1] \dicFlx{(‑u, ‑ur)}[13] \dicFieldCat{lit.} \dicDirectTranslationCS{epigram} \dicIndirectTranslationCS{(často od neznámého autora, velmi populární na Islandu)}
\dicEntry[lausblaðabók] \dicTerm{laus·blaða··|bók} \dicIPA{{l}{\oe i}{s}{\textsubring{b}}{l}{a}{ð}{a}{\textsubring{b}}{ou}{\r{g}}} \dicPos{f}[8] \dicFlx{(‑bókar, ‑bækur)}[5] \dicDirectTranslationCS{zápisník\,/\addthin sešit s~volnými listy}

\dicEntry[lausblaðamappa] \textls[15]{\dicTerm{laus·blaða··|mappa} \dicIPA{{l}{\oe i}{s}{\textsubring{b}}{l}{a}{ð}{a}{m}{a}{h}{\textsubring{b}}{a}} \dicPos{f}[1]} \dicFlx{(‑möppu, ‑möpp\-ur)}[8] \dicDirectTranslationCS{pořadač, šanon}

\dicEntry[lauslátur] \dicTerm{laus··látur} \dicIPA{{l}{\oe i}{s}{l}{au}{\textsubring{d}}{\textscy}{\textsubring{r}}} \dicPos{adj}[1]\dicFlx{}[-1] \dicDirectTranslationCS{nezávazný (flirt ap.), promiskuitní}
\dicEntry[lauslegur] \dicTerm{laus··legur} \dicIPA{{l}{\oe i}{s}{l}{\textepsilon}{\textbabygamma}{\textscy}{\textsubring{r}}} \dicPos{adj}[1]\dicFlx{}[-8] \textbf{1.} \dicSynonym{yfirborðslegur} \dicDirectTranslationCS{zběžný, letmý, hrubý (odhad ap.)} \dicExampleIS{lauslegur uppdráttur} \dicExampleCS{zběžný nákres}  \textbf{2.} \dicSynonym{óstöðugur} \dicDirectTranslationCS{volný, nepřipevněný}
\dicEntry[lauslæti] \dicTerm{laus··læti} \dicIPA{{l}{\oe i}{s}{l}{a}{i}{\textsubring{d}}{\textsci}} \dicPos{n}[2] \dicFlx{(‑s)}[20] \dicDirectTranslationCS{nezávaznost, promiskuita}
\dicEntry[lausmáll] \dicTerm{laus··máll} \dicIPA{{l}{\oe i}{s}{m}{au}{\textsubring{d}}{\textsubring{l}}} \dicPos{adj}[8]\dicFlx{}[-1] \textbf{1.} \dicSynonym{málgefinn} \dicDirectTranslationCS{upovídaný, užvaněný}  \textbf{2.} \dicSynonym*{óorðheldinn} \dicDirectTranslationCS{indiskrétní, netaktní}
\dicEntry[lausn] \dicTerm{lausn} \dicsymFrequent\  \dicIPA{{l}{\oe i}{s}{\textsubring{d}}{\textsubring{n}}} \dicPos{f}[7] \dicFlx{(‑ar, ‑ir)}[1] \textbf{1.} \dicSynonym{frelsun} \dicDirectTranslationCS{osvobození, vysvobození} \dicExampleIS{lausn úr þrældómi} \dicExampleCS{vysvobození z~otroctví}  \textbf{2.} \dicDirectTranslationCS{uvolnění, propuštění (z~práce ap.)} \dicExampleIS{lausn frá embætti} \dicExampleCS{uvolnění z~úřadu};  \dicPhraseIS{biðjast lausnar} \dicFlx{refl} \dicDirectTranslationCS{žádat o~uvolnění (z~pozice ap.)}  \textbf{3.} \dicSynonym{úrslit} \dicDirectTranslationCS{(vy)řešení, rozřešení, rozuzlení} \dicExampleIS{lausn á vandanum} \dicExampleCS{vyřešení problému}  \textbf{4.} \dicSynonym{endurlausn} \dicDirectTranslationCS{vykoupení, spasení}  \textbf{5.} \dicFieldCat{chem.} \dicDirectTranslationCS{roztok};  \dicPhraseIS{mettuð lausn} \dicFieldCat{chem.} \dicDirectTranslationCS{nasycený roztok}
\dicEntry[lausnarbeiðni] \dicTerm{lausnar··beiðn|i} \dicIPA{{l}{\oe i}{s}{\textsubring{d}}{n}{a}{r}{\textsubring{b}}{ei}{ð}{n}{\textsci}} \dicPos{f}[3] \dicFlx{(‑i, ‑ir)}[1] \dicDirectTranslationCS{výpověď, žádost o~uvolnění (z~práce ap.)};  \dicPhraseIS{senda inn lausnarbeiðni} \dicDirectTranslationCS{podat výpověď (z~práce ap.)}
\dicEntry[lausnargjald] \dicTerm{lausnar··|gjald} \dicIPA{{l}{\oe i}{s}{\textsubring{d}}{n}{a}{r}{\r{\textObardotlessj}}{a}{l}{\textsubring{d}}} \dicPos{n}[2] \dicFlx{(‑gjalds, ‑gjöld)}[8] \dicDirectTranslationCS{výkupné}
\dicEntry[lausnarhnappur] \dicTerm{lausnar··hnapp|ur} \dicIPA{{l}{\oe i}{s}{\textsubring{d}}{n}{a}{\textsubring{r}}{\textsubring{n}}{a}{h}{\textsubring{b}}{\textscy}{\textsubring{r}}} \dicPos{m}[6] \dicFlx{(‑s, ‑ar)}[25] \dicFieldCat{poč.} \dicDirectTranslationCS{klávesa escape}
\dicEntry[lausnari] \dicTerm{lausn··ar|i} \dicIPA{{l}{\oe i}{s}{\textsubring{d}}{n}{a}{r}{\textsci}} \dicPos{m}[1] \dicFlx{(‑a, ‑ar)}[13] \dicFieldCat{náb.} \dicSynonym{frelsari} \dicDirectTranslationCS{spasitel, vykupitel}
\dicEntry[laust] \dicTerm{laust\smash{\textsuperscript{1}}} \dicIPA{{l}{\oe i}{s}{\textsubring{d}}} \dicPos{v} \dicFlx{ind pf sg 1 pers} \dicLink{ljósta}
\dicEntry[laust] \dicTerm{laus|t\smash{\textsuperscript{2}}} \dicIPA{{l}{\oe i}{s}{\textsubring{d}}} \dicPos{adv} \dicFlx{(comp ‑ar, sup ‑ast)} \textbf{1.} \dicDirectTranslationCS{lehce, slabě, mírně} \dicExampleIS{berja laust að dyrum} \dicExampleCS{zaklepat lehce na dveře}  \textbf{2.} \dicDirectTranslationCS{krátce, těsně} \dicExampleIS{laust fyrir klukkan fimm} \dicExampleCS{krátce před pátou hodinou}
\dicEntry[lausung] \dicTerm{laus··ung} \dicIPA{{l}{\oe i}{\textlengthmark}{s}{u}{\ng}{\r{g}}} \dicPos{f}[4] \dicFlx{(‑ar)}[3] \textbf{1.} \dicSynonym*{staðfestuleysi} \dicDirectTranslationCS{nestabilita, nestálost}  \textbf{2.} \dicSynonym{lauslæti} \dicDirectTranslationCS{promiskuita};  \dicIdiom{lausung}{ \dicPhraseIS{gjalda lausung við lygi}} \dicLangCat{přen.} \dicDirectTranslationCS{oplatit faleš lží}
\dicEntry[laut] \dicTerm{laut\smash{\textsuperscript{1}}} \dicIPA{{l}{\oe i}{\textlengthmark}{\textsubring{d}}} \dicPos{f}[7] \dicFlx{(‑ar, ‑ir)}[6] \dicSynonym{lægð} \dicDirectTranslationCS{prohlubeň, dolík}
\dicEntry[laut] \dicTerm{laut\smash{\textsuperscript{2}}} \dicIPA{{l}{\oe i}{\textlengthmark}{\textsubring{d}}} \dicPos{v} \dicFlx{ind pf sg 1 pers} \dicLink{lúta}
\dicEntry[lautinant] \dicTerm{lautinant} \dicIPA{{l}{\oe i}{\textlengthmark}{\textsubring{d}}{\textsci}{n}{a}{\textsubring{n}}{\textsubring{d}}} \dicPos{m}[4] \dicFlx{(‑s, ‑ar)}[11] \dicFieldCat{voj.} \dicDirectTranslationCS{poručík, poručice}
\dicEntry[lax] \dicTerm{lax}\dicTerm{, laxfiskur} \dicIPA{{l}\-{a}\-{x}\-{s}\-} \dicPos{m}[4] \dicFlx{(‑, ‑ar)}[29] \dicFieldCat{zool.} \dicDirectTranslationCS{losos, losos obecný} \textit{(l.~{\textLA{Salmo salar}})}
\dicEntry[laxá] \dicTerm{lax··á} \dicIPA{{l}{a}{x}{s}{au}} \dicPos{f}[4] \dicFlx{(‑r, ‑r)}[18] \dicDirectTranslationCS{lososí řeka}
\dicEntry[laxerolía] \dicTerm{laxer··olí|a} \dicIPA{{l}{a}{x}{s}{\textepsilon}{r}{\textopeno}{l}{i}{j}{a}} \dicPos{f}[1] \dicFlx{(‑u)}[5] \dicDirectTranslationCS{ricinový olej}
\dicEntry[laxfiskur] \dicTerm{lax··fisk|ur} \dicIPA{{l}{a}{x}{s}{f}{\textsci}{s}{\r{g}}{\textscy}{\textsubring{r}}} \dicPos{m}[6] \dicFlx{(‑s, ‑ar)}[10] \dicLink{lax}
\dicEntry[laxveiði] \dicTerm{lax··veið|i} \dicIPA{{l}{a}{x}{s}{v}{ei}{ð}{\textsci}} \dicPos{f}[2] \dicFlx{(‑i\,/\addthin ‑ar, ‑ar)}[3] \textbf{1.} \dicDirectTranslationCS{lososí úlovek}  \textbf{2.} \dicDirectTranslationCS{chytání\,/\addthin lov lososů}
\dicEntry[lá] \dicTerm{lá\smash{\textsuperscript{1}}} \dicIPA{{l}{au}{\textlengthmark}} \dicPos{v}[2] \dicFlx{(‑ði, ‑ð)}[117] \dicFlx{dat} \dicSynonym{ásaka} \dicDirectTranslationCS{vyčíst, vyčítat, dávat za vinu} \dicIndirectTranslationCS{(používané v~otázce nebo záporu)} \dicExampleIS{Ég lái þér það ekki.} \dicExampleCS{Já ti to nevyčítám.};  \dicIdiom{lást}{ \dicPhraseIS{e‑m láist að (gera e‑ð)}} \dicFlx{refl impers} \dicDirectTranslationCS{(kdo) zapomíná (udělat (co)), (komu) vypadá z~paměti (udělat (co))} \dicExampleIS{Mér láðist að senda bréfið.} \dicExampleCS{Vypadlo mi (z~paměti), že mám poslat dopis.}
\dicEntry[lá] \dicTerm{lá\smash{\textsuperscript{2}}} \dicIPA{{l}{au}{\textlengthmark}} \dicPos{v} \dicFlx{ind pf sg 1 pers} \dicLink{liggja}
\dicEntry[láð] \dicTerm{láð} \dicIPA{{l}{au}{\textlengthmark}{\texttheta}} \dicPos{n}[2] \dicFlx{(‑s, ‑)}[5] \dicSynonym{land} \dicDirectTranslationCS{souš, země} \dicExampleIS{um láð og lög} \dicExampleCS{po souši a~po vodě}
\dicEntry[lág] \dicTerm{lág} \dicIPA{{l}{au}{\textlengthmark}} \dicPos{f}[4] \dicFlx{(‑ar, ‑ar)}[1] \dicSynonym{dæld} \dicDirectTranslationCS{dolina, kotlina};  \dicPhraseIS{láta e‑ð liggja í láginni} \dicLangCat{přen.} \dicDirectTranslationCS{nechat (co) ležet u~ledu}
\dicEntry[lágfiðla] \dicTerm{lág··fiðl|a} \dicIPA{{l}{au}{\textlengthmark}{f}{\textsci}{ð}{l}{a}} \dicPos{f}[1] \dicFlx{(‑u, ‑ur)}[7] \dicFieldCat{hud.} \dicDirectTranslationCS{viola}
\dicEntry[lágkúra] \dicTerm{lág··kúr|a} \dicIPA{{l}{au}{\textlengthmark}{k\smash{\textsuperscript{h}}}{u}{r}{a}} \dicPos{f}[1] \dicFlx{(‑u)}[5] \dicSynonym*{svipleysa} \dicDirectTranslationCS{banalita, přízemnost}
\dicEntry[lágkúrulegur] \textls[15]{\dicTerm{lág·kúru··legur} \dicIPA{{l}{au}{\textlengthmark}{k\smash{\textsuperscript{h}}}{u}{r}{\textscy}{l}{\textepsilon}{\textbabygamma}{\textscy}{\textsubring{r}}} \dicPos{adj}[1]\dicFlx{}[-8] \textbf{1.} \dicSynonym*{flatneskjulegur} \dicDirectTranslationCS{banální, všední, nezajímavý}} \dicExampleIS{lágkúrulegt hús} \dicExampleCS{všední dům}  \textbf{2.} \dicSynonym{lítilmótlegur} \dicDirectTranslationCS{přízemní, nízký, mrzký} \dicExampleIS{lágkúrulegur maður} \dicExampleCS{přízemní člověk}
\dicEntry[lágkveða] \dicTerm{lág··kveð|a} \dicIPA{{l}{au}{\textlengthmark}{k\smash{\textsuperscript{h}}}{v}{\textepsilon}{ð}{a}} \dicPos{f}[1] \dicFlx{(‑u, ‑ur)}[19] \dicFieldCat{lit.} \dicDirectTranslationCS{sudá stopa se slabším přízvukem}
\dicEntry[láglaunafólk] \dicTerm{lág·launa··fólk} \dicIPA{{l}{au}{\textlengthmark}{l}{\oe i}{n}{a}{f}{ou}{\textsubring{l}}{\r{g}}} \dicPos{n}[2] \dicFlx{(‑s)}[2] \dicDirectTranslationCS{osoby s~nízkým příjmem}
\dicEntry[láglendi] \dicTerm{lág··lendi} \dicIPA{{l}{au}{\textlengthmark}{l}{\textepsilon}{n}{\textsubring{d}}{\textsci}} \dicPos{n}[2] \dicFlx{(‑s)}[20] \dicDirectTranslationCS{nížina} \dicAntonym{hálendi}
\dicEntry[lágmark] \dicTerm{lág··|mark} \dicIPA{{l}{au}{\textlengthmark}{m}{a}{\textsubring{r}}{\r{g}}} \dicPos{n}[2] \dicFlx{(‑marks, ‑mörk)}[8] \dicDirectTranslationCS{minimum} \dicExampleIS{Kuldinn er í algjöru lágmarki.} \dicExampleCS{Zima je na úplném minimu.} \dicAntonym{hámark}
\dicEntry[lágmarksframfærsla] \dicTerm{lág·marks··fram·færsl|a} \dicIPA{{l}\-{au}\-{\textlengthmark}\-{m}\-{a}\-{\textsubring{r}}\-{x}\-{s}\-{f}\-{r}\-{a}\-{m}\-{f}\-{a}\-{i}\-{\textsubring{r}}\-{s}\-{\textsubring{d}}\-{l}\-{a}\-} \dicPos{f}[1] \dicFlx{(‑u)}[5] \dicFieldCat{ekon.} \dicDirectTranslationCS{životní minimum}
\dicEntry[lágmarksgjald] \dicTerm{lág·marks··|gjald} \dicIPA{{l}{au}{\textlengthmark}{m}{a}{\textsubring{r}}{x}{s}{\r{\textObardotlessj}}{a}{l}{\textsubring{d}}} \dicPos{n}[2] \dicFlx{(‑gjalds, ‑gjöld)}[8] \dicDirectTranslationCS{minimální sazba}
\begin{xtolerant}{}{1pt}
\dicEntry[lágmarkslaun] \dicTerm{lág·marks··laun}\addthinS\dicIPA{{l}{au}{\textlengthmark}{m}{a}{\textsubring{r}}{x}{s}{l}{\oe i}{\textsubring{n}}}\addthinS\dicPos{n}[2]\addthinS\dicFlx{pl}[1]\addthinS\dicDirectTranslationCS{minimální mzda\hspace{-0.06667em}\,/\addthin plat}
\end{xtolerant}
\dicEntry[lágmarkstímasókn] \dicTerm{lág·marks··tíma·sókn} \dicIPA{{l}{au}{\textlengthmark}{m}{a}{\textsubring{r}}{x}{s}{t\smash{\textsuperscript{h}}}{i}{m}{a}{s}{ou}{h}{\r{g}}{\textsubring{n}}} \dicPos{f}[7] \dicFlx{(‑ar)}[3] \dicDirectTranslationCS{minimální počet hodin školní docházky} \dicIndirectTranslationCS{(pro ukončení kurzu ap.)}
\dicEntry[lágmarksverð] \dicTerm{lág·marks··verð} \dicIPA{{l}{au}{\textlengthmark}{m}{a}{\textsubring{r}}{x}{s}{v}{\textepsilon}{r}{\texttheta}} \dicPos{n}[2] \dicFlx{(‑s)}[2] \dicDirectTranslationCS{minimální cena}
\dicEntry[lágmynd] \dicTerm{lág··mynd} \dicIPA{{l}{au}{\textlengthmark}{m}{\textsci}{n}{\textsubring{d}}} \dicPos{f}[7] \dicFlx{(‑ar, ‑ir)}[1] \dicDirectTranslationCS{reliéf}
\dicEntry[lágnætti] \dicTerm{lág··nætti} \dicIPA{{l}{au}{\textlengthmark}{n}{a}{i}{h}{\textsubring{d}}{\textsci}} \dicPos{n}[2] \dicFlx{(‑s, ‑)}[14] \dicDirectTranslationCS{půlnoc} \dicIndirectTranslationCS{(zvláště v~létě)}
\dicEntry[lágspenna] \dicTerm{lág··spenn|a} \dicIPA{{l}{au}{\textlengthmark}{s}{\textsubring{b}}{\textepsilon}{n}{a}} \dicPos{f}[1] \dicFlx{(‑u)}[5] \dicFieldCat{elek.} \dicDirectTranslationCS{nízké napětí} \dicAntonym{háspenna}
\dicEntry[lágstétt] \dicTerm{lág··stétt} \dicIPA{{l}{au}{\textlengthmark}{s}{\textsubring{d}}{j}{\textepsilon}{h}{\textsubring{d}}} \dicPos{f}[7] \dicFlx{(‑ar, ‑ir)}[1] \dicDirectTranslationCS{nižší (společenská) třída, proletariát} \dicAntonym{hástétt}
\dicEntry[lágt] \dicTerm{lágt} \dicsymFrequent\  \dicIPA{{l}{au}{x}{\textsubring{d}}} \dicPos{adv} \dicFlx{(comp lægra, sup lægst)} \textbf{1.} \dicSynonym*{lítill hljóðstyrkur} \dicDirectTranslationCS{tiše, potichu, tlumeně} \dicExampleIS{tala lágt} \dicExampleCS{mluvit potichu}  \textbf{2.} \dicSynonym{niðri} \dicDirectTranslationCS{nízko, dole}
\dicEntry[lágum] \dicTerm{lágum} \dicIPA{{l}{au}{\textlengthmark}{\textbabygamma}{\textscy}{\textsubring{m}}} \dicPos{v} \dicFlx{ind pf pl 1 pers} \dicLink{liggja}
\dicEntry[lágur] \dicTerm{lágur} \dicsymFrequent\  \dicIPA{{l}{au}{\textlengthmark}{\textscy}{\textsubring{r}}} \dicPos{adj}[10] \dicFlx{(comp lægri, sup lægstur)}[16] \textbf{1.} \dicSynonym{stuttur} \dicDirectTranslationCS{nízký, malý} \dicIndirectTranslationCS{(o~výšce)};  \dicPhraseIS{vera lágur vexti} \dicDirectTranslationCS{být malého vzrůstu}  \textbf{2.} \dicDirectTranslationCS{nízký} \dicIndirectTranslationCS{(o~čísle, ceně ap.)}  \textbf{3.} \dicDirectTranslationCS{tichý, tlumený} \dicExampleIS{lágur söngur} \dicExampleCS{tichý zpěv};  \dicPhraseIS{í lágum hljóðum} \dicFlx{adv} \dicDirectTranslationCS{tlumeně, tiše}  \textbf{4.} \dicDirectTranslationCS{nízký (tón ap.)} \dicExampleIS{lágur tónn} \dicExampleCS{nízký tón};  \dicIdiom{lágur}{ \dicPhraseIS{bera\,/\addthin bíða lægri hlut}} \dicLangCat{přen.} \dicDirectTranslationCS{tahat za kratší konec provazu}
\dicEntry[lágvaxinn] \dicTerm{lág··vaxinn} \dicIPA{{l}{au}{\textlengthmark}{v}{a}{x}{s}{\textsci}{\textsubring{n}}} \dicPos{adj}[6]\dicFlx{}[-3] \dicDirectTranslationCS{malý (vzrůstem), (jsoucí) malého vzrůstu} \dicAntonym{hávaxinn}
\dicEntry[lágþrýstingur] \dicTerm{lág··þrýst·ing|ur} \dicIPA{{l}{au}{\textlengthmark}{\texttheta}{r}{i}{s}{\textsubring{d}}{i}{\ng}{\r{g}}{\textscy}{\textsubring{r}}} \dicPos{m}[6] \dicFlx{(‑s)}[9] \textbf{1.} \dicFieldCat{med.} \dicDirectTranslationCS{nízký tlak}  \textbf{2.} \dicFieldCat{meteo.} \dicDirectTranslationCS{nízký tlak}
\dicEntry[lágþrýstisvæði] \dicTerm{lág·þrýsti··svæði} \dicIPA{{l}{au}{\textlengthmark}{\texttheta}{r}{i}{s}{\textsubring{d}}{\textsci}{s}{v}{a}{i}{ð}{\textsci}} \dicPos{n}[2] \dicFlx{(‑s, ‑)}[14] \dicFieldCat{meteo.} \dicDirectTranslationCS{oblast nízkého tlaku} \dicAntonym{háþrýstisvæði}
\dicEntry[lágþýska] \dicTerm{lág··þýsk|a} \dicIPA{{l}{au}{\textlengthmark}{\texttheta}{i}{s}{\r{g}}{a}} \dicPos{f}[1] \dicFlx{(‑u)}[5] \dicDirectTranslationCS{dolní němčina} \dicIndirectTranslationCS{(dialekt němčiny používaný v~severních částech Německa)} \dicAntonym{háþýska}
\dicEntry[lákum] \dicTerm{lákum} \dicIPA{{l}{au}{\textlengthmark}{\r{g}}{\textscy}{\textsubring{m}}} \dicPos{v} \dicFlx{ind pf pl 1 pers} \dicLink{leka}
\dicEntry[lán] \dicTerm{lán} \dicsymFrequent\  \dicIPA{{l}{au}{\textlengthmark}{\textsubring{n}}} \dicPos{n}[2] \dicFlx{(‑s, ‑)}[5] \textbf{1.} \dicSynonym{lánsfé} \dicDirectTranslationCS{(vý)půjčka, úvěr} \dicExampleIS{borga lán} \dicExampleCS{splatit půjčku};  \dicPhraseIS{taka lán} \dicDirectTranslationCS{vzít si půjčku\,/\addthin úvěr}  \textbf{2.} \dicSynonym*{það að lána} \dicDirectTranslationCS{(vy)půjčení};  \dicPhraseIS{fá e‑ð að láni} \dicDirectTranslationCS{(vy)půjčit si (co)};  \dicPhraseIS{hafa e‑ð að láni, vera með e‑ð í láni} \dicDirectTranslationCS{(vy)půjčit si (co), mít (co) půjčené}  \textbf{3.} \dicSynonym{gæfa} \dicDirectTranslationCS{štěstí, štěstěna};  \dicPhraseIS{lánið leikur við e‑n} \dicLangCat{přen.} \dicDirectTranslationCS{(kdo) má štěstí, na (koho) se usmívá štěstí};  \dicPhraseIS{það er lán í óláni að} \dicLangCat{přen.} \dicDirectTranslationCS{je to štěstí v~neštěstí, že}
\dicEntry[lána] \dicTerm{lán|a} \dicsymFrequent\  \dicIPA{{l}{au}{\textlengthmark}{n}{a}} \dicPos{v}[1] \dicFlx{(‑aði)}[1] \dicFlx{dat + acc} \dicDirectTranslationCS{(vy)půjčit, zapůjčit};  \dicPhraseIS{lána e‑m e‑ð} \dicDirectTranslationCS{půjčit (komu co)} \dicExampleIS{lána e‑m bókina} \dicExampleCS{půjčit (komu) knížku};  \dicPhraseIS{fá e‑ð lánað} \dicDirectTranslationCS{(vy)půjčit si (co)};  \dicIdiom{lána}[út]{ \dicPhraseIS{lána út e‑ð}} \dicDirectTranslationCS{(za)půjčovat (co)}
\dicEntry[lánardrottinn] \dicTerm{lánar··drott|inn} \dicIPA{{l}{au}{\textlengthmark}{n}{a}{r}{\textsubring{d}}{r}{\textopeno}{h}{\textsubring{d}}{\textsci}{\textsubring{n}}} \dicPos{m}[6] \dicFlx{(‑ins, ‑nar)}[39] \dicDirectTranslationCS{věřitel(ka)}
\dicEntry[lánasjóður] \dicTerm{lána··sjóð|ur} \dicIPA{{l}{au}{\textlengthmark}{n}{a}{s}{j}{ou}{ð}{\textscy}{\textsubring{r}}} \dicPos{m}[9] \dicFlx{(‑s, ‑ir)}[6] \dicFieldCat{ekon.} \dicDirectTranslationCS{úvěrový fond}
\dicEntry[lánastofnun] \dicTerm{lána··stofn|un} \dicIPA{{l}{au}{\textlengthmark}{n}{a}{s}{\textsubring{d}}{\textopeno}{\textsubring{b}}{n}{\textscy}{\textsubring{n}}} \dicPos{f}[7] \dicFlx{(‑unar, ‑anir)}[8] \dicFieldCat{ekon.} \dicDirectTranslationCS{úvěrová společnost}
\dicEntry[lánaveitandi] \dicTerm{lána··veit·|andi} \dicIPA{{l}{au}{\textlengthmark}{n}{a}{v}{ei}{\textsubring{d}}{a}{n}{\textsubring{d}}{\textsci}} \dicPos{m}[2] \dicFlx{(‑anda, ‑endur)}[1] \dicDirectTranslationCS{věřitel(ka)}
\dicEntry[lánlaus] \dicTerm{lán··laus} \dicIPA{{l}{au}{n}{l}{\oe i}{s}} \dicPos{adj}[5]\dicFlx{}[-1] \dicDirectTranslationCS{smolařský, nemající štěstí}
\dicEntry[lánsamur] \dicTerm{lán··|samur} \dicIPA{{l}{au}{n}{s}{a}{m}{\textscy}{\textsubring{r}}} \dicPos{adj}[1] \dicFlx{(f ‑söm)}[2] \dicSynonym{gæfusamur} \dicDirectTranslationCS{šťastný, mající štěstí}
\dicEntry[lánsfé] \dicTerm{láns··|fé} \dicIPA{{l}{au}{n}{s}{f}{j}{\textepsilon}} \dicPos{n}[3] \dicFlx{(‑fjár)}[5] \dicDirectTranslationCS{půjčka, půjčené peníze}
\dicEntry[lánsfjármögnun] \dicTerm{láns·fjár··mögn|un} \dicIPA{{l}{au}{n}{s}{f}{j}{au}{r}{m}{\oe}{\r{g}}{n}{\textscy}{\textsubring{n}}} \dicPos{f}[7] \dicFlx{(‑unar)}[12] \dicFieldCat{ekon.} \dicDirectTranslationCS{financování dluhu}
\begin{xtolerant}{}{1pt}
\dicEntry[lánskjaravísitala] \dicTerm{láns·kjara··vísi·|tala}\addthinS\dicIPA{{l}\-{au}\-{n}\-{s}\-{c\smash{\textsuperscript{h}}}\-{a}\-{r}\-{a}\-{v}\-{i}\-{s}\-{\textsci}\-{t\smash{\textsuperscript{h}}}\-{a}\-{l}\-{a}\-}\addthinS\dicPos{f}[1]\addthinS\dicFlx{(‑tölu, ‑tölur)}[14] \dicFieldCat{ekon.} \dicDirectTranslationCS{index úvěrových podmínek}
\end{xtolerant}
\dicEntry[lánskjör] \dicTerm{láns··kjör} \dicIPA{{l}{au}{n}{s}{c\smash{\textsuperscript{h}}}{\oe}{\textsubring{r}}} \dicPos{n}[2] \dicFlx{pl}[9] \dicFieldCat{ekon.} \dicDirectTranslationCS{úvěrové podmínky}
\dicEntry[lánstími] \dicTerm{láns··tím|i} \dicIPA{{l}{au}{n}{s}{t\smash{\textsuperscript{h}}}{i}{m}{\textsci}} \dicPos{m}[1] \dicFlx{(‑a)}[3] \textbf{1.} \dicFieldCat{ekon.} \dicDirectTranslationCS{doba splatnosti úvěru}  \textbf{2.} \dicDirectTranslationCS{výpůjční doba (v~knihovně ap.)}
\dicEntry[lánstraust] \dicTerm{láns··traust} \dicIPA{{l}{au}{n}{s}{t\smash{\textsuperscript{h}}}{r}{\oe i}{s}{\textsubring{d}}} \dicPos{n}[2] \dicFlx{(‑s)}[2] \dicFieldCat{ekon.} \dicDirectTranslationCS{úvěruschopnost}
\dicEntry[lántakandi] \dicTerm{lán··tak·|andi} \dicIPA{{l}{au}{n}{t\smash{\textsuperscript{h}}}{a}{\r{g}}{a}{n}{\textsubring{d}}{\textsci}} \dicPos{m}[2] \dicFlx{(‑anda, ‑endur)}[1] \dicSynonym{lánþegi} \dicDirectTranslationCS{vypůjčovatel(ka), příjemce\,/\addthin příjemkyně půjčky}
\dicEntry[lántaki] \dicTerm{lán··tak|i} \dicIPA{{l}{au}{n}{t\smash{\textsuperscript{h}}}{a}{\r{\textObardotlessj}}{\textsci}} \dicPos{m}[1] \dicFlx{(‑a, ‑ar)}[8] \dicSynonym{lánþegi} \dicDirectTranslationCS{vypůjčovatel(ka), příjemce\,/\addthin příjemkyně půjčky}
\dicEntry[lánveiting] \dicTerm{lán··veit·ing} \dicIPA{{l}{au}{n}{v}{ei}{\textsubring{d}}{i}{\ng}{\r{g}}} \dicPos{f}[4] \dicFlx{(‑ar, ‑ar)}[5] \dicFieldCat{ekon.} \dicDirectTranslationCS{poskytnutí\,/\addthin poskytování půjčky\,/\addthin úvěru}
\dicEntry[lánþegi] \dicTerm{lán··þeg|i} \dicIPA{{l}{au}{n}{\texttheta}{ei}{\textsci}} \dicPos{m}[1] \dicFlx{(‑a, ‑ar)}[1] \dicDirectTranslationCS{vypůjčovatel(ka), příjemce\,/\addthin příjemkyně půjčky}
\dicEntry[láréttur] \dicTerm{lá··réttur} \dicIPA{{l}{au}{\textlengthmark}{r}{j}{\textepsilon}{h}{\textsubring{d}}{\textscy}{\textsubring{r}}} \dicPos{adj}[1]\dicFlx{}[-10] \dicDirectTranslationCS{vodorovný, horizontální}
\dicEntry[lárpera] \dicTerm{lár··per|a} \dicIPA{{l}{au}{r}{p\smash{\textsuperscript{h}}}{\textepsilon}{r}{a}} \dicPos{f}[1] \dicFlx{(‑u, ‑ur)}[7] \dicDirectTranslationCS{avokádo} \dicIndirectTranslationCS{(plod)}
\dicEntry[lárperutré] \dicTerm{lár·peru··tré} \dicIPA{{l}{au}{r}{p\smash{\textsuperscript{h}}}{\textepsilon}{r}{\textscy}{t\smash{\textsuperscript{h}}}{r}{j}{\textepsilon}} \dicPos{n}[2] \dicFlx{(‑s, ‑)}[36] \dicFieldCat{bot.} \dicDirectTranslationCS{hruškovec přelahodný} \textit{(l.~{\textLA{Persea americana}})}  \dicsymPhoto\ 
\dicFigure{54128.jpg}{Lárperutré}{Lárperutré - Forest \& Kim [[p:2684;Starr]], Biolib, CC-BY}
\dicEntry[lárviðarlauf] \dicTerm{lár·viðar··lauf} \dicIPA{{l}{au}{r}{v}{\textsci}{ð}{a}{r}{l}{\oe i}{f}} \dicPos{n}[2] \dicFlx{(‑s, ‑)}[5] \dicDirectTranslationCS{bobkový list}
\dicEntry[lárviðarrós] \dicTerm{lár·viðar··rós} \dicIPA{{l}{au}{r}{v}{\textsci}{ð}{a}{r}{ou}{s}} \dicPos{f}[7] \dicFlx{(‑ar, ‑ir)}[1] \dicFieldCat{bot.} \dicSynonym{nería} \dicDirectTranslationCS{oleandr, oleandr obecný} \textit{(l.~{\textLA{Nerium oleander}})}  \dicsymPhoto\ 
\dicFigure{ds_image_larvidarros_0_1.jpg}{Lárviðarrós}{Lárviðarrós - Alvesgaspar, CC BY-SA 3.0}
\dicEntry[lárviður] \dicTerm{lár··við|ur} \dicIPA{{l}{au}{r}{v}{\textsci}{ð}{\textscy}{\textsubring{r}}} \dicPos{m}[10] \dicFlx{(‑ar, ‑ir)}[19] \dicFieldCat{bot.} \dicDirectTranslationCS{vavřín, vavřín vznešený} \textit{(l.~{\textLA{Laurus nobilis}})}  \dicsymPhoto\ 
\dicFigure{55905.jpg}{Lárviður}{Lárviður - Forest \& Kim [[p:2684;Starr]], Biolib, CC-BY}
\dicEntry[lás] \dicTerm{lás} \dicIPA{{l}{au}{\textlengthmark}{s}} \dicPos{m}[4] \dicFlx{(‑s, ‑ar)}[25] \textbf{1.} \dicDirectTranslationCS{zámek (u~dveří ap.)} \dicExampleIS{lásinn á hjólinu} \dicExampleCS{zámek na kolo};  \dicPhraseIS{skella e‑u í lás} \dicDirectTranslationCS{zabouchnout (si) (co) (dveře ap.)}  \textbf{2.} \dicDirectTranslationCS{pojistka} \dicIndirectTranslationCS{(součást střelné zbraně)} \dicExampleIS{lás á byssu} \dicExampleCS{pojistka pistole}
\dicEntry[lásasmiður] \dicTerm{lása··smið|ur} \dicIPA{{l}{au}{\textlengthmark}{s}{a}{s}{m}{\textsci}{ð}{\textscy}{\textsubring{r}}} \dicPos{m}[9] \dicFlx{(‑s, ‑ir)}[8] \dicDirectTranslationCS{zámečník, zámečnice}
\dicEntry[lásbogi] \dicTerm{lás··bog|i} \dicIPA{{l}{au}{\textlengthmark}{s}{\textsubring{b}}{\textopeno i}{j}{\textsci}} \dicPos{m}[1] \dicFlx{(‑a, ‑ar)}[1] \dicDirectTranslationCS{kuše, samostříl}
\dicEntry[lásum] \dicTerm{lásum} \dicIPA{{l}{au}{\textlengthmark}{s}{\textscy}{\textsubring{m}}} \dicPos{v} \dicFlx{ind pf pl 1 pers} \dicLink{lesa}
\dicEntry[lát] \dicTerm{lát} \dicIPA{{l}{au}{\textlengthmark}{\textsubring{d}}} \dicPos{n}[2] \dicFlx{(‑s, ‑)}[5] \textbf{1.} \dicSynonym{andlát} \dicDirectTranslationCS{skon, úmrtí} \dicExampleIS{auglýsa lát e‑rs} \dicExampleCS{oznámit (čí) úmrtí}  \textbf{2.} \dicSynonym{hlé} \dicDirectTranslationCS{přestávka, chvíle oddechu} \dicExampleIS{Það var ekkert lát á storminum alla nóttina.} \dicExampleCS{Bouře zuřila bez přestávky celou noc.}  \textbf{3.} \dicSynonym*{undanlát} \dicDirectTranslationCS{poddajnost} \dicExampleIS{Það er ekkert lát á honum.} \dicExampleCS{Není v~něm kousek poddajnosti.}
\dicEntry[láta] \dicTerm{láta} \dicsymFrequent\  \dicIPA{{l}{au}{\textlengthmark}{\textsubring{d}}{a}} \dicPos{v}[6] \dicFlx{(læt, lét, létum, léti, látið)}[84] \dicFlx{acc} \textbf{1.} \dicSynonym*{leyfa að e‑að gerist} \dicDirectTranslationCS{(po)nechat, (po)nechávat} \dicExampleIS{láta e‑n lesa} \dicExampleCS{nechat (koho) číst};  \dicPhraseIS{láta e‑n í friði} \dicDirectTranslationCS{nechat (koho) na pokoji};  \dicPhraseIS{láta e‑n vera} \dicDirectTranslationCS{nechat (koho) být}  \textbf{2.} \dicDirectTranslationCS{(po)nechat, (po)nechávat, zanechat, zanechávat} \dicExampleIS{láta bókina þar} \dicExampleCS{nechat tam knížku};  \dicPhraseIS{láta e‑ð af hendi} \dicDirectTranslationCS{dát (co) z~ruky};  \dicPhraseIS{láta e‑m e‑ð í té} \dicDirectTranslationCS{poskytnout (komu co), předat (komu co)}  \textbf{3.} \dicSynonym*{hegða sér} \dicDirectTranslationCS{chovat se, jednat} \dicExampleIS{láta eins og fífl} \dicExampleCS{chovat se jako hlupák};  \dicPhraseIS{láta sem ekkert sé} \dicDirectTranslationCS{dělat jakoby nic; dělat, jako by se nic nestalo};  \dicPhraseIS{láta öllum illum látum} \dicDirectTranslationCS{dělat problémy\,/\addthin potíže};  \dicPhraseIS{láttu ekki svona} \dicFlx{imper} \dicDirectTranslationCS{nech toho}  \textbf{4.} \dicPhraseIS{láta lífið} \dicDirectTranslationCS{přijít o~život}  \textbf{5.} \dicPhraseIS{láta í haf} \dicDirectTranslationCS{vyplout na moře};  \dicIdiom{láta}[að]{ \dicPhraseIS{láta vel að e‑n}} \dicDirectTranslationCS{být milý\,/\addthin laskavý ke (komu)};  \dicIdiom{láta}[af]{ \dicPhraseIS{láta af e‑u}} \dicDirectTranslationCS{(za)nechat (čeho) (práce ap.)};  \dicIdiom{láta}[aftur]{ \dicPhraseIS{láta aftur e‑ð}} \dicDirectTranslationCS{zavřít\,/\addthin zavírat (co)} \dicExampleIS{láta aftur hurðina\,/\addthin dyrnar} \dicExampleCS{zavřít dveře};  \dicIdiom{láta}[eftir]{ \dicPhraseIS{láta allt eftir sér}} \dicDirectTranslationCS{dopřávat si, nic si neodepřít}; { \dicPhraseIS{láta e‑ð eftir e‑m}} \dicDirectTranslationCS{nechat (co) na (kom), neodepřít (komu co)}; { \dicPhraseIS{láta eftir}} \dicDirectTranslationCS{podřídit se, podvolit se} \dicExampleIS{Hann varð að láta eftir.} \dicExampleCS{Musel se podřídit.};  \dicIdiom{láta}[nærri]{ \dicPhraseIS{það lætur nærri}} \dicFlx{impers} \dicDirectTranslationCS{moc tomu nechybí, téměř};  \dicIdiom{láta}[sér]{ \dicPhraseIS{láta sér e‑ð vel líka}} \dicDirectTranslationCS{být s~(čím) spokojený}; { \dicPhraseIS{láta sér ekki bregða}} \dicDirectTranslationCS{nenechat se překvapit\,/\addthin zaskočit}; { \dicPhraseIS{láta sér fátt um finnast}} \dicDirectTranslationCS{nedat najevo příliš velký zájem};  \dicIdiom{láta}[sig]{ \dicPhraseIS{láta sig ekki}} \dicDirectTranslationCS{nevzdat se, nevzdávat se};  \dicIdiom{láta}[undan]{ \dicPhraseIS{láta undan e‑m}} \dicSynonym{beygja\smash{\textsuperscript{2}}} \dicDirectTranslationCS{podvolit se (komu), podrobit se (komu)}; { \dicPhraseIS{láta undan}} \dicDirectTranslationCS{povolit, spadnout (střecha ap.)};  \dicIdiom{láta}[uppi]{ \dicPhraseIS{láta e‑ð uppi}} \dicDirectTranslationCS{vyjevit (co), dát (co) najevo};  \dicIdiom{láta}[út úr]{ \dicPhraseIS{láta e‑ð út úr sér}} \dicDirectTranslationCS{(dovolit si) říct (co)};  \dicIdiom{láta}[yfir]{ \dicPhraseIS{láta illa yfir e‑u}} \dicDirectTranslationCS{neschvalovat (co), odsuzovat (co)} \dicExampleIS{láta illa yfir ástandinu} \dicExampleCS{neschvalovat takovýto stav};  \dicIdiom{látast}{ \dicPhraseIS{látast}} \dicFlx{refl} {\textbf{a.}} \dicDirectTranslationCS{předstírat} \dicExampleIS{látast vera reiður} \dicExampleCS{předstírat rozčilení};  {\textbf{b.}} \dicDirectTranslationCS{skonat, zemřít} \dicExampleIS{látast úr elli} \dicExampleCS{zemřít v~důsledku vysokého věku};  \dicIdiom{láta}{ \dicPhraseIS{láta e‑ð eiga sig}} \dicDirectTranslationCS{odpustit si (co)}; { \dicPhraseIS{láta e‑ð ekki á sig fá}} \dicDirectTranslationCS{nepoddat se (čemu)}; { \dicPhraseIS{láta e‑ð í ljós}} \dicDirectTranslationCS{dát (co) najevo, vyjádřit (co)} \dicExampleIS{láta skoðun sína í ljós} \dicExampleCS{dát najevo své názory}; { \dicPhraseIS{láta fara vel um sig}} \dicDirectTranslationCS{udělat si pohodlí} \dicExampleIS{Afi lætur fara vel um sig í hægindastólnum.} \dicExampleCS{Děda si dělá pohodlí v~křesle.}; { \dicPhraseIS{láta ekki hugfallast}} \dicDirectTranslationCS{neztratit kuráž}; { \dicPhraseIS{láta ekki e‑ð hjá líða}} \dicDirectTranslationCS{nenechat si ujít (co)}; { \dicPhraseIS{láta til skarar skríða}} \dicLangCat{přen.} \dicDirectTranslationCS{jít na věc}
\dicEntry[látalæti] \dicTerm{láta··læti} \dicIPA{{l}{au}{\textlengthmark}{\textsubring{d}}{a}{l}{a}{i}{\textsubring{d}}{\textsci}} \dicPos{n}[2] \dicFlx{pl}[19] \dicSynonym{uppgerð} \dicDirectTranslationCS{předstírání, přetvářka, pokrytectví}
\dicEntry[látbragð] \dicTerm{lát··bragð} \dicIPA{{l}{au}{\textlengthmark}{\textsubring{d}}{\textsubring{b}}{r}{a}{\textbabygamma}{\texttheta}} \dicPos{n}[2] \dicFlx{(‑s)}[2] \textbf{1.} \dicDirectTranslationCS{posunek, gesto}  \textbf{2.} \dicDirectTranslationCS{pantomima}
\dicEntry[látbragðsleikur] \dicTerm{lát·bragðs··leik|ur} \dicIPA{{l}{au}{\textlengthmark}{\textsubring{d}}{\textsubring{b}}{r}{a}{\textbabygamma}{ð}{s}{l}{ei}{\r{g}}{\textscy}{\textsubring{r}}} \dicPos{m}[6] \dicFlx{(‑s)}[17] \dicDirectTranslationCS{němohra, pantomima}
\dicEntry[látið] \dicTerm{látið} \dicIPA{{l}{au}{\textlengthmark}{\textsubring{d}}{\textsci}{\texttheta}} \dicPos{v} \dicFlx{supin} \dicLink{láta}
\dicEntry[látinn] \dicTerm{látinn} \dicsymFrequent\  \dicIPA{{l}{au}{\textlengthmark}{\textsubring{d}}{\textsci}{\textsubring{n}}} \dicPos{adj}[6]\dicFlx{}[-6] \dicSynonym{dáinn} \dicDirectTranslationCS{zesnulý, zemřelý} \dicExampleIS{níu látnir og fjölmargir slasaðir} \dicExampleCS{devět mrtvých a~mnoho zraněných};  \dicPhraseIS{vera vant við látinn} \dicDirectTranslationCS{mít napilno, být zaneprázdněný};  \dicPhraseIS{vera vel látinn} \dicDirectTranslationCS{být populární\,/\addthin uznávaný}
\dicEntry[látlaus] \dicTerm{lát··laus} \dicIPA{{l}{au}{\textlengthmark}{\textsubring{d}}{l}{\oe i}{s}} \dicPos{adj}[5]\dicFlx{}[-1] \textbf{1.} \dicSynonym{tilgerðarlaus} \dicDirectTranslationCS{prostý, skromný, nenáročný}  \textbf{2.} \dicSynonym{linnulaus} \dicDirectTranslationCS{vytrvalý, soustavný, nepřetržitý} \dicExampleIS{látlaus rigning} \dicExampleCS{soustavný déšť}
\dicEntry[látleysi] \dicTerm{lát··leysi} \dicIPA{{l}{au}{\textlengthmark}{\textsubring{d}}{l}{ei}{s}{\textsci}} \dicPos{n}[2] \dicFlx{(‑s)}[20] \dicSynonym*{tilgerðarleysi} \dicDirectTranslationCS{prostota, přirozenost, jednoduchost}
\dicEntry[látún] \dicTerm{látún} \dicIPA{{l}{au}{\textlengthmark}{t\smash{\textsuperscript{h}}}{u}{\textsubring{n}}} \dicPos{n}[2] \dicFlx{(‑s)}[2] \dicDirectTranslationCS{mosaz}
\dicEntry[látæði] \dicTerm{lát··æði} \dicIPA{{l}{au}{\textlengthmark}{\textsubring{d}}{a}{i}{ð}{\textsci}} \dicPos{n}[2] \dicFlx{(‑s)}[20] \textbf{1.} \dicSynonym{hegðun} \dicDirectTranslationCS{chování, jednání}  \textbf{2.} \dicSynonym{látbragð} \dicDirectTranslationCS{posunek, gesto}
\dicEntry[lávarður] \dicTerm{lávarð|ur} \dicIPA{{l}{au}{\textlengthmark}{v}{a}{r}{ð}{\textscy}{\textsubring{r}}} \dicPos{m}[6] \dicFlx{(‑s\,/\addthin ‑ar, ‑ar)}[64] \dicDirectTranslationCS{lord}
\dicEntry[Lbs.] \dicTerm{Lbs.} \dicPos{zkr} \dicPhraseIS{Landsbókasafn} \dicDirectTranslationCS{Národní knihovna}
\dicEntry[ld.] \dicTerm{ld.} \dicPos{zkr} \dicPhraseIS{laugardagur} \dicDirectTranslationCS{sobota}
\dicEntry[leðja] \dicTerm{leðj|a} \dicIPA{{l}{\textepsilon}{ð}{j}{a}} \dicPos{f}[1] \dicFlx{(‑u)}[5] \dicSynonym{for\smash{\textsuperscript{1}}} \dicDirectTranslationCS{bláto, bahno} \dicExampleIS{sökkva ofan í leðjuna} \dicExampleCS{zapadnout do bláta}
\dicEntry[leður] \dicTerm{leður} \dicsymFrequent\  \dicIPA{{l}{\textepsilon}{\textlengthmark}{ð}{\textscy}{\textsubring{r}}} \dicPos{n}[2] \dicFlx{(‑s, ‑)}[25] \dicDirectTranslationCS{kůže} \dicIndirectTranslationCS{(materiál)} \dicExampleIS{vörur úr leðri} \dicExampleCS{výrobky z~kůže}
\dicEntry[leðurblaka] \dicTerm{leður··|blaka} \dicIPA{{l}{\textepsilon}{\textlengthmark}{ð}{\textscy}{r}{\textsubring{b}}{l}{a}{\r{g}}{a}} \dicPos{f}[1] \dicFlx{(‑blöku, ‑blökur)}[20] \textbf{1.} \dicPhraseIS{leðurblökur} \dicFlx{pl} \dicFieldCat{zool.} \dicDirectTranslationCS{letouni} \textit{(l.~{\textLA{Chiroptera}})}  \dicsymPhoto\   \textbf{2.} \dicFieldCat{zool.} \dicDirectTranslationCS{netopýr} \textit{(l.~{\textLA{Microchiroptera}})}
\dicFigure{62444.jpg}{Leðurblaka}{Leðurblaka - Krejčík Stanislav, Biolib, Copyright/CC-BY}
\dicEntry[leðurjakki] \dicTerm{leður··jakk|i} \dicIPA{{l}{\textepsilon}{\textlengthmark}{ð}{\textscy}{r}{j}{a}{h}{\r{\textObardotlessj}}{\textsci}} \dicPos{m}[1] \dicFlx{(‑a, ‑ar)}[8] \dicDirectTranslationCS{kožená bunda}
\dicEntry[leðurlíki] \dicTerm{leður··líki} \dicIPA{{l}{\textepsilon}{\textlengthmark}{ð}{\textscy}{r}{l}{i}{\r{\textObardotlessj}}{\textsci}} \dicPos{n}[2] \dicFlx{(‑s, ‑)}[16] \dicDirectTranslationCS{koženka}
\dicEntry[leg] \dicTerm{leg} \dicIPA{{l}{\textepsilon}{\textlengthmark}{x}} \dicPos{n}[2] \dicFlx{(‑s, ‑)}[5] \textbf{1.} \dicDirectTranslationCS{umístění, položení}  \textbf{2.} \dicFieldCat{anat.} \dicDirectTranslationCS{děloha}
\dicEntry[lega] \dicTerm{leg|a} \dicsymFrequent\  \dicIPA{{l}{\textepsilon}{\textlengthmark}{\textbabygamma}{a}} \dicPos{f}[1] \dicFlx{(‑u, ‑ur)}[19] \textbf{1.} \dicSynonym{staðsetning} \dicDirectTranslationCS{uložení, pozice} \dicExampleIS{landfræðileg lega Íslands} \dicExampleCS{geografická pozice Islandu}  \textbf{2.} \dicSynonym{rúmlega\smash{\textsuperscript{1}}} \dicDirectTranslationCS{ležení (na lůžku)}  \textbf{3.} \dicFieldCat{techn.} \dicDirectTranslationCS{ložisko}
\dicEntry[leggja] \dicTerm{leggja} \dicsymFrequent\  \dicIPA{{l}{\textepsilon}{\r{\textObardotlessj}}{\textlengthmark}{a}} \dicPos{v}[4] \dicFlx{(legg, lagði, lögðum, legði, lagt)}[20] \dicFlx{acc\,/\addthin dat} \textbf{1.} \dicFlx{acc} \dicDirectTranslationCS{položit, pokládat, umístit, umísťovat, uložit, ukládat} \dicExampleIS{leggja hlutinn á borðið} \dicExampleCS{položit věc na stůl};  \dicPhraseIS{leggja sig} \dicDirectTranslationCS{položit se, lehnout si, natáhnout se}  \textbf{2.} \dicFlx{acc} \dicDirectTranslationCS{(na)instalovat (potrubí ap.), (na)táhnout (elek. vedení ap.), položit, pokládat (dlaždice ap.)} \dicExampleIS{leggja nýjan veg} \dicExampleCS{položit novou silnici}  \textbf{3.} \dicDirectTranslationCS{vyjít, vyjet, vyrazit};  \dicPhraseIS{leggja af\,/\addthin á stað} \dicDirectTranslationCS{vyrazit (na cestu ap.), vydat se} \dicExampleIS{Hún leggur af stað til Akureyrar.} \dicExampleCS{Vyráží do Akureyri.}  \textbf{4.} \dicFlx{dat} \dicDirectTranslationCS{(za)parkovat} \dicExampleIS{leggja bílnum hjá búðinni} \dicExampleCS{zaparkovat auto u~obchodu}  \textbf{5.} \dicFlx{dat} \dicDirectTranslationCS{vyřadit, vyřazovat} \dicExampleIS{leggja gömlu kaffivélinni} \dicExampleCS{vyřadit starý kávovar}  \textbf{6.} \dicFlx{acc} \dicDirectTranslationCS{porazit, porážet (protivníka ap.)}  \textbf{7.} \dicFlx{acc} \dicDirectTranslationCS{(za)útočit, napadnout, napadat}  \textbf{8.} \dicPhraseIS{e‑ð leggur} \dicFlx{impers} {\textbf{a.}} \dicDirectTranslationCS{(co) se šíří, (co) se rozšiřuje (vůně ap.)};  {\textbf{b.}} \dicDirectTranslationCS{(co) zamrzá};  \dicIdiom{leggja}[að]{ \dicPhraseIS{leggja hart að sér}} \dicDirectTranslationCS{tvrdě pracovat, vyvinout veškeré úsilí}; { \dicPhraseIS{leggja fast að e‑m}} \dicDirectTranslationCS{vyvíjet tlak na (koho)}; { \dicPhraseIS{leggja e‑ð að veði}} \dicDirectTranslationCS{dát (co) v~sázku, vsadit (co)};  \dicIdiom{leggja}[af]{ \dicPhraseIS{leggja af}} \dicDirectTranslationCS{(z)hubnout} \dicExampleIS{Hún lagði af vegna þess að hún var veik.} \dicExampleCS{Zhubla, protože byla nemocná.};  \dicIdiom{leggja}[aftur]{ \dicPhraseIS{leggja aftur e‑ð}} \dicSynonym{loka\smash{\textsuperscript{2}}} \dicDirectTranslationCS{zavřít (co), zamknout (co)} \dicExampleIS{leggja aftur augun} \dicExampleCS{zavřít oči};  \dicIdiom{leggja}[á]{ \dicPhraseIS{leggja á borð}} \dicDirectTranslationCS{prostřít (stůl)}; { \dicPhraseIS{leggja ást á e‑n}} \dicDirectTranslationCS{zamilovat se do (koho), milovat (koho)}; { \dicPhraseIS{leggja áherslu á e‑ð}} \dicDirectTranslationCS{zdůraznit\,/\addthin zdůrazňovat (co), klást na (co) důraz}; { \dicPhraseIS{leggja mikið á sig}} \dicSynonym*{beita sig hörku} \dicDirectTranslationCS{opřít se do toho, vynasnažit se}; { \dicPhraseIS{leggja á}} \dicDirectTranslationCS{zavěsit, položit (telefon)} \dicExampleIS{Hann lagði símann á.} \dicExampleCS{Položil telefon.}; { \dicPhraseIS{leggja e‑ð á minnið, leggja e‑ð í minni}} \dicDirectTranslationCS{uložit (co) do paměti, zapamatovat si (co)}; { \dicPhraseIS{leggja skatt á e‑ð}} \dicDirectTranslationCS{(z)danit (co), uvalit daň na (co)}; { \dicPhraseIS{leggja hatur á e‑n}} \dicDirectTranslationCS{nenávidět (koho)}; { \dicPhraseIS{leggja dul á e‑ð}} \dicDirectTranslationCS{(u)tajit (co)}; { \dicPhraseIS{leggja stund á e‑ð}} \dicDirectTranslationCS{věnovat se (čemu)};  \dicIdiom{leggja}[fram]{ \dicPhraseIS{leggja sig allan fram}} \dicDirectTranslationCS{vynasnažit se, vyvinout úsilí}; { \dicPhraseIS{leggja fram e‑ð}} \dicDirectTranslationCS{předložit\,/\addthin předkládat (co) (k~diskuzi ap.)};  \dicIdiom{leggja}[fyrir]{ \dicPhraseIS{leggja e‑ð fyrir e‑n}} \dicDirectTranslationCS{předložit\,/\addthin předkládat (komu co) (otázku ap.)}; { \dicPhraseIS{leggja fyrir sig e‑ð}} \dicDirectTranslationCS{zabývat se (čím), být zapojen do (čeho)}; { \dicPhraseIS{leggja fyrir}} \dicDirectTranslationCS{odkládat, šetřit};  \dicIdiom{leggja}[inn]{ \dicPhraseIS{leggja e‑ð inn}} \dicDirectTranslationCS{uložit\,/\addthin ukládat (co)} \dicExampleIS{Ég legg peninga inn í bankann.} \dicExampleCS{Ukládám si peníze do banky.};  \dicIdiom{leggja}[í]{ \dicPhraseIS{leggja í e‑ð}} \dicDirectTranslationCS{odvážit se na (co)}; { \dicPhraseIS{leggja í hann}} \dicDirectTranslationCS{vydat se na cestu, vyrazit};  \dicIdiom{leggja}[niður]{ \dicPhraseIS{leggja niður e‑ð}} \dicDirectTranslationCS{skončit s~(čím), skoncovat s~(čím)} \dicExampleIS{leggja niður starf} \dicExampleCS{skončit s~prací}; { \dicPhraseIS{leggja niður vopn}} \dicDirectTranslationCS{složit zbraně}; \textls[15]{{ \dicPhraseIS{leggja sig niður við e‑ð}} \dicDirectTranslationCS{sní\nobreak žit se k~(čemu)};}  \dicIdiom{leggja}[saman]{ \dicPhraseIS{leggja saman e‑ð}} \dicFieldCat{mat.} \dicDirectTranslationCS{sčítat (co), přičítat (co)} \dicExampleIS{Barnið leggur saman tvo og þrjá og fær fimm.} \dicExampleCS{Dítě sečte dva plus tři a~vyjde mu pět.}; { \dicPhraseIS{leggja saman tvo og tvo}} \dicLangCat{přen.} \dicDirectTranslationCS{dát si dvě a~dvě dohromady};  \dicIdiom{leggja}[til]{ \dicPhraseIS{leggja e‑m til e‑ð}} \dicDirectTranslationCS{poskytnout\,/\addthin poskytovat (komu co)}; { \dicPhraseIS{leggja e‑ð til}} \dicDirectTranslationCS{navrhnout\,/\addthin navrhovat (co), předložit návrh o~(čem)} \dicExampleIS{Ég legg til að þessi plaköt verði lagfærð.} \dicExampleCS{Navrhuji, aby tyto plakáty byly upraveny.};  \dicPhraseIS{leggja sér e‑ð til munns} \dicDirectTranslationCS{jíst (co), pojídat (co)};  \dicIdiom{leggja}[undir]{ \dicPhraseIS{leggja undir sig e‑ð}} \dicDirectTranslationCS{podmanit\,/\addthin podmaňovat si (co), podrobit\,/\addthin podrobovat si (co)};  \dicIdiom{leggja}[upp]{ \dicPhraseIS{leggja upp}} \dicDirectTranslationCS{vyrazit, vydat se (na cestu ap.)} \dicExampleIS{leggja upp í ferðina} \dicExampleCS{vydat se na cestu};  \dicIdiom{leggja}[upp úr]{ \dicPhraseIS{leggja mikið\,/\addthin lítið upp úr e‑u}} \dicDirectTranslationCS{pokládat (co) za hodně\,/\addthin málo důležité};  \dicIdiom{leggja}[út]{ \dicPhraseIS{leggja út e‑ð}} {\textbf{a.}} \dicSynonym{borga} \dicDirectTranslationCS{vyplatit (co), přispět (čím) (částkou ap.)};  {\textbf{b.}} \dicSynonym{túlka} \dicDirectTranslationCS{přeložit (co), tlumočit (co)};  \dicIdiom{leggja}[við]{ \dicPhraseIS{leggja rækt við e‑ð}} \dicDirectTranslationCS{kultivovat (co), rozvíjet (co)}; { \dicPhraseIS{leggja e‑ð við e‑ð}} \dicDirectTranslationCS{přidat (co) k~(čemu)} \dicExampleIS{leggja tvo við þrjá} \dicExampleCS{přidat dvě ke třem};  \dicIdiom{leggjast}{ \dicPhraseIS{leggjast}} \dicFlx{refl} \dicDirectTranslationCS{lehnout si, natáhnout se};  \dicIdiom{leggjast}[að]{ \dicPhraseIS{leggjast að bryggju}} \dicFlx{refl} \dicDirectTranslationCS{přistát u~mola};  \dicIdiom{leggjast}[af]{ \dicPhraseIS{leggjast af}} \dicFlx{refl} \dicDirectTranslationCS{vymizet, zaniknout};  \dicIdiom{leggjast}[á]{ \dicPhraseIS{leggjast á eitt}} \dicFlx{refl} \dicDirectTranslationCS{spojit síly, táhnout za jeden provaz}; { \dicPhraseIS{e‑að leggst þungt á e‑n}} \dicFlx{refl} \dicLangCat{přen.} \dicDirectTranslationCS{(co koho) tíží (obavy ap.)};  \dicIdiom{leggjast}[gegn]{ \dicPhraseIS{leggjast gegn e‑u}} \dicFlx{refl} \dicDirectTranslationCS{postavit\,/\addthin stavět se proti (čemu)};  \dicIdiom{leggjast}[í]{ \dicPhraseIS{leggjast í drykkju}} \dicFlx{refl} \dicDirectTranslationCS{oddávat se pití (alkoholu)}; { \dicPhraseIS{leggjast í þunglyndi}} \dicFlx{refl} \dicDirectTranslationCS{upadnout do deprese};  \dicIdiom{leggjast}[með]{ \dicPhraseIS{leggjast með e‑m}} \dicFlx{refl} \dicDirectTranslationCS{vyspat se s~(kým), spát s~(kým)};  \dicIdiom{leggjast}[niður]{ \dicPhraseIS{leggjast niður}} \dicFlx{refl} \dicDirectTranslationCS{být přerušen, přestat};  \dicIdiom{leggjast}[yfir]{ \dicPhraseIS{e‑að leggst yfir e‑ð}} \dicFlx{refl} \dicDirectTranslationCS{(co) se usazuje nad (čím) (mlha nad městem ap.)}; { \dicPhraseIS{leggjast yfir e‑ð}} \dicFlx{refl} \dicDirectTranslationCS{pročíst (co), (důkladně) přečíst (co)}
\dicEntry[leggur] \dicTerm{legg|ur} \dicsymFrequent\  \dicIPA{{l}{\textepsilon}{\r{g}}{\textlengthmark}{\textscy}{\textsubring{r}}} \dicPos{m}[9] \dicFlx{(‑jar\,/\addthin ‑s, ‑ir)}[26] \textbf{1.} \dicSynonym{bein} \dicDirectTranslationCS{kost}  \textbf{2.} \dicSynonym{fótleggur} \dicDirectTranslationCS{noha, dolní končetina} \dicExampleIS{Ég er tilfinningalaus í leggnum.} \dicExampleCS{Nemám v~noze žádný cit.}  \textbf{3.} \dicFieldCat{bot.} \dicSynonym{stilkur} \dicDirectTranslationCS{stopka, stonek};  \dicIdiom{leggur}{ \dicPhraseIS{komast á legg}} \dicFlx{refl} \dicLangCat{přen.} \dicDirectTranslationCS{postavit se na vlastní nohy}
\dicEntry[leggöng] \dicTerm{leg··göng} \dicIPA{{l}{\textepsilon}{\textbabygamma}{\r{g}}{\oe i}{\ng}{\r{g}}} \dicPos{n}[2] \dicFlx{pl}[9] \dicFieldCat{anat.} \dicDirectTranslationCS{vagina, pochva}
\dicEntry[legið] \dicTerm{legið} \dicIPA{{l}{ei}{\textsci}{\texttheta}} \dicPos{v} \dicFlx{supin} \dicLink{liggja}
\dicEntry[legir] \dicTerm{legir} \dicIPA{{l}{ei}{\textsci}{\textsubring{r}}} \dicPos{m} \dicFlx{pl nom} \dicLink{lögur}
\dicEntry[legkaka] \dicTerm{leg··|kaka} \dicIPA{{l}{\textepsilon}{\textbabygamma}{k\smash{\textsuperscript{h}}}{a}{\r{g}}{a}} \dicPos{f}[1] \dicFlx{(‑köku, ‑kökur)}[20] \dicFieldCat{anat.} \dicDirectTranslationCS{placenta, plodový obal}
\dicEntry[legstaður] \dicTerm{leg··stað|ur} \dicIPA{{l}{\textepsilon}{x}{s}{\textsubring{d}}{a}{ð}{\textscy}{\textsubring{r}}} \dicPos{m}[10] \dicFlx{(‑ar, ‑ir)}[14] \dicSynonym{gröf} \dicDirectTranslationCS{hrob(ka)}
\dicEntry[legsteinn] \dicTerm{leg··stein|n} \dicIPA{{l}{\textepsilon}{x}{s}{\textsubring{d}}{ei}{\textsubring{d}}{\textsubring{n}}} \dicPos{m}[6] \dicFlx{(‑s, ‑ar)}[42] \dicDirectTranslationCS{náhrobní kámen}
\dicEntry[legubekkur] \dicTerm{legu··bekk|ur} \dicIPA{{l}{\textepsilon}{\textlengthmark}{\textbabygamma}{\textscy}{\textsubring{b}}{\textepsilon}{h}{\r{g}}{\textscy}{\textsubring{r}}} \dicPos{m}[9] \dicFlx{(‑s\,/\addthin ‑jar, ‑ir)}[26] \dicDirectTranslationCS{pohovka, lehátko}
\dicEntry[legudagur] \dicTerm{legu··dag|ur} \dicIPA{{l}{\textepsilon}{\textlengthmark}{\textbabygamma}{\textscy}{\textsubring{d}}{a}{\textbabygamma}{\textscy}{\textsubring{r}}} \dicPos{m}[6] \dicFlx{(‑s, ‑ar)}[62] \textbf{1.} \dicDirectTranslationCS{den na lůžku (v~nemocnici ap.)}  \textbf{2.} \dicDirectTranslationCS{den (lodi ap.) v~přístavu}
\dicEntry[legugjald] \dicTerm{legu··|gjald} \dicIPA{{l}{\textepsilon}{\textlengthmark}{\textbabygamma}{\textscy}{\r{\textObardotlessj}}{a}{l}{\textsubring{d}}} \dicPos{n}[2] \dicFlx{(‑gjalds, ‑gjöld)}[8] \dicFieldCat{nám.} \dicDirectTranslationCS{kotevné} \dicIndirectTranslationCS{(poplatek za kotvení)}
\dicEntry[legusár] \dicTerm{legu··sár} \dicIPA{{l}{\textepsilon}{\textlengthmark}{\textbabygamma}{\textscy}{s}{au}{\textsubring{r}}} \dicPos{n}[2] \dicFlx{(‑s, ‑)}[5] \dicFieldCat{med.} \dicDirectTranslationCS{proleženina}
\dicEntry[leið] \dicTerm{leið\smash{\textsuperscript{1}}} \dicsymFrequent\  \dicIPA{{l}{ei}{\textlengthmark}{\texttheta}} \dicPos{f}[7] \dicFlx{(‑ar, ‑ir)}[1] \textbf{1.} \dicSynonym{vegur} \dicDirectTranslationCS{cesta, stezka} \dicExampleIS{bein leið} \dicExampleCS{přímá cesta};  \dicPhraseIS{eiga leið hjá} \dicDirectTranslationCS{mít cestu kolem};  \dicPhraseIS{vera á leið(inni) (þangað)} \dicDirectTranslationCS{být na cestě (tam)}  \textbf{2.} \dicSynonym{háttur\smash{\textsuperscript{1}}} \dicDirectTranslationCS{způsob, postup};  \dicPhraseIS{það er engin leið að \dots{}} \dicDirectTranslationCS{není žádný způsob, jak\dots{}}  \textbf{3.} \dicSynonym{vegalengd} \dicDirectTranslationCS{cesta, vzdálenost} \dicExampleIS{dagleið} \dicExampleCS{jednodenní cesta}  \textbf{4.} \dicSynonym{áætlunarleið} \dicDirectTranslationCS{(dopravní) linka};  \dicIdiom{leið}{ \dicPhraseIS{á miðri leið}} \dicFlx{adv} \dicDirectTranslationCS{v~půli cesty}; { \dicPhraseIS{heim á leið}} \dicFlx{adv} \dicDirectTranslationCS{domů}; { \dicPhraseIS{leggja leið sína, leggja leiðir sínar}} \dicDirectTranslationCS{vydat se (do cizí země ap.)} \dicExampleIS{leggja leið sína til Íslands} \dicExampleCS{vydat se na Island};  \dicPhraseIS{um leið} \dicFlx{adv} \dicDirectTranslationCS{zároveň, současně, naráz}; { \dicPhraseIS{um leið og}} \dicFlx{conj} \dicSynonym*{strax og} \dicDirectTranslationCS{jakmile, hned jak} \dicExampleIS{Hann fór að sofa um leið og hann kom heim.} \dicExampleCS{Jakmile přišel domů, šel si lehnout.}; { \dicPhraseIS{vera komin (fimm) mánuði á leið}} \dicDirectTranslationCS{být v~(pátém) měsíci (těhotenství)}
\dicEntry[leið] \dicTerm{leið\smash{\textsuperscript{2}}} \dicIPA{{l}{ei}{\textlengthmark}{\texttheta}} \dicPos{v} \dicFlx{ind pf sg 1 pers} \dicLink{líða}
\dicEntry[leiða] \dicTerm{lei|ða} \dicsymFrequent\  \dicIPA{{l}{ei}{\textlengthmark}{ð}{a}} \dicPos{v}[2] \dicFlx{(‑ddi, ‑tt)}[169] \dicFlx{acc} \textbf{1.} \dicDirectTranslationCS{(při)vést, přivádět, vodit} \dicExampleIS{leiða e‑n þangað} \dicExampleCS{přivést (koho) tam}  \textbf{2.} \dicDirectTranslationCS{(při)vést, přivádět (elektřinu ap.)} \dicExampleIS{leiða rafmagn í hús} \dicExampleCS{přivádět elektřinu do domu};  \dicIdiom{leiða}[að]{ \dicPhraseIS{leiða hugann að e‑u}} \dicDirectTranslationCS{zavést myšlenku na (co)}; { \dicPhraseIS{leiða tal að e‑u}} \dicDirectTranslationCS{zavést řeč na (co), nadhodit (co)};  \dicIdiom{leiða}[af]{ \dicPhraseIS{leiða e‑ð af sér}} \dicDirectTranslationCS{mít (co) za následek, vést k~(čemu), skončit se (čím)};  \dicIdiom{leiða}[fyrir]{ \dicPhraseIS{leiða e‑ð e‑m fyrir sjónir}} \dicDirectTranslationCS{ukázat (komu co), vysvětlit (komu co)};  \dicIdiom{leiða}[hjá]{ \dicPhraseIS{leiða e‑ð hjá sér}} \dicDirectTranslationCS{nevšímat si (čeho), nebrat na vědomí (co)} \dicExampleIS{leiða orð hans hjá sér} \dicExampleCS{nevšímat si jeho slov};  \dicIdiom{leiða}[í]{ \dicPhraseIS{leiða e‑u í ljós}} \dicDirectTranslationCS{(po)odhalit (co)};  \dicIdiom{leiða}[til]{ \dicPhraseIS{leiða til e‑s}} \dicDirectTranslationCS{vést k~(čemu), vyústit v~(co)};  \dicIdiom{leiðast}{ \dicPhraseIS{leiðast}} \dicFlx{refl} \dicDirectTranslationCS{vést se za ruce} \dicExampleIS{Þau leiddust niður götuna.} \dicExampleCS{Šli po ulici a~vedli se za ruce.}; { \dicPhraseIS{leiðast út í drykkju}} \dicFlx{refl} \dicDirectTranslationCS{začít pít (alkohol)}; { \dicPhraseIS{láta til leiðast}} \dicDirectTranslationCS{nechat se přesvědčit};  \dicPhraseIS{e‑m leiðist} \dicFlx{refl impers} \dicDirectTranslationCS{(kdo) se nudí, (komu) je dlouhá chvíle} \dicExampleIS{Mér leiðist.} \dicExampleCS{Nudím se.}; { \dicPhraseIS{e‑m leiðist e‑að}} \dicFlx{refl impers} \dicDirectTranslationCS{(koho) nudí (co), (komu) přijde nudné (co)}
\dicEntry[leiðabók] \dicTerm{leiða··|bók} \dicIPA{{l}{ei}{\textlengthmark}{ð}{a}{\textsubring{b}}{ou}{\r{g}}} \dicPos{f}[8] \dicFlx{(‑bókar, ‑bækur)}[5] \dicDirectTranslationCS{jízdní řád}
\dicEntry[leiðangur] \dicTerm{leið··ang|ur} \dicsymFrequent\  \dicIPA{{l}{ei}{\textlengthmark}{ð}{au}{\ng}{\r{g}}{\textscy}{\textsubring{r}}} \dicPos{m}[5] \dicFlx{(‑urs, ‑rar)}[3] \dicSynonym{könnunarferð} \dicDirectTranslationCS{výprava, expedice, mise} \dicExampleIS{Leiðangurinn tekur tvær vikur.} \dicExampleCS{Expedice potrvá dva týdny.}
\dicEntry[leiðarendi] \dicTerm{leiðar··end|i} \dicIPA{{l}{ei}{\textlengthmark}{ð}{a}{r}{\textepsilon}{n}{\textsubring{d}}{\textsci}} \dicPos{m}[1] \dicFlx{(‑a, ‑ar)}[1] \dicSynonym{ákvörðunarstaður} \dicDirectTranslationCS{cíl cesty, destinace} \dicExampleIS{komast á leiðarenda} \dicExampleCS{dorazit do cíle cesty}
\dicEntry[leiðari] \dicTerm{leið··ar|i} \dicIPA{{l}{ei}{\textlengthmark}{ð}{a}{r}{\textsci}} \dicPos{m}[1] \dicFlx{(‑a, ‑ar)}[13] \textbf{1.} \dicSynonym{grein} \dicDirectTranslationCS{úvodník} \dicIndirectTranslationCS{(článek umístěný na první stránce novin)}  \textbf{2.} \dicFieldCat{fyz.} \dicDirectTranslationCS{vodič (elektřiny ap.)}
\dicEntry[leiðarljós] \dicTerm{leiðar··ljós} \dicIPA{{l}{ei}{\textlengthmark}{ð}{a}{r}{l}{j}{ou}{s}} \dicPos{n}[2] \dicFlx{(‑s, ‑)}[5] \dicLangCat{přen.} \dicDirectTranslationCS{záchytný bod, vodítko};  \dicPhraseIS{hafa e‑ð að leiðarljósi} \dicLangCat{přen.} \dicDirectTranslationCS{mít (co) jako záchytný bod}
\dicEntry[leiðarlok] \dicTerm{leiðar··lok} \dicIPA{{l}{ei}{\textlengthmark}{ð}{a}{r}{l}{\textopeno}{\r{g}}} \dicPos{n}[2] \dicFlx{pl}[1] \dicDirectTranslationCS{závěr\,/\addthin konec cesty (silnice i~putování)}
\dicEntry[leiðarlýsing] \dicTerm{leiðar··lýs·ing} \dicIPA{{l}{ei}{\textlengthmark}{ð}{a}{r}{l}{i}{s}{i}{\ng}{\r{g}}} \dicPos{f}[4] \dicFlx{(‑ar, ‑ar)}[5] \dicDirectTranslationCS{itinerář, cestovní plán}
\dicEntry[leiðarminni] \dicTerm{leiðar··minni} \dicIPA{{l}{ei}{\textlengthmark}{ð}{a}{r}{m}{\textsci}{n}{\textsci}} \dicPos{n}[2] \dicFlx{(‑s, ‑)}[14] \dicDirectTranslationCS{příznačný motiv, leitmotiv}
\dicEntry[leiðarvísir] \dicTerm{leiðar··vís|ir} \dicIPA{{l}{ei}{\textlengthmark}{ð}{a}{r}{v}{i}{s}{\textsci}{\textsubring{r}}} \dicPos{m}[7] \dicFlx{(‑is, ‑ar)}[1] \dicDirectTranslationCS{příručka, manuál, návod} \dicExampleIS{leiðarvísir um vél} \dicExampleCS{manuál k~zařízení}
\dicEntry[leiðbeina] \dicTerm{leið··bein|a} \dicIPA{{l}{ei}{ð}{\textsubring{b}}{ei}{n}{a}} \dicPos{v}[2] \dicFlx{(‑di, ‑t)}[152] \dicFlx{dat} \dicSynonym{kenna} \dicDirectTranslationCS{učit, vyučovat, cvičit, školit} \dicExampleIS{leiðbeina e‑m í e‑u} \dicExampleCS{školit (koho) v~(čem)}
\dicEntry[leiðbeinandi] \dicTerm{leið··bein·|andi} \dicIPA{{l}{ei}{ð}{\textsubring{b}}{ei}{n}{a}{n}{\textsubring{d}}{\textsci}} \dicPos{m}[2] \dicFlx{(‑anda, ‑endur)}[1] \textbf{1.} \dicDirectTranslationCS{školitel(ka), instruktor(ka), cvičitel(ka)}  \textbf{2.} \dicFieldCat{škol.} \dicDirectTranslationCS{nekvalifikovaný učitel, nekvalifikovaná učitelka}
\dicEntry[leiðbeining] \dicTerm{leið··bein·ing} \dicIPA{{l}{ei}{ð}{\textsubring{b}}{ei}{n}{i}{\ng}{\r{g}}} \dicPos{f}[4] \dicFlx{(‑ar, ‑ar)}[5] \textbf{1.} \dicSynonym{leiðsögn} \dicDirectTranslationCS{pokyn, instrukce}  \textbf{2.} \dicPhraseIS{leiðbeiningar} \dicFlx{pl} \dicSynonym{leiðarvísir} \dicDirectTranslationCS{pokyny, návod, manuál, instrukce} \dicExampleIS{leiðbeiningar um notkun} \dicExampleCS{návod k~použití}
\dicEntry[leiði] \dicTerm{leið|i\smash{\textsuperscript{1}}} \dicIPA{{l}{ei}{\textlengthmark}{ð}{\textsci}} \dicPos{m}[1] \dicFlx{(‑a)}[3] \dicSynonym{leiðindi} \dicDirectTranslationCS{nuda, dlouhá chvíle} \dicExampleIS{fá leiða á leikjum} \dicExampleCS{nudit se u~hry}
\dicEntry[leiði] \dicTerm{leiði\smash{\textsuperscript{2}}} \dicIPA{{l}{ei}{\textlengthmark}{ð}{\textsci}} \dicPos{n}[2] \dicFlx{(‑s, ‑)}[14] \textbf{1.} \dicSynonym{gröf} \dicDirectTranslationCS{náhrobek, hrobka} \dicExampleIS{setja kross á leiðið} \dicExampleCS{umístit kříž na náhrobek}  \textbf{2.} \dicSynonym{meðvindur} \dicDirectTranslationCS{příznivý vítr}  \textbf{3.} \dicSynonym{færð} \dicDirectTranslationCS{sjízdnost}
\dicEntry[leiðindamál] \dicTerm{leiðinda··mál} \dicIPA{{l}{ei}{\textlengthmark}{ð}{\textsci}{n}{\textsubring{d}}{a}{m}{au}{\textsubring{l}}} \dicPos{n}[2] \dicFlx{(‑s, ‑)}[5] \dicDirectTranslationCS{nepříjemnost, nepříjemná\,/\addthin otravná věc}
\dicEntry[leiðindi] \dicTerm{leið··indi} \dicIPA{{l}{ei}{\textlengthmark}{ð}{\textsci}{n}{\textsubring{d}}{\textsci}} \dicPos{n}[2] \dicFlx{pl}[19] \textbf{1.} \dicSynonym{leiði\smash{\textsuperscript{1}}} \dicDirectTranslationCS{nuda, nudnost, dlouhá chvíle}  \textbf{2.} \dicSynonym{óánægja} \dicDirectTranslationCS{protivnost, špatná nálada}
\dicEntry[leiðinlegur] \dicTerm{leiðin··legur} \dicsymFrequent\  \dicIPA{{l}{ei}{\textlengthmark}{ð}{\textsci}{n}{l}{\textepsilon}{\textbabygamma}{\textscy}{\textsubring{r}}} \dicPos{adj}[1]\dicFlx{}[-8] \textbf{1.} \dicDirectTranslationCS{nudný, fádní} \dicExampleIS{leiðinlegur kennari} \dicExampleCS{nudný učitel}  \textbf{2.} \dicSynonym{dapurlegur} \dicDirectTranslationCS{smutný, tristní}  \textbf{3.} \dicDirectTranslationCS{protivný, nepříjemný (počasí ap.)}
\dicEntry[leiðitamur] \textls[15]{\dicTerm{leiði··|tamur} \dicIPA{{l}{ei}{\textlengthmark}{ð}{\textsci}{t\smash{\textsuperscript{h}}}{a}{m}{\textscy}{\textsubring{r}}} \dicPos{adj}[1] \dicFlx{(f ‑töm)}[2] \dicSynonym{eftirlátur}} \dicDirectTranslationCS{svolný, povolný}
\dicEntry[leiðni] \dicTerm{leiðn|i} \dicIPA{{l}{ei}{ð}{n}{\textsci}} \dicPos{f}[3] \dicFlx{(‑i)}[3] \dicFieldCat{fyz.} \dicDirectTranslationCS{vodivost}
\dicEntry[leiðrétta] \dicTerm{leið··rétt|a} \dicIPA{{l}{ei}{ð}{r}{j}{\textepsilon}{h}{\textsubring{d}}{a}} \dicPos{v}[2] \dicFlx{(‑i, ‑)}[1] \dicFlx{acc} \dicDirectTranslationCS{opravit, opravovat, (z)korigovat (chyby ap.)} \dicExampleIS{leiðrétta verkefnið} \dicExampleCS{opravit cvičení}
\dicEntry[leiðrétting] \dicTerm{leið··rétt·ing} \dicIPA{{l}{ei}{ð}{r}{j}{\textepsilon}{h}{\textsubring{d}}{i}{\ng}{\r{g}}} \dicPos{f}[4] \dicFlx{(‑ar, ‑ar)}[5] \dicDirectTranslationCS{oprava, opravení, korekce, korektura} \dicExampleIS{leiðrétting  á mistökum} \dicExampleCS{oprava chyb}
\dicEntry[leiðsla] \dicTerm{leiðsl|a} \dicIPA{{l}{ei}{ð}{s}{\textsubring{d}}{l}{a}} \dicPos{f}[1] \dicFlx{(‑u, ‑ur)}[19] \textbf{1.} \dicSynonym{leiðsögn} \dicDirectTranslationCS{pokyn, instrukce}  \textbf{2.} \dicSynonym{rör} \dicDirectTranslationCS{roura, trubka, potrubí}  \textbf{3.} \dicSynonym{rafleiðsla} \dicDirectTranslationCS{(elektrické) vedení}  \textbf{4.} \dicSynonym*{ofurhrifning} \dicDirectTranslationCS{vytržení, trans} \dicExampleIS{vera í leiðslu} \dicExampleCS{být v~transu}
\dicEntry[leiðsögn] \dicTerm{leið··|sögn} \dicIPA{{l}{ei}{ð}{s}{\oe}{\r{g}}{\textsubring{n}}} \dicPos{f}[7] \dicFlx{(‑sagnar)}[19] \dicDirectTranslationCS{vedení, navádění, provádění}
\dicEntry[leiðsögumaður] \dicTerm{leið·sögu··|maður} \dicIPA{{l}{ei}{ð}{s}{\oe}{\textbabygamma}{\textscy}{m}{a}{ð}{\textscy}{\textsubring{r}}} \dicPos{m}[13] \dicFlx{(‑manns, ‑menn)}[2] \dicDirectTranslationCS{průvodce, průvodkyně} \dicExampleIS{leiðsögumaður ferðamanna} \dicExampleCS{průvodce turistů}
\dicEntry[leiðtogi] \dicTerm{leið··tog|i} \dicIPA{{l}{ei}{ð}{t\smash{\textsuperscript{h}}}{\textopeno i}{j}{\textsci}} \dicPos{m}[1] \dicFlx{(‑a, ‑ar)}[1] \dicSynonym{foringi} \dicDirectTranslationCS{vedoucí představitel(ka), lídr(yně)} \dicExampleIS{Flokkurinn hefur kosið sér leiðtoga.} \dicExampleCS{Strana si zvolila vedoucího představitele.}
\dicEntry[leiður] \dicTerm{leiður} \dicsymFrequent\  \dicIPA{{l}{ei}{\textlengthmark}{ð}{\textscy}{\textsubring{r}}} \dicPos{adj}[2]\dicFlx{}[-6] \textbf{1.} \dicSynonym{sorgmæddur} \dicDirectTranslationCS{zarmoucený, cítící lítost} \dicExampleIS{leiður út af e‑u} \dicExampleCS{zarmoucený kvůli čemu}  \textbf{2.} \dicSynonym{hvimleiður} \dicDirectTranslationCS{znuděný, unavený (z~poslouchání nářků ap.)} \dicExampleIS{vera orðinn leiður á verkinu} \dicExampleCS{být už unavený z~práce}
\dicEntry[leif] \dicTerm{leif} \dicsymFrequent\  \dicIPA{{l}{ei}{\textlengthmark}{f}} \dicPos{f}[4] \dicFlx{(‑ar, ‑ar)}[1] \textbf{1.} \dicLangCat{básn.} \dicSynonym{arfur} \dicDirectTranslationCS{dědictví, pozůstalost}  \textbf{2.} \dicPhraseIS{leifar} \dicFlx{pl} \dicDirectTranslationCS{zbytky, ostatky} \dicExampleIS{matarleifar} \dicExampleCS{zbytky jídla}  \textbf{3.} \dicFieldCat{mat.} \dicSynonym{afgangur} \dicDirectTranslationCS{zbytek}
\dicEntry[leifa] \dicTerm{leif|a} \dicIPA{{l}{ei}{\textlengthmark}{v}{a}} \dicPos{v}[2] \dicFlx{(‑ði, ‑t)}[99] \dicFlx{dat} \dicDirectTranslationCS{(za)nechat (jídlo nedojedené ap.)} \dicExampleIS{leifa matnum} \dicExampleCS{nechat jídlo}
\dicEntry[leiftra] \dicTerm{leiftr|a} \dicIPA{{l}{ei}{f}{\textsubring{d}}{r}{a}} \dicPos{v}[1] \dicFlx{(‑aði)}[1] \dicSynonym{glitra} \dicDirectTranslationCS{blýskat se, blyštět se, třpytit se} \dicExampleIS{Stjörnurnar leiftra á himninum.} \dicExampleCS{Hvězdy se blyští na obloze.}
\dicEntry[leiftur] \dicTerm{leift|ur\smash{\textsuperscript{1}}} \dicIPA{{l}{ei}{f}{\textsubring{d}}{\textscy}{\textsubring{r}}} \dicPos{m}[5] \dicFlx{(‑urs, ‑rar)}[1] \dicFieldCat{zool.} \dicDirectTranslationCS{plískavice běloboká} \textit{(l.~{\textLA{Lagenorhynchus acutus}})}  \dicsymPhoto\ 
\dicFigure{ds_image_leiftur_1_1.jpg}{Leiftur}{Leiftur - Twp, CC BY-SA 3.0}
\dicEntry[leiftur] \dicTerm{leiftur\smash{\textsuperscript{2}}} \dicIPA{{l}{ei}{f}{\textsubring{d}}{\textscy}{\textsubring{r}}} \dicPos{n}[2] \dicFlx{(‑s, ‑)}[25] \dicSynonym{blossi} \dicDirectTranslationCS{(za)blýsknutí, zá\-blesk}
\dicEntry[leifturhraði] \dicTerm{leiftur··hrað|i} \dicIPA{{l}{ei}{f}{\textsubring{d}}{\textscy}{\textsubring{r}}{a}{ð}{\textsci}} \dicPos{m}[1] \dicFlx{(‑a)}[3] \dicLangCat{přen.} \dicDirectTranslationCS{rychlost blesku, blesková rychlost}
\dicEntry[leiftursókn] \dicTerm{leiftur··sókn} \dicIPA{{l}{ei}{f}{\textsubring{d}}{\textscy}{\textsubring{r}}{s}{ou}{h}{\r{g}}{\textsubring{n}}} \dicPos{f}[7] \dicFlx{(‑ar, ‑ir)}[1] \dicDirectTranslationCS{bleskový útok}
\dicEntry[leiga] \dicTerm{leig|a} \dicsymFrequent\  \dicIPA{{l}{ei}{\textlengthmark}{\textbabygamma}{a}} \dicPos{f}[1] \dicFlx{(‑u, ‑ur)}[19] \textbf{1.} \dicSynonym*{það að leigja} \dicDirectTranslationCS{(pro)nájem} \dicExampleIS{leiga á íbúð} \dicExampleCS{pronájem bytu}  \textbf{2.} \dicSynonym*{leigugjald} \dicDirectTranslationCS{nájem, nájemné} \dicExampleIS{borga leiguna} \dicExampleCS{zaplatit nájem}
\dicEntry[leigja] \dicTerm{leig|ja} \dicsymFrequent\  \dicIPA{{l}{ei}{j}{\textlengthmark}{a}} \dicPos{v}[2] \dicFlx{(‑ði, ‑t)}[89] \dicFlx{acc} \dicDirectTranslationCS{(pro)najmout, (pro)najímat} \dicExampleIS{leigja íbúð af e‑m} \dicExampleCS{pronajmout si od (koho) byt}
\dicEntry[leigjandi] \dicTerm{leigj··|andi} \dicIPA{{l}{ei}{j}{\textlengthmark}{a}{n}{\textsubring{d}}{\textsci}} \dicPos{m}[2] \dicFlx{(‑anda, ‑endur)}[1] \dicDirectTranslationCS{(pod)nájemník, (pod)nájemnice}
\dicEntry[leigubifreið] \dicTerm{leigu··bif·reið} \dicIPA{{l}{ei}{\textlengthmark}{\textbabygamma}{\textscy}{\textsubring{b}}{\textsci}{v}{r}{ei}{\texttheta}} \dicPos{f}[4] \dicFlx{(‑ar, ‑ar\,/\addthin ‑ir)}[30] \dicLangCat{form.} \dicDirectTranslationCS{taxi, taxislužba}
\dicEntry[leigubílastöð] \textls[10]{\dicTerm{leigu·bíla··stöð} \dicIPA{{l}{ei}{\textlengthmark}{\textbabygamma}{\textscy}{\textsubring{b}}{i}{l}{a}{s}{\textsubring{d}}{\oe}{\texttheta}} \dicPos{f}[6] \dicFlx{(‑var, ‑var)}[1] \textbf{1.} \dicDirectTranslationCS{taxi\-služba}} \dicIndirectTranslationCS{(firma)}  \textbf{2.} \dicDirectTranslationCS{stanoviště taxi}
\dicEntry[leigubíll] \dicTerm{leigu··bíl|l} \dicIPA{{l}{ei}{\textlengthmark}{\textbabygamma}{\textscy}{\textsubring{b}}{i}{\textsubring{d}}{\textsubring{l}}} \dicPos{m}[6] \dicFlx{(‑s, ‑ar)}[48] \dicDirectTranslationCS{taxi, taxík}
\dicEntry[leigubílstjóri] \dicTerm{leigu··bíl·stjór|i} \dicIPA{{l}{ei}{\textlengthmark}{\textbabygamma}{\textscy}{\textsubring{b}}{i}{l}{s}{\textsubring{d}}{j}{ou}{r}{\textsci}} \dicPos{m}[1] \dicFlx{(‑a, ‑ar)}[1] \dicDirectTranslationCS{taxikář(ka), řidič(ka) taxi}
\dicEntry[leiguflug] \dicTerm{leigu··flug} \dicIPA{{l}{ei}{\textlengthmark}{\textbabygamma}{\textscy}{f}{l}{\textscy}{x}} \dicPos{n}[2] \dicFlx{(‑s, ‑)}[5] \dicDirectTranslationCS{charterový let}
\dicEntry[leiguíbúð] \dicTerm{leigu··í·búð} \dicIPA{{l}{ei}{\textlengthmark}{\textbabygamma}{\textscy}{i}{\textsubring{b}}{u}{\texttheta}} \dicPos{f}[7] \dicFlx{(‑ar, ‑ir)}[1] \dicDirectTranslationCS{pronajatý byt}
\dicEntry[leiguliði] \dicTerm{leigu··lið|i} \dicIPA{{l}{ei}{\textlengthmark}{\textbabygamma}{\textscy}{l}{\textsci}{ð}{\textsci}} \dicPos{m}[1] \dicFlx{(‑a, ‑ar)}[1] \dicDirectTranslationCS{nájemce, nájemkyně}
\dicEntry[leigumiðlari] \dicTerm{leigu··miðl·ar|i} \dicIPA{{l}{ei}{\textlengthmark}{\textbabygamma}{\textscy}{m}{\textsci}{ð}{l}{a}{r}{\textsci}} \dicPos{m}[1] \dicFlx{(‑a, ‑ar)}[13] \dicDirectTranslationCS{zprostředkovatel(ka) pronájmu}
\dicEntry[leigumorð] \dicTerm{leigu··morð} \dicIPA{{l}{ei}{\textlengthmark}{\textbabygamma}{\textscy}{m}{\textopeno}{r}{\texttheta}} \dicPos{n}[2] \dicFlx{(‑s, ‑)}[5] \dicDirectTranslationCS{nájemná vražda}
\dicEntry[leigumorðingi] \dicTerm{leigu··morð·ing|i} \dicIPA{{l}{ei}{\textlengthmark}{\textbabygamma}{\textscy}{m}{\textopeno}{r}{ð}{i}{\textltailn}{\r{\textObardotlessj}}{\textsci}} \dicPos{m}[1] \dicFlx{(‑ja, ‑jar)}[14] \dicDirectTranslationCS{nájemný vrah, nájemná vražedkyně}
\dicEntry[leigusali] \dicTerm{leigu··sal|i} \dicIPA{{l}{ei}{\textlengthmark}{\textbabygamma}{\textscy}{s}{a}{l}{\textsci}} \dicPos{m}[1] \dicFlx{(‑a, ‑ar)}[8] \dicDirectTranslationCS{pronajímatel(ka)}
\dicEntry[leigusamningur] \dicTerm{leigu··samn·ing|ur} \dicIPA{{l}{ei}{\textlengthmark}{\textbabygamma}{\textscy}{s}{a}{m}{n}{i}{\ng}{\r{g}}{\textscy}{\textsubring{r}}} \dicPos{m}[6] \dicFlx{(‑s, ‑ar)}[8] \dicDirectTranslationCS{nájem\-ní smlouva}
\dicEntry[leigutaki] \dicTerm{leigu··tak|i} \dicIPA{{l}{ei}{\textlengthmark}{\textbabygamma}{\textscy}{t\smash{\textsuperscript{h}}}{a}{\r{\textObardotlessj}}{\textsci}} \dicPos{m}[1] \dicFlx{(‑a, ‑ar)}[8] \dicDirectTranslationCS{nájemce, nájemkyně}
\dicEntry[leik] \dicTerm{leik} \dicIPA{{l}{ei}{\textlengthmark}{\r{g}}} \dicPos{v} \dicFlx{ind praes sg 1 pers} \dicLink{leika}
\dicEntry[leika] \dicTerm{leika} \dicsymFrequent\  \dicIPA{{l}{ei}{\textlengthmark}{\r{g}}{a}} \dicPos{v}[6] \dicFlx{(leik, lék, lékum, léki, leikið)}[63] \dicFlx{acc} \textbf{1.} \dicDirectTranslationCS{(za)hrát} \dicIndirectTranslationCS{(bavit se pro zábavu)};  \dicPhraseIS{leika sér} \dicDirectTranslationCS{hrát si} \dicExampleIS{Börnin leika sér saman.} \dicExampleCS{Děti si spolu hrají.};  \dicPhraseIS{leika sér að e‑u} \dicDirectTranslationCS{hrát si s~(čím)} \dicExampleIS{leika sér að dúkkum} \dicExampleCS{hrát si s~panenkami}  \textbf{2.} \dicDirectTranslationCS{(za)hrát (si)} \dicIndirectTranslationCS{(provozovat jako sport)} \dicExampleIS{leika knattspyrnu} \dicExampleCS{hrát fotbal}  \textbf{3.} \dicDirectTranslationCS{(za)hrát} \dicIndirectTranslationCS{(provozovat hudbu)};  \dicPhraseIS{leika á hljóðfæri} \dicDirectTranslationCS{hrát na hudební nástroj} \dicExampleIS{Hún lék á fiðlu sína.} \dicExampleCS{Hrála na housle.}  \textbf{4.} \dicDirectTranslationCS{(za)hrát} \dicIndirectTranslationCS{(předvádět divadelní hru nebo osobu v~ní)} \dicExampleIS{leika Hamlet} \dicExampleCS{hrát Hamleta}  \textbf{5.} \dicPhraseIS{leika e‑n grátt\,/\addthin hart\,/\addthin illa} \dicDirectTranslationCS{zatočit s~(kým), nemazlit se s~(kým)};  \dicIdiom{leika}[af]{ \dicPhraseIS{leika af sér}} \dicDirectTranslationCS{udělat špatný tah (v~šachu ap.)};  \dicIdiom{leika}[á]{ \dicPhraseIS{leika á e‑n}} \dicDirectTranslationCS{vyzrát na (koho), přelstít (koho)}; { \dicPhraseIS{leika á als oddi, leika á alsoddi}} \dicLangCat{přen.} \dicDirectTranslationCS{mít výbornou náladu, být v~dobrém rozmaru}; { \dicPhraseIS{það leikur ekki á tveim tungum}} \dicFlx{impers} \dicLangCat{přen.} \dicDirectTranslationCS{o~tom není pochyb};  \dicIdiom{leika}[eftir]{ \dicPhraseIS{leika e‑ð eftir e‑m}} \dicDirectTranslationCS{zopakovat (co) po (kom)};  \dicIdiom{leika}[í]{ \dicPhraseIS{leika í e‑u}} \dicDirectTranslationCS{hrát v~(čem) (ve filmu ap.)};  \dicIdiom{leika}[undir]{ \dicPhraseIS{leika undir (e‑ð)}} \dicFieldCat{hud.} \dicDirectTranslationCS{doprovodit\,/\addthin doprovázet ((co))};  \dicIdiom{leika}[við]{ \dicPhraseIS{leika við e‑n}} \dicDirectTranslationCS{hrát si s~(kým)} \dicExampleIS{Lilja leikur við mömmu sína.} \dicExampleCS{Lilja si hraje s~maminkou.}; { \dicPhraseIS{leika við hvern sinn fingur}} \dicLangCat{přen.} \dicDirectTranslationCS{být dobře naladěn, mít dobrou náladu}
\dicEntry[leikandi] \dicTerm{leik··|andi\smash{\textsuperscript{1}}} \dicIPA{{l}{ei}{\textlengthmark}{\r{g}}{a}{n}{\textsubring{d}}{\textsci}} \dicPos{m}[2] \dicFlx{(‑anda, ‑endur)}[1] \textbf{1.} \dicDirectTranslationCS{hráč(ka) (ve fotbale ap.)}  \textbf{2.} \dicDirectTranslationCS{herec, herečka}
\dicEntry[leikandi] \dicTerm{leik··andi\smash{\textsuperscript{2}}} \dicIPA{{l}{ei}{\textlengthmark}{\r{g}}{a}{n}{\textsubring{d}}{\textsci}} \dicPos{adv} \dicDirectTranslationCS{hravě, s~lehkostí} \dicExampleIS{stökkva leikandi yfir skurðinn} \dicExampleCS{přeskočit hravě (přes) příkop}
\dicEntry[leikaraskapur] \dicTerm{leikara··skap|ur} \dicIPA{{l}{ei}{\textlengthmark}{\r{g}}{a}{r}{a}{s}{\r{g}}{a}{\textsubring{b}}{\textscy}{\textsubring{r}}} \dicPos{m}[10] \dicFlx{(‑ar)}[15] \textbf{1.} \dicSynonym{fíflalæti} \dicDirectTranslationCS{šaškoviny, komediantství}  \textbf{2.} \dicSynonym{látalæti} \dicDirectTranslationCS{předstírání, hraní}
\dicEntry[leikari] \dicTerm{leik··ar|i} \dicsymFrequent\  \dicIPA{{l}{ei}{\textlengthmark}{\r{g}}{a}{r}{\textsci}} \dicPos{m}[1] \dicFlx{(‑a, ‑ar)}[13] \dicDirectTranslationCS{herec, herečka} \dicExampleIS{Leikarinn hlaut verðlaun.} \dicExampleCS{Herec obdržel ocenění.}
\dicEntry[leikbrúða] \dicTerm{leik··brúð|a} \dicIPA{{l}{ei}{\textlengthmark}{\r{g}}{\textsubring{b}}{r}{u}{ð}{a}} \dicPos{f}[1] \dicFlx{(‑u, ‑ur)}[19] \textbf{1.} \dicSynonym{brúða} \dicDirectTranslationCS{loutka} \dicIndirectTranslationCS{(v~loutkovém divadle)}  \textbf{2.} \dicSynonym{handbendi} \dicDirectTranslationCS{loutka} \dicIndirectTranslationCS{(zmanipulovaný člověk)}
\dicEntry[leikdómur] \dicTerm{leik··dóm|ur} \dicIPA{{l}{ei}{\textlengthmark}{\r{g}}{\textsubring{d}}{ou}{m}{\textscy}{\textsubring{r}}} \dicPos{m}[6] \dicFlx{(‑s, ‑ar)}[10] \dicDirectTranslationCS{divadelní kritika}
\dicEntry[leikfang] \dicTerm{leik··|fang} \dicIPA{{l}{ei}{\textlengthmark}{\r{g}}{f}{au}{\ng}{\r{g}}} \dicPos{n}[2] \dicFlx{(‑fangs, ‑föng)}[8] \dicDirectTranslationCS{hračka (na hraní)}
\dicEntry[leikfangabúð] \dicTerm{leik·fanga··búð} \dicIPA{{l}{ei}{\textlengthmark}{\r{g}}{f}{au}{\ng}{\r{g}}{a}{\textsubring{b}}{u}{\texttheta}} \dicPos{f}[7] \dicFlx{(‑ar, ‑ir)}[1] \dicDirectTranslationCS{hračkářství}
\dicEntry[leikfélag] \dicTerm{leik··fé·|lag} \dicIPA{{l}{ei}{\textlengthmark}{\r{g}}{f}{j}{\textepsilon}{l}{a}{x}} \dicPos{n}[2] \dicFlx{(‑lags, ‑lög)}[8] \dicDirectTranslationCS{divadelní společnost}
\dicEntry[leikfélagi] \dicTerm{leik··fé·lag|i} \dicIPA{{l}{ei}{\textlengthmark}{\r{g}}{f}{j}{\textepsilon}{l}{a}{i}{j}{\textsci}} \dicPos{m}[1] \dicFlx{(‑a, ‑ar)}[8] \dicDirectTranslationCS{kamarád(ka) z~dětství}
\dicEntry[leikfimi] \dicTerm{leik··fim|i} \dicIPA{{l}{ei}{\textlengthmark}{\r{g}}{f}{\textsci}{m}{\textsci}} \dicPos{f}[3] \dicFlx{(‑i)}[3] \dicDirectTranslationCS{tělocvik, (tělesné) cvičení}
\dicEntry[leikfimisalur] \dicTerm{leik·fimi··sal|ur} \dicIPA{{l}{ei}{\textlengthmark}{\r{g}}{f}{\textsci}{m}{\textsci}{s}{a}{l}{\textscy}{\textsubring{r}}} \dicPos{m}[10] \dicFlx{(‑ar, ‑ir)}[14] \dicDirectTranslationCS{tělocvična}
\dicEntry[leikfimiskennari] \dicTerm{leik·fimis··kenn·ar|i} \dicIPA{{l}\-{ei}\-{\textlengthmark}\-{\r{g}}\-{f}\-{\textsci}\-{m}\-{\textsci}\-{s}\-{c\smash{\textsuperscript{h}}}\-{\textepsilon}\-{n}\-{a}\-{r}\-{\textsci}\-} \dicPos{m}[1] \dicFlx{(‑a, ‑ar)}[13] \dicDirectTranslationCS{učitel(ka) tělesné výchovy, tělocvikář(ka)}
\dicEntry[leikgagnrýnandi] \dicTerm{leik··gagn·rýn·|andi} \dicIPA{{l}{ei}{\r{g}}{\textlengthmark}{a}{\r{g}}{n}{r}{i}{n}{a}{n}{\textsubring{d}}{\textsci}} \dicPos{m}[2] \dicFlx{(‑anda, ‑endur)}[1] \dicDirectTranslationCS{divadelní kritik\,/\addthin kritička}
\dicEntry[leikgerð] \dicTerm{leik··gerð} \dicIPA{{l}{ei}{\textlengthmark}{\r{g}}{\r{\textObardotlessj}}{\textepsilon}{r}{\texttheta}} \dicPos{f}[7] \dicFlx{(‑ar, ‑ir)}[1] \dicFieldCat{lit.} \dicDirectTranslationCS{dramatizace} \dicExampleIS{leikgerð Sölku Völku} \dicExampleCS{dramatizace Salky Valky}
\dicEntry[leikgrind] \dicTerm{leik··grind} \dicIPA{{l}{ei}{\r{g}}{\textlengthmark}{r}{\textsci}{n}{\textsubring{d}}} \dicPos{f}[8] \dicFlx{(‑ar, ‑ur)}[1] \dicDirectTranslationCS{(dětská) ohrádka}
\dicEntry[leikhús] \dicTerm{leik··hús} \dicIPA{{l}{ei}{\textlengthmark}{\r{g}}{h}{u}{s}} \dicPos{n}[2] \dicFlx{(‑s, ‑)}[5] \dicDirectTranslationCS{divadlo} \dicExampleIS{sýning í leikhúsinu} \dicExampleCS{představení v~divadle}
\dicEntry[leikhússtjóri] \dicTerm{leik·hús··stjór|i} \dicIPA{{l}{ei}{\textlengthmark}{\r{g}}{h}{u}{s}{\textsubring{d}}{j}{ou}{r}{\textsci}} \dicPos{m}[1] \dicFlx{(‑a, ‑ar)}[1] \dicDirectTranslationCS{divadelní ředitel(ka), ředitel(ka) divadla}
\dicEntry[leikið] \dicTerm{leikið} \dicIPA{{l}{ei}{\textlengthmark}{\r{\textObardotlessj}}{\textsci}{\texttheta}} \dicPos{v} \dicFlx{supin} \dicLink{leika}
\dicEntry[leikinn] \dicTerm{leikinn} \dicIPA{{l}{ei}{\textlengthmark}{\r{\textObardotlessj}}{\textsci}{\textsubring{n}}} \dicPos{adj}[6]\dicFlx{}[-2] \dicSynonym{snjall} \dicDirectTranslationCS{obratný, šikovný, dovedný} \dicExampleIS{vera leikinn í e‑u} \dicExampleCS{být šikovný v~(čem)};  \dicPhraseIS{vera grátt\,/\addthin illa leikinn} \dicDirectTranslationCS{nést známky špatného zacházení}
\dicEntry[leikjanámskeið] \dicTerm{leikja··nám·skeið} \dicIPA{{l}{ei}{\textlengthmark}{\r{\textObardotlessj}}{a}{n}{au}{m}{s}{\r{\textObardotlessj}}{ei}{\texttheta}} \dicPos{n}[2] \dicFlx{(‑s, ‑)}[5] \dicDirectTranslationCS{tábor, (dětský) prázdninový kurz}
\dicEntry[leikjatölva] \dicTerm{leikja··tölv|a} \dicIPA{{l}{ei}{\textlengthmark}{\r{\textObardotlessj}}{a}{t\smash{\textsuperscript{h}}}{\oe}{l}{v}{a}} \dicPos{f}[1] \dicFlx{(‑u, ‑ur)}[7] \dicFieldCat{poč.} \dicDirectTranslationCS{(hrací) konzole}
\dicEntry[leikkerahjól] \dicTerm{leik·kera··hjól} \dicIPA{{l}{ei}{\textlengthmark}{\r{g}}{c\smash{\textsuperscript{h}}}{\textepsilon}{r}{a}{\c{c}}{ou}{\textsubring{l}}} \dicPos{n}[2] \dicFlx{(‑s, ‑)}[5] \dicDirectTranslationCS{hrnčířské kolo}
\dicEntry[leikkerasmiður] \dicTerm{leik·kera··smið|ur} \dicIPA{{l}{ei}{\textlengthmark}{\r{g}}{c\smash{\textsuperscript{h}}}{\textepsilon}{r}{a}{s}{m}{\textsci}{ð}{\textscy}{\textsubring{r}}} \dicPos{m}[9] \dicFlx{(‑s, ‑ir)}[11] \dicDirectTranslationCS{hrnčíř(ka), keramik, keramička}
\dicEntry[leikkona] \dicTerm{leik··kon|a} \dicIPA{{l}{ei}{\textlengthmark}{\r{g}}{k\smash{\textsuperscript{h}}}{\textopeno}{n}{a}} \dicPos{f}[1] \dicFlx{(‑u, ‑ur)}[27] \dicDirectTranslationCS{herečka}
\dicEntry[leiklist] \dicTerm{leik··list} \dicIPA{{l}{ei}{\textlengthmark}{\r{g}}{l}{\textsci}{s}{\textsubring{d}}} \dicPos{f}[7] \dicFlx{(‑ar)}[3] \dicDirectTranslationCS{divadlo, divadelní umění}
\dicEntry[leiklistarskóli] \dicTerm{leik·listar··skól|i} \dicIPA{{l}{ei}{\textlengthmark}{\r{g}}{l}{\textsci}{s}{\textsubring{d}}{a}{\textsubring{r}}{s}{\r{g}}{ou}{l}{\textsci}} \dicPos{m}[1] \dicFlx{(‑a, ‑ar)}[1] \dicFieldCat{škol.} \dicDirectTranslationCS{divadelní škola}
\dicEntry[leikmaður] \dicTerm{leik··|maður} \dicIPA{{l}{ei}{\textlengthmark}{\r{g}}{m}{a}{ð}{\textscy}{\textsubring{r}}} \dicPos{m}[13] \dicFlx{(‑manns, ‑menn)}[2] \textbf{1.} \dicDirectTranslationCS{laik, laička, neodborník, neodbornice}  \textbf{2.} \dicDirectTranslationCS{hráč(ka)}  \textbf{3.} \dicSynonym{leikari} \dicDirectTranslationCS{herec, herečka}
\dicEntry[leikmunur] \dicTerm{leik··mun|ur} \dicIPA{{l}{ei}{\textlengthmark}{\r{g}}{m}{\textscy}{n}{\textscy}{\textsubring{r}}} \dicPos{m}[10] \dicFlx{(‑ar, ‑ir)}[11] \dicDirectTranslationCS{rekvizita}
\dicEntry[leikmynd] \dicTerm{leik··mynd} \dicIPA{{l}{ei}{\textlengthmark}{\r{g}}{m}{\textsci}{n}{\textsubring{d}}} \dicPos{f}[7] \dicFlx{(‑ar, ‑ir)}[1] \dicDirectTranslationCS{scénografie, scéna}
\dicEntry[leikmyndahönnuður] \dicTerm{leik·mynda··hönn·uð|ur} \dicIPA{{l}\-{ei}\-{\textlengthmark}\-{\r{g}}\-{m}\-{\textsci}\-{n}\-{\textsubring{d}}\-{a}\-{h}\-{\oe}\-{n}\-{\textscy}\-{ð}\-{\textscy}\-{\textsubring{r}}\-} \dicPos{m}[10] \dicFlx{(‑ar, ‑ir)}[4] \dicDirectTranslationCS{jevištní výtvarník\,/\addthin výtvarnice, scénograf(ka)}
\dicEntry[leikni] \dicTerm{leikn|i} \dicIPA{{l}{ei}{h}{\r{g}}{n}{\textsci}} \dicPos{f}[3] \dicFlx{(‑i)}[3] \dicDirectTranslationCS{šikovnost, obratnost, zručnost, dovednost} \dicExampleIS{leikni í íþróttum} \dicExampleCS{dovednost ve sportu}
\dicEntry[leikregla] \dicTerm{leik··regl|a} \dicIPA{{l}{ei}{\textlengthmark}{\r{g}}{r}{\textepsilon}{\r{g}}{l}{a}} \dicPos{f}[1] \dicFlx{(‑u, ‑ur)}[19] \dicDirectTranslationCS{pravidlo hry}
\dicEntry[leikrit] \dicTerm{leik··rit} \dicIPA{{l}{ei}{\textlengthmark}{\r{g}}{r}{\textsci}{\textsubring{d}}} \dicPos{n}[2] \dicFlx{(‑s, ‑)}[5] \dicDirectTranslationCS{divadelní hra}
\dicEntry[leikritahöfundur] \dicTerm{leik·rita··höf·und|ur} \dicIPA{{l}\-{ei}\-{\textlengthmark}\-{\r{g}}\-{r}\-{\textsci}\-{\textsubring{d}}\-{a}\-{h}\-{\oe}\-{v}\-{\textscy}\-{n}\-{\textsubring{d}}\-{\textscy}\-{\textsubring{r}}\-} \dicPos{m}[6] \dicFlx{(‑ar, ‑ar)}[56] \dicDirectTranslationCS{(divadelní) dramatik, dramatička} \dicExampleIS{frægur tékkneskur leikritahöfundur} \dicExampleCS{slavný český dramatik}
\dicEntry[leikrænn] \dicTerm{leik··rænn} \dicIPA{{l}{ei}{\textlengthmark}{\r{g}}{r}{a}{i}{\textsubring{d}}{\textsubring{n}}} \dicPos{adj}[7]\dicFlx{}[-1] \dicDirectTranslationCS{divadelní}
\dicEntry[leikskáld] \dicTerm{leik··skáld} \dicIPA{{l}{ei}{\textlengthmark}{\r{g}}{s}{\r{g}}{au}{l}{\textsubring{d}}} \dicPos{n}[2] \dicFlx{(‑s, ‑)}[5] \dicDirectTranslationCS{dramatik, dramatička, autor(ka) divadelních her}
\dicEntry[leikskóladeild] \dicTerm{leik·skóla··deild} \dicIPA{{l}{ei}{\textlengthmark}{\r{g}}{s}{\r{g}}{ou}{l}{a}{\textsubring{d}}{ei}{l}{\textsubring{d}}} \dicPos{f}[7] \dicFlx{(‑ar, ‑ir)}[1] \dicDirectTranslationCS{oddělení mateřské školky} \dicIndirectTranslationCS{(v~rámci školy ap.)}
\dicEntry[leikskóladvöl] \dicTerm{leik·skóla··|dvöl} \dicIPA{{l}{ei}{\textlengthmark}{\r{g}}{s}{\r{g}}{ou}{l}{a}{\textsubring{d}}{v}{\oe}{\textsubring{l}}} \dicPos{f}[7] \dicFlx{(‑dvalar, ‑dvalir)}[16] \dicDirectTranslationCS{pobyt ve školce}
\dicEntry[leikskólakennari] \dicTerm{leik·skóla··kenn·ar|i} \dicIPA{{l}\-{ei}\-{\textlengthmark}\-{\r{g}}\-{s}\-{\r{g}}\-{ou}\-{l}\-{a}\-{c\smash{\textsuperscript{h}}}\-{\textepsilon}\-{n}\-{a}\-{r}\-{\textsci}\-} \dicPos{m}[1] \dicFlx{(‑a, ‑ar)}[13] \dicDirectTranslationCS{učitel(ka) v~mateřské školce}
\dicEntry[leikskóli] \dicTerm{leik··skól|i} \dicIPA{{l}{ei}{\textlengthmark}{\r{g}}{s}{\r{g}}{ou}{l}{\textsci}} \dicPos{m}[1] \dicFlx{(‑a, ‑ar)}[1] \dicDirectTranslationCS{(mateřská) školka}
\dicEntry[leikskrá] \dicTerm{leik··skrá} \dicIPA{{l}{ei}{\textlengthmark}{\r{g}}{s}{\r{g}}{r}{au}} \dicPos{f}[4] \dicFlx{(‑r\,/\addthin ‑ar, ‑r)}[21] \dicDirectTranslationCS{divadelní program}
\dicEntry[leikslok] \dicTerm{leiks··lok} \dicIPA{{l}{ei}{\textlengthmark}{\r{g}}{s}{l}{\textopeno}{\r{g}}} \dicPos{n}[2] \dicFlx{pl}[1] \textbf{1.} \dicSynonym{niðurstaða} \dicDirectTranslationCS{závěr, rozhodnutí}  \textbf{2.} \dicDirectTranslationCS{výsledek (zápasu, utkání ap.)}
\dicEntry[leiksoppur] \dicTerm{leik··sopp|ur} \dicIPA{{l}{ei}{\textlengthmark}{\r{g}}{s}{\textopeno}{h}{\textsubring{b}}{\textscy}{\textsubring{r}}} \dicPos{m}[6] \dicFlx{(‑s, ‑ar)}[22] \dicDirectTranslationCS{loutka} \dicIndirectTranslationCS{(člověk vnitřně prázdný)}
\dicEntry[leikstjóri] \dicTerm{leik··stjór|i} \dicIPA{{l}{ei}{\textlengthmark}{\r{g}}{s}{\textsubring{d}}{j}{ou}{r}{\textsci}} \dicPos{m}[1] \dicFlx{(‑a, ‑ar)}[1] \textbf{1.} \dicDirectTranslationCS{režisér(ka)}  \textbf{2.} \dicFieldCat{sport.} \dicDirectTranslationCS{rozhodčí}
\dicEntry[leikstjórn] \dicTerm{leik··stjórn} \dicIPA{{l}{ei}{\textlengthmark}{\r{g}}{s}{\textsubring{d}}{j}{ou}{r}{\textsubring{d}}{\textsubring{n}}} \dicPos{f}[7] \dicFlx{(‑ar)}[3] \dicDirectTranslationCS{režie, režírování}
\dicEntry[leikstýra] \dicTerm{leik··stýr|a} \dicIPA{{l}{ei}{\textlengthmark}{\r{g}}{s}{\textsubring{d}}{i}{r}{a}} \dicPos{v}[2] \dicFlx{(‑ði, ‑t)}[105] \dicFlx{dat} \dicDirectTranslationCS{režírovat}
\dicEntry[leiksvið] \dicTerm{leik··svið} \dicIPA{{l}{ei}{\textlengthmark}{\r{g}}{s}{v}{\textsci}{\texttheta}} \dicPos{n}[2] \dicFlx{(‑s, ‑)}[5] \dicDirectTranslationCS{jeviště, scéna}
\dicEntry[leiksýning] \dicTerm{leik··sýn·ing} \dicIPA{{l}{ei}{\textlengthmark}{\r{g}}{s}{i}{n}{i}{\ng}{\r{g}}} \dicPos{f}[4] \dicFlx{(‑ar, ‑ar)}[5] \dicDirectTranslationCS{(divadelní) představení}
\dicEntry[leiktíð] \dicTerm{leik··tíð} \dicIPA{{l}{ei}{\textlengthmark}{\r{g}}{t\smash{\textsuperscript{h}}}{i}{\texttheta}} \dicPos{f}[7] \dicFlx{(‑ar, ‑ir)}[1] \dicFieldCat{sport.} \dicDirectTranslationCS{sezóna (fotbalová ap.)}
\dicEntry[leikur] \dicTerm{leik|ur} \dicsymFrequent\  \dicIPA{{l}{ei}{\textlengthmark}{\r{g}}{\textscy}{\textsubring{r}}} \dicPos{m}[6] \dicFlx{(‑s, ‑ir\,/\addthin ‑ar)}[67] \textbf{1.} \dicFlx{pl (-ir)} \dicSynonym{barnaleikur} \dicDirectTranslationCS{hra, hraní} \dicExampleIS{leikur barna} \dicExampleCS{hra dětí};  \dicPhraseIS{bregða á leik} \dicDirectTranslationCS{hrát si, skotačit}  \textbf{2.} \dicFlx{pl (-ar)} \dicSynonym{viðureign} \dicDirectTranslationCS{boj, střetnutí, utkání}  \textbf{3.} \dicFlx{pl (-ar)} \dicSynonym*{íþróttakeppni} \dicDirectTranslationCS{zápas, utkání, hra} \dicExampleIS{Ólympíuleikar} \dicExampleCS{olympijské hry};  \dicPhraseIS{fara á leikinn} \dicDirectTranslationCS{jít na zápas}  \textbf{4.} \dicSynonym{leikrit} \dicDirectTranslationCS{hra (divadelní ap.), představení}  \textbf{5.} \dicDirectTranslationCS{tah} \dicIndirectTranslationCS{(posunutí šachové figury ap.)};  \dicPhraseIS{eiga leik} \dicDirectTranslationCS{být na tahu};  \dicIdiom{leikur}{ \dicPhraseIS{á nýjan leik}} \dicFlx{adv} \dicDirectTranslationCS{znovu, opět}; { \dicPhraseIS{e‑að er leikur einn}} \dicLangCat{přen.} \dicDirectTranslationCS{(co) je hračka}; { \dicPhraseIS{með leik}} \dicFlx{adv} \dicDirectTranslationCS{snadno, lehce}; { \dicPhraseIS{vera úr leik}} \dicLangCat{přen.} \dicDirectTranslationCS{být mimo hru}; { \dicPhraseIS{við illan leik}} \dicFlx{adv} \dicDirectTranslationCS{horko těžko}
\dicEntry[leikvangur] \dicTerm{leik··vang|ur} \dicIPA{{l}{ei}{\textlengthmark}{\r{g}}{v}{au}{\ng}{\r{g}}{\textscy}{\textsubring{r}}} \dicPos{m}[6] \dicFlx{(‑s, ‑ar)}[25] \dicSynonym{íþróttavöllur} \dicDirectTranslationCS{stadión}
\dicEntry[leikvöllur] \dicTerm{leik··|völlur} \dicIPA{{l}{ei}{\textlengthmark}{\r{g}}{v}{\oe}{\textsubring{d}}{l}{\textscy}{\textsubring{r}}} \dicPos{m}[11] \dicFlx{(‑vallar, ‑vellir)}[5] \textbf{1.} \dicFieldCat{sport.} \dicSynonym{í\-þrótta\-völlur} \dicDirectTranslationCS{(sportovní) hřiště}  \textbf{2.} \dicSynonym*{barnaleikvöllur} \dicDirectTranslationCS{dětské hřiště}
\dicEntry[leir] \dicTerm{leir} \dicsymFrequent\  \dicIPA{{l}{ei}{\textlengthmark}{\textsubring{r}}} \dicPos{m}[4] \dicFlx{(‑s)}[18] \dicDirectTranslationCS{jíl, hlína} \dicExampleIS{hvít englabörn úr leir} \dicExampleCS{bílí andělíčci z~hlíny}
\dicEntry[leira] \dicTerm{leir|a\smash{\textsuperscript{1}}} \dicIPA{{l}{ei}{\textlengthmark}{r}{a}} \dicPos{f}[1] \dicFlx{(‑u, ‑ur)}[19] \dicDirectTranslationCS{jílovitá půda}
\dicEntry[Leira] \dicTerm{Leir|a\smash{\textsuperscript{2}}} \dicIPA{{l}{ei}{\textlengthmark}{r}{a}} \dicPos{f}[1] \dicFlx{(‑u)}[6] \dicFieldCat{geog.} \dicDirectTranslationCS{Loira} \dicIndirectTranslationCS{(řeka ve Francii)}
\dicEntry[leirburður] \dicTerm{leir··burð|ur} \dicIPA{{l}{ei}{r}{\textsubring{b}}{\textscy}{r}{ð}{\textscy}{\textsubring{r}}} \dicPos{m}[10] \dicFlx{(‑ar)}[7] \dicDirectTranslationCS{špatná poezie}
\dicEntry[leirker] \dicTerm{leir··ker} \dicIPA{{l}{ei}{\textsubring{r}}{c\smash{\textsuperscript{h}}}{\textepsilon}{\textsubring{r}}} \dicPos{n}[2] \dicFlx{(‑s, ‑)}[6] \dicDirectTranslationCS{keramika} \dicIndirectTranslationCS{(především z~archeologických nálezů)}
\dicEntry[leirlist] \dicTerm{leir··list} \dicIPA{{l}{ei}{r}{l}{\textsci}{s}{\textsubring{d}}} \dicPos{f}[7] \dicFlx{(‑ar)}[3] \dicDirectTranslationCS{keramika, keramické umění}
\dicEntry[leirmunur] \dicTerm{leir··mun|ur} \dicIPA{{l}{ei}{r}{m}{\textscy}{n}{\textscy}{\textsubring{r}}} \dicPos{m}[10] \dicFlx{(‑ar, ‑ir)}[11] \dicDirectTranslationCS{keramika, keramický výrobek}
\dicEntry[leit] \dicTerm{leit\smash{\textsuperscript{1}}} \dicsymFrequent\  \dicIPA{{l}{ei}{\textlengthmark}{\textsubring{d}}} \dicPos{f}[7] \dicFlx{(‑ar, ‑ir)}[1] \textbf{1.} \dicDirectTranslationCS{hledání, pátrání, vyhledávání} \dicExampleIS{Leitin er mjög umfangsmikil.} \dicExampleCS{Pátrání je velmi rozsáhlé.}  \textbf{2.} \dicPhraseIS{leitir} \dicFlx{pl} \dicSynonym{smölun} \dicDirectTranslationCS{shánění ovcí};  \dicPhraseIS{fara í leitir} \dicDirectTranslationCS{jet shánět ovce}
\dicEntry[leit] \dicTerm{leit\smash{\textsuperscript{2}}} \dicIPA{{l}{ei}{\textlengthmark}{\textsubring{d}}} \dicPos{v} \dicFlx{ind pf sg 1 pers} \dicLink{líta}
\dicEntry[leita] \dicTerm{leit|a} \dicsymFrequent\  \dicIPA{{l}{ei}{\textlengthmark}{\textsubring{d}}{a}} \dicPos{v}[1] \dicFlx{(‑aði)}[1] \dicFlx{gen} \textbf{1.} \dicDirectTranslationCS{(vy)hledat, (vy)pátrat} \dicIndirectTranslationCS{(snažit se nalézt)};  \dicPhraseIS{leita e‑s} \dicDirectTranslationCS{hledat (co)} \dicExampleIS{leita bátsins} \dicExampleCS{hledat člun};  \dicPhraseIS{leita að e‑u\,/\addthin e‑m} \dicDirectTranslationCS{hledat (co\,/\addthin koho)} \dicExampleIS{leita að lyklunum} \dicExampleCS{hledat klíče}  \textbf{2.} \dicDirectTranslationCS{vyhledat, navštívit (lékaře ap.)};  \dicPhraseIS{leita læknis} \dicDirectTranslationCS{vyhledat doktora};  \dicIdiom{leita}[á]{ \dicPhraseIS{leita á e‑m}} \dicDirectTranslationCS{prohledat (koho), prošacovat (koho)};  \dicIdiom{leita}[eftir]{ \dicPhraseIS{leita eftir e‑u við e‑n}} \dicDirectTranslationCS{požádat (koho) o~(co)};  \dicIdiom{leita}[fyrir]{ \dicPhraseIS{leita fyrir sér um e‑ð}} \dicDirectTranslationCS{shánět se po (čem), shánět (co) (práci ap.)};  \dicIdiom{leita}[til]{ \dicPhraseIS{leita til e‑rs}} \dicDirectTranslationCS{obrátit se na (koho)} \dicExampleIS{leita til hennar með vandamál sín} \dicExampleCS{obrátit se na ni se svými problémy};  \dicIdiom{leita}[uppi]{ \dicPhraseIS{leita e‑n uppi}} \dicDirectTranslationCS{vyhledat (koho), vypátrat (koho)};  \dicIdiom{leitast}[við]{ \dicPhraseIS{leitast við}} \dicFlx{refl} \dicDirectTranslationCS{pokoušet se, snažit se}
\dicEntry[leitandi] \dicTerm{leit··andi} \dicIPA{{l}{ei}{\textlengthmark}{\textsubring{d}}{a}{n}{\textsubring{d}}{\textsci}} \dicPos{adj}[13] \dicFlx{indecl}[1] \dicDirectTranslationCS{hledající, pátrající}
\dicEntry[leitarflokkur] \dicTerm{leitar··flokk|ur} \dicIPA{{l}{ei}{\textlengthmark}{\textsubring{d}}{a}{\textsubring{r}}{f}{l}{\textopeno}{h}{\r{g}}{\textscy}{\textsubring{r}}} \dicPos{m}[6] \dicFlx{(‑s, ‑ar)}[8] \dicDirectTranslationCS{pátrací skupina\,/\addthin četa}
\dicEntry[leitarvél] \dicTerm{leitar··vél} \dicIPA{{l}{ei}{\textlengthmark}{\textsubring{d}}{a}{r}{v}{j}{\textepsilon}{\textsubring{l}}} \dicPos{f}[4] \dicFlx{(‑ar, ‑ar)}[1] \dicFieldCat{poč.} \dicDirectTranslationCS{(internetový) vyhledávač}
\dicEntry[leiti] \dicTerm{leiti} \dicIPA{{l}{ei}{\textlengthmark}{\textsubring{d}}{\textsci}} \dicPos{n}[2] \dicFlx{(‑s, ‑)}[14] \dicSynonym{hæð} \dicDirectTranslationCS{vrch, vyvýšenina, kopec};  \dicPhraseIS{á næsta leiti} \dicFlx{adv} \dicDirectTranslationCS{blízko, nedaleko (časově), za dveřmi\,/\addthin rohem} \dicExampleIS{Hátíðin er á næsta leiti.} \dicExampleCS{Svátky jsou za dveřmi.}
\dicEntry[leitun] \dicTerm{leit|un} \dicIPA{{l}{ei}{\textlengthmark}{\textsubring{d}}{\textscy}{\textsubring{n}}} \dicPos{f}[7] \dicFlx{(‑unar)}[9] \dicDirectTranslationCS{hledání, vyhledávání, pátrání};  \dicPhraseIS{leitun að e‑u} \dicDirectTranslationCS{hledání (čeho)};  \dicPhraseIS{það er leitun á\,/\addthin að e‑u} \dicDirectTranslationCS{je vzácné\,/\addthin těžké nalézt (co)}
\dicEntry[lek] \dicTerm{lek} \dicIPA{{l}{\textepsilon}{\textlengthmark}{\r{g}}} \dicPos{v} \dicFlx{ind praes sg 1 pers} \dicLink{leka}
\dicEntry[leka] \dicTerm{leka} \dicsymFrequent\  \dicIPA{{l}{\textepsilon}{\textlengthmark}{\r{g}}{a}} \dicPos{v}[6] \dicFlx{(lek, lak, lákum, læki, lekið)}[17] \textbf{1.} \dicSynonym*{falla í dropum} \dicDirectTranslationCS{téct, stékat, kapat} \dicExampleIS{Vatnið lekur niður á gólfið.} \dicExampleCS{Voda stéká dolů na podlahu.}  \textbf{2.} \dicSynonym*{vera lekur} \dicDirectTranslationCS{téct, prosakovat} \dicIndirectTranslationCS{(propouštět kapalinu)}  \textbf{3.} \dicLangCat{přen.} \dicDirectTranslationCS{uniknout (na veřejnost), prosáknout (na veřejnost) (zpráva ap.)}
\dicEntry[lekandi] \dicTerm{lek··|andi} \dicIPA{{l}{\textepsilon}{\textlengthmark}{\r{g}}{a}{n}{\textsubring{d}}{\textsci}} \dicPos{m}[2] \dicFlx{(‑anda)}[2] \dicFieldCat{med.} \dicDirectTranslationCS{kapavka}
\dicEntry[leki] \dicTerm{lek|i} \dicIPA{{l}{\textepsilon}{\textlengthmark}{\r{\textObardotlessj}}{\textsci}} \dicPos{m}[1] \dicFlx{(‑a, ‑ar)}[1] \textbf{1.} \dicDirectTranslationCS{kapání, odkapávání} \dicExampleIS{leki úr þaki} \dicExampleCS{kapání ze střechy}  \textbf{2.} \dicDirectTranslationCS{prosakování, zatékání} \dicExampleIS{Það kom leki að bátnum.} \dicExampleCS{Loď začala zatékat.}  \textbf{3.} \dicDirectTranslationCS{únik, unikání (informací ap.)}
\dicEntry[lekið] \dicTerm{lekið} \dicIPA{{l}{\textepsilon}{\textlengthmark}{\r{\textObardotlessj}}{\textsci}{\texttheta}} \dicPos{v} \dicFlx{supin} \dicLink{leka}
\dicEntry[lektor] \dicTerm{lektor} \dicIPA{{l}{\textepsilon}{x}{\textsubring{d}}{\textopeno}{\textsubring{r}}} \dicPos{m}[4] \dicFlx{(‑s, ‑ar)}[14] \dicFieldCat{škol.} \dicDirectTranslationCS{odborný asistent, odborná asistentka}
\dicEntry[lekur] \dicTerm{lekur} \dicIPA{{l}{\textepsilon}{\textlengthmark}{\r{g}}{\textscy}{\textsubring{r}}} \dicPos{adj}[1]\dicFlx{}[-1] \dicDirectTranslationCS{zatékající, prosakující} \dicExampleIS{lekur bátur} \dicExampleCS{zatékající člun}
\dicEntry[lem] \dicTerm{lem} \dicIPA{{l}{\textepsilon}{\textlengthmark}{\textsubring{m}}} \dicPos{v} \dicFlx{ind praes sg 1 pers} \dicLink{lemja}
\dicEntry[lemdi] \dicTerm{lemdi} \dicIPA{{l}{\textepsilon}{m}{\textsubring{d}}{\textsci}} \dicPos{v} \dicFlx{con pf sg 1 pers} \dicLink{lemja}
\dicEntry[lemja] \dicTerm{lemja} \dicsymFrequent\  \dicIPA{{l}{\textepsilon}{m}{j}{a}} \dicPos{v}[4] \dicFlx{(lem, lamdi, lömdum, lemdi, lamið)}[3] \dicFlx{acc} \dicSynonym{slá\smash{\textsuperscript{3}}} \dicDirectTranslationCS{uhodit, udeřit, tlouct, mlátit} \dicExampleIS{lemja e‑n með barefli} \dicExampleCS{uhodit (koho) pálkou}
\dicEntry[lemstra] \dicTerm{lemstr|a} \dicIPA{{l}{\textepsilon}{m}{s}{\textsubring{d}}{r}{a}} \dicPos{v}[1] \dicFlx{(‑aði)}[1] \dicFlx{acc} \dicSynonym{limlesta} \dicDirectTranslationCS{zmrzačit}
\dicEntry[lend] \dicTerm{lend} \dicIPA{{l}{\textepsilon}{n}{\textsubring{d}}} \dicPos{f}[4] \dicFlx{(‑ar, ‑ar)}[1] \textbf{1.} \dicDirectTranslationCS{zadek (koně ap.)}  \textbf{2.} \dicPhraseIS{lendar} \dicFlx{pl} \dicSynonym{mjaðmir} \dicDirectTranslationCS{boky, bedra}
\dicEntry[lenda] \dicTerm{len|da} \dicsymFrequent\  \dicIPA{{l}{\textepsilon}{n}{\textsubring{d}}{a}} \dicPos{v}[2] \dicFlx{(‑ti, ‑t)}[41] \dicFlx{dat} \textbf{1.} \dicSynonym*{koma að landi} \dicDirectTranslationCS{přistát, přistávat, dosednout, dosedat} \dicExampleIS{lenda flugvélinni} \dicExampleCS{přistát s~letadlem}  \textbf{2.} \dicDirectTranslationCS{přistát, dopadnout (na zadek ap.)};  \dicIdiom{lenda}[í]{ \dicPhraseIS{lenda í e‑u}} \dicSynonym*{rata í e‑ð} \dicDirectTranslationCS{spadnout do (čeho), dostat se do (čeho)} \dicExampleIS{lenda í vandræðum} \dicExampleCS{dostat se do potíží};  \dicIdiom{lenda}[saman]{ \dicPhraseIS{e‑jum lendir saman}} \dicFlx{impers} \dicLangCat{přen.} \dicDirectTranslationCS{(kdo) se dostává do křížku} \dicExampleIS{Þeim lenti saman.} \dicExampleCS{Dostali se do křížku.}
\dicEntry[lending] \dicTerm{lend··ing} \dicIPA{{l}{\textepsilon}{n}{\textsubring{d}}{i}{\ng}{\r{g}}} \dicPos{f}[4] \dicFlx{(‑ar, ‑ar)}[5] \dicDirectTranslationCS{přistání, přistávání, dosednutí} \dicExampleIS{mjúk lending} \dicExampleCS{měkké přistání}
\dicEntry[lengd] \dicTerm{lengd} \dicsymFrequent\  \dicIPA{{l}{\textepsilon}{\ng}{\textsubring{d}}} \dicPos{f}[7] \dicFlx{(‑ar, ‑ir)}[1] \textbf{1.} \dicDirectTranslationCS{délka} \dicExampleIS{í fullri lengd} \dicExampleCS{v~plné délce}  \textbf{2.} \dicFieldCat{geog.} \dicDirectTranslationCS{zeměpisná délka};  \dicIdiom{lengd}{ \dicPhraseIS{til lengdar}} \dicFlx{adv} \dicDirectTranslationCS{dlouhodobě, v~dlouhém časovém období}
\dicEntry[lengdarbaugur] \dicTerm{lengdar··baug|ur} \dicIPA{{l}{\textepsilon}{\ng}{\textsubring{d}}{a}{r}{\textsubring{b}}{\oe i}{\textbabygamma}{\textscy}{\textsubring{r}}} \dicPos{m}[6] \dicFlx{(‑s, ‑ar)}[24] \dicFieldCat{geog.} \dicDirectTranslationCS{poledník}
\dicEntry[lengi] \dicTerm{lengi} \dicsymFrequent\  \dicIPA{{l}{ei}{\textltailn}{\r{\textObardotlessj}}{\textsci}} \dicPos{adv} \dicFlx{(comp lengur, sup lengst)} \dicSynonym*{langan tíma} \dicDirectTranslationCS{dlouho, dlouze} \dicExampleIS{lengi vetrar} \dicExampleCS{po většinu zimy};  \dicPhraseIS{lengi vel} \dicFlx{adv} \dicDirectTranslationCS{dlouho, dlouhou dobu}
\dicEntry[lenging] \dicTerm{leng··ing} \dicIPA{{l}{ei}{\textltailn}{\r{\textObardotlessj}}{i}{\ng}{\r{g}}} \dicPos{f}[4] \dicFlx{(‑ar, ‑ar)}[5] \textbf{1.} \dicSynonym*{það að lengja e‑ð} \dicDirectTranslationCS{prodloužení, prodlužování} \dicExampleIS{lenging á jólafríinu} \dicExampleCS{prodloužení vánočních prázdnin}  \textbf{2.} \dicSynonym{viðbót} \dicDirectTranslationCS{vylepšení, vylepšování} \dicExampleIS{lenging á bátnum} \dicExampleCS{vylepšení člunu}
\dicEntry[lengja] \dicTerm{lengj|a\smash{\textsuperscript{1}}} \dicIPA{{l}{ei}{\textltailn}{\r{\textObardotlessj}}{a}} \dicPos{f}[1] \dicFlx{(‑u, ‑ur)}[7] \textbf{1.} \dicSynonym{ræma} \dicDirectTranslationCS{pruh, proužek}  \textbf{2.} \dicDirectTranslationCS{karton (cigaret)}
\dicEntry[lengja] \dicTerm{leng|ja\smash{\textsuperscript{2}}} \dicIPA{{l}{ei}{\textltailn}{\r{\textObardotlessj}}{a}} \dicPos{v}[2] \dicFlx{(‑di, ‑t)}[130] \dicFlx{acc} \textbf{1.} \dicDirectTranslationCS{prodloužit, prodlužovat} \dicExampleIS{lengja vinnutímann} \dicExampleCS{prodloužit pracovní dobu}  \textbf{2.} \dicPhraseIS{e‑ð lengir} \dicFlx{impers} \dicDirectTranslationCS{(co) se prodlužuje} \dicExampleIS{Daginn lengir.} \dicExampleCS{Den se prodlužuje.};  \dicIdiom{lengja}[eftir]{ \dicPhraseIS{e‑n lengir eftir e‑m\,/\addthin e‑u}} \dicFlx{impers} \dicDirectTranslationCS{(kdo) touží po (kom\,/\addthin čem)} \dicExampleIS{Mig lengir eftir henni.} \dicExampleCS{Toužím po ní.};  \dicIdiom{lengjast}{ \dicPhraseIS{lengjast}} \dicFlx{refl} \dicDirectTranslationCS{stávat se delší, prodlužovat se} \dicExampleIS{Dagurinn lengist.} \dicExampleCS{Den se prodlužuje.}
\dicEntry[lengra] \dicTerm{lengra} \dicIPA{{l}{\textepsilon}{\ng}{\r{g}}{r}{a}} \dicPos{adv} \dicFlx{comp} \dicLink{langt}
\dicEntry[lengri] \dicTerm{lengri} \dicIPA{{l}{\textepsilon}{\ng}{\r{g}}{r}{\textsci}} \dicPos{adj} \dicFlx{comp m} \dicLink{langur\smash{\textsuperscript{2}}}
\dicEntry[lengst] \dicTerm{lengst} \dicIPA{{l}{\textepsilon}{\ng}{s}{\textsubring{d}}} \dicPos{adv} \dicFlx{sup} \textbf{1.} \dicLink{langt}  \textbf{2.} \dicLink{lengi}
\dicEntry[lengstur] \dicTerm{lengstur} \dicIPA{{l}{\textepsilon}{\ng}{s}{\textsubring{d}}{\textscy}{\textsubring{r}}} \dicPos{adj} \dicFlx{m sg nom sup} \dicLink{langur\smash{\textsuperscript{2}}}
\dicEntry[lengur] \dicTerm{lengur} \dicIPA{{l}{ei}{\ng}{\r{g}}{\textscy}{\textsubring{r}}} \dicPos{adv} \dicFlx{comp} \dicLink{lengi}
\dicEntry[lens] \dicTerm{lens} \dicIPA{{l}{\textepsilon}{n}{s}} \dicPos{n}[2] \dicFlx{(‑)}[24] \dicFieldCat{nám.} \dicSynonym{byr} \dicDirectTranslationCS{zadní vítr}
\dicEntry[lenska] \dicTerm{lensk|a} \dicIPA{{l}{\textepsilon}{n}{s}{\r{g}}{a}} \dicPos{f}[1] \dicFlx{(‑u)}[5] \dicSynonym*{landsvenja} \dicDirectTranslationCS{zvyk, obyčej, zvyklost}
\dicEntry[lenskur] \dicTerm{lenskur} \dicIPA{{l}{\textepsilon}{n}{s}{\r{g}}{\textscy}{\textsubring{r}}} \dicPos{adj}[1]\dicFlx{}[-1] \dicPhraseIS{hvers lenskur} \dicDirectTranslationCS{jaké národnosti} \dicExampleIS{Ég vissi ekki hvers lenskur hann var.} \dicExampleCS{Nevěděl jsem, jaké je národnosti.}
\dicEntry[lep] \dicTerm{lep} \dicIPA{{l}{\textepsilon}{\textlengthmark}{\textsubring{b}}} \dicPos{v} \dicFlx{ind praes sg 1 pers} \dicLink{lepja}
\dicEntry[lepja] \dicTerm{lepja} \dicIPA{{l}{\textepsilon}{\textlengthmark}{\textsubring{b}}{j}{a}} \dicPos{v}[4] \dicFlx{(lep, lapti, löptum, lepti, lapið)}[47] \dicFlx{acc} \textbf{1.} \dicDirectTranslationCS{(s)lízat, vylízat} \dicIndirectTranslationCS{(o~zvířatech)} \dicExampleIS{lepja úr skálinni} \dicExampleCS{lízat z~misky}  \textbf{2.} \dicDirectTranslationCS{papouškovat, opakovat (bez porozumění ap.)}
\dicEntry[leppa] \dicTerm{lepp|a} \dicIPA{{l}{\textepsilon}{h}{\textsubring{b}}{a}} \dicPos{v}[1] \dicFlx{(‑aði)}[1] \dicFlx{acc} \textbf{1.} \dicPhraseIS{leppa sig} \dicSynonym*{klæða sig} \dicDirectTranslationCS{obléct se, oblékat se}  \textbf{2.} \dicSynonym{bæta} \dicDirectTranslationCS{opravit, vyspravit} \dicExampleIS{leppa upp á e‑ð} \dicExampleCS{vyspravit (co)}  \textbf{3.} \dicDirectTranslationCS{vázat} \dicIndirectTranslationCS{(v~šachách)}
\dicEntry[Leppalúði] \dicTerm{Leppa··lúð|i} \dicIPA{{l}{\textepsilon}{h}{\textsubring{b}}{a}{l}{u}{ð}{\textsci}} \dicPos{m}[1] \dicFlx{(‑a)}[5] \dicFlx{prop} \dicFieldCat{pov.} \dicDirectTranslationCS{Leppalúði} \dicIndirectTranslationCS{(nesmírně líný obr, otec 13 vánočních skřítků)}
\dicEntry[leppríki] \dicTerm{lepp··ríki} \dicIPA{{l}{\textepsilon}{h}{\textsubring{b}}{r}{i}{\r{\textObardotlessj}}{\textsci}} \dicPos{n}[2] \dicFlx{(‑s, ‑)}[16] \dicFieldCat{pol.} \dicDirectTranslationCS{loutkový\,/\addthin satelitní stát}
\dicEntry[leppstjórn] \dicTerm{lepp··stjórn} \dicIPA{{l}{\textepsilon}{h}{\textsubring{b}}{s}{\textsubring{d}}{j}{ou}{r}{\textsubring{d}}{\textsubring{n}}} \dicPos{f}[7] \dicFlx{(‑ar, ‑ir)}[1] \dicDirectTranslationCS{loutková vláda}
\dicEntry[leppur] \dicTerm{lepp|ur} \dicIPA{{l}{\textepsilon}{h}{\textsubring{b}}{\textscy}{\textsubring{r}}} \dicPos{m}[6] \dicFlx{(‑s, ‑ar)}[22] \textbf{1.} \dicSynonym{bót} \dicDirectTranslationCS{záplata}  \textbf{2.} \dicDirectTranslationCS{vložka do bot}  \textbf{3.} \dicPhraseIS{leppar} \dicFlx{pl} \dicSynonym*{tötrar} \dicDirectTranslationCS{hadry} \dicIndirectTranslationCS{(o~šatech)}  \textbf{4.} \dicSynonym{handbendi} \dicDirectTranslationCS{loutka, nohsled}  \textbf{5.} \dicDirectTranslationCS{klapka (na oku ap.)}  \textbf{6.} \dicDirectTranslationCS{vazba} \dicIndirectTranslationCS{(v~šachách)}
\dicEntry[lepti] \dicTerm{lepti} \dicIPA{{l}{\textepsilon}{f}{\textsubring{d}}{\textsci}} \dicPos{v} \dicFlx{con pf sg 1 pers} \dicLink{lepja}
\dicEntry[lerki] \dicTerm{lerki} \dicIPA{{l}{\textepsilon}{\textsubring{r}}{\r{\textObardotlessj}}{\textsci}} \dicPos{n}[2] \dicFlx{(‑s)}[20] \dicFieldCat{bot.} \dicDirectTranslationCS{modřín} \textit{(l.~{\textLA{Larix}})}  \dicsymPhoto\ 
\dicFigure{286.jpg}{Lerki}{Lerki - Zicha Ondřej, Biolib, Copyright/CC-BY-NC}
\dicEntry[les] \dicTerm{les} \dicIPA{{l}{\textepsilon}{\textlengthmark}{s}} \dicPos{v} \dicFlx{ind praes sg 1 pers} \dicLink{lesa}
\dicEntry[lesa] \dicTerm{lesa} \dicsymFrequent\  \dicIPA{{l}{\textepsilon}{\textlengthmark}{s}{a}} \dicPos{v}[6] \dicFlx{(les, las, lásum, læsi, lesið)}[16] \dicFlx{acc} \textbf{1.} \dicDirectTranslationCS{(pře)číst} \dicExampleIS{lesa söguna upphátt} \dicExampleCS{číst příběh nahlas}  \textbf{2.} \dicSynonym{læra} \dicDirectTranslationCS{studovat, učit se} \dicExampleIS{Hann les sagnfræði í háskóla.} \dicExampleCS{Studuje historii na vysoké škole.}  \textbf{3.} \dicLangCat{zast.} \dicDirectTranslationCS{sbírat (květy ap.)} \dicExampleIS{lesa aldin af trjánum} \dicExampleCS{sbírat plody ze stromů};  \dicIdiom{lesa}[fyrir]{ \dicPhraseIS{lesa fyrir}} \dicDirectTranslationCS{diktovat}; { \dicPhraseIS{lesa fyrir e‑m e‑ð}} \dicDirectTranslationCS{(na)diktovat (komu co)} \dicExampleIS{lesa henni bréfið fyrir} \dicExampleCS{diktovat jí dopis};  \dicIdiom{lesa}[í]{ \dicPhraseIS{lesa í lófa}} \dicLangCat{přen.} \dicDirectTranslationCS{číst z~dlaně};  \dicIdiom{lesa}[saman]{ \dicPhraseIS{lesa saman e‑ð og e‑ð}} \dicDirectTranslationCS{srovnávat\,/\addthin porovnávat (co) s~(čím)} \dicExampleIS{lesa frumritið og afritið saman} \dicExampleCS{srovnávat původní text s~opisem};  \dicIdiom{lesa}[til]{ \dicPhraseIS{lesa sér til e‑ð}} \dicDirectTranslationCS{seznámit se s~(čím), dozvědět se (co) nového (čtením)};  \dicIdiom{lesa}[undir]{ \dicPhraseIS{lesa undir próf}} \dicDirectTranslationCS{učit se na zkoušku, studovat na zkoušku} \dicExampleIS{lesa undir skóla} \dicExampleCS{učit se do školy};  \dicIdiom{lesa}[upp]{ \dicPhraseIS{lesa upp e‑ð}} \dicDirectTranslationCS{recitovat (co), přednášet (co)} \dicExampleIS{lesa upp ljóð} \dicExampleCS{recitovat báseň}; { \dicPhraseIS{lesa sig upp}} \dicSynonym{klifra} \dicDirectTranslationCS{šplhat, vyšplhat (se)};  \dicIdiom{lesa}[yfir]{ \dicPhraseIS{lesa e‑ð yfir}} \dicDirectTranslationCS{pročíst si (co)}
\dicEntry[lesandi] \dicTerm{les··|andi} \dicsymFrequent\  \dicIPA{{l}{\textepsilon}{\textlengthmark}{s}{a}{n}{\textsubring{d}}{\textsci}} \dicPos{m}[2] \dicFlx{(‑anda, ‑endur)}[1] \dicDirectTranslationCS{čtenář(ka)} \dicExampleIS{Lesendurnir eru ánægðir með bókina.} \dicExampleCS{Čtenářům se knížka líbí.}
\dicEntry[lesbía] \dicTerm{lesbí|a} \dicIPA{{l}{\textepsilon}{s}{\textsubring{b}}{i}{j}{a}} \dicPos{f}[1] \dicFlx{(‑u, ‑ur)}[7] \dicDirectTranslationCS{lesba, lesbička}
\dicEntry[lesblinda] \dicTerm{les··blind|a} \dicIPA{{l}{\textepsilon}{\textlengthmark}{s}{\textsubring{b}}{l}{\textsci}{n}{\textsubring{d}}{a}} \dicPos{f}[1] \dicFlx{(‑u)}[5] \dicFieldCat{med.} \dicDirectTranslationCS{dyslexie}
\dicEntry[lesbók] \dicTerm{les··|bók} \dicIPA{{l}{\textepsilon}{\textlengthmark}{s}{\textsubring{b}}{ou}{\r{g}}} \dicPos{f}[8] \dicFlx{(‑bókar, ‑bækur)}[5] \dicDirectTranslationCS{čítanka}
\dicEntry[lesefni] \dicTerm{les··efni} \dicIPA{{l}{\textepsilon}{\textlengthmark}{s}{\textepsilon}{\textsubring{b}}{n}{\textsci}} \dicPos{n}[2] \dicFlx{(‑s, ‑)}[14] \dicDirectTranslationCS{čtení, četba, čtivo}
\dicEntry[lesgleraugu] \dicTerm{les··gler·augu} \dicIPA{{l}{\textepsilon}{\textlengthmark}{s}{\r{g}}{l}{\textepsilon}{r}{\oe i}{\textbabygamma}{\textscy}} \dicPos{n}[1] \dicFlx{pl}[3] \dicDirectTranslationCS{brýle na čtení}
\dicEntry[lesið] \dicTerm{lesið} \dicIPA{{l}{\textepsilon}{\textlengthmark}{s}{\textsci}{\texttheta}} \dicPos{v} \dicFlx{supin} \dicLink{lesa}
\dicEntry[lesinn] \dicTerm{lesinn} \dicIPA{{l}{\textepsilon}{\textlengthmark}{s}{\textsci}{\textsubring{n}}} \dicPos{adj}[6]\dicFlx{}[-2] \dicDirectTranslationCS{sečtělý}
\dicEntry[lesminni] \dicTerm{les··minni} \dicIPA{{l}{\textepsilon}{\textlengthmark}{s}{m}{\textsci}{n}{\textsci}} \dicPos{n}[2] \dicFlx{(‑s, ‑)}[14] \dicFieldCat{poč.} \dicDirectTranslationCS{ROM} \dicIndirectTranslationCS{(typ elektronické paměti používaný pro uložení firmware v~elektronických přístrojích)} \dicAntonym{vinnsluminni}
\dicEntry[lesning] \dicTerm{les··ning} \dicIPA{{l}{\textepsilon}{s}{\textsubring{d}}{n}{i}{\ng}{\r{g}}} \dicPos{f}[4] \dicFlx{(‑ar, ‑ar)}[5] \dicDirectTranslationCS{četba, čtení}
\dicEntry[Lesótó] \dicTerm{Lesótó} \dicIPA{{l}{\textepsilon}{\textlengthmark}{s}{ou}{\textsubring{d}}{ou}} \dicPos{n}[4] \dicFlx{indecl}[2] \dicFieldCat{geog.} \dicDirectTranslationCS{Lesotho}
\dicEntry[Lesótói] \dicTerm{Lesótó|i} \dicIPA{{l}{\textepsilon}{\textlengthmark}{s}{ou}{\textsubring{d}}{ou}{\textsci}} \dicPos{m}[1] \dicFlx{(‑a, ‑ar)}[1] \dicLink{Lesótómaður}
\dicEntry[Lesótómaður] \dicTerm{Lesótó··|maður}\dicTerm{, Lesótói} \dicIPA{{l}\-{\textepsilon}\-{\textlengthmark}\-{s}\-{ou}\-{\textsubring{d}}\-{ou}\-{m}\-{a}\-{ð}\-{\textscy}\-{\textsubring{r}}\-} \dicPos{m}[13] \dicFlx{(‑manns, ‑menn)}[2] \dicDirectTranslationCS{Lesothan(ka)}
\dicEntry[lesótóskur] \dicTerm{lesótóskur} \dicIPA{{l}{\textepsilon}{\textlengthmark}{s}{ou}{\textsubring{d}}{ou}{s}{\r{g}}{\textscy}{\textsubring{r}}} \dicPos{adj}[1]\dicFlx{}[-6] \dicDirectTranslationCS{lesothský}
\dicEntry[lesskilningur] \dicTerm{les··skiln·ing|ur} \dicIPA{{l}{\textepsilon}{\textlengthmark}{s}{\r{\textObardotlessj}}{\textsci}{l}{n}{i}{\ng}{\r{g}}{\textscy}{\textsubring{r}}} \dicPos{m}[6] \dicFlx{(‑s)}[9] \dicDirectTranslationCS{porozumění (psanému) textu}
\dicEntry[lest] \dicTerm{lest} \dicsymFrequent\  \dicIPA{{l}{\textepsilon}{s}{\textsubring{d}}} \dicPos{f}[7] \dicFlx{(‑ar, ‑ir)}[1] \textbf{1.} \dicDirectTranslationCS{vlak, vláček} \dicExampleIS{fara með lest} \dicExampleCS{jet vlakem};  \dicPhraseIS{missa af lestinni} {\textbf{a.}} \dicDirectTranslationCS{zmeškat vlak};  {\textbf{b.}} \dicLangCat{přen.} \dicDirectTranslationCS{nechat si ujet vlak} \dicIndirectTranslationCS{(promeškat příležitost)}  \textbf{2.} \dicFieldCat{nám.} \dicSynonym{farmur} \dicDirectTranslationCS{náklad}  \textbf{3.} \dicSynonym{tonn} \dicDirectTranslationCS{(lodní) tuna}  \textbf{4.} \dicFieldCat{nám.} \dicSynonym*{lestarrúm} \dicDirectTranslationCS{nákladní prostor}  \textbf{5.} \dicSynonym{röð} \dicDirectTranslationCS{řada, karavana} \dicExampleIS{löng lest af bílum} \dicExampleCS{dlouhá řada aut}
\dicEntry[lesta] \dicTerm{lest|a} \dicIPA{{l}{\textepsilon}{s}{\textsubring{d}}{a}} \dicPos{v}[1] \dicFlx{(‑aði)}[1] \dicSynonym{hlaða\smash{\textsuperscript{2}}} \dicDirectTranslationCS{naložit, nakládat} \dicExampleIS{lesta skipið} \dicExampleCS{nakládat loď}
\dicEntry[lestarstjóri] \dicTerm{lestar··stjór|i} \dicIPA{{l}{\textepsilon}{s}{\textsubring{d}}{a}{\textsubring{r}}{s}{\textsubring{d}}{j}{ou}{r}{\textsci}} \dicPos{m}[1] \dicFlx{(‑a, ‑ar)}[1] \dicDirectTranslationCS{strojvedoucí, strojvůdce, strojvůdkyně}
\dicEntry[lestir] \dicTerm{lestir} \dicIPA{{l}{\textepsilon}{s}{\textsubring{d}}{\textsci}{\textsubring{r}}} \dicPos{m} \dicFlx{pl nom} \dicLink{löstur}
\dicEntry[lestrarbók] \dicTerm{lestrar··|bók} \dicIPA{{l}{\textepsilon}{s}{\textsubring{d}}{r}{a}{r}{\textsubring{b}}{ou}{\r{g}}} \dicPos{f}[8] \dicFlx{(‑bókar, ‑bækur)}[5] \dicDirectTranslationCS{čítanka}
\dicEntry[lestrarsalur] \dicTerm{lestrar··sal|ur} \dicIPA{{l}{\textepsilon}{s}{\textsubring{d}}{r}{a}{\textsubring{r}}{s}{a}{l}{\textscy}{\textsubring{r}}} \dicPos{m}[10] \dicFlx{(‑ar, ‑ir)}[14] \dicDirectTranslationCS{čítárna}
\dicEntry[lestur] \dicTerm{lest|ur} \dicsymFrequent\  \dicIPA{{l}{\textepsilon}{s}{\textsubring{d}}{\textscy}{\textsubring{r}}} \dicPos{m}[5] \dicFlx{(‑rar\,/\addthin ‑urs, ‑rar)}[9] \textbf{1.} \dicSynonym*{það að lesa} \dicDirectTranslationCS{čtení, četba} \dicExampleIS{skemmtilegur lestur} \dicExampleCS{zábavné čtení}  \textbf{2.} \dicSynonym{nám} \dicDirectTranslationCS{studium, učení se}  \textbf{3.} \dicSynonym{klifur} \dicDirectTranslationCS{šplhání}
\dicEntry[let] \dicTerm{let} \dicIPA{{l}{\textepsilon}{\textlengthmark}{\textsubring{d}}} \dicPos{v} \dicFlx{ind praes sg 1 pers} \dicLink{letja}
\dicEntry[leti] \dicTerm{let|i} \dicIPA{{l}{\textepsilon}{\textlengthmark}{\textsubring{d}}{\textsci}} \dicPos{f}[3] \dicFlx{(‑i)}[3] \dicDirectTranslationCS{lenost, lenivost, zahálčivost}
\dicEntry[letidýr] \dicTerm{leti··dýr} \dicIPA{{l}{\textepsilon}{\textlengthmark}{\textsubring{d}}{\textsci}{\textsubring{d}}{i}{\textsubring{r}}} \dicPos{n}[2] \dicFlx{(‑s, ‑)}[5] \dicFieldCat{zool.} \dicDirectTranslationCS{lenochod} \textit{(l.~{\textLA{Bradypus}})}  \dicsymPhoto\ 
\dicFigure{ds_image_letidyr_0_1.jpg}{Letidýr}{Letidýr - Stefan Laube, PD}
\dicEntry[letihaugur] \dicTerm{leti··haug|ur} \dicIPA{{l}{\textepsilon}{\textlengthmark}{\textsubring{d}}{\textsci}{h}{\oe i}{\textbabygamma}{\textscy}{\textsubring{r}}} \dicPos{m}[6] \dicFlx{(‑s, ‑ar)}[22] \dicDirectTranslationCS{lenoch, lenoška, povaleč(ka)}
\dicEntry[letilegur] \dicTerm{leti··legur} \dicIPA{{l}{\textepsilon}{\textlengthmark}{\textsubring{d}}{\textsci}{l}{\textepsilon}{\textbabygamma}{\textscy}{\textsubring{r}}} \dicPos{adj}[1]\dicFlx{}[-8] \dicDirectTranslationCS{líný, lenivý}
\dicEntry[letingi] \dicTerm{let··ing|i} \dicIPA{{l}{\textepsilon}{\textlengthmark}{\textsubring{d}}{i}{\textltailn}{\r{\textObardotlessj}}{\textsci}} \dicPos{m}[1] \dicFlx{(‑ja, ‑jar)}[14] \dicDirectTranslationCS{lenoch, lenoška, povaleč(ka)}
\dicEntry[letja] \dicTerm{letja} \dicIPA{{l}{\textepsilon}{\textlengthmark}{\textsubring{d}}{j}{a}} \dicPos{v}[4] \dicFlx{(let, latti, löttum, letti, latt)}[48] \dicFlx{acc} \dicSynonym*{telja úr} \dicDirectTranslationCS{odradit, odrazovat};  \dicPhraseIS{letja e‑n e‑s} \dicDirectTranslationCS{odrazovat (koho) od (čeho)} \dicExampleIS{letja e‑n fararinnar} \dicExampleCS{odrazovat (koho) od cesty}
\dicEntry[letra] \dicTerm{letr|a} \dicIPA{{l}{\textepsilon}{\textlengthmark}{\textsubring{d}}{r}{a}} \dicPos{v}[1] \dicFlx{(‑aði)}[1] \dicFlx{acc} \textbf{1.} \dicSynonym{skrifa} \dicDirectTranslationCS{psát, napsat}  \textbf{2.} \dicSynonym*{höggva stafi í stein} \dicDirectTranslationCS{vysázet\,/\addthin vyrýt písmem}
\dicEntry[Letti] \dicTerm{Lett|i\smash{\textsuperscript{1}}} \dicIPA{{l}{\textepsilon}{h}{\textsubring{d}}{\textsci}} \dicPos{m}[1] \dicFlx{(‑a, ‑ar)}[1] \dicDirectTranslationCS{Lotyš(ka)}
\dicEntry[letti] \dicTerm{letti\smash{\textsuperscript{2}}} \dicIPA{{l}{\textepsilon}{h}{\textsubring{d}}{\textsci}} \dicPos{v} \dicFlx{con pf sg 1 pers} \dicLink{letja}
\dicEntry[Lettland] \dicTerm{Lett··land} \dicIPA{{l}{\textepsilon}{h}{\textsubring{d}}{l}{a}{n}{\textsubring{d}}} \dicPos{n}[2] \dicFlx{(‑s)}[4] \dicFieldCat{geog.} \dicDirectTranslationCS{Lotyšsko}
\dicEntry[lettneska] \dicTerm{lett··nesk|a} \dicIPA{{l}{\textepsilon}{h}{\textsubring{d}}{n}{\textepsilon}{s}{\r{g}}{a}} \dicPos{f}[1] \dicFlx{(‑u)}[5] \dicDirectTranslationCS{lotyština}
\dicEntry[lettneskur] \dicTerm{lett··neskur} \dicIPA{{l}{\textepsilon}{h}{\textsubring{d}}{n}{\textepsilon}{s}{\r{g}}{\textscy}{\textsubring{r}}} \dicPos{adj}[1]\dicFlx{}[-6] \dicDirectTranslationCS{lotyšský}
\dicEntry[letur] \dicTerm{letur} \dicIPA{{l}{\textepsilon}{\textlengthmark}{\textsubring{d}}{\textscy}{\textsubring{r}}} \dicPos{n}[2] \dicFlx{(‑s, ‑)}[25] \textbf{1.} \dicSynonym*{stafategund} \dicDirectTranslationCS{písmo, druh\,/\addthin typ písma, font}  \textbf{2.} \dicSynonym{áletrun} \dicDirectTranslationCS{nápis, epigraf} \dicExampleIS{letur á steini} \dicExampleCS{nápis na kameni}
\dicEntry[leturbreyting] \dicTerm{letur··breyt·ing} \dicIPA{{l}{\textepsilon}{\textlengthmark}{\textsubring{d}}{\textscy}{r}{\textsubring{b}}{r}{ei}{\textsubring{d}}{i}{\ng}{\r{g}}} \dicPos{f}[4] \dicFlx{(‑ar, ‑ar)}[5] \dicDirectTranslationCS{změna písma} \dicIndirectTranslationCS{(z~tučného písma na normální ap.)}
\dicEntry[leturgerð] \dicTerm{letur··gerð} \dicIPA{{l}{\textepsilon}{\textlengthmark}{\textsubring{d}}{\textscy}{r}{\r{\textObardotlessj}}{\textepsilon}{r}{\texttheta}} \dicPos{f}[7] \dicFlx{(‑ar, ‑ir)}[1] \dicSynonym*{leturtegund} \dicDirectTranslationCS{font, znaková sada}
\dicEntry[leturgrafari] \dicTerm{letur··graf·ar|i} \dicIPA{{l}{\textepsilon}{\textlengthmark}{\textsubring{d}}{\textscy}{r}{\r{g}}{r}{a}{v}{a}{r}{\textsci}} \dicPos{m}[1] \dicFlx{(‑a, ‑ar)}[10] \dicDirectTranslationCS{rytec\,/\addthin rytečka (písma)}
\dicEntry[leturstærð] \dicTerm{letur··stærð} \dicIPA{{l}{\textepsilon}{\textlengthmark}{\textsubring{d}}{\textscy}{\textsubring{r}}{s}{\textsubring{d}}{a}{i}{r}{\texttheta}} \dicPos{f}[7] \dicFlx{(‑ar, ‑ir)}[1] \dicDirectTranslationCS{velikost písma}
\dicEntry[lexía] \dicTerm{lexí|a} \dicIPA{{l}{\textepsilon}{x}{s}{i}{j}{a}} \dicPos{f}[1] \dicFlx{(‑u, ‑ur)}[7] \textbf{1.} \dicSynonym{námsefni} \dicDirectTranslationCS{lekce} \dicIndirectTranslationCS{(úsek učební látky)};  \dicPhraseIS{lesa lexíuna sína} \dicDirectTranslationCS{naučit se, udělat domácí úkoly}  \textbf{2.} \dicSynonym{áminning} \dicDirectTranslationCS{lekce, (nepříjemné) poučení}
\dicEntry[leyfa] \dicTerm{leyf|a} \dicsymFrequent\  \dicIPA{{l}{ei}{\textlengthmark}{v}{a}} \dicPos{v}[2] \dicFlx{(‑ði, ‑t)}[106] \dicFlx{dat + acc} \dicDirectTranslationCS{dovolit, dovolovat, povolit, povolovat, nechat, nechávat};  \dicPhraseIS{leyfa e‑m e‑ð} \dicDirectTranslationCS{dovolit (co komu)} \dicExampleIS{leyfa henni að halda á hestinum} \dicExampleCS{dovolit jí vést koně};  \dicIdiom{leyfast}{ \dicPhraseIS{e‑m leyfist að (gera e‑ð)}} \dicFlx{refl impers} \dicDirectTranslationCS{(komu) je dovoleno (udělat (co))} \dicExampleIS{ef mér leyfist að spyrja} \dicExampleCS{pokud se můžu zeptat}
\dicEntry[leyfður] \dicTerm{leyfður} \dicIPA{{l}{ei}{v}{ð}{\textscy}{\textsubring{r}}} \dicPos{adj}[2]\dicFlx{}[-1] \dicDirectTranslationCS{dovolený, povolený} \dicExampleIS{leyfður hraði} \dicExampleCS{povolená rychlost}
\dicEntry[leyfi] \dicTerm{leyfi} \dicsymFrequent\  \dicIPA{{l}{ei}{\textlengthmark}{v}{\textsci}} \dicPos{n}[2] \dicFlx{(‑s, ‑)}[14] \textbf{1.} \dicSynonym{heimild} \dicDirectTranslationCS{licence, oprávnění}  \textbf{2.} \dicDirectTranslationCS{povolení, dovolení, souhlas, svolení} \dicExampleIS{án leyfis} \dicExampleCS{bez povolení};  \dicPhraseIS{með leyfi} \dicFlx{adv} \dicDirectTranslationCS{s~dovolením}  \textbf{3.} \dicSynonym{frí} \dicDirectTranslationCS{dovolená, volno}
\dicEntry[leyfilegur] \dicTerm{leyfi··legur} \dicIPA{{l}{ei}{\textlengthmark}{v}{\textsci}{l}{\textepsilon}{\textbabygamma}{\textscy}{\textsubring{r}}} \dicPos{adj}[1]\dicFlx{}[-8] \dicSynonym{heimill} \dicDirectTranslationCS{povolený, dovolený, přípustný}
\dicEntry[leyfisleysi] \dicTerm{leyfis··leysi} \dicIPA{{l}{ei}{\textlengthmark}{v}{\textsci}{s}{l}{ei}{s}{\textsci}} \dicPos{n}[2] \dicFlx{(‑s)}[20] \dicPhraseIS{gera e‑ð í leyfisleysi} \dicDirectTranslationCS{udělat (co) bez povolení\,/\addthin dovolení\,/\addthin svolení}
\dicEntry[leyfisveiting] \dicTerm{leyfis··veit·ing} \dicIPA{{l}{ei}{\textlengthmark}{v}{\textsci}{s}{v}{ei}{\textsubring{d}}{i}{\ng}{\r{g}}} \dicPos{f}[4] \dicFlx{(‑ar, ‑ar)}[5] \dicDirectTranslationCS{udělení\,/\addthin vydání oprávnění\,/\addthin povolení}
\dicEntry[leyna] \dicTerm{leyn|a} \dicsymFrequent\  \dicIPA{{l}{ei}{\textlengthmark}{n}{a}} \dicPos{v}[2] \dicFlx{(‑di, ‑t)}[141] \dicFlx{dat} \dicDirectTranslationCS{ukrýt, ukrývat, skrýt, skrývat, (u)tajit, zatajovat};  \dicPhraseIS{leyna e‑n e‑u, leyna e‑u fyrir e‑m} \dicDirectTranslationCS{skrýt (co) před (kým), zatajit (co) před (kým), zatajit (komu co)} \dicExampleIS{leyna hana vitneskju um málið} \dicExampleCS{zatajit jí informace o~té záležitosti};  \dicIdiom{leyna}[á]{ \dicPhraseIS{leyna á sér}} \dicDirectTranslationCS{nezdát se} \dicIndirectTranslationCS{(neukazovat pravou skutečnost)};  \dicIdiom{leynast}{ \dicPhraseIS{leynast}} \dicFlx{refl} \dicDirectTranslationCS{skrýt se, skrývat se}
\dicEntry[leynd] \dicTerm{leynd} \dicIPA{{l}{ei}{n}{\textsubring{d}}} \dicPos{f}[7] \dicFlx{(‑ar, ‑ir)}[1] \textbf{1.} \dicSynonym{launung} \dicDirectTranslationCS{tajnost};  \dicPhraseIS{með leynd} \dicFlx{adv} \dicDirectTranslationCS{v~tajnosti, potají}  \textbf{2.} \dicSynonym*{leynistaður} \dicDirectTranslationCS{(místo) utajení}
\begin{xtolerant}{}{1pt}
\dicEntry[leyndardómsfullur] \dicTerm{leyndar·dóms··|fullur} \dicsymFrequent\addthinS\textls[-10]{\dicIPA{{l}\-{ei}\-{n}\-{\textsubring{d}}\-{a}\-{r}\-{\textsubring{d}}\-{ou}\-{m}\-{s}\-{f}\-{\textscy}\-{\textsubring{d}}\-{l}\-{\textscy}\-{\textsubring{r}}}} \dicPos{adj}[10]\addthinS\dicFlx{(comp ‑fyllri, sup ‑fyllstur)}[7] \dicDirectTranslationCS{tajuplný, tajemný, záhadný} \dicExampleIS{leyndardómsfullur á svip} \dicExampleCS{s~tajuplným výrazem ve tváři}
\end{xtolerant}
\dicEntry[leyndardómur] \dicTerm{leyndar··dóm|ur} \dicIPA{{l}{ei}{n}{\textsubring{d}}{a}{r}{\textsubring{d}}{ou}{m}{\textscy}{\textsubring{r}}} \dicPos{m}[6] \dicFlx{(‑s, ‑ar)}[10] \dicDirectTranslationCS{tajemství, tajemnost, záhada} \dicExampleIS{leyndardómur lífsins} \dicExampleCS{tajemství života}
\dicEntry[leyndarmál] \dicTerm{leyndar··mál} \dicsymFrequent\  \dicIPA{{l}{ei}{n}{\textsubring{d}}{a}{r}{m}{au}{\textsubring{l}}} \dicPos{n}[2] \dicFlx{(‑s, ‑)}[5] \dicDirectTranslationCS{tajemství} \dicExampleIS{opinbert leyndarmál} \dicExampleCS{veřejné tajemství}
\dicEntry[leyndur] \dicTerm{leyndur} \dicsymFrequent\  \dicIPA{{l}{ei}{n}{\textsubring{d}}{\textscy}{\textsubring{r}}} \dicPos{adj}[2]\dicFlx{}[-14] \textbf{1.} \dicDirectTranslationCS{skrytý, schovaný} \dicExampleIS{leyndir hæfileikar} \dicExampleCS{skryté schopnosti}  \textbf{2.} \dicDirectTranslationCS{tajný, utajený};  \dicPhraseIS{fara leynt með e‑ð} \dicDirectTranslationCS{tajit (co), držet (co) v~tajnosti};  \dicPhraseIS{halda e‑ð leyndum} \dicDirectTranslationCS{tajit (co), držet (co) v~tajnosti}
\dicEntry[leyni] \dicTerm{leyni} \dicIPA{{l}{ei}{\textlengthmark}{n}{\textsci}} \dicPos{n}[2] \dicFlx{(‑s, ‑)}[14] \dicDirectTranslationCS{skrýš, úkryt};  \dicPhraseIS{í leyni} \dicFlx{adv} \dicDirectTranslationCS{tajně, potají};  \dicPhraseIS{liggja í leyni} \dicDirectTranslationCS{skrývat se, ukrývat se}
\dicEntry[leynifélag] \dicTerm{leyni··fé·|lag} \dicIPA{{l}{ei}{\textlengthmark}{n}{\textsci}{f}{j}{\textepsilon}{l}{a}{x}} \dicPos{n}[2] \dicFlx{(‑lags, ‑lög)}[8] \dicDirectTranslationCS{tajné společenství, tajná společnost}
\dicEntry[leynilegur] \dicTerm{leyni··legur} \dicIPA{{l}{ei}{\textlengthmark}{n}{\textsci}{l}{\textepsilon}{\textbabygamma}{\textscy}{\textsubring{r}}} \dicPos{adj}[1]\dicFlx{}[-8] \dicSynonym{dulinn} \dicDirectTranslationCS{tajný, utajený}
\dicEntry[leynilögregla] \dicTerm{leyni··lög·regl|a} \dicIPA{{l}{ei}{\textlengthmark}{n}{\textsci}{l}{\oe}{\textbabygamma}{r}{\textepsilon}{\r{g}}{l}{a}} \dicPos{f}[1] \dicFlx{(‑u)}[5] \dicDirectTranslationCS{tajná policie}

\dicEntry[leynilögreglumaður] \textls[23]{\dicTerm{leyni··lög·reglu·|maður} \dicIPA{{l}\-{ei}\-{\textlengthmark}\-{n}\-{\textsci}\-{l}\-{\oe}\-{\textbabygamma}\-{r}\-{\textepsilon}\-{\r{g}}\-{l}\-{\textscy}\-{m}\-{a}\-{ð}\-{\textscy}\-{\textsubring{r}}} \dicPos{m}[13] \dicFlx{(‑manns, ‑menn)}[2]} \dicDirectTranslationCS{detektiv, tajný policista, tajná policistka}

\dicEntry[leynilögreglusaga] \dicTerm{leyni·lög·reglu··|saga} \dicIPA{{l}\-{ei}\-{\textlengthmark}\-{n}\-{\textsci}\-{l}\-{\oe}\-{\textbabygamma}\-{r}\-{\textepsilon}\-{\r{g}}\-{l}\-{\textscy}\-{s}\-{a}\-{\textbabygamma}\-{a}\-} \dicPos{f}[1] \dicFlx{(‑sögu, ‑sög\-ur)}[14] \dicDirectTranslationCS{detektivka, detektivní příběh}
\dicEntry[leyniskytta] \dicTerm{leyni··skytt|a} \dicIPA{{l}{ei}{\textlengthmark}{n}{\textsci}{s}{\r{\textObardotlessj}}{\textsci}{h}{\textsubring{d}}{a}} \dicPos{f}[1] \dicFlx{(‑u, ‑ur)}[19] \dicDirectTranslationCS{odstřelovač(ka), snajper(ka)}
\dicEntry[leyniþjónusta] \textls[15]{\dicTerm{leyni··þjón·ust|a} \dicIPA{{l}{ei}{\textlengthmark}{n}{\textsci}{\texttheta}{j}{ou}{n}{\textscy}{s}{\textsubring{d}}{a}} \dicPos{f}[1] \dicFlx{(‑u)}[5] \dicDirectTranslationCS{tajná\,/\addthin vý\-zvěd\-ná služba, rozvědka}}
\dicEntry[leysa] \dicTerm{leys|a} \dicsymFrequent\  \dicIPA{{l}{ei}{\textlengthmark}{s}{a}} \dicPos{v}[2] \dicFlx{(‑ti, ‑t)}[64] \dicFlx{acc} \textbf{1.} \dicSynonym{losa} \dicDirectTranslationCS{rozvázat, rozvazovat, rozmotat, rozmotávat, uvolnit, uvolňovat} \dicExampleIS{leysa bandið} \dicExampleCS{rozvázat provázek}  \textbf{2.} \dicSynonym*{finna lausn} \dicDirectTranslationCS{vyřešit, rozřešit, rozuzlit}  \textbf{3.} \dicPhraseIS{e‑ð leysir} \dicFlx{impers} \dicDirectTranslationCS{(co) taje, (co) roztává} \dicExampleIS{Ísa leysir.} \dicExampleCS{Ledy tají.};  \dicIdiom{leysa}[af]{ \dicPhraseIS{leysa e‑n af}} \dicDirectTranslationCS{zastoupit (koho)}; { \dicPhraseIS{leysa e‑ð af hólmi}} \dicDirectTranslationCS{nahradit (co), zastoupit (co)};  \dicIdiom{leysa}[niður]{ \dicPhraseIS{leysa niður um sig}} \dicDirectTranslationCS{stáhnout si kalhoty};  \dicIdiom{leysa}[undan]{ \dicPhraseIS{leysa e‑n undan e‑u}} \dicDirectTranslationCS{oprostit (koho) od (čeho), osvobodit (koho) od (čeho)};  \dicIdiom{leysa}[upp]{ \dicPhraseIS{leysa e‑ð upp}} \dicDirectTranslationCS{rozpustit\,/\addthin rozpouštět (co) (sůl ve vodě ap.)};  \dicIdiom{leysa}[úr]{ \dicPhraseIS{leysa úr e‑u}} \dicDirectTranslationCS{vyřešit (co), nalézt řešení (čeho)} \dicExampleIS{leysa úr vanda} \dicExampleCS{vyřešit problém};  \dicIdiom{leysa}[út]{ \dicPhraseIS{leysa e‑ð út}} \dicDirectTranslationCS{zaplatit za (co) (za objednané zboží)};  \dicPhraseIS{leysa e‑n út með e‑ð} \dicDirectTranslationCS{obdarovat (koho čím) na odchodu};  \dicIdiom{leysast}{ \dicPhraseIS{leysast}} \dicFlx{refl} {\textbf{a.}} \dicDirectTranslationCS{rozvázat se, rozmotat se};  {\textbf{b.}} \dicDirectTranslationCS{vyřešit se, rozuzlit se} \dicExampleIS{Deilan leystist með samkomulagi.} \dicExampleCS{Spor se vyřešil dohodou.}
\dicEntry[leysanlegur] \dicTerm{leysan··legur} \dicIPA{{l}{ei}{\textlengthmark}{s}{a}{n}{l}{\textepsilon}{\textbabygamma}{\textscy}{\textsubring{r}}} \dicPos{adj}[1]\dicFlx{}[-8] \textbf{1.} \dicFieldCat{chem.} \dicDirectTranslationCS{rozpustný}  \textbf{2.} \dicDirectTranslationCS{(vy)řešitelný}
\dicEntry[leysiefni] \dicTerm{leysi··efni} \dicIPA{{l}{ei}{\textlengthmark}{s}{\textsci}{\textepsilon}{\textsubring{b}}{n}{\textsci}} \dicPos{n}[2] \dicFlx{(‑s, ‑)}[14] \dicFieldCat{chem.} \dicDirectTranslationCS{rozpouštědlo}
\dicEntry[leysigeisli] \dicTerm{leysi··geisl|i} \dicIPA{{l}{ei}{\textlengthmark}{s}{\textsci}{\r{\textObardotlessj}}{ei}{s}{\textsubring{d}}{l}{\textsci}} \dicPos{m}[1] \dicFlx{(‑a, ‑ar)}[1] \dicFieldCat{fyz.} \dicDirectTranslationCS{laserový paprsek}
\dicEntry[leysing] \dicTerm{leys··ing} \dicIPA{{l}{ei}{\textlengthmark}{s}{i}{\ng}{\r{g}}} \dicPos{f}[4] \dicFlx{(‑ar, ‑ar)}[5] \textbf{1.} \dicSynonym{lausn} \dicDirectTranslationCS{(vy)řešení, rozuzlení}  \textbf{2.} \dicSynonym*{snjóbráð} \dicDirectTranslationCS{tání, obleva}  \textbf{3.} \dicSynonym{upplausn} \dicDirectTranslationCS{rozpuštění, rozpouštění (látky ve vodě ap.)} \dicExampleIS{leysing sykurs í vatni} \dicExampleCS{rozpuštění cukru ve vodě}
\dicEntry[leysingi] \dicTerm{leys··ing|i} \dicIPA{{l}{ei}{\textlengthmark}{s}{i}{\textltailn}{\r{\textObardotlessj}}{\textsci}} \dicPos{m}[1] \dicFlx{(‑ja, ‑jar)}[14] \dicFieldCat{hist.} \dicDirectTranslationCS{osvobozený otrok, osvobozená otrokyně}
\dicEntry[leysinn] \dicTerm{leysinn} \dicIPA{{l}{ei}{\textlengthmark}{s}{\textsci}{\textsubring{n}}} \dicPos{adj}[6]\dicFlx{}[-6] \dicDirectTranslationCS{rozpustný}
\dicEntry[leysiprentari] \dicTerm{leysi··prent·ar|i} \dicIPA{{l}{ei}{\textlengthmark}{s}{\textsci}{p\smash{\textsuperscript{h}}}{r}{\textepsilon}{\textsubring{n}}{\textsubring{d}}{a}{r}{\textsci}} \dicPos{m}[1] \dicFlx{(‑a, ‑ar)}[13] \dicFieldCat{poč.} \dicDirectTranslationCS{laserová tiskárna}
\dicEntry[leysir] \dicTerm{leys|ir} \dicIPA{{l}{ei}{\textlengthmark}{s}{\textsci}{\textsubring{r}}} \dicPos{m}[7] \dicFlx{(‑is, ‑ar)}[1] \dicFieldCat{fyz.} \dicDirectTranslationCS{laser}
\dicEntry[leysni] \dicTerm{leysn|i} \dicIPA{{l}{ei}{s}{\textsubring{d}}{n}{\textsci}} \dicPos{f}[3] \dicFlx{(‑i)}[3] \dicFieldCat{chem.} \dicDirectTranslationCS{rozpustnost (materiálu ap.)}
\dicEntry[leyti] \dicTerm{leyti} \dicsymFrequent\  \dicIPA{{l}{ei}{\textlengthmark}{\textsubring{d}}{\textsci}} \dicPos{n}[2] \dicFlx{(‑s)}[20] \textbf{1.} \dicSynonym{tillit} \dicDirectTranslationCS{ohled, zřetel, míra};  \dicPhraseIS{að mestu leyti} \dicFlx{adv} \dicDirectTranslationCS{převážně, většinou};  \dicPhraseIS{að miklu leyti} \dicFlx{adv} \dicDirectTranslationCS{do velké míry, z~velké části};  \dicPhraseIS{að mörgu leyti} \dicFlx{adv} \dicDirectTranslationCS{do značné míry};  \dicPhraseIS{að nokkru leyti} \dicFlx{adv} \dicDirectTranslationCS{do určité míry, částečně};  \dicPhraseIS{að vissu leyti} \dicFlx{adv} \dicDirectTranslationCS{do jisté míry};  \dicPhraseIS{að öðru leyti} \dicFlx{adv} \dicDirectTranslationCS{jinak};  \dicPhraseIS{að þessu leyti} \dicFlx{adv} \dicDirectTranslationCS{v~tomto ohledu\,/\addthin směru}  \textbf{2.} \dicSynonym{tími} \dicDirectTranslationCS{doba, čas} \dicExampleIS{um sjöleytið} \dicExampleCS{kolem sedmé hodiny};  \dicPhraseIS{um hvaða leyti} \dicDirectTranslationCS{v~kolik (hodin)};  \dicPhraseIS{um sama leyti} \dicFlx{adv} \dicDirectTranslationCS{(přibližně) ve stejný čas};  \dicPhraseIS{um það leyti} \dicFlx{adv} \dicDirectTranslationCS{tehdy, tenkrát}
\dicEntry[léð] \dicTerm{léð} \dicIPA{{l}{j}{\textepsilon}{\textlengthmark}{\texttheta}} \dicPos{v} \dicFlx{supin} \dicLink{ljá}
\dicEntry[léði] \dicTerm{léði} \dicIPA{{l}{j}{\textepsilon}{\textlengthmark}{ð}{\textsci}} \dicPos{v} \dicFlx{ind\,/\addthin con pf sg 1 pers} \dicLink{ljá}
\dicEntry[léðum] \dicTerm{léðum} \dicIPA{{l}{j}{\textepsilon}{\textlengthmark}{ð}{\textscy}{\textsubring{m}}} \dicPos{v} \dicFlx{ind pf pl 1 pers} \dicLink{ljá}
\dicEntry[lék] \dicTerm{lék} \dicIPA{{l}{j}{\textepsilon}{\textlengthmark}{\r{g}}} \dicPos{v} \dicFlx{ind pf sg 1 pers} \dicLink{leika}
\dicEntry[léki] \dicTerm{léki} \dicIPA{{l}{j}{\textepsilon}{\textlengthmark}{\r{\textObardotlessj}}{\textsci}} \dicPos{v} \dicFlx{con pf sg 1 pers} \dicLink{leika}
\dicEntry[lékum] \dicTerm{lékum} \dicIPA{{l}{j}{\textepsilon}{\textlengthmark}{\r{g}}{\textscy}{\textsubring{m}}} \dicPos{v} \dicFlx{ind pf pl 1 pers} \dicLink{leika}
\dicEntry[lélega] \dicTerm{lé··lega} \dicIPA{{l}{j}{\textepsilon}{\textlengthmark}{l}{\textepsilon}{\textbabygamma}{a}} \dicPos{adv} \dicDirectTranslationCS{špatně, bídně, mizerně}
\dicEntry[lélegur] \dicTerm{lé··legur} \dicsymFrequent\  \dicIPA{{l}{j}{\textepsilon}{\textlengthmark}{l}{\textepsilon}{\textbabygamma}{\textscy}{\textsubring{r}}} \dicPos{adj}[1]\dicFlx{}[-8] \dicSynonym{slæmur} \dicDirectTranslationCS{špatný, bídný, mizerný, chabý} \dicExampleIS{benda á lélegan árangur} \dicExampleCS{upozornit na špatný výsledek}
\dicEntry[lén] \dicTerm{lén} \dicIPA{{l}{j}{\textepsilon}{\textlengthmark}{\textsubring{n}}} \dicPos{n}[2] \dicFlx{(‑s, ‑)}[5] \textbf{1.} \dicSynonym{ríki} \dicDirectTranslationCS{léno}  \textbf{2.} \dicFieldCat{poč.} \dicSynonym{umdæmi} \dicDirectTranslationCS{doména}  \textbf{3.} \dicFieldCat{biol.} \dicDirectTranslationCS{doména}
\dicEntry[lénskur] \dicTerm{lénskur} \dicIPA{{l}{j}{\textepsilon}{n}{s}{\r{g}}{\textscy}{\textsubring{r}}} \dicPos{adj}[1]\dicFlx{}[-1] \dicDirectTranslationCS{feudální}
\dicEntry[lénsmaður] \dicTerm{léns··|maður} \dicIPA{{l}{j}{\textepsilon}{n}{s}{m}{a}{ð}{\textscy}{\textsubring{r}}} \dicPos{m}[13] \dicFlx{(‑manns, ‑menn)}[2] \dicDirectTranslationCS{leník, vazal(ka)}
\dicEntry[léreft] \dicTerm{léreft} \dicIPA{{l}{j}{\textepsilon}{\textlengthmark}{r}{\textepsilon}{f}{\textsubring{d}}} \dicPos{n}[2] \dicFlx{(‑s, ‑)}[5] \dicDirectTranslationCS{(lněné) plátno}
\dicEntry[lét] \dicTerm{lét} \dicIPA{{l}{j}{\textepsilon}{\textlengthmark}{\textsubring{d}}} \dicPos{v} \dicFlx{ind pf sg 1 pers} \dicLink{láta}
\dicEntry[léti] \dicTerm{léti} \dicIPA{{l}{j}{\textepsilon}{\textlengthmark}{\textsubring{d}}{\textsci}} \dicPos{v} \dicFlx{con pf sg 1 pers} \dicLink{láta}
\dicEntry[létt] \dicTerm{létt} \dicsymFrequent\  \dicIPA{{l}{j}{\textepsilon}{h}{\textsubring{d}}} \dicPos{adv} \dicFlx{(comp ‑ar, sup ‑ast)} \dicDirectTranslationCS{lehce, mírně, jemně, zlehka} \dicExampleIS{Hann klappaði létt á bakið á mér.} \dicExampleCS{Lehce mě poklepal po zádech.}
\dicEntry[létta] \dicTerm{létt|a} \dicsymFrequent\  \dicIPA{{l}{j}{\textepsilon}{h}{\textsubring{d}}{a}} \dicPos{v}[2] \dicFlx{(‑i, ‑)}[2] \dicFlx{acc\,/\addthin dat} \textbf{1.} \dicDirectTranslationCS{odlehčit, odlehčovat, ulehčit, ulehčovat, (u)činit lehčím} \dicExampleIS{létta byrðina} \dicExampleCS{odlehčit břemeno}  \textbf{2.} \dicPhraseIS{e‑m léttir} \dicFlx{impers} \dicDirectTranslationCS{(komu) se ulehčuje, (komu) se ulevuje}  \textbf{3.} \dicPhraseIS{e‑u léttir} \dicFlx{impers} \dicDirectTranslationCS{(čeho) ubývá (mlhy ap.), (co) mizí};  \dicIdiom{létta}[af]{ \dicPhraseIS{létta af e‑m e‑u}} \dicDirectTranslationCS{odlehčit (komu) od (čeho)}; { \dicPhraseIS{það léttir af e‑m e‑u}} \dicFlx{impers} \dicDirectTranslationCS{(komu) se ulevuje od (čeho)};  \dicIdiom{létta}[á]{ \dicPhraseIS{létta á sér}} \dicDirectTranslationCS{ulevit si} \dicIndirectTranslationCS{(vykonat tělesnou potřebu)};  \dicIdiom{létta}[sér]{ \dicPhraseIS{létta sér lífið}} \dicDirectTranslationCS{pobavit se, povyrazit si};  \dicIdiom{létta}[til]{ \dicPhraseIS{létta til}} \dicDirectTranslationCS{vyčasit se, vyjasnit se, vyjasňovat se} \dicExampleIS{Það er að létta til.} \dicExampleCS{Začíná se vyjasňovat.};  \dicIdiom{létta}[undir]{ \dicPhraseIS{létta undir með e‑m}} \dicDirectTranslationCS{ulehčit (komu), učinit to (komu) lehčím};  \dicIdiom{léttast}{ \dicPhraseIS{léttast}} \dicFlx{refl} \dicDirectTranslationCS{(z)hubnout};  \dicIdiom{létta}{ \dicPhraseIS{létta akkerum}} \dicFieldCat{nám.} \dicDirectTranslationCS{zvednout kotvu}
\dicEntry[léttasótt] \dicTerm{létta··sótt} \dicIPA{{l}{j}{\textepsilon}{h}{\textsubring{d}}{a}{s}{ou}{h}{\textsubring{d}}} \dicPos{f}[7] \dicFlx{(‑ar)}[3] \dicDirectTranslationCS{porodní bolesti, kontrakce} \dicExampleIS{taka léttasóttina} \dicExampleCS{mít porodní bolesti}
\dicEntry[létteind] \dicTerm{létt··eind} \dicIPA{{l}{j}{\textepsilon}{h}{\textsubring{d}}{ei}{n}{\textsubring{d}}} \dicPos{f}[7] \dicFlx{(‑ar, ‑ir)}[1] \dicFieldCat{fyz.} \dicDirectTranslationCS{lepton}
\dicEntry[létti] \dicTerm{létt|i} \dicIPA{{l}{j}{\textepsilon}{h}{\textsubring{d}}{\textsci}} \dicPos{m}[1] \dicFlx{(‑a, ‑ar)}[1] \textbf{1.} \dicSynonym{léttir} \dicDirectTranslationCS{ulehčení, úleva}  \textbf{2.} \dicSynonym{axlaband} \dicDirectTranslationCS{popruh, šle}  \textbf{3.} \dicSynonym*{fúsleiki} \dicDirectTranslationCS{ochota};  \dicPhraseIS{segja allt af létta} \dicDirectTranslationCS{ochotně všechno povědět}
\dicEntry[léttir] \dicTerm{létt|ir} \dicIPA{{l}{j}{\textepsilon}{h}{\textsubring{d}}{\textsci}{\textsubring{r}}} \dicPos{m}[7] \dicFlx{(‑is, ‑ar)}[1] \textbf{1.} \dicSynonym{huggun} \dicDirectTranslationCS{útěcha, úleva}  \textbf{2.} \dicSynonym{hjálp} \dicDirectTranslationCS{pomoc};  \dicPhraseIS{vera e‑m til léttis} \dicDirectTranslationCS{pomoci (komu), ulehčit (komu)}  \textbf{3.} \dicFieldCat{zool.} \dicDirectTranslationCS{delfín obecný} \textit{(l.~{\textLA{Delphinus delphis}})}  \dicsymPhoto\ 
\dicFigure{26462.jpg}{Léttir}{Léttir - Shane Anderson, Biolib, PD}
\dicEntry[léttklæddur] \dicTerm{létt··klæddur} \dicIPA{{l}{j}{\textepsilon}{h}{\textsubring{d}}{k\smash{\textsuperscript{h}}}{l}{a}{i}{\textsubring{d}}{\textscy}{\textsubring{r}}} \dicPos{adj}[2]\dicFlx{}[-18] \dicDirectTranslationCS{lehce oblečený}
\dicEntry[léttleiki] \dicTerm{létt··leik|i} \dicIPA{{l}{j}{\textepsilon}{h}{\textsubring{d}}{l}{ei}{\r{\textObardotlessj}}{\textsci}} \dicPos{m}[1] \dicFlx{(‑a)}[3] \dicDirectTranslationCS{lehkost} \dicExampleIS{óbærilegur léttleiki tilverunnar} \dicExampleCS{nesnesitelná lehkost bytí}
\dicEntry[léttlyndi] \dicTerm{létt··lyndi} \dicIPA{{l}{j}{\textepsilon}{h}{\textsubring{d}}{l}{\textsci}{n}{\textsubring{d}}{\textsci}} \dicPos{n}[2] \dicFlx{(‑s)}[20] \dicDirectTranslationCS{dobrá nálada, uvolněnost, bezstarostnost}
\dicEntry[léttlyndur] \dicTerm{létt··lyndur} \dicIPA{{l}{j}{\textepsilon}{h}{\textsubring{d}}{l}{\textsci}{n}{\textsubring{d}}{\textscy}{\textsubring{r}}} \dicPos{adj}[2]\dicFlx{}[-14] \dicDirectTranslationCS{uvolněný, bezstarostný}
\dicEntry[léttmjólk] \dicTerm{létt··mjólk} \dicIPA{{l}{j}{\textepsilon}{h}{\textsubring{d}}{m}{j}{ou}{\textsubring{l}}{\r{g}}} \dicPos{f}[10] \dicFlx{(‑ur)}[2] \dicDirectTranslationCS{nízkotučné mléko}
\dicEntry[léttskýjaður] \dicTerm{létt··skýj·|aður} \dicIPA{{l}{j}{\textepsilon}{h}{\textsubring{d}}{s}{\r{\textObardotlessj}}{i}{j}{a}{ð}{\textscy}{\textsubring{r}}} \dicPos{adj}[3] \dicFlx{(f ‑uð)}[3] \dicPhraseIS{það er léttskýjað} \dicDirectTranslationCS{je polojasno}
\dicEntry[léttur] \dicTerm{léttur} \dicsymFrequent\  \dicIPA{{l}{j}{\textepsilon}{h}{\textsubring{d}}{\textscy}{\textsubring{r}}} \dicPos{adj}[1]\dicFlx{}[-10] \textbf{1.} \dicDirectTranslationCS{lehký, snadný} \dicExampleIS{létt verk} \dicExampleCS{lehká práce} \dicAntonym{erfiður}  \textbf{2.} \dicDirectTranslationCS{lehký} \dicIndirectTranslationCS{(vážící málo)} \dicAntonym{þungur}  \textbf{3.} \dicSynonym{glaður} \dicDirectTranslationCS{uvolněný, klidný, bezstarostný} \dicExampleIS{Hún er létt í lund.} \dicExampleCS{Je bezstarostná.}  \textbf{4.} \dicDirectTranslationCS{dietní, lehký (strava ap.)} \dicExampleIS{léttar veitingar} \dicExampleCS{lehké občerstvení}  \textbf{5.} \dicDirectTranslationCS{lehký, slabý (alkohol ap.)} \dicExampleIS{létt vín} \dicExampleCS{slabý alkohol}
\dicEntry[léttúð] \dicTerm{létt··úð} \dicIPA{{l}{j}{\textepsilon}{h}{\textsubring{d}}{u}{\texttheta}} \dicPos{f}[4] \dicFlx{(‑ar)}[3] \textbf{1.} \dicSynonym{alvöruleysi} \dicDirectTranslationCS{lehkovážnost}  \textbf{2.} \dicSynonym{lauslæti} \dicDirectTranslationCS{rozpustilost, frivolnost}
\dicEntry[léttúðugur] \dicTerm{létt·úð··ugur} \dicIPA{{l}{j}{\textepsilon}{h}{\textsubring{d}}{u}{ð}{\textscy}{\textbabygamma}{\textscy}{\textsubring{r}}} \dicPos{adj}[1]\dicFlx{}[-8] \dicSynonym{lauslátur} \dicDirectTranslationCS{frivolní, prostopášný}
\dicEntry[léttvín] \dicTerm{létt··vín} \dicIPA{{l}{j}{\textepsilon}{h}{\textsubring{d}}{v}{i}{\textsubring{n}}} \dicPos{n}[2] \dicFlx{(‑s, ‑)}[5] \dicDirectTranslationCS{víno}
\dicEntry[léttvægur] \dicTerm{létt··vægur} \dicIPA{{l}{j}{\textepsilon}{h}{\textsubring{d}}{v}{a}{i}{\textbabygamma}{\textscy}{\textsubring{r}}} \dicPos{adj}[1]\dicFlx{}[-1] \dicDirectTranslationCS{nedůležitý, nepodstat\-ný, nevýznamný}
\dicEntry[létum] \dicTerm{létum} \dicIPA{{l}{j}{\textepsilon}{\textlengthmark}{\textsubring{d}}{\textscy}{\textsubring{m}}} \dicPos{v} \dicFlx{ind pf pl 1 pers} \dicLink{láta}
\dicEntry[lh. nt.] \dicTerm{lh. nt.} \dicPos{zkr} \dicPhraseIS{lýsingarháttur nútíðar} \dicFieldCat{jaz.} \dicDirectTranslationCS{příčestí přítomné}
\dicEntry[lh. þt.] \dicTerm{lh. þt.} \dicPos{zkr} \dicPhraseIS{lýsingarháttur þátíðar} \dicFieldCat{jaz.} \dicDirectTranslationCS{příčestí minulé}
\dicEntry[lið] \dicTerm{lið} \dicsymFrequent\  \dicIPA{{l}{\textsci}{\textlengthmark}{\texttheta}} \dicPos{n}[2] \dicFlx{(‑s, ‑)}[5] \textbf{1.} \dicSynonym{hópur} \dicDirectTranslationCS{družstvo, oddíl, četa};  \dicPhraseIS{ganga í lið við e‑n} \dicDirectTranslationCS{spojit síly s~(kým)};  \dicPhraseIS{safna liði} \dicDirectTranslationCS{shromažďovat podporu}  \textbf{2.} \dicFieldCat{sport.} \dicDirectTranslationCS{družina, tým, mužstvo}  \textbf{3.} \dicSynonym{aðstoð} \dicDirectTranslationCS{výpomoc, pomocná ruka};  \dicPhraseIS{koma til liðs við e‑n} \dicDirectTranslationCS{vypomoct (komu), přijít (komu) na pomoc};  \dicPhraseIS{verða e‑m að liði} \dicDirectTranslationCS{vypomoct (komu), asistovat (komu)}  \textbf{4.} \dicLangCat{hovor.} \dicSynonym{fólk} \dicDirectTranslationCS{lidi}
\dicEntry[liða] \dicTerm{lið|a} \dicIPA{{l}{\textsci}{\textlengthmark}{ð}{a}} \dicPos{v}[1] \dicFlx{(‑aði)}[1] \dicFlx{acc} \textbf{1.} \dicSynonym{skipta} \dicDirectTranslationCS{(roz)dělit (do skupin)}  \textbf{2.} \dicSynonym*{gera hrokkið} \dicDirectTranslationCS{natočit, natáčet, (na)kadeřit (vlasy ap.)} \dicExampleIS{liða hár} \dicExampleCS{natáčet vlasy};  \dicIdiom{liðast}{ \dicPhraseIS{liðast}} \dicFlx{refl} \dicDirectTranslationCS{vlnit se, kroutit se, kadeřit se}
\dicEntry[liðaður] \dicTerm{lið··|aður} \dicIPA{{l}{\textsci}{\textlengthmark}{ð}{a}{ð}{\textscy}{\textsubring{r}}} \dicPos{adj}[3] \dicFlx{(f ‑uð)}[3] \dicSynonym*{bylgjaður} \dicDirectTranslationCS{natočený, nakadeřený, navlněný}
\dicEntry[liðagigt] \dicTerm{liða··gigt} \dicIPA{{l}{\textsci}{\textlengthmark}{ð}{a}{\r{\textObardotlessj}}{\textsci}{x}{\textsubring{d}}} \dicPos{f}[4] \dicFlx{(‑ar)}[3] \dicFieldCat{med.} \dicDirectTranslationCS{revmatoidní artritida}
\dicEntry[liðamót] \dicTerm{liða··mót} \dicIPA{{l}{\textsci}{\textlengthmark}{ð}{a}{m}{ou}{\textsubring{d}}} \dicPos{n}[2] \dicFlx{pl}[1] \dicFieldCat{anat.} \dicDirectTranslationCS{kloub}
\dicEntry[liðband] \dicTerm{lið··|band} \dicIPA{{l}{\textsci}{ð}{\textsubring{b}}{a}{n}{\textsubring{d}}} \dicPos{n}[2] \dicFlx{(‑bands, ‑bönd)}[8] \dicFieldCat{anat.} \dicDirectTranslationCS{vazivo}
\dicEntry[liðdýr] \dicTerm{lið··dýr} \dicIPA{{l}{\textsci}{ð}{\textsubring{d}}{i}{\textsubring{r}}} \dicPos{n}[2] \dicFlx{(‑s, ‑)}[5] \dicFieldCat{zool.} \dicDirectTranslationCS{členovec} \textit{(l.~{\textLA{Arthropoda}})}
\dicEntry[liðfall] \dicTerm{lið··|fall} \dicIPA{{l}{\textsci}{ð}{f}{a}{\textsubring{d}}{\textsubring{l}}} \dicPos{n}[2] \dicFlx{(‑falls, ‑föll)}[8] \dicFieldCat{jaz.} \dicSynonym{brottfall} \dicDirectTranslationCS{elipsa, výpustka}
\dicEntry[liðfelldur] \dicTerm{lið··felldur} \dicIPA{{l}{\textsci}{ð}{f}{\textepsilon}{l}{\textsubring{d}}{\textscy}{\textsubring{r}}} \dicPos{adj}[2]\dicFlx{}[-14] \dicFieldCat{jaz.} \dicDirectTranslationCS{eliptický} \dicExampleIS{liðfellt orðtak} \dicExampleCS{eliptický výraz}
\dicEntry[liðhlaupi] \dicTerm{lið··hlaup|i} \dicIPA{{l}{\textsci}{ð}{\textsubring{l}}{\oe i}{\textsubring{b}}{\textsci}} \dicPos{m}[1] \dicFlx{(‑a, ‑ar)}[1] \dicFieldCat{voj.} \dicDirectTranslationCS{zběh, dezertér(ka)}
\dicEntry[liði] \dicTerm{lið|i\smash{\textsuperscript{1}}} \dicIPA{{l}{\textsci}{\textlengthmark}{ð}{\textsci}} \dicPos{m}[1] \dicFlx{(‑a, ‑ar)}[1] \dicDirectTranslationCS{člen(ka) týmu (vojenského, zdravotního ap.)} \dicIndirectTranslationCS{(vyskytuje se jako druhá část složeného slova)} \dicExampleIS{skæruliði} \dicExampleCS{partyzán}
\dicEntry[liði] \dicTerm{liði\smash{\textsuperscript{2}}} \dicIPA{{l}{\textsci}{\textlengthmark}{ð}{\textsci}} \dicPos{v} \dicFlx{con pf sg 1 pers} \dicLink{líða}
\dicEntry[liðið] \dicTerm{liðið} \dicIPA{{l}{\textsci}{\textlengthmark}{ð}{\textsci}{\texttheta}} \dicPos{v} \dicFlx{supin} \dicLink{líða}
\dicEntry[liðinn] \dicTerm{liðinn} \dicsymFrequent\  \dicIPA{{l}{\textsci}{\textlengthmark}{ð}{\textsci}{\textsubring{n}}} \dicPos{adj}[6]\dicFlx{}[-6] \textbf{1.} \dicDirectTranslationCS{minulý, uplynulý} \dicExampleIS{í liðnum mánuði} \dicExampleCS{v~minulém měsíci}  \textbf{2.} \dicDirectTranslationCS{zesnulý};  \dicPhraseIS{vera vel liðinn} \dicDirectTranslationCS{být populární\,/\addthin oblíbený}
\dicEntry[liðka] \dicTerm{liðk|a} \dicIPA{{l}{\textsci}{\texttheta}{\r{g}}{a}} \dicPos{v}[1] \dicFlx{(‑aði)}[1] \dicFlx{acc} \dicDirectTranslationCS{rozhýbat, rozpohybovat} \dicExampleIS{liðka handlegginn} \dicExampleCS{rozhýbat ruku};  \dicPhraseIS{liðka sig} \dicDirectTranslationCS{rozhýbat se};  \dicIdiom{liðka}[til]{ \dicPhraseIS{liðka e‑ð til}} \dicDirectTranslationCS{usnadnit (co)};  \dicIdiom{liðkast}{ \dicPhraseIS{liðkast}} \dicFlx{refl} \dicDirectTranslationCS{rozhýbat se}
\dicEntry[liðlangur] \dicTerm{lið··|langur} \dicIPA{{l}{\textsci}{ð}{l}{au}{\ng}{\r{g}}{\textscy}{\textsubring{r}}} \dicPos{adj}[10] \dicFlx{(f ‑löng, comp ‑lengri, sup ‑leng\-st\-ur)}[2] \dicPhraseIS{allan liðlangan daginn} \dicFlx{adv} \dicDirectTranslationCS{(po) celý celičký den}
\dicEntry[liðlegur] \dicTerm{lið··legur} \dicIPA{{l}{\textsci}{ð}{l}{\textepsilon}{\textbabygamma}{\textscy}{\textsubring{r}}} \dicPos{adj}[1]\dicFlx{}[-8] \textbf{1.} \dicSynonym{fimur} \dicDirectTranslationCS{obratný}  \textbf{2.} \dicSynonym{hjálpsamur} \dicDirectTranslationCS{vstřícný, úslužný, ochotný pomoci}
\dicEntry[liðleiki] \dicTerm{lið··leik|i} \dicIPA{{l}{\textsci}{ð}{l}{ei}{\r{\textObardotlessj}}{\textsci}} \dicPos{m}[1] \dicFlx{(‑a)}[3] \textbf{1.} \dicDirectTranslationCS{pohyblivost, obratnost}  \textbf{2.} \dicDirectTranslationCS{vstřícnost, ochota pomoci}
\dicEntry[liðsauki] \dicTerm{liðs··auk|i} \dicIPA{{l}{\textsci}{ð}{s}{\oe i}{\r{\textObardotlessj}}{\textsci}} \dicPos{m}[1] \dicFlx{(‑a)}[3] \dicDirectTranslationCS{posila, posílení\,/\addthin rozšíření řad}
\dicEntry[liðsemd] \textls[30]{\dicTerm{lið··semd} \dicIPA{{l}{\textsci}{ð}{s}{\textepsilon}{m}{\textsubring{d}}} \dicPos{f}[7] \dicFlx{(‑ar)}[3] \dicSynonym{aðstoð} \dicDirectTranslationCS{výpomoc, asis\nobreak tence}}
\dicEntry[liðsforingi] \dicTerm{liðs··for·ing|i} \dicIPA{{l}{\textsci}{ð}{s}{f}{\textopeno}{r}{i}{\textltailn}{\r{\textObardotlessj}}{\textsci}} \dicPos{m}[1] \dicFlx{(‑ja, ‑jar)}[14] \dicFieldCat{voj.} \dicDirectTranslationCS{poručík, poručice}
\dicEntry[liðsinna] \dicTerm{lið··sinn|a} \dicIPA{{l}{\textsci}{ð}{s}{\textsci}{n}{a}} \dicPos{v}[2] \dicFlx{(‑ti, ‑t)}[78] \dicFlx{dat} \dicSynonym{hjálpa} \dicDirectTranslationCS{asistovat, (vy)pomáhat} \dicExampleIS{liðsinna e‑m við e‑ð} \dicExampleCS{asistovat (komu) při (čem)}
\dicEntry[liðsinni] \dicTerm{lið··sinni} \dicIPA{{l}{\textsci}{ð}{s}{\textsci}{n}{\textsci}} \dicPos{n}[2] \dicFlx{(‑s)}[20] \dicSynonym{aðstoð} \dicDirectTranslationCS{asistence, (vý)pomoc}
\dicEntry[liðsmaður] \dicTerm{liðs··|maður} \dicIPA{{l}{\textsci}{ð}{s}{m}{a}{ð}{\textscy}{\textsubring{r}}} \dicPos{m}[13] \dicFlx{(‑manns, ‑menn)}[2] \textbf{1.} \dicSynonym{hermaður} \dicDirectTranslationCS{voják(yně)}  \textbf{2.} \dicSynonym{stuðningsmaður} \dicDirectTranslationCS{přívrženec, přívrženkyně, stoupenec, stoupenkyně}  \textbf{3.} \dicDirectTranslationCS{(spolu)hráč(ka), člen(ka) družiny\,/\addthin týmu}
\dicEntry[liðtækur] \dicTerm{lið··tækur} \dicIPA{{l}{\textsci}{ð}{t\smash{\textsuperscript{h}}}{a}{i}{\r{g}}{\textscy}{\textsubring{r}}} \dicPos{adj}[1]\dicFlx{}[-6] \dicSynonym{hæfur} \dicDirectTranslationCS{kompetentní, způsobilý}
\dicEntry[liðugur] \dicTerm{lið··ugur} \dicIPA{{l}{\textsci}{\textlengthmark}{ð}{\textscy}{\textbabygamma}{\textscy}{\textsubring{r}}} \dicPos{adj}[1]\dicFlx{}[-8] \dicSynonym{fimur} \dicDirectTranslationCS{hbitý, mrštný} \dicExampleIS{vera liðugur í hreyfingum} \dicExampleCS{pohybovat se hbitě};  \dicPhraseIS{vera laus og liðugur} \dicDirectTranslationCS{být volný, užívat si svobody}
\dicEntry[liðum] \dicTerm{liðum} \dicIPA{{l}{\textsci}{\textlengthmark}{ð}{\textscy}{\textsubring{m}}} \dicPos{v} \dicFlx{ind pf pl 1 pers} \dicLink{líða}
\dicEntry[liður] \dicTerm{lið|ur} \dicsymFrequent\  \dicIPA{{l}{\textsci}{\textlengthmark}{ð}{\textscy}{\textsubring{r}}} \dicPos{m}[10] \dicFlx{(‑ar\,/\addthin ‑s, ‑ir)}[37] \textbf{1.} \dicFieldCat{anat.} \dicSynonym{liðamót} \dicDirectTranslationCS{kloub} \dicExampleIS{axlarliður} \dicExampleCS{ramenní kloub};  \dicPhraseIS{geta hvorki hrært legg né lið} \dicLangCat{přen.} \dicDirectTranslationCS{nemoci se vůbec pohnout}  \textbf{2.} \dicSynonym{bylgja} \dicDirectTranslationCS{kudrna, vlnka (ve vlasech)} \dicExampleIS{liður í hári} \dicExampleCS{kudrna ve vlasech}  \textbf{3.} \dicSynonym{atriði} \dicDirectTranslationCS{položka, bod} \dicExampleIS{Þessi liður er rangur.} \dicExampleCS{Tato položka je špatně.}  \textbf{4.} \dicSynonym{ættliður} \dicDirectTranslationCS{koleno, pokolení, generace}  \textbf{5.} \dicFieldCat{mat.} \dicDirectTranslationCS{člen}
\dicEntry[liðveisla] \dicTerm{lið··veisl|a} \dicIPA{{l}{\textsci}{ð}{v}{ei}{s}{\textsubring{d}}{l}{a}} \dicPos{f}[1] \dicFlx{(‑u, ‑ur)}[13] \dicDirectTranslationCS{asistence, (vý)pomoc}
\dicEntry[liðþjálfi] \dicTerm{lið··þjálf|i} \dicIPA{{l}{\textsci}{ð}{\texttheta}{j}{au}{l}{v}{\textsci}} \dicPos{m}[1] \dicFlx{(‑a, ‑ar)}[1] \dicFieldCat{voj.} \dicDirectTranslationCS{seržant(ka)}
\dicEntry[Liechtenstein] \dicTerm{Liechtenstein} \dicIPA{{l}{i}{x}{\textsubring{d}}{\textepsilon}{n}{s}{\textsubring{d}}{a}{i}{\textsubring{n}}} \dicPos{n}[4] \dicFlx{indecl}[2] \dicFieldCat{geog.} \dicDirectTranslationCS{Lichtenštejnsko}
\dicEntry[Liechtensteini] \dicTerm{Liechtenstein|i} \dicIPA{{l}{i}{x}{\textsubring{d}}{\textepsilon}{n}{s}{\textsubring{d}}{a}{i}{n}{\textsci}} \dicPos{m}[1] \dicFlx{(‑a, ‑ar)}[1] \dicDirectTranslationCS{Lichtenštejňan(ka), Lichtenštejnec, Lichtenštejnka}
\dicEntry[liechtensteinskur] \dicTerm{liechtensteinskur} \dicIPA{{l}{i}{x}{\textsubring{d}}{\textepsilon}{n}{s}{\textsubring{d}}{a}{i}{n}{s}{\r{g}}{\textscy}{\textsubring{r}}} \dicPos{adj}[1]\dicFlx{}[-6] \dicDirectTranslationCS{lichtenštejnský}
\dicEntry[lifa] \dicTerm{lif|a} \dicsymFrequent\  \dicIPA{{l}{\textsci}{\textlengthmark}{v}{a}} \dicPos{v}[2] \dicFlx{(‑ði, ‑að)}[124] \dicFlx{acc} \textbf{1.} \dicDirectTranslationCS{žít} \dicExampleIS{lifa langa ævi} \dicExampleCS{žít dlouhý život}  \textbf{2.} \dicDirectTranslationCS{existovat}  \textbf{3.} \dicDirectTranslationCS{zažít, prožít} \dicExampleIS{lifa þessa atburði} \dicExampleCS{zažít tyto události}  \textbf{4.} \dicSynonym*{vera logandi} \dicDirectTranslationCS{planout, svítit};  \dicIdiom{lifa}[af]{ \dicPhraseIS{lifa af}} \dicSynonym{endast} \dicDirectTranslationCS{přežít}; { \dicPhraseIS{lifa af e‑u}} \dicDirectTranslationCS{vyžít z~(čeho)};  \dicIdiom{lifa}[á]{ \dicPhraseIS{lifa á e‑u}} \dicDirectTranslationCS{živit se (čím)};  \dicIdiom{lifa}[eftir]{ \dicPhraseIS{lifa eftir e‑u}} \dicDirectTranslationCS{žít podle (čeho), žít v~souladu s~(čím)};  \dicIdiom{lifa}[fyrir]{ \dicPhraseIS{lifa fyrir e‑n}} \dicDirectTranslationCS{žít pro (koho)};  \dicIdiom{lifa}[inn í]{ \dicPhraseIS{lifa sig inn í e‑ð}} \dicDirectTranslationCS{vžít se do (čeho)} \dicExampleIS{lifa sig inn í starfið} \dicExampleCS{vžít se do práce};  \dicIdiom{lifa}[við]{ \dicPhraseIS{lifa við e‑ð}} \dicDirectTranslationCS{žít s~(čím) (pocitem viny ap.)}
\dicEntry[lifandi] \dicTerm{lif··|andi\smash{\textsuperscript{1}}} \dicIPA{{l}{\textsci}{\textlengthmark}{v}{a}{n}{\textsubring{d}}{\textsci}} \dicPos{m}[2] \dicFlx{(‑anda, ‑endur)}[1] \dicDirectTranslationCS{živý člověk, živá bytost} \dicIndirectTranslationCS{(zvláště v~pl)}
\dicEntry[lifandi] \dicTerm{lif··andi\smash{\textsuperscript{2}}} \dicsymFrequent\  \dicIPA{{l}{\textsci}{\textlengthmark}{v}{a}{n}{\textsubring{d}}{\textsci}} \dicPos{adj}[13] \dicFlx{indecl}[1] \textbf{1.} \dicSynonym*{á lífi} \dicDirectTranslationCS{živý, (jsoucí) naživu} \dicExampleIS{ná birninum lifandi} \dicExampleCS{chytit medvěda živého}  \textbf{2.} \dicSynonym{líflegur} \dicDirectTranslationCS{živý, čilý};  \dicIdiom{lifandi}{ \dicPhraseIS{ekki nokkur\,/\addthin enginn lifandi maður}} \dicLangCat{přen.} \dicDirectTranslationCS{ani živáček};  \dicPhraseIS{lifandi tónlist} \dicDirectTranslationCS{živá hudba}
\dicEntry[lifna] \dicTerm{lifn|a} \dicsymFrequent\  \dicIPA{{l}{\textsci}{\textsubring{b}}{n}{a}} \dicPos{v}[1] \dicFlx{(‑aði)}[1] \textbf{1.} \dicDirectTranslationCS{oživit se, rozhořet se}  \textbf{2.} \dicPhraseIS{lifna við} \dicSynonym*{vakna til lífs} \dicDirectTranslationCS{oživnout, ožít} \dicExampleIS{Fjöllin lifnuðu við eftir fimmtíu ár.} \dicExampleCS{Hory ožily po padesáti letech.}  \textbf{3.} \dicPhraseIS{það lifnar yfir e‑m} \dicFlx{impers} \dicLangCat{přen.} \dicDirectTranslationCS{(kdo) ožívá, (kdo) se stává živějším}
\dicEntry[lifnaður] \dicTerm{lif··nað|ur} \dicIPA{{l}{\textsci}{\textsubring{b}}{n}{a}{ð}{\textscy}{\textsubring{r}}} \dicPos{m}[10] \dicFlx{(‑ar)}[7] \dicSynonym{líferni} \dicDirectTranslationCS{způsob života, životní styl}
\dicEntry[lifrarkæfa] \dicTerm{lifrar··kæf|a} \dicIPA{{l}{\textsci}{v}{r}{a}{\textsubring{r}}{c\smash{\textsuperscript{h}}}{a}{i}{v}{a}} \dicPos{f}[1] \dicFlx{(‑u)}[5] \dicFieldCat{kulin.} \dicDirectTranslationCS{játrová paštika, játrovka}
\dicEntry[lifrarpylsa] \dicTerm{lifrar··pyls|a} \dicIPA{{l}{\textsci}{v}{r}{a}{\textsubring{r}}{p\smash{\textsuperscript{h}}}{\textsci}{l}{s}{a}} \dicPos{f}[1] \dicFlx{(‑u, ‑ur)}[13] \dicFieldCat{kulin.} \dicDirectTranslationCS{jitrnice}
\dicEntry[lifur] \dicTerm{lif|ur} \dicIPA{{l}{\textsci}{\textlengthmark}{v}{\textscy}{\textsubring{r}}} \dicPos{f}[6] \dicFlx{(‑rar, ‑rar)}[4] \dicFieldCat{anat.} \dicDirectTranslationCS{játra}
\dicEntry[ligg] \dicTerm{ligg} \dicIPA{{l}{\textsci}{\r{g}}{\textlengthmark}} \dicPos{v} \dicFlx{ind praes sg 1 pers} \dicLink{liggja}
\dicEntry[liggja] \dicTerm{liggja} \dicsymFrequent\  \dicIPA{{l}{\textsci}{\r{\textObardotlessj}}{\textlengthmark}{a}} \dicPos{v}[7] \dicFlx{(ligg, lá, lágum, lægi, legið)}[9] \textbf{1.} \dicDirectTranslationCS{ležet} \dicIndirectTranslationCS{(spočívat na něčem podélně)} \dicExampleIS{Blaðið liggur á gólfinu.} \dicExampleCS{Noviny leží na podlaze.}  \textbf{2.} \dicDirectTranslationCS{ležet} \dicIndirectTranslationCS{(být uložen ke spánku)}  \textbf{3.} \dicDirectTranslationCS{ležet, rozkládat se, prostírat se} \dicExampleIS{Vegurinn liggur meðfram ánni.} \dicExampleCS{Cesta vede podél řeky.};  \dicIdiom{liggja}[að]{ \dicPhraseIS{liggja að e‑u}} \dicDirectTranslationCS{sousedit s~(čím)};  \dicIdiom{liggja}[á]{ \dicPhraseIS{ekkert liggur á}} \dicDirectTranslationCS{všechno popořadě, vše ve správný čas}; { \dicPhraseIS{e‑m liggur á}} \dicFlx{impers} \dicDirectTranslationCS{(kdo) má naspěch, (kdo) pospíchá, (kdo) chvátá} \dicExampleIS{Honum lá á.} \dicExampleCS{Měl naspěch.}; { \dicPhraseIS{e‑m liggur á e‑u}} \dicFlx{impers} \dicDirectTranslationCS{(kdo) má naspěch s~(čím), (komu) záleží na (čem)} \dicExampleIS{Honum liggur mikið á sendingunni.} \dicExampleCS{Na té zásilce mu moc záleží.}; { \dicPhraseIS{liggja á eggjum}} \dicDirectTranslationCS{sedět na vejcích}; { \dicPhraseIS{það liggur vel á e‑m}} \dicFlx{impers} \dicDirectTranslationCS{(kdo) se dobře cítí, (kdo) je v~dobrém rozpoložení};  \dicIdiom{liggja}[fyrir]{ \dicPhraseIS{e‑að á fyrir e‑m að liggja}} \dicDirectTranslationCS{(co) na (koho) čeká, (co) má (kdo) před sebou}; { \dicPhraseIS{e‑að liggur fyrir e‑m}} \dicDirectTranslationCS{(co) na (koho) čeká, (co) leží před (kým) (rozhodnutí ap.)}; { \dicPhraseIS{liggja fyrir}} \dicDirectTranslationCS{ležet v~posteli};  \dicIdiom{liggja}[í]{ \dicPhraseIS{liggja í e‑u}} \dicDirectTranslationCS{ležet s~(čím) (chřipkou ap.)} \dicExampleIS{liggja í flensu} \dicExampleCS{ležet s~chřipkou}; { \dicPhraseIS{e‑að liggur í loftinu}} \dicLangCat{přen.} \dicDirectTranslationCS{(co) visí ve vzduchu, (co) se vznáší ve vzduchu (napětí ap.)}; { \dicPhraseIS{liggja í því}} \dicDirectTranslationCS{prohrát, dostat nakládačku};  \dicIdiom{liggja}[með]{ \dicPhraseIS{liggja með e‑m}} \dicDirectTranslationCS{spát s~(kým)};  \dicIdiom{liggja}[nærri]{ \dicPhraseIS{það lá nærri e‑u}} \dicFlx{impers} \dicDirectTranslationCS{schylovalo se k~(čemu), (co) bylo na spadnutí};  \dicIdiom{liggja}[undir]{ \dicPhraseIS{liggja undir grun}} \dicDirectTranslationCS{čelit podezření, být v~podezření};  \dicIdiom{liggja}[úti]{ \dicPhraseIS{liggja úti}} \dicDirectTranslationCS{spát venku, spát pod širým nebem};  \dicIdiom{liggja}[við]{ \dicPhraseIS{e‑m liggur við e‑u}} \dicFlx{impers} \dicDirectTranslationCS{(kdo) málem (co) (upadl ap.)}; { \dicPhraseIS{e‑að liggur beint við}} \dicDirectTranslationCS{(co) je přirozené}; { \dicPhraseIS{ef\,/\addthin þegar mikið liggur við}} \dicDirectTranslationCS{pokud to je nutné}; { \dicPhraseIS{það lá við e‑u}} \dicFlx{impers} \dicDirectTranslationCS{schylovalo se k~(čemu)};  \dicIdiom{liggja}[yfir]{ \dicPhraseIS{liggja yfir e‑u}} \dicDirectTranslationCS{ležet nad (čím), ležet v~(čem) (v~učení ap.)};  \dicIdiom{liggja}{ \dicPhraseIS{liggja í augum uppi}} \dicDirectTranslationCS{být jasný\,/\addthin zřejmý\,/\addthin očividný, bít do očí}
\dicEntry[lilja] \dicTerm{lilj|a} \dicIPA{{l}{\textsci}{l}{j}{a}} \dicPos{f}[1] \dicFlx{(‑u, ‑ur)}[7] \dicFieldCat{bot.} \dicDirectTranslationCS{lilie} \textit{(l.~{\textLA{Lilium}})}  \dicsymPhoto\ 
\dicFigure{1249.jpg}{Lilja}{Lilja - Zicha Ondřej, Biolib, Copyright/CC-BY-NC}
\dicEntry[lim] \dicTerm{lim} \dicIPA{{l}{\textsci}{\textlengthmark}{\textsubring{m}}} \dicPos{f}[4] \dicFlx{(‑ar, ‑ar)}[1] \dicDirectTranslationCS{listoví}
\dicEntry[limaburður] \dicTerm{lima··burð|ur} \dicIPA{{l}{\textsci}{\textlengthmark}{m}{a}{\textsubring{b}}{\textscy}{r}{ð}{\textscy}{\textsubring{r}}} \dicPos{m}[10] \dicFlx{(‑ar)}[7] \dicDirectTranslationCS{postoj, póza} \dicIndirectTranslationCS{(držení těla při stání)}
\dicEntry[limgerði] \dicTerm{lim··gerði} \dicIPA{{l}{\textsci}{m}{\r{\textObardotlessj}}{\textepsilon}{r}{ð}{\textsci}} \dicPos{n}[2] \dicFlx{(‑s, ‑)}[14] \dicDirectTranslationCS{živý plot}
\dicEntry[limlesta] \dicTerm{lim··lest|a} \dicIPA{{l}{\textsci}{m}{l}{\textepsilon}{s}{\textsubring{d}}{a}} \dicPos{v}[2] \dicFlx{(‑i, ‑)}[12] \dicFlx{acc} \dicDirectTranslationCS{zmrzačit, polámat kosti}
\dicEntry[limra] \dicTerm{limr|a} \dicIPA{{l}{\textsci}{m}{r}{a}} \dicPos{f}[1] \dicFlx{(‑u, ‑ur)}[7] \dicFieldCat{lit.} \dicDirectTranslationCS{limerik} \dicIndirectTranslationCS{(zábavný poetický útvar)}
\dicEntry[limur] \dicTerm{lim|ur} \dicIPA{{l}{\textsci}{\textlengthmark}{m}{\textscy}{\textsubring{r}}} \dicPos{m}[9] \dicFlx{(‑s, ‑ir)}[8] \textbf{1.} \dicSynonym{útlimur} \dicDirectTranslationCS{úd, končetina} \dicExampleIS{limir líkamans} \dicExampleCS{údy těla}  \textbf{2.} \dicSynonym{getnaðarlimur} \dicDirectTranslationCS{pyj}
\dicEntry[lina] \dicTerm{lin|a} \dicIPA{{l}{\textsci}{\textlengthmark}{n}{a}} \dicPos{v}[1] \dicFlx{(‑aði)}[1] \dicFlx{acc} \dicDirectTranslationCS{(z)mírnit, (z)tlumit, (z)měkčit} \dicExampleIS{Kærleikurinn linar sorg hennar.} \dicExampleCS{Láska mírní její žal.};  \dicIdiom{linast}{ \dicPhraseIS{linast}} \dicFlx{refl} {\textbf{a.}} \dicSynonym*{mýkjast} \dicDirectTranslationCS{(z)měknout} \dicIndirectTranslationCS{(stát\,/\addthin stávat se měkkým)};  {\textbf{b.}} \dicSynonym*{slá af fyrri skoðun} \dicDirectTranslationCS{obměkčit se, ustoupit};  {\textbf{c.}} \dicSynonym*{missa krafta} \dicDirectTranslationCS{(ze)slábnout, ztratit\,/\addthin ztrácet sílu}
\dicEntry[lind] \dicTerm{lind} \dicsymFrequent\  \dicIPA{{l}{\textsci}{n}{\textsubring{d}}} \dicPos{f}[7] \dicFlx{(‑ar, ‑ir)}[1] \textbf{1.} \dicSynonym{uppspretta} \dicDirectTranslationCS{zřídlo, pramen} \dicExampleIS{Hún vísaði þyrstum hermönnum á lind.} \dicExampleCS{Ukázala žíznivým vojákům cestu k~pramenu.}  \textbf{2.} \dicLangCat{přen.} \dicSynonym{upptök} \dicDirectTranslationCS{zřídlo, původ}  \textbf{3.} \dicFieldCat{bot.} \dicSynonym{linditré} \dicDirectTranslationCS{lípa} \textit{(l.~{\textLA{Tilia}})}  \dicsymPhoto\ 
\dicFigure{16367.jpg}{Lind}{Lind - Hanzlík Václav, Biolib, Copyright/CC-BY-NC}
\dicEntry[lindá] \dicTerm{lind··á} \dicIPA{{l}{\textsci}{n}{\textsubring{d}}{au}} \dicPos{f}[4] \dicFlx{(‑r, ‑r)}[18] \dicFieldCat{geol.} \dicDirectTranslationCS{řeka napájená zřídly}
\dicEntry[lindi] \dicTerm{lind|i} \dicIPA{{l}{\textsci}{n}{\textsubring{d}}{\textsci}} \dicPos{m}[1] \dicFlx{(‑a, ‑ar)}[1] \dicSynonym{band} \dicDirectTranslationCS{opasek, pás}
\dicEntry[linditré] \dicTerm{lindi··tré} \dicIPA{{l}{\textsci}{n}{\textsubring{d}}{\textsci}{t\smash{\textsuperscript{h}}}{r}{j}{\textepsilon}} \dicPos{n}[2] \dicFlx{(‑s, ‑)}[36] \dicFieldCat{bot.} \dicSynonym{lind} \dicDirectTranslationCS{lípa} \textit{(l.~{\textLA{Tilia}})}
\dicEntry[lindýr] \dicTerm{lin··dýr} \dicIPA{{l}{\textsci}{n}{\textsubring{d}}{i}{\textsubring{r}}} \dicPos{n}[2] \dicFlx{(‑s, ‑)}[5] \dicFieldCat{zool.} \dicDirectTranslationCS{měkkýš} \textit{(l.~{\textLA{Mollusca}})}
\dicEntry[lingerður] \dicTerm{lin··gerður} \dicIPA{{l}{\textsci}{n}{\r{\textObardotlessj}}{\textepsilon}{r}{ð}{\textscy}{\textsubring{r}}} \dicPos{adj}[2]\dicFlx{}[-1] \textbf{1.} \dicSynonym*{kveifarlegur} \dicDirectTranslationCS{změkčilý}  \textbf{2.} \dicSynonym{veikburða} \dicDirectTranslationCS{zesláblý, neduživý}
\dicEntry[linka] \dicTerm{link|a} \dicIPA{{l}{i}{\r{\ng}}{\r{g}}{a}} \dicPos{f}[1] \dicFlx{(‑u)}[5] \textbf{1.} \dicSynonym{lasleiki} \dicDirectTranslationCS{churavost, neduživost}  \textbf{2.} \dicSynonym{dugleysi} \dicDirectTranslationCS{neschopnost, nezpůsobilost}
\dicEntry[linmæli] \dicTerm{lin··mæli} \dicIPA{{l}{\textsci}{n}{m}{a}{i}{l}{\textsci}} \dicPos{n}[2] \dicFlx{(‑s)}[20] \dicFieldCat{jaz.} \dicIndirectTranslationCS{druh výslovnosti -- používání neaspirované plozivy ([\textsubring{b}], [\textsubring{d}], [\r{g}], [\r{\textObardotlessj}]) uprostřed slova po dlouhé samohlásce (standardní výslovnost)} \dicAntonym{harðmæli}
\dicEntry[linna] \dicTerm{linn|a} \dicIPA{{l}{\textsci}{n}{\textlengthmark}{a}} \dicPos{v}[2] \dicFlx{(‑ti, ‑t)}[80] \dicPhraseIS{e‑u linnir} \dicFlx{impers} \dicDirectTranslationCS{(co) přestává\,/\addthin ustává\,/\addthin končí} \dicExampleIS{Óveðrinu er tekið að linna.} \dicExampleCS{Bouře pomalu ustává.};  \dicPhraseIS{linna ekki látum} \dicDirectTranslationCS{nepřestat, neustat (s~prosbami ap.)}
\dicEntry[linnulaus] \dicTerm{linnu··laus} \dicIPA{{l}{\textsci}{n}{\textlengthmark}{\textscy}{l}{\oe i}{s}} \dicPos{adj}[5]\dicFlx{}[-1] \dicSynonym{sífelldur} \dicDirectTranslationCS{nepřetržitý, neustálý, trvalý} \dicExampleIS{linnulausar árásir} \dicExampleCS{nepřetržité útoky}
\dicEntry[linsa] \dicTerm{lins|a} \dicIPA{{l}{\textsci}{n}{s}{a}} \dicPos{f}[1] \dicFlx{(‑u, ‑ur)}[7] \textbf{1.} \dicFieldCat{fyz.} \dicDirectTranslationCS{čočka}  \textbf{2.} \dicSynonym*{augnlinsa} \dicDirectTranslationCS{(kontaktní) čočka}
\dicEntry[linsoðinn] \dicTerm{lin··soðinn} \dicIPA{{l}{\textsci}{n}{s}{\textopeno}{ð}{\textsci}{\textsubring{n}}} \dicPos{adj}[6]\dicFlx{}[-2] \dicDirectTranslationCS{(jsoucí) naměkko} \dicExampleIS{linsoðið egg} \dicExampleCS{vejce uvařené naměkko}
\dicEntry[linsubaun] \dicTerm{linsu··baun} \dicIPA{{l}{\textsci}{n}{s}{\textscy}{\textsubring{b}}{\oe i}{\textsubring{n}}} \dicPos{f}[7] \dicFlx{(‑ar, ‑ir)}[1] \dicFieldCat{bot.} \dicDirectTranslationCS{čočka} \textit{(l.~{\textLA{Lens culinaris}})}
\dicEntry[linun] \dicTerm{lin|un} \dicIPA{{l}{\textsci}{\textlengthmark}{n}{\textscy}{\textsubring{n}}} \dicPos{f}[7] \dicFlx{(‑unar)}[9] \textbf{1.} \dicSynonym{fró} \dicDirectTranslationCS{uklidnění, úleva}  \textbf{2.} \dicSynonym{mildun} \dicDirectTranslationCS{zmírnění} \dicExampleIS{linun refsingar} \dicExampleCS{zmírnění trestu}
\dicEntry[linur] \dicTerm{linur} \dicIPA{{l}{\textsci}{\textlengthmark}{n}{\textscy}{\textsubring{r}}} \dicPos{adj}[1]\dicFlx{}[-1] \textbf{1.} \dicSynonym{mjúkur} \dicDirectTranslationCS{hebký, měkký, jemný}  \textbf{2.} \dicSynonym{slakur} \dicDirectTranslationCS{slabý, chabý} \dicExampleIS{linur í fótbolta} \dicExampleCS{slabý ve fotbale}  \textbf{3.} \dicSynonym{slappur} \dicDirectTranslationCS{indisponovaný, slabý}
\dicEntry[lipur] \dicTerm{lipur} \dicIPA{{l}{\textsci}{\textlengthmark}{\textsubring{b}}{\textscy}{\textsubring{r}}} \dicPos{adj}[9] \dicFlx{(f ‑)}[1] \textbf{1.} \dicSynonym*{flinkur} \dicDirectTranslationCS{zručný, šikovný} \dicExampleIS{Hún er lipur í höndunum.} \dicExampleCS{Má šikovné ruce.}  \textbf{2.} \dicSynonym{þægilegur} \dicDirectTranslationCS{vstřícný, přívětivý}  \textbf{3.} \dicSynonym{meðfærilegur} \dicDirectTranslationCS{(uživatelsky) přívětivý}
\dicEntry[lipurð] \dicTerm{lipurð} \dicIPA{{l}{\textsci}{\textlengthmark}{\textsubring{b}}{\textscy}{r}{\texttheta}} \dicPos{f}[4] \dicFlx{(‑ar)}[3] \textbf{1.} \dicSynonym*{fimleiki} \dicDirectTranslationCS{obratnost, zručnost}  \textbf{2.} \dicSynonym*{lipur í viðmóti} \dicDirectTranslationCS{vstřícnost, přívětivost}
\dicEntry[lirfa] \dicTerm{lirf|a} \dicIPA{{l}{\textsci}{r}{v}{a}} \dicPos{f}[1] \dicFlx{(‑u, ‑ur)}[19] \dicDirectTranslationCS{larva}
\dicEntry[list] \dicTerm{list} \dicsymFrequent\  \dicIPA{{l}{\textsci}{s}{\textsubring{d}}} \dicPos{f}[7] \dicFlx{(‑ar, ‑ir)}[1] \textbf{1.} \dicDirectTranslationCS{umění} \dicIndirectTranslationCS{(odvětví lidské tvorby)} \dicExampleIS{unna listum} \dicExampleCS{zbožňovat umění}  \textbf{2.} \dicSynonym{færni} \dicDirectTranslationCS{umění, um, dovednost}
\dicEntry[listaháskóli] \dicTerm{lista··há·skól|i} \dicIPA{{l}{\textsci}{s}{\textsubring{d}}{a}{h}{au}{s}{\r{g}}{ou}{l}{\textsci}} \dicPos{m}[1] \dicFlx{(‑a, ‑ar)}[1] \dicFieldCat{škol.} \dicDirectTranslationCS{akademie umění}
\dicEntry[listahátíð] \dicTerm{lista··há·tíð} \dicIPA{{l}{\textsci}{s}{\textsubring{d}}{a}{h}{au}{t\smash{\textsuperscript{h}}}{i}{\texttheta}} \dicPos{f}[7] \dicFlx{(‑ar, ‑ir)}[1] \dicDirectTranslationCS{umělecký festival}
\dicEntry[listakona] \dicTerm{lista··kon|a} \dicIPA{{l}{\textsci}{s}{\textsubring{d}}{a}{k\smash{\textsuperscript{h}}}{\textopeno}{n}{a}} \dicPos{f}[1] \dicFlx{(‑u, ‑ur)}[27] \dicDirectTranslationCS{umělkyně}
\dicEntry[listamaður] \dicTerm{lista··|maður} \dicsymFrequent\  \dicIPA{{l}{\textsci}{s}{\textsubring{d}}{a}{m}{a}{ð}{\textscy}{\textsubring{r}}} \dicPos{m}[13] \dicFlx{(‑manns, ‑menn)}[2] \dicDirectTranslationCS{umělec, umělkyně} \dicExampleIS{Það er eitt dáðasta verk listamannsins.} \dicExampleCS{Je to jedno z~nejvíce obdivovaných děl umělce.}
\dicEntry[listasafn] \dicTerm{lista··|safn} \dicIPA{{l}{\textsci}{s}{\textsubring{d}}{a}{s}{a}{\textsubring{b}}{\textsubring{n}}} \dicPos{n}[2] \dicFlx{(‑safns, ‑söfn)}[8] \dicDirectTranslationCS{galerie (umění)}
\dicEntry[listasaga] \dicTerm{lista··|saga} \dicIPA{{l}{\textsci}{s}{\textsubring{d}}{a}{s}{a}{\textbabygamma}{a}} \dicPos{f}[1] \dicFlx{(‑sögu)}[2] \dicDirectTranslationCS{dějiny umění}
\dicEntry[listaskóli] \dicTerm{lista··skól|i} \dicIPA{{l}{\textsci}{s}{\textsubring{d}}{a}{s}{\r{g}}{ou}{l}{\textsci}} \dicPos{m}[1] \dicFlx{(‑a, ‑ar)}[1] \dicFieldCat{škol.} \dicDirectTranslationCS{umělecká škola}
\dicEntry[listaverk] \dicTerm{lista··verk} \dicIPA{{l}{\textsci}{s}{\textsubring{d}}{a}{v}{\textepsilon}{\textsubring{r}}{\r{g}}} \dicPos{n}[2] \dicFlx{(‑s, ‑)}[5] \dicDirectTranslationCS{umělecké dílo}
\dicEntry[listdans] \dicTerm{list··dans} \dicIPA{{l}{\textsci}{s}{\textsubring{d}}{a}{n}{s}} \dicPos{m}[4] \dicFlx{(‑, ‑ar)}[29] \dicSynonym{ballett} \dicDirectTranslationCS{umělecký tanec, balet}
\dicEntry[listfengur] \dicTerm{list··fengur} \dicIPA{{l}{\textsci}{s}{\textsubring{d}}{f}{ei}{\ng}{\r{g}}{\textscy}{\textsubring{r}}} \dicPos{adj}[1]\dicFlx{}[-1] \dicDirectTranslationCS{umělecky nadaný}
\dicEntry[listfræðingur] \dicTerm{list·fræð··ing|ur} \dicIPA{{l}{\textsci}{s}{\textsubring{d}}{f}{r}{a}{i}{ð}{i}{\ng}{\r{g}}{\textscy}{\textsubring{r}}} \dicPos{m}[6] \dicFlx{(‑s, ‑ar)}[8] \dicDirectTranslationCS{historik\,/\addthin historička umění, kunsthistorik, kunsthistorička}
\dicEntry[listgáfa] \dicTerm{list··gáf|a} \dicIPA{{l}{\textsci}{s}{\textsubring{d}}{\r{g}}{au}{v}{a}} \dicPos{f}[1] \dicFlx{(‑u, ‑ur)}[13] \dicDirectTranslationCS{umělecký talent}
\dicEntry[listi] \dicTerm{list|i} \dicsymFrequent\  \dicIPA{{l}{\textsci}{s}{\textsubring{d}}{\textsci}} \dicPos{m}[1] \dicFlx{(‑a, ‑ar)}[1] \textbf{1.} \dicSynonym{skrá\smash{\textsuperscript{1}}} \dicDirectTranslationCS{seznam, soupis, listina, rejstřík, žebříček} \dicExampleIS{óskalisti} \dicExampleCS{seznam přání};  \dicPhraseIS{listi yfir e‑n\,/\addthin e‑ð} \dicDirectTranslationCS{seznam (koho\,/\addthin čeho)}  \textbf{2.} \dicSynonym{ræma} \dicDirectTranslationCS{lišta, laťka}
\dicEntry[listiðn] \dicTerm{list··iðn} \dicIPA{{l}{\textsci}{s}{\textsubring{d}}{\textsci}{ð}{\textsubring{n}}} \dicPos{f}[7] \dicFlx{(‑ar, ‑ir)}[1] \dicDirectTranslationCS{umělecké řemeslo}
\dicEntry[listiðnaður] \dicTerm{list··iðn·að|ur} \dicIPA{{l}{\textsci}{s}{\textsubring{d}}{\textsci}{ð}{n}{a}{ð}{\textscy}{\textsubring{r}}} \dicPos{m}[10] \dicFlx{(‑ar)}[9] \dicDirectTranslationCS{užité umění}
\dicEntry[listmálari] \dicTerm{list··mál·ar|i} \dicIPA{{l}{\textsci}{s}{\textsubring{d}}{m}{au}{l}{a}{r}{\textsci}} \dicPos{m}[1] \dicFlx{(‑a, ‑ar)}[13] \dicDirectTranslationCS{malíř(ka)} \dicIndirectTranslationCS{(umělec malující obrazy)}
\dicEntry[listnám] \dicTerm{list··nám} \dicIPA{{l}{\textsci}{s}{\textsubring{d}}{n}{au}{\textsubring{m}}} \dicPos{n}[2] \dicFlx{(‑s)}[2] \dicDirectTranslationCS{studium umění}
\dicEntry[listrænn] \dicTerm{list··rænn} \dicIPA{{l}{\textsci}{s}{\textsubring{d}}{r}{a}{i}{\textsubring{d}}{\textsubring{n}}} \dicPos{adj}[7]\dicFlx{}[-1] \dicDirectTranslationCS{umělecký}
\dicEntry[listskautar] \dicTerm{list··skautar} \dicIPA{{l}{\textsci}{s}{\textsubring{d}}{s}{\r{g}}{\oe i}{\textsubring{d}}{a}{\textsubring{r}}} \dicPos{m}[1] \dicFlx{pl}[2] \dicFieldCat{sport.} \dicDirectTranslationCS{krasobruslení}
\dicEntry[listunnandi] \dicTerm{list··unn·|andi} \dicIPA{{l}{\textsci}{s}{\textsubring{d}}{\textscy}{n}{a}{n}{\textsubring{d}}{\textsci}} \dicPos{m}[2] \dicFlx{(‑anda, ‑endur)}[1] \dicDirectTranslationCS{milovník\,/\addthin milovnice umění}
\dicEntry[lita] \dicTerm{lit|a} \dicsymFrequent\  \dicIPA{{l}{\textsci}{\textlengthmark}{\textsubring{d}}{a}} \dicPos{v}[1] \dicFlx{(‑aði)}[1] \dicFlx{acc} \textbf{1.} \dicSynonym{mála} \dicDirectTranslationCS{(na)barvit, nabarvovat, obarvit, obarvovat} \dicExampleIS{lita fötin} \dicExampleCS{nabarvit oblečení}  \textbf{2.} \dicDirectTranslationCS{zabarvit, zkreslit (skutečnost ap.)} \dicExampleIS{Skáldskapurinn litar frá sér.} \dicExampleCS{Literární dílo zkresluje.};  \dicIdiom{litast}{ \dicPhraseIS{litast}} \dicFlx{refl} \dicDirectTranslationCS{zbarvit se, zabarvit se, obarvit se};  \dicIdiom{litast}[um]{ \dicPhraseIS{litast um}} \dicFlx{refl} \dicDirectTranslationCS{(po)rozhlédnout se};  \dicIdiom{litast}[upp]{ \dicPhraseIS{litast upp}} \dicFlx{refl} \dicSynonym*{tapa lit} \dicDirectTranslationCS{(vy)blednout, ztratit\,/\addthin ztrácet barvu}
\dicEntry[litaraft] \dicTerm{litar··aft} \dicIPA{{l}{\textsci}{\textlengthmark}{\textsubring{d}}{a}{r}{a}{f}{\textsubring{d}}} \dicPos{n}[2] \dicFlx{(‑s)}[2] \dicDirectTranslationCS{barva\,/\addthin zbarvení obličeje}
\dicEntry[litarháttur] \dicTerm{litar··hátt|ur} \dicIPA{{l}{\textsci}{\textlengthmark}{\textsubring{d}}{a}{\textsubring{r}}{h}{au}{h}{\textsubring{d}}{\textscy}{\textsubring{r}}} \dicPos{m}[12] \dicFlx{(‑ar)}[8] \dicDirectTranslationCS{barva pleti} \dicExampleIS{fordómar vegna litarháttar} \dicExampleCS{předsudky k~barvě pleti}
\dicEntry[litasamsetning] \dicTerm{lita··sam·set·ning} \dicIPA{{l}{\textsci}{\textlengthmark}{\textsubring{d}}{a}{s}{a}{m}{s}{\textepsilon}{h}{\textsubring{d}}{n}{i}{\ng}{\r{g}}} \dicPos{f}[4] \dicFlx{(‑ar, ‑ar)}[5] \dicDirectTranslationCS{barevná kombinace, barevné schéma}
\dicEntry[litaspjald] \dicTerm{lita··|spjald} \dicIPA{{l}{\textsci}{\textlengthmark}{\textsubring{d}}{a}{s}{\textsubring{b}}{j}{a}{l}{\textsubring{d}}} \dicPos{n}[2] \dicFlx{(‑spjalds, ‑spjöld)}[8] \textbf{1.} \dicDirectTranslationCS{paleta} \dicIndirectTranslationCS{(malířská deska k~míchání a~roztírání barev)}  \textbf{2.} \dicFieldCat{poč.} \dicDirectTranslationCS{paleta} \dicIndirectTranslationCS{(na výběr barev)}
\dicEntry[litblinda] \dicTerm{lit··blind|a} \dicIPA{{l}{\textsci}{\textlengthmark}{\textsubring{d}}{\textsubring{b}}{l}{\textsci}{n}{\textsubring{d}}{a}} \dicPos{f}[1] \dicFlx{(‑u)}[5] \dicDirectTranslationCS{barvoslepost}
\dicEntry[litblindur] \dicTerm{lit··blindur} \dicIPA{{l}{\textsci}{\textlengthmark}{\textsubring{d}}{\textsubring{b}}{l}{\textsci}{n}{\textsubring{d}}{\textscy}{\textsubring{r}}} \dicPos{adj}[2]\dicFlx{}[-14] \dicFieldCat{med.} \dicDirectTranslationCS{barvoslepý}
\dicEntry[litbrigði] \dicTerm{lit··brigði} \dicIPA{{l}{\textsci}{\textlengthmark}{\textsubring{d}}{\textsubring{b}}{r}{\textsci}{\textbabygamma}{ð}{\textsci}} \dicPos{n}[2] \dicFlx{(‑s, ‑)}[14] \dicDirectTranslationCS{odstín (barvy)}
\dicEntry[litgreining] \dicTerm{lit··grein·ing} \dicIPA{{l}{\textsci}{\textlengthmark}{\textsubring{d}}{\r{g}}{r}{ei}{n}{i}{\ng}{\r{g}}} \dicPos{f}[4] \dicFlx{(‑ar)}[7] \textbf{1.} \dicFieldCat{polygr.} \dicDirectTranslationCS{separace barev}  \textbf{2.} \dicDirectTranslationCS{barevná analýza} \dicIndirectTranslationCS{(osobní kombinace barev)}
\dicEntry[Litháen] \dicTerm{Litháen} \dicIPA{{l}{\textsci}{\textlengthmark}{t\smash{\textsuperscript{h}}}{au}{\textepsilon}{\textsubring{n}}} \dicPos{n}[2] \dicFlx{(‑s)}[32] \dicFieldCat{geog.} \dicDirectTranslationCS{Litva}
\dicEntry[Lithái] \dicTerm{Lithá|i} \dicIPA{{l}{\textsci}{\textlengthmark}{t\smash{\textsuperscript{h}}}{au}{\textsci}} \dicPos{m}[1] \dicFlx{(‑a, ‑ar)}[1] \dicDirectTranslationCS{Litevec, Litevka}
\dicEntry[litháíska] \dicTerm{litháísk|a} \dicIPA{{l}{\textsci}{\textlengthmark}{t\smash{\textsuperscript{h}}}{au}{i}{s}{\r{g}}{a}} \dicPos{f}[1] \dicFlx{(‑u)}[5] \dicLink{litháska}
\dicEntry[litháískur] \dicTerm{litháískur} \dicIPA{{l}{\textsci}{\textlengthmark}{t\smash{\textsuperscript{h}}}{au}{i}{s}{\r{g}}{\textscy}{\textsubring{r}}} \dicPos{adj}[1]\dicFlx{}[-6] \dicLink{litháskur}
\dicEntry[litháska] \dicTerm{lithásk|a}\dicTerm{, litháíska} \dicIPA{{l}\-{\textsci}\-{\textlengthmark}\-{t\smash{\textsuperscript{h}}}\-{au}\-{s}\-{\r{g}}\-{a}\-} \dicPos{f}[1] \dicFlx{(‑u)}[5] \dicDirectTranslationCS{litevština}
\dicEntry[litháskur] \dicTerm{litháskur}\dicTerm{, litháískur} \dicIPA{{l}\-{\textsci}\-{\textlengthmark}\-{t\smash{\textsuperscript{h}}}\-{au}\-{s}\-{\r{g}}\-{\textscy}\-{\textsubring{r}}\-} \dicPos{adj}[1]\dicFlx{}[-6] \dicDirectTranslationCS{litevský}
\dicEntry[lithimna] \dicTerm{lit··himn|a} \dicIPA{{l}{\textsci}{\textlengthmark}{\textsubring{d}}{h}{\textsci}{m}{n}{a}} \dicPos{f}[1] \dicFlx{(‑u, ‑ur)}[7] \dicFieldCat{anat.} \dicDirectTranslationCS{duhovka}
\dicEntry[liti] \dicTerm{liti} \dicIPA{{l}{\textsci}{\textlengthmark}{\textsubring{d}}{\textsci}} \dicPos{v} \dicFlx{con pf sg 1 pers} \dicLink{líta}
\dicEntry[litið] \dicTerm{litið} \dicIPA{{l}{\textsci}{\textlengthmark}{\textsubring{d}}{\textsci}{\texttheta}} \dicPos{v} \dicFlx{supin} \dicLink{líta}
\dicEntry[litlafingurs] \dicTerm{litla··fingurs} \dicIPA{{l}{\textsci}{h}{\textsubring{d}}{l}{a}{f}{\textsci}{\ng}{\r{g}}{\textscy}{\textsubring{r}}{s}} \dicPos{m} \dicFlx{sg gen} \dicLink{litlifingur}
\dicEntry[litlatá] \dicTerm{litla··tá} \dicIPA{{l}{\textsci}{h}{\textsubring{d}}{l}{a}{t\smash{\textsuperscript{h}}}{au}} \dicPos{f}[9] \dicFlx{(litlutáar, litlutær)}[2] \dicDirectTranslationCS{malík\,/\addthin malíček (u~nohy)}
\dicEntry[litlaus] \dicTerm{lit··laus} \dicIPA{{l}{\textsci}{\textlengthmark}{\textsubring{d}}{l}{\oe i}{s}} \dicPos{adj}[5]\dicFlx{}[-1] \textbf{1.} \dicSynonym*{ekki í lit} \dicDirectTranslationCS{bezbarvý}  \textbf{2.} \dicSynonym{dauflegur} \dicDirectTranslationCS{mdlý, fádní, bezbarvý} \dicExampleIS{litlaus frásögn} \dicExampleCS{bezbarvé vyprávění}
\dicEntry[litlifingur] \dicTerm{litli··fingur} \dicIPA{{l}{\textsci}{h}{\textsubring{d}}{l}{\textsci}{f}{i}{\ng}{\r{g}}{\textscy}{\textsubring{r}}} \dicPos{m}[5] \dicFlx{(litlafingurs, litlufingur)}[14] \dicDirectTranslationCS{malík\,/\addthin malíček (u~ruky)}
\dicEntry[litlufingur] \dicTerm{litlu··fingur} \dicIPA{{l}{\textsci}{h}{\textsubring{d}}{l}{\textscy}{f}{\textsci}{\ng}{\r{g}}{\textscy}{\textsubring{r}}} \dicPos{m} \dicFlx{pl nom} \dicLink{litlifingur}
\dicEntry[litmynd] \dicTerm{lit··mynd} \dicIPA{{l}{\textsci}{\textlengthmark}{\textsubring{d}}{m}{\textsci}{n}{\textsubring{d}}} \dicPos{f}[7] \dicFlx{(‑ar, ‑ir)}[1] \dicDirectTranslationCS{barevný obraz, barevná fotografie}
\dicEntry[litningur] \dicTerm{litn··ing|ur} \dicIPA{{l}{\textsci}{h}{\textsubring{d}}{n}{i}{\ng}{\r{g}}{\textscy}{\textsubring{r}}} \dicPos{m}[6] \dicFlx{(‑s, ‑ar)}[8] \dicFieldCat{biol.} \dicDirectTranslationCS{chromozom}
\dicEntry[litprentun] \dicTerm{lit··prent|un} \dicIPA{{l}{\textsci}{\textlengthmark}{\textsubring{d}}{p\smash{\textsuperscript{h}}}{r}{\textepsilon}{\textsubring{n}}{\textsubring{d}}{\textscy}{\textsubring{n}}} \dicPos{f}[7] \dicFlx{(‑unar)}[9] \dicFieldCat{polygr.} \dicDirectTranslationCS{barevný tisk}
\dicEntry[litríkur] \dicTerm{lit··ríkur} \dicIPA{{l}{\textsci}{\textlengthmark}{\textsubring{d}}{r}{i}{\r{g}}{\textscy}{\textsubring{r}}} \dicPos{adj}[1]\dicFlx{}[-1] \textbf{1.} \dicDirectTranslationCS{pestrobarevný, pestrý} \dicExampleIS{litrík föt} \dicExampleCS{pestrobarevné oblečení}  \textbf{2.} \dicDirectTranslationCS{pestrý, rozmanitý} \dicExampleIS{Söngkonan á litríkan feril að baki.} \dicExampleCS{Zpěvačka má za sebou pestrou kariéru.}
\dicEntry[litróf] \dicTerm{lit··róf} \dicIPA{{l}{\textsci}{\textlengthmark}{\textsubring{d}}{r}{ou}{f}} \dicPos{n}[2] \dicFlx{(‑s, ‑)}[5] \textbf{1.} \dicFieldCat{fyz.} \dicDirectTranslationCS{barevné spektrum}  \textbf{2.} \dicLangCat{přen.} \dicDirectTranslationCS{spektrum (názorové ap.)}
\dicEntry[litskrúð] \dicTerm{lit··skrúð} \dicIPA{{l}{\textsci}{\textlengthmark}{\textsubring{d}}{s}{\r{g}}{r}{u}{\texttheta}} \dicPos{n}[2] \dicFlx{(‑s)}[2] \dicDirectTranslationCS{pestrobarevnost}
\dicEntry[litskyggna] \dicTerm{lit··skyggn|a} \dicIPA{{l}{\textsci}{\textlengthmark}{\textsubring{d}}{s}{\r{\textObardotlessj}}{\textsci}{\r{g}}{n}{a}} \dicPos{f}[1] \dicFlx{(‑u, ‑ur)}[7] \dicDirectTranslationCS{barevný diapozitiv}
\dicEntry[litum] \dicTerm{litum} \dicIPA{{l}{\textsci}{\textlengthmark}{\textsubring{d}}{\textscy}{\textsubring{m}}} \dicPos{v} \dicFlx{ind pf pl 1 pers} \dicLink{líta}
\dicEntry[litun] \dicTerm{lit|un} \dicIPA{{l}{\textsci}{\textlengthmark}{\textsubring{d}}{\textscy}{\textsubring{n}}} \dicPos{f}[7] \dicFlx{(‑unar)}[9] \dicDirectTranslationCS{(na)barvení}
\dicEntry[litur] \dicTerm{lit|ur} \dicsymFrequent\  \dicIPA{{l}{\textsci}{\textlengthmark}{\textsubring{d}}{\textscy}{\textsubring{r}}} \dicPos{m}[10] \dicFlx{(‑ar, ‑ir)}[13] \textbf{1.} \dicDirectTranslationCS{barva, barevný odstín} \dicExampleIS{Hvernig er það á litinn?} \dicExampleCS{Jakou to má barvu?}  \textbf{2.} \dicSynonym*{litarefni} \dicDirectTranslationCS{barva, nátěr}  \textbf{3.} \dicSynonym*{tegund spila} \dicDirectTranslationCS{barva} \dicIndirectTranslationCS{(v~kartách)}
\dicEntry[Líbani] \dicTerm{Líban|i} \dicIPA{{l}{i}{\textlengthmark}{\textsubring{b}}{a}{n}{\textsci}} \dicPos{m}[3] \dicFlx{(‑a, ‑ar\,/\addthin ‑ir)}[3] \dicDirectTranslationCS{Libanonec, Libanonka}
\dicEntry[Líbanon] \dicTerm{Líbanon} \dicIPA{{l}{i}{\textlengthmark}{\textsubring{b}}{a}{n}{\textopeno}{\textsubring{n}}} \dicPos{n}[2] \dicFlx{(‑s)}[32] \dicFieldCat{geog.} \dicDirectTranslationCS{Libanon}
\dicEntry[líbanskur] \dicTerm{líbanskur} \dicIPA{{l}{i}{\textlengthmark}{\textsubring{b}}{a}{n}{s}{\r{g}}{\textscy}{\textsubring{r}}} \dicPos{adj}[1] \dicFlx{(f líbönsk)}[3] \dicDirectTranslationCS{libanonský}
\dicEntry[Líbería] \dicTerm{Líberí|a} \dicIPA{{l}{i}{\textlengthmark}{\textsubring{b}}{\textepsilon}{r}{i}{j}{a}} \dicPos{f}[1] \dicFlx{(‑u)}[6] \dicFieldCat{geog.} \dicDirectTranslationCS{Libérie}
\dicEntry[líberískur] \dicTerm{líberískur} \dicIPA{{l}{i}{\textlengthmark}{\textsubring{b}}{\textepsilon}{r}{i}{s}{\r{g}}{\textscy}{\textsubring{r}}} \dicPos{adj}[1]\dicFlx{}[-6] \dicDirectTranslationCS{liberijský}
\dicEntry[Líberíumaður] \dicTerm{Líberíu··|maður} \dicIPA{{l}{i}{\textlengthmark}{\textsubring{b}}{\textepsilon}{r}{i}{j}{\textscy}{m}{a}{ð}{\textscy}{\textsubring{r}}} \dicPos{m}[13] \dicFlx{(‑manns, ‑menn)}[2] \dicDirectTranslationCS{Liberijec, Liberijka}
\dicEntry[Líbía] \dicTerm{Líbí|a}\dicTerm{, Líbýa} \dicIPA{{l}\-{i}\-{\textlengthmark}\-{\textsubring{b}}\-{i}\-{j}\-{a}\-} \dicPos{f}[1] \dicFlx{(‑u)}[6] \dicFieldCat{geog.} \dicDirectTranslationCS{Libye}
\dicEntry[líbískur] \dicTerm{líbískur}\dicTerm{, líbýskur} \dicIPA{{l}\-{i}\-{\textlengthmark}\-{\textsubring{b}}\-{i}\-{s}\-{\r{g}}\-{\textscy}\-{\textsubring{r}}\-} \dicPos{adj}[1]\dicFlx{}[-6] \dicDirectTranslationCS{libyjský}
\dicEntry[Líbíumaður] \dicTerm{Líbíu··|maður} \dicIPA{{l}{i}{\textlengthmark}{\textsubring{b}}{i}{j}{\textscy}{m}{a}{ð}{\textscy}{\textsubring{r}}} \dicPos{m}[13] \dicFlx{(‑manns, ‑menn)}[2] \dicDirectTranslationCS{Libyjec, Libyjka}
\dicEntry[Líbýa] \dicTerm{Líbý|a} \dicIPA{{l}{i}{\textlengthmark}{\textsubring{b}}{i}{j}{a}} \dicPos{f}[1] \dicFlx{(‑u)}[6] \dicLink{Líbía}
\dicEntry[líbýskur] \dicTerm{líbýskur} \dicIPA{{l}{i}{\textlengthmark}{\textsubring{b}}{i}{s}{\r{g}}{\textscy}{\textsubring{r}}} \dicPos{adj}[1]\dicFlx{}[-6] \dicLink{líbískur}
\dicEntry[líð] \dicTerm{líð} \dicIPA{{l}{i}{\textlengthmark}{\texttheta}} \dicPos{v} \dicFlx{ind praes sg 1 pers} \dicLink{líða}
\dicEntry[líða] \dicTerm{líða} \dicsymFrequent\  \dicIPA{{l}{i}{\textlengthmark}{ð}{a}} \dicPos{v}[6] \dicFlx{(líð, leið, liðum, liði, liðið)}[66] \dicFlx{acc\,/\addthin dat} \textbf{1.} \dicSynonym*{um tíma} \dicDirectTranslationCS{plynout, ubíhat, utíkat} \dicIndirectTranslationCS{(o~čase)} \dicExampleIS{Árið líður hratt.} \dicExampleCS{Rok rychle ubíhá.}  \textbf{2.} \dicFlx{dat} \dicDirectTranslationCS{pokračovat, postupovat} \dicExampleIS{Hvernig líður verkinu?} \dicExampleCS{Jak pokračuje práce?}  \textbf{3.} \dicFlx{acc} \dicSynonym{þola} \dicDirectTranslationCS{snést, snášet, tolerovat} \dicIndirectTranslationCS{(často se záporem)} \dicExampleIS{líða þetta ekki} \dicExampleCS{nesnést to}  \textbf{4.} \dicPhraseIS{e‑m líður} \dicFlx{impers} \dicDirectTranslationCS{(kdo) se cítí (dobře ap.)} \dicExampleIS{Mér líður vel.} \dicExampleCS{Cítím se dobře.\,/\addthin Je mi dobře.};  \dicPhraseIS{hvernig líður þér?} \dicDirectTranslationCS{jak se máš?\,/\addthin jak se cítíš?};  \dicIdiom{líða}[að]{ \dicPhraseIS{það líður að e‑u}} \dicFlx{impers} \dicDirectTranslationCS{(co) se blíží (svátky ap.)};  \dicIdiom{líða}[hjá]{ \dicPhraseIS{e‑að líður hjá}} \dicDirectTranslationCS{(co) míjí, (co) pomine};  \dicIdiom{líða}[í]{ \dicPhraseIS{líða í ómegin}} \dicDirectTranslationCS{omdlít, upadnout do mdlob};  \dicIdiom{líða}[undir]{ \dicPhraseIS{e‑að líður undir lok}} \dicDirectTranslationCS{(co) se schyluje\,/\addthin chýlí ke konci};  \dicIdiom{líða}[út af]{ \dicPhraseIS{líða út af}} {\textbf{a.}} \dicDirectTranslationCS{omdlít};  {\textbf{b.}} \dicDirectTranslationCS{usnout} \dicExampleIS{líða út af í hægindastólnum} \dicExampleCS{usnout v~křesle};  \dicIdiom{líða}[yfir]{ \dicPhraseIS{það líður yfir e‑n}} \dicFlx{impers} \dicDirectTranslationCS{(kdo) omdlívá, (kdo) ztrácí vědomí} \dicExampleIS{Það leið yfir mig.} \dicExampleCS{Omdlel jsem.};  \dicIdiom{líðast}{ \dicPhraseIS{e‑m líðst e‑að}} \dicFlx{refl impers} \dicDirectTranslationCS{(co) se (komu) toleruje} \dicIndirectTranslationCS{(často se záporem)};  \dicIdiom{líða}{ \dicPhraseIS{hvað sem öðru líður}} \dicDirectTranslationCS{beztak, když se to vezme kolem a~kolem}
\dicEntry[líðan] \dicTerm{líðan} \dicIPA{{l}{i}{\textlengthmark}{ð}{a}{\textsubring{n}}} \dicPos{f}[4] \dicFlx{(‑ar)}[3] \dicDirectTranslationCS{(zdravotní) stav} \dicExampleIS{Líðan sjúklingsins fer batnandi.} \dicExampleCS{Zdravotní stav pacienta se zlepšuje.}
\dicEntry[líðandi] \dicTerm{líð··andi} \dicIPA{{l}{i}{\textlengthmark}{ð}{a}{n}{\textsubring{d}}{\textsci}} \dicPos{adj}[13] \dicFlx{indecl}[1] \dicPhraseIS{á líðandi stund} \dicFlx{adv} \dicSynonym{núna} \dicDirectTranslationCS{teď, v~současné chvíli}
\dicEntry[líf] \dicTerm{líf} \dicsymFrequent\  \dicIPA{{l}{i}{\textlengthmark}{f}} \dicPos{n}[2] \dicFlx{(‑s, ‑)}[5] \textbf{1.} \dicSynonym*{það að lifa} \dicDirectTranslationCS{život, žití, bytí};  \dicPhraseIS{vera á lífi} \dicDirectTranslationCS{být naživu}  \textbf{2.} \dicSynonym*{lífverur} \dicDirectTranslationCS{život, existence (organismu)} \dicExampleIS{lífið á jörðinni} \dicExampleCS{život na Zemi}  \textbf{3.} \dicSynonym*{dagleg tilvera} \dicDirectTranslationCS{život (všední ap.)}  \textbf{4.} \dicSynonym{fjör} \dicDirectTranslationCS{živo, rušno};  \dicIdiom{líf}{ \dicPhraseIS{af lífi og sál}} \dicFlx{adv} \dicDirectTranslationCS{z~celých sil}; { \dicPhraseIS{blása lífið í e‑ð}} \dicLangCat{přen.} \dicDirectTranslationCS{vdechnout život (čemu)}; { \dicPhraseIS{elska e‑n út af lífinu}} \dicDirectTranslationCS{milovat (koho) z~celé duše};  \dicPhraseIS{láta lífið} \dicDirectTranslationCS{přijít o~život}; { \dicPhraseIS{svipta sig lífi}} \dicDirectTranslationCS{vzít si život}; { \dicPhraseIS{taka e‑n af lífi}} \dicDirectTranslationCS{vzít (komu) život}; { \dicPhraseIS{vekja til lífsins}} \dicLangCat{přen.} \dicDirectTranslationCS{probudit k~životu}
\dicEntry[lífbátur] \dicTerm{líf··bát|ur} \dicIPA{{l}{i}{v}{\textsubring{b}}{au}{\textsubring{d}}{\textscy}{\textsubring{r}}} \dicPos{m}[6] \dicFlx{(‑s, ‑ar)}[22] \dicDirectTranslationCS{záchranný člun}
\dicEntry[lífbein] \dicTerm{líf··bein} \dicIPA{{l}{i}{v}{\textsubring{b}}{ei}{\textsubring{n}}} \dicPos{n}[2] \dicFlx{(‑s, ‑)}[5] \dicFieldCat{anat.} \dicDirectTranslationCS{stydká kost}
\dicEntry[lífeðlisfræði] \dicTerm{líf··eðlis·fræð|i} \dicIPA{{l}{i}{\textlengthmark}{v}{\textepsilon}{ð}{l}{\textsci}{s}{f}{r}{a}{i}{ð}{\textsci}} \dicPos{f}[3] \dicFlx{(‑i)}[3] \dicDirectTranslationCS{fyziologie}
\dicEntry[lífeðlisfræðingur] \dicTerm{líf·eðlis·fræð··ing|ur} \dicIPA{{l}\-{i}\-{\textlengthmark}\-{v}\-{\textepsilon}\-{ð}\-{l}\-{\textsci}\-{s}\-{f}\-{r}\-{a}\-{i}\-{ð}\-{i}\-{\ng}\-{\r{g}}\-{\textscy}\-{\textsubring{r}}\-} \dicPos{m}[6] \dicFlx{(‑s, ‑ar)}[8] \dicDirectTranslationCS{fyziolog, fyzioložka}
\dicEntry[lífefnafræði] \dicTerm{líf··efna·fræð|i} \dicIPA{{l}{i}{\textlengthmark}{v}{\textepsilon}{\textsubring{b}}{n}{a}{f}{r}{a}{i}{ð}{\textsci}} \dicPos{f}[3] \dicFlx{(‑i)}[3] \dicDirectTranslationCS{biochemie}
\dicEntry[lífefnafræðingur] \dicTerm{líf·efna·fræð··ing|ur} \dicIPA{{l}\-{i}\-{\textlengthmark}\-{v}\-{\textepsilon}\-{\textsubring{b}}\-{n}\-{a}\-{f}\-{r}\-{a}\-{i}\-{ð}\-{i}\-{\ng}\-{\r{g}}\-{\textscy}\-{\textsubring{r}}\-} \dicPos{m}[6] \dicFlx{(‑s, ‑ar)}[8] \dicDirectTranslationCS{biochemik, biochemička}
\dicEntry[líferni] \dicTerm{líf··erni} \dicIPA{{l}{i}{\textlengthmark}{v}{\textepsilon}{r}{\textsubring{d}}{n}{\textsci}} \dicPos{n}[2] \dicFlx{(‑s)}[20] \dicDirectTranslationCS{životospráva, životní styl}
\dicEntry[lífeyrir] \dicTerm{líf··eyr|ir} \dicIPA{{l}{i}{\textlengthmark}{v}{ei}{r}{\textsci}{\textsubring{r}}} \dicPos{m}[7] \dicFlx{(‑is)}[2] \dicDirectTranslationCS{důchod, penze} \dicIndirectTranslationCS{(pravidelný příjem osob plynoucí z~pojištění)} \dicExampleIS{taka lífeyris} \dicExampleCS{příjem důchodu}
\dicEntry[lífeyrisgreiðsla] \dicTerm{líf·eyris··greiðsl|a} \dicIPA{{l}{i}{\textlengthmark}{v}{ei}{r}{\textsci}{s}{\r{g}}{r}{ei}{ð}{s}{\textsubring{d}}{l}{a}} \dicPos{f}[1] \dicFlx{(‑u, ‑ur)}[13] \dicDirectTranslationCS{výplata důchodu}
\dicEntry[lífeyrisréttindi] \dicTerm{líf·eyris··rétt·indi} \dicIPA{{l}{i}{\textlengthmark}{v}{ei}{r}{\textsci}{s}{r}{j}{\textepsilon}{h}{\textsubring{d}}{\textsci}{n}{\textsubring{d}}{\textsci}} \dicPos{n}[2] \dicFlx{pl}[19] \dicDirectTranslationCS{nárok na důchod}
\dicEntry[lífeyrissjóðstekjur] \dicTerm{líf·eyris·sjóðs··tekjur} \dicIPA{{l}\-{i}\-{\textlengthmark}\-{v}\-{ei}\-{r}\-{\textsci}\-{s}\-{j}\-{ou}\-{ð}\-{s}\-{t\smash{\textsuperscript{h}}}\-{\textepsilon}\-{\r{\textObardotlessj}}\-{\textscy}\-{\textsubring{r}}\-} \dicPos{f}[1] \dicFlx{pl}[18] \dicDirectTranslationCS{příjmy z~penzijního fondu}
\dicEntry[lífeyrissjóður] \dicTerm{líf·eyris··sjóð|ur} \dicIPA{{l}{i}{\textlengthmark}{v}{ei}{r}{\textsci}{s}{j}{ou}{ð}{\textscy}{\textsubring{r}}} \dicPos{m}[9] \dicFlx{(‑s, ‑ir)}[6] \dicFieldCat{ekon.} \dicDirectTranslationCS{penzijní fond}
\dicEntry[lífeyrissparnaður] \dicTerm{líf·eyris··spar·nað|ur} \dicIPA{{l}\-{i}\-{\textlengthmark}\-{v}\-{ei}\-{r}\-{\textsci}\-{s}\-{\textsubring{b}}\-{a}\-{r}\-{\textsubring{d}}\-{n}\-{a}\-{ð}\-{\textscy}\-{\textsubring{r}}\-} \dicPos{m}[10] \dicFlx{(‑ar)}[9] \dicDirectTranslationCS{důchodové spoření}
\dicEntry[lífeyristryggingar] \dicTerm{líf··eyris··trygg·ingar} \dicIPA{{l}\-{i}\-{\textlengthmark}\-{v}\-{ei}\-{r}\-{\textsci}\-{s}\-{t\smash{\textsuperscript{h}}}\-{r}\-{\textsci}\-{\r{\textObardotlessj}}\-{i}\-{\ng}\-{\r{g}}\-{a}\-{\textsubring{r}}\-} \dicPos{f}[4] \dicFlx{pl}[6] \dicDirectTranslationCS{důchodové pojištění}
\dicEntry[lífeyrisþegi] \dicTerm{líf·eyris··þeg|i} \dicIPA{{l}{i}{\textlengthmark}{v}{ei}{r}{\textsci}{s}{\texttheta}{ei}{\textsci}} \dicPos{m}[1] \dicFlx{(‑a, ‑ar)}[1] \dicDirectTranslationCS{důchodce, důchodkyně, penzista, penzistka}
\dicEntry[líffræði] \dicTerm{líf··fræð|i} \dicIPA{{l}{i}{f}{\textlengthmark}{r}{a}{i}{ð}{\textsci}} \dicPos{f}[3] \dicFlx{(‑i)}[3] \dicDirectTranslationCS{biologie}
\dicEntry[líffræðingur] \dicTerm{líf·fræð··ing|ur} \dicIPA{{l}{i}{f}{\textlengthmark}{r}{a}{i}{ð}{i}{\ng}{\r{g}}{\textscy}{\textsubring{r}}} \dicPos{m}[6] \dicFlx{(‑s, ‑ar)}[8] \dicDirectTranslationCS{biolog, bioložka}
\dicEntry[líffærafræði] \dicTerm{líf·færa··fræð|i} \dicIPA{{l}{i}{f}{\textlengthmark}{a}{i}{r}{a}{f}{r}{a}{i}{ð}{\textsci}} \dicPos{f}[3] \dicFlx{(‑i)}[3] \dicDirectTranslationCS{anatomie}
\dicEntry[líffæri] \dicTerm{líf··færi} \dicIPA{{l}{i}{f}{\textlengthmark}{a}{i}{r}{\textsci}} \dicPos{n}[2] \dicFlx{(‑s, ‑)}[14] \dicFieldCat{anat.} \dicDirectTranslationCS{orgán, ústrojí}
\dicEntry[lífga] \dicTerm{lífg|a} \dicIPA{{l}{i}{v}{\r{g}}{a}} \dicPos{v}[1] \dicFlx{(‑aði)}[1] \dicFlx{acc} \textbf{1.} \dicSynonym{fjörga} \dicDirectTranslationCS{oživit, oživovat, vzkřísit} \dicIndirectTranslationCS{(činit živějším)} \dicExampleIS{lífga eldinn} \dicExampleCS{oživit oheň};  \dicPhraseIS{lífga upp á e‑ð} \dicDirectTranslationCS{oživit\,/\addthin osvěžit (co) (pokoj barvami ap.)}  \textbf{2.} \dicSynonym*{vekja til lífs} \dicDirectTranslationCS{oživit, oživovat, vzkřísit, resuscitovat} \dicIndirectTranslationCS{(probudit člověka k~životu)};  \dicPhraseIS{lífga e‑n úr dauðadái} \dicDirectTranslationCS{probudit (koho) z~kómatu};  \dicPhraseIS{lífga e‑n við} \dicDirectTranslationCS{oživit (koho), probudit (koho) k~životu}
\dicEntry[lífgjafi] \dicTerm{líf··gjaf|i} \dicIPA{{l}{i}{v}{\r{\textObardotlessj}}{a}{v}{\textsci}} \dicPos{m}[1] \dicFlx{(‑a, ‑ar)}[8] \dicDirectTranslationCS{zachránce\,/\addthin zachránkyně (života)}
\dicEntry[lífgjöf] \dicTerm{líf··|gjöf} \dicIPA{{l}{i}{v}{\r{\textObardotlessj}}{\oe}{f}} \dicPos{f}[7] \dicFlx{(‑gjafar)}[19] \dicDirectTranslationCS{záchrana života}
\dicEntry[lífgun] \dicTerm{lífg|un} \dicIPA{{l}{i}{v}{\r{g}}{\textscy}{\textsubring{n}}} \dicPos{f}[7] \dicFlx{(‑unar, ‑anir)}[8] \dicDirectTranslationCS{resuscitace, reanimace, oživování}
\dicEntry[lífhimna] \dicTerm{líf··himn|a} \dicIPA{{l}{i}{\textlengthmark}{f}{h}{\textsci}{m}{n}{a}} \dicPos{f}[1] \dicFlx{(‑u, ‑ur)}[7] \dicFieldCat{anat.} \dicDirectTranslationCS{pobřišnice}
\dicEntry[lífklukka] \dicTerm{líf··klukk|a} \dicIPA{{l}{i}{v}{k\smash{\textsuperscript{h}}}{l}{\textscy}{h}{\r{g}}{a}} \dicPos{f}[1] \dicFlx{(‑u, ‑ur)}[13] \dicDirectTranslationCS{biologické hodiny}
\dicEntry[líflaus] \dicTerm{líf··laus} \dicIPA{{l}{i}{v}{l}{\oe i}{s}} \dicPos{adj}[5]\dicFlx{}[-1] \textbf{1.} \dicDirectTranslationCS{(jsoucí) bez života, mrtvý}  \textbf{2.} \dicFieldCat{biol.} \dicLangCat{zast.} \dicSynonym{ólífrænn} \dicDirectTranslationCS{neživý}
\dicEntry[líflát] \dicTerm{líf··lát} \dicIPA{{l}{i}{v}{l}{au}{\textsubring{d}}} \dicPos{n}[2] \dicFlx{(‑s, ‑)}[5] \dicSynonym{aftaka\smash{\textsuperscript{1}}} \dicDirectTranslationCS{poprava} \dicExampleIS{hóta e‑m  líflát} \dicExampleCS{hrozit (komu) popravou}
\dicEntry[lífláta] \dicTerm{líf··|láta} \dicIPA{{l}{i}{v}{l}{au}{\textsubring{d}}{a}} \dicPos{v}[6] \dicFlx{(‑læt, ‑lét, ‑létum, ‑léti, ‑látið)}[87] \dicFlx{acc} \dicDirectTranslationCS{popravit} \dicExampleIS{lífláta e‑n} \dicExampleCS{popravit (koho)}
\dicEntry[líflátið] \dicTerm{líf··látið} \dicIPA{{l}{i}{v}{l}{au}{\textsubring{d}}{\textsci}{\texttheta}} \dicPos{v} \dicFlx{supin} \dicLink{lífláta}
\dicEntry[líflegur] \dicTerm{líf··legur} \dicIPA{{l}{i}{v}{l}{\textepsilon}{\textbabygamma}{\textscy}{\textsubring{r}}} \dicPos{adj}[1]\dicFlx{}[-8] \dicSynonym{fjörugur} \dicDirectTranslationCS{živý, svěží, plný života}
\dicEntry[líflét] \dicTerm{líf··lét} \dicIPA{{l}{i}{v}{l}{j}{\textepsilon}{\textsubring{d}}} \dicPos{v} \dicFlx{ind pf sg 1 pers} \dicLink{lífláta}
\dicEntry[lífléti] \dicTerm{líf··léti} \dicIPA{{l}{i}{v}{l}{j}{\textepsilon}{\textsubring{d}}{\textsci}} \dicPos{v} \dicFlx{con pf sg 1 pers} \dicLink{lífláta}
\dicEntry[líflétum] \dicTerm{líf··létum} \dicIPA{{l}{i}{v}{l}{j}{\textepsilon}{\textsubring{d}}{\textscy}{\textsubring{m}}} \dicPos{v} \dicFlx{ind pf pl 1 pers} \dicLink{lífláta}
\dicEntry[líflæt] \dicTerm{líf··læt} \dicIPA{{l}{i}{v}{l}{a}{i}{\textsubring{d}}} \dicPos{v} \dicFlx{ind praes sg 1 pers} \dicLink{lífláta}
\dicEntry[lífríki] \dicTerm{líf··ríki} \dicIPA{{l}{i}{v}{r}{i}{\r{\textObardotlessj}}{\textsci}} \dicPos{n}[2] \dicFlx{(‑s, ‑)}[16] \dicDirectTranslationCS{biosystém, biologický systém}
\dicEntry[lífrænn] \dicTerm{líf··rænn} \dicIPA{{l}{i}{v}{r}{a}{i}{\textsubring{d}}{\textsubring{n}}} \dicPos{adj}[7]\dicFlx{}[-1] \textbf{1.} \dicDirectTranslationCS{živý, živoucí, rušný} \dicExampleIS{lífrænt starf} \dicExampleCS{živá práce}  \textbf{2.} \dicSynonym{náttúrulegur} \dicDirectTranslationCS{organický, ekologický} \dicExampleIS{lífrænn búskapur} \dicExampleCS{organické zemědělství} \dicAntonym{ólífrænn}
\dicEntry[lífsbarátta] \dicTerm{lífs··bar·átt|a} \dicIPA{{l}{i}{f}{s}{\textsubring{b}}{a}{r}{au}{h}{\textsubring{d}}{a}} \dicPos{f}[1] \dicFlx{(‑u)}[5] \dicDirectTranslationCS{existenční\,/\addthin životní boj}
\dicEntry[lífseigur] \dicTerm{líf··seigur} \dicIPA{{l}{i}{f}{s}{ei}{\textbabygamma}{\textscy}{\textsubring{r}}} \dicPos{adj}[1]\dicFlx{}[-1] \dicDirectTranslationCS{houževnatý, (jsoucí) s~tuhým kořínkem}
\dicEntry[lífsflótti] \dicTerm{lífs··flótt|i} \dicIPA{{l}{i}{f}{s}{f}{l}{ou}{h}{\textsubring{d}}{\textsci}} \dicPos{m}[1] \dicFlx{(‑a)}[3] \dicDirectTranslationCS{útěk před realitou, eskapismus}
\dicEntry[lífsglaður] \dicTerm{lífs··|glaður} \dicIPA{{l}{i}{f}{s}{\r{g}}{l}{a}{ð}{\textscy}{\textsubring{r}}} \dicPos{adj}[2] \dicFlx{(f ‑glöð)}[7] \dicSynonym{glaðlyndur} \dicDirectTranslationCS{radostný, šťastný, radující se ze života}
\dicEntry[lífsgæði] \dicTerm{lífs··gæði} \dicIPA{{l}{i}{f}{s}{\r{\textObardotlessj}}{a}{i}{ð}{\textsci}} \dicPos{n}[2] \dicFlx{pl}[19] \dicDirectTranslationCS{životní úroveň, kvalita života}
\dicEntry[lífsháski] \dicTerm{lífs··hásk|i} \dicIPA{{l}{i}{f}{s}{h}{au}{s}{\r{\textObardotlessj}}{\textsci}} \dicPos{m}[1] \dicFlx{(‑a)}[3] \dicDirectTranslationCS{ohrožení života}
\dicEntry[lífshætta] \dicTerm{lífs··hætt|a} \dicIPA{{l}{i}{f}{s}{h}{a}{i}{h}{\textsubring{d}}{a}} \dicPos{f}[1] \dicFlx{(‑u)}[5] \dicDirectTranslationCS{nebezpečí života}
\dicEntry[lífshættulegur] \dicTerm{lífs··hættu·legur} \dicIPA{{l}{i}{f}{s}{h}{a}{i}{h}{\textsubring{d}}{\textscy}{l}{\textepsilon}{\textbabygamma}{\textscy}{\textsubring{r}}} \dicPos{adj}[1]\dicFlx{}[-8] \dicDirectTranslationCS{životu nebezpečný}
\dicEntry[lífskjör] \dicTerm{lífs··kjör} \dicIPA{{l}{i}{f}{s}{c\smash{\textsuperscript{h}}}{\oe}{\textsubring{r}}} \dicPos{n}[2] \dicFlx{pl}[9] \dicDirectTranslationCS{životní podmínky}
\dicEntry[lífsleiði] \dicTerm{lífs··leið|i} \dicIPA{{l}{i}{f}{s}{l}{ei}{ð}{\textsci}} \dicPos{m}[1] \dicFlx{(‑a)}[3] \dicDirectTranslationCS{životní nuda, omrzelost životem}
\dicEntry[lífsleikni] \dicTerm{lífs··leikn|i} \dicIPA{{l}{i}{f}{s}{l}{ei}{h}{\r{g}}{n}{\textsci}} \dicPos{f}[3] \dicFlx{(‑i)}[3] \dicFieldCat{škol.} \dicDirectTranslationCS{praktická výchova} \dicIndirectTranslationCS{(předmět ve škole)}
\dicEntry[lífsnauðsyn] \dicTerm{lífs··nauð·syn} \dicIPA{{l}{i}{f}{s}{n}{\oe i}{ð}{s}{\textsci}{\textsubring{n}}} \dicPos{f}[4] \dicFlx{(‑jar, ‑jar)}[9] \dicDirectTranslationCS{životní nutnost}
\dicEntry[lífsorka] \dicTerm{lífs··ork|a} \dicIPA{{l}{i}{f}{s}{\textopeno}{\textsubring{r}}{\r{g}}{a}} \dicPos{f}[1] \dicFlx{(‑u)}[5] \dicDirectTranslationCS{vitalita, životní energie}
\dicEntry[lífsregla] \dicTerm{lífs··regl|a} \dicIPA{{l}{i}{f}{s}{r}{\textepsilon}{\r{g}}{l}{a}} \dicPos{f}[1] \dicFlx{(‑u, ‑ur)}[19] \dicDirectTranslationCS{životní zásada}
\dicEntry[lífsreynsla] \dicTerm{lífs··reynsl|a} \dicIPA{{l}{i}{f}{s}{r}{ei}{n}{s}{\textsubring{d}}{l}{a}} \dicPos{f}[1] \dicFlx{(‑u)}[5] \dicDirectTranslationCS{životní zkušenost}
\dicEntry[lífsskilyrði] \dicTerm{lífs··skil·yrði} \dicIPA{{l}{i}{f}{s}{\r{\textObardotlessj}}{\textsci}{l}{\textsci}{r}{ð}{\textsci}} \dicPos{n}[2] \dicFlx{pl}[19] \dicDirectTranslationCS{životní podmínky}
\dicEntry[lífsskoðun] \dicTerm{lífs··skoð|un} \dicIPA{{l}{i}{f}{s}{\r{g}}{\textopeno}{ð}{\textscy}{\textsubring{n}}} \dicPos{f}[7] \dicFlx{(‑unar, ‑anir)}[8] \dicDirectTranslationCS{životní postoj\,/\addthin názor}
\dicEntry[lífsspeki] \dicTerm{lífs··spek|i} \dicIPA{{l}{i}{f}{s}{\textsubring{b}}{\textepsilon}{\r{\textObardotlessj}}{\textsci}} \dicPos{f}[3] \dicFlx{(‑i)}[3] \dicDirectTranslationCS{životní moudrost}
\dicEntry[lífsstarf] \dicTerm{lífs··|starf} \dicIPA{{l}{i}{f}{s}{\textsubring{d}}{a}{r}{f}} \dicPos{n}[2] \dicFlx{(‑starfs, ‑störf)}[8] \dicDirectTranslationCS{životní poslání}
\dicEntry[lífsstíll] \dicTerm{lífs··stíl|l} \dicIPA{{l}{i}{f}{s}{\textsubring{d}}{i}{\textsubring{d}}{\textsubring{l}}} \dicPos{m}[6] \dicFlx{(‑s)}[50] \dicDirectTranslationCS{životní styl}
\dicEntry[lífstykki] \dicTerm{líf··stykki} \dicIPA{{l}{i}{f}{s}{\textsubring{d}}{\textsci}{h}{\r{\textObardotlessj}}{\textsci}} \dicPos{n}[2] \dicFlx{(‑s, ‑)}[16] \dicDirectTranslationCS{korzet, šněrovačka}
\dicEntry[lífsveig] \dicTerm{lífs··veig} \dicIPA{{l}{i}{f}{s}{v}{ei}{x}} \dicPos{f}[4] \dicFlx{(‑ar, ‑ar)}[1] \dicDirectTranslationCS{elixír života}
\dicEntry[lífsviðurværi] \dicTerm{lífs··viður·væri} \dicIPA{{l}{i}{f}{s}{v}{\textsci}{ð}{\textscy}{r}{v}{a}{i}{r}{\textsci}} \dicPos{n}[2] \dicFlx{(‑s)}[20] \dicDirectTranslationCS{obživa (jídlo ap.)}
\dicEntry[lífsþægindi] \dicTerm{lífs··þæg·indi} \dicIPA{{l}{i}{f}{s}{\texttheta}{a}{i}{j}{\textsci}{n}{\textsubring{d}}{\textsci}} \dicPos{n}[2] \dicFlx{pl}[19] \dicDirectTranslationCS{životní komfort\,/\addthin pohodlí}
\dicEntry[líftóra] \dicTerm{líf··tór|a} \dicIPA{{l}{i}{f}{t\smash{\textsuperscript{h}}}{ou}{r}{a}} \dicPos{f}[1] \dicFlx{(‑u)}[5] \dicDirectTranslationCS{jiskra života};  \dicPhraseIS{hræða úr e‑m líftóruna} \dicLangCat{přen.} \dicDirectTranslationCS{vyděsit (koho) k~smrti}
\dicEntry[líftrygging] \dicTerm{líf··trygg·ing} \dicIPA{{l}{i}{f}{t\smash{\textsuperscript{h}}}{r}{\textsci}{\r{\textObardotlessj}}{i}{\ng}{\r{g}}} \dicPos{f}[4] \dicFlx{(‑ar, ‑ar)}[5] \dicDirectTranslationCS{životní pojištění}
\dicEntry[líftryggja] \textls[15]{\dicTerm{líf··trygg|ja} \dicIPA{{l}{i}{f}{t\smash{\textsuperscript{h}}}{r}{\textsci}{\r{\textObardotlessj}}{a}} \dicPos{v}[2] \dicFlx{(‑ði, ‑t)}[92] \dicFlx{acc} \dicDirectTranslationCS{uzavřít životní pojiš\nobreak tění}}
\dicEntry[lífvera] \dicTerm{líf··ver|a} \dicIPA{{l}{i}{v}{\textlengthmark}{\textepsilon}{r}{a}} \dicPos{f}[1] \dicFlx{(‑u, ‑ur)}[7] \dicDirectTranslationCS{živá bytost, živý organismus}
\dicEntry[lífvænlegur] \dicTerm{líf··væn·legur} \dicIPA{{l}{i}{v}{\textlengthmark}{a}{i}{n}{l}{\textepsilon}{\textbabygamma}{\textscy}{\textsubring{r}}} \dicPos{adj}[1]\dicFlx{}[-8] \textbf{1.} \dicDirectTranslationCS{rentabilní} \dicExampleIS{lífvænleg atvinna} \dicExampleCS{rentabilní zaměstnání}  \textbf{2.} \dicDirectTranslationCS{příznivý pro život} \dicExampleIS{lífvænlegur staður} \dicExampleCS{místo příznivé pro život}
\begin{xtolerant}{}{1pt}
\dicEntry[lífvörður] \dicTerm{líf··|vörður}\addthinS\dicIPA{{l}{i}{v}{\textlengthmark}{\oe}{r}{ð}{\textscy}{\textsubring{r}}}\addthinS\dicPos{m}[11]\addthinS\dicFlx{(‑varðar, ‑verðir)}[5] \dicDirectTranslationCS{osobní strážce\,/\addthin strážkyně, tělesná stráž, bodyguard}
\end{xtolerant}
\dicEntry[lík] \dicTerm{lík} \dicsymFrequent\  \dicIPA{{l}{i}{\textlengthmark}{\r{g}}} \dicPos{n}[2] \dicFlx{(‑s, ‑)}[5] \textbf{1.} \dicSynonym*{dauður líkami} \dicDirectTranslationCS{mrtvola} \dicExampleIS{líkið af e‑m} \dicExampleCS{(čí) mrtvola}  \textbf{2.} \dicLangCat{hovor.} \dicSynonym*{tóm vínflaska} \dicDirectTranslationCS{prázdná láhev (po alkoholu)}
\dicEntry[líka] \dicTerm{lík|a\smash{\textsuperscript{1}}} \dicsymFrequent\  \dicIPA{{l}{i}{\textlengthmark}{\r{g}}{a}} \dicPos{v}[1] \dicFlx{(‑aði)}[82] \dicFlx{impers} \dicPhraseIS{e‑m líkar e‑að} \dicDirectTranslationCS{(co) se (komu) líbí, (co) má (kdo) rád} \dicExampleIS{Mér líkar það ekki.} \dicExampleCS{Mně se to nelíbí.};  \dicPhraseIS{e‑m líkar vel við e‑n} \dicDirectTranslationCS{(kdo) má (koho) rád}
\dicEntry[líka] \dicTerm{líka\smash{\textsuperscript{2}}} \dicsymFrequent\  \dicIPA{{l}{i}{\textlengthmark}{\r{g}}{a}} \dicPos{adv} \dicSynonym{einnig} \dicDirectTranslationCS{také, taky, rovněž} \dicExampleIS{Dagblaðið er líka á netinu.} \dicExampleCS{Deník je také na internetu.}
\dicEntry[líkami] \dicTerm{líkam|i} \dicsymFrequent\  \dicIPA{{l}{i}{\textlengthmark}{\r{g}}{a}{m}{\textsci}} \dicPos{m}[1] \dicFlx{(‑a, ‑ar)}[8] \dicDirectTranslationCS{tělo} \dicExampleIS{æfingar fyrir líkama} \dicExampleCS{cvičení pro tělo}
\dicEntry[líkamlegur] \dicTerm{líkam··legur} \dicsymFrequent\  \dicIPA{{l}{i}{\r{g}}{a}{m}{l}{\textepsilon}{\textbabygamma}{\textscy}{\textsubring{r}}} \dicPos{adj}[1]\dicFlx{}[-8] \dicDirectTranslationCS{tělesný, fyzický} \dicExampleIS{líkamleg atvinna} \dicExampleCS{fyzická práce}
\dicEntry[líkamsárás] \dicTerm{líkams··á·rás} \dicIPA{{l}{i}{\textlengthmark}{\r{g}}{a}{m}{s}{au}{r}{au}{s}} \dicPos{f}[7] \dicFlx{(‑ar, ‑ir)}[1] \dicDirectTranslationCS{tělesný útok, fyzické napadení}
\dicEntry[líkamsástand] \dicTerm{líkams··á·stand} \dicIPA{{l}{i}{\textlengthmark}{\r{g}}{a}{m}{s}{au}{s}{\textsubring{d}}{a}{n}{\textsubring{d}}} \dicPos{n}[2] \dicFlx{(‑s)}[2] \dicDirectTranslationCS{tělesný stav}
\dicEntry[líkamsburðir] \dicTerm{líkams··burðir} \dicIPA{{l}{i}{\textlengthmark}{\r{g}}{a}{m}{s}{\textsubring{b}}{\textscy}{r}{ð}{\textsci}{\textsubring{r}}} \dicPos{m}[10] \dicFlx{pl}[1] \dicDirectTranslationCS{tělesná\,/\addthin fyzická síla}
\dicEntry[líkamshiti] \dicTerm{líkams··hit|i} \dicIPA{{l}{i}{\textlengthmark}{\r{g}}{a}{m}{s}{h}{\textsci}{\textsubring{d}}{\textsci}} \dicPos{m}[1] \dicFlx{(‑a)}[3] \dicDirectTranslationCS{tělesná teplota}
\dicEntry[líkamsrækt] \dicTerm{líkams··rækt} \dicIPA{{l}{i}{\textlengthmark}{\r{g}}{a}{m}{s}{r}{a}{i}{x}{\textsubring{d}}} \dicPos{f}[4] \dicFlx{(‑ar)}[3] \dicDirectTranslationCS{tělesné cvičení, posilování, fitness}
\dicEntry[líkamsræktarstöð] \dicTerm{líkams·ræktar··stöð} \dicIPA{{l}{i}{\textlengthmark}{\r{g}}{a}{m}{s}{r}{a}{i}{x}{\textsubring{d}}{a}{\textsubring{r}}{s}{\textsubring{d}}{\oe}{\texttheta}} \dicPos{f}[6] \dicFlx{(‑var, ‑var)}[1] \dicDirectTranslationCS{posilovna, fitness centrum}
\dicEntry[líkamsstærð] \dicTerm{líkams··stærð} \dicIPA{{l}{i}{\textlengthmark}{\r{g}}{a}{m}{s}{\textsubring{d}}{a}{i}{r}{\texttheta}} \dicPos{f}[7] \dicFlx{(‑ar)}[3] \dicDirectTranslationCS{životní velikost} \dicExampleIS{í fullri líkamsstærð} \dicExampleCS{v~životní velikosti}
\dicEntry[líkan] \dicTerm{lík|an} \dicIPA{{l}{i}{\textlengthmark}{\r{g}}{a}{\textsubring{n}}} \dicPos{n}[2] \dicFlx{(‑ans, ‑ön)}[8] \dicDirectTranslationCS{model, maketa} \dicIndirectTranslationCS{(zmenšené provedení předmětu)} \dicExampleIS{líkan af húsi} \dicExampleCS{model domu}
\dicEntry[líkbrennsla] \dicTerm{lík··brennsl|a} \dicIPA{{l}{i}{\textlengthmark}{\r{g}}{\textsubring{b}}{r}{\textepsilon}{n}{s}{\textsubring{d}}{l}{a}} \dicPos{f}[1] \dicFlx{(‑u, ‑ur)}[19] \dicSynonym{bálför} \dicDirectTranslationCS{kremace}
\dicEntry[líkbörur] \dicTerm{lík··börur} \dicIPA{{l}{i}{\textlengthmark}{\r{g}}{\textsubring{b}}{\oe}{r}{\textscy}{\textsubring{r}}} \dicPos{f}[12] \dicFlx{pl}[4] \dicDirectTranslationCS{máry}
\dicEntry[líkfylgd] \dicTerm{lík··fylgd} \dicIPA{{l}{i}{\textlengthmark}{\r{g}}{f}{\textsci}{l}{\textsubring{d}}} \dicPos{f}[7] \dicFlx{(‑ar, ‑ir)}[1] \dicDirectTranslationCS{pohřební průvod}
\dicEntry[líkhús] \dicTerm{lík··hús} \dicIPA{{l}{i}{\textlengthmark}{\r{g}}{h}{u}{s}} \dicPos{n}[2] \dicFlx{(‑s, ‑)}[5] \dicDirectTranslationCS{márnice}
\dicEntry[líki] \dicTerm{lík|i\smash{\textsuperscript{1}}} \dicIPA{{l}{i}{\textlengthmark}{\r{\textObardotlessj}}{\textsci}} \dicPos{m}[1] \dicFlx{(‑a, ‑ar)}[1] \dicSynonym{jafningi} \dicDirectTranslationCS{sobě rovný, podobný (člověk)} \dicExampleIS{hann og hans líkar} \dicExampleCS{on a~jemu podobní};  \dicPhraseIS{eiga engan sinn líka} \dicDirectTranslationCS{nemít sobě rovného}
\dicEntry[líki] \dicTerm{líki\smash{\textsuperscript{2}}} \dicIPA{{l}{i}{\textlengthmark}{\r{\textObardotlessj}}{\textsci}} \dicPos{n}[2] \dicFlx{(‑s, ‑)}[16] \dicSynonym{sköpulag} \dicDirectTranslationCS{podoba, forma};  \dicPhraseIS{í heilu líki} \dicFlx{adj} \dicDirectTranslationCS{celý, netknutý}
\dicEntry[líkindi] \dicTerm{lík··indi} \dicsymFrequent\  \dicIPA{{l}{i}{\textlengthmark}{\r{\textObardotlessj}}{\textsci}{n}{\textsubring{d}}{\textsci}} \dicPos{n}[2] \dicFlx{pl}[19] \dicDirectTranslationCS{pravděpodobnost} \dicExampleIS{Það eru mikil líkindi til þess að \dots{}} \dicExampleCS{Je velká pravděpodobnost, že\dots{}};  \dicPhraseIS{að öllum líkindum} \dicFlx{adv} \dicDirectTranslationCS{s~největší pravděpodobností}
\dicEntry[líking] \dicTerm{lík··ing} \dicIPA{{l}{i}{\textlengthmark}{\r{\textObardotlessj}}{i}{\ng}{\r{g}}} \dicPos{f}[4] \dicFlx{(‑ar, ‑ar)}[5] \textbf{1.} \dicSynonym{líki\smash{\textsuperscript{2}}} \dicDirectTranslationCS{podoba, podobnost, obdoba, analogie}  \textbf{2.} \dicDirectTranslationCS{přirovnání, podobenství} \dicExampleIS{tala í líkingum} \dicExampleCS{mluvit v~přirovnáních}  \textbf{3.} \dicFieldCat{lit.} \dicSynonym{samlíking} \dicDirectTranslationCS{metafora} \dicExampleIS{líkingar í skáldskap} \dicExampleCS{metafora v~literatuře}  \textbf{4.} \dicFieldCat{mat.} \dicSynonym{jafna\smash{\textsuperscript{1}}} \dicDirectTranslationCS{rovnost}
\dicEntry[líkja] \dicTerm{lík|ja} \dicsymFrequent\  \dicIPA{{l}{i}{\textlengthmark}{\r{\textObardotlessj}}{a}} \dicPos{v}[2] \dicFlx{(‑ti, ‑t)}[25] \dicFlx{dat} \textbf{1.} \dicDirectTranslationCS{přirovnat, přirovnávat};  \dicPhraseIS{líkja e‑u saman} \dicDirectTranslationCS{srovnávat (co)};  \dicPhraseIS{líkja e‑u við e‑ð} \dicDirectTranslationCS{přirovnávat (co) k~(čemu)} \dicExampleIS{líkja ævi mannsins við blóm} \dicExampleCS{přirovnávat život člověka ke květině}  \textbf{2.} \dicDirectTranslationCS{napodobit, napodobovat};  \dicPhraseIS{líkja eftir e‑u} \dicDirectTranslationCS{napodobit (co), imitovat (co)} \dicExampleIS{líkja eftir hljóði fuglsins} \dicExampleCS{napodobovat zvuk ptáčka};  \dicIdiom{líkjast}{ \dicPhraseIS{líkjast e‑m}} \dicFlx{refl} \dicDirectTranslationCS{připomínat (koho), podobat se (komu)} \dicExampleIS{Hverjum líkist ég?} \dicExampleCS{Komu se podobám?}
\dicEntry[líkjör] \dicTerm{líkjör} \dicIPA{{l}{i}{\textlengthmark}{\r{\textObardotlessj}}{\oe}{\textsubring{r}}} \dicPos{m}[4] \dicFlx{(‑s, ‑ar)}[14] \dicDirectTranslationCS{likér}
\dicEntry[líkkista] \dicTerm{lík··kist|a} \dicIPA{{l}{i}{\textlengthmark}{\r{g}}{c\smash{\textsuperscript{h}}}{\textsci}{s}{\textsubring{d}}{a}} \dicPos{f}[1] \dicFlx{(‑u, ‑ur)}[19] \dicDirectTranslationCS{rakev}
\dicEntry[líkklæði] \dicTerm{lík··klæði} \dicIPA{{l}{i}{\textlengthmark}{\r{g}}{k\smash{\textsuperscript{h}}}{l}{a}{i}{ð}{\textsci}} \dicPos{n}[2] \dicFlx{pl}[19] \dicDirectTranslationCS{pohřební roucho}
\dicEntry[líkl.] \dicTerm{líkl.} \dicPos{zkr} \dicPhraseIS{líklega} \dicFlx{adv} \dicDirectTranslationCS{pravděpodobně}
\dicEntry[líklega] \dicTerm{lík··leg|a} \dicsymFrequent\  \dicIPA{{l}{i}{\textlengthmark}{\r{g}}{l}{\textepsilon}{\textbabygamma}{a}} \dicPos{adv} \dicFlx{(comp ‑ar, sup ‑ast)} \dicSynonym{sennilega} \dicDirectTranslationCS{pravděpodobně, nejspíš, asi} \dicExampleIS{Þau segja líklega nei.} \dicExampleCS{Pravděpodobně řeknou ne.}
\dicEntry[líklegur] \dicTerm{lík··legur} \dicsymFrequent\  \dicIPA{{l}{i}{\textlengthmark}{\r{g}}{l}{\textepsilon}{\textbabygamma}{\textscy}{\textsubring{r}}} \dicPos{adj}[1]\dicFlx{}[-8] \textbf{1.} \dicSynonym{efnilegur} \dicDirectTranslationCS{slibný, nadějný}  \textbf{2.} \dicSynonym{sennilegur} \dicDirectTranslationCS{pravděpodobný, předpokládaný} \dicExampleIS{Þeir eru líklegir til að feta í sömu fótspor.} \dicExampleCS{Pravděpodobně budou kráčet ve stejných stopách.}
\dicEntry[líkn] \dicTerm{líkn} \dicIPA{{l}{i}{h}{\r{g}}{\textsubring{n}}} \dicPos{f}[7] \dicFlx{(‑ar, ‑ir)}[1] \textbf{1.} \dicSynonym{miskunn} \dicDirectTranslationCS{milost, smilování}  \textbf{2.} \dicSynonym{hjúkrun} \dicDirectTranslationCS{opatrování, péče}
\dicEntry[líkna] \dicTerm{líkn|a} \dicIPA{{l}{i}{h}{\r{g}}{n}{a}} \dicPos{v}[1] \dicFlx{(‑aði)}[1] \dicFlx{dat} \dicSynonym{hjúkra} \dicDirectTranslationCS{pečovat, poskytnout péči, opatrovat} \dicExampleIS{líkna e‑m} \dicExampleCS{pečovat o~(koho)}
\dicEntry[líknarbelgur] \dicTerm{líknar··belg|ur} \dicIPA{{l}{i}{h}{\r{g}}{n}{a}{r}{\textsubring{b}}{\textepsilon}{l}{\r{g}}{\textscy}{\textsubring{r}}} \dicPos{m}[9] \dicFlx{(‑s\,/\addthin ‑jar, ‑ir)}[26] \textbf{1.} \dicFieldCat{anat.} \dicSynonym*{fósturfylgja} \dicDirectTranslationCS{amnion, vnitřní plodový obal}  \textbf{2.} \dicSynonym*{loftpúði í bíl} \dicDirectTranslationCS{airbag} \dicIndirectTranslationCS{(zařízení pasivní bezpečnosti v~autě)}
\dicEntry[líknardeild] \dicTerm{líknar··deild} \dicIPA{{l}{i}{h}{\r{g}}{n}{a}{r}{\textsubring{d}}{ei}{l}{\textsubring{d}}} \dicPos{f}[7] \dicFlx{(‑ar, ‑ir)}[1] \dicDirectTranslationCS{hospic}
\dicEntry[líknardráp] \dicTerm{líknar··dráp} \dicIPA{{l}{i}{h}{\r{g}}{n}{a}{r}{\textsubring{d}}{r}{au}{\textsubring{b}}} \dicPos{n}[2] \dicFlx{(‑s, ‑)}[5] \dicDirectTranslationCS{eutanazie}
\dicEntry[líknarmál] \dicTerm{líknar··mál} \dicIPA{{l}{i}{h}{\r{g}}{n}{a}{r}{m}{au}{\textsubring{l}}} \dicPos{n}[2] \dicFlx{(‑s, ‑)}[5] \dicDirectTranslationCS{charita, dobročinnost}
\dicEntry[líknarstofnun] \dicTerm{líknar··stofn|un} \dicIPA{{l}{i}{h}{\r{g}}{n}{a}{\textsubring{r}}{s}{\textsubring{d}}{\textopeno}{\textsubring{b}}{n}{\textscy}{\textsubring{n}}} \dicPos{f}[7] \dicFlx{(‑unar, ‑anir)}[8] \dicSynonym{góðgerðastofnun} \dicDirectTranslationCS{dobročinná\,/\addthin charitativní organizace}
\dicEntry[líkneski] \dicTerm{lík··neski} \dicIPA{{l}{i}{h}{\r{g}}{n}{\textepsilon}{s}{\r{\textObardotlessj}}{\textsci}} \dicPos{n}[2] \dicFlx{(‑s, ‑)}[16] \dicSynonym{höggmynd} \dicDirectTranslationCS{podobizna}
\dicEntry[líknsamur] \dicTerm{líkn··|samur} \dicIPA{{l}{i}{h}{\r{g}}{\textsubring{n}}{s}{a}{m}{\textscy}{\textsubring{r}}} \dicPos{adj}[1] \dicFlx{(f ‑söm)}[2] \dicDirectTranslationCS{milosrdný, milostivý, soucitný}
\dicEntry[líkræða] \dicTerm{lík··ræð|a} \dicIPA{{l}{i}{\textlengthmark}{\r{g}}{r}{a}{i}{ð}{a}} \dicPos{f}[1] \dicFlx{(‑u, ‑ur)}[19] \dicDirectTranslationCS{smuteční proslov}
\dicEntry[líkskoðun] \dicTerm{lík··skoð|un} \dicIPA{{l}{i}{\textlengthmark}{\r{g}}{s}{\r{g}}{\textopeno}{ð}{\textscy}{\textsubring{n}}} \dicPos{f}[7] \dicFlx{(‑unar)}[9] \dicDirectTranslationCS{ohledání mrtvoly}
\dicEntry[líkskurður] \dicTerm{lík··skurð|ur} \dicIPA{{l}{i}{\textlengthmark}{\r{g}}{s}{\r{g}}{\textscy}{r}{ð}{\textscy}{\textsubring{r}}} \dicPos{m}[10] \dicFlx{(‑ar, ‑ir)}[4] \dicDirectTranslationCS{pitva}
\dicEntry[líkur] \dicTerm{líkur\smash{\textsuperscript{1}}} \dicsymFrequent\  \dicIPA{{l}{i}{\textlengthmark}{\r{g}}{\textscy}{\textsubring{r}}} \dicPos{f}[12] \dicFlx{pl}[5] \dicSynonym{líkindi} \dicDirectTranslationCS{šance, pravděpodobnost} \dicExampleIS{líkur á e‑u\,/\addthin líkur til e‑s} \dicExampleCS{šance na (co)};  \dicPhraseIS{ef að líkum lætur} \dicFlx{adv} \dicDirectTranslationCS{pravděpodobně, patrně}
\dicEntry[líkur] \dicTerm{líkur\smash{\textsuperscript{2}}} \dicsymFrequent\  \dicIPA{{l}{i}{\textlengthmark}{\r{g}}{\textscy}{\textsubring{r}}} \dicPos{adj}[1]\dicFlx{}[-1] \textbf{1.} \dicDirectTranslationCS{podobný} \dicExampleIS{lík vandamál} \dicExampleCS{podobné problémy};  \dicPhraseIS{vera líkur e‑m} \dicDirectTranslationCS{být (komu) podobný, podobat se (komu)}  \textbf{2.} \dicDirectTranslationCS{obdobný, analogický};  \dicIdiom{líkur}{ \dicPhraseIS{líkt og}} \dicFlx{conj} \dicDirectTranslationCS{jak, jako} \dicExampleIS{líta út líkt og Íslendingur} \dicExampleCS{vypadat jako Islanďan}
\dicEntry[líkþorn] \dicTerm{lík··þorn} \dicIPA{{l}{i}{\textlengthmark}{\r{g}}{\texttheta}{\textopeno}{r}{\textsubring{d}}{\textsubring{n}}} \dicPos{n}[2] \dicFlx{(‑s, ‑)}[5] \dicDirectTranslationCS{kuří oko}
\dicEntry[líkþrá] \dicTerm{lík··þrá} \dicIPA{{l}{i}{\textlengthmark}{\r{g}}{\texttheta}{r}{au}} \dicPos{f}[4] \dicFlx{(‑r)}[19] \dicFieldCat{med.} \dicDirectTranslationCS{lepra, malomocenství}
\dicEntry[líkþrár] \dicTerm{lík··þrár} \dicIPA{{l}{i}{\textlengthmark}{\r{g}}{\texttheta}{r}{au}{\textsubring{r}}} \dicPos{adj}[4]\dicFlx{}[-3] \dicDirectTranslationCS{malomocný}
\dicEntry[lím] \dicTerm{lím} \dicIPA{{l}{i}{\textlengthmark}{\textsubring{m}}} \dicPos{n}[2] \dicFlx{(‑s, ‑)}[5] \dicDirectTranslationCS{lepidlo, lep, klih}
\dicEntry[líma] \dicTerm{lím|a} \dicsymFrequent\  \dicIPA{{l}{i}{\textlengthmark}{m}{a}} \dicPos{v}[2] \dicFlx{(‑di, ‑t)}[141] \dicFlx{acc} \dicDirectTranslationCS{(na)lepit, nalepovat, přilepit, přilepovat} \dicExampleIS{líma mynd á blað} \dicExampleCS{nalepit obrázek na papír};  \dicIdiom{límast}{ \dicPhraseIS{límast}} \dicFlx{refl} \dicDirectTranslationCS{(při)lepit se}
\dicEntry[límband] \dicTerm{lím··|band} \dicIPA{{l}{i}{m}{\textsubring{b}}{a}{n}{\textsubring{d}}} \dicPos{n}[2] \dicFlx{(‑bands, ‑bönd)}[8] \dicDirectTranslationCS{lepicí páska, izolepa}
\dicEntry[límmiði] \dicTerm{lím··mið|i} \dicIPA{{l}{i}{m}{\textlengthmark}{\textsci}{ð}{\textsci}} \dicPos{m}[1] \dicFlx{(‑a, ‑ar)}[1] \dicDirectTranslationCS{nálepka, štítek, etiketa}
\dicEntry[lín] \dicTerm{lín} \dicIPA{{l}{i}{\textlengthmark}{\textsubring{n}}} \dicPos{n}[2] \dicFlx{(‑s)}[2] \textbf{1.} \dicFieldCat{bot.} \dicSynonym{hör} \dicDirectTranslationCS{len setý} \textit{(l.~{\textLA{Linum usitatissimum}})}  \dicsymPhoto\   \textbf{2.} \dicSynonym{léreft} \dicDirectTranslationCS{len, lněné plátno} \dicExampleIS{spinna lín} \dicExampleCS{tkát len}
\dicFigure{ds_image_lin_0_2.jpg}{Lín}{Lín - Franz Eugen Kohler, PD}
\dicEntry[lína] \dicTerm{lín|a} \dicsymFrequent\  \dicIPA{{l}{i}{\textlengthmark}{n}{a}} \dicPos{f}[1] \dicFlx{(‑u, ‑ur)}[7] \textbf{1.} \dicSynonym{strik} \dicDirectTranslationCS{čára, linka, linie}  \textbf{2.} \dicSynonym*{textalína} \dicDirectTranslationCS{řádek, řádka} \dicExampleIS{nokkrar línur úr textanum} \dicExampleCS{několik řádek z~textu}  \textbf{3.} \dicSynonym{band} \dicDirectTranslationCS{lano, šňůra, provázek}  \textbf{4.} \dicSynonym{lóð\smash{\textsuperscript{2}}} \dicDirectTranslationCS{lanová síť}  \textbf{5.} \dicSynonym{stefna\smash{\textsuperscript{1}}} \dicDirectTranslationCS{politický směr, politická linie}  \textbf{6.} \dicFieldCat{mat.} \dicDirectTranslationCS{přímka};  \dicIdiom{lína}{ \dicPhraseIS{lesa milli línanna}} \dicLangCat{přen.} \dicDirectTranslationCS{číst mezi řádky} \dicIndirectTranslationCS{(porozumět více, než je napsáno)}; { \dicPhraseIS{passa upp á línurnar}} \dicLangCat{přen.} \dicDirectTranslationCS{dávat pozor na linii} \dicIndirectTranslationCS{(snažit se neztloustnout)}
\dicEntry[línolía] \dicTerm{lín··olí|a} \dicIPA{{l}{i}{\textlengthmark}{n}{\textopeno}{l}{i}{j}{a}} \dicPos{f}[1] \dicFlx{(‑u)}[5] \dicDirectTranslationCS{lněný olej}
\dicEntry[línubil] \dicTerm{línu··bil} \dicIPA{{l}{i}{\textlengthmark}{n}{\textscy}{\textsubring{b}}{\textsci}{\textsubring{l}}} \dicPos{n}[2] \dicFlx{(‑s, ‑)}[5] \dicFieldCat{poč.} \dicDirectTranslationCS{řádkování}
\dicEntry[línumaður] \dicTerm{línu··|maður} \dicIPA{{l}{i}{\textlengthmark}{n}{\textscy}{m}{a}{ð}{\textscy}{\textsubring{r}}} \dicPos{m}[13] \dicFlx{(‑manns, ‑menn)}[2] \dicFieldCat{sport.} \dicDirectTranslationCS{krajní útočník\,/\addthin útočnice, křídlo}
\dicEntry[línurit] \dicTerm{línu··rit} \dicIPA{{l}{i}{\textlengthmark}{n}{\textscy}{r}{\textsci}{\textsubring{d}}} \dicPos{n}[2] \dicFlx{(‑s, ‑)}[5] \dicDirectTranslationCS{graf, diagram} \dicExampleIS{töflur og línurit} \dicExampleCS{tabulky a~grafy}
\dicEntry[línuskauti] \dicTerm{línu··skaut|i} \dicIPA{{l}{i}{\textlengthmark}{n}{\textscy}{s}{\r{g}}{\oe i}{\textsubring{d}}{\textsci}} \dicPos{m}[1] \dicFlx{(‑a, ‑ar)}[1] \dicPhraseIS{línuskautar} \dicFlx{pl} \dicFieldCat{sport.} \dicDirectTranslationCS{kolečkové brusle}
\dicEntry[línuskil] \dicTerm{línu··skil} \dicIPA{{l}{i}{\textlengthmark}{n}{\textscy}{s}{\r{\textObardotlessj}}{\textsci}{\textsubring{l}}} \dicPos{n}[2] \dicFlx{pl}[1] \dicFieldCat{poč.} \dicDirectTranslationCS{zalomení řádku}
\dicEntry[líra] \dicTerm{lír|a} \dicIPA{{l}{i}{\textlengthmark}{r}{a}} \dicPos{f}[1] \dicFlx{(‑u, ‑ur)}[7] \textbf{1.} \dicDirectTranslationCS{lira} \dicIndirectTranslationCS{(bývalá italská měna)}  \textbf{2.} \dicFieldCat{hud.} \dicDirectTranslationCS{lyra}
\dicEntry[lít] \dicTerm{lít} \dicIPA{{l}{i}{\textlengthmark}{\textsubring{d}}} \dicPos{v} \dicFlx{ind praes sg 1 pers} \dicLink{líta}
\dicEntry[líta] \dicTerm{líta} \dicsymFrequent\  \dicIPA{{l}{i}{\textlengthmark}{\textsubring{d}}{a}} \dicPos{v}[6] \dicFlx{(lít, leit, litum, liti, litið)}[76] \dicFlx{acc} \dicDirectTranslationCS{(po)dívat se, hledět, pohlédnout, pohlížet} \dicExampleIS{líta á klukkuna} \dicExampleCS{podívat se na hodiny};  \dicPhraseIS{líta e‑n} \dicDirectTranslationCS{pohlédnout na (koho)};  \dicPhraseIS{líta á e‑n} \dicDirectTranslationCS{podívat se na (koho)} \dicExampleIS{líta spyrjandi á hana} \dicExampleCS{podívat se na ni tázavě};  \dicIdiom{líta}[af]{ \dicPhraseIS{líta af e‑u\,/\addthin e‑m}} \dicDirectTranslationCS{moct spustit oči z~(čeho\,/\addthin koho)} \dicIndirectTranslationCS{(většinou v~záporu)};  \dicIdiom{líta}[á]{ \dicPhraseIS{líta á e‑ð}} \dicDirectTranslationCS{podívat se na (co), kouknout na (co), prozkoumat (co)} \dicExampleIS{líta á vélina} \dicExampleCS{kouknout na motor};  \dicPhraseIS{líta stórt á sig} \dicDirectTranslationCS{mít o~sobě vysoké mínění};  \dicIdiom{líta}[eftir]{ \dicPhraseIS{líta eftir e‑m}} \dicDirectTranslationCS{dohlédnout\,/\addthin dohlížet na (koho)} \dicExampleIS{líta eftir barninu} \dicExampleCS{dohlížet na dítě};  \dicIdiom{líta}[fram hjá]{ \dicPhraseIS{líta fram hjá e‑u}} \dicDirectTranslationCS{přehlédnout\,/\addthin přehlížet (co)} \dicExampleIS{líta fram hjá mikilvægum staðreyndum} \dicExampleCS{přehlížet důležitá fakta};  \dicIdiom{líta}[inn]{ \dicPhraseIS{líta inn hjá e‑m}} \dicDirectTranslationCS{(na)kouknout ke (komu), (za)stavit se u~(koho) (na návštěvu ap.)};  \dicIdiom{líta}[í]{ \dicPhraseIS{líta í e‑ð}} \dicDirectTranslationCS{podívat se do (čeho), nahlédnout do (čeho) (novin ap.)};  \dicIdiom{líta}[niður]{ \dicPhraseIS{líta niður á e‑n}} \dicDirectTranslationCS{dívat se na (koho) spatra, shlížet na (koho) povýšeně};  \dicIdiom{líta}[til]{ \dicPhraseIS{líta til e‑rs}} \dicDirectTranslationCS{zastavit se u~(koho), podívat se ke (komu)};  \dicIdiom{líta}[undan]{ \dicPhraseIS{líta undan}} \dicDirectTranslationCS{sklopit zrak, podívat se dolů};  \dicIdiom{líta}[upp]{ \dicPhraseIS{líta upp til e‑rs}} \dicDirectTranslationCS{obdivovat (koho), vzhlížet ke (komu) (s~úctou ap.)};  \dicIdiom{líta}[út]{ \dicPhraseIS{líta út}} \dicDirectTranslationCS{vypadat, vyhlížet} \dicExampleIS{líta vel út} \dicExampleCS{vypadat dobře}; { \dicPhraseIS{líta út fyrir að (vera ríkur)}} \dicDirectTranslationCS{vypadat na (boháče), vypadat jako (bohatý)};  \dicIdiom{líta}[við]{ \dicPhraseIS{líta ekki við e‑u}} \dicDirectTranslationCS{netknout se (čeho), nejevit zájem o~(co)}; { \dicPhraseIS{líta við (þar)}} \dicDirectTranslationCS{stavit se (tam) na skok, zaskočit (tam)};  \dicIdiom{líta}[yfir]{ \dicPhraseIS{líta yfir e‑ð}} \dicDirectTranslationCS{projít (co), prohlédnout (co)} \dicExampleIS{líta yfir textann} \dicExampleCS{projít text};  \dicIdiom{lítast}[á]{ \dicPhraseIS{e‑m líst vel á e‑ð}} \dicFlx{refl impers} \dicDirectTranslationCS{(co) se (komu) líbí}
\dicEntry[lítið] \dicTerm{lítið} \dicsymFrequent\  \dicIPA{{l}{i}{\textlengthmark}{\textsubring{d}}{\textsci}{\texttheta}} \dicPos{adv} \dicFlx{(comp minna, sup minnst)} \dicDirectTranslationCS{málo, trochu} \dicExampleIS{vera lítið áberandi} \dicExampleCS{být trochu nápadný}
\dicEntry[lítilfjörlegur] \dicTerm{lítil··fjör·legur} \dicIPA{{l}{i}{\textlengthmark}{\textsubring{d}}{\textsci}{l}{f}{j}{\oe}{r}{l}{\textepsilon}{\textbabygamma}{\textscy}{\textsubring{r}}} \dicPos{adj}[1]\dicFlx{}[-8] \textbf{1.} \dicSynonym{ómerkilegur} \dicDirectTranslationCS{nezajímavý, nezáživný, všední}  \textbf{2.} \dicSynonym{lasinn} \dicDirectTranslationCS{indisponovaný, slabý}
\dicEntry[lítill] \dicTerm{lítill} \dicsymFrequent\  \dicIPA{{l}{i}{\textlengthmark}{\textsubring{d}}{\textsci}{\textsubring{d}}{\textsubring{l}}} \dicPos{adj}[11] \dicFlx{(comp minni, sup minnstur)}[10] \textbf{1.} \dicSynonym{smár} \dicDirectTranslationCS{malý, drobný} \dicExampleIS{lítill maður} \dicExampleCS{malý člověk} \dicAntonym{stór}  \textbf{2.} \dicSynonym*{þann yngri til aðgreiningar frá þeim eldri} \dicDirectTranslationCS{malý, mladší} \dicExampleIS{Sölvi litli} \dicExampleCS{malý Sölvi}  \textbf{3.} \dicSynonym{óverulegur} \dicDirectTranslationCS{malý, zanedbatelný, nepatrný}  \textbf{4.} \dicSynonym{óduglegur} \dicDirectTranslationCS{zahálčivý, málo snaživý} \dicExampleIS{lítill íþróttamaður} \dicExampleCS{malý sportovec};  \dicIdiom{lítill}{ \dicPhraseIS{vera lítill í sér}} \dicDirectTranslationCS{být zbabělý}
\dicEntry[lítillátur] \dicTerm{lítil··látur} \dicIPA{{l}{i}{\textlengthmark}{\textsubring{d}}{\textsci}{l}{au}{\textsubring{d}}{\textscy}{\textsubring{r}}} \dicPos{adj}[1]\dicFlx{}[-1] \textbf{1.} \dicSynonym*{hrokalaus} \dicDirectTranslationCS{pokorný}  \textbf{2.} \dicSynonym{auðmjúkur} \dicDirectTranslationCS{nenáročný, skromný}
\dicEntry[lítillækka] \dicTerm{lítil··lækk|a} \dicIPA{{l}{i}{\textlengthmark}{\textsubring{d}}{\textsci}{l}{a}{i}{h}{\r{g}}{a}} \dicPos{v}[1] \dicFlx{(‑aði)}[1] \dicFlx{acc} \dicDirectTranslationCS{pokořit, pokořovat, ponížit, ponižovat} \dicExampleIS{lítillækka sig} \dicExampleCS{ponížit se}
\dicEntry[lítillækkun] \dicTerm{lítil··lækk|un} \dicIPA{{l}{i}{\textlengthmark}{\textsubring{d}}{\textsci}{l}{a}{i}{h}{\r{g}}{\textscy}{\textsubring{n}}} \dicPos{f}[7] \dicFlx{(‑unar)}[9] \dicDirectTranslationCS{pokoření, ponížení}
\dicEntry[lítillæti] \dicTerm{lítil··læti} \dicIPA{{l}{i}{\textlengthmark}{\textsubring{d}}{\textsci}{l}{a}{i}{\textsubring{d}}{\textsci}} \dicPos{n}[2] \dicFlx{(‑s)}[20] \textbf{1.} \dicSynonym{yfirlætisleysi} \dicDirectTranslationCS{pokora}  \textbf{2.} \dicSynonym{auðmýkt} \dicDirectTranslationCS{skromnost}
\dicEntry[lítilmagni] \dicTerm{lítil··magn|i} \dicIPA{{l}{i}{\textsubring{d}}{\textsci}{l}{m}{a}{\r{g}}{n}{\textsci}} \dicPos{m}[1] \dicFlx{(‑a, ‑ar)}[8] \dicDirectTranslationCS{padavka, outsider(ka)}
\dicEntry[lítilmenni] \dicTerm{lítil··menni} \dicIPA{{l}{i}{\textlengthmark}{\textsubring{d}}{\textsci}{l}{m}{\textepsilon}{n}{\textsci}} \dicPos{n}[2] \dicFlx{(‑s, ‑)}[14] \textbf{1.} \dicDirectTranslationCS{malý člověk (vzrůstem)}  \textbf{2.} \dicDirectTranslationCS{malicherný\,/\addthin nanicovatý člověk}
\dicEntry[lítilmótlegur] \dicTerm{lítil··mót·legur} \dicIPA{{l}{i}{\textlengthmark}{\textsubring{d}}{\textsci}{l}{m}{ou}{\textsubring{d}}{l}{\textepsilon}{\textbabygamma}{\textscy}{\textsubring{r}}} \dicPos{adj}[1]\dicFlx{}[-8] \dicSynonym{ómerkilegur} \dicDirectTranslationCS{ubohý, mizerný, mrzký}
\dicEntry[lítilræði] \dicTerm{lítil··ræði} \dicIPA{{l}{i}{\textlengthmark}{\textsubring{d}}{\textsci}{l}{r}{a}{i}{ð}{\textsci}} \dicPos{n}[2] \dicFlx{(‑s)}[20] \dicDirectTranslationCS{maličkost, maličko}
\dicEntry[lítilsháttar] \dicTerm{lítils··háttar\smash{\textsuperscript{1}}} \dicIPA{{l}{i}{\textlengthmark}{\textsubring{d}}{\textsci}{l}{s}{h}{au}{h}{\textsubring{d}}{a}{\textsubring{r}}} \dicPos{adj}[13] \dicFlx{indecl}[1] \dicDirectTranslationCS{malý, drobný}
\dicEntry[lítilsháttar] \dicTerm{lítils··háttar\smash{\textsuperscript{2}}} \dicIPA{{l}{i}{\textlengthmark}{\textsubring{d}}{\textsci}{l}{s}{h}{au}{h}{\textsubring{d}}{a}{\textsubring{r}}} \dicPos{adv} \dicDirectTranslationCS{málo, trochu}
\dicEntry[lítilsigldur] \dicTerm{lítil··sigldur} \dicIPA{{l}{i}{\textlengthmark}{\textsubring{d}}{\textsci}{l}{s}{\textsci}{l}{\textsubring{d}}{\textscy}{\textsubring{r}}} \dicPos{adj}[2]\dicFlx{}[-14] \dicSynonym{lítilfjörlegur} \dicDirectTranslationCS{nevýznamný, ubohý, nicotný}
\dicEntry[lítilsvirða] \dicTerm{lítils··vir|ða} \dicIPA{{l}{i}{\textlengthmark}{\textsubring{d}}{\textsci}{l}{s}{v}{\textsci}{r}{ð}{a}} \dicPos{v}[2] \dicFlx{(‑ti, ‑t)}[52] \dicFlx{acc} \dicDirectTranslationCS{pohrdat, opovrhovat}
\dicEntry[lítilsvirðing] \dicTerm{lítils··virð·ing} \dicIPA{{l}{i}{\textlengthmark}{\textsubring{d}}{\textsci}{l}{s}{v}{\textsci}{r}{ð}{i}{\ng}{\r{g}}} \dicPos{f}[4] \dicFlx{(‑ar)}[7] \dicSynonym{óvirðing} \dicDirectTranslationCS{opovržení, pohrdání}
\dicEntry[lítilvægur] \dicTerm{lítil··vægur} \dicIPA{{l}{i}{\textlengthmark}{\textsubring{d}}{\textsci}{l}{v}{a}{i}{\textbabygamma}{\textscy}{\textsubring{r}}} \dicPos{adj}[1]\dicFlx{}[-1] \dicSynonym*{þýðingarlítill} \dicDirectTranslationCS{nedůležitý, nepodstatný, bezvýznamný}
\dicEntry[lítilþægur] \dicTerm{lítil··þægur} \dicIPA{{l}{i}{\textlengthmark}{\textsubring{d}}{\textsci}{l}{\texttheta}{a}{i}{\textbabygamma}{\textscy}{\textsubring{r}}} \dicPos{adj}[1]\dicFlx{}[-1] \dicDirectTranslationCS{skromný, nenáročný} \dicIndirectTranslationCS{(často používané v~záporu)}
\dicEntry[lítri] \dicTerm{lítr|i} \dicIPA{{l}{i}{\textlengthmark}{\textsubring{d}}{r}{\textsci}} \dicPos{m}[1] \dicFlx{(‑a, ‑ar)}[1] \dicDirectTranslationCS{litr}
\dicEntry[lítt] \dicTerm{lítt} \dicsymFrequent\  \dicIPA{{l}{i}{h}{\textsubring{d}}} \dicPos{adv} \dicFlx{(comp miður, sup minnst)} \textbf{1.} \dicSynonym{lítið} \dicDirectTranslationCS{trochu, málo} \dicExampleIS{lítt þekktur ráðherra} \dicExampleCS{málo známý ministr}  \textbf{2.} \dicSynonym{illa} \dicDirectTranslationCS{špatně}
\dicEntry[ljá] \dicTerm{ljá} \dicIPA{{l}{j}{au}{\textlengthmark}} \dicPos{v}[5] \dicFlx{(ljæ, léði, léðum, léði, léð)}[15] \dicFlx{dat + acc} \textbf{1.} \dicLangCat{zast.} \dicDirectTranslationCS{(vy)půjčit} \dicExampleIS{ljá e‑m e‑ð} \dicExampleCS{půjčit (komu co)}  \textbf{2.} \dicPhraseIS{ljá e‑m eyra} \dicLangCat{přen.} \dicDirectTranslationCS{vyslyšet (koho), dopřát (komu) sluchu}  \textbf{3.} \dicDirectTranslationCS{vnést, vnášet, přidat, přidávat} \dicIndirectTranslationCS{(způsobit uplatnění, začlenění něčeho)}  \textbf{4.} \dicPhraseIS{ljá máls á e‑u} \dicDirectTranslationCS{zavést téma na (co)}
\dicEntry[ljár] \dicTerm{ljá|r} \dicIPA{{l}{j}{au}{\textlengthmark}{\textsubring{r}}} \dicPos{m}[9] \dicFlx{(‑s, ‑ir)}[12] \dicDirectTranslationCS{kosa} \dicIndirectTranslationCS{(ruční nástroj)};  \dicPhraseIS{slá með orfi og ljá} \dicDirectTranslationCS{sekat kosou};  \dicPhraseIS{vera e‑m óþægur ljár í þúfu} \dicLangCat{přen.} \dicDirectTranslationCS{být (komu) trnem v~oku}
\dicEntry[ljóð] \dicTerm{ljóð} \dicsymFrequent\  \dicIPA{{l}{j}{ou}{\textlengthmark}{\texttheta}} \dicPos{n}[2] \dicFlx{(‑s, ‑)}[5] \dicDirectTranslationCS{báseň} \dicExampleIS{ljóð eftir e‑n} \dicExampleCS{(čí) báseň}
\dicEntry[ljóðabók] \dicTerm{ljóða··|bók} \dicIPA{{l}{j}{ou}{\textlengthmark}{ð}{a}{\textsubring{b}}{ou}{\r{g}}} \dicPos{f}[8] \dicFlx{(‑bókar, ‑bækur)}[5] \dicDirectTranslationCS{sbírka\,/\addthin kniha poezie}
\dicEntry[ljóðaháttur] \dicTerm{ljóða··|háttur} \dicIPA{{l}{j}{ou}{\textlengthmark}{ð}{a}{h}{au}{h}{\textsubring{d}}{\textscy}{\textsubring{r}}} \dicPos{m}[12] \dicFlx{(‑háttar, ‑hættir)}[7] \dicFieldCat{lit.} \dicDirectTranslationCS{ed\nobreak dické metrum (šestiverší)}
\dicEntry[ljóðlist] \dicTerm{ljóð··list} \dicIPA{{l}{j}{ou}{ð}{l}{\textsci}{s}{\textsubring{d}}} \dicPos{f}[7] \dicFlx{(‑ar)}[3] \dicDirectTranslationCS{básnictví, básnické umění}
\dicEntry[ljóðmæli] \dicTerm{ljóð··mæli} \dicIPA{{l}{j}{ou}{ð}{m}{a}{i}{l}{\textsci}} \dicPos{n}[2] \dicFlx{pl}[19] \dicDirectTranslationCS{sbírka poezie}
\dicEntry[ljóðræna] \dicTerm{ljóð··ræn|a} \dicIPA{{l}{j}{ou}{ð}{r}{a}{i}{n}{a}} \dicPos{f}[1] \dicFlx{(‑u)}[5] \dicFieldCat{lit.} \dicDirectTranslationCS{lyrika}
\dicEntry[ljóðrænn] \dicTerm{ljóð··rænn} \dicIPA{{l}{j}{ou}{ð}{r}{a}{i}{\textsubring{d}}{\textsubring{n}}} \dicPos{adj}[7]\dicFlx{}[-1] \dicDirectTranslationCS{poetický, lyrický}
\dicEntry[ljóðskáld] \dicTerm{ljóð··skáld} \dicIPA{{l}{j}{ou}{ð}{s}{\r{g}}{au}{l}{\textsubring{d}}} \dicPos{n}[2] \dicFlx{(‑s, ‑)}[5] \dicDirectTranslationCS{básník, básnířka, poeta, poetka}
\dicEntry[ljóður] \dicTerm{ljóð|ur} \dicIPA{{l}{j}{ou}{\textlengthmark}{ð}{\textscy}{\textsubring{r}}} \dicPos{m}[6] \dicFlx{(‑s)}[26] \dicSynonym{lýti} \dicDirectTranslationCS{kaz, vada}
\dicEntry[ljóma] \dicTerm{ljóm|a} \dicIPA{{l}{j}{ou}{\textlengthmark}{m}{a}} \dicPos{v}[1] \dicFlx{(‑aði)}[1] \dicDirectTranslationCS{zářit, vyzařovat} \dicExampleIS{Sólin ljómar á himninum.} \dicExampleCS{Slunce září na obloze.};  \dicPhraseIS{ljóma af gleði} \dicLangCat{přen.} \dicDirectTranslationCS{zářit štěstím}
\dicEntry[ljómandi] \dicTerm{ljóm··andi\smash{\textsuperscript{1}}} \dicIPA{{l}{j}{ou}{\textlengthmark}{m}{a}{n}{\textsubring{d}}{\textsci}} \dicPos{adj}[13] \dicFlx{indecl}[1] \dicDirectTranslationCS{zářivý, zářící} \dicExampleIS{ljómandi sól} \dicExampleCS{zářící slunce}
\dicEntry[ljómandi] \dicTerm{ljóm··andi\smash{\textsuperscript{2}}} \dicIPA{{l}{j}{ou}{\textlengthmark}{m}{a}{n}{\textsubring{d}}{\textsci}} \dicPos{adv} \dicDirectTranslationCS{úžasně, fantasticky} \dicExampleIS{ljómandi vel} \dicExampleCS{fantasticky dobře}
\dicEntry[ljómi] \dicTerm{ljóm|i} \dicIPA{{l}{j}{ou}{\textlengthmark}{m}{\textsci}} \dicPos{m}[1] \dicFlx{(‑a)}[3] \textbf{1.} \dicSynonym{birta\smash{\textsuperscript{1}}} \dicDirectTranslationCS{zář(e), svit}  \textbf{2.} \dicFieldCat{fyz.} \dicDirectTranslationCS{jas}
\dicEntry[ljón] \dicTerm{ljón} \dicsymFrequent\  \dicIPA{{l}{j}{ou}{\textlengthmark}{\textsubring{n}}} \dicPos{n}[2] \dicFlx{(‑s, ‑)}[5] \textbf{1.} \dicFieldCat{zool.} \dicDirectTranslationCS{lev} \textit{(l.~{\textLA{Panthera leo}})}  \dicsymPhoto\  \dicExampleIS{Ljón hleypur í átt að manni.} \dicExampleCS{Směrem k~člověku běží lev.}  \textbf{2.} \dicDirectTranslationCS{Lev} \dicIndirectTranslationCS{(znamení zvěrokruhu)} \textit{(l.~{\textLA{Leo}})}
\dicFigure{20786.jpg}{Ljón}{Ljón - Stansell, Ken, Biolib, PD}
\dicEntry[ljónatemjari] \dicTerm{ljóna··temj·ar|i} \dicIPA{{l}{j}{ou}{\textlengthmark}{n}{a}{t\smash{\textsuperscript{h}}}{\textepsilon}{m}{j}{a}{r}{\textsci}} \dicPos{m}[1] \dicFlx{(‑a, ‑ar)}[13] \dicDirectTranslationCS{krotitel(ka) lvů\,/\addthin šelem\,/\addthin dravé zvěře}
\dicEntry[ljónslappi] \dicTerm{ljóns··lapp|i} \dicIPA{{l}{j}{ou}{n}{s}{l}{a}{h}{\textsubring{b}}{\textsci}} \dicPos{m}[1] \dicFlx{(‑a)}[3] \dicFieldCat{bot.} \dicDirectTranslationCS{kontryhel alpinský} \textit{(l.~{\textLA{Alchemilla alpina}})}  \dicsymPhoto\ 
\dicFigure{88632.jpg}{Ljónslappi}{Ljónslappi - Pleva František, Biolib, PD}
\dicEntry[ljónynja] \dicTerm{ljón··ynj|a} \dicIPA{{l}{j}{ou}{\textlengthmark}{n}{\textsci}{n}{j}{a}} \dicPos{f}[1] \dicFlx{(‑u, ‑ur)}[7] \dicDirectTranslationCS{lvice}
\dicEntry[ljóri] \dicTerm{ljór|i} \dicIPA{{l}{j}{ou}{\textlengthmark}{r}{\textsci}} \dicPos{m}[1] \dicFlx{(‑a, ‑ar)}[1] \dicDirectTranslationCS{světlík}
\dicEntry[ljós] \dicTerm{ljós\smash{\textsuperscript{1}}} \dicsymFrequent\  \dicIPA{{l}{j}{ou}{\textlengthmark}{s}} \dicPos{n}[2] \dicFlx{(‑s, ‑)}[5] \textbf{1.} \dicSynonym{birta\smash{\textsuperscript{1}}} \dicDirectTranslationCS{světlo, záře, jas, svit} \dicExampleIS{ljósið frá lampanum} \dicExampleCS{světlo lampy};  \dicPhraseIS{kveikja ljósið} \dicDirectTranslationCS{rozsvítit, zapnout\,/\addthin zapálit světlo};  \dicPhraseIS{slökkva ljósið} \dicDirectTranslationCS{zhasnout, vypnout světlo}  \textbf{2.} \dicSynonym*{þægt barn} \dicDirectTranslationCS{zlatíčko, andílek} \dicIndirectTranslationCS{(o~poslušném dítěti)};  \dicIdiom{ljós}{ \dicPhraseIS{í nýju ljósi}} \dicFlx{adv} \dicLangCat{přen.} \dicDirectTranslationCS{v~novém světle}; { \dicPhraseIS{það kom í ljós að}} \dicDirectTranslationCS{ukázalo se, že; vyšlo najevo, že}
\dicEntry[ljós] \dicTerm{ljós\smash{\textsuperscript{2}}} \dicsymFrequent\  \dicIPA{{l}{j}{ou}{\textlengthmark}{s}} \dicPos{adj}[5]\dicFlx{}[-1] \dicDirectTranslationCS{jasný, světlý} \dicExampleIS{vera ljós á húð} \dicExampleCS{mít světlou pokožku};  \dicPhraseIS{gera sér e‑ð ljóst} \dicDirectTranslationCS{uvědomit si (co)};  \dicPhraseIS{e‑m er ljóst e‑að} \dicFlx{impers} \dicDirectTranslationCS{(co) je (komu) jasné};  \dicPhraseIS{e‑að er deginum ljósara} \dicLangCat{přen.} \dicDirectTranslationCS{(co) je nad slunce jasnější, (co) je jasné jak facka}
\dicEntry[ljósakróna] \dicTerm{ljósa··krón|a} \dicIPA{{l}{j}{ou}{\textlengthmark}{s}{a}{k\smash{\textsuperscript{h}}}{r}{ou}{n}{a}} \dicPos{f}[1] \dicFlx{(‑u, ‑ur)}[7] \dicDirectTranslationCS{(mnohoramenný) lustr}
\dicEntry[ljósapera] \dicTerm{ljósa··per|a} \dicIPA{{l}{j}{ou}{\textlengthmark}{s}{a}{p\smash{\textsuperscript{h}}}{\textepsilon}{r}{a}} \dicPos{f}[1] \dicFlx{(‑u, ‑ur)}[7] \dicFieldCat{elek.} \dicDirectTranslationCS{žárovka}
\dicEntry[ljósaskilti] \dicTerm{ljósa··skilti} \dicIPA{{l}{j}{ou}{\textlengthmark}{s}{a}{s}{\r{\textObardotlessj}}{\textsci}{\textsubring{l}}{\textsubring{d}}{\textsci}} \dicPos{n}[2] \dicFlx{(‑s, ‑)}[14] \dicDirectTranslationCS{neonová reklama}
\dicEntry[ljósaskipti] \dicTerm{ljósa··skipti} \dicIPA{{l}{j}{ou}{\textlengthmark}{s}{a}{s}{\r{\textObardotlessj}}{\textsci}{f}{\textsubring{d}}{\textsci}} \dicPos{n}[2] \dicFlx{pl}[19] \dicSynonym{rökkur} \dicDirectTranslationCS{soumrak}
\dicEntry[ljósastaur] \dicTerm{ljósa··staur} \dicIPA{{l}{j}{ou}{\textlengthmark}{s}{a}{s}{\textsubring{d}}{\oe i}{\textsubring{r}}} \dicPos{m}[4] \dicFlx{(‑s, ‑ar)}[14] \dicDirectTranslationCS{kandelábr, stojan veřejného osvětlení}
\dicEntry[ljósastæði] \dicTerm{ljósa··stæði} \dicIPA{{l}{j}{ou}{\textlengthmark}{s}{a}{s}{\textsubring{d}}{a}{i}{ð}{\textsci}} \dicPos{n}[2] \dicFlx{(‑s, ‑)}[14] \dicDirectTranslationCS{přípojka na světlo}
\dicEntry[Ljósálfheimur] \dicTerm{Ljós··álf·heim|ur}\dicTerm{, Álfheimur} \dicIPA{{l}\-{j}\-{ou}\-{\textlengthmark}\-{s}\-{au}\-{l}\-{f}\-{h}\-{ei}\-{m}\-{\textscy}\-{\textsubring{r}}\-} \dicPos{m}[6] \dicFlx{(‑s)}[3] \dicFieldCat{myt.} \dicDirectTranslationCS{Álfheim} \dicIndirectTranslationCS{(říše světla, světlých elfů)} \dicAntonym{Svartálfheimur}
\dicEntry[ljósár] \dicTerm{ljós··ár} \dicIPA{{l}{j}{ou}{\textlengthmark}{s}{au}{\textsubring{r}}} \dicPos{n}[2] \dicFlx{(‑s, ‑)}[5] \dicFieldCat{astro.} \dicDirectTranslationCS{světelný rok}
\dicEntry[ljósberi] \dicTerm{ljós··ber|i} \dicIPA{{l}{j}{ou}{s}{\textsubring{b}}{\textepsilon}{r}{\textsci}} \dicPos{m}[1] \dicFlx{(‑a, ‑ar)}[1] \dicFieldCat{bot.} \dicDirectTranslationCS{kohoutek alpský} \textit{(l.~{\textLA{Lychnis alpina}})}  \dicsymPhoto\ 
\dicFigure{72890.jpg}{Ljósberi}{Ljósberi - Kesl Michael, Biolib, Copyright/CC-BY-NC}
\dicEntry[ljóseind] \dicTerm{ljós··eind} \dicIPA{{l}{j}{ou}{\textlengthmark}{s}{ei}{n}{\textsubring{d}}} \dicPos{f}[7] \dicFlx{(‑ar, ‑ir)}[1] \dicFieldCat{fyz.} \dicDirectTranslationCS{foton}
\dicEntry[ljósfræði] \dicTerm{ljós··fræð|i} \dicIPA{{l}{j}{ou}{s}{f}{r}{a}{i}{ð}{\textsci}} \dicPos{f}[3] \dicFlx{(‑i)}[3] \dicFieldCat{fyz.} \dicDirectTranslationCS{optika}
\dicEntry[ljósgeisli] \dicTerm{ljós··geisl|i} \dicIPA{{l}{j}{ou}{s}{\r{\textObardotlessj}}{ei}{s}{\textsubring{d}}{l}{\textsci}} \dicPos{m}[1] \dicFlx{(‑a, ‑ar)}[1] \dicDirectTranslationCS{světelný paprsek, paprsek světla}
\dicEntry[ljósgulur] \dicTerm{ljós··gulur} \dicIPA{{l}{j}{ou}{s}{\r{g}}{\textscy}{l}{\textscy}{\textsubring{r}}} \dicPos{adj}[1]\dicFlx{}[-1] \dicDirectTranslationCS{světle žlutý, nažloutlý}
\dicEntry[ljóshraði] \dicTerm{ljós··hrað|i} \dicIPA{{l}{j}{ou}{\textlengthmark}{s}{\textsubring{r}}{a}{ð}{\textsci}} \dicPos{m}[1] \dicFlx{(‑a)}[3] \dicFieldCat{fyz.} \dicDirectTranslationCS{rychlost světla}
\dicEntry[ljóshærður] \dicTerm{ljós··hærður} \dicsymFrequent\  \dicIPA{{l}{j}{ou}{\textlengthmark}{s}{h}{a}{i}{r}{ð}{\textscy}{\textsubring{r}}} \dicPos{adj}[2]\dicFlx{}[-1] \dicDirectTranslationCS{blonďatý, světlovlasý} \dicExampleIS{Hann er ljóshærður.} \dicExampleCS{Má světlé vlasy.}
\dicEntry[ljóshöfðaönd] \dicTerm{ljós·höfða··|önd} \dicIPA{{l}{j}{ou}{\textlengthmark}{s}{h}{\oe}{v}{ð}{a}{\oe}{n}{\textsubring{d}}} \dicPos{f}[8] \dicFlx{(‑andar, ‑endur)}[2] \dicFieldCat{zool.} \dicDirectTranslationCS{hvízdák americký} \textit{(l.~{\textLA{Anas americana}})}  \dicsymPhoto\ 
\dicFigure{ds_image_ljoshofdaond_0_2.jpg}{Ljóshöfðaönd}{Ljóshöfðaönd - Donna Dewhurst, USFWS, PD}
\dicEntry[ljóska] \dicTerm{ljósk|a} \dicIPA{{l}{j}{ou}{s}{\r{g}}{a}} \dicPos{f}[1] \dicFlx{(‑u, ‑ur)}[19] \dicDirectTranslationCS{blondýna, blondýnka}
\dicEntry[ljóskastari] \dicTerm{ljós··kast·ar|i} \dicIPA{{l}{j}{ou}{s}{k\smash{\textsuperscript{h}}}{a}{s}{\textsubring{d}}{a}{r}{\textsci}} \dicPos{m}[1] \dicFlx{(‑a, ‑ar)}[10] \dicDirectTranslationCS{světlomet, reflektor}
\dicEntry[ljósker] \dicTerm{ljós··ker} \dicIPA{{l}{j}{ou}{s}{c\smash{\textsuperscript{h}}}{\textepsilon}{\textsubring{r}}} \dicPos{n}[2] \dicFlx{(‑s, ‑)}[6] \dicDirectTranslationCS{lampa, svítilna, lucerna}
\dicEntry[ljóslifandi] \dicTerm{ljós··lif·andi} \dicIPA{{l}{j}{ou}{s}{l}{\textsci}{v}{a}{n}{\textsubring{d}}{\textsci}} \dicPos{adj}[13] \dicFlx{indecl}[1] \dicDirectTranslationCS{živý, barvitý} \dicExampleIS{e‑að stendur e‑m ljóslifandi fyrir hugskotssjónum} \dicExampleCS{(co) má (kdo) v~živé paměti}
\dicEntry[ljósmerki] \dicTerm{ljós··merki} \dicIPA{{l}{j}{ou}{s}{m}{\textepsilon}{\textsubring{r}}{\r{\textObardotlessj}}{\textsci}} \dicPos{n}[2] \dicFlx{(‑s, ‑)}[16] \dicDirectTranslationCS{světelný signál, světelné znamení}
\dicEntry[ljósmóðir] \dicTerm{ljós··|móðir} \dicIPA{{l}{j}{ou}{s}{m}{ou}{ð}{\textsci}{\textsubring{r}}} \dicPos{f}[11] \dicFlx{(‑móður, ‑mæður)}[1] \dicDirectTranslationCS{porodní asistentka}
\dicEntry[ljósmynd] \dicTerm{ljós··mynd} \dicsymFrequent\  \dicIPA{{l}{j}{ou}{s}{m}{\textsci}{n}{\textsubring{d}}} \dicPos{f}[7] \dicFlx{(‑ar, ‑ir)}[1] \dicDirectTranslationCS{fotografie, fotka} \dicExampleIS{svarthvít ljósmynd af ungri konu} \dicExampleCS{černobílá fotografie mladé ženy}
\dicEntry[ljósmynda] \dicTerm{ljós··mynd|a} \dicIPA{{l}{j}{ou}{s}{m}{\textsci}{n}{\textsubring{d}}{a}} \dicPos{v}[1] \dicFlx{(‑aði)}[1] \dicDirectTranslationCS{(vy)fotit, (vy)fotografovat, (u)dělat fotografii}
\dicEntry[ljósmyndari] \dicTerm{ljós··mynd·ar|i} \dicIPA{{l}{j}{ou}{s}{m}{\textsci}{n}{\textsubring{d}}{a}{r}{\textsci}} \dicPos{m}[1] \dicFlx{(‑a, ‑ar)}[13] \dicDirectTranslationCS{fotograf(ka)}
\dicEntry[ljósmyndavél] \dicTerm{ljós··mynda·vél} \dicIPA{{l}{j}{ou}{s}{m}{\textsci}{n}{\textsubring{d}}{a}{v}{j}{\textepsilon}{\textsubring{l}}} \dicPos{f}[4] \dicFlx{(‑ar, ‑ar)}[1] \dicDirectTranslationCS{fotoaparát}
\dicEntry[ljósmyndun] \dicTerm{ljós··mynd|un} \dicIPA{{l}{j}{ou}{s}{m}{\textsci}{n}{\textsubring{d}}{\textscy}{\textsubring{n}}} \dicPos{f}[7] \dicFlx{(‑unar, ‑anir)}[8] \textbf{1.} \dicDirectTranslationCS{fotografování}  \textbf{2.} \dicDirectTranslationCS{fotografie} \dicIndirectTranslationCS{(umělecký obor)}
\dicEntry[ljósmælir] \dicTerm{ljós··mæl|ir} \dicIPA{{l}{j}{ou}{s}{m}{a}{i}{l}{\textsci}{\textsubring{r}}} \dicPos{m}[7] \dicFlx{(‑is, ‑ar)}[1] \dicFieldCat{fyz.} \dicDirectTranslationCS{expozimetr}
\dicEntry[ljósnemi] \dicTerm{ljós··nem|i} \dicIPA{{l}{j}{ou}{s}{n}{\textepsilon}{m}{\textsci}} \dicPos{m}[1] \dicFlx{(‑a, ‑ar)}[1] \dicFieldCat{fyz.} \dicDirectTranslationCS{fotočlánek}
\dicEntry[ljósnæmur] \dicTerm{ljós··næmur} \dicIPA{{l}{j}{ou}{s}{n}{a}{i}{m}{\textscy}{\textsubring{r}}} \dicPos{adj}[1]\dicFlx{}[-1] \dicDirectTranslationCS{citlivý na světlo, fotosenzitivní}
\dicEntry[ljósrit] \dicTerm{ljós··rit} \dicIPA{{l}{j}{ou}{s}{r}{\textsci}{\textsubring{d}}} \dicPos{n}[2] \dicFlx{(‑s, ‑)}[5] \dicDirectTranslationCS{fotokopie}
\dicEntry[ljósrita] \dicTerm{ljós··rit|a} \dicIPA{{l}{j}{ou}{s}{r}{\textsci}{\textsubring{d}}{a}} \dicPos{v}[1] \dicFlx{(‑aði)}[1] \dicFlx{acc} \dicDirectTranslationCS{(o)kopírovat (pomocí fotokopírky)}
\dicEntry[ljósritunarvél] \dicTerm{ljós·ritunar··vél} \dicIPA{{l}{j}{ou}{s}{r}{\textsci}{\textsubring{d}}{\textscy}{n}{a}{r}{v}{j}{\textepsilon}{\textsubring{l}}} \dicPos{f}[4] \dicFlx{(‑ar, ‑ar)}[1] \dicDirectTranslationCS{(foto)kopírka}
\dicEntry[ljósta] \dicTerm{ljósta} \dicsymFrequent\  \dicIPA{{l}{j}{ou}{s}{\textsubring{d}}{a}} \dicPos{v}[6] \dicFlx{(lýst, laust, lustum, lysti, lostið)}[97] \dicFlx{acc} \dicLangCat{zast.} \dicDirectTranslationCS{udeřit, uhodit} \dicExampleIS{ljósta e‑n í andlitið með e‑u} \dicExampleCS{uhodit (koho) do tváře (čím)};  \dicIdiom{ljósta}[niður]{ \dicPhraseIS{e‑u lýstur niður}} \dicFlx{impers} \dicDirectTranslationCS{(co) udeří, (co) zasáhne (blesk ap.)} \dicExampleIS{Eldingu laust niður í húsið.} \dicExampleCS{Blesk udeřil do domu.};  \dicIdiom{ljósta}[saman]{ \dicPhraseIS{e‑m lýstur saman}} \dicFlx{impers} \dicDirectTranslationCS{(kdo) se střetne} \dicExampleIS{Herjunum laust saman.} \dicExampleCS{Armády se střetly.};  \dicIdiom{ljósta}[upp]{ \dicPhraseIS{ljósta upp ópi}} \dicDirectTranslationCS{vykřiknout}
\dicEntry[ljóstillífun] \dicTerm{ljós··til·líf|un} \dicIPA{{l}{j}{ou}{s}{t\smash{\textsuperscript{h}}}{\textsci}{l}{i}{v}{\textscy}{\textsubring{n}}} \dicPos{f}[7] \dicFlx{(‑unar)}[9] \dicFieldCat{biol.} \dicDirectTranslationCS{fotosyntéza}
\dicEntry[ljóstra] \dicTerm{ljóstr|a} \dicIPA{{l}{j}{ou}{s}{\textsubring{d}}{r}{a}} \dicPos{v}[1] \dicFlx{(‑aði)}[1] \dicFlx{dat} \dicPhraseIS{ljóstra e‑u upp} \dicDirectTranslationCS{odhalit\,/\addthin odhalovat (co), odkrýt\,/\addthin odkrývat (co)} \dicExampleIS{ljóstra upp leyndarmáli} \dicExampleCS{odkrýt tajemství}
\dicEntry[ljósvaki] \dicTerm{ljós··vak|i} \dicIPA{{l}{j}{ou}{\textlengthmark}{s}{v}{a}{\r{\textObardotlessj}}{\textsci}} \dicPos{m}[1] \dicFlx{(‑a, ‑ar)}[8] \dicDirectTranslationCS{éter} \dicIndirectTranslationCS{(podle starých představ látka zaplňující vesmírný prostor)};  \dicPhraseIS{á öldum ljósvakans} \dicFlx{adv} \dicDirectTranslationCS{na vlnách éteru}
\dicEntry[ljótleiki] \dicTerm{ljót··leik|i} \dicIPA{{l}{j}{ou}{\textlengthmark}{\textsubring{d}}{l}{ei}{\r{\textObardotlessj}}{\textsci}} \dicPos{m}[1] \dicFlx{(‑a)}[3] \dicSynonym{lýti} \dicDirectTranslationCS{šerednost, ohyzdnost (vzhled ap.)}
\dicEntry[ljótur] \dicTerm{ljótur} \dicsymFrequent\  \dicIPA{{l}{j}{ou}{\textlengthmark}{\textsubring{d}}{\textscy}{\textsubring{r}}} \dicPos{adj}[1]\dicFlx{}[-1] \dicDirectTranslationCS{šeredný, ošklivý, škaredý} \dicExampleIS{ljótar athugasemdir} \dicExampleCS{ošklivé poznámky}
\dicEntry[Ljubljana] \dicTerm{Ljubljana} \dicIPA{{l}{j}{\textscy}{\textsubring{b}}{l}{j}{a}{n}{a}} \dicPos{subs} \dicFlx{indecl} \dicFieldCat{geog.} \dicDirectTranslationCS{Lublaň} \dicIndirectTranslationCS{(hlavní město Slovinska)}
\dicEntry[ljúffengur] \dicTerm{ljúf··fengur} \dicIPA{{l}{j}{u}{f}{\textlengthmark}{ei}{\ng}{\r{g}}{\textscy}{\textsubring{r}}} \dicPos{adj}[1]\dicFlx{}[-1] \dicDirectTranslationCS{chutný, lahodný, delikátní} \dicIndirectTranslationCS{(o~jídle)}
\dicEntry[ljúflingur] \dicTerm{ljúf··ling|ur} \dicIPA{{l}{j}{u}{v}{l}{i}{\ng}{\r{g}}{\textscy}{\textsubring{r}}} \dicPos{m}[6] \dicFlx{(‑s, ‑ar)}[8] \textbf{1.} \dicSynonym{huldumaður} \dicDirectTranslationCS{elf}  \textbf{2.} \dicSynonym{eftirlæti} \dicDirectTranslationCS{miláček, zlato}
\dicEntry[ljúflyndi] \dicTerm{ljúf··lyndi} \dicIPA{{l}{j}{u}{v}{l}{\textsci}{n}{\textsubring{d}}{\textsci}} \dicPos{n}[2] \dicFlx{(‑s)}[20] \dicDirectTranslationCS{přívětivost, laskavost}
\dicEntry[ljúfmenni] \dicTerm{ljúf··menni} \dicIPA{{l}{j}{u}{v}{m}{\textepsilon}{n}{\textsci}} \dicPos{n}[2] \dicFlx{(‑s, ‑)}[14] \dicSynonym{góðmenni} \dicDirectTranslationCS{milý\,/\addthin přívětivý člověk}
\dicEntry[ljúfmennska] \dicTerm{ljúf··mennsk|a} \dicIPA{{l}{j}{u}{v}{m}{\textepsilon}{n}{s}{\r{g}}{a}} \dicPos{f}[1] \dicFlx{(‑u)}[5] \dicDirectTranslationCS{dobrota, přívětivost, vlídnost}
\dicEntry[ljúfur] \dicTerm{ljúf|ur\smash{\textsuperscript{1}}} \dicIPA{{l}{j}{u}{\textlengthmark}{v}{\textscy}{\textsubring{r}}} \dicPos{m}[6] \dicFlx{(‑s, ‑ar)}[24] \dicDirectTranslationCS{poklad, drahoušek, miláček} \dicExampleIS{ljúfur(inn) (minn)} \dicExampleCS{ty můj poklade}
\dicEntry[ljúfur] \dicTerm{ljúfur\smash{\textsuperscript{2}}} \dicsymFrequent\  \dicIPA{{l}{j}{u}{\textlengthmark}{v}{\textscy}{\textsubring{r}}} \dicPos{adj}[1]\dicFlx{}[-1] \textbf{1.} \dicSynonym{blíður} \dicDirectTranslationCS{laskavý, přívětivý, vlídný}  \textbf{2.} \dicSynonym{þægilegur} \dicDirectTranslationCS{příjemný, milý} \dicExampleIS{ljúfar minningar} \dicExampleCS{příjemné vzpomínky};  \dicPhraseIS{e‑m er ljúft að (gera e‑ð)} \dicFlx{impers} \dicDirectTranslationCS{(komu) je příjemné (udělat (co))}
\dicEntry[ljúga] \dicTerm{ljúga} \dicsymFrequent\  \dicIPA{{l}{j}{u}{\textlengthmark}{a}} \dicPos{v}[6] \dicFlx{(lýg, laug, lugum, lygi, logið)}[52] \dicFlx{dat} \dicDirectTranslationCS{lhát, zalhat} \dicExampleIS{ljúga e‑u} \dicExampleCS{zalhat o~(čem)};  \dicPhraseIS{ljúga að e‑m} \dicDirectTranslationCS{lhát (komu)};  \dicPhraseIS{ljúga til um e‑ð} \dicDirectTranslationCS{lhát o~(čem)} \dicExampleIS{ljúga til um aldur} \dicExampleCS{lhát o~(svém) věku}
\dicEntry[ljúgvitni] \dicTerm{ljúg··vitni} \dicIPA{{l}{j}{u}{\textlengthmark}{v}{\textsci}{h}{\textsubring{d}}{n}{\textsci}} \dicPos{n}[2] \dicFlx{(‑s, ‑)}[14] \dicDirectTranslationCS{křivá přísaha, křivé doznání};  \dicPhraseIS{bera ljúgvitni gegn e‑m} \dicDirectTranslationCS{křivě svědčit proti (komu)}
\dicEntry[ljúka] \dicTerm{ljúka} \dicsymFrequent\  \dicIPA{{l}{j}{u}{\textlengthmark}{\r{g}}{a}} \dicPos{v}[6] \dicFlx{(lýk, lauk, lukum, lyki, lokið)}[47] \dicFlx{dat} \textbf{1.} \dicSynonym{klára} \dicDirectTranslationCS{(u)končit, ukončovat, zakončit, zakončovat, dokončit, dokončovat};  \dicPhraseIS{ljúka e‑u} \dicDirectTranslationCS{ukončit (co)};  \dicPhraseIS{ljúka við e‑ð} \dicDirectTranslationCS{ukončit (co), (s)končit s~(čím)}  \textbf{2.} \dicPhraseIS{e‑u lýkur} \dicFlx{impers} \dicDirectTranslationCS{(co) končí} \dicExampleIS{Skólanum lýkur í júní.} \dicExampleCS{Škola končí v~červnu.};  \dicIdiom{ljúka}[af]{ \dicPhraseIS{ljúka e‑u af}} \dicDirectTranslationCS{ukončit (co), skončit (co)};  \dicIdiom{ljúka}[upp]{ \dicPhraseIS{ljúka upp e‑u}} \dicLangCat{zast.} \dicSynonym{opna\smash{\textsuperscript{2}}} \dicDirectTranslationCS{otevřít (co)} \dicExampleIS{ljúka dyrunum upp} \dicExampleCS{otevřít dveře};  \dicIdiom{ljúka}[yfir]{ \dicPhraseIS{þar til yfir lýkur}} \dicFlx{adv} \dicDirectTranslationCS{až do samého konce}; { \dicPhraseIS{áður en yfir lýkur}} \dicFlx{adv} \dicDirectTranslationCS{dříve než bude po všem};  \dicIdiom{ljúkast}[saman]{ \dicPhraseIS{ljúkast saman}} \dicFlx{refl} \dicLangCat{zast.} \dicDirectTranslationCS{zavřít se};  \dicIdiom{ljúkast}[upp]{ \dicPhraseIS{ljúkast upp}} \dicFlx{refl} \dicLangCat{zast.} \dicDirectTranslationCS{otevřít se, rozevřít se}; { \dicPhraseIS{e‑að lýkst upp fyrir e‑m}} \dicFlx{refl} \dicLangCat{přen.} \dicDirectTranslationCS{(co) se otevírá před (kým) (nová myšlenka ap.)}
\dicEntry[ljæ] \dicTerm{ljæ} \dicIPA{{l}{j}{a}{i}{\textlengthmark}} \dicPos{v} \dicFlx{ind praes sg 1 pers} \dicLink{ljá}
\dicEntry[lo.] \dicTerm{lo.} \dicPos{zkr} \dicPhraseIS{lýsingarorð} \dicFieldCat{jaz.} \dicDirectTranslationCS{přídavné jméno}
\dicEntry[lobbíisti] \dicTerm{lobbí··ist|i} \dicIPA{{l}{\textopeno}{\textsubring{b}}{\textlengthmark}{i}{j}{\textsci}{s}{\textsubring{d}}{\textsci}} \dicPos{m}[1] \dicFlx{(‑a, ‑ar)}[1] \dicDirectTranslationCS{lobbista, lobbistka}
\dicEntry[loddari] \dicTerm{lodd··ar|i} \dicIPA{{l}{\textopeno}{\textsubring{d}}{\textlengthmark}{a}{r}{\textsci}} \dicPos{m}[1] \dicFlx{(‑a, ‑ar)}[13] \textbf{1.} \dicSynonym{svikari} \dicDirectTranslationCS{šarlatán(ka)}  \textbf{2.} \dicSynonym{trúður} \dicDirectTranslationCS{šašek, klaun(ka)}
\dicEntry[loða] \dicTerm{lo|ða} \dicIPA{{l}{\textopeno}{\textlengthmark}{ð}{a}} \dicPos{v}[2] \dicFlx{(‑ddi, ‑ðað)}[182] \dicSynonym{tolla\smash{\textsuperscript{2}}} \dicDirectTranslationCS{přiléhat, lnout, přimknout se, přimykat};  \dicIdiom{loða}[saman]{ \dicPhraseIS{loða saman}} \dicDirectTranslationCS{přilnout k~sobě};  \dicIdiom{loða}[við]{ \dicPhraseIS{loða við e‑ð}} \dicDirectTranslationCS{přilnout k~(čemu)} \dicExampleIS{Málningin loðir við vegginn.} \dicExampleCS{Nátěr přilne ke zdi.}
\dicEntry[loðdýr] \dicTerm{loð··dýr} \dicIPA{{l}{\textopeno}{ð}{\textsubring{d}}{i}{\textsubring{r}}} \dicPos{n}[2] \dicFlx{(‑s, ‑)}[5] \dicDirectTranslationCS{kožešinové zvíře}
\dicEntry[loðfeldur] \dicTerm{loð··feld|ur} \dicIPA{{l}{\textopeno}{ð}{f}{\textepsilon}{l}{\textsubring{d}}{\textscy}{\textsubring{r}}} \dicPos{m}[10] \dicFlx{(‑ar, ‑ir)}[4] \dicSynonym{loðskinn} \dicDirectTranslationCS{kožešina}
\dicEntry[loðfíll] \dicTerm{loð··fíl|l} \dicIPA{{l}{\textopeno}{ð}{f}{i}{\textsubring{d}}{\textsubring{l}}} \dicPos{m}[6] \dicFlx{(‑s, ‑ar)}[48] \dicFieldCat{zool.} \dicDirectTranslationCS{mamut} \textit{(l.~{\textLA{Mammuthus primigenius}})}
\dicEntry[loðhúfa] \dicTerm{loð··húf|a} \dicIPA{{l}{\textopeno}{\textlengthmark}{\texttheta}{h}{u}{v}{a}} \dicPos{f}[1] \dicFlx{(‑u, ‑ur)}[19] \dicDirectTranslationCS{kožešinová čepice, beranice}
\dicEntry[loðinn] \dicTerm{loðinn} \dicIPA{{l}{\textopeno}{\textlengthmark}{ð}{\textsci}{\textsubring{n}}} \dicPos{adj}[6]\dicFlx{}[-2] \textbf{1.} \dicDirectTranslationCS{chlupatý, srstnatý, huňatý}  \textbf{2.} \dicSynonym{óhreinn} \dicDirectTranslationCS{vyhýbavý, mlhavý (odpověď ap.)}
\dicEntry[loðkápa] \dicTerm{loð··káp|a} \dicIPA{{l}{\textopeno}{ð}{k\smash{\textsuperscript{h}}}{au}{\textsubring{b}}{a}} \dicPos{f}[1] \dicFlx{(‑u, ‑ur)}[19] \dicDirectTranslationCS{kožich}
\dicEntry[loðna] \dicTerm{loðn|a} \dicIPA{{l}{\textopeno}{ð}{n}{a}} \dicPos{f}[1] \dicFlx{(‑u, ‑ur)}[7] \textbf{1.} \dicDirectTranslationCS{ochlupení}  \textbf{2.} \dicFieldCat{zool.} \dicDirectTranslationCS{huňáček severní} \textit{(l.~{\textLA{Mallotus villosus}})}
\dicEntry[loðskinn] \dicTerm{loð··skinn} \dicIPA{{l}{\textopeno}{ð}{s}{\r{\textObardotlessj}}{\textsci}{\textsubring{n}}} \dicPos{n}[2] \dicFlx{(‑s, ‑)}[5] \dicDirectTranslationCS{kožešina, kožka}
\dicEntry[loðvíðir] \dicTerm{loð··víð|ir} \dicIPA{{l}{\textopeno}{ð}{v}{i}{ð}{\textsci}{\textsubring{r}}} \dicPos{m}[7] \dicFlx{(‑is)}[2] \dicFieldCat{bot.} \dicDirectTranslationCS{vrba vlnatá} \textit{(l.~{\textLA{Salix lanata}})}  \dicsymPhoto\ 
\dicFigure{66869.jpg}{Loðvíðir}{Loðvíðir - Hroneš Michal, Biolib, Copyright/CC-BY-NC}
\dicEntry[lof] \dicTerm{lof} \dicIPA{{l}{\textopeno}{\textlengthmark}{f}} \dicPos{n}[2] \dicFlx{(‑s)}[2] \dicSynonym{hrós} \dicDirectTranslationCS{chvála, pochvala};  \dicPhraseIS{bera lof á e‑n} \dicDirectTranslationCS{chválit (koho)};  \dicPhraseIS{guði sé lof} \dicDirectTranslationCS{díky bohu};  \dicPhraseIS{klappa lof í lófa} \dicDirectTranslationCS{tleskat, aplaudovat}
\dicEntry[lofa] \dicTerm{lof|a} \dicsymFrequent\  \dicIPA{{l}{\textopeno}{\textlengthmark}{v}{a}} \dicPos{v}[1] \dicFlx{(‑aði)}[1] \dicFlx{acc\,/\addthin dat} \textbf{1.} \dicFlx{acc} \dicSynonym{hrósa} \dicDirectTranslationCS{(po)chválit, vychvalovat};  \dicPhraseIS{lofa e‑n\,/\addthin e‑ð} \dicDirectTranslationCS{chválit (koho\,/\addthin co)} \dicExampleIS{lofa listaverkið} \dicExampleCS{pochválit umělecké dílo}  \textbf{2.} \dicFlx{dat + dat} \dicSynonym{heita} \dicDirectTranslationCS{(při)slíbit, slibovat};  \dicPhraseIS{lofa e‑m e‑u} \dicDirectTranslationCS{slíbit (komu co)}  \textbf{3.} \mbox{\dicFlx{dat + acc}} \dicSynonym{leyfa} \dicDirectTranslationCS{dovolit, dovolovat, povolit, povolovat};  \dicPhraseIS{lofa e‑m e‑ð} \dicDirectTranslationCS{dovolit (komu co)} \dicExampleIS{Ég lofaði henni að dansa.} \dicExampleCS{Dovolil jsem jí tančit.};  \dicIdiom{lofast}{ \dicPhraseIS{lofast e‑m}} \dicFlx{refl} \dicLangCat{zast.} \dicDirectTranslationCS{zaslíbit se (komu)}
\dicEntry[lofgerð] \dicTerm{lof··gerð} \dicIPA{{l}{\textopeno}{v}{\r{\textObardotlessj}}{\textepsilon}{r}{\texttheta}} \dicPos{f}[7] \dicFlx{(‑ar, ‑ir)}[1] \dicSynonym{lof} \dicDirectTranslationCS{chvála, pochvala}
\dicEntry[loforð] \dicTerm{lof··orð} \dicIPA{{l}{\textopeno}{\textlengthmark}{v}{\textopeno}{r}{\texttheta}} \dicPos{n}[2] \dicFlx{(‑s, ‑)}[5] \textbf{1.} \dicDirectTranslationCS{slib, slovo};  \dicPhraseIS{standa við loforð} \dicDirectTranslationCS{držet slovo}  \textbf{2.} \dicDirectTranslationCS{přísaha}
\dicEntry[lofsamlegur] \dicTerm{lof·sam··legur} \dicIPA{{l}{\textopeno}{f}{s}{a}{m}{l}{\textepsilon}{\textbabygamma}{\textscy}{\textsubring{r}}} \dicPos{adj}[1]\dicFlx{}[-8] \dicDirectTranslationCS{pochvalný}
\dicEntry[lofsorð] \dicTerm{lofs··orð} \dicIPA{{l}{\textopeno}{f}{s}{\textopeno}{r}{\texttheta}} \dicPos{n}[2] \dicFlx{(‑s, ‑)}[5] \dicPhraseIS{ljúka lofsorði á e‑ð} \dicDirectTranslationCS{chválit (co), vyjádřit se o~(čem) pochvalně}
\dicEntry[lofsverður] \dicTerm{lofs··verður} \dicIPA{{l}{\textopeno}{f}{s}{v}{\textepsilon}{r}{ð}{\textscy}{\textsubring{r}}} \dicPos{adj}[2]\dicFlx{}[-1] \dicDirectTranslationCS{chvályhodný, záslužný} \dicExampleIS{lofsverð viðleitni} \dicExampleCS{chvályhodné úsilí}
\dicEntry[lofsöngur] \dicTerm{lof··söng|ur} \dicIPA{{l}{\textopeno}{f}{s}{\oe i}{\ng}{\r{g}}{\textscy}{\textsubring{r}}} \dicPos{m}[8] \dicFlx{(‑s, ‑var)}[2] \dicDirectTranslationCS{chvalozpěv}
\dicEntry[loft] \dicTerm{loft} \dicsymFrequent\  \dicIPA{{l}{\textopeno}{f}{\textsubring{d}}} \dicPos{n}[2] \dicFlx{(‑s, ‑)}[5] \textbf{1.} \dicSynonym{andrúmsloft} \dicDirectTranslationCS{vzduch, povětří, ovzduší} \dicExampleIS{ferskt loft} \dicExampleCS{čerstvý vzduch}  \textbf{2.} \dicSynonym{himinn} \dicDirectTranslationCS{vzduch, obloha} \dicExampleIS{kasta e‑u upp í loftið} \dicExampleCS{vyhodit (co) do vzduchu}  \textbf{3.} \dicDirectTranslationCS{poschodí, podlaží};  \dicPhraseIS{uppi á lofti} \dicFlx{adv} \dicDirectTranslationCS{v~horním patře}  \textbf{4.} \dicSynonym{háaloft} \dicDirectTranslationCS{podkroví};  \dicIdiom{loft}{ \dicPhraseIS{það er allt upp í loft}} \dicLangCat{přen.} \dicDirectTranslationCS{všechno je vzhůru nohama};  \dicPhraseIS{e‑að er úr lausu lofti gripið} \dicLangCat{přen.} \dicDirectTranslationCS{(co) je neopodstatněné}
\dicEntry[loftárás] \dicTerm{loft··á·rás} \dicIPA{{l}{\textopeno}{f}{\textsubring{d}}{au}{r}{au}{s}} \dicPos{f}[7] \dicFlx{(‑ar, ‑ir)}[1] \dicDirectTranslationCS{nálet, letecký útok}
\dicEntry[loftbelgur] \dicTerm{loft··belg|ur} \dicIPA{{l}{\textopeno}{f}{\textsubring{d}}{\textsubring{b}}{\textepsilon}{l}{\r{g}}{\textscy}{\textsubring{r}}} \dicPos{m}[9] \dicFlx{(‑s\,/\addthin ‑jar, ‑ir)}[26] \dicFieldCat{let.} \dicDirectTranslationCS{(horkovzdušný) balón}
\dicEntry[loftbor] \dicTerm{loft··bor} \dicIPA{{l}{\textopeno}{f}{\textsubring{d}}{\textsubring{b}}{\textopeno}{\textsubring{r}}} \dicPos{m}[4] \dicFlx{(‑s, ‑ar)}[14] \dicDirectTranslationCS{sbíječka, sbíjecí kladivo}
\dicEntry[loftbréf] \dicTerm{loft··bréf} \dicIPA{{l}{\textopeno}{f}{\textsubring{d}}{\textsubring{b}}{r}{j}{\textepsilon}{f}} \dicPos{n}[2] \dicFlx{(‑s, ‑)}[5] \dicDirectTranslationCS{letecké psaní}
\dicEntry[loftbrú] \dicTerm{loft··|brú} \dicIPA{{l}{\textopeno}{f}{\textsubring{d}}{\textsubring{b}}{r}{u}} \dicPos{f}[9] \dicFlx{(‑brúar, ‑brýr)}[4] \dicDirectTranslationCS{letecký most}
\dicEntry[loftfar] \dicTerm{loft··|far} \dicIPA{{l}{\textopeno}{f}{\textsubring{d}}{f}{a}{\textsubring{r}}} \dicPos{n}[2] \dicFlx{(‑fars, ‑för)}[8] \dicDirectTranslationCS{létající prostředek (letadlo, balón, vzducholoď ap.)}
\dicEntry[loftfimleikar] \dicTerm{loft··fim·leikar} \dicIPA{{l}{\textopeno}{f}{\textsubring{d}}{f}{\textsci}{m}{l}{ei}{\r{g}}{a}{\textsubring{r}}} \dicPos{m}[6] \dicFlx{pl}[2] \dicDirectTranslationCS{letecká akrobacie}
\dicEntry[lofthræddur] \dicTerm{loft··hræddur} \dicIPA{{l}{\textopeno}{f}{\textsubring{d}}{\textsubring{r}}{a}{i}{\textsubring{d}}{\textscy}{\textsubring{r}}} \dicPos{adj}[2]\dicFlx{}[-18] \dicDirectTranslationCS{mající strach z~výšek}
\dicEntry[lofthræðsla] \dicTerm{loft··hræðsl|a} \dicIPA{{l}{\textopeno}{f}{\textsubring{d}}{\textsubring{r}}{a}{i}{ð}{s}{\textsubring{d}}{l}{a}} \dicPos{f}[1] \dicFlx{(‑u)}[5] \dicDirectTranslationCS{strach z~výšek}
\dicEntry[loftkastali] \dicTerm{loft··kastal|i} \dicIPA{{l}{\textopeno}{f}{\textsubring{d}}{k\smash{\textsuperscript{h}}}{a}{s}{\textsubring{d}}{a}{l}{\textsci}} \dicPos{m}[1] \dicFlx{(‑a, ‑ar)}[12] \dicDirectTranslationCS{vzdušný zámek};  \dicPhraseIS{byggja loftkastala} \dicLangCat{přen.} \dicDirectTranslationCS{stavět vzdušné zámky}
\dicEntry[loftkenndur] \dicTerm{loft··kenndur} \dicIPA{{l}{\textopeno}{f}{\textsubring{d}}{c\smash{\textsuperscript{h}}}{\textepsilon}{n}{\textsubring{d}}{\textscy}{\textsubring{r}}} \dicPos{adj}[2]\dicFlx{}[-14] \textbf{1.} \dicDirectTranslationCS{vzdušný}  \textbf{2.} \dicDirectTranslationCS{plynný} \dicExampleIS{loftkennd gosefni} \dicExampleCS{plynné magma}
\dicEntry[loftkældur] \dicTerm{loft··kældur} \dicIPA{{l}{\textopeno}{f}{\textsubring{d}}{c\smash{\textsuperscript{h}}}{a}{i}{l}{\textsubring{d}}{\textscy}{\textsubring{r}}} \dicPos{adj}[2]\dicFlx{}[-17] \dicDirectTranslationCS{klimatizovaný} \dicExampleIS{loftkæld skrifstofa} \dicExampleCS{klimatizovaná kancelář}
\dicEntry[loftkæling] \dicTerm{loft··kæl·ing} \dicIPA{{l}{\textopeno}{f}{\textsubring{d}}{c\smash{\textsuperscript{h}}}{a}{i}{l}{i}{\ng}{\r{g}}} \dicPos{f}[4] \dicFlx{(‑ar, ‑ar)}[5] \dicDirectTranslationCS{klimatizace}
\begin{xtolerant}{}{1pt}
\dicEntry[loftmengun] \dicTerm{loft··meng|un}\addthinS\dicIPA{{l}{\textopeno}{f}{\textsubring{d}}{m}{ei}{\ng}{\r{g}}{\textscy}{\textsubring{n}}}\addthinS\dicPos{f}[7]\addthinS\dicFlx{(‑unar)}[9] \dicDirectTranslationCS{znečištění ovzduší}
\end{xtolerant}
\dicEntry[loftnet] \dicTerm{loft··net} \dicIPA{{l}{\textopeno}{f}{\textsubring{d}}{n}{\textepsilon}{\textsubring{d}}} \dicPos{n}[2] \dicFlx{(‑s, ‑)}[5] \dicDirectTranslationCS{anténa}
\dicEntry[loftræsta] \dicTerm{loft··ræst|a} \dicIPA{{l}{\textopeno}{f}{\textsubring{d}}{r}{a}{i}{s}{\textsubring{d}}{a}} \dicPos{v}[2] \dicFlx{(‑i, ‑)}[16] \dicFlx{acc} \dicDirectTranslationCS{(vy)větrat, (vy)ventilovat} \dicExampleIS{loftræsta sal} \dicExampleCS{vyvětrat sál}
\dicEntry[loftræsting] \dicTerm{loft··ræst·ing} \dicIPA{{l}{\textopeno}{f}{\textsubring{d}}{r}{a}{i}{s}{\textsubring{d}}{i}{\ng}{\r{g}}} \dicPos{f}[4] \dicFlx{(‑ar, ‑ar)}[5] \dicDirectTranslationCS{větrání, ventilace}
\begin{xtolerant}{}{1pt}
\dicEntry[loftskeytamaður] \dicTerm{loft·skeyta··|maður}\addthinSS\dicIPA{{l}{\textopeno}{f}{\textsubring{d}}{s}{c\smash{\textsuperscript{h}}}{ei}{\textsubring{d}}{a}{m}{a}{ð}{\textscy}{\textsubring{r}}}\addthinSS\dicPos{m}[1\hspace{-0.10667em}3]\addthinS\dicFlx{(‑manns,\addthinS‑menn)}[2] \dicDirectTranslationCS{radista, radistka, rádiový operátor, rádiová operátorka}
\end{xtolerant}
\dicEntry[loftskip] \dicTerm{loft··skip} \dicIPA{{l}{\textopeno}{f}{\textsubring{d}}{s}{\r{\textObardotlessj}}{\textsci}{\textsubring{b}}} \dicPos{n}[2] \dicFlx{(‑s, ‑)}[5] \dicFieldCat{let.} \dicDirectTranslationCS{vzducholoď}
\dicEntry[loftslag] \dicTerm{lofts··lag} \dicIPA{{l}{\textopeno}{f}{s}{l}{a}{x}} \dicPos{n}[2] \dicFlx{(‑s)}[2] \dicSynonym{veðurfar} \dicDirectTranslationCS{podnebí, klima}
\dicEntry[loftslagsbelti] \dicTerm{lofts·lags··belti} \dicIPA{{l}{\textopeno}{f}{s}{l}{a}{x}{s}{\textsubring{b}}{\textepsilon}{\textsubring{l}}{\textsubring{d}}{\textsci}} \dicPos{n}[2] \dicFlx{(‑s, ‑)}[14] \dicFieldCat{meteo.} \dicDirectTranslationCS{podnebný pás}
\dicEntry[loftslagsfræði] \dicTerm{lofts·lags··fræð|i} \dicIPA{{l}{\textopeno}{f}{s}{l}{a}{x}{s}{f}{r}{a}{i}{ð}{\textsci}} \dicPos{f}[3] \dicFlx{(‑i)}[3] \dicSynonym{veðurfarsfræði} \dicDirectTranslationCS{klimatologie}
\dicEntry[loftsteinn] \dicTerm{loft··stein|n} \dicIPA{{l}{\textopeno}{f}{\textsubring{d}}{s}{\textsubring{d}}{ei}{\textsubring{d}}{\textsubring{n}}} \dicPos{m}[6] \dicFlx{(‑s, ‑ar)}[42] \dicFieldCat{astro.} \dicDirectTranslationCS{meteorit}
\dicEntry[lofttegund] \dicTerm{loft··tegund} \dicIPA{{l}{\textopeno}{f}{t\smash{\textsuperscript{h}}}{\textepsilon}{\textbabygamma}{\textscy}{n}{\textsubring{d}}} \dicPos{f}[7] \dicFlx{(‑ar, ‑ir)}[1] \dicFieldCat{chem.} \dicDirectTranslationCS{plyn}
\dicEntry[lofttóm] \textls[15]{\dicTerm{loft··tóm} \dicIPA{{l}{\textopeno}{f}{t\smash{\textsuperscript{h}}}{ou}{\textsubring{m}}} \dicPos{n}[2] \dicFlx{(‑s)}[2] \dicFieldCat{fyz.} \dicDirectTranslationCS{vakuum, vzduchoprázdno}}
\dicEntry[loftvarnir] \dicTerm{loft··varnir} \dicIPA{{l}{\textopeno}{f}{\textsubring{d}}{v}{a}{r}{\textsubring{d}}{n}{\textsci}{\textsubring{r}}} \dicPos{f}[7] \dicFlx{pl}[18] \dicFieldCat{voj.} \dicDirectTranslationCS{protiletecká\,/\addthin protiletadlová obrana}
\dicEntry[loftveiki] \dicTerm{loft··veik|i} \dicIPA{{l}{\textopeno}{f}{\textsubring{d}}{v}{ei}{\r{\textObardotlessj}}{\textsci}} \dicPos{f}[3] \dicFlx{(‑i)}[3] \dicFieldCat{med.} \dicDirectTranslationCS{výšková\,/\addthin horská nemoc}
\dicEntry[loftvog] \dicTerm{loft··vog} \dicIPA{{l}{\textopeno}{f}{\textsubring{d}}{v}{\textopeno}{x}} \dicPos{f}[7] \dicFlx{(‑ar, ‑ir)}[1] \dicDirectTranslationCS{barometr}
\dicEntry[loftþéttur] \dicTerm{loft··þéttur} \dicIPA{{l}{\textopeno}{f}{\textsubring{d}}{\texttheta}{j}{\textepsilon}{h}{\textsubring{d}}{\textscy}{\textsubring{r}}} \dicPos{adj}[1]\dicFlx{}[-10] \dicDirectTranslationCS{vzduchotěsný, hermetický}
\dicEntry[loftþrýstingur] \dicTerm{loft··þrýst·ing|ur} \dicIPA{{l}{\textopeno}{f}{\textsubring{d}}{\texttheta}{r}{i}{s}{\textsubring{d}}{i}{\ng}{\r{g}}{\textscy}{\textsubring{r}}} \dicPos{m}[6] \dicFlx{(‑s)}[9] \dicDirectTranslationCS{tlak vzduchu, atmosférický tlak}
\dicEntry[loga] \dicTerm{log|a} \dicsymFrequent\  \dicIPA{{l}{\textopeno}{\textlengthmark}{\textbabygamma}{a}} \dicPos{v}[1] \dicFlx{(‑aði)}[44] \textbf{1.} \dicDirectTranslationCS{hořet, planout};  \dicPhraseIS{það logar á e‑u} \dicFlx{impers} \dicDirectTranslationCS{(co) hoří} \dicExampleIS{Það logar vel á kertinu.} \dicExampleCS{Svíčka dobře hoří.}  \textbf{2.} \dicDirectTranslationCS{hořet, zuřit (město ap.)} \dicIndirectTranslationCS{(o~nepokojích nebo válce)};  \dicPhraseIS{e‑að logar í e‑u} \dicDirectTranslationCS{v~(čem) zuří (co) (ve městě nepokoje ap.)}
\dicEntry[logandi] \dicTerm{log··andi\smash{\textsuperscript{1}}} \dicIPA{{l}{\textopeno}{\textlengthmark}{\textbabygamma}{a}{n}{\textsubring{d}}{\textsci}} \dicPos{adj}[13] \dicFlx{indecl}[1] \dicDirectTranslationCS{hořící, planoucí} \dicExampleIS{logandi kerti} \dicExampleCS{hořící svíce}
\dicEntry[logandi] \dicTerm{log··andi\smash{\textsuperscript{2}}} \dicIPA{{l}{\textopeno}{\textlengthmark}{\textbabygamma}{a}{n}{\textsubring{d}}{\textsci}} \dicPos{adv} \dicDirectTranslationCS{hrozně, strašně} \dicIndirectTranslationCS{(k~zdůraznění)} \dicExampleIS{Hann var logandi hræddur.} \dicExampleCS{Strašně se bál.};  \dicPhraseIS{hvert þó í logandi!} \dicDirectTranslationCS{panebože!, svatá dobroto!}
\dicEntry[logga] \dicTerm{logg|a} \dicIPA{{l}{\textopeno}{\r{g}}{\textlengthmark}{a}} \dicPos{v}[1] \dicFlx{(‑aði)}[1] \dicFieldCat{poč.} \dicDirectTranslationCS{logovat, zaznamenat do souboru};  \dicIdiom{logga}[inn]{ \dicPhraseIS{logga sig inn}} \dicFieldCat{poč.} \dicDirectTranslationCS{(za)logovat se, nalogovat se, přihlásit se};  \dicIdiom{logga}[út]{ \dicPhraseIS{logga sig út}} \dicFieldCat{poč.} \dicDirectTranslationCS{vylogovat se, odlogovat se, odhlásit se}
\dicEntry[logi] \dicTerm{log|i} \dicsymFrequent\  \dicIPA{{l}{\textopeno i}{j}{\textlengthmark}{\textsci}} \dicPos{m}[1] \dicFlx{(‑a, ‑ar)}[1] \dicDirectTranslationCS{plamen, planutí} \dicExampleIS{Loginn hefur dofnað.} \dicExampleCS{Plamen zeslábl.};  \dicPhraseIS{fara eins og logi yfir akur} \dicLangCat{přen.} \dicDirectTranslationCS{přehnat se jako bouře}
\dicEntry[logið] \dicTerm{logið} \dicIPA{{l}{\textopeno i}{j}{\textlengthmark}{\textsci}{\texttheta}} \dicPos{v} \dicFlx{supin} \dicLink{ljúga}
\dicEntry[logn] \dicTerm{logn} \dicsymFrequent\  \dicIPA{{l}{\textopeno}{\r{g}}{\textsubring{n}}} \dicPos{n}[2] \dicFlx{(‑s, ‑)}[5] \dicSynonym*{vindleysa} \dicDirectTranslationCS{bezvětří} \dicIndirectTranslationCS{(Beaufort 0) \mbox{(0--0,2~m/s)}} \dicExampleIS{Það er logn í dag.} \dicExampleCS{Dnes je bezvětří.}
\dicEntry[lognmolla] \dicTerm{logn··moll|a} \dicIPA{{l}{\textopeno}{\r{g}}{n}{m}{\textopeno}{\textsubring{d}}{l}{a}} \dicPos{f}[1] \dicFlx{(‑u)}[5] \textbf{1.} \dicSynonym{andleysi} \dicDirectTranslationCS{nečinnost, neaktivita}  \textbf{2.} \dicDirectTranslationCS{poklidné počasí}
\dicEntry[logsauð] \dicTerm{log··sauð} \dicIPA{{l}{\textopeno}{x}{s}{\oe i}{\texttheta}} \dicPos{v} \dicFlx{ind pf sg 1 pers} \dicLink{logsjóða}
\dicEntry[logsjóða] \dicTerm{log··|sjóða} \dicIPA{{l}{\textopeno}{x}{s}{j}{ou}{ð}{a}} \dicPos{v}[6] \dicFlx{(‑sýð, ‑sauð, ‑suðum, ‑syði, ‑soðið)}[104] \dicFlx{acc} \dicFieldCat{techn.} \dicDirectTranslationCS{svářet, svařovat}
\dicEntry[logsoðið] \dicTerm{log··soðið} \dicIPA{{l}{\textopeno}{x}{s}{\textopeno}{ð}{\textsci}{\texttheta}} \dicPos{v} \dicFlx{supin} \dicLink{logsjóða}
\dicEntry[logsuða] \dicTerm{log··suð|a} \dicIPA{{l}{\textopeno}{x}{s}{\textscy}{ð}{a}} \dicPos{f}[1] \dicFlx{(‑u)}[5] \dicSynonym*{það að logsjóða} \dicDirectTranslationCS{sváření, svařování}
\dicEntry[logsuðum] \dicTerm{log··suðum} \dicIPA{{l}{\textopeno}{x}{s}{\textscy}{ð}{\textscy}{\textsubring{m}}} \dicPos{v} \dicFlx{ind pf pl 1 pers} \dicLink{logsjóða}
\dicEntry[logsyði] \dicTerm{log··syði} \dicIPA{{l}{\textopeno}{x}{s}{\textsci}{ð}{\textsci}} \dicPos{v} \dicFlx{con pf sg 1 pers} \dicLink{logsjóða}
\dicEntry[logsýð] \dicTerm{log··sýð} \dicIPA{{l}{\textopeno}{x}{s}{i}{\texttheta}} \dicPos{v} \dicFlx{ind praes sg 1 pers} \dicLink{logsjóða}
\dicEntry[lok] \dicTerm{lok} \dicsymFrequent\  \dicIPA{{l}{\textopeno}{\textlengthmark}{\r{g}}} \dicPos{n}[2] \dicFlx{(‑s, ‑)}[5] \textbf{1.} \dicDirectTranslationCS{uzávěr, víčko, víko, poklička, poklice} \dicExampleIS{lokið á dósinni} \dicExampleCS{víčko plechovky}  \textbf{2.} \dicPhraseIS{lok} \dicFlx{pl} \dicSynonym{endir} \dicDirectTranslationCS{konec, závěr};  \dicPhraseIS{að lokum} \dicFlx{adv} \dicDirectTranslationCS{nakonec, na závěr};  \dicPhraseIS{e‑að líður undir lok} \dicDirectTranslationCS{(co) se chýlí k~závěru};  \dicPhraseIS{í lokin} \dicFlx{adv} \dicDirectTranslationCS{nakonec}
\dicEntry[loka] \dicTerm{lok|a\smash{\textsuperscript{1}}} \dicIPA{{l}{\textopeno}{\textlengthmark}{\r{g}}{a}} \dicPos{f}[1] \dicFlx{(‑u, ‑ur)}[19] \textbf{1.} \dicSynonym{hespa\smash{\textsuperscript{1}}} \dicDirectTranslationCS{západka, petlice, závora}  \textbf{2.} \dicFieldCat{anat.} \dicDirectTranslationCS{chlopeň}
\dicEntry[loka] \dicTerm{lok|a\smash{\textsuperscript{2}}} \dicsymFrequent\  \dicIPA{{l}{\textopeno}{\textlengthmark}{\r{g}}{a}} \dicPos{v}[1] \dicFlx{(‑aði)}[1] \dicFlx{dat} \dicDirectTranslationCS{(u)zavřít, zavírat} \dicExampleIS{loka glugganum} \dicExampleCS{zavřít okno};  \dicIdiom{loka}[að]{ \dicPhraseIS{loka að sér}} \dicDirectTranslationCS{zavřít za sebou (dveře)};  \dicIdiom{loka}[fyrir]{ \dicPhraseIS{loka fyrir e‑ð}} \dicDirectTranslationCS{zavřít (co), vypnout (co) (elektřinu ap.)};  \dicIdiom{loka}[inni]{ \dicPhraseIS{loka e‑n inni}} \dicDirectTranslationCS{zavřít (koho) (do blázince ap.)};  \dicIdiom{loka}[úti]{ \dicPhraseIS{loka e‑n úti}} \dicDirectTranslationCS{vyloučit (koho) (z~party ap.), nevpustit (koho) dovnitř};  \dicIdiom{lokast}{ \dicPhraseIS{lokast}} \dicFlx{refl} \dicDirectTranslationCS{zavřít se, zavírat se}
\dicEntry[lokaáfangi] \dicTerm{loka··á·fang|i} \dicIPA{{l}{\textopeno}{\textlengthmark}{\r{g}}{a}{au}{f}{au}{\textltailn}{\r{\textObardotlessj}}{\textsci}} \dicPos{m}[1] \dicFlx{(‑a, ‑ar)}[8] \dicDirectTranslationCS{poslední etapa}
\dicEntry[lokaður] \dicTerm{lok··|aður} \dicsymFrequent\  \dicIPA{{l}{\textopeno}{\textlengthmark}{\r{g}}{a}{ð}{\textscy}{\textsubring{r}}} \dicPos{adj}[3] \dicFlx{(f ‑uð)}[3] \textbf{1.} \dicSynonym*{ekki opinn} \dicDirectTranslationCS{zavřený, uzavřený} \dicExampleIS{lokaður vegur} \dicExampleCS{uzavřená cesta}  \textbf{2.} \dicLangCat{přen.} \dicDirectTranslationCS{uzavřený, odměřený, zdrženlivý, rezervovaný}
\dicEntry[lokafrestur] \dicTerm{loka··frest|ur} \dicIPA{{l}{\textopeno}{\textlengthmark}{\r{g}}{a}{f}{r}{\textepsilon}{s}{\textsubring{d}}{\textscy}{\textsubring{r}}} \dicPos{m}[6] \dicFlx{(‑s)}[7] \dicDirectTranslationCS{uzávěrka, (konečný) termín, deadline}
\dicEntry[lokaorð] \dicTerm{loka··orð} \dicIPA{{l}{\textopeno}{\textlengthmark}{\r{g}}{a}{\textopeno}{r}{\texttheta}} \dicPos{n}[2] \dicFlx{(‑s, ‑)}[5] \dicDirectTranslationCS{konečné\,/\addthin závěrečné slovo, doslov};  \dicPhraseIS{eiga lokaorðið} \dicLangCat{přen.} \dicDirectTranslationCS{mít poslední slovo}
\dicEntry[lokapróf] \dicTerm{loka··próf} \dicIPA{{l}{\textopeno}{\textlengthmark}{\r{g}}{a}{p\smash{\textsuperscript{h}}}{r}{ou}{f}} \dicPos{n}[2] \dicFlx{(‑s, ‑)}[5] \dicFieldCat{škol.} \dicDirectTranslationCS{závěrečná zkouška}
\dicEntry[lokaritgerð] \dicTerm{loka··rit·gerð} \dicIPA{{l}{\textopeno}{\textlengthmark}{\r{g}}{a}{r}{\textsci}{\textsubring{d}}{\r{\textObardotlessj}}{\textepsilon}{r}{\texttheta}} \dicPos{f}[7] \dicFlx{(‑ar, ‑ir)}[1] \dicFieldCat{škol.} \dicDirectTranslationCS{závěrečná práce}
\dicEntry[lokasjóður] \dicTerm{loka··sjóð|ur} \dicIPA{{l}{\textopeno}{\textlengthmark}{\r{g}}{a}{s}{j}{ou}{ð}{\textscy}{\textsubring{r}}} \dicPos{m}[9] \dicFlx{(‑s, ‑ir)}[6] \dicFieldCat{bot.} \dicDirectTranslationCS{kokrhel menší} \textit{(l.~{\textLA{Rhinanthus minor}})}  \dicsymPhoto\ 
\dicFigure{71849.jpg}{Lokasjóður}{Lokasjóður - ŠARŽÍK František, Biolib, Copyright/PD}
\dicEntry[lokaæfing] \dicTerm{loka··æf·ing} \dicIPA{{l}{\textopeno}{\textlengthmark}{\r{g}}{a}{a}{i}{v}{i}{\ng}{\r{g}}} \dicPos{f}[4] \dicFlx{(‑ar, ‑ar)}[5] \dicDirectTranslationCS{generální zkouška, generálka}
\dicEntry[lokhljóð] \dicTerm{lok··hljóð} \dicIPA{{l}{\textopeno}{\textlengthmark}{\r{g}}{\textsubring{l}}{j}{ou}{\texttheta}} \dicPos{n}[2] \dicFlx{(‑s, ‑)}[5] \dicFieldCat{jaz.} \dicDirectTranslationCS{ražená souhláska, ploziva}
\dicEntry[loki] \dicTerm{lok|i} \dicIPA{{l}{\textopeno}{\textlengthmark}{\r{\textObardotlessj}}{\textsci}} \dicPos{m}[1] \dicFlx{(‑a, ‑ar)}[1] \textbf{1.} \dicSynonym{hnútur} \dicDirectTranslationCS{uzel, zauzlení}  \textbf{2.} \dicFieldCat{techn.} \dicSynonym*{öryggisloki} \dicDirectTranslationCS{ventil, ventilek}
\dicEntry[lokið] \dicTerm{lokið} \dicIPA{{l}{\textopeno}{\textlengthmark}{\r{\textObardotlessj}}{\textsci}{\texttheta}} \dicPos{v} \dicFlx{supin} \dicLink{ljúka}
\dicEntry[lokinn] \dicTerm{lokinn} \dicsymFrequent\  \dicIPA{{l}{\textopeno}{\textlengthmark}{\r{\textObardotlessj}}{\textsci}{\textsubring{n}}} \dicPos{adj}[6]\dicFlx{}[-6] \dicSynonym{búinn} \dicDirectTranslationCS{skončený, uzavřený} \dicExampleIS{Rannsókn er lokið.} \dicExampleCS{Vyšetřování je uzavřené.};  \dicPhraseIS{e‑m er öllum lokið} \dicFlx{impers} \dicLangCat{přen.} \dicDirectTranslationCS{(kdo) toho má dost, (kdo) je hotový}
\dicEntry[lokka] \dicTerm{lokk|a} \dicIPA{{l}{\textopeno}{h}{\r{g}}{a}} \dicPos{v}[1] \dicFlx{(‑aði)}[1] \dicFlx{acc} \textbf{1.} \dicSynonym{ginna} \dicDirectTranslationCS{(na)lákat, zlákat} \dicExampleIS{lokka e‑n til sín} \dicExampleCS{nalákat (koho) k~sobě}  \textbf{2.} \dicDirectTranslationCS{vylákat, vymámit, vyloudit} \dicExampleIS{lokka e‑ð út úr e‑m} \dicExampleCS{vylákat (co) od (koho)}
\dicEntry[lokkur] \dicTerm{lokk|ur} \dicIPA{{l}{\textopeno}{h}{\r{g}}{\textscy}{\textsubring{r}}} \dicPos{m}[6] \dicFlx{(‑s, ‑ar)}[22] \dicDirectTranslationCS{lokna, kudrna}
\dicEntry[lokræsi] \dicTerm{lok··ræsi} \dicIPA{{l}{\textopeno}{\textlengthmark}{\r{g}}{r}{a}{i}{s}{\textsci}} \dicPos{n}[2] \dicFlx{(‑s, ‑)}[14] \dicLangCat{zast.} \dicDirectTranslationCS{stoka, odpad}
\dicEntry[loks] \dicTerm{loks} \dicsymFrequent\  \dicIPA{{l}{\textopeno}{\r{g}}{s}} \dicPos{adv} \dicSynonym*{að lokum} \dicDirectTranslationCS{nakonec, konečně} \dicExampleIS{Hurðin var loks opnuð.} \dicExampleCS{Dveře se nakonec otevřely.}
\dicEntry[loksins] \dicTerm{loksins} \dicsymFrequent\  \dicIPA{{l}{\textopeno}{\r{g}}{s}{\textsci}{n}{s}} \dicPos{adv} \dicSynonym*{á endanum} \dicDirectTranslationCS{konečně} \dicExampleIS{Fundur er loksins að hefjast.} \dicExampleCS{Schůze konečně začíná.}
\dicEntry[lokun] \dicTerm{lok|un} \dicIPA{{l}{\textopeno}{\textlengthmark}{\r{g}}{\textscy}{\textsubring{n}}} \dicPos{f}[7] \dicFlx{(‑unar, ‑anir)}[8] \textbf{1.} \dicSynonym*{það að loka} \dicDirectTranslationCS{zavírání, (u)zavření} \dicExampleIS{lokun búða} \dicExampleCS{zavření obchodu}  \textbf{2.} \dicSynonym*{það að setja hindrun} \dicDirectTranslationCS{uzavírka (cesty ap.)} \dicExampleIS{lokun vegar} \dicExampleCS{uzavírka cesty}
\dicEntry[lokunartími] \textls[15]{\dicTerm{lokunar··tím|i} \dicIPA{{l}{\textopeno}{\r{g}}{\textscy}{n}{a}{\textsubring{r}}{t\smash{\textsuperscript{h}}}{i}{m}{\textsci}} \dicPos{m}[1] \dicFlx{(‑a)}[3] \dicDirectTranslationCS{zavírací doba}} \dicAntonym{opnunartími}
\dicEntry[lopapeysa] \dicTerm{lopa··peys|a} \dicIPA{{l}{\textopeno}{\textlengthmark}{\textsubring{b}}{a}{p\smash{\textsuperscript{h}}}{ei}{s}{a}} \dicPos{f}[1] \dicFlx{(‑u, ‑ur)}[7] \dicIndirectTranslationCS{pletený vlněný svetr s~typickými islandskými vzory}
\dicEntry[lopi] \dicTerm{lop|i} \dicIPA{{l}{\textopeno}{\textlengthmark}{\textsubring{b}}{\textsci}} \dicPos{m}[1] \dicFlx{(‑a, ‑ar)}[1] \textbf{1.} \dicSynonym*{ullarþráður} \dicDirectTranslationCS{mykaná vlna}  \textbf{2.} \dicFieldCat{med.} \dicSynonym{bjúgur} \dicDirectTranslationCS{edém, otok};  \dicIdiom{lopi}{ \dicPhraseIS{spinna\,/\addthin teygja lopann}} \dicLangCat{přen.} \dicDirectTranslationCS{pustit si pusu na špacír}
\dicEntry[loppa] \dicTerm{lopp|a} \dicIPA{{l}{\textopeno}{h}{\textsubring{b}}{a}} \dicPos{f}[1] \dicFlx{(‑u, ‑ur)}[19] \dicDirectTranslationCS{tlapa, tlapka}
\dicEntry[los] \dicTerm{los} \dicIPA{{l}{\textopeno}{\textlengthmark}{s}} \dicPos{n}[2] \dicFlx{(‑s, ‑)}[5] \dicDirectTranslationCS{nestabilita, nestabilnost, nesoudržnost} \dicExampleIS{það er los á e‑u} \dicExampleCS{(co) je nestabilní}
\dicEntry[losa] \dicTerm{los|a} \dicsymFrequent\  \dicIPA{{l}{\textopeno}{\textlengthmark}{s}{a}} \dicPos{v}[1] \dicFlx{(‑aði)}[1] \dicFlx{acc} \textbf{1.} \dicDirectTranslationCS{uvolnit, uvolňovat, povolit, povolovat (šroubek ap.)} \dicExampleIS{Hann losar hjólið af bílnum.} \dicExampleCS{Uvolňuje kolo u~auta.}  \textbf{2.} \dicDirectTranslationCS{uvolnit, uvolňovat (pokoj ap.)} \dicExampleIS{losa hótelherbergi} \dicExampleCS{uvolnit hotelový pokoj}  \textbf{3.} \dicDirectTranslationCS{vyprázdnit, vyprazdňovat} \dicExampleIS{losa flöskuna} \dicExampleCS{vyprázdnit láhev}  \textbf{4.} \dicPhraseIS{e‑að losar (20 tonn)} \dicDirectTranslationCS{(co) má něco málo přes (20 tun), (co) přesahuje (20 tun)};  \dicIdiom{losa}[um]{ \dicPhraseIS{losa um e‑ð}} \dicDirectTranslationCS{uvolnit (co), zmírnit (co) (o\nobreak mezení ap.)};  \dicIdiom{losa}[úr]{ \dicPhraseIS{losa e‑n úr e‑u}} \dicDirectTranslationCS{uvolnit (koho) z~(čeho), vymanit (koho) z~(čeho)};  \dicIdiom{losa}[við]{ \dicPhraseIS{losa e‑n við e‑ð}} \dicDirectTranslationCS{pomoct (komu) s~(čím)}; { \dicPhraseIS{losa sig við e‑ð}} \dicDirectTranslationCS{zbavit se (čeho)}
\dicEntry[losaralegur] \dicTerm{losara··legur} \dicIPA{{l}{\textopeno}{\textlengthmark}{s}{a}{r}{a}{l}{\textepsilon}{\textbabygamma}{\textscy}{\textsubring{r}}} \dicPos{adj}[1]\dicFlx{}[-8] \dicSynonym{sundurlaus} \dicDirectTranslationCS{náhod\-ný, nahodilý, neplánovaný}
\dicEntry[losna] \dicTerm{losn|a} \dicsymFrequent\  \dicIPA{{l}{\textopeno}{s}{\textsubring{d}}{n}{a}} \dicPos{v}[1] \dicFlx{(‑aði)}[1] \textbf{1.} \dicDirectTranslationCS{uvolnit se, uprázdnit se (místo ap.)} \dicExampleIS{Starfið losnar bráðlega.} \dicExampleCS{Zaměstnání se brzy uvolní.}  \textbf{2.} \dicDirectTranslationCS{uvolnit\,/\addthin uvolňovat se, povolit, povolovat (šroubek ap.)} \dicExampleIS{Það hefur losnað hjól af bílnum.} \dicExampleCS{U~auta se uvolnilo kolo.}  \textbf{3.} \dicDirectTranslationCS{uvolnit\,/\addthin uvolňovat se (z~práce ap.)} \dicExampleIS{losna úr vinnunni} \dicExampleCS{uvolnit se z~práce};  \dicIdiom{losna}[undan]{ \dicPhraseIS{losna undan e‑u}} \dicDirectTranslationCS{vyhnout\,/\addthin vyhýbat se (čemu), uniknout\,/\addthin unikat (čemu)} \dicExampleIS{losna undan mörgum skyldum} \dicExampleCS{vyhýbat se mnoha povinnostem};  \dicIdiom{losna}[við]{ \dicPhraseIS{losna við e‑ð}} \dicSynonym*{koma e‑u af sér} \dicDirectTranslationCS{zbavit se (čeho)}
\dicEntry[lost] \dicTerm{lost} \dicIPA{{l}{\textopeno}{s}{\textsubring{d}}} \dicPos{n}[2] \dicFlx{(‑s, ‑)}[5] \textbf{1.} \dicFieldCat{med.} \dicDirectTranslationCS{šok}  \textbf{2.} \dicLangCat{hovor.} \dicDirectTranslationCS{šok, překvapení} \dicExampleIS{Ég var alveg í losti á eftir.} \dicExampleCS{Byl jsem po tom v~šoku.}
\dicEntry[lostafullur] \dicTerm{losta··|fullur} \dicIPA{{l}{\textopeno}{s}{\textsubring{d}}{a}{f}{\textscy}{\textsubring{d}}{l}{\textscy}{\textsubring{r}}} \dicPos{adj}[10] \dicFlx{(comp ‑fyllri, sup ‑fyllstur)}[7] \dicSynonym*{munaðargjarn} \dicDirectTranslationCS{chlípný, žádostivý, chtivý, vilný}
\dicEntry[losti] \dicTerm{lost|i} \dicIPA{{l}{\textopeno}{s}{\textsubring{d}}{\textsci}} \dicPos{m}[1] \dicFlx{(‑a)}[3] \dicSynonym{girnd} \dicDirectTranslationCS{chlípnost, žádostivost, chtíč, vilnost}
\dicEntry[lostið] \dicTerm{lostið} \dicIPA{{l}{\textopeno}{s}{\textsubring{d}}{\textsci}{\texttheta}} \dicPos{v} \dicFlx{supin} \dicLink{ljósta}
\dicEntry[lostinn] \dicTerm{lostinn} \dicIPA{{l}{\textopeno}{s}{\textsubring{d}}{\textsci}{\textsubring{n}}} \dicPos{adj}[6]\dicFlx{}[-6] \dicPhraseIS{skelfingu lostinn} \dicFlx{adj} \dicDirectTranslationCS{vyděšený, vystrašený};  \dicPhraseIS{vera furðu lostinn yfir e‑u} \dicDirectTranslationCS{oněmět úžasem nad (čím)};  \dicPhraseIS{vera þrumu lostinn yfir e‑u} \dicDirectTranslationCS{být šokován (čím)}
\dicEntry[lostæti] \dicTerm{lost··æti} \dicIPA{{l}{\textopeno}{s}{\textsubring{d}}{a}{i}{\textsubring{d}}{\textsci}} \dicPos{n}[2] \dicFlx{(‑s, ‑)}[14] \dicSynonym{góðgæti} \dicDirectTranslationCS{pochoutka, lahůdka, delikatesa}
\dicEntry[losun] \dicTerm{los|un} \dicIPA{{l}{\textopeno}{\textlengthmark}{s}{\textscy}{\textsubring{n}}} \dicPos{f}[7] \dicFlx{(‑unar)}[9] \textbf{1.} \dicSynonym*{losun sorps o.fl.} \dicDirectTranslationCS{vypouštění, vyprazdňování (odpadu, odpadků ap.)}  \textbf{2.} \dicSynonym*{afferming} \dicDirectTranslationCS{vykládka, vyložení} \dicExampleIS{losun á farminum} \dicExampleCS{vyložení nákladu}
\dicEntry[lota] \dicTerm{lot|a} \dicIPA{{l}{\textopeno}{\textlengthmark}{\textsubring{d}}{a}} \dicPos{f}[1] \dicFlx{(‑u, ‑ur)}[19] \dicDirectTranslationCS{období, cyklus, perioda, (časový) úsek} \dicExampleIS{í einni lotu} \dicExampleCS{na jeden zátah}
\dicEntry[lotið] \dicTerm{lotið} \dicIPA{{l}{\textopeno}{\textlengthmark}{\textsubring{d}}{\textsci}{\texttheta}} \dicPos{v} \dicFlx{supin} \dicLink{lúta}
\dicEntry[lotinn] \dicTerm{lotinn} \dicIPA{{l}{\textopeno}{\textlengthmark}{\textsubring{d}}{\textsci}{\textsubring{n}}} \dicPos{adj}[6]\dicFlx{}[-2] \dicSynonym{boginn} \dicDirectTranslationCS{shrbený} \dicExampleIS{lotinn í herðum} \dicExampleCS{shrbený v~zádech}
\dicEntry[lotning] \dicTerm{lot··ning} \dicIPA{{l}{\textopeno}{h}{\textsubring{d}}{n}{i}{\ng}{\r{g}}} \dicPos{f}[4] \dicFlx{(‑ar)}[7] \dicSynonym{virðing} \dicDirectTranslationCS{úcta, uctívání} \dicExampleIS{sýna e‑m lotningu} \dicExampleCS{projevovat ke (komu) úctu}
\begin{xtolerant}{}{1pt}
\dicEntry[lotningarhreimur] \dicTerm{lotningar··hreim|ur}\addthinS\dicIPA{{l}{\textopeno}{h}{\textsubring{d}}{n}{i}{\ng}{\r{g}}{a}{\textsubring{r}}{ei}{m}{\textscy}{\textsubring{r}}}\addthinS\dicPos{m}[6]\addthinS\dicFlx{(‑s, ‑ar)}[22] \dicDirectTranslationCS{pochlebovačný tón}
\end{xtolerant}
\dicEntry[lottó] \dicTerm{lottó} \dicIPA{{l}{\textopeno}{h}{\textsubring{d}}{ou}} \dicPos{n}[2] \dicFlx{(‑s, ‑)}[34] \dicDirectTranslationCS{loterie, sportka}
\dicEntry[lotugræðgi] \dicTerm{lotu··græðg|i} \dicIPA{{l}{\textopeno}{\textlengthmark}{\textsubring{d}}{\textscy}{\r{g}}{r}{a}{i}{ð}{\r{g}}{\textsci}} \dicPos{f}[3] \dicFlx{(‑i)}[3] \dicFieldCat{med.} \dicDirectTranslationCS{bulimie}
\dicEntry[lotukerfi] \dicTerm{lotu··kerfi} \dicIPA{{l}{\textopeno}{\textlengthmark}{\textsubring{d}}{\textscy}{c\smash{\textsuperscript{h}}}{\textepsilon}{r}{v}{\textsci}} \dicPos{n}[2] \dicFlx{(‑s)}[20] \dicFieldCat{chem.} \dicDirectTranslationCS{periodická tabulka\,/\addthin soustava prvků}
\dicEntry[lotutúlkun] \dicTerm{lotu··túlk|un} \dicIPA{{l}{\textopeno}{\textlengthmark}{\textsubring{d}}{\textscy}{t\smash{\textsuperscript{h}}}{u}{\textsubring{l}}{\r{g}}{\textscy}{\textsubring{n}}} \dicPos{f}[7] \dicFlx{(‑unar, ‑anir)}[8] \dicDirectTranslationCS{konsekutivní tlumočení}
\dicEntry[lóa] \dicTerm{ló|a} \dicIPA{{l}{ou}{\textlengthmark}{a}} \dicPos{f}[1] \dicFlx{(‑u, ‑ur)}[7] \dicFieldCat{zool.} \dicSynonym{heiðlóa} \dicDirectTranslationCS{kulík, kulík zlatý} \textit{(l.~{\textLA{Pluvialis apricaria}})}  \dicsymPhoto\ 
\dicFigure{21284.jpg}{Lóa}{Lóa - Karney, Lee, Biolib, PD}
\dicEntry[lóð] \dicTerm{lóð\smash{\textsuperscript{1}}} \dicIPA{{l}{ou}{\textlengthmark}{\texttheta}} \dicPos{f}[7] \dicFlx{(‑ar, ‑ir)}[1] \dicSynonym*{landspilda} \dicDirectTranslationCS{parcela, pozemek} \dicExampleIS{Húsið stendur á lóðinni.} \dicExampleCS{Dům stojí na pozemku.}
\dicEntry[lóð] \dicTerm{lóð\smash{\textsuperscript{2}}} \dicIPA{{l}{ou}{\textlengthmark}{\texttheta}} \dicPos{f}[7] \dicFlx{(‑ar, ‑ir)}[1] \dicFieldCat{nám.} \dicSynonym{fiskilína} \dicDirectTranslationCS{síť s~mnoha háčky};  \dicPhraseIS{leggja lóðirnar} {\textbf{a.}} \dicDirectTranslationCS{položit síť} \dicIndirectTranslationCS{(síť s~mnoha háčky)};  {\textbf{b.}} \dicSynonym*{selja upp af sjóveiki} \dicDirectTranslationCS{zvracet} \dicIndirectTranslationCS{(kvůli mořské nemoci)}
\dicEntry[lóð] \dicTerm{lóð\smash{\textsuperscript{3}}} \dicIPA{{l}{ou}{\textlengthmark}{\texttheta}} \dicPos{n}[2] \dicFlx{(‑s, ‑)}[5] \dicDirectTranslationCS{závaží} \dicExampleIS{lóð á vog} \dicExampleCS{závaží na váze};  \dicPhraseIS{það er lóðið} \dicLangCat{přen.} \dicDirectTranslationCS{to je jádro věci}
\dicEntry[lóða] \dicTerm{lóð|a} \dicIPA{{l}{ou}{\textlengthmark}{ð}{a}} \dicPos{v}[1] \dicFlx{(‑aði)}[1] \dicFlx{acc} \textbf{1.} \dicFieldCat{techn.} \dicDirectTranslationCS{letovat, pájet} \dicExampleIS{lóða e‑ð saman} \dicExampleCS{sletovat (co)}  \textbf{2.} \dicSynonym*{mæla dýpt} \dicDirectTranslationCS{sondovat (dno nebo moře pro výskyt ryb)};  \dicPhraseIS{lóða á fisk(i)} \dicFieldCat{nám.} \dicDirectTranslationCS{hledat ryby sonarem}
\dicEntry[lóðbolti] \dicTerm{lóð··bolt|i} \dicIPA{{l}{ou}{ð}{\textsubring{b}}{\textopeno}{\textsubring{l}}{\textsubring{d}}{\textsci}} \dicPos{m}[1] \dicFlx{(‑a, ‑ar)}[1] \dicFieldCat{techn.} \dicDirectTranslationCS{páječka, pájedlo}
\dicEntry[lóðréttur] \dicTerm{lóð··réttur} \dicIPA{{l}{ou}{ð}{r}{j}{\textepsilon}{h}{\textsubring{d}}{\textscy}{\textsubring{r}}} \dicPos{adj}[1]\dicFlx{}[-10] \dicDirectTranslationCS{svislý, vertikální, kolmý}
\dicEntry[lóðun] \dicTerm{lóð|un} \dicIPA{{l}{ou}{\textlengthmark}{ð}{\textscy}{\textsubring{n}}} \dicPos{f}[7] \dicFlx{(‑unar)}[9] \dicFieldCat{techn.} \dicSynonym*{það að lóða} \dicDirectTranslationCS{letování, pájení}
\dicEntry[lófaklapp] \dicTerm{lófa··klapp} \dicIPA{{l}{ou}{\textlengthmark}{v}{a}{k\smash{\textsuperscript{h}}}{l}{a}{h}{\textsubring{b}}} \dicPos{n}[2] \dicFlx{(‑s)}[2] \dicDirectTranslationCS{potlesk, aplaus}
\dicEntry[lófatak] \dicTerm{lófa··|tak} \dicIPA{{l}{ou}{\textlengthmark}{v}{a}{t\smash{\textsuperscript{h}}}{a}{\r{g}}} \dicPos{n}[2] \dicFlx{(‑taks, ‑tök)}[8] \dicSynonym{lófaklapp} \dicDirectTranslationCS{potlesk, aplaus}
\dicEntry[lófatölva] \dicTerm{lófa··tölv|a} \dicIPA{{l}{ou}{\textlengthmark}{v}{a}{t\smash{\textsuperscript{h}}}{\oe}{l}{v}{a}} \dicPos{f}[1] \dicFlx{(‑u, ‑ur)}[7] \dicFieldCat{poč.} \dicDirectTranslationCS{palmtop, PDA, kapesní počítač}
\dicEntry[lófi] \dicTerm{lóf|i} \dicsymFrequent\  \dicIPA{{l}{ou}{\textlengthmark}{v}{\textsci}} \dicPos{m}[1] \dicFlx{(‑a, ‑ar)}[1] \dicDirectTranslationCS{dlaň} \dicExampleIS{geyma hring í lófa sínum} \dicExampleCS{držet prsten ve své dlani}
\dicEntry[lófótur] \textls[15]{\dicTerm{ló··|fótur} \dicIPA{{l}{ou}{\textlengthmark}{f}{ou}{\textsubring{d}}{\textscy}{\textsubring{r}}} \dicPos{m}[13] \dicFlx{(‑fótar, ‑fætur)}[5] \dicFieldCat{bot.} \dicDirectTranslationCS{prustka obecná}} \textit{(l.~{\textLA{Hippuris vulgaris}})}  \dicsymPhoto\ 
\dicFigure{3847.jpg}{Lófótur}{Lófótur - Zicha Ondřej, Biolib, Copyright/CC-BY-NC}
\dicEntry[lóga] \dicTerm{lóg|a} \dicIPA{{l}{ou}{\textlengthmark}{a}} \dicPos{v}[1] \dicFlx{(‑aði)}[34] \dicFlx{dat} \textbf{1.} \dicSynonym{slátra} \dicDirectTranslationCS{porazit, porážet (ovci ap.)} \dicExampleIS{lóga fé} \dicExampleCS{porážet ovce}  \textbf{2.} \dicSynonym*{aflífa} \dicDirectTranslationCS{utratit, uspat (psa ap.)} \dicExampleIS{lóga hundi} \dicExampleCS{uspat psa}
\dicEntry[lógaritmi] \dicTerm{lógaritm|i} \dicIPA{{l}{ou}{\textlengthmark}{\r{g}}{a}{r}{\textsci}{h}{\textsubring{d}}{m}{\textsci}} \dicPos{m}[1] \dicFlx{(‑a, ‑ar)}[1] \dicFieldCat{mat.} \dicDirectTranslationCS{logaritmus}
\dicEntry[lómur] \dicTerm{lóm|ur} \dicIPA{{l}{ou}{\textlengthmark}{m}{\textscy}{\textsubring{r}}} \dicPos{m}[6] \dicFlx{(‑s, ‑ar)}[24] \textbf{1.} \dicFieldCat{zool.} \dicDirectTranslationCS{potáplice malá} \textit{(l.~{\textLA{Gavia stellata}})}  \dicsymPhoto\   \textbf{2.} \dicSynonym{kveinstafir} \dicDirectTranslationCS{nářky, stížnosti};  \dicPhraseIS{berja\,/\addthin lemja lóminn} \dicDirectTranslationCS{stěžovat si, naříkat}
\dicFigure{21071.jpg}{Lómur}{Lómur - Bowman, Tim, Biolib, PD}
\dicEntry[lón] \dicTerm{lón} \dicIPA{{l}{ou}{\textlengthmark}{\textsubring{n}}} \dicPos{n}[2] \dicFlx{(‑s, ‑)}[5] \dicDirectTranslationCS{laguna};  \dicPhraseIS{Bláa lónið} \dicFieldCat{geog.} \dicDirectTranslationCS{Modrá laguna} \dicIndirectTranslationCS{(laguna v~blízkosti Keflavíku)}
\dicEntry[lóuþræll] \dicTerm{lóu··þræl|l} \dicIPA{{l}{ou}{\textlengthmark}{\textscy}{\texttheta}{r}{a}{i}{\textsubring{d}}{\textsubring{l}}} \dicPos{m}[6] \dicFlx{(‑s, ‑ar)}[51] \dicFieldCat{zool.} \dicDirectTranslationCS{jespák, jespák obecný} \textit{(l.~{\textLA{Calidris alpina}})}  \dicsymPhoto\ 
\dicFigure{21101.jpg}{Lóuþræll}{Lóuþræll - Bowman, Tim, Biolib, PD}
\dicEntry[luggarasegl] \dicTerm{luggara··segl} \dicIPA{{l}{\textscy}{\r{g}}{\textlengthmark}{a}{r}{a}{s}{\textepsilon}{\r{g}}{\textsubring{l}}} \dicPos{n}[2] \dicFlx{(‑s, ‑)}[5] \dicFieldCat{nám.} \dicDirectTranslationCS{lugrová plachta}
\dicEntry[lugt] \dicTerm{lugt} \dicIPA{{l}{\textscy}{x}{\textsubring{d}}} \dicPos{f}[7] \dicFlx{(‑ar, ‑ir)}[1] \dicLink{lukt}
\dicEntry[lugum] \dicTerm{lugum} \dicIPA{{l}{\textscy}{\textlengthmark}{\textbabygamma}{\textscy}{\textsubring{m}}} \dicPos{v} \dicFlx{ind pf pl 1 pers} \dicLink{ljúga}
\dicEntry[lukið] \dicTerm{lukið} \dicIPA{{l}{\textscy}{\textlengthmark}{\r{\textObardotlessj}}{\textsci}{\texttheta}} \dicPos{v} \dicFlx{supin} \dicLink{lykja}
\dicEntry[lukka] \dicTerm{lukk|a} \dicIPA{{l}{\textscy}{h}{\r{g}}{a}} \dicPos{f}[1] \dicFlx{(‑u)}[5] \dicSynonym{hamingja} \dicDirectTranslationCS{štěstí, zdar} \dicExampleIS{óska e‑m til lukku með afmælið} \dicExampleCS{popřát (komu) štěstí k~narozeninám}
\dicEntry[lukkast] \dicTerm{lukk|ast} \dicIPA{{l}{\textscy}{h}{\r{g}}{a}{s}{\textsubring{d}}} \dicPos{v}[1] \dicFlx{(‑aðist)}[97] \dicFlx{refl} \dicSynonym{heppnast} \dicDirectTranslationCS{(po)dařit se, zdařit se} \dicExampleIS{Hjónabandið lukkast vel.} \dicExampleCS{Manželství se daří.}
\dicEntry[lukkulegur] \dicTerm{lukku··legur} \dicIPA{{l}{\textscy}{h}{\r{g}}{\textscy}{l}{\textepsilon}{\textbabygamma}{\textscy}{\textsubring{r}}} \dicPos{adj}[1]\dicFlx{}[-8] \dicSynonym{ánægður} \dicDirectTranslationCS{šťastný, spokojený, potěšený}
\dicEntry[lukt] \dicTerm{lukt}\dicTerm{, lugt} \dicIPA{{l}\-{\textscy}\-{x}\-{\textsubring{d}}\-} \dicPos{f}[7] \dicFlx{(‑ar, ‑ir)}[1] \textbf{1.} \dicDirectTranslationCS{lucerna, svítilna} \dicExampleIS{vasalukt} \dicExampleCS{kapesní svítilna}  \textbf{2.} \dicDirectTranslationCS{přední světlo (u~auta, kola ap.)}
\dicEntry[lukti] \dicTerm{lukti} \dicIPA{{l}{\textscy}{x}{\textsubring{d}}{\textsci}} \dicPos{v} \dicFlx{ind pf sg 1 pers} \dicLink{lykja}
\dicEntry[luktum] \dicTerm{luktum} \dicIPA{{l}{\textscy}{x}{\textsubring{d}}{\textscy}{\textsubring{m}}} \dicPos{v} \dicFlx{ind pf pl 1 pers} \dicLink{lykja}
\dicEntry[luktur] \dicTerm{luktur} \dicIPA{{l}{\textscy}{x}{\textsubring{d}}{\textscy}{\textsubring{r}}} \dicPos{adj}[1]\dicFlx{}[-10] \dicDirectTranslationCS{(u)zavřený, obklopený}
\dicEntry[lukum] \dicTerm{lukum} \dicIPA{{l}{\textscy}{\textlengthmark}{\r{g}}{\textscy}{\textsubring{m}}} \dicPos{v} \dicFlx{ind pf pl 1 pers} \dicLink{ljúka}
\dicEntry[luma] \dicTerm{lum|a} \dicIPA{{l}{\textscy}{\textlengthmark}{m}{a}} \dicPos{v}[1] \dicFlx{(‑aði)}[1] \dicPhraseIS{luma á e‑u} \dicDirectTranslationCS{ukrýt si (co), schovat si (co), ulít si (co)} \dicExampleIS{Hann lumar á peningum.} \dicExampleCS{Ukrývá si peníze.}
\dicEntry[lumbra] \dicTerm{lumbr|a} \dicIPA{{l}{\textscy}{m}{\textsubring{b}}{r}{a}} \dicPos{v}[1] \dicFlx{(‑aði)}[1] \dicDirectTranslationCS{zašantročit, zastrčit} \dicExampleIS{lumbra e‑u e‑s staðar} \dicExampleCS{zašantročit (co kde)};  \dicIdiom{lumbra}[á]{ \dicPhraseIS{lumbra á e‑m}} \dicSynonym{berja} \dicDirectTranslationCS{(z)bít (koho), (z)mlátit (koho)}
\dicEntry[lumma] \dicTerm{lumm|a} \dicIPA{{l}{\textscy}{m}{\textlengthmark}{a}} \dicPos{f}[1] \dicFlx{(‑u, ‑ur)}[7] \dicFieldCat{kulin.} \dicDirectTranslationCS{lívanec, lívaneček};  \dicPhraseIS{e‑að selst eins og heitar lummur} \dicFlx{refl} \dicLangCat{přen.} \dicDirectTranslationCS{(co) jde na dračku};  \dicPhraseIS{það er gömul lumma} \dicLangCat{přen.} \dicDirectTranslationCS{to je ohraná písnička}
\dicEntry[lummulegur] \dicTerm{lummu··legur} \dicIPA{{l}{\textscy}{m}{\textlengthmark}{\textscy}{l}{\textepsilon}{\textbabygamma}{\textscy}{\textsubring{r}}} \dicPos{adj}[1]\dicFlx{}[-8] \dicLangCat{hovor.} \dicSynonym{tilkomulítill} \dicDirectTranslationCS{nijaký, nezáživný, nudný}
\dicEntry[lund] \dicTerm{lund} \dicsymFrequent\  \dicIPA{{l}{\textscy}{n}{\textsubring{d}}} \dicPos{f}[7] \dicFlx{(‑ar, ‑ir)}[1] \textbf{1.} \dicSynonym{geð} \dicDirectTranslationCS{povaha, temperament} \dicExampleIS{vera léttur í lund} \dicExampleCS{být šťastný}  \textbf{2.} \dicSynonym{háttur\smash{\textsuperscript{1}}} \dicDirectTranslationCS{způsob}
\dicEntry[Lundey] \dicTerm{Lund··ey} \dicIPA{{l}{\textscy}{n}{\textsubring{d}}{ei}} \dicPos{f}[4] \dicFlx{(‑jar)}[17] \dicFieldCat{geog.} \dicDirectTranslationCS{Lundey} \dicIndirectTranslationCS{(ostrov severně od Islandu)}
\dicEntry[lundi] \dicTerm{lund|i} \dicIPA{{l}{\textscy}{n}{\textsubring{d}}{\textsci}} \dicPos{m}[1] \dicFlx{(‑a, ‑ar)}[1] \dicFieldCat{zool.} \dicDirectTranslationCS{papuchálek, papuchalk, papuchalk ploskozobý\,/\addthin bělobradý} \textit{(l.~{\textLA{Fratercula arctica}})}  \dicsymPhoto\ 
\dicFigure{ds_image_lundi_0_1.jpg}{Lundi}{Lundi - Tgo2002, CC BY-SA 2.5}
\dicEntry[lundir] \dicTerm{lundir} \dicIPA{{l}{\textscy}{n}{\textsubring{d}}{\textsci}{\textsubring{r}}} \dicPos{f}[7] \dicFlx{pl}[2] \dicDirectTranslationCS{jemné filety, svíčková}
\dicEntry[lundur] \dicTerm{lund|ur} \dicIPA{{l}{\textscy}{n}{\textsubring{d}}{\textscy}{\textsubring{r}}} \dicPos{m}[10] \dicFlx{(‑ar, ‑ir)}[4] \dicDirectTranslationCS{háj, lesík}
\dicEntry[Lundúnir] \dicTerm{Lundúnir} \dicIPA{{l}{\textscy}{n}{\textsubring{d}}{u}{n}{\textsci}{\textsubring{r}}} \dicPos{f}[12] \dicFlx{pl}[6] \dicFieldCat{geog.} \dicDirectTranslationCS{Londýn} \dicIndirectTranslationCS{(hlavní město Spojeného království)}
\dicEntry[lunga] \dicTerm{lung|a} \dicIPA{{l}{u}{\ng}{\r{g}}{a}} \dicPos{n}[1] \dicFlx{(‑a, ‑u)}[2] \dicFieldCat{anat.} \dicDirectTranslationCS{plíce}
\dicEntry[lungnabólga] \dicTerm{lungna··bólg|a} \dicIPA{{l}{u}{\ng}{n}{a}{\textsubring{b}}{ou}{l}{\r{g}}{a}} \dicPos{f}[1] \dicFlx{(‑u)}[5] \dicFieldCat{med.} \dicDirectTranslationCS{zápal plic, pneumonie}
\dicEntry[lungnakvef] \dicTerm{lungna··kvef} \dicIPA{{l}{u}{\ng}{n}{a}{k\smash{\textsuperscript{h}}}{v}{\textepsilon}{f}} \dicPos{n}[2] \dicFlx{(‑s)}[2] \dicFieldCat{med.} \dicDirectTranslationCS{zánět průdušek, bronchitida}
\dicEntry[lungnapípa] \dicTerm{lungna··píp|a} \dicIPA{{l}{u}{\ng}{n}{a}{p\smash{\textsuperscript{h}}}{i}{\textsubring{b}}{a}} \dicPos{f}[1] \dicFlx{(‑u, ‑ur)}[19] \dicFieldCat{anat.} \dicDirectTranslationCS{průduška}
\dicEntry[luntalegur] \dicTerm{lunta··legur} \dicIPA{{l}{\textscy}{\textsubring{n}}{\textsubring{d}}{a}{l}{\textepsilon}{\textbabygamma}{\textscy}{\textsubring{r}}} \dicPos{adj}[1]\dicFlx{}[-8] \dicSynonym{ólundarlegur} \dicDirectTranslationCS{mrzutý, nevrlý}
\dicEntry[luralegur] \dicTerm{lura··legur} \dicIPA{{l}{\textscy}{\textlengthmark}{r}{a}{l}{\textepsilon}{\textbabygamma}{\textscy}{\textsubring{r}}} \dicPos{adj}[1]\dicFlx{}[-8] \dicSynonym*{þyngslalegur} \dicDirectTranslationCS{nemotorný, neohrabaný}
\dicEntry[lustum] \dicTerm{lustum} \dicIPA{{l}{\textscy}{s}{\textsubring{d}}{\textscy}{\textsubring{m}}} \dicPos{v} \dicFlx{ind pf pl 1 pers} \dicLink{ljósta}
\dicEntry[lutum] \dicTerm{lutum} \dicIPA{{l}{\textscy}{\textlengthmark}{\textsubring{d}}{\textscy}{\textsubring{m}}} \dicPos{v} \dicFlx{ind pf pl 1 pers} \dicLink{lúta}
\dicEntry[lúalegur] \dicTerm{lúa··legur} \dicIPA{{l}{u}{\textlengthmark}{a}{l}{\textepsilon}{\textbabygamma}{\textscy}{\textsubring{r}}} \dicPos{adj}[1]\dicFlx{}[-8] \textbf{1.} \dicSynonym{þreytulegur} \dicDirectTranslationCS{vyčerpaný, unavený}  \textbf{2.} \dicSynonym{ódrengilegur} \dicDirectTranslationCS{hanebný, sprostý} \dicExampleIS{lúalegar getsakir} \dicExampleCS{hanebné narážky}
\dicEntry[lúð] \dicTerm{lúð} \dicIPA{{l}{u}{\textlengthmark}{\texttheta}} \dicPos{v} \dicFlx{supin} \dicLink{lýja}
\dicEntry[lúða] \dicTerm{lúð|a} \dicIPA{{l}{u}{\textlengthmark}{ð}{a}} \dicPos{f}[1] \dicFlx{(‑u, ‑ur)}[19] \dicFieldCat{zool.} \dicSynonym{flyðra} \dicDirectTranslationCS{platýs, platýs obecný} \textit{(l.~{\textLA{Hippoglossus hippoglossus}})}
\dicEntry[lúðalegur] \dicTerm{lúða··legur} \dicIPA{{l}{u}{\textlengthmark}{ð}{a}{l}{\textepsilon}{\textbabygamma}{\textscy}{\textsubring{r}}} \dicPos{adj}[1]\dicFlx{}[-8] \dicDirectTranslationCS{slizký, odporný}
\dicEntry[lúði] \dicTerm{lúð|i\smash{\textsuperscript{1}}} \dicIPA{{l}{u}{\textlengthmark}{ð}{\textsci}} \dicPos{m}[1] \dicFlx{(‑a, ‑ar)}[1] \dicDirectTranslationCS{křupan, balík}
\dicEntry[lúði] \dicTerm{lúði\smash{\textsuperscript{2}}} \dicIPA{{l}{u}{\textlengthmark}{ð}{\textsci}} \dicPos{v} \dicFlx{ind pf sg 1 pers} \dicLink{lýja}
\dicEntry[lúðrasveit] \dicTerm{lúðra··sveit} \dicIPA{{l}{u}{ð}{r}{a}{s}{v}{ei}{\textsubring{d}}} \dicPos{f}[7] \dicFlx{(‑ar, ‑ir)}[1] \dicFieldCat{hud.} \dicDirectTranslationCS{dechová kapela, dechovka}
\dicEntry[lúðum] \dicTerm{lúðum} \dicIPA{{l}{u}{\textlengthmark}{ð}{\textscy}{\textsubring{m}}} \dicPos{v} \dicFlx{ind pf pl 1 pers} \dicLink{lýja}
\dicEntry[lúður] \dicTerm{lúð|ur} \dicIPA{{l}{u}{\textlengthmark}{ð}{\textscy}{\textsubring{r}}} \dicPos{m}[5] \dicFlx{(‑urs, ‑rar)}[1] \dicFieldCat{hud.} \dicDirectTranslationCS{trumpeta} \dicExampleIS{blása í lúður} \dicExampleCS{foukat na trumpetu}
\dicEntry[lúðurþeytari] \dicTerm{lúður··þeyt·ar|i} \dicIPA{{l}{u}{\textlengthmark}{ð}{\textscy}{\textsubring{r}}{\texttheta}{ei}{\textsubring{d}}{a}{r}{\textsci}} \dicPos{m}[1] \dicFlx{(‑a, ‑ar)}[13] \dicDirectTranslationCS{trubač(ka)}
\dicEntry[lúga] \dicTerm{lúg|a} \dicIPA{{l}{u}{\textlengthmark}{a}} \dicPos{f}[1] \dicFlx{(‑u, ‑ur)}[7] \dicDirectTranslationCS{poklop, příklop}
\dicEntry[lúi] \dicTerm{lú|i} \dicIPA{{l}{u}{\textlengthmark}{\textsci}} \dicPos{m}[1] \dicFlx{(‑a)}[3] \dicSynonym{þreyta\smash{\textsuperscript{1}}} \dicDirectTranslationCS{znavení, vyčerpanost, vysílení} \dicExampleIS{finna fyrir lúa eftir ferðina} \dicExampleCS{cítit vysílení po cestě}
\dicEntry[lúið] \dicTerm{lúið} \dicIPA{{l}{u}{\textlengthmark}{\textsci}{\texttheta}} \dicPos{v} \dicFlx{supin} \dicLink{lýja}
\dicEntry[lúinn] \dicTerm{lúinn} \dicIPA{{l}{u}{\textlengthmark}{\textsci}{\textsubring{n}}} \dicPos{adj}[6]\dicFlx{}[-2] \textbf{1.} \dicSynonym{þreyttur} \dicDirectTranslationCS{znavený, vyčerpaný, vysílený}  \textbf{2.} \dicSynonym{slitinn} \dicDirectTranslationCS{opotřebovaný, obnošený}
\dicEntry[lúka] \dicTerm{lúk|a} \dicIPA{{l}{u}{\textlengthmark}{\r{g}}{a}} \dicPos{f}[1] \dicFlx{(‑u, ‑ur)}[19] \textbf{1.} \dicSynonym{lófi} \dicDirectTranslationCS{dlaň};  \dicPhraseIS{vera með lífið í lúkunum} \dicLangCat{přen.} \dicDirectTranslationCS{mít srdce v~kalhotách}  \textbf{2.} \dicSynonym{handfylli} \dicDirectTranslationCS{hrst} \dicExampleIS{lúka af salti} \dicExampleCS{hrst soli}
\dicEntry[lúkar] \dicTerm{lúkar} \dicIPA{{l}{u}{\textlengthmark}{\r{g}}{a}{\textsubring{r}}} \dicPos{m}[4] \dicFlx{(‑s, ‑ar)}[15] \dicFieldCat{nám.} \dicSynonym{káeta} \dicDirectTranslationCS{důstojnická kabina\,/\addthin kajuta}
\dicEntry[lúlla] \dicTerm{lúll|a} \dicIPA{{l}{u}{l}{\textlengthmark}{a}} \dicPos{v}[1] \dicFlx{(‑aði)}[44] \dicLangCat{dět.} \dicSynonym{sofa} \dicDirectTranslationCS{spinkat, hajat} \dicExampleIS{Farðu nú að lúlla.} \dicExampleCS{Jdi spinkat.}
\dicEntry[lúmskur] \dicTerm{lúmskur} \dicIPA{{l}{u}{m}{s}{\r{g}}{\textscy}{\textsubring{r}}} \dicPos{adj}[1]\dicFlx{}[-1] \dicSynonym{undirförull} \dicDirectTranslationCS{podlý, zákeřný};  \dicPhraseIS{hafa lúmskt gaman af að (gera e‑ð)} \dicDirectTranslationCS{mít škodolibou radost z~(dělání (čeho))}
\dicEntry[lúpína] \dicTerm{lúpín|a} \dicIPA{{l}{u}{\textlengthmark}{p\smash{\textsuperscript{h}}}{i}{n}{a}} \dicPos{f}[1] \dicFlx{(‑u, ‑ur)}[7] \dicFieldCat{bot.} \dicDirectTranslationCS{lupina nutkajská, vlčí bob} \textit{(l.~{\textLA{Lupinus nootkatensis}})}  \dicsymPhoto\ 
\dicFigure{ds_image_lupina_0_1.jpg}{Lúpína}{Lúpína - Jutta234, GFDL}
\dicEntry[lúpulegur] \dicTerm{lúpu··legur} \dicIPA{{l}{u}{\textlengthmark}{\textsubring{b}}{\textscy}{l}{\textepsilon}{\textbabygamma}{\textscy}{\textsubring{r}}}\addthinS\dicPos{adj}[1]\dicFlx{}[-8] \dicSynonym{skömmustulegur} \dicDirectTranslationCS{zahanbený, (jsoucí) v~rozpacích}
\dicEntry[lúr] \dicTerm{lúr} \dicIPA{{l}{u}{\textlengthmark}{\textsubring{r}}} \dicPos{m}[4] \dicFlx{(‑s, ‑ar)}[14] \dicSynonym{blundur} \dicDirectTranslationCS{zdřímnutí, šlofík} \dicExampleIS{fá sér svolítinn lúr} \dicExampleCS{trochu si zdřímnout}
\dicEntry[lúra] \dicTerm{lúr|a} \dicIPA{{l}{u}{\textlengthmark}{r}{a}} \dicPos{v}[2] \dicFlx{(‑ði, ‑að\,/\addthin ‑t)}[184] \dicSynonym{sofa} \dicDirectTranslationCS{zdřímnout si, dřímat, podřimovat};  \dicIdiom{lúra}[á]{ \dicPhraseIS{lúra á e‑u}} \dicDirectTranslationCS{schovávat (co), ukrýt (si) (co)} \dicExampleIS{lúra á góðgæti} \dicExampleCS{schovávat dobroty}
\dicEntry[lús] \dicTerm{lús} \dicIPA{{l}{u}{\textlengthmark}{s}} \dicPos{f}[9] \dicFlx{(‑ar, lýs)}[6] \dicFieldCat{zool.} \dicDirectTranslationCS{veš} \textit{(l.~{\textLA{Anoplura}})};  \dicPhraseIS{læðast eins og lús með saum(i)} \dicFlx{refl} \dicLangCat{přen.} \dicDirectTranslationCS{loudat se jako hlemýžď};  \dicPhraseIS{lús kviknaði milli þeirra} \dicLangCat{přen.} \dicDirectTranslationCS{dostali se do křížku};  \dicPhraseIS{nú detta mér allar dauðar lýs úr höfði} \dicLangCat{přen.} \dicDirectTranslationCS{to jsem z~toho jelen}
\dicEntry[lúsarlegur] \dicTerm{lúsar··legur} \dicIPA{{l}{u}{\textlengthmark}{s}{a}{r}{l}{\textepsilon}{\textbabygamma}{\textscy}{\textsubring{r}}} \dicPos{adj}[1]\dicFlx{}[-8] \dicSynonym{smásálarlegur} \dicDirectTranslationCS{ubohý, tristní, mizerný (plat ap.)}
\dicEntry[lúta] \dicTerm{lúta} \dicsymFrequent\  \dicIPA{{l}{u}{\textlengthmark}{\textsubring{d}}{a}} \dicPos{v}[6] \dicFlx{(lýt, laut, lutum, lyti, lotið)}[106] \dicFlx{dat} \textbf{1.} \dicSynonym{beygja\smash{\textsuperscript{2}}} \dicDirectTranslationCS{sklonit, sklánět, sehnout, shýbat};  \dicPhraseIS{lúta höfði} \dicDirectTranslationCS{sklonit hlavu}  \textbf{2.} \dicDirectTranslationCS{sklonit se, podrobit se};  \dicPhraseIS{lúta e‑m} \dicDirectTranslationCS{sklonit se před (kým), podrobit se (komu)};  \dicIdiom{lúta}[að]{ \dicPhraseIS{lúta að e‑u}} \dicDirectTranslationCS{vztahovat se na (co), týkat se (čeho)};  \dicIdiom{lúta}{ \dicPhraseIS{lúta (svo) lágt}} \dicLangCat{přen.} \dicDirectTranslationCS{snížit se, klesnout} \dicIndirectTranslationCS{(propůjčit se k~něčemu špatnému)}; { \dicPhraseIS{lúta í lægra haldi fyrir e‑m}} \dicLangCat{přen.} \dicDirectTranslationCS{sklonit se před (kým)} \dicIndirectTranslationCS{(přijmout prohru)}
\dicEntry[lúterskur] \dicTerm{lúterskur} \dicIPA{{l}{u}{\textlengthmark}{t\smash{\textsuperscript{h}}}{e}{\textsubring{r}}{s}{\r{g}}{\textscy}{\textsubring{r}}} \dicPos{adj}[1]\dicFlx{}[-1] \dicFieldCat{náb.} \dicDirectTranslationCS{luteránský}
\dicEntry[lúterstrú] \dicTerm{lúters··trú} \dicIPA{{l}{u}{\textlengthmark}{t\smash{\textsuperscript{h}}}{\textepsilon}{\textsubring{r}}{s}{t\smash{\textsuperscript{h}}}{r}{u}} \dicPos{f}[4] \dicFlx{(‑ar)}[27] \dicFieldCat{náb.} \dicDirectTranslationCS{luteránství}
\dicEntry[lútur] \dicTerm{lút|ur} \dicIPA{{l}{u}{\textlengthmark}{\textsubring{d}}{\textscy}{\textsubring{r}}} \dicPos{m}[6] \dicFlx{(‑s\,/\addthin ‑ar, ‑ar)}[65] \dicFieldCat{chem.} \dicDirectTranslationCS{louh}
\dicEntry[Lúxemborg] \dicTerm{Lúxem··borg} \dicIPA{{l}{u}{x}{s}{\textepsilon}{m}{\textsubring{b}}{\textopeno}{r}{\r{g}}} \dicPos{f}[7] \dicFlx{(‑ar)}[4] \textbf{1.} \dicFieldCat{geog.} \dicDirectTranslationCS{Lucembursko}  \textbf{2.} \dicFieldCat{geog.} \dicDirectTranslationCS{Lucemburk} \dicIndirectTranslationCS{(hlavní město Lucemburska)}
\dicEntry[Lúxemborgari] \dicTerm{Lúxem··borg·ar|i} \dicIPA{{l}{u}{x}{s}{\textepsilon}{m}{\textsubring{b}}{\textopeno}{r}{\r{g}}{a}{r}{\textsci}} \dicPos{m}[1] \dicFlx{(‑a, ‑ar)}[13] \dicDirectTranslationCS{Lucemburčan(ka)}
\dicEntry[lúxemborgíska] \dicTerm{lúxem··borgísk|a} \dicIPA{{l}{u}{x}{s}{\textepsilon}{m}{\textsubring{b}}{\textopeno}{r}{\r{\textObardotlessj}}{i}{s}{\r{g}}{a}} \dicPos{f}[1] \dicFlx{(‑u)}[5] \dicDirectTranslationCS{lucemburština}
\dicEntry[lúxemborgskur] \dicTerm{lúxem··borgskur} \dicIPA{{l}{u}{x}{s}{\textepsilon}{m}{\textsubring{b}}{\textopeno}{r}{\r{g}}{s}{\r{g}}{\textscy}{\textsubring{r}}} \dicPos{adj}[1]\dicFlx{}[-6] \dicDirectTranslationCS{lucemburský}
\dicEntry[lúxus] \dicTerm{lúxus} \dicIPA{{l}{u}{x}{s}{\textscy}{s}} \dicPos{m}[4] \dicFlx{(‑s)}[21] \dicDirectTranslationCS{luxus, přepych}
\dicEntry[lyf] \dicTerm{lyf} \dicsymFrequent\  \dicIPA{{l}{\textsci}{\textlengthmark}{f}} \dicPos{n}[2] \dicFlx{(‑s, ‑)}[13] \dicDirectTranslationCS{lék} \dicExampleIS{lyf gegn lyfseðli} \dicExampleCS{lék na předpis};  \dicPhraseIS{lyf gegn e‑u} \dicDirectTranslationCS{lék proti (čemu), lék na (co)} \dicExampleIS{lyf gegn kvefi} \dicExampleCS{lék na rýmu};  \dicPhraseIS{lyf við e‑ð} \dicDirectTranslationCS{lék na (co)};  \dicPhraseIS{taka inn lyf} \dicDirectTranslationCS{brát léky};  \dicPhraseIS{verja á lyfjum} \dicDirectTranslationCS{být na prášcích}
\dicEntry[lyfjabúð] \dicTerm{lyfja··búð} \dicIPA{{l}{\textsci}{v}{j}{a}{\textsubring{b}}{u}{\texttheta}} \dicPos{f}[7] \dicFlx{(‑ar, ‑ir)}[1] \dicDirectTranslationCS{lékárna}
\dicEntry[lyfjafræði] \dicTerm{lyfja··fræð|i} \dicIPA{{l}{\textsci}{v}{j}{a}{f}{r}{a}{i}{ð}{\textsci}} \dicPos{f}[3] \dicFlx{(‑i)}[3] \textbf{1.} \dicDirectTranslationCS{lékárnictví, farmacie}  \textbf{2.} \dicDirectTranslationCS{farmakologie}
\dicEntry[lyfjafræðingur] \dicTerm{lyfja·fræð··ing|ur} \dicIPA{{l}{\textsci}{v}{j}{a}{f}{r}{a}{i}{ð}{i}{\ng}{\r{g}}{\textscy}{\textsubring{r}}} \dicPos{m}[6] \dicFlx{(‑s, ‑ar)}[8] \textbf{1.} \dicDirectTranslationCS{lékárník, lékárnice, farmaceut(ka)}  \textbf{2.} \dicDirectTranslationCS{farmakolog, farmakoložka}
\dicEntry[lyfjagras] \dicTerm{lyfja··|gras} \dicIPA{{l}{\textsci}{v}{j}{a}{\r{g}}{r}{a}{s}} \dicPos{n}[2] \dicFlx{(‑grass, ‑grös)}[8] \dicFieldCat{bot.} \dicDirectTranslationCS{tučnice obecná} \textit{(l.~{\textLA{Pinguicula vulgaris}})}  \dicsymPhoto\ 
\dicFigure{31254.jpg}{Lyfjagras}{Lyfjagras - Dvořák Václav, Biolib, Copyright/CC-BY-NC}
\dicEntry[lyfjakort] \dicTerm{lyfja··kort} \dicIPA{{l}{\textsci}{v}{j}{a}{k\smash{\textsuperscript{h}}}{\textopeno}{\textsubring{r}}{\textsubring{d}}} \dicPos{n}[2] \dicFlx{(‑s, ‑)}[5] \dicIndirectTranslationCS{karta pacienta s~informacemi o~užívaných lécích}
\dicEntry[lyfjakostnaður] \dicTerm{lyfja··kost·nað|ur} \dicIPA{{l}{\textsci}{v}{j}{a}{k\smash{\textsuperscript{h}}}{\textopeno}{s}{\textsubring{d}}{n}{a}{ð}{\textscy}{\textsubring{r}}} \dicPos{m}[10] \dicFlx{(‑ar)}[9] \dicDirectTranslationCS{náklady na léky}
\dicEntry[lyfjaskírteini] \dicTerm{lyfja··skír·teini} \dicIPA{{l}{\textsci}{v}{j}{a}{s}{\r{\textObardotlessj}}{i}{\textsubring{r}}{t\smash{\textsuperscript{h}}}{ei}{n}{\textsci}} \dicPos{n}[2] \dicFlx{(‑s, ‑)}[14] \dicIndirectTranslationCS{průkaz opravňující ke slevě na léky}
\dicEntry[lyfjaverð] \dicTerm{lyfja··verð} \dicIPA{{l}{\textsci}{v}{j}{a}{v}{\textepsilon}{r}{\texttheta}} \dicPos{n}[2] \dicFlx{(‑s)}[2] \dicDirectTranslationCS{cena léku}
\dicEntry[lyfjaverslun] \dicTerm{lyfja··versl|un} \dicIPA{{l}{\textsci}{v}{j}{a}{v}{\textepsilon}{\textsubring{r}}{s}{\textsubring{d}}{l}{\textscy}{\textsubring{n}}} \dicPos{f}[7] \dicFlx{(‑unar, ‑anir)}[8] \dicSynonym{lyfjabúð} \dicDirectTranslationCS{lékárna}
\dicEntry[lyflækningar] \dicTerm{lyf··lækn·ingar} \dicIPA{{l}{\textsci}{v}{l}{a}{i}{h}{\r{g}}{n}{i}{\ng}{\r{g}}{a}{\textsubring{r}}} \dicPos{f}[4] \dicFlx{pl}[6] \dicFieldCat{med.} \dicDirectTranslationCS{klinická medicína}
\dicEntry[lyfsali] \dicTerm{lyf··sal|i} \dicIPA{{l}{\textsci}{f}{s}{a}{l}{\textsci}} \dicPos{m}[1] \dicFlx{(‑a, ‑ar)}[8] \dicDirectTranslationCS{lékárník, lékárnice}
\dicEntry[lyfseðill] \dicTerm{lyf··seð|ill} \dicIPA{{l}{\textsci}{f}{s}{\textepsilon}{ð}{\textsci}{\textsubring{d}}{\textsubring{l}}} \dicPos{m}[6] \dicFlx{(‑ils, ‑lar)}[35] \dicDirectTranslationCS{(lékařský) předpis}
\dicEntry[lyfta] \dicTerm{lyft|a\smash{\textsuperscript{1}}} \dicIPA{{l}{\textsci}{f}{\textsubring{d}}{a}} \dicPos{f}[1] \dicFlx{(‑u, ‑ur)}[7] \textbf{1.} \dicDirectTranslationCS{výtah}  \textbf{2.} \dicDirectTranslationCS{zdviž}  \textbf{3.} \dicSynonym{skíðalyfta} \dicDirectTranslationCS{vlek (lyžařský ap.)}
\dicEntry[lyfta] \dicTerm{lyft|a\smash{\textsuperscript{2}}} \dicsymFrequent\  \dicIPA{{l}{\textsci}{f}{\textsubring{d}}{a}} \dicPos{v}[2] \dicFlx{(‑i, ‑)}[14] \dicFlx{dat} \dicSynonym{hefja} \dicDirectTranslationCS{zvednout, zvedat, zdvihat, nadzvednout, nadzdvihnout} \dicExampleIS{lyfta steini} \dicExampleCS{zvednout kámen};  \dicPhraseIS{lyfta brúnum} \dicDirectTranslationCS{zvednout obočí};  \dicIdiom{lyfta}[sér]{ \dicPhraseIS{lyfta sér}} \dicDirectTranslationCS{(na)kynout (těsto ap.)};  \dicIdiom{lyfta}[upp]{ \dicPhraseIS{lyfta sér upp}} \dicSynonym*{skemmta sér} \dicDirectTranslationCS{povyrazit si, odvázat se};  \dicIdiom{lyftast}{ \dicPhraseIS{lyftast}} \dicFlx{refl} \dicDirectTranslationCS{zvednout\,/\addthin zvedat se (letadlo ap.)}
\dicEntry[lyftari] \dicTerm{lyft··ar|i} \dicIPA{{l}{\textsci}{f}{\textsubring{d}}{a}{r}{\textsci}} \dicPos{m}[1] \dicFlx{(‑a, ‑ar)}[13] \dicFieldCat{techn.} \dicDirectTranslationCS{vysokozdvižný vozík}
\dicEntry[lyftiduft] \dicTerm{lyfti··duft} \dicIPA{{l}{\textsci}{f}{\textsubring{d}}{\textsci}{\textsubring{d}}{\textscy}{f}{\textsubring{d}}} \dicPos{n}[2] \dicFlx{(‑s)}[2] \dicDirectTranslationCS{prášek do pečiva}
\dicEntry[lyfting] \dicTerm{lyft··ing} \dicIPA{{l}{\textsci}{f}{\textsubring{d}}{i}{\ng}{\r{g}}} \dicPos{f}[4] \dicFlx{(‑ar, ‑ar)}[5] \textbf{1.} \dicSynonym*{það að lyfta} \dicDirectTranslationCS{zvednutí, zvedání, (vy)zdvižení, nadzvednutí}  \textbf{2.} \dicSynonym*{fjörgun} \dicDirectTranslationCS{pozvednutí, povznesení (na duchu ap.)} \dicExampleIS{lyfting andans} \dicExampleCS{povznesení na duchu}  \textbf{3.} \dicPhraseIS{lyftingar} \dicFlx{pl} \dicFieldCat{sport.} \dicDirectTranslationCS{vzpírání}
\begin{xtolerant}{}{1pt}
\dicEntry[lyftingamaður] \dicTerm{lyftinga··|maður}\addthinS\dicIPA{{l}{\textsci}{f}{\textsubring{d}}{i}{\ng}{\r{g}}{a}{m}{a}{ð}{\textscy}{\textsubring{r}}}\addthinS\dicPos{m}[13]\addthinS\dicFlx{(‑manns, ‑menn)}[2] \dicDirectTranslationCS{vzpě\-rač(ka)}
\end{xtolerant}
\dicEntry[lygamælir] \dicTerm{lyga··mæl|ir} \dicIPA{{l}{\textsci}{\textlengthmark}{\textbabygamma}{a}{m}{a}{i}{l}{\textsci}{\textsubring{r}}} \dicPos{m}[7] \dicFlx{(‑is, ‑ar)}[1] \dicDirectTranslationCS{detektor lži}
\dicEntry[lygar] \dicTerm{lygar} \dicIPA{{l}{\textsci}{\textlengthmark}{\textbabygamma}{a}{\textsubring{r}}} \dicPos{f} \dicFlx{pl nom} \dicLink{lygi\smash{\textsuperscript{1}}}
\dicEntry[lygari] \dicTerm{lyg··ar|i} \dicIPA{{l}{\textsci}{\textlengthmark}{\textbabygamma}{a}{r}{\textsci}} \dicPos{m}[1] \dicFlx{(‑a, ‑ar)}[13] \dicDirectTranslationCS{lhář(ka)}
\dicEntry[lygasaga] \dicTerm{lyga··|saga} \dicIPA{{l}{\textsci}{\textlengthmark}{\textbabygamma}{a}{s}{a}{\textbabygamma}{a}} \dicPos{f}[1] \dicFlx{(‑sögu, ‑sögur)}[14] \dicDirectTranslationCS{báchorka, pohádka}
\dicEntry[lygi] \dicTerm{lyg|i\smash{\textsuperscript{1}}} \dicsymFrequent\  \dicIPA{{l}{i}{j}{\textlengthmark}{\textsci}} \dicPos{f}[2] \dicFlx{(‑i, ‑ar)}[1] \dicDirectTranslationCS{lež, nepravda} \dicExampleIS{Bæjarstjórinn segir lygar.} \dicExampleCS{Starosta lže.};  \dicPhraseIS{fara með lygar} \dicDirectTranslationCS{lhát}
\dicEntry[lygi] \dicTerm{lygi\smash{\textsuperscript{2}}} \dicIPA{{l}{i}{j}{\textlengthmark}{\textsci}} \dicPos{v} \dicFlx{con pf sg 1 pers} \dicLink{ljúga}
\dicEntry[lygilegur] \textls[15]{\dicTerm{lygi··legur} \dicIPA{{l}{i}{j}{\textlengthmark}{\textsci}{l}{\textepsilon}{\textbabygamma}{\textscy}{\textsubring{r}}} \dicPos{adj}[1]\dicFlx{}[-8] \textbf{1.} \dicSynonym{ótrúlegur} \dicDirectTranslationCS{nevěrohodný, nepravděpodobný}}  \textbf{2.} \dicSynonym{ósannur} \dicDirectTranslationCS{nepravdivý, lživý}
\dicEntry[lyginn] \dicTerm{lyginn} \dicIPA{{l}{i}{j}{\textlengthmark}{\textsci}{\textsubring{n}}} \dicPos{adj}[6]\dicFlx{}[-2] \dicDirectTranslationCS{ulhaný, prolhaný}
\dicEntry[lygn] \dicTerm{lygn} \dicIPA{{l}{\textsci}{\r{g}}{\textsubring{n}}} \dicPos{adj}[5]\dicFlx{}[-5] \textbf{1.} \dicSynonym*{vindlaus} \dicDirectTranslationCS{bezvětrný, klidný} \dicIndirectTranslationCS{(o~počasí)}  \textbf{2.} \dicSynonym*{straumlaus} \dicDirectTranslationCS{nehybný, nezčeřený, klidný} \dicIndirectTranslationCS{(o~vodě)}
\dicEntry[lygna] \dicTerm{lygn|a\smash{\textsuperscript{1}}} \dicIPA{{l}{\textsci}{\r{g}}{n}{a}} \dicPos{f}[1] \dicFlx{(‑u, ‑ur)}[7] \textbf{1.} \dicSynonym*{kyrrt vatn} \dicDirectTranslationCS{klidná voda}  \textbf{2.} \dicSynonym*{kyrrt loft} \dicDirectTranslationCS{bezvětří}
\dicEntry[lygna] \dicTerm{lygn|a\smash{\textsuperscript{2}}} \dicIPA{{l}{\textsci}{\r{g}}{n}{a}} \dicPos{v}[2] \dicFlx{(‑di, ‑t)}[150] \dicFlx{dat} \textbf{1.} \dicPhraseIS{það lygnir} \dicFlx{impers} \dicSynonym*{kyrrast} \dicDirectTranslationCS{utišuje se (vítr), přestává foukat}  \textbf{2.} \dicPhraseIS{lygna aftur augunum} \dicDirectTranslationCS{zavřít\,/\addthin přivřít oči}
\dicEntry[lyk] \dicTerm{lyk} \dicIPA{{l}{\textsci}{\textlengthmark}{\r{g}}} \dicPos{v} \dicFlx{ind praes sg 1 pers} \dicLink{lykja}
\dicEntry[lyki] \dicTerm{lyki} \dicIPA{{l}{\textsci}{\textlengthmark}{\r{\textObardotlessj}}{\textsci}} \dicPos{v} \dicFlx{con pf sg 1 pers} \dicLink{ljúka}
\dicEntry[lykilaðstaða] \dicTerm{lykil··að·|staða} \dicIPA{{l}{\textsci}{\textlengthmark}{\r{\textObardotlessj}}{\textsci}{l}{a}{ð}{s}{\textsubring{d}}{a}{ð}{a}} \dicPos{f}[1] \dicFlx{(‑stöðu)}[2] \dicDirectTranslationCS{klíčové postavení}
\dicEntry[lykilatriði] \dicTerm{lykil··at·riði} \dicIPA{{l}{\textsci}{\textlengthmark}{\r{\textObardotlessj}}{\textsci}{l}{a}{\textsubring{d}}{r}{\textsci}{ð}{\textsci}} \dicPos{n}[2] \dicFlx{(‑s, ‑)}[14] \dicDirectTranslationCS{klíčový bod, klíčová otázka}
\dicEntry[lykill] \dicTerm{lyk|ill} \dicsymFrequent\  \dicIPA{{l}{\textsci}{\textlengthmark}{\r{\textObardotlessj}}{\textsci}{\textsubring{d}}{\textsubring{l}}} \dicPos{m}[6] \dicFlx{(‑ils, ‑lar)}[35] \textbf{1.} \dicDirectTranslationCS{klíč (od dveří ap.)};  \dicPhraseIS{lykill að e‑u} \dicDirectTranslationCS{klíč od (čeho)} \dicExampleIS{lykill að húsinu} \dicExampleCS{klíč od domu}  \textbf{2.} \dicFieldCat{techn.} \dicDirectTranslationCS{klíč} \dicIndirectTranslationCS{(ruční nástroj)}  \textbf{3.} \dicFieldCat{hud.} \dicDirectTranslationCS{klíč} \dicIndirectTranslationCS{(hudební značka)} \dicExampleIS{nótnalykill} \dicExampleCS{hudební klíč}  \textbf{4.} \dicSynonym{skýring} \dicDirectTranslationCS{klíč, nápověda} \dicIndirectTranslationCS{(pomůcka podstatná pro vysvětlení něčeho)} \dicExampleIS{orðalykill} \dicExampleCS{konkordance}
\dicEntry[lykilmaður] \dicTerm{lykil··|maður} \dicIPA{{l}{\textsci}{\textlengthmark}{\r{\textObardotlessj}}{\textsci}{l}{m}{a}{ð}{\textscy}{\textsubring{r}}} \dicPos{m}[13] \dicFlx{(‑manns, ‑menn)}[2] \dicDirectTranslationCS{klíčová postava}
\dicEntry[lykilorð] \dicTerm{lykil··orð} \dicIPA{{l}{\textsci}{\textlengthmark}{\r{\textObardotlessj}}{\textsci}{l}{\textopeno}{r}{\texttheta}} \dicPos{n}[2] \dicFlx{(‑s, ‑)}[5] \textbf{1.} \dicSynonym{aðalatriði} \dicDirectTranslationCS{klíčové slovo}  \textbf{2.} \dicFieldCat{poč.} \dicSynonym*{aðgangsorð} \dicDirectTranslationCS{(přístupové) heslo}
\dicEntry[lykja] \dicTerm{lykja} \dicIPA{{l}{\textsci}{\textlengthmark}{\r{\textObardotlessj}}{a}} \dicPos{v}[4] \dicFlx{(lyk, lukti, luktum, lykti, lukið)}[61] \dicDirectTranslationCS{obklopit, obklopovat, obstoupit, obstupovat} \dicExampleIS{Vatn lukti um hana.} \dicExampleCS{Voda ji obklopila.};  \dicIdiom{lykjast}{ \dicPhraseIS{lykjast}} \dicFlx{refl} \dicDirectTranslationCS{obklopit, obklopovat, obstoupit, obstupovat}
\dicEntry[lykkja] \dicTerm{lykkj|a} \dicIPA{{l}{\textsci}{h}{\r{\textObardotlessj}}{a}} \dicPos{f}[1] \dicFlx{(‑u, ‑ur)}[25] \textbf{1.} \dicSynonym{hanki} \dicDirectTranslationCS{smyčka, oko, klička (na provázku ap.)} \dicExampleIS{lykkja á bandi} \dicExampleCS{smyčka na provázku}  \textbf{2.} \dicDirectTranslationCS{oko} \dicIndirectTranslationCS{(v~pletení)}  \textbf{3.} \dicSynonym{getnaðarvörn} \dicDirectTranslationCS{spirála} \dicIndirectTranslationCS{(druh antikoncepce)}
\dicEntry[lykkjufall] \dicTerm{lykkju··|fall} \dicIPA{{l}{\textsci}{h}{\r{\textObardotlessj}}{\textscy}{f}{a}{\textsubring{d}}{\textsubring{l}}} \dicPos{n}[2] \dicFlx{(‑falls, ‑föll)}[8] \dicDirectTranslationCS{puštěné oko} \dicIndirectTranslationCS{(v~pletení)} \dicExampleIS{lykkjufall á sokkabuxurnar} \dicExampleCS{oko na punčoše}
\dicEntry[lyklaborð] \dicTerm{lykla··borð} \dicsymFrequent\  \dicIPA{{l}{\textsci}{h}{\r{g}}{l}{a}{\textsubring{b}}{\textopeno}{r}{\texttheta}} \dicPos{n}[2] \dicFlx{(‑s, ‑)}[5] \dicFieldCat{poč.} \dicDirectTranslationCS{klávesnice} \dicExampleIS{íslenskt lyklaborð} \dicExampleCS{islandská klávesnice}
\dicEntry[lyklakippa] \dicTerm{lykla··kipp|a} \dicIPA{{l}{\textsci}{h}{\r{g}}{l}{a}{c\smash{\textsuperscript{h}}}{\textsci}{h}{\textsubring{b}}{a}} \dicPos{f}[1] \dicFlx{(‑u, ‑ur)}[19] \dicDirectTranslationCS{svazek klíčů}
\dicEntry[lykt] \dicTerm{lykt\smash{\textsuperscript{1}}} \dicIPA{{l}{\textsci}{x}{\textsubring{d}}} \dicPos{f}[7] \dicFlx{(‑ar, ‑ir)}[1] \dicSynonym{endir} \dicDirectTranslationCS{konec, závěr};  \dicPhraseIS{að lyktum} \dicFlx{adv} \dicSynonym*{að lokum} \dicDirectTranslationCS{nakonec, na závěr};  \dicPhraseIS{leiða e‑ð til lykta} \dicSynonym*{ljúka e‑u} \dicDirectTranslationCS{dovést (co) do konce, dokončit (co)};  \dicPhraseIS{gera lyktir á e‑u} \dicSynonym*{ljúka við e‑ð} \dicDirectTranslationCS{dokončit\,/\addthin dokončovat (co)}
\dicEntry[lykt] \dicTerm{lykt\smash{\textsuperscript{2}}} \dicsymFrequent\  \dicIPA{{l}{\textsci}{x}{\textsubring{d}}} \dicPos{f}[7] \dicFlx{(‑ar, ‑ir)}[1] \dicSynonym{ilmur} \dicDirectTranslationCS{(zá)pach, vůně, aroma} \dicExampleIS{lykt af blómunum} \dicExampleCS{vůně květin}
\dicEntry[lykta] \dicTerm{lykt|a\smash{\textsuperscript{1}}} \dicIPA{{l}{\textsci}{x}{\textsubring{d}}{a}} \dicPos{v}[1] \dicFlx{(‑aði)}[5] \dicPhraseIS{e‑u lyktar} \dicFlx{impers} \dicDirectTranslationCS{(co) končí} \dicExampleIS{Leiknum lyktaði með jafntefli.} \dicExampleCS{Zápas skončil remízou.}
\dicEntry[lykta] \dicTerm{lykt|a\smash{\textsuperscript{2}}} \dicIPA{{l}{\textsci}{x}{\textsubring{d}}{a}} \dicPos{v}[1] \dicFlx{(‑aði)}[34] \textbf{1.} \dicSynonym{ilma} \dicDirectTranslationCS{páchnout, být cítit, vonět, vydávat vůni} \dicExampleIS{lykta vel} \dicExampleCS{vonět dobře};  \dicPhraseIS{lykta af e‑u} \dicDirectTranslationCS{páchnout (čím), být cítit (čím), vonět (čím)} \dicExampleIS{lykta af ilmvatni} \dicExampleCS{vonět parfémem}  \textbf{2.} \dicSynonym{þefa} \dicDirectTranslationCS{přivonět (si), přičichnout (si)};  \dicPhraseIS{lykta af e‑u} \dicSynonym{þefa} \dicDirectTranslationCS{přivonět (si) k~(čemu)} \dicExampleIS{lykta af blóminu} \dicExampleCS{přivonět si ke květině}
\dicEntry[lyktandi] \dicTerm{lykt··andi} \dicIPA{{l}{\textsci}{x}{\textsubring{d}}{a}{n}{\textsubring{d}}{\textsci}} \dicPos{adj}[13] \dicFlx{indecl}[1] \dicDirectTranslationCS{vonící, šířící vůni};  \dicPhraseIS{e‑að er illa lyktandi} \dicDirectTranslationCS{(co) smrdí\,/\addthin páchne}
\dicEntry[lykti] \dicTerm{lykti} \dicIPA{{l}{\textsci}{x}{\textsubring{d}}{\textsci}} \dicPos{v} \dicFlx{con pf sg 1 pers} \dicLink{lykja}
\dicEntry[lyktnæmur] \dicTerm{lykt··næmur} \dicIPA{{l}{\textsci}{x}{\textsubring{d}}{n}{a}{i}{m}{\textscy}{\textsubring{r}}} \dicPos{adj}[1]\dicFlx{}[-1] \dicDirectTranslationCS{mající dobrý čich}
\dicEntry[lymska] \dicTerm{lymsk|a} \dicIPA{{l}{\textsci}{m}{s}{\r{g}}{a}} \dicPos{f}[1] \dicFlx{(‑u)}[5] \dicSynonym*{lævísi} \dicDirectTranslationCS{záludnost, zákeřnost, proradnost}
\dicEntry[lymskur] \dicTerm{lymskur} \dicIPA{{l}{\textsci}{m}{s}{\r{g}}{\textscy}{\textsubring{r}}} \dicPos{adj}[1]\dicFlx{}[-1] \dicSynonym{lævís} \dicDirectTranslationCS{záludný, zákeřný, proradný}
\dicEntry[lynda] \dicTerm{lyn|da} \dicIPA{{l}{\textsci}{n}{\textsubring{d}}{a}} \dicPos{v}[2] \dicFlx{(‑ti, ‑t)}[44] \dicPhraseIS{e‑m lyndir við e‑n} \dicFlx{impers} \dicDirectTranslationCS{(kdo) si rozumí s~(kým), (kdo) má pochopení pro (koho)} \dicExampleIS{Mér lyndir ekki við hana.} \dicExampleCS{Nerozumím si s~ní.};  \dicPhraseIS{e‑jum lyndir} \dicFlx{impers} \dicDirectTranslationCS{(kdo) si rozumí, (kdo) má pochopení jeden pro druhého};  \dicPhraseIS{láta sér e‑ð lynda} \dicDirectTranslationCS{spokojit se s~(čím)}
\dicEntry[lyndi] \dicTerm{lyndi} \dicIPA{{l}{\textsci}{n}{\textsubring{d}}{\textsci}} \dicPos{n}[2] \dicFlx{(‑s)}[20] \dicSynonym{skap} \dicDirectTranslationCS{nálada, rozpoložení};  \dicPhraseIS{e‑að leikur í lyndi} \dicDirectTranslationCS{(co) klape, (co) hraje do karet};  \dicPhraseIS{e‑m leikur allt í lyndi} \dicFlx{impers} \dicDirectTranslationCS{(komu) vše vychází, (komu) vše hraje do noty}
\dicEntry[lyndiseinkunn] \dicTerm{lyndis··ein·kunn} \dicIPA{{l}{\textsci}{n}{\textsubring{d}}{\textsci}{s}{ei}{\r{\ng}}{\r{g}}{\textscy}{\textsubring{n}}} \dicPos{f}[7] \dicFlx{(‑ar, ‑ir)}[1] \dicSynonym*{skaphöfn} \dicDirectTranslationCS{charakter}
\dicEntry[lyng] \dicTerm{lyng} \dicIPA{{l}{i}{\ng}{\r{g}}} \dicPos{n}[2] \dicFlx{(‑s)}[2] \dicFieldCat{bot.} \dicDirectTranslationCS{vřesovec} \textit{(l.~{\textLA{Erica}})}  \dicsymPhoto\ 
\dicFigure{74636.jpg}{Lyng}{Lyng - Kesl Michael, Biolib, Copyright/CC-BY-NC}
\dicEntry[lyst] \dicTerm{lyst} \dicIPA{{l}{\textsci}{s}{\textsubring{d}}} \dicPos{f}[7] \dicFlx{(‑ar, ‑ir)}[1] \textbf{1.} \dicSynonym{matarlyst} \dicDirectTranslationCS{chuť (k~jídlu), apetit};  \dicPhraseIS{hafa lyst á e‑u} \dicDirectTranslationCS{mít chuť na (co)} \dicExampleIS{hafa góða lyst á matnum} \dicExampleCS{mít dobrý apetit}  \textbf{2.} \dicSynonym{löngun} \dicDirectTranslationCS{chuť, přání};  \dicPhraseIS{af hjartans lyst} \dicFlx{adv} \dicDirectTranslationCS{z~celého srdce}
\dicEntry[lystarlaus] \dicTerm{lystar··laus} \dicIPA{{l}{\textsci}{s}{\textsubring{d}}{a}{r}{l}{\oe i}{s}} \dicPos{adj}[5]\dicFlx{}[-1] \dicDirectTranslationCS{(jsoucí) bez apetitu, nemající chuť (na jídlo)}
\dicEntry[lystarleysi] \dicTerm{lystar··leysi} \dicIPA{{l}{\textsci}{s}{\textsubring{d}}{a}{r}{l}{ei}{s}{\textsci}} \dicPos{n}[2] \dicFlx{(‑s)}[20] \dicDirectTranslationCS{nechutenství}
\dicEntry[lystarstol] \dicTerm{lystar··stol} \dicIPA{{l}{\textsci}{s}{\textsubring{d}}{a}{\textsubring{r}}{s}{\textsubring{d}}{\textopeno}{\textsubring{l}}} \dicPos{n}[2] \dicFlx{(‑s)}[2] \dicFieldCat{med.} \dicDirectTranslationCS{(mentální) anorexie}
\dicEntry[lysti] \dicTerm{lysti} \dicIPA{{l}{\textsci}{s}{\textsubring{d}}{\textsci}} \dicPos{v} \dicFlx{con pf sg 1 pers} \dicLink{ljósta}
\dicEntry[lystibátur] \dicTerm{lysti··bát|ur} \dicIPA{{l}{\textsci}{s}{\textsubring{d}}{\textsci}{\textsubring{b}}{au}{\textsubring{d}}{\textscy}{\textsubring{r}}} \dicPos{m}[6] \dicFlx{(‑s, ‑ar)}[22] \dicFieldCat{nám.} \dicDirectTranslationCS{jachta}
\dicEntry[lystigarður] \dicTerm{lysti··garð|ur} \dicIPA{{l}{\textsci}{s}{\textsubring{d}}{\textsci}{\r{g}}{a}{r}{ð}{\textscy}{\textsubring{r}}} \dicPos{m}[6] \dicFlx{(‑s, ‑ar)}[5] \dicSynonym{skemmtigarður} \dicDirectTranslationCS{(zábavní) park}
\dicEntry[lystiskip] \dicTerm{lysti··skip} \dicIPA{{l}{\textsci}{s}{\textsubring{d}}{\textsci}{s}{\r{\textObardotlessj}}{\textsci}{\textsubring{b}}} \dicPos{n}[2] \dicFlx{(‑s, ‑)}[5] \dicDirectTranslationCS{jachta, výletní loď}
\dicEntry[lystugur] \dicTerm{lyst··ugur} \dicIPA{{l}{\textsci}{s}{\textsubring{d}}{\textscy}{\textbabygamma}{\textscy}{\textsubring{r}}} \dicPos{adj}[1]\dicFlx{}[-8] \textbf{1.} \dicSynonym{sólginn} \dicDirectTranslationCS{mající dobrý apetit} \dicExampleIS{vera lystugur á matinn} \dicExampleCS{mít dobrý apetit}  \textbf{2.} \dicSynonym*{sem vekur lyst} \dicDirectTranslationCS{chutný, vábný} \dicIndirectTranslationCS{(o~jídle)}
\dicEntry[lyti] \dicTerm{lyti} \dicIPA{{l}{\textsci}{\textlengthmark}{\textsubring{d}}{\textsci}} \dicPos{v} \dicFlx{con pf sg 1 pers} \dicLink{lúta}
\dicEntry[lýðfrjáls] \dicTerm{lýð··frjáls} \dicIPA{{l}{i}{ð}{f}{r}{j}{au}{l}{s}} \dicPos{adj}[5]\dicFlx{}[-1] \dicDirectTranslationCS{demokratický, svobodný} \dicExampleIS{lýðfrjálst land} \dicExampleCS{demokratická země}
\dicEntry[lýðháskóli] \dicTerm{lýð··há·skól|i} \dicIPA{{l}{i}{\textlengthmark}{\texttheta}{h}{au}{s}{\r{g}}{ou}{l}{\textsci}} \dicPos{m}[1] \dicFlx{(‑a, ‑ar)}[1] \dicFieldCat{škol.} \dicDirectTranslationCS{veřejná vysoká škola}
\dicEntry[lýðhylli] \dicTerm{lýð··hyll|i} \dicIPA{{l}{i}{\textlengthmark}{\texttheta}{h}{\textsci}{\textsubring{d}}{l}{\textsci}} \dicPos{f}[3] \dicFlx{(‑i)}[3] \dicDirectTranslationCS{přízeň lidu, popularita}
\dicEntry[lýði] \dicTerm{lýði} \dicIPA{{l}{i}{\textlengthmark}{ð}{\textsci}} \dicPos{v} \dicFlx{con pf sg 1 pers} \dicLink{lýja}
\dicEntry[lýðréttindi] \dicTerm{lýð··rétt·indi} \dicIPA{{l}{i}{ð}{r}{j}{\textepsilon}{h}{\textsubring{d}}{\textsci}{n}{\textsubring{d}}{\textsci}} \dicPos{n}[2] \dicFlx{pl}[19] \dicDirectTranslationCS{občanská práva}
\dicEntry[lýðræði] \dicTerm{lýð··ræði} \dicIPA{{l}{i}{ð}{r}{a}{i}{ð}{\textsci}} \dicPos{n}[2] \dicFlx{(‑s)}[20] \dicDirectTranslationCS{demokracie}
\dicEntry[lýðræðislegur] \dicTerm{lýð·ræðis··legur} \dicIPA{{l}{i}{ð}{r}{a}{i}{ð}{\textsci}{s}{l}{\textepsilon}{\textbabygamma}{\textscy}{\textsubring{r}}} \dicPos{adj}[1]\dicFlx{}[-8] \dicDirectTranslationCS{demokratický}
\dicEntry[lýðræðissinnaður] \dicTerm{lýð·ræðis··sinn·|aður} \dicIPA{{l}\-{i}\-{ð}\-{r}\-{a}\-{i}\-{ð}\-{\textsci}\-{s}\-{\textsci}\-{n}\-{a}\-{ð}\-{\textscy}\-{\textsubring{r}}\-} \dicPos{adj}[3] \dicFlx{(f ‑uð)}[3] \dicDirectTranslationCS{(pro)demokratický}
\dicEntry[lýðræðissinni] \dicTerm{lýð·ræðis··sinn|i} \dicIPA{{l}{i}{ð}{r}{a}{i}{ð}{\textsci}{s}{\textsci}{n}{\textsci}} \dicPos{m}[1] \dicFlx{(‑a, ‑ar)}[1] \dicDirectTranslationCS{demokrat(ka)}
\dicEntry[lýðskrum] \dicTerm{lýð··skrum} \dicIPA{{l}{i}{ð}{s}{\r{g}}{r}{\textscy}{\textsubring{m}}} \dicPos{n}[2] \dicFlx{(‑s)}[2] \dicDirectTranslationCS{demagogie, populismus}
\dicEntry[lýðskrumari] \dicTerm{lýð··skrum·ar|i} \dicIPA{{l}{i}{ð}{s}{\r{g}}{r}{\textscy}{m}{a}{r}{\textsci}} \dicPos{m}[1] \dicFlx{(‑a, ‑ar)}[13] \dicDirectTranslationCS{demagog, demagožka, populista, populistka}
\dicEntry[lýður] \dicTerm{lýð|ur} \dicIPA{{l}{i}{\textlengthmark}{ð}{\textscy}{\textsubring{r}}} \dicPos{m}[9] \dicFlx{(‑s, ‑ir)}[8] \textbf{1.} \dicSynonym{almenningur} \dicDirectTranslationCS{lid(é), občané}  \textbf{2.} \dicSynonym{skríll} \dicDirectTranslationCS{lůza, holota};  \dicPhraseIS{e‑að er enn við lýði} \dicDirectTranslationCS{(co) stále existuje, (co) se stále vyskytuje (zvyk ap.)}
\dicEntry[lýðveldi] \dicTerm{lýð··veldi} \dicIPA{{l}{i}{ð}{v}{\textepsilon}{l}{\textsubring{d}}{\textsci}} \dicPos{n}[2] \dicFlx{(‑s, ‑)}[14] \dicDirectTranslationCS{republika}
\dicEntry[lýg] \dicTerm{lýg} \dicIPA{{l}{i}{\textlengthmark}{x}} \dicPos{v} \dicFlx{ind praes sg 1 pers} \dicLink{ljúga}
\dicEntry[lýi] \dicTerm{lýi} \dicIPA{{l}{i}{j}{\textsci}} \dicPos{v} \dicFlx{ind praes sg 1 pers} \dicLink{lýja}
\dicEntry[lýja] \dicTerm{lýja} \dicIPA{{l}{i}{j}{\textlengthmark}{a}} \dicPos{v}[4] \dicFlx{(lýi, lúði, lúðum, lýði, lúið\,/\addthin lúð)}[54] \dicFlx{acc} \textbf{1.} \dicDirectTranslationCS{kovat, kout} \dicExampleIS{lýja járn} \dicExampleCS{kovat železo}  \textbf{2.} \dicSynonym{þreyta\smash{\textsuperscript{2}}} \dicDirectTranslationCS{unavit, utahat} \dicExampleIS{Þetta lýir mann.} \dicExampleCS{To člověka unaví.};  \dicIdiom{lýjast}{ \dicPhraseIS{lýjast}} \dicFlx{refl} \dicSynonym*{þreytast} \dicDirectTranslationCS{vyčerpat se, unavit se} \dicExampleIS{Strákurinn er farinn að lýjast.} \dicExampleCS{Kluk se začíná unavovat.}
\dicEntry[lýjandi] \dicTerm{lýj··andi} \dicIPA{{l}{i}{j}{\textlengthmark}{a}{n}{\textsubring{d}}{\textsci}} \dicPos{adj}[13] \dicFlx{indecl}[1] \dicDirectTranslationCS{vyčerpávající, vysilující} \dicExampleIS{lýjandi starf} \dicExampleCS{vyčerpávající práce}
\dicEntry[lýk] \dicTerm{lýk} \dicIPA{{l}{i}{\textlengthmark}{\r{g}}} \dicPos{v} \dicFlx{ind praes sg 1 pers} \dicLink{ljúka}
\dicEntry[lýrik] \dicTerm{lýrik} \dicIPA{{l}{i}{\textlengthmark}{r}{\textsci}{\r{g}}} \dicPos{f}[10] \dicFlx{(‑ur)}[2] \dicFieldCat{lit.} \dicSynonym{ljóðræna} \dicDirectTranslationCS{lyrika}
\dicEntry[lýrískur] \dicTerm{lýrískur} \dicIPA{{l}{i}{\textlengthmark}{r}{i}{s}{\r{g}}{\textscy}{\textsubring{r}}} \dicPos{adj}[1]\dicFlx{}[-1] \dicFieldCat{lit.} \dicSynonym{ljóðrænn} \dicDirectTranslationCS{lyrický}
\dicEntry[lýs] \dicTerm{lýs} \dicIPA{{l}{i}{\textlengthmark}{s}} \dicPos{f} \dicFlx{pl nom} \dicLink{lús}
\dicEntry[lýsa] \dicTerm{lýs|a\smash{\textsuperscript{1}}} \dicIPA{{l}{i}{\textlengthmark}{s}{a}} \dicPos{f}[1] \dicFlx{(‑u, ‑ur)}[7] \textbf{1.} \dicSynonym{ljómi} \dicDirectTranslationCS{světlo, paprsek} \dicExampleIS{fyrsta lýsa morgunsins} \dicExampleCS{první ranní paprsek}  \textbf{2.} \dicFieldCat{zool.} \dicDirectTranslationCS{treska bezvousá} \textit{(l.~{\textLA{Merlangius merlangus}})}  \dicsymPhoto\ 
\dicFigure{ds_image_lysa_1_2.jpg}{Lýsa}{Lýsa - Georges Jansoone (JoJan), CC BY 3.0}
\dicEntry[lýsa] \dicTerm{lýs|a\smash{\textsuperscript{2}}} \dicsymFrequent\  \dicIPA{{l}{i}{\textlengthmark}{s}{a}} \dicPos{v}[2] \dicFlx{(‑ti, ‑t)}[67] \dicFlx{dat} \textbf{1.} \dicSynonym*{bera birtu} \dicDirectTranslationCS{osvětlit, osvětlovat, osvítit, ozářit, ozařovat} \dicExampleIS{Tunglið lýsir mér.} \dicExampleCS{Měsíc mě osvětluje.};  \dicPhraseIS{það lýsir (af degi)} \dicFlx{impers} \dicSynonym*{það birtir} \dicDirectTranslationCS{rozednívá se}  \textbf{2.} \dicSynonym*{lýsa mynd} \dicDirectTranslationCS{vyvolat, vyvolávat (film ap.)} \dicExampleIS{lýsa myndina} \dicExampleCS{vyvolat film}  \textbf{3.} \dicSynonym{sýna} \dicDirectTranslationCS{projevit, projevovat, vyjádřit, vyjadřovat}  \textbf{4.} \dicSynonym*{gefa lýsingu} \dicDirectTranslationCS{popsat, popisovat, (vy)líčit, vykreslit, vykreslovat} \dicExampleIS{lýsa reynslu sinni} \dicExampleCS{popsat svou zkušenost}  \textbf{5.} \dicSynonym{tilkynna} \dicDirectTranslationCS{prohlásit, prohlašovat, deklarovat};  \dicIdiom{lýsa}[eftir]{ \dicPhraseIS{lýsa eftir e‑m}} \dicDirectTranslationCS{vyhlásit pátrání po (kom)};  \dicIdiom{lýsa}[upp]{ \dicPhraseIS{lýsa upp e‑ð}} \dicDirectTranslationCS{osvětlit (co) (místnost ap.)};  \dicIdiom{lýsa}[yfir]{ \dicPhraseIS{lýsa yfir e‑u}} \dicDirectTranslationCS{ohlásit (co), oznámit (co)}
\dicEntry[lýsi] \dicTerm{lýsi} \dicIPA{{l}{i}{\textlengthmark}{s}{\textsci}} \dicPos{n}[2] \dicFlx{(‑s, ‑)}[14] \dicDirectTranslationCS{rybí tuk} \dicIndirectTranslationCS{(medicína)}
\dicEntry[lýsing] \dicTerm{lýs··ing} \dicsymFrequent\  \dicIPA{{l}{i}{\textlengthmark}{s}{i}{\ng}{\r{g}}} \dicPos{f}[4] \dicFlx{(‑ar, ‑ar)}[5] \textbf{1.} \dicDirectTranslationCS{osvětlení, nasvícení, iluminace} \dicExampleIS{lýsing í eldhúsinu} \dicExampleCS{osvětlení v~kuchyni}  \textbf{2.} \dicSynonym{frásögn} \dicDirectTranslationCS{popis, vylíčení, vykreslení} \dicExampleIS{lýsing á ferðinni} \dicExampleCS{vylíčení cesty}  \textbf{3.} \dicDirectTranslationCS{osvit, expozice}  \textbf{4.} \dicSynonym{skreyting} \dicDirectTranslationCS{ilustrace, ilustrování (rukopisu, knihy ap.)}
\dicEntry[lýsingarháttur] \dicTerm{lýsingar··|háttur} \dicIPA{{l}{i}{\textlengthmark}{s}{i}{\ng}{\r{g}}{a}{\textsubring{r}}{h}{au}{h}{\textsubring{d}}{\textscy}{\textsubring{r}}} \dicPos{m}[12] \dicFlx{(‑háttar, ‑hættir)}[7] \dicFieldCat{jaz.} \dicDirectTranslationCS{příčestí};  \dicPhraseIS{lýsingarháttur nútíðar} \dicFieldCat{jaz.} \dicDirectTranslationCS{příčestí přítomné};  \dicPhraseIS{lýsingarháttur þátíðar} \dicFieldCat{jaz.} \dicDirectTranslationCS{příčestí minulé}
\dicEntry[lýsingarorð] \dicTerm{lýsingar··orð} \dicIPA{{l}{i}{\textlengthmark}{s}{i}{\ng}{\r{g}}{a}{r}{\textopeno}{r}{\texttheta}} \dicPos{n}[2] \dicFlx{(‑s, ‑)}[5] \dicFieldCat{jaz.} \dicDirectTranslationCS{přídavné jméno, adjektivum} \textit{(l.~{\textLA{nomen adiectivum}})}
\dicEntry[lýst] \dicTerm{lýst} \dicIPA{{l}{i}{s}{\textsubring{d}}} \dicPos{v} \dicFlx{ind praes sg 1 pers} \dicLink{ljósta}
\dicEntry[lýt] \dicTerm{lýt} \dicIPA{{l}{i}{\textlengthmark}{\textsubring{d}}} \dicPos{v} \dicFlx{ind praes sg 1 pers} \dicLink{lúta}
\dicEntry[lýtalaus] \textls[15]{\dicTerm{lýta··laus} \dicIPA{{l}{i}{\textlengthmark}{\textsubring{d}}{a}{l}{\oe i}{s}} \dicPos{adj}[5]\dicFlx{}[-1] \dicSynonym{gallalaus} \dicDirectTranslationCS{bezchybný, (jsoucí) bez poskvrnky}}
\dicEntry[lýtalækningar] \dicTerm{lýta··lækn·ingar} \dicIPA{{l}{i}{\textlengthmark}{\textsubring{d}}{a}{l}{a}{i}{h}{\r{g}}{n}{i}{\ng}{\r{g}}{a}{\textsubring{r}}} \dicPos{f}[4] \dicFlx{pl}[6] \dicFieldCat{med.} \dicDirectTranslationCS{plastická chirurgie}
\dicEntry[lýtalæknir] \dicTerm{lýta··lækn|ir} \dicIPA{{l}{i}{\textlengthmark}{\textsubring{d}}{a}{l}{a}{i}{h}{\r{g}}{n}{\textsci}{\textsubring{r}}} \dicPos{m}[7] \dicFlx{(‑is, ‑ar)}[1] \dicDirectTranslationCS{plastický chirurg, plastická chirurgyně\,/\addthin chiruržka}
\dicEntry[lýti] \dicTerm{lýti} \dicIPA{{l}{i}{\textlengthmark}{\textsubring{d}}{\textsci}} \dicPos{n}[2] \dicFlx{(‑s, ‑)}[14] \textbf{1.} \dicSynonym{ljóður} \dicDirectTranslationCS{poskvrna, kaz, vada} \dicExampleIS{lýti á bókinni} \dicExampleCS{vada na knížce}  \textbf{2.} \dicSynonym{ljótleiki} \dicDirectTranslationCS{šerednost, ošklivost (vzhled člověka ap.)}
\dicEntry[læða] \dicTerm{læð|a\smash{\textsuperscript{1}}} \dicIPA{{l}{a}{i}{\textlengthmark}{ð}{a}} \dicPos{f}[1] \dicFlx{(‑u, ‑ur)}[19] \dicDirectTranslationCS{kočka} \dicIndirectTranslationCS{(samice)}
\dicEntry[læða] \dicTerm{læ|ða\smash{\textsuperscript{2}}} \dicsymFrequent\  \dicIPA{{l}{a}{i}{\textlengthmark}{ð}{a}} \dicPos{v}[2] \dicFlx{(‑ddi, ‑tt)}[172] \dicFlx{dat} \dicPhraseIS{læða e‑u að e‑m} \dicDirectTranslationCS{propašovat (komu co), vstrčit (komu co) skrytě} \dicExampleIS{læða þessari frétt að henni} \dicExampleCS{nenápadně jí sdělit zprávu};  \dicIdiom{læðast}{ \dicPhraseIS{læðast}} \dicFlx{refl} \dicDirectTranslationCS{(v)plížit se, plazit se, krást se} \dicExampleIS{læðast burt} \dicExampleCS{odplížit se pryč}
\dicEntry[læðupokast] \dicTerm{læðu··pok|ast} \dicIPA{{l}{a}{i}{\textlengthmark}{ð}{\textscy}{p\smash{\textsuperscript{h}}}{\textopeno}{\r{g}}{a}{s}{\textsubring{d}}} \dicPos{v}[1] \dicFlx{(‑aðist)}[91] \dicFlx{refl} \dicDirectTranslationCS{(při)krást se, plížit se}
\dicEntry[lægð] \dicTerm{lægð} \dicIPA{{l}{a}{i}{\textbabygamma}{\texttheta}} \dicPos{f}[7] \dicFlx{(‑ar, ‑ir)}[1] \textbf{1.} \dicSynonym{dæld} \dicDirectTranslationCS{prohlubeň, proláklina} \dicExampleIS{Það myndast djúp lægð í götuna.} \dicExampleCS{Na ulici se tvoří hluboká prohlubeň.}  \textbf{2.} \dicFieldCat{meteo.} \dicSynonym{lágþrýstisvæði} \dicDirectTranslationCS{tlaková níže} \dicAntonym{hægð}  \textbf{3.} \dicSynonym{deyfð} \dicDirectTranslationCS{útlum};  \dicPhraseIS{vera í lægð} \dicDirectTranslationCS{být v~útlumu}
\dicEntry[lægi] \dicTerm{lægi\smash{\textsuperscript{1}}} \dicIPA{{l}{a}{i}{j}{\textlengthmark}{\textsci}} \dicPos{n}[2] \dicFlx{(‑s, ‑)}[16] \dicSynonym{skipalægi} \dicDirectTranslationCS{kotviště}
\dicEntry[lægi] \dicTerm{lægi\smash{\textsuperscript{2}}} \dicIPA{{l}{a}{i}{j}{\textlengthmark}{\textsci}} \dicPos{v} \dicFlx{con pf sg 1 pers} \dicLink{liggja}
\dicEntry[lægja] \dicTerm{læg|ja} \dicIPA{{l}{a}{i}{j}{\textlengthmark}{a}} \dicPos{v}[2] \dicFlx{(‑ði, ‑t)}[90] \dicFlx{acc} \dicSynonym{auðmýkja} \dicDirectTranslationCS{zklidnit, zklidňovat, zmírnit, zmírňovat} \dicExampleIS{lægja æsing} \dicExampleCS{zklidnit vztek};  \dicPhraseIS{lægja rostann í e‑m} \dicDirectTranslationCS{zklidnit (koho)};  \dicPhraseIS{e‑ð lægir} \dicFlx{impers} {\textbf{a.}} \dicDirectTranslationCS{(co) klesá, (co) zapadá (slunce ap.)} \dicExampleIS{Sólina lægir.} \dicExampleCS{Slunce klesá.};  {\textbf{b.}} \dicDirectTranslationCS{(co) ustává (o~větru ap.), (co) se uklidňuje} \dicExampleIS{Veðrið lægir.} \dicExampleCS{Počasí se uklidňuje.}
\dicEntry[lægra] \dicTerm{lægra} \dicIPA{{l}{a}{i}{\textbabygamma}{r}{a}} \dicPos{adv} \dicFlx{comp} \dicLink{lágt}
\dicEntry[lægri] \dicTerm{lægri} \dicIPA{{l}{a}{i}{\textbabygamma}{r}{\textsci}} \dicPos{adj} \dicFlx{comp m} \dicLink{lágur}
\dicEntry[lægst] \dicTerm{lægst} \dicIPA{{l}{a}{i}{x}{s}{\textsubring{d}}} \dicPos{adv} \dicFlx{sup} \dicLink{lágt}
\dicEntry[lægstur] \dicTerm{lægstur} \dicIPA{{l}{a}{i}{x}{s}{\textsubring{d}}{\textscy}{\textsubring{r}}} \dicPos{adj} \dicFlx{m sg nom sup} \dicLink{lágur}
\dicEntry[læki] \dicTerm{læki} \dicIPA{{l}{a}{i}{\textlengthmark}{\r{\textObardotlessj}}{\textsci}} \dicPos{v} \dicFlx{con pf sg 1 pers} \dicLink{leka}
\dicEntry[lækka] \dicTerm{lækk|a} \dicsymFrequent\  \dicIPA{{l}{a}{i}{h}{\r{g}}{a}} \dicPos{v}[1] \dicFlx{(‑aði)}[48] \dicFlx{acc} \textbf{1.} \dicSynonym{minnka} \dicDirectTranslationCS{snížit, snižovat, (z)redukovat} \dicExampleIS{lækka skatta} \dicExampleCS{snížit daně};  \dicPhraseIS{lækka í e‑u} \dicDirectTranslationCS{zeslabit (co), ztišit (co), ztlumit (co), zmírnit (co)} \dicExampleIS{lækka í útvarpinu} \dicExampleCS{ztišit rádio};  \dicPhraseIS{lækka röddina\,/\addthin róminn} \dicDirectTranslationCS{snížit hlas}  \textbf{2.} \dicDirectTranslationCS{klesnout, klesat, sestoupit, sestupovat, snášet se} \dicExampleIS{lækka flugið} \dicExampleCS{klesat s~letadlem};  \dicPhraseIS{e‑að lækkar} \dicDirectTranslationCS{(co) klesá, (co) ustupuje (voda ap.)};  \dicPhraseIS{það lækkar í e‑u} \dicFlx{impers} \dicDirectTranslationCS{(co) klesá, (co) opadává} \dicExampleIS{Það lækkar í ánni.} \dicExampleCS{Voda v~řece opadává.}
\dicEntry[lækkun] \dicTerm{lækk|un} \dicIPA{{l}{a}{i}{h}{\r{g}}{\textscy}{\textsubring{n}}} \dicPos{f}[7] \dicFlx{(‑unar)}[9] \textbf{1.} \dicSynonym*{það að lækka} \dicDirectTranslationCS{snížení, snižování, redukce, redukování} \dicExampleIS{lækkun skatta} \dicExampleCS{snížení daní}  \textbf{2.} \dicSynonym*{það að e‑að lækkar} \dicDirectTranslationCS{klesání, opadání, ubývání} \dicExampleIS{lækkun á yfirborði sjávar} \dicExampleCS{klesání hladiny moře}
\dicEntry[lækna] \dicTerm{lækn|a} \dicsymFrequent\  \dicIPA{{l}{a}{i}{h}{\r{g}}{n}{a}} \dicPos{v}[1] \dicFlx{(‑aði)}[1] \dicFlx{acc} \dicDirectTranslationCS{(vy)léčit};  \dicPhraseIS{lækna e‑n af e‑u} \dicDirectTranslationCS{léčit (koho) z~(čeho)} \dicExampleIS{lækna e‑n af berklum} \dicExampleCS{léčit (koho) z~tuberkulózy}
\dicEntry[læknadeild] \dicTerm{lækna··deild} \dicIPA{{l}{a}{i}{h}{\r{g}}{n}{a}{\textsubring{d}}{ei}{l}{\textsubring{d}}} \dicPos{f}[7] \dicFlx{(‑ar, ‑ir)}[1] \dicFieldCat{škol.} \dicDirectTranslationCS{lékařská fakulta}
\dicEntry[læknandi] \dicTerm{lækn··andi} \dicIPA{{l}{a}{i}{h}{\r{g}}{n}{a}{n}{\textsubring{d}}{\textsci}} \dicPos{adj}[13] \dicFlx{indecl}[1] \dicDirectTranslationCS{léčivý, léčebný}
\dicEntry[læknanemi] \dicTerm{lækna··nem|i} \dicIPA{{l}{a}{i}{h}{\r{g}}{n}{a}{n}{\textepsilon}{m}{\textsci}} \dicPos{m}[1] \dicFlx{(‑a, ‑ar)}[1] \dicDirectTranslationCS{student(ka) medicíny\,/\addthin lékařství}
\dicEntry[læknanlegur] \dicTerm{læknan··legur} \dicIPA{{l}{a}{i}{h}{\r{g}}{n}{a}{n}{l}{\textepsilon}{\textbabygamma}{\textscy}{\textsubring{r}}} \dicPos{adj}[1]\dicFlx{}[-8] \dicDirectTranslationCS{(vy)léčitelný}
\dicEntry[læknastofa] \dicTerm{lækna··stof|a} \dicIPA{{l}{a}{i}{h}{\r{g}}{n}{a}{s}{\textsubring{d}}{\textopeno}{v}{a}} \dicPos{f}[1] \dicFlx{(‑u, ‑ur)}[7] \dicDirectTranslationCS{ordinace}
\dicEntry[lækning] \dicTerm{lækn··ing} \dicsymFrequent\  \dicIPA{{l}{a}{i}{h}{\r{g}}{n}{i}{\ng}{\r{g}}} \dicPos{f}[4] \dicFlx{(‑ar, ‑ar)}[5] \dicDirectTranslationCS{léčení, léčba};  \dicPhraseIS{lækning við e‑u} \dicDirectTranslationCS{léčba (čeho)} \dicExampleIS{lækning við sjúkdómi} \dicExampleCS{léčba nemoci}
\dicEntry[lækningastofa] \dicTerm{lækninga··stof|a} \dicIPA{{l}{a}{i}{h}{\r{g}}{n}{i}{\ng}{\r{g}}{a}{s}{\textsubring{d}}{\textopeno}{v}{a}} \dicPos{f}[1] \dicFlx{(‑u, ‑ur)}[7] \dicDirectTranslationCS{ordinace}
\dicEntry[læknir] \dicTerm{lækn|ir} \dicsymFrequent\  \dicIPA{{l}{a}{i}{h}{\r{g}}{n}{\textsci}{\textsubring{r}}} \dicPos{m}[7] \dicFlx{(‑is, ‑ar)}[1] \dicDirectTranslationCS{doktor(ka), lékař(ka)} \dicExampleIS{panta tíma hjá lækni} \dicExampleCS{objednat se k~lékaři}
\dicEntry[læknisfræði] \dicTerm{læknis··fræð|i} \dicIPA{{l}{a}{i}{h}{\r{g}}{n}{\textsci}{s}{f}{r}{a}{i}{ð}{\textsci}} \dicPos{f}[3] \dicFlx{(‑i)}[3] \dicDirectTranslationCS{lékařství, medicína} \dicIndirectTranslationCS{(nauka)}
\dicEntry[læknisskoðun] \dicTerm{læknis··skoð|un} \dicIPA{{l}{a}{i}{h}{\r{g}}{n}{\textsci}{s}{\r{g}}{\textopeno}{ð}{\textscy}{\textsubring{n}}} \dicPos{f}[7] \dicFlx{(‑unar)}[9] \dicDirectTranslationCS{zdravotní\,/\addthin lékařská prohlídka}
\dicEntry[læknisvottorð] \dicTerm{læknis··vott·orð} \dicIPA{{l}{a}{i}{h}{\r{g}}{n}{\textsci}{s}{v}{\textopeno}{h}{\textsubring{d}}{\textopeno}{r}{\texttheta}} \dicPos{n}[2] \dicFlx{(‑s, ‑)}[5] \dicDirectTranslationCS{lékařská zpráva}
\dicEntry[læknisþjónusta] \dicTerm{læknis··þjón·ust|a} \dicIPA{{l}{a}{i}{h}{\r{g}}{n}{\textsci}{s}{\texttheta}{j}{ou}{n}{\textscy}{s}{\textsubring{d}}{a}} \dicPos{f}[1] \dicFlx{(‑u)}[5] \dicDirectTranslationCS{zdravot\-ní\,/\addthin lékařská služba}
\dicEntry[lækur] \dicTerm{læk|ur} \dicsymFrequent\  \dicIPA{{l}{a}{i}{\textlengthmark}{\r{g}}{\textscy}{\textsubring{r}}} \dicPos{m}[10] \dicFlx{(‑jar, ‑ir)}[23] \textbf{1.} \dicDirectTranslationCS{potok, říčka} \dicExampleIS{ár og lækir} \dicExampleCS{řeky a~potoky}  \textbf{2.} \dicFieldCat{anat.} \dicSynonym{naflastrengur} \dicDirectTranslationCS{pupeční šňůra}
\dicEntry[læmingi] \dicTerm{læm··ing|i} \dicIPA{{l}{a}{i}{\textlengthmark}{m}{i}{\textltailn}{\r{\textObardotlessj}}{\textsci}} \dicPos{m}[1] \dicFlx{(‑ja, ‑jar)}[14] \dicFieldCat{zool.} \dicDirectTranslationCS{lumík norský} \textit{(l.~{\textLA{Lemmus lemmus}})}  \dicsymPhoto\ 
\dicFigure{ds_image_laemingi_0_1.jpg}{Læmingi}{Læmingi - Argus fin, PD}
\dicEntry[læra] \dicTerm{lær|a} \dicsymFrequent\  \dicIPA{{l}{a}{i}{\textlengthmark}{r}{a}} \dicPos{v}[2] \dicFlx{(‑ði, ‑t)}[103] \dicFlx{acc} \dicDirectTranslationCS{(na)učit se, studovat} \dicExampleIS{læra íslensku} \dicExampleCS{učit se islandštinu};  \dicPhraseIS{læra e‑ð utanbókar} \dicDirectTranslationCS{učit se (co) zpaměti, memorovat (co)} \dicExampleIS{læra utanbókar ljóð} \dicExampleCS{učit se básničku nazpaměť};  \dicIdiom{læra}[af]{ \dicPhraseIS{læra af e‑m}} \dicDirectTranslationCS{učit se od (koho)};  \dicIdiom{læra}[á]{ \dicPhraseIS{læra á e‑ð}} \dicDirectTranslationCS{učit se na (co) (na hudební nástroj ap.)} \dicExampleIS{Ég lærði á píanó.} \dicExampleCS{Učil jsem se na piano.};  \dicIdiom{læra}[fyrir]{ \dicPhraseIS{læra fyrir e‑ð}} \dicDirectTranslationCS{učit se na (co)} \dicExampleIS{læra fyrir próf} \dicExampleCS{učit se na zkoušku};  \dicIdiom{læra}[undir]{ \dicPhraseIS{læra undir e‑ð}} \dicDirectTranslationCS{učit se na (co) (na zkoušku ap.)} \dicExampleIS{læra undir skóla} \dicExampleCS{učit se do školy};  \dicIdiom{læra}[utan að]{ \dicPhraseIS{læra utan að}} \dicDirectTranslationCS{učit se zpaměti, memorovat};  \dicIdiom{lærast}{ \dicPhraseIS{e‑að lærist vel}} \dicFlx{refl} \dicDirectTranslationCS{(co) se učí dobře}; { \dicPhraseIS{e‑m lærist að (gera e‑ð)}} \dicFlx{refl impers} \dicDirectTranslationCS{(kdo) se učí (dělat (co)), (kdo) si osvojuje (dělat (co))};  \dicProverb\  \dicPhraseIS{Svo lengi lærir sem lifir.} \dicLangCat{přís.} \dicDirectTranslationCS{Člověk se učí po celý život.}
\dicEntry[lærbein] \dicTerm{lær··bein} \dicIPA{{l}{a}{i}{r}{\textsubring{b}}{ei}{\textsubring{n}}} \dicPos{n}[2] \dicFlx{(‑s, ‑)}[5] \dicFieldCat{anat.} \dicDirectTranslationCS{stehenní kost}
\dicEntry[lærdómur] \dicTerm{lær··dóm|ur} \dicIPA{{l}{a}{i}{r}{\textsubring{d}}{ou}{m}{\textscy}{\textsubring{r}}} \dicPos{m}[6] \dicFlx{(‑s, ‑ar)}[10] \dicSynonym{þekking} \dicDirectTranslationCS{učení, učenost, vzdělání};  \dicPhraseIS{draga lærdóma af e‑u} \dicDirectTranslationCS{poučit se z~(čeho)}
\dicEntry[lærður] \dicTerm{lærður} \dicsymFrequent\  \dicIPA{{l}{a}{i}{r}{ð}{\textscy}{\textsubring{r}}} \dicPos{adj}[2]\dicFlx{}[-1] \dicDirectTranslationCS{učený, vzdělaný} \dicExampleIS{mjög lærður maður} \dicExampleCS{velmi vzdělaný člověk}
\dicEntry[læri] \dicTerm{læri} \dicsymFrequent\  \dicIPA{{l}{a}{i}{\textlengthmark}{r}{\textsci}} \dicPos{n}[2] \dicFlx{(‑s, ‑)}[14] \textbf{1.} \dicFieldCat{anat.} \dicDirectTranslationCS{stehno}  \textbf{2.} \dicDirectTranslationCS{kýta} \dicExampleIS{læri af lambi} \dicExampleCS{jehněčí kýta}  \textbf{3.} \dicSynonym{nám} \dicDirectTranslationCS{učení} \dicIndirectTranslationCS{(učňovská léta)};  \dicPhraseIS{vera í læri hjá e‑m} \dicDirectTranslationCS{být u~(koho) v~učení}
\dicEntry[lærifaðir] \dicTerm{læri··|faðir} \dicIPA{{l}{a}{i}{\textlengthmark}{r}{\textsci}{f}{a}{ð}{\textsci}{\textsubring{r}}} \dicPos{m}[13] \dicFlx{(‑föður, ‑feður)}[3] \dicSynonym{kennari} \dicDirectTranslationCS{duchovní otec, mistr, guru}
\dicEntry[lærisveinn] \dicTerm{læri··svein|n} \dicIPA{{l}{a}{i}{\textlengthmark}{r}{\textsci}{s}{v}{ei}{\textsubring{d}}{\textsubring{n}}} \dicPos{m}[6] \dicFlx{(‑s, ‑ar)}[42] \dicSynonym{nemandi} \dicDirectTranslationCS{učeň(ka), učnice, učedník, učednice, žák(yně), žačka} \dicExampleIS{lærisveinar Jesú} \dicExampleCS{učedníci Ježíše}
\dicEntry[lærleggur] \dicTerm{lær··legg|ur} \dicIPA{{l}{a}{i}{r}{l}{\textepsilon}{\r{g}}{\textscy}{\textsubring{r}}} \dicPos{m}[9] \dicFlx{(‑jar\,/\addthin ‑s, ‑ir)}[26] \dicFieldCat{anat.} \dicDirectTranslationCS{stehenní kost}
\dicEntry[lærlingur] \dicTerm{lær··ling|ur} \dicIPA{{l}{a}{i}{r}{l}{i}{\ng}{\r{g}}{\textscy}{\textsubring{r}}} \dicPos{m}[6] \dicFlx{(‑s, ‑ar)}[8] \dicSynonym{nemandi} \dicDirectTranslationCS{učedník, učednice, učeň(ka), učnice}
\dicEntry[læs] \dicTerm{læs} \dicIPA{{l}{a}{i}{\textlengthmark}{s}} \dicPos{adj}[5]\dicFlx{}[-1] \dicDirectTranslationCS{gramotný, umějící číst}
\dicEntry[læsa] \dicTerm{læs|a} \dicsymFrequent\  \dicIPA{{l}{a}{i}{\textlengthmark}{s}{a}} \dicPos{v}[2] \dicFlx{(‑ti, ‑t)}[63] \dicFlx{dat\,/\addthin acc} \textbf{1.} \dicFlx{dat} \dicSynonym*{loka með lás} \dicDirectTranslationCS{(u)zamknout, (u)zamykat, (u)zavřít, (u)zavírat} \dicExampleIS{Dyrunum var læst.} \dicExampleCS{Dveře byly zamčeny.};  \dicPhraseIS{læsa að sér} \dicDirectTranslationCS{zamknout (za sebou) dveře};  \dicPhraseIS{læsa á eftir sér} \dicDirectTranslationCS{zavřít za sebou (dveře)};  \dicPhraseIS{læsa e‑n inni} \dicDirectTranslationCS{zavřít (koho), dát (koho) pod zámek (do vězení ap.)}  \textbf{2.} \dicFlx{dat} \dicSynonym{innsigla} \dicDirectTranslationCS{zalepit, zalepovat, zapečetit, zapečeťovat (obálku ap.)}  \textbf{3.} \dicFlx{acc} \dicSynonym*{breiðast út} \dicDirectTranslationCS{(roz)šířit se, rozšiřovat se};  \dicPhraseIS{læsa sig um e‑ð} \dicDirectTranslationCS{rozšířit se po (čem)}  \textbf{4.} \dicFlx{dat} \dicDirectTranslationCS{sevřít, svírat (v~drápech\,/\addthin zubech)};  \dicIdiom{læsast}{ \dicPhraseIS{læsast}} \dicFlx{refl} {\textbf{a.}} \dicDirectTranslationCS{zavřít se (klíče v~domě ap.)};  {\textbf{b.}} \dicDirectTranslationCS{sevřít, chopit se}
\dicEntry[læsi] \dicTerm{læsi\smash{\textsuperscript{1}}} \dicIPA{{l}{a}{i}{\textlengthmark}{s}{\textsci}} \dicPos{n}[2] \dicFlx{(‑s)}[20] \dicDirectTranslationCS{gramotnost, dovednost čtení}
\dicEntry[læsi] \dicTerm{læsi\smash{\textsuperscript{2}}} \dicIPA{{l}{a}{i}{\textlengthmark}{s}{\textsci}} \dicPos{v} \dicFlx{con pf sg 1 pers} \dicLink{lesa}
\dicEntry[læsilegur] \dicTerm{læsi··legur} \dicIPA{{l}{a}{i}{\textlengthmark}{s}{\textsci}{l}{\textepsilon}{\textbabygamma}{\textscy}{\textsubring{r}}} \dicPos{adj}[1]\dicFlx{}[-8] \textbf{1.} \dicDirectTranslationCS{čitelný} \dicExampleIS{læsileg skrift} \dicExampleCS{čitelné písmo}  \textbf{2.} \dicDirectTranslationCS{čtivý} \dicExampleIS{læsileg bók} \dicExampleCS{čtivá knížka}
\dicEntry[læsing] \dicTerm{læs··ing} \dicIPA{{l}{a}{i}{\textlengthmark}{s}{i}{\ng}{\r{g}}} \dicPos{f}[4] \dicFlx{(‑ar, ‑ar)}[5] \textbf{1.} \dicSynonym{lás} \dicDirectTranslationCS{zámek, zamykání (u~dveří ap.)} \dicExampleIS{læsing á dyrunum} \dicExampleCS{zámek u~dveří}  \textbf{2.} \dicSynonym{lokun} \dicDirectTranslationCS{(u)zamčení, (u)zamknutí}
\dicEntry[læstur] \dicTerm{læstur} \dicsymFrequent\  \dicIPA{{l}{a}{i}{s}{\textsubring{d}}{\textscy}{\textsubring{r}}} \dicPos{adj}[1]\dicFlx{}[-10] \dicSynonym*{aflokaður} \dicDirectTranslationCS{(u)zamčený, (u)zamknutý} \dicExampleIS{læstar dyr} \dicExampleCS{zamknuté dveře}
\dicEntry[læt] \dicTerm{læt} \dicIPA{{l}{a}{i}{\textlengthmark}{\textsubring{d}}} \dicPos{v} \dicFlx{ind praes sg 1 pers} \dicLink{láta}
\dicEntry[læti] \dicTerm{læti} \dicsymFrequent\  \dicIPA{{l}{a}{i}{\textlengthmark}{\textsubring{d}}{\textsci}} \dicPos{n}[2] \dicFlx{pl}[19] \dicDirectTranslationCS{rozruch, povyk} \dicExampleIS{vera með læti} \dicExampleCS{dělat povyk};  \dicPhraseIS{kunna sér ekki læti} \dicLangCat{přen.} \dicDirectTranslationCS{být štěstím bez sebe}
\dicEntry[lævirki] \dicTerm{læ··virk|i} \dicIPA{{l}{a}{i}{\textlengthmark}{v}{\textsci}{\textsubring{r}}{\r{\textObardotlessj}}{\textsci}} \dicPos{m}[1] \dicFlx{(‑ja, ‑jar)}[14] \dicFieldCat{zool.} \dicDirectTranslationCS{skřivan} \textit{(l.~{\textLA{Alauda}})}  \dicsymPhoto\ 
\dicFigure{ds_image_laevirki_0_2.jpg}{Lævirki}{Lævirki - Daniel Pettersson, CC BY-SA 2.5}
\dicEntry[lævís] \dicTerm{læ··vís} \dicIPA{{l}{a}{i}{\textlengthmark}{v}{i}{s}} \dicPos{adj}[5]\dicFlx{}[-1] \dicSynonym{slóttugur} \dicDirectTranslationCS{záludný, zákeřný}
\dicEntry[löðrungur] \dicTerm{löðr··ung|ur} \dicIPA{{l}{\oe}{ð}{r}{u}{\ng}{\r{g}}{\textscy}{\textsubring{r}}} \dicPos{m}[6] \dicFlx{(‑s, ‑ar)}[8] \dicSynonym{kinnhestur} \dicDirectTranslationCS{facka, políček} \dicExampleIS{gefa e‑m löðrungur} \dicExampleCS{dát (komu) facku}
\dicEntry[löður] \dicTerm{löður} \dicIPA{{l}{\oe}{\textlengthmark}{ð}{\textscy}{\textsubring{r}}} \dicPos{n}[2] \dicFlx{(‑s)}[28] \dicSynonym{froða} \dicDirectTranslationCS{(vodní) pěna}
\dicEntry[lög] \dicTerm{lög} \dicIPA{{l}{\oe}{\textlengthmark}{x}} \dicPos{n}[2] \dicFlx{pl}[9] \dicFieldCat{práv.} \dicDirectTranslationCS{zákon};  \dicPhraseIS{leiða e‑ð í lög} \dicDirectTranslationCS{uzákonit (co)};  \dicPhraseIS{setja lög} \dicDirectTranslationCS{vydat zákon}
\dicEntry[lögaldur] \dicTerm{lög··ald|ur} \dicIPA{{l}{\oe}{\textlengthmark}{\textbabygamma}{a}{l}{\textsubring{d}}{\textscy}{\textsubring{r}}} \dicPos{m}[5] \dicFlx{(‑urs, ‑rar)}[3] \dicFieldCat{práv.} \dicDirectTranslationCS{plnoletost, zletilost, zákonný věk};  \dicPhraseIS{ná lögaldri} \dicDirectTranslationCS{dosáhnout plnoletosti}
\dicEntry[lögbann] \dicTerm{lög··bann} \dicIPA{{l}{\oe}{\textbabygamma}{\textsubring{b}}{a}{\textsubring{n}}} \dicPos{n}[2] \dicFlx{(‑s)}[2] \dicFieldCat{práv.} \dicDirectTranslationCS{soudní zákaz}
\dicEntry[lögbrot] \dicTerm{lög··brot} \dicIPA{{l}{\oe}{\textbabygamma}{\textsubring{b}}{r}{\textopeno}{\textsubring{d}}} \dicPos{n}[2] \dicFlx{(‑s, ‑)}[5] \dicFieldCat{práv.} \dicDirectTranslationCS{porušení zákona, delikt}
\dicEntry[lögfesta] \dicTerm{lög··fest|a} \dicIPA{{l}{\oe}{x}{f}{\textepsilon}{s}{\textsubring{d}}{a}} \dicPos{v}[2] \dicFlx{(‑i, ‑)}[12] \dicFlx{acc} \dicFieldCat{práv.} \dicDirectTranslationCS{uzákonit, legalizovat, kodifikovat}
\dicEntry[lögformlegur] \dicTerm{lög··form·legur} \dicIPA{{l}{\oe}{x}{f}{\textopeno}{r}{m}{l}{\textepsilon}{\textbabygamma}{\textscy}{\textsubring{r}}} \dicPos{adj}[1]\dicFlx{}[-8] \dicFieldCat{práv.} \dicDirectTranslationCS{zákonný}
\dicEntry[lögfræði] \dicTerm{lög··fræð|i} \dicIPA{{l}{\oe}{x}{f}{r}{a}{i}{ð}{\textsci}} \dicPos{f}[3] \dicFlx{(‑i)}[3] \dicDirectTranslationCS{právní věda, práva}
\dicEntry[lögfræðideild] \dicTerm{lög·fræði··deild} \dicIPA{{l}{\oe}{x}{f}{r}{a}{i}{ð}{\textsci}{\textsubring{d}}{ei}{l}{\textsubring{d}}} \dicPos{f}[7] \dicFlx{(‑ar, ‑ir)}[1] \dicFieldCat{škol.} \dicDirectTranslationCS{právnická fakulta}
\dicEntry[lögfræðilegur] \dicTerm{lög·fræði··legur} \dicIPA{{l}{\oe}{x}{f}{r}{a}{i}{ð}{\textsci}{l}{\textepsilon}{\textbabygamma}{\textscy}{\textsubring{r}}} \dicPos{adj}[1]\dicFlx{}[-8] \dicDirectTranslationCS{právní} \dicExampleIS{lögfræðilegur ráðgjafi} \dicExampleCS{právní poradce}
\dicEntry[lögfræðingur] \dicTerm{lög·fræð··ing|ur} \dicsymFrequent\  \dicIPA{{l}{\oe}{x}{f}{r}{a}{i}{ð}{i}{\ng}{\r{g}}{\textscy}{\textsubring{r}}} \dicPos{m}[6] \dicFlx{(‑s, ‑ar)}[8] \dicDirectTranslationCS{právník, právnička} \dicExampleIS{aðstoð lögfræðings} \dicExampleCS{pomoc právníka}
\dicEntry[lögg] \dicTerm{lögg} \dicIPA{{l}{\oe}{\r{g}}{\textlengthmark}} \dicPos{f}[7] \dicFlx{(laggar\,/\addthin ‑var, laggir\,/\addthin ‑ir\,/\addthin ‑var)}[30] \dicDirectTranslationCS{kapka, doušek (v~láhvi ap.)} \dicExampleIS{kaffilögg} \dicExampleCS{doušek kávy};  \dicPhraseIS{setja e‑ð á laggirnar} \dicDirectTranslationCS{založit (co), zřídit (co)} \dicExampleIS{setja fyrirtæki á laggirnar} \dicExampleCS{zřídit společnost};  \dicPhraseIS{koma e‑u á laggirnar} \dicDirectTranslationCS{založit (co), zřídit (co)}
\dicEntry[lögga] \dicTerm{lögg|a} \dicsymFrequent\  \dicIPA{{l}{\oe}{\r{g}}{\textlengthmark}{a}} \dicPos{f}[1] \dicFlx{(‑u, ‑ur)}[7] \dicLangCat{hovor.} \dicSynonym{lögreglumaður} \dicDirectTranslationCS{polda, policajt(ka)} \dicExampleIS{löggur í kvikmyndum} \dicExampleCS{poldové ve filmech}
\dicEntry[löggilda] \dicTerm{lög··gil|da} \dicIPA{{l}{\oe}{\textbabygamma}{\r{\textObardotlessj}}{\textsci}{l}{\textsubring{d}}{a}} \dicPos{v}[2] \dicFlx{(‑ti, ‑t)}[37] \dicFlx{acc} \dicFieldCat{práv.} \dicDirectTranslationCS{oprávnit, zmocnit, autorizovat}
\dicEntry[löggilding] \dicTerm{lög··gild·ing} \dicIPA{{l}{\oe}{\textbabygamma}{\r{\textObardotlessj}}{\textsci}{l}{\textsubring{d}}{i}{\ng}{\r{g}}} \dicPos{f}[4] \dicFlx{(‑ar, ‑ar)}[5] \dicFieldCat{práv.} \dicDirectTranslationCS{oprávnění, zmocnění, autorizace}
\dicEntry[löggiltur] \dicTerm{lög··giltur} \dicIPA{{l}{\oe}{\textbabygamma}{\r{\textObardotlessj}}{\textsci}{\textsubring{l}}{\textsubring{d}}{\textscy}{\textsubring{r}}} \dicPos{adj}[1]\dicFlx{}[-13] \dicFieldCat{práv.} \dicDirectTranslationCS{oprávněný, autorizovaný};  \dicPhraseIS{löggiltur skjalaþýðandi} \dicDirectTranslationCS{úřední překladatel(ka)}
\dicEntry[löggjafarvald] \dicTerm{lög·gjafar··vald} \dicIPA{{l}{\oe}{\textbabygamma}{\r{\textObardotlessj}}{a}{v}{a}{r}{v}{a}{l}{\textsubring{d}}} \dicPos{n}[2] \dicFlx{(‑s)}[2] \dicDirectTranslationCS{zákonodárná moc}
\dicEntry[löggjafarþing] \dicTerm{lög·gjafar··þing} \dicIPA{{l}{\oe}{\textbabygamma}{\r{\textObardotlessj}}{a}{f}{a}{\textsubring{r}}{\texttheta}{i}{\ng}{\r{g}}} \dicPos{n}[2] \dicFlx{(‑s, ‑)}[5] \dicFieldCat{práv.} \dicDirectTranslationCS{zákonodárný sbor}
\dicEntry[löggjafi] \dicTerm{lög··gjaf|i} \dicIPA{{l}{\oe}{\textbabygamma}{\r{\textObardotlessj}}{a}{v}{\textsci}} \dicPos{m}[1] \dicFlx{(‑a, ‑ar)}[8] \dicFieldCat{pol.} \dicDirectTranslationCS{zákonodárce, zákonodárkyně}
\dicEntry[löggjöf] \dicTerm{lög··|gjöf} \dicIPA{{l}{\oe}{\textbabygamma}{\r{\textObardotlessj}}{\oe}{f}} \dicPos{f}[7] \dicFlx{(‑gjafar)}[19] \dicFieldCat{práv.} \dicDirectTranslationCS{legislativa, zákonodárství}
\dicEntry[löggæsla] \dicTerm{lög··gæsl|a} \dicIPA{{l}{\oe}{\textbabygamma}{\r{\textObardotlessj}}{a}{i}{s}{\textsubring{d}}{l}{a}} \dicPos{f}[1] \dicFlx{(‑u)}[5] \dicDirectTranslationCS{prosazování práva}
\dicEntry[löghald] \dicTerm{lög··hald} \dicIPA{{l}{\oe}{\textlengthmark}{x}{h}{a}{l}{\textsubring{d}}} \dicPos{n}[2] \dicFlx{(‑s)}[2] \dicFieldCat{práv.} \dicDirectTranslationCS{(úřední) zabavení, konfiskace};  \dicPhraseIS{leggja löghald á e‑u} \dicDirectTranslationCS{zkonfiskovat (co)}
\dicEntry[lögheimili] \dicTerm{lög··heimili} \dicIPA{{l}{\oe}{\textlengthmark}{x}{h}{ei}{m}{\textsci}{l}{\textsci}} \dicPos{n}[2] \dicFlx{(‑s, ‑)}[14] \dicFieldCat{práv.} \dicDirectTranslationCS{trvalé bydliště}
\dicEntry[löghlýðinn] \dicTerm{lög··hlýðinn} \dicIPA{{l}{\oe}{x}{\textsubring{l}}{i}{ð}{\textsci}{\textsubring{n}}} \dicPos{adj}[6]\dicFlx{}[-2] \dicDirectTranslationCS{dodržující zákony}
\dicEntry[löghyggja] \dicTerm{lög··hyggj|a} \dicIPA{{l}{\oe}{\textlengthmark}{x}{h}{\textsci}{\r{\textObardotlessj}}{a}} \dicPos{f}[1] \dicFlx{(‑u)}[5] \dicFieldCat{filos.} \dicDirectTranslationCS{determinismus}
\dicEntry[löglegur] \dicTerm{lög··legur} \dicIPA{{l}{\oe}{\textbabygamma}{l}{\textepsilon}{\textbabygamma}{\textscy}{\textsubring{r}}} \dicPos{adj}[1]\dicFlx{}[-8] \dicDirectTranslationCS{legální, právoplatný}
\dicEntry[lögmaður] \dicTerm{lög··|maður} \dicsymFrequent\  \dicIPA{{l}{\oe}{\textbabygamma}{m}{a}{ð}{\textscy}{\textsubring{r}}} \dicPos{m}[13] \dicFlx{(‑manns, ‑menn)}[2] \dicDirectTranslationCS{advokát(ka)} \dicExampleIS{Vantar þig lögmann?} \dicExampleCS{Potřebuješ advokáta?}
\dicEntry[lögmál] \dicTerm{lög··mál} \dicIPA{{l}{\oe}{\textbabygamma}{m}{au}{\textsubring{l}}} \dicPos{n}[2] \dicFlx{(‑s, ‑)}[5] \textbf{1.} \dicSynonym{lög} \dicDirectTranslationCS{zákon}  \textbf{2.} \dicDirectTranslationCS{zákon, řád (přírody ap.)} \dicExampleIS{lögmál lífsins} \dicExampleCS{zákon života}  \textbf{3.} \dicFieldCat{filos.} \dicSynonym{grundvallarregla} \dicDirectTranslationCS{princip, zákonitost}
\dicEntry[lögmæti] \dicTerm{lög··mæti} \dicIPA{{l}{\oe}{\textbabygamma}{m}{a}{i}{\textsubring{d}}{\textsci}} \dicPos{n}[2] \dicFlx{(‑s)}[20] \dicDirectTranslationCS{legálnost, zákonnost}
\dicEntry[lögmætur] \dicTerm{lög··mætur} \dicIPA{{l}{\oe}{\textbabygamma}{m}{a}{i}{\textsubring{d}}{\textscy}{\textsubring{r}}} \dicPos{adj}[1]\dicFlx{}[-1] \dicSynonym{gildur} \dicDirectTranslationCS{legální, zákonný}
\dicEntry[lögn] \dicTerm{lögn} \dicIPA{{l}{\oe}{\r{g}}{\textsubring{n}}} \dicPos{f}[7] \dicFlx{(lagnar, lagnir)}[16] \textbf{1.} \dicSynonym{leiðsla} \dicDirectTranslationCS{vedení, instalace} \dicExampleIS{raflögn} \dicExampleCS{elektrické vedení}  \textbf{2.} \dicSynonym*{það að leggja net} \dicDirectTranslationCS{pokládání sítě}
\dicEntry[lögregla] \dicTerm{lög··regl|a} \dicsymFrequent\  \dicIPA{{l}{\oe}{\textbabygamma}{r}{\textepsilon}{\r{g}}{l}{a}} \dicPos{f}[1] \dicFlx{(‑u, ‑ur)}[19] \textbf{1.} \dicSynonym*{stofnun til löggæslu} \dicDirectTranslationCS{policie} \dicExampleIS{Lögreglan leitar að vitnum að árekstri.} \dicExampleCS{Policie hledá svědky srážky vozidel.}  \textbf{2.} \dicSynonym{lögreglumaður} \dicDirectTranslationCS{policista, policistka}
\dicEntry[lögreglubíll] \dicTerm{lög·reglu··bíl|l} \dicsymFrequent\  \dicIPA{{l}{\oe}{\textbabygamma}{r}{\textepsilon}{\r{g}}{l}{\textscy}{\textsubring{b}}{i}{\textsubring{d}}{\textsubring{l}}} \dicPos{m}[6] \dicFlx{(‑s, ‑ar)}[48] \dicDirectTranslationCS{policejní auto\,/\addthin vůz} \dicExampleIS{Lögreglubílar lentu í árekstri.} \dicExampleCS{Policejní auta se srazila.}
\dicEntry[lögreglumaður] \dicTerm{lög·reglu··|maður} \dicsymFrequent\  \dicIPA{{l}{\oe}{\textbabygamma}{r}{\textepsilon}{\r{g}}{l}{\textscy}{m}{a}{ð}{\textscy}{\textsubring{r}}} \dicPos{m}[13] \dicFlx{(‑manns, ‑menn)}[2] \dicDirectTranslationCS{policista, policistka} \dicExampleIS{Lögreglumenn réðust inn í húsið.} \dicExampleCS{Policisté vtrhli do domu.}
\dicEntry[lögreglusamþykkt] \dicTerm{lög·reglu··sam·þykkt} \dicIPA{{l}{\oe}{\textbabygamma}{r}{\textepsilon}{\r{g}}{l}{\textscy}{s}{a}{m}{\texttheta}{\textsci}{x}{\textsubring{d}}} \dicPos{f}[7] \dicFlx{(‑ar, ‑ir)}[1] \dicDirectTranslationCS{policejní nařízení}
\dicEntry[lögreglustjóri] \dicTerm{lög·reglu··stjór|i} \dicIPA{{l}{\oe}{\textbabygamma}{r}{\textepsilon}{\r{g}}{l}{\textscy}{s}{\textsubring{d}}{j}{ou}{r}{\textsci}} \dicPos{m}[1] \dicFlx{(‑a, ‑ar)}[1] \dicDirectTranslationCS{policejní prezident(ka)}
\dicEntry[lögreglustöð] \dicTerm{lög·reglu··stöð} \dicIPA{{l}{\oe}{\textbabygamma}{r}{\textepsilon}{\r{g}}{l}{\textscy}{s}{\textsubring{d}}{\oe}{\texttheta}} \dicPos{f}[6] \dicFlx{(‑var, ‑var)}[1] \dicDirectTranslationCS{policejní stanice}
\dicEntry[lögregluvarðstjóri] \dicTerm{lög·reglu··varð·stjór|i} \dicIPA{{l}\-{\oe}\-{\textbabygamma}\-{r}\-{\textepsilon}\-{\r{g}}\-{l}\-{\textscy}\-{v}\-{a}\-{r}\-{ð}\-{s}\-{\textsubring{d}}\-{j}\-{ou}\-{r}\-{\textsci}\-} \dicPos{m}[1] \dicFlx{(‑a, ‑ar)}[1] \dicDirectTranslationCS{policejní komisař(ka)}
\dicEntry[lögregluvarðstofa] \dicTerm{lög·reglu··varð·stof|a} \dicIPA{{l}\-{\oe}\-{\textbabygamma}\-{r}\-{\textepsilon}\-{\r{g}}\-{l}\-{\textscy}\-{v}\-{a}\-{r}\-{ð}\-{s}\-{\textsubring{d}}\-{\textopeno}\-{v}\-{a}\-} \dicPos{f}[1] \dicFlx{(‑u, ‑ur)}[7] \dicDirectTranslationCS{okrsková služebna, policejní stanice}
\dicEntry[lögregluþjónn] \dicTerm{lög·reglu··þjón|n} \dicsymFrequent\  \dicIPA{{l}{\oe}{\textbabygamma}{r}{\textepsilon}{\r{g}}{l}{\textscy}{\texttheta}{j}{ou}{\textsubring{d}}{\textsubring{n}}} \dicPos{m}[6] \dicFlx{(‑s, ‑ar)}[42] \dicDirectTranslationCS{policista, policistka} \dicExampleIS{lögregluþjónar á næturvakt} \dicExampleCS{policisté na noční službě}
\dicEntry[Lögrétta] \dicTerm{Lög··rétt|a} \dicIPA{{l}{\oe}{\textbabygamma}{r}{j}{\textepsilon}{h}{\textsubring{d}}{a}} \dicPos{f}[1] \dicFlx{(‑u)}[6] \dicFieldCat{geog., hist.} \dicDirectTranslationCS{Lögrétta} \dicIndirectTranslationCS{(zákonodárný a~soudní sněm nacházející se v~Altingu)}
\dicEntry[lögræði] \dicTerm{lög··ræði} \dicIPA{{l}{\oe}{\textbabygamma}{r}{a}{i}{ð}{\textsci}} \dicPos{n}[2] \dicFlx{(‑s)}[20] \dicFieldCat{práv.} \dicDirectTranslationCS{svéprávnost};  \dicPhraseIS{svipta e‑n lögræði} \dicDirectTranslationCS{zbavit (koho) svéprávnosti}
\dicEntry[lögsótt] \dicTerm{lög··sótt} \dicIPA{{l}{\oe}{x}{s}{ou}{h}{\textsubring{d}}} \dicPos{v} \dicFlx{supin} \dicLink{lögsækja}
\dicEntry[lögsótti] \dicTerm{lög··sótti} \dicIPA{{l}{\oe}{x}{s}{ou}{h}{\textsubring{d}}{\textsci}} \dicPos{v} \dicFlx{ind pf sg 1 pers} \dicLink{lögsækja}
\dicEntry[lögsóttum] \dicTerm{lög··sóttum} \dicIPA{{l}{\oe}{x}{s}{ou}{h}{\textsubring{d}}{\textscy}{\textsubring{m}}} \dicPos{v} \dicFlx{ind pf pl 1 pers} \dicLink{lögsækja}
\dicEntry[lögsæki] \dicTerm{lög··sæki} \dicIPA{{l}{\oe}{x}{s}{a}{i}{\r{\textObardotlessj}}{\textsci}} \dicPos{v} \dicFlx{ind praes sg 1 pers} \dicLink{lögsækja}
\dicEntry[lögsækja] \dicTerm{lög··|sækja} \dicIPA{{l}{\oe}{x}{s}{a}{i}{\r{\textObardotlessj}}{a}} \dicPos{v}[5] \dicFlx{(‑sæki, ‑sótti, ‑sóttum, ‑sækti, ‑sótt)}[7] \dicFlx{acc} \dicFieldCat{práv.} \dicDirectTranslationCS{žalovat, podat\,/\addthin podávat žalobu} \dicExampleIS{lögsækja e‑n fyrir e‑ð} \dicExampleCS{žalovat (koho) za (co)}
\dicEntry[lögsækti] \dicTerm{lög··sækti} \dicIPA{{l}{\oe}{x}{s}{a}{i}{x}{\textsubring{d}}{\textsci}} \dicPos{v} \dicFlx{con pf sg 1 pers} \dicLink{lögsækja}
\dicEntry[lögtak] \dicTerm{lög··|tak} \dicIPA{{l}{\oe}{x}{t\smash{\textsuperscript{h}}}{a}{\r{g}}} \dicPos{n}[2] \dicFlx{(‑taks, ‑tök)}[8] \dicFieldCat{práv.} \dicDirectTranslationCS{exekuce}
\dicEntry[lögun] \dicTerm{lögun} \dicsymFrequent\  \dicIPA{{l}{\oe}{\textlengthmark}{\textbabygamma}{\textscy}{\textsubring{n}}} \dicPos{f}[7] \dicFlx{(‑ar, laganir)}[16] \textbf{1.} \dicSynonym{form} \dicDirectTranslationCS{tvar, podoba} \dicExampleIS{breyta löguninni} \dicExampleCS{změnit podobu}  \textbf{2.} \dicSynonym*{kaffilögun o.þ.h.} \dicDirectTranslationCS{příprava (kávy ap.)} \dicExampleIS{lögun á kaffi} \dicExampleCS{příprava kávy}
\dicEntry[lögur] \dicTerm{lögur} \dicIPA{{l}{\oe}{\textlengthmark}{\textbabygamma}{\textscy}{\textsubring{r}}} \dicPos{m}[11] \dicFlx{(lagar, legir)}[5] \textbf{1.} \dicDirectTranslationCS{kapalina, tekutina}  \textbf{2.} \dicSynonym{vatn} \dicDirectTranslationCS{voda, vodní plocha} \dicExampleIS{á láði og legi} \dicExampleCS{na souši a~na vodě}
\dicEntry[lök] \dicTerm{lök\smash{\textsuperscript{1}}} \dicIPA{{l}{\oe}{\textlengthmark}{\r{g}}} \dicPos{n} \dicFlx{pl nom} \dicLink{lak\smash{\textsuperscript{1}}}
\dicEntry[lök] \dicTerm{lök\smash{\textsuperscript{2}}} \dicIPA{{l}{\oe}{\textlengthmark}{\r{g}}} \dicPos{adj} \dicFlx{f sg nom pos} \dicLink{lakur}
\dicEntry[lökk] \dicTerm{lökk} \dicIPA{{l}{\oe}{h}{\r{g}}} \dicPos{n} \dicFlx{pl nom} \dicLink{lakk}
\dicEntry[lömb] \dicTerm{lömb} \dicIPA{{l}{\oe}{m}{\textsubring{b}}} \dicPos{n} \dicFlx{pl nom} \dicLink{lamb}
\dicEntry[lömdum] \dicTerm{lömdum} \dicIPA{{l}{\oe}{m}{\textsubring{d}}{\textscy}{\textsubring{m}}} \dicPos{v} \dicFlx{ind pf pl 1 pers} \dicLink{lemja}
\dicEntry[lömuð] \dicTerm{lömuð} \dicIPA{{l}{\oe}{\textlengthmark}{m}{\textscy}{\texttheta}} \dicPos{adj} \dicFlx{f sg nom pos} \dicLink{lamaður}
\dicEntry[lömun] \dicTerm{löm|un} \dicIPA{{l}{\oe}{\textlengthmark}{m}{\textscy}{\textsubring{n}}} \dicPos{f}[7] \dicFlx{(‑unar)}[12] \dicFieldCat{med.} \dicDirectTranslationCS{ochrnutí, paralýza} \dicExampleIS{fá lömun í fótinn} \dicExampleCS{ochrnout na nohu}
\dicEntry[lömunarveiki] \dicTerm{lömunar··veik|i} \dicIPA{{l}{\oe}{\textlengthmark}{m}{\textscy}{n}{a}{r}{v}{ei}{\r{\textObardotlessj}}{\textsci}} \dicPos{f}[3] \dicFlx{(‑i)}[3] \dicFieldCat{med.} \dicDirectTranslationCS{obrna}
\dicEntry[lönd] \dicTerm{lönd} \dicIPA{{l}{\oe}{n}{\textsubring{d}}} \dicPos{n} \dicFlx{pl nom} \dicLink{land}
\dicEntry[löndun] \dicTerm{löndun} \dicIPA{{l}{\oe}{n}{\textsubring{d}}{\textscy}{\textsubring{n}}} \dicPos{f}[7] \dicFlx{(‑ar, landanir)}[16] \dicSynonym*{uppskipun} \dicDirectTranslationCS{vykládka, vylodění}
\dicEntry[löng] \dicTerm{löng} \dicIPA{{l}{\oe i}{\ng}{\r{g}}} \dicPos{adj} \dicFlx{f sg nom pos} \dicLink{langur\smash{\textsuperscript{2}}}
\dicEntry[löngu] \dicTerm{löngu\smash{\textsuperscript{1}}} \dicIPA{{l}{\oe i}{\ng}{\r{g}}{\textscy}} \dicPos{f} \dicFlx{sg gen} \dicLink{langa\smash{\textsuperscript{1}}}
\dicEntry[löngu] \dicTerm{löngu\smash{\textsuperscript{2}}} \dicIPA{{l}{\oe i}{\ng}{\r{g}}{\textscy}} \dicPos{adv} \dicDirectTranslationCS{dlouho, dávno};  \dicPhraseIS{endur fyrir löngu} \dicFlx{adv} \dicDirectTranslationCS{dávno, před dávnými dobami};  \dicPhraseIS{fyrir (langa\,/\addthin margt) löngu} \dicFlx{adv} \dicDirectTranslationCS{dávno, před dávnou dobou};  \dicPhraseIS{fyrir lifandi löngu} \dicFlx{adv} \dicDirectTranslationCS{už\,/\addthin již (velmi) dávno};  \dicPhraseIS{löngu áður} \dicFlx{adv} \dicDirectTranslationCS{dlouho předtím};  \dicPhraseIS{löngu fyrr} \dicFlx{adv} \dicDirectTranslationCS{mnohem dříve, dlouho předtím};  \dicPhraseIS{löngu síðar\,/\addthin seinna} \dicFlx{adv} \dicDirectTranslationCS{dlouho poté, o~dost později}
\dicEntry[löngum] \dicTerm{löngum} \dicsymFrequent\  \dicIPA{{l}{\oe i}{\ng}{\r{g}}{\textscy}{\textsubring{m}}} \dicPos{adv} \dicDirectTranslationCS{dlouho} \dicExampleIS{Það hefur löngum verið vitað að sólböð eru hættuleg.} \dicExampleCS{Už se dlouho ví, že opalování je nebezpečné.}
\dicEntry[löngun] \dicTerm{löngun} \dicsymFrequent\  \dicIPA{{l}{\oe i}{\ng}{\r{g}}{\textscy}{\textsubring{n}}} \dicPos{f}[7] \dicFlx{(‑ar, langanir)}[16] \dicDirectTranslationCS{touha, toužení, tužba} \dicExampleIS{löngun í nýjan bíl} \dicExampleCS{touha po novém autě}
\dicEntry[löngutangar] \dicTerm{löngu··tangar} \dicIPA{{l}{\oe i}{\ng}{\r{g}}{\textscy}{t\smash{\textsuperscript{h}}}{au}{\ng}{\r{g}}{a}{\textsubring{r}}} \dicPos{f} \dicFlx{sg gen} \dicLink{langatöng}
\dicEntry[löngutengur] \dicTerm{löngu··tengur} \dicIPA{{l}{\oe i}{\ng}{\r{g}}{\textscy}{t\smash{\textsuperscript{h}}}{\textepsilon}{i}{\ng}{\r{g}}{\textscy}{\textsubring{r}}} \dicPos{f} \dicFlx{pl nom} \dicLink{langatöng}
\dicEntry[löngutöng] \dicTerm{löngu··|töng} \dicIPA{{l}{\oe i}{\ng}{\r{g}}{\textscy}{t\smash{\textsuperscript{h}}}{\oe i}{\ng}{\r{g}}} \dicPos{f}[8] \dicFlx{(‑tangar, ‑tengur\,/\addthin ‑tangir)}[9] \dicLink{langatöng}
\dicEntry[löpp] \dicTerm{löpp} \dicsymFrequent\  \dicIPA{{l}{\oe}{h}{\textsubring{b}}} \dicPos{f}[7] \dicFlx{(lappar, lappir)}[16] \dicDirectTranslationCS{noha, chodidlo} \dicExampleIS{vöðvar í löppum} \dicExampleCS{svaly na nohou};  \dicPhraseIS{koma sér á lappir} \dicDirectTranslationCS{vstát z~postele}
\dicEntry[löptum] \dicTerm{löptum} \dicIPA{{l}{\oe}{f}{\textsubring{d}}{\textscy}{\textsubring{m}}} \dicPos{v} \dicFlx{ind pf pl 1 pers} \dicLink{lepja}
\dicEntry[löstur] \dicTerm{löstur} \dicIPA{{l}{\oe}{s}{\textsubring{d}}{\textscy}{\textsubring{r}}} \dicPos{m}[11] \dicFlx{(lastar, lestir)}[5] \dicDirectTranslationCS{zlozvyk, neřest} \dicExampleIS{drykkjuskapur og aðrir lestir} \dicExampleCS{opilství a~jiné neřesti}
\dicEntry[löt] \dicTerm{löt} \dicIPA{{l}{\oe}{\textlengthmark}{\textsubring{d}}} \dicPos{adj} \dicFlx{f sg nom pos} \dicLink{latur}
\dicEntry[lötra] \dicTerm{lötr|a} \dicIPA{{l}{\oe}{\textlengthmark}{\textsubring{d}}{r}{a}} \dicPos{v}[1] \dicFlx{(‑aði)}[44] \dicDirectTranslationCS{loudat se, šinout se, courat se} \dicExampleIS{Hann lötraði á eftir hópnum.} \dicExampleCS{Coural se za skupinou.}
\dicEntry[löttum] \dicTerm{löttum} \dicIPA{{l}{\oe}{h}{\textsubring{d}}{\textscy}{\textsubring{m}}} \dicPos{v} \dicFlx{ind pf pl 1 pers} \dicLink{letja}
}
\fi
\restoregeometry
\pagestyle{plain}

\fi

% ================================================================ PHONETICS =

\ifinputphon

\cleardoublepage

\chapter{Seznam islandských fonémů}              \label{sec:phon_phonems}

\begin{table}[h]
\begin{tabular}{lllll} \toprule

\textbf{foném} & \textbf{vysvětlivka} & \textbf{příklad} & \textbf{fonetický přepis} & \textbf{překlad} \\
\midrule
{\textipa{[{a}]}} & otevřená přední nezaokrouhlená samohl. & rata & {\textipa{[{r}{a}{\textlengthmark}{\textsubring{d}}{a}]}} & trefit \\ 
{\textipa{[{\textepsilon}]}} & polootevřená přední nezaokrouhlená samohl. & meta & {\textipa{[{m}{\textepsilon}{\textlengthmark}{\textsubring{d}}{a}]}} & ocenit \\ 
{\textipa{[{e}]}} & polozavřená přední nezaokrouhlená samohl. & heiður & {\textipa{[{h}{ei}{\textlengthmark}{ð}{\textscy}{\textsubring{r}}]}} & čest \\
{\textipa{[{\textsci}]}} & téměř zavřená téměř přední nezaokrouhlená samohl. & vinna & {\textipa{[{v}{\textsci}{n}{\textlengthmark}{a}]}} & pracovat \\ 
{\textipa{[{i}]}} & zavřená přední nezaokrouhlená samohl. & fínn & {\textipa{[{f}{i}{\textsubring{d}}{\textsubring{n}}]}} & jemný \\ 
{\textipa{[{\textopeno}]}} & polootevřená zadní zaokrouhlená samohl. & lofa & {\textipa{[{l}{\textopeno}{\textlengthmark}{v}{a}]}} & slibovat \\ 
{\textipa{[{o}]}} & polozavřená zadní zaokrouhlená samohl. & bóndi & {\textipa{[{\textsubring{b}}{ou}{n}{\textsubring{d}}{\textsci}]}} & farmář \\ 
{\textipa{[{\textscy}]}} & téměř zavřená téměř přední zaokrouhlená samohl. & hundur & {\textipa{[{h}{\textscy}{n}{\textsubring{d}}{\textscy}{\textsubring{r}}]}} & pes \\ 
{\textipa{[{u}]}} & zavřená zadní zaokrouhlená samohl. & súkkulaði & {\textipa{[{s}{u}{h}{\r{g}}{\textscy}{l}{a}{ð}{\textsci}]}} & čokoláda \\ 
{\textipa{[{\oe}]}} & téměř otevřená přední nezaokrouhlená samohl. & lönd & {\textipa{[{l}{\oe}{n}{\textsubring{d}}]}} & země \\ 
{\textipa{[{m}]}} & bilabiální nazála & mýri & {\textipa{[{m}{i}{\textlengthmark}{r}{\textsci}]}} & mokřina \\ 
{\textipa{[{\textsubring{m}}]}} & neznělá bilabiální nazála & þyrmt & {\textipa{[{\texttheta}{\textsci}{\textsubring{r}}{\textsubring{m}}{\textsubring{d}}]}} & {\textit{$\shortrightarrow$ þyrma}} \\  
{\textipa{[{n}]}} & alveolární nazála & nemandi & {\textipa{[{n}{\textepsilon}{\textlengthmark}{m}{a}{n}{\textsubring{d}}{\textsci}]}} & žák \\ 
{\textipa{[{\textsubring{n}}]}} & neznělá alveolární nazála & mennt & {\textipa{[{m}{\textepsilon}{\textsubring{n}}{\textsubring{d}}]}} & vzdělání \\ 
{\textipa{[{\textltailn}]}} & palatální nazála & angi & {\textipa{[{au}{\textltailn}{\r{\textObardotlessj}}{\textsci}]}} & drobek \\ 
{\textipa{[{\r{\textltailn}}]}} & neznělá palatální nazála & banki & {\textipa{[{\textsubring{b}}{au}{\r{\textltailn}}{\r{\textObardotlessj}}{\textsci}]}} & banka \\
{\textipa{[{\ng}]}} & velární nazála & skyggndur & {\textipa{[{s}{\r{\textObardotlessj}}{\textsci}{\ng}{\textsubring{d}}{\textscy}{\textsubring{r}}]}} & {\textit{$\shortrightarrow$ skyggna}} \\ 
{\textipa{[{\r{\ng}}]}} & neznělá velární nazála & skyggnt & {\textipa{[{s}{\r{\textObardotlessj}}{\textsci}{\r{\ng}}{\textsubring{d}}]}} & {\textit{$\shortrightarrow$ skyggna}} \\ 
{\textipa{[{\textsubring{b}}]}} & neznělá bilabiální ploziva & bara & {\textipa{[{\textsubring{b}}{a}{\textlengthmark}{r}{a}]}} & jenom \\ 
{\textipa{[{p\textsuperscript{h}}]}} & neznělá bilabiální ploziva s přídechem & prestur & {\textipa{[{p\textsuperscript{h}}{r}{\textepsilon}{s}{\textsubring{d}}{\textscy}{\textsubring{r}}]}} & kněz \\ 
{\textipa{[{\textsubring{d}}]}} & neznělá alveolární ploziva & dæmi & {\textipa{[{\textsubring{d}}{a}{i}{\textlengthmark}{m}{\textsci}]}} & příklad \\ 
{\textipa{[{t\textsuperscript{h}}]}} & neznělá alveolární ploziva s přídechem & tómur & {\textipa{[{t\textsuperscript{h}}{ou}{\textlengthmark}{m}{\textscy}{\textsubring{r}}]}} & prázdný \\ 
{\textipa{[{\r{\textObardotlessj}}]}} & neznělá palatální ploziva & gefa & {\textipa{[{\r{\textObardotlessj}}{\textepsilon}{\textlengthmark}{v}{a}]}} & dát \\ 
{\textipa{[{c\textsuperscript{h}}]}} & neznělá palatální ploziva s přídechem & kenna & {\textipa{[{c\textsuperscript{h}}{\textepsilon}{n}{\textlengthmark}{a}]}} & učit \\
{\textipa{[{\r{g}}]}} & neznělá velární ploziva & gata & {\textipa{[{\r{g}}{a}{\textlengthmark}{\textsubring{d}}{a}]}} & ulice \\ 
{\textipa{[{k\textsuperscript{h}}]}} & neznělá velární ploziva s přídechem & koma & {\textipa{[{k\textsuperscript{h}}{\textopeno}{\textlengthmark}{m}{a}]}} & přijít \\
{\textipa{[{v}]}} & znělá labiodentální frikativa & hafa & {\textipa{[{h}{a}{\textlengthmark}{v}{a}]}} & mít \\ 
{\textipa{[{f}]}} & neznělá labiodentální frikativa & finna & {\textipa{[{f}{\textsci}{n}{\textlengthmark}{a}]}} & najít \\ 
{\textipa{[{ð}]}} & znělá dentální frikativa & borða & {\textipa{[{\textsubring{b}}{\textopeno}{r}{ð}{a}]}} & jíst \\ 
{\textipa{[{\texttheta}]}} & neznělá dentální frikativa & þola & {\textipa{[{\texttheta}{\textopeno}{\textlengthmark}{l}{a}]}} & snést \\
{\textipa{[{s}]}} & neznělá alveolární frikativa & systir & {\textipa{[{s}{\textsci}{s}{\textsubring{d}}{\textsci}{\textsubring{r}}]}} & sestra \\ 
{\textipa{[{\c{c}}]}} & neznělá palatální frikativa & hjóla & {\textipa{[{\c{c}}{ou}{\textlengthmark}{l}{a}]}} & jet na kole \\ 
{\textipa{[{\textbabygamma}]}} & znělá velární frikativa & sagði & {\textipa{[{s}{a}{\textbabygamma}{ð}{\textsci}]}} & řekl \\ 
{\textipa{[{x}]}} & neznělá velární frikativa & lag & {\textipa{[{l}{a}{\textlengthmark}{x}]}} & píseň \\
{\textipa{[{h}]}} & neznělá glotální frikativa & heima & {\textipa{[{h}{ei}{\textlengthmark}{m}{a}]}} & doma \\ 
{\textipa{[{l}]}} & alveolární laterální aproximanta & loforð & {\textipa{[{l}{\textopeno}{\textlengthmark}{v}{\textopeno}{r}{\texttheta}]}} & slib \\ 
{\textipa{[{\textsubring{l}}]}} & neznělá alveolární laterální aproximanta & hlutur & {\textipa{[{\textsubring{l}}{\textscy}{\textlengthmark}{\textsubring{d}}{\textscy}{\textsubring{r}}]}} & věc \\  
{\textipa{[{j}]}} & palatální aproximanta & jól & {\textipa{[{j}{ou}{\textlengthmark}{\textsubring{l}}]}} & Vánoce \\ 
{\textipa{[{r}]}} & alveolární vibranta & reynsla & {\textipa{[{r}{ei}{n}{s}{\textsubring{d}}{l}{a}]}} & zkušenost \\ 
{\textipa{[{\textsubring{r}}]}} & neznělá alveolární vibranta & hraun & {\textipa{[{\textsubring{r}}{\oe i}{\textlengthmark}{\textsubring{n}}]}} & láva \\ 
\bottomrule
\end{tabular}
\end{table}

\clearpage

%\section{Výslovnost islandských samohlásek}      \label{sec:phon_vowels}
%\LTXtable{\columnwidth}{phon/vowels}
%\clearpage

%\section{Výslovnost islandských souhlásek}       \label{sec:phon_consonants}
%\LTXtable{\columnwidth}{phon/consonants}
%\clearpage
\fi

% =============================================================== MORPHOLOGY =

\ifinputmorpho

\cleardoublepage

\chapter{Morfologický klíč}                      \label{sec:morpho}

\section{Podstatné jméno, rod mužský}            \label{sec:morpho_m}
{\small\LTXfw{morpho/m}}
\clearpage

\section{Podstatné jméno, rod ženský}            \label{sec:morpho_f}
{\small\LTXfw{morpho/f}}
\clearpage

\section{Podstatné jméno, rod střední}           \label{sec:morpho_n}
{\small\LTXfw{morpho/n}}
\clearpage

\section{Přídavné jméno}                         \label{sec:morpho_adj}
{\small\LTXfw{morpho/adj}}
\clearpage

\section{Zájmeno}                                \label{sec:morpho_pron}
{\small\LTXfw{morpho/pron}}
 \begin{center}
\begin{minipage}[t]{.45\textwidth}
\begin{tabular}{l>{\footnotesize\itshape}lll}
 {\textbf{\textit{třída}}} & {\textit{pád}}   & \textit{sg} & \textit{pl}  \\ 
 \hline
\multirow{3}{*}{{{\textbf{pron} \Large{\textbf{17}}}}}  &  nom & ég & við   \\*
 & acc &  mig  & okkur  \\*
 & dat & mér & okkur   \\*
 & gen & mín  & okkar  \\
\hline

\multirow{3}{*}{{{\textbf{pron} \Large{\textbf{17}}}}}  &  nom & hann & þeir   \\*
 & acc &  hann  & þá  \\*
 & dat & honum & þeim   \\*
 & gen & hans  & þeirra  \\
\hline

\multirow{3}{*}{{{\textbf{pron} \Large{\textbf{17}}}}}  &  nom & hún & þær   \\*
 & acc &  hana  & þær  \\*
 & dat & henni & þeim   \\*
 & gen & hennar  & þeirra  \\
\hline

\multirow{3}{*}{{{\textbf{pron} \Large{\textbf{17}}}}}  &  nom &  & vér   \\*
 & acc &    & oss  \\*
 & dat &  & oss   \\*
 & gen &   & vor  \\
\hline

\multirow{3}{*}{{{\textbf{pron} \Large{\textbf{17}}}}}  &  nom &  & þér   \\*
 & acc &    & yður  \\*
 & dat &  & yður   \\*
 & gen &   & yðar  \\
\hline

\multirow{3}{*}{{{\textbf{pron} \Large{\textbf{17}}}}}  &  nom & þú & þið   \\*
 & acc &  þig  & ykkur  \\*
 & dat & þér & ykkur   \\*
 & gen & þín  & ykkar  \\
\hline

\multirow{3}{*}{{{\textbf{pron} \Large{\textbf{17}}}}}  &  nom & það & þau   \\*
 & acc &  það  & þau  \\*
 & dat & því & þeim   \\*
 & gen & þess  & þeirra  \\
\hline
\end{tabular}
\end{minipage}\hfil
\begin{minipage}[t]{.45\textwidth}
\begin{tabular}{l>{\footnotesize\itshape}lll}
\hline
\multirow{3}{*}{{{\textbf{pron} \Large{\textbf{18}}}}}  &  nom &  &    \\*
 & acc &  sig  &   \\*
 & dat & sér &    \\*
 & gen & sín  &   \\
 \hline
\end{tabular}
\end{minipage}
\end{center}

\clearpage

\section{Číslovka}                               \label{sec:morpho_num}
 {\small\LTXfw{morpho/num}}
\clearpage

\begin{landscape}
\section{Sloveso}                                \label{sec:morpho_v}
{\small\LTXfw{morpho/v}}
\end{landscape}

\clearpage

\fi

% ============================================================== BACK MATTER =

\ifmakebackmatter

% ------------------------------------------------ Photographs authors index -

\cleardoublepage
%\pagestyle{empty}
%\pagestyle{indexstyle}
%\phantomsection
%\addcontentsline{toc}{chapter}
%  {Seznam autorů fotografií a ilustrací}        
\label{sec:photo}
\printindex[figures]

% ------------------------------------------------------------- Bibliography -
\cleardoublepage
\chapter{Bibliografie}
%\twocolumn
%\phantomsection
%\addcontentsline{toc}{chapter}{Bibliografie}
\begin{multicols}{2}
\nocite{*}\printbibliography[heading=none]
\end{multicols}

% ------------------------------------------------------------------ Apendix -
\cleardoublepage
\onecolumn
\ifPDF
\else
\chapter{Apendix: Islandská gramatika}
\clearpage
\covergeometry  
\makeatletter
\includepdf[
  pages={1,2,3,4,5,6,7,8,9,10,11,12,13,14,15,16,17,18,19,20,21,22,23,24,25,26,27,28,29,30,31,32,33,34,35,36,37,38,39,40,41,42,43,44,45,46,47,48,49,50,51,52,53,54,55,56,57,58,59,60,61,62,63,64,65,66,67,68,69,70,71,72,73},
  turn=false,
  scale=0.79,
  keepaspectratio,
  offset={\dimexpr(\paperwidth-\csname Gm@layoutwidth\endcsname)/2-\csname Gm@layouthoffset\endcsname\relax}
  {\dimexpr(\paperheight-\csname  Gm@layoutheight\endcsname)/2-\csname Gm@layoutvoffset\endcsname\relax}
   ]{gram.pdf}
\makeatother
\restoregeometry  
\fi

% ------------------------------------------------------------------ Barcode -
\pagenumbering{gobble}
\ifPDF
\else
\cleardoublepage\null\clearpage % clear to verso
\vspace*\fill
\EANisbn[SC3]
\fi
\fi

% --------------------------------------------------------------- Back cover -

\ifmakecovers
  \cleardoublepage\null\clearpage % inner side of the cover
\covergeometry  
  \makecoverwith\backcoverimages
\restoregeometry  
\fi

\end{document}

